\documentclass[11pt]{article}
\usepackage[T1]{fontenc}
\usepackage[margin=1.02in]{geometry}
\usepackage{amsmath,amssymb,amsthm,mathtools,bm,mathrsfs}
\usepackage{booktabs,array,longtable,enumitem,microtype}
\usepackage[colorlinks=true,linkcolor=blue,citecolor=blue,urlcolor=blue]{hyperref}
\setlength{\emergencystretch}{3em}
\allowdisplaybreaks[2]
\newtheorem{theorem}{Theorem}[section]
\newtheorem{proposition}[theorem]{Proposition}
\newtheorem{lemma}[theorem]{Lemma}
\newtheorem{corollary}[theorem]{Corollary}
\theoremstyle{definition}
\newtheorem{definition}[theorem]{Definition}
\newtheorem{example}[theorem]{Example}
\theoremstyle{remark}
\newtheorem{remark}[theorem]{Remark}
\newcommand{\C}{\mathbb C}
\newcommand{\R}{\mathbb R}
\newcommand{\A}{\mathcal A}
\newcommand{\B}{\mathcal B}
\newcommand{\Hh}{\mathcal H}
\newcommand{\E}{\mathcal E}
\newcommand{\W}{\mathcal W}
\newcommand{\Nn}{\mathcal N}
\newcommand{\End}{\operatorname{End}}
\newcommand{\Hom}{\operatorname{Hom}}
\newcommand{\Span}{\operatorname{span}}
\newcommand{\Ran}{\operatorname{Ran}}
\newcommand{\Ker}{\operatorname{Ker}}
\newcommand{\Tr}{\operatorname{Tr}}
\newcommand{\rank}{\operatorname{rank}}
\newcommand{\supp}{\operatorname{supp}}
\newcommand{\diag}{\operatorname{diag}}
\newcommand{\id}{\mathrm{id}}
\newcommand{\HS}{\mathrm{HS}}
\newcommand{\op}{\mathrm{op}}
\newcommand{\norm}[1]{\lVert#1\rVert}
\newcommand{\ip}[2]{\langle#1,#2\rangle}
\newcommand{\src}[1]{\texttt{\small\nolinkurl{#1}}}
\newcolumntype{L}[1]{>{\raggedright\arraybackslash}p{#1}}
\title{\textbf{Historical Standard-Model Source Selection}\\[0.35em]
\large Word modules, reciprocal returns, and internal saturation\\[0.35em]
\normalsize Research Notes V1}
\author{Renewal Geometry research programme}
\date{}
\begin{document}
\maketitle
\begin{abstract}
The structural Standard-Model branch already has exact colour-bridge and weak-router criteria.  The remaining question is whether the required operations belong to the actual source rather than to a larger algebra of kinematically available multiplicity endomorphisms.  This note begins a separate source-selection series.  It first identifies word-induced invariant actions by compatibility with right multiplication and gives an exact defect formula.  A twelve-root example shows that even a full noncommutative history algebra can have large invariant word multiplicities without a physical internal colour factor.  The positive result is a reciprocal-return theorem: two nonzero operational ports through the same irreducible mediator force a colour bridge among their finite return words.  With mediator multiplicity and dynamics, a finite positive Krylov Gram gives a necessary-and-sufficient return test and a quantitative sufficient reachability condition.  A separate subdirect-product theorem isolates the central separation needed when blockwise matrix saturation is assembled into a complete internal algebra.  These results provide a source-level replacement for an independently asserted bridge on a precisely specified route class.  They do not assert that the historical Store already supplies the required reciprocal ports, nor do they force the determinant or matter tensor types.
\end{abstract}
\setcounter{tocdepth}{2}
\tableofcontents

\section{Context, audited inputs, and the precise target}
\label{sel1:context}

\subsection{Why a new V1 is needed}
The duality manuscript classifies a finite joint spacetime--gauge--incidence packet and reconstructs its opposite multiplicity algebra.  The selection objection concerns a different implication: whether the literal operational source supplies the colour-eligible carrier, the colour bridge, the weak operations, and the later determinant/matter incidence.  A proof that an added bridge generates $M_3(\C)$ is not a proof that the bridge occurred.

This is a new note lineage, distinct from the previous spacetime--gauge rigidity notes.  V1 records the context and develops only the source-selection questions not settled by the audited inputs.  Subsequent versions are to preserve the V1 mathematical body, retain its labels, and append new development blocks before \verb|\end{document}|.  Any correction is to be explicitly identified.  The filenames use the suffixes \texttt{\_V1.tex}, \texttt{\_V2.tex}, and so forth.

The source audit includes the supplied Grand Tensor V076, the updated SM--Einstein monograph, and the two CMP manuscripts \cite{Grand,Continuum,Duality,Classical}.  It also follows the previously cited V106 result to the accessible cumulative \emph{SM Unconditional Emergence Attack} through V128 \cite{Attack}, especially V106--V110.  This extra check changes the starting frontier: V108 already constructs the weak reset packet on the repaired primitive, and V109 distinguishes a positive type-line Schur floor from an unnecessary nonzero type--private overlap.  Those results must not be rediscovered here.

\begin{center}
\small
\begin{tabular}{L{.37\textwidth}L{.55\textwidth}}
\toprule
Audited result retained as input & Source and scope\\
\midrule
External versus internal multiplicities; flat-depth obstruction & \cite{Attack}, \src{thm:V106-relative-commutant-census}, \src{thm:V106-flat-colour-obstruction}.\\[.35em]
One-router and full-block matrix generation & \cite{Attack}, \src{thm:V106-connected-router}, \src{thm:V107-full-block-router}. Not reproved as new generation theorems.\\[.35em]
Profile words versus operation words; external twirl & \cite{Attack}, \src{cth:V107-profile-no-go}, \src{thm:V107-twirl}.\\[.35em]
Weak reset occurrence on the repaired primitive and its colour-support obstruction & \cite{Attack}, \src{thm:V108-weak-closure}, \src{thm:V108-block-separation}. Strict old-alphabet admission is separately qualified there.\\[.35em]
Type-line Schur survival and metric-covariant colour-bridge criterion & \cite{Attack}, \src{thm:V109-type-short}, \src{thm:V109-one-bridge}.\\[.35em]
Determinant projection and its independent occurrence & \cite{Attack}, \src{thm:V110-det-incidence-iso}, \src{cth:V110-three-blocks-no-det}.\\[.35em]
Howe positivity, genericity, incidence rank and singular rigidity & \cite{Rigidity,Duality}. Their existing proofs are not repeated.\\
\bottomrule
\end{tabular}
\end{center}

\subsection{The actual advance attempted here}
There are three tasks.  First, an increasing history-word Hilbert space must not silently replace the exposed physical system: its invariant endomorphisms need not be word operations.  Second, a bridge may already be a \emph{return product} of two actual routes through a common mediator; in that case no independent bridge primitive is required.  Third, blockwise surjectivity does not always assemble into a full direct sum without separating repeated simple blocks.

The new proofs below address these tasks.  The earlier finite bimodule-alignment panel in V2 of \cite{Attack} and the principal matter-module theorem in \cite{Grand} remain separate inherited results: the present module test concerns endomorphisms of the entire regular word algebra and their descent to the exposed system.  The representation-theoretic tools are standard finite-dimensional methods; background is available in \cite{Etingof}.  The claims of development are relative to the audited programme, not claims of priority for regular representations, Schur averaging, or finite spectral interpolation.  Every specialized assertion is proved in the note.

\subsection{What is, and is not, called an operation}
An \emph{actual operation word} is an admissible composition of operators supplied by the source on specified input and output cuts.  A linear combination in their represented algebra is not automatically one newly executable event.  In particular, formal adjoints belong to the $*$-algebra; an executable reverse route requires physical admission of that reverse.  The reciprocal-return theorems state this admission explicitly.  Their positive sums certify that at least one already admissible labelled word is nonzero, rather than treating the sum as an invented physical operation.

A positive source Gram determines a source span and its inner product.  It does not, without multiplication and action data, determine which endomorphisms of that span are executable.  Neither a large trivial-isotypic word rank nor a nonzero source overlap substitutes for the latter data.

\section{The multiplication-compatible internal algebra}
\label{sel1:module}

Let $\A\subseteq\End(\Hh)$ be the finite unital $C^*$-algebra generated by the declared system operations on one fixed exposed carrier.  Let a finite group $G$ act by unitaries $u_g$ on $\Hh$.  The historically generated internal algebra on that carrier is
\begin{equation}
 \B_{\rm hist}:=\A\cap\{u_g:g\in G\}'.
 \label{sel1:eq:actual-internal}
\end{equation}
It is not, by definition, the entire kinematic commutant $\{u_g\}'$.

For the regular-word formulation assume additionally that $u_g\A u_g^*=\A$ and write $\alpha_g(a)=u_gau_g^*$.  Equip $\A$ with the normalized regular trace
\begin{equation}
 \tau_{\rm reg}(a)=\frac1D\sum_\beta n_\beta\Tr(a_\beta),
 \qquad \A\simeq\bigoplus_\beta M_{n_\beta}(\C),\qquad D=\dim_\C\A=\sum_\beta n_\beta^2.
 \label{sel1:eq:trace}
\end{equation}
Automorphisms preserve this trace.  Set $\W=L^2(\A,\tau_{\rm reg})$ and define
\[
 L_a x=ax,\qquad R_bx=xb,\qquad U_gx=\alpha_g(x).
\]
Right multiplication is an opposite-algebra representation; no commutativity of $\A$ is assumed.

\begin{theorem}[Word-module recognition of internal actions]
\label{sel1:thm:module}
On the regular word carrier,
\begin{equation}
 \boxed{R(\A)'\cap U(G)'=L(\A^G)=L(\B_{\rm hist}).}
 \label{sel1:eq:module-intersection}
\end{equation}
More explicitly, an endomorphism $T\in\End(\W)$ is induced by left multiplication by an invariant system word-algebra element if and only if it commutes with every right multiplication and every $U_g$.  When this holds, its unique coefficient is $a=T1\in\A^G$.
\end{theorem}
\begin{proof}
Every $L_a$ commutes with every $R_b$.  Conversely, if $T$ commutes with all $R_b$, then for each $b\in\A$,
\[
 T(b)=T(R_b1)=R_b(T1)=(T1)b.
\]
Thus $T=L_a$ for $a=T1$, uniquely.  Since $U_gL_aU_g^{-1}=L_{\alpha_g(a)}$, faithfulness of left multiplication shows that $L_a$ commutes with $U_g$ exactly when $\alpha_g(a)=a$.  The latter is equivalent to $a\in\A\cap\{u_g\}'$ on the original exposed carrier.
\end{proof}

This theorem is a representation check, not a postulate forbidding other kinds of physical operation.  A genuinely supplied superoperator or a new instrument on $\W$ can fail right-module compatibility.  It must then be counted as a separate supplied operation, not inferred from multiplication of the old system words.

\begin{proposition}[Exact module defect and an invariant approximation bound]
\label{sel1:prop:defect}
Let $(e_\nu)_{\nu=1}^D$ be a $\tau_{\rm reg}$-orthonormal basis of $\A$.  For arbitrary $T\in\End(\W)$, put $a=T1$ and
\[
 \delta_R(T)^2=\sum_{\nu=1}^D\norm{[T,R_{e_\nu}]1}_{\tau}^2,
 \qquad
 \delta_G(T)^2=\frac1{|G|}\sum_g\norm{[U_g,T]1}_{\tau}^2.
\]
For $E_G(a)=|G|^{-1}\sum_g\alpha_g(a)$,
\begin{align}
 \norm{T-L_a}_{\HS(\W)}^2&=\delta_R(T)^2,\label{sel1:eq:module-defect}\\
 \norm{T-L_{E_G(a)}}_{\HS(\W)}
 &\le\delta_R(T)+\sqrt{D/2}\,\delta_G(T).
 \label{sel1:eq:invariant-defect}
\end{align}
In particular, both defects vanish exactly for the algebra in \eqref{sel1:eq:module-intersection}.
\end{proposition}
\begin{proof}
For each basis vector,
$(T-L_a)e_\nu=T(R_{e_\nu}1)-R_{e_\nu}T1=[T,R_{e_\nu}]1$.
Summing squared norms gives \eqref{sel1:eq:module-defect}.  Since $U_g1=1$,
$[U_g,T]1=\alpha_g(a)-a$.  Orthogonal group averaging in the tracial Hilbert space gives
\[
 \frac1{|G|}\sum_g\norm{\alpha_g(a)-a}_\tau^2
 =2\norm{a-E_G(a)}_\tau^2.
\]
On each matrix summand left multiplication has $n_\beta$ copies of its defining representation, so
$\norm{L_b}_{\HS(\W)}^2=\sum_\beta n_\beta\Tr(b_\beta^*b_\beta)=D\norm b_\tau^2$.
The triangle inequality now gives \eqref{sel1:eq:invariant-defect}.
\end{proof}

\begin{corollary}[Word-rank growth is not a physical multiplicity-birth theorem]
\label{sel1:cor:word-birth}
If the exposed carrier and its external representation $u(G)$ are fixed, increasing the evaluated word span cannot change the multiplicities of $u(G)$ on that carrier.  It can reveal a larger subalgebra of \eqref{sel1:eq:actual-internal}.  Multiplicities of the induced conjugation action $U(G)$ on $L^2(\A)$ describe a different representation; its full commutant need not consist of word-induced system operations.
\end{corollary}
\begin{proof}
The multiplicities of the fixed matrices $u_g$ do not depend on the choice of a generating word list.  Their internal subalgebras are contained in the fixed ambient commutant.  On the regular carrier, Theorem~\ref{sel1:thm:module} identifies precisely the additional right-module constraints required of word-induced operations.  Dropping those constraints enlarges the allowed endomorphism algebra, as the explicit example in Section~\ref{sel1:root} demonstrates.
\end{proof}

\begin{remark}[Application to the V106 word-depth census]
The V106 census remains correct when $\mathcal V_r$ is the actual reachable--observable physical carrier carrying the stated external action.  When it is represented instead by a vector space of algebra words or by a regular trace quotient, its commutant cannot be identified with executable internal operations without the action/module comparison above.  This note does not silently reinterpret all earlier uses of $\mathcal V_r$ in either way: the representation map must be specified for the particular source being audited.
\end{remark}

\section{A twelve-root test with a fully noncommutative history algebra}
\label{sel1:root}

Let $\Omega=\{(i,j):1\le i,j\le4,\ i\ne j\}$ and $\Hh=\C^\Omega$.  The simultaneous relabeling representation is
$u_g|i,j\rangle=|g(i),g(j)\rangle$, $g\in S_4$.
Consider a benchmark operation bank containing all diagonal point projections and the two three-point fibre expectations
\begin{equation}
 (T f)(i,j)=\frac13\sum_{k\ne i}f(i,k),\qquad
 (H f)(i,j)=\frac13\sum_{k\ne j}f(k,j).
 \label{sel1:eq:root-resets}
\end{equation}
The reset definitions are the established root operations; adjoining the complete diagonal profile bank specifies the benchmark.  Its full admission is not inferred for an unprovided historical alphabet.

\begin{theorem}[Large invariant word multiplicities without internal colour]
\label{sel1:thm:root-benchmark}
For this benchmark the system operation algebra is $\A=M_{12}(\C)$, while
\begin{equation}
 \B_{\rm hist}=\A^{S_4}
 \cong\C\oplus M_2(\C)\oplus\C\oplus\C.
 \label{sel1:eq:root-system-algebra}
\end{equation}
In the order $(\mathbf1,\operatorname{sgn},W,W\otimes\operatorname{sgn},V_2)$, the system multiplicities and the regular-word multiplicities are respectively
\begin{equation}
 (1,0,2,1,1),\qquad (7,5,19,17,12).
 \label{sel1:eq:root-multiplicities}
\end{equation}
Thus the full commutant of external conjugation on $L^2(M_{12})$ has complex dimension
\begin{equation}
 7^2+5^2+19^2+17^2+12^2=868
 \label{sel1:eq:868}
\end{equation}
and contains an $M_7(\C)$ block on invariant word vectors.  The word-induced invariant system algebra nevertheless has dimension only seven and no $M_3(\C)$ corner.
\end{theorem}
\begin{proof}
The support graph of $T+H$ joins pairs with equal tail or equal head.  It is connected.  Indeed, to connect $(i,j)$ to $(k,l)$, choose $q$ distinct from $i,k$ and use
$(i,j)\to(i,q)\to(k,q)\to(k,l)$, omitting repeated vertices when necessary.  Each transition is an allowed support edge.  The existing diagonal-plus-router theorem therefore gives $C^*(\text{diagonal projections},T,H)=M_{12}$.

For the conjugacy classes $1,(12),(12)(34),(123),(1234)$, the class sizes are $1,6,3,8,6$.  The permutation character on ordered unequal pairs is the number $f(g)(f(g)-1)$, where $f(g)$ is the number of fixed vertices, and is therefore $(12,2,0,0,0)$.  The five irreducible characters in the stated order are
\[
\begin{array}{c|rrrrr}
 &1&(12)&(12)(34)&(123)&(1234)\\\hline
\mathbf1&1&1&1&1&1\\
\operatorname{sgn}&1&-1&1&1&-1\\
W&3&1&-1&0&-1\\
W\otimes\operatorname{sgn}&3&-1&-1&0&1\\
V_2&2&0&2&-1&0
\end{array}
\]
Here $W$ is the vertex permutation representation minus its invariant line; $V_2$ is the permutation representation on the three pair partitions minus its invariant line.  Character inner products give $(1,0,2,1,1)$, hence \eqref{sel1:eq:root-system-algebra} by the finite commutant census.

Conjugation on $M_{12}\cong\Hh\otimes\overline\Hh$ has character $|\chi_\Hh|^2=(144,4,0,0,0)$.  Its character inner products give $(7,5,19,17,12)$, and its commutant dimension is the sum of their squares.  In contrast, Theorem~\ref{sel1:thm:module} keeps only $L(\A^{S_4})$.  Each simple factor of the latter has size at most two.  A nonzero representation of $M_3$ in a matrix factor of size at most two is impossible, even allowing a supported nonunital corner, so no such colour corner exists.
\end{proof}

\begin{remark}[What the example rules out]
This is stronger than the already known profile-algebra obstruction: $\A$ here is the full noncommutative matrix algebra.  The false inference is from a large invariant \emph{word-space} endomorphism algebra to a new internal \emph{system} factor.  The example neither changes the external-versus-internal census nor rules out colour on an actual larger source-minimal system carrier.  It requires that any claimed enlargement be accompanied by that carrier and its action map.
\end{remark}

\section{A reciprocal-return mechanism that forces the colour bridge}
\label{sel1:returns}

\subsection{The source-native route class}
The imported colour theorem needs a faithful shorted line $T\cong\C$, a private plane $H\cong\C^2$, the supported private algebra $M_2(H)$, and one nonzero word from $T$ to $H$ or back.  We retain the first three inputs.  Instead of assuming the last word, we derive it from a specified family of reciprocal routes.

Let $G$ act trivially on $T$ and $H$.  At a mediator cut let
\[
 \E=V\otimes\Nn,\qquad \rho(g)=v_g\otimes I_{\Nn},
 \qquad d=\dim V,\quad n=\dim\Nn\ge1,
\]
where $v$ is an irreducible unitary representation.  The same mediator cut, common identifications, and common group labeling are part of the supplied route table.  Assume the actual entry and reverse-exit families have the evaluated form
\begin{equation}
 A_g=(v_g\otimes I)A:T\to\E,
 \qquad B_h^*=B^*(v_h^*\otimes I):\E\to H,
 \label{sel1:eq:ports}
\end{equation}
for fixed maps $A:T\to\E$ and $B:H\to\E$.  The labels $g,h$ in this equation denote admitted operations, not an assumed ability to synthesize an arbitrary group average.  The composable products $B_h^*A_g$ must belong to the declared grammar.  A physical reverse realizing $B_h^*$ is a hypothesis; a source vector $B$ without a reverse operation is not sufficient.

Let the actual mediator operation be $D=I_V\otimes h$ with $h=h^*\in\End(\Nn)$; the identity word $D^0$ is included.  The finite returned words relevant below are
\begin{equation}
 r_{g,k}=B^*(v_g\otimes h^k)A:T\to H,
 \qquad g\in G,\quad 0\le k\le n-1.
 \label{sel1:eq:return-words}
\end{equation}
They are the products $B_e^*D^kA_g$.  General $B_h^*D^kA_g$ gives the same list with the relative label $h^{-1}g$.  Each returned operator, extended by zero on other final-cut sectors, commutes with the external action because its source and target are externally trivial.

Define the multiplicity source densities, using unnormalized partial traces,
\begin{equation}
 \rho_A=\Tr_V(AA^*),\qquad \rho_B=\Tr_V(BB^*)\succeq0.
 \label{sel1:eq:densities}
\end{equation}
They measure the two ports in the \emph{same} mediator multiplicity frame.

\begin{theorem}[Exact positive return mass]
\label{sel1:thm:return-mass}
For any $h=h^*$ and every $k\ge0$,
\begin{equation}
 \frac1{|G|}\sum_g\norm{r_{g,k}}_{\HS}^2
 =\frac1d\Tr(\rho_Bh^k\rho_Ah^k).
 \label{sel1:eq:return-mass-k}
\end{equation}
Consequently, for the positive multiplicity reachability Gram
\begin{equation}
 W_A=\sum_{k=0}^{n-1}h^k\rho_Ah^k,
 \qquad
 \mathfrak m_{\rm ret}:=\frac1{|G|}\sum_{g,k<n}\norm{r_{g,k}}_{\HS}^2,
 \label{sel1:eq:WA}
\end{equation}
one has
\begin{equation}
 \boxed{\mathfrak m_{\rm ret}=\frac1d\Tr(\rho_BW_A)\ge0.}
 \label{sel1:eq:positive-return}
\end{equation}
This mass is positive if and only if at least one of the admitted returned words is a nonzero colour corner.
\end{theorem}
\begin{proof}
Finite group averaging on the irreducible factor gives
\[
 \frac1{|G|}\sum_g(v_g\otimes I)AA^*(v_g^*\otimes I)
 =\frac{I_V}{d}\otimes\rho_A.
\]
Indeed the average commutes with every $v_g$; Schur's lemma determines its first factor and partial trace determines the second.  Since $D$ commutes with the group action,
\begin{align*}
 \frac1{|G|}\sum_g\norm{B^*(v_g\otimes h^k)A}_{\HS}^2
 &=\Tr_{\E}\!\left(BB^*\left[\frac{I_V}{d}\otimes h^k\rho_Ah^k\right]\right)\\
 &=\frac1d\Tr_{\Nn}(\rho_Bh^k\rho_Ah^k).
\end{align*}
The right side is nonnegative, also directly as a squared Hilbert--Schmidt norm.  Summing proves \eqref{sel1:eq:positive-return}.  A finite sum of nonnegative squared norms vanishes exactly when every summand vanishes.
\end{proof}

\begin{corollary}[A shared irreducible mediator removes the separate bridge assumption]
\label{sel1:cor:one-copy}
If the mediator has multiplicity one, so $\Nn=\C$, then
\begin{equation}
 \boxed{\mathfrak m_{\rm ret}=
 \frac{\norm A_{\HS}^2\norm B_{\HS}^2}{d}.}
 \label{sel1:eq:one-copy-mass}
\end{equation}
Thus two nonzero reciprocal operational ports through the same irreducible mediator force a nonzero two-leg return word.  On the already faithful $T\oplus H$, that word and the existing private algebra generate $M_3(\C)$ by the imported colour-bridge theorem.  No independent direct bridge operation is needed.
\end{corollary}
\begin{proof}
For $n=1$ the panel contains only $k=0$, and the two partial traces are the scalar squared Hilbert--Schmidt norms of the ports.  Positivity and the finite witness conclusion follow from Theorem~\ref{sel1:thm:return-mass}.  The last step is exactly the existing one-bridge criterion, not an additional matrix-generation argument.
\end{proof}

\begin{remark}[What has been reduced]
The sufficient source condition is now reciprocal access to one \emph{common irreducible mediator}, together with nonzero ports.  This is a check on upstream route structure, not a restatement that a colour corner is nonzero.  It still has content: separate ports, identical marginal dimensions, or an unrecorded choice of their relative frame do not satisfy the hypothesis.  Nor does it imply that the given historical Store already has those ports.
\end{remark}

\subsection{Exact finite-horizon exclusion, including multiple excursions}

Define
\begin{equation}
 \Nn_A=\Span\{h^k\Ran\rho_A:0\le k<n\},
 \qquad \Nn_B=\Span\{h^k\Ran\rho_B:0\le k<n\}.
 \label{sel1:eq:cyclic}
\end{equation}
Cayley--Hamilton makes these spaces invariant under $h$, and self-adjointness makes them reducing.  They need not intersect when they are nonorthogonal; those two notions are not identified.

\begin{theorem}[Necessary-and-sufficient return criterion for the declared route bank]
\label{sel1:thm:complete-return}
Consider the algebra generated on $T\oplus H\oplus\E$ by the three cut projections, the supported private $M_2(H)$, the mediator operator $D$, and all route operators $A_g,B_h$ and their adjoints.  Then the following are equivalent:
\begin{enumerate}[label=\textup{(\roman*)},leftmargin=2.6em]
\item some word in this bank has a nonzero $H\leftarrow T$ corner;
\item some $r_{g,k}$ in \eqref{sel1:eq:return-words} is nonzero;
\item $\Tr(\rho_BW_A)>0$;
\item $\Nn_A$ and $\Nn_B$ are not orthogonal.
\end{enumerate}
If the conditions hold, a returned witness has at most $n+1$ letters, counting entry, mediator steps, and exit.  If they fail, a reducing projection separates $T$ from $H$ for the entire route bank, so arbitrary repeated excursions cannot create a bridge within that bank.
\end{theorem}
\begin{proof}
The equivalence of (ii) and (iii), and (ii)$\Rightarrow$(i), follow from Theorem~\ref{sel1:thm:return-mass}.  The range of $W_A$ is $\Nn_A$, since its kernel consists of vectors orthogonal to every $h^k\Ran\rho_A$.  Hence
\[
 \Tr(\rho_BW_A)=0\quad\Longleftrightarrow\quad
 \Ran\rho_B\perp\Nn_A.
\]
As $\Nn_A$ reduces $h$, this is equivalent to $\Nn_B\perp\Nn_A$, proving (iii)$\Leftrightarrow$(iv).

It remains to exclude longer words when the mass is zero.  Partial-trace positivity implies $\Ran A\subseteq V\otimes\Ran\rho_A$ and $\Ran B\subseteq V\otimes\Ran\rho_B$.  To see this, a vector $z\in\Ker\rho_A$ satisfies
$\sum_i\norm{A^*(e_i\otimes z)}^2=0$, so every $e_i\otimes z$ is orthogonal to $\Ran A$; the proof for $B$ is identical.  Therefore each $A_g$ connects $T$ only to $V\otimes\Nn_A$.  When the mass vanishes, every $B_h$ connects $H$ only to $V\otimes\Nn_A^\perp$.  The projection
\[
 Q=I_T\oplus0_H\oplus(I_V\otimes P_{\Nn_A})
\]
commutes with every generator, including $D$ and all adjoints.  It thus commutes with every word.  Since $Qp_T=p_T$ and $Qp_H=0$, every $p_Hwp_T$ vanishes.  This proves (i)$\Rightarrow$(iii) by contraposition and the all-excursion assertion.  The finite panel uses $k\le n-1$, hence at most $k+2\le n+1$ letters.
\end{proof}

\begin{corollary}[Port-relative minimality forces a bridge quantitatively]
\label{sel1:cor:reach-force}
If the actual type-port orbit and mediator dynamics are cyclic on the retained multiplicity space, namely $\Nn_A=\Nn$, and $B\ne0$, then the colour bridge is forced.  More precisely, with $\kappa=\lambda_{\min}(W_A)>0$,
\begin{equation}
 \mathfrak m_{\rm ret}\ge\frac{\kappa}{d}\norm B_{\HS}^2,
 \qquad
 \max_{g,k<n}\norm{r_{g,k}}_{\HS}^2
 \ge\frac{\kappa}{dn}\norm B_{\HS}^2.
 \label{sel1:eq:reach-margin}
\end{equation}
\end{corollary}
\begin{proof}
$W_A\succeq\kappa I$ and $\Tr\rho_B=\norm B_{\HS}^2$ give the first inequality.  The maximum of $|G|n$ nonnegative summands is at least their average, giving the second.  Positive mass supplies a word by the preceding theorem.
\end{proof}

\begin{remark}[Port-relative, not joint, source minimality]
The condition $\Nn_A=\Nn$ is reachability from the type port alone, using the admitted orbit and mediator dynamics.  Reachability from the union of type and private ports is weaker and does not force a bridge.  This distinction is not dropped when describing the result as assumption reduction.  The displayed lower bound is a route-word amplitude margin; it is not automatically the full calibrated Howe gap.
\end{remark}

\subsection{A spectral form of the test}

\begin{proposition}[Finite pole separation and the source-overlap panel]
\label{sel1:prop:poles}
Let $h=\sum_{\ell=1}^s\lambda_\ell P_\ell$ with distinct real $\lambda_\ell$, $s\le n$.  The return condition is equivalent to
\begin{equation}
 \exists\ell:\quad
 \Tr(\rho_BP_\ell\rho_AP_\ell)>0.
 \label{sel1:eq:pole-test}
\end{equation}
It suffices to use $0\le k<s$.  If $V_{k\ell}=\lambda_\ell^k$, $0\le k<s$, and
$\sigma_-,\sigma_+$ are its least and greatest singular values, then
\begin{equation}
 \frac{\sigma_-^2}{d}\sum_\ell\Tr(\rho_BP_\ell\rho_AP_\ell)
 \le\frac1{|G|}\sum_{g,k<s}\norm{r_{g,k}}_{\HS}^2
 \le\frac{\sigma_+^2}{d}\sum_\ell\Tr(\rho_BP_\ell\rho_AP_\ell).
 \label{sel1:eq:vandermonde}
\end{equation}
\end{proposition}
\begin{proof}
For fixed $g$ write $R_{g\ell}=B^*(v_g\otimes P_\ell)A$.  Then $r_{g,k}=\sum_\ell V_{k\ell}R_{g\ell}$.  The Vandermonde matrix is invertible, so all these moments vanish exactly when all the residues vanish.  Its singular-value bounds apply to the stacked matrix-valued residues as well.  Schur averaging, with $P_\ell$ on both sides of $\rho_A$, gives
\[
 \frac1{|G|}\sum_g\norm{R_{g\ell}}_{\HS}^2
 =\frac1d\Tr(\rho_BP_\ell\rho_AP_\ell).
\]
Sum over $\ell$ to prove the bounds and the criterion.  Distinct poles may be close; no regulator-uniform lower bound on $\sigma_-$ is inferred without a separate separation estimate.
\end{proof}

\begin{proposition}[Direct error interval for a reconstructed return panel]
\label{sel1:prop:panel-error}
Suppose the reconstructed return operators have a certified error
\[
 \frac1{|G|}\sum_{g,k<n}\norm{\widehat r_{g,k}-r_{g,k}}_{\HS}^2\le\epsilon^2,
 \qquad
 \widehat{\mathfrak m}=\frac1{|G|}\sum_{g,k<n}\norm{\widehat r_{g,k}}_{\HS}^2.
\]
Then
\begin{equation}
 \bigl(\sqrt{\widehat{\mathfrak m}}-\epsilon\bigr)_+^2
 \le\mathfrak m_{\rm ret}\le
 \bigl(\sqrt{\widehat{\mathfrak m}}+\epsilon\bigr)^2.
 \label{sel1:eq:return-error}
\end{equation}
In particular, $\sqrt{\widehat{\mathfrak m}}>\epsilon$ certifies an actual nonzero returned colour word.  A failed positive test with nonzero error does not certify exact absence.
\end{proposition}
\begin{proof}
Apply the triangle and reverse triangle inequalities in the Hilbert direct sum of the returned Hom spaces with weights $1/|G|$.  Its true and measured norms are the two square roots in \eqref{sel1:eq:return-error}.
\end{proof}

\section{Exact examples and the boundary of the return mechanism}
\label{sel1:examples}

\subsection{Tetrahedral return forced despite a vanishing unrotated product}

\begin{example}[An exact $S_4$ mediator calculation]
\label{sel1:ex:S4-return}
Let $V=W_0=\{x\in\C^4:\sum_i x_i=0\}$ with the usual $S_4$ action and $\Nn=\C$.  Let $T=\C$, $H=\C^2$, and set, in the ambient vertex coordinates,
\[
 A=\begin{pmatrix}1\\-1\\0\\0\end{pmatrix},\qquad
 B=\begin{pmatrix}1&1\\1&1\\-2&1\\0&-3\end{pmatrix}.
\]
Both maps take values in $W_0$, and $B^*A=0$.  Nevertheless for the transposition $g=(23)$,
\[
 B^*v_gA=\begin{pmatrix}3\\0\end{pmatrix}\ne0.
\]
The complete positive return calculation is
\begin{equation}
 \frac1{24}\sum_{g\in S_4}\norm{B^*v_gA}^2
 =\frac{2\cdot18}{3}=12,
 \qquad\sum_g\norm{B^*v_gA}^2=288.
 \label{sel1:eq:tetrahedral-return}
\end{equation}
Thus inspecting only the unrotated entry/exit pair would miss the bridge.  An admitted complete orbit makes its occurrence unavoidable.  This is a symbolic benchmark using displayed port matrices, not a claim that those matrices were extracted from the historical Store.
\end{example}

\begin{proposition}[The benchmark admits normalized finite instruments]
\label{sel1:prop:normalization}
The displayed entry and exit families can be implemented as finite completely positive instrument branches without changing which returned products are nonzero.  In particular, let every entry Kraus operator be $A_g/8$.  Their summed effect is $3I_T/4$ and a no-entry branch completes it to $I_T$.  Let every exit Kraus operator be $B_g^*/12$.  Their summed effect on $W_0$ is $I_{W_0}$.  Every nonzero returned product is scaled by $1/96$ and remains nonzero.
\end{proposition}
\begin{proof}
Each $A_g^*A_g=2I_T$, so the sum over 24 labels after division by $8^2$ is $48/64=3/4$.  Schur averaging on the irreducible $W_0$ gives
$\sum_g B_gB_g^*=(24\norm B_{\HS}^2/3)I_{W_0}=144I_{W_0}$.
The positive residual entry effect is supplied by its positive square root as a no-entry Kraus operator; the exit list is already normalized.  Composing any successful entry and exit gives the claimed scalar rescaling.  This is an explicit active instrument realization, not evidence of historical admission.
\end{proof}

\subsection{All ports can occur while the bridge remains absent}

\begin{example}[Jointly minimal but mutually dark mediator copies]
\label{sel1:ex:dark}
Let $V$ be any irreducible unitary $G$-module, take $\Nn=\C^2$ and $h=\diag(0,1)$.  Choose nonzero $a,b\in V$ and a nonzero functional $\ell:H\to\C$.  Define
\[
 A(t)=t\,a\otimes e_1,\qquad B(x)=\ell(x)b\otimes e_2.
\]
Both operational port orbits are nonzero.  Their combined $G$-orbit span is all of $V\otimes\C^2$, so the mediator is source minimal relative to the \emph{union} of the two port queries.  But
\[
 \rho_A=\norm a^2|e_1\rangle\langle e_1|,\qquad
 \rho_B=\norm b^2\norm\ell^2|e_2\rangle\langle e_2|,
 \qquad \Tr(\rho_BW_A)=0.
\]
No return word in the declared route bank bridges the colour components.  Nonzero ports, common external irrep name, complete covariance, and joint source minimality do not remove the need for a common multiplicity-overlap or reachability condition.
\end{example}

\begin{example}[One occurring mediator transition opens the return channel]
\label{sel1:ex:dynamic}
In the preceding example replace $h$ by
$h(t)=\begin{psmallmatrix}0&t\\t&0\end{psmallmatrix}$, $t\in\R$, keeping both ports fixed.  Write $a_0=\norm a^2$ and $b_0=\norm b^2\norm\ell^2$.  The fixed two-layer panel gives
\[
 W_A(t)=a_0\diag(1,t^2),\qquad
 \mathfrak m_{\rm ret}(t)=\frac{a_0b_0}{d}t^2.
\]
At $t=0$ the bank is dark.  At every $t\ne0$ a three-letter entry/mediator/exit word gives colour, with an explicitly vanishing raw margin at the boundary.  A nonzero entry of the \emph{actual} mediator dynamics, not a change of external irrep dimensions, creates the needed route.  This does not claim generic dynamics will select $t\ne0$ in an unspecified microscopic family.
\end{example}

\section{Admissible invariant extraction and central separation}
\label{sel1:assembly}

\subsection{A colour corner does not require an extra twirl control}

\begin{proposition}[Support-native extraction of an internal return]
\label{sel1:prop:corner}
Let $\A\subseteq\End(\Hh)$ be the actual word algebra, and suppose $p_T,p_H\in\A$ project onto externally trivial subspaces.  For every $a\in\A$,
\begin{equation}
 p_Hap_T\in\A\cap u(G)',
 \qquad p_HE_G(a)p_T=p_Hap_T,
 \label{sel1:eq:corner-twirl}
\end{equation}
where $E_G$ denotes averaging in $\End(\Hh)$.  The first conclusion does not require $u_g\A u_g^*=\A$, and no separate physical implementation of $E_G$ is needed to extract this corner in the generated algebra.
\end{proposition}
\begin{proof}
The corner belongs to $\A$ because its three factors do.  The group acts as identity on its input and output supports and the corner vanishes elsewhere, so it commutes with $u_g$.  Compression of each summand $u_gau_g^*$ by $p_H,p_T$ gives the same $p_Hap_T$, proving the second identity.
\end{proof}

The support projections are retained inputs here, not arbitrary kinematic projectors.  A source construction proving that a projection exists as a subspace does not by itself show that the projection belongs to the executed algebra.

\begin{proposition}[When an external twirl stays inside the historical algebra]
\label{sel1:prop:twirl}
If $u_g\A u_g^*=\A$, then $E_G(\A)=\A^G$.  If an evaluated word basis spans $\A$, its twirls linearly span $\A^G$.  Neither twirling only the primitive generators nor twirling an algebra not stable under the external action has this guarantee.
\end{proposition}
\begin{proof}
Under invariance every summand of $E_G(a)$ lies in $\A$, the average is invariant, and every invariant element is fixed by the average.  Apply the linear map to a basis to obtain the spanning statement.

For the first limitation let $X,Z$ be Pauli matrices on $\C^2$, let $u=Z\otimes I$, and take actual Hermitian generators
\[
 a_0=X\otimes I,\qquad a_1=X\otimes X,\qquad a_2=X\otimes Z.
\]
They generate $\A=C^*(X)\otimes M_2(\C)$, since $a_0a_1=I\otimes X$ and $a_0a_2=I\otimes Z$.  The algebra is invariant under conjugation by $u$.  Every primitive generator has zero twirl, but $\A^{\langle u\rangle}=I\otimes M_2(\C)$, supplied by the even products.  Twirling before word composition loses these operations.

For the second limitation take the diagonal algebra on $\C^2$ and the order-two unitary $F=2^{-1/2}\begin{psmallmatrix}1&1\\1&-1\end{psmallmatrix}$.  The average of $Z$ is $(Z+X)/2$, which is not diagonal.  Thus averaging in the ambient endomorphism algebra can introduce an operator absent from the original algebra when normalization by the group is not established.
\end{proof}

These are not new proofs of the external-twirl formula.  They specify the source-algebra and word-composition conditions under which that existing formula is usable for an occurrence argument.

\subsection{Blockwise fullness versus fullness of the assembled algebra}

\begin{theorem}[Subdirect matrix-algebra classification]
\label{sel1:thm:subdirect}
Let $\B\subseteq\bigoplus_{i=1}^r M_{n_i}(\C)$ be a unital $C^*$-subalgebra whose projection onto each displayed factor is surjective.  There is a partition of $\{1,\ldots,r\}$ into classes, each containing only equal matrix sizes, such that $\B$ consists of one independent full matrix per class, embedded diagonally across that class up to fixed inner automorphisms.  In particular, correlations between different simple factors are possible only between isomorphic factors.
\end{theorem}
\begin{proof}
Write the intrinsic finite-dimensional decomposition of $\B$ as $\bigoplus_\beta M_{m_\beta}(\C)$.  A surjective $*$-homomorphism from this direct sum to the simple algebra $M_{n_i}$ annihilates every summand except one, and on that remaining summand is an isomorphism.  To justify uniqueness, the images of the central units of the summands are central orthogonal projections in $M_{n_i}$ summing to identity; exactly one can be nonzero.  Its kernel is an ideal in a simple algebra and the map is nonzero, so the kernel vanishes and $m_\beta=n_i$.

Group the coordinates $i$ according to which intrinsic summand they see.  Every intrinsic summand must be seen somewhere because the inclusion of $\B$ into the product is faithful.  Within a class the coordinate isomorphisms are $*$-automorphisms of a full complex matrix algebra, hence inner.  The innerness can be verified by transporting a matrix-unit system and choosing an orthonormal basis of the transported rank-one projections.  Coordinates in different classes depend on independent intrinsic summands.  This proves the asserted form.
\end{proof}

\begin{corollary}[Exact fourteen-or-fifteen alternative on the proposed internal carrier]
\label{sel1:cor:14-15}
Suppose the actually generated internal algebra is a unital subalgebra
\[
 \B\subseteq M_3(\C)\oplus M_2(\C)\oplus\C\oplus\C
\]
with surjective projections onto each of the four displayed blocks.  Then exactly one of the following holds:
\begin{align}
 \B&=M_3(\C)\oplus M_2(\C)\oplus\C\oplus\C,
 &\dim_\C\B&=15;\label{sel1:eq:15}\\
 \B&=\{(A,B,z,z):A\in M_3,\ B\in M_2,\ z\in\C\},
 &\dim_\C\B&=14.\label{sel1:eq:14}
\end{align}
The full case holds if and only if the two scalar characters $\chi_1,\chi_2$ of $\B$ differ.  For any generating bank $(b_j)$, it is equivalent to the finite condition
\begin{equation}
 \boxed{\eta_{\rm cen}:=\sum_j|\chi_1(b_j)-\chi_2(b_j)|^2>0.}
 \label{sel1:eq:central-separator}
\end{equation}
\end{corollary}
\begin{proof}
The nonisomorphic $M_3$ and $M_2$ coordinates cannot be identified with each other or with a scalar coordinate by Theorem~\ref{sel1:thm:subdirect}.  Only the two scalar coordinates can be identified, and a complex scalar $*$-automorphism is the identity.  This gives exactly the two possibilities.  Characters agreeing on every generator agree on adjoints, products and linear combinations, hence on the generated algebra.  Thus they differ if and only if one of the summands in \eqref{sel1:eq:central-separator} is positive.
\end{proof}

\begin{remark}[Relation to the existing $9/11/13/15$ deficit table]
That table applies to independently supported block operations.  Its independent scalar blocks are justified when the corresponding central projections occur in the source algebra.  The present corollary does not replace that table; it identifies what must be checked if one has established only four surjective compressed actions and has not established independent central supports.  Before colour closure, two isomorphic $M_2$ blocks can also be diagonally correlated, so componentwise dimensions cannot simply be added in an arbitrary source family.
\end{remark}

\begin{proposition}[When the scalar-separation check is already discharged]
\label{sel1:prop:central-controls}
If the external representation matrices $u_g$ themselves belong to the actual system word algebra $\A$, then every external isotypic central projector belongs to $\A\cap u(G)'$.  In particular, two inequivalent external scalar-multiplicity sectors are independently separated; the alternative \eqref{sel1:eq:14} is impossible once both are present.
\end{proposition}
\begin{proof}
For an irreducible character $\chi_\lambda$, the standard central projector is
\[
 p_\lambda=\frac{d_\lambda}{|G|}\sum_g\overline{\chi_\lambda(g)}u_g.
\]
It belongs to $\A$ by linear closure and commutes with every $u_g$ by character centrality.  Its action is identity on the $\lambda$-isotypic sector and zero on inequivalent sectors.  Thus it separates their scalar characters.
\end{proof}

The proposition is an assumption-reuse result: an independently established physical external-control bank already supplies this part of global internal saturation.  An external symmetry used only as an abstract relabeling need not supply those projectors as historical operations.

\section{The combined source-level internal landing theorem}
\label{sel1:landing}

We now state exactly what the new return mechanism adds to the inherited structural branch.  It does not postulate the desired commutant as an equality.

\begin{theorem}[Internal landing from reciprocal source reachability]
\label{sel1:thm:landing}
On one exposed final-cut carrier suppose the verified external $S_4$ decomposition is
\begin{equation}
 \Hh_0=(T\oplus H)\oplus(W\otimes M)\oplus V_2\oplus W^{\rm sgn},
 \qquad \dim T=1,\quad\dim H=\dim M=2,
 \label{sel1:eq:final-carrier}
\end{equation}
where $T,H$ are externally trivial and $W,V_2,W^{\rm sgn}$ are inequivalent irreducibles of dimensions $3,2,3$.  Assume the following independently checked source inputs.
\begin{enumerate}[label=\textup{(P\arabic*)},leftmargin=3em]
\item The type line survives the full source short, and $T\oplus H$ is the corresponding faithful common carrier.  Its support projections and supported private $M_2(H)$ are in the actual word algebra.
\item The already constructed weak packet supplies $I_W\otimes M_2(M)$ on the indicated multiplicity block.
\item Through an actual common mediator, the reciprocal route bank of Section~\ref{sel1:returns} is admitted, with $B\ne0$ and $\Nn_A=\Nn$.
\item The two remaining scalar characters are separated by one actual internal word, or by an external-control bank as in Proposition~\ref{sel1:prop:central-controls}.
\end{enumerate}
Let $\A_0$ be the actual final-cut word algebra including returned paths, and put $\B_0=\A_0\cap u(S_4)'$.  Then
\begin{equation}
 \boxed{\B_0=u(S_4)'\cong M_3(\C)\oplus M_2(\C)\oplus\C\oplus\C.}
 \label{sel1:eq:landing}
\end{equation}
The colour bridge is a derived word with at most $n+1$ letters, not an independent primitive.  A returned word has the quantitative lower bound in \eqref{sel1:eq:reach-margin}.  In the one-copy mediator case, (P3) is replaced simply by the two nonzero reciprocal ports, and a two-leg word suffices.
\end{theorem}
\begin{proof}
The external decomposition in \eqref{sel1:eq:final-carrier} fixes its kinematic commutant to the product algebra on the right of \eqref{sel1:eq:landing}; hence $\B_0$ is contained in that algebra.  Corollary~\ref{sel1:cor:reach-force} produces a nonzero admitted $T\to H$ return.  It is internally invariant and supported on the colour carrier, so belongs to $\B_0$.  Together with (P1), the imported metric-covariant bridge theorem generates the full colour $M_3$.  Assumption (P2) already gives the weak $M_2$ and the unital algebra projects onto each scalar coordinate.  Corollary~\ref{sel1:cor:14-15} and (P4) now give equality with the full product.  The word length and mass bounds come from the return theorem.  Corollary~\ref{sel1:cor:one-copy} gives the last claim.
\end{proof}

\begin{remark}[What this theorem has not assumed and has not proved]
No independent direct colour bridge, nonzero type--private source angle, second weak router, or new persistent final-cut dimension is assumed.  The mediator is a typed intermediate cut, not a secretly adjoined final-cut colour line.  Its existence and reciprocal ports remain operational inputs.  The theorem does assume the checked carrier \eqref{sel1:eq:final-carrier}; it does not prove its universal selection.  It neither forces the alternating determinant line nor selects the matter tensor types.  Those remain the separate source tests already isolated in V110 and in the duality manuscript.
\end{remark}

\begin{corollary}[Exact alternative on a complete declared mediator bank]
\label{sel1:cor:alternative}
With (P1), (P2), and (P4) held fixed, replace (P3) by the complete declared single-mediator bank of Theorem~\ref{sel1:thm:complete-return}, with no other routes coupling $T$ and $H$.  A positive value of $\Tr(\rho_BW_A)$ supplies the historical colour bridge.  A zero value gives the explicit reducing projection in that theorem and excludes every longer route in the declared bank.  Missing port or mediator entries give neither conclusion: they are uninstantiated source inputs, not zero values.
\end{corollary}
\begin{proof}
The positive branch is Theorem~\ref{sel1:thm:complete-return} followed by the inherited bridge theorem.  The zero branch is its reducing-subspace construction.  The final statement is the logical distinction between an evaluated zero and an absent entry.
\end{proof}

\section{Finite source audit and the V2 frontier}
\label{sel1:status}

\subsection{A terminating audit once the actual route table is supplied}
The following computation does not fit operators to a target gauge group.  At fixed cutoff it takes their actually evaluated matrices and their action/cut maps as inputs.

First identify the exposed system representation and the history multiplication representation.  If using a regular word model, retain the module test \eqref{sel1:eq:module-intersection}.  Compute the internal algebra by simultaneous linear equations inside the actual word algebra, not by declaring the whole ambient multiplicity commutant to be executed.  The existing flatness theorem supplies a finite algebra basis when the complete source-word table is known.

Next retain the established weak reset packet rather than charging it again.  Compute the type-line Schur floor in the common source metric, using the inherited V109 formula.  Its positivity decides the faithful line, not a nonzero overlap angle.

For every independently identified common mediator, extract its actual $G$ action, reciprocal entry/exit maps, and mediator operation.  Test whether it is of the class $V\otimes\Nn$ with $D=I\otimes h$ used here.  Only then form $\rho_A,\rho_B,W_A$ and evaluate \eqref{sel1:eq:positive-return}.  The panel is finite, with at most $|G|n$ displayed returned operators.  A positive result returns a literal labelled witness; a zero result excludes longer paths only when the declared bank is complete.  A different or incomplete mediator grammar requires its own reachability analysis, not an unsupported application of the single-mediator theorem.

Finally test global central separation and compute the actual centre and simple-block dimensions.  Only after these checks should the determinant/Clifford/Dirac incidence tests be applied on the same source.  Their positivity and same-history provenance are not implied by the internal matrix-algebra test.

\subsection{Source availability and what has actually been evaluated}
The supplied monographs provide general source-reconstruction mechanisms, source-short formulas, structural carriers and conditional finite-Dirac constructions.  The audited V106--V110 continuation further supplies the repaired weak resets, the positive active type-line short, and exact historical alternatives.  In the material inspected for this note, no common-frame numerical or symbolic \emph{historical} reciprocal-port panel $(A,B,h)$ has been identified for which the present colour-return mass can be evaluated as a Store output.  The matrices in Section~\ref{sel1:examples} are explicitly constructed benchmarks, not substitutes for that missing historical panel.

Accordingly the gain is a new source-level sufficient theorem and an exact exclusion test for a nontrivial route class.  It is not a declaration that historical colour or unconditional Standard-Model selection is now closed.

\begin{center}
\small
\begin{tabular}{L{.35\textwidth}L{.57\textwidth}}
\toprule
Question & Status in this V1\\
\midrule
Repeat weak-router or private-plane generation? & No. V108--V109 are retained as audited inputs.\\[.3em]
When does an invariant word-space endomorphism come from the old system algebra? & Exactly when it satisfies the right-module compatibility test; Theorem~\ref{sel1:thm:module}.\\[.3em]
Can large noncommutative word ranks falsely suggest colour capacity? & Yes. The exact twelve-root benchmark has a formal invariant $M_7$ word block but no physical internal $M_3$; Theorem~\ref{sel1:thm:root-benchmark}.\\[.3em]
Can a bridge follow from upstream operations rather than being added? & Yes. Shared irreducible reciprocal access forces a two-leg return. Port-relative cyclicity gives the multiplicity-valued extension; Section~\ref{sel1:returns}.\\[.3em]
Can repeated excursions overcome a zero complete return panel? & Not in the declared self-adjoint single-mediator route class; an explicit reducing projection excludes them.\\[.3em]
Does local block fullness imply the full dimension fifteen? & Only after separating repeated scalar characters, or verifying their central projectors; Corollary~\ref{sel1:cor:14-15}.\\[.3em]
Historical Store evaluation of the new return test & Not completed: the actual common-frame reciprocal-port panel has not been identified in this audit.\\[.3em]
Determinant inevitability, matter-type uniqueness, and microscopic attraction & Not claimed; these are distinct from the proved internal return criterion.\\
\bottomrule
\end{tabular}
\end{center}

\subsection{The next focused mathematical step}
V2 should first seek an actual typed pair of reciprocal routes in the pre-release Store, keeping the already constructed weak packet fixed.  The immediate test is whether their common mediator is source minimal from the type port alone.  When it is, the bridge follows without a separate occurrence assumption.  When it is not, the returned reducing-space defect specifies exactly which mediator multiplicity direction is dark.

If the available mediator dynamics has several noncommuting controls, the genuinely new extension is a \emph{word-valued} version of the return criterion with the same source and module bookkeeping.  That would replace the single Hermitian $h$ by the actual control algebra, not by a freely chosen full matrix algebra.  The absence of a historical panel is not a reason to invent one; it is a reason to isolate its precise source domain and admission rules.

The later determinant and tensor-type selection questions should remain downstream.  No new abstract proof of the one-bridge matrix theorem, Howe genericity, or the profile-only obstruction is needed for that continuation.

\subsection{Verification scope}
The proofs in this note are finite-dimensional and self-contained apart from the explicitly identified standard representation decompositions and inherited source/generation results.  An accompanying exact-arithmetic script verifies the $S_4$ character census, the twelve-root orbital algebra and its centre, the return mass $288/24=12$, the normalized instrument effects, the dark and dynamically opened mediator examples, the early-twirl counterexample, the right-module defect identity, and the fourteen/fifteen scalar-separation alternatives.  These calculations check the examples and implementation identities.  They are not a formal proof-assistant verification or an independent review of the general theorems.

\begin{thebibliography}{99}
\bibitem{Grand}
A.~P\'elissier, \emph{Renewal Grand Tensor Monograph}, V076,
supplied source \src{Renewal_Grand_Tensor_Monograph_V076(1).tex}.
The history algebra, regular-trace word hierarchy, source shorts and typed source interpretations are retained as source inputs.

\bibitem{Continuum}
A.~P\'elissier, \emph{Renewal Standard Model--Einstein Continuum Emergence},
supplied source \src{SM_Einstein_Continuum_Monograph_Updated(3)(1).tex}.

\bibitem{Duality}
A.~P\'elissier, \emph{Spacetime--Gauge Commutant Duality: Finite Rigidity and Cofinal Stability}, supplied source
\src{Spacetime_Gauge_Commutant_Duality_CMP_Analytic_Strengthened(1).tex}.

\bibitem{Classical}
A.~P\'elissier, \emph{Renewal Geometry and the Emergence of the Classical Einstein--Standard-Model System}, supplied source
\src{Renewal_Geometry_Classical_Einstein_SM_CMP_Submission(1).tex}.

\bibitem{Attack}
A.~P\'elissier, \emph{Renewal SM Unconditional Emergence Attack}, cumulative accessible source
\src{Renewal_SM_Unconditional_Emergence_Attack_V128.tex}, especially the V106--V110 development blocks and their scope statements.  Later blocks concern separate action, marginal-flow and source-escape problems.

\bibitem{Rigidity}
\emph{Renewal Spacetime--Gauge Duality Research Notes}, cumulative V33,
supplied source \src{Renewal_Spacetime_Gauge_Duality_Research_Notes_V33(2).tex}.
Used to exclude repetition of the finite rigidity, support, genericity and operational-identifiability programme.

\bibitem{Etingof}
P.~Etingof, O.~Golberg, S.~Hensel, T.~Liu, A.~Schwendner, D.~Vaintrob and E.~Yudovina,
\emph{Introduction to Representation Theory}, arXiv:0901.0827;
\href{https://arxiv.org/abs/0901.0827}{author preprint}.
Standard representation-theoretic background; no claim of new priority is made for the finite algebra methods used here.
\end{thebibliography}

% END OF SOURCE-SELECTION V1 DEVELOPMENT BLOCK.
% Subsequent versions append new sections here, preserving this mathematical body.

\clearpage
\hypersetup{pdftitle={Historical Standard-Model Source Selection: Research Notes V1--V2}}
\begin{center}
{\Large\bfseries V2: Recurrence before reversal}\par\medskip
{\large Forward operational words and endpoint-certified colour returns}
\end{center}
\addcontentsline{toc}{part}{V2: Recurrence before reversal}

\section{The upstream audit: do not manufacture a reverse operation}
\label{sel2:context}

The positive result of V1 needs actual reverse-exit routes.  The next source
question is whether that hypothesis can be removed without adjoining the
missing reverse as a new control.  It is not enough to observe that every
operator has an adjoint, or that the reconstructed history object is a
$C^*$-algebra.

The Grand Tensor makes this distinction explicitly in
\src{sec:process-history-algebra}: its primitive maps are reduced
Schr\"odinger completely positive (CP) superoperators, and the formal reverse
is the Hilbert--Schmidt adjoint, not an asserted executable reverse channel.
The preceding assessment separates the forward predictive algebra from this
positive adjoint completion.  These are source statements, not omissions to be
filled by identifying the two algebras.

A second nonduplication check is essential.  V6--V7 of \cite{Rigidity} already
prove the faithful-state fixed-observable/commutant identity, the finite
quantum hitting test, maximal recurrent support, and its semidefinite
faithfulness certificate.  V10 of \cite{Attack} already uses physically
available inverse unitary pulses to implement a reversible monitored kernel.
None of those results is presented here as a new recurrence or
first-return theorem.  The new question concerns the \emph{unital algebra of
forward Kraus words}, and whether a nonzero corner obtained using formal
adjoints must also occur in a genuinely forward, declared CP history.

\begin{center}
\small
\begin{tabular}{L{.37\textwidth}L{.55\textwidth}}
\toprule
Inherited input & Additional question treated in V2\\
\midrule
Faithful-state fixed-observable identity, \src{v6:thm:dirichlet} in \cite{Rigidity}
& Does faithfulness also identify the forward coefficient algebra with its
adjoint closure?\\[.35em]
Maximal stationary support, \src{v7:thm:recurrent-support}
& On which actual carrier is such a forward-word identification legitimate?\\[.35em]
Quantum hitting-or-trap test, \src{v6:thm:hitting}
& Can one extract a nonzero finite forward word, rather than assume a reverse
pulse or prove almost-sure first hitting?\\[.35em]
The V1 reciprocal-mediator theorem
& Can two forward entries into a common mediator force a return without
physical admission of either reverse entry?\\[.35em]
Source-minimal CP realization and Kraus-span invariance
& When does a nonzero CP transfer supply a uniquely defined amplitude branch,
without choosing an unobserved Kraus label?\\
\bottomrule
\end{tabular}
\end{center}

The answer is positive on a specified recurrent operational class.  It is not
unconditional: the complete actual scheduler, its recurrent support, the
exposed ports, and any endpoint filters must be verified.  Two explicit
negative evaluations below show why the six-edge normalization and the
one-way V1 benchmark do not by themselves supply this recurrent class.
General faithful-channel structure is established theory
\cite{sel2-Carbone}; the new role developed here is its finite
forward-witness consequence for the source-selection problem.  No priority
claim is made for standard finite channel or representation theory.

\subsection{Three algebras on two different spaces}
Fix a finite physical Hilbert space $\mathcal F$ of dimension $N$ and an
actually specified instrument $(\mathcal J_a)_a$.  Its nonselective state map
is
\begin{equation}
 \Phi(X)=\sum_a\mathcal J_a(X)
       =\sum_{a,\mu}K_{a\mu}XK_{a\mu}^*,
 \qquad \sum_{a,\mu}K_{a\mu}^*K_{a\mu}=I_{\mathcal F}.
 \label{sel2:eq:channel}
\end{equation}
Kraus indices $\mu$ need not be physically resolved.  They are coordinates
for the actual CP branch.  Let $\mathcal K$ be the span of all its Kraus
coefficients and define
\begin{equation}
 \mathcal A_+:=\operatorname{alg}(I,\mathcal K),\qquad
 \mathcal A_*:=C^*(\mathcal K),\qquad
 \mathcal S_+:=\operatorname{alg}(I_{\End(\mathcal F)},\mathcal J_a:a).
 \label{sel2:eq:three-algebras}
\end{equation}
The first two act on $\mathcal F$; the last acts on $\End(\mathcal F)$.
Kraus changes of coordinates preserve $\mathcal K$ and $\mathcal A_+$.
The equality $\mathcal A_+=\mathcal A_*$ will \emph{not} imply closure of
$\mathcal S_+$ under Hilbert--Schmidt adjoints, and no linear combination in
$\mathcal A_+$ is called a new executable instrument.

The scheduler is required to be complete on these cuts.  If future
admission depends on a controller state, that state is included in
$\mathcal F=\bigoplus_x\mathcal F_x$ and every primitive coefficient has its
actual $y\leftarrow x$ block.  Products of incompatible blocks vanish.
Hidden classical or quantum memory must not be discarded before taking
powers of $\Phi$.  These hypotheses identify one fixed finite source; they
do not authorize arbitrary rescheduling of separately supplied controls.

\section{Faithful recurrence is exactly forward adjoint completeness}
\label{sel2:forward}

A stationary density is $\sigma\succeq0$ with $\Tr\sigma=1$ and
$\Phi(\sigma)=\sigma$.  It is faithful when $\sigma\succ0$ on the entire
stated carrier.  Neither detailed balance nor uniqueness of $\sigma$ is
assumed.

\begin{lemma}[No one-way enclosure on a faithful stationary carrier]
\label{sel2:lem:reducing}
Suppose $\Phi$ has a faithful stationary density.  Every subspace invariant
under all $K_{a\mu}$ is reducing under all those coefficients.  In particular,
if $P$ is its orthogonal projection, then
\begin{equation}
 (I-P)K_{a\mu}P=0\ \text{for all }a,\mu
 \quad\Longrightarrow\quad
 PK_{a\mu}(I-P)=0\ \text{for all }a,\mu.
 \label{sel2:eq:enclosure}
\end{equation}
\end{lemma}
\begin{proof}
Write a coefficient relative to $P\mathcal F\oplus(I-P)\mathcal F$ as
$K_j=\begin{psmallmatrix}A_j&B_j\\0&D_j\end{psmallmatrix}$, combining
$a,\mu$ into $j$.  Trace preservation gives
$\sum_j A_j^*A_j=I$ and $\sum_j A_j^*B_j=0$.  Thus the unital Heisenberg
dual satisfies
\begin{equation}
 \Phi^*(P)-P=0\oplus\sum_j B_j^*B_j\succeq0.
 \label{sel2:eq:subharmonic-defect}
\end{equation}
Stationarity gives
$\Tr\sigma[\Phi^*(P)-P]=\Tr P[\Phi(\sigma)-\sigma]=0$.
A positive operator with zero expectation in a faithful state is zero.
Hence each $B_j$ is zero.  No commutation of $\sigma$ with $P$ was used.
\end{proof}

\begin{theorem}[Faithful recurrence and the forward coefficient algebra]
\label{sel2:thm:forward-star}
For a finite trace-preserving CP map, the following are equivalent:
\begin{enumerate}[label=\textup{(\roman*)},leftmargin=2.6em]
\item $\Phi$ admits a faithful stationary density;
\item $\mathcal A_+$ is closed under adjoints;
\item $\mathcal A_+=\mathcal A_*$.
\end{enumerate}
In the positive case, every formal coefficient-adjoint word is a linear
combination of forward coefficient words.  If $D=\dim_\C\mathcal A_+$,
forward words of length at most $D-1\le N^2-1$ suffice to span that algebra.
This is an algebraic span statement, not physical implementation of a reverse
CP channel or of coherent linear combinations of histories.
\end{theorem}
\begin{proof}
To prove (i)$\Rightarrow$(ii), amplify the channel to
$\operatorname{id}_{M_N}\otimes\Phi$ on $\mathcal F^{\oplus N}$.
Its coefficients are $I_N\otimes K_j$ and it has faithful stationary density
$(I_N/N)\otimes\sigma$.  For an orthonormal basis $e_1,\ldots,e_N$ of
$\mathcal F$, the subspace
\[
 \mathcal L=\{(ae_1,\ldots,ae_N):a\in\mathcal A_+\}
\]
is invariant under every amplified coefficient.  By
Lemma~\ref{sel2:lem:reducing}, it is also invariant under their adjoints.
The vector $(e_1,\ldots,e_N)$ belongs to $\mathcal L$; applying
$I_N\otimes K_j^*$ shows that some $a\in\mathcal A_+$ agrees with $K_j^*$
on every basis vector.  Thus $K_j^*\in\mathcal A_+$, which proves adjoint
closure.  This proof does not assume that every invariant subspace of an
arbitrary nonselfadjoint algebra has a complement.

For (ii)$\Rightarrow$(i), take the represented finite $C^*$-decomposition
\[
 \mathcal F\cong\bigoplus_\beta\C^{d_\beta}\otimes\C^{m_\beta},\qquad
 \mathcal A_+=\bigoplus_\beta M_{d_\beta}(\C)\otimes I_{m_\beta}.
\]
Write $K_j=\bigoplus_\beta k_{j\beta}\otimes I$.
Each map $\Phi_\beta(X)=\sum_j k_{j\beta}Xk_{j\beta}^*$ is trace preserving
and has a stationary density, by finite Cesaro averaging.  The support of
any stationary density is invariant under every $k_{j\beta}$: positivity of
the stationary equation makes every coefficient preserve that support.
The $k_{j\beta}$ generate $M_{d_\beta}$, so a nonzero invariant support is
the full block.  Choose a faithful stationary density $\sigma_\beta$ in
each block.  A strictly positive weighted direct sum of
$\sigma_\beta\otimes I_{m_\beta}/m_\beta$ is faithful and stationary on
$\mathcal F$.  Statements (ii) and (iii) are equivalent by the definition
of $\mathcal A_*$.

Finally let $V_\ell$ be the span of $I$ and all forward words of length at
most $\ell$.  Then $V_{\ell+1}=V_\ell+\sum_jK_jV_\ell$.  A plateau makes
this space invariant under every left multiplication, so it already contains
all later words.  Starting at dimension one, at most $D-1$ strict dimension
increases are possible.  This gives the bound.
\end{proof}

\begin{corollary}[Reuse of the existing recurrent-support compiler]
\label{sel2:cor:recurrent-reuse}
Let $p$ be the canonical maximal stationary support reconstructed in V7 of
\cite{Rigidity}.  The actual recurrent coefficients $pK_jp$ have
forward algebra equal to their $C^*$-algebra.  This conclusion applies on
$p\mathcal F$ only.  A type, private, or mediator direction removed by $p$
must not be declared recurrent merely to apply
Theorem~\ref{sel2:thm:forward-star}.
\end{corollary}
\begin{proof}
The inherited recurrent-support theorem gives a normalized compressed
instrument and a faithful stationary density on $p\mathcal F$.  Apply
Theorem~\ref{sel2:thm:forward-star} there.
\end{proof}

\begin{proposition}[A quantitative witness against a proposed faithful state]
\label{sel2:prop:stationary-error}
Let $P\mathcal F$ be exactly invariant under all coefficients, and suppose
$\sigma\succeq sI$, $s>0$, $\Tr\sigma=1$, is an estimated stationary density.
Then
\begin{equation}
 \sum_j\norm{PK_j(I-P)}_{\HS}^2
 \le \frac{\norm{\Phi(\sigma)-\sigma}_1}{2s}.
 \label{sel2:eq:stationary-error}
\end{equation}
Thus a measured one-way enclosure and positive reverse leakage give a
lower bound on the stationarity residual of any density with the stated
faithfulness floor.
\end{proposition}
\begin{proof}
Equation~\eqref{sel2:eq:subharmonic-defect} has trace equal to the sum on the
left.  Its expectation in $\sigma$ is at least $s$ times that trace and is
$\Tr P[\Phi(\sigma)-\sigma]$.  The bracket is Hermitian and trace zero,
so its expectation in a projection has absolute value at most half its trace
norm.  Combine these facts.  Approximate invariance of $P$ is not covered by
this particular estimate.
\end{proof}

\section{Extracting an actual forward history, without a reverse pulse}
\label{sel2:words}

For a nonzero projection $P$ of rank $r$, define
\begin{equation}
 W_P^{(\ell)}=\sum_{k=0}^{\ell}\Phi^k(P),\qquad
 \mathcal R_P^{(\ell)}=\Ran W_P^{(\ell)},\qquad
 W_P=W_P^{(N-r)}.
 \label{sel2:eq:reachable-gram}
\end{equation}
The initial operator $P$ is not normalized; dividing by $r$ gives the
maximally mixed state on its support.  These are unmonitored forward
iterates of the actual channel, not a newly inserted sequence of projective
measurements.

\begin{theorem}[Finite forward-transfer test]
\label{sel2:thm:forward-transfer}
For any second projection $Q$,
\begin{equation}
 \boxed{
 \Gamma_{Q\leftarrow P}:=\Tr(QW_P)
 =\sum_{k=0}^{N-r}\ \sum_{|w|=k}\norm{QK_wP}_{\HS}^2.}
 \label{sel2:eq:transfer-mass}
\end{equation}
The following are equivalent:
\begin{enumerate}[label=\textup{(\roman*)},leftmargin=2.6em]
\item some forward coefficient word has a nonzero $Q\leftarrow P$ corner;
\item $\Gamma_{Q\leftarrow P}>0$;
\item some actual, possibly unresolved, CP outcome word $w$ of length at most
$N-r$ satisfies $\Tr[Q\mathcal J_w(P)]>0$.
\end{enumerate}
If the mass vanishes, $\mathcal R_P=\Ran W_P$ is an explicitly reconstructed
invariant subspace containing $P\mathcal F$ and orthogonal to
$Q\mathcal F$, excluding every longer forward history.  If the channel has a
faithful stationary density, that subspace is reducing and
\begin{equation}
 \Gamma_{Q\leftarrow P}=0
 \quad\Longleftrightarrow\quad
 \Gamma_{P\leftarrow Q}=0,
 \label{sel2:eq:qualitative-reciprocity}
\end{equation}
where each mass uses its own rank-dependent horizon.  No equality of
transition probabilities or detailed balance is implied.
\end{theorem}
\begin{proof}
For a positive matrix $X$, $\Ran\Phi(X)$ is the span of the spaces
$K_j\Ran X$.  This follows either from a positive square root of $X$, or
from the identity
$\langle v,\Phi(X)v\rangle=\sum_j\norm{X^{1/2}K_j^*v}^2$.
Consequently
\[
 \mathcal R_P^{(\ell+1)}=\mathcal R_P^{(\ell)}+
                      \sum_jK_j\mathcal R_P^{(\ell)}.
\]
Once dimensions plateau, this space is invariant and no further increase is
possible.  Its initial dimension is $r$ and its maximum is $N$, so it has
stabilized by $\ell=N-r$.  It is therefore the span of all forward images
of $P\mathcal F$.

Expanding $\Phi^k$ into Kraus words and taking trace gives
\eqref{sel2:eq:transfer-mass}.  A positive sum is nonzero precisely when
some summand is nonzero.  Regrouping fine Kraus words by their actual outcome
labels gives
$\Gamma=\sum_{k\le N-r}\sum_{|w|=k}\Tr[Q\mathcal J_w(P)]$; positivity
proves (iii) without declaring the Kraus labels to be observed.  Zero mass
means that $\Ran W_P$ is orthogonal to $Q\mathcal F$, giving the invariant
separator.  Under faithfulness, Lemma~\ref{sel2:lem:reducing} makes it
reducing.  Equivalently, Theorem~\ref{sel2:thm:forward-star} gives
$Q\mathcal A_+P=0$ if and only if $P\mathcal A_+Q=0$ by taking adjoints.
This proves the last equivalence.
\end{proof}

\begin{remark}[Difference from the previously proved quantum hitting theorem]
Theorem~\ref{sel2:thm:forward-transfer} tests \emph{existence} of a forward
history with positive terminal transfer, not probability-one first arrival
from every initial state.  Its $N-r$ bound comes from a reachable
\emph{Hilbert-subspace} chain.  The inherited V6 first-hitting theorem uses
a killed CP map and a different operator-space horizon.  Neither result is
substituted for the other, and no unrecorded intermediate monitor is added.
\end{remark}

\begin{corollary}[A finite conversion of a formal corner to forward histories]
\label{sel2:cor:formal-to-forward}
On a faithful stationary carrier, every nonzero $Q\mathcal A_*P$ corner
contains a nonzero corner of a forward Kraus word of length at most $N-r$.
For a general linear source-shadow map $F:\mathcal A_*\to\mathcal Z$,
$F(\mathcal A_*)\ne0$ implies $F(K_w)\ne0$ for some forward word of length
at most $N^2-1$.
\end{corollary}
\begin{proof}
Use $\mathcal A_* =\mathcal A_+$ from
Theorem~\ref{sel2:thm:forward-star}.  The first assertion follows from
Theorem~\ref{sel2:thm:forward-transfer}; the second follows by applying the
linear map to the finite spanning bank of
Theorem~\ref{sel2:thm:forward-star}.
\end{proof}

The second assertion does not imply that one history has several separately
nonzero shadows or satisfies a nonlinear same-history provenance relation.
The existing determinant--matter synchronization obstruction is retained.
For unresolved outcomes, a coefficient-shadow witness is not automatically
an independently resolved event; the following endpoint construction avoids
that ambiguity for the colour corner.

\subsection{A pure returned branch without resolving the environment}

\begin{lemma}[A repeatable rank-one Read is its L\"uders filter]
\label{sel2:lem:read}
Let an actual CP outcome $\mathcal R$ have effect
$\mathcal R^*(I)=p=|t\rangle\langle t|$, with $\norm t=1$.  If its
outcome is repeatable, in the sense
$\Tr[(I-p)\mathcal R(p)]=0$, then
\begin{equation}
 \mathcal R(X)=pXp\quad\text{for every }X.
 \label{sel2:eq:repeatable-filter}
\end{equation}
A rank-one effect alone, without repeatability, does not imply this equality.
\end{lemma}
\begin{proof}
For Kraus operators $R_j$, $\sum R_j^*R_j=p$ implies $R_j=R_jp$.
Repeatability and positivity imply $(I-p)R_jp=0$ for every $j$.
Thus $R_j=c_jp$, with $\sum_j|c_j|^2=1$, giving
\eqref{sel2:eq:repeatable-filter}.  Without repeatability the map can instead
prepare any normalized output density from the input weight $\Tr(pX)$.
\end{proof}

\begin{theorem}[Endpoint certification of an unresolved CP return]
\label{sel2:thm:endpoint}
Let $p=|t\rangle\langle t|$ and $q=|h\rangle\langle h|$ be actual
repeatable rank-one Reads on a common source and target frame.  For any
admitted CP history $\mathcal J_w$, the selected composite
\begin{equation}
 \mathcal C_w(X):=q\mathcal J_w(pXp)q
 \label{sel2:eq:filtered-return}
\end{equation}
has Choi rank at most one.  More precisely, with
$c_w=\Tr[q\mathcal J_w(p)]\ge0$,
\begin{equation}
 \boxed{\mathcal C_w(X)=C_wXC_w^*,\qquad
 C_w=\sqrt{c_w}\,|h\rangle\langle t|.}
 \label{sel2:eq:pure-return}
\end{equation}
When $c_w>0$, this coefficient is unique up to phase.  It is a physical
conditional amplitude for the displayed composite and not an arbitrarily
chosen Kraus component of $\mathcal J_w$.
\end{theorem}
\begin{proof}
Lemma~\ref{sel2:lem:read} identifies the actual Read maps with their L\"uders
filters.  Since $pXp=\langle t,Xt\rangle p$, and
$q\mathcal J_w(p)q=c_wq$, equation~\eqref{sel2:eq:pure-return} follows
for all matrices, not only states.  The Choi matrix is
$c_w|h\otimes\overline t\rangle\langle h\otimes\overline t|$.
Its unique nonzero eigenspace proves the phase uniqueness.
\end{proof}

\begin{corollary}[Colour from a measured transfer rather than a chosen exit]
\label{sel2:cor:colour-transfer}
Keep the faithful shorted type line $T$ and the actually supported private
$M_2(H)$ of V1.  Suppose repeatable rank-one Reads on $T$ and on a finite bank of rays spanning $H$ are actually admitted.  If
$\Gamma_{p_H\leftarrow p_T}>0$, there is an actual filtered history of at
most $N+1$ elementary operations whose uniquely defined amplitude is a
nonzero $H\leftarrow T$ corner.  The colour amplitude algebra is then $M_3$
by the inherited bridge theorem.  Unresolved Kraus outcomes along the route
need not be individually measured.
\end{corollary}
\begin{proof}
Here $\rank p_T=1$.  The positive transfer sum has at least one CP outcome
word of length at most $N-1$ with positive $H$ output.  The sum of the admitted private rank-one Read projections is
strictly positive on $H$, so at least one gives $c_w>0$.  Theorem~\ref{sel2:thm:endpoint} supplies the amplitude.
The two endpoint filters add at most two operations.  Apply the inherited
colour-bridge criterion on $T\oplus H$.
\end{proof}

If the private Read projections are $q_1,\ldots,q_s$, put
$\nu_H=\lambda_{\min}(\sum_jq_j|_H)>0$.  Their total tested transfer obeys
\begin{equation}
 \sum_{j=1}^s\Tr(q_jW_{p_T})\ge\nu_H\Gamma_{p_H\leftarrow p_T}.
 \label{sel2:eq:read-frame}
\end{equation}
In particular the two already retained, distinct nonorthogonal private rays
suffice when their Reads are operationally repeatable; no extra orthogonal
basis measurement is needed.  This inequality is simply the positive form
$\sum_jq_j\succeq\nu_Hp_H$ applied to $W_{p_T}$.

The output in this corollary is an algebra of \emph{physical conditional
amplitudes on the declared Hilbert carrier}.  It is not an assertion that the
superoperator algebra $\mathcal S_+$ is isomorphic to $M_3$, nor that the
nonselective channel performs a coherent unitary colour rotation.  A private
projection belonging only to a formal algebra is not an admitted repeatable
Read; its operational availability must be checked separately.

\section{A common mediator with forward entries only}
\label{sel2:mediator}

We now replace the reciprocal route hypothesis of V1 by conditions on the
actual nonselective source dynamics.  The mediator controls may be
noncommuting and nonselfadjoint.  No single-Hermitian-$h$ normal form and no
physical reverse-pulse bank are required.

\begin{theorem}[Recurrent common-mediator forcing]
\label{sel2:thm:mediator}
Let the actual instrument \eqref{sel2:eq:channel} have a faithful stationary
density on $\mathcal F$.  Let $T,H,E\subseteq\mathcal F$ be mutually
orthogonal, with $\dim T=1$, and let $p_T,p_H,p_E$ be their projections.
Suppose the following two forward entry rows belong to the instrument:
\begin{enumerate}[label=\textup{(E\arabic*)},leftmargin=2.8em]
\item a branch, or a finite family of actual branches, from $T$ to $E$ whose
combined output on $p_T$ is strictly positive on $E$;
\item a nonzero actual branch from $H$ to $E$.
\end{enumerate}
All input and output supports in these rows are exact CP branch supports, not
merely nonzero marginal effects.  Then
\begin{equation}
 \boxed{\Gamma_{p_H\leftarrow p_T}>0.}
 \label{sel2:eq:forced-forward}
\end{equation}
Some actual forward CP history of length at most $N-1$ transfers positive
weight from $T$ to $H$.  Neither an exit $E\to H$ nor the adjoint of the
$H\to E$ entry is assumed to be physically executable.

In particular, (E1) follows from one nonzero \emph{covariant} CP entry from
the trivial type line into an irreducible unitary $G$-mediator $E$: for its
actual branch $\mathcal J_A$,
\begin{equation}
 \mathcal J_A(p_T)=\frac{\Tr\mathcal J_A(p_T)}{\dim E}\,p_E\succ0
 \quad\hbox{on }E.
 \label{sel2:eq:covariant-entry}
\end{equation}
An actually admitted full orbit of nonzero entry branches gives the same
support conclusion by Schur averaging, without requiring the average itself
to be an additional physical control.
\end{theorem}
\begin{proof}
Let $\mathcal R_T=\Ran W_{p_T}$.  By
Theorem~\ref{sel2:thm:forward-transfer}, it is invariant under all
coefficients and contains every one-step image of $T$.  Condition (E1)
therefore gives $E\subseteq\mathcal R_T$.  By faithful stationarity and
Lemma~\ref{sel2:lem:reducing}, $\mathcal R_T$ is reducing for all
coefficients.  Choose a nonzero coefficient $B$ of the branch in (E2).
It has exact support $B=p_EBp_H$.  A vector $v\in E$ with $B^*v\ne0$
exists, and reduction gives
\[
 0\ne B^*v\in H\cap\mathcal R_T.
\]
Consequently $\mathcal R_T$ is not orthogonal to $H$.
Theorem~\ref{sel2:thm:forward-transfer} proves positivity and extracts a
forward CP word within the stated horizon.  The vector $B^*v$ is used only
inside a proof of reachability; it is not declared an executed reverse.

For the covariance assertion, the input $p_T$ is invariant, so its image
commutes with the irreducible $G$ action on $E$.  Schur's lemma makes that
positive image scalar on $E$, with scalar determined by the trace.
Nonzero entry makes the trace positive.  For an admitted orbit, the sum of
its positive images has the same full support; positive scheduling weights
do not change the support of this sum.
\end{proof}

\begin{remark}[The entry may be mixed]
If $\dim E>1$, a covariant CP entry from an invariant line normally has
mixed output of full rank; no rank-one Kraus resolution of that branch is
asserted.  This is why Theorem~\ref{sel2:thm:mediator} first concludes a
positive \emph{CP history}.  The endpoint theorem, when its actual
rank-one Reads are supplied, then gives a pure selected amplitude.  The two
steps must not be compressed into an invented coherent entry vector.
\end{remark}

\begin{corollary}[Noncommuting mediator control and one-sided reachability]
\label{sel2:cor:general-mediator}
In Theorem~\ref{sel2:thm:mediator}, replace (E1) by the weaker condition that
the actual forward reachable space from $T$ contains $E$.  The conclusion
is unchanged.  Thus any supplied collection of noncommuting mediator controls
can be used; its reachability is tested by \eqref{sel2:eq:reachable-gram},
not by declaring a freely controllable matrix algebra on the mediator.
\end{corollary}
\begin{proof}
The proof of Theorem~\ref{sel2:thm:mediator} uses (E1) only to show
$E\subseteq\mathcal R_T$.  Retain that conclusion and repeat its final
three lines.
\end{proof}

\begin{theorem}[Source-level internal saturation without an independent reverse]
\label{sel2:thm:landing}
Retain the verified final-cut carrier, private amplitude algebra, weak packet,
and central separation of Theorem~\ref{sel1:thm:landing}.  In place of its
reciprocal-port condition (P3), assume a complete actual instrument on the
cut/controller carrier satisfying Theorem~\ref{sel2:thm:mediator} or
Corollary~\ref{sel2:cor:general-mediator}.  Assume in addition that the
rank-one type Read and a spanning bank of rank-one private Reads are actual
repeatable outcomes, with the endpoint selection grammar admitted.
Then a nonzero colour amplitude is furnished by an actual forward selected
history.  The resulting internal amplitude algebra on the final-cut carrier
is
\begin{equation}
 \boxed{M_3(\C)\oplus M_2(\C)\oplus\C\oplus\C,}
 \label{sel2:eq:internal-landing}
\end{equation}
equal to the stated external relative commutant.  No independent direct
colour bridge, inverse pulse, or resolved environment Kraus label is needed.
A selected bridge uses at most $N+1$ operations, where $N$ is the dimension
of the full retained cut/controller carrier, not merely the final-cut
colour carrier.
\end{theorem}
\begin{proof}
Theorem~\ref{sel2:thm:mediator} gives a positive forward transfer.  The
actual endpoint Reads then give a nonzero, phase-unique amplitude by
Corollary~\ref{sel2:cor:colour-transfer}.  Its support is $H\leftarrow T$;
both spaces are externally trivial, so this amplitude belongs to the
internal algebra.  The inherited colour theorem closes the private plane
and type line to $M_3$.  The weak and central parts are the retained inputs.
The assembly proof of Theorem~\ref{sel1:thm:landing} now applies without
change.  The operation count is that in
Corollary~\ref{sel2:cor:colour-transfer}.
\end{proof}

\begin{remark}[What is reduced, and what is not]
The V1 requirement of a physically admitted $B^*$ is replaced by faithful
recurrence of the actual complete channel.  This property is independently
testable using the already established V6--V7 source compiler; it is not a
choice of a favourable inverse instrument.  The shared mediator and actual
entry rows still have to be supplied.  The repeatable endpoint Reads have
been stated explicitly rather than inferred from kinematic projections.
The external multiplicity pattern, determinant incidence, matter tensor types,
and dynamical attraction are not selected by this theorem.
\end{remark}

\section{Exact tests that prevent a circular recurrence argument}
\label{sel2:tests}

\subsection{A directed recurrent benchmark: no reverse entry in the alphabet}

\begin{example}[Forward return with an absent inverse primitive]
\label{sel2:ex:directed}
On $\mathcal F=\C^4$ take $T=\C e_0$, $H=\Span\{e_1,e_2\}$ and
$E=\C e_3$.  Let $0<\epsilon<1$ and use the five actual Kraus outcomes
\begin{equation}
 K_0=E_{30},\quad K_1=E_{31},\quad K_2=E_{12},\quad
 K_3=\sqrt{1-\epsilon}\,E_{03},\quad
 K_4=\sqrt\epsilon\,E_{23}.
 \label{sel2:eq:directed-bank}
\end{equation}
Their effects sum to identity.  The stationary density is
\begin{equation}
 \sigma_\epsilon=\frac1{2+\epsilon}
        \diag(1-\epsilon,\epsilon,\epsilon,1)\succ0.
 \label{sel2:eq:directed-state}
\end{equation}
The entries $K_0:T\to E$ and $K_1:H\to E$ are nonzero.  There is no
primitive $K_1^*=E_{13}$, and the classical dynamics is not reversible:
$1\to3$ has positive probability but $3\to1$ has probability zero in one
step.  Nevertheless
\begin{equation}
 K_4K_0=\sqrt\epsilon\,E_{20},\qquad
 K_2K_4K_0=\sqrt\epsilon\,E_{10},\qquad
 K_1^*=\epsilon^{-1/2}K_2K_4.
 \label{sel2:eq:directed-returns}
\end{equation}
Thus both required return amplitudes arise from forward histories, not from
execution of $K_1^*$.

For the exact three-step horizon $N-\rank p_T=3$,
\begin{equation}
 \Gamma_{p_H\leftarrow p_T}=2\epsilon.
 \label{sel2:eq:directed-mass}
\end{equation}
At $\epsilon=16/25$ all primitive coefficients are rational and
$\sigma=\diag(9,16,16,25)/66$.  The two nonzero returns in
\eqref{sel2:eq:directed-returns} have coefficient $4/5$ and the transfer mass
is $32/25$.  The mass is a sum over several observation times, not one
probability, so it can exceed one.
\end{example}
\begin{proof}
The five input-column effects are respectively $p_0,p_1,p_2$,
$(1-\epsilon)p_3,\epsilon p_3$, proving normalization.
The state update kills all off-diagonal entries and maps diagonal weights to
$((1-\epsilon)x_3,x_2,\epsilon x_3,x_0+x_1)$.
Solving the stationary equations and normalizing gives
\eqref{sel2:eq:directed-state}.  Matrix-unit multiplication proves
\eqref{sel2:eq:directed-returns}.  Starting at $p_0$, steps one, two and
three have $H$ weights $0,\epsilon,\epsilon$ respectively, proving
\eqref{sel2:eq:directed-mass}.
\end{proof}

\begin{remark}[The calibration must not be lost when replacing a reverse]
As $\epsilon\downarrow0$, both entry amplitudes remain fixed, but the actual
forward return mass tends to zero.  The stationary faithfulness floor also
vanishes.  The formal coefficient $\epsilon^{-1/2}$ converting the missing
reverse into forward words diverges.  It would therefore be incorrect to
import the unscaled two-leg V1 return margin after replacing a reverse by
such a linear combination.  V2 retains the actual forward probabilities.
At zero, $T\oplus E$ is recurrent while $H$ is transient, and no forward
$T\to H$ history exists.
\end{remark}

\subsection{The existing V1 one-way benchmark is not recurrent on its full carrier}

\begin{proposition}[Exact recurrence audit of the V1 normalized route example]
\label{sel2:prop:V1-audit}
Complete the normalized instrument of
Proposition~\ref{sel1:prop:normalization} on
$\mathcal F=T\oplus E\oplus H$ by the no-entry coefficient $p_T/2$ and an
absorbing coefficient $p_H$ on $H$.  Keep its actual entries $A_g/8$ and
its actually supplied exits $B_g^*/12$.  The maximal stationary support is
exactly $H$.  Therefore the V1 benchmark does not satisfy the full-carrier
faithful recurrence hypothesis of Theorem~\ref{sel2:thm:mediator}.
\end{proposition}
\begin{proof}
The inherited entry and exit effects give
\[
 \frac14p_T+\frac34p_T+p_E+p_H=I.
\]
At a stationary density, let $t,e,h$ denote the three cut populations.
The type equation is $t=t/4$, so $t=0$.  The mediator equation is
$e=3t/4$, so $e=0$.  Positivity then makes all rows and columns involving
$T$ or $E$ vanish.  Every density supported on $H$ is stationary because
of the absorbing coefficient, proving the maximal-support assertion.
\end{proof}

This is a negative evaluation of an already displayed construction, not a
newly chosen adversarial source.  Its V1 bridge remains valid because its
exit was explicitly supplied.  Recurrence is a different sufficient mechanism;
its failure does not invalidate a directly supplied exit, and adding a new
return branch merely to make the state faithful would change the source.

\subsection{The locked six-edge identities do not imply recurrence}

\begin{proposition}[A transient completion of the exact six-edge normalization]
\label{sel2:prop:six-transient}
On $\C^3$ define
\begin{equation}
 B_1=\diag(1,1/3,1/3),\quad
 B_2=\tfrac23E_{01},\quad B_3=\tfrac23E_{02},\quad
 B_4=\tfrac23E_{12},\quad B_5=\tfrac23E_{21}.
 \label{sel2:eq:transient-five}
\end{equation}
Let $O\in\R^{6\times5}$ satisfy
$O^{\mathsf T}O=I$ and $O^{\mathsf T}\mathbf1=0$, as in the six-edge
normal form of V33 of \cite{Rigidity}, and put
\begin{equation}
 K_e=\frac1{\sqrt{18}}I+
              \sqrt{\frac23}\sum_{\mu=1}^5O_{e\mu}B_\mu.
 \label{sel2:eq:transient-six}
\end{equation}
These six matrices are linearly independent and satisfy
\begin{equation}
 \sum_eK_e^*K_e=I,\qquad \sum_eK_e=\sqrt2 I.
 \label{sel2:eq:locked-normalizations}
\end{equation}
Their channel has only one stationary density, $E_{00}$, and no faithful
one.  Its forward coefficient algebra has dimension seven,
\begin{equation}
 \mathcal A_+=
 \left\{\begin{pmatrix}a&b_1&b_2\\0&c_{11}&c_{12}\\0&c_{21}&c_{22}
 \end{pmatrix}\right\},
 \qquad \mathcal A_*=M_3(\C).
 \label{sel2:eq:seven-nine}
\end{equation}
Thus the literal six-edge constraints and identity coefficient pivot do not
supply the faithful-recurrence premise or equality of forward and formal
adjoint word algebras.
\end{proposition}
\begin{proof}
The source-column effects of \eqref{sel2:eq:transient-five} add to identity,
since $1/9+4/9+4/9=1$ on the two nonzero columns.  The six matrices
$I,B_1,\ldots,B_5$ are independent: four are separate off-diagonal units and
$I,B_1$ have independent diagonals.  The inherited orthogonal coefficient
change gives the independence and identities in
\eqref{sel2:eq:locked-normalizations}.  Direct expansion gives
\[
 \Phi=\tfrac13\operatorname{id}+\tfrac23\Psi,\qquad
 \Psi(X)=\sum_{\mu=1}^5B_\mu X B_\mu^*.
\]
Both channels have the same stationary densities.  The $00$ entry of the
stationary equation for $\Psi$ is
$\rho_{00}=\rho_{00}+\frac49(\rho_{11}+\rho_{22})$.
Positivity forces the other diagonal entries and their rows and columns to
vanish, leaving only $E_{00}$.

The forward algebra contains $B_1$, hence $E_{00}$ and $E_{11}+E_{22}$ by
spectral interpolation.  The products $B_4B_5$ and $B_5B_4$ give the two
remaining diagonal units.  Together with the four displayed off-diagonal
units these are exactly the seven matrix units in
\eqref{sel2:eq:seven-nine}; this triangular space is closed under
multiplication.  Adjoining the missing adjoints supplies all matrix units.
No faithful stationary state exists, consistently with
Theorem~\ref{sel2:thm:forward-star}.
\end{proof}

The example does not prove absence of every colour corner; some one-way
corners are present.  It proves the precise failure needed here: the locked
normalization does not justify inferring a missing forward reverse from a
formal adjoint.  Appending the original lazy no-response branch also does
not repair recurrence, since $(1-\vartheta)\operatorname{id}+\vartheta\Phi$
has exactly the same stationary densities for $0<\vartheta\le1$.

\subsection{Faithful populations cannot be transferred to the wrong operator level}

\begin{example}[A faithful reset channel whose forward superoperator algebra is not adjoint closed]
\label{sel2:ex:superoperator}
Let $\sigma=\diag(2/3,1/3)$ and $\mathcal J(X)=\Tr(X)\sigma$ on $M_2$,
with complete input/output operator tomography retained.  It is a
trace-preserving channel with faithful stationary density $\sigma$.
Its Kraus-span algebra is $M_2$, in agreement with
Theorem~\ref{sel2:thm:forward-star}.  But $\mathcal J^2=\mathcal J$, so its
forward superoperator algebra is only
\[
 \operatorname{alg}(I_{M_2},\mathcal J)
   =\Span\{I_{M_2},\mathcal J\}.
\]
The Hilbert--Schmidt adjoint is
$\mathcal J^*(X)=\Tr(\sigma X)I$, and is not in this two-dimensional span.
Indeed, off-diagonal inputs force the coefficient of $I_{M_2}$ to vanish in
any putative expression for $\mathcal J^*$, while evaluation on $I$ would
then require $I$ to be proportional to the nonscalar $\sigma$.
Consequently the source's CP-superoperator history construction cannot be
replaced silently by the Kraus-amplitude construction of this note.
\end{example}

\begin{example}[A common incoming mediator without recurrence can be an irreversible sink]
\label{sel2:ex:sink}
On $\C^3$, let $T=\C e_0$, $H=\C e_1$, $E=\C e_2$ and take
$K_0=E_{20},K_1=E_{21},K_2=E_{22}$.  Both forward ports enter the same
one-dimensional mediator, but every state is reset to $E$ after one step.
There is no forward word from $T$ to $H$ or from $H$ to $T$.
Nevertheless the formal product $K_1^*K_0=E_{10}$ is nonzero.
The unique stationary density is $E_{22}$, so recurrence on the exposed
$T\oplus H\oplus E$ carrier fails.  This shows why the new theorem cannot
be obtained from normalization and shared entry supports alone.
\end{example}

\section{What the upstream step has resolved}
\label{sel2:status}

\subsection{A source audit, not a new freely loaded controller}
The mathematical advance is the following chain:
\begin{equation}
 \boxed{\begin{aligned}
 &\text{actual normalized source with faithful recurrence}\\
 &\quad\Longrightarrow\ \text{forward coefficient algebra is adjoint complete}\\
 &\quad\Longrightarrow\ \text{shared forward mediator access has a forward return}\\
 &\quad\Longrightarrow\ \text{an endpoint-certified pure colour branch}.
 \end{aligned}}
 \label{sel2:eq:chain}
\end{equation}
This is a replacement for one independent reverse-operation hypothesis, not
a proof that every normalized source has faithful recurrence or shares the
required mediator.  The old V6--V7 compiler already determines the recurrent
support and tests a positive stationary floor.  Applying it to the \emph{actual}
channel is legitimate; adjoining an artificial reset or reverse control to
force that test to pass would be a different source.

Two specific source audits have now been completed.  The natural absorbing
completion of V1's displayed route instrument has stationary support $H$
only, so the new recurrence theorem cannot be retroactively used on its
entire prepared route.  Independently, the exact six-edge normalization from
the locked-source analysis admits the transient realization
\eqref{sel2:eq:transient-six}.  Thus those identities do not determine the
missing stationary-support row.  Neither conclusion requires evaluating an
unprovided historical matrix entry as zero.

The targeted inspection of the supplied Grand Tensor, the cumulative
V106--V128 internal construction, and the older duality operational notes did
not identify an evaluated historical complete CP panel having both the
specified type/private incoming rows into one mediator and the required
faithful support on all of those cuts.  General process-tensor
reconstruction is a procedure for recovering such a panel when the full
operational table is provided, not the numerical or symbolic panel itself.
No historical colour occurrence is announced on the basis of the new
benchmark.

\subsection{Minimal next audit}
The next computation should start with the actual forward CP branch table,
not with another invented mediator Hamiltonian.  Its cut/controller memory
must be explicit.  On that same table:
\begin{enumerate}[label=\textup{(A\arabic*)},leftmargin=2.8em]
\item use the inherited maximal-stationary-support calculation and retain its
actual source projection; check whether the type, private, and candidate
mediator supports survive;
\item test the two incoming CP supports, and use their covariance or their
actual positive reachable Gram to establish full access to the mediator;
\item compute $W_{p_T}=\sum_{k=0}^{N-1}\Phi^k(p_T)$ and its measured private
Read responses; positive mass returns a literal outcome history and zero mass
returns the invariant separator for this complete grammar;
\item verify repeatability of the already declared endpoint rank-one Reads,
or instead use a directly resolved Choi-rank-one outcome if one is already
supplied; do not refine an environment merely to select a favourable Kraus
coefficient;
\item only after the returned amplitude is identified on the actual final
carrier, apply the inherited private/weak/central assembly and the separate
determinant--matter incidence tests.
\end{enumerate}

If the actual route lives partly outside the recurrent support, the V1 direct
reciprocal-return criterion and the general forward-transfer test remain
available; failure of this sufficient recurrence route is not a general
no-colour theorem.  Likewise, if the endpoint Read is not repeatable, the
positive CP transfer remains proved but its pure amplitude identification is
not supplied by the rank-one effect alone.

\begin{center}
\small
\begin{tabular}{L{.37\textwidth}L{.55\textwidth}}
\toprule
Question & Status after the V2 development\\
\midrule
Does faithful recurrence merely reproduce the old fixed-observable theorem?
& No: Theorem~\ref{sel2:thm:forward-star} identifies the full unital
\emph{forward coefficient algebra} with its adjoint closure.\\[.35em]
Must every inverse Kraus operator be executed separately?
& No: a nonzero formal corner yields an actual forward CP word; no inverse
channel is synthesized.\\[.35em]
Do arbitrary noncommuting mediator controls fit?
& Yes: their actual reachable support is tested by the forward CP Gram,
without a one-Hermitian-generator assumption.\\[.35em]
Must unresolved Kraus labels be promoted to physical records?
& No for the colour return with actual repeatable rank-one endpoints: the
selected composite has Choi rank one intrinsically.\\[.35em]
Does the V1 prepared benchmark already pass the recurrence test?
& No on its full route carrier; its maximal stationary support is $H$.  Its
explicit-exit V1 conclusion is unaffected.\\[.35em]
Do the locked six-edge identities force recurrence?
& No: the exact transient completion has a seven-dimensional forward algebra
and a nine-dimensional formal adjoint closure.\\[.35em]
Does faithful recurrence identify the CP-superoperator and amplitude levels?
& No: the faithful reset example distinguishes them.\\[.35em]
Is historical Standard-Model selection now closed?
& No: the actual common source panel and endpoint admission remain to be
identified; determinant and tensor-type occurrence remain separate.\\
\bottomrule
\end{tabular}
\end{center}

\subsection{Verification and append-only record}
The V1 source is retained byte for byte through its final append marker; the
only removed bytes are its final \verb|\end{document}| so that this V2 block
can be appended.  The cumulative file closes with a new
\verb|\end{document}|.  A verification manifest records the retained-prefix
hash.  V1 remains a historical development block rather than being silently
rewritten to satisfy the recurrence assumptions of V2.

The accompanying exact SymPy checks verify the directed nonreversible
benchmark, its stationary density and forward algebra, the finite transfer
mass, the missing-reverse word identity, the transient six-edge construction
through its orthogonal-frame identities, the sharp stationary-error bound,
the endpoint Choi-rank reduction, the faithful-reset superoperator
counterexample, the irreversible sink, and the V1 cut-population audit.
These are exact regression checks for explicit examples.  The general
statements are established by the written proofs; no proof-assistant
formalization or independent referee verification is claimed.

\begin{thebibliography}{99}
\bibitem[8]{sel2-Carbone}
R.~Carbone and A.~Jen\v{c}ov\'a,
\emph{On period, cycles and fixed points of a quantum channel},
2019, \href{https://arxiv.org/abs/1905.00857}{arXiv:1905.00857}.
Cited for the established faithful-invariant-state channel framework; the
specialized finite forward-word proofs are provided above.
\end{thebibliography}

% END OF SOURCE-SELECTION V2 DEVELOPMENT BLOCK.
% Append the next development here. Preserve all V1 and V2 statements;
% record any necessary correction as an explicit amendment.


\clearpage
\hypersetup{pdftitle={Historical Standard-Model Source Selection: Research Notes V1--V3}}
\begin{center}
{\Large\bfseries V3: Read the actual reset before completing its algebra}\\[.4em]
{\large Regenerative write frames, coherent source domains, and retained provenance}
\end{center}
\addcontentsline{toc}{part}{V3: Regenerative writes and retained provenance}

\section{Source audit and the new question}
\label{sel3:context}

V2 removes an independently executable reverse from a sufficient colour
mechanism.  It does not remove the distinction between an actual CP history,
a coefficient in a Kraus span, and a phase-unique conditional amplitude.  Its
last step uses actual repeatable rank-one endpoint Reads.  The present block
goes upstream of that step, to the concrete completed-private reset already
written in the Grand Tensor, rather than choosing another mediator Hamiltonian.

There are two useful outcomes.  An actual atomic writer with a rank-one input
effect and a pure private output already supplies a coherent typed amplitude:
no terminal private filter is required.  Two such histories with a common
type-line input and spanning private outputs also supply the private matrix
algebra, so that algebra need not be inferred from state support projectors.
Conversely, a rank-one \emph{predictive reset} can have Choi rank greater than
one.  In an exact regenerative grammar, that input-rank obstruction survives
arbitrarily many subsequent resets.  A faithful recurrent window can therefore
have full forward coefficient algebra while none of its positive-length
resolved histories is a single-amplitude operation on the whole window.

\subsection{Audited source rows, and what is not inferred from them}
The following upstream statements are used as supplied inputs.
\begin{center}
\small
\begin{tabular}{L{.37\textwidth}L{.55\textwidth}}
\toprule
Source statement & Its role here\\
\midrule
\cite{Grand}, \src{def:completed-private-packet} and
\src{thm:fresh-endpoint-retained-record}
& The completed hub reset writes the two private pure states with overlap
$1/5$ and a fresh fair hub endpoint.\\[.35em]
\cite{Grand}, \src{thm:atomic-reset-characterization}
& Predictive reset rank one means $X\mapsto\Tr(EX)\omega$; it is not a
statement that the Choi rank is one.\\[.35em]
\cite{Grand}, \src{thm:two-state-active-renewal}
& The active phase matrix and its stationary law are explicit.  They do not
identify the active phases with the colour type/private supports.\\[.35em]
\cite{Attack}, \src{thm:V33-private-qutrit}
& The two private \emph{support projectors} generate $M_2$.  Their physical
admission as repeatable operations is a different question.\\[.35em]
Theorems~\ref{sel2:thm:forward-star}, \ref{sel2:thm:forward-transfer}, and
\ref{sel2:thm:endpoint}
& Recurrence, finite forward transfer, and endpoint purification are inherited,
not reproved as new results.\\
\bottomrule
\end{tabular}
\end{center}

The supplied completed-boundary map has input the hub qubit $U$ and output a
fresh hub together with a private record $R$.  It is not supplied as an
endomorphism of the colour space $T\oplus H$.  Retaining the last record and
retaining one fixed older record are also different observation policies.
We state explicitly the fresh-register renewal composition needed to obtain
a last-record channel.  A single one-boundary map does not, by itself, prove
that composition rule for every possible hidden-memory realization.

The general measure--prepare formalism is established channel theory
\cite{sel3-HSR}.  Nondisturbance of nonorthogonal records is also established
\cite{sel3-KI}.  The contributions here are their source-specific combination:
the all-history input-rank obstruction, the reset-writer source criterion,
the exact completed-private window calculation, and the common-domain test
needed before its write amplitudes can be called colour.  No claim of
literature priority is made for elementary matrix or channel identities.

\section{Atomic histories retain the first input effect}
\label{sel3:atomic}

Use unnormalized matrix traces throughout.  Let an actually admitted CP branch
between two exposed ports be atomic in the supplied sense,
\begin{equation}
 \mathcal J(X)=\Tr(EX)\omega,
 \qquad E\succeq0,\quad\omega\succeq0,\quad\Tr\omega=1.
 \label{sel3:eq:atomic}
\end{equation}
It is a trace-nonincreasing outcome precisely when $E\preceq I$.
We use the input-first Choi convention
$J(\mathcal J)=\sum_{i,j}|i\rangle\langle j|\otimes
\mathcal J(|i\rangle\langle j|)$.
The inherited atomic characterization directly gives
\begin{equation}
 J(\mathcal J)=E^{\mathsf T}\otimes\omega,
 \qquad \rank J(\mathcal J)=\rank E\,\rank\omega.
 \label{sel3:eq:atomic-choi}
\end{equation}
This formula is recalled, not claimed as a new characterization of reset.

\begin{theorem}[All-history input-rank persistence]
\label{sel3:thm:atomic-history}
Let $w=(e_1,\ldots,e_k)$ be an admissible sequence of atomic outcomes, with
$\mathcal J_{e_j}(X)=\Tr(E_{e_j}X)\omega_{e_j}$ and matching physical cuts.
Set
\begin{equation}
 a_w=\prod_{j=2}^k\Tr(E_{e_j}\omega_{e_{j-1}}),
 \quad a_{(e_1)}=1.
 \label{sel3:eq:atomic-weight}
\end{equation}
Then
\begin{equation}
 \boxed{\mathcal J_w(X)=a_w\Tr(E_{e_1}X)\omega_{e_k}.}
 \label{sel3:eq:atomic-word}
\end{equation}
If $a_w>0$ and $E_{e_1}\ne0$, its Choi rank is
$\rank E_{e_1}\,\rank\omega_{e_k}$.  In particular, no length of
subsequent atomic continuation removes the first effect's rank.
More generally, for any admitted terminal CP operation $\mathcal F$, every
nonzero $\mathcal F\mathcal J_w$ has Choi rank at least
$\rank E_{e_1}$.  For an admitted initial CP operation $\mathcal L$,
\begin{equation}
 J(\mathcal F\mathcal J_w\mathcal L)
   =a_w\,[\mathcal L^*(E_{e_1})]^{\mathsf T}
                     \otimes\mathcal F(\omega_{e_k}).
 \label{sel3:eq:atomic-prepost}
\end{equation}
Consequently a nonzero single-amplitude composite requires and is implied by
rank one of both factors in \eqref{sel3:eq:atomic-prepost}.
\end{theorem}
\begin{proof}
Applying $\mathcal J_{e_{j+1}}$ to a multiple of $\omega_{e_j}$ multiplies
that multiple by $\Tr(E_{e_{j+1}}\omega_{e_j})$ and replaces the output by
$\omega_{e_{j+1}}$.  Induction proves \eqref{sel3:eq:atomic-word}.
The Choi formula then follows from \eqref{sel3:eq:atomic-choi}.
Postprocessing changes only the second tensor factor, so cannot lower the
nonzero first factor's rank.  Preprocessing replaces the input functional by
$\Tr[E_{e_1}\mathcal L(X)]=\Tr[\mathcal L^*(E_{e_1})X]$.
The rank of a tensor product is the product of ranks.  A nonzero positive
Choi matrix has rank one exactly when the CP map has one coefficient, unique
up to its phase.  This proves the last assertion.
\end{proof}

\begin{remark}[The obstruction is at the input, not an unresolved phase]
A pure written output does not repair an input effect of rank two or more.
Nor does selecting a later private output ray: that is a terminal operation
in \eqref{sel3:eq:atomic-prepost}.  Selecting a favourable Kraus vector of
$E^{\mathsf T}\otimes\omega$ would add a finer environment outcome not
present in the stipulated instrument.  Conversely, an actual initial filter
can remove the obstruction if its pulled-back effect has rank one.  The
criterion is a condition on the complete selected CP map, not a prohibition
against physically supplied filters.
\end{remark}

\begin{proposition}[Orthogonal-source alternatives and common-source amplitudes]
\label{sel3:prop:domain}
The statement that several output states are pure does not identify their
input ports with one common coherent line.  In particular, let
$U_0=\C u_0$, $U_1=\C u_1$ be orthogonal input lines and let $H=\C^2$
contain independent unit vectors $r_0,r_1$ with nonzero overlap.  The two
actual arrows
\[
 b_i=|r_i\rangle\langle u_i|:U_i\to H
\]
on $U_0\oplus U_1\oplus H$ have forward algebra
$\Span\{I,b_0,b_1\}$ when no other nonidentity operations are supplied,
because every product of two arrows is zero.  Their formal adjoint closure
is $M_4(\C)$, not the desired $M_3$ on a common $T\oplus H$.
\end{proposition}
\begin{proof}
Each arrow vanishes on $H$, which proves the forward product claim.
In the adjoint closure,
$b_i^*b_i=|u_i\rangle\langle u_i|$ and
$b_0^*b_1=\langle r_0,r_1\rangle|u_0\rangle\langle u_1|$.
Hence the full matrix algebra on $U_0\oplus U_1$ occurs.  Multiplying the
arrows by its matrix units supplies every map from that input plane to
$\Span\{r_0,r_1\}=H$; adjoints and products give the remaining blocks.
These are all matrix units on the four-dimensional carrier.
\end{proof}

A common scalar dimension is therefore not a common source-domain
identification.  This is relevant to the hub labels of the completed-private
writer: neither hub input line is silently renamed the categorical colour
type contrast.

\section{Recurrence reconstructed from a regenerative write graph}
\label{sel3:regenerative}

Consider a finite typed carrier $\mathcal F=\bigoplus_{a\in V}F_a$ with
$d_a=\dim F_a$.  Each actually resolved edge $e:a\to b$ is a scalar-effect
reset,
\begin{equation}
 \mathcal J_e(X)=p_e\Tr(X)\omega_e,
 \qquad p_e>0,\quad\Tr\omega_e=1,\quad
 \supp\omega_e\subseteq F_b,
 \qquad\sum_{s(e)=a}p_e=1.
 \label{sel3:eq:reset-graph}
\end{equation}
Here $X\in\B(F_a)$, and the CP extension to the full Hilbert carrier uses
its actual $b\leftarrow a$ block and annihilates other input blocks.
The last identity is verified operator normalization.  A scalar probability
measured on just one input state does not establish the effect $p_eI_{F_a}$.
Let
\[
 P_{ab}=\sum_{e:a\to b}p_e,
 \qquad (\Phi X)_b=\sum_{e:a\to b}p_e\Tr(X_a)\omega_e.
\]
All controller memory relevant to edge admission is retained in the vertices.

\begin{theorem}[Explicit stationary support of a regenerative writer]
\label{sel3:thm:reset-stationary}
Stationary densities of $\Phi$ are in one-to-one affine correspondence with
stationary probability rows $\pi=\pi P$.  The corresponding density is
\begin{equation}
 \boxed{\sigma_b=\sum_{e:a\to b}\pi_a p_e\omega_e.}
 \label{sel3:eq:reset-sigma}
\end{equation}
Let $C$ range over the closed communicating classes of the finite graph and
let $\pi^C$ be the strictly positive stationary probability on $C$.  The
maximal stationary Hilbert support is the direct sum over those classes of
\begin{equation}
 \bigoplus_{b\in C}
 \Span\{\Ran\omega_e:e:a\to b,\ a\in C\}.
 \label{sel3:eq:reset-support}
\end{equation}
In particular, on an irreducible graph a faithful stationary density exists
if and only if the incoming write states span each target Hilbert space.
On that branch the density is unique, and its faithfulness floor is exactly
\begin{equation}
 \min_b\lambda_{\min}
       \left(\sum_{e:a\to b}\pi_a p_e\omega_e\right)>0.
 \label{sel3:eq:reset-floor}
\end{equation}
There is no assumption of detailed balance or of a physically executable
reverse edge.
\end{theorem}
\begin{proof}
Every output of the full extension is block diagonal.  Taking block traces
of its fixed-point equation gives $\pi_b=\sum_a\pi_aP_{ab}$.  Once the
traces are fixed, the same equation gives \eqref{sel3:eq:reset-sigma};
conversely its block traces are $\pi_b$, so that formula is stationary.
Finite Markov stationarity is supported on closed communicating classes.
Each such class has a stationary row positive at every vertex, whereas
transient vertices have zero weight in every stationary row.  These facts
follow, for example, by decomposing the graph into strongly connected
components: any stationary mass flowing out of a nonclosed component forces
an incoming component, and following the finite condensation graph leaves
stationary mass only on closed components.  On a closed irreducible finite
class a stationary row exists by Cesaro averaging, its support is a nonempty
closed subset and hence the full class, and uniqueness is the usual maximum
principle for the finite irreducible transition matrix.

The support of a finite positive sum is the span of the supports of its
positive-weight summands.  Apply this to \eqref{sel3:eq:reset-sigma} with
$\pi^C$.  It proves \eqref{sel3:eq:reset-support}, the faithfulness
criterion, and the displayed eigenvalue floor.
\end{proof}

\begin{corollary}[A finite source criterion that supplies V2's recurrence]
\label{sel3:cor:reset-recurrence}
In the irreducible, full-incoming-frame branch of
Theorem~\ref{sel3:thm:reset-stationary}, the actual forward coefficient
algebra is $M_N(\C)$ on $\mathcal F$, where $N=\sum_a d_a$.
In particular, its adjoint completeness is a conclusion of graph and
write-frame data rather than an additional recurrence premise.
\end{corollary}
\begin{proof}
The preceding theorem supplies a faithful stationary state, so
Theorem~\ref{sel2:thm:forward-star} makes the forward coefficient algebra
a represented $C^*$-algebra.  If this algebra were not $M_N$, its represented
Wedderburn decomposition would contain at least two irreducible invariant
subspaces (counting equivalent copies separately).  Trace preservation on
each such reducing subspace gives a stationary density supported there by
Cesaro averaging.  Their supports can be chosen orthogonal, contradicting
uniqueness in Theorem~\ref{sel3:thm:reset-stationary}.  Hence the algebra is
$M_N$.  This proof does not identify the algebra with a physically executable
set of all unitary controls.
\end{proof}

\begin{proposition}[Recurrence does not purify regenerative histories]
\label{sel3:prop:reset-no-purification}
For a path $w=(e_1,\ldots,e_k)$ in \eqref{sel3:eq:reset-graph},
\begin{equation}
 \mathcal J_w(X)=\left(\prod_{j=1}^kp_{e_j}\right)
                    \Tr(X)\omega_{e_k}.
 \label{sel3:eq:scalar-reset-word}
\end{equation}
Its Choi rank on the original source vertex is
$d_{s(e_1)}\rank\omega_{e_k}$, independent of the number of intermediate
resets.  Thus even on the faithful recurrent branch there is no nonzero
single-amplitude history from a source vertex of dimension greater than one,
unless an additional input-selective operation is admitted before the first
reset.  Terminal selection alone cannot remedy this.
\end{proposition}
\begin{proof}
Here each transition overlap in \eqref{sel3:eq:atomic-weight} is the next
$p_e$, because its effect is scalar on the whole matching source space.
Apply Theorem~\ref{sel3:thm:atomic-history}.  The first effect is a nonzero
scalar multiple of the identity on that original vertex, so its rank is
$d_{s(e_1)}$.  The postprocessing assertion is the same theorem.
\end{proof}

\section{Pure writes can replace the private-filter and private-algebra assumptions}
\label{sel3:writes}

The preceding obstruction has a useful positive converse on a genuine
one-dimensional type port.  This does not require an extra physical
measurement of that port: it requires that the actual incoming CP branch
has that \emph{exact input support} on the exposed carrier.

\begin{theorem}[Source-native pure-write amplitude]
\label{sel3:thm:pure-write}
Let $T=\C t$ and $H$ be orthogonal, physically identified subspaces of one
exposed final-cut carrier.  Suppose an admitted history begins with an atomic
branch whose actual effect is $\alpha p_T$, $\alpha>0$, and continues by
atomic branches whose transition factors in \eqref{sel3:eq:atomic-weight}
are positive.  Suppose its last branch writes the pure state
$|h\rangle\langle h|$ on $H$, with $\norm h=1$.  Then the full selected
history, extended by zero on the other source supports, is
\begin{equation}
 \boxed{\mathcal J_w(X)=b_wXb_w^*,\qquad
 b_w=\sqrt{\alpha a_w}\,|h\rangle\langle t|.}
 \label{sel3:eq:write-amplitude}
\end{equation}
The amplitude is intrinsic up to a single phase.  No reverse route,
faithful-recurrence premise, terminal private Read, or resolution of hidden
Kraus labels is needed for this conclusion.
\end{theorem}
\begin{proof}
The first effect has rank one and the last output is pure.  Formula
\eqref{sel3:eq:atomic-word} therefore has the displayed coefficient, and
its positive Choi matrix has rank one.  Its unique Choi range determines
that coefficient line.  All products used are admitted histories; their
positive probability factors are not chosen coherent superpositions.
\end{proof}

\begin{corollary}[Full colour from a spanning bank of common-origin writes]
\label{sel3:cor:write-colour}
Suppose the actually identified histories of
Theorem~\ref{sel3:thm:pure-write} have the same input line $T$ and output
vectors $h_j\in H$, with weights $c_j>0$.  Define their phase-independent
output frame
\begin{equation}
 F_H=\sum_j\mathcal J_{w_j}(p_T)
     =\sum_j c_j|h_j\rangle\langle h_j|.
 \label{sel3:eq:write-frame}
\end{equation}
Let $H_0=\Ran F_H$ and $k=\dim H_0$.  For a nonempty bank,
\begin{equation}
 C^*(I,p_T,b_j:j)
 =\B(T\oplus H_0)\oplus\C I_{H\ominus H_0},
 \label{sel3:eq:write-algebra}
\end{equation}
with the last summand omitted if $H_0=H$.  In particular, for $\dim H=2$,
\begin{equation}
 \boxed{F_H\succ0\ \Longleftrightarrow\
        C^*(I,p_T,b_j:j)=M_3(\C).}
 \label{sel3:eq:write-colour}
\end{equation}
The private $M_2(H)$ is then generated by products of actual amplitude lines;
it is not assumed from the mathematical support projectors of prepared states.
Two independent output rays are sufficient and, absent other operations on
$H$, necessary.
\end{corollary}
\begin{proof}
The amplitude span contains every $|v\rangle\langle t|$, $v\in H_0$.
Its products with adjoints contain every $|v\rangle\langle w|$ on $H_0$,
as well as $p_T$.  These are all matrix units on $T\oplus H_0$.
Every nonidentity generator vanishes on $H\ominus H_0$, so only the scalar
identity remains there.  This is the matrix-unit step of the inherited
bridge/router theorem, applied to source-certified writes.  It proves
\eqref{sel3:eq:write-algebra} and the rank criterion.  One ray cannot span
a plane; two independent rays do.
\end{proof}

\begin{proposition}[The calibrated private overlap gives an exact write floor]
\label{sel3:prop:private-floor}
For the private pair with $|\langle r_0,r_1\rangle|=1/5$ and two actual
common-origin write masses $c_0,c_1>0$, the frame has
\begin{align}
 \det F_H&=\frac{24}{25}c_0c_1,\
 \lambda_{\min}(F_H)
 &=\frac{c_0+c_1-
       \sqrt{(c_0-c_1)^2+4c_0c_1/25}}{2}>0.
 \label{sel3:eq:private-floor}
\end{align}
For $c_0=c_1=1/2$, its eigenvalues are $3/5$ and $2/5$.
For an evaluated Hermitian frame $\widehat F_H$ with
$\norm{\widehat F_H-F_H}_{\op}\le\epsilon$, the strict inequality
$\lambda_{\min}(\widehat F_H)>\epsilon$ certifies the colour conclusion.
\end{proposition}
\begin{proof}
The trace is $c_0+c_1$ and the determinant of the two-vector Gram is
$c_0c_1(1-|\langle r_0,r_1\rangle|^2)$.  Its two nonzero eigenvalues are
those of $F_H$, yielding the formula.  The last assertion is the min--max
bound $\lambda_{\min}(F_H)\ge
\lambda_{\min}(\widehat F_H)-\epsilon$.
\end{proof}

\begin{remark}[Exact scope of the assumption reduction]
The actual branch data must certify the input effect, pure output, common
physical $T\oplus H$ identification, and positive path factors.  A source
state $r_i$ alone certifies none of its incoming effect.  Nor do two unrelated
hub origins satisfy the common-origin condition; see
Proposition~\ref{sel3:prop:domain}.  The conclusion concerns the
$C^*$-algebra of those physical conditional amplitudes, not coherent
execution of every unitary in $U(3)$ and not selection of its matter tensor
types or determinant relation.
\end{remark}

\begin{example}[Normalized two-write realization with no private Read]
\label{sel3:ex:two-write}
On $T\oplus H=\C t\oplus\C^2$, use actual resolved write branches
\[
 \mathcal J_i(X)=\tfrac12\langle t,Xt\rangle
                         |r_i\rangle\langle r_i|,
 \qquad i=0,1,
\]
and one unresolved return branch
$\mathcal R(X)=\Tr(p_HX)p_T$.  Their effects sum to $p_T+p_H=I$.
The two writes are already single-amplitude branches; the return has Choi
rank two and is not refined.  The state
\begin{equation}
 \sigma=\tfrac12p_T+\tfrac14
             (|r_0\rangle\langle r_0|+|r_1\rangle\langle r_1|)
 \label{sel3:eq:two-write-state}
\end{equation}
is the unique stationary density and has eigenvalues
$1/2,3/10,1/5$.  Both colour and faithful recurrence follow without an
actual private projective Read.  This is an explicit normalized realization
of the sufficient write criterion, not an evaluation of a historical
categorical type port.
\end{example}
\begin{proof}
The classical graph alternates between $T$ and $H$, with the two $T\to H$
outcomes weighted $1/2$ and the $H\to T$ outcome weighted one.  Its stationary
populations are $(1/2,1/2)$; Theorem~\ref{sel3:thm:reset-stationary} gives
\eqref{sel3:eq:two-write-state}.  The private mean state has eigenvalues
$3/5,2/5$, proving the spectrum.  The output frame of the two resolved
writes is that same positive private mean; apply
Corollary~\ref{sel3:cor:write-colour}.  This periodic graph is recurrent,
not a claim of convergence to stationarity at every integer time.
\end{proof}

\section{Evaluate the supplied completed-private writer on a declared record window}
\label{sel3:window}

\subsection{One boundary, one window, and one older record are different objects}
The source supplies the completed-boundary channel
\begin{equation}
 \mathcal R(X)=\frac{I_U}{2}\otimes
      \left(\langle0,X0\rangle\rho_0+
            \langle1,X1\rangle\rho_1\right),
 \quad U=\C^2,\quad \rho_x=|r_x\rangle\langle r_x|,
 \quad\langle r_0,r_1\rangle=\frac15.
 \label{sel3:eq:source-reset}
\end{equation}
It is the source's
\src{thm:fresh-endpoint-retained-record}.  Here the input is the hub $U$;
the output is a new hub and a two-dimensional written record $R$.
The private output states span $R$.  The input of this displayed map is not
$T\oplus H$ and not $U\otimes R$.

To test a natural candidate for V2's recurrent panel, declare the following
\emph{fresh-register composition rule}.  At every completed boundary the
same map \eqref{sel3:eq:source-reset} acts on the hub and a fresh record is
written; previously written records have no interaction with this operation.
An observer retaining only the hub and the most recently written record
traces out the older records.  This is a specified observation window of that
writer realization.  It is not an extra reset applied to a fixed older record.
The one-boundary source formula alone does not establish this composition
rule in a realization with additional unobserved record interactions.

\begin{proposition}[Exact closure of the fresh last-record window]
\label{sel3:prop:window-lift}
Under the declared fresh-register composition, the closed channel on the
last-record window $F=U\otimes R$ is
\begin{equation}
 \Psi(X)=\mathcal R(\Tr_R X)
 =\frac{I_U}{2}\otimes
    \sum_{x=0}^1\Tr[(|x\rangle\langle x|\otimes I_R)X]\rho_x.
 \label{sel3:eq:window-channel}
\end{equation}
This holds even for correlations between the hub and older records.  It is
therefore an exact reduced window dynamics of the declared realization, not
a fitted channel with a selected stationary state.
\end{proposition}
\begin{proof}
The full boundary applies $\mathcal R\otimes\id$ to the hub and the
retained old records.  Taking the partial trace of those old records commutes
with a map acting only on the hub.  Thus the reduced new hub--record state is
$\mathcal R(\Tr_{\rm old}X)$.  Keeping just the immediately preceding
record in the input gives \eqref{sel3:eq:window-channel}.  Linearity of the
partial trace proves the identity on correlated inputs as well as products.
\end{proof}

\begin{theorem}[An exact two-step recurrent window]
\label{sel3:thm:window-recurrence}
Put
\begin{equation}
 \bar\rho=\tfrac12(\rho_0+\rho_1),\qquad
 \sigma_*=\tfrac12 I_U\otimes\bar\rho.
 \label{sel3:eq:window-stationary}
\end{equation}
Then on the whole $M_4$ window,
\begin{equation}
 \boxed{\Psi^2(X)=\Tr(X)\sigma_*,\qquad
        \Psi^3=\Psi^2,\qquad \Psi(\sigma_*)=\sigma_*.}
 \label{sel3:eq:two-step-window}
\end{equation}
The stationary density is unique and faithful, with eigenvalues
\begin{equation}
 (3/10,3/10,1/5,1/5).
 \label{sel3:eq:window-floor}
\end{equation}
The unital algebra of its forward Kraus coefficients is $M_4$ and its
commutant is scalar.  However, as a CP superoperator, $\Psi$ has ordinary
linear rank two; its minimal polynomial is $z^2(z-1)$ and
$\operatorname{alg}(I,\Psi)$ has dimension three.  Its one-step and two-step
Choi ranks are respectively eight and sixteen.
\end{theorem}
\begin{proof}
The partial trace over the just-written record is
$\Tr_R\Psi(X)=\Tr(X)I_U/2$.  Applying \eqref{sel3:eq:source-reset} again
proves \eqref{sel3:eq:two-step-window}.  Every stationary state is its own
second iterate, so must equal $\sigma_*$.  Since the two unit vectors have
overlap $1/5$, $\bar\rho$ has eigenvalues $3/5,2/5$, giving
\eqref{sel3:eq:window-floor}.

The output window's hub values may be used as the two vertices of the
regenerative writer of Section~\ref{sel3:regenerative}, each with private
fibre dimension two.  The edge $x\to h$ has weight $1/2$ and writes
$\rho_x$.  Its graph is complete and each incoming private frame spans $R$.
Corollary~\ref{sel3:cor:reset-recurrence} gives $M_4$; alternatively use
V2's forward-adjoint theorem and the unique faithful stationary state.

The image of $\Psi$ is the span of $I_U\otimes\rho_0$ and
$I_U\otimes\rho_1$, which are independent, so its rank is two.  The
operator $X=(|0\rangle\langle0|-|1\rangle\langle1|)\otimes\tau_R$,
for any unit-trace $\tau_R$, has
$\Psi(X)=I_U\otimes(\rho_0-\rho_1)/2\ne0$ and $\Psi^2(X)=0$.
There is therefore a length-two zero-eigenvalue Jordan chain.  There is also
the eigenvalue one on $\sigma_*$.  Together with
$\Psi^2(\Psi-I)=0$ these facts give the claimed minimal polynomial and
three-dimensional generated superoperator algebra.

The Choi matrix of \eqref{sel3:eq:window-channel} is a sum of two positive
terms with orthogonal input-hub supports:
\[
 J(\Psi)=\sum_{x=0}^1
 (|x\rangle\langle x|\otimes I_R)^{\mathsf T}
             \otimes\left(\tfrac12I_U\otimes\rho_x\right).
\]
Each term has rank $2\cdot2=4$, hence total rank eight.  The two-step reset
has Choi matrix $I_F\otimes\sigma_*$, of rank sixteen.
\end{proof}

\begin{remark}[The first actual recurrence calculation is not a colour identification]
Theorem~\ref{sel3:thm:window-recurrence} uses the source's displayed overlap,
fair endpoint, and reset map; the fresh-register composition rule is stated
separately.  It supplies an exact recurrent \emph{window} candidate without
altering the displayed one-boundary law.  It does not identify its two hub
values with the categorical type line and a private plane, or convert its
$M_4$ amplitude algebra into $M_3$.  The stationary floor is measured per
completed boundary.  No claim about a primitive-tick or continuum relaxation
rate is obtained by dropping the actual renewal durations.
\end{remark}

\subsection{Complete hub histories do not purify the lost private input}
To give the window the most informative hub record, allow both the current
hub input $x$ and the newly written hub $h$ to be resolved.  These are the
atomic branches
\begin{equation}
 \mathcal J_{hx}(X)=\tfrac12
       \Tr[(|x\rangle\langle x|\otimes I_R)X]
       (|h\rangle\langle h|\otimes\rho_x).
 \label{sel3:eq:hub-refinement}
\end{equation}
Their sum is $\Psi$.  This refinement is an actual subprotocol only when
those hub labels are physically available on the same history; it is also a
valid mathematical refinement for a no-purification upper-bound argument.

\begin{theorem}[All-length Choi-rank obstruction on the recurrent window]
\label{sel3:thm:window-rank}
Consider a word $w=((h_1,x_1),\ldots,(h_k,x_k))$ of positive length.
The branch vanishes if $x_{j+1}\ne h_j$ at any intermediate step.
Otherwise its exact action is
\begin{equation}
 \boxed{\mathcal J_w(X)=2^{-k}
       \Tr[(|x_1\rangle\langle x_1|\otimes I_R)X]
       (|h_k\rangle\langle h_k|\otimes\rho_{x_k}).}
 \label{sel3:eq:window-word}
\end{equation}
Every such nonzero branch has ordinary linear rank one and Choi rank exactly
two on the original window input.  Every nonempty coarse graining of these
positive-length histories has Choi rank at least two.  No terminal CP
postprocessing after the first reset can produce a nonzero Choi-rank-one
window history.  Initial input selection can do so only by reducing the
pulled-back first effect to rank one, as in
Theorem~\ref{sel3:thm:atomic-history}.
\end{theorem}
\begin{proof}
The effect of the next branch on the preceding pure output is
$\frac12\delta_{x_{j+1},h_j}\Tr\rho_{x_j}
=\frac12\delta_{x_{j+1},h_j}$.  The product law gives
\eqref{sel3:eq:window-word}.  Its first effect has rank two and its final
state is pure, so its Choi rank is two.  For positive matrices,
$\Ker(\sum_iA_i)=\bigcap_i\Ker A_i$; therefore a nonempty sum cannot have
rank smaller than an included nonzero summand.  This proves the
coarse-graining assertion.  Terminal postprocessing preserves the first
rank-two factor by \eqref{sel3:eq:atomic-prepost}.  This conclusion covers
arbitrary CP processing after that first discard, not just further reset
steps.
\end{proof}

\begin{remark}[Pure on a smaller preparation domain is a different statement]
If the old private input is restricted to one \emph{actually prepared} ray,
the domain is smaller and a selected write can have Choi rank one there.
This is not Choi rank one of the operation on $F$.  Replacing the original
window input by a rank-one code without recording that preparation would
silently remove the lost private information from the question.  Likewise,
at the very first boundary the source map has hub-only input; selecting its
old hub value and fresh endpoint gives a pure write on that hub line.  The
common-domain test in Proposition~\ref{sel3:prop:domain} is still needed
before such a line is identified with the colour type source.
\end{remark}

\section{A retained private record cannot be replaced by the newest-record slot}
\label{sel3:retained}

The latest-record window and a chosen older record can both be retained in a
finite experiment.  Under the same explicitly archive-preserving composition
rule, let $R_{\rm old}$ denote one such fixed record.  The observed map is
\begin{equation}
 \widetilde\Psi=\Psi\otimes\id_{\B(R_{\rm old})}.
 \label{sel3:eq:archived-channel}
\end{equation}
No-signalling from a private record to the active hub is weaker than this
identity action on that record; the source's more general sealed-provenance
extensions allow other private propagation.  Equation~\eqref{sel3:eq:archived-channel}
is used only for the stated archive-preserving realization.

\begin{proposition}[Exact algebra and stationary states on the archived realization]
\label{sel3:prop:archive}
If $\dim R_{\rm old}=r$, the forward coefficient algebra and its commutant are
\begin{equation}
 \mathcal A_+(\widetilde\Psi)=M_4\otimes I_r,
 \qquad
 \mathcal A_+(\widetilde\Psi)'=I_4\otimes M_r.
 \label{sel3:eq:archive-algebras}
\end{equation}
All stationary states are exactly $\sigma_*\otimes\eta$, where $\eta$ is
any density on $R_{\rm old}$.  Thus a faithful stationary state on the joint
carrier exists, but there is still no executed nonscalar operator on the
old record in the coefficient algebra.  In particular,
$I_4\otimes|r_i\rangle\langle r_i|$ does not belong to that algebra for
an old private plane of dimension two.
\end{proposition}
\begin{proof}
The coefficients are precisely $K_j\otimes I_r$ for any Kraus list of
$\Psi$.  Their forward algebra is the first algebra in
\eqref{sel3:eq:archive-algebras} by
Theorem~\ref{sel3:thm:window-recurrence}; taking its commutant gives the
second.  Since $\Psi^2$ is the full reset to $\sigma_*$,
\[
 (\Psi^2\otimes\id)(X)=\sigma_*\otimes\Tr_F X.
\]
Every fixed state therefore has the asserted product form, and every such
product is fixed.  Finally, the only elements of $M_4\otimes I_r$ which
also have the form $I_4\otimes B$ have scalar $B$.  The private support
projectors are nonscalar.
\end{proof}

\begin{remark}[Relation to the physical duality statement]
The old record contributes a \emph{multiplicity commutant}; the newest-record
window contributes the forward acted-on algebra.  These roles are not
interchangeable.  In particular, the manuscript's mutual-centralizer theorem
classifies the action and its opposite multiplicity algebra on a specified
carrier; it does not declare every mathematical element of the latter to be
an independently executed operation.  V1's right-module and action-map audit
is relevant at exactly this point.  A terminal probability Read of the old
record need not be a composable system operation back into that same record.
\end{remark}

\section{What nondisturbing access to the private pair would imply}
\label{sel3:nondisturbance}

The Grand Tensor's adjective ``passive'' means that reading a boundary mark
does not trigger or change that boundary event.  It is not a general assertion
that every quantum state in a private record is left unchanged by every
possible measurement.  The following statement tests the stronger
nondisturbance interpretation only.  Its general setting belongs to the
Koashi--Imoto theory \cite{sel3-KI}; a direct proof of the pure-state case is
short enough to include.

\begin{proposition}[Connected nonorthogonal records admit only scalar nondisturbing actions]
\label{sel3:prop:nondisturbance}
Let normalized pure states $r_i$ span a Hilbert space $R$.  Join $i,j$ when
$\langle r_i,r_j\rangle\ne0$, and suppose the resulting graph is connected.
Let an actual normalized instrument $(\mathcal I_a)_a$ have nonselective
map $\Lambda$ with
\begin{equation}
 \Lambda(|r_i\rangle\langle r_i|)=|r_i\rangle\langle r_i|
 \quad\text{for every }i.
 \label{sel3:eq:preserve-records}
\end{equation}
Then, on $\B(R)$,
\begin{equation}
 \boxed{\mathcal I_a(X)=p_aX,\qquad p_a\ge0,\quad\sum_ap_a=1.}
 \label{sel3:eq:trivial-instrument}
\end{equation}
The outcome probabilities contain no information distinguishing the input
$r_i$, and the coefficient algebra on $R$ is scalar.  Thus prepared
nonorthogonal states and state-preserving access do not supply their
noncommuting projectors as actual filter operations.
\end{proposition}
\begin{proof}
Choose Kraus operators $L_{a\mu}$ and combine $(a,\mu)$ into an index $j$.
Positivity and the pure output in \eqref{sel3:eq:preserve-records} imply
$L_jr_i=c_{ji}r_i$.  Trace preservation gives
$\sum_j|c_{ji}|^2=1$ and, for two states,
\[
 \langle r_i,r_k\rangle
 =\sum_j\langle L_jr_i,L_jr_k\rangle
 =\langle r_i,r_k\rangle\sum_j\overline{c_{ji}}c_{jk}.
\]
When the overlap is nonzero, the two unit vectors $(c_{ji})_j$ and
$(c_{jk})_j$ have inner product one and are equal.  Connectivity makes these
coefficient vectors the same for all $i$.  Since the states span $R$,
$L_j=c_jI_R$.  Grouping by the actual outcome $a$ gives
\eqref{sel3:eq:trivial-instrument} with
$p_a=\sum_\mu|c_{a\mu}|^2$.
\end{proof}

\begin{proposition}[Exact disturbance of a private repeatable binary Read]
\label{sel3:prop:read-disturbance}
On the private plane with $|\langle r_0,r_1\rangle|=1/5$, suppose an
actual binary instrument has effects $p_0=|r_0\rangle\langle r_0|$ and
$I_R-p_0$, and both outcomes are repeatable.  Its nonselective map is the
L\"uders pinching
\[
 \mathcal D_0(X)=p_0Xp_0+(I_R-p_0)X(I_R-p_0).
\]
It changes the other private state by the exact amounts
\begin{equation}
 1-\Tr[p_1\mathcal D_0(p_1)]
   =\frac{48}{625}=\norm{[p_0,p_1]}_{\HS}^2,
 \qquad
 \norm{\mathcal D_0(p_1)-p_1}_1=\frac{4\sqrt6}{25}.
 \label{sel3:eq:read-disturbance}
\end{equation}
In particular, the source's nonzero private projector commutator is also an
explicit obstruction to interpreting this Read as nondisturbing on both
prepared private states.
\end{proposition}
\begin{proof}
Both effects are rank one on the plane.  The inherited
Lemma~\ref{sel2:lem:read} makes each repeated outcome its L\"uders filter.
Choose a basis with $r_0=(1,0)$ and, after phases,
$r_1=(c,\sqrt{1-c^2})$, $c=1/5$.  Pinching removes the two off-diagonal
entries $c\sqrt{1-c^2}$ of $p_1$.  Its overlap with $p_1$ is
$c^4+(1-c^2)^2$, giving the loss $2c^2(1-c^2)=48/625$.
The difference has eigenvalues $\pm c\sqrt{1-c^2}$, giving the trace norm.
Direct multiplication gives the same $2c^2(1-c^2)$ for the commutator norm.
\end{proof}

\begin{remark}[No prohibition on informative private experiments]
The last two propositions do not show that a private Read is forbidden, and
passive no-feedback access need not satisfy
\eqref{sel3:eq:preserve-records}.  An actual measurement can alter a private
state while leaving all active hub statistics unchanged.  The propositions
identify the finite response that must be retained if a repeatable filter is
used in V2.  In particular, the optimal distinguishability of the two prepared
states does not itself certify that a particular repeatable filter belongs
to the source's admitted operation alphabet.
\end{remark}

\section{Source-level conclusion and the next noncircular input}
\label{sel3:status}

\subsection{A positive mechanism and a completed source audit}
Two mechanisms are now distinct and available.
On a recurrent common-mediator source, V2 turns formal corners into forward
CP histories, with actual endpoint filters supplying their pure amplitudes.
On an atomic-write source, Theorem~\ref{sel3:thm:pure-write} instead uses the
first input effect and last output state.  A pure write from a physically
identified type line is already a single amplitude.  Two independent such
writes give the full private plane and colour algebra through
Corollary~\ref{sel3:cor:write-colour}, without an assumed private Read,
preloaded private matrix algebra, or independent reverse.

The supplied completed-private reset has been evaluated rather than replaced
by an unspecified example.  Under the explicitly declared fresh-register
renewal composition, its newest-record window has a faithful stationary floor
$1/5$, exact two-step reset, and full $M_4$ forward coefficient algebra.
That statement is not the historical colour conclusion.  The source input is
a hub qubit, not the checked common colour type line; all positive-length
resolved histories on the full window retain input effect rank two; and one
fixed older record, if archived unchanged, remains a multiplicity factor
rather than an acted-on private algebra.  If the historical multitime source
has a different record propagation rule, the displayed one-boundary map does
not decide that rule and the window calculation is a stated conditional
realization, not a numerical certificate of it.

The precise source-domain distinction is consistent with the finite joint
packet definition in the duality manuscript: mixed operators, their cuts,
and their actual common carrier must be supplied together.  A scalar
preparation domain is not automatically the categorical type contrast, just
as a newest-register identification is not an operation on a retained older
register.  No theorem here assigns missing historical entries the value zero.

\begin{center}
\small
\begin{tabular}{L{.35\textwidth}L{.57\textwidth}}
\toprule
Question & Status after V3\\
\midrule
Does predictive reset rank one imply one amplitude?
& No. The inherited Choi factorization and the all-history theorem retain the
first effect's rank.\\[.3em]
Can a physical private filter be avoided?
& Yes on an actual common-origin pure-write branch; the input effect and output
purity replace it, not an inferred Kraus refinement.\\[.3em]
Can the private $M_2$ also be derived?
& Yes from two source-certified common-origin writes spanning the plane.\\[.3em]
Can recurrence be reconstructed without the generic stationary SDP?
& Yes for a complete scalar-effect regenerative writer: its classical recurrent
classes and incoming write frames give the exact support and state.\\[.3em]
Does the source reset have a faithful finite realization?
& Its declared fresh last-record window does, with exact floor $1/5$; this is
not a claim about every historical multitime completion.\\[.3em]
Does that window already contain a pure coherent history on its whole input?
& No positive-length hub-resolved word does; every one has Choi rank two,
regardless of recurrence and algebra saturation.\\[.3em]
Does retaining two private states execute their projectors?
& No. Archive-preserving evolution acts trivially on a fixed old record;
nondisturbing instruments on the nonorthogonal pair are scalar.\\[.3em]
Is the complete historical SM source now selected?
& Not established. A common colour-domain write panel or the V2 admitted
Read/mediator panel is still required; determinant and matter types remain
separate.\\
\bottomrule
\end{tabular}
\end{center}

\subsection{The minimal next calculation}
The next source audit should seek the earliest actual atomic branch whose
input effect survives, on the full geometry-shorted system, as a positive
multiple of the categorical type projection.  Its effect, output density,
physical input/output cut, and composability labels must come from the same
CP experiment or symbolic source map.  Trace-one output purity is the finite
test $\Tr\omega^2=1$; effect rank one is independent.  Two such branches
or admitted continuations are sufficient when their private output frame is
positive definite.  The source overlap $1/5$ then gives the explicit floor in
Proposition~\ref{sel3:prop:private-floor}.

If no such common-origin writes occur, the alternative is to verify an
actual input-selective private Read before the first discarding reset and
then use V2's forward transfer.  The all-history rank theorem proves that
waiting through more copies of the same reset cannot replace this input
operation.  A purely terminal Read is too late to recover already discarded
input selectivity.  The algebraic support projectors, an optimal
state-discrimination formula, and full recurrence of a relabeled record
window are not substitutes for either source panel.

The narrow remaining task is therefore an \emph{input-effect and typed-domain
identification}, not another weak-router proof, another proof of the bridge
matrix algebra, or another freely completed recurrent channel.  Resolving that
input is required before the determinant/Clifford/matter occurrence tests can
be applied to the same historical packet.

\subsection{Verification scope and cumulative preservation}
V1 and V2 are retained byte for byte before their final document terminator.
The attached exact-arithmetic checks verify the atomic path formula and Choi
rank, the two-write frame and stationary state, distinct-origin algebra ranks,
the completed-private window's superoperator and Choi ranks, its two-step
reset and stationary floor, finite cases of the all-length hub-word formula,
the archive commutant, and the exact private-Read disturbance.  The all-length
statements are proved above, not inferred from those finite tests.  There is
no claim of formal proof-assistant verification or independent peer review.

\begin{thebibliography}{99}
\bibitem[9]{sel3-HSR}
M.~Horodecki, P.~W.~Shor and M.~B.~Ruskai,
\emph{Entanglement Breaking Channels}, Rev.\ Math.\ Phys.\ \textbf{15}
(2003), 629--641;
\href{https://arxiv.org/abs/quant-ph/0302031}{arXiv:quant-ph/0302031}.
The measure--prepare form is established background, not a new result of this
programme.
\bibitem[10]{sel3-KI}
M.~Koashi and N.~Imoto,
\emph{Operations that do not disturb partially known quantum states},
Phys.\ Rev.\ A \textbf{66} (2002), 022318;
\href{https://doi.org/10.1103/PhysRevA.66.022318}{doi:10.1103/PhysRevA.66.022318};
\href{https://arxiv.org/abs/quant-ph/0101144}{arXiv:quant-ph/0101144}.
The pure-state nondisturbance argument is specialized and proved directly in
Section~\ref{sel3:nondisturbance}.
\end{thebibliography}

% END OF SOURCE-SELECTION V3 DEVELOPMENT BLOCK.
% Append future development before the final document terminator.
% Preserve all V1--V3 bodies; identify corrections explicitly.


\clearpage
\hypersetup{pdftitle={Historical Standard-Model Source Selection: Research Notes V1--V4}}
\begin{center}
{\Large\bfseries V4: Pull the categorical line back to the actual input}\\[.4em]
{\large Coefficient-domain separation, admissible pure subchannels, and record erasure}
\end{center}
\addcontentsline{toc}{part}{V4: The actual type input and its refinement boundary}

\section{The upstream source question is a map, not a dimension}
\label{sel4:context}

V3 replaces endpoint filtering by common-origin pure writes, but leaves their
input domain uninstantiated: the input effect must be a positive multiple of
one actual type-line projection.  The present continuation follows the
categorical type line back to its definition in the accepted coefficient
space.  It evaluates that line in the locked Stinespring realization before
asking whether a rank-one physical input has been obtained.  This is an audit
of the supplied source formula, not another choice of a mediator or a reset.

Two levels of refinement must then be distinguished.  A positive subchannel
may be mathematically dominated by a known channel and therefore be obtainable
by a suitable environment effect in a minimal dilation.  Its availability in
the actual record alphabet is a separate condition.  We determine the
complete finite common-input feasibility test, and evaluate it on V3's
completed-private window.  A coarse two-step channel permits pure
common-input subchannels that a preserved first record can forbid.  The exact
record-erasure obstruction has a quantitative channel-norm version.

\subsection{Source rows used and results not repeated}
\begin{center}
\small
\begin{tabular}{L{.38\textwidth}L{.54\textwidth}}
\toprule
Audited source & Content used here\\
\midrule
\cite{Grand}, \src{def:locked-opportunity-instrument},
\src{thm:locked-opportunity-instrument}
& The actual factors $K_e\otimes\Pi_\sigma\otimes P_r$, their coherent sum,
normalization, and the coefficient matrix $M_\vartheta$.\\[.35em]
\cite{Grand}, \src{thm:SM-active-residual-short}
& The source type line is the root-uniform $c-\ell$ coefficient contrast.
On the stated orthogonal active completion it survives the source short.\\[.35em]
\cite{Rigidity}, V32 type-contrast compression and V33 pullback reconstruction
& Environment-effect compression is not system action; a complete
operator-valued complementary pullback reconstructs the latter.  Those general
results are not reproved.\\[.35em]
\cite{Grand}, \src{thm:fresh-endpoint-retained-record};
Theorems~\ref{sel3:thm:window-recurrence}, \ref{sel3:thm:window-rank}
& The completed-private writer, its declared fresh-window iteration, exact
two-step reset, and fixed-record Choi-rank obstruction are inherited.\\[.35em]
Theorem~\ref{sel3:thm:pure-write} and
Corollary~\ref{sel3:cor:write-colour}
& Once actual common-origin pure writes are present, their positive private
frame generates the colour block.  No additional generation theorem is needed.\\
\bottomrule
\end{tabular}
\end{center}

The finite CP-order/environment correspondence is already part of the Grand
Tensor's minimal-realization analysis.  Standard background is the
Radon--Nikodym theorem for completely positive maps
\cite{sel4-Raginsky}.  We use that correspondence as an inherited tool, while
proving the specialized rank, common-domain, and sharp-success formulas below
by finite matrix arguments.  No priority claim is made for CP order,
vectorization, or rank-one positive updates.

\subsection{Conventions and provenance}
All spaces are finite-dimensional and all traces are unnormalized.  For a
CP map $\mathcal J:\B(S)\to\B(O)$ use the input-first Choi convention of V3.
The vectorized coefficient $b:S\to O$ is
$|b\rangle\!\rangle=\sum_i|i\rangle\otimes b|i\rangle$.
Thus a rank-one write $b=|h\rangle\langle t|$ has Choi vector
$|\overline t\rangle\otimes|h\rangle$.  The notation
$\mathcal K\le_{\rm CP}\mathcal J$ means that $\mathcal J-\mathcal K$ is CP;
it does not declare $\mathcal K$ to be an executed branch.

The system, coefficient, environment, private-record, and Choi Hilbert spaces
retain their distinct types.  A line in any one of them is not identified
with a line in another without its displayed source map.  All statements
about the private window retain V3's explicitly declared fresh-register
composition rule.  They are not asserted for every memoryful completion of
one boundary map.

\section{A coefficient ray has two different physical ranks}
\label{sel4:ranks}

Let $V:S\to O\otimes E$ be a Stinespring operator for an actual CP branch,
with coefficient $C(z)=(I\otimes\langle z|)V$ for $z\in E$.  An environment
ray $|z\rangle\langle z|$ produces the single-coefficient CP map
\[
 \mathcal J_z(X)=C(z)XC(z)^*.
\]
Its Choi rank is one whenever $C(z)\ne0$.  Its input effect, however, is
$E_z=C(z)^*C(z)$ and can have any rank up to $\min\{\dim S,\dim O\}$.

\begin{proposition}[Exact input-line test for a coefficient source]
\label{sel4:prop:rank-test}
For $C=C(z)\ne0$, put
\[
 \mathfrak d_{\rm in}(C)
   :=(\Tr C^*C)^2-\Tr[(C^*C)^2].
\]
Then $\mathfrak d_{\rm in}(C)\ge0$, and the following are equivalent:
\begin{enumerate}[label=\textup{(\roman*)},leftmargin=2.5em]
\item the full input effect $C^*C$ has rank one;
\item $C=|h\rangle\langle t|$ for nonzero vectors $h,t$;
\item the vector $|C\rangle\!\rangle$ has Schmidt rank one;
\item $\mathfrak d_{\rm in}(C)=0$.
\end{enumerate}
A unitary change of minimal environment coordinates, accompanied by the same
change of the physical ray, or unitary changes of system input/output frames
preserve these statements.  For a unitary $U$ on an $N$-dimensional space and
$a>0$, the coefficient $C=\sqrt a\,U$ instead has
\begin{equation}
 E_z=aI_N,\qquad
 \mathfrak d_{\rm in}(C)=a^2N(N-1),\qquad
 \frac{\Tr E_z^2}{(\Tr E_z)^2}=\frac1N.
 \label{sel4:eq:unitary-rank}
\end{equation}
\end{proposition}
\begin{proof}
If the eigenvalues of $C^*C$ are $\lambda_i\ge0$, then
$\mathfrak d_{\rm in}(C)=2\sum_{i<j}\lambda_i\lambda_j$.
It vanishes exactly when at most one eigenvalue is positive.  The singular
value decomposition of $C$ gives both the rank-one factorization and the
Schmidt decomposition of its vectorization.  The stipulated coordinate
changes preserve the coefficient up to left/right unitaries; hence they
preserve its singular values.  Substitution of $C^*C=aI_N$ gives
\eqref{sel4:eq:unitary-rank}.
\end{proof}

\begin{proposition}[A robust exclusion of a rank-one donor effect]
\label{sel4:prop:rank-distance}
If $N\ge2$ and $E=aI_N$, then among positive operators $F$ of rank at most
one,
\begin{equation}
 \inf_F\norm{E-F}_{\op}=a,
 \qquad
 \inf_F\norm{E-F}_{\HS}=a\sqrt{N-1}.
 \label{sel4:eq:rank-distance}
\end{equation}
If a reconstructed effect obeys $\norm{\widehat E-aI_N}_{\op}<a$, it is
strictly positive on the complete input, not rank one.  Multiplying the
source frequency by a small positive number changes these absolute scales
but does not change the ranks.
\end{proposition}
\begin{proof}
Every such $F$ has a kernel of dimension at least $N-1$, on which $E-F$ acts
by $a$.  This gives both lower bounds.  Taking $F=ap$ for any rank-one
projection attains them.  The eigenvalue estimate
$\widehat E\succeq(a-\norm{\widehat E-aI})I$ proves the second statement.
\end{proof}

These tests do not require treating the input effect as an independently
executable Read.  They diagnose the input domain of an already supplied
coefficient or CP branch.  Conversely, a positive norm of a coefficient
source ray only proves that its coefficient is nonzero; it does not prove
any of the equivalent input-line statements above.

\section{Exact system pullback of the locked categorical type line}
\label{sel4:locked}

For the literal binary-factor realization write
\[
 S=H_E\otimes\C^2_\sigma\otimes\C^2_r,
 \qquad D_E=\dim H_E,\qquad N=4D_E,
\]
with rank-one pointer projections $\Pi_\pm$ and type projections
$P_c,P_\ell$.  Put
\[
 Q_r=I_{H_E}\otimes I_\sigma\otimes P_r,
 \qquad Z=Q_c-Q_\ell.
\]
The supplied accepted coefficients are
$A_{e\sigma r}=K_e\otimes\Pi_\sigma\otimes P_r$ and satisfy
\begin{equation}
 \sum_eK_e^*K_e=I_{H_E},\qquad
 \sum_eK_e=\sqrt2 I_{H_E}.
 \label{sel4:eq:locked-K}
\end{equation}
There are six $e$ labels.  Their independence implies $D_E\ge3$.
Nothing below selects a particular completion of the unevaluated $K_e$.

On the accepted coefficient environment $E_{24}=\C^{24}$ define real unit
vectors
\begin{equation}
 u_{e\sigma r}=\frac1{\sqrt{24}},\qquad
 t_{e\sigma c}=\frac1{\sqrt{24}},\qquad
 t_{e\sigma\ell}=-\frac1{\sqrt{24}}.
 \label{sel4:eq:ut}
\end{equation}
The vector $t$ is precisely the normalized root-uniform type contrast
$12^{-1/2}\mathbf1_\Phi\otimes t_1$.
Let $c=\mathbf1/\sqrt2=\sqrt{12}\,u$ and
\[
 M_\vartheta=\vartheta I+(1-\vartheta)cc^*,
 \qquad h_\vartheta=12-11\vartheta.
\]
If $B$ synthesizes the accepted Choi vectors, the supplied minimal
Stinespring coefficients are the columns of $BM_\vartheta^{1/2}$.
The corresponding complementary pullback is denoted $\mathcal P^*$.

\begin{theorem}[The categorical contrast pulls back to a full-domain phase]
\label{sel4:thm:type-pullback}
For the exact locked source, independently of the unspecified individual
$K_e$,
\begin{equation}
 \boxed{
 C(u)=\sqrt{\frac{h_\vartheta}{12}}I_S,
 \qquad
 C(t)=\sqrt{\frac\vartheta{12}}Z.}
 \label{sel4:eq:type-coefficient}
\end{equation}
Consequently
\begin{align}
 \mathcal P^*(|t\rangle\langle t|)&=\frac\vartheta{12}I_S,
 \label{sel4:eq:type-effect}\\
 \mathcal P^*(|u\rangle\langle t|)&=
 \frac{\sqrt{\vartheta h_\vartheta}}{12}Z.
 \label{sel4:eq:cross-effect}
\end{align}
A hypothetical execution of the coherent environment ray $t$ would therefore
have state-independent probability $\vartheta/12$ and normalized conditional
system action $X\mapsto ZXZ$.  Its Choi rank is one, its full input effect
has rank $N\ge12$, and its normalized Choi vector is maximally entangled
between the input and output copies of $S$.  In particular this ray does not
supply the rank-one type-input effect required by V3's common-write theorem.
\end{theorem}
\begin{proof}
The vectors $u,t$ are orthogonal, $M_\vartheta u=h_\vartheta u$, and
$M_\vartheta t=\vartheta t$.  The coherent sum in
\eqref{sel4:eq:locked-K} gives
\[
 \sum_{e,\sigma,r}u_{e\sigma r}A_{e\sigma r}=I_S/\sqrt{12},
 \qquad
 \sum_{e,\sigma,r}t_{e\sigma r}A_{e\sigma r}=Z/\sqrt{12}.
\]
Multiplying the coefficient vectors by $M_\vartheta^{1/2}$ proves
\eqref{sel4:eq:type-coefficient}.  The standard complementary identity
$\mathcal P^*(|z\rangle\langle w|)=C(z)^*C(w)$ gives the two pullbacks.
Since $Z=Z^*=Z^{-1}$, all singular values of $C(t)$ equal
$\sqrt{\vartheta/12}$.  Proposition~\ref{sel4:prop:rank-test} gives the
remaining ranks and Schmidt statement.  The dimension bound follows from
six independent matrices in $M_{D_E}(\C)$.
\end{proof}

\begin{remark}[What is evaluated, and what is not identified]
The coefficient source metric gives $t^*M_\vartheta t=\vartheta$.
The full system input effect is instead \eqref{sel4:eq:type-effect}.
These two positive objects live in different spaces.  The exact source
Schur theorem is not contradicted: it identifies a line in its stipulated
common source Hilbert space.  The present calculation shows that its literal
coefficient direction, mapped through the locked Stinespring system
synthesis, does not become a one-dimensional system input.  A different
historical common-carrier realization must supply its different physical
identification map; it is not ruled out merely by this calculation.
On the orthogonal active short, where the retained type contrast is $t$
itself, this particular mismatch cannot be removed by additional precision
on the $K_e$.
\end{remark}

\begin{proposition}[A record-conditioned type is also not this source ray]
\label{sel4:prop:resolved-type}
The resolved accepted type-$r$ event has system effect
\begin{equation}
 \sum_{e,\sigma}L_{e\sigma r}^*L_{e\sigma r}=\vartheta Q_r.
 \label{sel4:eq:resolved-type-effect}
\end{equation}
Its input support has rank $2D_E\ge6$.  After conditioning on acceptance,
the type Read has effects $Q_c,Q_\ell$, whereas the coherent contrast event
of Theorem~\ref{sel4:thm:type-pullback} has a scalar effect.  Neither is the
rank-one projection $|t\rangle\langle t|$ on the coefficient environment.
\end{proposition}
\begin{proof}
Sum $\vartheta K_e^*K_e\otimes\Pi_\sigma\otimes P_r$ and use
\eqref{sel4:eq:locked-K}.  The ranks and the distinction follow from the
explicit factor spaces and \eqref{sel4:eq:type-effect}.
\end{proof}

\subsection{All root-uniform coherent readouts, not just one ray}
Let $L=\Span\{u,t\}\subset E_{24}$.  To obtain a strong finite obstruction,
allow provisionally every positive environment effect $0\preceq F\preceq I$
supported on $L$, not merely the old observed outcomes.  This is a larger
comparison class, not an assertion of physical availability.

\begin{theorem}[Root-uniform refinements cannot create an input line]
\label{sel4:thm:uniform-class}
Every induced system CP map in that class has the form
\begin{equation}
 \mathcal J_F(X)=\sum_{r,s\in\{c,\ell\}}H_{rs}Q_rXQ_s,
 \qquad H\succeq0.
 \label{sel4:eq:Schur-type}
\end{equation}
Its effect is $H_{cc}Q_c+H_{\ell\ell}Q_\ell$.  Every nonzero such map has
input-effect rank at least $2D_E$.  Any finite composition, randomization,
or coarse graining of these maps has the same form and rank alternative.
All their coefficient operators lie in $\C Q_c\oplus\C Q_\ell$.
Thus even arbitrary coherent readouts restricted to the root-uniform
coefficient plane cannot supply a one-dimensional physical donor or a
noncommutative colour action.
\end{theorem}
\begin{proof}
A positive effect decomposes into positive ray effects with vectors in $L$.
By \eqref{sel4:eq:type-coefficient} every resulting Kraus operator has the
form $a_\nu I+b_\nu Z=k_{\nu c}Q_c+k_{\nu\ell}Q_\ell$.
Set $H_{rs}=\sum_\nu k_{\nu r}\overline{k_{\nu s}}$.  It is a Gram
matrix, and expansion gives \eqref{sel4:eq:Schur-type}.  Orthogonality of
the $Q_r$ gives the effect formula.  If both diagonal entries vanish,
positivity forces the whole $H$ to vanish.  Otherwise at least one
$2D_E$-dimensional input sector remains.
Composition replaces $H$ by the entrywise product of the two Gram matrices:
its positivity follows by taking tensor products of their Gram vectors.
Randomization and coarse graining add positive matrices.  The same effect
and coefficient statements therefore persist through the entire indicated
comparison grammar.
\end{proof}

The theorem does not include a later independently supplied rank-one system
filter, nonuniform edge controls, or an intervention coupling the retained
private register.  Such operations change the declared comparison grammar.
In contrast to V3's irreversible atomic reset, the phase action $ZXZ$ loses
no system information; a genuinely available later input-selective filter
could therefore select a line.  No such filter is inferred here.

\section{The maximal CP subchannel on a prescribed input line}
\label{sel4:Choi-short}

The previous calculation excludes one proposed way of obtaining the donor.
A complementary constructive question is now finite and exact.  Given an
actual CP map and an independently identified system line, which output
frames can be carried by CP subchannels supported on that line?  This asks
what the complete channel permits under environment refinement.  It must be
answered before, and separately from, deciding whether that refinement was
physically admitted.

Here $t$ denotes an independently supplied physical input vector, not the
coefficient vector $t$ of Section~\ref{sel4:locked}.
Let $J\succeq0$ be the input-first Choi matrix of
$\mathcal J:\B(S)\to\B(O)$, let $t\in S$ be a unit vector, and fix a
candidate target subspace $H\subseteq O$.  Define the isometry
\[
 V_t:H\longrightarrow\overline S\otimes O,
 \qquad V_th=\overline t\otimes h.
\]
Let $\Pi_J$ be the support projection of $J$.  Form the positive matrix
\begin{equation}
 N_J(t)=V_t^*(I-\Pi_J)V_t,\qquad
 H_J(t)=\Ker N_J(t)\subseteq H.
 \label{sel4:eq:fibre-support}
\end{equation}
On this subspace let $V_t^0=V_t|_{H_J(t)}$ and
\begin{equation}
 F_J^{\max}(t)
 =\left[(V_t^0)^*J^\dagger V_t^0\right]^{-1}
 \quad\hbox{on }H_J(t),
 \qquad F_J^{\max}(t)=0\quad\hbox{on }H_J(t)^\perp.
 \label{sel4:eq:input-short}
\end{equation}
If $H_J(t)=0$, this definition means the zero operator.  There is no inverse
on a discarded nullspace.

\begin{theorem}[Exact common-input Choi short]
\label{sel4:thm:Choi-short}
A CP map of the form
\begin{equation}
 \mathcal K_F(X)=\langle t,Xt\rangle F,
 \qquad F\succeq0,\quad\Ran F\subseteq H,
 \label{sel4:eq:common-input-map}
\end{equation}
is dominated by $\mathcal J$ if and only if
\begin{equation}
 \boxed{0\preceq F\preceq F_J^{\max}(t).}
 \label{sel4:eq:input-short-order}
\end{equation}
The matrix in brackets in \eqref{sel4:eq:input-short} is strictly positive
on $H_J(t)$.  Thus $F_J^{\max}(t)$ is the unique largest output frame of a
CP subchannel on the prescribed full input line.  In particular,
\[
 \rank F_J^{\max}(t)=\dim H_J(t),
\]
and that rank is the maximal number of linearly independent pure output
vectors that can occur in a CP refinement on this one input.
For a unit $h\in H$, a positive subchannel
$X\mapsto c\langle t,Xt\rangle|h\rangle\langle h|$ exists exactly when
$h\in H_J(t)$, and its sharp coefficient bound is
\begin{equation}
 \boxed{
 0<c\le c_{\max}(t,h):=
 \left[\langle\overline t\otimes h,
                 J^\dagger(\overline t\otimes h)\rangle\right]^{-1}.}
 \label{sel4:eq:single-budget}
\end{equation}
\end{theorem}
\begin{proof}
The Choi matrix of \eqref{sel4:eq:common-input-map} is $V_tFV_t^*$.
By the finite Choi characterization, CP domination is precisely
$V_tFV_t^*\preceq J$.  A positive operator dominated by $J$ has range in
$\Ran J$: if $x\in\Ker J$, its quadratic form and then its positive square
root annihilate $x$.  The identity
\[
 \langle h,N_J(t)h\rangle=\norm{(I-\Pi_J)V_th}^2
\]
therefore shows that a dominated $F$ is supported on $H_J(t)$.
On that space $V_t^0$ is injective with image in $\Ran J$, so
$M=(V_t^0)^*J^\dagger V_t^0\succ0$.
Whitening by $J^{\dagger/2}$ reduces the inequality to
\[
 WFW^*\preceq I_{\Ran J},\qquad
 W=J^{\dagger/2}V_t^0,\quad W^*W=M.
\]
Write $W=UM^{1/2}$ with $U$ an isometry.  The last inequality is equivalent
to $M^{1/2}FM^{1/2}\preceq I$, hence to $F\preceq M^{-1}$.
This proves \eqref{sel4:eq:input-short-order}, including singular $F$.
A spectral decomposition of $F_J^{\max}(t)$ gives rank-one CP subchannels
whose sum is dominated by $\mathcal J$, with independent outputs spanning
$H_J(t)$.  No larger output support can occur by the range condition.
Finally, whitening $c|V_th\rangle\langle V_th|\preceq J$ gives
$c\norm{J^{\dagger/2}V_th}^2\le1$, yielding
\eqref{sel4:eq:single-budget}.
\end{proof}

\begin{remark}[Why simple compression is insufficient]
The admissible output support is defined by
$V_th\in\Ran J$, not merely by
$\langle V_th,JV_th\rangle>0$.
The positive compression $V_t^*JV_t$ can be nonzero even when
$H_J(t)=0$.  For example, if $J=|I_N\rangle\!\rangle
\langle\!\langle I_N|$ with $N>1$, then
$V_t^*JV_t=|t\rangle\langle t|$ in matching coordinates, but the sole
vectorized identity in $\Ran J$ is not a product vector.  Consequently
$F_J^{\max}(t)=0$.  A nonzero response on a prepared input is not by itself
a subchannel whose \emph{entire} input effect is supported on that line.
\end{remark}

\begin{corollary}[Exact result for an atomic branch]
\label{sel4:cor:atomic-short}
For the inherited atomic map $\mathcal J(X)=\Tr(EX)\omega$, with
$E\succeq0$, $\omega\succeq0$, the maximal frame on a unit input $t$,
with unrestricted output target, is
\begin{equation}
 F_J^{\max}(t)=
 \begin{cases}
 \omega/\langle t,E^\dagger t\rangle,&t\in\Ran E,\\
 0,&t\notin\Ran E.
 \end{cases}
 \label{sel4:eq:atomic-short}
\end{equation}
In particular an atomic branch with pure output can be refined to a
common-input pure write precisely on lines in its effect support.  Such a
refinement is not present merely because the effect support is known.
\end{corollary}
\begin{proof}
Use $J=E^{\mathsf T}\otimes\omega$ and
$J^\dagger=(E^\dagger)^{\mathsf T}\otimes\omega^\dagger$.
A product $\overline t\otimes h$ belongs to the support exactly when
$t\in\Ran E$ and $h\in\Ran\omega$.  The restricted inverse in
\eqref{sel4:eq:input-short} is then
$[\langle t,E^\dagger t\rangle\,\omega^\dagger]^{-1}$ on that
support, proving the formula.
\end{proof}

\begin{proposition}[The environment effect required for the refinement]
\label{sel4:prop:RN-effect}
Choose a minimal Kraus list $(C_\nu)$ for $\mathcal J$, and let $W$ be the
full-column-rank matrix with columns $|C_\nu\rangle\!\rangle$; thus
$J=WW^*$.  If $V_tFV_t^*\preceq J$, the unique coefficient-side
Radon--Nikodym matrix is
\begin{equation}
 R_F=W^\dagger V_tFV_t^*(W^\dagger)^*,\qquad 0\preceq R_F\preceq I.
 \label{sel4:eq:required-effect}
\end{equation}
For the dilation $\sum_\nu C_\nu\otimes|\nu\rangle$, the physical
environment effect has coordinate matrix $R_F^{\mathsf T}$.
For a pure write with $F=c|h\rangle\langle h|$ this is a rank-one effect.
A declared environment effect bank can therefore be tested against
\eqref{sel4:eq:required-effect}; the CP-order test alone does not supply that
bank.
\end{proposition}
\begin{proof}
The inherited CP Radon--Nikodym correspondence gives uniqueness on a minimal
coefficient space.  Directly, $WR_FW^*=V_tFV_t^*$ because its range is
in $\Ran W$.  Whitening by $W^\dagger$ makes domination equivalent to
$0\preceq R_F\preceq I$; injectivity of $W$ proves uniqueness.
An environment effect $F_E$ in the stated dilation gives Choi matrix
$WF_E^{\mathsf T}W^*$, by expanding the partial trace over the environment.
This accounts for the transpose and avoids equating unconjugated coefficient
and physical-environment coordinates.  Rank one follows from the displayed
formula for a pure $F$ and from $W^\dagger V_th\ne0$.
\end{proof}

\begin{corollary}[The locked type-phase outcome has no donor refinement]
\label{sel4:cor:phase-short}
For the hypothetical coherent type outcome of
Theorem~\ref{sel4:thm:type-pullback},
$F_J^{\max}(x)=0$ for every physical unit input $x\in S$.
The same conclusion holds for any root-uniform outcome in
Theorem~\ref{sel4:thm:uniform-class} when $D_E\ge1$ and both type sectors
have dimension greater than one.  No CP refinement of such an outcome,
even with unrestricted access to its own minimal environment, creates a
nonzero common-input pure write.
\end{corollary}
\begin{proof}
A rank-one write coefficient is a product vector in Choi space.
The phase outcome has Choi support $\C|Z\rangle\!\rangle$, whose
nonzero coefficients have rank $N>1$.  The general root-uniform outcome has
coefficient support contained in $\Span\{Q_c,Q_\ell\}$, where every
nonzero coefficient has rank at least $2D_E>1$.  Neither support contains
a nonzero rank-one write.  The support criterion of
Theorem~\ref{sel4:thm:Choi-short} gives the result.
\end{proof}

\section{Preserved histories and forgotten histories have different input shorts}
\label{sel4:marked}

Let $(\mathcal J_a)_a$ be the actually resolved branches of an instrument,
with Choi matrices $J_a$, and $J=\sum_aJ_a$ the coarse channel.  A
\emph{record-respecting refinement} consists of positive subchannels beneath
each $\mathcal J_a$, possibly followed by discarding the newly added success
labels.  The old record $a$ is still available; no coherent mixing of its
orthogonal record values is declared.

\begin{theorem}[The exact cost of retaining the source record]
\label{sel4:thm:marked-short}
For a prescribed physical input $t$ and target $H$, the largest output frame
of a common-input subchannel obtainable by such a refinement is
\begin{equation}
 \boxed{
 F_{\rm marked}^{\max}(t)=\sum_aF_{J_a}^{\max}(t)
 \preceq F_J^{\max}(t)=:F_{\rm coarse}^{\max}(t).}
 \label{sel4:eq:marked-short}
\end{equation}
More precisely, a frame $F$ is feasible with the record retained exactly
when the finite semidefinite decomposition problem
\begin{equation}
 F=\sum_aF_a,\qquad 0\preceq F_a\preceq F_{J_a}^{\max}(t)
 \label{sel4:eq:marked-feasibility}
\end{equation}
has a solution.  Every such frame is bounded above by the sum in
\eqref{sel4:eq:marked-short}, and that sum is attained.
It supports two independent private write rays if and only if it is
positive definite on the specified private plane.  The positive matrix
$F_{\rm coarse}^{\max}(t)-F_{\rm marked}^{\max}(t)$ measures a genuine
CP-order enlargement when only the record-forgotten channel is specified.
Realizing that enlargement may require coherent access to the erased record;
classical discard is not sufficient.  It need not vanish and can have positive rank even when
$F_{\rm marked}^{\max}(t)=0$.
\end{theorem}
\begin{proof}
Suppose positive refined Choi operators $J'_a\preceq J_a$ sum to
$V_tFV_t^*$.  Each $J'_a$ is positive and bounded above by that sum; hence
its range is contained in $V_tH$.  Thus $J'_a=V_tF_aV_t^*$ for
$F_a\succeq0$, and Theorem~\ref{sel4:thm:Choi-short} gives
$F_a\preceq F_{J_a}^{\max}(t)$.  Summing proves the upper bound.
Conversely, the maximal subchannel from each recorded branch may be selected
in a formal CP refinement and the selections summed, attaining the sum in
\eqref{sel4:eq:marked-short}.  It is a subchannel of the coarse channel,
so the coarse maximality theorem gives the remaining inequality.
The same argument in both directions gives \eqref{sel4:eq:marked-feasibility}.
Spectral rank-one refinements of the individual $F_{J_a}^{\max}(t)$ have
output span equal to the range of their sum.  This proves the private-frame
criterion; it does not assume a phase-coherent sum of different outcomes.
The strict possibility is established by the completed-private calculation
below.
\end{proof}

\begin{remark}[A maximum frame does not determine the entire marked interval]
The condition $F\preceq\sum_aF_{J_a}^{\max}(t)$ alone need not solve
\eqref{sel4:eq:marked-feasibility}.  On a scalar input take two recorded
preparations $\mathcal J_0(x)=x|0\rangle\langle0|/2$ and
$\mathcal J_1(x)=x|1\rangle\langle1|/2$.  Both marked and coarse maximal
frames equal $I_2/2$.  Nevertheless $|+\rangle\langle+|/2\preceq I_2/2$
is a coarse pure subchannel and cannot be a sum of positive refinements of
the two recorded preparations, whose output supports are separate lines.
The semidefinite decomposition, not just the maximal rank, retains this
additional constraint.
\end{remark}

Attainability here is in CP order with the requisite environment effects.
It is not an assertion that the extra effects were observed historically.
For a finite physically supplied refinement bank, one must still compare its
actual Radon--Nikodym matrices with
Proposition~\ref{sel4:prop:RN-effect}.

\subsection{A closed formula for orthogonal input classes}

\begin{proposition}[Harmonic output frame across unresolved input sectors]
\label{sel4:prop:harmonic}
Let $S=\bigoplus_a S_a$ have orthogonal projections $P_a$, and suppose
$E_a$ is strictly positive on $S_a$ and zero elsewhere.  Let
\[
 \mathcal J(X)=\sum_a\Tr(E_aX)\omega_a,
 \qquad\omega_a\succeq0.
\]
For a unit $t=\sum_at_a$ put $I_t=\{a:t_a\ne0\}$ and
\[
 H_t=\bigcap_{a\in I_t}\Ran\omega_a,
 \qquad \alpha_a=\langle t_a,E_a^{-1}t_a\rangle.
\]
With full output target $O$, the common-input short is zero if $H_t=0$;
otherwise it is supported on $H_t$ and equals
\begin{equation}
 \boxed{
 F_J^{\max}(t)=
 \left[P_{H_t}\left(\sum_{a\in I_t}\alpha_a\omega_a^\dagger\right)
           P_{H_t}\bigm|_{H_t}\right]^{-1}.}
 \label{sel4:eq:harmonic-frame}
\end{equation}
If the old record resolves $a$ and $|I_t|>1$, its marked common-input short
is zero even when \eqref{sel4:eq:harmonic-frame} is positive.
\end{proposition}
\begin{proof}
The coarse Choi operator is block diagonal on the orthogonal input factors:
$J=\bigoplus_a E_a^{\mathsf T}\otimes\omega_a$.
The vector $\overline t\otimes h$ lies in its range exactly when
$h\in\Ran\omega_a$ for every nonzero $t_a$.  Its inverse quadratic form
on that intersection is
$\sum_a\alpha_a\langle h,\omega_a^\dagger h\rangle$.
Theorem~\ref{sel4:thm:Choi-short} gives
\eqref{sel4:eq:harmonic-frame}.  In an individual recorded branch the
input support is only $S_a$.  A vector with nonzero components in two input
classes belongs to none of those supports, so
$F_{J_a}^{\max}(t)=0$ for every $a$ by
Corollary~\ref{sel4:cor:atomic-short}.  Apply
Theorem~\ref{sel4:thm:marked-short}.
\end{proof}

\subsection{A sharp channel-norm obstruction to restoring an erased input phase}
For two orthogonal input classes let
$\mathcal D(X)=P_0XP_0+P_1XP_1$.  The norm $\|\cdot\|_\diamond$ is the
completely bounded trace norm on linear maps.  It is used on unnormalized
outcome maps; one is not permitted to divide away a small success probability.

\begin{theorem}[Exact distance from a common write to a record-dephased map]
\label{sel4:thm:dephasing-distance}
Let $t$ be unit, $p_a=\|P_at\|^2$, and $\omega$ a density operator.
For $c>0$ define
$\mathcal K(X)=c\langle t,Xt\rangle\omega$.
Among CP maps $\mathcal L$ satisfying
$\mathcal L=\mathcal L\mathcal D$,
\begin{equation}
 \boxed{
 \inf_{\mathcal L=\mathcal L\mathcal D}
       \|\mathcal K-\mathcal L\|_\diamond
 =c\sqrt{p_0p_1}.}
 \label{sel4:eq:dephasing-distance}
\end{equation}
The same minimum holds if maps are required trace-nonincreasing and
$0<c\le1$.  The minimizer $\mathcal K\mathcal D$ is CP (and in that case
trace-nonincreasing).  A balanced donor across the two input classes therefore
has an unavoidable distance $c/2$ from every such record-dephased operation.
\end{theorem}
\begin{proof}
If either $p_a$ is zero, $\mathcal K=\mathcal K\mathcal D$ and the claim is
immediate.  Otherwise set $u_a=P_at/\sqrt{p_a}$.
The input $X=|u_0\rangle\langle u_1|$ has trace norm one and
$\mathcal D(X)=0$.  Hence every admissible $\mathcal L$ gives
$\mathcal L(X)=0$, whereas
$\|\mathcal K(X)\|_1=c\sqrt{p_0p_1}$, proving the lower bound.
For the upper bound take $\mathcal L=\mathcal K\mathcal D$.
Their difference is $X\mapsto\Tr(AX)\omega$ with
\[
 A=c\bigl(|t\rangle\langle t|-\mathcal D(|t\rangle\langle t|)\bigr),
 \qquad \|A\|_{\op}=c\sqrt{p_0p_1}.
\]
The completely bounded trace norm of this measure--prepare linear map is
$\|A\|_{\op}\|\omega\|_1=\|A\|_{\op}$.
To check this directly, at any ancillary dimension the functional part
$X\mapsto\Tr_S[(A\otimes I)X]$ has trace norm at most
$\|A\|_{\op}\|X\|_1$ by trace-norm duality against all ancillary
contractions.  A pure input eigenvector for a largest absolute eigenvalue of
$A$ attains the bound.  Tensoring its output with the fixed density $\omega$
preserves the trace norm.  Finally the effect of $\mathcal K\mathcal D$ is
$c\mathcal D(|t\rangle\langle t|)\preceq I$ when $c\le1$.
\end{proof}

\section{Evaluate the new input short on the completed-private source}
\label{sel4:window}

Retain exactly V3's source and its declared fresh-register composition.
The physical window input and output are $F=U\otimes R$, with
$\dim U=\dim R=2$.  Let
\[
 P_a=|a\rangle\langle a|\otimes I_R,\qquad
 \rho_a=|r_a\rangle\langle r_a|,\qquad
 \langle r_0,r_1\rangle=1/5.
\]
The known formulas are
\[
 \Psi(X)=\sum_{a=0}^1\Tr(P_aX)\omega_a,
 \quad \omega_a=\tfrac12 I_U\otimes\rho_a,
 \quad \Psi^2(X)=\Tr(X)\sigma_*,
 \quad \sigma_*=\tfrac12 I_U\otimes\bar\rho,
 \quad \bar\rho=\tfrac12(\rho_0+\rho_1).
\]
The recurrence, spectrum, and all-history rank statements are not proved
again here.  The new question is the maximal \emph{full-input-line}
subchannel under different record access policies.

\begin{theorem}[One boundary has no common private-write frame]
\label{sel4:thm:one-step-short}
For a unit physical input $t\in F$ and unrestricted window output,
\begin{equation}
 \boxed{
 F_{J(\Psi)}^{\max}(t)=
 \begin{cases}
 \tfrac12I_U\otimes\rho_a,&t\in P_aF\text{ for one }a,\\
 0,&P_0t\ne0\text{ and }P_1t\ne0.
 \end{cases}}
 \label{sel4:eq:one-step-short}
\end{equation}
Resolving both the input hub and the fresh output hub gives the same maximal
frame after summing record-respecting refinements.  For a fixed fresh endpoint
$|h\rangle$, the maximal private frame is $\rho_a/2$ on a one-sector
input, and zero on a two-sector input.  Thus one boundary cannot produce two
linearly independent private outputs from one full input line, even under
arbitrary CP refinement of the known coarse channel.
\end{theorem}
\begin{proof}
The two output supports are
$O_a=U\otimes\C r_a$.  Because $r_0,r_1$ are linearly independent,
$O_0\cap O_1=0$; their nonorthogonality does not change this intersection.
Apply Proposition~\ref{sel4:prop:harmonic} with $E_a=P_a$.  On a
one-sector unit input its coefficient $\alpha_a$ is one and the maximal
frame is $\omega_a$; on a two-sector input the common output support is
zero.  The hub-resolved maps have Choi operators
\[
 J_{ha}=\tfrac12P_a^{\mathsf T}\otimes
                (|h\rangle\langle h|\otimes\rho_a).
\]
Their individual shorts are $|h\rangle\langle h|\otimes\rho_a/2$ if
$t\in P_aF$, and zero otherwise.  Summing over the admitted hub labels
recovers \eqref{sel4:eq:one-step-short}.  Restricting the target to one
fresh endpoint gives the private statement.
\end{proof}

The output frame in \eqref{sel4:eq:one-step-short} can have rank two on the
\emph{hub--record product}.  Its private factor still has rank one.  Those
two dimensions are not two independent private writes and are not identified
with colour.

\begin{theorem}[Two-step common-input frames depend sharply on record retention]
\label{sel4:thm:two-step-short}
For the two-step coarse window channel and every physical unit input $t$,
\begin{equation}
 \boxed{F_{J(\Psi^2)}^{\max}(t)=\sigma_*.}
 \label{sel4:eq:coarse-two-short}
\end{equation}
If the first input hub record is preserved, the exact record-respecting
maximum, even allowing arbitrary further CP refinements within each old
record, is instead
\begin{equation}
 \boxed{
 F_{\rm marked,2}^{\max}(t)=
 \begin{cases}
 \sigma_*,&t\in P_0F\text{ or }t\in P_1F,\\
 0,&P_0t\ne0\text{ and }P_1t\ne0.
 \end{cases}}
 \label{sel4:eq:marked-two-short}
\end{equation}
The same formula holds when every intermediate hub label is retained.
Without any added refinement, the actual positive-length resolved histories
continue to have Choi rank two, as proved in V3.  The positive maxima here
are therefore conditional on the extra environment effects, not automatic
operations of the original hub alphabet.
\end{theorem}
\begin{proof}
The coarse Choi operator is $I_F^{\mathsf T}\otimes\sigma_*$ and is
strictly positive.  Theorem~\ref{sel4:thm:Choi-short} gives
$F^{\max}=\sigma_*$ for every unit input.
If just the first input hub $a$ is preserved, its two-step CP branch is
$X\mapsto\Tr(P_aX)\sigma_*$.  Its individual short is $\sigma_*$ for
$t\in P_aF$ and zero otherwise by
Corollary~\ref{sel4:cor:atomic-short}.  The two cases in
\eqref{sel4:eq:marked-two-short} follow by summing those maxima.

For the fully resolved two-step history, write its labels as
$(a,b,h)$, where $a$ is the first input hub, $b$ the next hub, and $h$ the
final hub.  The inherited path formula gives
\begin{equation}
 \mathcal J_{a,b,h}^{(2)}(X)
  =\tfrac14\Tr(P_aX)
            (|h\rangle\langle h|\otimes\rho_b).
 \label{sel4:eq:two-histories}
\end{equation}
Its maximal frame on a one-sector input $t\in P_aF$ is exactly the output
on the right divided by four.  Summing over $b,h$ gives $\sigma_*$; the
other initial sector contributes zero.  For a two-sector input every
individual maximum is zero.  This proves that further old hub labels do not
alter the displayed maximum.  The last distinction follows from V3's
all-length Choi-rank theorem: a proper CP refinement is not a coarse graining
or a continuation within that original alphabet.
\end{proof}

\begin{corollary}[Sharp two-write success budget at a fixed fresh endpoint]
\label{sel4:cor:pair-budget}
Fix a unit physical input $t$ and one fresh endpoint $|h\rangle$.
For the coarse two-step channel, two hypothetical refined writes with
coefficients
\[
 b_j=\sqrt{c_j}\,|h\rangle\otimes|r_j\rangle\langle t|,
 \qquad j=0,1,
\]
can simultaneously be CP subchannels if and only if
\begin{equation}
 \boxed{0\le c_0\le\tfrac14,\qquad0\le c_1\le\tfrac14.}
 \label{sel4:eq:pair-budget}
\end{equation}
At the maximum $c_0=c_1=1/4$, their total success probability on $p_t$ is
$1/2$, and their private output frame has
\begin{equation}
 F_R=\tfrac14(\rho_0+\rho_1),\qquad
 \operatorname{spec}(F_R)=\{3/10,1/5\},\qquad
 \det F_R=3/50.
 \label{sel4:eq:private-optimum}
\end{equation}
On a one-sector input these refinements can respect all old hub records;
on a two-sector input they cannot respect even the first hub record.
No identification of $t$ with the categorical source ray is made.
\end{corollary}
\begin{proof}
The maximal output frame restricted to the fixed fresh endpoint is
$|h\rangle\langle h|\otimes\bar\rho/2$ by
\eqref{sel4:eq:coarse-two-short}.  Hence simultaneous domination is
\[
 c_0\rho_0+c_1\rho_1\preceq\tfrac14(\rho_0+\rho_1).
\]
Let $R_0$ be the invertible two-column matrix $(r_0,r_1)$.
Congruence by $R_0^{-1}$ changes this inequality to
$\diag(c_0,c_1)\preceq I_2/4$, proving the sharp rectangle
\eqref{sel4:eq:pair-budget}.  The eigenvalues use the inherited overlap
$1/5$, and their product gives the determinant.
For a one-sector $t$, the two individual histories
\eqref{sel4:eq:two-histories} with fixed $a,h$ and $b=0,1$ have exactly
these maximal subchannels.  For a two-sector $t$,
\eqref{sel4:eq:marked-two-short} excludes every nonzero such refinement.
\end{proof}

\begin{remark}[Classical forgetting is not quantum erasure]
The coarse two-step CP-order interval is larger because the completely
positive map by itself no longer retains the old record decomposition.
Classically deleting an already dephased record does not implement the new
pure subchannels.  A physical implementation for a two-sector $t$ needs
coherent access to the degrees of freedom carrying the erased initial hub
information.  The required minimal-environment effect is given by
Proposition~\ref{sel4:prop:RN-effect}.  An actual permanently classical
record cannot be kept intact while that same input phase is restored.
For a one-sector input, a new selection of the initially discarded private
input is still required; the old hub labels alone are not such a selection.
\end{remark}

\begin{corollary}[A nonzero quantitative barrier for a balanced donor]
\label{sel4:cor:window-distance}
If $\|P_0t\|^2=\|P_1t\|^2=1/2$, each of the optimal pure-write maps in
Corollary~\ref{sel4:cor:pair-budget}, of weight $c_j=1/4$, has diamond
distance at least $1/8$ from every CP operation that respects the first
hub dephasing.  This is an absolute unnormalized operation error, not an
error divided by its success probability.
\end{corollary}
\begin{proof}
Every record-respecting refinement in \eqref{sel4:eq:two-histories} is
supported on one initial hub sector and hence annihilates the off-diagonal
input blocks.  The same is true after summing such maps.  Apply
Theorem~\ref{sel4:thm:dephasing-distance} with $c=1/4$ and $p_0=p_1=1/2$.
\end{proof}

\section{What is now decided and the actual next source row}
\label{sel4:status}

\subsection{A source evaluation rather than a new favourable completion}
The categorical type ray has been evaluated through the source's own locked
coefficient synthesis.  It gives a full-domain phase with scalar effect,
not the pure input effect required by V3.  The result does not depend on
choosing the six unknown edge coefficients.  All root-uniform coherent
readouts satisfy a stronger finite rank-floor obstruction.  Thus neither
reconstructing those same coefficients more accurately nor reading the
uniform contrast more coherently closes this particular input-domain gap.

The positive mathematical addition is the common-input Choi short.  Given an
independently identified physical line, it returns the maximal possible
private output frame, sharp weights for pure writes, and the exact
environment effect required.  The marked version retains physical history
labels and quantifies what a record-forgotten channel comparison would add.
It is a complete finite feasibility calculation, not another declaration that
suitable writes exist.

For the inherited completed-private window, one boundary cannot provide two
independent private outputs on a common full input line.  The two-step coarse
channel can admit such refinements, with the exact success budget
\eqref{sel4:eq:pair-budget}; whether the first record is preserved changes
the common-input answer from a positive full frame to zero for a donor
spanning both hub classes.  Even when the marked order interval is positive,
the extra environmental refinement remains an input.  These assertions use
V3's already stated fresh-register observation policy; the original
one-boundary formula alone does not certify that policy for other sources.

\begin{center}
\small
\begin{tabular}{L{.38\textwidth}L{.54\textwidth}}
\toprule
Question & V4 result\\
\midrule
Does the normalized locked coefficient contrast supply a rank-one system input?
& No under its canonical system pullback: its effect is
$\vartheta I/12$ and its input rank is the full system dimension.\\[.35em]
Can the whole root-uniform coherent readout plane repair that rank?
& No: every nonzero effect retains an entire type sector of dimension at least
six.\\[.35em]
Can a general channel's pure-input refinement possibilities be computed?
& Yes: $F_J^{\max}(t)$ is an exact support-restricted inverse and a unique
maximal positive output frame.\\[.35em]
Does a feasible positive subchannel certify a historical Read?
& No: its unique Radon--Nikodym effect must also be in the physically admitted
readout bank.\\[.35em]
Can forgetting a classical record change that mathematical feasibility?
& Yes: the two-step source gives a full coarse frame but zero marked frame on
a two-sector input.  Classical discard alone does not implement it.\\[.35em]
What is the exact source-specific two-write bound?
& With fixed new hub endpoint: $c_0,c_1\le1/4$, private floor $1/5$, and
success $1/2$ at the joint maximum.\\[.35em]
Is historical colour, determinant occurrence, or SM selection now proved?
& No.  A physical input-domain map and the relevant actual refinement or
nonuniform system operation are still needed.\\
\bottomrule
\end{tabular}
\end{center}

\subsection{The noncircular continuation}
The next audit must not identify the source-vector projection
$|t\rangle\langle t|$ with its system pullback: the latter is now explicitly
known and has the wrong rank.  It must locate an actual physical input map
for the retained categorical line, on the same Hilbert representation as the
private carrier.  There are two precise possibilities to inspect.

A source may provide an exposed input-selective branch outside the
root-uniform coefficient plane.  Its full Choi matrix, not just a state
prepared on a smaller test domain, can be fed to
Theorem~\ref{sel4:thm:Choi-short}.  The output frame and the source's actual
Radon--Nikodym effect then decide whether V3's common-origin writes occurred.
Alternatively, it may supply coherent access to degrees of freedom erased at
the first reset.  In that case the source must specify which old records are
retained and whether the required effects in
Proposition~\ref{sel4:prop:RN-effect} can be executed without violating that
record policy.  A new rank-one label fitted to a coarse Choi matrix is not a
physical identification.

The general common-carrier theorem and the positive type-line Schur floor
remain valid on their declared realization.  The new obstruction concerns
one concrete attempted passage from its coefficient line to its system
input.  It does not rule out an independently reconstructed common source
with a different action map, nor any colour proof using nonuniform operations
outside the tested grammar.  Once such a source is actually supplied,
V1--V3 and the inherited weak/central tests still give the internal assembly;
the determinant and matter-type conditions remain downstream and separate.

\subsection{Verification scope}
The V3 prefix is retained byte for byte before its terminal document command.
The exact-arithmetic checks accompanying V4 evaluate the locked contrast
without fitting individual edge coefficients, the two rank obstructions,
the support-restricted inverse on positive and singular Choi examples, the
one- and two-step completed-window shorts, the marked/coarse difference,
the simultaneous two-write rectangle, the private eigenvalues and determinant,
and the dephasing-distance witness.  The V3 regression suite is rerun.
These are checks of explicit matrix identities, not a proof-assistant
formalization or an independent audit of every inherited theorem.  The
historical missing inputs are not set to zero in those computations.

\begin{thebibliography}{99}
\bibitem[11]{sel4-Raginsky}
M.~Raginsky, \emph{Radon--Nikodym derivatives of quantum operations},
J.\ Math.\ Phys.\ \textbf{44} (2003), 5003--5020;
\href{https://arxiv.org/abs/math-ph/0303056}{arXiv:math-ph/0303056};
\href{https://doi.org/10.1063/1.1615697}{doi:10.1063/1.1615697}.
The CP-order/environment correspondence is established background and an
inherited tool; the specialized source-domain and marked-short calculations
are proved in this development block.
\end{thebibliography}

% END OF SOURCE-SELECTION V4 DEVELOPMENT BLOCK.
% Append subsequent developments before the final document terminator.
% Preserve the V1--V4 bodies and state necessary corrections explicitly.


\clearpage
\hypersetup{pdftitle={Historical Standard-Model Source Selection: Research Notes V1--V5}}
\begin{center}
{\Large\bfseries V5: What the source can actually Read}\\[.4em]
{\large Access-constrained input shorts, erased-input identification,\\
weak-read amplification, and exact-algebra instability}
\end{center}
\addcontentsline{toc}{part}{V5: Readout access and the finite/cofinal selection boundary}

\section{The next upstream datum is an operational readout cone}
\label{sel5:context}

V4 evaluated the canonical type-contrast pullback and ruled it out as the
rank-one physical donor.  It also computed every common-input subchannel
permitted by unrestricted CP order.  Repeating either calculation cannot
establish that the required environment effect occurred.  We now keep the
channel and its input/output cuts fixed and ask what follows from the
\emph{actually exposed readout operations}.

The upstream definition \src{def:bare-predictive-renewal} in \cite{Grand}
requires a finite physical CP frame spanning each exposed Hermitian
Choi-coordinate space, finite classical control, composition, and coarse
graining.  It does not say that every positive element of the linear span is
an admitted outcome.  Its \src{def:completed-private-packet} supplies the
completed-private record and overlap $1/5$, not a selective measurement of the
input register erased by that completion.  These distinctions are consistent
with, rather than corrections to, the source's stated operational entrance.

This block has three positive outputs.  First, a finite algebra-relative
support computation returns the largest common-input frame allowed by a
specified readout algebra, and a separate test identifies its pure-write
part.  Second, on the completed-private history, the missing readout is
identified exactly: it is a selective Read of the discarded \emph{input}
private register in a calibrated frame, not a new Read of the written output.
Third, an actually admitted weak nondemolition Read can reconstruct a
candidate donor ray and approximate its selection with a controlled error
and nonvanishing total success, without a sharp projector primitive;
forgetting the new Read outcomes recovers the old writer exactly.  The last statement is a cofinal approximation,
not exact finite algebraic selection.  A four-dimensional calculation proves
why that distinction cannot be dropped.

\subsection{Audit and inherited inputs}
The type pullback, unrestricted common-input Choi short, marked CP-order
interval, two-boundary private writer, and its sharp $1/4$ write budgets are
V4 inputs.  The private-frame colour theorem is an input from V3.  The
forward-algebra/recurrence theorem of V2 and the earlier Howe convergence
programme are not proved again.  Targeted searches of the available source
notes located no algebra-constrained common-input short or the weak-read
source calculation developed below.  This is a nonduplication statement
relative to the material inspected, not a claim of literature priority.

The CP Radon--Nikodym correspondence is inherited from
\cite{sel4-Raginsky}.  Repeated quantum nondemolition measurement and its
collapse limit are established topics; see \cite{sel5-BB}.  The specialized
finite coarse-graining formula, its success/error budgets on the specified
private history, and the algebraic warning are proved explicitly here.

\subsection{Conventions that prevent an implicit source enlargement}
An \emph{access algebra} below is a represented finite $C^*$-algebra
containing the coefficient-side transposes of the admitted environment
effects.  An upper-bound theorem can use such an algebra without assuming all
its effects are physical.  A constructive attainability assertion for its
entire effect interval requires that those effects be admitted.  A finite
measurement bank which merely generates the same algebra need not have that
property.  Section~\ref{sel5:tomography} supplies an exact counterexample.

Environment coordinates in a specified dilation are kept fixed.  Changing
coordinates conjugates the access algebra together with the dilation; it does
not turn an inaccessible effect into an accessible one.  A redundant dilation
is allowed in the first theorem and is not reduced in a way that loses a
physical record.  Claims that compare effect rank with Choi rank explicitly
require an injective Kraus synthesis.

\section{A common-input short constrained by the readout algebra}
\label{sel5:access}

Let an actual CP map $\mathcal J:\B(S)\to\B(O)$ have a specified Kraus
synthesis
\[
 W:E\longrightarrow\overline S\otimes O,
 \qquad W e_\nu=|C_\nu\rangle\!\rangle,
 \qquad J=WW^*.
\]
The synthesis need not be injective.  If the physical environment effect is
$F_E$, the selected Choi operator is $WF_E^{\mathsf T}W^*$.  Write
$R=F_E^{\mathsf T}$, and let $\mathcal D\subseteq\B(E)$ be the
coefficient-side access algebra.  Its comparison effect interval is
$0\preceq R\preceq I_E$, $R\in\mathcal D$.

Fix a unit physical input $t\in S$, an output target $H\subseteq O$, and
$V_t:H\to\overline S\otimes O$, $V_th=\overline t\otimes h$ as in V4.
Put
\begin{equation}
 \Lambda_t=(I-V_tV_t^*)W,\qquad \mathcal S_t=\Ker\Lambda_t.
 \label{sel5:eq:access-null}
\end{equation}
In a verified represented decomposition of the access algebra write
\begin{equation}
 E=\bigoplus_\alpha\C^{m_\alpha}\otimes\C^{d_\alpha},
 \qquad
 \mathcal D=\bigoplus_\alpha I_{m_\alpha}\otimes M_{d_\alpha}(\C).
 \label{sel5:eq:access-type}
\end{equation}
Let $\iota_\alpha$ denote the inclusion of a block into $E$.  On
$\C^{d_\alpha}$ define the positive matrices
\begin{equation}
 N_\alpha(t)=\sum_{i=1}^{m_\alpha}
   \Lambda_{\alpha i}^*\Lambda_{\alpha i},
 \quad
 \Lambda_{\alpha i}x=\Lambda_t\iota_\alpha(e_i\otimes x),
 \quad L_\alpha(t)=\Ker N_\alpha(t),
 \label{sel5:eq:block-null}
\end{equation}
and the access projection
\begin{equation}
 q_{\mathcal D,t}=\bigoplus_\alpha
                    I_{m_\alpha}\otimes P_{L_\alpha(t)}\in\mathcal D.
 \label{sel5:eq:access-projection}
\end{equation}

\begin{theorem}[Exact access-constrained common-input maximum]
\label{sel5:thm:access-short}
The projection $q_{\mathcal D,t}$ is the largest projection in $\mathcal D$
whose range lies in $\mathcal S_t$.  For $0\preceq R\preceq I_E$ in
$\mathcal D$, the selected map is a common-input map
$X\mapsto\langle t,Xt\rangle F$ with $\Ran F\subseteq H$ if and only if
\begin{equation}
 0\preceq R\preceq q_{\mathcal D,t}.
 \label{sel5:eq:access-interval}
\end{equation}
Consequently the unique largest output frame in this comparison class is
\begin{equation}
 \boxed{
 F_{\mathcal D}^{\max}(t)
   =V_t^*Wq_{\mathcal D,t}W^*V_t.}
 \label{sel5:eq:access-max}
\end{equation}
It is attained by the coefficient effect $q_{\mathcal D,t}$.  Every
admitted common-input outcome whose environment effect is in $\mathcal D$
is bounded above by this frame, even when the whole interval of
$\mathcal D$ has not been admitted.
\end{theorem}
\begin{proof}
The range condition $\Ran(WRW^*)\subseteq\Ran V_t$ is equivalent to
$\Lambda_tR^{1/2}=0$.  This follows because
$WRW^*=(WR^{1/2})(WR^{1/2})^*$ and the range of $AA^*$ equals the range
of $A$ in finite dimension.  Equivalently, $\Ran R\subseteq\mathcal S_t$.
Every spectral support projection of an element of $\mathcal D$ is in
$\mathcal D$.  In the decomposition \eqref{sel5:eq:access-type}, a projection
has the form $\bigoplus I_{m_\alpha}\otimes p_\alpha$.  Its range lies in
$\mathcal S_t$ precisely when $\Lambda_{\alpha i}p_\alpha=0$ for all
$\alpha,i$.  Positivity of the sum in \eqref{sel5:eq:block-null} says this
is equivalent to $p_\alpha\le P_{L_\alpha(t)}$.  This proves maximality
of $q_{\mathcal D,t}$ and \eqref{sel5:eq:access-interval}.

For every such $R$, its output frame is $V_t^*WRW^*V_t$.  Order preservation
and $R\preceq q_{\mathcal D,t}$ give the upper bound in
\eqref{sel5:eq:access-max}; setting $R=q_{\mathcal D,t}$ attains it.
The proof never assumes injectivity of $W$.  In particular, directions in
$\Ker W$ do not contribute spurious nonzero output even if the maximizing
projection contains them.
\end{proof}

\begin{remark}[This is not a claim that all smaller frames are accessible]
For a specified target frame $F$, its exact feasibility problem remains
\[
 R\in\mathcal D,\quad0\preceq R\preceq q_{\mathcal D,t},
 \quad V_t^*WRW^*V_t=F.
\]
It is a finite semidefinite problem.  The condition
$F\preceq F_{\mathcal D}^{\max}(t)$ alone need not solve it, just as the
marked maximum in V4 did not determine the whole marked order interval.
For a genuinely finite effect bank, membership and outcome budgets must be
imposed in addition.  Linear spanning or generation of $\mathcal D$ is not
that membership condition.
\end{remark}

\subsection{A positive mixed frame is not automatically a pure-write frame}
Assume in this subsection that the specified Kraus synthesis $W$ is
injective, as for a minimal dilation of one fixed recorded CP branch.  Define
\begin{equation}
 q_{\mathcal D,t}^{(1)}=
   \bigoplus_{\alpha:m_\alpha=1}P_{L_\alpha(t)},
 \qquad
 F_{\mathcal D}^{(1)}(t)=V_t^*Wq_{\mathcal D,t}^{(1)}W^*V_t,
 \label{sel5:eq:pure-projection}
\end{equation}
with zero blocks in the other summands.

\begin{theorem}[Pure-write capacity of an exposed effect algebra]
\label{sel5:thm:pure-capacity}
Under injectivity of $W$, an accessible nonzero rank-one Choi outcome has
coefficient effect $R=c|x\rangle\langle x|$ supported in one
$m_\alpha=1$ block.  It is a pure common-input write from $t$ exactly when
$x\in L_\alpha(t)$ in that block.  If the effect interval of $\mathcal D$
is available, then
\begin{equation}
 \boxed{
 \dim\Span\{\text{accessible pure output vectors on }t\}
   =\rank F_{\mathcal D}^{(1)}(t)
   =\sum_{\alpha:m_\alpha=1}\dim L_\alpha(t).}
 \label{sel5:eq:pure-capacity}
\end{equation}
The frame $F_{\mathcal D}^{(1)}(t)$ is the largest sum of output frames of
simultaneously selected pure common-input outcomes.  A positive mixed maximum
$F_{\mathcal D}^{\max}(t)$ need not have any accessible pure refinement.
\end{theorem}
\begin{proof}
Injectivity gives $\rank(WRW^*)=\rank R$.  A positive $R\in\mathcal D$
has rank $\sum_\alpha m_\alpha\rank R_\alpha$.  Rank one is possible
exactly in a single multiplicity-one block, with a rank-one $R_\alpha$.
The common-input condition is \eqref{sel5:eq:access-interval}, namely
$x\in L_\alpha(t)$.  Rank-one projections onto orthonormal bases of those
spaces sum to $q_{\mathcal D,t}^{(1)}$ and are a valid simultaneous effect
selection.  Their images are linearly independent: $W$ is injective and
$V_t^*$ is isometric on the containing subspace $\Ran V_t$.
This proves \eqref{sel5:eq:pure-capacity} and attainability.  Every sum of
such effects is at most $q_{\mathcal D,t}^{(1)}$, so order preservation proves
maximality of the pure frame.
\end{proof}

\begin{example}[A full mixed output with no accessible pure component]
\label{sel5:ex:mixed-no-pure}
Take scalar input $S=\C$, output $O=\C^2$, $W=I_2/\sqrt2$ and access
algebra $\mathcal D=\C I_2$.  Then $q_{\mathcal D,t}=I_2$ and
$F_{\mathcal D}^{\max}=I_2/2$, but $m_\alpha=2$ and
$q_{\mathcal D,t}^{(1)}=0$.  The only selected maps prepare scalar multiples
of $I_2/2$.  Its mathematical spectral decomposition is not a readout of
its inaccessible environment.
\end{example}

\begin{proposition}[A quantitative inaccessible-multiplicity floor]
\label{sel5:prop:rank-floor}
Suppose $W$ is injective, with least and greatest singular values
$s_->0$ and $s_+$, and every block multiplicity in $\mathcal D$ is at
least $m_0\ge2$.  Every nonzero selected Choi matrix $J_R=WRW^*$ obeys
\begin{equation}
 \frac{\lambda_1(J_R)}{\Tr J_R}
 \le \left[1+(m_0-1)\frac{s_-^2}{s_+^2}\right]^{-1}<1,
 \label{sel5:eq:rank-floor}
\end{equation}
where eigenvalues are in decreasing order.  Thus scaling an outcome to small
probability does not make its trace-normalized Choi matrix pure.
\end{proposition}
\begin{proof}
The largest eigenvalue of $R$ is repeated at least $m_0$ times.
The nonzero eigenvalues of $J_R$ are those of
$R^{1/2}W^*WR^{1/2}$, which lies between $s_-^2R$ and $s_+^2R$.
The min--max principle gives
$\lambda_j(J_R)\ge s_-^2\lambda_1(R)$ for $1\le j\le m_0$ and
$\lambda_1(J_R)\le s_+^2\lambda_1(R)$.
Summing the first $m_0$ eigenvalues proves the bound.
\end{proof}

\begin{remark}[Why actual redundant environments cannot be discarded silently]
The injectivity qualification is necessary.  If
$W(z_1,z_2)=(z_1+z_2)|b\rangle\!\rangle/\sqrt2$ and
$\mathcal D=\C I_2$, the allowed effect $I_2$ has rank two but its Choi
image is $|b\rangle\!\rangle\langle\!\langle b|$ of rank one.
The general maximum theorem still applies.  The multiplicity rank test does
not.  A physical dilation must be reduced together with its readout map, not
replaced by a minimal environment to which unrestricted new measurements
are then granted.
\end{remark}

\section{Identify the missing environment Read in the actual reset history}
\label{sel5:erased-input}

Retain V3--V4's fresh-register window, its original hub records, and its
physical spaces $F=U\otimes R$, $\dim U=\dim R=2$.  No new window dynamics
is chosen.  Its already evaluated two-step recorded branch is
\begin{equation}
 \mathcal J_{a,b,h}(X)=\tfrac14\Tr(P_aX)\,\omega_{b,h},
 \quad P_a=|a\rangle\langle a|\otimes I_R,
 \quad \omega_{b,h}=|h\rangle\langle h|\otimes\rho_b.
 \label{sel5:eq:original-history}
\end{equation}
Here $a$ is the initial hub, $b$ the intervening hub, and $h$ the final hub.
The pure vectors $r_0,r_1$ spanning $R$ have overlap $1/5$ as supplied by
\src{def:completed-private-packet} in \cite{Grand}.  The iteration policy
remains the conditional fresh-register realization already declared in V3.

\begin{theorem}[The discarded private input is the minimal history environment]
\label{sel5:thm:erased-identification}
For a fixed nonzero recorded branch \eqref{sel5:eq:original-history}, define
$V_{a,b,h}:F\to F\otimes R_{\rm er}$ by
\begin{equation}
 V_{a,b,h}x
 =\tfrac12(|h\rangle\otimes|r_b\rangle)
       \otimes ((\langle a|\otimes I_R)x),
 \label{sel5:eq:erased-dilation}
\end{equation}
where $R_{\rm er}$ is a copy of the initial private input with the displayed
identification.  This is a minimal Stinespring operator for the branch.
For a physical effect $0\preceq F_{\rm er}\preceq I$ on that environment,
its selected map is exactly
\begin{equation}
 \boxed{
 \mathcal J_{a,b,h;F_{\rm er}}(X)
  =\tfrac14\Tr[(|a\rangle\langle a|\otimes F_{\rm er})X]\,
       \omega_{b,h}.}
 \label{sel5:eq:effect-transport}
\end{equation}
Its Choi rank and input-effect rank both equal $\rank F_{\rm er}$.
For $t=|a\rangle\otimes s$, $\|s\|=1$, a nonzero pure common-input
refinement is possible precisely with $F_{\rm er}=\eta p_s$,
$0<\eta\le1$.  The selected amplitude is then
\begin{equation}
 b_{b,h}=\frac{\sqrt\eta}{2}
                  |h\rangle\otimes|r_b\rangle\langle t|.
 \label{sel5:eq:actual-effect-write}
\end{equation}
\end{theorem}
\begin{proof}
Tracing $R_{\rm er}$ in \eqref{sel5:eq:erased-dilation} gives
\eqref{sel5:eq:original-history}.  Its two Kraus operators in an orthonormal
basis $(e_i)$ of $R$ are
$C_i=\tfrac12|h,r_b\rangle\langle a,e_i|$; they are linearly independent
and $W^*W=I_2/4$, so the dilation is minimal.  Inserting $F_{\rm er}$ in
its partial trace gives \eqref{sel5:eq:effect-transport}.  Its Choi operator
is
$\tfrac14(|a\rangle\langle a|\otimes F_{\rm er})^{\mathsf T}
 \otimes\omega_{b,h}$, proving the two rank statements.
The full input effect is supported on $\C t$ exactly when
$F_{\rm er}=\eta p_s$.  Substitution gives the displayed amplitude.
\end{proof}

This is an identification in a minimal dilation, not a claim that a hidden
input register can be accessed physically.  Any equivalent minimal dilation
has the same conclusion after its calibrated environment unitary.  The new
source requirement is precisely access to that registered input effect; an
arbitrary measurement on the newly written $r_b$ is a different operation.

\begin{corollary}[Exact readout requirement for the two private writes]
\label{sel5:cor:readout-floor}
Suppose the old hub labels are retained and an independently supplied access
algebra $\mathcal D_a\subseteq\B(R_{\rm er})$ is available in the common
input-register frame for both branches $b=0,1$ at fixed $a,h$.
Assume its full effect interval is admitted.  The two maximal common-origin
writes on $t=|a\rangle\otimes s$ occur in this refinement class if and only
if
\begin{equation}
 \boxed{p_s\in\mathcal D_a.}
 \label{sel5:eq:readout-floor}
\end{equation}
Their weights are $1/4$ each, their joint success on $p_t$ is $1/2$, and the
private output frame has spectrum $\{3/10,1/5\}$.
On the two-dimensional erased input, the complete algebraic alternatives are
\[
 \begin{array}{c|c}
 \mathcal D_a&\text{possible exact donor rays}\\ \hline
 \C I_2&\text{none}\\
 \C p\oplus\C(I-p)&\Ran p\text{ or }\Ran(I-p)\text{ only}\\
 M_2(\C)&\text{every ray}.
 \end{array}
\]
Thus a binary sharp selection on the erased private input is algebraically
sufficient for a specified donor; full qubit tomography or full matrix
control is not necessary.  A spanning measurement bank alone is not the
full-effect hypothesis.
\end{corollary}
\begin{proof}
By Theorem~\ref{sel5:thm:erased-identification}, every nonzero pure refinement
from this donor needs the positive effect $\eta p_s$; membership in the
linear algebra is equivalent to membership of $p_s$.  If it is available,
use $F_{\rm er}=p_s$ in the two actually recorded branches.  The frame and
weights are the inherited V4 calculation.  A unital $C^*$-subalgebra on
$\C^2$ is scalar, diagonal in one orthonormal basis, or all of $M_2$.
Its rank-one projections give the table.  This can also be seen directly:
a nonscalar Hermitian element supplies its two spectral projections, and
an additional nondiagonal element together with them generates $M_2$.
\end{proof}

For the unrefined hub-record alphabet, the environment effect of each retained
history is $I_{R_{\rm er}}$.  Its access algebra for further selections of
that unexposed register is therefore scalar, and the first row applies.
This is a precise statement about that restricted alphabet.  The general
source framework allows additional declared interventions; their actual maps
must be examined rather than assigned this scalar algebra.

\begin{proposition}[No cross-hub donor with the first hub record preserved]
\label{sel5:prop:hub-rule}
Even arbitrary readout access on the erased private register cannot produce
an exact common-input write on $t$ with $P_0t\ne0$ and $P_1t\ne0$ while
preserving the first hub record.  The obstruction survives arbitrary later
record-respecting processing.
\end{proposition}
\begin{proof}
Every selected Choi operator in a fixed old-hub branch has input support in
$\overline{P_aF}$.  If positive branch operators sum to a common-input
Choi operator supported on $\C\overline t\otimes H$, each summand must
have its range in that same subspace.  Its intersection with each individual
old-hub support is zero.  Every summand therefore vanishes.  Later
record-respecting processing does not restore the missing cross-hub input
blocks.  This is V4's marked support obstruction applied after the explicit
erased-input identification, not a new bypass of it.
\end{proof}

\section{Tomographic completeness does not supply the missing refinement}
\label{sel5:tomography}

The input datum in \cite{Grand} assumes a \emph{spanning physical CP frame}.
The following counterexample retains that strong linear property, not just
an informationally incomplete classical Read.

\begin{theorem}[A complete CP frame with no finite pure-write histories]
\label{sel5:thm:interior-frame}
For every $d\ge2$ there is a finite normalized instrument on $M_d(\C)$
whose branch Choi matrices span all Hermitian Choi coordinates and are all
strictly positive.  Every nonzero positive-length history generated by those
branches, finite classical control, randomization and coarse graining has
strictly positive Choi matrix.  None is a pure common-input write.
Moreover all normalized output states of such histories satisfy a common
strictly positive eigenvalue floor.  After the first frame letter, even
arbitrary terminal system processing cannot remove a corresponding strict
floor from the trace-normalized full input effect.  The full algebraic span
of the Kraus coefficients is nevertheless $M_d(\C)$.
\end{theorem}
\begin{proof}
Let $D=d^2$ and choose a basis $F_1,\ldots,F_{D^2-1}$ of traceless
Hermitian operators on $\C^D$, rescaled to $\|F_j\|_{\op}\le1$.
Choose $0<a<1$ and put
\begin{equation}
 J_{j,\pm}=c(I_D\pm aF_j),\qquad
 c=\frac1{2d(D^2-1)}.
 \label{sel5:eq:interior-frame}
\end{equation}
They are strictly positive and their differences and sums span all Hermitian
matrices.  Their total is $I_D/d$, whose output partial trace is $I_d$;
thus they are the branches of a normalized instrument.

If $J(\Phi)\succeq\epsilon I$ then
$\Phi\ge_{\rm CP}\epsilon\mathcal T$ for
$\mathcal T(X)=\Tr(X)I_d$.  If also $J(\Psi)\succeq\eta I$,
composition preserves CP order and gives
\[
 \Psi\Phi\ge_{\rm CP}\epsilon\eta\mathcal T^2
      =d\epsilon\eta\mathcal T.
\]
Hence every word has positive definite Choi matrix.  A nonzero positive sum
of such histories has the same property.  Classical control chooses words,
not signed combinations, so does not invalidate the argument.
For a positive input $X\ne0$ and any last branch,
\[
 \Phi_{j,\pm}(X)\succeq c(1-a)\Tr(X)I_d,
 \quad
 \Tr\Phi_{j,\pm}(X)\le cd(1+a)\Tr(X).
\]
The normalized output is therefore at least
\begin{equation}
 \kappa I_d,\qquad \kappa=\frac{1-a}{d(1+a)}>0,
 \label{sel5:eq:uniform-output-floor}
\end{equation}
independently of the earlier word.  Mixtures preserve this floor.
For the input-effect assertion, every terminal selected system operation has
an effect $F\succeq0$ at the output of the first branch.  The first-branch
Choi bounds imply, by taking dual CP maps,
\[
 c(1-a)\Tr(F)I_d\preceq\Phi_{j,\pm}^*(F)
                \preceq c(1+a)\Tr(F)I_d.
\]
If the resulting effect is nonzero, divide by its trace and obtain the same
floor $\kappa I_d$.  This remains true after positive sums of histories,
including adaptive choices of terminal effects.  No inverse map or
measurement of the first branch's environment is included in terminal system
processing.  Finally,
any positive definite Choi matrix has full coefficient support
$\B(\C^d)$, so the Kraus linear span is already $M_d$.  That linear span
is not a list of its individually executed pure outcomes.
\end{proof}

\begin{remark}[Scope of the counterexample]
This theorem does not assign the displayed interior frame to the historical
Store.  It shows why the general tomographic entrance cannot logically supply
a particular rank-one refinement without specifying its operation bank.
Additional preparation, intervention, ancillary coupling, or a sharper Read
may enlarge that bank.  Those are additional actual maps, not consequences of
linear reconstruction.  The zero-length identity operation has full input
effect and does not supply the missing donor either.
\end{remark}

\begin{proposition}[Even a full effect algebra can hide a finite-depth rank floor]
\label{sel5:prop:weak-IC}
On a qubit the six effects
\begin{equation}
 E_{j,\pm}=\frac16(I\pm a\sigma_j),\qquad j=x,y,z,
 \quad0<a<1,
 \label{sel5:eq:IC-effects}
\end{equation}
form an informationally complete POVM.  If their actual instrument has
L\"uders coefficients $L_{j,\pm}=E_{j,\pm}^{1/2}$, every finite nonempty
outcome word has an invertible coefficient and a positive definite input
effect.  A nonzero coarse-grained finite selection has a positive definite
effect too, despite $C^*(E_{j,\pm})=M_2$.
\end{proposition}
\begin{proof}
The six effects sum to identity, and their sums and differences span the
identity and Pauli matrices.  Every eigenvalue is strictly positive, so every
$L_{j,\pm}$ is invertible; the same holds for any product.
The effect of a coarse-grained event is a nonzero positive sum of the word
effects, so remains positive definite.  This includes adaptive choices among
a finite bank: each realized word is still a product of invertible operators.
\end{proof}

The last proposition is a finite-depth statement.  A nonzero-probability
\emph{limit} of increasingly many readout outcomes can approach a sharp Read.
The next section proves precisely such a positive alternative, rather than
incorrectly extending the finite obstruction to all asymptotic protocols.

\section{Weak-read amplification with a nonvanishing success budget}
\label{sel5:weak-amplification}

Let $p=|s\rangle\langle s|$ be a physically identified qubit direction and
$q=I-p$.  Assume the actual exposed readout instrument, before input erasure
or on its coherently preserved erased-input copy, has two coefficients
\begin{equation}
 L_+=\sqrt{\frac{1+a}{2}}p+\sqrt{\frac{1-a}{2}}q,
 \qquad
 L_-=\sqrt{\frac{1-a}{2}}p+\sqrt{\frac{1+a}{2}}q,
 \quad0<a<1.
 \label{sel5:eq:QND}
\end{equation}
Repeated use means the same qubit is retained, no unmodelled dynamics
intervene, and these two state-update maps recur.  The assumption is about
the instrument, not just two POVM probabilities.  Such an operation is
nondemolition in the indicated eigenbasis but can disturb other private
states.  Neither its occurrence nor that eigenbasis is inferred from the
source contrast line.

\begin{proposition}[The candidate donor is reconstructed from the admitted Read]
\label{sel5:prop:read-native-ray}
Suppose an actual binary qubit instrument has effects $E_+$ and $I-E_+$,
with $0\prec E_+\prec I$, $\Tr E_+=1$, and $E_+$ nonscalar.  Suppose its
state updates are the L\"uders maps for these effects.  Then the parameters
in \eqref{sel5:eq:QND} are reconstructed by
\[
 a=\lambda_{\max}(E_+)-\lambda_{\min}(E_+),\qquad
 p=\frac{E_+-\lambda_{\min}(E_+)I}{a},\qquad q=I-p.
\]
They satisfy $0<a<1$ and $p$ has rank one.  Thus the physical candidate
donor ray can be an output of the measured Read, not a direction selected to
fit the colour target.  Its identification with the categorical source line
is still a different, required physical comparison.
\end{proposition}
\begin{proof}
The two eigenvalues are positive, distinct, sum to one, and lie below one;
they are therefore $(1+a)/2$ and $(1-a)/2$.  Spectral calculus gives the
rank-one projector and the two L\"uders coefficients in
\eqref{sel5:eq:QND}.  Conjugating the actual input frame conjugates $p$,
so the reconstruction has the appropriate unitary covariance.
\end{proof}

For an odd $n$, retain all length-$n$ readout strings and declare success for
the classical event with more $+$ than $-$ outcomes.  This is one allowed
coarse graining of actually recorded strings, not a coherent sum of them.
Put $p_+=(1+a)/2$, $p_-=(1-a)/2$ and
\begin{equation}
 \delta_n=
 \sum_{k=(n+1)/2}^{n}\binom nk p_-^k p_+^{n-k}.
 \label{sel5:eq:binomial-error}
\end{equation}

\begin{theorem}[An exact finite majority Read and its exponential error]
\label{sel5:thm:majority}
The input effect of the majority-success event is
\begin{equation}
 \boxed{E_{n,+}=(1-\delta_n)p+\delta_n q,\qquad
     0<\delta_n<\tfrac12,
     \delta_n\le(1-a^2)^{n/2}.}
 \label{sel5:eq:majority-effect}
\end{equation}
In particular
$\|E_{n,+}-p\|_{\op}=\delta_n$, while its success on $p$ tends to one.
Every finite $n$ still has input-effect rank two.  If $a=a_n$ varies with
$n$, the sufficient limit condition $na_n^2\to\infty$ makes
$\delta_n\to0$.
\end{theorem}
\begin{proof}
For a string with $k$ plus outcomes the product coefficient is
$p_+^{k/2}p_-^{(n-k)/2}p+
 p_-^{k/2}p_+^{(n-k)/2}q$.
Its effect is the sum of the squared coefficients times $p,q$.
Summing over the successful strings gives the two binomial tail
probabilities.  Swapping plus and minus bijects success and failure for odd
$n$, so the two coefficients sum to one and the $q$ coefficient is
\eqref{sel5:eq:binomial-error}.  Pairing each upper-tail binomial term with
its lower-tail reflection shows $\delta_n<1/2$ because $p_-<p_+$;
positivity is strict since $p_->0$.

For the exponential bound let $K$ be binomial with parameters $n,p_-$ and
choose $z=p_+/p_->1$.  Markov's inequality gives
\[
 \Pr(K\ge n/2)\le z^{-n/2}(p_++p_-z)^n
       =(2\sqrt{p_+p_-})^n=(1-a^2)^{n/2}.
\]
The majority tail is no larger.  The operator-norm formula follows from the
two eigenvalues of $E_{n,+}-p=\delta_n(q-p)$.
Finally $(1-a_n^2)^{n/2}\le e^{-na_n^2/2}$ gives the stated sufficient
variable-strength condition.
\end{proof}

\begin{remark}[Why one long postselected word is not the same protocol]
The all-plus word alone has effect $p_+^np+p_-^nq$.
Its trace-normalized effect approaches $p$, but its success on $p$ is
$p_+^n\to0$.  Majority coarse graining retains many distinguishable strings
and has success $1-\delta_n\to1$.  The improvement uses only classical
processing of recorded outcomes; it does not coherently recombine histories.
It does require the stated repeated instrument to be physically executable.
\end{remark}

\subsection{Apply the majority Read to the evaluated private write}
Fix the first hub $a$ and final hub $h$, put $t=|a\rangle\otimes s$, and
keep the original two branches $b=0,1$ of
\eqref{sel5:eq:original-history}.  Insert the repeated weak instrument on the
old private input \emph{before it is discarded}, or equivalently on the
specified surviving erased-input copy if that access is actually available.
Retain the old hub labels throughout.

\begin{theorem}[Approximate common-origin private writes at fixed total success]
\label{sel5:thm:weak-writes}
For each $b$ the actually coarse-grained success map is
\begin{equation}
 \mathcal J_{b,n}^{+}(X)
  =\tfrac14\Tr[(|a\rangle\langle a|\otimes E_{n,+})X]\,
       \omega_{b,h}.
 \label{sel5:eq:approximate-write}
\end{equation}
Let $\mathcal K_b(X)=\tfrac14\langle t,Xt\rangle\omega_{b,h}$ be its
ideal common-input comparison map.  With the unnormalized diamond norm,
\begin{equation}
 \boxed{
 \|\mathcal J_{b,n}^{+}-\mathcal K_b\|_\diamond=\delta_n/4,
 \qquad
 \left\|\sum_b\mathcal J_{b,n}^{+}-\sum_b\mathcal K_b\right\|_\diamond
       =\delta_n/2.}
 \label{sel5:eq:diamond-writes}
\end{equation}
The same $\delta_n/2$ holds for the direct-sum output retaining $b$.
The combined success probability on the desired input $p_t$ is
$(1-\delta_n)/2$, tending to V4's optimal $1/2$, and the false success on
$|a\rangle\otimes s^\perp$ is $\delta_n/2$.
Thus the ideal two-write packet is a diamond-norm limit of actual
record-respecting operations under the repeated weak-read hypothesis.
It is not an exact finite-$n$ outcome.
\end{theorem}
\begin{proof}
The two-step atomic writer sees only the input effect of the preceding
measurement.  The effect is \eqref{sel5:eq:majority-effect}, so composition
or \eqref{sel5:eq:effect-transport} gives
\eqref{sel5:eq:approximate-write}.
The difference from $\mathcal K_b$ is
$X\mapsto\Tr(A_nX)\omega_{b,h}$ with
$A_n=\tfrac14\delta_n|a\rangle\langle a|\otimes(q-p)$.
As proved in V4's diamond-distance argument, the diamond norm of
$X\mapsto\Tr(AX)B$ is $\|A\|_{\op}\|B\|_1$ for Hermitian $A$
and positive $B$.  Here $\|A_n\|=\delta_n/4$ and
$\|\omega_{b,h}\|_1=1$.  Summing the outputs gives a positive operator of
trace two, whether or not they occupy orthogonal record blocks; this proves
all three norm equalities.  Direct evaluation on $p_t$ and its orthogonal
private donor gives the success and false-success probabilities.
\end{proof}

\begin{corollary}[The added Read is conservative for the old source records]
\label{sel5:cor:conservative-read}
For any normalized instrument on the initial private register, its insertion
before the discarded-input reset, followed by forgetting its new outcomes,
leaves every old branch \eqref{sel5:eq:original-history} exactly unchanged.
In particular the majority-success and majority-failure branches sum to the
original private writer.  The weak-read construction changes neither its
unconditioned two-step CP map nor its old hub-record probabilities.
\end{corollary}
\begin{proof}
Let $M_\nu$ be the inserted private coefficients, with
$\sum_\nu M_\nu^*M_\nu=I_R$.  The pulled-back old input effect is
\[
 \sum_\nu(I_U\otimes M_\nu)^*P_a(I_U\otimes M_\nu)=P_a.
\]
The old branch's output is fixed, so its complete CP map is unchanged.
Iterate the same identity for the normalized weak instrument and partition
its strings into majority success and failure.  Equivalently, their effects
$E_{n,+}$ and $I-E_{n,+}$ sum to identity.
\end{proof}

This is a conservative \emph{readable enrichment} conditional on the actual
extra intervention, not a proof that the original alphabet contained it.
Knowing its forgetful channel cannot determine which new readout, if any, was
performed.  The construction supplies the exact incremental operation to audit
without altering the previously evaluated reset law.

The limiting private frame is still the already evaluated
$(\rho_0+\rho_1)/4$, with eigenvalues $3/10,1/5$ and determinant $3/50$.
The theorem does not rederive these constants or identify $t$ with the
categorical source ray.  It removes the \emph{sharpness} of a readout from a
sufficient cofinal protocol, not the need to supply its physical input cut,
its nondemolition direction, and its executable repeated state update.

\begin{proposition}[A finite accumulated implementation budget]
\label{sel5:prop:execution-error}
Suppose the full two-outcome readout channel, including its classical flag,
is implemented at each step with diamond error at most $\eta_n$, while
between steps the declared nondemolition system is otherwise retained.
Let $\widehat{\mathcal J}_n$ be the resulting selected private-write map,
with $b$ retained.  If all exact and implemented full instruments are
normalized, then the sufficient estimate is
\begin{equation}
 \left\|\widehat{\mathcal J}_n-
              \bigoplus_b\mathcal K_b\right\|_\diamond
 \le n\eta_n+\delta_n/2.
 \label{sel5:eq:implementation-budget}
\end{equation}
Consequently $na_n^2\to\infty$ and $n\eta_n\to0$ suffice for the same
nonvanishing-success limit.  Fixed per-step errors are not silently divided
by the selected success probability.
\end{proposition}
\begin{proof}
Telescope the difference of the composed flagged instruments, keeping their
classical histories in the output.  A CPTP map has diamond norm one, so each
of the $n$ terms is bounded by $\eta_n$.  Majority selection and the
subsequent normalized source followed by selection are CP trace-nonincreasing
maps and have diamond norm at most one.  They cannot increase the bound.
Add the exact comparison error from
Theorem~\ref{sel5:thm:weak-writes}.
\end{proof}

\section{A vanishing channel error need not lock the finite internal algebra}
\label{sel5:algebra-limit}

The preceding positive approximation must not be mistaken for exact source
selection.  Its residual wrong-donor amplitude can alter the generated algebra
at every finite stage.  The following comparison retains that amplitude
rather than projecting it away.

Let $T=\C t$, $T'=\C u$, and $H\cong\C^2$ be mutually orthogonal
subspaces of a four-dimensional declared linking carrier.  Let
$r_0,r_1\in H$ be unit independent vectors.  For $0\le\delta<1$ put
\begin{equation}
 b_{j,t}(\delta)=\frac{\sqrt{1-\delta}}2|r_j\rangle\langle t|,
 \qquad
 b_{j,u}(\delta)=\frac{\sqrt\delta}2|r_j\rangle\langle u|,
 \qquad j=0,1.
 \label{sel5:eq:linking-coefficients}
\end{equation}
These are the coefficient supports of
\eqref{sel5:eq:approximate-write} under this explicit linking-carrier
identification.  That identification is a testbed hypothesis, not an
identification of the physical input and output cuts by their equal dimension.

\begin{theorem}[Exact algebra jump despite fixed-success write convergence]
\label{sel5:thm:algebra-jump}
The unital coefficient $C^*$-algebra generated by
\eqref{sel5:eq:linking-coefficients} is
\begin{equation}
 \boxed{
 \A_\delta=M_4(\C)\quad(0<\delta<1),\qquad
 \A_0=M_3(T\oplus H)\oplus\C p_{T'}.}
 \label{sel5:eq:algebra-jump}
\end{equation}
Thus its commutant has dimension one for every positive $\delta$ and
dimension two at zero.  For the calibrated form
\[
 q_\delta(X)=\sum_{j=0}^1\sum_{v\in\{t,u\}}
    \left(\|[b_{j,v}(\delta),X]\|_{\HS}^2+
          \|[b_{j,v}(\delta)^*,X]\|_{\HS}^2\right),
\]
the least eigenvalue on the traceless space satisfies
\begin{equation}
 \boxed{0<\gamma_\delta\le\frac43\delta\quad(0<\delta<1).}
 \label{sel5:eq:gap-collapse}
\end{equation}
In particular the majority protocol has $\gamma_{\delta_n}\to0$,
even though its write error tends to zero with success tending to $1/2$.
\end{theorem}
\begin{proof}
For $\delta>0$, the two $t$-donor coefficients span all $H\leftarrow T$
rank-one maps, and the two $u$-donor coefficients span all
$H\leftarrow T'$ maps.  Products with adjoints give all matrix units on
$H$, both donor projections, and the donor--donor rectangular units.
Together these span $M_4$.  At zero the $u$-donor coefficients vanish.
The same argument on $T\oplus H$ gives its full $M_3$; the remaining
unit supplies the independent scalar on $T'$.  Taking commutants proves the
dimension statements and positivity of the traceless gap for $\delta>0$.

For its upper bound use the traceless Hermitian vector
$X=p_{T'}-\tfrac13p_{T\oplus H}$, whose squared HS norm is $4/3$.
It commutes with both strong $t$-donor coefficients.  Across either weak
$u$-donor edge the eigenvalue difference is $4/3$.  Each weak coefficient
has squared HS norm $\delta/4$ and its adjoint contributes the same amount.
Therefore
\[
 q_\delta(X)=2\cdot2\cdot\frac\delta4\left(\frac43\right)^2
             =\frac{16\delta}{9}.
\]
Its Rayleigh quotient is $4\delta/3$, proving
\eqref{sel5:eq:gap-collapse}.
\end{proof}

\begin{remark}[Where a cofinal locking theorem would need an extra hypothesis]
This does not contradict the compact-Mosco or finite-anchor results in the
duality manuscript.  Protecting only the scalar commutant leaves an extra
limiting zero mode, so the exact-limit-kernel hypothesis fails.  Protecting
the two-dimensional limiting commutant instead fails its exact inclusion in
the finite kernels.  Approximate channel convergence alone supplies neither
missing condition.  A cutoff-dependent comparison with a larger limiting
algebra is meaningful, but it cannot be described as exact finite-kernel
locking without a separate protected-support mechanism.
\end{remark}

\begin{remark}[No fictitious Kraus outcome is introduced]
The algebra in \eqref{sel5:eq:algebra-jump} is the algebra of the complete
Kraus span of the selected CP maps.  Its computation does not assert that
the two coefficients of a rank-two map were resolved as separate historical
outcomes.  Selecting only the leading Choi eigenvector would be a further
refinement and would erase the very wrong-donor term whose algebraic effect
is being tested.
\end{remark}

\section{Source status and the remaining finite test}
\label{sel5:status}

The continuation has replaced an unspecified environment refinement by an
exact access question.  On each completed two-step history the minimal
refinement environment is the discarded private input, with effect transport
\eqref{sel5:eq:effect-transport}.  For a fixed input-hub sector the exact
common donor needs precisely its private input projector in the physically
admitted effect cone.  For a donor crossing the input-hub sectors the old
record obstruction still excludes it.  The locked coefficient contrast from
V4 is neither of these physical input objects.

The general constrained short is a positive finite compiler once the actual
dilation and access algebra are supplied.  The pure-write capacity theorem
then distinguishes a mixed maximum from genuinely rank-one outcomes.  Its
positive attainability statements require an effect-complete access bank;
its negative bounds already apply to any smaller bank contained in that
algebra.  Tomographic spanning is not effect completeness, as the normalized
interior-frame counterexample proves at the CP-map level.

There is also a controlled positive route without an ideal sharp input Read.
If the source actually admits repeated nondemolition weak access to the
identified erased input, classical majority selection approaches the desired
write without vanishing success.  The explicit error is $\delta_n/2$ on
the selected two-write instrument, plus the accumulated implementation
budget.  This is not an automatic property of normalization, recurrence,
nonorthogonal output states, or the scalar completed-return law.

\begin{center}
\small
\begin{tabular}{L{.37\textwidth}L{.55\textwidth}}
\toprule
Question & Status after V5\\
\midrule
Which missing readout implements a common donor in the evaluated history?
& A rank-one effect on the discarded initial private input; the transfer map
is explicit.\\[.35em]
Can the actual access algebra be included in the Choi short?
& Yes. A blockwise kernel computation gives the maximal allowed frame and,
for a faithful Kraus synthesis, its pure-write capacity.\\[.35em]
Does a spanning physical CP frame supply the required outcome?
& No. A normalized interior frame spans every Choi coordinate yet has no
finite pure-write history or asymptotically pure output in that grammar.\\[.35em]
Must a sufficient asymptotic protocol assume a sharp primitive?
& No. An actual informative nondemolition weak Read, retained input, and
classical majority processing give a nonvanishing-success limit.\\[.35em]
Does that imply exact finite Standard-Model selection?
& No. The explicit small wrong-donor term changes $M_3\oplus\C$ to $M_4$
and closes the relative gap as it vanishes.\\[.35em]
Are the requisite erased-input access and categorical identification in the
historical source now verified?
& Not identified in the inspected source rows. They are not assigned zero,
nor supplied by choosing the comparison instrument.\\
\bottomrule
\end{tabular}
\end{center}

The next noncircular source task is to inspect the \emph{actual state-update
map} of an exposed pre-erasure or erased-input Read, if the source declares
one, together with its cut identification and retained records.  Its
Radon--Nikodym effect can be tested directly against the constrained short.
If only a weak instrument occurs, its repeatability, information strength,
duration and error budget must be established before the majority theorem
applies.  If no such access is part of the declared source, the present
readout route stops there; a proof by a different nonuniform system route
would require that route's own evaluated maps.  There is no reason to repeat
the completed reset powers or to introduce another favorable recurrent
completion.

The determinant and matter tensor types remain separate source conditions.
The note proves neither that all admissible microscopic sources choose the
internal Standard-Model branch nor that this access class exhausts every
possible selection mechanism.

\subsection{Verification and cumulative integrity}
The V4 byte prefix before its terminal document command is preserved by
appending this block immediately before that command.  The accompanying
exact-arithmetic script checks access projections in faithful and redundant
dilations, the pure/mixed-frame distinction, the erased-input effect map,
the complete positive frame, finite majority probabilities, exact write
errors and success, the algebra dimensions $16$ and $10$, and the
$4\delta/3$ Rayleigh quotient.  It also reruns the V4 example checks.
General all-depth and limiting statements are proved above; they are not
inferred from finitely many numerical samples.  No physical experiment,
formal proof-assistant certification, or independent peer review is claimed.

\begin{thebibliography}{99}
\bibitem[12]{sel5-BB}
M.~Bauer and D.~Bernard,
\emph{Convergence of repeated quantum nondemolition measurements and
wave-function collapse}, Phys.\ Rev.\ A \textbf{84} (2011), 044103;
\href{https://doi.org/10.1103/PhysRevA.84.044103}{doi:10.1103/PhysRevA.84.044103};
\href{https://arxiv.org/abs/1106.4953}{arXiv:1106.4953}.
The general nondemolition-collapse phenomenon is established background.
The finite majority-selection budget and its use for the present recorded
private writer are proved in this block.
\end{thebibliography}

% END OF SOURCE-SELECTION V5 DEVELOPMENT BLOCK.
% Append further work before the terminal document command.
% Preserve the V1--V5 mathematical bodies and identify corrections explicitly.


\clearpage
\hypersetup{pdftitle={Historical Standard-Model Source Selection: Research Notes V1--V6}}
\begin{center}
{\Large\bfseries V6: Put the missing Read at a certified source cut}\\[.4em]
{\large Complete-future conservation, arbitrary informative updates,\\
identity-branch disturbance, and the rare-opportunity budget}
\end{center}
\addcontentsline{toc}{part}{V6: Read location, complete futures, and opportunity budgets}

\section{The upstream question is where the Read can occur}
\label{sel6:context}

V5 proves that a physically supplied weak Read can approach a common-input
private write with nonvanishing conditional success.  Its conservation
statement is exact for the displayed completed-private reset: after the new
outcomes are forgotten, that reset and its retained hub records are unchanged.
It does not identify the insertion with every earlier physical cut.  At an
earlier cut, an old admitted future can still carry quantum input information
which the same Read changes.  We therefore inspect the source's complete
future interface and, in particular, its already specified no-response branch.

The new positive result removes V5's special symmetric L\"uders update from
one sufficient amplification theorem.  A source-native conserved binary
future projection forces any compatible actual qubit instrument to be
nondemolition; a single informative record distribution then supplies a
likelihood-selection protocol.  The negative and quantitative results locate
its scope.  The actual locked source has a positive identity branch.  No
informative transparent insertion can preserve that complete quantum branch,
and every exact refinement of a fully resolved history is only a classical
splitting of that history.  The calibrated no-response probability also gives
a sharp universal upper budget on information and pure-write success.  These are
source-evaluated constraints, not another favourable controller completion.

\subsection{Audit and nonduplication}
\begin{center}
\small
\begin{tabular}{L{.38\textwidth}L{.54\textwidth}}
\toprule
Inherited source row & Its role here\\
\midrule
\cite{Grand}, \src{def:bare-predictive-renewal}
& Actual typed intervention and Read alphabets, complete probabilities on
exposed quantum ports, and their causal composition.  A spanning frame does
not assert that every desired operation is admitted.\\[.3em]
\cite{Grand}, \src{def:locked-opportunity-instrument}
& The exact branch $L_\varnothing=\sqrt{1-\vartheta}I$ and the calibrated
$\vartheta=576851/417942208512$.  No individual unknown $K_e$ is fitted.\\[.3em]
V4's common-input and marked Choi shorts
& Feasible refined write frames and their record constraints remain the
existing CP-order tools.  They are not reproved.\\[.3em]
V5's erased-input identification, majority law and old-reset conservation
& The reset, input-copy meaning, conditional write normalization and
special symmetric majority protocol are retained exactly.\\[.3em]
\cite{Grand} future saturation and \cite{Rigidity} fixed-algebra arguments
& Finite Heisenberg saturation and the multiplicative-domain identity are
standard inherited mechanisms; below they are applied to insertion
conservation, not advertised as new abstract reconstruction theorems.\\
\bottomrule
\end{tabular}
\end{center}

The searched source rows do not supply an evaluated historical pre-erasure
Read on the categorical type input.  They do supply the locked identity
branch, so that branch can be evaluated without resolving the missing
Read.  Nondisturbance and commutation are established subjects
\cite{sel6-BS}; the general information/coherence inequality used below has
its standard interferometric counterpart in \cite{sel6-Englert}.  Our
contribution here is the complete-cut test, its combination with the actual
source branch and rare-event calibration, and the update-independent
private-write implication.  No claim of literature priority is made for
those general measurement principles.

\subsection{A scope clarification, not a replacement of V5}
\label{sel6:scope}
Three changes of protocol must be distinguished.
An \emph{available intervention} is a new command which can change old
future probabilities when it is deliberately executed.  A \emph{transparent
insertion} is an intervention at a fixed old cut whose new outcomes may be
forgotten without changing any specified old future there.  A
\emph{record-respecting refinement} decomposes each old CP branch into
positive subbranches whose sum is exactly that same old branch.  The last
two are different from merely enlarging the alphabet.

No theorem below forbids an actual disturbing Read as an independently
admitted control.  Nor does it forbid a refinement inside a genuinely
erasing branch.  It forbids using conservation of a later reset as proof of
transparency at an earlier, coherently observable cut.  The old V5 statement
in Corollary~\ref{sel5:cor:conservative-read} remains true on its explicitly
specified reset interface.

All spaces are finite-dimensional, traces are unnormalized, and diamond
norms have no factor of $1/2$.  Quantum ports, old classical records,
retained input copies and geometric coefficient spaces are not identified
without a displayed source map.  A physical reference is never acted on by
the inserted instrument.  Where complete channel equality is translated into
an operational test, the input preparations and output effect frames are
assumed to be actually exposed and spanning, as in the declared tomographic
entrance.  Otherwise the smaller actual future interface must be used.

\section{The complete-future conservation test at one physical cut}
\label{sel6:futures}

At an exposed port $S$, let $\mathcal F_w$ be the CP map of an admitted
old future, including its retained branch labels and quantum terminal output.
If $E$ is one of its actual terminal effects, put
\begin{equation}
 F_{w,E}=\mathcal F_w^*(E),\qquad
 \mathcal V_{\rm fut}=\Span_\C\{I,F_{w,E}:w,E\}\subseteq\B(S).
 \label{sel6:eq:future-space}
\end{equation}
It is a self-adjoint operator space, not automatically an algebra.  An
inserted finite instrument has branches $\mathcal I_y$, total channel
$\Lambda=\sum_y\mathcal I_y$, and Kraus operators $M_{y\alpha}$ satisfying
$\sum_{y,\alpha}M_{y\alpha}^*M_{y\alpha}=I$.

\begin{proposition}[Exact future test and a finite witness]
\label{sel6:prop:future-test}
The insertion is transparent at this cut if and only if
\begin{equation}
 \boxed{\Lambda^*F=F\qquad(F\in\mathcal V_{\rm fut}).}
 \label{sel6:eq:future-fixed}
\end{equation}
For a basis $(F_\ell)$ chosen from the actual pulled-back effects and the
identity, the finite residual
\begin{equation}
 \mathfrak c_{\rm fut}(\Lambda)
 =\sum_\ell\|\Lambda^*F_\ell-F_\ell\|_{\HS}^2
 \label{sel6:eq:future-residual}
\end{equation}
vanishes exactly under \eqref{sel6:eq:future-fixed}.  If it is positive,
some basis future and some member of an actual spanning preparation frame
have changed probability.
For a finite closed-carrier grammar of old CP letters $\Phi_a$, the future
space is computed by the iteration
\[
 V_0=\Span_\C\{I,\text{terminal effects}\},\qquad
 V_{k+1}=V_k+\sum_a\Phi_a^*(V_k).
\]
A first exact plateau is permanent; at most $(\dim S)^2-\dim V_0$ strict
rank increases occur.  With a finite scheduling memory, the same calculation
is performed on the actually retained typed carrier, not on a freely
composed alphabet with the scheduling restrictions discarded.
\end{proposition}
\begin{proof}
The difference in a future probability on input $\rho$ is
$\Tr\rho(\Lambda^*F_{w,E}-F_{w,E})$.  Vanishing on a spanning preparation
frame is equivalent to vanishing of that Hermitian operator.  Linear span
then gives \eqref{sel6:eq:future-fixed}; its basis version is
\eqref{sel6:eq:future-residual}.  If a residual operator is nonzero, its
pairing cannot vanish on every member of a spanning preparation frame,
which gives the claimed actual test row.  The saturation assertion is the
usual finite forward/Heisenberg span argument: $V_{k+1}=V_k$ implies
$\Phi_a^*V_k\subseteq V_k$ for every letter.  Induction excludes all later
increases; every strict inclusion increases a dimension bounded by
$(\dim S)^2$.  Typed source and target blocks or a retained finite
controller implement the same recursion without inserting forbidden words.
\end{proof}

\begin{remark}[No unsupported algebra completion]
In general, fixing $\mathcal V_{\rm fut}$ does not imply fixing the ordinary
algebra it generates.  The fixed-operator-system examples already supplied
in \cite{Rigidity} are a reason not to make that replacement.  If the
actual determining future space is an algebra, or contains the needed
Hermitian generators together with their squares, the familiar
multiplicative-domain argument does apply.  On a qubit the complete
classification below needs no additional square observations: a nonscalar
Hermitian operator has only two eigenvalues, so its spectral projections
are affine functions of it and the identity.
\end{remark}

For a local private Read $\id_U\otimes\Lambda_R$ at a physically identified
port $U\otimes R$, take the coefficient slices of every $F\in
\mathcal V_{\rm fut}$ in a Hermitian basis on $U$.  Their span, together
with $I_R$, is denoted $\mathcal V_R$.  The local insertion is transparent
exactly when $\Lambda_R^*$ fixes $\mathcal V_R$ pointwise; linear
independence of the $U$-basis proves this assertion.  This construction does
not replace a tensor-product private slot by an unrelated type-line
summand.

\begin{lemma}[A conserved sharp effect constrains every branch coefficient]
\label{sel6:lem:projection}
For a normalized channel $\Lambda$ and a projection $p$,
\begin{equation}
 \Lambda^*(p)=p
 \quad\Longleftrightarrow\quad
 [M_{y\alpha},p]=0\quad\text{for all }y,\alpha.
 \label{sel6:eq:projection-fixed}
\end{equation}
More generally, if a Hermitian $A$ and $A^2$ are fixed, every coefficient
commutes with $A$.
\end{lemma}
\begin{proof}
For $x\in\Ker p$, fixedness gives
$\sum_{y,\alpha}\|pM_{y\alpha}x\|^2=0$, hence no coefficient maps
$\Ker p$ into $\Ran p$.  Normalization and fixedness of $I-p$ give the
opposite block condition.  This proves commutation; the converse follows by
normalization.  For the last statement the exact positive identity is
\[
 \sum_{y,\alpha}[M_{y\alpha},A]^*[M_{y\alpha},A]
 =\Lambda^*(A^2)-\Lambda^*(A)A-A\Lambda^*(A)+A^2.
\]
Its right side is zero, so every summand vanishes.  This is the standard
finite multiplicative-domain identity, recorded to fix its precise use.
\end{proof}

\begin{theorem}[Complete qubit alternatives for transparent access]
\label{sel6:thm:qubit-futures}
Let $R\cong\C^2$ and let $\mathcal V_R$ be its actual self-adjoint future
space containing $I$.  Let $d$ be the real dimension of its traceless
Hermitian part.  Every normalized inserted instrument is classified as
follows.
\begin{enumerate}[label=\textup{(\roman*)},leftmargin=2.5em]
\item If $d=0$, every normalized instrument is transparent relative to this
future interface.
\item If $d=1$, write $\mathcal V_R=\Span_\C\{p,I-p\}$, where $p$ is
rank one and is reconstructed from any nonscalar actual future effect.
Transparency is equivalent to
\begin{equation}
 M_{y\alpha}=a_{y\alpha}p+b_{y\alpha}(I-p)
 \quad\text{for every coefficient}.
 \label{sel6:eq:diagonal-updates}
\end{equation}
\item If $d\ge2$, transparency forces
$M_{y\alpha}=a_{y\alpha}I$ for every coefficient.  Thus
$\mathcal I_y=c_y\id$, $\sum_yc_y=1$, and the new outcome record is
independent of the input.
\end{enumerate}
In particular, two noncommuting qubit future effects exclude every
informative transparent insertion, even if their span is not
informationally complete.  A single nonscalar future effect instead
identifies a native binary axis on which an actual informative compatible
instrument can operate.
\end{theorem}
\begin{proof}
The first case is trace preservation.  In the second case, if $A=A^*$ is a
nonscalar future operator with eigenvalues $\lambda_+>\lambda_-$, then
$p=(A-\lambda_-I)/(\lambda_+-\lambda_-)$ belongs to the same span.
Apply Lemma~\ref{sel6:lem:projection}; the commutant of a rank-one qubit
projection is precisely its diagonal algebra.  If $d\ge2$, choose two
linearly independent traceless Hermitian members.  Their Pauli vectors are
nonparallel, so they do not commute.  The preceding argument makes every
coefficient commute with both spectral projections, whose common commutant
on $\C^2$ is scalar.  Alternatively, a diagonal nonscalar coefficient in
the first basis could not commute with the nondiagonal part of the second
projection.  The scalar coefficients give the asserted branch maps.
\end{proof}

This classifies compatibility, not availability.  In the second case the
allowed full comparison class contains a sharp Read of $p$, but the source
need not execute that Read.  Conversely, if its actual operation table does
contain an informative transparent instrument, the special update assumption
can now be removed as follows.

\section{Amplification from a measured conservation law, not a chosen L\"uders update}
\label{sel6:general-amplification}

Assume the $d=1$ case of Theorem~\ref{sel6:thm:qubit-futures}, and suppose
one actual finite instrument is transparent there and can be repeated on the
same retained qubit without additional unmodelled dynamics or memory.  Its
coefficients have \eqref{sel6:eq:diagonal-updates}.  The outcome laws on the
two native pure states are
\begin{equation}
 P_y=\sum_\alpha|a_{y\alpha}|^2,\qquad
 Q_y=\sum_\alpha|b_{y\alpha}|^2,\qquad
 \sum_yP_y=\sum_yQ_y=1.
 \label{sel6:eq:native-record-laws}
\end{equation}
They are measured probabilities; hidden Kraus labels $\alpha$ need not be
resolved.  Put $b_{\rm cl}=\sum_y\sqrt{P_yQ_y}$.

\begin{theorem}[General record-likelihood amplification]
\label{sel6:thm:likelihood}
For each $n$, let $A_n$ be the actual classical length-$n$ record event on
which $P^{\otimes n}(w)\ge Q^{\otimes n}(w)$, with any fixed assignment
of equality cases.  Write
$\epsilon_p=P^{\otimes n}(A_n^c)$ and
$\epsilon_q=Q^{\otimes n}(A_n)$.  Its effect is
\begin{equation}
 \boxed{E_n=(1-\epsilon_p)p+\epsilon_q(I-p),\qquad
 \epsilon_p+\epsilon_q
 =\sum_w\min\{P^{\otimes n}(w),Q^{\otimes n}(w)\}
 \le b_{\rm cl}^{\,n}.}
 \label{sel6:eq:likelihood-effect}
\end{equation}
If $P\ne Q$ then $b_{\rm cl}<1$ and this is a nonvanishing-success
approximation to the native sharp effect $p$.  In particular
$\|E_n-p\|_{\op}=\max(\epsilon_p,\epsilon_q)\le b_{\rm cl}^{\,n}$.
If $P=Q$, no finite or limiting processing of these repeated classical
records distinguishes the two native inputs.  Phases in
\eqref{sel6:eq:diagonal-updates} and unobserved within-outcome Kraus
multiplicity do not alter either assertion.
\end{theorem}
\begin{proof}
The pure states $p,I-p$ remain the same pure states after any nonzero
outcome, although off-diagonal input coherences can change.  Repetition
therefore gives the displayed product laws, including when an outcome has
more than one unobserved coefficient.  Every word effect is diagonal in
$p,I-p$ with the corresponding product probabilities as its two entries.
Sum over $A_n$ to obtain \eqref{sel6:eq:likelihood-effect}.  On each record
word the sum of the two decision-error contributions is the smaller of the
two probabilities; ties do not change that sum.  Since
$\min(x,y)\le\sqrt{xy}$, summation factorizes and gives
$b_{\rm cl}^{\,n}$.  Cauchy--Schwarz gives $b_{\rm cl}\le1$, with equality
exactly when the two normalized probability vectors agree.  In that case
all word laws agree and classical postprocessing cannot distinguish them.
\end{proof}

\begin{corollary}[Update-independent sufficient private-write limit]
\label{sel6:cor:general-writes}
At the same completed-private reset cut and with the same physical input
identification $t=|a\rangle\otimes s$ as in V5, suppose $p=|s\rangle
\langle s|$ is the axis in Theorem~\ref{sel6:thm:likelihood}.  Inserting its
actual record event before erasure gives the two selected writes
\[
 \mathcal J_{b,n}(X)=\tfrac14
 \Tr[(|a\rangle\langle a|\otimes E_n)X]\,\omega_{b,h}.
\]
Their flagged distance to V5's two ideal common-input writes is exactly
\begin{equation}
 \boxed{\left\|\bigoplus_{b=0}^1\mathcal J_{b,n}
       -\bigoplus_{b=0}^1\mathcal K_b\right\|_\diamond
       =\tfrac12\max(\epsilon_p,\epsilon_q)
       \le\tfrac12 b_{\rm cl}^{\,n}.}
 \label{sel6:eq:general-write-error}
\end{equation}
Their success on $p_t$ is $(1-\epsilon_p)/2\to1/2$ if $P\ne Q$.
No symmetric two-outcome law, L\"uders state update, resolved Kraus frame,
or inverse pulse is required.  The actual conserved future effect,
nonidentical record laws, repeatable execution and common input map remain
source conditions.
\end{corollary}
\begin{proof}
Only the full effect $E_n$ enters the subsequent atomic reset.  The
measure--prepare difference has Hermitian input functional of norm
$\max(\epsilon_p,\epsilon_q)/4$ and a flagged positive output of trace
two.  V4's elementary diamond-norm identity for such a map gives
\eqref{sel6:eq:general-write-error}.  The success formula follows by
substitution of $p_t$.  The private frame and its colour consequence, when
an actual common-carrier identification is supplied, are inherited from
V3--V5 and are not new generation claims.
\end{proof}

\begin{remark}[How the update premise has been reduced]
One can test $\Lambda^*A=A$ for a nonscalar actual qubit future effect $A$
using the already declared input/output frames.  That conservation identity
forces the diagonal form for \emph{every} coefficient, rather than imposing
the positive square roots used in V5.  Outcome frequencies then determine
$P,Q$.  This replaces a particular microscopic state-update model by a
finite conservation test and measured informativeness.  It does not infer
that either test passes in an unpopulated historical table, nor does it
infer repeated memoryless execution from a single measurement panel.
\end{remark}

\section{An informative record has a sharp quantum-disturbance cost}
\label{sel6:disturbance}

The preceding theorem needs transparency only on the specified binary future
space.  It need not preserve a quantum continuation capable of testing
coherence between those two inputs.  The following bound applies to any
finite instrument on one retained carrier, not only a nondemolition or a
L\"uders model.  It also applies to an adaptive finite tree once all its
outcomes and actual updates are collected into its complete instrument.

For two orthogonal unit inputs $t,u\in S$, let
\[
 p_y=\Tr\mathcal I_y(|t\rangle\langle t|),\qquad
 q_y=\Tr\mathcal I_y(|u\rangle\langle u|),\qquad
 D_{\rm rec}=\tfrac12\sum_y|p_y-q_y|.
\]
This is total variation distance of the normalized classical record laws,
not a trace distance after postselecting away a rare outcome.

\begin{theorem}[Sharp record-information versus unconditioned disturbance]
\label{sel6:thm:information-disturbance}
For every normalized finite instrument with total channel $\Lambda$,
\begin{equation}
 \boxed{\|\Lambda-\id\|_\diamond
 \ge 1-\sqrt{1-D_{\rm rec}^2}.}
 \label{sel6:eq:information-disturbance}
\end{equation}
The constant and function are sharp already on a qubit.  If a retained
binary decision has false-negative and false-positive probabilities
$\epsilon_p,\epsilon_q$ on $t,u$ and
$\epsilon_p+\epsilon_q\le1$, then
\begin{equation}
 \|\Lambda-\id\|_\diamond
 \ge 1-\sqrt{1-(1-\epsilon_p-\epsilon_q)^2}.
 \label{sel6:eq:decision-disturbance}
\end{equation}
In particular, two errors at most $\delta\le1/2$ require disturbance at
least $1-2\sqrt{\delta(1-\delta)}$.
\end{theorem}
\begin{proof}
Take any Kraus decomposition $M_{y\alpha}$ of the actual instrument and set
\[
 z=\langle t,\Lambda(|t\rangle\langle u|)u\rangle
   =\sum_{y,\alpha}\langle t,M_{y\alpha}t\rangle
                 \overline{\langle u,M_{y\alpha}u\rangle}.
\]
Within each outcome, Cauchy--Schwarz and
$|\langle t,Mt\rangle|\le\|Mt\|$ imply
\[
 |z|\le b:=\sum_y\sqrt{p_yq_y}.
\]
A second Cauchy--Schwarz inequality gives
\[
 \sum_y|p_y-q_y|
 =\sum_y|\sqrt{p_y}-\sqrt{q_y}|(\sqrt{p_y}+\sqrt{q_y})
 \le2\sqrt{1-b^2}.
\]
Thus $|z|\le\sqrt{1-D_{\rm rec}^2}$.
To turn this into a valid quantum experiment rather than a formal
non-Hermitian input, use a reference qubit and the state
$|\Omega\rangle=(|0\rangle\otimes t+|1\rangle\otimes u)/\sqrt2$.
The Hermitian contraction
\[
 W=|0,t\rangle\langle1,u|+|1,u\rangle\langle0,t|
\]
has expectation one before $\Lambda$ and $\operatorname{Re}z$ afterwards.
Trace-norm duality therefore bounds the amplified output difference below
by $1-\operatorname{Re}z\ge1-|z|$, proving
\eqref{sel6:eq:information-disturbance}.

For sharpness, on $\C t\oplus\C u$ take
$E_\pm=(I\pm DZ)/2$ with $Z=p_t-p_u$ and use their L\"uders coefficients.
Their record distance is $D$, and their total channel multiplies the
coherences by $\sqrt{1-D^2}$.  It equals
$\tfrac{1+\sqrt{1-D^2}}2\id+
 \tfrac{1-\sqrt{1-D^2}}2\operatorname{Ad}_Z$.
The diamond distance from identity is $1-\sqrt{1-D^2}$: the triangle
inequality gives the upper bound, and the input
$(t+u)/\sqrt2$ attains it.  Finally, a binary decision has probability
difference $1-\epsilon_p-\epsilon_q$ on the two inputs.  Total variation
cannot increase on classical coarse graining, so substituting this lower
bound for $D_{\rm rec}$ proves \eqref{sel6:eq:decision-disturbance}.
\end{proof}

The record/coherence inequality in this proof is established measurement
theory, closely related to the visibility--distinguishability relation
\cite{sel6-Englert}.  Its role here is an unnormalized, actual-source error
budget and an obstruction valid beyond the special symmetric weak update.
It does not assert that an arbitrary private preparation is preserved.

\begin{proposition}[Exact old-future defect of V5's repeated weak Read]
\label{sel6:prop:majority-future-defect}
For V5's $n$ repeated weak Reads on a fixed qubit with strength $a$, let
$\kappa=\sqrt{1-a^2}$ and $\eta_n=\kappa^n$.  Forgetting all the new
outcomes gives the exact channel
\begin{equation}
 \Lambda_n(X)=pXp+qXq+\eta_n(pXq+qXp),\qquad q=I-p.
 \label{sel6:eq:dephasing-channel}
\end{equation}
For any Hermitian old future effect $F$,
\begin{equation}
 \|\Lambda_n^*F-F\|_{\op}
 =(1-\eta_n)\|pFq+qFp\|_{\op}.
 \label{sel6:eq:future-disturbance}
\end{equation}
On a determining basis of old future effects,
\begin{equation}
 \boxed{\mathfrak c_{\rm fut}(\Lambda_n)
 =(1-\eta_n)^2\sum_\ell\|[p,F_\ell]\|_{\HS}^2.}
 \label{sel6:eq:future-commutator}
\end{equation}
The insertion is transparent at every finite $n\ge1$ if and only if all
those old effects commute with $p$.  If $na_n^2\to\infty$, as in the
sufficient weak-read limit of V5, then $\eta_n\to0$; the disturbance of
any fixed noncommuting future effect does not tend to zero.
\end{proposition}
\begin{proof}
One use multiplies each off-diagonal matrix unit by
$2\sqrt{(1+a)(1-a)}/2=\kappa$ and leaves the two diagonal blocks fixed.
Composition gives \eqref{sel6:eq:dephasing-channel}, independently of the
later majority grouping.  This channel is self-adjoint for the HS pairing,
so its action on $F$ gives \eqref{sel6:eq:future-disturbance}.
The HS norms of $pFq+qFp$ and $[p,F]=pFq-qFp$ agree, since their two
rectangular blocks are orthogonal.  Sum the squared identities to obtain
\eqref{sel6:eq:future-commutator}.  Finally
$(1-a_n^2)^{n/2}\le e^{-na_n^2/2}$.
\end{proof}

\begin{remark}[The tradeoff is relative to the actual old futures]
The noncommuting effects in this proposition must occur as old terminal
queries at the tested cut.  A private \emph{state} $\rho_b$ is not thereby
an admitted effect $F=\rho_b$.  Conversely, a complete quantum output at an
exposed port has an actual spanning terminal frame, so some such effect
exists whenever the inserted channel is not identity.  Quantum tomography
performed only \emph{after} an erasing reset cannot reconstruct those
pre-erasure coherences by reversing the reset.
\end{remark}

\section{The actual no-response branch locates the conservation boundary}
\label{sel6:idle}

An informative Read can be hidden from a reset which discards its input.
The locked opportunity is different: one of its old resolved branches
preserves the entire acted-on quantum input.
For a normalized old instrument $\mathcal O=(\mathcal O_0,
\mathcal O_1,\ldots)$ on $S$, suppose an old branch is
\begin{equation}
 \mathcal O_0=r\,\operatorname{Ad}_U,
 \qquad 0<r\le1,
 \qquad U:S\to S\ \text{unitary}.
 \label{sel6:eq:unitary-old-branch}
\end{equation}
The source application has $U=I$.  The comparison below keeps that branch
label and its quantum output, rather than just its total probability.

\begin{theorem}[A retained coherent branch detects the inserted Read]
\label{sel6:thm:idle-witness}
For a normalized inserted instrument with total channel $\Lambda$, forgetting
its new outcomes but keeping all the old output flags gives
\begin{equation}
 \boxed{\left\|\bigoplus_a\mathcal O_a\Lambda-
                     \bigoplus_a\mathcal O_a\right\|_\diamond
 \ge r\|\Lambda-\id\|_\diamond.}
 \label{sel6:eq:idle-distance}
\end{equation}
Thus exact transparency for the complete old branch forces $\Lambda=\id$.
Every outcome map of such an insertion is then $c_y\id$, with
$c_y\ge0$ and $\sum_yc_y=1$; it supplies no input-dependent record.
If the inserted record distinguishes orthogonal inputs with distance
$D_{\rm rec}$, the old flagged output changes by at least
\begin{equation}
 \boxed{r\bigl(1-\sqrt{1-D_{\rm rec}^2}\bigr).}
 \label{sel6:eq:idle-information}
\end{equation}
This lower bound holds for arbitrary finite adaptive implementations, not
only for V5's fixed weak Read.
\end{theorem}
\begin{proof}
Selecting the old flag zero is CP and trace-nonincreasing, hence is a
contraction for the diamond norm.  On that branch the difference is
$r\operatorname{Ad}_U(\Lambda-\id)$, of diamond norm
$r\|\Lambda-\id\|_\diamond$ by unitary invariance.  This proves
\eqref{sel6:eq:idle-distance}.  If it is zero, $\Lambda=\id$.
The positive Choi matrices of its outcome maps sum to the rank-one matrix
$|I\rangle\!\rangle\langle\!\langle I|$; every summand must be a
nonnegative scalar multiple of that matrix.  This gives the scalar branch
maps.  Theorem~\ref{sel6:thm:information-disturbance} now yields
\eqref{sel6:eq:idle-information}.  An adaptive finite implementation is still
one normalized instrument when all its history labels are retained, so the
same argument applies to it.
\end{proof}

\begin{corollary}[Exact evaluated locked-source obstruction]
\label{sel6:cor:locked-idle}
For the supplied locked instrument, set
\begin{equation}
 \vartheta=\frac{576851}{417942208512},\qquad
 r=1-\vartheta=\frac{417941631661}{417942208512}.
 \label{sel6:eq:source-calibration}
\end{equation}
A claimed transparent insertion on its acted-on system, placed before the
old opportunity while preserving its full resolved CP output, is
input-uninformative.  More quantitatively, an inserted donor decision with
both errors at most $\delta\le1/2$ changes that output by at least
\begin{equation}
 \boxed{r\bigl(1-2\sqrt{\delta(1-\delta)}\bigr).}
 \label{sel6:eq:locked-cost}
\end{equation}
For V5's repeated symmetric weak Read on an actually identified subsystem,
the no-response contribution alone is exactly
\begin{equation}
 \boxed{r\left[1-(1-a^2)^{n/2}\right].}
 \label{sel6:eq:locked-majority-cost}
\end{equation}
Thus the limiting private-write error may tend to zero while this old
quantum branch changes by a quantity tending to $r$, which is greater than
$0.9999986$ for the stated source calibration.
\end{corollary}
\begin{proof}
The source gives $L_\varnothing=\sqrt r I$, so
Theorem~\ref{sel6:thm:idle-witness} applies independently of all individual
$K_e$.  The decision bound follows from
\eqref{sel6:eq:decision-disturbance}.  For the weak Read,
Proposition~\ref{sel6:prop:majority-future-defect} gives a dephasing channel
with parameter $\eta_n=(1-a^2)^{n/2}$ and diamond distance $1-\eta_n$
from identity.  Tensoring with an untouched subsystem leaves that norm
unchanged.  Multiply by $r$ to obtain the exact flagged contribution.
The rational value in \eqref{sel6:eq:source-calibration} implies the stated
decimal lower bound.
\end{proof}

\begin{remark}[Which source cut has actually been tested]
The acted-on system in this corollary is the system of
\src{def:locked-opportunity-instrument}, not automatically the completed
writer's private register or its categorical coefficient ray.  The
corollary applies to any proposed common-carrier identification placing the
Read on that system before its identity branch.  If the Read acts instead
on a separately retained erased-input copy while the old system has already
been decoupled, that is a different cut and the identity-branch hypothesis
for the copy has not been supplied.  The source does not identify those two
ports merely because both carry a Hilbert representation.

Also, preserving just the probability of the no-response flag does not
suffice: its effect is $rI$, which is unchanged by every normalized
insertion.  It is the branch's full quantum output, or a determining set of
its actually exposed future queries, that detects the disturbance.
\end{remark}

\begin{proposition}[A finite calibrated tomography witness]
\label{sel6:prop:finite-holdout}
At a fixed old coherent branch suppose actual states $(\rho_i)$ and effects
$(F_j)$ form Hermitian bases on its input and output.  Let $A_i,B_j$ be the
trace-dual bases, so $\Tr(A_i\rho_k)=\delta_{ik}$ and
$\Tr(F_jB_k)=\delta_{jk}$.  For the branch-map difference
$\Delta=\mathcal O_0\Lambda-\mathcal O_0$, put
$d_{ij}=\Tr[F_j\Delta(\rho_i)]$ and
\[
 C_{\rm frame}=
 \left(\sum_i\|A_i\|_{\op}\right)
 \left(\sum_j\|B_j\|_1\right).
\]
Then
\begin{equation}
 \|\Delta\|_\diamond\le C_{\rm frame}\max_{ij}|d_{ij}|.
 \label{sel6:eq:tomographic-budget}
\end{equation}
Consequently the inserted record distance enforces at least one old
probability-row defect of size
$r(1-\sqrt{1-D_{\rm rec}^2})/C_{\rm frame}$.
\end{proposition}
\begin{proof}
Dual reconstruction gives the exact map expansion
$\Delta(X)=\sum_{ij}d_{ij}\Tr(A_iX)B_j$.
The diamond norm of $X\mapsto\Tr(AX)B$ is at most
$\|A\|_{\op}\|B\|_1$: apply trace-norm duality to its amplified
functional and then tensor with $B$.  Triangle inequality proves
\eqref{sel6:eq:tomographic-budget}.  The previous lower bound on the coherent
branch gives the final assertion.  The frame conditioning is retained;
no uniform probability-error constant is inferred merely from spanning.
\end{proof}

\begin{proposition}[Why the completed-private reset remains a permitted hiding cut]
\label{sel6:prop:reset-futures}
For the completed-private window and every old future forced to begin with
its displayed reset, all subsequent output-sensitive effects pull back to
$\Span\{P_0,P_1\}$ on the window input.  A local normalized insertion on
the old private factor is therefore transparent for \emph{every} such
future, not just for the coarse reset count.  This statement does not include
old interventions or terminal Reads allowed before that first reset.
\end{proposition}
\begin{proof}
The actual first reset is
$\Psi(X)=\sum_a\Tr(P_aX)\omega_a$.  For any terminal effect $F$ after
any further admitted CP continuation $\mathcal G$,
\[
 (\mathcal G\Psi)^*(F)
 =\sum_a\Tr[\mathcal G(\omega_a)F]P_a.
\]
The retained old hub flag only replaces the scalar coefficients by separate
known scalar rows and does not change their input supports.  A local
normalized private insertion fixes both $P_a=|a\rangle\langle a|\otimes I$.
Proposition~\ref{sel6:prop:future-test} proves transparency on the complete
specified future family.  A pre-reset query has a different pullback and was
not used in the argument.
\end{proof}

\section{The fully resolved locked instrument has no hidden refinement direction}
\label{sel6:efficient-record}

The identity branch is not the only rigid part of the locked record.  Each
accepted branch also has one Kraus coefficient.  The extremality of a
rank-one Choi ray, already used throughout V3--V5, now applies to the
\emph{entire resolved history instrument} and gives a stronger stopping
statement than analysis of its coarse environment algebra.

\begin{theorem}[Every exact refinement of a single-coefficient history is classical]
\label{sel6:thm:efficient-refinement}
Let an old finite instrument have branches
$\mathcal O_a(X)=L_aXL_a^*$.  For any record-respecting CP refinement
$\mathcal O_{a,y}$ with $\sum_y\mathcal O_{a,y}=\mathcal O_a$, every
nonzero old branch satisfies
\begin{equation}
 \boxed{\mathcal O_{a,y}=p(y\mid a)\mathcal O_a,
 \qquad p(y\mid a)\ge0,\qquad\sum_yp(y\mid a)=1.}
 \label{sel6:eq:efficient-refinement}
\end{equation}
The same statement holds for each nonzero resolved history
$\mathcal O_w(X)=L_wXL_w^*$ of any finite length.  Consequently:
\begin{enumerate}[label=\textup{(\roman*)},leftmargin=2.5em]
\item given the old history label, the new label is independent of the input;
\item retaining both labels gives exactly the same total variation distance
between any two input-dependent record laws as retaining the old label alone;
\item the refined coefficient span on the original acted-on carrier and its
generated algebra are unchanged when all new outcomes are retained;
\item a nonzero pure common-input write obtained by such refinements and
positive classical selection requires an old contributing history with a
rank-one \emph{operator} $L_w=|h\rangle\langle t|$ up to scalar.
\end{enumerate}
Thus refinement cannot turn a Choi-rank-one, full-input-rank history into a
one-dimensional donor operation.
\end{theorem}
\begin{proof}
The positive Choi matrices $J_{a,y}$ sum to
$|L_a\rangle\!\rangle\langle\!\langle L_a|$.  Positivity puts each
$J_{a,y}$ on this same line, so it is a nonnegative scalar multiple;
normalization of their sum gives \eqref{sel6:eq:efficient-refinement}.
Products of single coefficients are single coefficients, including zero
products, proving the all-history assertion.
For inputs $\rho,\sigma$, let $P(w),Q(w)$ be the old record probabilities.
Then
\[
 \frac12\sum_{w,y}|p(y\mid w)P(w)-p(y\mid w)Q(w)|
 =\frac12\sum_w|P(w)-Q(w)|.
\]
Each nonzero old coefficient is still spanned by the nonzero scalar
multiples in its complete refinement.  Finally, a positive sum with Choi
rank one can contain only Choi vectors on the final line.  If its amplitude
is a rank-one operator $|h\rangle\langle t|$, every old contributing
coefficient must be proportional to that operator.  Neither signed
cancellation nor coherent recombination of distinct old records is permitted
in this argument.
\end{proof}

\begin{corollary}[Evaluate the whole locked history grammar]
\label{sel6:cor:locked-efficient}
Every exact record-respecting refinement of the supplied locked opportunity
instrument, at any fixed finite history length, is of the form
\eqref{sel6:eq:efficient-refinement}.  In particular its coarse Choi rank
of twenty-four does not furnish an extra quantum refinement beneath a
\emph{fixed old history}.
For a length-$n$ history containing $k\ge1$ accepted labels, its coefficient
is zero unless all accepted pointer and type values are the same.  If they
are the same $(\sigma,r_0)$, then
\begin{equation}
 \boxed{L_w=r^{(n-k)/2}\vartheta^{k/2}
 (K_{e_k}\cdots K_{e_1})\otimes\Pi_\sigma\otimes P_{r_0},
 \qquad \rank L_w=\rank(K_{e_k}\cdots K_{e_1}).}
 \label{sel6:eq:locked-word-rank}
\end{equation}
The all-no-response word has coefficient $r^{n/2}I$ and full input rank.
On the literal binary pointer/type realization of V4, a pure donor in this
unchanged recorded grammar must therefore already occur as a rank-one edge
product, in a single locked pointer/type block.  If every $K_e$ is invertible,
no positive-length old history or its record-respecting refinement supplies
a pure donor.
\end{corollary}
\begin{proof}
Each of the twenty-five source outcomes has a single displayed coefficient,
so Theorem~\ref{sel6:thm:efficient-refinement} applies.  The no-response
coefficient is scalar identity.  Orthogonality of $\Pi_\sigma$ and $P_{r_0}$
makes a product with inconsistent accepted values zero.  Otherwise the
pointer/type factors remain the same rank-one projections and the edge
factors multiply, giving \eqref{sel6:eq:locked-word-rank}.  If every edge
coefficient is invertible its product has rank $D_E\ge3$, while the empty
accepted word has rank $4D_E$.  Theorem~\ref{sel6:thm:efficient-refinement}
then excludes all the claimed pure refinements and positive selections.
\end{proof}

\begin{example}[An existing six-edge completion has no pure finite history]
\label{sel6:ex:invertible-six}
Use the normalized three-dimensional completion of
Proposition~\ref{sel2:prop:six-transient}, with its same five matrices
$B_1,\ldots,B_5$.  Specialize its free real orthonormal contrast matrix to
\[
 O_{e\mu}=\begin{cases}
 1/\sqrt{\mu(\mu+1)},&1\le e\le\mu,\\
 -\mu/\sqrt{\mu(\mu+1)},&e=\mu+1,\\
 0,&e>\mu+1,
 \end{cases}
 \quad 1\le\mu\le5.
\]
It obeys the required $O^{\mathsf T}O=I$ and
$O^{\mathsf T}\mathbf1=0$.  Direct evaluation of the already defined
$K_e=I/\sqrt{18}+\sqrt{2/3}\sum_\mu O_{e\mu}B_\mu$ gives
\[
 \begin{array}{c|c}
 e&4860\det K_e\\\hline
 1&231\sqrt2+202\sqrt3\\
 2&279\sqrt2-218\sqrt3\\
 3,4&45\sqrt2-8\sqrt3\\
 5&45\sqrt2+32\sqrt3\\
 6&45\sqrt2
 \end{array}
\]
Every displayed value is strictly positive.  For example
$2\cdot279^2>3\cdot218^2$ and $2\cdot45^2>3\cdot8^2$ verify the only
two differences.  Hence all six coefficients are invertible.  Every nonzero
locked history containing an acceptance has input rank three, and the
all-no-response history has input rank twelve, at every finite length.
No exact record refinement supplies a pure donor.
The forward and adjoint-closed edge algebras still have dimensions seven
and nine, as established in V2.  Thus an algebra containing rank-one
matrix units does not imply that an individual recorded history is such a
matrix unit.  This is a specified specialization of the existing normalized
countermodel, not a claimed reconstruction of the historical $K_e$.
\end{example}

\begin{remark}[What this closes, and what it does not decide]
This is a stopping theorem for searching for a \emph{new} donor hidden in
an exact refinement of the already fully resolved locked instrument.  It is
not a proof that all original $K_e$ or all their products are invertible.
Those operators remain unevaluated.  Existing rank-one edge words, should
any occur, are not ruled out.  Nor does a rank-one edge word identify its
physical domain with V4's categorical coefficient contrast.  That separate
map and the private output alignment still have to be supplied.

The completed-private history of V5 is different: it has Choi rank two
because the old private input has been discarded, and its refined erased-input
effect is a genuine additional direction.  Replacing this history by one of
the locked opportunity's Choi-rank-one branches silently changes the cut and
the CP map.  A mathematical algebra of the full twenty-four-dimensional
coarse environment is not evidence for additional refinements preserving
each one-dimensional old outcome support.
\end{remark}

\section{The source calibration imposes an absolute opportunity budget}
\label{sel6:budget}

There is a quantitative constraint even before the unknown edge products
have been examined.  This constraint also applies to more general accepted
instruments than the fully resolved one above.

\begin{theorem}[No-response histories bound all record-respecting donor success]
\label{sel6:thm:rare-budget}
Suppose a fixed normalized source instrument on $S$, $\dim S>1$, has
no-response branch $r\id$, $0<r<1$.  Repeat that source for $n$ opportunities
without adding unrecorded operations, and retain every old outcome.  Allow
arbitrary exact CP refinements inside each old history, followed by classical
selection.  Then:
\begin{enumerate}[label=\textup{(\roman*)},leftmargin=2.5em]
\item the all-no-response history contributes no nonzero pure common-input
subchannel;
\item for two arbitrary input states, the total variation distance of the
record laws of the old flags and their new refinement labels is at most
$1-r^n$;
\item a collection of selected pure writes with the same full unit input
$t$ has total success probability on that input at most $1-r^n$.
\end{enumerate}
In particular, a target success $s\in(0,1)$ requires
\begin{equation}
 \boxed{n\ge\left\lceil\frac{\log(1-s)}{\log r}\right\rceil,}
 \label{sel6:eq:necessary-opportunities}
\end{equation}
and equal-prior discrimination error of any classical donor decision is at
least $r^n/2$.  These are necessary opportunity counts, not assertions of
attainability by the unspecified accepted branches.  The decision uses only
those classical record labels; a new terminal measurement of the surviving
quantum output is a different operation and is not included as a free
refinement preserving that output.
\end{theorem}
\begin{proof}
The all-no-response CP branch is $r^n\id$.  Its Choi matrix has one line,
spanned by the full-rank identity coefficient.  Every exact refinement is a
scalar multiple of it, so none is a rank-one operator write on $S$.
Moreover its old probability is $r^n$ for every input and its new sublabels
have the same input-independent scalar distribution.  The remaining records
carry mass $1-r^n$ for each input; their total variation contribution is
therefore at most $1-r^n$.  For pure-write success, no nonzero portion can
come from the all-no-response branch.  The sum of all other old history
effects equals $(1-r^n)I$ by normalization.  CP domination makes the sum of
selected effects no larger.  A common-input collection has effect
$(\sum_jc_j)p_t$, so evaluation on $t$ gives $\sum_jc_j\le1-r^n$.
Solving this inequality gives \eqref{sel6:eq:necessary-opportunities}.
For equal priors the optimal classical testing error is
$(1-D_{\rm rec})/2$, yielding the last bound.
\end{proof}

\begin{example}[Sharpness of the general opportunity bound]
\label{sel6:ex:rare-sharp}
On a qubit with orthogonal rank-one projections $p,q$, take the normalized
instrument with coefficients $\sqrt r I$, $\sqrt{1-r}\,p$ and
$\sqrt{1-r}\,q$.  A nonzero history with an accepted letter has a
coefficient proportional to $p$ or to $q$; histories containing both
accepted types vanish.  The all-no-response history has mass $r^n$.
On inputs $p,q$, all the other record mass has disjoint support, so the
record distance is exactly $1-r^n$.  Selecting the histories with an
accepted $p$ gives total map $(1-r^n)(X\mapsto pXp)$, a pure write on
that donor with success $1-r^n$.  Thus the universal bounds are sharp.
This example establishes sharpness over the declared general source class;
it does not assign those accepted projectors to the locked historical source.
\end{example}

\begin{corollary}[Evaluate the locked success scale]
\label{sel6:cor:numerical-budget}
For \eqref{sel6:eq:source-calibration}, any exact record-respecting pure
common-input success at one locked opportunity is at most
\[
 \vartheta=\frac{576851}{417942208512}<1.381\times10^{-6}.
\]
A fixed-horizon scheme with success at least $1/2$ requires at least
\begin{equation}
 \boxed{502202\ \text{locked opportunities}.}
 \label{sel6:eq:half-budget}
\end{equation}
This does not contradict V4--V5's two-write success $1/2$ at the separate
completed-private reset cut.  Those writes are not record-respecting
refinements of one locked opportunity with its no-response branch retained.
\end{corollary}
\begin{proof}
Apply Theorem~\ref{sel6:thm:rare-budget} with
$r=417941631661/417942208512$.  The integer in
\eqref{sel6:eq:half-budget} is certified without exponentiating enormous
rational powers.  For $x=\vartheta$, use
\[
 x+x^2/2< -\log(1-x)
 <x+x^2/2+\frac{x^3}{3(1-x)}.
\]
For $L_m=2\sum_{j=0}^{m}3^{-(2j+1)}/(2j+1)$,
\[
 L_m<\log2< L_m+
 \frac{2\,3^{-(2m+3)}}{(2m+3)(1-1/9)}.
\]
At $m=12$, exact rational division of these bounds puts
$\log2/[-\log(1-\vartheta)]$ strictly between $502201$ and $502202$.
The ceiling is therefore \eqref{sel6:eq:half-budget}.
\end{proof}

\begin{proposition}[Expected-cost bound for an actually stopped protocol]
\label{sel6:prop:stopping-budget}
In the same frozen instrument, let $T$ be an adapted nonnegative
integer-valued stopping time with $\mathbb ET<\infty$.  Refinement updates are causal and $T$ is determined before the next
original opportunity.  Do not insert other input-sensitive controls or
replace the original CP coefficients.  If a
selected pure common-input output has probability $s$ on its specified
input, then
\begin{equation}
 \boxed{\mathbb ET\ge\frac{s}{1-r}.}
 \label{sel6:eq:expected-cost}
\end{equation}
For the locked calibration and $s=1/2$ this is more than $362261$ expected
opportunities.  It is an opportunity budget, not a physical-time estimate;
converting it into time requires the actual joint duration record.
\end{proposition}
\begin{proof}
A successful pure write requires at least one accepted opportunity before
stopping, by the same positive support argument as in
Theorem~\ref{sel6:thm:rare-budget}.  Conditional on every previous history,
the next original acceptance probability is $1-r$, since the total accepted
effect is $(1-r)I$; the state need not be stationary.  Let $N_T$ count
acceptances before $T$.  Tonelli and adapted stopping give
\[
 \mathbb EN_T=\sum_{k\ge1}\mathbb E[
       \mathbf1_{\{T\ge k\}}\mathbf1_{\{k\text{ accepted}\}}]
 =(1-r)\sum_{k\ge1}\Pr(T\ge k)=(1-r)\mathbb ET.
\]
Since $\mathbf1_{\{\mathrm{success}\}}\le N_T$, the bound follows.  Randomized
classical stopping and outcome refinements which sum to each original branch
can be retained in the filtration; they do not change the identity accepted
effect of the next unchanged opportunity.  The numerical comparison is direct
rational evaluation of $1/(2\vartheta)$.
\end{proof}

\begin{remark}[Two distinct ways of exceeding these bounds]
A coarse-channel comparison which coherently recombines the no-response
record with accepted records is not a record-respecting refinement; V4 already
exhibits the general marked/coarse distinction.  Alternatively, a new
input-selective intervention can change the old resolved instrument.  Neither
change is prohibited as a new source operation.  Neither can be concealed
inside the unchanged record alphabet or described as having no source cost.
\end{remark}

\section{Outcome of the upstream audit and the next nonduplicative task}
\label{sel6:status}

The most decisive evaluated source result is now stronger than another
conditional readout construction.  Every nonzero fully resolved locked
history has one coefficient, so every exact refinement beneath it is
classical.  There is no additional donor direction waiting to be exposed by
a more elaborate measurement of that history's own minimal environment.
A donor in that grammar must already be a rank-one original edge-product
operation with the required physical domain.  The whole root-uniform phase
obstruction of V4 remains valid; this new full-record statement also covers
nonuniform selections while keeping their unknown edge products explicit.

V5's weak-read route remains valid at a genuinely erasing reset cut.  The
complete-future compiler now tests that cut rather than assuming every earlier
future is erased as well.  On an exposed qubit, one nonscalar conserved future
effect forces any actually informative compatible instrument into the needed
nondemolition form.  Likelihood selection therefore removes the special
symmetric L\"uders update from the sufficient write limit.  Two noncommuting
old effects, by contrast, forbid an informative transparent insertion.

For the locked acted-on system, the explicit no-response branch already
carries the full quantum input.  Informative insertion before it has a
nonzero old-output cost, approaching nearly one at high readout fidelity for
the stated $\vartheta$.  Exact within-record refinements face a different
constraint: all new information is already in old accepted histories, whose
first-occurrence probability imposes the calibrated opportunity bounds.  The
two mechanisms must not be combined as if they described the same experiment.

\begin{center}
\small
\begin{tabular}{L{.37\textwidth}L{.55\textwidth}}
\toprule
Source question & V6 result\\
\midrule
Does V5's reset conservation imply transparency at an arbitrary earlier cut?
& No. The actual future-effect span supplies the exact insertion test; old
quantum coherences can remain visible there.\\[.3em]
Must the sufficient weak Read have the stipulated symmetric L\"uders update?
& No. One conserved native qubit projection plus an actually informative,
repeatable instrument implies a general likelihood-amplification route.\\[.3em]
Can a Read reveal a donor while preserving the full locked no-response map?
& No. Its positive identity branch forces a transparent total channel to be
identity, with input-independent new outcomes.\\[.3em]
Can a full resolved locked history hide a nonclassical exact refinement?
& No. Every positive subbranch is a scalar multiple of that old history map.\\[.3em]
Can rare outcomes be renormalized away in a source-selection success claim?
& No. Success by $n$ locked opportunities is at most $1-r^n$; the source
rate gives the explicit count and stopping lower bounds.\\[.3em]
Are original rank-one edge words ruled out?
& Not in general. Their ranks depend on the actual $K_e$ products; this
source row is still unevaluated.\\[.3em]
Are historical colour and the full SM selected?
& Not established. The same-carrier domain map, actual original or additional
write operations, determinant incidence and matter types remain distinct
source requirements.\\
\bottomrule
\end{tabular}
\end{center}

The next task should not repeat ideal environment filtering, majority
amplification, source covariance or finite commutant generation.  There are
two genuinely different source rows to evaluate.  One is an actual original
accepted word $K_{e_k}\cdots K_{e_1}$ together with its full input/output
identifications: its rank, donor support and private output span can be
tested without adding a readout.  The other is a separately exposed erased
input register with its actual operation table, located after decoupling
from the futures being preserved.  A map supplied only at the locked
coefficient or record level cannot be substituted for either row.

For any candidate additional Read, the finite future residual should now be
run first.  A failed transparency test does not imply that the operation is
unphysical; it means that its execution changes the complete old process and
must be represented as that change.  This is the required information for a
source-selection claim, rather than an assumption that every conservative
extension of a marginal already occurred historically.

\subsection{Proof status and cumulative preservation}
The V5 source prefix before its terminal document command is preserved
byte for byte.  The exact-arithmetic verification checks the future-space
classification examples, conserved non-L\"uders likelihood laws, the
record/coherence identities and sharp qubit witness, V5 majority channel
multipliers, the source rational idle factor, refinement classicality,
locked word-rank examples, finite rare-budget saturation examples, and the
rational logarithm certificate for \eqref{sel6:eq:half-budget}.  The V5
regression suite is rerun.  General all-history, adaptive-instrument and
stopping assertions are proved above rather than inferred from finite tests.
No new historical $K_e$, private Read, common-domain identification,
physical duration, formal proof-assistant certification, or independent
peer review is claimed.

\begin{thebibliography}{99}
\bibitem[13]{sel6-BS}
P.~Busch and J.~Singh, \emph{L\"uders theorem for unsharp quantum measurements},
Phys.\ Lett.\ A \textbf{249} (1998), 10--12;
\href{https://doi.org/10.1016/S0375-9601(98)00704-X}{doi:10.1016/S0375-9601(98)00704-X};
\href{https://arxiv.org/abs/1304.0054}{arXiv:1304.0054}.
Nondisturbance and commutation are established background; the finite
source-cut use and qubit classification are proved explicitly here.
\bibitem[14]{sel6-Englert}
B.-G.~Englert, \emph{Fringe visibility and which-way information: an inequality},
Phys.\ Rev.\ Lett.\ \textbf{77} (1996), 2154--2157;
\href{https://doi.org/10.1103/PhysRevLett.77.2154}{doi:10.1103/PhysRevLett.77.2154}.
The general record/coherence tradeoff is not a new principle.  Its
unnormalized channel bound, source-branch calibration and finite readout
application are proved in this block.
\end{thebibliography}

% END OF SOURCE-SELECTION V6 DEVELOPMENT BLOCK.
% Append subsequent work before the final document terminator.
% Preserve the V1--V6 bodies and state necessary amendments explicitly.


\clearpage
\hypersetup{pdftitle={Historical Standard-Model Source Selection: cumulative V1--V7}}
\begin{center}
{\Large\bfseries V7: Native six-edge histories, coherent phase locking,\\[.2em]
rank-minimal sources and a three-letter occurrence test}\par\medskip
{14 September 2026}
\end{center}
\addcontentsline{toc}{part}{V7: Original six-edge histories and a finite source-record test}

\section{Go upstream of readout access: classify the original coefficients}
\label{sel7:context}

V6 proves that a fully resolved single-coefficient history has no
nonclassical exact refinement.  Its remaining forward-history question is
therefore not answered by proposing a more complete environment Read.  One
must inspect the original coefficients and their actual products.  This
continuation does so on an entire, independently specified covariant source
class, rather than assigning a desired colour bridge.

The source supplies six linearly independent edge operators on a common
acted-on space $S$ and the two identities
\begin{equation}
 \sum_{e=1}^6 K_e^*K_e=I_S,
 \qquad \sum_{e=1}^6K_e=\sqrt2\,I_S.
 \label{sel7:eq:source-identities}
\end{equation}
These are exactly the edge identities of
\cite{Grand}, \src{def:locked-opportunity-instrument}, retained in
V4--V6.  The physical opportunity coefficients additionally carry the original
pointer/type projections and the factor $\sqrt\vartheta$; the no-response
coefficient is $\sqrt{1-\vartheta}I$.  That calibration will not be divided
out of a success or error budget.

\paragraph{What is inherited.}
The source's coherent normalization, the six-to-five contrast normal form,
V6's efficient-refinement theorem, the earlier root-dyadic independence,
V1's distinction between an external action and its internal commutant, and
V3's common-input write theorem are inputs.  None is claimed as a new general
result.  The relation between tight frames and rank-one quantum measurements
is established background \cite{sel7-EldarForney}; the new statements here
concern a fixed coherent identity sum, its support-rank defect, the complete
covariant six-coefficient family, and products in its original record grammar.
No priority claim over frame or representation theory is made.

\paragraph{The additional source class is explicit.}
The detailed classification assumes that $\dim S=3$ and that the actual
resolved edge CP maps are covariant under a verified unitary action of $S_4$
which permutes their six unoriented-edge records.  Three is the \emph{smallest
possible} acted-on dimension compatible with six independent matrices; the
historical source is not declared to attain it.  Likewise an edge label
$e\in E(K_4)$ does not by itself verify covariance on the acted-on system.
These two qualifications distinguish a complete classification of a source
class from identification of that class with a particular historical writer.

\begin{center}
\small
\begin{tabular}{L{.36\textwidth}L{.56\textwidth}}
\toprule
Input or result & Logical status in this block\\
\midrule
Six independent coefficients and \eqref{sel7:eq:source-identities}
& Supplied locked-source normalization; no individual $K_e$ is numerically assigned.\\[.3em]
Three-dimensional acted-on carrier
& The minimal possible dimension, not a consequence that the historical carrier is minimal.\\[.3em]
Resolved system CP covariance
& An explicit test on the actual system representation and branch maps, not a consequence of record relabelling.\\[.3em]
A coherent covariant Kraus phase convention
& Derived from the identity sum and independence, rather than imposed separately.\\[.3em]
Original rank-one history
& Classified on the whole stated family; existence is decided by two- and three-letter record masses.\\[.3em]
Internal colour, private state alignment, determinant and matter occurrence
& Not inferred from an external three-dimensional coefficient algebra.\\
\bottomrule
\end{tabular}
\end{center}

All products below are consecutive applications of the same original edge
instrument on its complete retained carrier.  If a controller restricts which
edge can follow which other edge, that controller belongs in the carrier and
the theorem must be applied to the corresponding complete typed maps.  An
unadmitted record word is not a vanishing matrix product on a smaller carrier.

\section{The coherent identity sum gives an exact operator-rank defect}
\label{sel7:rank}

\begin{theorem}[Coherent normalization and the total support-rank floor]
\label{sel7:thm:rank-defect}
Let $K_1,\ldots,K_m\in\B(\C^D)$ satisfy
$\sum_jK_j^*K_j=I$ and $\sum_jK_j=\lambda I$, with $\lambda>0$.
Let $P_j$ be the orthogonal projection onto $\Ran K_j$, and put
$R=\sum_j\rank K_j$.  Then
\begin{equation}
 \boxed{\sum_{j=1}^m
   \norm{K_j-\lambda^{-1}P_j}_{\HS}^{2}
   =\frac{R}{\lambda^2}-D.}
 \label{sel7:eq:rank-defect}
\end{equation}
In particular $R\ge\lambda^2D$.  Equality holds precisely when
\begin{equation}
 K_j=\lambda^{-1}P_j\quad\hbox{for every }j,
 \qquad \sum_jP_j=\lambda^2 I.
 \label{sel7:eq:rank-saturation}
\end{equation}
The same defect identity holds with the initial-support projections
$\supp(K_j^*K_j)$ in place of $P_j$.
\end{theorem}
\begin{proof}
Since $P_jK_j=K_j$,
$\Tr(K_j^*P_j)=\overline{\Tr K_j}$.  Expanding the squared norms gives
\[
 \sum_j\Tr K_j^*K_j+
 \lambda^{-2}\sum_j\Tr P_j-
 2\lambda^{-1}\operatorname{Re}\Tr\sum_jK_j
 =D+\lambda^{-2}R-2D.
\]
This is \eqref{sel7:eq:rank-defect}.  Equality in its nonnegative left side
means each difference vanishes, and then the coherent sum gives
\eqref{sel7:eq:rank-saturation}.  Conversely those identities give equality.
For the initial projection $Q_j$, use $K_jQ_j=K_j$ and cyclicity of trace
in the same expansion.
\end{proof}

\begin{corollary}[The minimum six-edge source is a projector instrument]
\label{sel7:cor:minimal-six}
For six independent coefficients satisfying
\eqref{sel7:eq:source-identities}, one has $D\ge3$ and $R\ge2D$.
If every coefficient has operator rank one, then necessarily $D=3$ and
\begin{equation}
 \boxed{K_e=\frac1{\sqrt2}P_e,
 \qquad \rank P_e=1,
 \qquad \sum_{e=1}^6P_e=2I_3.}
 \label{sel7:eq:projector-instrument}
\end{equation}
Conversely, any six independent rank-one projections satisfying the last
identity give a normalized source of this kind.  The bound $R=6$ at $D=3$
is attained.  Thus no separate L\"uders update or outcome phase is needed
\emph{on the saturating branch}; the coherent normalization forces it.
\end{corollary}
\begin{proof}
Independence gives $6\le D^2$, so $D\ge3$.
The theorem gives $R\ge2D$.  Six rank-one coefficients have $R=6$, which
forces $D\le3$ and hence $D=3$.  Equality in the rank floor gives
\eqref{sel7:eq:projector-instrument}.  The converse is direct.  The six
normalized $A_3$ root projections below have sum $2I_3$ and are independent,
providing an attaining example.
\end{proof}

\begin{remark}[An exact characterization, not a free minimality axiom]
The rank $R$ counts physical input/output support dimensions of the original
operators, not their number of coefficient directions or the rank of a source
Hilbert Gram.  A positive source Gram does not force $R=6$.
No minimization of $R$ is added to the historical law.  The corollary explains
what would follow if this measurable rank saturation actually held and what
its absence means.
\end{remark}

\section{System CP covariance fixes the phases and the entire minimal family}
\label{sel7:classification}

Write $\mathscr E=\{\{i,j\}:1\le i<j\le4\}$ and let $\bar e$ denote
the disjoint complementary edge.  A unitary representation $U_g$ on $S$
implements resolved covariance when
\begin{equation}
 K_{ge}XK_{ge}^*
 =U_gK_eU_g^*XU_gK_e^*U_g^*
 \quad\text{for all }X,e,g.
 \label{sel7:eq:CP-covariance}
\end{equation}
This statement concerns the actual branch maps, not only their traces.

\begin{lemma}[An independent identity expansion locks the Kraus phases]
\label{sel7:lem:phase-lock}
For a linearly independent nonzero list with
$\sum_eK_e=\lambda I$, $\lambda\ne0$,
\eqref{sel7:eq:CP-covariance} implies
\begin{equation}
 \boxed{K_{ge}=U_gK_eU_g^*.}
 \label{sel7:eq:amplitude-covariance}
\end{equation}
No further phase choice is required in the already calibrated list.
\end{lemma}
\begin{proof}
Equality of two nonzero single-coefficient CP maps makes their coefficients
unit scalar multiples, by equality of their rank-one Choi rays.  For fixed
$g$ write $U_gK_eU_g^*=z_{g,e}K_{ge}$, $|z_{g,e}|=1$.
Summing gives $\lambda I=\sum_ez_{g,e}K_{ge}=\sum_eK_{ge}$.
Independence forces every $z_{g,e}=1$.
\end{proof}

\begin{lemma}[The minimal covariant acted-on representation is a triplet]
\label{sel7:lem:triplet}
Suppose $\dim S=3$ and six independent coefficients obey
\eqref{sel7:eq:CP-covariance} and the coherent identity sum.  Then $U$
is the standard three-dimensional representation of $S_4$ or its sign twist,
up to unitary equivalence.  Their conjugation actions are identical.
\end{lemma}
\begin{proof}
The five complex irreducible dimensions of $S_4$ are $1,1,2,3,3$; the two
triplets are the standard representation and its sign twist.  This finite
representation list is already used in the source decomposition; see
\cite{Etingof}.  Both one-dimensional representations and the two-dimensional
representation are trivial on the normal Klein four-group of double
transpositions.  Hence any reducible three-dimensional representation is
trivial there.  But, for example, $(12)(34)$ exchanges the edge records
$\{1,3\}$ and $\{2,4\}$.  Equation~\eqref{sel7:eq:amplitude-covariance}
would identify those two independent coefficients, a contradiction.
Thus the representation is irreducible, yielding the stated list.
Multiplication of $U_g$ by the sign scalar does not change conjugation.
\end{proof}

On
\[
 W=\{x\in\C^4:\sum_ix_i=0\}
\]
define $u_e=(e_i-e_j)/\sqrt2$, with an arbitrary fixed sign for each
unoriented edge, and let
\begin{equation}
 P_e=|u_e\rangle\langle u_e|,
 \qquad Q_e=I_W-P_e-P_{\bar e}.
 \label{sel7:eq:root-projectors}
\end{equation}
The three terms are mutually orthogonal rank-one projections.
$Q_e=Q_{\bar e}$ is the pair-sum contrast; the three different $Q_e$ form
an orthogonal resolution of $I_W$.  The root frame satisfies
\begin{equation}
 \sum_eP_e=2I_W,
 \qquad\sum_eQ_e=2I_W.
 \label{sel7:eq:root-frame}
\end{equation}
These follow directly from the complete-graph Laplacian, or from the three
pair-contrast coordinates.  The six $P_e$ are independent: a vanishing
combination is a graph Laplacian with zero restriction to $W$ and zero action
on the uniform line, so its off-diagonal entries force every edge coefficient
to vanish.  This is the source's existing root-dyadic independence argument.

\begin{theorem}[Complete minimal-carrier coefficient family]
\label{sel7:thm:normal-form}
Under the hypotheses of Lemma~\ref{sel7:lem:triplet}, all six coefficients
have, in one source-compatible unitary frame, the form
\begin{equation}
 \boxed{K_e=\frac1{\sqrt2}
   (\alpha P_e+\beta P_{\bar e}+\gamma Q_e),}
 \label{sel7:eq:normal-form}
\end{equation}
where
\begin{equation}
 \boxed{\alpha+\beta+\gamma=1,
 \qquad |\alpha|^2+|\beta|^2+|\gamma|^2=1.}
 \label{sel7:eq:parameter-sphere}
\end{equation}
The six operators are independent exactly when
\begin{equation}
 x:=\alpha-\beta\ne0,
 \qquad y:=\alpha+\beta-2\gamma\ne0.
 \label{sel7:eq:independence}
\end{equation}
Conversely every parameter satisfying
\eqref{sel7:eq:parameter-sphere}--\eqref{sel7:eq:independence} gives an
independent normalized covariant six-edge source.  All its coefficients are
normal, the coarse channel is unital, and its forward coefficient algebra is
\begin{equation}
 \boxed{\operatorname{alg}(I,K_e:e)=C^*(K_e:e)=M_3(\C).}
 \label{sel7:eq:full-algebra}
\end{equation}
In particular the stationary state $I_3/3$ is faithful throughout this family.
\end{theorem}
\begin{proof}
The stabilizer of $e=\{1,2\}$ is generated by $(12)$ and $(34)$.
On $W$ its three distinct character lines are the ranges of
$P_{12},P_{34},Q_{12}$, with respective characters $(-,+),(+,-),(+,+)$.
By Lemma~\ref{sel7:lem:phase-lock}, $K_{12}$ commutes with this stabilizer
and is diagonal on those three lines.  Transport gives
\eqref{sel7:eq:normal-form} with the same three scalars for all edges.
The two source identities and \eqref{sel7:eq:root-frame} give
\eqref{sel7:eq:parameter-sphere}; conversely those identities normalize the
whole family.

Since $I=\tfrac12\sum_fP_f$, the coefficient transformation relative to the
independent $P_f$ has matrix
\[
 T=\frac\gamma2\mathbf1\mathbf1^{\mathsf T}
       +(\alpha-\gamma)I_6+(\beta-\gamma)C,
\]
where $C$ exchanges complementary edges.  On the uniform line its eigenvalue
is $1$; on the complementary-edge-even zero-sum plane it is $y$; on the
complementary-edge-odd three-space it is $x$.
This proves \eqref{sel7:eq:independence}.  If the transformation is invertible,
the forward algebra contains every $P_e$ as a linear combination.  Nonzero
overlaps connect the six root rays, and their products give all dyads between
a spanning set of $W$.  Thus their algebra is $M_3$.
This is an algebraic span statement, not an assertion that those linear
combinations are newly executable histories.  Normality is manifest in
\eqref{sel7:eq:normal-form}; consequently the trace-preserving identity also
gives $\sum_eK_eK_e^*=I$, proving unitality and the stationary claim.
\end{proof}

\begin{proposition}[The one-letter coefficient Gram]
\label{sel7:prop:Gram}
The labelled Hilbert--Schmidt Gram $G_{ef}=\Tr(K_e^*K_f)$ is
\begin{equation}
 \boxed{G=\Pi_1+\frac{|y|^2}{4}\Pi_2+
                     \frac{|x|^2}{2}\Pi_3,}
 \label{sel7:eq:Gram}
\end{equation}
where $\Pi_1,\Pi_2,\Pi_3$ are the uniform, complement-even zero-sum, and
complement-odd projections, respectively.  Their ranks are $1,2,3$.
In particular all $G_{ee}=1/2$, and
$|y|^2+3|x|^2=4$.  The raw operator determinant is instead
\begin{equation}
 \boxed{\det K_e=\frac{\alpha\beta\gamma}{2\sqrt2}.}
 \label{sel7:eq:det-coefficient}
\end{equation}
\end{proposition}
\begin{proof}
The root-projector Gram has entries $1$ on the diagonal, zero on complementary
edges, and $1/4$ on the four adjacent edges.  Its eigenvalues on the three
stated subspaces are $2,1/2,1$.  Conjugate that Gram by $T/\sqrt2$ from the
previous proof.  This gives \eqref{sel7:eq:Gram}.  Taking its trace gives
$1+|y|^2/2+3|x|^2/2=3$, hence the constraint.  The three diagonal
eigenvalues of one raw coefficient are
$\alpha/\sqrt2,\beta/\sqrt2,\gamma/\sqrt2$, proving the determinant formula.
\end{proof}

\section{All finite rank-one histories on this source class have a depth-two test}
\label{sel7:histories}

No projection, adjoint, or coherent combination of distinct records is inserted
into the histories in this section.  A word is precisely
$K_w=K_{e_k}\cdots K_{e_1}$.

\begin{theorem}[Complete native pure-history alternative]
\label{sel7:thm:pure-word}
On the family of Theorem~\ref{sel7:thm:normal-form}, the following are
equivalent:
\begin{enumerate}[label=\textup{(\roman*)},leftmargin=2.5em]
\item a nonzero original forward word has operator rank one;
\item such a word exists with length at most two;
\item $\alpha\beta\gamma=0$;
\item one original edge effect $K_e^*K_e$ is singular.
\end{enumerate}
More precisely, exactly one of the following applies:
\begin{center}
\begin{tabular}{L{.36\textwidth}L{.54\textwidth}}
\toprule
Parameters & Exact native history result\\
\midrule
$\alpha\beta\gamma\ne0$ & Every word is invertible; no finite pure donor exists.\\[.25em]
$(\alpha,\beta,\gamma)=(1,0,0)$ or $(0,1,0)$
& Every one-letter coefficient is rank one.\\[.25em]
$\alpha=0$ or $\beta=0$, with exactly two nonzero parameters
& $K_{\bar e}K_e=(\gamma^2/2)Q_e\ne0$ is rank one.\\[.25em]
$\gamma=0$, $\alpha\beta\ne0$
& Any pair of coefficients from different complementary-edge pairs has a
rank-one product.\\
\bottomrule
\end{tabular}
\end{center}
The tuple $(0,0,1)$ is excluded by coefficient independence.  When the
one-letter rank is two, the minimal pure-word length is exactly two.
\end{theorem}
\begin{proof}
If all three parameters are nonzero, every factor is invertible by
\eqref{sel7:eq:det-coefficient}, and so is every product.  This proves the
negative branch without a bound on word length.
If two parameters vanish, normalization sets the remaining one to $1$;
independence excludes the $Q_e$-only family, leaving the two rank-one cases.
For a complementary pair, direct multiplication in their common three-line
basis gives
\begin{equation}
 K_{\bar e}K_e
 =\frac12\bigl[\alpha\beta(P_e+P_{\bar e})+\gamma^2Q_e\bigr].
 \label{sel7:eq:opposite-product}
\end{equation}
Thus either $\alpha=0$ or $\beta=0$, with rank two, gives the stated
nonzero rank-one product.

If $\gamma=0$ and $\alpha\beta\ne0$, the range and support of $K_e$ are
$Q_e^\perp$.  For a different complementary pair $f$, $Q_e,Q_f$ are
orthogonal rank-one projections.  Thus $\Ker K_f=\Ran Q_f$ lies in
$\Ran K_e=Q_e^\perp$.  The restriction of $K_f$ to that two-dimensional
range loses exactly one dimension, so $\rank(K_fK_e)=1$.
All singular one-letter cases have now been exhausted.  Since their one-letter
rank is two unless already listed as rank one, the asserted minimal lengths
follow.
\end{proof}

\begin{proposition}[Pure finite words occupy explicit exceptional circles]
\label{sel7:prop:generic}
Before the independence exclusions, \eqref{sel7:eq:parameter-sphere} is a
real three-sphere in the complex affine plane
$\alpha+\beta+\gamma=1$.  The pure-word locus is the union of its three
circles $\alpha=0$, $\beta=0$, $\gamma=0$, with the inadmissible points
removed.  Thus the no-pure-finite-history branch is open, dense and of full
three-dimensional surface measure in the admissible parameter set, even
though \eqref{sel7:eq:full-algebra} and faithful stationarity hold throughout.
\end{proposition}
\begin{proof}
Translate by $(1/3,1/3,1/3)$.  The constraint becomes a sphere of radius
$\sqrt{2/3}$ in the two-complex-dimensional zero-sum plane.
For example, $\alpha=0$ reduces the other constraints to
$\gamma=1-\beta$ and $|\beta-1/2|=1/2$, a circle.  The other two cases
are identical.  These closed one-dimensional subsets have empty interior
and zero surface measure in the three-sphere.  Restricting to the open
independence set preserves the conclusion.  Use
Theorem~\ref{sel7:thm:pure-word} for its history interpretation.
\end{proof}

\begin{remark}[What this genericity statement does not forbid]
It excludes \emph{exact finite rank-one coefficient histories} on the stated
open branch.  It neither gives a uniform lower bound on the singular values
of arbitrarily long products nor rules out approximate conditional
purification.  It also does not forbid a nonzero algebraic corner relative
to already supplied physical supports.  Such a corner and a full-input pure
write are different operational claims.
\end{remark}

\begin{example}[Rank two letters with an exactly pure two-letter output]
\label{sel7:ex:rank-two}
Take $\gamma=0$, $\alpha=(1+z)/2$, $\beta=(1-z)/2$,
$|z|=1$, $z\ne\pm1$.  These are all the rank-two parameters on this
circle.  In orthonormal pair-contrast coordinates, the two coefficients
acting on the $(1,2)$-plane have the form
\[
 K_{12}^{\epsilon}=\frac1{2\sqrt2}
 \begin{pmatrix}1&\epsilon z&0\\\epsilon z&1&0\\0&0&0\end{pmatrix},
 \qquad \epsilon=\pm1,
\]
and similarly $K_{23}^{\eta}$ acts on the $(2,3)$-plane.
These plane subscripts describe the pair-contrast coordinates, not a new edge
alphabet.  Direct multiplication gives
\begin{equation}
 K_{12}^{\epsilon}K_{23}^{\eta}
 =\frac14|h_\epsilon\rangle\langle t_\eta|,
 \quad h_\epsilon=\frac{\epsilon z e_1+e_2}{\sqrt2},
 \quad t_\eta=\frac{e_2+\eta\overline z e_3}{\sqrt2}.
 \label{sel7:eq:rank-two-example}
\end{equation}
The two last-edge choices have a common physical input and orthogonal output
rays; each has success $1/16$ on that input in the normalized edge process.
The fixed two-opportunity locked realization carries the additional success
factor $\vartheta^2$.  No terminal filter or environment refinement is used.
The displayed input is not declared to be the categorical source contrast,
and these outputs are not declared to be the historical private pair.
\end{example}

\section{A three-letter original record decides the singular-history branch}
\label{sel7:record-test}

The matrices need not be reconstructed individually to run the preceding
alternative once the source class has been independently verified.
For $k\ge1$, let $s_k$ be the probability, in $k$ consecutive uses of the
\emph{normalized edge instrument}, that all $k$ edge labels are identical.
The event is a positive classical selection of old records:
\[
 s_k(\rho)=\sum_{e\in\mathscr E}
   \Tr(K_e^k\rho(K_e^*)^k).
\]
It has no coherent recombination.  The notation does not assert independent
edge probabilities at successive uses.

\begin{theorem}[Input-independent repeated-edge masses and the exact cubic test]
\label{sel7:thm:record-test}
For every input density operator on the classified source,
\begin{equation}
 \boxed{s_k=2^{1-k}
 (|\alpha|^{2k}+|\beta|^{2k}+|\gamma|^{2k}).}
 \label{sel7:eq:repeated-mass}
\end{equation}
In particular, define
\begin{equation}
 d_{\rm word}:=1-6s_2+8s_3.
 \label{sel7:eq:record-defect}
\end{equation}
Then
\begin{equation}
 \boxed{d_{\rm word}=6|\alpha\beta\gamma|^2\ge0,
 \qquad \det(K_e^*K_e)=\frac{d_{\rm word}}{48}.}
 \label{sel7:eq:record-determinant}
\end{equation}
Thus, within the stated source class,
\begin{equation}
 \boxed{d_{\rm word}=0
 \quad\Longleftrightarrow\quad
 \text{a nonzero original pure history of length at most two exists}.}
 \label{sel7:eq:record-decision}
\end{equation}
If $d_{\rm word}>0$, every finite original word is invertible.
Moreover $1/6\le s_2\le1/2$, and $s_2=1/2$ is equivalent to the
rank-minimal projector branch of Corollary~\ref{sel7:cor:minimal-six}.
\end{theorem}
\begin{proof}
Normality and the three orthogonal spectral lines give
\[
 (K_e^*)^kK_e^k=2^{-k}
 (|\alpha|^{2k}P_e+|\beta|^{2k}P_{\bar e}
                         +|\gamma|^{2k}Q_e).
\]
Summation using \eqref{sel7:eq:root-frame} gives
$\sum_e(K_e^*)^kK_e^k=s_kI$, proving both the formula and input independence.
Put $x_1=|\alpha|^2,x_2=|\beta|^2,x_3=|\gamma|^2$; these are nonnegative
with sum one.  The elementary identity
\[
 1-3\sum_ix_i^2+2\sum_ix_i^3=6x_1x_2x_3
\]
and $s_2=\tfrac12\sum_ix_i^2$, $s_3=\tfrac14\sum_ix_i^3$ prove
\eqref{sel7:eq:record-determinant}.  Apply
Theorem~\ref{sel7:thm:pure-word} for the decision statement.
Finally $1/3\le\sum_ix_i^2\le1$, with upper equality precisely when
one $x_i$ is one and the others zero.  Normalization then gives one of the
three single-parameter points; independence excludes $(0,0,1)$.
\end{proof}

\begin{corollary}[The independent uniform-record endpoint is a no-donor branch]
\label{sel7:cor:uniform-record}
Within the same verified source class, the condition $s_2=1/6$ is equivalent
to $|\alpha|^2=|\beta|^2=|\gamma|^2=1/3$ and therefore to
\[
 K_e^*K_e=I/6\quad\text{for every }e.
\]
Every $\sqrt6K_e$ is then unitary.  Every length-$k$ old edge word has
effect $6^{-k}I$, the record is independent uniform on all inputs, and no
finite word is a donor.  In particular $s_3=1/36$ and
$d_{\rm word}=2/9>0$.  It suffices to verify the two-letter repeated-edge
mass, not to assume an entire infinite iid history law.
\end{corollary}
\begin{proof}
Equality in the lower bound $\sum_i x_i^2\ge1/3$ with $\sum_i x_i=1$
forces $x_1=x_2=x_3=1/3$.  The coefficient spectral decomposition gives
the stated effects.  A square matrix with $K_e^*K_e=I/6$ is a scaled
unitary, and products preserve that property.  The word probabilities and
remaining constants follow directly.
\end{proof}

\begin{remark}[Chronological labels are not automatically the same local writer]
The independent chronological edge records in the source must not simply be
inserted into this corollary as the local system history law.
The source explicitly distinguishes the transitive ordered-pair action, with
its orientation cocycle, from a locked local-edge representation with separate
sign orbits; see \cite{Attack},
\src{fire:chronological-local-distinction}.  The same-history identification
and mixed operational comparison remain required.  Once that identification
and the source-class hypotheses are verified, however, a genuinely measured
uniform two-letter endpoint would rule out native finite donors rather than
supply them.
\end{remark}

\begin{proposition}[An opposite-edge event detects the minimum-rank branch]
\label{sel7:prop:opposite-test}
Let $o_2$ be the probability that two consecutive edge records are complementary,
with both orders and all six first labels included.  Its effect is scalar:
\begin{equation}
 \boxed{\sum_e(K_{\bar e}K_e)^*(K_{\bar e}K_e)
       =o_2I,
 \qquad o_2=|\alpha\beta|^2+\tfrac12|\gamma|^4.}
 \label{sel7:eq:opposite-test}
\end{equation}
It follows that $o_2=0$ exactly on the two independent rank-one projector
orbits $(1,0,0)$ and $(0,1,0)$.  This is equivalent to $R=6$ and to $s_2=1/2$.
\end{proposition}
\begin{proof}
Square \eqref{sel7:eq:opposite-product} on its mutually orthogonal supports,
and sum the result over all edges.  The two root terms sum to $4I$ and the
$Q_e$ terms to $2I$, giving \eqref{sel7:eq:opposite-test}.  Positivity makes
$o_2=0$ equivalent to $\alpha\beta=0$ and $\gamma=0$; the remaining
normalization sets the nonzero parameter to one.  The prior rank theorem
and repeated-edge formula identify the other conditions.
\end{proof}

\begin{corollary}[The original locked calibration is retained]
\label{sel7:cor:locked-test}
Let $p_k$ be the probability of $k$ consecutive \emph{accepted} locked
opportunities with identical edge label.  Sum over the original pointer/type
labels, retaining their consistent values.  For any input on the full
locked carrier,
\begin{equation}
 p_k=\vartheta^k s_k,
 \qquad
 p_{\rm opp}=\vartheta^2o_2.
 \label{sel7:eq:locked-probabilities}
\end{equation}
Consequently the source-normalized test is
\begin{equation}
 \boxed{d_{\rm word}
 =1-6\frac{p_2}{\vartheta^2}
       +8\frac{p_3}{\vartheta^3}.}
 \label{sel7:eq:locked-cubic}
\end{equation}
This is a classification statistic from old recorded cylinders, not an
unnormalized probability of successful writes.  Their actual masses retain
all factors of $\vartheta$.
\end{corollary}
\begin{proof}
For a consistent pointer/type block, the corresponding word has coefficient
$\vartheta^{k/2}K_e^k$ tensored with its original projections.
Inconsistent block sequences vanish.  The sum of the retained projections is
identity, and the edge event effect is the scalar in
\eqref{sel7:eq:repeated-mass}.  This gives the assertion even for entangled
edge--pointer/type inputs.  The complementary-pair calculation is the same.
\end{proof}

\begin{proposition}[A certified negative decision with finite probability errors]
\label{sel7:prop:record-errors}
Suppose $\vartheta$ is the exact declared calibration and measurements satisfy
$|\widehat p_k-p_k|\le\epsilon_k$ for $k=2,3$.  Set
\[
 \widehat d=1-6\widehat p_2/\vartheta^2+8\widehat p_3/\vartheta^3,
 \qquad e_d=6\epsilon_2/\vartheta^2+8\epsilon_3/\vartheta^3.
\]
On a verified source class, $\widehat d-e_d>0$ certifies that no finite
original word is pure.  If $\widehat d+e_d<0$, the data and error budget
are inconsistent with the asserted class.  An interval containing zero does
not certify exact purity, and no statistical observation of zero counts is
identified with an exact zero probability.
\end{proposition}
\begin{proof}
The triangle inequality gives $|\widehat d-d_{\rm word}|\le e_d$.
Use nonnegativity and Theorem~\ref{sel7:thm:record-test}.
\end{proof}

\begin{remark}[Rare-event conditioning is not free]
The source has $\vartheta=576851/417942208512$.  The large factors
$\vartheta^{-2}$ and $\vartheta^{-3}$ in the error propagation are retained
explicitly.  Exact symbolic cylinder values would decide the alternatives;
finite experimental data generally certify the strict negative branch more
readily than an exact rank drop.  V6's opportunity and stopping-time bounds
remain applicable and have not been relaxed by writing a normalized statistic.
An unavailable word in a differently gated protocol cannot be entered as
$p_{\rm opp}=0$ without verifying that it is the same complete carrier and
same repeated instrument used here.
\end{remark}

\section{Identical one-letter Grams can conceal opposite native-history answers}
\label{sel7:Gram-counterexample}

\begin{theorem}[A complete covariant Gram-preserving separation]
\label{sel7:thm:Gram-counterexample}
On the same three-dimensional $W$, with the same six labels and the same
external $S_4$ action, compare
\begin{equation}
 K_e^{\rm p}=\frac1{\sqrt2}P_e,
 \qquad
 K_e^{\rm i}=\frac1{3\sqrt2}(2I-3P_{\bar e}).
 \label{sel7:eq:counterexample-pair}
\end{equation}
Both satisfy all source identities, resolved covariance and coefficient
independence.  Both have forward and adjoint-closed algebra $M_3$ and faithful
stationary state $I_3/3$.  Their \emph{entire labelled one-letter}
Hilbert--Schmidt Grams are equal:
\begin{equation}
 G_{ef}=\begin{cases}
 1/2,&e=f,\\
 0,&f=\bar e,\\
 1/8,&e,f\text{ adjacent and distinct}.
 \end{cases}
 \label{sel7:eq:equal-Gram}
\end{equation}
Their nonzero one-step coarse Choi spectra are therefore also equal.
Nevertheless every $K_e^{\rm p}$ is rank one, whereas every original word in
$K_e^{\rm i}$ is invertible.  The recorded two- and three-letter statistics
separate the two sources:
\begin{center}
\begin{tabular}{lcccc}
\toprule
Source & $s_2$ & $s_3$ & $o_2$ & $d_{\rm word}$\\
\midrule
$K^{\rm p}$ & $1/2$ & $1/4$ & $0$ & $0$\\
$K^{\rm i}$ & $11/54$ & $43/972$ & $4/27$ & $32/243$\\
\bottomrule
\end{tabular}
\end{center}
Thus the result does not claim equality of the two full operational laws.
\end{theorem}
\begin{proof}
The parameter points are $(1,0,0)$ and $(2/3,-1/3,2/3)$.
Both have $x=1$; their respective $y$ values are $1$ and $-1$.
Theorems~\ref{sel7:thm:normal-form} and~\ref{sel7:thm:pure-word} prove
the source and history assertions.  Proposition~\ref{sel7:prop:Gram} shows
equality of the labelled Grams, not just their spectra; the displayed entries
are the root-projector overlaps divided by two.  If $W_K$ synthesizes the
vectorized coefficients, the positive matrices $W_K^*W_K$ and $W_KW_K^*$
have the same nonzero spectrum, giving the Choi statement.
At the second point the squared eigenvalue parameters are
$(4/9,1/9,4/9)$.  Substitution into
\eqref{sel7:eq:repeated-mass} and~\eqref{sel7:eq:opposite-test} gives the
table, including $d_{\rm word}=6(16/729)=32/243$.
\end{proof}

\begin{proposition}[An arbitrarily small same-Gram perturbation removes every exact pure word]
\label{sel7:prop:deformation}
For $z=e^{i\phi}$ set
\begin{equation}
 \alpha=\frac{5+z}{6},\qquad
 \beta=\frac{z-1}{6},\qquad
 \gamma=\frac{1-z}{3}.
 \label{sel7:eq:deformation}
\end{equation}
This is a normalized independent covariant family with constant labelled Gram
\eqref{sel7:eq:equal-Gram}.  At $\phi=0$ it is $K^{\rm p}$, and for
$\phi\notin2\pi\mathbb Z$ every coefficient is invertible.
For the complete labelled edge instrument $\mathcal I_\phi$,
\begin{equation}
 \boxed{\|\mathcal I_\phi-\mathcal I_0\|_\diamond
       \le\frac{2}{\sqrt6}|e^{i\phi}-1|.}
 \label{sel7:eq:deformation-distance}
\end{equation}
The corresponding complete locked opportunity maps differ by at most
$2\vartheta|e^{i\phi}-1|/\sqrt6$.  Repeating them $k$ times and retaining
all flags gives the upper bound $2k\vartheta|e^{i\phi}-1|/\sqrt6$.
These are comparison estimates, not an implementation of either source from
the other.
\end{proposition}
\begin{proof}
Here $x=1,y=z$, so the equivalent constraint $|y|^2+3|x|^2=4$ and the
fixed sum imply normalization.  Neither $x$ nor $y$ vanishes.  For $z\ne1$,
$\beta,\gamma$ are nonzero and $5+z$ never vanishes, proving invertibility.
The Gram statement follows from \eqref{sel7:eq:Gram}.

Let $V_\phi x=\sum_e K_e(\phi)x\otimes|e\rangle$ be the channel isometry.
Covariance, or direct summation of the three projector components, gives
\[
 (V_\phi-V_0)^*(V_\phi-V_0)
 =\sum_e(K_e(\phi)-K_e(0))^*(K_e(\phi)-K_e(0))
 =\frac{|z-1|^2}{6}I.
\]
For any ancillary trace-class input, expanding the difference between
$V_\phi X V_\phi^*$ and $V_0XV_0^*$ into two terms bounds its trace norm
by $2\|V_\phi-V_0\|\|X\|_1$.
Tracing or dephasing the environment is a channel and cannot increase this
bound.  Dephasing the displayed edge basis retains the full classical record,
so \eqref{sel7:eq:deformation-distance} follows.
The locked no-response maps agree and the accepted difference has the common
factor $\vartheta$; the pointer/type projection instrument is a fixed channel.
Finally telescope the $k$-fold sequential channels, each of diamond norm one.
\end{proof}

\begin{remark}[The relevant source datum is not another Gram rank]
These two sources and the interpolating family already have the same maximal
six-dimensional coefficient span, the same positive one-letter Gram and the
same full forward algebra.  Rank saturation in that Hilbert Gram does not
select the operator-rank-minimal point.  The actual old two- and three-letter
record tests do distinguish it.  Neither the positive alternative nor its
strict negative counterpart is inferred by assigning unknown historical
coefficients to one of the displayed points.
\end{remark}

\section{Native donor rays, private alignment, and the external-action distinction}
\label{sel7:alignment}

\begin{proposition}[Every actual word ray on the projector branch]
\label{sel7:prop:word-rays}
For $K_e=P_e/\sqrt2$, every original length-$k$ word has the exact form
\begin{equation}
 \boxed{K_{e_k}\cdots K_{e_1}
 =2^{-k/2}\left(\prod_{j=1}^{k-1}
               \langle u_{e_{j+1}},u_{e_j}\rangle\right)
                |u_{e_k}\rangle\langle u_{e_1}|.}
 \label{sel7:eq:word-rays}
\end{equation}
It is nonzero precisely when no consecutive pair is complementary.
If $a(w)$ counts distinct adjacent transitions, its squared operator norm is
$2^{-k}4^{-a(w)}$.  The graph of allowed distinct transitions is connected,
so every pair of edge rays is linked by a nonzero word, with length at most
three.  All input and pure output rays still belong to the original six-ray
set.  The complementary projector orbit has the same conclusions after
complementing all labels.
\end{proposition}
\begin{proof}
Multiply rank-one dyads.  For root rays the absolute overlap is one on the
same ray, zero on complementary edges, and $1/2$ on distinct edges with one
shared vertex.  This gives the zero and norm statements.  Any complementary
pair has a common adjacent edge, proving the length bound.  The initial and
terminal dyads in \eqref{sel7:eq:word-rays} identify the rays for every
length; longer products do not create additional linear combinations of
physical outcomes.
\end{proof}

\begin{corollary}[Minimum operational rank does not realize the supplied private pair]
\label{sel7:cor:private-mismatch}
On the rank-minimal covariant branch:
\begin{enumerate}[label=\textup{(\roman*)},leftmargin=2.5em]
\item two distinct native pure output rays have absolute overlap either zero
or $1/2$, never $1/5$;
\item for a fixed native donor line $T=\C u_e$, the only native output ray in
$T^\perp$ is $\C u_{\bar e}$; hence native pure writes into that orthogonal
private plane have output frame rank at most one;
\item if a proposed donor is not a native root ray, no nonzero original word
has precisely that entire input line.
\end{enumerate}
The same restrictions survive exact record-respecting refinements and positive
classical selection of original histories.  They do not exclude a different
source representation or independently executed input/output maps.
\end{corollary}
\begin{proof}
The first three claims follow from the ray list and overlaps in
Proposition~\ref{sel7:prop:word-rays}.  V6's
Theorem~\ref{sel6:thm:efficient-refinement} says that an exact refinement of
an old history only scales its ray.  If a positive sum has a pure Choi output,
all nonzero summands lie on that output ray; selection therefore cannot
manufacture a new coherent output vector.  An isometric identification of two
physical rays preserves their absolute overlap, so the source's private
value $1/5$ is not repaired by a unitary change of coordinates.
\end{proof}

The source already gives the private overlap, whereas the geometric type
contrast is a line in a different coefficient/source construction.  The
corollary is a test of a proposed identification of these objects with the
native edge rays, not a proof that the private source does not exist.
A source-positive type-line Schur residual does not change either of these
physical ray lists without a specified action map.

\begin{remark}[The landed $M_3$ here is external, not an internal colour proof]
On this minimal acted-on carrier, the verified external triplet is irreducible,
so
\[
 C^*(K_e:e)\cap\{U_g:g\in S_4\}'=\C I_3.
\]
Thus \eqref{sel7:eq:full-algebra} cannot be renamed an internal colour factor.
It is an exact source-coefficient algebra on an external triplet.  This is
precisely the external-versus-multiplicity distinction already established in
V106 and retained in V1.  A colour identification must occur on the appropriate
relative multiplicity carrier of the common physical packet, not by equating
matrix sizes.
\end{remark}

\section{The classification has an operator-valued extension on multiplicities}
\label{sel7:multiplicity}

The preceding external-action qualification does not end the upstream
calculation.  It identifies where a genuinely internal source must be tested.
The same stabilizer argument gives a compact operator-valued source panel
without choosing an internal target algebra.

\begin{theorem}[Three actual multiplicity coefficients and their calibration]
\label{sel7:thm:multiplicity-panel}
Let $S=W\otimes M$ and let the independently verified external action be
$U_g=u_g\otimes I_M$, where $u_g$ is the standard triplet or its sign twist.
Suppose six independent original coefficients satisfy
\eqref{sel7:eq:source-identities} and resolved CP covariance.
Then, in the same common frame, there are operators $A,B,C\in\B(M)$ with
\begin{equation}
 \boxed{K_e=\frac1{\sqrt2}
 (P_e\otimes A+P_{\bar e}\otimes B+Q_e\otimes C),}
 \label{sel7:eq:matrix-family}
\end{equation}
subject to the exact identities
\begin{equation}
 \boxed{A+B+C=I_M,
 \qquad A^*A+B^*B+C^*C=I_M.}
 \label{sel7:eq:matrix-normalization}
\end{equation}
Conversely these identities give a normalized covariant source, and its six
coefficients are independent exactly when
$A-B\ne0$ and $A+B-2C\ne0$ as operators.
The operators are reconstructed by partial matrix elements, for example
\[
 A=\sqrt2(\langle u_e|\otimes I)K_e(|u_e\rangle\otimes I),
\]
with the analogous formulas on the other two stabilizer lines.
If the external matrix algebra is included among the independently occurring
base actions, then
\begin{equation}
 \boxed{C^*(\B(W)\otimes I,K_e:e)
       =\B(W)\otimes C^*(A,B,C),}
 \label{sel7:eq:relative-panel}
\end{equation}
and its intersection with the commutant of that external algebra is
$I_W\otimes C^*(A,B,C)$.
\end{theorem}
\begin{proof}
Lemma~\ref{sel7:lem:phase-lock} applies unchanged.  The three inequivalent
stabilizer characters now have common multiplicity space $M$, so a commuting
coefficient is block diagonal with three arbitrary operators on $M$.
Transport proves \eqref{sel7:eq:matrix-family}.  Summing its coefficients and
positive squares with \eqref{sel7:eq:root-frame} gives
\eqref{sel7:eq:matrix-normalization}.
For independence, the edge coefficient space decomposes as $1\oplus2\oplus3$.
The induced maps on those three summands have coefficients proportional to
$I_M$, $A+B-2C$, and $A-B$, respectively.  Each nonzero equivariant map
from these irreducible summands is injective; distinct types cannot cancel.
This proves the criterion and its converse.

The partial matrix elements recover the displayed operators directly.
Every generator lies in the right side of \eqref{sel7:eq:relative-panel},
proving one containment.  The actual external base provides all its matrix
units.  Compressing $K_e$ on a stabilizer line yields, for instance,
$P_e\otimes A/\sqrt2$; external matrix units transport and sum this to
$I_W\otimes A$.  The same holds for $B,C$.  This proves the other containment
as an algebraic statement.  Intersect with
$(\B(W)\otimes I)'=I_W\otimes\B(M)$ for the final assertion.
\end{proof}

\begin{proposition}[The three-operator calibration leaves a two-Kraus freedom]
\label{sel7:prop:two-Kraus}
Choose a real $3\times2$ contrast frame $O$ with
$O^{\mathsf T}O=I_2$ and $O^{\mathsf T}\mathbf1=0$.
Writing $(A_1,A_2,A_3)=(A,B,C)$, equations
\eqref{sel7:eq:matrix-normalization} are equivalent to
\begin{equation}
 A_j=\frac13 I+\sqrt{\frac23}\sum_{\mu=1}^2O_{j\mu}B_\mu,
 \qquad \sum_{\mu=1}^2B_\mu^*B_\mu=I.
 \label{sel7:eq:two-Kraus}
\end{equation}
In particular $C^*(A,B,C)=C^*(B_1,B_2)$.
This does not select its matrix size or represented algebra type.
\end{proposition}
\begin{proof}
Make the orthogonal coefficient change with columns $\mathbf1/\sqrt3,O$.
The uniform coefficient becomes $I/\sqrt3$ and the other two coefficients
$D_\mu$ satisfy $\sum_\mu D_\mu^*D_\mu=2I/3$.
Set $B_\mu=\sqrt{3/2}D_\mu$ and invert the change.  Both directions and the
algebra equality follow.  This is the inherited coherent-sum contrast method,
now applied to the multiplicity panel forced by the stabilizer decomposition;
it is not claimed as a new general channel normal form.
\end{proof}

\begin{example}[The multiplicity panel is not forced to be noncommutative]
\label{sel7:ex:matrix-freedom}
For the explicit contrast frame
\[
 O=\begin{pmatrix}1/\sqrt2&1/\sqrt6\\
                 -1/\sqrt2&1/\sqrt6\\0&-2/\sqrt6\end{pmatrix},
\]
take $B_1=I/\sqrt2$, $B_2=I/\sqrt2$ on $M=\C^2$.  The resulting
$A,B,C$ are scalar and the internal algebra is $\C I_2$.
Alternatively take $B_1=X/\sqrt2$, $B_2=Z/\sqrt2$ with Pauli matrices
$X,Z$; their algebra is $M_2$.  Both choices satisfy the same normalized
operator identities and give independent six-edge covariant lists, since
$A-B=\sqrt{4/3}B_1\ne0$ and
$A+B-2C=2B_2\ne0$.  This example is a mathematical separation, not an
assignment of either multiplicity panel to the historical source.
\end{example}

\begin{remark}[The scalar history test is not applied to matrix multiplicities]
The cubic test in Theorem~\ref{sel7:thm:record-test} uses three scalar
spectral parameters on a three-dimensional acted-on carrier.  It is not
asserted for noncommuting operator coefficients $A,B,C$ or for additional
external isotypic sectors.  Likewise a partial matrix element used in
reconstruction is not automatically an executed source intervention.
The operator-valued panel must be retained when those additional dimensions
actually occur.
\end{remark}

\section{V7 outcome and the source row which can now be evaluated}
\label{sel7:status}

The upstream question has changed from an unrestricted search through
unknown products to a finite alternative on a stated minimal covariant
class.  The coherent identity sum removes the projective phase ambiguity;
the stabilizer gives the entire three-parameter family; its original
rank-one histories are exactly the singular branch and are witnessed by
depth two; two- and three-letter original record masses test that branch
without fitting a Kraus frame or adding a Read.

There are two distinct positive results.  On rank saturation, the state update
is forced to be a tight projector instrument by an exact support-rank defect.
On the rank-two singular strata, pure two-letter histories already occur even
though no one-letter coefficient is pure.  Neither result inserts the missing
amplitude through a coherent recombination of old outcomes.

There are also concrete source-selection restrictions.  The same full
one-letter Gram and full coefficient algebra can describe either pure
one-letter histories or no exact pure history at any finite length.  A
generic minimal covariant coefficient list is in the latter class.  Even the
rank-minimal positive branch has a specific native input/output ray system
which does not match the supplied private overlap $1/5$.  Finally its raw
$M_3$ is external.  The operator-valued extension identifies the actual
multiplicity panel $A,B,C$ which must be reconstructed to make an internal
claim; normalization alone leaves that panel nonunique.

\begin{center}
\small
\begin{tabular}{L{.37\textwidth}L{.55\textwidth}}
\toprule
Question & Result and exact boundary\\
\midrule
Can the original six coefficients be constrained without an extra Read?
& Yes. Their coherent sum gives the exact total support-rank defect.\\[.3em]
Does CP covariance require a separately fitted Kraus phase law?
& No on an independent identity-calibrated list: the phases are forced.\\[.3em]
Is native rank-one-word existence finitely decided here?
& Yes on the complete three-dimensional covariant class: the cubic record
statistic vanishes exactly when a witness of depth at most two exists.\\[.3em]
Does a positive one-letter Gram decide that property?
& No. An explicit same-Gram continuous family has opposite answers at its
singular endpoint and every nearby nonzero parameter.\\[.3em]
Does operational rank minimality land the old private pair or internal colour?
& No. The native overlap list and the external scalar commutant exclude those
particular identifications.\\[.3em]
Is historical SM selection closed?
& No. The actual acted-on dimension, system covariance, corresponding original
record rows, common-domain/private map and internal multiplicity panel remain
to be identified; determinant and matter conditions remain separate.\\
\bottomrule
\end{tabular}
\end{center}

\paragraph{Next source audit.}
Before reconstructing every unknown $K_e$, test whether the actual retained
system has dimension three and whether the resolved \emph{system} covariance
holds, including the given coherent identity expansion.  These facts cannot
be read from the six edge names alone.  If they hold, evaluate the original
same-edge cylinders $p_2,p_3$ and, when needed, the opposite-edge cylinder.
An exact symbolic zero of the cubic test returns the explicit one- or
two-letter pure histories; a strictly positive lower bound stops the exact
finite-history search.  Neither test by itself identifies the categorical
type ray or private source state.

If the source instead has a verified external triplet with nontrivial
multiplicity, retain it and reconstruct $A,B,C$ from the actual mixed
system panel of Theorem~\ref{sel7:thm:multiplicity-panel}.  Do not compress
that multiplicity to a scalar model or interpret external Burnside blocks as
gauge factors.  The earlier return and common-input theorems can then be
applied to actual source-domain maps.  This is a source calculation with
checkable alternatives, not permission to select a favourable point of the
three-sphere or a favourable internal channel.

\paragraph{Proof, audit and preservation status.}
The V6 prefix is preserved byte for byte before its final document command.
The new exact-arithmetic checks verify the phase/stabilizer examples, root
frame and coefficient Gram, source normalization, rank-one products on all
singular cases, the input-independent repeated-edge and opposite-edge effects,
the same-Gram counterexample, the phase deformation, the private-ray overlap
obstruction and explicit multiplicity examples.  V6's regression script is
also rerun.  Finite tests support the displayed identities; the all-word
negative statements and complete parameter classification are established by
the proofs, not by finite enumeration.  No unknown historical source row is
filled by an example, no proof-assistant formalization is claimed, and the
results have not had independent peer review.

\begin{thebibliography}{99}
\bibitem[15]{sel7-EldarForney}
Y.~C.~Eldar and G.~D.~Forney, Jr.,
\emph{Optimal tight frames and quantum measurement},
IEEE Trans.\ Inform.\ Theory \textbf{48} (2002), 599--610;
\href{https://doi.org/10.1109/18.985949}{doi:10.1109/18.985949};
\href{https://arxiv.org/abs/quant-ph/0106070}{arXiv:quant-ph/0106070}.
\end{thebibliography}

% END OF SOURCE-SELECTION V7 DEVELOPMENT BLOCK.
% Append subsequent developments before the final document terminator.
% Preserve the V1--V7 bodies and state necessary amendments explicitly.


\clearpage
\begin{center}
{\Large\bfseries Source-selection continuation V8}\\[.4em]
{\large Internal extraction from original covariant coefficients}\\[.25em]
Quadratic--cubic recovery, exact algebra decomposition, and the reduced source
\end{center}

\section{The remaining external-control hypothesis can be removed}
\label{sel8:context}

V7 reached the multiplicity-valued source class
\[
 K_e=2^{-1/2}(P_e\otimes A+P_{\bar e}\otimes B+Q_e\otimes C),
 \qquad A+B+C=I,
 \qquad A^*A+B^*B+C^*C=I.
\]
Its internal-algebra assertion still adjoined independently occurring
external matrix controls.  The present continuation removes that extra
algebra-generation hypothesis.  All operators used to generate the algebra
below are the six original coefficients and their formal adjoints; the
external representation is retained to identify invariant operators, not
added as a control bank.  Quadratic and cubic combinations of these original
coefficients recover $I_W\otimes A$, $I_W\otimes B$, and $I_W\otimes C$,
even at singular coefficient supports.  This is an assertion about the
represented original source algebra, not a prescription for executing a
coherent sum, an adjoint channel, or a Moore--Penrose inverse.

The calculation gives a second, output-sensitive route.  The original edge
channel has an autonomous reduced channel on the actual multiplicity factor.
The support of one rank-at-most-two positive residual Choi matrix determines
the entire internal algebra and its calibrated commutator form.  This reduces
the needed operator data without assuming a favourable internal factor.

\paragraph{Audit and nonduplication.}
The V7 multiplicity panel and its two-Kraus normal form, V1's
multiplication-compatible word-action test, the older V107 external twirl,
and the positive Choi--commutator construction in V3 of \cite{Rigidity} are
inputs.  They do not by themselves prove that the original six coefficients
contain their own multiplicity operators without an adjoined external action.
The new results are the explicit quadratic--cubic extraction, the exact raw
source algebra including its abelian exceptional sector, the autonomous
multiplicity-channel specialization, and its sharply delimited selection
consequences.  The source normalization is the one in
\cite{Grand}, \src{def:locked-opportunity-instrument} and
\src{thm:GT76-locked-normalform}; it does not assign a numerical multiplicity
panel.  No such assignment is made here.

The general finite functional-calculus and Kraus/Choi facts are standard;
\cite{sel8-Watrous} supplies background for the latter.  Every specialized
assertion is proved below.  The claims of development concern the audited
programme, not mathematical priority over the operator-algebra literature.

\paragraph{Fixed source class.}
Assume the actual exposed system is $S=W\otimes M$, where $W$ is the
three-dimensional standard tetrahedral representation or its sign twist,
$m=\dim M\ge1$, and the source satisfies precisely the hypotheses of
Theorem~\ref{sel7:thm:multiplicity-panel}.  These include resolved system
covariance and the identity-calibrated six-coefficient source, not merely
an action on six names.  Put
\[
 \mathcal A_K=C^*(I,K_e:e),\qquad
 \mathcal R=C^*(A,B,C)\subseteq\B(M),\qquad
 \mathcal I_K=\mathcal A_K\cap(u(S_4)\otimes I_M)'.
\]
The prime in the last expression refers to the external representation,
not the algebra generated by its physically executed controls.  The
cofinal regulator, the type/private identification, and determinant and
matter incidences are not assumed by specifying this finite class.

\section{Quadratic and cubic source words recover the multiplicity operators}
\label{sel8:extraction}

Use the orthonormal pair-contrast basis of $W$.  The four tetrahedral vectors
may be written as the rows of
\[
 \frac12\begin{pmatrix}1&1&1\\1&-1&-1\\-1&1&-1\\-1&-1&1\end{pmatrix}.
\]
The six unoriented root rays have coordinates
$(e_i\pm e_j)/\sqrt2$.  Relabel a complementary edge pair by $(k,+),(k,-)$,
where $\{i,j,k\}=\{1,2,3\}$.  Write $E_{ij}=|i\rangle\langle j|$ and
\begin{equation}
 T_{ij}=E_{ij}+E_{ji},\qquad H_k=E_{kk}-I_3/3,\qquad
 P_k^\pm=\tfrac12(I_3-E_{kk}\pm T_{ij}),\quad Q_k=E_{kk}.
 \label{sel8:eq:external-coordinates}
\end{equation}
These are changes of coordinates on the given root representation; they
are not additional source operations.

Define two intrinsic multiplicity contrasts
\begin{equation}
 D=A-B,\qquad F=A+B-2C.
 \label{sel8:eq:DF}
\end{equation}
Then
\begin{equation}
 \boxed{A=(2I+F+3D)/6,\quad B=(2I+F-3D)/6,\quad C=(I-F)/3,}
 \qquad 3D^*D+F^*F=4I_M.
 \label{sel8:eq:contrast-calibration}
\end{equation}
The last equation is an equivalent form of the positive normalization;
there is no commutativity or normality assumption on $D,F$.  In particular
$\|D\|\le2/\sqrt3$ and $\|F\|\le2$.

The following linear combinations of actual source coefficients are useful:
\begin{equation}
 X_{ij}=X_{ji}:=\sqrt2(K_{k,+}-K_{k,-})=T_{ij}\otimes D,
 \quad
 Y_k:=\tfrac23I_S-\sqrt2(K_{k,+}+K_{k,-})=H_k\otimes F.
 \label{sel8:eq:source-contrasts}
\end{equation}
They belong to the source algebra at word degree at most one, with the
identity counted as degree zero.  They need not be single executed letters.

\begin{theorem}[Four low-degree invariant operator combinations]
\label{sel8:thm:cubic-extraction}
Form, solely from \eqref{sel8:eq:source-contrasts},
\begin{align}
 Q_D&=\tfrac12\sum_{i<j}X_{ij}^*X_{ij},&
 V_D&=\tfrac12\sum_{i,j,k\ {
m distinct}}X_{ij}X_{jk}^*X_{ki},\nonumber\\
 Q_F&=\tfrac32\sum_kY_k^*Y_k,&
 V_F&=\tfrac92\sum_kY_kY_k^*Y_k.
 \label{sel8:eq:four-moments}
\end{align}
They are externally invariant, and obey the exact identities
\begin{equation}
 \boxed{Q_D=I_3\otimes D^*D,\quad V_D=I_3\otimes DD^*D,
 \quad Q_F=I_3\otimes F^*F,\quad V_F=I_3\otimes FF^*F.}
 \label{sel8:eq:four-identities}
\end{equation}
In particular
\begin{equation}
 \boxed{I_3\otimes D=V_DQ_D^\dagger,\qquad
        I_3\otimes F=V_FQ_F^\dagger.}
 \label{sel8:eq:invariant-recovery}
\end{equation}
All identities remain valid on singular supports.  The recovered operators
are polynomials in the original star-closed coefficient bank of degree at
most $2m+1$.  The calibration also gives $3Q_D+Q_F=4I_S$, so one of the
two quadratic quantities is redundant when that exact source identity has
already been verified.
\end{theorem}
\begin{proof}
The external identities are
\[
 \sum_{i<j}T_{ij}^2=2I_3,\qquad
 T_{ij}T_{jk}T_{ki}=E_{ii}\quad(i,j,k\text{ distinct}),
 \qquad
 \sum_kH_k^2=\tfrac23I_3,\quad\sum_kH_k^3=\tfrac29I_3.
\]
There are two ordered triples beginning with each $i$, so substitution into
\eqref{sel8:eq:four-moments} gives \eqref{sel8:eq:four-identities}.
For any matrix $Z$, putting $Q=Z^*Z$ gives
$ZZ^*Z=ZQ$ and $ZQ Q^\dagger=Z$, because $Z$ vanishes on $\Ker Q$.
This proves \eqref{sel8:eq:invariant-recovery}.  Each positive matrix $Q_D$
or $Q_F$ has at most $m$ distinct eigenvalues.  Interpolation of the function
$\lambda\mapsto1/\lambda$ on its positive eigenvalues, with value zero at
zero when present, represents its pseudoinverse by a polynomial of degree
at most $m-1$.  Multiplication by the corresponding cubic $V$ gives word
degree at most $3+2(m-1)$.  When the matrix is zero the recovered operator
is zero and no inverse is taken.
\end{proof}

\begin{theorem}[Internal algebra of the original six-coefficient source]
\label{sel8:thm:internal-algebra}
Without adjoining any external matrix control to $\mathcal A_K$,
\begin{equation}
 \boxed{\mathcal I_K=I_3\otimes\mathcal R
       =C^*(I_S,Q_D,V_D,Q_F,V_F).}
 \label{sel8:eq:internal-algebra}
\end{equation}
Thus the complete invariant system algebra is generated by the multiplicity
contrasts already encoded in the original source.  This strengthens the
intersection conclusion of Theorem~\ref{sel7:thm:multiplicity-panel} by
removing its extra algebra-generation assumption.
\end{theorem}
\begin{proof}
Every $K_e$ lies in $M_3(\C)\otimes\mathcal R$, hence so does
$\mathcal A_K$.  Theorem~\ref{sel8:thm:cubic-extraction} puts
$I_3\otimes D$ and $I_3\otimes F$ in $\mathcal A_K$; by
\eqref{sel8:eq:contrast-calibration} they generate $I_3\otimes\mathcal R$.
The external representation on $W$ is irreducible, so the externally
invariant part of $M_3\otimes\mathcal R$ is exactly
$I_3\otimes\mathcal R$.  This proves both containments.  The four
combinations in \eqref{sel8:eq:four-identities} belong to that algebra and
recover its generators, proving the last equality.
\end{proof}

\begin{remark}[What was removed]
The theorem does not assume $I_3\otimes\mathcal R$ as an initial gauge
bank.  It exhibits it inside the original generated system algebra.  Nor
does it identify $\mathcal R$ with $M_3$ by a dimension match.  Its actual
represented type is still determined by the source matrices or their
output-sensitive Choi data.  Polynomial functional calculus is an algebraic
recovery, not physical availability of an inverse filter or of a star word as
one amplitude.  V6's efficient-record restriction remains in force.
\end{remark}

\section{The complete raw algebra and its precisely located diagonal sector}
\label{sel8:raw-algebra}

Let $p\in Z(\mathcal R)$ be the central support of the ideal generated by
$D$ in $\mathcal R$, and put $q=I_M-p$.  Equivalently, $qM$ is the largest
$\mathcal R$-reducing subspace on which $D$ vanishes.  Write
$\mathcal R_p=p\mathcal R|_{pM}$ and $\mathcal R_q=q\mathcal R|_{qM}$,
with zero-dimensional summands omitted.  The calibration gives
$F_q^*F_q=4I_{qM}$; finite dimensionality makes $F_q/2$ unitary.
Consequently $\mathcal R_q=C^*(F_q)$ is abelian.  Denote by
$\mathcal D_3$ the diagonal algebra in the basis
\eqref{sel8:eq:external-coordinates}.

\begin{theorem}[Exact raw-algebra decomposition]
\label{sel8:thm:raw-algebra}
On $S=(W\otimes pM)\oplus(W\otimes qM)$,
\begin{equation}
 \boxed{\mathcal A_K=(M_3(\C)\otimes\mathcal R_p)
                \oplus(\mathcal D_3\otimes\mathcal R_q),}
 \label{sel8:eq:raw-algebra}
\end{equation}
and, in the corresponding representations,
\begin{equation}
 \boxed{\mathcal A_K'=(I_3\otimes\mathcal R_p')
                \oplus(\mathcal D_3\otimes\mathcal R_q').}
 \label{sel8:eq:raw-commutant}
\end{equation}
In particular the original source contains $M_3(\C)\otimes I_M$ exactly
when $q=0$.  If $\mathcal R$ is simple and $D\ne0$, then $q=0$ automatically.
The invariant internal algebra remains $I_3\otimes\mathcal R$ in both cases.
\end{theorem}
\begin{proof}
By Theorem~\ref{sel8:thm:internal-algebra}, any $X\in\mathcal A_K'$ has
external matrix entries $X_{ij}\in\mathcal R'$.  Commutation with
$H_k\otimes F$ implies $FX_{ij}=0$ for $i\ne j$: at least one of the
corresponding diagonal differences of $H_k$ is nonzero.  Given distinct
$i,j,k$, commutation with $T_{jk}\otimes D$ implies $X_{ij}D=0$, by comparing
the $(i,k)$ block.  Since the entries commute with $\mathcal R$, also
$DX_{ij}=0$.  Thus
\[
 4X_{ij}=(3D^*D+F^*F)X_{ij}=0\qquad(i\ne j).
\]
The commutant is therefore externally diagonal, say
$X=\operatorname{diag}(X_1,X_2,X_3)$ with $X_i\in\mathcal R'$.
Its remaining conditions are
$(X_i-X_j)D=0$ for each $i\ne j$ and the adjoint equations.  Since every
difference commutes with $\mathcal R$, this is equivalent to annihilating
the whole ideal $\mathcal R D\mathcal R$, hence to
$(X_i-X_j)p=0$.  The $p$ parts must agree; the $q$ parts can be chosen
independently in $\mathcal R_q'$.  The central projections $p,q$ belong to
$\mathcal R$, so no intertwiner between them remains.  This proves
\eqref{sel8:eq:raw-commutant}.  Finite bicommutant closure gives
\eqref{sel8:eq:raw-algebra}.

If $q\ne0$, a nondiagonal external matrix tensored with $I_{qM}$ does not
belong to the $q$ summand; if $q=0$, the full external matrix algebra is
present.  A nonzero ideal in a simple algebra is the whole algebra, proving
the simple-block assertion.  The final invariant statement also follows
from Theorem~\ref{sel8:thm:internal-algebra}.
\end{proof}

\begin{proposition}[A finite positive certificate for the exceptional sector]
\label{sel8:prop:ideal-certificate}
Let $r=\rank D$ and let $\mathcal G=\{D,D^*,F,F^*\}$.  If $D\ne0$, set
\begin{equation}
 H_D=\sum_{|w|\le m-r}wDD^*w^*,
 \label{sel8:eq:ideal-gram}
\end{equation}
where the sum includes the empty word and every word in $\mathcal G$ of the
indicated lengths.  Then $\supp H_D=p$, and $q=0$ exactly when $H_D\succ0$.
If $D=0$, set $H_D=0$ and $q=I_M$.  A zero direction of $H_D$ belongs to the
explicit diagonal sector of Theorem~\ref{sel8:thm:raw-algebra}.
\end{proposition}
\begin{proof}
The range of $H_D$ is the span of all $w\Ran D$ in the indicated list.
Starting from the $r$-dimensional range of $D$, each nonflat word extension
increases dimension by at least one.  If one step is flat, invariance under
the star-closed generating bank prevents all further increases.  Hence
length $m-r$ suffices to reach $\mathcal R\Ran D$.  Its projection commutes
with $\mathcal R$ by reducing invariance and belongs to $\mathcal R$ by
functional calculus of $H_D$.  It is the central support of the ideal of
$D$, as claimed.  This uses the existing finite source-cyclicity mechanism
only to identify the newly computed exceptional sector; it is not a new
operational word-execution rule.
\end{proof}

\begin{example}[An independent covariant source can retain a diagonal sector]
\label{sel8:ex:dark}
On $M=\C^2$ choose
\[
 D=\operatorname{diag}(1,0),\qquad F=\operatorname{diag}(-1,2).
\]
They satisfy $3D^*D+F^*F=4I$.  The reconstructed coefficient triples on the
two multiplicity lines are respectively
\[
 (A,B,C)=(2/3,-1/3,2/3),\qquad (2/3,2/3,-1/3).
\]
All six full $K_e$ are invertible, and the six are independent because both
$D$ and $F$ are nonzero.  The source is unital as well as trace preserving.
Nevertheless
\[
 \mathcal R=\C^2,\quad p=\operatorname{diag}(1,0),\quad
 \mathcal A_K\cong M_3\oplus\mathcal D_3,
 \quad\dim\mathcal A_K=12,\quad\dim\mathcal A_K'=4.
\]
Thus invertibility, independence, covariance and faithful stationarity do
not imply that the source executes the full external matrix algebra on
every multiplicity summand.  The invariant internal algebra is exactly the
expected $\C^2$, not a larger algebra manufactured by adding controls.
\end{example}

\section{The actual source induces an autonomous multiplicity channel}
\label{sel8:reduced-source}

Let $\Phi_*(X)=\sum_eK_eXK_e^*$ be the normalized accepted edge channel,
and let $\Phi^*$ denote its Heisenberg dual.  The following calculation
removes the need to choose partial coefficient phases when only the
internal algebra and its rigidity form are sought.

\begin{theorem}[Autonomous reduced source and a rank-two positive Choi residual]
\label{sel8:thm:reduced-channel}
The original channel satisfies, for every $x\in\B(M)$ and every
$X\in\B(W\otimes M)$,
\begin{align}
 \Phi^*(I_3\otimes x)&=I_3\otimes\Psi^*(x),\nonumber\\
 \operatorname{Tr}_W\Phi_*(X)&=\Psi_*(\operatorname{Tr}_W X),
 \label{sel8:eq:autonomy}
\end{align}
where the trace-preserving map determined by the source is
\begin{equation}
 \boxed{\Psi_*(\rho)=A\rho A^*+B\rho B^*+C\rho C^*
 =\tfrac13\rho+\tfrac12D\rho D^*+\tfrac16F\rho F^*.}
 \label{sel8:eq:reduced-channel}
\end{equation}
Also
\[
 \Phi_*(I_3/3\otimes\rho)=I_3/3\otimes\Psi_*(\rho).
\]
If $J_\Psi$ is the unnormalized Choi matrix, the data-defined residual
\begin{equation}
 \boxed{J_\Xi:=\tfrac32(J_\Psi-\tfrac13|I\rangle\!\rangle
                                      \langle\!\langle I|)
   =\tfrac34|D\rangle\!\rangle\langle\!\langle D|
      +\tfrac14|F\rangle\!\rangle\langle\!\langle F|\succeq0}
 \label{sel8:eq:rank-two-Choi}
\end{equation}
has rank at most two.  Its unvectorized support determines
\begin{equation}
 \boxed{\mathcal I_K=I_3\otimes
 C^*(I,\operatorname{unvec}\operatorname{Ran}J_\Xi).}
 \label{sel8:eq:Choi-internal}
\end{equation}
The full edge source has a faithful stationary state if and only if
$\Psi_*$ has one.
\end{theorem}
\begin{proof}
On one edge the three external projections are mutually orthogonal, so
\[
 K_e^*(I_3\otimes x)K_e=\tfrac12(P_e\otimes A^*xA+
 P_{\bar e}\otimes B^*xB+Q_e\otimes C^*xC).
\]
The three projection sums over all six edges are $2I_3$.  This gives the
Heisenberg identity; duality gives the partial-trace identity, including
entangled $X$.  The same projection calculation in Schr\"odinger form
gives the product-state identity.  Substituting
\eqref{sel8:eq:contrast-calibration} into $\sum A_j\rho A_j^*$ cancels
all mixed terms and yields \eqref{sel8:eq:reduced-channel}.
The matrix in \eqref{sel8:eq:rank-two-Choi} is the Choi matrix of the
trace-preserving channel with two coefficients $\sqrt3D/2,F/2$.
Its support is their vectorized span, independent of the Kraus factorization.
Unvectorizing and taking a unital star algebra recovers $\mathcal R$,
and Theorem~\ref{sel8:thm:internal-algebra} gives
\eqref{sel8:eq:Choi-internal}.

A faithful stationary $\rho$ of $\Psi_*$ lifts to the faithful stationary
$I_3/3\otimes\rho$ of $\Phi_*$.  Conversely, a faithful stationary state
of $\Phi_*$ has a faithful partial trace, which is stationary by
\eqref{sel8:eq:autonomy}.  No detailed-balance assumption is needed.
\end{proof}

\begin{remark}[Reduced quantum output is not the classical edge record]
The map $\Psi_*$ is derived from the actual channel and the independently
identified tensor factor.  Its input and output operator data are required;
an edge-count table is not that channel.  The two terms in $J_\Xi$ need not
be physically resolved outcomes.  Subtracting the known identity component
is a positive data calculation justified by the proven normalization, not
an instruction to execute a subtracted event.  No original record is
coherently recombined in asserting the algebra equality.
\end{remark}

\begin{theorem}[The reduced Choi packet preserves the original internal calibration]
\label{sel8:thm:calibrated-form}
Use the star-closed commutator form of the original accepted coefficients,
\[
 q_K(Z)=\sum_e\bigl(\|[K_e,Z]\|_{\HS}^2+
                         \|[K_e^*,Z]\|_{\HS}^2\bigr),
\]
and the Hilbert--Schmidt isometric embedding $\iota(x)=I_3\otimes x/\sqrt3$.
Then
\begin{equation}
 \boxed{q_K(\iota x)=q_{\rm int}(x)
 =\tfrac12\bigl(\|[D,x]\|_{\HS}^2+\|[D^*,x]\|_{\HS}^2\bigr)
  +\tfrac16\bigl(\|[F,x]\|_{\HS}^2+\|[F^*,x]\|_{\HS}^2\bigr).}
 \label{sel8:eq:internal-calibration}
\end{equation}
This is also the star-closed commutator form of any Kraus representation of
$\Psi_*$.  Its kernel is $\mathcal R'$, and it is determined by $J_\Psi$
or, equivalently with the known weights, by $J_\Xi$.
\end{theorem}
\begin{proof}
On one edge, Hilbert--Schmidt orthogonality of the three external rank-one
projections gives $1/6$ times the sum of the three multiplicity commutator
squares for $[K_e,\iota x]$.  Sum six edges and their adjoints.  The affine
formulas for $A,B,C$ cancel the $D,F$ cross terms and give
\eqref{sel8:eq:internal-calibration}.  A Kraus isometry preserves the sums
of commutator squares, including the adjoint family, so the form is
factorization independent.  Its nonnegative summands vanish exactly on
$C^*(D,F)'$.  Equivalently the ordinary Choi commutator formula expresses
it linearly in the positive Choi packet.  This specializes the inherited
positive-CP certificate to the actual source normalization.
\end{proof}

For the locked physical opportunity channel with idle coefficient
$\sqrt{1-\vartheta}I$ and accepted coefficients $\sqrt\vartheta K_e$,
the idle term commutes with every test and the form above is multiplied by
$\vartheta$.  Pointer/type blocks and their own commutation constraints
must be retained when they are part of the exposed system.  Conditional
normalization of accepted events does not remove this physical scale.

\section{Stable extraction does not need a fixed nonzero singular-value floor}
\label{sel8:stability}

Exact algebra recovery can use a pseudoinverse on the actual finite support.
For noisy operator data that operation is poorly conditioned near support
birth.  The specific cubic numerators above provide a uniform alternative.
All norms in this section are operator norms unless explicitly marked.

\begin{proposition}[Regularized recovery and a quantitative data error]
\label{sel8:prop:regularized}
For a matrix $Z$, set $Q=Z^*Z$ and $V=ZQ$.  For every $\eta>0$,
\begin{equation}
 \boxed{\|Z-V(Q+\eta I)^{-1}\|\le\tfrac12\sqrt\eta.}
 \label{sel8:eq:bias}
\end{equation}
Suppose $\widehat Q\succeq0$ and $\widehat V$ are supplied with certified
errors $\|\widehat Q-Q\|\le\epsilon_Q$ and
$\|\widehat V-V\|\le\epsilon_V$.  If $\|Z\|\le M_Z$, then
\begin{equation}
 \boxed{\|\widehat V(\widehat Q+\eta I)^{-1}-Z\|
 \le\tfrac12\sqrt\eta+
                (\epsilon_V+M_Z\epsilon_Q)/\eta.}
 \label{sel8:eq:noisy-recovery}
\end{equation}
For $c=\epsilon_V+M_Z\epsilon_Q>0$, the choice
$\eta=(4c)^{2/3}$ makes this bound $3\,2^{-4/3}c^{1/3}$.
These estimates apply to both contrasts in
Theorem~\ref{sel8:thm:cubic-extraction}, using
$M_D=2/\sqrt3$ and $M_F=2$.
\end{proposition}
\begin{proof}
On a right singular vector with singular value $s$, the bias magnitude is
$s\eta/(s^2+\eta)\le\sqrt\eta/2$.  This proves
\eqref{sel8:eq:bias}, including $s=0$.
For the data perturbation write the exact identity
\begin{align*}
 \widehat V(\widehat Q+\eta I)^{-1}-V(Q+\eta I)^{-1}
 &=(\widehat V-V)(\widehat Q+\eta I)^{-1}\\
 &\quad+V(Q+\eta I)^{-1}(Q-\widehat Q)(\widehat Q+\eta I)^{-1}.
\end{align*}
The inverse norm is at most $1/\eta$, and
$\|V(Q+\eta I)^{-1}\|=\|ZQ(Q+\eta I)^{-1}\|\le M_Z$.
The triangle inequality proves \eqref{sel8:eq:noisy-recovery}.
Elementary minimization gives the stated choice and value.
\end{proof}

\begin{corollary}[A calibrated internal Howe stopping test]
\label{sel8:cor:robust-Howe}
Suppose a represented candidate algebra $\mathcal N\subseteq\B(M)$ is
independently known to commute with $D,F$ and their adjoints.  Recover
$\widehat D,\widehat F$ with errors $e_D,e_F$, for example by
Proposition~\ref{sel8:prop:regularized}.  On $\mathcal N^\perp$ let
$\widehat s$ be the least singular value of the measured stacked map with
weights $1/\sqrt2$ on $\operatorname{ad}_{\widehat D}$ and its adjoint
coefficient, and $1/\sqrt6$ on the two $\widehat F$ commutators.  If
\begin{equation}
 \boxed{\widehat s>2\sqrt{e_D^2+e_F^2/3},}
 \label{sel8:eq:Howe-stopping}
\end{equation}
then $\mathcal R'=\mathcal N$ and the true gap of
\eqref{sel8:eq:internal-calibration} on this complement is at least
$(\widehat s-2\sqrt{e_D^2+e_F^2/3})^2$.
\end{corollary}
\begin{proof}
The bound $\|[Z,x]\|_{\HS}\le2\|Z\|\|x\|_{\HS}$ gives
$2\sqrt{e_D^2+e_F^2/3}$ for the norm difference of the two weighted
stacked maps.  Apply the least-singular-value perturbation inequality and
the exact kernel identity in Theorem~\ref{sel8:thm:calibrated-form}.
This is the inherited Howe certificate with the actual source weights,
not an assumption that a fitted candidate commutant is protected.
\end{proof}

\begin{remark}[Data positivity and execution remain separate]
The hypothesis $\widehat Q\succeq0$ must be checked, and its error bound
must refer to the actual positive $Q$; an unbudgeted spectral fit is not
an exact measurement.  A mathematical external twirl can remove noninvariant
operator reconstruction noise without increasing its norm, but is not a new
executed control.  Regularization constructs estimates of the original
algebraic generators.  It does not turn a small wrong-donor amplitude into
an exactly absent physical history, as V5 already cautions.
\end{remark}

\section{Cubic information is genuinely needed on an exactly calibrated family}
\label{sel8:order-sharpness}

For the external representation define its mathematical conditional
expectation
\[
 \mathsf T(Z)=\frac1{24}\sum_{g\in S_4}(u_g\otimes I)Z(u_g^*\otimes I)
            =\frac{I_3}{3}\otimes\operatorname{Tr}_W Z.
\]
This is the inherited external twirl.  It is used to compare invariant
operator data, not as a newly available history operation.

\begin{theorem}[Complete invariant degree-two data can hide internal noncommutativity]
\label{sel8:thm:quadratic-blindness}
On the same multiplicity qubit compare the two panels
\[
 (D_0,F_0)=(Z,Z),\qquad (D_1,F_1)=(X,Z),
\]
where $X,Z$ are the Hermitian Pauli involutions.  Reconstruct $A,B,C$ and
then the original six $K_e$ by the stated formulas.  Both panels are
normalized, covariant and six-coefficient independent.  For every labelled
star word $w$ of length at most two,
\begin{equation}
 \boxed{\mathsf T(w(\bm K^{(0)}))=\mathsf T(w(\bm K^{(1)}))}
 \label{sel8:eq:quadratic-blindness}
\end{equation}
as operators on the same $W\otimes\C^2$.  Their complete labelled
one-letter Hilbert--Schmidt Grams are also equal.  Nevertheless,
\begin{equation}
 \boxed{\mathcal I_{K^{(0)}}=I_3\otimes\C^2,
 \qquad\mathcal I_{K^{(1)}}=I_3\otimes M_2(\C).}
 \label{sel8:eq:two-internal-types}
\end{equation}
The cubic combinations distinguish the panels: $V_D=I_3\otimes D_i$ and
$V_F=I_3\otimes Z$.  Thus degree two is insufficient in general for the
invariant operator panel, whereas the degree-two and degree-three data in
Theorem~\ref{sel8:thm:cubic-extraction} are sufficient with finite
functional calculus.
\end{theorem}
\begin{proof}
Both panels obey $3D_i^2+F_i^2=4I$ and $D_i,F_i\ne0$, proving the stated
normalization and independence.  In the original coefficient expression,
the nonconstant external parts lie in the span of the off-diagonal $T_{ij}$
and the traceless diagonal $H_k$.  Every such matrix is traceless, and
\[
 \Tr(T_{ij}H_k)=0.
\]
Consequently a twirl of a word of degree at most two can contain only an
identity coefficient and terms proportional to $D_i^2$, $F_i^2$, or,
for star words, $D_i^*D_i,D_iD_i^*,F_i^*F_i,F_iF_i^*$.  All those
multiplicity products are identity for both panels.  The external trace
coefficients depend only on the same labels, proving
\eqref{sel8:eq:quadratic-blindness}.  Taking the full trace proves the Gram
identity.  The contrasts generate $\C^2$ in the first case and $M_2$ in
the second, giving \eqref{sel8:eq:two-internal-types} by
Theorem~\ref{sel8:thm:internal-algebra}.  The cubic identities follow from
$D_i^*D_i=F_i^*F_i=I$.
\end{proof}

\begin{proposition}[Exact calibrated internal spectra for the separating pair]
\label{sel8:prop:Pauli-gaps}
On the Hilbert--Schmidt normalized Pauli basis of the traceless multiplicity
operators, the form \eqref{sel8:eq:internal-calibration} has spectrum
\[
 \{0,16/3,16/3\}\quad\text{for }(D,F)=(Z,Z),\qquad
 \{4/3,4,16/3\}\quad\text{for }(D,F)=(X,Z).
\]
The first panel has the additional diagonal commutant direction $Z$; the
second has only the scalar commutant and gap $4/3$.  The original raw source
algebras have dimensions $18$ and $36$, respectively.
\end{proposition}
\begin{proof}
Hermiticity reduces the form to
$\|[D,x]\|_{\HS}^2+\|[F,x]\|_{\HS}^2/3$.
Two distinct normalized Pauli directions have commutator squared norm four;
a direction commutes with itself.  This gives the two spectra directly.
Here $D$ is invertible, so $p=I$ and
Theorem~\ref{sel8:thm:raw-algebra} gives
$M_3\otimes\C^2$ and $M_3\otimes M_2$.
\end{proof}

\begin{remark}[Why there is no conflict with a quadratic commutator certificate]
The comparison above concerns invariant words made only from the original
coefficients, with no inserted system test.  A commutator certificate is
quadratic in the coefficients but also contains an independently specified
operator test $x$ between them.  Those output-sensitive matrix elements are
not the same data as the untested degree-two twirls.  Nor are the cubic
operator combinations the scalar repeated-edge probabilities of V7.
An ordinary classical record table does not automatically supply them.
\end{remark}

\section{The normalized source class permits every internal algebra without pure histories}
\label{sel8:freedom}

The positive extraction results answer what the source algebra contains
once its actual multiplicity panel is known.  They should not be converted
into a uniqueness assertion about that panel.  The following realization
stays within the very same six-edge normalization and covariance class and
shows how much freedom still remains.

\begin{lemma}[Two arbitrarily small unitary perturbations generate any finite algebra]
\label{sel8:lem:two-unitaries}
Let $\mathcal R\subseteq\B(M)$ be any represented finite-dimensional unital
complex $C^*$-algebra.  For every $\delta>0$, there are unitaries
$U,V\in\mathcal R$ with
\[
 C^*(U,V)=\mathcal R,\qquad
 \|U-I\|<\delta,\quad\|V-I\|<\delta.
\]
\end{lemma}
\begin{proof}
Write $\mathcal R=\bigoplus_aM_{n_a}(\C)\otimes I_{r_a}$ in its actual
representation.  Choose a self-adjoint diagonal $H$ whose eigenvalues at all
pairs $(a,i)$ are distinct, and a self-adjoint $L$ with nonzero nearest-neighbor
entries along a chain inside each $M_{n_a}$ block and no off-block entries.
Spectral projections of $H$ and compressions of $L$ generate every matrix
unit in every simple block, so $C^*(H,L)=\mathcal R$.  For sufficiently
small $\epsilon>0$, $U=e^{i\epsilon H}$ and $V=e^{i\epsilon L}$ are
$\delta$-close to identity.  Finite spectral interpolation recovers $H,L$
from $U,V$ when their eigenvalue phases do not wrap or collide.  Hence
$C^*(U,V)=\mathcal R$.  Repeated representation multiplicities are retained
as identity factors throughout.
\end{proof}

\begin{theorem}[Exact internal realization with all original letters invertible]
\label{sel8:thm:all-algebras}
For every such represented $\mathcal R$ and every $0<\delta<1/2$, take
$U,V$ as in Lemma~\ref{sel8:lem:two-unitaries} and set
\begin{equation}
 D=U,\quad F=-V,\qquad
 A=(2I-V+3U)/6,\quad B=(2I-V-3U)/6,\quad C=(I+V)/3.
 \label{sel8:eq:unitary-family}
\end{equation}
The resulting six original edge coefficients satisfy all hypotheses of
V7's multiplicity-valued class, and
\begin{equation}
 \boxed{\mathcal A_K=M_3(\C)\otimes\mathcal R,
 \quad\mathcal I_K=I_3\otimes\mathcal R,
 \quad\sigma_*=I_{3m}/(3m)\text{ is faithful stationary}.}
 \label{sel8:eq:arbitrary-internal}
\end{equation}
Every original coefficient is invertible, with
\begin{equation}
 \boxed{s_{\min}(K_e)\ge\frac{1-2\delta}{3\sqrt2}>0.}
 \label{sel8:eq:invertible-floor}
\end{equation}
Therefore no nonzero finite original history has operator rank one, and
V6's exact record refinements cannot supply one.  The coefficient forward
span nevertheless equals the full star algebra by the faithful-recurrence
theorem proved in V2.

For fixed $M$, these sources can approach the same scalar-multiplicity
source $(A,B,C)=(2I/3,-I/3,2I/3)$.  If $\mathfrak K$ denotes the complete
classically flagged six-edge instrument, then
\begin{equation}
 \boxed{\|\mathfrak K_{U,V}-\mathfrak K_{I,I}\|_\diamond
       \le2\sqrt{2/3}\,\delta.}
 \label{sel8:eq:source-close}
\end{equation}
Hence neither the exact normalization nor arbitrarily close source
probabilities and channels uniquely select an internal represented type.
\end{theorem}
\begin{proof}
Unitary $U,V$ give both $3D^*D+F^*F=4I$ and
$3DD^*+FF^*=4I$.  The latter, by the same contrast expansion, gives
$AA^*+BB^*+CC^*=I$.  Thus the edge channel is both trace preserving and
unital.  Since $D,F$ are nonzero, the six coefficients are independent.
The central support of the invertible $D$ is full and
$C^*(D,F)=C^*(U,V)=\mathcal R$, proving
\eqref{sel8:eq:arbitrary-internal}.

Relative to the scalar reference, the changes of $A,B$ have norm at most
$2\delta/3$, and that of $C$ at most $\delta/3$.  Their least singular
values are therefore at least $2(1-\delta)/3$, $(1-2\delta)/3$, and
$(2-\delta)/3$, respectively.  Each $K_e$ is an orthogonal three-block
operator with these blocks divided by $\sqrt2$, proving
\eqref{sel8:eq:invertible-floor}.  Every finite product is invertible;
its exact within-record refinement is a scalar multiple by the inherited
efficient-record theorem.

For the channel estimate put $\Delta D=U-I$, $\Delta F=-(V-I)$ and use
orthogonality of the external projections.  Direct expansion gives
\[
 \sum_e(\Delta K_e)^*\Delta K_e
 =I_3\otimes\left(\tfrac12\Delta D^*\Delta D+
                             \tfrac16\Delta F^*\Delta F\right)
 \preceq\tfrac23\delta^2 I_S.
\]
The Stinespring isometries with identical output and environment copies of
each classical flag therefore differ in norm by at most
$\sqrt{2/3}\,\delta$.  Expanding the difference of the two dilated output
maps gives at most twice that norm, also with any ancilla and after partial
trace, proving \eqref{sel8:eq:source-close}.
\end{proof}

\begin{corollary}[The spectral cost of approaching the scalar source]
\label{sel8:cor:small-source-gap}
Fix self-adjoint $H,L$ generating $\mathcal R$, set
$U_\epsilon=e^{i\epsilon H}$ and $V_\epsilon=e^{i\epsilon L}$ in
\eqref{sel8:eq:unitary-family}, and let $\gamma_\epsilon$ be the gap of
\eqref{sel8:eq:internal-calibration} on $\mathcal R'^\perp$.  Suppose this
complement is nonzero.  Define
\[
 \gamma_0^{HL}:=\lambda_{\min}\left(
   \operatorname{ad}_H^*\operatorname{ad}_H+
   \tfrac13\operatorname{ad}_L^*\operatorname{ad}_L
        \bigm|_{\mathcal R'^\perp}\right).
\]
This number is strictly positive, and for all sufficiently small nonzero $\epsilon$,
\begin{equation}
 \boxed{\left|\frac{\gamma_\epsilon}{\epsilon^2}-\gamma_0^{HL}\right|
 \le\frac{\epsilon^2}{12}
  \left(\|\operatorname{ad}_H\|^4+
                         \tfrac13\|\operatorname{ad}_L\|^4\right).}
 \label{sel8:eq:small-source-gap}
\end{equation}
In particular $\gamma_\epsilon=\epsilon^2\gamma_0^{HL}+O(\epsilon^4)$
while every finite original letter remains uniformly invertible.
\end{corollary}
\begin{proof}
Self-adjoint $H$ makes $\operatorname{ad}_H$ self-adjoint on the
Hilbert--Schmidt operator space.  The internal Laplacian is exactly
\[
 2I-2\cos(\epsilon\operatorname{ad}_H)
 +\tfrac13\bigl(2I-2\cos(\epsilon\operatorname{ad}_L)\bigr).
\]
The scalar remainder bound
$|2-2\cos t-t^2|\le t^4/12$ and spectral calculus bound the norm difference
after division by $\epsilon^2$.  Restriction to the fixed orthogonal
complement and the min--max principle give
\eqref{sel8:eq:small-source-gap}.  For small nonzero $\epsilon$ the
unitaries generate the same $\mathcal R$ by the interpolation argument,
so the kernel on that interval is precisely $\mathcal R'$.  Its limiting
rescaled gap is strictly positive because $H,L$ generate $\mathcal R$.
\end{proof}

\begin{example}[A fully explicit internal qubit with no pure original history]
\label{sel8:ex:rational-unitaries}
On $M=\C^2$ put
\[
 U=\frac1{101}\begin{pmatrix}99+20i&0\\0&99-20i\end{pmatrix},
 \qquad
 V=\frac1{101}\begin{pmatrix}99&20\\-20&99\end{pmatrix}.
\]
Both are unitary and have distance $2/\sqrt{101}<1/2$ from identity.
The distinct eigenvalues of $U$ force its commutant to be diagonal, and the
nonzero off-diagonal entries of $V$ make the common commutant scalar.
Thus $C^*(U,V)=M_2$.  The exact multiplicity determinants in
\eqref{sel8:eq:unitary-family} are
\[
 \det A=\frac{41000}{91809},\qquad
 \det B=\frac{10409}{91809},\qquad
 \det C=\frac{400}{909}.
\]
They are nonzero, so every $K_e$ and every finite original word has full
operator rank six.  Yet the source algebra is $M_6$, its original invariant
internal algebra is $I_3\otimes M_2$, and its stationary state is faithful.
The companion exact-arithmetic script checks these determinants and the
quadratic--cubic extraction.  This specified example is not a proposed
numerical value of the historical source panel.
\end{example}

\section{What has been discharged and which source calculation remains}
\label{sel8:status}

The continuation has removed one genuine assumption from the multiplicity
route.  The internal generators no longer have to be recovered after adding
a bank of external matrix controls.  The six original coefficients already
contain them as invariant elements, and the four quadratic--cubic operator
combinations reconstruct them explicitly.  The exact raw source algebra is
also classified: it is a full external matrix factor on the central support
of the off-diagonal contrast, with a precisely identified diagonal sector
on the complement.  This complement is computed by a finite positive
source-ideal Gram rather than discarded.

There are now two equivalent source data routes to the internal algebra.
One uses the common-frame quadratic and cubic operator combinations.  The
other uses the actual reduced multiplicity channel, whose positive residual
Choi support has dimension at most two.  The latter also preserves the
original calibrated internal commutator form and transfers faithful
stationarity in both directions.  It needs quantum output data, not a
favourably interpreted classical record.

The distinction from the pure-donor programme is important.  Theorem
\ref{sel8:thm:all-algebras} shows that an internal $M_3$ can occur in the
original invariant source algebra while every actual finite history is
invertible.  A global operator-rank-one donor is therefore a sufficient
construction used in V1--V7, not a necessary condition for algebraic colour.
Searching exclusively for such histories would miss legitimate internal
matrix algebras.  Conversely, the existence of an internal matrix algebra
does not select an executed invariant bridge, the physical type/private
identifications, or a determinant-bearing same-history tensor.  Formal
matrix units do not supply those additional occurrence statements.

\begin{center}
\small
\begin{tabular}{L{.36\textwidth}L{.56\textwidth}}
\toprule
Question & V8 result and boundary\\
\midrule
Are separately adjoined external controls needed to identify the internal
algebra of V7's multiplicity-valued source?
& No. The original invariant algebra equals $I_3\otimes C^*(A,B,C)$ by
quadratic--cubic extraction.\\[.3em]
Does the source necessarily generate all external controls on every sector?
& No. The exact exceptional diagonal sector is the central ideal complement
where the off-diagonal contrast vanishes.\\[.3em]
Can the required internal panel be made smaller than a complete reconstruction
of every edge coefficient?
& Yes. The reduced multiplicity channel supplies one rank-at-most-two positive
residual Choi support and the exact calibrated form.\\[.3em]
Do complete invariant quadratic words settle the internal type?
& No. An exactly calibrated pair has identical degree-two data and different
commutative versus noncommutative internal algebras; cubic data distinguish it.\\[.3em]
Are pure original histories necessary for an internal nonabelian algebra?
& No. Every represented finite internal algebra has an invertible,
faithfully recurrent six-edge realization in the same class.\\[.3em]
Is historical Standard-Model selection established?
& No. The actual multiplicity factor and reduced source row remain to be
identified, and the source-specific type/private, determinant and matter
conditions remain separate.\\
\bottomrule
\end{tabular}
\end{center}

\paragraph{The next nonduplicative calculation.}
On a verified external-triplet multiplicity factor, reconstruct the actual
reduced channel $\Psi_*$, retaining the system input/output identification
and the opportunity calibration.  Check the positive residual
\eqref{sel8:eq:rank-two-Choi}, determine its supported operator space, and
compute its invariant source algebra and calibrated commutant.  Alternatively
populate the four operator combinations in
\eqref{sel8:eq:four-moments}.  Neither route presupposes which internal
block sizes will appear.  If the actual factor $M$ or this output-sensitive
row has not been identified, the status is a missing source row, not failure
of a numerical example to be persuasive.

Only after the internal algebra is known should the physically named type
line, private plane and determinant incidence be tested on that same source.
The literal source definitions used in the audit still supply normalization
and independent coefficients rather than evaluated matrices $A,B,C$ or the
reduced channel.  The new theorems do not fill those entries by choosing one
of their realizations.  They do show that another search for a sharp input
Read is not required merely to determine the internal algebra of this
already specified source class.

\paragraph{Proof and preservation status.}
The V7 source prefix before its final document command is retained byte for
byte.  The accompanying exact-arithmetic checks verify the tetrahedral
contrasts, four extraction identities for nonnormal and singular matrices,
normalization, raw-algebra dimension examples, autonomous partial traces,
Choi support recovery, the degree-two indistinguishable pair, the calibrated
Pauli spectra, and the rational invertible realization.  V7's regression
suite is rerun.  The all-dimension algebra decomposition, finite extraction
bound, arbitrary-algebra realization and error estimates are proved in the
text, not inferred from sampled matrices.  No new historical operator row,
proof-assistant formalization, or independent referee verification is claimed.

\begin{thebibliography}{99}
\bibitem[16]{sel8-Watrous}
J.~Watrous, \emph{The Theory of Quantum Information}, Cambridge University
Press, 2018; \href{https://doi.org/10.1017/9781316848142}{doi:10.1017/9781316848142}.
Kraus/Choi representation and channel-norm estimates are standard background;
the normalized source-extraction and decomposition results are proved above.
\end{thebibliography}

% END OF SOURCE-SELECTION V8 DEVELOPMENT BLOCK.
% Append subsequent developments before the final document terminator.
% Preserve the V1--V8 bodies and state any necessary amendments explicitly.


\clearpage
\begin{center}
{\Large\bfseries Source-selection continuation V9}\\[.4em]
{\large Native conserved sectors and normalized source geometry}\\[.25em]
Represented-type codimensions, exact leakage costs, and a scalar colour test
\end{center}

\section{Go upstream of a chosen internal block structure}
\label{sel9:context}

V8 identifies the internal algebra of a verified covariant source with the
operator algebra of its actual reduced multiplicity channel.  It also proves
that the source normalization permits every finite represented algebra, even
with all original histories invertible.  The next issue is therefore not how
to manufacture another pure donor.  It is which source information can protect
a multi-block internal algebra instead of a single full matrix algebra.

This continuation answers three finite questions.  It computes the dimensions
of the represented-algebra strata inside the \emph{exactly normalized} source
family.  It derives the least normalized coefficient change needed to make
independently specified sectors invariant, directly from their measured
leakage operator.  Finally, on an already identified type--private carrier,
it reduces the remaining algebraic colour test to one positive population
transfer scalar, without a Kraus-frame reconstruction or a pure-history
assumption.  The last statement uses the inherited full-block bridge theorem;
it is its reduced-source specialization, not another proof of matrix-unit
generation presented as a new principle.

\paragraph{Audit and nonduplication.}
The V8 extraction, autonomous channel, calibration and arbitrary-algebra
realization are inputs.  The earlier rigidity programme already proves
analytic Howe genericity with a fixed protected algebra, and the older
SM-emergence programme already uses type-transfer graphs and positive
occurrence masses.  Neither is repeated as a new general theorem.  The new
geometric content is the exact represented-type dimension inside the
normalized two-coefficient source manifold, including the unrestricted
versus sector-conserving distinction.  The new optimization content is an
exact constrained source-distance formula, including singular compressed
supports and physical normalization.  The fixed-projection and transfer
identities below are proved briefly to specify precisely the data needed by
those two results.

The source audit uses the literal normalization and record factors of
\cite{Grand}, under the labels
\begin{center}
\small
\src{def:locked-opportunity-instrument}\\
\src{thm:GT76-locked-normalform}.
\end{center}
These explicitly leave the individual system coefficients undetermined.  Their existing tensor-algebra formula is
used in Section~\ref{sel9:actual-records}, not rediscovered.  A targeted search
of \cite{Grand,Attack,Rigidity} and V1--V8 located related source-innovation,
fixed-point and polar-completion results, but not the type-stratum count or
the exact nearest normalized sector-conserving source formula proved here.
This is a programme-level audit, not an exhaustive priority claim.

Stiefel descriptions of fixed-Kraus-rank channels and polar best approximation
are established tools; see \cite{sel9-ItenColbeck,sel9-Higham}.  The specialized
proofs below state the required finite hypotheses.  No distribution over
microscopic sources is postulated by using the word \emph{generic}.

\subsection{The inherited channel is the parameter space, not the target algebra}

Fix the physically identified multiplicity space $M\cong\C^m$, $m\ge2$, of
V8's external-triplet source class.  In this block only, set
\begin{equation}
 E_1=\frac{\sqrt3}{2}D,\qquad E_2=\frac12F,\qquad
 V_E=\begin{pmatrix}E_1\\E_2\end{pmatrix},\qquad V_E^*V_E=I_m.
 \label{sel9:eq:source-coordinates}
\end{equation}
The symbols $E_1,E_2$ are coordinates on V8's positive residual Choi packet;
they are not asserted to be newly resolved physical outcomes.  Retain
\begin{equation}
 \Xi_{E,*}(\rho)=\sum_{j=1}^2 E_j\rho E_j^*,\qquad
 \Psi_{E,*}=\tfrac13\id+\tfrac23\Xi_{E,*},\qquad
 \mathcal R_E=C^*(I,E_1,E_2).
 \label{sel9:eq:residual-channel}
\end{equation}
V8 proves that the original invariant internal algebra is
$I_W\otimes\mathcal R_E$.  Recover the original $A,B,C$ and $K_e$ by
\eqref{sel8:eq:contrast-calibration} and
\eqref{sel7:eq:matrix-family}.  Six-coefficient independence is exactly
$E_1\ne0$, $E_2\ne0$.  We first allow their closed normalized parameter
space
\begin{equation}
 \mathfrak S_m=\{(E_1,E_2)\in M_m(\C)^2:
                         E_1^*E_1+E_2^*E_2=I_m\}.
 \label{sel9:eq:Stiefel}
\end{equation}
All generic conclusions restrict to the open independent-source subset.
Source-distance minimizers may lie on its boundary; when this occurs, the
notes say so explicitly.

\begin{lemma}[Exact normalized source manifold]
\label{sel9:lem:manifold}
The space $\mathfrak S_m$ is the connected compact complex Stiefel manifold
of isometries $\C^m\to\C^{2m}$, of real dimension $3m^2$.  If a represented
unital algebra
\[
 \mathcal R_0=\bigoplus_{a=1}^k M_{n_a}(\C)\otimes I_{r_a},
 \qquad \sum_a n_ar_a=m,
\]
is held fixed, its normalized pairs form a connected smooth compact manifold
of real dimension $3\sum_a n_a^2$.
\end{lemma}
\begin{proof}
Stacking gives the first identification.  The unitary group on $\C^{2m}$
acts transitively on isometries with stabilizer $\mathrm U(m)$, so the real
dimension is $(2m)^2-m^2=3m^2$ and the space is connected and compact.
Alternatively the derivative of $V\mapsto V^*V$ is onto the Hermitian
matrices at every isometry, since $\delta V=VH/2$ maps to $H$.
Inside $\mathcal R_0$ the normalization is an independent Stiefel equation
on each $n_a\times n_a$ matrix block; the representation repetitions $r_a$
do not create coefficient parameters.  Taking the product proves the claim.
\end{proof}

\section{What is generic in the actual normalized source class}
\label{sel9:strata}

A represented algebra type is an unordered list
$\tau=((n_a,r_a))_{a=1}^k$ with $\sum_a n_ar_a=m$.  Different representation
multiplicities are not identified.  Let $\mathfrak S_\tau$ consist of pairs
whose generated unital $C^*$-algebra is unitarily conjugate to
$\mathcal R_\tau=\bigoplus_a M_{n_a}\otimes I_{r_a}$.

\begin{theorem}[Dimension of each represented internal-algebra stratum]
\label{sel9:thm:type-dimension}
Every $\mathfrak S_\tau$ is nonempty and is a smooth semialgebraic stratum
of real dimension
\begin{equation}
 \boxed{\dim_{\R}\mathfrak S_\tau
    =m^2+2\sum_a n_a^2-\sum_a r_a^2+k.}
 \label{sel9:eq:type-dimension}
\end{equation}
Consequently its codimension in the normalized source manifold is
\begin{equation}
 \boxed{\operatorname{codim}\mathfrak S_\tau
    =2m^2-2\sum_a n_a^2+\sum_a r_a^2-k.}
 \label{sel9:eq:type-codimension}
\end{equation}
On a fixed represented algebra $\mathcal R_\tau$, the pairs generating
that entire algebra, including separation of distinct central blocks, form
an open dense full-measure subset of its normalized-pair manifold.
Here full measure is with respect to smooth volume on the stated manifold,
not a law of historical source occurrence.
\end{theorem}
\begin{proof}
Put $a_\tau=\sum n_a^2$, $b_\tau=\sum r_a^2$.  The normalized pairs in
$\mathcal R_\tau$ have dimension $3a_\tau$ by
Lemma~\ref{sel9:lem:manifold}.  Generation of $\mathcal R_\tau$ is a finite
rank condition on star words: a unital word span either grows at the next
step or is already invariant under all generators.  Depth $a_\tau-1$
therefore suffices.  Full rank is open and is detected by nonvanishing of a
finite sum of squared minors.  A pair of unitaries generating
$\mathcal R_\tau$, divided by $\sqrt2$, is a normalized witness; such a pair
was constructed in Lemma~\ref{sel8:lem:two-unitaries}.  The rank detector is
real analytic and not identically zero on the connected product manifold.
The real-analytic identity principle, in local coordinate charts, implies
that its zero set has empty interior and zero smooth volume.  This proves
the last assertion and nonemptiness.  The usual measure-zero fact follows
locally by the first nonvanishing derivative and induction on the number of
real coordinates; no ambient Lebesgue measure on a constraint surface is used.

Let $N_\tau$ be the unitary normalizer of $\mathcal R_\tau$ in $\mathrm U(m)$.
Its identity component is the product of the unitary groups of
$\mathcal R_\tau$ and $\mathcal R_\tau'$, with their common central torus
identified.  Indeed, a continuous automorphism fixes the minimal central
projections and is inner on each full matrix block; removing that inner
unitary leaves a commutant unitary.  Thus
\[
 \dim N_\tau=a_\tau+b_\tau-k.
\]
Possible permutations of equivalent represented blocks are finite and do
not change this dimension.

Let $X_\tau$ be the open set of normalized generating pairs inside
$\mathcal R_\tau$.  The map
\[
 \mathrm U(m)\mathbin{\times}_{N_\tau}X_\tau
 \longrightarrow\mathfrak S_\tau,\qquad
 [U,(E_j)]\longmapsto (UE_jU^*)
\]
is bijective: equality of two pairs identifies the algebras they generate
and hence makes their coordinate change a normalizer element.  It is a
smooth identification.  To check the local assertion, select a nonzero
word minor at a pair.  Its selected words give a smoothly varying basis of
the generated algebra on this type stratum, hence a smooth map into the
unitary orbit of $\mathcal R_\tau$ in the Grassmannian.  A local section of
the compact-group orbit map lifts that algebra back to $\mathcal R_\tau$;
the resulting pair is in $X_\tau$.  Equivalently, an infinitesimal conjugacy
$Z^*=-Z$ with $[Z,E_j]\in\mathcal R_\tau$ for all $j$ also satisfies
$[Z,\mathcal R_\tau]\subseteq\mathcal R_\tau$, by the product and adjoint
rules, so precisely the normalizer tangent has been quotiented out.
The dimension is therefore
$m^2+3a_\tau-(a_\tau+b_\tau-k)$, giving
\eqref{sel9:eq:type-dimension}.

Finally, the fixed algebra, normalization, unitarity and finite rank
conditions are real polynomial conditions in real and imaginary matrix
entries.  Existentially quantifying the conjugating unitary gives a
semialgebraic set by the projection theorem already used in the rigidity
programme.  The local description above establishes smoothness of this
represented-type set, not a claim that its closure is smooth.
\end{proof}

\begin{theorem}[Generic source saturation and the codimension of every reduction]
\label{sel9:thm:generic-full}
For $m\ge2$, the set
\[
 \{E\in\mathfrak S_m:\mathcal R_E=M_m(\C)\}
\]
is open, dense and of full smooth volume.  Its complement has real dimension
\begin{equation}
 \boxed{3m^2-4(m-1).}
 \label{sel9:eq:reducible-dimension}
\end{equation}
On an open dense full-measure subset, the coefficient forward span
$\operatorname{alg}(I,E_1,E_2)$ also equals $M_m$, the residual Choi rank is
two, the source has a faithful
stationary state, and every original six-edge coefficient is invertible.
In particular, generic algebraic internal saturation does not imply the
existence of any rank-one original finite history.
\end{theorem}
\begin{proof}
A proper unital $C^*$-subalgebra has a nontrivial reducing projection.
For a fixed projection of rank $r$, the preserving normalized pairs have
dimension $3r^2+3(m-r)^2$.  The possible projections form a complex
Grassmannian of real dimension $2r(m-r)$.  Their incidence space therefore
has dimension
\[
 3m^2-4r(m-r)\le3m^2-4(m-1).
\]
The images for $1\le r<m$ are compact semialgebraic sets.  Their finite
union is exactly the reducible locus, so it is closed, has empty interior,
and has zero smooth volume.  The type
$M_{m-1}(\C)\oplus\C$ (two scalar blocks when $m=2$) has the stated maximal
dimension by Theorem~\ref{sel9:thm:type-dimension}.  This gives equality in
\eqref{sel9:eq:reducible-dimension}.

Forward generation is another finite word-rank condition, now without
adjoint letters, of depth at most $m^2-1$.  A generating unitary pair is a
witness, since each unitary inverse is a polynomial in that unitary on its
finite spectrum.  The same analytic-minor argument gives an open dense
full-measure forward-full locus.  Any stationary state of a finite
trace-preserving channel exists by Cesaro averaging; its support is forward
invariant, since the positive image has no mass outside that support.
Forward fullness makes that support all of $M$, so the stationary state is
faithful.  This is the source-support argument used in V2, not a new
recurrence criterion.

Linear independence of $E_1,E_2$, and invertibility of the original $A,B,C$,
are also analytic nonvanishing conditions.  V8's near-identity unitary
construction, with two unitaries generating $M_m$, witnesses all of them
simultaneously.  Their exceptional sets have measure zero.  Intersecting
these finitely many open dense sets proves the remaining assertions.
An invertible product cannot have rank one on $W\otimes M$, whose dimension
is $3m>1$.
\end{proof}

\begin{corollary}[The count descends to the reduced Choi data]
\label{sel9:cor:Choi-dimension}
Restrict to linearly independent two-coefficient pairs.  The map
$E\mapsto J_{\Xi_E}$ identifies exactly the free $\mathrm U(2)$ Kraus
mixing orbits.  The manifold of such channels has dimension $3m^2-4$.
Every represented-algebra stratum with nonempty rank-two part loses the
same four dimensions on this quotient; its codimension is still
\eqref{sel9:eq:type-codimension}.  The known affine identity addition
$\Psi=\id/3+2\Xi/3$ does not change these dimensions.
\end{corollary}
\begin{proof}
Two full-column Choi syntheses with the same positive product differ by a
unitary on their two coefficient coordinates.  Linear independence makes
this action free, and compactness makes it proper.  The quotient therefore
subtracts $\dim\mathrm U(2)=4$.  Mixing changes neither the coefficient
span nor its generated algebra.  The quotient construction is the standard
fixed-Kraus-rank manifold description \cite{sel9-ItenColbeck}; the present
codimensions concern its represented source-algebra strata.  The affine
addition is injective.  The rank-one scalar-algebra stratum is not included
in this rank-two count.
\end{proof}

\begin{example}[Exact low-dimensional counts and the structural target]
\label{sel9:ex:counts}
For a multiplicity qutrit the complete represented list is
\[
\begin{array}{c|r|r}
 \mathcal R_E\subset M_3&\dim\mathfrak S_\tau&\operatorname{codim}\mathfrak S_\tau\\\hline
 M_3&27&0\\
 M_2\oplus\C&19&8\\
 \C^3&15&12\\
 \C I_2\oplus\C&10&17\\
 \C I_3&3&24
\end{array}
\]
For the \emph{minimal faithful representation} of
$M_3\oplus M_2\oplus\C\oplus\C$ on $M=\C^7$, the corresponding source
stratum has dimension $79$ in an ambient manifold of dimension $147$;
its codimension is $68$.  Allowing all orientations of those sectors is
already included in that number.  On rank-two reduced Choi data the
dimensions are $75$ and $143$, with the same codimension.

This seven-dimensional multiplicity space is an explicitly stated
representation-theoretic comparison, not an identification with the
historical fourteen-dimensional source carrier or an experimental
fine-tuning probability.  Other faithful representations, or physically
restricted parameter families, have different counts.  In particular, once
native sector conservation is imposed, the same block algebra is generic
\emph{inside that restricted family}, as shown next.
\end{example}

\section{Native conservation, not a fitted block partition}
\label{sel9:native-sectors}

Let $P=(p_1,\ldots,p_k)$ be nonzero orthogonal projections on the
\emph{same physical multiplicity input and output}, with $\sum_a p_a=I$ and
$r_a=\rank p_a$.  The projections are supplied by independently identified
source observables or tested as candidate physical sectors.  They are not
inferred from the desired block sizes.  Define
\begin{align}
 t_{ba}&=\Tr[p_b\Xi_{E,*}(p_a)]
        =\sum_j\|p_bE_jp_a\|_{\HS}^2,\nonumber\\
 \ell_a&=\sum_{b\ne a}t_{ba},\qquad
 \ell_P=\sum_a\ell_a,\qquad
 L_P=\sum_a p_a\Xi_E^*(I-p_a)p_a.
 \label{sel9:eq:leakage}
\end{align}
Then $L_P\succeq0$ and $\Tr L_P=\ell_P$.  The actual normalized state for a
sector-uniform preparation is $p_a/r_a$, so its exit probability is
$\ell_a/r_a$, not $\ell_a$.  These experiments need their preparation and
terminal sector Read to be admitted; the matrix formula alone does not
supply those operations.

\begin{proposition}[Exact one-step conservation test]
\label{sel9:prop:conservation}
The following are equivalent:
\[
 \ell_P=0;\qquad [E_j,p_a]=0\ \text{for every }j,a;\qquad
 \Xi_E^*(p_a)=p_a\ \text{for every }a.
\]
Thus the normalized source pairs preserving $P$ form precisely
\begin{equation}
 \mathfrak S_P=\{E\in\mathfrak S_m:E_j\in\bigoplus_a\B(p_aM)\},
 \qquad\dim_{\R}\mathfrak S_P=3\sum_a r_a^2.
 \label{sel9:eq:fixed-sectors}
\end{equation}
No stationary-faithfulness assumption is required.
\end{proposition}
\begin{proof}
Zero trace of the positive sum in \eqref{sel9:eq:leakage} makes every
cross-block coefficient zero.  The source is block diagonal, and
normalization gives $\Xi_E^*(p_a)=p_a$.  Conversely, if a projection $p$ is
fixed, then
\[
 (I-p)\Xi_E^*(p)(I-p)=0,\qquad
 p\Xi_E^*(I-p)p=0.
\]
Their positive Kraus sums imply $pE_j(I-p)=0$ and $(I-p)E_jp=0$.
Apply this to every $p_a$.  The dimension is the product Stiefel count of
Lemma~\ref{sel9:lem:manifold}.  This is the sharp-projection specialization
of the already used multiplicative-domain argument; it is included to fix
the source-level input of the new geometry and distance theorems.
\end{proof}

\begin{theorem}[A native conserved observable selects the sector-constrained family]
\label{sel9:thm:native-observable}
Suppose a native Hermitian observable has spectral resolution
$Z=\sum_a z_ap_a$, with all $z_a$ distinct.  Its positive squared-change
operator is
\begin{equation}
 V_Z=\Xi_E^*(Z^2)-\Xi_E^*(Z)Z-Z\Xi_E^*(Z)+Z^2
     =\sum_j[Z,E_j]^*[Z,E_j]\succeq0.
 \label{sel9:eq:charge-variance}
\end{equation}
The scalar data obey
\begin{equation}
 \boxed{\Tr V_Z=\sum_{a,b}(z_b-z_a)^2t_{ba},\qquad
  \delta_Z^2\ell_P\le\Tr V_Z\le\Delta_Z^2\ell_P,}
 \label{sel9:eq:variance-leakage}
\end{equation}
where $\delta_Z=\min_{a\ne b}|z_a-z_b|>0$ and
$\Delta_Z=\max_{a,b}|z_a-z_b|$ for $k>1$.
Thus $V_Z=0$ is equivalent to exact preservation of every native sector.
It is also equivalent to simultaneous conservation of $Z$ and $Z^2$.
Within this exactly conserved family, the generic internal algebra is
\begin{equation}
 \boxed{\mathcal R_E=\bigoplus_a M_{r_a}(\C).}
 \label{sel9:eq:native-generic-algebra}
\end{equation}
In particular native ranks $(3,2,1,1)$ give the structural internal algebra
on an open dense full-measure subset of that conservation family, but the
ranks and conservation have not been derived by this statement.
For an actually admitted sector-uniform experiment with initial state $I/m$,
the mean squared change of the recorded native values is $\Tr V_Z/m$; the
trace in \eqref{sel9:eq:variance-leakage} retains that normalization explicitly.
\end{theorem}
\begin{proof}
Expand the Kraus commutator squares, using $\sum E_j^*E_j=I$.
On the $(b,a)$ block, $[Z,E_j]$ equals $(z_b-z_a)p_bE_jp_a$.
Hilbert--Schmidt orthogonality gives \eqref{sel9:eq:variance-leakage}.
Its zero set is precisely \eqref{sel9:eq:fixed-sectors}.  Commuting with
$Z$ makes both displayed moments conserved; conversely, if both moments
are conserved, the first formula for $V_Z$ is zero.  The generic-algebra
claim is Theorem~\ref{sel9:thm:type-dimension} for the fixed algebra
$\bigoplus_a\B(p_aM)$.  Its generating-pair rank condition includes central
separation, so two scalar sectors have not silently been identified.
\end{proof}

\begin{remark}[A conditional selection improvement and its boundary]
The conservation equation above is upstream of separately selecting the
matrix entries of every block.  It gives a native-family route to generic
fullness and central separation.  It is not a prediction of the spectrum or
degeneracies of $Z$, nor a reason why that observable is conserved.  Replacing
an unknown physical $Z$ by projectors fitted to the intended Standard-Model
answer would simply reinsert the original selection assumption.  Likewise,
with no such conserved observable, Theorem~\ref{sel9:thm:generic-full}
predicts generic \emph{simple} internal algebra in this mathematical family,
not the multi-block structural target.
\end{remark}

\begin{proposition}[A source-native positive statistic with a stable error bound]
\label{sel9:prop:stable-leakage}
Let $J_\Xi$ use the input-first vectorization convention of V4 and define
\[
 W_P=\sum_a p_a^{\mathsf T}\otimes(I-p_a).
\]
Then $W_P$ is an orthogonal projection and
\begin{equation}
 \boxed{\ell_P=\Tr(J_\Xi W_P).}
 \label{sel9:eq:Choi-leakage}
\end{equation}
Hence $\|\widehat J_\Xi-J_\Xi\|_1\le\epsilon$ implies
$|\widehat\ell_P-\ell_P|\le\epsilon$.  Moreover
\[
 \ell_P(\Psi_E)=\tfrac23\ell_P(\Xi_E),\qquad
 \ell_P((1-\vartheta)\id+\vartheta\Psi_E)
       =\tfrac{2\vartheta}{3}\ell_P(\Xi_E).
\]
Leakage can therefore be tested directly in the actual reduced source,
without subtracting or diagonalizing its Choi matrix.  A certified positive
lower bound excludes exact sector conservation.  An error interval meeting
zero does not certify it.
\end{proposition}
\begin{proof}
The terms of $W_P$ are orthogonal projections with mutually orthogonal
input supports.  The Choi pairing gives each $t_{ba}$, hence the identity.
Trace duality and $\|W_P\|\le1$ give the error bound.  The identity channel
has no cross-sector transition, proving both scale relations.  The known
calibration factors are retained even when only original event frequencies
are used.
\end{proof}

\section{The exact cost of imposing an unverified sector split}
\label{sel9:nearest}

A nonzero leakage should not be set to zero by an unbudgeted block fit.
There is an exact answer to how much a normalized source must change.
For this section $P$ is a fixed candidate partition; it need not already
be conserved.  Set
\begin{equation}
 B_j=\sum_a p_aE_jp_a,\qquad
 S_P=\sum_jB_j^*B_j=I-L_P.
 \label{sel9:eq:pinched-coefficients}
\end{equation}
The equality is blockwise: on $p_aM$, the retained same-sector mass plus the
outgoing mass equals $p_a$.  Therefore $0\preceq L_P\preceq I$.

\begin{theorem}[Nearest normalized sector-conserving source]
\label{sel9:thm:nearest}
The least squared coefficient Hilbert--Schmidt change to a normalized pair
preserving $P$ is
\begin{equation}
 \boxed{
 d_P(E)^2:=\min_{\widehat E\in\mathfrak S_P}
     \sum_{j=1}^2\|E_j-\widehat E_j\|_{\HS}^2
   =2\Tr\bigl(I-\sqrt{I-L_P}\bigr).}
 \label{sel9:eq:nearest-distance}
\end{equation}
If $S_P\succ0$, the unique minimizer is
\begin{equation}
 \boxed{\widehat E_j=B_jS_P^{-1/2}.}
 \label{sel9:eq:nearest-source}
\end{equation}
If $S_P$ is singular, blockwise isometric extensions of the polar partial
isometry of $\begin{psmallmatrix}B_1\\B_2\end{psmallmatrix}$ give minimizers
with the same distance.  No spectral-floor assumption is needed for the
minimum value.  In every case
\begin{equation}
 \boxed{\ell_P\le d_P(E)^2\le2\ell_P.}
 \label{sel9:eq:distance-leakage-bounds}
\end{equation}
The value is invariant under unitary Kraus mixing and under a simultaneous
unitary change of physical system and sector coordinates.  It depends only
on $\Xi_E$ and $P$, not on a chosen factorization of the residual Choi matrix.
\end{theorem}
\begin{proof}
Write $V=\begin{psmallmatrix}E_1\\E_2\end{psmallmatrix}$,
$B=\begin{psmallmatrix}B_1\\B_2\end{psmallmatrix}$ and similarly
$\widehat V$.  Both normalized columns have squared Hilbert--Schmidt norm
$m$, so
\[
 \|V-\widehat V\|_{\HS}^2
     =2m-2\operatorname{Re}\Tr(\widehat V^*V).
\]
A sector-preserving $\widehat V$ is Hilbert--Schmidt orthogonal to all
cross-sector blocks of $V$, hence the last trace is
$\Tr(\widehat V^*B)$.  On sector $a$, $B$ is a rectangular map
$p_aM\to p_aM\oplus p_aM$.  If its singular values are $s_i$, then for
any isometry $U_a$ on that sector,
\[
 \operatorname{Re}\Tr(U_a^*B_a)\le\sum_i s_i.
\]
To verify the inequality directly, take a singular-value decomposition and
sum $s_i$ times the real diagonal entries of a contraction, each at most one.
The polar isometry attains equality.  When $B_a$ is rank deficient, extend
its polar partial isometry on the missing input directions to orthonormal
vectors in the output complement; its output dimension $2r_a$ always permits
this.  If it has full column rank, equality forces the polar isometry, proving
uniqueness.  Summing over sectors gives the maximum
$\Tr\sqrt{B^*B}=\Tr\sqrt{I-L_P}$ and proves
\eqref{sel9:eq:nearest-distance}--\eqref{sel9:eq:nearest-source}.
This is a constrained use of the standard polar best-approximation method
\cite{sel9-Higham}; the source normalization and leakage formula are retained.

For $0\le t\le1$, one has
$t\le2(1-\sqrt{1-t})\le2t$.  Apply this to the eigenvalues of $L_P$.
Unitary mixing of the two coefficient coordinates mixes the $B_j$ in the
same way, leaving $S_P$ and $L_P$ unchanged.  Any two length-two Kraus
representations of the same rank-at-most-two channel are related by such a
unitary after zero padding.  Alternatively the formula for $L_P$ already
uses only $\Xi_E^*$ and $P$.  Simultaneous system conjugation preserves its
spectrum and the stated distance.
\end{proof}

\begin{remark}[What the minimizer means]
Equation~\eqref{sel9:eq:nearest-source} constructs a different normalized
source unless $\ell_P=0$.  It is not a state update performed by the old
source, an exact refinement of its recorded branches, or permission to erase
an observed leakage.  The minimizer may make one multiplicity contrast zero
or may identify two scalar blocks.  The theorem minimizes over the closed
normalization/conservation class.  If strict six-coefficient independence or
full block generation is also demanded, the same value is the \emph{infimum},
by density in that class, but need not be attained there.  This distinction
prevents a best fit from being mislabeled a source-selection proof.
\end{remark}

\begin{corollary}[Sharp geometric distance to any reducing decomposition]
\label{sel9:cor:distance-to-reducible}
The distance in the metric of \eqref{sel9:eq:nearest-distance} to the full
reducible locus in $\mathfrak S_m$ is
\begin{equation}
 \boxed{d_{\rm red}(E)^2=
   \min_{\substack{1\le r<m\\p=p^*=p^2,\ \rank p=r}}
        2\Tr\bigl(I-\sqrt{I-L_{(p,I-p)}}\bigr).}
 \label{sel9:eq:distance-reducible}
\end{equation}
The minimum is attained; it is positive exactly when $\mathcal R_E=M_m$.
For any fixed target rank list $(r_a)$, the analogous minimization over all
partitions with these ranks is the distance to its closed sector-preserving
class.  It equals the infimum distance to the exact, independently represented
full-block type $\bigoplus_aM_{r_a}$, though an exact-type nearest point may
not exist.
\end{corollary}
\begin{proof}
Every reducible finite unital $C^*$-algebra has a nontrivial reducing
projection, and every pair preserving such a projection is reducible.
The Grassmannians and normalized source spaces are compact; the displayed
objective is continuous, including at singular $S_P$.  The minimum is thus
attained and zero precisely on the closed reducible locus.
The rank-list statement follows from the same compactness and the density
of generating, centrally separated pairs in each fixed block algebra,
proved in Theorem~\ref{sel9:thm:type-dimension}.
\end{proof}

\begin{proposition}[Channel cost and the original opportunity scale]
\label{sel9:prop:channel-cost}
For a minimizing pair $\widehat E$, let
$\eta=\|V_E-V_{\widehat E}\|_{\op}\le d_P(E)$.  Then
\begin{equation}
 \|\Xi_E-\Xi_{\widehat E}\|_\diamond\le2\eta,
 \qquad
 \|\Psi_E-\Psi_{\widehat E}\|_\diamond\le\tfrac43\eta.
 \label{sel9:eq:channel-upper}
\end{equation}
For \emph{every} sector-conserving trace-preserving comparison $\widehat\Xi$,
regardless of its Kraus rank,
\begin{equation}
 \boxed{\|\Xi_E-\widehat\Xi\|_\diamond
           \ge2\max_a\frac{\ell_a}{r_a}.}
 \label{sel9:eq:channel-lower}
\end{equation}
If $\mathfrak K_E$ is the complete classically flagged six-edge instrument,
\begin{equation}
 \|\mathfrak K_E-\mathfrak K_{\widehat E}\|_\diamond
       \le2\sqrt{2/3}\,\eta.
 \label{sel9:eq:original-cost}
\end{equation}
With the same locked idle branch and acceptance probability $\vartheta$,
the last upper bound is multiplied by $\vartheta$; the reduced locked
channel lower bound is
$\frac{4\vartheta}{3}\max_a\ell_a/r_a$.
\end{proposition}
\begin{proof}
For normalized Stinespring columns the difference of the two output maps
expands into two terms, each bounded by the column difference, also with
any ancillary input.  This proves the first inequality, and the fixed
identity admixture proves the second.  On input $p_a/r_a$, the true channel
puts mass $\ell_a/r_a$ outside that sector whereas a preserving channel puts
none.  The output difference is Hermitian with trace zero, so its trace norm
is at least twice that probability difference.  Taking the best sector gives
\eqref{sel9:eq:channel-lower}.

The original contrast formulas give the exact operator identity
\[
 \sum_e(\Delta K_e)^*\Delta K_e
   =\tfrac23 I_3\otimes\sum_j(\Delta E_j)^*\Delta E_j.
\]
The flagged Stinespring construction therefore has column difference
$\sqrt{2/3}\,\eta$, proving \eqref{sel9:eq:original-cost}.
The old idle branch is identical in both comparisons.  The difference of
accepted flagged maps is thus multiplied by $\vartheta$, rather than
renormalized to an accepted experiment.  Partial trace and forgetting flags
recover the reduced locked map
$(1-2\vartheta/3)\id+(2\vartheta/3)\Xi_E$; applying the lower bound proves
the last statement.  None of these estimates equates a Frobenius nearest
source with a diamond-nearest channel.
\end{proof}

\section{A finite colour test and a quantitative obstruction to fitting sectors}
\label{sel9:certificates}

\begin{proposition}[A positive Howe gap bounds the cost of any sector fit]
\label{sel9:prop:Howe-cost}
Suppose $\mathcal R_E=M_m$ and the actually calibrated internal form of V8
has traceless gap $\gamma>0$.  Then for every partition $P$ as above,
\begin{equation}
 \boxed{\ell_P\ge\frac{3\gamma}{8}
       \left(m-\frac{\sum_a r_a^2}{m}\right),\qquad
 d_P(E)^2\ge\frac{3\gamma}{8}
       \left(m-\frac{\sum_a r_a^2}{m}\right).}
 \label{sel9:eq:gap-distance}
\end{equation}
In particular
\begin{equation}
 \boxed{d_{\rm red}(E)^2\ge\frac{3\gamma(m-1)}{4m}.}
 \label{sel9:eq:gap-any-reduction}
\end{equation}
Thus a certified full-algebra gap gives an explicit exclusion radius for all
proper sector decompositions.  It is not consistent to claim that gap and
simultaneously replace a positive cross-sector coupling by an exact zero
without paying the displayed source error.
\end{proposition}
\begin{proof}
In the residual coordinates \eqref{sel9:eq:source-coordinates},
\[
 q_{\rm int}(x)=\tfrac23\sum_j
  (\|[E_j,x]\|_{\HS}^2+\|[E_j^*,x]\|_{\HS}^2).
\]
For Hermitian $p_a$, the two norms agree.  Every cross block is counted twice
when the projection commutators are summed, so
\[
 \sum_aq_{\rm int}(p_a)=\tfrac83\ell_P.
\]
Subtract the scalar $r_aI/m$ from $p_a$, apply the gap, and use
$\|p_a-r_aI/m\|_{\HS}^2=r_a-r_a^2/m$.  This proves the leakage bound;
Theorem~\ref{sel9:thm:nearest} proves the distance bound.  For two sectors
of ranks $r,m-r$, minimize $r(m-r)$ over $1\le r<m$ to obtain
\eqref{sel9:eq:gap-any-reduction}.  The commutator coercivity step is the
inherited Howe mechanism; its consequence here is a distance to the
\emph{normalized} reducing-source locus.
\end{proof}

\begin{theorem}[Algebraic colour from one original population-transfer scalar]
\label{sel9:thm:scalar-colour}
Suppose the actual multiplicity carrier is $M=T\oplus H$, with
$\dim T=1$, $\dim H=2$, and their same-input/output projections $p_T,p_H$
are independently identified.  Assume the already occurring source algebra
contains
\[
 \mathcal B_0=\C p_T\oplus\B(H)\subseteq\mathcal R_E.
\]
No rank-one original history is assumed.  Define from the actual reduced
channel of V8
\begin{equation}
 c_{\rm col}=\Tr[p_H\Psi_{E,*}(p_T)]
              +\Tr[p_T\Psi_{E,*}(p_H)].
 \label{sel9:eq:colour-scalar}
\end{equation}
Then
\begin{equation}
 \boxed{c_{\rm col}>0\iff\mathcal R_E=M_3(\C),\qquad
 c_{\rm col}=0\iff\mathcal R_E=\mathcal B_0.}
 \label{sel9:eq:scalar-colour-test}
\end{equation}
For every one of the six original covariant coefficients, putting
$P_T=I_3\otimes p_T$, $P_H=I_3\otimes p_H$, one has
\begin{equation}
 \boxed{\|P_HK_eP_T\|_{\HS}^2+\|P_TK_eP_H\|_{\HS}^2
                 =\tfrac12c_{\rm col}.}
 \label{sel9:eq:all-original-corners}
\end{equation}
Thus a positive scalar already witnesses nonzero original source corners;
no long word search or ideal sharp-input refinement is required for this
algebraic conclusion.  The scalar is $2/3$ times the two-sector residual
leakage of \eqref{sel9:eq:leakage}.  Its physical locked probability scale
retains the additional factor $\vartheta$.
\end{theorem}
\begin{proof}
The scalar is the sum of the squared $T$--$H$ corners of $A,B,C$.  If it
vanishes, all three are block diagonal, so $\mathcal R_E\subseteq\mathcal B_0$;
the assumed reverse containment gives equality.  If it is positive, at
least one actual coefficient has a nonzero corner.  The inherited
necessary-and-sufficient colour-bridge theorem, applied to $\mathcal B_0$,
gives $M_3$.  This invokes that theorem rather than inferring an executed
matrix unit from a formal linear combination.  V8 identifies this generated
multiplicity algebra with the original invariant internal algebra.

Each $K_e$ is an orthogonal three-block external sum with coefficient blocks
$A,B,C$ divided by $\sqrt2$.  The external projections have rank one, so
their squared Hilbert--Schmidt corners add to precisely one half of the
reduced sum, proving \eqref{sel9:eq:all-original-corners}.  The scale
statements are Proposition~\ref{sel9:prop:stable-leakage}.
\end{proof}

\begin{remark}[Operational and data scope of the shortcut]
The two normalized population experiments in
\eqref{sel9:eq:colour-scalar} are a type-line input and the private input
$p_H/2$, with the latter measured probability multiplied by two.  If these
native input and output queries already occur, full reduced-channel
tomography is unnecessary for this specific bridge decision.  The
independent private $M_2$ and the same-carrier projectors must not be fitted
from the desired answer.  A positive raw corner need not be a globally
rank-one amplitude or an externally invariant individually executed word.
The theorem closes the internal \emph{algebra} on its stated entrance, not
determinant-bearing common-history occurrence or matter tensor selection.
V4's categorical coefficient contrast is still not automatically $p_T$.
\end{remark}

\begin{example}[An exact nearest reducing source and a sharp gap bound]
\label{sel9:ex:nearest-qubit}
Let $M=\C^2$, let $X,Y,Z$ be Pauli matrices, and put
\[
 E_1=Z/\sqrt2,\qquad
 E_2=R/\sqrt2,\qquad
 R=\begin{pmatrix}3/5&4/5\\-4/5&3/5\end{pmatrix}.
\]
This is a normalized, unital rank-two source with $\mathcal R_E=M_2$.
For the native coordinate projections $p_0,p_1$,
\[
 L_P=\frac8{25}I_2,\qquad \ell_P=\frac{16}{25},\qquad
 S_P=\frac{17}{25}I_2.
\]
The unique nearest preserving source is
\begin{equation}
 \widehat E_1=\frac5{\sqrt{34}}Z,\qquad
 \widehat E_2=\frac3{\sqrt{34}}I,
 \qquad d_P(E)^2=4-\frac{4\sqrt{17}}5.
 \label{sel9:eq:nearest-qubit}
\end{equation}
It generates exactly the diagonal algebra, not a collapsed scalar algebra.
For a general rank-one qubit projection
$p=(I+n_xX+n_yY+n_zZ)/2$, $|n|=1$, direct evaluation gives
\[
 \ell_{(p,I-p)}=1+\frac{16}{25}-n_z^2-\frac{16}{25}n_y^2.
\]
Because the channel is unital on a qubit, the two sector exit probabilities
are equal, so $L_{(p,I-p)}=\ell_{(p,I-p)}I_2/2$.
The expression is minimized at $n_z=\pm1$.  Therefore
\eqref{sel9:eq:nearest-qubit} is also the exact squared distance to
\emph{all} reducible normalized pairs, not just to this coordinate split.

On the normalized Pauli basis the calibrated internal spectrum is
\[
 \{128/75,\ 8/3,\ 328/75\}.
\]
Thus $\gamma=128/75$ and
$3\gamma/8=16/25=\ell_P$: the leakage lower bound in
\eqref{sel9:eq:gap-distance} is sharp.  The nearest-source formula gives the
stronger exact distance after normalization has been enforced.
All these are exact computations on a declared comparison source, not
historical numerical coefficients.
\end{example}

\begin{example}[Singular compressed support and an independence boundary]
\label{sel9:ex:singular-nearest}
With the same coordinate partition, take first
$E_1=E_2=X/\sqrt2$.  Then $B_1=B_2=0$, $L_P=I_2$, and
$d_P(E)^2=4=2\ell_P$.  Every preserving normalized pair has zero overlap
with the original stacked source and is a minimizer.  The singular case
cannot be treated by taking an inverse of $S_P$.

For a different normalized source set
\[
 E_1=\begin{pmatrix}3/5&0\\0&4/5\end{pmatrix},\qquad
 E_2=\begin{pmatrix}0&3/5\\4/5&0\end{pmatrix}.
\]
Then $L_P=\operatorname{diag}(16/25,9/25)$ and $d_P(E)^2=6/5$.
The unique nearest preserving pair is $(I,0)$, which makes the $F$ contrast
zero and the original six-edge list dependent.  Independent, centrally
separated preserving sources approach that point, but no such source attains
its minimum.  This verifies the infimum qualification following
Theorem~\ref{sel9:thm:nearest}.
\end{example}

\section{Evaluate the conserved sectors which the locked source actually supplies}
\label{sel9:actual-records}

The new conservation test can be run without knowing $K_e$ on one row that
\emph{is} explicitly present.  The supplied source has four orthogonal
sign--type record projections $Q_a$ and coefficient maps
\[
 L_\varnothing=\sqrt{1-\vartheta}\,I,\qquad
 L_{e,a}=\sqrt\vartheta\,K_e\otimes Q_a,
 \qquad \sum_eK_e=\sqrt2I.
\]
This is \cite{Grand}, \src{def:locked-opportunity-instrument}.  Write
$\mathcal D_R=C^*(Q_a)$ and $\mathcal B=C^*(K_e)$.

\begin{proposition}[The recorded sectors are conserved but are not the internal target]
\label{sel9:prop:actual-records}
The projectors $q_a=I_{H_E}\otimes Q_a$ have exactly zero leakage under
the full locked instrument and its unconditioned channel.  Their normalized
sector-enforcement cost is zero.  The existing source-algebra formula is
\begin{equation}
 \mathcal A_{\rm sys}=\mathcal B\otimes\mathcal D_R
             \cong\underbrace{\mathcal B\oplus\cdots\oplus\mathcal B}_{4}.
 \label{sel9:eq:record-repetitions}
\end{equation}
Thus those four actual sectors carry four copies of the same abstract edge
algebra.  Their existence does not supply independently identified internal
blocks of sizes $(3,2,1,1)$.  In particular the algebra in
\eqref{sel9:eq:record-repetitions} is not itself isomorphic to
$M_3\oplus M_2\oplus\C\oplus\C$.
\end{proposition}
\begin{proof}
Every displayed source coefficient commutes with every $q_a$, proving
conservation and zero cost.  Summing over $e$ at fixed $a$ gives
$\sqrt2 I\otimes Q_a$, and summing over $a$ gives $K_e\otimes I$.
These are exactly the generators of the tensor algebra, reproducing the
already proved formula in \cite{Grand}, \src{thm:GT76-locked-normalform}.
The $Q_a$ may have nontrivial ranks in their representation; their abstract
commutative algebra still has four minimal central projections.  Every simple
factor of $\mathcal B$ consequently occurs four times in the tensor algebra.
The indicated target has one $M_3$, one $M_2$, and two scalar factors, which
cannot have that repetition pattern.  Abstract finite $C^*$-algebra
classification proves nonisomorphism.
\end{proof}

\begin{remark}[What this evaluated row does not exclude]
This calculation does not rule out an internal Standard-Model algebra on an
appropriately reconstructed residual source.  The record factor is not to be
renamed an internal particle label merely because both descriptions have
four summands.  If the external action also permutes record sectors, its
invariant algebra must additionally be calculated; it is not silently assumed
to act trivially here.  The proposition concerns the explicitly supplied raw
record tensor algebra and the conservation it actually certifies.  The
internal multiplicity observables, their same-carrier locations, and their
transition data remain separate from this already evaluated row.
\end{remark}

\section{V9 outcome and the next source-dependent calculation}
\label{sel9:status}

The source-selection issue now has a quantitative split.  The full exactly
normalized two-coefficient class generically gives one simple internal matrix
algebra.  Multi-block source algebras occupy explicitly counted lower-dimensional
strata, even after their orientations and Kraus-coordinate ambiguities have
been allowed.  This is a statement about the displayed family, not a prior
probability assigned to the microscopic theory.  Once independently occurring
sector conservation is imposed, however, full block algebras and central
separation are generic within that native family.  Thus genericity cannot be
invoked until the actual allowed variations and conserved observables are
specified.

A proposed block projection can now be tested before a complete multiplicity
Choi reconstruction.  Its one-step leakage, or the squared change of one
native observable with that spectral decomposition, is a positive source
statistic.  A nonzero value gives both an operational exclusion and an exact
least coefficient-change cost for enforcing the partition.  Pinching followed
by a normalization is not an innocuous reinterpretation of the same source:
it is a precisely measured change unless that cost is zero.  The exact
minimization also handles singular supports instead of assuming their
nondegeneracy.

For an independently established type line and private $M_2$, the positive
algebraic colour decision is smaller still: the single scalar
\eqref{sel9:eq:colour-scalar} suffices.  V8's reduced channel need not be fully
reconstructed for that particular query.  This is a conditional data reduction,
not evidence that the physical type/private embedding or the required
population experiments have already been supplied.  No pure-history
determinant or matter tensor is inferred from the algebraic result.

\begin{center}
\small
\begin{tabular}{L{.37\textwidth}L{.55\textwidth}}
\toprule
Question & V9 result and exact remaining condition\\
\midrule
What does genericity select without protected internal sectors?
& A full $M_m$ on the fixed multiplicity carrier; the reducible locus has
codimension $4(m-1)$.\\[.3em]
How exceptional is a specified represented algebra?
& Its exact codimension is \eqref{sel9:eq:type-codimension}; the count is
relative to the normalized source family, not an experimental probability.\\[.3em]
Can a source-native conservation law replace a chosen block fit?
& Yes as a finite criterion: zero squared native-observable change forces
exact block preservation.  The observable and its degeneracies must occur
independently.\\[.3em]
Can measured leakage be normalized away for free?
& No. Equation~\eqref{sel9:eq:nearest-distance} is the exact minimum source
change, and \eqref{sel9:eq:channel-lower} gives a channel-level lower bound.\\[.3em]
Must algebraic colour require full tomography or a pure donor?
& Not once the native line/private carrier and its base $M_2$ are established:
one positive original population-transfer scalar suffices.\\[.3em]
Do the four supplied sign--type records constitute the desired internal split?
& No. Their source algebra repeats the same edge algebra four times; their
zero leakage does not identify heterogeneous internal blocks.\\[.3em]
Is historical SM selection closed?
& No. Actual multiplicity source rows, native internal-sector identifications,
and the separate determinant/matter incidence conditions remain unpopulated.\\
\bottomrule
\end{tabular}
\end{center}

\paragraph{The next noncircular input.}
Inspect the pre-matter, geometry-shorted source for an \emph{actual system}
observable or same-carrier projection panel, rather than a coefficient-space
contrast or a classical label count.  Record its spectral supports and evaluate
its one-step squared-change statistic under the actual reduced channel.  If
it is conserved, restrict the subsequent algebra audit to that physically
justified family.  If it is not, retain the nonzero leakage and its correction
cost rather than deleting cross-block entries to obtain the target.
When a type line and a faithful private base algebra are independently
available, test \eqref{sel9:eq:colour-scalar} directly.  If none of those
physical input/output identifications is supplied, the reduced Choi panel of
V8 remains the appropriate missing source row.  No additional construction
inside the unrestricted Stiefel family can turn that missing row into a
historical selection theorem.

\paragraph{Proof and cumulative preservation.}
The complete V8 source prefix before its terminal document command is
preserved byte for byte.  The accompanying exact-arithmetic checks verify
normalization, represented-type tangent dimensions on all five qutrit types,
the codimension counts, Kraus mixing invariance, the variance/leakage
identities, positive and singular polar minimizers, the explicit all-partition
qubit minimum, the calibrated gap, original six-edge error scaling, and the
scalar colour test.  The V8 regression suite is rerun.  Finite examples test
stated formulas; the all-dimension genericity, stratum and optimization
statements are proved above, not inferred from numerical sampling.  The
source examples are not fitted historical matrices, and no formal proof
assistant or independent referee verification is claimed.

\begin{thebibliography}{99}
\bibitem[17]{sel9-ItenColbeck}
R.~Iten and R.~Colbeck, \emph{Smooth manifold structure for extreme channels},
J.\ Math.\ Phys.\ \textbf{59} (2018), 012202;
\href{https://doi.org/10.1063/1.5019837}{doi:10.1063/1.5019837};
\href{https://arxiv.org/abs/1610.02513}{arXiv:1610.02513}.
The Stiefel/Kraus quotient is established background; the represented-algebra
strata of the normalized source are calculated here.
\bibitem[18]{sel9-Higham}
N.~J.~Higham, \emph{Computing the polar decomposition---with applications},
SIAM J.\ Sci.\ Stat.\ Comput.\ \textbf{7} (1986), 1160--1174;
\href{https://doi.org/10.1137/0907079}{doi:10.1137/0907079}.
Polar best approximation is standard; the sector-conserving normalized-source
minimum is proved explicitly above.
\end{thebibliography}

% END OF SOURCE-SELECTION V9 DEVELOPMENT BLOCK.
% Append subsequent work before the final document terminator.
% Preserve the V1--V9 bodies and state any necessary amendments explicitly.


\clearpage
\begin{center}
{\Large\bfseries Source-selection continuation V10}\\[.4em]
{\large Full tetrahedral covariance, capacity saturation, and native returns}\\[.25em]
A source-derived line--plane split and generic invariant-algebra saturation
\end{center}

\section{The upstream distinction: covariance is not conservation}
\label{sel10:context}

V9 concerns the normalized residual channel on \emph{one} external isotypic
factor $W\otimes M$.  Its unrestricted generic algebra is $M_m$, and an
additional sector-conservation condition is sufficient to retain a direct sum
inside that single multiplicity space.  That result must not be extrapolated
to a physical carrier with several inequivalent external representations.
On such a carrier, the invariant part of a full source algebra already has
several summands because of the external representation.  Their central
projections need not commute with the original, noninvariant source letters.

This continuation returns to the externally typed carrier used in the
structural construction rather than adding a fitted conserved observable.
It asks what the six original resolved coefficients, their coherent identity
sum, and their actual system covariance imply on that carrier.  Three results
are constructive.  The normalization bounds the trivial multiplicity by
three.  At saturation the actual input/output channel reconstructs a native
rank-one projection and its rank-two complement.  Quadratic invariant words
then furnish a finite return criterion, and an explicit six-letter source
proves generic saturation of the full invariant algebra on the entire
connected normalized covariance family.  No independent colour bridge,
private $M_2$ control, or weak-projector bank is inserted in that last result.

\paragraph{Audited inputs and nonduplication.}
The source is the identity-calibrated locked coefficient class of
\cite{Grand}, \src{def:locked-opportunity-instrument} and
\src{thm:GT76-locked-normalform}.  The carrier census, external/internal
distinction, and external twirl of V106--V107 of \cite{Attack} are inherited.
V7--V8 classify and extract multiplicity operators on a \emph{single}
external triplet; V9 counts the algebra strata on that single-factor source
space.  None of those statements supplies the normalized intertwiners between
all the external isotypes considered here.  The general covariant
Stinespring mechanism and finite representation decomposition are standard
background \cite{Etingof,sel10-Verdon}.  We prove their calibrated specialization
and all selection consequences used below.  The source-cyclicity and
subdirect-product steps are used as established mechanisms, not advertised
as new general theorems.

\paragraph{The physical scope cannot be shortened.}
The carrier below must be the acted-on finite Hilbert space, not merely a
coefficient-space contrast or a regular word-module multiplicity.  System
covariance, the coherent sum, and the full record policy must be verified on
that same space.  The larger symmetry obtained by retaining root reversal
requires an additional compatibility test; Section~\ref{sel10:reversal}
proves that this can obstruct the six-letter class.  The explicit positive
source is a benchmark in the stated class, not an evaluated historical Store.
In particular a positive geometry-source Schur quotient does not by itself
transport the six-coefficient CP normalization to that quotient.  Such an
application requires an actual normalized coefficient descent.  For a literal
retained system projection $p$, the familiar required trace-preservation test is
\[
 p\sum_eK_e^*(I-p)K_ep=0;
\]
otherwise the compressed bank is trace decreasing and cannot be inserted into
\eqref{sel10:eq:general-source} without changing its source.  The projection
must also preserve the declared system symmetry.  This is an application
condition, not a claim that the source Schur short has already passed it.

\section{The coherent identity fixes an equivariant Stinespring capacity}
\label{sel10:capacity}

Let $G$ be finite, let $U$ be its unitary action on a finite space $S$, and
let $\pi$ be a real permutation representation on a transitive set of $n$
resolved coefficient labels.  Write $u=n^{-1/2}\mathbf1$ and
$\mathcal E_0=u^\perp$.  Assume
\begin{equation}
 \sum_eK_e^*K_e=I_S,\qquad \sum_eK_e=cI_S,\qquad
 K_{ge}=U_gK_eU_g^*,\qquad 0<|c|^2<n.
 \label{sel10:eq:general-source}
\end{equation}
The covariance is that of the actual coherent coefficient frame.  When only
resolved CP covariance is initially given, the inherited phase-lock lemma
\ref{sel7:lem:phase-lock} applies after independence and the coherent identity
expansion have been verified.

\begin{theorem}[Calibrated equivariant contrast isometry]
\label{sel10:thm:contrast-isometry}
Put $a=|c|^2/n$.  Equations~\eqref{sel10:eq:general-source} are equivalent to
\begin{equation}
 V_K=\sum_eK_e\otimes|e\rangle
 =\frac{c}{\sqrt n}\,I_S\otimes|u\rangle+
             \sqrt{1-a}\,B,
 \qquad B^*B=I_S,
 \label{sel10:eq:column}
\end{equation}
where $B:S\to S\otimes\mathcal E_0$ is equivariant.  If
\[
 S=\bigoplus_\lambda V_\lambda\otimes\C^{m_\lambda},\qquad
 S\otimes\mathcal E_0
   =\bigoplus_\lambda V_\lambda\otimes\C^{n_\lambda},
\]
then such sources exist exactly when $n_\lambda\ge m_\lambda$ for every
occurring $\lambda$.  Their normalized coefficient space is a connected
compact real-analytic manifold
\begin{equation}
 \prod_{\lambda:m_\lambda>0}
        V_{m_\lambda}(\C^{n_\lambda}),\qquad
 \dim_{\R}=\sum_\lambda(2n_\lambda m_\lambda-m_\lambda^2).
 \label{sel10:eq:equivariant-Stiefel}
\end{equation}
Here $V_r(\C^s)$ denotes the complex Stiefel space of $r$-frames in $\C^s$.
The accepted system channel has the exact decomposition
\begin{equation}
 \Phi_* = a\,\id+(1-a)\,\Xi_*,\qquad
 \Xi_*(X)=\operatorname{Tr}_{\mathcal E_0}(BXB^*).
 \label{sel10:eq:identity-fraction}
\end{equation}
\end{theorem}
\begin{proof}
The component of $V_K$ along $u$ is
$n^{-1/2}\sum_eK_e=cI/\sqrt n$.  Orthogonality of $u$ and $\mathcal E_0$
and $V_K^*V_K=I$ give $B^*B=I$.  Covariance of the column means
$(U_g\otimes\pi_g)V_K=V_KU_g$; its two reducing environment components are
therefore equivariant.  Conversely the displayed column reconstructs every
identity in \eqref{sel10:eq:general-source} by taking its coefficients.
Schur's lemma makes an intertwiner $B$ equal to
$\bigoplus_\lambda I_{V_\lambda}\otimes b_\lambda$, with
$b_\lambda:\C^{m_\lambda}\to\C^{n_\lambda}$.
Isometry is precisely $b_\lambda^*b_\lambda=I$.  This proves existence,
the product description, and the dimension using
$\dim_\R V_r(\C^s)=2sr-r^2$.  Each factor is a connected compact homogeneous
unitary space.  Partial trace over the orthogonal environment components
removes their cross terms and proves \eqref{sel10:eq:identity-fraction}.
\end{proof}

We now use $G=S_4$, the six unoriented edge labels, $c=\sqrt2$, and the
actual external representation class
\begin{equation}
 S_r=(\mathbf1\otimes\C^r)\oplus(W\otimes\C^2)
                       \oplus V_2\oplus W^{\rm s},\qquad
 N=\dim S_r=r+11.
 \label{sel10:eq:carrier}
\end{equation}
Here $\dim W=3$, $\dim V_2=2$, and $W^{\rm s}=W\otimes\operatorname{sgn}$.
The degree-one categorical census supplies the same nontrivial multiplicities
and $r=1$; admitting a coherent private plane on the same acted-on carrier
would give $r=3$.  The theorem below does not assert that this admission has
occurred.  The six-edge contrast environment is
$\mathcal E_0=W\oplus V_2$, not a freely chosen auxiliary carrier.

\begin{theorem}[Exact tetrahedral multiplicity capacity]
\label{sel10:thm:capacity}
On \eqref{sel10:eq:carrier}, the output multiplicities in
$S_r\otimes(W\oplus V_2)$, in the order
$(\mathbf1,W,V_2,W^{\rm s},\operatorname{sgn})$, are
\begin{equation}
 (n_{\mathbf1},n_W,n_{V_2},n_{W^{\rm s}},n_{\operatorname{sgn}})
                 =(3,r+7,r+4,7,2).
 \label{sel10:eq:multiplicities}
\end{equation}
Consequently the exactly normalized covariant six-letter class exists if
and only if $0\le r\le3$.  Its real dimension is
\begin{equation}
 \boxed{44+12r-r^2.}
 \label{sel10:eq:source-dimension}
\end{equation}
In particular $r=1,2,3$ give dimensions $55,64,71$.  At $r=3$, every
coefficient list in this class has six linearly independent coefficients.
\end{theorem}
\begin{proof}
The needed tensor decompositions follow from the $S_4$ character products:
\begin{align*}
 W\otimes W&=\mathbf1\oplus V_2\oplus W\oplus W^{\rm s},&
 W\otimes V_2&=W\oplus W^{\rm s},\\
 W^{\rm s}\otimes W&=\operatorname{sgn}\oplus V_2\oplus W\oplus W^{\rm s},&
 V_2\otimes V_2&=\mathbf1\oplus\operatorname{sgn}\oplus V_2,\\
 W^{\rm s}\otimes V_2&=W\oplus W^{\rm s}.&&
\end{align*}
Adding the contributions with their actual input multiplicities gives
\eqref{sel10:eq:multiplicities}.  The only restrictive isometry inequality
is $r\le3$.  Substitution into \eqref{sel10:eq:equivariant-Stiefel} gives
$(6r-r^2)+(4r+24)+(2r+7)+13$.
For independence at $r=3$, the trivial-input block of $B$ is a unitary
$\C^3\to\C^2\oplus\C$.  Its projections onto both target summands are
surjective, so both $W$ and $V_2$ environment coefficient subspaces are
nonzero.  The coefficient map from $\mathbf1\oplus W\oplus V_2$ is
nonzero on each inequivalent irreducible summand and therefore injective.
Its mean component is already nonzero by the coherent identity.
\end{proof}

\begin{remark}[Capacity is not occurrence or minimization]
The source identities exclude a fourth trivial copy while the other
multiplicities are held fixed.  They do not choose between the existing
classes $r=1,2,3$, and maximizing the capacity is not the same as minimizing
physical source innovation.  The dimension formula also assumes no additional
reversal or superselection condition has silently reduced the allowed class.
\end{remark}

\section{At saturation the source reconstructs a native line and plane}
\label{sel10:split}

Let $p_0,p_W,p_2,p_{\rm s}$ be the external isotypic projections on $S_r$.
They identify system subspaces; at this stage they are not asserted to be
executed interventions or elements of $\mathcal A_K=C^*(I,K_e:e)$.
Choose real orthonormal bases of $W,V_2$ and write the coefficients of the
contrast isometry as $L_{\lambda a}$, where $\lambda=W,2$ and
$1\le a\le d_\lambda$, with $d_W=3$, $d_2=2$.
The $L_{\lambda a}$ are fixed linear combinations of the six original
coefficients, not newly resolved physical outcomes.

\begin{theorem}[Native transition effects and an exact line--plane split]
\label{sel10:thm:native-split}
There are maps $T_W:\C^r\to\C^2$ and $T_2:\C^r\to\C$ such that
\begin{equation}
 L_{\lambda a}p_0
       =\frac{|a\rangle\otimes T_\lambda}{\sqrt{d_\lambda}},
 \qquad T_W^*T_W+T_2^*T_2=I_r.
 \label{sel10:eq:entries}
\end{equation}
In particular, for every operator $\rho$ supported in $p_0S_r$,
\begin{equation}
 \Phi_*(\rho)=\frac13\rho\ \oplus\
 \frac23\left(\frac{I_W}{3}\otimes T_W\rho T_W^*\right)
 \ \oplus\ \frac23\left(\frac{I_{V_2}}2\otimes T_2\rho T_2^*\right)
 \ \oplus\ 0.
 \label{sel10:eq:trivial-channel}
\end{equation}
At $r=3$ the stacked map $T=(T_W,T_2)^{\mathsf T}$ is unitary.  Thus
\begin{equation}
 \boxed{P_H:=T_W^*T_W,\qquad P_T:=T_2^*T_2,\qquad
 P_HP_T=0,\quad P_H+P_T=I_3,\quad
 \rank P_H=2,\quad\rank P_T=1.}
 \label{sel10:eq:native-projectors}
\end{equation}
These projections are recovered from the original accepted channel by
\begin{equation}
 \boxed{P_H=\frac32p_0\Phi^*(p_W)p_0,\qquad
        P_T=\frac32p_0\Phi^*(p_2)p_0,}
 \label{sel10:eq:native-effects}
\end{equation}
where the right sides are restricted to $p_0S_3$.
\end{theorem}
\begin{proof}
A trivial input vector can map equivariantly into $S_r\otimes\mathcal E_0$
only through the unique invariant vectors of $W\otimes W$ and
$V_2\otimes V_2$.  In real orthonormal bases these are
$d_\lambda^{-1/2}\sum_a|a\rangle\otimes|a\rangle$.
The remaining coefficients are exactly $T_\lambda$, giving
\eqref{sel10:eq:entries}; isometry gives their summed Gram identity.
Tracing out the two orthogonal environment summands yields
\eqref{sel10:eq:trivial-channel}, including the identity component.
For $r=3$, an isometry between two three-dimensional spaces is unitary;
its two target coordinate projections pull back to the orthogonal
rank-two and rank-one projections displayed above.  Computing the probability
of each output isotypic sector for arbitrary input $\rho$ gives the two
Heisenberg effects.
\end{proof}

\begin{corollary}[Unavoidable departure from the trivial external sector]
\label{sel10:cor:leakage}
For every normalized input supported in $p_0S_r$, $r>0$,
\begin{equation}
 \boxed{\operatorname{Tr}[(I-p_0)\Phi_*(\rho)]=\frac23.}
 \label{sel10:eq:forced-leakage}
\end{equation}
Hence the six accepted source coefficients cannot conserve $p_0$.
For the original locked opportunity with idle probability $1-\vartheta$,
the corresponding departure probability is $2\vartheta/3$, not $2/3$.
\end{corollary}
\begin{proof}
Take the trace of the two nontrivial output blocks in
\eqref{sel10:eq:trivial-channel} and use $T^*T=I_r$.
The idle branch does not leave $p_0$.
\end{proof}

\begin{remark}[Which conclusion this improves]
The native $P_T$ is a projection on an \emph{acted-on system} input, recovered
from output-sensitive transition effects.  It is not the root-uniform
categorical coefficient ray whose V4 pullback had full input rank.  Neither
is its equality with the manuscript's geometrically named type line automatic.
That equality, and identification of $P_H$ with the supplied private states,
remain same-carrier source comparisons.  Nor does the existence of the
projections imply that their sharp measurement is physically available.
\end{remark}

\begin{proposition}[The source already carries a conditional quantum transfer]
\label{sel10:prop:transfer}
On $P_H\C^3$, the map $T_W$ is unitary onto the actual $W$ multiplicity
plane.  The $p_W$ output of one accepted step is a state-independent
success branch of weight $2/3$, with output
$(I_W/3)\otimes T_W\rho T_W^*$.  Its normalized multiplicity output
retains the entire private-plane density matrix, including entanglement with
an untouched auxiliary system.  The complementary native line goes to the
$V_2$ multiplicity line with the same success weight.
\end{proposition}
\begin{proof}
Restrict \eqref{sel10:eq:trivial-channel} to each native input support.
The restriction of $T_\lambda$ is an isometry onto its target, so partial
trace over the external factor followed by its inverse gives a mathematical
left inverse.  The formula tensors with the identity on any auxiliary
system.  Physical implementation of that inverse or of the output-sector
Read is not part of the assertion.
\end{proof}

\section{Quadratic source returns replace independently loaded internal controls}
\label{sel10:returns}

Set $r=3$.  The reverse-direction \emph{blocks of the same forward
coefficients} have the form
\begin{equation}
 p_0 L_{\lambda a}p_\lambda
   =S_\lambda(\langle a|\otimes I),
 \qquad S_W:\C^2\to\C^3,\quad S_2:\C\to\C^3.
 \label{sel10:eq:exits}
\end{equation}
This is a representation-theoretic form for actually present forward
coefficients, not an assumption that an adjoint of an entry is executable.
The maps $S_\lambda$ may initially be zero or rank deficient.
Write
\begin{equation}
 \widehat S=(\sqrt3S_W,\sqrt2S_2),\qquad R=\widehat S T.
 \label{sel10:eq:return-matrix}
\end{equation}

\begin{proposition}[Faithful stationarity forces a nonsingular native return]
\label{sel10:prop:recurrent-return}
If the accepted channel has a faithful stationary state, then $\widehat S$
and $R$ are invertible, and each $S_\lambda$ is injective.  More precisely,
twirl a faithful stationary state to write it as
\[
 \sigma=\sigma_0\oplus(I_W/3\otimes\sigma_W)
               \oplus(I_{V_2}/2\otimes\sigma_2)
               \oplus(I_{W^{\rm s}}/3\otimes\sigma_{\rm s}),
\]
where every displayed multiplicity density is positive definite.  Then
\begin{equation}
 \boxed{R\,\mathcal B_0R^*=\sigma_0,\qquad
 \mathcal B_0=T^*\diag(\sigma_W/3,\sigma_2/2)T\succ0,}
 \label{sel10:eq:stationary-return}
\end{equation}
and
$s_{\min}(R)^2\ge\lambda_{\min}(\sigma_0)/\lambda_{\max}(\mathcal B_0)$.
If the channel is unital, then $\widehat S$ and $R$ are unitary.
\end{proposition}
\begin{proof}
Covariance preserves the group-twirl of a stationary state and twirling
preserves strict positivity.  Since
$\Phi_*=(\id+2\Xi_*)/3$, stationarity for $\Phi_*$ is equivalent to
stationarity for $\Xi_*$.  Only inputs of type $W$ or $V_2$ can reach the
trivial output through its contrast coefficients.  Using
\eqref{sel10:eq:exits}, the trivial output stationarity equation is
$\sigma_0=S_W\sigma_WS_W^*+S_2\sigma_2S_2^*$.
This is \eqref{sel10:eq:stationary-return}.  Its right side is faithful, so
the square matrix $\widehat S$ is surjective and invertible.  The singular
value bound follows from
$\sigma_0\preceq\lambda_{\max}(\mathcal B_0)RR^*$.
For unitality take the unnormalized stationary operator $I_{S_3}$; its
multiplicity blocks in this convention are $I_3,3I_2,2,3$.
Then $\widehat S\widehat S^*=I_3$, proving the unitary statement.
\end{proof}

Let $\Pi_W,\Pi_2$ be the real orthogonal projections of the six-label
permutation space onto its two contrast irreducibles.  The following
operators are fixed, basis-independent combinations of the original
coefficient words:
\begin{align}
 \mathsf Q_\lambda
 &=\sum_aL_{\lambda a}^*L_{\lambda a}
   =\frac32\sum_{e,f}(\Pi_\lambda)_{ef}K_e^*K_f,\nonumber\\
 \mathsf H_\lambda
 &=\sum_aL_{\lambda a}L_{\lambda a}
   =\frac32\sum_{e,f}(\Pi_\lambda)_{ef}K_eK_f,
 \qquad \mathsf Q=\mathsf Q_2,\quad
 \mathsf H=\mathsf H_W+\mathsf H_2.
 \label{sel10:eq:quadratic-invariants}
\end{align}
The real structure in these formulas is inherited from the actual edge
permutation representation.  No arbitrary complex Kraus rephasing is made.

\begin{theorem}[Source-contained return operators]
\label{sel10:thm:quadratic-returns}
All operators in \eqref{sel10:eq:quadratic-invariants} belong to
\[
 \mathcal I_K:=\mathcal A_K\cap U(S_4)'
 \ \subseteq\ U(S_4)'=M_3\oplus(I_W\otimes M_2)\oplus\C\oplus\C.
\]
Their restrictions to the trivial external sector satisfy
\begin{equation}
 \boxed{\mathsf Q|_{p_0}=P_T,\qquad
 \mathsf Q_W|_{p_0}=P_H,\qquad
 \mathsf H_\lambda|_{p_0}=\sqrt{d_\lambda}S_\lambda T_\lambda,
 \qquad \mathsf H|_{p_0}=R.}
 \label{sel10:eq:return-restrictions}
\end{equation}
Thus the native splitter and the return matrix are compressed quadratic
invariant words from the original source, without adjoining external controls.
\end{theorem}
\begin{proof}
The contrast coefficients transform in real orthogonal irreducible
representations.  Contracting two equal real representation indices therefore
makes both their sesquilinear and bilinear sums invariant.  They are
polynomials in the original bank, hence in $\mathcal A_K$.
The first restrictions follow from \eqref{sel10:eq:entries}.
For the bilinear sum, each path from the trivial sector through a contrast
coefficient goes to its matching type $\lambda$, and the next coefficient
returns by \eqref{sel10:eq:exits}.  Summing over $a$ gives
$d_\lambda S_\lambda T_\lambda/\sqrt{d_\lambda}$.
Adding the two components proves the return formula.
\end{proof}

Define the finite positive return-cyclicity Gram on the \emph{actual native
trivial sector} by
\begin{equation}
 \Omega_R=\sum_{|w|\le2,\ w\in\{R,R^*\}^*}wP_Tw^*.
 \label{sel10:eq:return-cyclicity}
\end{equation}
The sum includes the empty word.  The word star in the indexing set means
all words in the two indicated letters, not an additional physical grammar.

\begin{theorem}[Native-return saturation of the full internal algebra]
\label{sel10:thm:return-saturation}
Suppose $R$ is invertible and $\Omega_R\succ0$.  Invertibility is automatic
under the faithful-stationarity condition of
Proposition~\ref{sel10:prop:recurrent-return}.  Then
\begin{equation}
 \boxed{\mathcal I_K=U(S_4)'
 \cong M_3(\C)\oplus M_2(\C)\oplus\C\oplus\C.}
 \label{sel10:eq:full-internal}
\end{equation}
All four central supports and the full private and weak matrix blocks belong
to the algebra generated by the original coefficients.  No independently
loaded private $M_2$, weak router, colour bridge, or central separator is
required for this conclusion.
\end{theorem}
\begin{proof}
The range of $\Omega_R$ is the cyclic space generated from the native line
by $R,R^*$.  In dimension three, words of length two suffice: any nonflat
extension increases the dimension, while a flat extension is invariant under
both generators and stays flat.  On a full cyclic space, the operators
$wP_Tv^*$ span all matrix units between a basis of cyclic vectors.  Hence
$C^*(P_T,R)=M_3$.
By Theorem~\ref{sel10:thm:quadratic-returns}, the trivial-sector image of
$\mathcal I_K$ is therefore $M_3$.

The other external multiplicity images act on spaces of dimensions $2,1,1$.
They cannot support a nonzero unital representation of this simple
$M_3$ summand.  The finite subdirect-product argument of
Theorem~\ref{sel1:thm:subdirect} therefore separates the donor summand:
$p_0\in\mathcal I_K$ and $p_0\mathcal I_Kp_0=M_3$.
For any $x\in M_3$ use this actual algebra element to form
\[
 \mathsf Z_\lambda(x)
       =\sum_a L_{\lambda a}^*(p_0xp_0)L_{\lambda a}
       \in\mathcal I_K.
\]
The operator is supported only on the input sector $\lambda$ and equals
$I_{V_\lambda}\otimes S_\lambda^*xS_\lambda$ there.
Invertibility of $R=\widehat ST$ makes each $S_\lambda$ injective.
Consequently $S_W^*M_3S_W=M_2$ and $S_2^*M_3S_2=\C$.
This puts the entire separately supported $W$ block and the $V_2$ scalar
block in $\mathcal I_K$.  Subtracting their identities and $p_0$ from the
full identity gives $p_{\rm s}$, whose multiplicity is one.  These are all
summands of $U(S_4)'$, proving equality and central separation.
\end{proof}

\begin{remark}[A failed return panel is not a general no-go]
The condition $\Omega_R\succ0$ is necessary and sufficient for saturation by
$P_T,R$ on the trivial sector.  If it fails, other invariant source words
can still enlarge the donor algebra.  The theorem gives a sufficient
finite source test for the full internal result, not an all-word exclusion
from a singular $\Omega_R$ alone.  The recovered algebra is not the raw
source commutant, a noiseless algebra, or the group preserving every source
channel.  Determinant and matter-incidence preservation remain later tests.
\end{remark}

\begin{corollary}[The native return exposes original two-letter coupling]
\label{sel10:cor:two-letter}
For the native projections on $p_0S_3$, embedded in the full system,
\begin{align}
 \sum_{e,f}\|P_HK_fK_eP_T\|_{\HS}^2
   &=\frac29\|P_HRP_T\|_{\HS}^2,\nonumber\\
 \sum_{e,f}\|P_TK_fK_eP_H\|_{\HS}^2
   &=\frac4{27}\|P_TRP_H\|_{\HS}^2.
 \label{sel10:eq:two-letter-mass}
\end{align}
Thus a positive return-cyclicity Gram forces an actual original history of
length two with a nonzero native line--plane corner.  If $R$ is unitary,
the sum of the two masses is
$\frac{10}{27}(1-|\langle t,Rt\rangle|^2)$ for a unit vector $t$ on the
native line.  At the cyclic witness below it is exactly $10/27$.
For two locked opportunities these accepted-history masses are multiplied
by $\vartheta^2$.
\end{corollary}
\begin{proof}
The trivial-input/output compression of a two-letter product is
\[
 p_0K_fK_ep_0=\frac1{18}I_3+
 \frac23\sum_{\lambda=W,2}
 \frac{(\Pi_\lambda)_{fe}}{d_\lambda}
                    \sqrt{d_\lambda}S_\lambda T_\lambda.
\]
At saturation $\sqrt{d_\lambda}S_\lambda T_\lambda=RP_\lambda$.
The identity has no off-diagonal native corner.  Orthogonality of the
real environment projectors gives
$\sum_{e,f}(\Pi_\lambda)_{ef}(\Pi_\mu)_{ef}
=\delta_{\lambda\mu}d_\lambda$.  Squaring the corners and summing proves
\eqref{sel10:eq:two-letter-mass}.  If both right sides vanish, the native
line reduces $R$ and cannot be cyclic.  A unitary has equal squared
cross-block Hilbert--Schmidt masses across a projection and its complement;
for the rank-one projection that value is
$1-|\langle t,Rt\rangle|^2$.  The original accepted amplitude per letter
is $\sqrt\vartheta K_e$, which gives the last scaling.  This is a source
corner witness; preparing or reading its endpoints is a separate admission
question.
\end{proof}

\section{An exact covariant witness on the same fourteen-dimensional carrier}
\label{sel10:witness}

The following family proves that the native-return condition is nonempty.
It is an explicitly declared source realization, not a numerical assignment
to the historical writer.  In the pair-contrast basis of the external
triplet put
\begin{equation}
 Z_1=\frac{\diag(1,-1,0)}{\sqrt2},\quad
 Z_2=\frac{\diag(1,1,-2)}{\sqrt6},\quad
 X=\begin{pmatrix}0&1\\1&0\end{pmatrix},\quad
 Z=\diag(1,-1).
 \label{sel10:eq:witness-basis}
\end{equation}
The tetrahedral triplet consists of signed permutation matrices in this
basis.  Their conjugation action on $\Span\{Z_1,Z_2\}$ is $V_2$.
The pair $(X,Z)/\sqrt2$ transforms by that same $V_2$ representation under
conjugation on the $V_2$ carrier.  These identities can also be checked on
the two permutations $(12)$ and $(1234)$, which generate $S_4$.

Use the block order $(\C^3,W\otimes\C^2,V_2,W^{\rm s})$ and put
\[
 T_W=\begin{pmatrix}1&0&0\\0&1&0\end{pmatrix},\qquad
 T_2=\begin{pmatrix}0&0&1\end{pmatrix}.
\]
Choose any unitary $R\in U(3)$ and define
$S_W=RT_W^*/\sqrt3$, $S_2=RT_2^*/\sqrt2$.
For each triplet index $i=1,2,3$, define $L_{Wi}$ by its only nonzero blocks
\begin{equation}
 (L_{Wi})_{W,0}=|i\rangle\otimes T_W/\sqrt3,
 \qquad
 (L_{Wi})_{0,W}=S_W(\langle i|\otimes I_2).
 \label{sel10:eq:witness-W}
\end{equation}
For $a=1,2$, the nonzero blocks of $L_{2a}$ are
\begin{equation}
 \begin{gathered}
 (L_{2a})_{2,0}=|a\rangle T_2/\sqrt2,
 \qquad (L_{2a})_{0,2}=S_2\langle a|,\\
 (L_{2a})_{W,W}=Z_a\otimes I_2,\qquad
 (L_{2a})_{2,2}=\begin{cases}X/2,&a=1,\\Z/2,&a=2,\end{cases}\qquad
 (L_{2a})_{{\rm s},{\rm s}}=\sqrt{3/2}\,Z_a.
 \end{gathered}
 \label{sel10:eq:witness-2}
\end{equation}
All unspecified blocks are zero.  The five symbols describe the contrast
coordinates of six actual letters.  They are not added outcome labels.

\begin{proposition}[Complete source realization for every unitary return]
\label{sel10:prop:witness}
The coefficients in \eqref{sel10:eq:witness-W}--\eqref{sel10:eq:witness-2}
are equivariant and satisfy
\begin{equation}
 \sum_{\lambda,a}L_{\lambda a}^*L_{\lambda a}=I_{14},\qquad
 \sum_{\lambda,a}L_{\lambda a}L_{\lambda a}^*=I_{14}.
 \label{sel10:eq:witness-normalization}
\end{equation}
Writing the six edges as the complementary pairs $(k,+),(k,-)$, their
original coefficients are
\begin{equation}
 \boxed{K_{k,\pm}=\frac{I_{14}}{3\sqrt2}
       +\frac1{\sqrt3}\left(\pm L_{Wk}
                      +\sum_{a=1}^2(Z_a)_{kk}L_{2a}\right).}
 \label{sel10:eq:witness-six}
\end{equation}
They are independent, obey the original coherent and positive normalizations,
and have faithful stationary state $I_{14}/14$.  Their native return matrix
is exactly the chosen $R$.  The accepted instrument is the normalized
six-outcome instrument with these coefficients, not a postselected subset.
\end{proposition}
\begin{proof}
The column families in \eqref{sel10:eq:witness-W} transform as $W$ because
both their input/output contraction and the external representation are
real.  The stated transformation of $Z_a$ and $(X,Z)$ proves the $V_2$
covariance of \eqref{sel10:eq:witness-2}; the sign twist cancels in
conjugation on $W^{\rm s}$.  Direct identities are
\[
 \sum_a Z_a^2=\frac23I_3,\qquad
 X^2=Z^2=I_2,\qquad
 \sum_{a=1}^2\langle a|\begin{cases}X,&a=1,\\Z,&a=2\end{cases}=0.
\]
For $\sum L^*L$, the trivial input has total Gram
$T_W^*T_W+T_2^*T_2=I_3$.  A $W$ input receives $I_2/3$ from its exit
and $2I_2/3$ from its diagonal self-blocks; a $V_2$ input receives $1/2$
from each contribution; the $W^{\rm s}$ input receives $I_3$ from its
self-blocks.  Off-isotypic sums vanish by covariance (or the last displayed
identity).  For $\sum LL^*$, the trivial output is
$3S_WS_W^*+2S_2S_2^*=I_3$ by unitarity of $R$, and the other diagonal
sums have the same complementary fractions.  This proves both identities.

An orthonormal real contrast frame on the six edges is
\[
 O_{(k,\pm),(W,i)}=\frac{\pm\delta_{ki}}{\sqrt2},\qquad
 O_{(k,\pm),(2,a)}=\frac{(Z_a)_{kk}}{\sqrt2}.
\]
Together with $\mathbf1/\sqrt6$ it is an orthogonal coefficient change.
Substituting it into Theorem~\ref{sel10:thm:contrast-isometry} gives
\eqref{sel10:eq:witness-six}.  The second normalization makes the source
unital as well as trace preserving, giving the stationary state.  Independence
follows from Theorem~\ref{sel10:thm:capacity}.  The entry and exit definitions
give $T=I_3$ and $\widehat S=R$.
\end{proof}

\begin{example}[A cyclic native return with no added internal control]
\label{sel10:ex:cycle}
Take
\begin{equation}
 R=\begin{pmatrix}0&0&1\\1&0&0\\0&1&0\end{pmatrix}.
 \label{sel10:eq:cycle}
\end{equation}
The native line is $P_T=|3\rangle\langle3|$.  Its first three forward
return vectors $|3\rangle,R|3\rangle,R^2|3\rangle$ are orthonormal.
Thus $\Omega_R\succ0$ and
Theorem~\ref{sel10:thm:return-saturation} gives the complete internal
algebra.  In this example the two quadratic invariants already have the
block forms
\begin{align}
 \mathsf Q&=\diag\left(P_T,\frac23I_W\otimes I_2,I_{V_2},I_{W^{\rm s}}\right),
 \nonumber\\
 \mathsf H&=\diag\left(R,I_W\otimes
          \begin{pmatrix}2/3&0\\1/3&2/3\end{pmatrix},
                     \frac12I_{V_2},I_{W^{\rm s}}\right).
 \label{sel10:eq:two-invariant-example}
\end{align}
Consequently $C^*(I,\mathsf Q,\mathsf H)=U(S_4)'$ directly.  The donor
block is full by cyclicity, the weak Jordan block generates $M_2$ with its
adjoint, and the two scalar characters have different values $1/2$ and $1$
on $\mathsf H$.  Different simple matrix sizes cannot be diagonally tied.
In particular the four central projectors are generated rather than added.
\end{example}

\begin{remark}[The raw action need not be a full matrix algebra]
For the displayed cyclic witness, the $W^{\rm s}$ sector is reducing for
every source coefficient and carries a diagonal raw algebra.  Thus raw
irreducibility on $\C^{14}$ is not the mechanism.  The invariant-algebra
result is an intersection inside the original source algebra, not an
identification of its raw $M_{14}$ commutant with a gauge algebra.  It is
also not a source-level detailed-balance assumption: the forward return
in \eqref{sel10:eq:cycle} is not its own adjoint.
\end{remark}

\section{A degree-six original-source panel proves generic internal saturation}
\label{sel10:generic}

Keep the fixed carrier $S_3$ and its normalized covariant source manifold
of dimension $71$.  No stationary, unital, conserved-sector, or independently
executed external-control hypothesis is imposed in this section.
Set $q=\mathsf Q$ and $h=\mathsf H$, and take the following fixed list of
fifteen invariant words:
\begin{equation}
 \mathcal W=\bigl(I,q,h,h^*,q^2,hq,h^*q,qh,h^2,h^*h,
                    qh^*,hh^*,(h^*)^2,qhq,h^*qh\bigr).
 \label{sel10:eq:fifteen-words}
\end{equation}
Each word has degree at most three in the quadratic invariants and hence
at most six in the original star-closed six-letter bank.  On the actual
physical Hilbert--Schmidt metric define
\begin{equation}
 \mathbb G_{10}(K)_{ij}=\operatorname{Tr}(w_i(K)^*w_j(K)),\qquad
 \Delta_{10}(K)=\det\mathbb G_{10}(K)\ge0.
 \label{sel10:eq:invariant-panel}
\end{equation}
The operator trace is on the fourteen-dimensional acted-on carrier, not an
unweighted seven-dimensional list of multiplicity coordinates.

\begin{theorem}[Generic full internal algebra in the complete covariance class]
\label{sel10:thm:generic}
If $\Delta_{10}(K)>0$, then
\begin{equation}
 \boxed{C^*(I,q,h)=\mathcal I_K=U(S_4)'
      \cong M_3\oplus M_2\oplus\C\oplus\C.}
 \label{sel10:eq:generic-full}
\end{equation}
The exact cyclic witness of Example~\ref{sel10:ex:cycle} gives
\begin{equation}
 \boxed{\Delta_{10}=\frac{98}{177147}=\frac{2\cdot7^2}{3^{11}}>0.}
 \label{sel10:eq:det-anchor}
\end{equation}
Therefore the positive-panel condition holds on an open dense full-smooth-volume
subset of the entire connected $71$-dimensional normalized source manifold.
In particular full invariant-algebra saturation is generic there without
adjoining any external matrix controls or imposing the V9 conservation
conditions on the raw source.
\end{theorem}
\begin{proof}
Every $w_i$ belongs to $\mathcal I_K\subseteq U(S_4)'$, whose complex
dimension is $3^2+2^2+1+1=15$.  A positive determinant says that fifteen
such operators are linearly independent.  They therefore span the whole
external commutant and prove \eqref{sel10:eq:generic-full}.

For an explicit rational evaluation of the determinant, replace the repeated
external blocks in \eqref{sel10:eq:two-invariant-example} by the reduced
seven-dimensional matrices
\[
 q_7=\diag(0,0,1,2/3,2/3,1,1),\quad
 h_7=R\oplus\begin{pmatrix}2/3&0\\1/3&2/3\end{pmatrix}
                         \oplus(1/2)\oplus1.
\]
For all invariant operators the physical trace is exactly the weighted trace
$\operatorname{Tr}(W_7\,\cdot)$ with
$W_7=\diag(1,1,1,3,3,2,3)$.
Thus the entries of the fifteen-by-fifteen Gram are explicitly
$\operatorname{Tr}(W_7w_i(q_7,h_7)^*w_j(q_7,h_7))$.
In row-major matrix-unit coordinates on $M_3\oplus M_2\oplus\C^2$,
the fifteen word columns have determinant of absolute value $7/3^8$.
The Gram of that coordinate basis has determinant $3^4\cdot2\cdot3$.
Multiplying these determinants gives \eqref{sel10:eq:det-anchor}; the
accompanying rational check supplies both exact elimination calculations.
Positivity also follows independently from the block generation argument
of Example~\ref{sel10:ex:cycle} and its word-span ranks $1,4,13,15$ at
lengths zero through three.

The entries of this Gram are polynomial in the real and imaginary entries
of the original coefficients, and the determinant is a sum of squared
absolute minors of their word synthesis.  Restricted to the connected
real-analytic Stiefel product it is analytic and not identically zero.
A nonzero real-analytic function on a connected finite-dimensional analytic
manifold has a zero set of empty interior and of measure zero in smooth
coordinate charts.  This standard local identity/zero-set argument can be
applied in a countable atlas; it is the analytic genericity mechanism already
used in the rigidity programme.  Continuity gives openness.  No distribution
over physical sources has been assumed.
\end{proof}

\begin{corollary}[A positive noisy panel is a one-sided source certificate]
\label{sel10:cor:noise}
Under the exact carrier and covariance hypotheses above, suppose a Hermitian measured Gram $\widehat{\mathbb G}_{10}$ has certified
entrywise errors at most $\varepsilon$ from the actual Gram, and
\begin{equation}
 \lambda_{\min}(\widehat{\mathbb G}_{10})>15\varepsilon.
 \label{sel10:eq:noise}
\end{equation}
Then the exact internal algebra is the full external commutant.  A vanishing
or unresolved determinant of this particular panel is not an all-word
exclusion; it may require a larger invariant word bank.
\end{corollary}
\begin{proof}
The spectral norm of a fifteen-by-fifteen matrix whose entries are bounded
by $\varepsilon$ is at most $15\varepsilon$.  The least-eigenvalue
perturbation inequality makes the true Gram positive definite.  Apply
Theorem~\ref{sel10:thm:generic}.  Only its forward implication is used.
\end{proof}

\begin{remark}[Why this does not contradict V9]
V9's codimension $68$ compares a direct-sum algebra with $M_7$ on one
multiplicity space in a source carried by $W\otimes\C^7$.
The present carrier has dimension fourteen and four inequivalent external
types.  Its source coefficients can transfer between those types; only the
invariant subalgebra is block diagonal.  Indeed
Corollary~\ref{sel10:cor:leakage} proves that the trivial-sector projector
cannot be conserved by this raw channel.  Thus centrality inside
$\mathcal I_K$ does not imply centrality inside $\mathcal A_K$ or a fixed
point of $\Phi^*$.  The two genericity statements concern different,
explicitly identified parameter spaces.  The external representation supplies
the block restriction here; an additional fitted charge does not.
\end{remark}

\begin{remark}[Which data the finite test uses]
The invariant words are algebraic combinations of original coefficients and
their formal adjoints.  Their faithful word Gram can be reconstructed from
appropriate common-frame quantum input/output or mixed-moment data.  A bare
classical edge-frequency table need not contain these entries, and coherent
sums of words are not promoted to executed operations.  The accepted
normalization in \eqref{sel10:eq:invariant-panel} is declared explicitly;
converting physical opportunity coefficients $\sqrt\vartheta K_e$ to it
uses the measured nonzero $\vartheta$ and propagates its error.  No physical
waiting or acceptance factor is discarded in claiming an operational rate.
\end{remark}

\section{The reversal record can obstruct the six-letter hypothesis}
\label{sel10:reversal}

The preceding theorem concerns the actual group $S_4$ and its six unoriented
outcomes.  The source manuscripts also retain root reversal and distinguish
$W_+$, $W_-$, and $(W^{\rm s})_-$.  An additional involution acting on the
same physical system must not be forgotten to obtain a larger commutant.
The next statement is a direct, evaluable compatibility test, not a claim
that every source has the same parity assignment.

\begin{theorem}[Capacity obstruction for an even locked alphabet]
\label{sel10:thm:parity}
Suppose the independently identified acted-on symmetry is
$S_4\times\langle\iota\rangle$, the six outcome labels are reversal even,
and every coefficient is even.  On the proposed carrier
\begin{equation}
 S^{\rm even}_r=(\mathbf1_+\otimes\C^r)
                   \oplus W_+\oplus W_-\oplus(V_2)_+
                   \oplus(W^{\rm s})_-,
 \label{sel10:eq:parity-carrier}
\end{equation}
the positive normalization and coherent identity sum force
\begin{equation}
 \boxed{r\le2.}
 \label{sel10:eq:parity-bound}
\end{equation}
In particular a coherent three-dimensional even trivial carrier cannot be
realized by this unchanged even six-letter source.

If the trivial multiplicity is instead split between the two reversal signs,
the capacities are $r_+\le2$ and $r_-\le1$.  Such a split may restore
normalization feasibility, but an internal algebra required also to commute
with $\iota$ cannot join those two trivial sectors into $M_3$.
\end{theorem}
\begin{proof}
The contrast environment is now $W_+\oplus(V_2)_+$.
Apply Theorem~\ref{sel10:thm:contrast-isometry} to the full product group.
A trivial even input can enter an invariant output/environment pair only
through $W_+\otimes W_+$ or $(V_2)_+\otimes(V_2)_+$.
Each contributes one multiplicity coordinate, so its output capacity is two.
For a trivial odd input only $W_-\otimes W_+$ contributes, giving capacity
one.  Commutation with the involution makes its two eigenspaces reducing,
so a full matrix algebra mixing them is excluded on the invariant side.
This uses a physically specified involution, not a relabeling of $W$ copies.
\end{proof}

\begin{proposition}[Keeping oriented outcomes changes the capacity]
\label{sel10:prop:oriented}
For the twelve ordered roots with their actual reversal permutation, the
contrast environment is
\begin{equation}
 \mathcal E_{0,12}=W_+\oplus(V_2)_+\oplus W_-\oplus(W^{\rm s})_-.
 \label{sel10:eq:oriented-environment}
\end{equation}
On \eqref{sel10:eq:parity-carrier}, the trivial-even capacity becomes four.
Thus the specific obstruction \eqref{sel10:eq:parity-bound} is absent for
that richer environment.  This does not assert that the normalized
\emph{original} six-letter source can acquire the missing odd directions
by changing the names of its outcomes, or that a full oriented source has
passed its other isometry inequalities and incidence tests.
\end{proposition}
\begin{proof}
The inherited ordered-root decomposition gives
\eqref{sel10:eq:oriented-environment} after removing its single mean.
Each of the four matching nontrivial system types supplies one invariant
output/environment pairing of even total parity.  Their multiplicities
sum to four.  The two environments are different representations; a
mathematical refinement of their labels is not an identity of physical
source operations.
\end{proof}

\begin{remark}[Read the group from the source, not from the target]
The group and environmental representation in the capacity test are physical
inputs to be identified.  If reversal acts only on an independent recorded
factor, or if actual source coefficients transform nontrivially under it,
then \eqref{sel10:eq:parity-carrier} or the even-alphabet hypothesis need not
hold.  Conversely, if those hypotheses do hold, the $r=3$ $S_4$ witness is
not an admissible realization of that stricter source.  The audit must keep
the relevant involution and its output action, rather than weakening them
silently to use the positive construction.
\end{remark}

\section{V10 status and the next source row}
\label{sel10:status}

There is now a positive source-level statement on the full external isotypic
carrier.  The internal matrix algebra is not selected by imposing a conserved
$3+2+1+1$ split on an otherwise single-factor channel.  On the explicitly
specified $S_4$ carrier $S_3$, it is the generic invariant part of the original
normalized six-coefficient source algebra.  A fixed positive Gram of fifteen
words, each of original source degree at most six, certifies that result.
The exact covariant benchmark has determinant $98/177147$ in the stated
physical trace convention.

The source normalization also supplies more than a genericity argument.
It bounds the trivial multiplicity and, at saturation, reconstructs an actual
line--plane split from output transition effects.  Faithful recurrence forces
a nonsingular same-source return, while its finite cyclicity test supplies
the complete internal algebra without independent private/weak controls.
These are specific assumption reductions.  They do not reconstruct a missing
acted-on carrier or its physical source matrices by selecting a benchmark.

The capacity obstruction with reversal is equally substantive.  The
six-even-letter normalization cannot act on three even trivial copies with
the stated nontrivial parity allocation.  A claim involving the full physical
reversal must therefore identify the actual environmental representation or
retain this obstruction.  Merely forgetting a recorded grading to recover
the $S_4$ covariance class would be circular.

\begin{center}
\small
\begin{tabular}{L{.37\textwidth}L{.55\textwidth}}
\toprule
Question & V10 result and remaining condition\\
\midrule
Can the internal target be generic without an added conserved charge?
& Yes on the verified full $S_3$ covariance family, not on V9's different
single-isotypic parameter space.\\[.3em]
Does the six-letter normalization constrain the trivial multiplicity?
& It imposes the exact capacity $r\le3$ on the stated $S_4$ carrier;
it does not force $r=3$ historically.\\[.3em]
Must the type line and private plane be chosen independently?
& At capacity saturation a native rank-one/rank-two split is reconstructed
from actual channel effects.  Identification with the manuscript's named
source supports remains a separate test.\\[.3em]
Are independent weak and private controls needed for the positive theorem?
& No. Nonsingular cyclic returns, or the positive fifteen-word panel,
generate the full original invariant algebra and its central supports.\\[.3em]
Can an even reversal constraint be ignored?
& No. On the displayed parity-resolved carrier it lowers the trivial-even
capacity to two and excludes the $r=3$ six-letter branch.\\[.3em]
Is historical Standard-Model selection proved?
& No. Actual system representation, complete symmetry/output action,
source-moment panel, type/private identification, and determinant/matter
occurrence still need their historical source rows.\\
\bottomrule
\end{tabular}
\end{center}

\paragraph{Next noncircular calculation.}
First determine the \emph{full acted-on} external representation and the
action of reversal on its actual output environment.  Apply the capacity
test before attempting algebra saturation.  On a verified $S_4$ $r=3$
source, populate the two quadratic invariant operators or their faithful
fifteen-word Gram; alternatively evaluate the native transition effects,
stationary return, and cyclicity Gram.  On a parity-obstructed six-letter
source, search for an actually supplied odd environmental sector or use the
source's correct larger alphabet.  Do not insert it as an unobserved Kraus
label.  Determinant stabilizers, anomaly-free matter types and common-history
odd incidence are downstream of this internal-algebra calculation and are
not consequences of its successful rank test.

\paragraph{Verification and cumulative preservation.}
The V9 prefix is retained byte for byte before its terminal document command.
The new exact-arithmetic checks verify all $S_4$ character products and
capacity counts, the full isometry and bistochastic normalization of the
fourteen-dimensional witness, covariance under all twenty-four group elements,
independence and normalization of the six original coefficients, the native
transition effects and return formulas, the fifteen-word rank and exact
physical Gram determinant, and the reversal capacity obstruction.
The V9 regression is rerun.  These computations test explicit finite
identities; the general capacity, recurrence, saturation and genericity
claims are proved above rather than extrapolated from samples.  The witness
is not a measured historical source, and no proof-assistant certification,
independent peer review, or mathematical priority claim is made.

\begin{thebibliography}{99}
\raggedright
\bibitem[19]{sel10-Verdon}
D.~Verdon, \emph{A covariant Stinespring theorem},
J.\ Math.\ Phys.\ \textbf{63} (2022), 091705;
\href{https://doi.org/10.1063/5.0071215}{doi:10.1063/5.0071215};
\href{https://arxiv.org/abs/2108.09872}{arXiv:2108.09872}.
Covariant dilation is established background; the original-frame decomposition,
finite multiplicity counts, native return and source-algebra consequences
used here are proved explicitly.
\end{thebibliography}

% END OF SOURCE-SELECTION V10 DEVELOPMENT BLOCK.
% Preserve V1--V10 and append later developments before the final terminator.
% State amendments explicitly instead of silently changing inherited results.


\clearpage
\begin{center}
{\Large\bfseries Source Selection V11}\\[.3em]
{\large Reversal-compatible source classes and native colour from four effects}\\[.2em]
An exact record-lift obstruction, quantitative parity cost, and a three-scalar certificate
\end{center}

\section{Audit the reversal map before importing its representation}
\label{sel11:context}

V10 proves generic internal saturation for the normalized six-edge source on
an explicitly verified $S_4$ carrier.  Its last section also proves a
capacity obstruction if an additional acted-on parity is imposed and all
six letters are even.  The present continuation evaluates which of these
symmetry statements can be connected to the displayed upstream writer.
It does not remove a reversal mark to obtain a favourable matrix dimension.

Three inherited objects must remain distinct.  The chronological ordered-pair
record has an edge-dependent orientation cocycle under $S_4$.  The literal
locked writer instead has square coefficient operators of product-pointer form
$K_e\otimes\Pi_\sigma\otimes P_r$.  Finally, reversal of chronological
matrix-unit histories is transpose, an order-reversing algebra map, rather
than conjugation by a system involution.  These statements occur in
\cite{Attack}, under
\src{fire:chronological-local-distinction} and the chronological transition
algebra calculation, and in \cite{Grand}, under
\src{def:locked-opportunity-instrument} and
\src{thm:GT76-locked-normalform}.  They do not identify one common physical
system representation merely because all three descriptions have twelve
oriented names.

The main negative result below turns the inherited warning into an exact
finite obstruction: the original product-pointer coefficients with coherent
sum $\sqrt2 I$ cannot carry chronological oriented-edge covariance on the
same system, for any choice of the unknown edge matrices.  A mixed-product
or Hilbert--Schmidt moment residual detects the obstruction.  A separate
calculation gives the sharp minimum parity-changing coefficient mass on
V10's fourteen-dimensional carrier.

The main positive result uses a \emph{genuinely oriented} source when that
is independently supplied.  All of its covariance and normalization
inequalities are solved, not only its trivial-sector capacity.  At trivial
multiplicity three, four native transition effects form a redundant Parseval
frame.  Three independent transition probabilities determine whether these
original-source quadratic effects already generate a supported internal
$M_3$.  No private matrix controls, independent bridge, physical adjoints, or
return-cyclicity assumption is needed for that statement.  The representation,
original record policy, and available preparations/Reads are still source
inputs.

\paragraph{Nonduplication and proof scope.}
The calibrated contrast-isometry theorem \ref{sel10:thm:contrast-isometry},
the earlier efficient-record refinement theorem, the external/internal census,
and the finite subdirect-product mechanism are inherited.  The chronological
and local label actions were already distinguished in \cite{Attack}; they
are not newly discovered representations.  The additions are the exact
product-pointer incompatibility and moment floors, its symmetry-resolved
source-short cost, the sharp parity-mismatch optimization, the full oriented
normalization class, and the probability-only native-effect criterion.
Finite covariant dilation and representation theory remain standard tools
\cite{sel10-Verdon,Etingof}.  These are programme-level advances with proofs,
not claims of literature priority or evaluated historical operator data.

\section{The locked product pointer cannot implement chronological edge covariance}
\label{sel11:cocycle}

Fix reference orientations on the six edges $e$ of $K_4$.  For $g\in S_4$
let $\epsilon_g(e)\in\{1,-1\}$ record whether $g$ preserves the reference
orientation after carrying $e$ to $ge$.  Chronological labels transform as
\[
 (e,\sigma)\longmapsto(ge,\epsilon_g(e)\sigma).
\]
On an actual edge system $H\cong\C^d$ suppose
\begin{equation}
 \sum_{e=1}^6 K_e^*K_e=I_d,\qquad \sum_{e=1}^6 K_e=cI_d,
 \qquad c\ne0.
 \label{sel11:eq:edge-normalization}
\end{equation}
Let $\Pi_+,\Pi_-$ be nonzero complementary orthogonal projections on a
pointer space, and define $A_{e,\sigma}=K_e\otimes\Pi_\sigma$.
A fixed additional type projection can be tensored throughout.  The
question is covariance on the \emph{same acted-on algebra}, not covariance
of an independently written classical label register.

\begin{lemma}[Orientation separation of two edges]
\label{sel11:lem:separation}
For any distinct target edges $e,f$ there is $g\in S_4$ such that
$\epsilon_g(g^{-1}e)\ne\epsilon_g(g^{-1}f)$.  This holds for every fixed
choice of the six reference orientations.
\end{lemma}
\begin{proof}
For disjoint edges, transpose the endpoints of $e$ and fix both endpoints
of $f$.  For adjacent edges, regard the reference orientations as a tournament
on four vertices.  At any vertex, two of its three incident edges point in
the same direction relative to that vertex.  Some vertex also has both an
incoming and an outgoing edge: otherwise every vertex would be a source or
a sink, while a tournament can have at most one source and at most one sink.
Thus adjacent pairs of both relative orientation types occur.  Choose a
preimage adjacent pair of the type opposite to the given target pair, and
map its common vertex and two neighbours to the target triple.  The product
of the two orientation signs is then minus one.  Extend this bijection by
the fourth vertex.  The argument applies to every reference tournament.
\end{proof}

\begin{theorem}[Exact product-pointer lift obstruction]
\label{sel11:thm:pointer-no-go}
Suppose that for each $g\in S_4$ there is a unital $*$-automorphism
$\alpha_g$ of
$C^*(I,A_{e,\sigma})$ satisfies
\begin{equation}
 \alpha_g(A_{e,\sigma})=A_{ge,\epsilon_g(e)\sigma}.
 \label{sel11:eq:chronological-lift}
\end{equation}
Then the six edge coefficients must have the form
\begin{equation}
 \boxed{K_e=cP_e,\qquad P_eP_f=0\ (e\ne f),\qquad
        P_e=P_e^*=P_e^2,\quad\sum_eP_e=I_d,\quad |c|=1.}
 \label{sel11:eq:projective-boundary}
\end{equation}
In particular the literal locked coherent sum $c=\sqrt2$ is incompatible
with \eqref{sel11:eq:chronological-lift}, independently of dimension,
faithful stationarity, or invertibility of the $K_e$.

Conversely, \eqref{sel11:eq:projective-boundary} with all $P_e\ne0$
admits the chronological permutation as an abstract algebra automorphism.
If the six $P_e$ have equal ranks and the two pointer ranks agree, it also
has a unitary implementation on the represented system.
\end{theorem}
\begin{proof}
Initially $A_{a,+}^*A_{b,-}=0$ for all $a,b$, by pointer orthogonality.
Given distinct target edges $e,f$, choose $g$ by
Lemma~\ref{sel11:lem:separation} and put $a=g^{-1}e$, $b=g^{-1}f$.
The two opposite input pointer labels then have the same output pointer
label.  Applying $\alpha_g$ to their zero product gives
$K_e^*K_f\otimes\Pi_\sigma=0$.  Nonzero pointer support implies
$K_e^*K_f=0$.  Thus all off-diagonal edge products vanish.

Multiplying the coherent sum by its adjoint gives
$|c|^2I=\sum_eK_e^*K_e=I$, so $|c|=1$.
Also $K_e^*\sum_fK_f=K_e^*K_e=cK_e^*$.
Taking adjoints gives the same positive matrix as $\overline cK_e$.
Consequently $P_e=\overline cK_e$ is Hermitian and $P_e^2=P_e$;
all pairwise products vanish and the coherent sum gives $\sum P_e=I$.
The converse is a permutation of the twelve minimal central projections
$P_e\otimes\Pi_\sigma$.  Equal representation ranks allow the corresponding
permutation of orthonormal block bases to implement it unitarily.
\end{proof}

\begin{corollary}[Resolved CP covariance does not evade the obstruction]
\label{sel11:cor:CP-lift}
For the independent locked coefficient bank, the same no-go applies if only
unitary covariance of each resolved rank-one Choi branch is initially
postulated.  It is not removed by choosing different coefficient phases.
\end{corollary}
\begin{proof}
Rank-one Choi equality makes the transformed coefficient a phase multiple
of the requested coefficient.  Summing all transformed coefficients gives
$cI$, as does the requested coherent coefficient sum.  Independence forces
every phase to be one, exactly as in the inherited phase-lock lemma
\ref{sel7:lem:phase-lock}.  Theorem~\ref{sel11:thm:pointer-no-go} applies.
\end{proof}

\subsection{A finite moment witness with a calibration floor}

Put
\begin{equation}
 \mathfrak C_K:=\sum_{e\ne f}\|K_e^*K_f\|_{\HS}^2.
 \label{sel11:eq:product-defect}
\end{equation}
This is a positive degree-four source moment.  It is not a probability of
executing an adjoint history.  For one rank-one pointer of each sign also
define the physical normalized one-letter Gram
\[
 G_{(e,\sigma),(f,\tau)}=\frac1{2d}\Tr(A_{e,\sigma}^*A_{f,\tau})
       =\frac{\delta_{\sigma\tau}}{2d}\Tr(K_e^*K_f).
\]
Let $\pi_g$ be the chronological permutation on these twelve labels.

\begin{proposition}[Uniform product and Gram-covariance residuals]
\label{sel11:prop:moment-floor}
For every bank satisfying \eqref{sel11:eq:edge-normalization},
\begin{equation}
 \boxed{\mathfrak C_K\ge\frac{d(|c|^2-1)^2}{30},\qquad
 \max_{g,i,j}|G_{\pi_g i,\pi_g j}-G_{i,j}|
             \ge\frac{\bigl||c|^2-1\bigr|}{60}.}
 \label{sel11:eq:moment-floor}
\end{equation}
Thus the locked $c=\sqrt2$ frame has a dimensionless normalized Gram
incompatibility of at least $1/60$.  With fixed $c$ but normalization errors
\[
 \|\sum_eK_e^*K_e-I\|\le\eta_N,\qquad
 \|\sum_eK_e-cI\|\le\eta_C,
\]
the second bound remains valid with numerator
$[\,\bigl||c|^2-1\bigr|-\eta_N-2|c|\eta_C-\eta_C^2\,]_+$.
If every normalized Gram entry is measured with error at most $\eta_G$,
the measured maximum discrepancy is at least the exact bound minus
$2\eta_G$.
\end{proposition}
\begin{proof}
The exact source identities give
\[
 \sum_{e\ne f}K_e^*K_f=(|c|^2-1)I_d.
\]
Cauchy--Schwarz for the thirty Hilbert--Schmidt vectors proves the first
bound.  Taking normalized trace shows that at least one of the thirty
numbers $\Tr(K_e^*K_f)/d$ has modulus at least
$\bigl||c|^2-1\bigr|/30$.  The orientation-separation lemma maps a
cross-pointer zero Gram entry to this same-pointer entry, proving the
second bound.  For errors, write $\sum K_e=cI+E$ and
$\sum K_e^*K_e=I+N$.  The extra operator in the summed off-diagonal
products is $\overline cE+cE^*+E^*E-N$, of norm at most
$2|c|\eta_C+\eta_C^2+\eta_N$.  Normalized trace cannot increase that
bound.  The final measurement statement is the triangle inequality.
\end{proof}

\begin{remark}[What this rules out]
The result rules out one identification of \emph{system} branches, not the
existence of an ordered classical record.  The source may write chronological
labels on another carrier, or admit additional controls, or use genuinely
non-product oriented coefficients.  Those are distinct data.  The accepted
probability-score Gram with scalar metric $\vartheta I$ is not the physical
Hilbert--Schmidt coefficient Gram used here.  The two may not be substituted
for one another.  At physical opportunity amplitudes $\sqrt\vartheta A_j$,
the Gram discrepancy is multiplied by $\vartheta$ and the product residual
by $\vartheta^2$; conditional normalization does not erase this scale.
\end{remark}

\section{The label-module mismatch has a nonzero source-short cost}
\label{sel11:modules}

The preceding no-go is stronger than a mismatch of label permutations.
The label calculation is nevertheless useful for auditing a claimed common
source synthesis before attempting a system lift.
Let $\mathcal E^{\rm loc}$ be the twelve-dimensional edge--pointer module
where $S_4$ acts on edges only and reversal swaps the pointer sign.
Let $\mathcal E^{\rm chr}$ be the ordered-root module where $S_4$ acts on
both endpoints and reversal interchanges them.  Both are representations
of $S_4\times\mathbb Z_2$.

\begin{proposition}[Nine common coefficient dimensions, not twelve]
\label{sel11:prop:module-deficit}
Their decompositions are
\begin{align}
 \mathcal E^{\rm loc}
 &=\mathbf1_+\oplus W_+\oplus(V_2)_+
                   \oplus\mathbf1_-\oplus W_-\oplus(V_2)_-,\\
 \mathcal E^{\rm chr}
 &=\mathbf1_+\oplus W_+\oplus(V_2)_+
                   \oplus W_-\oplus(W^{\rm s})_-.
 \label{sel11:eq:two-environments}
\end{align}
Their largest common subrepresentation has dimension nine.  Every
intertwiner between them has rank at most nine, and rank nine is attained.
After removing their uniform means, the respective dimensions are eleven
and the maximal common rank is eight.

Suppose actual equivariant syntheses $S_{\rm loc},S_{\rm chr}$ into one
Hilbert source satisfy $S_i^*S_i=\kappa I_{12}$ with $\kappa>0$.
Each directional Schur residual has rank at least three, trace at least
$3\kappa$, and three eigenvalues exactly $\kappa$ on its unmatched
irreducible sector.  Both lower bounds can be attained.
\end{proposition}
\begin{proof}
The local representation is the tensor product of the six-edge permutation
module $\mathbf1\oplus W\oplus V_2$ and the two parity characters.
The chronological decomposition is inherited from the ordered-root
calculation.  The common summands are
$\mathbf1_+,W_+,(V_2)_+,W_-$, of dimensions $1,3,2,3$.
Schur's lemma proves the rank assertion and provides an isomorphism on
these common summands.  Removing the mean deletes the common trivial line.

The mixed Gram $C=S_{\rm loc}^*S_{\rm chr}$ is an intertwiner and vanishes
on the unmatched irreducibles.  Hence
$\kappa I-\kappa^{-1}C^*C$ equals $\kappa I$ on the three unmatched
chronological dimensions, and the reverse short has the same property.
Both shorts are positive.  Identifying the common summands isometrically
and placing the two unmatched triples orthogonally in a common source
attains the bounds.
\end{proof}

\begin{remark}[Transpose reversal is a different operation]
The chronological transition calculation in \cite{Attack} realizes
reversal as $X\mapsto X^{\mathsf T}$ on its matrix units.  It reverses
multiplication: $(XY)^{\mathsf T}=Y^{\mathsf T}X^{\mathsf T}$.
It therefore cannot be inserted as $X\mapsto JXJ$ on the same noncommutative
system algebra.  For example $E_{12}E_{23}=E_{13}$, but
$E_{21}E_{32}=0$.  More generally a bijection that is both an automorphism
and an anti-automorphism forces its algebra to be commutative.  Transpose
is a Hilbert--Schmidt unitary on the \emph{operator space}; regarding that
larger space as a new physical input is exactly the word-module substitution
excluded in V1.  The source's chronological transpose statement is retained,
not corrected into a different involution.
\end{remark}

\section{A separate recorded reversal is fully compatible with the local source}
\label{sel11:local-lift}

The no-go does not apply to the actual local label action
$(e,\sigma)\mapsto(ge,\sigma)$ followed by an independent global pointer
swap.  On an equal-rank binary pointer choose an identification of its two
fibres and let $X_R$ exchange them.  For an already supplied $S_4$-covariant
edge system $U_g$ define
\[
 \widetilde U_g=U_g\otimes I_R,\qquad
 \widetilde J=I_H\otimes X_R.
\]
The identification is a declared record-frame convention; the statement
below is about represented covariance, not execution of a new swap gate.

\begin{proposition}[Local reversal lift with its record multiplicity retained]
\label{sel11:prop:local-lift}
The unchanged product coefficients satisfy
\[
 \widetilde U_gA_{e,\sigma}\widetilde U_g^*=A_{ge,\sigma},\qquad
 \widetilde JA_{e,\sigma}\widetilde J=A_{e,-\sigma},\qquad
 [\widetilde J,\widetilde U_g]=0.
\]
Writing $\mathcal A_K=C^*(I,K_e)$,
$\mathcal D_R=C^*(\Pi_+,\Pi_-)$ and
$\mathcal I_K=\mathcal A_K\cap U(S_4)'$, the actual invariant algebra is
\begin{equation}
 \boxed{
 (\mathcal A_K\otimes\mathcal D_R)
       \cap(\widetilde U(S_4),\widetilde J)'
       =\mathcal I_K\otimes I_R.}
 \label{sel11:eq:local-invariants}
\end{equation}
If an independent two-type record is also retained but is not acted on by
these symmetries, its diagonal algebra remains as a tensor factor.  Thus a
V10 internal algebra of dimension fifteen gives dimension thirty in the
complete type-marked invariant algebra, not a silent deletion of one copy.
\end{proposition}
\begin{proof}
The covariance identities are direct products.  The original coherent sum
recovers $I\otimes\Pi_\sigma$ from $\sum_e A_{e,\sigma}/c$, and summing
the two pointer branches gives $K_e\otimes I$.  Hence their source algebra
is exactly $\mathcal A_K\otimes\mathcal D_R$, as in the inherited tensor
formula.  Commutation with $\widetilde U_g$ puts both pointer blocks in
$\mathcal I_K$; commutation with the swap makes them equal.  The separate
type factor is unaffected and has dimension two.
\end{proof}

This evaluates the local realization of the reversal mark without changing
any $K_e$.  It also explains why V10's parity obstruction cannot simply be
asserted for this different acted-on representation: every payload irrep is
repeated in both swap-parity sectors of $H\otimes R$.  The chronological
cocycle is still not implemented, by Theorem~\ref{sel11:thm:pointer-no-go}.

\section{A sharp quantitative form of the parity-capacity obstruction}
\label{sel11:parity-cost}

A parity assignment on the payload, if independently verified, has a
stronger consequence than infeasibility of exactly even coefficients.
Let a finite-group system $S$ carry a commuting involution $J$.  In the
calibrated contrast-isometry theorem \ref{sel10:thm:contrast-isometry}, let
the contrast environment be parity even.  For each group irrep $\lambda$
write $m_{\lambda,+},m_{\lambda,-}$ for its input multiplicities and
$n_{\lambda,+},n_{\lambda,-}$ for its multiplicities in the output
$S\otimes\mathcal E_0$.  Assume the ordinary group-covariant normalized
source class is nonempty, so $m_{\lambda,+}+m_{\lambda,-}
\le n_{\lambda,+}+n_{\lambda,-}$.
Put $a=|c|^2/n$ for its fixed identity fraction.

\begin{theorem}[Exact minimum parity-changing coefficient mass]
\label{sel11:thm:parity-distance}
Over all such normalized group-covariant coefficient lists,
\begin{equation}
 \boxed{
 \min\sum_e\left\|\frac{K_e-JK_eJ}{2}\right\|_{\HS}^2
 =(1-a)\sum_\lambda \dim V_\lambda
   \left[(m_{\lambda,+}-n_{\lambda,+})_+
        +(m_{\lambda,-}-n_{\lambda,-})_+\right].}
 \label{sel11:eq:parity-distance}
\end{equation}
The minimum is attained by an explicit blockwise isometry.  It is a
coefficient-mass minimization on the specified closed normalization class;
it does not assert that an additional source constraint, such as a prescribed
stationary density, survives the minimizing construction.
\end{theorem}
\begin{proof}
By the inherited contrast decomposition, the odd coefficient mass is
$(1-a)$ times the mismatched-parity squared Hilbert--Schmidt mass of an
equivariant isometry.  On one irrep it is an isometry
$b:\C^{m_+}\oplus\C^{m_-}\to\C^{n_+}\oplus\C^{n_-}$.
The even-input matching block has rank at most $n_+$ and is a contraction,
so its squared Hilbert--Schmidt norm is at most $\min(m_+,n_+)$.
The unmatched even-input mass is at least $(m_+-n_+)_+$; similarly for
the odd input.  Multiply by $\dim V_\lambda$ and add.

For attainment, match $\min(m_+,n_+)$ positive input basis vectors to
positive output basis vectors and similarly for negative parity.  At most
one input parity has a deficit, because total output capacity is sufficient.
Send the unmatched input vectors to unused opposite-parity output vectors.
The resulting column is an isometry with exactly the stated mass.
Orthogonal coefficient synthesis reconstructs the original calibrated list.
\end{proof}

\begin{theorem}[A forced parity-changing input ray on V10's saturated carrier]
\label{sel11:thm:forced-parity}
Retain V10's actual $S_4$ source class on
\[
 S_3=\mathbf1^{\oplus3}\oplus W^{\oplus2}\oplus V_2\oplus W^{\rm s},
\]
and independently specify the parity allocation
$\mathbf1_+^{\oplus3},W_+,W_-,(V_2)_+,(W^{\rm s})_-$.
Let $p_0$ be the trivial-system projection and $p_-$ the total odd-system
projection.  For \emph{every} source in this class there is a rank-one
projection $q_-\le p_0$ determined by its entry isometry such that
\begin{equation}
 \boxed{p_0\Phi^*(p_-)p_0=\frac23 q_-,\qquad
 \Tr[p_-\Phi_*(p_0/3)]=\frac29.}
 \label{sel11:eq:forced-parity}
\end{equation}
In particular
\[
 \sum_e\|\tfrac12(K_e-JK_eJ)p_0\|_{\HS}^2=\frac23,
 \qquad
 \sum_e\|[J,K_e]p_0\|_{\HS}^2=\frac83.
\]
Every parity-preserving channel $\widehat\Phi_*$ on this same system obeys
\begin{equation}
 \boxed{\|\Phi_* -\widehat\Phi_*\|_\diamond\ge\frac43.}
 \label{sel11:eq:parity-channel-distance}
\end{equation}
For locked opportunities, comparing with the same idle branch multiplies
this last lower bound by $\vartheta$.
\end{theorem}
\begin{proof}
At saturation the three entry rows $T_W$ and $T_2$ in
Theorem~\ref{sel10:thm:native-split} are a unitary matrix.
One row of $T_W$ lands in $W_-$ and the other in $W_+$; $T_2$ lands in
even $V_2$.  The odd row therefore has a rank-one input Gram $q_-$,
which gives \eqref{sel11:eq:forced-parity} by the inherited channel
formula.  The mean component is even, and $p_0$ is even, so odd coefficient
norms measure precisely this outgoing transition mass.  The commutator
identity has the factor four because $[J,K]$ is twice its odd part up to
multiplication by $J$.

On the input state $q_-$, the old channel places probability $2/3$ in the
odd sector.  A parity-preserving channel places probability zero there.
Measuring the two parity outcomes bounds the output trace norm below by
$2(2/3)$, proving the diamond lower bound without an ancilla.
The two locked channels differ by $\vartheta(\Phi_*-\widehat\Phi_*)$.
\end{proof}

For completeness the output parity multiplicities, in the order
$(\mathbf1,W,V_2,W^{\rm s},\operatorname{sgn})$, are
\[
 n_+=(2,6,5,3,1),\qquad n_-=(1,4,2,4,1).
\]
The input counts are $(3,1,1,0,0)$ and $(0,1,0,1,0)$.  Only the even
trivial block is deficient.  The sharp global minimum in
\eqref{sel11:eq:parity-distance} is consequently $2/3$.
This is a positive, calibrated obstruction; an approximately even source
cannot satisfy all the exact hypotheses with arbitrarily small odd mass.

\section{A genuinely oriented source has a complete admissible normalization class}
\label{sel11:oriented}

Now suppose twelve \emph{actual} square system coefficients $C_{ij}$,
$i\ne j$, are supplied, and their resolved system covariance is with respect
to the chronological action and a commuting reversal involution $J$:
\[
 U_gC_{ij}U_g^*=C_{g(i)g(j)},\qquad JC_{ij}J=C_{ji}.
\]
Assume $\sum C_{ij}^*C_{ij}=I$ and $\sum C_{ij}=c_{12}I$, with
$0<a:=|c_{12}|^2/12<1$.  This is a new \emph{class to test}, not a
refinement of the unchanged product-pointer source.  Its coherent calibration
must actually be supplied.  In particular $a=1/3$ means $|c_{12}|=2$,
whereas keeping $c_{12}=\sqrt2$ would mean $a=1/6$.

Use the parity-resolved system
\begin{equation}
 S_r^J=\mathbf1_+^{\oplus r}\oplus W_+\oplus W_-
                           \oplus(V_2)_+\oplus(W^{\rm s})_-.
 \label{sel11:eq:oriented-system}
\end{equation}
Its contrast environment is precisely
$\mathcal E_{0,12}=W_+\oplus(V_2)_+\oplus W_-\oplus(W^{\rm s})_-$.

\begin{theorem}[Complete oriented capacity and source dimension]
\label{sel11:thm:oriented-capacity}
For $r\ge1$, the normalized covariant twelve-coefficient class on
\eqref{sel11:eq:oriented-system} exists exactly when $r\le4$.
The output multiplicities for the five \emph{occurring input} types
$(\mathbf1_+,W_+,W_-,(V_2)_+,(W^{\rm s})_-)$ are
\begin{equation}
 \boxed{(4,\ r+7,\ r+8,\ r+6,\ r+8).}
 \label{sel11:eq:oriented-capacities}
\end{equation}
Its normalized source manifold is connected and real analytic of dimension
\begin{equation}
 \boxed{54+16r-r^2.}
 \label{sel11:eq:oriented-dimension}
\end{equation}
In particular $r=3$ gives dimension $93$, and $r=4$ gives dimension $102$.
Twelve independent coefficients occur on an open dense full-volume subset;
at $r=4$ independence is automatic.
\end{theorem}
\begin{proof}
Tensor the displayed system with the four contrast irreps, retaining the
product of their parity signs.  The $S_4$ fusion rules already used in V10
give \eqref{sel11:eq:oriented-capacities}; for example a trivial-even input
has one matching output/environment contraction for each of the four
nontrivial types, hence capacity four.  Every other occurring input has
multiplicity one and strictly positive output capacity.  The inherited
contrast-isometry theorem therefore gives exactly $r\le4$ and the product
of one $V_r(\C^4)$ with four one-column Stiefel spaces.  Its real dimension
is
\[
 (8r-r^2)+(2r+13)+(2r+15)+(2r+11)+(2r+15).
\]
A $4\times r$ isometry whose four rows are nonzero exists for every
$1\le r\le4$.  Its trivial-input block makes each inequivalent contrast
environment irrep occur nontrivially in the coefficient synthesis.  The
nonzero mean completes the twelve-dimensional permutation module, so the
coefficient map is injective.  Zero rows or vanished coefficient components
are proper analytic conditions; the connected-manifold argument gives the
generic assertion.  At $r=4$ the trivial-input block is unitary and has no
zero row, giving automatic independence.
\end{proof}

\begin{proposition}[Native effects and unavoidable odd source activity]
\label{sel11:prop:four-effects}
Let $\alpha$ run over the four contrast types, with system dimensions
$d_\alpha=(3,2,3,3)$ and system projections $p_\alpha$.
The trivial-input block of the contrast isometry is a matrix
$T=(t_\alpha)_\alpha:\C^r\to\C^4$, $T^*T=I_r$.
Set $E_\alpha=t_\alpha^*t_\alpha$.  The actual channel satisfies
\begin{equation}
 \boxed{
 \Phi_*(\rho)=a\rho\ \oplus\
 (1-a)\bigoplus_\alpha\frac{I_{d_\alpha}}{d_\alpha}
                             t_\alpha\rho t_\alpha^*,
 \qquad \sum_\alpha E_\alpha=I_r,
 \quad
 p_0\Phi^*(p_\alpha)p_0=(1-a)E_\alpha}
 \label{sel11:eq:oriented-trivial-channel}
\end{equation}
for $\rho$ supported on $p_0$.
At $r=3$, at least one input ray transfers to the odd system sector with
probability exactly $1-a$.  The uniform trivial input has odd transition
probability at least $(1-a)/3$.
\end{proposition}
\begin{proof}
Each matching irreducible contraction has the unique normalized form
$L_{\alpha b}p_0=|b\rangle t_\alpha/\sqrt{d_\alpha}$.
Tracing its environment gives the displayed formula.  The two even target
rows define an effect $E_{\rm even}$ of rank at most two on $\C^3$.
A unit vector in its kernel obeys
$\langle E_{\rm odd}\rangle=1$ because the four effects sum to identity.
Also $\Tr E_{\rm odd}=3-\Tr E_{\rm even}\ge1$, since $E_{\rm even}$ is
a contraction of rank at most two.  These give the two probability claims.
\end{proof}

A nonzero odd environmental representation is therefore not merely additional
label capacity on this branch: it must carry positive source amplitude.
The old six-even-letter source cannot acquire it through scalar refinements
of efficient outcomes.  Nor does an arbitrary parity-covariant source conserve
parity; covariance permits an odd output paired with an odd environment.

\section{Four native transition effects give a probability-only colour certificate}
\label{sel11:colour}

Fix $r=3$ in the genuine oriented class.  Define contrast coefficients
$L_{\alpha b}$ in the verified orthonormal environment frame and the four
quadratic original-source invariants
\begin{equation}
 Q_\alpha=\sum_{b=1}^{d_\alpha}L_{\alpha b}^*L_{\alpha b}.
 \label{sel11:eq:Qalpha}
\end{equation}
Every $Q_\alpha$ belongs to $C^*(I,C_{ij})$ and commutes with the full
$S_4\times\mathbb Z_2$ action.  Its trivial block is $E_\alpha$ from
\eqref{sel11:eq:oriented-trivial-channel}.  No external projection has
been added to the generated source algebra in this definition.

\begin{theorem}[Exact four-effect alternative]
\label{sel11:thm:four-frame}
Let $T:\C^3\to\C^4$ be the trivial-input isometry.  Choose a unit
$w\in\Ker T^*$ and put $s=|\supp w|$.  Then
\begin{equation}
 \boxed{
 C^*(I_3,E_1,E_2,E_3,E_4)
 \cong
 \begin{cases}
 M_3,&s=4,\\
 M_2\oplus\C,&s=3,\\
 \C^3,&s=1\text{ or }2.
 \end{cases}}
 \label{sel11:eq:four-frame-algebra}
\end{equation}
On the $s=4$ branch, the \emph{original system coefficient algebra} contains
the supported internal summand $\B(p_0S_3^J)\cong M_3(\C)$.
This conclusion needs no reverse operation, faithful stationary state,
private $M_2$ bank, or separately assumed colour bridge.
\end{theorem}
\begin{proof}
Put $v_\alpha=T^*e_\alpha$, so
$E_\alpha=|v_\alpha\rangle\langle v_\alpha|$ and
\[
 \langle v_\alpha,v_\beta\rangle
       =\delta_{\alpha\beta}-w_\alpha\overline w_\beta.
\]
If $w_\alpha=0$, $v_\alpha$ is a unit vector orthogonal to every other
$v_\beta$.  It gives an independent scalar summand.  The remaining $s$
vectors span an $(s-1)$-dimensional space.  For $s\ge2$ they are nonzero
and every pair has nonzero overlap, so their rank-one operators generate
all matrix units on that span: their products give nonzero multiples of
$|v_\alpha\rangle\langle v_\beta|$, and the vectors span.  For $s=1$
the corresponding vector is zero and the other three are an orthonormal
basis.  This proves \eqref{sel11:eq:four-frame-algebra}.

On the whole system the algebra $\mathcal Q=C^*(I,Q_1,\ldots,Q_4)$ is
contained in the full symmetry commutant $M_3\oplus\C^4$.  Its trivial
compression is $M_3$ for $s=4$.  The ideal of $\mathcal Q$ generated by
commutators vanishes on all four scalar blocks and projects onto the entire
simple $M_3$ block.  It therefore equals the supported $M_3$ summand.
In particular its unit is $p_0$, obtained inside the original source
algebra rather than loaded as a control.  This is the earlier
subdirect-product mechanism applied to the newly computed effects.
\end{proof}

For an admitted preparation $p_0/3$ and terminal sector readout, let
\[
 \pi_\alpha=\Tr[p_\alpha\Phi_*(p_0/3)].
\]
Otherwise these quantities denote reconstructed channel functionals, not
claimed experiments.  Define
\begin{equation}
 \delta_\alpha=1-\frac{3\pi_\alpha}{1-a},\qquad
 \kappa_{\rm col}=\prod_{\alpha=1}^4\delta_\alpha.
 \label{sel11:eq:probability-certificate}
\end{equation}
These defects are reconstructed from the source transition masses; they
are not fitted entries of an internal Hamiltonian.

\begin{theorem}[Three independent scalar probabilities certify native colour]
\label{sel11:thm:probability-colour}
The source normalization implies
\begin{equation}
 \boxed{\delta_\alpha=|w_\alpha|^2\ge0,\qquad
        \sum_\alpha\delta_\alpha=1,\qquad
        \sum_\alpha\pi_\alpha=1-a.}
 \label{sel11:eq:defects}
\end{equation}
Consequently
\begin{equation}
 \boxed{\kappa_{\rm col}>0
 \iff C^*(E_\alpha:\alpha)=M_3
 \ \Longrightarrow\ 
 \B(p_0S_3^J)\subseteq C^*(I,C_{ij})\cap(U(S_4),J)'.}
 \label{sel11:eq:probability-colour}
\end{equation}
The other two cases are classified exactly by the number of positive
$\delta_\alpha$ as in \eqref{sel11:eq:four-frame-algebra}.  The positive
condition holds on an open dense full-volume subset of the complete
93-dimensional normalized covariance manifold.  If it fails, other source
words may still generate colour; failure excludes only saturation by the
four native effects.
\end{theorem}
\begin{proof}
$TT^*=I_4-ww^*$ gives
$\Tr E_\alpha=1-|w_\alpha|^2$.  Substitute
$\pi_\alpha=(1-a)\Tr E_\alpha/3$ from
\eqref{sel11:eq:oriented-trivial-channel} to obtain
\eqref{sel11:eq:defects}.  The preceding theorem gives the equivalence and
inclusion.  Each equation $w_\alpha=0$, equivalently
$\Tr E_\alpha=1$, is a proper real-analytic condition on
$V_3(\C^4)$.  A frame with all four complement coordinates nonzero exists,
so their finite union has empty interior and measure zero.  The other
Stiefel factors are independent, proving the assertion on the full source
manifold.  The four transition masses have one fixed sum, leaving three
independent scalars; no coefficient phase reconstruction is used.
\end{proof}

\begin{corollary}[One-sided certification with probability error]
\label{sel11:cor:probability-noise}
If $|\widehat\pi_\alpha-\pi_\alpha|\le\epsilon_\alpha$ and
\[
 1-\frac{3\widehat\pi_\alpha}{1-a}
                >\frac{3\epsilon_\alpha}{1-a}
 \quad\text{for every }\alpha,
\]
then the original source has the supported colour summand in
\eqref{sel11:eq:probability-colour}.  If an independently calibrated opportunity source has raw masses
$p_\alpha=\vartheta\pi_\alpha$, replace $1-a$ throughout by
$\vartheta(1-a)$.  Exact capacity, covariance and input-dimension identities
must be verified independently of these statistical errors.
\end{corollary}
\begin{proof}
The inequalities make every true $\delta_\alpha$ strictly positive.
Theorem~\ref{sel11:thm:probability-colour} applies.  Multiplication of all
accepted transition masses by $\vartheta$ gives the raw formula.
\end{proof}

\subsection{An exact frame, its calibrated gap, and a dimension ambiguity}

\begin{example}[A balanced native frame]
\label{sel11:ex:tetra-frame}
Take the actual trivial-input isometry, as a declared benchmark, to be
\[
 T=\frac12\begin{pmatrix}
 1&1&1\\1&-1&-1\\-1&1&-1\\-1&-1&1
 \end{pmatrix}.
\]
It has complement $w=(1,1,1,1)^{\mathsf T}/2$, so
$\delta_\alpha=1/4$, $\kappa_{\rm col}=1/256$, and
$\pi_\alpha=(1-a)/4$ for every sector.  The other equivariant isometry
blocks may be chosen arbitrarily from the admissible Stiefel factors;
Theorem~\ref{sel11:thm:oriented-capacity} then supplies a complete normalized
source, and the four nonzero rows ensure all twelve coefficients are
independent.  This is an existence witness, not the historical channel.

For the four Hermitian effects $E_\alpha$ counted once each, the commutator
form $q_E(x)=\sum_\alpha\|[E_\alpha,x]\|_{\HS}^2$ has spectrum
\begin{equation}
 \boxed{\{0^{(1)},(1/2)^{(3)},(3/2)^{(5)}\}.}
 \label{sel11:eq:effect-spectrum}
\end{equation}
Indeed $\sum E_\alpha^2=3I/4$, while
$x\mapsto\sum E_\alpha xE_\alpha$ has eigenvalues $3/4$ on $I$,
$1/2$ on the three symmetric off-diagonal matrices, and zero on their
five-dimensional traceless orthogonal complement.  Thus the commutator
Laplacian is $(3/2)\id-2\sum E_\alpha(\cdot)E_\alpha$.
For the unrescaled channel effects $(1-a)E_\alpha$, the diagnostic gap is
$(1-a)^2/2$; raw locked effects give $\vartheta^2(1-a)^2/2$.
This is a transition-effect diagnostic, not an assertion of an identical
spectral gap for the original coefficient Laplacian.
\end{example}

\begin{proposition}[The same sector probabilities do not select dimension three]
\label{sel11:prop:dimension-ambiguity}
At the other admissible saturation $r=4$, the trivial entry matrix is
unitary and the four native effects are orthogonal rank-one projections.
Their generated algebra is $\C^4$, not $M_3$.  For the maximally mixed
trivial input their four transition probabilities are again
$(1-a)/4$, exactly those in Example~\ref{sel11:ex:tetra-frame}.
\end{proposition}
\begin{proof}
A square isometry is unitary, so its four rows are an orthonormal basis.
Each effect has trace one, and dividing by the actual input dimension four
gives $(1-a)/4$.  The source-capacity theorem admits this class.
\end{proof}

The positive colour test is therefore conditional on the actual rank-three
trivial carrier.  It cannot use its own transition counts to manufacture
that rank or identify the private overlap $1/5$.  This explicit ambiguity
prevents a population-only dimension-selection claim.

\section{Reversal-normalized symmetry is not reversal-commuting symmetry}
\label{sel11:status}

Suppose an actual involution $J$ commutes with $U(S_4)$ and normalizes the
source algebra $\mathcal A$.  Then it normalizes
$\mathcal I=\mathcal A\cap U(S_4)'$.  If that internal algebra saturates
the full $S_4$ commutant, write
$J=\bigoplus_\lambda I_{V_\lambda}\otimes j_\lambda$.
The fixed algebra for this additional action is
\begin{equation}
 \mathcal I\cap J'
 =\bigoplus_\lambda I_{V_\lambda}\otimes
          \bigl(M_{m_{\lambda,+}}\oplus M_{m_{\lambda,-}}\bigr).
 \label{sel11:eq:parity-fixed-internal}
\end{equation}
This follows by diagonalizing each Hermitian involution $j_\lambda$.
It is an elementary commutant calculation, included to specify the physical
claim rather than advertised as a new classification principle.

On \eqref{sel11:eq:oriented-system} with $r=3$, the largest internal
algebra commuting with the \emph{whole} product group is
\[
 M_3\oplus\C\oplus\C\oplus\C\oplus\C,
 \qquad\dim_\C=13.
\]
The two $W$ copies have opposite parities.  The weak off-diagonal matrix
units, if they occur in $\mathcal A\cap U(S_4)'$, are odd under $J$;
they do not commute with it.  One may retain $J$ as a normalizing automorphism
of the larger internal algebra, but that is a different statement from
requiring all internal operations to commute with $J$.  Which statement is
physically required must be read from the source and the joint Clifford
interface.  A determinant seed and its transformation law require a further
independent check; neither interpretation selects it here.

\paragraph{The evaluated upstream alternatives.}
For the literal locked product-pointer writer, the chronological cocycle
cannot be lifted to its same-system coefficients at $c=\sqrt2$.
The scalar-type record does not remove this obstruction.  A separate local
pointer swap is instead an exact compatible reversal and retains the payload
internal algebra, with all remaining type copies explicitly included.
If the fourteen-dimensional payload itself has the stated reversal parity,
the six-letter source has a forced parity-changing ray and cannot be
arbitrarily close to a parity-conserving channel.  A genuinely oriented
source has a different, fully solvable normalization class.  On its verified
rank-three trivial carrier, the four native effects generically generate
colour and three independent scalar transition masses certify that fact.

\paragraph{The next source row is a symmetry-labelled operation, not another router.}
The historical source still does not specify the same-cut physical map
identifying its categorical reversal allocation with the acted-on payload,
its independently stored pointer, or a genuine oriented amplitude bank.
The notes supply exact tests distinguishing those cases and do not replace
that missing identification by a benchmark.  For a verified oriented
rank-three source, evaluate $\pi_\alpha$ or the quadratic $Q_\alpha$.
For the original local product source, retain its actual action and V10's
source-moment test instead of importing the chronological label decomposition.
The existing source-short theorem does not assert normalized CP descent on
its quotient.  Determinant occurrence, same-history matter incidence, and the
physical generation/type/private comparison remain separate.

\paragraph{Verification and append-only record.}
The complete V10 source prefix is retained before its terminal document
command.  The companion script verifies the orientation cocycle, the
pair-separation property, exact module multiplicities and maximal common
rank, calibrated product and Gram defects, all oriented capacity counts,
parity-flip formulas on the inherited V10 witness, the four-effect
classification examples, the transition diagnostic spectrum, and the
rank-three/rank-four probability ambiguity.  V10's regression is rerun.
General all-source assertions are established by the written proofs rather
than by sampled frames.  No new measured historical operators, physical
symmetry lift, proof-assistant verification, or independent peer review is
claimed.

% END OF SOURCE-SELECTION V11 DEVELOPMENT BLOCK.
% Preserve V1--V11 and append subsequent work before the final terminator.


\clearpage
\begin{center}
{\Large\bfseries Source-selection continuation V12}\\[.4em]
{\large Normalized descent, retained records, and external memory}\\[.25em]
A harmonic-metric test and the exact edge-blindness boundary
\end{center}

\section{Go upstream of the reduced carrier: what process has actually descended?}
\label{sel12:context}

The last source row in V11 concerns the acted-on carrier, its full symmetry,
and the physical map implementing the claimed reduction.  Its warning about
geometry shorting is substantive: a positive source quotient has a well-defined
Hilbert metric, but does not automatically carry the original normalized
instrument.  Conversely, covariance does supply a different, canonical
normalized reduction onto the external multiplicity algebra.  The missing
question is whether this reduction also preserves the original resolved
histories.

This block separates those operations.  First, an exact harmonic-metric
identity decides normalized amplitude descent through a specified positive
quotient.  Second, the external conditional expectation gives an autonomous
multiplicity channel with no hypothesis of a physical inverse or a chosen
internal gauge factor.  Third, the original record-respecting reduction has a
strictly stronger test.  On V10's full six-edge source with both native return
ports present, that test forces every retained classical outcome to be
independent of the edge label.  A smaller three-label coarsening suffices for
an isolated return block, but not for the full coherent source.  The difference
is quantified and evaluated on the already constructed V10 source.

\paragraph{Audit and nonduplication.}
The source short itself is inherited from \cite{Duality},
\src{thm:geometry-short}; it is not reconstructed here.  The fixed-algebra
identification under a faithful stationary state is an input from V6 of
\cite{Rigidity}.  The full source representation, contrast coefficients,
entry and return blocks are those of Sections~\ref{sel10:capacity}--\ref{sel10:returns}.
The single-isotype partial trace of V8 and the record/refinement results of
V4--V6 remain valid.  They do not already give the multi-isotype
record-preserving descent test or the calibrated six-edge obstruction below.
A targeted audit of the supplied monograph, \cite{Attack,Rigidity}, and
V1--V11 did not locate those specialized results.  This is not a claim of
exhaustive literature priority.

Exact quantum model reduction by conditional expectations is established
operator-theoretic machinery; in particular, reduction of averaged and
conditioned dynamics must be distinguished \cite{sel12-GT,sel8-Watrous}.
We give the finite proofs needed here.  The programme-specific additions are
the harmonic test on a proposed source short, the return-output dimension
bound, the exact tetrahedral coarsening classification, the edge-blindness
and error formulas, and the two-step evaluation of the unchanged V10
benchmark.  No new historical matrix or intervention is postulated.

\paragraph{Scope.}
All state reductions below are linear completely positive maps, acting also
on states entangled with an untouched auxiliary system.  A reduced channel is
a description of an actual source under a stated observation map, not an
instruction to perform that map physically between events.  Exact descent is
required for every input on the declared carrier.  A smaller reachable input
class, an enlarged classical/quantum memory, or a history-dependent decoder is
a different reduction problem and is not excluded by a failure of this test.

\section{An exact normalization test for a positive source quotient}
\label{sel12:metric}

Let $K_1,\ldots,K_s$ be the actual operators on a finite Hilbert space $H$,
with $\sum_jK_j^*K_j=I$.  Write
$\Phi^*(X)=\sum_jK_j^*XK_j$.  Suppose a proposed source quotient has metric
$G\succeq0$, $G\ne0$, and let $p=\supp G$.  In canonical coordinates its
synthesis is $q_G=G^{1/2}:H\to pH$; any other synthesis with the same Gram
is unitarily equivalent on its image.

The demand that the old coefficients descend means that there are operators
$\widetilde K_j$ on $pH$ satisfying
\begin{equation}
 \widetilde K_jG^{1/2}=G^{1/2}K_j
 \quad\text{on all of }H.
 \label{sel12:eq:descent}
\end{equation}
This requires an actual action of the $K_j$ on the domain of $G$.  A source
Gram on a coefficient space, without that action, is not an instance of this
problem.

\begin{theorem}[Harmonic-metric characterization of normalized descent]
\label{sel12:thm:metric-descent}
There exist operators satisfying \eqref{sel12:eq:descent} and
$\sum_j\widetilde K_j^*\widetilde K_j=I_{pH}$ if and only if
\begin{equation}
 \boxed{\Phi^*(G)=G.}
 \label{sel12:eq:harmonic-metric}
\end{equation}
They are then unique and are given by
\begin{equation}
 \boxed{\widetilde K_j=G^{1/2}K_jG^{\dagger/2}\big|_{pH}.}
 \label{sel12:eq:metric-coefficients}
\end{equation}
In particular, \eqref{sel12:eq:harmonic-metric} forces
$\Ker G$ to be invariant under every $K_j$.
\end{theorem}
\begin{proof}
For $v\in\Ker G$, harmonicity gives
$0=\langle v,Gv\rangle=\sum_j\|G^{1/2}K_jv\|^2$.
Each summand vanishes, proving invariance.  Equation
\eqref{sel12:eq:metric-coefficients} therefore satisfies
\eqref{sel12:eq:descent}, and
\[
 \sum_j\widetilde K_j^*\widetilde K_j
 =G^{\dagger/2}\Phi^*(G)G^{\dagger/2}\big|_{pH}=I_{pH}.
\]
Conversely, multiply the normalization of the descended coefficients on both
sides by $G^{1/2}$ and use \eqref{sel12:eq:descent} to obtain harmonicity.
Surjectivity of $G^{1/2}:H\to pH$ proves uniqueness.
\end{proof}

\begin{proposition}[Separate well-definedness and normalization defects]
\label{sel12:prop:metric-defect}
Put
\[
 b_G=\sum_j\|G^{1/2}K_j(I-p)\|_{\HS}^2.
\]
The old coefficients descend as linear maps exactly when $b_G=0$.  On that
branch their normalization defect is exactly
\begin{equation}
 \boxed{\sum_j\widetilde K_j^*\widetilde K_j-I_{pH}
 =G^{\dagger/2}(\Phi^*(G)-G)G^{\dagger/2}\big|_{pH}.}
 \label{sel12:eq:metric-defect}
\end{equation}
If $G\succeq g p$ with $g>0$, its operator norm is at most
$\|\Phi^*(G)-G\|/g$.  Without $b_G=0$, the same displayed matrices still
exist, but do not implement \eqref{sel12:eq:descent} on the old null space.
\end{proposition}
\begin{proof}
The nonnegative sum is zero exactly when each old null vector has a null
image.  This is the usual well-definedness condition for a quotient map.
The sum of the squared matrices in \eqref{sel12:eq:metric-coefficients} gives
\eqref{sel12:eq:metric-defect}.  The spectral bound on $G$ bounds its inverse
square root by $g^{-1/2}$ on the support.
\end{proof}

\begin{corollary}[Recurrence restricts normalized source shorts]
\label{sel12:cor:recurrent-metric}
If $\Phi_*$ has a faithful stationary state, normalized descent through $G$
is possible exactly when
\begin{equation}
 \boxed{G\in C^*(I,K_1,\ldots,K_s)'.}
 \label{sel12:eq:metric-commutant}
\end{equation}
Consequently $p$ reduces the original coefficient bank, and its descended
operators are simply the restrictions $K_j|_{pH}$.  If the original coefficient
algebra is $\B(H)$, the only nonzero admissible metrics are $G=cI$, $c>0$;
no proper rank-reducing quotient satisfies the criterion.
\end{corollary}
\begin{proof}
Under faithful stationarity the fixed space of the unital adjoint channel is
its Kraus commutant, as proved in the inherited faithful Dirichlet theorem.
Apply that result to the Hermitian fixed operator $G$ in
\eqref{sel12:eq:harmonic-metric}.  Functional calculus makes $p$ and
$G^{\dagger/2}$ commute with every $K_j$, yielding the restriction formula.
The final case uses the scalar commutant of a full matrix algebra.
\end{proof}

\begin{remark}[What this does not change in the geometry theorem]
The manuscript's $G_{\rm int\mid ST}=S_{\rm int}^*(I-P_{\rm ST})S_{\rm int}$
is a positive source metric and has its stated rank regardless of
\eqref{sel12:eq:harmonic-metric}.  The latter is an additional test only when
that quotient is used as the domain of the \emph{same normalized amplitude
process}.  Even then, $\rho\mapsto G^{1/2}\rho G^{1/2}$ is not generally a
trace-preserving preparation map: it tracks $\Tr(G\rho)$.  An identification
with the original probability law must retain its physical source metric and
preparation map.  A harmonic reweighting is not automatically a lossless
operational quotient.
\end{remark}

\section{Covariance gives a different, genuinely trace-preserving reduction}
\label{sel12:covariant-reduction}

Let a finite group act on the actual system as
\[
 H=\bigoplus_\lambda V_\lambda\otimes M_\lambda,
 \qquad U_g=\bigoplus_\lambda u_g^\lambda\otimes I,
 \qquad d_\lambda=\dim V_\lambda.
\]
The inequivalent $V_\lambda$ are external irreducibles.  Define the abstract
multiplicity operator algebra and two state maps by
\begin{align}
 \mathfrak M&=\bigoplus_\lambda\B(M_\lambda),\nonumber\\
 \mathsf R_*(X)&=\bigoplus_\lambda\Tr_{V_\lambda}(p_\lambda Xp_\lambda),\nonumber\\
 \mathsf S_*(x)&=\bigoplus_\lambda\frac{I_{V_\lambda}}{d_\lambda}\otimes x_\lambda.
 \label{sel12:eq:RS}
\end{align}
We use the ordinary sum of traces on $\mathfrak M$.
Both maps are CPTP, $\mathsf R_*\mathsf S_*=\id$, and
$\mathsf S_*\mathsf R_*=\mathfrak T$, the group twirl on $\B(H)$.
The adjoint of $\mathsf R_*$ is the unital embedding
\[
 \mathsf i(x)=\bigoplus_\lambda I_{V_\lambda}\otimes x_\lambda,
 \qquad\Ran\mathsf i=U(G)'.
\]
These are standard conditional-expectation factors, included to fix the
physical normalization and the erased variables \cite{sel12-GT}.

\begin{theorem}[Autonomous averaged multiplicity channel]
\label{sel12:thm:coarse-descent}
For every covariant CPTP map $\Phi_*$ on $H$, the map
\begin{equation}
 \boxed{\Psi_*:=\mathsf R_*\Phi_*\mathsf S_*}
 \label{sel12:eq:reduced-channel}
\end{equation}
is CPTP on $\mathfrak M$ and obeys
\begin{equation}
 \boxed{\mathsf R_*\Phi_*^n=\Psi_*^n\mathsf R_*\quad(n\ge0).}
 \label{sel12:eq:coarse-autonomy}
\end{equation}
The identities hold for all inputs and all untouched auxiliary systems.  In
particular no resetting intervention is required to reproduce the averaged
multiplicity evolution.  The independently calibrated opportunity map
$(1-\vartheta)\id+\vartheta\Phi_*$ descends to
$(1-\vartheta)\id+\vartheta\Psi_*$ with the same $\vartheta$.
\end{theorem}
\begin{proof}
Covariance gives $\mathfrak T\Phi_*=\Phi_*\mathfrak T$.  Since
$\mathsf R_*\mathfrak T=\mathsf R_*$, one has
\[
 \mathsf R_*\Phi_*
 =\mathsf R_*\mathfrak T\Phi_*
 =\mathsf R_*\Phi_*\mathsf S_*\mathsf R_*
 =\Psi_*\mathsf R_*.
\]
Induction proves the power identity.  Composition of the displayed CPTP maps
proves normalization and complete positivity, including every amplification.
Linearity proves the opportunity statement.
\end{proof}

This theorem does not identify the entire algebra $\mathfrak M$ with the
historically generated internal action.  The latter is still
$\mathcal I_K=C^*(I,K_e)\cap U(G)'$, and may be smaller than
$\Ran\mathsf i$.  A source-algebra saturation test remains necessary.
Nor is \eqref{sel12:eq:RS} the geometry Schur quotient of the preceding
section: it erases external representation factors by a CP state map,
retaining all multiplicity blocks as a direct-sum state algebra.

\subsection{The exact additional condition for retaining the original outcomes}

Let $\{\Phi_{a,*}\}_a$ be an actual instrument with sum $\Phi_*$.  Say that
its original record descends \emph{autonomously through $\mathsf R_*$} when
there are CP maps $\Psi_{a,*}$ such that
\begin{equation}
 \mathsf R_*\Phi_{a,*}=\Psi_{a,*}\mathsf R_*\quad\text{for every }a.
 \label{sel12:eq:record-autonomy}
\end{equation}
This is a fixed-state, memoryless reduction at the declared cut, for all
inputs.  If it holds, composition preserves every original labelled cylinder
and its terminal multiplicity state.

\begin{theorem}[Record-respecting descent and a finite observable defect]
\label{sel12:thm:record-descent}
The following conditions are equivalent:
\begin{enumerate}[label=\textup{(\roman*)},leftmargin=2.7em]
\item \eqref{sel12:eq:record-autonomy} holds;
\item $\Phi_a^*(\mathsf i(\mathfrak M))\subseteq\mathsf i(\mathfrak M)$ for every $a$;
\item for a Hermitian basis $b_1,\ldots,b_q$ of $\mathfrak M$,
\begin{equation}
 \boxed{\mathfrak d_{\rm rec}
  :=\sum_{a,j}\|(I-\mathfrak T)\Phi_a^*(\mathsf i(b_j))\|_{\HS}^2=0.}
 \label{sel12:eq:record-defect}
\end{equation}
\end{enumerate}
On that branch the descended maps are unique,
$\Psi_{a,*}=\mathsf R_*\Phi_{a,*}\mathsf S_*$, and form a normalized
instrument.  If a summand of \eqref{sel12:eq:record-defect} is nonzero, there
are two input states with the same reduced state but different branch-weighted
terminal multiplicity expectation.  Thus no choice of memoryless reduced maps
can reproduce all such data.
\end{theorem}
\begin{proof}
Take trace duals in \eqref{sel12:eq:record-autonomy} to obtain
$\Phi_a^*\mathsf i=\mathsf i\Psi_a^*$.  Conversely the invariant-range
condition and the left inverse of $\mathsf i$ give this identity and its
state dual.  Composing on the right with $\mathsf S_*$ proves uniqueness.
Complete positivity follows from the explicit composition.  Its sum is CPTP
by Theorem~\ref{sel12:thm:coarse-descent}.  The twirl is the HS orthogonal
projection onto the invariant algebra, giving equivalence with the finite
sum.

For a nonzero summand let
$Y=(I-\mathfrak T)\Phi_a^*(\mathsf i(b_j))$; it is Hermitian, traceless,
and $\mathsf R_*(Y)=0$.  For small $\epsilon>0$,
$\rho_\pm=I_H/\dim H\pm\epsilon Y$ are density matrices with identical
reduction.  Their expectation difference is
$2\epsilon\Tr[Y\Phi_a^*(\mathsf i(b_j))]=2\epsilon\|Y\|_{\HS}^2>0$.
\end{proof}

\begin{corollary}[Transitive labels look uniformly random after resetting]
\label{sel12:cor:uniform-projected}
If a group acts transitively on $n$ original coefficient labels and their
instrument is covariant, then
\begin{equation}
 \boxed{\mathsf R_*\Phi_{e,*}\mathsf S_*=\Psi_*/n\quad\text{for every }e.}
 \label{sel12:eq:uniform-projected}
\end{equation}
These maps reproduce the reduced one-step output for invariant input states,
but their compositions reproduce the original resolved histories only if the
corresponding descent identities hold on the relevant inputs.
\end{corollary}
\begin{proof}
Both $\mathsf R_*$ and $\mathsf S_*$ intertwine the trivial group action on
$\mathfrak M$ with the system conjugation action.  Covariance therefore makes
the left side constant along the label orbit.  Summing gives $\Psi_*$.
Iteration without \eqref{sel12:eq:record-autonomy} inserts
$\mathsf S_*\mathsf R_*=\mathfrak T$ between actual events, which need not
leave their resolved probabilities invariant.
\end{proof}

\section{A return cannot erase an external quantum factor in an efficient branch}
\label{sel12:rank-bound}

Let $\mathcal J_*$ be a nonzero CP branch from $V\otimes M$ to
$W\otimes N$, of Choi rank $k$, where $d=\dim V$ and $b=\dim W$.
Suppose its reduced output depends only on the reduced input:
\begin{equation}
 \Tr_W\mathcal J_*(X)=\mathcal T_*(\Tr_VX).
 \label{sel12:eq:return-autonomy}
\end{equation}
Let $r=\rank J(\mathcal T_*)$.

\begin{theorem}[Sharp branch-rank capacity for autonomous returns]
\label{sel12:thm:rank-capacity}
Equation~\eqref{sel12:eq:return-autonomy} implies
\begin{equation}
 \boxed{kb\ge dr.}
 \label{sel12:eq:rank-capacity}
\end{equation}
The inequality is sharp: for any CP trace-nonincreasing $\mathcal T_*$ with
Choi rank $r$, such a branch with at most $k$ Kraus operators can be constructed
whenever $kb\ge dr$.
In particular a nonzero one-coefficient return from $d>1$ to an external
trivial output ($b=1$) cannot descend autonomously after erasing the incoming
external factor.
\end{theorem}
\begin{proof}
A $k$-coefficient representation of $\mathcal J_*$ followed by tracing $W$
has at most $kb$ coefficients.  In input ordering $V\otimes M$, the Choi
matrix of the right side of \eqref{sel12:eq:return-autonomy} is, up to a
fixed tensor permutation, $I_V\otimes J(\mathcal T_*)$, of rank $dr$.
Choi rank equals minimum coefficient number, proving the bound.

For sharpness take a minimal list $T_\ell$ of $\mathcal T_*$ and orthonormal
vectors $v_{i\ell}$ in $\C^k\otimes W$, possible because $kb\ge dr$.
Define
\[
 A(|i\rangle\otimes x)=\sum_{\ell=1}^{r}v_{i\ell}\otimes T_\ell x
       \quad\text{in }\C^k\otimes W\otimes N.
\]
Taking its $k$ components produces the required branch.  Orthogonality gives
$A^*A=I_V\otimes\mathcal T^*(I_N)\preceq I$, and tracing the first two
output factors proves \eqref{sel12:eq:return-autonomy}.
\end{proof}

The bound concerns the rank of the actual coarse CP branch, not an instruction
to introduce unobserved physical outcomes.  Coarsening original recorded
outcomes can increase that Choi rank.  An enlarged retained memory can also
change $b$; neither operation is the same as claiming that the old efficient
branch already acts on the smaller multiplicity state alone.

\section{The coherent six-edge source admits only edge-blind autonomous records}
\label{sel12:edge-blindness}

Return to the fully verified $S_4$ source of V10 with $r=3$; do not substitute
the genuinely oriented source or its different reversal action.  Write
$f_{\lambda a}(e)$ for the real orthonormal contrast columns of its original
six-label permutation representation, $\lambda=W,2$, and retain
\begin{equation}
 K_e=\frac{I}{\sqrt{18}}+
           \sqrt{\frac23}\sum_{\lambda,a}f_{\lambda a}(e)L_{\lambda a}.
 \label{sel12:eq:source-frame}
\end{equation}
From \eqref{sel10:eq:entries} and \eqref{sel10:eq:exits},
\begin{align}
 p_0K_ep_0&=I_{p_0}/\sqrt{18},\nonumber\\
 p_0K_ep_\lambda
 &=\sqrt{\frac23}\,S_\lambda
       (\langle f_\lambda(e)|\otimes I),\qquad\lambda=W,2.
 \label{sel12:eq:source-corners}
\end{align}
The other nontrivial input type cannot enter the trivial output through the
five-dimensional contrast environment.  The $S_\lambda$ here are precisely
the inherited forward return blocks; they are not adjoints added as controls.

For a classical postprocessing row $x=(x_e)\in[0,1]^6$, define
$\Phi_{x,*}=\sum_ex_eK_e(\cdot)K_e^*$.  A complete postprocessing consists
of such rows $x^{(y)}$ with $\sum_yx_e^{(y)}=1$ for every old edge $e$.

\begin{theorem}[Exact edge-blindness of multiplicity-only records]
\label{sel12:thm:edge-blindness}
Assume $S_W\ne0$ and $S_2\ne0$.  The coarse branch $\Phi_{x,*}$ descends
through the external-multiplicity state map $\mathsf R_*$ if and only if
\begin{equation}
 \boxed{x_e=\overline x:=\frac16\sum_fx_f\quad\text{for every }e.}
 \label{sel12:eq:edge-blind}
\end{equation}
Consequently every autonomous classical postprocessing of the original
six-edge instrument is edge blind:
$\Phi_{y,*}=p_y\Phi_*$, $p_y\ge0$, $\sum_yp_y=1$.
The assumptions on the two return blocks follow in particular from faithful
stationarity on the complete V10 carrier.
\end{theorem}
\begin{proof}
Let $e_0$ denote the identity in the trivial summand of $\mathfrak M$ and
zero in the others, so $\mathsf i(e_0)=p_0$.  Define
$\xi_\lambda(x)=\sum_ex_ef_\lambda(e)$.  Direct substitution in
\eqref{sel12:eq:source-corners} gives the cross input blocks
\begin{equation}
 \boxed{p_0\Phi_x^*(p_0)p_\lambda
 =\frac1{3\sqrt3}S_\lambda(\langle\xi_\lambda(x)|\otimes I).}
 \label{sel12:eq:blind-cross}
\end{equation}
If the branch descends, Theorem~\ref{sel12:thm:record-descent} requires these
blocks to vanish: an invariant observable has no cross terms between
inequivalent irreducibles.  Since $S_\lambda\ne0$, its tensor product with
a nonzero row cannot vanish; therefore $\xi_W=\xi_2=0$.  The five contrast
columns form an orthonormal basis of $\mathbf1^\perp$, so this is exactly
\eqref{sel12:eq:edge-blind}.  Conversely a constant row is a scalar multiple
of the covariant total channel and descends by
Theorem~\ref{sel12:thm:coarse-descent}.  Apply this rowwise to a complete
postprocessing.  The stationarity assertion is the inherited
Proposition~\ref{sel10:prop:recurrent-return}.
\end{proof}

\begin{remark}[This is stronger than a Kraus-rank obstruction]
An isolated return can be made autonomous by grouping several efficient
outcomes.  The full source also contains its coherently calibrated identity
amplitude.  Its interference with the same return blocks is exactly
\eqref{sel12:eq:blind-cross}; it forces \emph{all five} edge contrasts to
vanish in any surviving postprocessed label.  Removing the identity fraction
as if it were a separately resolved event would change the original
instrument.  The actual no-response event is separate and can remain recorded;
the result also applies separately on a retained sign--type record block.
It does not cancel the accepted branch's cross blocks.
\end{remark}

\begin{proposition}[A calibrated lower bound on a failed descent]
\label{sel12:prop:blind-error}
For any real row $x$, put
\[
 \beta(x)^2=2\sum_{\lambda=W,2}
               \|p_0\Phi_x^*(p_0)p_\lambda\|_{\HS}^2.
\]
Then
\begin{equation}
 \boxed{\beta(x)^2=\frac2{27}\sum_{\lambda=W,2}
              \|S_\lambda\|_{\HS}^2\|\xi_\lambda(x)\|^2
 \ge\frac{2s_*}{27}\|x-\overline x\mathbf1\|^2,}
 \label{sel12:eq:blind-HS-bound}
\end{equation}
where $s_*:=\min(\|S_W\|_{\HS}^2,\|S_2\|_{\HS}^2)>0$.
For every candidate linear reduced branch $\widehat\Psi_{x,*}$,
\begin{equation}
 \boxed{\|\mathsf R_*\Phi_{x,*}-\widehat\Psi_{x,*}\mathsf R_*\|_\diamond
 \ge\frac1{3\sqrt3}\max_{\lambda=W,2}
                         \|S_\lambda\|_\op\|\xi_\lambda(x)\|.}
 \label{sel12:eq:blind-diamond-bound}
\end{equation}
At actual accepted opportunity weight $\vartheta$, the map bound is
multiplied by $\vartheta$, and the squared moment defect by $\vartheta^2$.
\end{proposition}
\begin{proof}
The HS norm of a tensor product factors.  The orthonormal contrast frame gives
$\|\xi_W\|^2+\|\xi_2\|^2=\|x-\overline x\mathbf1\|^2$; this proves
\eqref{sel12:eq:blind-HS-bound}.  Take the dual difference at $e_0$.
The term $\mathsf i\widehat\Psi_x^*(e_0)$ has zero cross-isotypic blocks,
so its $0,\lambda$ compression remains \eqref{sel12:eq:blind-cross}.
The operator norm of a compression does not exceed the whole norm.  Duality
of induced trace and operator norms, bounded above by the diamond norm,
proves \eqref{sel12:eq:blind-diamond-bound}.  Scaling the original CP branch
proves the final assertion.
\end{proof}

\subsection{The minimum coherent information forced by one-step readbacks}

Let
\[
 \mathcal O_1=C^*\bigl(U(S_4)',\Phi_e^*(U(S_4)'):e=1,\ldots,6\bigr).
\]
This is an observable algebra generated by native terminal multiplicity
queries and their one-step old-record pullbacks, not an added bank of controls.
An autonomous reduction through a represented observable algebra containing
all multiplicity queries must contain $\mathcal O_1$ by
Theorem~\ref{sel12:thm:record-descent}.

\begin{theorem}[Native returns force an eleven-dimensional coherent predictor sector]
\label{sel12:thm:predictor-core}
For every V10 source with $S_W\ne0$ and $S_2\ne0$, set
$H_c=(p_0+p_W+p_2)H$, of dimension $3+6+2=11$.  Then
\begin{equation}
 \boxed{\B(H_c)\oplus\C p_{\rm s}\ \subseteq\ \mathcal O_1.}
 \label{sel12:eq:predictor-core}
\end{equation}
Thus a multiplicity-only state on the abstract algebra
$M_3\oplus M_2\oplus\C^2$ cannot preserve all original edge-conditioned
internal expectations.  Any exact reduction of this kind must retain a
faithful representation of the full coherent $M_{11}$ sector.
\end{theorem}
\begin{proof}
The central system projections $p_\lambda$ belong to $U(S_4)'$ and hence
can be used inside the observable algebra.  By
\eqref{sel12:eq:blind-cross}, the cross operators
$p_0\Phi_e^*(p_0)p_\lambda$ have the form
$cS_\lambda(\langle f_\lambda(e)|\otimes I)$, with $c\ne0$.
The rows $f_\lambda(e)$ span $V_\lambda^*$.
Left multiplication by $\B(p_0H)$ and right multiplication by
$I_{V_\lambda}\otimes\B(M_\lambda)$ turn any nonzero $S_\lambda$ into
all maps $M_\lambda\to p_0H$: to see this, sandwich one nonzero matrix entry
between matrix units.  Hence all maps $p_\lambda H\to p_0H$ belong to
$\mathcal O_1$.  Taking their adjoints and products gives every matrix unit
on $H_c$.  The remaining unit $p_{\rm s}$ is already one of the external
central projections.  This proves the supported inclusion.
\end{proof}

This is a restriction on ordinary algebraic state reduction preserving the
specified internal outputs for all input states.  It is not a lower bound
for arbitrary classical simulation, for nonlinear history-dependent
estimators, or for a deliberately restricted preparation class.

\section{What can be retained when only a return is queried?}
\label{sel12:return-groups}

The preceding no-go must not be applied to a more limited query without
checking its hypotheses.  Isolate a nonzero return from type $\lambda=W,2$
to the trivial output.  Such a component is an actual selected operation only
if the input-domain and terminal sector selection are physically available.
Otherwise the following is an exact component calculation of the channel.

Set $d=d_\lambda$, $R_\lambda=\sqrt d\,S_\lambda$, and
$v_{\lambda,e}=\sqrt{6/d}\,f_\lambda(e)$, which is a unit vector.  The
six original return coefficients are
\begin{equation}
 B_e=p_0K_ep_\lambda
       =\frac13R_\lambda(\langle v_{\lambda,e}|\otimes I).
 \label{sel12:eq:return-coeff}
\end{equation}
For a row $x\ge0$, put
$F_{\lambda,x}=\sum_ex_e|v_{\lambda,e}\rangle\langle v_{\lambda,e}|$.

\begin{theorem}[Exact return-only coarsenings]
\label{sel12:thm:return-groups}
For $R_\lambda\ne0$, the return map
$\sum_ex_eB_e(\cdot)B_e^*$ depends only on $\Tr_{V_\lambda}X$ if and
only if $F_{\lambda,x}$ is scalar.  In the six-edge tetrahedral frame this
is equivalent, for either $\lambda=W$ or $\lambda=2$, to
\begin{equation}
 \boxed{x_{e_1}+x_{\bar e_1}
       =x_{e_2}+x_{\bar e_2}
       =x_{e_3}+x_{\bar e_3},}
 \label{sel12:eq:pair-sum}
\end{equation}
where $(e_k,\bar e_k)$ are the three complementary edge pairs.
A nonempty proper deterministic subset therefore works exactly when it
contains one edge from each pair.  There are eight such subsets, the four
vertex stars and their four complementary triangles.  They yield four
unordered two-bin partitions.  The minimum nonzero deterministic bin has
three original edges.
\end{theorem}
\begin{proof}
For a product input $Z\otimes\rho$, the output is
$(\Tr F_{\lambda,x}Z)R_\lambda\rho R_\lambda^*/9$.
If it depends only on $\Tr Z\rho$, a $\rho$ with nonzero output shows that
$\Tr(F_{\lambda,x}Z)$ is a fixed multiple of $\Tr Z$ for every $Z$.
Thus $F_{\lambda,x}$ is scalar.  Conversely a scalar frame gives the
claimed dependence by linearity, including entangled inputs.

For $W$, the two vectors of each complementary pair are opposite vectors
on one of three orthogonal axes, so $F_{W,x}$ is diagonal with the three
pair sums.  For $V_2$, complementary edges have the same rank-one
projection and the three distinct directions form a real trine.  Their three
rank-one projectors are linearly independent in the real symmetric
$2\times2$ matrices and have sum $3I_2/2$.  Hence their weighted sum is
scalar exactly when the three weights are equal.  This proves
\eqref{sel12:eq:pair-sum}.  For $0$--$1$ rows the common pair sum is
$0,1$, or $2$; the only proper nonempty case is $1$, proving all counts.
\end{proof}

For a three-edge bin the return Choi rank is $d$, saturating
\eqref{sel12:eq:rank-capacity} for its reduced rank-one map.  In the $V_2$
case the three resolved coefficients are linearly dependent; the physical
number of grouped outcomes and the minimum Choi coefficient number are
therefore different.  On an unweighted bin the reduced return is
$R_W(\cdot)R_W^*/9$ for $W$ and
$R_2(\cdot)R_2^*/6$ for $V_2$.

\begin{remark}[A covariant coarse record without a chosen vertex]
A fixed star--triangle partition singles out a vertex.  It need not be
introduced as a new source mark: one can instead choose a vertex uniformly
by independent classical randomization and report that vertex together with
whether the old edge belongs to its star.  There are eight reported outcomes,
each with row $x_e=1/4$ on its star or triangle and zero otherwise.  The rows
sum to one for every original edge and transform covariantly under $S_4$.
Every \emph{return component} descends.  The full source still fails
\eqref{sel12:eq:edge-blind}: return-only autonomy does not preserve the
coherent identity--return interference.
\end{remark}

\begin{proposition}[Exact cost of erasing the external input of one return]
\label{sel12:prop:single-return-error}
For one original return \eqref{sel12:eq:return-coeff}, its canonical
multiplicity approximation is
$\overline{\mathcal J}_{e,*}(\sigma)=R_\lambda\sigma R_\lambda^*/(9d)$.
For $d>1$,
\begin{equation}
 \boxed{\|\mathcal J_{e,*}
       -\overline{\mathcal J}_{e,*}\Tr_{V_\lambda}\|_\diamond
 =\frac19\left(1-\frac1d\right)\|R_\lambda\|_\op^2.}
 \label{sel12:eq:single-return-error}
\end{equation}
The best possible approximation by any map depending only on
$\Tr_{V_\lambda}X$ has diamond error
$\|R_\lambda\|_\op^2/18$.  This optimum is attained by the CP map
$\sigma\mapsto R_\lambda\sigma R_\lambda^*/18$, but generally does not
assemble with the other bins into the original normalized reduced instrument.
\end{proposition}
\begin{proof}
The difference is $1/9$ times the functional
$Z\mapsto\Tr[(|v\rangle\langle v|-I/d)Z]$ tensored with
$\rho\mapsto R_\lambda\rho R_\lambda^*$.  A Hermitian functional has
completely bounded trace norm equal to the operator norm of its representing
matrix.  The CP factor has diamond norm $\|R_\lambda\|^2$; a product
input achieving both norms proves equality.  For the optimization, compare
$|v\rangle\langle v|\otimes\rho$ with an orthogonal external pure state
and the same maximizing $\rho$.  Their reduced inputs agree, while the
true outputs differ by trace norm $\|R_\lambda\|^2/9$.  The triangle
inequality gives half that difference as a lower bound for any approximation.
The functional $|v\rangle\langle v|-I/2$ has operator norm $1/2$, proving
attainment by the displayed CP map.
\end{proof}

\section{Evaluate the unchanged full-source benchmark}
\label{sel12:example}

No new favourable channel is required to see the distinction.  Use exactly
V10's source \eqref{sel10:eq:witness-six} with
\[
 R=\begin{pmatrix}0&0&1\\1&0&0\\0&1&0\end{pmatrix},
 \qquad T_W=(I_2\ \ 0),\quad T_2=(0\ \ 0\ \ 1),
 \qquad S_W=R_{[:,1:2]}/\sqrt3,\quad S_2=R_{[:,3]}/\sqrt2.
\]
The source is normalized, unital and covariant, and its invariant algebra
saturates the target.  These inherited properties are not re-established as
new selection theorems.

\begin{proposition}[The native trivial projection is not a normalized source short]
\label{sel12:prop:benchmark-metric}
For this source,
\begin{equation}
 \Phi^*(p_0)-p_0
 =-\tfrac23 I_3\oplus\tfrac29 I_6\oplus\tfrac13 I_2\oplus0_3,
 \qquad\|\Phi^*(p_0)-p_0\|_{\HS}^2=\frac{50}{27}.
 \label{sel12:eq:benchmark-metric}
\end{equation}
Thus $G=p_0$ fails normalized quotient descent.  Simple compression onto the
trivial system input loses probability $2/3$ in one accepted step, whereas
keeping all multiplicity sectors by \eqref{sel12:eq:RS} gives a normalized
channel.
\end{proposition}
\begin{proof}
The mean coefficient contributes $p_0/3$ to $\Phi^*(p_0)$.
The inherited return blocks contribute
$(2/3)I_W\otimes S_W^*S_W=(2/9)I_6$ and
$(2/3)I_{V_2}\otimes S_2^*S_2=(1/3)I_2$; the remaining sector contributes
zero.  This proves the displayed matrix and its squared norm.  The input
$p_0$ retention probability is $1/3$ by the native transition theorem.
This projection is a diagnostic example, not an identification with the
historical geometry Schur projection.
\end{proof}

\begin{theorem}[A two-step old-record witness against repeated internal resetting]
\label{sel12:thm:two-step}
Prepare the native trivial-sector state $\rho=|3\rangle\langle3|$ in the
benchmark.  Label the edges in three complementary pairs as in V10's
orthonormal contrast frame.  Its original two-edge probabilities are
\begin{equation}
 \boxed{\Pr(e,f\mid\rho)=
 \begin{cases}
 23/324,&e,f\text{ belong to the same complementary pair},\\
 1/162,&\text{otherwise}.
 \end{cases}}
 \label{sel12:eq:old-pairs}
\end{equation}
The canonical memoryless instrument
$\Psi_{e,*}=\mathsf R_*\Phi_{e,*}\mathsf S_*=
\Psi_*/6$ instead predicts $1/36$ for every ordered pair.  The two classical
laws therefore have total variation distance
\begin{equation}
 \boxed{\operatorname{TV}(\Pr,\widehat\Pr)=\frac{14}{27}.}
 \label{sel12:eq:TV}
\end{equation}
For any fixed $e$, the additional terminal trivial-sector event has respective
probabilities $5/324$ and $1/108$ after $(e,e)$.
At two actual locked opportunities, all displayed pair masses carry
$\vartheta^2$.  Summing over old edge labels restores the exact averaged
multiplicity evolution, as required by Theorem~\ref{sel12:thm:coarse-descent}.
\end{theorem}
\begin{proof}
For completeness the calculation needs only the $V_2$ contrasts of the
unchanged benchmark.  On complementary pair $k$, write the corresponding row
as
\[
 f^{(1)}=(1/2,1/\sqrt{12}),\quad
 f^{(2)}=(-1/2,1/\sqrt{12}),\quad
 f^{(3)}=(0,-1/\sqrt3).
\]
Both edges in a pair have this same row.  Set $a_0=1/\sqrt{18}$,
$H_f=f_1X+f_2Z$ in the $V_2$ block, and let $t=|3\rangle$,
$u=Rt=|1\rangle$.  The only nonzero blocks of $K_fK_et$ are
\begin{align*}
 p_0K_fK_et&=a_0^2t+\frac{\langle f,e\rangle}{3}u,\\
 p_2K_fK_et&=\frac{a_0}{\sqrt3}(e+f)
                      +\frac{\sqrt2}{6}H_fe.
\end{align*}
Here $e,f$ in these two formulas denote the two contrast rows, not their label
vectors.  Their squared norms sum to $23/324$ when the rows agree and
$1/162$ for two different rows, by direct substitution of the three displayed
rows and the Pauli matrices.  When the rows agree the first norm is $5/324$.
There are twelve same-pair ordered outcomes and twenty-four other outcomes;
$12(23/324)+24(1/162)=1$.

Equation~\eqref{sel12:eq:uniform-projected} and trace preservation of
$\Psi_*$ give $1/36$ for all projected pairs.  The sum of the twelve positive
probability differences is
$12(23/324-1/36)=14/27$, proving the TV formula.  After one projected step,
the multiplicity state is $\rho/18$ in the trivial block and $1/9$ in the
$V_2$ block.  Its next same-edge trivial output has trace
$1/324+1/162=1/108$.  Finally, each actual accepted letter carries
$\sqrt\vartheta$, so a two-letter probability carries $\vartheta^2$.
\end{proof}

\begin{proposition}[Exact minimal ordinary predictor algebra of the benchmark]
\label{sel12:prop:benchmark-predictor}
For the same V10 benchmark,
\begin{equation}
 \boxed{\mathcal O_1=M_{11}(\C)\oplus\mathcal D_3,\qquad
       \dim_\C\mathcal O_1=124.}
 \label{sel12:eq:benchmark-predictor}
\end{equation}
It is invariant under every original $\Phi_e^*$ and is the smallest unital
$*$-subalgebra containing $U(S_4)'$ with that invariance property.  Its
conditional-expectation reduction is therefore an exact model of every old
edge record and terminal internal query.  A faithful Hilbert representation
has dimension at least $11+1+1+1=14$, compared with seven for the abstract
multiplicity algebra; within this reduction class there is no Hilbert-dimension
saving for the benchmark.
\end{proposition}
\begin{proof}
Theorem~\ref{sel12:thm:predictor-core} supplies the supported $M_{11}$.
In V10's explicit frame, on the remaining $W^{\rm s}$ block the two edges
in complementary pair $k$ have
$K_e|_{W^{\rm s}}=E_{kk}/\sqrt2$.
Indeed, with the traceless diagonal basis $z_a$ of that construction,
$\sum_a(z_a)_{kk}(z_a)_{jj}=\delta_{kj}-1/3$, and its scalar term cancels
the $-1/3$ part.  Therefore
$p_{\rm s}\Phi_e^*(p_{\rm s})p_{\rm s}=E_{kk}/2$.
These three old-record pullbacks generate $\mathcal D_3$, proving the lower
inclusion.  Every original $K_e$ is block diagonal between $H_c$ and
$W^{\rm s}$ and diagonal on the latter.  It follows that
$M_{11}\oplus\mathcal D_3$ contains the initial internal algebra and all
one-step pullbacks and is invariant under every $\Phi_e^*$.  This proves
equality, minimality and exact reducibility.  The smallest faithful
representation of a direct sum of full matrix algebras has dimension the sum
of their matrix sizes, giving fourteen.
\end{proof}

\section{What has advanced and the next upstream source comparison}
\label{sel12:status}

There are now two exact but different descent statements.  A specified
positive source short carries the old normalized coefficients precisely when
its metric is harmonic for those coefficients.  This requires their action on
the source domain to have been identified; the positivity/rank theorem alone
is not that datum.  Separately, actual covariance always gives a trace-preserving
multiplicity \emph{state} reduction.  It does not require extra sharp Reads,
a full-matrix physical compression, or a reconstructed gauge group.

The new limitation is also exact.  On the full six-edge class with both native
return ports, all autonomous postprocessings compatible with that multiplicity
state are edge blind.  An original record-preserving instrument on a smaller
state must therefore retain additional predictive information or pass a more
restricted reachable-input analysis.  This is not a contradiction to the
original invariant-algebra saturation theorem: the internal action algebra,
the averaged state algebra, and the predictive state needed for resolved
histories are different objects.

For the benchmark this extra information has been determined, not merely
postulated: its smallest invariant ordinary observable algebra is
$M_{11}\oplus\mathcal D_3$, reconstructed from one-step native readbacks.
The next source calculation is to evaluate $b_G$ and $\Phi^*(G)-G$ on an
actually identified quotient, or $\mathfrak d_{\rm rec}$ on the retained
record policy, and determine which map the historical construction supplies: a normalized
amplitude descent, an averaged multiplicity readout, or a larger
record-conditioned predictor.  If a larger predictor is required, its
additional rows must come from the original future table rather than a new
chosen control.  Algebraic colour can remain valid while such information is
needed for the source-complete mixed histories.  Determinant incidence,
charged matter types, and their physical source identification are not
consequences of either descent theorem, and historical Standard-Model
selection is not asserted.

\paragraph{Verification and preservation.}
The V11 prefix is retained byte for byte.  The accompanying exact checks test
the quotient identities in singular and nonsingular cases, the full
multi-isotype averaged reduction, the return Choi-rank bounds, all sixty-four
deterministic edge coarsenings, the edge-blindness moment matrix, and the
unchanged V10 two-step record probabilities.  V11's regression is rerun.
These finite examples check the formulas; the all-input descent, all-source
coarsening and sharp rank claims are proved in the text.  No physical source
admission, historical symmetry lift, formal proof-assistant certification, or
independent peer review is claimed.

\begin{thebibliography}{99}
\raggedright
\bibitem[20]{sel12-GT}
T.~Grigoletto and F.~Ticozzi,
\emph{Exact model reduction for discrete-time conditional quantum dynamics},
IEEE Control Systems Letters \textbf{8} (2024), 550--555;
\href{https://doi.org/10.1109/LCSYS.2024.3399100}{doi:10.1109/LCSYS.2024.3399100}.
\end{thebibliography}

% END OF SOURCE-SELECTION V12 DEVELOPMENT BLOCK.
% Preserve V1--V12 and append later development before the final terminator.


\clearpage
\begin{center}
{\Large\bfseries Source-selection continuation V13}\\[.4em]
{\large Original-record rank, capacity selection, and positive memory}\\[.25em]
A finite probability certificate for the full internal algebra
\end{center}

% The appended section numbers reach three digits; widen only these TOC labels.
\addtocontents{toc}{\protect\begingroup}
\addtocontents{toc}{\protect\renewcommand{\protect\numberline}{\protect\makebox[2.1em][l]}}

\section{Go upstream to the actual future table, not another control}
\label{sel13:context}

V12 separates the internal action algebra from the state needed to predict
resolved histories.  Its minimal ordinary predictor is not automatically a
minimum among all completely positive encodings.  More importantly, the
source's defining entrance is a table of conditional futures: before asking
for another operator-valued intervention, one should determine what the
\emph{original record probabilities themselves} can certify.

The main result of this continuation is positive.  In the verified
calibrated tetrahedral carrier category of V10, a finite original-record
Hankel rank of at least 171 forces the three-dimensional trivial multiplicity,
the full original system algebra, and hence the internal algebra
$M_3\oplus M_2\oplus\C\oplus\C$.  No separate bridge, weak projector,
private control, coherent readout, or stationary-state hypothesis enters that
implication.  A normalized point in the same source manifold gives an exact
196-by-196 certificate using one initial ray and words of length at most twelve.
It proves generic full record rank on that manifold.  It is a mathematical
witness, not a replacement for an unevaluated historical probability table.

\paragraph{Audit and nonduplication.}
The future-equivalence and Hankel constructions are already part of the
Grand-Tensor programme.  The supplied spacetime manuscript explicitly
separates deterministic predictive minimization, linear Hankel minimization,
and contextual $C^*$-history separation.  V12 of \cite{Rigidity} already
proves linear-capacity bounds for positive realizations and robust sampled
Hankel tests.  V10 of the present series already supplies the calibrated
covariant Stiefel manifold, its multiplicity capacity, and an invariant
source-algebra certificate; V12 supplies record-descent and ordinary
predictor-algebra results.  Those general mechanisms are inputs, not new
principles claimed here.  The additions are the original-record selector,
its explicit same-class full-rank witness, and the exact pointed and
positive-realization costs of the inherited benchmark.  The finite
conditional-model and positive-realization distinctions also have established
literature \cite{sel12-GT,sel13-MW,sel13-BRW}.

\paragraph{Scope retained.}
The external representation must act on the actual system, not only on a
coefficient Gram or regular word module.  A positive geometry short must pass
V12's normalized descent test before the resulting quotient is used as that
system.  The $S_4$ source class below is not the stronger reversal-even class
excluded in V10--V11.  The source chooses its physical symmetry; the rank test
does not authorize forgetting an additional constraint.  All statements
concern the previously identified edge-system factor.  Independent retained
pointer/type factors remain present in the full instrument.

\section{A finite original-record selector on the capacity-bounded carrier}
\label{sel13:selector}

Let $K_1,\ldots,K_6$ be the original resolved coefficients of one actual
source on
\begin{equation}
 S_r=\mathbf1^{\oplus r}\oplus(W\otimes\C^2)
                        \oplus V_2\oplus W^{\rm s},\qquad N=r+11,
 \label{sel13:eq:carrier}
\end{equation}
with the independently verified $S_4$ representation and calibration
\begin{equation}
 \sum_eK_e^*K_e=I,
 \qquad \sum_eK_e=\sqrt2 I,
 \qquad K_{ge}=U_gK_eU_g^*.
 \label{sel13:eq:calibration}
\end{equation}
The inherited capacity theorem \ref{sel10:thm:capacity} gives $r\le3$ and
$N\le14$.  This bound does not itself select $r=3$.

For a fixed normalized initial state $\rho$, write a word chronologically as
$w=(e_1,\ldots,e_k)$ and set
\[
 K_w=K_{e_k}\cdots K_{e_1},\quad K_\emptyset=I,\qquad
 p_\rho(w)=\Tr(K_w\rho K_w^*).
\]
Given finite past and future word lists $\mathcal V,\mathcal U$, define
\begin{equation}
 H_{u,v}=p_\rho(vu)
        =\Tr\bigl(K_u^*K_u\,K_v\rho K_v^*\bigr).
 \label{sel13:eq:record-Hankel}
\end{equation}
These are unconditional word probabilities.  They are recovered from the
complete conditional-future table by multiplying by the prefix probability.
No amplitude interference experiment or terminal quantum Read is used.

\begin{lemma}[Record rank bounds both state coordinates and the source algebra]
\label{sel13:lem:rank}
For every such panel,
\begin{equation}
 \rank H\le \min\{N^2,\dim_\C C^*(I,K_e:e)\}.
 \label{sel13:eq:rank-source}
\end{equation}
If the same probabilities have a finite time-homogeneous linear realization
of dimension $q$, then $\rank H\le q$.  A realization on an ordinary quantum
memory algebra $\mathcal B_{\boldsymbol n}=\bigoplus_bM_{n_b}(\C)$ therefore
satisfies
\begin{equation}
 \rank H\le\sum_bn_b^2\le D^2,\qquad D=\sum_bn_b.
 \label{sel13:eq:positive-capacity}
\end{equation}
Every retained classical memory flag is included in this cost.
\end{lemma}
\begin{proof}
Factor \eqref{sel13:eq:record-Hankel} through the state vectors
$K_v\rho K_v^*$ in $M_N$ and the dual observable vectors $K_u^*K_u$.
The latter belong to $C^*(I,K_e)$, giving both bounds in
\eqref{sel13:eq:rank-source}.  For a linear realization with transition
matrices $T_e$, initial vector $b$ and terminal row $a$, the same matrix
factors as $H_{u,v}=aT_uT_vb$ through its $q$ coordinates.  Vectorizing an
ordinary quantum-memory algebra gives $q=\sum_bn_b^2$, proving the last
assertion.  Complete positivity is not needed for the rank inequality, but
is part of the positive-realization class to which it is applied.
\end{proof}

\begin{lemma}[Dimension of a proper represented matrix subalgebra]
\label{sel13:lem:proper}
A proper unital $C^*$-subalgebra of $M_N(\C)$ has complex dimension at most
$(N-1)^2+1$.
\end{lemma}
\begin{proof}
A proper finite $C^*$-algebra cannot act irreducibly, by the finite
bicommutant theorem.  A nontrivial reducing projection of rank $k$ embeds it
in $M_k\oplus M_{N-k}$.  Thus its dimension is at most
$k^2+(N-k)^2\le(N-1)^2+1$ for $1\le k\le N-1$.
\end{proof}

\begin{theorem}[Original-record capacity-and-algebra selection]
\label{sel13:thm:selection}
Assume \eqref{sel13:eq:carrier}--\eqref{sel13:eq:calibration} on the actual
system.  If an original-record panel satisfies $\rank H\ge170$, then
$r=3$.  If
\begin{equation}
 \boxed{\rank H\ge171,}
 \label{sel13:eq:selection-threshold}
\end{equation}
then
\begin{equation}
 \boxed{r=3,\quad C^*(I,K_e)=M_{14}(\C),\quad
 C^*(I,K_e)\cap U(S_4)'\cong M_3\oplus M_2\oplus\C^2.}
 \label{sel13:eq:selection}
\end{equation}
Moreover the least total Hilbert dimension of any finite ordinary
quantum-memory realization of these original records is exactly fourteen.
At that minimum its ordinary algebra must be $M_{14}$.
\end{theorem}
\begin{proof}
The source capacity bounds $N$ by fourteen.  If $r\le2$, then $N\le13$ and
$\rank H\le169$.  Hence rank at least 170 forces $r=3$.  A proper
$C^*$-subalgebra of $M_{14}$ has dimension at most 170 by
Lemma~\ref{sel13:lem:proper}.  Rank at least 171 therefore forces the full
source algebra.  Its intersection with the verified external commutant is
that whole commutant; the external multiplicities $(3,2,1,1)$ give
\eqref{sel13:eq:selection}.  Any memory of total Hilbert dimension at most
thirteen has at most 169 linear coordinates.  The original source gives a
fourteen-dimensional positive realization, so the bound is attained.
A direct sum with total dimension fourteen but more than one nonzero block
has at most $13^2+1=170$ operator coordinates.  It cannot realize the panel.
\end{proof}

\begin{remark}[What the theorem selects]
This is a finite, source-category-relative selection theorem.  It replaces
an independent multiplicity-three occurrence assertion and an additional
internal saturation panel by an original-record rank test, once the
nontrivial external multiplicities, same-system normalization, and covariance
are established.  It does not derive that category from arbitrary microscopic
processes.  It also does not identify the native line and plane with the
manuscript's named type/private sources, create a determinant tensor, or
select charged matter types.  The raw algebra $M_{14}$ is not the gauge
algebra: the internal algebra is its \emph{intersection} with the external
commutant, as throughout this programme.
\end{remark}

\section{A normalized witness inside the original six-edge covariance class}
\label{sel13:witness}

Use the unchanged V10 benchmark
\eqref{sel10:eq:witness-six}, with the cyclic return
\eqref{sel10:eq:cycle}.  Denote its five contrast coefficients by
$L_{Wa}$ ($a=1,2,3$) and $L_{2b}$ ($b=1,2$), and its real six-by-five
contrast isometry by $O$.  The basis blocks are $\C^3$, $W\otimes\C^2$,
$V_2$, $W^{\rm s}$, in that order.  Let $h=(1,0)^{\mathsf T}$ be the first
existing $W$-multiplicity basis vector, and define
\[
 P=(I_W\otimes |h\rangle\langle h|)\oplus I_{W^{\rm s}},
\]
with zero action on the other two blocks.  On $W\cong\R^3$ put
$(J_a)_{bc}=\epsilon_{bac}$.  Define $C_a$ to be the Hermitian operator
whose only nonzero blocks are
\begin{equation}
 p_WC_ap_{\rm s}=\frac1{\sqrt2}J_a\otimes |h\rangle,
 \qquad p_{\rm s}C_ap_W=(p_WC_ap_{\rm s})^*,
 \qquad C_4=C_5=0.
 \label{sel13:eq:C}
\end{equation}
The tensor in the first formula means the map
$x\mapsto(J_ax)\otimes h/\sqrt2$.

\begin{proposition}[An orthogonal equivariant deformation, with no added outcome]
\label{sel13:prop:deformation}
The columns satisfy
\begin{equation}
 \sum_{a=1}^5L_a^*C_a=0,\qquad
 \sum_{a=1}^5C_a^*C_a=P.
 \label{sel13:eq:orthogonal}
\end{equation}
For real $u,v$ with $u^2+v^2=1$ and any scalar $|\zeta|=1$, let
\begin{align}
 U_\zeta&=I+(\zeta-1)p_2,\nonumber\\
 L_a(u,v,\zeta)&=\bigl[L_a(I-P+uP)+vC_a\bigr]U_\zeta,\nonumber\\
 K_e(u,v,\zeta)&=\frac{I}{\sqrt{18}}
       +\sqrt{\frac23}\sum_{a=1}^5O_{ea}L_a(u,v,\zeta).
 \label{sel13:eq:deformation}
\end{align}
These coefficients satisfy the original normalized six-outcome covariance
and coherent-sum identities for every displayed parameter.  At $(u,v,\zeta)
=(1,0,1)$ they are the unchanged benchmark.
\end{proposition}
\begin{proof}
The cross-product matrices satisfy $\sum_aJ_a^*J_a=2I$ and
$\sum_a\langle a|J_a=0$.  The first identity proves the second formula in
\eqref{sel13:eq:orthogonal}.  For the first, an original $W$-contrast can
output into $W$ only from the trivial sector, with output component
$|a\rangle\otimes T_W/\sqrt3$.  Its contraction with the new $W$ output
therefore sums to zero by the second identity.  Its other output is trivial,
where $C_a$ has no output.  There is no original $W^{\rm s}$ output in a
$W$-contrast, and the two $V_2$-contrast components of $C$ are zero.
This proves complete column orthogonality.

The old contrast column is isometric.  Thus the squared norm of the new
column before $U_\zeta$ is
$(I-P+uP)^2+v^2P=I$; right multiplication by the unitary $U_\zeta$ preserves
it.  For an orthogonal tetrahedral matrix $u_g$,
$u_gJ_au_g^{\mathsf T}=\det(u_g)\sum_b(u_g)_{ba}J_b$.
The sign twist in $W^{\rm s}$ cancels this determinant, so $C_a$ transforms
as the actual $W$ contrast.  The projections $P,p_2$ commute with the
external action.  Consequently the whole new column is equivariant in the
\emph{same} environment $W\oplus V_2$.  The inherited calibrated
contrast-isometry theorem gives the original six coefficients and both
normalizations.  No extra outcome or runtime control has been adjoined.
\end{proof}

The deformation changes the mathematical source coefficients when its
parameters change.  It is not asserted to preserve the old benchmark's
probabilities or to be an executed intervention on the historical source.
Its purpose is to establish a nonzero point of a fixed admissible parameter
space.  We use the exact point
\begin{equation}
 (u,v,\zeta)=\left(\frac35,\frac45,i\right),\qquad
 \rho_0=|3\rangle\langle3|\text{ in the trivial block}.
 \label{sel13:eq:full-point}
\end{equation}
The initial state is the same ray used in V12; no additional terminal Read
will be required.

\begin{theorem}[An exact full-rank original-record panel]
\label{sel13:thm:full-panel}
At \eqref{sel13:eq:full-point}, there are fixed lists of 196 past words and
196 future words, each of length at most six, for which
\begin{equation}
 \boxed{H_{u,v}=p_{\rho_0}(vu),\qquad \rank H=196.}
 \label{sel13:eq:full-panel}
\end{equation}
Every entry is an original six-edge word probability of length at most twelve.
The past states span $M_{14}$ and the future effects $K_u^*K_u$ span
$M_{14}$.  In particular the source meets
Theorem~\ref{sel13:thm:selection} without a mixed operator-moment readout.
\end{theorem}
\begin{proof}
A completely specified exact determinant certificate is given in
Section~\ref{sel13:certificate}; it defines the two word lists by finite
ordered elimination, so their selection contains no fitted probability or
numerical threshold.  In the residue homomorphism
\[
 \sqrt2\mapsto14,\quad\sqrt3\mapsto10,\quad i\mapsto22
 \quad\text{in }\mathbb F_{97},
\]
the determinant of the indicated \emph{physical probability} minor is 37,
which is nonzero.  All denominators are units at 97, so a vanishing
characteristic-zero determinant would have vanished under this homomorphism.
Thus the 196 columns are independent over the original number field and
over $\C$.  Factorization \eqref{sel13:eq:record-Hankel}, through the
196-dimensional space $M_{14}$, gives equality and both span assertions.
\end{proof}

\begin{corollary}[Generic original-record selection, with one fixed preparation]
\label{sel13:cor:generic}
On V10's connected 71-dimensional normalized $S_4$ source manifold on $S_3$,
keep the fixed preparation $\rho_0$ and the fixed lists from
Theorem~\ref{sel13:thm:full-panel}.  The locus $\det H\ne0$ is open, dense,
and has full smooth volume.  On it the minimum ordinary quantum-memory
realization has dimension fourteen and algebra $M_{14}$, while the invariant
internal algebra is $M_3\oplus M_2\oplus\C^2$.
\end{corollary}
\begin{proof}
The word probabilities are real polynomials in the real and imaginary entries
of the coefficients.  For the fixed lists the determinant has total degree
2956, twice the sum of all row- and column-word lengths.  Its restriction to
the source manifold is real analytic and is nonzero at
\eqref{sel13:eq:full-point}.  A nonzero real-analytic function on a connected
real-analytic manifold has a zero set of empty interior and zero smooth
volume.  For completeness, the identity theorem excludes an identically
zero germ at any point.  At a zero, choose a nonzero partial derivative of
least total order $k$.  A derivative of order $k-1$ then vanishes there and
has nonzero gradient.  Thus the zero set of the original function is
contained in the union, over all derivative multiindices, of the regular
zero sets of those derivatives.  Each is locally a hypersurface by the
implicit-function theorem and has zero volume.  A countable analytic atlas
and the countably many multiindices give the result.  Continuity gives
openness of the nonzero locus.  Theorem~\ref{sel13:thm:selection} gives the conclusions.
\end{proof}

\begin{remark}[Genericity does not populate the historical table]
This is genericity in the actual normalized covariance family under discussion,
not a probability law for microscopic Renewal sources or a dynamical basin
theorem.  In particular the lower $r=1,2$ carrier classes remain admissible,
and a special $r=3$ source can fail the record test.  The theorem proves that
an independently observed sufficient record rank selects the saturated branch
\emph{within the capacity-bounded category}.  It does not assume that the
historical record ranks have already been evaluated.
\end{remark}

\section{The unchanged benchmark: linear reduction is not algebraic reduction}
\label{sel13:benchmark}

Return now to the original V10/V12 benchmark $K_e$, not the deformation.
Write $H_c=(p_0+p_W+p_2)H\cong\C^{11}$ and use zero-based coordinate
indices $0,\ldots,13$ only in the following matrix units.  Its $V_2$ block
is at indices 9 and 10.  Define
\begin{equation}
 Y=i(E_{9,10}-E_{10,9}),\qquad
 \mathcal S_c=\{X\in\B(H_c):\Tr(YX)=0\}.
 \label{sel13:eq:Sc}
\end{equation}
Thus $\dim_\C\mathcal S_c=120$.  This is an adjoint-stable linear space,
not a subalgebra.

\begin{lemma}[A coordinate annihilated by every original branch]
\label{sel13:lem:killed}
For every original benchmark coefficient,
\begin{equation}
 \boxed{K_eYK_e^*=0,\qquad K_e^*YK_e=0.}
 \label{sel13:eq:killed}
\end{equation}
The core $H_c$ is reducing for every $K_e$.  Both the forward and backward
branch maps preserve $\mathcal S_c$.
\end{lemma}
\begin{proof}
The two columns of $K_e$ indexed by $V_2$ are collinear, and its two
$V_2$ rows are collinear.  Here is an explicit check covering all six
letters.  Put $A_e=18\sqrt2K_e$.  On the first complementary edge pair,
the $V_2$ diagonal block is
$3\left(\begin{smallmatrix}3&\sqrt3\\\sqrt3&1\end{smallmatrix}\right)$;
on the second it is
$3\left(\begin{smallmatrix}3&-\sqrt3\\-\sqrt3&1\end{smallmatrix}\right)$;
on the third it is $\diag(0,12)$.
Their coupling rows into the trivial sector are respectively proportional
to $(\sqrt6,\sqrt2)$, $(-\sqrt6,\sqrt2)$ and $(0,-2\sqrt2)$.
The input couplings have the same two-component ratios.  Hence each of the
two relevant rectangular blocks has rank one, giving
\eqref{sel13:eq:killed} by expansion of $Y$.  There are no core--$W^{\rm s}$
corners in the inherited source.  Finally
$\Tr(YK_eXK_e^*)=\Tr(K_e^*YK_eX)=0$, and the opposite identity proves
the backward assertion.
\end{proof}

Let $\mathcal N=U(S_4)'$ and consider the two actual linear word spaces
\begin{align}
 \mathcal V_{\rm int}
  &=\Span\{K_w^*FK_w:F\in\mathcal N,\ w\text{ an original word}\},
       \nonumber\\
 \mathcal R_0
  &=\Span\{K_wXK_w^*:X\in\B(p_0H),\ w\text{ an original word}\}.
 \label{sel13:eq:linear-spaces}
\end{align}

\begin{theorem}[Exact source-specific future and reachable spaces]
\label{sel13:thm:spaces}
For the unchanged benchmark,
\begin{equation}
 \boxed{\mathcal V_{\rm int}=\mathcal S_c\oplus\mathcal D_3,
 \quad\dim\mathcal V_{\rm int}=123,
 \qquad \mathcal R_0=\mathcal S_c,
 \quad\dim\mathcal R_0=120.}
 \label{sel13:eq:spaces}
\end{equation}
Words of length at most three span both displayed spaces.  The complete
record-and-internal-output response from arbitrary native-qutrit inputs has
linear Hankel rank exactly 120, witnessed by a 120-by-120 probability panel
whose past and future words have length at most three.
\end{theorem}
\begin{proof}
The initial internal observables lie in
$\mathcal S_c\oplus\mathcal D_3$.  On the last three coordinates the
original coefficient in complementary pair $k$ is $E_{kk}/\sqrt2$.
Lemma~\ref{sel13:lem:killed} therefore gives the first upper inclusion for
all words.  The same lemma and the initial $p_0$ support give
$\mathcal R_0\subseteq\mathcal S_c$.  The exact certificate of
Section~\ref{sel13:certificate} supplies respectively 123 independent
future vectors and 120 independent reachable vectors at depth three, proving
equality.  It uses nine positive qutrit preparations and fifteen positive
terminal effects spanning $\mathcal N$, rather than treating matrix units as
physical preparations.  Their selected 120-by-120 probability minor has
residue 20 modulo 97.  Hence the Hankel rank is at least 120; the displayed
reachable space bounds it above by 120.
\end{proof}

The 123-dimensional future space is not the 124-dimensional ordinary algebra
of V12.  For example $E_{0,9}$ and $E_{0,10}$ belong to $\mathcal S_c$,
but their adjoint product $E_{9,10}$ does not.  Its ordinary algebraic hull
is $M_{11}$, and adjoining the three final classical coordinates gives the
inherited $M_{11}\oplus\mathcal D_3$.  V12's algebra theorem is retained,
not contradicted.

\begin{proposition}[The linear representative need not be a positive state]
\label{sel13:prop:nonpositive}
The Hilbert--Schmidt projection onto $\mathcal S_c$ is
\[
 P_c(X)=X-\frac12\Tr(YX)Y.
\]
It is not even positive.  On the pure state with vector
$(|0\rangle+|9\rangle+i|10\rangle)/\sqrt3$, its three nonzero-block
eigenvalues are
\begin{equation}
 \boxed{\frac13,\quad\frac{1+\sqrt2}{3},\quad
                     \frac{1-\sqrt2}{3}.}
 \label{sel13:eq:negative}
\end{equation}
\end{proposition}
\begin{proof}
In those three coordinates, the image is
$\frac13\left(\begin{smallmatrix}1&1&-i\\1&1&0\\i&0&1\end{smallmatrix}\right)$.
Its characteristic polynomial is
$(3t-1)(9t^2-6t-1)/27$, proving the assertion.  The negative eigenvalue
excludes positivity and therefore complete positivity.  Thus deletion of the
invisible linear coordinate is not automatically a physical quotient map.
The actual reachable states are still positive matrices in $\mathcal S_c$;
this fact does not extend $P_c$ to a positive map on all original states.
\end{proof}

\section{Positive memory costs without a conditional-expectation ansatz}
\label{sel13:positive}

An ordinary finite quantum predictor here consists of a CPTP encoder, a fixed
normalized instrument on $\mathcal B_{\boldsymbol n}$, and a CPTP decoder
to the declared terminal multiplicity-state algebra.  It must reproduce
unnormalized conditional outputs for every word and every declared input.
The maps are arbitrary CPTP maps: they need not be compressions, algebra
homomorphisms, or conditional expectations.  The retained cost is
$D=\sum_bn_b$.  A history-dependent external controller or stored prefix
that influences subsequent prediction is part of the memory, not free side
information.

\begin{theorem}[Exact positive minimum from native-qutrit preparation]
\label{sel13:thm:native-minimum}
For the unchanged benchmark with arbitrary inputs in $\B(p_0H)$ and all
original records and terminal internal queries retained, the minimum total
Hilbert dimension is exactly eleven.  At that minimum the ordinary algebra
must be $M_{11}$, although the linear minimum is 120.
\end{theorem}
\begin{proof}
The rank-120 panel in Theorem~\ref{sel13:thm:spaces} and
\eqref{sel13:eq:positive-capacity} exclude $D\le10$.
The original $H_c$ is reducing; restricting its actual coefficients to
$H_c$ gives a normalized instrument on $M_{11}$, with the original initial
embedding and terminal multiplicity decoder.  This attains eleven.
A nontrivial direct sum with total Hilbert dimension eleven has at most
$10^2+1=101$ operator coordinates, so it cannot realize rank 120.
No claim about the globally least \emph{operator-space dimension} at a larger
Hilbert cost is needed for this conclusion.
\end{proof}

\begin{theorem}[One native preparation and only the original edge record]
\label{sel13:thm:single-start}
For the unchanged benchmark started at the single state
$|3\rangle\langle3|$, with no terminal quantum Read, the scalar original
record process has linear Hankel rank exactly 66.  A finite 66-by-66 panel
with past depth at most six and future depth at most four witnesses it.
Every ordinary quantum predictor consequently has total Hilbert dimension at
least nine.  The original core gives an eleven-dimensional realization;
existence of a nine- or ten-dimensional realization is not decided here.
\end{theorem}
\begin{proof}
Every coefficient and the initial state are real, so all reachable states are
real symmetric and supported in $H_c$.  Their complex span has dimension
at most $11\cdot12/2=66$.  The exact original-record certificate in
Section~\ref{sel13:certificate} gives a rank-66 Hankel minor, of residue
26 modulo 97, proving equality without an additional output observable.
The lower bound follows from $8^2<66$.  Restriction to the reducing core
supplies the stated upper bound.  This uses just one initial state, not a
hypothetical ability to prepare all states of the fourteen-dimensional system.
\end{proof}

For comparison, there is also an exact extension of V12's all-input minimum
beyond its ordinary-subalgebra ansatz.  We state the packing mechanism because
it distinguishes a genuine positivity restriction from the weaker global
Hankel bound.

\begin{lemma}[Exact sector readout forces orthogonal encoded supports]
\label{sel13:lem:packing}
Suppose a CPTP encoder $\mathcal E_*$ and decoder preserve the sharp sector
probabilities $\Tr(p_\lambda\rho)$ at the initial cut.  Let
$q_\lambda=\supp\mathcal E_*(p_\lambda/\Tr p_\lambda)$ and
$d_\lambda=\rank q_\lambda$.  These supports are pairwise orthogonal.
If the old future observables compressed to $p_\lambda H$ span $r_\lambda$
complex dimensions, every exact predictor satisfies
\begin{equation}
 d_\lambda^2\ge r_\lambda,\qquad
 D\ge\sum_\lambda\lceil\sqrt{r_\lambda}\rceil.
 \label{sel13:eq:packing}
\end{equation}
If $k$ input states in one sector are perfectly distinguished by some
normalized old future experiment, that sector additionally needs
$d_\lambda\ge k$.
\end{lemma}
\begin{proof}
The decoded sector effects form a POVM $F_\lambda$ on the encoded memory.
For the encoded sector density $\sigma_\lambda$, one has
$\Tr(F_\lambda\sigma_\lambda)=1$.  Positivity forces its support into the
unit-eigenspace of $F_\lambda$.  Unit-eigenspaces of distinct members of a
POVM are orthogonal, proving the first assertion.  Every encoded operator
from $p_\lambda\B(H)p_\lambda$ is supported on $q_\lambda$: bound a
Hermitian input above and below by multiples of $p_\lambda$ and use positivity,
then extend linearly.  The compressed future functionals must factor through
this $d_\lambda^2$-dimensional image, giving the rank bound.  A normalized
experiment perfectly distinguishing $k$ states gives, by the same
unit-eigenspace argument, $k$ orthogonal nonzero encoded supports within
$q_\lambda$.
\end{proof}

\begin{corollary}[The all-input positive benchmark minimum is fourteen]
\label{sel13:cor:all-input}
For the unchanged benchmark with all original system inputs, records and
terminal multiplicity outputs retained, the minimum total Hilbert dimension
among arbitrary ordinary CPTP predictors is fourteen.
\end{corollary}
\begin{proof}
The compressed future dimensions from
Theorem~\ref{sel13:thm:spaces} are $(9,36,3,3)$ on
$(p_0,p_W,p_2,p_{\rm s})$.  Thus the first three encoded supports require
at least $3,6,2$ dimensions.  In $W^{\rm s}$, summing the two edges in
complementary pair $k$ gives exactly the effect $E_{kk}$; the three native
axes are perfectly distinguished by the original record.  That support
requires three dimensions.  Lemma~\ref{sel13:lem:packing} gives
$D\ge3+6+2+3=14$.  The original system, or the exact
$M_{11}\oplus\C^3$ reduction of V12, attains fourteen.  No common
conditional-expectation coordinates were assumed in the lower bound.
\end{proof}

\section{Finite exact probability certificates, with a reproducible minor}
\label{sel13:certificate}

The following specification makes the source calculations finite arithmetic
proofs rather than numerical rank estimates.  All source entries used here
belong to $\mathbb Z[1/30,\sqrt2,\sqrt3,i]$ after multiplying each
$K_e$ by $18\sqrt2$.  There is a ring homomorphism to $\mathbb F_{97}$
sending
\begin{equation}
 \sqrt2\mapsto14,\quad\sqrt3\mapsto10,\quad
 \sqrt6\mapsto43,\quad i\mapsto22.
 \label{sel13:eq:residue}
\end{equation}
The squares of these residues are $2,3,6,-1$ respectively, and
$14\cdot10=43$ modulo 97.  All denominators, including
$648=(18\sqrt2)^2$, are invertible.  Symbolic adjoints are taken
\emph{before} applying the residue homomorphism; it is not asserted to
preserve the usual complex conjugation as an operation on the finite field.

Here is a deterministic definition of every minor reported below.  Edge
labels are in the inherited order $(1,+),(1,-),(2,+),(2,-),(3,+),(3,-)$.
Flatten matrices in row-major order.  Starting with the ordered seed list,
scan candidates in order and keep a column precisely when finite-field
elimination against previously kept columns leaves a nonzero residue.  Its
pivot is the first remaining nonzero coordinate; normalize it to one and
record the removed pivot factor.  Retain the \emph{original} matrix and word
label as well as its normalized elimination vector.  At the next depth scan
only the newly retained matrices, in retained order, and then the six edge
labels in order.  Future matrices are updated by $A_e^*FA_e$ and prepend
$e$ to the chronological word; reachable matrices are updated by $A_eXA_e^*$
and append $e$, where $A_e=18\sqrt2K_e$.  Thus every retained column is an
actual scaled word state or effect, not an invented control.

For an observable list $(F_i)$ and reachable list $(X_j)$ obtained this way,
form the matrix $\Tr(F_iX_j)$ modulo 97 and apply the same ordered column
elimination.  Its selected columns and ordered pivot rows define a square
minor.  The product of its recorded pivot factors is its determinant.
Finally divide row $i$ by $648^{|u_i|}$ and column $j$ by $648^{|v_j|}$
to obtain the determinant of the physical probability minor.  The lists,
pivots and residues are exported explicitly in the companion
\src{Renewal_SM_Source_Selection_Research_Notes_V13_certificate.json};
the self-contained checking script regenerates them.

For internal-output tests use the positive basis of $M_m$ consisting of
$|a\rangle\langle a|$ and, for each $a<b$, the projectors onto
$(|a\rangle+|b\rangle)/\sqrt2$ and
$(|a\rangle+i|b\rangle)/\sqrt2$, in that order.  It has $m^2$ members.
Embed these as $I_{V_\lambda}\otimes F$ in the four external sectors,
ordered $(p_0,p_W,p_2,p_{\rm s})$.  This supplies fifteen actual-effect
candidates; their use as laboratory Reads is an explicitly separate
availability condition.  The main record-only certificate uses just $F=I$
and so has no such terminal-input requirement.

\begin{center}
\small
\begin{tabular}{L{.29\textwidth}rrrr}
\toprule
Source and query & Size & Past & Future & $\det H\bmod97$\\
\midrule
Unchanged; native qutrit, internal outputs & 120 & 3 & 3 & 20\\
Unchanged; one ray, original records only & 66 & 6 & 4 & 26\\
\eqref{sel13:eq:full-point}; one ray, original records only
 & 196 & 6 & 6 & 37\\
\bottomrule
\end{tabular}
\end{center}

For the last row the record-effect elimination retains respectively
$1,6,35,144,172,187,196$ columns through depths zero to six.  Reachable
state elimination retains $1,4,10,37,96,181,196$.  These intermediate
finite-field counts are only lower bounds on the corresponding
characteristic-zero ranks; no upper bound at an intermediate depth is inferred
from them.  At the last depth the ambient dimension 196 supplies the matching
upper bound.  The selected \emph{scaled} probability minor has residue 46;
after the physical row and column factors it has residue 37.  The sums of
the selected word lengths total 1478, giving determinant degree 2956.

For the unchanged benchmark the final native-qutrit and internal-future spans
have certified dimensions 120 and 123, with the analytic upper spaces
\eqref{sel13:eq:spaces}.  For the record-only single start, the forward
upper space is the complex span of real symmetric matrices on $H_c$, of
dimension 66.  Its nonzero probability minor proves the matching lower
bound directly.  These certificates establish the assertions used in the
proofs; they do not extrapolate all-length claims from a sequence of
floating-point ranks.  The general all-length upper inclusions are supplied
by the invariant-space arguments above.

\section{The physical opportunity clock and one-sided error tests}
\label{sel13:clock}

The coefficients $K_e$ refer to the accepted six-edge source.  In the actual
locked opportunity process the payload branches are
\[
 \mathcal T_{0,*}=(1-\vartheta)\id,
 \qquad\mathcal T_{e,*}=\vartheta K_e(\cdot)K_e^*,
 \qquad \vartheta=\frac{576851}{417942208512}.
\]
This edge marginal follows by summing the original sign/type record branches,
not by modifying them.  When a whole source packet is studied those other
record factors remain in its state and in any additional requested outputs.

\begin{proposition}[Raw-time and accepted-time panels have the same exact rank]
\label{sel13:prop:clock}
For a word of $k$ accepted letters at specified consecutive opportunities,
$p_{\rm raw}(w)=\vartheta^kp(w)$.  Hence the corresponding raw probability
panel satisfies
\begin{equation}
 H_{\rm raw}
  =\diag(\vartheta^{|u|})H\diag(\vartheta^{|v|}),
 \qquad \rank H_{\rm raw}=\rank H.
 \label{sel13:eq:raw-H}
\end{equation}
Alternatively, retain all intervening no-response records and stop after the
$k$th acceptance.  The accepted edge sequence has exactly the law $p(w)$,
without an exponentially small success postselection.  Its mean number of
opportunities is $k/\vartheta$.
\end{proposition}
\begin{proof}
Each prescribed accepted amplitude acquires a factor $\sqrt\vartheta$,
proving the first formula and the invertible row/column scaling.  For waiting
counts $g_1,\ldots,g_k\ge1$, the full stopped word weight is
\[
 \left[\prod_{j=1}^k\vartheta(1-\vartheta)^{g_j-1}\right]p(w).
\]
The identity no-response branch does not change the state.  Summing each
geometric factor gives one; the waiting counts are independent geometric
variables of mean $1/\vartheta$.  Thus the accepted law and stated mean
follow, with all timing records accounted for.
\end{proof}

For the main twelve-letter panel the mean budget is
$12/\vartheta=5015306502144/576851$ opportunities.  This is an event-count
budget, not an elapsed-time estimate and not a sample-complexity assertion.
Rare individual edge sequences and estimation of a large probability panel
remain substantive statistical requirements.  Raw and stopped panels have
identical exact ranks, but not identical singular values or error bars.

\begin{corollary}[A finite noisy probability certificate]
\label{sel13:cor:error}
Let $H$ be the fixed 196-by-196 record panel and let every estimated entry
have absolute error at most $\varepsilon$.  Then
\[
 \|\widehat H-H\|_\op\le196\varepsilon.
\]
If $\sigma_{171}(\widehat H)>196\varepsilon$, with singular values ordered
decreasingly, Theorem~\ref{sel13:thm:selection} applies to the true source.
If $\sigma_{196}(\widehat H)>196\varepsilon$, the full rank 196 is
certified.  More generally, any candidate predictor whose retained operator
dimension is less than $j$ must make a uniform probability error at least
$\sigma_j(H)/196$ on this panel.
\end{corollary}
\begin{proof}
The Frobenius norm bounds the operator norm of the entrywise error by
$\sqrt{196^2}\varepsilon$.  The singular-value perturbation inequality
therefore bounds the $j$th true singular value below by
$\sigma_j(\widehat H)-196\varepsilon$.  A rank-less-than-$j$ matrix is
at operator-norm distance at least $\sigma_j(H)$, proving the last statement.
The dimension bound on every candidate panel is
Lemma~\ref{sel13:lem:rank}.  The general mechanism is the inherited
finite-Hankel stability argument, specialized here to actual scalar word
probabilities and the source-selection threshold.
\end{proof}

Only certified error bounds justify this one-sided conclusion.  The modular
nonzero determinant proves an exact witness; its residue is not a lower bound
on its real smallest singular value.  Dividing raw probabilities by powers of
$\vartheta$ requires propagating the corresponding error amplification.

\section{What the new certificate closes, and what remains upstream}
\label{sel13:status}

There is now a source-native \emph{record} route to the algebraic structural
branch.  Within the verified capacity-bounded external category, a sufficiently
large finite original-record rank establishes the trivial multiplicity,
full raw system algebra, and full internal intersection.  Full source
matrices, a supplied bridge, a private control bank, stationary faithfulness,
and a separately queried quantum-output moment panel are unnecessary for
that implication.  Its nonempty generic branch is established by an exactly
normalized deformation in the original six-outcome source manifold, not by
adding an intervention to the old alphabet.

The exact positive-memory calculation also resolves a limitation in V12.
Its all-input fourteen-dimensional minimum holds for arbitrary CPTP encoders
and decoders, not just ordinary-subalgebra reductions.  On the unchanged
benchmark, genuinely restricting the preparation to the native qutrit removes
the unreachable three-dimensional sector and gives the sharp positive minimum
eleven.  A single initial state and just the old edge record already rule
out every memory of total Hilbert dimension at most eight.  Linear rank,
ordinary algebra dimension, and physical Hilbert cost remain different
invariants; the explicit numbers $120,121,11$ are one evaluated instance.

The remaining historical step has not been filled by these calculations.
One must identify the actual same-cut normalized carrier and its full symmetry
first.  In the valid $S_4$ category, the next source row can now be the finite
original-record table \eqref{sel13:eq:record-Hankel}, with its timing and
initial-preparation policy retained.  If its certified rank exceeds 170, the
internal algebra follows by Theorem~\ref{sel13:thm:selection}.  If it does
not, neither colour nor the structural branch is excluded: this is a sufficient
selector, not a necessary-and-sufficient classification of every possible
source with that internal algebra.  The unchanged V10 benchmark, whose
internal algebra is already full while its single-start record rank is 66,
is an explicit reminder of that boundary.

No historical word probabilities have been invented or measured here.  The
chronological/local reversal obstruction and the normalized-descent tests
remain in force.  Even a successful rank selector does not supply the
categorical type/private identification, determinant incidence, charged
matter modules, or their common-history provenance.  It strengthens the
finite algebraic selection part of the CMP programme without claiming an
unconditional microscopic Standard-Model derivation.

\paragraph{Verification and append-only preservation.}
The V12 source prefix is retained byte for byte before its final document
terminator.  The companion checks verify the inherited benchmark,
normalization of the orthogonal-column deformation, all tetrahedral
intertwiner identities, the missing-coordinate identities, the positivity
counterexample, the exact finite-field minors, sector-packing data and
physical clock factors.  The source-specific rank equalities use nonzero
finite minors plus analytic upper bounds, not numerical thresholds.
The V12 regression is rerun.  The results are written proofs with exact
computational certificates, not proof-assistant formalization, independent
peer review, or a claim of literature priority.

\begin{thebibliography}{99}
\raggedright
\bibitem[21]{sel13-MW}
A.~Monr\`as and A.~Winter,
\emph{Quantum learning of classical stochastic processes: The completely
positive realization problem}, J.\ Math.\ Phys.\ \textbf{57} (2016), 015219;
\href{https://doi.org/10.1063/1.4936935}{doi:10.1063/1.4936935}.
\bibitem[22]{sel13-BRW}
A.~Bluhm, L.~Rauber and M.~M.~Wolf,
\emph{Quantum compression relative to a set of measurements},
Ann.\ Henri Poincar\'e \textbf{19} (2018), 1891--1937;
\href{https://doi.org/10.1007/s00023-018-0660-z}{doi:10.1007/s00023-018-0660-z}.
These references concern the established positive-realization and
measurement-relative compression questions; the finite source-specific
selectors and certificates above are proved here.
\end{thebibliography}

\addtocontents{toc}{\protect\endgroup}

% END OF SOURCE-SELECTION V13 DEVELOPMENT BLOCK.
% Preserve V1--V13 and append later work before the final terminator.


\clearpage
\begin{center}
{\Large\bfseries Source-selection continuation V14}\\[.4em]
{\large Recover the internal type from symmetry-resolved records}\\[.25em]
Two probability ranks, projective symmetry, and the remaining representation ambiguity
\end{center}
\addtocontents{toc}{\protect\begingroup}
\addtocontents{toc}{\protect\renewcommand{\protect\numberline}{\protect\makebox[2.2em][l]}}

\section{The next assumption to remove is the external inventory}
\label{sel14:context}

V13 obtains internal saturation from sufficiently large original-record rank,
but first fixes the nontrivial external multiplicities in
\eqref{sel13:eq:carrier}.  This continuation goes upstream of that inventory.
It asks what the relabelling symmetry of the \emph{same recorded process}
can determine, without placing its system in a preselected list of
irreducible representations.

The main result is an alternative selector.  Suppose an independently audited
ordinary quantum memory has total Hilbert dimension at most fourteen, and
that the actual record covariance is implemented by reversible completely
positive symmetries of that memory.  No external multiplicity list, coherent
six-coefficient sum, private plane, colour bridge, or weak control is assumed
in the selector.  A full-rank 196-by-196 original-record panel and a
rank-fifteen \emph{averaged relabelled-past} panel force the abstract internal
algebra
\[
 M_3(\C)\oplus M_2(\C)\oplus\C\oplus\C.
\]
The proof initially allows a projective unitary implementation of the physical
$S_4$ action and excludes a nontrivial multiplier at invariant dimension
fifteen.  For the inherited full-rank witness, four class traces computed
from the same probabilities leave exactly two represented algebra types,
each with this abstract internal algebra.

\paragraph{What is traded, rather than silently eliminated.}
The independent scalar memory bound replaces V13's specific external
inventory and its derived bound.  It is not deduced from six outcome names or
from a large Hankel rank.  The required record rank is the full 196, rather
than V13's sufficient 171.  Physical complete-positive symmetry must also be
verified; symmetry of a scalar probability law alone does not equip its
minimum linear predictor with a positive matrix cone.  These are different
hypotheses, made explicit below.  The new result is not claimed to dominate
V13 in every source category.

The bound concerns the actual memory used by the queried process.  In the
worked example this is the independently defined edge-system factor.
Retained pointer/type registers of a larger source are not silently included
in fourteen dimensions or treated as free memory.  Applying the result to a
factor first requires its valid autonomous state reduction; V12's distinction
between coarse and record-preserving descent remains in force.

\paragraph{Audit and nonduplication.}
The ordinary Hankel factorization, finite-rank reconstruction, and positive
memory capacity bounds are inherited from Section~\ref{sel13:selector} and
\cite{Rigidity}.  The supplied spacetime manuscript separates predictive,
linear, and contextual algebraic minimization; that distinction is preserved.
The character-projector compiler in \cite{Attack},
\src{thm:V17-character-compiler}, starts with an inserted, already
reconstructed action on a residual multiplicity carrier.  Here the matrices
of the \emph{state-space symmetry} are instead obtained by relabelling
original past words in a probability panel.  These are different source
inputs and different representations.  The additions are the two-rank
selector, its projective extension, the finite record-symmetry audit, and the
exact character fibre and representation boundary.  Finite character theory
and the distinction between linear and completely positive realizations are
standard background \cite{Etingof,sel13-MW}; no priority claim is made for
those underlying methods.

\section{Reconstruct the symmetry in the coordinates of actual histories}
\label{sel14:records}

Let $(\Phi_{e,*})_{e\in\mathscr E}$ be a normalized finite instrument on an
ordinary quantum memory, with a fixed initial density $\rho$.  Its branch
maps are completely positive and their sum is trace preserving.  They need
not have Choi rank one.  For a chronological word $w=(e_1,\ldots,e_k)$ put
\[
 \Phi_{w,*}=\Phi_{e_k,*}\cdots\Phi_{e_1,*},\qquad
 p(w)=\Tr\Phi_{w,*}(\rho).
\]
The empty word is retained.  All probabilities are unnormalized word
probabilities, obtained from conditional futures by their prefix weights.
For finite lists $\mathcal U=(u_i)_{i=1}^q$ and
$\mathcal V=(v_j)_{j=1}^q$, define
\begin{equation}
 H_{ij}=p(v_ju_i),\qquad
 (H_e)_{ij}=p(v_je u_i),\qquad
 (H_g)_{ij}=p((gv_j)u_i).
 \label{sel14:eq:panels}
\end{equation}
Here a finite group $G$ acts on the actual event labels and hence on words.
Only the \emph{past} is relabelled in $H_g$; the future word stays fixed.
Relabelling the entire concatenated word would instead leave an invariant
probability unchanged and would not measure the action required here.

\begin{proposition}[Probability-coordinate reconstruction]
\label{sel14:prop:coordinates}
Suppose the actual linear state space has dimension $q$ and $H$ is
invertible.  Put
\begin{equation}
 A_e=H^{-1}H_e,\qquad
 b=H^{-1}(p(u_i))_i,\qquad a=(p(v_j))_j,
 \qquad T_g=H^{-1}H_g.
 \label{sel14:eq:coordinates}
\end{equation}
Then $A_e,b,a$ are the matrices of the actual branch maps, initial state and
trace functional in the basis of states $\Phi_{v_j,*}(\rho)$.  In particular
$p(w)=aA_{e_k}\cdots A_{e_1}b$ for every word.  The $j$th column of $T_g$
is the coordinate vector of the actual state $\Phi_{gv_j,*}(\rho)$.
\end{proposition}
\begin{proof}
Let $\mathsf B$ synthesize the $q$ past states and let $\mathsf O$ map a
state $x$ to $(\Tr\Phi_{u_i,*}(x))_i$.  The measured factorization is
$H=\mathsf O\mathsf B$.  Invertibility and the dimension bound make both
maps invertible.  Also $H_e=\mathsf O\Phi_{e,*}\mathsf B$ and
$(p(u_i))_i=\mathsf O\rho$.  Thus
$A_e=\mathsf B^{-1}\Phi_{e,*}\mathsf B$, $b=\mathsf B^{-1}\rho$, and
$a=\Tr\circ\mathsf B$.  If $\mathsf B_g$ synthesizes the relabelled
past states, then $H_g=\mathsf O\mathsf B_g$ and
$T_g=\mathsf B^{-1}\mathsf B_g$.  This proves all assertions.  The use of
invertible Hankel coordinates is standard; no positive structure is yet
assigned to arbitrary coordinate matrices.
\end{proof}

\begin{theorem}[A finite audit of complete record symmetry]
\label{sel14:thm:finite-symmetry}
Under Proposition~\ref{sel14:prop:coordinates}, let $\mathscr S$ generate
$G$.  The original law is invariant under all relabellings,
$p(gw)=p(w)$, if and only if for every $s\in\mathscr S$,
\begin{equation}
 \boxed{T_sb=b,\qquad aT_s=a,\qquad
             T_sA_e=A_{se}T_s\quad(e\in\mathscr E).}
 \label{sel14:eq:finite-symmetry}
\end{equation}
When these conditions hold, $g\mapsto T_g$ is the unique representation on
the minimum linear state space which fixes $b$ and intertwines the relabelled
branch maps.  No additional group-relation assumption is required.
\end{theorem}
\begin{proof}
First assume invariance of the entire law.  A relation
$\sum_jc_j\Phi_{v_j,*}(\rho)=0$ implies
$\sum_jc_jp(v_ju)=0$ for every future $u$.  Invariance gives
$\sum_jc_jp((gv_j)u)=\sum_jc_jp(v_j(g^{-1}u))=0$.
Future separation therefore carries the relation to the relabelled states.
The rule $\Phi_{v,*}(\rho)\mapsto\Phi_{gv,*}(\rho)$ is consequently a
well-defined linear map, with inverse given by $g^{-1}$.  Word composition
proves its group law and the three identities in
\eqref{sel14:eq:finite-symmetry}.  Its matrix on the chosen basis is $T_g$.

Conversely, commute $T_s$ through a chronological product of the $A_e$ and
use $T_sb=b$ and $aT_s=a$.  It follows that $p(sw)=p(w)$ for every word.
Every product of these maps intertwines the corresponding product
relabeling and fixes $b$.  If that relabeling is the identity, the product
commutes with every $A_e$ and fixes every reachable vector $A_wb$.
Those vectors span the state space, so the product is the identity.
Thus all defining relations of $G$ hold and the representation is unique.
The first part then identifies its matrices with $H^{-1}H_g$.
\end{proof}

The audit is finite: if the longest row and column words have lengths $l_U$
and $l_V$, its entries have length at most $l_U+l_V+1$.  It proves an
all-future statement because the actual dimension is already saturated, not
because a sequence of numerical ranks appears to have stabilized.

\subsection{The indispensable completely positive lift}

Suppose the same source supplies reversible CPTP maps $\alpha_g$ on its
actual quantum state algebra, forming an action of $G$ and obeying
\begin{equation}
 \alpha_g\Phi_{e,*}=\Phi_{ge,*}\alpha_g.
 \label{sel14:eq:physical-covariance}
\end{equation}
No runtime application of these symmetries is needed to evaluate
\eqref{sel14:eq:panels}; covariance is an independently checked property of
the source operations.  If the record law passes
Theorem~\ref{sel14:thm:finite-symmetry}, then
\begin{equation}
 \alpha_g\rho=\rho,\qquad
 T_g=\mathsf B^{-1}\alpha_g\mathsf B.
 \label{sel14:eq:physical-lift}
\end{equation}
Indeed, covariance gives
$p(gw)=\Tr\Phi_{w,*}(\alpha_g^{-1}\rho)$, and the full future-effect
span distinguishes $\alpha_g^{-1}\rho$ from $\rho$.  Thus the invariance of
the initial state can itself be checked by the finite record audit.

\begin{proposition}[The quantum symmetry is unique if it exists]
\label{sel14:prop:CP-lift}
On a matrix memory $M_N$, the lifted maps in
\eqref{sel14:eq:physical-lift} are uniquely determined by the original
probability panels and the physical past-state synthesis $\mathsf B$.
If only the record audit is initially known, define
$\alpha_s=\mathsf BT_s\mathsf B^{-1}$ for three adjacent-transposition
generators of $S_4$.  They give the required reversible CPTP action if and
only if their Choi matrices are positive semidefinite.  Trace preservation
is already supplied by $aT_s=a$.
\end{proposition}
\begin{proof}
The past states are a basis, so their images determine each map uniquely.
The finite audit implies $T_s^2=I$ and all Coxeter relations.  Thus a
completely positive $\alpha_s$ has completely positive inverse, namely
itself, and is trace preserving.  Its adjoint and inverse adjoint are unital
completely positive maps.  Applying their Schwarz inequalities successively
forces equality in both inequalities.  The multiplicative-domain identity
then makes them inverse $*$-automorphisms.  A $*$-automorphism of $M_N$ is
unitary conjugation, as follows by transporting a complete matrix-unit
system.  Products of the generators give the stated action.  The converse
is Choi positivity of each actual CPTP symmetry.
\end{proof}

This is a test on the \emph{actual operator coordinates}.  A reconstructed
linear similarity matrix is not automatically a quantum operation.
Without an independently specified positive realization, the columns of
$H$ do not supply $\mathsf B$ as matrices, and the Choi tests cannot be
performed by declaring those columns to be density operators.  This is the
same positive-realization boundary distinguished in
\cite{sel13-MW}.  Anti-automorphic transposition of chronological histories
is not a reversible CPTP map on a full noncommutative system, so it cannot
be inserted into \eqref{sel14:eq:physical-covariance} as a substitute.

\section{Two original-probability ranks select the abstract internal type}
\label{sel14:two-ranks}

Average only the past relabellings:
\begin{equation}
 \overline H=\frac1{|G|}\sum_{g\in G}H_g,
 \qquad P=H^{-1}\overline H.
 \label{sel14:eq:averaged-panel}
\end{equation}
Every entry of $\overline H$ is a classical arithmetic average of original
word probabilities.  It is not an amplitude superposition or a newly
executed symmetry pulse.

\begin{theorem}[Invariant algebra dimension is a probability rank]
\label{sel14:thm:twirl-rank}
Assume the complete record symmetry and the physical lift
\eqref{sel14:eq:physical-covariance} on $M_N$, with an invertible panel
of size $q=N^2$.  Then
\begin{equation}
 \boxed{P^2=P,\qquad
 \rank\overline H=\Tr P
 =\dim_\C\{x\in M_N:\alpha_g(x)=x\ \forall g\}.}
 \label{sel14:eq:twirl-rank}
\end{equation}
The last space is an ordinary unital $C^*$-algebra.  If
$\mathcal A=C^*(K_{e\mu}:e,\mu)$ is the algebra of the actual Kraus
coefficient spans, the full record rank also gives $\mathcal A=M_N$.
Thus \eqref{sel14:eq:twirl-rank} is the dimension of the historically
generated internal intersection, not just of a kinematically available
commutant.
\end{theorem}
\begin{proof}
The averaged physical map
$\mathsf E=|G|^{-1}\sum_g\alpha_g$ is the trace-preserving conditional
expectation onto the fixed algebra.  By
\eqref{sel14:eq:physical-lift}, $P=\mathsf B^{-1}\mathsf E\mathsf B$.
It is idempotent with trace equal to its rank.  Multiplication by the
invertible $H$ gives $\rank\overline H=\rank P$.  The fixed space is
closed under products because every $\alpha_g$ is an automorphism.
For the last assertion, every future effect is a sum of products of the
actual coefficients and their adjoints, hence belongs to $\mathcal A$.
Full Hankel rank makes those effects span $M_N$, proving $\mathcal A=M_N$.
\end{proof}

The coordinate projector $P$ need not be Hermitian or norm one: the past
basis is not orthonormal.  Idempotence and its integer trace, rather than
an unjustified Euclidean orthogonal projection, are used below.

\begin{lemma}[A fifteen-dimensional invariant algebra under $S_4$]
\label{sel14:lem:fifteen}
Let $S_4$ act by automorphisms on any $M_N$.  If its fixed
algebra has complex dimension fifteen, its abstract algebra type is
\begin{equation}
 \boxed{M_3(\C)\oplus M_2(\C)\oplus\C\oplus\C.}
 \label{sel14:eq:abstract-type}
\end{equation}
An ordinary unitary lift of the action is not required as an assumption.
\end{lemma}
\begin{proof}
Choose unitary implementers $U_g$.  They form a projective representation
with a common multiplier $\omega$, since the centre of $M_N$ is scalar.
Decompose it into inequivalent irreducible representations with this
multiplier:
\[
 \C^N=\bigoplus_\lambda V_\lambda\otimes\C^{m_\lambda},
 \quad N=\sum_\lambda d_\lambda m_\lambda,
 \quad M_N^{S_4}\cong\bigoplus_\lambda M_{m_\lambda}.
\]
Consequently $\sum_\lambda m_\lambda^2=15$.

If $\omega$ is cohomologically trivial, an ordinary lift exists and $S_4$
has five inequivalent irreducibles.  There are at most five nonzero
$m_\lambda$.  No $m_\lambda$ exceeds three.  If $a,b,c$ count the entries
three, two, and one, then $9a+4b+c=15$ and $a+b+c\le5$.
If $a=0$, then $b\le3$ and $a+b+c=15-3b\ge6$, a contradiction.
Hence $a=1$, and $b=0$ would give seven entries.  The only possibility is
$(a,b,c)=(1,1,2)$.

If $\omega$ is nontrivial, Lemma~\ref{sel14:lem:twisted-S4} below shows
that it has at most three inequivalent irreducible projective
representations.  Thus $15$ would be a sum of at most three nonzero
integer squares.  Each entry is at most three; if one entry is three,
the remaining six is not a sum of two squares from $\{0,1,4,9\}$.
If none is three, three entries contribute at most twelve.  Both are
impossible.  Hence the multiplier is trivial, and the ordinary calculation
already gives \eqref{sel14:eq:abstract-type}.  In particular this algebra-type
conclusion itself needs no dimension-fourteen assumption; that bound is used
to identify the full physical state space from the specified probability panel.
\end{proof}

\begin{theorem}[Two-rank source selector without an external multiplicity inventory]
\label{sel14:thm:selection}
Suppose the actual ordinary finite quantum memory has independently
verified total Hilbert cost at most fourteen.  Assume its record alphabet
has an $S_4$ action implemented by reversible CPTP source symmetries and
that the initial record law is invariant, for example as verified by
Theorem~\ref{sel14:thm:finite-symmetry}.  If a square original-record panel
and its averaged relabelled-past panel satisfy
\begin{equation}
 \boxed{\rank H=196,\qquad\rank\overline H=15,}
 \label{sel14:eq:two-rank-selector}
\end{equation}
then the actual ordinary state algebra is $M_{14}$, the original coefficient
algebra is all of $M_{14}$, and its invariant internal intersection has
abstract type \eqref{sel14:eq:abstract-type}.
No particular external irreducible inventory, resolved Kraus rank,
six-edge coherent sum, private action, weak projector, or bridge is assumed.
\end{theorem}
\begin{proof}
An ordinary memory $\bigoplus_bM_{n_b}$ with total Hilbert cost
$D=\sum_bn_b\le14$ has at most $D^2\le196$ linear coordinates.
Full rank 196 forces equality throughout and hence a single block
$M_{14}$.  Apply Theorem~\ref{sel14:thm:twirl-rank} to identify its full
original coefficient algebra and its fifteen-dimensional fixed algebra.
Lemma~\ref{sel14:lem:fifteen} gives the asserted abstract type.
\end{proof}

\begin{remark}[Information content and remaining scope]
Two ranks are two \emph{summaries of finite probability panels}, not two
single-shot measurements.  The source dimension cap is an independent
resource or reconstruction statement.  Without it, invisible spectators
and larger positive realizations are not excluded by an observed rank.
The theorem determines the algebra's abstract matrix sizes.  It does not
yet specify their represented support ranks or the geometric names of their
carriers.  Those distinctions are resolved as far as the same records allow
in the next section.
\end{remark}

\section{Projective symmetry and the exact representation fibre}
\label{sel14:character}

For any $g\in G$, the same panels give the basis-independent class scalar
\begin{equation}
 \kappa(g):=\Tr(H^{-1}H_g)=\Tr\alpha_g=|\Tr U_g|^2.
 \label{sel14:eq:character}
\end{equation}
The last equality follows by vectorizing unitary conjugation.  It holds for
projective implementers and is independent of their scalar phases.

\begin{lemma}[The nontrivial projective $S_4$ branch has only three irreducibles]
\label{sel14:lem:twisted-S4}
For a nontrivial unitary multiplier $\omega$ of $S_4$, the finite twisted
group algebra has exactly three simple summands, of matrix sizes $2,2,4$.
In particular it has at most three inequivalent irreducible projective
representations.  This standard projective representation fact admits the
following elementary proof sufficient for the selection argument.
\end{lemma}
\begin{proof}
In the unitary twisted group algebra, choose lifts $S_1,S_2,S_3$ of the
adjacent transpositions and multiply them by phases so that $S_i^2=I$.
The distant relation is $S_1S_3=\varepsilon S_3S_1$, where
$\varepsilon^2=1$.  The adjacent product cubes are scalars $z_1,z_2$;
conjugating each product by one of its involutions inverts it, so
$z_i\in\{1,-1\}$.  Signs of the three generators can make both cubes
one.  If $\varepsilon=1$, the resulting lifts satisfy the ordinary
$S_4$ Coxeter presentation.  Every resulting group word is a scalar
multiple of the original twisted basis unit, so these words give a section
whose multiplier is identically one.  This would make $\omega$ trivial.
Therefore a nontrivial multiplier forces $\varepsilon=-1$.

Conjugation by a basis unit sends each twisted group basis vector to a
scalar multiple of the vector in its ordinary conjugacy class.  A central
vector consequently has at most one free coefficient per conjugacy class.
The transposition class contributes none: conjugation by the disjoint
transposition sends $S_1$ to $-S_1$.  The double-transposition class
also contributes none, since conjugation by $S_1$ sends $S_1S_3$ to
$-S_1S_3$.  Of the five ordinary classes, only three remain.
Thus the centre of the twisted algebra has dimension at most three.

The twisted algebra is a finite-dimensional $C^*$-algebra of complex
dimension $24$ in its faithful unitary regular representation.  It has no
one-dimensional representation, because such a representation would make
$\omega$ a coboundary.  Its simple block dimensions therefore satisfy
$\sum d_i^2=24$, with at most three $d_i$ and every $d_i\ge2$.
The only possibility is $(d_i)=(2,2,4)$: each $d_i\le4$, and direct
substitution into $4a+9b+16c=24$, $a+b+c\le3$, gives
$(a,b,c)=(2,0,1)$.  This proves the assertion without importing a
projective character table.
\end{proof}

\begin{corollary}[A dimension-free abstract-type diagnostic]
\label{sel14:cor:dimension-free}
On any verified matrix memory whose original record panel saturates all
$N^2$ state coordinates, actual reversible $S_4$ covariance and
$\rank\overline H=15$ force an ordinary unitary symmetry lift and the
abstract internal type $M_3\oplus M_2\oplus\C^2$.  The number fourteen
in Theorem~\ref{sel14:thm:selection} is the physical memory bound for its
particular $196\times196$ panel, not a prerequisite of this algebra-type
classification.
\end{corollary}
\begin{proof}
Theorem~\ref{sel14:thm:twirl-rank} identifies the invariant algebra dimension.
Lemma~\ref{sel14:lem:fifteen}, including its exclusion of a nontrivial
multiplier, gives the conclusion.  Full record rank makes the original
coefficient algebra the entire matrix algebra.
\end{proof}

\begin{theorem}[A nonzero transposition trace removes the projective lift obstruction]
\label{sel14:thm:projective-lift}
Let $S_n$, $n\ge3$, act by $*$-automorphisms on a full matrix algebra.
If $\Tr\alpha_s>0$ for a transposition $s$, the action admits an ordinary
unitary representation lift.  In particular the condition applies to the
$S_4$ character recovered in \eqref{sel14:eq:character}.
\end{theorem}
\begin{proof}
Choose unitary implementers $S_i$ of the adjacent transpositions
$(i,i+1)$ and rescale them so that $S_i^2=I$.  They are then Hermitian.
All transpositions are conjugate at the automorphism level, so
$\Tr\alpha_{s_i}=|\Tr S_i|^2>0$ for every $i$.
For $|i-j|>1$, projective commutation gives
$S_iS_j=\varepsilon_{ij}S_jS_i$.  The involution identities force
$\varepsilon_{ij}^2=1$.  Taking the trace after conjugation by $S_j$
shows $\Tr S_i=\varepsilon_{ij}\Tr S_i$; hence
$\varepsilon_{ij}=1$.

For adjacent generators, put $P_i=S_iS_{i+1}$.  Its cube is a scalar
$z_iI$ because the corresponding automorphism has order dividing three.
Since $S_iP_iS_i=P_i^{-1}$, one has $z_i=z_i^{-1}$ and $z_i\in\{1,-1\}$.
Choose signs $\delta_1=1$ and recursively
$\delta_{i+1}=\delta_i z_i$.  Replacing $S_i$ by $\delta_iS_i$ makes
all adjacent cubes equal to $I$ and preserves the distant commutation
relations and involution squares.  These are the Coxeter relations for
$S_n$.  They therefore define a genuine unitary representation implementing
the original automorphisms.  The generator presentation is standard; the
argument exhibits the required phase choice explicitly.
\end{proof}

\begin{example}[Why nonzero trace is a substantive condition]
\label{sel14:ex:projective}
With Pauli matrices $X,Y,Z$, take
\[
 S_1=X,\qquad S_2=-\tfrac12(X+Y)+2^{-1/2}Z,\qquad S_3=Y.
\]
Their squares and adjacent product cubes equal $I$, while
$S_1S_3=-S_3S_1$.  Their conjugations consequently satisfy the $S_4$
Coxeter relations, but no scalar rephasing can make the distant lifts
commute.  All three lifts are traceless, so the transposition character of
the conjugation action is zero.  This is a finite algebraic comparison, not
an additional source operation assigned to the Store.
\end{example}

For $S_4$, order the five conjugacy classes as
$1,(12),(12)(34),(123),(1234)$, with sizes $1,6,3,8,6$.
The ordinary character table, already used in \cite{Attack}, is
\[
\begin{array}{c|rrrrr}
 &1&(12)&(12)(34)&(123)&(1234)\\\hline
 \mathbf1&1&1&1&1&1\\
 \mathrm{sgn}&1&-1&1&1&-1\\
 V_2&2&0&2&-1&0\\
 W&3&1&-1&0&-1\\
 W^{\rm s}&3&-1&-1&0&1
\end{array}
\]
Its orthogonality and its interpretation as the five irreducibles are
standard finite representation theory \cite{Etingof}.  The inverse fibre
calculation below is the new source-identification question.

\begin{theorem}[Complete external inventory fibre of the recorded character]
\label{sel14:thm:character-fibre}
On a fourteen-dimensional matrix memory, suppose the physical symmetry
character reconstructed from the original probability panels is
\begin{equation}
 \boxed{\kappa=(196,16,4,4,4).}
 \label{sel14:eq:target-character}
\end{equation}
Then an ordinary lift exists.  Its irreducible multiplicity vector in the
order $(\mathbf1,\mathrm{sgn},V_2,W,W^{\rm s})$ is exactly one of
\begin{equation}
 \boxed{(3,0,1,2,1),\quad(0,3,1,1,2),\quad
        (2,1,1,3,0),\quad(1,2,1,0,3).}
 \label{sel14:eq:character-fibre}
\end{equation}
All four have abstract fixed algebra $M_3\oplus M_2\oplus\C^2$.
Up to sign twist, however, they have two different represented types.
\end{theorem}
\begin{proof}
The transposition scalar is positive, so
Theorem~\ref{sel14:thm:projective-lift} supplies the ordinary lift.
Let its multiplicities be $(a,b,c,d,e)$, and write
$x=a+b$, $y=a-b$, $z=d+e$, $t=d-e$.
The defining-system character is
\[
 (x+2c+3z,\ y+t,\ x+2c-z,\ x-c,\ y-t).
\]
Its first value is fourteen and the squared values are
\eqref{sel14:eq:target-character}.  Put $v=x+2c-z\in\{2,-2\}$ and
$w=x-c\in\{2,-2\}$.  Then
$z=(14-v)/4$ and $x+2c=(14+3v)/4$.
For $v=2$, integrality and $w=\pm2$ force $(x,c,z)=(3,1,3)$.
For $v=-2$, the only integral nonnegative option is $(x,c,z)=(2,0,4)$.
But $y=((y+t)+(y-t))/2$ is odd because $y+t=\pm4$ and $y-t=\pm2$;
it cannot have the parity of $x=2$.  The second option is impossible.

For $(x,c,z)=(3,1,3)$, the possible pairs $(y,t)$ are
$(3,1)$, $(1,3)$, $(-1,-3)$, and $(-3,-1)$.
Solving $a,b=(x\pm y)/2$ and $d,e=(z\pm t)/2$ gives exactly
\eqref{sel14:eq:character-fibre}.  Twisting by sign exchanges the first
pair and the last pair.  Their positive multiplicities are all
$(3,2,1,1)$, proving the abstract assertion.
\end{proof}

\begin{corollary}[Exactly what the character fails to identify]
\label{sel14:cor:represented}
Let $p_{(3)}$ denote the central unit of the unique $M_3$ summand of the
invariant algebra.  For the four inventories in
\eqref{sel14:eq:character-fibre}, precisely two physical support lists occur:
\begingroup\small
\begin{equation}
\begin{array}{c|c|c|c}
\text{inventory up to sign twist}&\rank p_{(3)}&
 \text{ranks of }(M_3,M_2,\C,\C)&\dim C^*(U(S_4))\\\hline
 (3,0,1,2,1)&3&(3,6,2,3)&23\\
 (2,1,1,3,0)&9&(9,2,1,2)&15
\end{array}
\label{sel14:eq:represented-types}
\end{equation}
\endgroup
Thus the class character alone does not identify the geometric colour carrier
with the native trivial sector.  Once the actual invariant \emph{product}
and represented trace are reconstructed, the single scalar
$\Tr_Hp_{(3)}\in\{3,9\}$ distinguishes the two types.
\end{corollary}
\begin{proof}
A summand $M_m$ attached to an external irrep of dimension $d$ acts on a
support of rank $dm$.  The algebra generated by the external action has
one $M_d$ for each occurring distinct irrep, of dimension $\sum d^2$.
Apply these formulas to \eqref{sel14:eq:character-fibre}.
Sign twist changes the irrep names but not their dimensions or commutant.
The $M_3$ summand is unique by its matrix size, so its central unit is
intrinsic after multiplication and the physical trace have been identified.
\end{proof}

This is an ambiguity of the extracted \emph{class-character summary}, not
an assertion that all four types realize the same complete probability
process.  Full product-sensitive source data may distinguish them.  The
scalar trace in the corollary is not automatically an admitted sharp Read.
Even its value three would not identify particular type/private rays or
supply determinant and charged-matter incidence.

\section{Evaluate the unchanged full-rank source and its relabelled probabilities}
\label{sel14:certificate}

Use exactly the source, initial ray and word lists of
\eqref{sel13:eq:full-point} and Theorem~\ref{sel13:thm:full-panel}.
No source parameter, instrument branch, or preparation is changed.
The initial ray belongs to the trivial sector, so its density is invariant.
The original six coefficients already have the required physical covariance.
In V13's explicit frame, the external matrices are
\begin{equation}
 U_g=I_3\oplus(u_g\otimes I_2)\oplus v_g
                   \oplus(\det u_g)u_g,
 \label{sel14:eq:witness-action}
\end{equation}
where $u_g$ is the tetrahedral standard matrix and $v_g$ is its action on
traceless diagonal matrices.  These matrices verify the hypotheses of the
probability theorem; the probability panels themselves are computed by
original-word recursion, not fabricated from a desired character.

\begin{theorem}[An exact two-rank probability witness with no additional Read]
\label{sel14:thm:witness}
For that unchanged source and V13's fixed lists, all words in $H_g$ have
length at most twelve and
\begin{equation}
 \boxed{\rank H=196,\qquad\rank\overline H=15,
       \qquad\Tr(H^{-1}H_g)=(196,16,4,4,4)_{[g]}.}
 \label{sel14:eq:witness}
\end{equation}
The finite record-symmetry audit additionally uses original words of length
at most thirteen.  No quantum terminal effect other than the trace is needed
for any of these entries.
\end{theorem}
\begin{proof}
The normalized covariance and full-rank word lists are inherited from V13.
Here the additional probability calculation is as follows.  For each of the
24 tetrahedral label permutations, relabel the selected past word letter by
letter, apply its original Kraus maps to the same initial ray, and pair with
the unchanged future effect.  The companion script computes all these
$196\times196$ matrices over V13's residue ring, with the physical word
scales retained.  The base minor has residue 37 modulo 97.  Thus its inverse
exists in the localized characteristic-zero ring and modulo 97.
The recovered class-trace residues are
\begin{equation}
 (2,16,4,4,4)\pmod{97},
 \label{sel14:eq:character-residues}
\end{equation}
where $196\equiv2$.  The script also checks the initial/terminal and
transition intertwinings of Theorem~\ref{sel14:thm:finite-symmetry}.

There is an exact integer lift of the residues, not a floating-point trace
estimate.  For a projective involution on $\C^{14}$, choose a Hermitian
unitary lift.  Its trace is an even integer of absolute value at most 14,
so its conjugation trace is that integer's square.  Among these squares,
residue 16 forces absolute trace 4, since $97>2\cdot14$.
Theorem~\ref{sel14:thm:projective-lift} now gives an ordinary $S_4$ lift.
Every $S_4$ character value is integral.  If $0\le r,s\le14$ and
$r^2\equiv s^2\pmod{97}$, primality and $97>28$ imply $r=s$.
Thus the remaining residue 4 values have unique squared-character lift 4.
This proves the exact class values in \eqref{sel14:eq:witness} without
assuming the desired external inventory in the inversion.

Finally the rank of the averaged action equals its trace,
\[
 \frac1{24}(196+6\cdot16+3\cdot4+8\cdot4+6\cdot4)=15.
\]
Theorem~\ref{sel14:thm:twirl-rank} gives the exact rank of $\overline H$.
The companion calculation independently supplies a nonzero 15-by-15 minor
and all modular idempotence identities.  The analytic trace computation
provides the characteristic-zero upper bound; a modular rank plateau by
itself would not do so.  The word-length bounds follow from relabelling
without adding letters and from the one inserted original event in $H_e$.
\end{proof}

The exported certificate retains the past/future lists, all 24 label
permutations and recovered class residues, the averaged minor, and the four
inventory candidates.  Its integer lifting uses representation and dimension
bounds.  The residue 37 is not a real determinant size or a lower singular
value bound.  This calculation evaluates the existing mathematical witness;
it supplies no measurements of the historical Store.

\begin{proposition}[The same panel is generic in the inherited normalized family]
\label{sel14:prop:generic}
On V10's connected rank-three normalized covariance manifold, keep the
invariant initial ray and V13's fixed word lists.  On the open dense
full-volume set where $H$ is invertible, $\overline H$ has rank exactly
fifteen and the new selector applies.  No further nonzero polynomial or
new source deformation is needed to establish this set.
\end{proposition}
\begin{proof}
V13 already proves the full-rank assertion for the fixed $H$.  Covariance
and the same invariant preparation hold throughout that family.
The external fixed algebra in that already verified family has dimension
fifteen.  Theorem~\ref{sel14:thm:twirl-rank} identifies this number with the
averaged probability rank wherever $H$ is invertible.
\end{proof}

This proposition reuses V13's nonzero point and genericity result; it is not
presented as a second discovery of their existence.  The new application is
the removal of their external inventory from the \emph{inference rule}
when a scalar physical memory cap and symmetry audit are available instead.

\section{Certified errors for integer symmetry dimensions}
\label{sel14:error}

A positive singular value can certify a rank lower bound, but by itself it
cannot certify that an averaged panel has rank \emph{exactly} fifteen.
The integer trace of the symmetry projector supplies the missing upper
identification.

\begin{proposition}[Inverse-panel and trace error budget]
\label{sel14:prop:error}
Let $H$ be invertible of size $q$ and let $C$ denote either $H_g$ or
$\overline H$.  Suppose certified estimates satisfy
\[
 \|\widehat H-H\|_\op\le\delta,
 \qquad\|\widehat C-C\|_\op\le\delta_C,
 \qquad s=\sigma_{\min}(\widehat H)>\delta.
\]
Set $\widehat T=\widehat H^{-1}\widehat C$.  Then
\begin{equation}
 \boxed{\|H^{-1}C-\widehat T\|_\op
 \le \frac{\delta_C+\delta\|\widehat T\|_\op}{s-\delta},
 \qquad
 |\Tr(H^{-1}C)-\Tr\widehat T|
 \le q\frac{\delta_C+\delta\|\widehat T\|_\op}{s-\delta}.}
 \label{sel14:eq:error}
\end{equation}
Under the exact physical symmetry hypotheses, with $C=\overline H$,
\begin{equation}
 \boxed{|\Tr\widehat T-15|
       +q\frac{\delta_C+\delta\|\widehat T\|_\op}{s-\delta}
       <\frac12}
 \label{sel14:eq:integer-test}
\end{equation}
certifies $\rank\overline H=15$.  Together with $q=196$ and the independent
memory bound, this is a one-sided noisy version of the two-rank selector.
\end{proposition}
\begin{proof}
The least singular value inequality gives
$\|H^{-1}\|\le(s-\delta)^{-1}$.  The exact identity
\[
 H^{-1}C-\widehat T
 =H^{-1}\bigl[(C-\widehat C)-(H-\widehat H)\widehat T\bigr]
\]
gives the first bound.  The estimate $|\Tr X|\le q\|X\|$ gives the
second.  For $C=\overline H$, the true trace is an integer by
Theorem~\ref{sel14:thm:twirl-rank}.  Condition
\eqref{sel14:eq:integer-test} places it strictly within one half of fifteen,
so it must equal fifteen.  Invertibility of $H$ is already certified by
$s>\delta$.
\end{proof}

For entrywise errors bounded by $\varepsilon$ in each $q\times q$ panel,
one may use $\delta,\delta_C\le q\varepsilon$.  Averaging group-related
entries does not justify a further statistical improvement unless their
error dependence is analyzed.  The matrices $\widehat T$ can be poorly
conditioned in the past-state basis; the displayed norm and denominator
must be retained.  The bounds do not convert approximate record covariance
into exact quantum symmetry.  That is a separate hypothesis or
positive-lift audit.  Individual class traces can be treated by the same
formula once their integral-character interpretation has been verified.

\paragraph{The original opportunity calibration.}
For a fixed source factor and relabelling that preserves word length, the
raw opportunity panels have the same diagonal scaling as V13:
\[
 H^{\rm raw}=D_UHD_V,\qquad
 H_g^{\rm raw}=D_UH_gD_V,\qquad
 \overline H^{\rm raw}=D_U\overline HD_V,
\]
where $(D_U)_{ii}=\vartheta^{|u_i|}$ and
$(D_V)_{jj}=\vartheta^{|v_j|}$.  Consequently
\[
 (H^{\rm raw})^{-1}H_g^{\rm raw}=D_V^{-1}T_gD_V.
\]
Both ranks and exact traces are invariant under this scaling; singular
values and error bounds are not.  V13's accepted-time stopping construction
therefore applies unchanged, with mean budget $k/\vartheta$ opportunities
for $k$ accepted events.  It is not a sample-complexity bound for estimating
these probability panels.

\section{The upstream conclusion and the next unresolved source row}
\label{sel14:status}

The independent external \emph{multiplicity inventory} is no longer necessary
for one sufficient algebraic selector.  A verified small quantum carrier,
actual reversible $S_4$ covariance, and the two original-probability ranks
\eqref{sel14:eq:two-rank-selector} reconstruct the abstract internal matrix
sizes.  The physical action is not freely fitted afterward: when the
record state space is full, its lift on that particular matrix realization
is unique.  A nontrivial projective multiplier is excluded at invariant dimension
fifteen.  Independently, a positive transposition trace removes the lift
obstruction before the more detailed character inversion.

The remaining distinctions are not hidden in the word ``selection''.
A scalar probability symmetry must still be compared with the actual
positive system action.  The bound fourteen must come from a real memory
or a validated positive source reduction, not from the label count.
The two represented types in \eqref{sel14:eq:represented-types} show why
the abstract algebra does not by itself identify its geometric carriers.
Neither the determinant tensor, charged matter types, nor the common-history
Dirac incidence follows from the symmetry ranks.  The original normalization
and reversal obstructions of V11--V12 are not removed by changing the
representation language.

For the historical source, the next finite row is therefore either its
symmetry-resolved probability panel on an independently bounded positive
carrier, or the positive-lift check for the unique linear action already
reconstructed from that panel.  If the abstract type passes, a source-native
represented trace and product distinguish the remaining carrier allocations.
An absent physical panel is not filled by the existing benchmark.  No new
historical probabilities, source admission, dynamical basin, or unconditional
microscopic Standard-Model derivation is asserted.

\paragraph{Exact checks and preservation.}
The V13 source prefix is retained byte for byte before its terminal document
command.  The checking script reuses the unchanged full-rank coefficient
family, recomputes relabelled words for all 24 group elements, verifies the
probability-coordinate and projector identities, enumerates all 291
ordinary dimension-fourteen $S_4$ inventories, and checks the projective
Coxeter example and exact inverse-error identity.  The V13 regression is
rerun.  General selector and ambiguity assertions are proved in the text;
modular calculations supply explicit finite witnesses with stated integer
lifting bounds.  These are written proofs and reproducible exact checks,
not proof-assistant formalization or independent peer review.

\addtocontents{toc}{\protect\endgroup}
% END OF SOURCE-SELECTION V14 DEVELOPMENT BLOCK.
% Preserve V1--V14 and append later development before the final terminator.



\clearpage
\begin{center}
{\Large\bfseries Source-selection continuation V15}\\[.4em]
{\large Operational products, quantum symmetry, and the represented carrier}\\[.25em]
Finite overlap tests replace a memory cap and an assumed positive lift
\end{center}
\addtocontents{toc}{\protect\begingroup}
\addtocontents{toc}{\protect\renewcommand{\protect\numberline}{\protect\makebox[2.2em][l]}}

\section{Go upstream from record coordinates to their physical product}
\label{sel15:context}

V14 removes an independently stipulated external multiplicity inventory, but
retains two substantive premises: an upper bound on the actual quantum memory
and complete positivity of the symmetry reconstructed from scalar records.
It also leaves two represented types with the same conjugation character.
Repeating the Hankel-rank argument cannot remove any of these distinctions.
The present continuation asks which additional \emph{scalar operational
comparisons} determine the required product, rather than choosing matrix
coordinates for a linear predictor.

There are two positive routes.  First, on a verified matrix memory, pair and
\emph{symmetric cubic} overlaps of actual past states determine their Jordan
product.  Invariance of those overlaps makes the record symmetry a Jordan
automorphism.  An involutive anti-automorphism has state-space trace $\pm N$;
the already recorded transposition trace $16$ on $M_{14}$ excludes that
alternative and therefore derives the completely positive lift.  A quadratic
trace invariant of the recovered internal product then separates the two
remaining represented types.

Second, selected fourth-order overlaps measure the failure of the chosen
past-state span to be an ordinary algebra.  Together with two-copy tests of
its invariance under each original event, their vanishing constructs an exact
positive predictor.  A one-dimensional centre and a rank-196 record panel
then derive, rather than assume, a minimal $M_{14}$ memory.  Ordered cubic
covariance certifies its physical symmetry.  This route needs richer
same-cut, multiple-copy comparison data; it is not a claim that the original
single-copy record law alone fixes a quantum cone.

\paragraph{Audit and scope.}
The finite record reconstruction and symmetry audit, projective $S_4$
classification, and four-inventory character fibre are inherited from
Sections~\ref{sel14:records}--\ref{sel14:character}.  The positive-realization
and source-short distinctions in V12--V13 are retained.  The supplied spacetime
manuscript already separates predictive, linear, and contextual algebraic
minimization; its full operational entrance includes physical intervention
frames, not just an unlabelled scalar sequence.  The added overlap identities
must be supplied on the \emph{same acted-on states}; neither a categorical
source Gram nor the regular trace of a different history representation is a
substitute.

The Jordan automorphism classification and cyclic-permutation estimation of
state moments are established methods \cite{sel15-Jordan,sel15-Ekert}.  Their
finite proofs and the specialized reconstruction below do not claim general
literature priority.  New here relative to the audited development are the
combined symmetry/capacity selector, the represented-trace discriminator, the
same-rank nonpositive symmetry countermodel, and the explicit application to
the unchanged full-rank source.  Determinant incidence, charged matter, and
historical admission remain separate.

\section{The physical second, third, and fourth overlap panels}
\label{sel15:panels}

Let $\Phi_{e,*}$ be the original normalized instrument on a finite Hilbert
space $H$, initially in a density $\rho$.  Its dimension is not bounded in
this section.  Choose original pasts $v_1,\ldots,v_q$ with $v_1=\emptyset$
and write
\[
 X_i=\Phi_{v_i,*}(\rho)\succeq0,\qquad
 \mathcal V=\Span_\C\{X_1,\ldots,X_q\},\qquad
 a_i=\Tr_H X_i=p(v_i).
\]
Retain an invertible original probability panel $H_{ui}=p(v_i u)$ of size
$q$.  This already proves independence of the $X_i$, but supplies no upper
bound on other physical operator coordinates.  Define
\begin{align}
 G_{ij}&=\Tr_H(X_iX_j),\label{sel15:eq:G}\\
 C_{kij}&=\Tr_H(X_kX_iX_j),\label{sel15:eq:C}\\
 C^+_{kij}&=\operatorname{Re}C_{kij}
       =\Tr_H\!\left(X_k\frac{X_iX_j+X_jX_i}{2}\right),
       \label{sel15:eq:Csym}\\
 M_{ij}&=\Tr_H(X_i^2X_j^2).
       \label{sel15:eq:M}
\end{align}
Here $G$ is real symmetric positive definite, $C^+$ is fully symmetric and
real, and $M_{ij}\ge0$.  The imaginary part of $C$ is not discarded in the
ordinary-product route.

\begin{proposition}[Meaning of the comparison data]
\label{sel15:prop:copies}
The entries of $G,C,M$ are bounded permutation-observable expectations on
respectively two, three, and four independently run copies of the same-cut
process, weighted by the original past probabilities.  Specifically, cyclic
permutation operators give $\Tr(A_1\cdots A_k)$, up to the harmless choice
of cycle orientation.  Hermitian and anti-Hermitian parts give its real and
imaginary parts.  No unknown input state is cloned and no inverse event is
used.
\end{proposition}
\begin{proof}
In a matrix-unit basis, contraction of the output indices of a cyclic
permutation with the next input indices gives
$\Tr[V_k(A_1\otimes\cdots\otimes A_k)]=\Tr(A_1\cdots A_k)$ after fixing
the orientation of $V_k$.  The real and imaginary parts of a unitary
permutation have operator norm at most one.  On a copy with recorded past
$v_i$, the normalized conditional state is $X_i/a_i$; multiplying the
conditional comparison expectation by the product of the $a_i$ recovers
\eqref{sel15:eq:G}--\eqref{sel15:eq:M}.  Independence is established by
separate preparations, not by copying one unknown state.  Independence of the
$X_i$ implies $a_i>0$.  Using $X_i,X_i,X_j,X_j$ gives $M_{ij}$.
\end{proof}

This proposition specifies a sufficient operational realization of the data,
not permission to execute SWAP or cycle comparisons in every source.  A
complete same-cut state tomography panel can instead determine the same
numbers.  Multiple-copy access, or that tomography, is additional to the
single-copy original-record experiment.  The comparison device and its copies
are not uncounted persistent memory in a simulation of the original source.

\begin{proposition}[Exact product and event-closure residuals]
\label{sel15:prop:residuals}
Put $c_{ij}=(C_{kij})_k$.  The nonnegative numbers
\begin{equation}
 \boxed{b_{ij}=M_{ij}-c_{ij}^*G^{-1}c_{ij}
       =\|(I-P_{\mathcal V})X_iX_j\|_{\HS}^2}
 \label{sel15:eq:product-defect}
\end{equation}
vanish for all $i,j$ exactly when $\mathcal V$ is closed under ordinary
multiplication.  For $Y_{e j}=\Phi_{e,*}(X_j)$ and
$g_{ej}=(\Tr_H X_iY_{ej})_i$, the nonnegative numbers
\begin{equation}
 \boxed{d_{ej}=\Tr_H Y_{ej}^2-g_{ej}^{\mathsf T}G^{-1}g_{ej}
       =\|(I-P_{\mathcal V})Y_{ej}\|_{\HS}^2}
 \label{sel15:eq:event-defect}
\end{equation}
vanish for all $e,j$ exactly when every original branch preserves
$\mathcal V$.  All event residuals require only two-copy comparisons of
original pasts and their one-event extensions.
\end{proposition}
\begin{proof}
The orthogonal projection of any matrix $Z$ onto the span of the Hermitian
basis $X_i$ has coordinates $G^{-1}(\Tr_H X_iZ)_i$ and squared norm
$c(Z)^*G^{-1}c(Z)$.  Also
$\|X_iX_j\|_{\HS}^2=\Tr_H(X_jX_i^2X_j)=M_{ij}$.
Orthogonal Pythagoras proves both identities.  Checking products of basis
elements and branch images of basis elements proves the two equivalences.
\end{proof}

\section{An autonomous positive memory from zero operational defects}
\label{sel15:memory}

\begin{theorem}[Finite ordinary-algebra and positive-predictor certificate]
\label{sel15:thm:memory}
Assume the invertible size-$q$ original panel, physical overlaps above, and
\[
 b_{ij}=0\quad\hbox{for every }i,j,\qquad
 d_{ej}=0\quad\hbox{for every original }e,j.
\]
Then $\mathcal V$ is an ordinary finite $C^*$-algebra on its support, and
every original recorded word from $\rho$ is reproduced by an exact
normalized CP instrument on an algebra isomorphic to $\mathcal V$.
The multiplication and unit are reconstructed by
\begin{equation}
 \boxed{\mu_{ij}=G^{-1}c_{ij},\qquad e=G^{-1}a,
       \qquad \sum_i e_iX_i=p,}
 \label{sel15:eq:product-unit}
\end{equation}
where $p$ is the common support of the chosen positive states.
If the reconstructed centre has dimension one and $q=n^2$, then
$\mathcal V\cong M_n(\C)$.  The minimum total Hilbert dimension of an
ordinary quantum predictor for the original record process is exactly $n$.
\end{theorem}
\begin{proof}
The span is adjoint stable.  Vanishing of the $b_{ij}$ makes it a finite
self-adjoint subalgebra, possibly without the ambient identity.  Set
$S=\sum_iX_i\succeq0$.  Its support $p$ is the union of the supports of
all $X_i$.  Each $X_i=pX_ip$, so $p$ is an identity on the span.  Since
$S$ has finite spectrum, a polynomial which vanishes at zero and equals one
at every positive eigenvalue gives $p$ as a polynomial in $S$ without a
constant term.  Thus $p\in\mathcal V$.
The inner products $\Tr_H X_ip=\Tr_H X_i=a_i$ give the unit formula;
the product formula is the projection formula, now without a residual.

The event residuals imply $\Phi_{e,*}(\mathcal V)\subseteq\mathcal V$.
Every restriction is completely positive for the inherited matrix order,
and the sum preserves the physical trace.  Since $X_1=\rho$, induction
places every original unnormalized state in this same algebra and proves
all-word equality.  No intermediate replacement state or conditional
renormalization is inserted.

Write a faithful representation of $\mathcal V$ on $pH$ as
$\bigoplus_b I_{r_b}\otimes M_{n_b}$.
A state-coordinate encoder and decoder are
\[
 \mathsf S_*(x)=\bigoplus_b(I_{r_b}/r_b)\otimes x_b,
 \qquad
 \mathsf R_*(X)=\bigoplus_b\Tr_{r_b}(p_bXp_b).
\]
They are CPTP on the indicated support, and $\mathsf R_*\mathsf S_*=\id$.
The maps $\Psi_{e,*}=\mathsf R_*\Phi_{e,*}\mathsf S_*$ form a normalized
instrument and reproduce the original process because $\mathsf S_*\mathsf
R_*$ is the identity on $\mathcal V$.  An extension of the decoder outside
$pH$ is irrelevant to this prepared-source statement.

The centre is found by the finite linear equations
$\sum_i z_i(\mu_{ij}-\mu_{ji})=0$ for every $j$.  If it has dimension one,
finite $C^*$-structure gives $\mathcal V\cong M_n$ with $q=n^2$.
The construction supplies an $n$-dimensional quantum realization.  Conversely,
every ordinary memory of total Hilbert cost $D$ has at most $D^2$ linear
coordinates, whereas the original panel has rank $n^2$.  Hence $D\ge n$.
\end{proof}

\begin{corollary}[Represented support and its redundant multiplicity]
\label{sel15:cor:support}
In the simple case of Theorem~\ref{sel15:thm:memory},
\begin{equation}
 \boxed{\Tr_Hp=a^{\mathsf T}G^{-1}a=rn,\qquad
        r=\frac{a^{\mathsf T}G^{-1}a}{n}\in\mathbb N.}
 \label{sel15:eq:physical-rank}
\end{equation}
Here $r$ is the physical representation multiplicity of the retained $M_n$,
not its number of quantum-memory coordinates.  If the initial state obeys
$\Tr_H\rho^2>1/2$, then $r=1$.  In particular a pure initial state has
no such redundancy on the retained support.
\end{corollary}
\begin{proof}
The unit coefficients give $\Tr_Hp=a^{\mathsf T}e$.
On a simple represented algebra every state has the form
$(I_r/r)\otimes\sigma$, so its purity is $\Tr\sigma^2/r\le1/r$.
The integer $r$ must therefore equal one under the stated strict bound.
\end{proof}

The certificate concerns the original preparation and its reachable retained
algebra.  It does not erase unrelated original input sectors or promise the
same reduction for arbitrary states outside that algebra.  This is why it
does not contradict V12's all-input predictor obstruction.  The matrices
$G,C,M$ are genuine physical overlaps: assigning a convenient product to
Hankel coordinates would not prove any of the residual identities.

\section{Derive completely positive symmetry from operational products}
\label{sel15:lift}

Apply V14's finite record audit on the now certified $q$-dimensional process.
It reconstructs $T_g=H^{-1}H_g$, with $T_gb=b$, $a^{\mathsf T}T_g=a^{\mathsf
T}$, and the event intertwinings.  The real matrix $T_g$ acts on the Hermitian
past-state span.  Extend it complex linearly and write
$\alpha_g(\sum_i x_iX_i)=\sum_i(T_gx)_iX_i$.

\begin{theorem}[Ordered cubic covariance certifies the physical CP lift]
\label{sel15:thm:ordered-lift}
Under Theorem~\ref{sel15:thm:memory}, suppose for every generator $s$ of
the record-symmetry group that
\begin{equation}
 \boxed{T_s^{\mathsf T}GT_s=G,\qquad
 C(T_sx,T_sy,T_sz)=C(x,y,z)\quad(x,y,z\in\R^q).}
 \label{sel15:eq:ordered-invariance}
\end{equation}
It suffices to check the displayed tensor identity on basis triples.
Then $\alpha_s$ is a trace-preserving $*$-automorphism of the retained
algebra and intertwines the original reduced instruments.  On a simple
retained $M_n$, it is unitary conjugation and hence reversible CPTP.
No separately reconstructed Choi matrix of the symmetry is needed.
\end{theorem}
\begin{proof}
For Hermitian $X,Y,Z\in\mathcal V$, pair invariance and ordered cubic
invariance give
\[
 \Tr_H\!\left[\alpha_s(X)
   \bigl(\alpha_s(Y)\alpha_s(Z)-\alpha_s(YZ)\bigr)\right]=0.
\]
The product is in $\mathcal V$ by its already certified closure.  Since
$\alpha_s$ is bijective and the trace pairing on $\mathcal V$ is
nondegenerate, the expression in parentheses is zero.  Complex linearity
extends multiplicativity to all matrices.  The Hermitian basis and reality of
$T_s$ give adjoint preservation.  A bijective algebra map preserves its unit.
The record trace identity supplies trace preservation.  Thus $\alpha_s$ is
a $*$-automorphism.  Transport of a matrix-unit system proves that every
$*$-automorphism of $M_n$ is unitary conjugation.  The intertwining already
holds in record coordinates and hence on these same physical states.
\end{proof}

\begin{theorem}[A source selector without an assumed memory cap or CP symmetry]
\label{sel15:thm:selector}
Suppose one actual finite normalized source satisfies the following finite
tests on its original past states:
\begin{enumerate}[label=\textup{(P\arabic*)},leftmargin=3em]
\item an original probability panel of size $196$ is invertible;
\item the physical product and event residuals
      \eqref{sel15:eq:product-defect}--\eqref{sel15:eq:event-defect} vanish;
\item the centre of the reconstructed multiplication has dimension one;
\item the original record law passes the finite $S_4$ symmetry audit and
      the physical pair and ordered cubic tensors satisfy
      \eqref{sel15:eq:ordered-invariance};
\item the averaged relabelled-past probability panel has rank fifteen.
\end{enumerate}
Then the original record process has an exact ordinary quantum predictor of
minimum total Hilbert dimension fourteen, with retained algebra $M_{14}$.
Its reduced original coefficient algebra is all of $M_{14}$, and its internal
invariant intersection is
\begin{equation}
 \boxed{\mathcal I\cong M_3(\C)\oplus M_2(\C)\oplus\C\oplus\C.}
 \label{sel15:eq:selection}
\end{equation}
Neither a preassigned external inventory, an upper bound on the original
ambient Hilbert dimension, nor an independently assumed CP symmetry lift is
required.  The multiple-copy product comparisons in \textup{(P2)} and
\textup{(P4)} are substantive source inputs.
\end{theorem}
\begin{proof}
Theorem~\ref{sel15:thm:memory} supplies the exact $M_{14}$ predictor and its
sharp minimum.  Theorem~\ref{sel15:thm:ordered-lift} supplies reversible CP
$S_4$ covariance on that predictor.  The rank-196 original record panel makes
its future effects span $M_{14}$; those effects lie in the algebra generated
by its original Kraus coefficient spans.  That coefficient algebra is
therefore full.  V14's twirl-rank theorem identifies the rank-fifteen average
with the invariant algebra dimension.  Its projective $S_4$ classification,
Lemma~\ref{sel14:lem:fifteen}, gives \eqref{sel15:eq:selection}.
\end{proof}

This selector replaces inherited assumptions by finite \emph{operational
closure and product tests}; it does not show that every source passes them.
A future-invisible linear direction is not removed unless the positive
predictor has actually been certified.  No algebraic reconstruction in this
section makes a sharp Read, a determinant incidence, or a symmetry pulse
physically executable.

\section{A smaller third-order route when the matrix memory is already verified}
\label{sel15:Jordan}

The fourth-order closure tests are unnecessary if the actual past-state span
is independently known to be a full matrix algebra.  There is then a useful
way to avoid the imaginary part of the cubic tensor as well.

\begin{theorem}[Pair and symmetric cubic data reconstruct the Jordan product]
\label{sel15:thm:Jordan}
Let the $X_i$ be a Hermitian basis of a verified $M_n$.  Pair and symmetric
cubic data determine the Jordan product $X\circ Y=(XY+YX)/2$ by
\begin{equation}
 \boxed{X_i\circ X_j=\sum_k\nu_{ij}^kX_k,\qquad
       \nu_{ij}=G^{-1}(C^+_{kij})_k.}
 \label{sel15:eq:Jordan}
\end{equation}
A real invertible record symmetry preserving the trace is a Jordan
$*$-automorphism if and only if it preserves $G$ and $C^+$.
\end{theorem}
\begin{proof}
The coefficient formula follows by pairing with each basis element.
If $T$ preserves $G$ and $C^+$, the same nondegenerate pairing proves
$\alpha(X\circ Y)=\alpha(X)\circ\alpha(Y)$.  The identity is fixed:
its coordinate vector is $G^{-1}a$, and trace preservation and $G$-orthogonality
force that vector to be fixed.  The map and its inverse preserve Hermitian
squares, hence the positive cone.  Conversely a Jordan $*$-automorphism of
$M_n$ has the form established in the next lemma and preserves trace, $G$
and $C^+$.
\end{proof}

\begin{lemma}[Finite Jordan alternatives and the involution trace]
\label{sel15:lem:Jordan-alternatives}
Every complex-linear unital Jordan $*$-automorphism of $M_n$, $n\ge2$, is
of exactly one of the forms
\[
 X\longmapsto UXU^*,\qquad X\longmapsto UX^{\mathsf T}U^*.
\]
If the second form is an involution, its complex linear trace on $M_n$ is
$+n$ or $-n$.  The negative case requires even $n$.  The second form is not
completely positive.
\end{lemma}
\begin{proof}
We include the finite matrix argument behind the standard classification.
Images of the diagonal matrix units are pairwise orthogonal rank-one
projections summing to the identity: positivity of the map and its inverse
preserves the minimal nonzero projections.  After a unitary conjugation they
are the original diagonal units.  Their Jordan multiplication eigenspaces
then put the image of $E_{ij}$ in $\Span\{E_{ij},E_{ji}\}$.
Write it as $a_{ij}E_{ij}+b_{ij}E_{ji}$.  Its square is zero, so
$a_{ij}b_{ij}=0$.  Its Jordan product with its adjoint is
$(E_{ii}+E_{jj})/2$, so the nonzero coefficient has modulus one.
For distinct indices, $E_{ij}\circ E_{jk}=E_{ik}/2$ forces a consistent
orientation on every pair: a mixture of the two orientations would make the
left side zero.  Connectivity of these triples gives one common orientation
when $n\ge3$; for $n=2$ there is only one pair.  The same identity makes
the phases multiplicative, hence of the form $a_{ij}=u_i\overline u_j$
or the corresponding reversed form.  A diagonal unitary removes those phases,
proving the two alternatives.

For $\alpha(X)=UX^{\mathsf T}U^*$, the identity $\alpha^2=\id$ gives
$U\overline U=\zeta I$.  Equivalently $U=\zeta U^{\mathsf T}$, which
forces $\zeta^2=1$.  The operator-space trace, calculated on matrix units,
is $\Tr(U\overline U)=\zeta n$.  An invertible antisymmetric matrix has
even size, proving the restriction.  Finally transposition applied to one
half of a normalized maximally entangled state gives the flip divided by $n$,
which has a negative eigenvalue on the antisymmetric subspace.  Unitary
conjugation does not remove that negative eigenvalue.
\end{proof}

\begin{theorem}[The recorded character removes the anti-automorphism alternative]
\label{sel15:thm:trace-lift}
On a verified matrix memory $M_n$, assume the finite $S_4$ record audit and
pair/symmetric-cubic invariance for the adjacent-transposition generators.
If, for one transposition $s$,
\begin{equation}
 \boxed{\Tr(H^{-1}H_s)\notin\{n,-n\},}
 \label{sel15:eq:anti-trace}
\end{equation}
then all the reconstructed symmetries are unitary conjugations.  In
particular the recorded value $16$ on $M_{14}$ derives the CP lift needed
by V14, without an imaginary cubic comparison or an entangled-input Choi test.
\end{theorem}
\begin{proof}
The preceding theorem makes every generator a Jordan automorphism.  The
record audit gives its involution relation.  The lemma excludes the
anti-automorphism alternative for $s$.  All transpositions are conjugate
inside $S_4$, and conjugating an automorphism by any Jordan automorphism is
still an automorphism.  Thus every transposition is an automorphism, and
transpositions generate the group.  Their unitary implementations are
reversible CPTP maps and intertwine the original instrument by the record
identities.
\end{proof}

Condition \eqref{sel15:eq:anti-trace} is sufficient, not necessary.  When it
fails, one alternative is a nonzero ordered imaginary cubic: an
anti-automorphism sends $\operatorname{Im}\Tr(XYZ)$ to its negative, while
an automorphism preserves it.  A preserved nonzero such comparison for one
transposition also excludes the anti-automorphism branch.  Pair and symmetric
cubic data alone cannot distinguish transposition from unitary conjugation.

\begin{proposition}[A positive Jordan-covariance residual]
\label{sel15:prop:Jordan-defect}
Assume $T^{\mathsf T}GT=G$ and define
$\Delta_{ijk}=C^+(Te_i,Te_j,Te_k)-C^+_{ijk}$.
Its squared tensor norm with the three inverse-Gram metrics is exactly the
sum, over orthonormal Hermitian input pairs, of the squared Hilbert--Schmidt
Jordan-intertwining errors
\[
 \alpha(X)\circ\alpha(Y)-\alpha(X\circ Y).
\]
Thus $\Delta=0$ is a finite positive zero test, not a count of symmetric
coordinates.  A nonzero certified residual disproves the proposed product
symmetry on the same physical source.
\end{proposition}
\begin{proof}
The inner product of the displayed error with $\alpha(X_k)$ is
$\Delta_{ijk}$.  Expand the input pairs and output in an orthonormal basis.
Parseval gives the asserted squared sum; transforming to the original basis
introduces one $G^{-1}$ for each tensor index.
\end{proof}

\section{A full-rank same-character countermodel for pairwise-only inference}
\label{sel15:countermodel}

The need for a product test is not an appeal to an unspecified hidden cone.
The next construction passes the exact V14 probability ranks and class trace,
and preserves physical pair overlaps, but has no positive symmetry lift on
its own minimal matrix realization.

\begin{theorem}[Original record symmetry need not preserve quantum positivity]
\label{sel15:thm:countermodel}
There exists a finite normalized instrument on $M_{14}$, from one invariant
initial state, with
\[
 \rank H=196,\qquad \rank\overline H=15,\qquad
 \Tr(H^{-1}H_g)=(196,16,4,4,4)_{[g]},
\]
whose unique record-symmetry maps preserve the physical trace and
Hilbert--Schmidt pairing but are not positive.  Its fixed linear space is
not an ordinary algebra.  All recorded branches have strictly positive Choi
matrices.  This is a countermodel in V14's arbitrary finite-instrument class,
not a six-letter calibrated historical source.
\end{theorem}
\begin{proof}
Start with the actual ordinary representation type
$(3,0,1,2,1)$ from V14 and let $\beta_g=\operatorname{Ad}U_g$ on
$M_{14}^{\rm sa}$.  For one transposition $s$, choose orthonormal vectors
$f_1,f_2$ in its trivial isotypic sector and a $-1$ eigenvector $f_3$ in
its $W$ sector.  Put
\[
 F=\frac{|f_1\rangle\langle f_2|+|f_2\rangle\langle f_1|}{\sqrt2},
 \qquad
 Q=\frac{|f_1\rangle\langle f_3|+|f_3\rangle\langle f_1|}{\sqrt2}.
\]
These are orthonormal traceless Hermitian matrices.  Let the real orthogonal
map $O$ exchange $F,Q$ and fix their orthogonal complement, including the
identity.  Extend it complex linearly and define
$\gamma_g=O\beta_gO^{-1}$.  This is a trace-preserving orthogonal
representation with exactly the same class character and fixed-space
dimension as $\beta$.

On the density with vector $(f_1+f_2+if_3)/\sqrt3$, the $3\times3$
nonzero block of $\gamma_s(\rho)$ is
\begin{equation}
 \frac13\begin{pmatrix}1&-1&i\\-1&1&i\\-i&-i&1\end{pmatrix},
 \qquad\operatorname{spec}=\{-1/3,2/3,2/3\}.
 \label{sel15:eq:counternegative}
\end{equation}
Thus $\gamma_s$ is not positive.  Its cubic trace on this state is $5/9$,
whereas $\Tr\rho^3=1$.  Also $Q=O(F)$ belongs to the fixed space,
but $Q^2$ does not: it contains a rank-one external $W$ projection, which
is not invariant, and $O$ fixes that square.  Hence the fixed space is not
closed even under squares.

It remains to realize this linear symmetry as an actual normalized quantum
record process.  Let $d=14$, $h=d^2-1=195$, and choose any orthonormal
traceless Hermitian basis $F_1,\ldots,F_h$.  Write $M=2|S_4|h=9360$.
For $g\in S_4$, $1\le a\le h$, and $\sigma\in\{+1,-1\}$, define
\begin{align}
 F_{ga}&=\gamma_g(F_a),\nonumber\\
 E_{ga\sigma}&=\frac{I+\sigma\epsilon F_{ga}}{M},
 &\omega_{ga\sigma}&=\frac{I}{d}+\sigma\eta F_{ga},\nonumber\\
 \Phi_{ga\sigma,*}(X)&=\Tr(E_{ga\sigma}X)\,\omega_{ga\sigma},
 &0<\epsilon<1,&\quad0<\eta<1/d.
 \label{sel15:eq:measureprepare}
\end{align}
The operator norm of every $F_{ga}$ is at most its HS norm one.  Hence both
$E_{ga\sigma}$ and $\omega_{ga\sigma}$ are strictly positive, their
traces have the required normalization, and $\sum E_{ga\sigma}=I$.
Every branch is CP with positive-definite Choi matrix
$E_{ga\sigma}^{\mathsf T}\otimes\omega_{ga\sigma}$.
Start at $I/d$ and relabel $(g,a,\sigma)$ by left multiplication of $g$.
Orthogonality of $\gamma$ gives
$\gamma_k\Phi_{ga\sigma,*}=\Phi_{kg,a,\sigma,*}\gamma_k$ as linear
maps.  This identity does not require positivity of $\gamma_k$.
The whole scalar record law is therefore invariant.

Take pasts and futures consisting of the empty word and the $h$ positive
outcomes $(1,a,+)$.  Their probability panel is
\begin{equation}
 H=\begin{pmatrix}
 1&M^{-1}\mathbf1^{\mathsf T}\\
 M^{-1}\mathbf1&M^{-2}(\mathbf1\mathbf1^{\mathsf T}+\epsilon\eta I_h)
 \end{pmatrix},\qquad
 \boxed{\det H=(\epsilon\eta/M^2)^h>0.}
 \label{sel15:eq:counter-H}
\end{equation}
The corresponding past states span $M_d$.  Relabelling implements precisely
$\gamma_g$ on their coordinates, so the averaged panel has rank fifteen and
all class traces are those claimed.  Uniqueness of the full-rank record
symmetry excludes a different positive map with the same action on the pasts.
All pair overlaps are preserved, while the cubic counterexample proves that
the next tensor fails.  The rank $196$ also makes fourteen the minimum
positive Hilbert-memory cost, so this is not an invisible-spectator example.
\end{proof}

The construction does not claim equality of its full probabilities with
V13's witness: it has the \emph{same selection summaries} and still fails the
missing physical hypothesis.  It also has more outcomes than the original
six-edge writer.  Its purpose is the exact logical boundary of a selector
whose stated class allows general finite instruments.

\section{Recover the represented internal type from a third-order contraction}
\label{sel15:represented}

Once one of the positive-lift routes passes, the invariant space is an
ordinary algebra.  Let $Y_1,\ldots,Y_r$ be any Hermitian basis of it obtained
from independent averaged past-state columns.  Use the trace $\tau$ of the
\emph{minimal retained matrix memory}.  In Theorem~\ref{sel15:thm:memory},
this is $\tau=\Tr_H/r_0$ with $r_0$ determined by
\eqref{sel15:eq:physical-rank}; in the multiplicity-free benchmark it is
simply $\Tr_H$.  Put
\[
 g_{ab}=\tau(Y_aY_b),\qquad
 c^+_{kab}=\tau(Y_k(Y_a\circ Y_b)).
\]
Define the coordinate vector $t_k=\sum_{a,b}(g^{-1})_{ab}c^+_{kab}$,
$z=g^{-1}t$, and the represented operator
\begin{equation}
 \boxed{Z_{\mathcal I}=\sum_a z_aY_a
     =\sum_{a,b}(g^{-1})_{ab}(Y_a\circ Y_b),
     \qquad Q_{\mathcal I}=z^{\mathsf T}gz.}
 \label{sel15:eq:quadratic-central}
\end{equation}
Both are basis independent.  Only pair and symmetric cubic overlaps are
needed; $Q_{\mathcal I}$ does not require a new fourth-copy measurement.

\begin{theorem}[Trace-Casimir discriminator of the represented multiplicities]
\label{sel15:thm:Casimir}
If the internal algebra is represented as
$\bigoplus_b I_{d_b}\otimes M_{m_b}$ on the retained matrix memory, then
\begin{equation}
 \boxed{Z_{\mathcal I}=\sum_b\frac{m_b}{d_b}p_b,
 \qquad\tau(Z_{\mathcal I})=\sum_bm_b^2,
 \qquad Q_{\mathcal I}=\sum_b\frac{m_b^3}{d_b}.}
 \label{sel15:eq:Casimir-formula}
\end{equation}
For V14's recorded character $(196,16,4,4,4)$,
\begin{equation}
 \boxed{
 Q_{\mathcal I}=\begin{cases}
 61/2,&\text{support ranks }(3,6,2,3),\\
 37/2,&\text{support ranks }(9,2,1,2).
 \end{cases}}
 \label{sel15:eq:Casimir-two}
\end{equation}
In particular this scalar resolves the complete represented-type ambiguity
left by the class-character summary.  It does not distinguish an ordinary
lift from its sign twist, which implements the same automorphisms.
\end{theorem}
\begin{proof}
The inverse Gram contracts two copies of a basis covariantly, so it is enough
to use an orthonormal Hermitian basis.  On $I_d\otimes M_m$ this is
$I_d\otimes B_a/\sqrt d$, where $(B_a)$ is HS orthonormal in $M_m^{\rm sa}$.
The diagonal units and the two normalized Hermitian off-diagonal units for
each pair give $\sum_aB_a^2=mI_m$.  Therefore their square sum is
$(m/d)p$.  Adding the blocks proves the operator identity and its first and
second trace moments.  V14's two inventories give respectively
\[
 27+\frac83+\frac12+\frac13=\frac{61}{2},\qquad
 9+8+1+\frac12=\frac{37}{2}.
\]
There are no other inventories in that recorded-character fibre by
Theorem~\ref{sel14:thm:character-fibre}.  Thus the two different scalar values
resolve its two represented types.
\end{proof}

On the first type the four eigenvalues of $Z_{\mathcal I}$ are
$3,2/3,1/2,1/3$.  Its unique colour-block central unit is reconstructed
without first choosing the colour support by
\begin{equation}
 p_{(3)}=
 \frac{(Z_{\mathcal I}-\tfrac23p)(Z_{\mathcal I}-\tfrac12p)
                   (Z_{\mathcal I}-\tfrac13p)}{140/9}.
 \label{sel15:eq:colour-central-projector}
\end{equation}
Products of this one element may equivalently be computed with its Jordan
product.  The result has retained trace three.  This is a recovered algebra
projection, not automatically an executable sharp Read; it still does not
specify particular private rays or a determinant-bearing tensor.

\section{Evaluate the unchanged full-rank source}
\label{sel15:checks}

Use exactly the V13 source at $(u,v,\zeta)=(3/5,4/5,i)$ and the original
196 past words retained through V14.  No coefficient, preparation, word
alphabet, or source parameter is changed.  The present calculation concerns
extra comparisons of the resulting physical states, not newly inserted
single-copy controls.

\begin{theorem}[Exact product and represented-type certificate for the inherited witness]
\label{sel15:thm:witness}
For this source, all product and event-closure residuals in
\eqref{sel15:eq:product-defect}--\eqref{sel15:eq:event-defect} vanish.
Its recovered centre has dimension one, its support trace is fourteen, and
its moment tensors satisfy both the ordered and symmetric cubic covariance
identities.  On the fifteen-dimensional invariant span,
\begin{equation}
 \boxed{Q_{\mathcal I}=\frac{61}{2},\qquad
      Z_{\mathcal I}=3I_3\oplus\frac23I_6\oplus\frac12I_2
                              \oplus\frac13I_3.}
 \label{sel15:eq:witness-Casimir}
\end{equation}
Thus its $M_3$ summand has physical support rank three, not nine.  Both the
bound-free ordered-product selector and the bounded symmetric-cubic selector
apply to this same mathematical witness.
\end{theorem}
\begin{proof}
The inherited nonzero original probability minor makes the physical past
states a basis of the actual $M_{14}$.  Hence the residuals vanish exactly,
not just modulo a prime; its centre is scalar.  Normalized source covariance
and unitary invariance of traces give every cubic and pair identity.
The initial state is pure, so the physical support has no representation
redundancy.  These analytic facts provide exact upper identifications for the
finite calculations, not new assumptions about historical data.

Here is a product-sensitive reconstruction of the final scalar.  Regenerate
the physical past matrices with V13's exact number-ring arithmetic, and
average each over the 24 actual symmetry conjugations.  Select the first
fifteen independent averaged columns by the inherited ordered elimination
modulo 97.  Compute their full $15\times15$ physical Gram and the 680
unordered symmetric cubic entries.  The inverse-Gram contraction
\eqref{sel15:eq:quadratic-central} gives $\tau Z_{\mathcal I}=15$ and
\[
 Q_{\mathcal I}\equiv79\pmod{97}.
\]
The two candidate rational values in \eqref{sel15:eq:Casimir-two} reduce to
$79$ and $67$ respectively, so the residue identifies $61/2$ exactly within
the already proved character fibre.  It is not a floating-point inverse or
an attempt to lift an unrestricted rational from its residue.

The companion certificate records the actual selected averaged pasts,
Gram, cubic entries, inverse-Gram contraction and central-projector trace.
Direct substitution in the inherited represented blocks gives the displayed
operator in \eqref{sel15:eq:witness-Casimir}; its minimal polynomial also
verifies \eqref{sel15:eq:colour-central-projector}.  These independent
checks agree with the reconstruction rather than replacing it.
\end{proof}

The maximum old past depth is six.  Event-closure comparisons use pasts of
depth at most seven.  Up to four \emph{separate} copies are compared for the
ordinary-product closure test; this is not a length-24 single-system word.
The restricted invariant contraction uses only 120 unordered pairs and 680
unordered symmetric triples after its fifteen basis states have been
identified.  Certifying the full symmetry lift on the 196-dimensional space
requires the corresponding full tensor identities or an equivalent exact
structural proof; the restricted 15-dimensional table alone does not prove
that full lift.

\section{Error budgets and the retained physical calibration}
\label{sel15:errors}

All zero-defect theorems above are exact.  An approximate small residual does
not prove exact closure or a CP automorphism of unchanged algebraic type.
Conversely, a certified positive lower bound on a residual is a genuine
obstruction.  The following estimate is useful after the exact structural
hypotheses and character fibre have independently been established.

\begin{proposition}[Certified interval for the represented-type scalar]
\label{sel15:prop:scalar-error}
On one fixed invariant basis of size $r$, suppose
$\|\widehat g-g\|_\op\le\delta_g$ and
$\|\widehat c^+-c^+\|_F\le\delta_c$, with
$s=\lambda_{\min}(\widehat g)>\delta_g$.
Set
\[
 u=(s-\delta_g)^{-1},\quad
 d_u=\frac{\delta_g}{s(s-\delta_g)},\quad
 \widehat t_k=\sum_{a,b}(\widehat g^{-1})_{ab}\widehat c^+_{kab},
 \quad\widehat z=\widehat g^{-1}\widehat t.
\]
Define
\begin{align}
 e_t&=\sqrt r\bigl(\delta_c u+\|\widehat c^+\|_F d_u\bigr),\nonumber\\
 e_z&=u e_t+d_u\|\widehat t\|,\nonumber\\
 e_Q&=\delta_g(\|\widehat z\|+e_z)^2
       +\|\widehat g\|_\op e_z(2\|\widehat z\|+e_z).
 \label{sel15:eq:Qerror}
\end{align}
Then $|Q_{\mathcal I}-\widehat z^{\mathsf T}\widehat g\widehat z|\le e_Q$.
Within the exact two-value fibre \eqref{sel15:eq:Casimir-two},
\[
 |\widehat z^{\mathsf T}\widehat g\widehat z-61/2|+e_Q<6
\]
certifies the first represented type.  The analogous test around $37/2$
certifies the second.
\end{proposition}
\begin{proof}
The resolvent identity for inverses gives
$\|g^{-1}-\widehat g^{-1}\|\le d_u$ and $\|g^{-1}\|\le u$.
Contracting a tensor with a matrix is bounded by the product of their
Frobenius norms.  Since $\|g^{-1}\|_F\le\sqrt r\,u$, the two changes
in $t=c^+:g^{-1}$ give $\|t-\widehat t\|\le e_t$.
The inverse identity then gives $\|z-\widehat z\|\le e_z$.
Expand $z^{\mathsf T}gz-\widehat z^{\mathsf T}\widehat g\widehat z$ by
first changing $g$ and then $z$; the triangle inequality gives
\eqref{sel15:eq:Qerror}.  The two possible exact values are twelve apart,
so an interval strictly within distance six of one cannot contain the other.
\end{proof}

If the invariant basis itself is estimated from $H^{-1}\overline H$, its
coordinate errors must first be propagated into $\delta_g,\delta_c$.
An arbitrary coordinate projector is not assumed orthogonal.  The V14
inverse-panel estimate remains applicable.  This proposition does not turn
noisy agreement into exact covariance, exact product closure, or exact zero
residuals.

For raw opportunity histories of length $l_i$, the old factors are
$X_i^{\rm raw}=\vartheta^{l_i}X_i$.  With $D=\diag(\vartheta^{l_i})$, the
physical panels therefore scale as
\[
 G^{\rm raw}=DGD,\qquad
 C^{\rm raw}_{kij}=D_{kk}D_{ii}D_{jj}C_{kij},\qquad
 M^{\rm raw}_{ij}=D_{ii}^2D_{jj}^2M_{ij}.
\]
The reconstructed product is related by the matching basis change, and
$Z_{\mathcal I}$ and its correctly calibrated trace scalar are invariant.
Their conditioning and statistical errors are not.  Accepted-time stopping
retains V13's opportunity budget; it does not make multi-copy comparisons or
rare past preparations free.

\section{The remaining upstream source task}
\label{sel15:status}

This development closes neither a historical measurement gap nor the full
Standard-Model programme.  It gives a finite route in which two formerly
independent assumptions are replaced by actual scalar tests.  The ordinary
product and event residuals can establish an autonomous positive memory and
its dimension; ordered cubic covariance can establish its CP symmetry.
When the matrix memory is already known, the symmetric cubic route and the
recorded transposition trace avoid the imaginary part and fourth-order
closure tests.  The additional central trace contraction removes V14's
represented-type ambiguity for the actual full-rank witness.

The countermodel shows why further single-copy record ranks, or pairwise
state overlaps alone, cannot replace this product information in the general
finite-instrument class.  It has the exact same selector ranks and symmetry
character while its uniquely reconstructed linear symmetry is nonpositive.
The obstruction is not a missing numerical tolerance or an invisible
spectator.

The next historical input should therefore be a same-cut operational
\emph{product-comparison panel} on the original past states, or an already
supplied positive process realization from which that panel follows.  The
source must justify tensor-copy identification and readout availability.
The product of a categorical coefficient quotient is not silently substituted
for that of the acted-on quantum state algebra.  Once the physical retained
algebra and its symmetry pass, the determinant/matter tests must still be
performed on the same source.  Abstract matrix sizes, represented colour
support, named type/private rays, a protected determinant tensor, and the
charged matter packet remain distinct assertions.

\paragraph{Verification and preservation.}
V14 is retained byte for byte before its final document terminator.  The new
script regenerates the inherited states and probability certificate, verifies
the restricted overlap contraction and its exact residue separation, checks
the transposition alternatives, the nonpositive same-character construction,
finite product and event residual identities, and the perturbation formulas.
General selection claims are proved above; the computational checks are
reproducible exact witnesses, not proof-assistant certification, independent
peer review, or fabricated historical data.

\begin{thebibliography}{99}
\raggedright
\bibitem[23]{sel15-Jordan}
I.~Gogi\'c and M.~Toma\v sevi\'c,
\emph{Classification of Jordan multiplicative maps on matrix algebras},
Aequationes Mathematicae (2025),
\href{https://doi.org/10.1007/s00010-025-01208-y}{doi:10.1007/s00010-025-01208-y};
\href{https://arxiv.org/abs/2503.24094}{arXiv:2503.24094}.
Cited for established Jordan-preserver background; the finite linear
$*$-preserving argument needed here is included explicitly.
\bibitem[24]{sel15-Ekert}
A.~K.~Ekert, C.~Moura Alves, D.~K.~L.~Oi, M.~Horodecki, P.~Horodecki and
L.~C.~Kwek,
\emph{Direct estimations of linear and non-linear functionals of a quantum
state}, Phys.\ Rev.\ Lett.\ \textbf{88} (2002), 217901;
\href{https://doi.org/10.1103/PhysRevLett.88.217901}{doi:10.1103/PhysRevLett.88.217901}.
Cited for physical multi-copy moment comparisons, not as an occurrence
assertion about this source.
\end{thebibliography}
\addtocontents{toc}{\protect\endgroup}
% END OF SOURCE-SELECTION V15 DEVELOPMENT BLOCK.
% Preserve V1--V15 and append subsequent work before the final terminator.


\clearpage
\begin{center}
{\Large\bfseries Source-selection continuation V16}\\[.4em]
{\large Determinant incidence, coherent character differences,\\
and the tensor-order boundary of the retained memory}
\end{center}
\addtocontents{toc}{\protect\begingroup}
\addtocontents{toc}{\protect\renewcommand{\protect\numberline}{\protect\makebox[2.1em][l]}}

\section{Move upstream of the determinant handoff}
\label{sel16:context}

V15 identifies a finite operational route to the internal algebra and resolves
its represented colour-support ambiguity.  Repeating that reconstruction does
not decide the next structural question: what precisely does a physical source
have to retain for a determinant incidence to restrict its unitary group?  This
question has two distinct parts.  A fixed multilinear tensor is not the same
datum as the completely positive map defined by its Kraus ray.  In addition,
an endomorphism of the retained single-system memory is not automatically a
map between the tensor types used by the determinant incidence.

We keep the internal-algebra and product reconstruction of
Theorems~\ref{sel15:thm:selector} and~\ref{sel15:thm:witness} as inputs.
The determinant-line isomorphism, its fixed-tensor stabilizer, and its possible
coexistence with a neutral Clifford shadow are also inherited:
\cite{Attack}, \src{thm:V110-det-incidence-iso},
\src{thm:V110-central-character}, and \src{thm:V110-one-ray}.
Those results do not need another proof of colour, weak, or tensor-line
existence.  What is added below is the exact criterion for extracting their
\emph{fixed-tensor conclusion from a recorded CP source}, including the global
disconnected alternatives, an error margin, and a same-carrier tensor-order
obstruction.  A known protected tensor remains sufficient; nothing below
invalidates its stabilizer theorem.

\paragraph{Scope of the group being tested.}
Throughout this block, $G_0=\mathrm U(C)\times\mathrm U(W)$, where
$\dim C=3$ and $\dim W=2$, is the candidate block-unitary group acting on
already identified internal carriers.  We test the symmetry of an
\emph{additional typed incidence packet}, with its original outcome labels
fixed.  We do not identify this stabilizer with the commutant of every raw
source coefficient.  The earlier raw source algebra can be all of $M_{14}$
while its internal invariant intersection is reducible.  Nor do we identify a
finite determinant-incidence phase with the determinant/Pfaffian phase of a
continuum fermionic measure.

\paragraph{Established background and new specialization.}
Phase references, symmetry modes of CP maps, and the loss of coherent-control
information under passage to a channel are established subjects
\cite{sel16-MS,sel16-GS,sel16-AFCB}.  The proofs below use their elementary
finite-dimensional content, not a claim of priority for a general resource
theory.  The new development relative to the audited programme is the marked
character-lattice test for the determinant handoff, its transfer to the
inherited common Clifford--determinant ray, the exact normalized comparisons,
and the tensor-order calculation on V15's actual represented type.

\section{A protected tensor and its CP ray have different stabilizers}
\label{sel16:rays}

Let $R_{\rm in},R_{\rm out}$ be specified unitary representations of $G_0$
on actual input and output spaces.  For an amplitude $B:H_{\rm in}\to
H_{\rm out}$ write
\[
 \mathcal J_B(X)=BXB^*,\qquad
 |B\rangle\!\rangle\in H_{\rm out}\otimes\overline{H_{\rm in}},\qquad
 J_B=|B\rangle\!\rangle\langle\!\langle B|.
\]
The Choi vectorization is output-first; its covariance representation is
$\mathcal R(g)=R_{\rm out}(g)\otimes\overline{R_{\rm in}(g)}$.
Set
\[
 \chi(g)=\det U\det V,\qquad g=(U,V)\in G_0.
\]

\begin{proposition}[Fixed amplitude versus phase-insensitive channel]
\label{sel16:prop:ray}
Suppose $B\ne0$ and
\begin{equation}
 R_{\rm out}(g) B R_{\rm in}(g)^*=\chi(g)^qB,
 \qquad q\in\mathbb Z.
 \label{sel16:eq:weight-B}
\end{equation}
The stabilizer of the \emph{fixed tensor} $B$ is $\ker\chi^q$.
For $q\ne0$ this is generally smaller than $G_0$.  In contrast,
\begin{equation}
 \boxed{\mathcal R(g)J_B\mathcal R(g)^*=J_B
 \quad(g\in G_0).}
 \label{sel16:eq:ray-invariance}
\end{equation}
Thus the CP map is covariant under all of $G_0$.  Its complete Choi matrix,
all output states for arbitrary inputs, and any tensor-copy experiments that
use only that CP map do not specify a coherent reference to another amplitude.
\end{proposition}
\begin{proof}
Fixing the nonzero vector in \eqref{sel16:eq:weight-B} requires
$\chi(g)^q=1$.  Its rank-one outer product instead multiplies by
$\chi(g)^q\overline{\chi(g)^q}=1$.  Equality of Choi matrices is equality
of the CP maps on every input and after every ancillary amplification.
Tensoring or composing identical maps cannot distinguish their Kraus phases.
A controlled superposition with another amplitude is additional data, not a
function of $J_B$ alone.
\end{proof}

For the inherited determinant-incidence tensor
\[
 \Theta_\tau:C\otimes W\otimes W\longrightarrow\Lambda^2C^*,
 \qquad \tau\in(\det C\otimes\det W)^*\setminus\{0\},
\]
the character in \eqref{sel16:eq:weight-B} is $\chi^{-1}$.  Its fixed-tensor
stabilizer is therefore $S(\mathrm U(3)\times\mathrm U(2))$, as already
proved in \cite{Attack,Duality}.  But the isolated map
$X\mapsto\Theta_\tau X\Theta_\tau^*$ is $G_0$-covariant.  Its positive
alternating mass certifies a nonzero \emph{character ray}, not by itself a
smaller CP-covariance group.

\begin{remark}[Overall phase freedom is not projective invariance]
Multiplying a fixed tensor by a constant $c\ne0$ does not change its
fixed-vector stabilizer.  This is the phase independence stated by the
inherited determinant theorem.  Allowing a \emph{group-dependent} phase when
comparing the tensor with its transform is a different equivalence relation;
it gives the ray stabilizer.  We keep these statements separate rather than
changing the meaning of the inherited protected relation.
\end{remark}

\begin{proposition}[No primitive determinant reference from covariant processing]
\label{sel16:prop:no-free-phase}
Fix a subgroup $G_1\subseteq G_0$.  If all admitted CP event maps are
$G_1$-covariant, all prepared ancillary states are $G_1$-invariant, and all
terminal contractions are $G_1$-covariant, then every finite adaptive protocol
formed using these operations, tensor products, and classical outcome
postprocessing remains $G_1$-covariant on each retained label.  In particular,
character-homogeneous Kraus-ray primitives do not acquire a determinant phase
reference merely by increasing temporal word depth or comparing more copies.
\end{proposition}
\begin{proof}
Covariance is preserved under composition and tensor product.  Invariant
ancillary preparation and equivariant discard preserve it as well.  At a
classical decision node each conditional path is a composition of covariant
maps.  A retained outcome is a sum with nonnegative, input-independent
classical weights, which preserves the covariance equation.  Induction on
the finite protocol proves the assertion.  An asymmetric control state,
phase-sensitive contraction, or coherent control not specified by those maps
would violate a displayed premise and is not included in the induction.
\end{proof}

\section{The exact marked character-lattice compiler}
\label{sel16:lattice}

The preceding distinction has a finite positive test.  Let
$\pi:G_0\to\mathbb T^2$ be $(U,V)\mapsto(\det U,\det V)$.
For $w=(w_1,w_2)\in\mathbb Z^2$ put
$\chi_w(g)=(\det U)^{w_1}(\det V)^{w_2}$.
Consider a specified finite character subspace of the physical amplitude
space,
\[
 \mathscr V=\bigoplus_{w\in\mathcal W}\mathscr V_w,
 \qquad \mathcal R(g)|_{\mathscr V_w}=\chi_w(g)I.
\]
Thus $\mathrm{SU}(3)\times\mathrm{SU}(2)$ acts trivially on these
intertwiner spaces; their labels are verified tensor transformation laws, not
weights assigned to fit the desired subgroup.  Let $P_w$ be their orthogonal
projections.  Every original retained outcome $y$ has a CP Choi operator
$J_y\succeq0$ supported in $\mathscr V$.  Define
\begin{align}
 J_{wz}^{(y)}&=P_wJ_yP_z,\nonumber\\
 L_{\rm rec}
 &=\operatorname{span}_{\mathbb Z}
    \{w-z:J_{wz}^{(y)}\ne0\text{ for some }y\}
       \subseteq\mathbb Z^2.
 \label{sel16:eq:lattice}
\end{align}

\begin{theorem}[Exact operational determinant-character stabilizer]
\label{sel16:thm:lattice}
The subgroup preserving every recorded CP operation is
\begin{equation}
 \boxed{G_{\rm inc}
 =\{g:\mathcal R(g)J_y\mathcal R(g)^*=J_y\ \forall y\}
 =\pi^{-1}(\operatorname{Ann}L_{\rm rec}).}
 \label{sel16:eq:stabilizer}
\end{equation}
Here $\operatorname{Ann}L=\{z\in\mathbb T^2:z^\ell=1\ \forall\ell\in L\}$.
If the nonzero Smith invariants of $L$ are $s_1,\ldots,s_k$, with
$k\le2$ and $s_i\mid s_{i+1}$, then
\begin{equation}
 G_{\rm inc}/(\mathrm{SU}(3)\times\mathrm{SU}(2))
 \cong\mathbb T^{2-k}\times\prod_{i=1}^k\mu_{s_i}.
 \label{sel16:eq:smith}
\end{equation}
In particular,
\begin{equation}
 \boxed{G_{\rm inc}=S(\mathrm U(3)\times\mathrm U(2))
 \quad\Longleftrightarrow\quad
 L_{\rm rec}=\mathbb Z(1,1).}
 \label{sel16:eq:primitive}
\end{equation}
\end{theorem}
\begin{proof}
The $(w,z)$ block of the transformed Choi operator is
$\chi_{w-z}(g)J_{wz}^{(y)}$.  Orthogonal block decomposition precludes
cancellation between different pairs.  A nonzero block is fixed exactly when
that character equals one.  Requiring this on all labels yields
\eqref{sel16:eq:stabilizer} and its integer-span formulation.
Integer row and column operations put a generator matrix of $L$ into Smith
normal form.  Row operations change generators; unimodular column operations
are automorphisms of $\mathbb T^2$.  Its annihilator then has equations
$z_i^{s_i}=1$ for $i\le k$ and no other equations, proving
\eqref{sel16:eq:smith}.  The map $\pi$ is onto with kernel
$\mathrm{SU}(3)\times\mathrm{SU}(2)$.
Finally the integer characters vanishing on $\{z_1z_2=1\}$ are exactly
$\mathbb Z(1,1)$: substitute $(z,z^{-1})$ into $z_1^az_2^b$.
The Smith description also shows that a subgroup has precisely the original
lattice as the set of characters vanishing on its annihilator.  This proves
both directions of \eqref{sel16:eq:primitive}.
\end{proof}

\begin{corollary}[Powers of the determinant retain a discrete ambiguity]
\label{sel16:cor:gcd}
Suppose all weights are $q(1,1)$ with $q\in\mathbb Z$.  Let $d$ be the
positive greatest common divisor of all differences $q-r$ whose marked Choi
blocks are nonzero.  If there are none, $G_{\rm inc}=G_0$.  Otherwise
\[
 \boxed{G_{\rm inc}=\{g:\chi(g)^d=1\}.}
\]
It has $d$ connected components, with identity component
$S(\mathrm U(3)\times\mathrm U(2))$.  Equality with that connected group
requires $d=1$; a nonzero quadratic phase test can leave $d=2$.
\end{corollary}
\begin{proof}
The lattice is $d\mathbb Z(1,1)$ by Bezout's identity.  Each allowed value
of $\chi$ is a distinct $d$th root of unity.  The kernel of $\chi$ is
connected, since it is the image of the connected group
$\mathrm{SU}(3)\times\mathrm{SU}(2)\times\mathbb T$ under the surjection
$(A,B,z)\mapsto(z^2A,z^{-3}B)$.  Its fibres are translates of that kernel,
which proves the component count.  The sixfold kernel of this latter covering
map is the inherited $\mathbb Z_6$ quotient; it is \emph{not} the extra
component group $\mu_d$ in the present statement.
\end{proof}

For example, nonzero coherences at differences $2(1,1)$ and $3(1,1)$ together
select the primitive constraint, whereas either difference alone does not.
Replacing an integer lattice by its real span would lose precisely this
finite-group information.

\begin{proposition}[Forgetting labels can erase the constraint]
\label{sel16:prop:marked}
For the coarse Choi operator $J=\sum_yJ_y$, its character-difference lattice
satisfies $L_{\rm coarse}\subseteq L_{\rm rec}$, and its covariance group
contains $G_{\rm inc}$.  The containment can be strict even when each marked
outcome is a single coherent ray.  Exact refinements of an already
rank-one \emph{homogeneous} Choi outcome cannot enlarge its zero difference
lattice.
\end{proposition}
\begin{proof}
A nonzero block of a sum has at least one nonzero summand.  Thus every coarse
difference is present in the marked lattice.  Cancellation of nonzero blocks
between labels is possible and is exhibited below.  For the last assertion,
if positive matrices sum to $|B\rangle\!\rangle\langle\!\langle B|$, each
has range in $\mathbb CB$ and is a nonnegative scalar multiple of that
matrix.  Homogeneous character support is unchanged.  This uses, rather than
replaces, V6's rank-one refinement obstruction.
\end{proof}

\section{A neutral shadow makes the determinant condition observable}
\label{sel16:anchor}

Suppose an actual incidence packet has only two character spaces:
$\mathscr V_0$ with trivial character and $\mathscr V_1$ with character
$\chi^{-1}$.  With the original labels retained, write
\begin{equation}
 J_y=\begin{pmatrix}A_y&C_y\\C_y^*&B_y\end{pmatrix}\succeq0,
 \qquad h_{\det}=\sum_y\|C_y\|_{\HS}^2.
 \label{sel16:eq:two-block}
\end{equation}
The block projections are analysis maps on an identified physical Choi space;
using them in a calculation is not a claim that new projected outcomes occur.

\begin{theorem}[Phase-free scalar test for an anchored determinant constraint]
\label{sel16:thm:anchor}
On this two-character packet,
\begin{equation}
 \boxed{h_{\det}>0\iff G_{\rm inc}=\ker\chi,
 \qquad h_{\det}=0\iff G_{\rm inc}=G_0.}
 \label{sel16:eq:anchor}
\end{equation}
The scalar is reconstructed without choosing a phase of either shadow:
\begin{equation}
 \boxed{h_{\det}
 =\frac12\sum_y\bigl(\Tr J_y^2-\Tr A_y^2-\Tr B_y^2\bigr).}
 \label{sel16:eq:purity}
\end{equation}
It is invariant under changes of Kraus frame and unitary changes of the
identified character coordinates.  For a general packet with more weights,
a nonzero neutral--determinant block establishes only
$G_{\rm inc}\subseteq\ker\chi$; the complete lattice decides equality.
\end{theorem}
\begin{proof}
The only possible nonzero differences are $\pm(1,1)$.  Their occurrence is
exactly $C_y\ne0$ for some $y$, so Theorem~\ref{sel16:thm:lattice} gives
\eqref{sel16:eq:anchor}.  Squaring the block matrix and taking its trace gives
$\Tr J_y^2=\Tr A_y^2+\Tr B_y^2+2\Tr C_yC_y^*$.
Choi operators are invariant under Kraus changes; orthogonal block norms are
invariant under the stated coordinate changes.  Additional weight
differences can impose extra equations, proving the final qualification.
\end{proof}

\begin{proposition}[Coherence budget and the common-ray equality case]
\label{sel16:prop:budget}
For one nonzero two-character packet put $\mu=\Tr A$, $\nu=\Tr B$.
Then
\begin{equation}
 0\le\|C\|_{\HS}^2\le\mu\nu.
 \label{sel16:eq:budget}
\end{equation}
When $\mu,\nu>0$, equality at the upper bound holds if and only if $J$ has
rank one.  In particular one actually coherent neutral--determinant source
ray with nonzero projections automatically satisfies the primitive anchor
criterion.  No separately added phase reference is required in that case.
\end{proposition}
\begin{proof}
Factor $J=LL^*$ and write its two block rows as $L_0,L_1$.  Then
$A=L_0L_0^*$, $B=L_1L_1^*$ and $C=L_0L_1^*$.
With $P=L_0^*L_0$, $Q=L_1^*L_1$,
\[
 \|C\|_{\HS}^2=\Tr(PQ)
 \le\lambda_{\max}(P)\Tr Q\le\Tr P\Tr Q=\mu\nu.
\]
If $\mu,\nu>0$, equality forces $P$ to have rank one and $Q$ to be
supported on its same one-dimensional range.  Thus both $L_i$ have that
single column support and $J$ has rank one.  Conversely a single column
$L=(b,t)^{\mathsf T}$ gives equality immediately.
\end{proof}

\begin{corollary}[Transfer of the inherited shared Clifford--determinant ray]
\label{sel16:cor:shared}
Suppose the common source ray of \cite{Attack}, \src{thm:V110-one-ray}, is
realized by one physical Choi ray on coherently identified typed input/output
ports, with a gauge-neutral normalized Clifford amplitude $b$ and a normalized
determinant-character amplitude $t$.  For
$k=\sqrt\mu\,b+\sqrt\nu\,t$, where $\mu,\nu>0$, its marked Choi packet
satisfies
\[
 \boxed{C=\sqrt{\mu\nu}\,|b\rangle\!\rangle\langle\!\langle t|,
 \quad h_{\det}=\mu\nu>0.}
\]
It therefore supplies the determinant constraint operationally as well as
as a fixed-tensor statement.  A direct sum of \emph{classically recorded}
CP shadows with the same two masses instead has $C=0$.
\end{corollary}
\begin{proof}
The two characters are distinct, so the Choi vectors are orthogonal.  Expand
$|k\rangle\!\rangle\langle\!\langle k|$ in those two orthogonal spaces
and apply Theorem~\ref{sel16:thm:anchor}.  Recording the two contributions
as different classical alternatives keeps only their diagonal Choi blocks.
\end{proof}

The corollary requires a \emph{physical common Choi ray}, not merely a
one-dimensional domain for an abstract source synthesis.  The inherited
common-ray construction remains valid.  Its transfer to an operational channel
uses the stated coherent input/output identifications; two marginal CP masses
alone do not supply them.  In particular a neutral Clifford amplitude must
really be neutral under the internal group on that same carrier.

If each diagonal block is itself a rank-one ray, every positive joint packet
with those fixed diagonal blocks has the form
\begin{equation}
 J_\eta=\begin{pmatrix}\mu&\eta\sqrt{\mu\nu}\\
                     \overline\eta\sqrt{\mu\nu}&\nu\end{pmatrix},
 \qquad |\eta|\le1,
 \label{sel16:eq:visibility}
\end{equation}
in their normalized ray basis.  Thus identical positive determinant and
Clifford masses admit both $h_{\det}=0$ and
$h_{\det}=\mu\nu|\eta|^2>0$.  Positivity gives the bound on $\eta$ by the
two-by-two determinant.  These are CP packets; a trace-nonincreasing scale or
complete-instrument normalization must still be checked before treating one
as a source event.  The next example includes that normalization explicitly.

\section{Global phase rigidity and the exact cost of deleting its coherence}
\label{sel16:rigidity}

\begin{theorem}[Calibrated covariance residual and its central null direction]
\label{sel16:thm:residual}
For the general character packet define
$\mathfrak d(g)^2=\sum_y\|\mathcal R(g)J_y\mathcal R(g)^*-J_y\|_{\HS}^2$.
Then
\begin{equation}
 \boxed{\mathfrak d(g)^2
 =\sum_{y,w,z}|\chi_{w-z}(g)-1|^2\|J_{wz}^{(y)}\|_{\HS}^2.}
 \label{sel16:eq:residual}
\end{equation}
For \eqref{sel16:eq:two-block}, if $\chi(g)=e^{i\phi}$, this becomes
\begin{equation}
 \boxed{\mathfrak d(g)^2=8h_{\det}\sin^2(\phi/2).}
 \label{sel16:eq:phase-margin}
\end{equation}
On central matrices $(e^{i\alpha}I_3,e^{i\beta}I_2)$ its quadratic part is
$2h_{\det}(3\alpha+2\beta)^2$.  Its rank-one quadratic-form matrix is
\[
 2h_{\det}\begin{pmatrix}9&6\\6&4\end{pmatrix},
\]
with positive eigenvalue $26h_{\det}$ and kernel spanned by $(-2,3)$.
This is a source-calibrated symmetry diagnostic, not a physical gauge-boson
mass or a normalization of a continuum kinetic term.
\end{theorem}
\begin{proof}
The blocks $\mathscr V_z\to\mathscr V_w$ are mutually orthogonal for the
Hilbert--Schmidt product.  Subtract the original block after multiplying by
its character and sum squared norms.  For two weights the two adjoint cross
blocks contribute $2|e^{i\phi}-1|^2h_{\det}$, proving the formula.
Taylor expansion at zero and $\phi=3\alpha+2\beta$ give the quadratic
matrix; its trace and null vector give the stated eigenvalue and kernel.
\end{proof}

\begin{theorem}[Exact marked Choi distance to a phase-blind comparison]
\label{sel16:thm:cost}
Suppose \eqref{sel16:eq:two-block} is a complete normalized instrument on an
input of dimension $n$, with its labels fixed.  Among all normalized
$G_0$-covariant comparison instruments on these same input, output, and label
spaces,
\begin{equation}
 \boxed{\inf_{\widehat{\mathcal J}}
       \sum_y\|J_y-\widehat J_y\|_1
       =2\sum_y\|C_y\|_1.}
 \label{sel16:eq:cost}
\end{equation}
The infimum is attained by group averaging each Choi operator, which removes
its cross blocks.  For the flagged channels one also has
\begin{equation}
 \inf_{\widehat{\mathcal J}}
 \|\mathcal J_{\rm flag}-\widehat{\mathcal J}_{\rm flag}\|_\diamond
 \ge\frac{2}{n}\sum_y\|C_y\|_1.
 \label{sel16:eq:diamond-cost}
\end{equation}
The trace-norm optimum is not asserted to be the diamond-norm optimum in
all dimensions.
\end{theorem}
\begin{proof}
Let $C_y=U_y|C_y|$ be its polar decomposition.  The Hermitian off-diagonal
operator with blocks $U_y,U_y^*$ has operator norm at most one.  Its trace
pairing with $J_y-\widehat J_y$ is $2\Tr|C_y|$, since a covariant comparison
has zero cross block between distinct characters.  Trace-norm duality gives
the lower bound, even if that comparison has other character sectors.
The off-diagonal Hermitian difference after pinching has trace norm
$2\|C_y\|_1$, giving equality.  Averaging physical unitary conjugations of
a normalized instrument preserves complete positivity and total trace
preservation, so it is an admissible comparison.  Applying a flagged channel
difference to a normalized maximally entangled state produces its flagged
Choi operator divided by $n$, proving \eqref{sel16:eq:diamond-cost}.
\end{proof}

\begin{proposition}[One-sided certificates retain the physical scale]
\label{sel16:prop:errors}
If the reconstructed marked cross blocks have certified errors
$\sum_y\|\widehat C_y-C_y\|_{\HS}^2\le\epsilon^2$, then
\[
 \left|\sqrt{h_{\det}}-
     \Bigl(\sum_y\|\widehat C_y\|_{\HS}^2\Bigr)^{1/2}\right|
 \le\epsilon.
\]
Consequently a measured cross-block norm exceeding $\epsilon$ certifies the
primitive determinant constraint on the exact two-character class.  If instead
\eqref{sel16:eq:purity} is estimated with absolute squared-trace errors
$e_{J,y},e_{A,y},e_{B,y}$, then
\[
 |\widehat h-h_{\det}|
 \le\frac12\sum_y(e_{J,y}+e_{A,y}+e_{B,y}).
\]
A positive lower endpoint certifies $h_{\det}>0$.  A small interval containing
zero does not certify an exactly phase-blind source.
\end{proposition}
\begin{proof}
The first bound is the reverse triangle inequality in the direct-sum
Hilbert--Schmidt space.  The second is the triangle inequality in the exact
squared-trace identity.  No rank rounding or unverified sign of a noisy
residual is used.
\end{proof}

A physical factor $s_y$ in an unnormalized event multiplies its Choi operator
by $s_y$, its cross block by $s_y$, and its contribution to $h_{\det}$ by
$s_y^2$.  Thus an $l$-letter accepted event at specified locked opportunities
retains $s_y=\vartheta^l$.  The Choi-distance and diamond lower bounds have
one power of that event weight, not two.  Comparisons of normalized successful
states must restore the probability factors before these source claims are
made.

\section{A normalized phase record with the same forgotten reset}
\label{sel16:reset}

This comparison is not a new historical source assignment.  It shows that the
preceding marked/coarse distinction is compatible with complete normalization.
On the two actual character lines $\mathbb C|0\rangle\oplus
\mathbb C|1\rangle$ take the group action $\operatorname{diag}(1,\chi(g))$.
For $0\le v\le1$ let
\begin{equation}
 F_\pm(v)=\frac12(I\pm vX),\qquad
 \mathcal J_{\pm,v}(\rho)=\Tr(F_\pm(v)\rho)|0\rangle\langle0|.
 \label{sel16:eq:reset}
\end{equation}
The output line is neutral.  The two amplitude-coordinate lines are spanned
by $|0\rangle\langle0|$ and $|0\rangle\langle1|$, of characters $1$ and
$\chi^{-1}$.

\begin{theorem}[Exact normalized marked/coarse alternative]
\label{sel16:thm:reset}
The maps in \eqref{sel16:eq:reset} form a normalized instrument.  Its forgotten
channel is the \emph{same} reset $\rho\mapsto\Tr(\rho)|0\rangle\langle0|$
for every $v$.  In the amplitude basis its marked Choi matrices are
\[
 J_\pm(v)=\frac12\begin{pmatrix}1&\pm v\\\pm v&1\end{pmatrix},
 \qquad h_{\det}=\frac{v^2}{2}.
\]
For fixed outcome labels, its covariance group is $G_0$ at $v=0$ and
$\ker\chi$ at every $v>0$.  Moreover
\begin{equation}
 \boxed{\|\mathcal J_{{\rm flag},v}
            -\mathcal J_{{\rm flag},0}\|_\diamond=v.}
 \label{sel16:eq:reset-distance}
\end{equation}
This attains the general lower bound \eqref{sel16:eq:diamond-cost}.
If relabelling the two outcomes is permitted, its symmetry group at $v>0$
instead contains exactly the extra alternative $\chi=-1$ and becomes
$\ker\chi^2$.  That is a different record policy.
\end{theorem}
\begin{proof}
The effects have eigenvalues $(1\pm v)/2$ and sum to $I$.  Their transpose
is their Choi matrix on the displayed two amplitude directions, giving the
normalization and cross-mass assertions.  The cross blocks have opposite signs
and cancel when labels are forgotten.  With labels fixed, their transformation
requires $\chi=1$.  Allowing their exchange permits $\chi=-1$ as well and
no other phase for nonzero real off-diagonal entries.

The flagged difference is the functional $\rho\mapsto\Tr(X\rho)$ times
$v/2$ in one output block and its negative in the other.  A Hermitian
functional represented by $X$ has completely bounded trace norm
$\|X\|_\op=1$: this follows by dualizing to the map $z\mapsto zX$.
The two flags sum their trace norms.  This proves the upper bound $v$, and an
$X$ eigenstate attains it.  Here $n=2$ and
$\sum_\pm\|C_\pm\|_1=v$, giving equality with
\eqref{sel16:eq:diamond-cost}.
\end{proof}

At $v=1$ the two actual amplitudes are
$k_\pm=|0\rangle(\langle0|\pm\langle1|)/\sqrt2$.
Each includes a neutral and a determinant-character contribution coherently,
while the complete reset is phase blind.  This does not evade V6: it refines a
rank-two reset Choi operator, not a previously resolved rank-one homogeneous
outcome.  On an unchanged instrument the admissibility of this refined Read
must be proved, not assumed from the CP domination relation.

\section{The inherited determinant incidence in a full coherent cell}
\label{sel16:dirac}

The previous criterion can also be evaluated on the precise tensor already
used in the earlier determinant handoff, rather than on an arbitrary router.
For a unit determinant functional the inherited map obeys
\[
 \Theta\Theta^*=2I_3,\qquad\|\Theta\|_{\HS}^2=6.
\]
Put $V=\Theta/\sqrt2$, so $VV^*=I_3$, and write
$P=V^*V$ on $S=C\otimes W\otimes W$ of dimension twelve.
Its rank is three.  The target $T=\Lambda^2C^*$ has dimension three.
On $S\oplus T$ define
\begin{equation}
 D=s\begin{pmatrix}0&V^*\\V&0\end{pmatrix},\quad s>0,
 \qquad P_D=P\oplus I_T,\quad Q=I-P_D,
 \label{sel16:eq:Dirac}
\end{equation}
so $\rank Q=9$ and $D^2=s^2P_D$.
This is the Hermitian completion of that \emph{existing typed incidence}.
The unitary dynamics below is an additional source specification only when
it is asserted to be physically executed; its existence is not inferred from
formal functional calculus.

\begin{theorem}[Isolated transition, odd CP square, and full-cell covariance]
\label{sel16:thm:Dirac}
For the fixed input/output representation inherited from $C,W$:
\begin{enumerate}[label=\textup{(\roman*)},leftmargin=3em]
\item The map $X\mapsto VXV^*$ is $G_0$-covariant.
\item The CP map $X\mapsto DXD$ on $S\oplus T$ has covariance group
      $\ker\chi^2$.
\item Let $U_t=e^{-itD}$ and $\mathcal U_t(X)=U_tXU_t^*$.  Then
\begin{equation}
 \boxed{\operatorname{Cov}(\mathcal U_t)=
 \begin{cases}
 G_0,&\sin(st)=0,\\
 \ker\chi,&\sin(st)\ne0.
 \end{cases}}
 \label{sel16:eq:unitary-group}
\end{equation}
      Hence every sufficiently small positive time has exactly the primitive
      determinant constraint, if this full coherent evolution occurs.
\end{enumerate}
The three mutually orthogonal character components of $|U_t\rangle\!\rangle$
have weights $0,1,-1$ and squared norms
\begin{equation}
 \mu_0=9+6\cos^2(st),\qquad
 \mu_+=\mu_-=3\sin^2(st).
 \label{sel16:eq:unitary-masses}
\end{equation}
Each neutral--charged Choi cross block therefore has squared HS norm
$(9+6\cos^2(st))3\sin^2(st)$.
\end{theorem}
\begin{proof}
Part (i) is Proposition~\ref{sel16:prop:ray}.  The two off-diagonal terms of
$D$ are nonzero and transform by opposite characters $\chi$ and $\chi^{-1}$.
A single-Kraus CP map is unchanged precisely when its coefficient changes by
a common scalar phase.  For $D$ this requires $\chi=\chi^{-1}$, or
$\chi^2=1$, proving (ii).

The identity $D^2=s^2P_D$ gives the exact exponential
\begin{equation}
 U_t=Q+\cos(st)P_D-i\sin(st)D/s.
 \label{sel16:eq:unitary}
\end{equation}
Its neutral part $Q+\cos(st)P_D$ never vanishes, since $\rank Q=9$.
When $\sin(st)\ne0$, both charged components are present.  Equality of the
unitary channels requires the transformed $U_t$ to equal a scalar multiple
of $U_t$.  The nonzero neutral component fixes that scalar to one; either
charged component then requires $\chi=1$.  When $\sin(st)=0$, the whole
unitary is neutral.  The ranks of $Q,P_D$ and $V$ give
\eqref{sel16:eq:unitary-masses}, and the single Choi ray gives the cross norms.
\end{proof}

\begin{corollary}[An exact full-cell experiment distinguishes the missing sign]
\label{sel16:cor:Dirac-witness}
At $st=\pi/2$, choose unit vectors $a\in\Ran P$, $b=Va$, and
$k\in\Ran Q$.  The input $(k+a)/\sqrt2$ has output
$(k-ib)/\sqrt2$.  A transform with $\chi=-1$ instead gives
$(k+ib)/\sqrt2$.  These outputs are orthogonal.  Thus the two full unitary
channels have diamond distance two although their isolated transition CP maps
and their odd CP squares agree under that same sign transformation.
\end{corollary}
\begin{proof}
Use \eqref{sel16:eq:unitary}, $Vk=0$ on the kernel, and orthogonality of
$k,b$.  The output inner product is zero.  A pair of orthogonal pure outputs
has trace distance two, which is also the universal upper bound for the
diamond distance between channels.  The invariances of the isolated maps are
parts (i)--(ii) of Theorem~\ref{sel16:thm:Dirac}.
\end{proof}

This experiment requires the stated coherent input and full output access.
It is an exact conditional calculation of the inherited tensor, not a claim
that those preparations occur in the locked writer.  It shows why keeping
the unsplit full-cell source can recover a determinant condition that an
isolated second-order transition square loses.

\section{The retained fourteen-dimensional memory has a tensor-order boundary}
\label{sel16:arity}

There is a second reason not to search indefinitely among single-system word
products.  Use V15's \emph{already evaluated} represented type, with no change
to its source coefficients.  For the canonical block-unitary embedding of the
two nonabelian factors, and with the other scalar blocks held neutral, the
candidate action on the retained memory is
\begin{equation}
 R_{\rm nat}(U,V)=U\oplus(I_3\otimes V)\oplus I_2\oplus I_3.
 \label{sel16:eq:native-group}
\end{equation}
The neutral five-dimensional sector in this formula may retain its two
original scalar labels; joining them in a charge count does not erase either
label from the source.

\begin{theorem}[All-depth single-system exclusion and sharp five-copy floor]
\label{sel16:thm:arity}
For \eqref{sel16:eq:native-group}, the central-weight support of
$\End(H_{14})$ is
\begin{equation}
 \boxed{(0,0),\quad\pm(1,0),\quad\pm(0,1),\quad\pm(1,-1).}
 \label{sel16:eq:weight-support}
\end{equation}
Here $(e^{i\alpha}I_3,e^{i\beta}I_2)$ acts on a weight $(a,b)$ by
$e^{i(a\alpha+b\beta)}$.  None is the determinant-incidence weight
$(-3,-2)$.  Consequently no single-system operator, and hence no temporal
word of any length on that same memory, can be a nonzero determinant-character
amplitude under this specified action.

On the parallel carrier $H_{14}^{\otimes k}$ with diagonal group action,
that character is absent for $k<5$ and is present for $k=5$.
The bound is sharp as an algebraic occurrence statement.  It does not assert
physical admission of a five-copy preparation, contraction, or outcome.
\end{theorem}
\begin{proof}
The physical system has central weights $(1,0),(0,1),(0,0)$ on subspaces of
dimensions $3,6,5$.  An endomorphism has target-minus-source weight, giving
exactly \eqref{sel16:eq:weight-support}.  Their dimensions are respectively
\[
 70;\quad15,15;\quad30,30;\quad18,18,
 \qquad70+30+60+36=196.
\]
This is a statement about the entire ambient $M_{14}$, so increasing word
depth inside that algebra cannot evade it.

In $k$ copies a Hilbert-vector weight is $(a,b)$ with $a,b\ge0$ and
$a+b\le k$.  A target-minus-source weight therefore has sum at least $-k$.
The required $(-3,-2)$ has sum $-5$, excluding $k<5$.
At $k=5$, take a fixed multiplicity copy of $W$ and the subspace
$C^{\otimes3}\otimes W^{\otimes2}$.  Its normalized alternating vector
\[
 \Omega=\Omega_C\otimes\Omega_W,\qquad
 \Omega_C=\frac1{\sqrt6}\sum_{\sigma\in S_3}
          \operatorname{sgn}(\sigma)e_{\sigma(1)}\otimes
                      e_{\sigma(2)}\otimes e_{\sigma(3)},\qquad
 \Omega_W=\frac{f_1\otimes f_2-f_2\otimes f_1}{\sqrt2}
\]
transforms by $\det U\det V$.  Let $r$ be a unit vector in the original
neutral sector.  Then
$B=|r\rangle^{\otimes5}\langle\Omega|$ is a nonzero rank-one operator
on the five-copy carrier and transforms by $\chi^{-1}$.
This proves presence and sharpness.  Its isolated Choi ray remains phase blind
by Proposition~\ref{sel16:prop:ray}; tensor availability alone is not the
neutral-reference condition.
\end{proof}

\begin{remark}[What has been held fixed]
The floor refers to the defining block-unitary embedding
\eqref{sel16:eq:native-group}.  Assigning determinant powers to a previously
neutral scalar sector, adjoining charged dual carriers, or changing the
central action changes this premise.  Such a change can lower a parallel-copy
requirement and must be separately sourced; it is not determined by the
abstract algebra $M_3\oplus M_2\oplus\C^2$.  The inherited up-incidence
$\Hom(C\otimes W^{\otimes2},\Lambda^2C^*)$ already has a different typed
source and target, with five elementary charged legs in total.  There is no
contradiction: it was never an endomorphism of the fundamental-only memory
specified in \eqref{sel16:eq:native-group}.  Five tensor copies here means
neither five temporal steps, five spatial dimensions, nor five generations.
\end{remark}

\section{The finite determinant handoff and its remaining source row}
\label{sel16:status}

\begin{corollary}[Operational completion of the determinant-group step]
\label{sel16:cor:handoff}
Suppose the existing algebra selector has established the physical internal
factors $M_3,M_2$, and actual typed input/output data identify a complete
incidence subpacket with verified character spaces $1,\chi^{-1}$.  If its
marked Choi cross mass satisfies $h_{\det}>0$, its exact stabilizer inside
$\mathrm U(3)\times\mathrm U(2)$ is
\[
 S(\mathrm U(3)\times\mathrm U(2))
 \cong(\mathrm{SU}(3)\times\mathrm{SU}(2)\times\mathrm U(1))/\mathbb Z_6.
\]
An actually coherent realization of the inherited common neutral--determinant
ray supplies this cross mass automatically.  If other character blocks occur,
the complete integer difference lattice replaces the two-block test.
\end{corollary}
\begin{proof}
Apply Theorem~\ref{sel16:thm:anchor} or
Theorem~\ref{sel16:thm:lattice}.  The quotient of the connected kernel is the
already established determinant stabilizer.  Corollary~\ref{sel16:cor:shared}
gives the automatic cross mass for one common physical ray.
\end{proof}

The new step is not another proof that $M_3$ can be generated.  It establishes
what a \emph{recorded positive source} must retain before the determinant
stabilizer follows.  Character-ray mass and phase-sensitive coupling are
different source quantities.  Their difference survives arbitrary additional
single-copy record data, because a global Kraus phase does not change the CP
map at all.  A native neutral Clifford component or the identity component of
an actual full coherent cell can provide the needed reference; an externally
invented phase control is not logically necessary on those branches.

The fixed-memory charge count supplies a stopping rule.  On the defining
internal action of the unchanged V15 witness, no longer temporal word can
create a missing determinant-character operator.  The next object must be an
actually identified tensor-incidence port or a separately verified central
character on a new source sector.  After that port is identified, evaluate its
original marked Choi cross blocks, rather than adding a determinant line or
calling a classical pair of shadows one coherent source.

None of these calculations populates that historical tensor port.  V15's
physical product-comparison requirement is retained, and the determinant
character test does not determine named type/private rays, anomaly-free matter
types, their joint grading, or continuum measure phases.  A successful
determinant test is one step of the structural programme, not an unconditional
microscopic Standard-Model derivation.  The general source-complete joint
packet continues to require all declared mixed maps on one actual carrier, as
in \cite{Duality}, \src{def:joint-packet}.

\paragraph{Verification and preservation.}
The complete V15 prefix is preserved.  The companion exact-arithmetic script
checks the marked charge lattices, primitive and disconnected examples,
positivity and purity identities, the sharp normalized reset distance,
phase-residual calibration, the inherited $12$-to-$3$ alternating incidence,
its full $15$-dimensional unitary, and the single-copy/five-copy weight boundary.
It reuses the V15 certificate and reruns the V15 checks rather than changing
the full-rank source.  General all-word and all-parameter statements are
proved above, not inferred from a numerical sample.  The scripts supply
reproducible algebraic checks, not proof-assistant certification or independent
peer review.

\begin{thebibliography}{99}
\raggedright
\bibitem[25]{sel16-MS}
I.~Marvian and R.~W.~Spekkens,
\emph{Modes of asymmetry: The application of harmonic analysis to symmetric
quantum dynamics and quantum reference frames},
Phys.\ Rev.\ A \textbf{90} (2014), 062110;
\href{https://doi.org/10.1103/PhysRevA.90.062110}{doi:10.1103/PhysRevA.90.062110}.
\bibitem[26]{sel16-GS}
G.~Gour and R.~W.~Spekkens,
\emph{The resource theory of quantum reference frames: manipulations and
monotones}, New J.\ Phys.\ \textbf{10} (2008), 033023;
\href{https://arxiv.org/abs/0711.0043}{arXiv:0711.0043}.
\bibitem[27]{sel16-AFCB}
M.~Ara\'{u}jo, A.~Feix, F.~Costa and \v C.~Brukner,
\emph{Quantum circuits cannot control unknown operations},
New J.\ Phys.\ \textbf{16} (2014), 093026;
\href{https://arxiv.org/abs/1309.7976}{arXiv:1309.7976}.
These references establish the surrounding symmetry-mode and coherent-control
frameworks; no physical reference resource is imported from them into the
historical source.
\end{thebibliography}
\addtocontents{toc}{\protect\endgroup}
% END OF SOURCE-SELECTION V16 DEVELOPMENT BLOCK.
% Preserve V1--V16 and append subsequent work before the final terminator.


\clearpage
\hypersetup{pdftitle={Historical Standard-Model Source Selection: V1--V17}}
\begin{center}
{\Large\bfseries Source-selection continuation V17}\\[.4em]
{\large The first operational determinant mode: sharp tensor-slot bounds,\\
primitive global selection, and lower-order invisibility}
\end{center}
\addtocontents{toc}{\protect\begingroup}
\addtocontents{toc}{\protect\renewcommand{\protect\numberline}{\protect\makebox[2.1em][l]}}

\section{The next source question is the order of the CP operation}
\label{sel17:context}

V16 proves that a determinant-character \emph{amplitude} cannot occur in the
endomorphisms of fewer than five parallel copies of the specified native
memory.  It separately proves that a CP symmetry is fixed by differences of
amplitude characters.  These two results leave a smaller upstream question:
what is the minimum tensor-port order of a \emph{CP operation} whose covariance
group is the primitive determinant subgroup?  A neutral amplitude and an
amplitude of the full determinant character are sufficient, but need not be
necessary.  Two different nonneutral amplitudes can have the required relative
character.

The distinction is substantive.  We prove below that fewer than five total
native input/output slots cannot distinguish the determinant subgroup from
the full product group.  Five slots suffice, with two inputs and three outputs;
on equal input and output carriers three parallel copies suffice.  The
five-copy \emph{amplitude} theorem is retained unchanged.  At the first
allowed operational orders, a nonzero symmetry signal is automatically
primitive: the higher character harmonics responsible for a disconnected
alternative do not fit in the available slots.

\paragraph{Audit and dependence.}
The inherited results are the character-lattice theorem
\ref{sel16:thm:lattice}, the neutral-anchor criterion
\ref{sel16:thm:anchor}, and the amplitude floor
\ref{sel16:thm:arity}.  The older determinant/up-incidence isomorphism and
shared Clifford--determinant source ray are \cite{Attack},
\src{thm:V110-det-incidence-iso} and \src{thm:V110-one-ray}.  We do not
reprove their determinant-line construction, claim that its physical port
already occurs, or add another colour router.  The additions are a
whole-Choi bandwidth theorem that does not require character-scalar Kraus
spaces, its sharp normalized realization, an exact lower-order
indistinguishability theorem, and a direct audit of the unchanged V13 source.
Symmetry-mode decomposition and quantum-network composition are established
methods \cite{sel16-MS,sel17-CDP}; all bounds and specialized claims used
here are proved explicitly.

\paragraph{Physical scope.}
A tensor slot below means a factor carrying the \emph{defining} native action,
not a time label or a freely available copy operation.  Each proposed port,
its representation, its input preparation, and its output readout must be
identified in the actual typed grammar.  Neutral registers may carry their
original provenance labels.  No particle-statistics assignment to these
internal factors is made: a physical fermionic grading, when required, is an
additional compatibility test, not silently replaced by tensor degree.

\section{A sharp bandwidth bound for the complete Choi operator}
\label{sel17:bandwidth}

Keep the actual candidate action of V16,
\begin{equation}
 R_0(U,V)=U\oplus(I_3\otimes V)\oplus I_2\oplus I_3,
 \qquad G_0=\mathrm U(3)\times\mathrm U(2).
 \label{sel17:eq:native}
\end{equation}
Put $G_{\det}=\ker\chi$, where $\chi(U,V)=\det U\det V$.
An input port $H_a$ and an output port $H_b$ are invariant subspaces of
$H_{14}^{\otimes a}$ and $H_{14}^{\otimes b}$, respectively.  Arbitrary
finite neutral multiplicity factors are allowed.  The nonnegative integers
$a,b$ count native factors and $s=a+b$ is their total slot budget.  A scalar
port has zero slots.  Retain a marked CP instrument with Choi operators
$J_y$ on $H_b\otimes\overline{H_a}$ and representation
$\mathcal R=R_b\otimes\overline{R_a}$.  For this block suppose
\begin{equation}
 \mathcal R(h)J_y\mathcal R(h)^*=J_y
       \quad(h\in G_{\det},\ \hbox{every original }y).
 \label{sel17:eq:Hcov}
\end{equation}
This is a condition on the original maps, not permission to twirl them into
satisfying it.  A finite test for it is given in
Section~\ref{sel17:finite-test}.

\begin{theorem}[Determinant bandwidth of a native tensor port]
\label{sel17:thm:bandwidth}
Every Choi operator satisfying \eqref{sel17:eq:Hcov} has a unique orthogonal
mode expansion
\begin{equation}
 J_y=\sum_{q=-m}^{m}J_y^{[q]},\qquad
 m=\left\lfloor\frac{s}{5}\right\rfloor,
 \qquad
 \mathcal R(g)J_y^{[q]}\mathcal R(g)^*
       =\chi(g)^q J_y^{[q]}.
 \label{sel17:eq:modes}
\end{equation}
Here $(J_y^{[q]})^*=J_y^{[-q]}$.  The modes are not individually asserted to
be CP maps; $J_y^{[0]}$ is the group average and is positive.  If there is a
nonzero nonconstant mode and $d$ is the positive greatest common divisor of
its occurring indices, then the fixed-label covariance group inside $G_0$ is
\begin{equation}
 \boxed{G_{\rm inc}=\ker\chi^d,\qquad 1\le d\le m.}
 \label{sel17:eq:group}
\end{equation}
If there is no such mode, $G_{\rm inc}=G_0$.
\end{theorem}
\begin{proof}
The finite-dimensional subspace of Choi operators fixed by $G_{\det}$ is
invariant under $G_0$, since $G_{\det}$ is normal.  Its representation factors
through the compact circle $G_0/G_{\det}$, which is identified with
$\mathbb T$ by $\chi$.  A finite-dimensional unitary circle representation
is an orthogonal sum of its integer-character spaces.  One elementary proof
is to simultaneously diagonalize its commuting unitary matrices; a continuous
scalar character of the circle has the form $e^{it}\mapsto e^{iqt}$ with
$q\in\mathbb Z$.  This proves the mode expansion before the bound on $q$.

For the central subgroup $g_t=(e^{it}I_3,e^{it}I_2)$, each native Hilbert
factor has only charges zero and one.  Thus $R_a(g_t)$ has charges between
zero and $a$, and $R_b(g_t)$ has charges between zero and $b$.  The amplitude
representation $\mathcal R(g_t)$ has charges between $-a$ and $b$.
Conjugation on its operator space has charge differences between $-s$ and
$s$.  But the mode $\chi^q$ has charge $5q$ under $g_t$.  Therefore
$5|q|\le s$.  Invariant subspaces and neutral multiplicities cannot enlarge
this range.  Adjoint and orthogonality give the other expansion properties.

A mode is unchanged by $g$ exactly when $\chi(g)^q=1$ or its coefficient
is zero.  Orthogonality prevents cancellation between different modes.  The
simultaneous nonzero-mode constraints are $\chi^d=1$ by Bezout's identity.
The gcd is no larger than the largest occurring positive absolute index,
which proves \eqref{sel17:eq:group}.  The zero mode is the Haar average of
$J_y$; averaging preserves positivity and the normalization of a complete
instrument.
\end{proof}

The theorem concerns the complete Choi operator, including nonscalar
$\mathrm{SU}(3)\times\mathrm{SU}(2)$ Kraus modules.  It does not first project
onto a hand-picked scalar intertwiner space.  In particular it applies where
V16's two-character amplitude hypothesis has not been supplied.

\begin{corollary}[The first allowed operational order is necessarily primitive]
\label{sel17:cor:primitive}
Under \eqref{sel17:eq:Hcov}:
\begin{enumerate}[label=\textup{(\roman*)},leftmargin=2.8em]
\item If $s<5$, every original CP branch is $G_0$-covariant.
\item If $5\le s\le9$, the fixed-label covariance group is either $G_0$ or
      exactly $G_{\det}$.  It is $G_{\det}$ precisely when a nonconstant mode
      occurs.
\item A group $\ker\chi^d$ with $d\ge2$ requires $s\ge5d$.  Thus an extra
      determinant-square component requires at least ten slots in this native
      category.
\end{enumerate}
\end{corollary}
\begin{proof}
For $s<5$, $m=0$.  For $5\le s\le9$, $m=1$, and any nonempty set of
nonzero indices has gcd one.  Finally $d\le\lfloor s/5\rfloor$ implies
$s\ge5d$.
\end{proof}

These assertions compare groups \emph{inside} the fixed $G_0$ and retain the
original marks.  They do not exclude an outer charge-conjugation symmetry or
a different discrete symmetry that is permitted to permute those marks.
The connectedness and the familiar sixfold covering kernel of $G_{\det}$
remain the inherited group facts, not additional tensor-count conclusions.

\section{A normalized two-input, three-output source attains the bound}
\label{sel17:sharp}

The lower bound has a simple exact realization on invariant subspaces of the
\emph{same native representation}; it does not need charged duals or a
previously assigned determinant power on a scalar sector.  This is a declared
comparison instrument, not an unevaluated historical source being filled in.
Choose a unit vector $r$ in one original neutral summand and choose one of the
three actual defining weak multiplicity copies.  Define
\begin{align}
 n_i&=r\otimes r, &
 w_i&=(f_1\otimes f_2-f_2\otimes f_1)/\sqrt2
                   &&\text{in }H_{14}^{\otimes2},\nonumber\\
 n_o&=r^{\otimes3}, &
 c_o&=\frac1{\sqrt6}\sum_{\sigma\in S_3}\operatorname{sgn}(\sigma)
       e_{\sigma(1)}\otimes e_{\sigma(2)}\otimes e_{\sigma(3)}
                   &&\text{in }H_{14}^{\otimes3}.
 \label{sel17:eq:vectors}
\end{align}
The two $n$ vectors are neutral; $w_i$ transforms by $\det V$ and $c_o$
by $\det U$.  The rays and their orthogonal projections are canonical
alternating subspaces once the physical factors are fixed.  Choosing unit
vectors in those rays chooses comparison coordinates, not a new claim of
historical accessibility.
Put
\begin{equation}
 B_0=|c_o\rangle\langle n_i|,\qquad
 B_1=|n_o\rangle\langle w_i|,\qquad
 P=|n_i\rangle\langle n_i|+|w_i\rangle\langle w_i|.
 \label{sel17:eq:arms}
\end{equation}
For $0<\lambda\le1$ and $|\eta|\le1$, define the two marked maps
\begin{align}
 \mathcal J_{{\rm s},\eta}(X)
 &=\lambda\bigl(B_0XB_0^*+B_1XB_1^*
              +\eta B_0XB_1^*+\overline\eta B_1XB_0^*\bigr),\nonumber\\
 \mathcal J_{\rm f}(X)
 &=\Tr[(I-\lambda P)X]|n_o\rangle\langle n_o|,
 \qquad \Lambda_\eta=\mathcal J_{{\rm s},\eta}+\mathcal J_{\rm f}.
 \label{sel17:eq:instrument}
\end{align}
The letters s and f are retained labels.  Both maps act on the entire
specified input, not only its two-dimensional success support.

\begin{theorem}[Sharp operational realization without a neutral amplitude]
\label{sel17:thm:sharp}
The maps \eqref{sel17:eq:instrument} form a normalized instrument.  Its
success probability on $P$ is the state-independent $\lambda$.  The covariance
groups of both the marked instrument and its forgotten channel are
\begin{equation}
 \boxed{\operatorname{Cov}(\Lambda_\eta)
   =\begin{cases}G_0,&\eta=0,\\G_{\det},&\eta\ne0.\end{cases}}
 \label{sel17:eq:sharp-group}
\end{equation}
For $\eta=1$, the success branch has one amplitude
$\sqrt\lambda(B_0+B_1)$, but neither arm is neutral or has character
$\chi^{\pm1}$.  They have characters $(\det U)$ and $(\det V)^{-1}$,
respectively.  Their relative character is $\chi$.

For the complete family,
\begin{equation}
 \boxed{\|\Lambda_\eta-\Lambda_\zeta\|_\diamond
       =\lambda|\eta-\zeta|.}
 \label{sel17:eq:sharp-distance}
\end{equation}
The same equality holds with the two labels retained.  In particular the
exact distance from $\Lambda_\eta$ to the set of normalized $G_0$-covariant
channels on the same ports is $\lambda|\eta|$.
\end{theorem}
\begin{proof}
The Choi operator of the success branch is the image, in the orthonormal
amplitude basis $(B_0,B_1)$, of
$\lambda\left(\begin{smallmatrix}1&\eta\\\overline\eta&1\end{smallmatrix}\right)$.
This is positive exactly when $|\eta|\le1$.  The cross products
$B_0^*B_1=B_1^*B_0=0$ show that its input effect is $\lambda P$.
The other effect is $I-\lambda P\succeq0$, giving normalization and the
success statement.

Both individual arm maps are $G_0$-covariant.  Their Choi cross block transforms
by $\det U\det V$, so its nonzero coefficient is fixed exactly on
$G_{\det}$.  The failure effect is invariant and its output is neutral, so
the failure map is $G_0$-covariant.  It cannot cancel a nonzero nonconstant
mode in the sum.  This proves \eqref{sel17:eq:sharp-group} with or without
labels.

On the success support, identify $(n_i,w_i)$ with a logical input qubit and
$(c_o,n_o)$ with a logical output qubit.  The channel difference multiplies
the two off-diagonal coordinates by $\lambda(\eta-\zeta)$ and its conjugate,
and kills diagonal coordinates.  The real unit off-diagonal map is
$(\id-\operatorname{Ad}_Z)/2$, of diamond norm at most one; a unitary phase
change accounts for a complex coefficient.  Compression to the input support
and isometric output embedding are complete trace-norm contractions.  This
gives the upper bound in \eqref{sel17:eq:sharp-distance}, also at all ancillary
amplifications.  Input $(n_i+w_i)/\sqrt2$ has a difference of output density
matrices with eigenvalues $\pm\lambda|\eta-\zeta|/2$, proving equality.

For any $G_0$-covariant comparison, its output matrix element between $c_o$
and $n_o$ must vanish on every input supported in $P$.  Indeed that output
functional has character $\det U$, whereas the input operator space on $P$
has only characters $1,(\det V)^{\pm1}$.  Equivariance forbids such an
intertwiner.  Pairing the output of the above balanced input with a norm-one
off-diagonal Hermitian observable of the appropriate phase therefore gives
trace-distance lower bound $\lambda|\eta|$.  The group average
$\Lambda_0$ attains it.  The failure branches agree between all displayed
channels, proving the flagged equality too.
\end{proof}

\begin{corollary}[Three different exact tensor costs]
\label{sel17:cor:costs}
On the native fundamental-only representation, the following minima are sharp:
\begin{center}
\small
\begin{tabular}{L{.64\textwidth}L{.27\textwidth}}
\toprule
Object whose occurrence is required & Minimum native order\\
\midrule
A determinant-character amplitude in $\End(H_{14}^{\otimes k})$
 & $k=5$ (inherited V16).\\[.25em]
A CP instrument on $H_{14}^{\otimes k}$ with fixed-label covariance group
exactly $G_{\det}$ & $k=3$.\\[.25em]
A CP instrument from $a$ native inputs to $b$ native outputs with that group
 & $a+b=5$, attained at $(a,b)=(2,3)$.\\
\bottomrule
\end{tabular}
\end{center}
A state preparation alone requires five output factors; a measurement alone
requires five input factors.  Neutral spectator factors do not change these
lower bounds.
\end{corollary}
\begin{proof}
The amplitude statement is Theorem~\ref{sel16:thm:arity}.  The total-slot
lower bound is Corollary~\ref{sel17:cor:primitive};
Theorem~\ref{sel17:thm:sharp} attains it.  For an equal-port realization,
replace $n_i$ by $r^{\otimes3}$ and $w_i$ by $w_i\otimes r$ in the same
construction.  It gives a normalized channel on three copies.  Two copies
have only four slots and are excluded.  For state preparation use
$(r^{\otimes5}+\Omega_C\otimes\Omega_W)/\sqrt2$; its density has relative
character $\chi$ and group $G_{\det}$.  The binary projective measurement
onto that ray and its complement gives the measurement endpoint.  Their
existence is a comparison construction, not an admission of copy or
measurement operations in the old source.
\end{proof}

\begin{proposition}[The disconnected-component bound is also sharp]
\label{sel17:prop:higher}
For every integer $d\ge1$, a normalized instrument with covariance group
$\ker\chi^d$ exists with $a=2d$, $b=3d$.  Thus the minimum total-slot cost
of exactly $d$ determinant components is $5d$.
\end{proposition}
\begin{proof}
In \eqref{sel17:eq:vectors} replace $w_i,c_o$ by
$\Omega_W^{\otimes d},\Omega_C^{\otimes d}$ and use the corresponding
neutral tensor powers.  The two arms have characters $(\det U)^d$ and
$(\det V)^{-d}$.  The proof of Theorem~\ref{sel17:thm:sharp} applies
unchanged and gives $\ker\chi^d$ for any nonzero visibility.  The lower bound
is Corollary~\ref{sel17:cor:primitive}(iii).
\end{proof}

\section{Lower-order process data can be identical at the two group choices}
\label{sel17:marginals}

The sharp realization is not an instruction to replace the original source.
It instead gives a complete positivity-compatible separating pair for the
\emph{missing order of data}.  In this section a marginal comparison means
preparing removed inputs in invariant states, discarding removed outputs, or
more generally applying fixed $G_0$-covariant input and output channels.  No
asymmetric reference is introduced in the discarded legs.

\begin{theorem}[Every subcritical native marginal is phase blind]
\label{sel17:thm:marginals}
Let $\Lambda$ be a $G_{\det}$-covariant channel on a native tensor port and
let $\Lambda^{(0)}$ be its $G_0$ group average.  Suppose a physical marginal
construction has the form
\[
 \mathfrak F(\Lambda)=\mathcal D\Lambda\mathcal E,
\]
where $\mathcal E,\mathcal D$ are $G_0$-covariant CPTP maps and the resulting
channel has $c$ native input factors and $d$ native output factors.  If
$c+d<5$, then
\begin{equation}
 \boxed{\mathfrak F(\Lambda)=\mathfrak F(\Lambda^{(0)}).}
 \label{sel17:eq:marginal-equality}
\end{equation}
The assertion holds branch by branch for covariant CP contractions of a
marked source.  In particular every proper marginal of the five-slot
instrument \eqref{sel17:eq:instrument}, formed by such leg removal, is the
same for all $\eta$.  Its complete five-slot operation nevertheless has group
$G_{\det}$ at $\eta\ne0$ and $G_0$ at zero, and their diamond distance is
$\lambda|\eta|$.
\end{theorem}
\begin{proof}
Composition makes $\mathfrak F(\Lambda)$ $G_{\det}$-covariant.  Its slot
budget is below five, so Theorem~\ref{sel17:thm:bandwidth} makes it
$G_0$-covariant.  The covariance of $\mathcal E,\mathcal D$ allows them to
pass through the Haar average, giving
\[
 \mathfrak F(\Lambda^{(0)})
  =\int_{G_0} g\cdot\mathfrak F(\Lambda)\,dg
  =\mathfrak F(\Lambda).
\]
The argument uses linearity, not trace preservation of the individual branch,
so it also proves the marked CP statement.  The average of
\eqref{sel17:eq:instrument} is its $\eta=0$ member.  The group and distance
are Theorem~\ref{sel17:thm:sharp}.
\end{proof}

This is equality of the reduced channels on \emph{all} their retained quantum
inputs, not merely equality of scalar populations on one preparation.  In
particular the three-copy equal-port construction and its average have the
same one-copy and two-copy input/output marginal channels when the omitted
inputs are supplied invariantly.  Their joint three-copy operation differs.
No inference from those marginals can decide which joint group occurs.

For the explicit five-slot pair, the reason is visible without group theory.
Its only varying Choi term is a scalar multiple of
$|B_0\rangle\!\rangle\langle\!\langle B_1|$ and its adjoint.  Tracing any
output leg pairs a colour vector with $r$ and gives zero.  An invariant
contraction of either input leg pairs $r$ with a weak vector and also gives
zero.  Thus each missing physical leg removes the entire phase signal.

\begin{corollary}[No temporal workaround through covariant narrow-memory steps]
\label{sel17:cor:width}
Consider a finite adaptive protocol in which each executed CP branch has at
most two native input and two native output factors, possibly together with
arbitrary neutral quantum registers.  Suppose each branch is
$G_{\det}$-covariant, preparations and terminal contractions obey the same
port restriction, and classical decision records transform trivially.
Then every resulting branch is $G_0$-covariant.  It cannot supply an operational
determinant selector, irrespective of temporal depth.
\end{corollary}
\begin{proof}
Each primitive has fewer than five slots and is therefore $G_0$-covariant by
Theorem~\ref{sel17:thm:bandwidth}.  Composition, tensoring with neutral
registers, and classical sums preserve covariance, by the inherited
Proposition~\ref{sel16:prop:no-free-phase}.  Induction on the protocol tree
proves the assertion.
\end{proof}

This corollary concerns the actual memory-bearing operations, not just their
observed low-order marginals.  A correlated source with an additional charged
memory can have invariant single-step marginals without meeting its
hypotheses.  A charged dual module, a wider memory, or an asymmetric reference
must be counted in the physical port rather than hidden in an unobserved
controller.  The full-causal process and its lower-order marginals are not
interchangeable.  Nor is varying $\eta$ a conservative enrichment of the
\emph{complete} five-slot source: only the stated lower-order images stay
unchanged.

\section{A finite global selector from covariance and one phase comparison}
\label{sel17:finite-test}

The hypothesis \eqref{sel17:eq:Hcov} need not be restated as an unexplained
group constraint.  On the verified native ports it is a finite set of
operator equations.  Let $N_C,N_W$ be the native projections onto the colour
and weak summands.  Sum them over the input or output factors and write
\[
 \widehat N_C=N_C^{\rm out}\otimes I-I\otimes\overline{N_C^{\rm in}},
 \qquad
 \widehat N_W=N_W^{\rm out}\otimes I-I\otimes\overline{N_W^{\rm in}}.
\]
For eleven Hermitian Lie-algebra basis generators $T_j$ of
$\mathfrak{su}(3)\oplus\mathfrak{su}(2)$, use the analogous Choi generators
$\widehat T_j$.  Define the positive covariance residual
\begin{equation}
 r_{\det}=\sum_y\left[
   \sum_{j=1}^{11}\|[\widehat T_j,J_y]\|_{\HS}^2
   +\|[-2\widehat N_C+3\widehat N_W,J_y]\|_{\HS}^2\right].
 \label{sel17:eq:Lie-test}
\end{equation}
Each term uses the full original Choi map at the identified cut.

\begin{theorem}[Finite primitive determinant selector below the second harmonic]
\label{sel17:thm:finite}
On a port with $5\le a+b\le9$, $r_{\det}=0$ is equivalent to
\eqref{sel17:eq:Hcov}.  If it vanishes, put
\[
 g_\phi=(e^{i\phi/3}I_3,I_2),\quad
 J_y(\phi)=\mathcal R(g_\phi)J_y\mathcal R(g_\phi)^*,\quad
 \mathfrak h=\sum_y\|J_y^{[1]}\|_{\HS}^2.
\]
Then
\begin{equation}
 \boxed{\mathfrak h>0
  \iff G_{\rm inc}=G_{\det}
  \iff\sum_y\|J_y-J_y(\pi)\|_{\HS}^2>0.}
 \label{sel17:eq:finite-selector}
\end{equation}
The exact formulas are
\begin{align}
 J_y^{[1]}&=\frac13\sum_{j=0}^{2}e^{-2\pi i j/3}J_y(2\pi j/3),
                         \label{sel17:eq:Fourier}\\
 \mathfrak h&=\frac18\sum_y\|J_y-J_y(\pi)\|_{\HS}^2
  =\frac14\sum_y\bigl[\Tr J_y^2-\Tr(J_yJ_y(\pi))\bigr],
                         \label{sel17:eq:phase-purity}\\
 \sum_y\|[\widehat N_C/3,J_y]\|_{\HS}^2&=2\mathfrak h.
                         \label{sel17:eq:derivative}
\end{align}
Thus no independent primitive-lattice assumption, absolute determinant phase,
or scalar-Kraus-space decomposition is needed at these orders.
\end{theorem}
\begin{proof}
Vanishing of \eqref{sel17:eq:Lie-test} is commutation with a Lie generating
set.  Exponentiation yields invariance under
$\mathrm{SU}(3)\times\mathrm{SU}(2)$ and the central circle
$(e^{-2it}I_3,e^{3it}I_2)$.  These generate the connected $G_{\det}$, as
in Corollary~\ref{sel16:cor:gcd}; the converse follows by differentiation.
For this exact symmetry class, Theorem~\ref{sel17:thm:bandwidth} leaves
only indices $-1,0,1$.  The three-point discrete Fourier sum in
\eqref{sel17:eq:Fourier} therefore has no aliasing.  Although $g_{2\pi}$
need not equal the group identity, it lies in $G_{\det}$, so the Choi orbit is
$2\pi$ periodic.  At $\phi=\pi$ its two nonconstant modes change sign.
Their orthogonality and equal squared norms give the first equality in
\eqref{sel17:eq:phase-purity}.  The two Choi operators have the same squared
trace, giving the second.  Differentiation of the orbit gives
$i[\widehat N_C/3,J_y]$ and multiplies each mode by $iq$; this proves
\eqref{sel17:eq:derivative}.  A nonconstant mode has gcd one, proving
\eqref{sel17:eq:finite-selector}.
\end{proof}

The finite covariance equations are tests, not a claim that full source
covariance follows from group labels or an algebra with the desired block
sizes.  A nonzero $r_{\det}$ excludes this sufficient selector.  Replacing the
channel by its $G_{\det}$ average would change the source and is not a proof
that the residual was zero.

\begin{proposition}[Nontriviality is generic within the verified minimal-port class]
\label{sel17:prop:generic}
On the full $(a,b)=(2,3)$ native ports, within the finite-dimensional convex
set of normalized $G_{\det}$-covariant channels, the subset with covariance
group exactly $G_{\det}$ is relatively open and dense.  Its complement
within the feasible set has zero Lebesgue measure in the affine hull of that
set.  The exceptional $G_0$-covariant set is a proper affine section.  In particular nontriviality is not forced by covariance and
normalization, since the averaged channel is always allowed.
\end{proposition}
\begin{proof}
The squared nonconstant-mode norm is continuous.  Its zero locus is the
intersection with a linear subspace.  The channel of
Theorem~\ref{sel17:thm:sharp} with $\eta=1$ lies outside that subspace,
so it is proper in the affine hull of the feasible set.  Mixing any feasible
point with this witness gives arbitrarily nearby points with nonzero mode;
a nonzero affine-line mode can vanish at at most one mixing parameter.
A proper affine subspace has zero Lebesgue measure in the affine hull.
Corollary~\ref{sel17:cor:primitive} identifies every nonzero point with the
claimed group.  Averaging produces the explicit exceptional points.
\end{proof}

This is source-category-relative genericity.  It is not a probability law for
historical sources or a derivation of the twelve covariance equations from
arbitrary microscopic dynamics.

\begin{proposition}[An interference panel and a one-sided error certificate]
\label{sel17:prop:readout}
For an identified two-line input $(n_i,w_i)$ and output $(c_o,n_o)$, an actual
channel determines
\[
 z=\langle c_o|\mathcal J(|n_i\rangle\langle w_i|)|n_o\rangle.
\]
It is reconstructed by the four input states with vectors
$(n_i\pm w_i)/\sqrt2$, $(n_i\pm i w_i)/\sqrt2$ and two output Pauli
readouts.  In the convention $Y=-iE_{01}+iE_{10}$, let
$A_X,A_Y$ be the output $X,Y$ expectation differences for the two $X$ inputs,
and $B_X,B_Y$ the corresponding differences for the two $Y$ inputs.  Then
\begin{equation}
 \boxed{z=\frac{A_X+B_Y+i(B_X-A_Y)}4.}
 \label{sel17:eq:readout}
\end{equation}
For \eqref{sel17:eq:instrument}, $z=\lambda\eta$ on either the success
branch or the forgotten channel, and $\mathfrak h=\lambda^2|\eta|^2$.

More generally, under the exact hypotheses of
Theorem~\ref{sel17:thm:finite}, suppose the direct-sum HS errors in the
estimated packets at phases zero and $\pi$ are at most $\epsilon_0$ and
$\epsilon_\pi$.  If
\begin{equation}
 \left(\sum_y\|\widehat J_y-\widehat J_y(\pi)\|_{\HS}^2\right)^{1/2}
       >\epsilon_0+\epsilon_\pi,
 \label{sel17:eq:noise}
\end{equation}
the true fixed-label group is $G_{\det}$.
\end{proposition}
\begin{proof}
The differences of the two input state pairs are $X$ and $Y$, and
$E_{01}=(X+iY)/2$.  For a Hermitian output $Z$,
$Z_{01}=(\Tr(XZ)-i\Tr(YZ))/2$.  Apply these two identities successively to
obtain \eqref{sel17:eq:readout}.  The explicit arm formula gives $z$ and its
single off-diagonal Choi mode gives $\mathfrak h$.  The last statement is the
reverse triangle inequality in the direct-sum HS space, followed by
Theorem~\ref{sel17:thm:finite}.
\end{proof}

The displayed preparations are entangled alternating-sector comparisons on
actual tensor ports.  They are not guaranteed by single-copy scalar future
statistics.  Equivalent physical process tomography can supply the same
matrix element without this particular measurement implementation.  Small
nonzero estimates of $r_{\det}$ do not certify exact symmetry: a robust exact
group conclusion still requires its exact source hypotheses.  If an event has
an unnormalized physical weight $w$, then $z$ scales by $w$, $\mathfrak h$
by $w^2$, and the diamond separation by $w$.  Fixed-opportunity accepted
histories retain their original powers of $\vartheta$.

\section{Audit the unchanged six-letter source rather than assign it the port}
\label{sel17:old-source}

The preceding existence theorem must not be attached to the V13/V15
full-record witness simply because that witness already has the correct
internal invariant algebra.  V16 explicitly distinguishes an additional typed
incidence packet from covariance of all raw source coefficients.  Here this
separation has an exact evaluation.

\begin{corollary}[The identity-calibrated full source has no nontrivial fixed-label symmetry]
\label{sel17:cor:old-covariance}
Let the original coefficients $K_1,\ldots,K_s$ be linearly independent with
$\sum_eK_e=cI$, $c\ne0$.  A unitary $U$ preserves each resolved CP branch
$K_e(\cdot)K_e^*$ if and only if $U$ commutes with every $K_e$.  Hence if
$C^*(I,K_e)=M_N$, every such $U$ is scalar.

For the unchanged V13/V15 full-record source, no nontrivial connected internal
$\mathrm{SU}(3)$ or $\mathrm{SU}(2)$ action is a symmetry of that complete
fixed-label instrument.  Allowing a \emph{continuous} action on its finite
outcome set does not change the conclusion for a connected group.
\end{corollary}
\begin{proof}
Equality of the two rank-one Choi operators gives
$UK_eU^*=z_eK_e$ with $|z_e|=1$.  Conjugate the independent identity
expansion to obtain $\sum_e(z_e-1)K_e=0$, and hence $z_e=1$ for every $e$.
This is the fixed-label specialization of the inherited phase-lock lemma
\ref{sel7:lem:phase-lock}, not a new phase-choice principle.
The converse is immediate.  The full matrix algebra has scalar commutant,
so the first statement follows.  Theorem~\ref{sel13:thm:full-panel} establishes
that full algebra for the retained witness.  A continuous homomorphism from
a connected group to a finite permutation group is constant, giving the final
statement.
\end{proof}

\begin{proposition}[An exact nonabelian covariance defect in the existing witness]
\label{sel17:prop:old-defect}
Keep exactly the V13 point
$(u,v,\zeta)=(3/5,4/5,i)$, the six coefficients in their original order, and
$K=K_{1,+}$.  Let $T=\diag(1,-1,0)$ on the native trivial sector and zero
elsewhere, an internal colour Lie generator.  Then
\begin{align}
 \|K\|_{\HS}^2&=\frac73,&
 \|[T,K]\|_{\HS}^2&=\frac{28}{75},&
 \Tr(K^*[T,K])&=\frac{16}{225},\nonumber\\
 &\boxed{\|[\widehat T,J_K]\|_{\HS}^2=\frac{87688}{50625}>0.}
 \label{sel17:eq:old-defect}
\end{align}
Here $J_K=|K\rangle\!\rangle\langle\!\langle K|$ and
$\widehat T=T\otimes I-I\otimes\overline T$.  This is the unnormalized
accepted-branch defect.  Its squared HS value at one locked opportunity is
multiplied by $\vartheta^2$.
\end{proposition}
\begin{proof}
The first three identities follow by substitution of the original V13
orthogonal-column deformation; the accompanying exact script uses those same
matrices.  For $k=|K\rangle\!\rangle$ and
$l=|[T,K]\rangle\!\rangle$, the commutator is
$lk^*-kl^*$.  The generator is Hermitian, so $k^*l$ is real, and
\[
 \|lk^*-kl^*\|_{\HS}^2
 =2\|k\|^2\|l\|^2-2(k^*l)^2.
\]
Inserting the three rational numbers gives the displayed defect.
The opportunity weight multiplies $J_K$ by $\vartheta$.
\end{proof}

The inherited internal intersection is not lost:
$\mathcal A_K\cap U(S_4)'$ still has its proved structural type.  The new
calculation says that this algebra must not be reinterpreted as a symmetry
of the complete original efficient instrument.  A proposed gauge-covariant
incidence port needs its own actual source maps or a verified covariant
reduction.  It cannot be obtained merely by applying the target name to the
old full matrix algebra, by taking its tensor powers, or by analytically
projecting away the above covariance defect.

\section{Source-level conclusion and reproducible verification}
\label{sel17:status}

The upstream task is now smaller and more sharply delimited.  An operational
determinant selector does not require a neutral amplitude and a full
five-copy determinant amplitude.  A coherent comparison between a
colour-determinant \emph{write} and a weak-determinant \emph{read} can impose
the same primitive condition on a two-input, three-output native port.  The
need is still coherent source data, but the neutral-reference assumption and
the five-parallel-copy construction are not necessary for this CP-level
conclusion.

The exact lower bound also removes two unproductive routes.  No
$G_{\det}$-covariant operation below five total slots can carry the missing
mode, and arbitrary depth of actual covariant two-copy memory steps does not
change that fact.  Complete lower-order marginals can remain identical between
a phase-blind source and a primitive selector, so increasing their precision
cannot reconstruct the omitted five-slot correlation.  The source must
supply a larger physical port, a counted charged memory, or a different
verified charge inventory.

At five through nine total slots, exact nonabelian and hypercharge covariance
plus one nonzero central phase response give the \emph{connected} determinant
group; the unwanted higher-harmonic component ambiguity is excluded by the
finite representation budget itself.  These covariance equations are
finite and testable.  Their vanishing is not derived from arbitrary renewal
normalization, and the new genericity statement remains relative to that
verified class.

The next historical comparison should therefore inspect the earliest actual
typed tensor-source panel with at least five native input/output slots and
its complete retained cross correlations.  The tomographically complete
multitime entrance in \cite{sel17-Spacetime} explains how a supplied panel is
reconstructed; it does not assert that every proposed parallel port or
coherent contraction is admitted.  On a valid low-order port, evaluate
\eqref{sel17:eq:Lie-test} and the phase response
\eqref{sel17:eq:finite-selector}.  A positive result closes the operational
determinant-group step.  A missing port, a nonzero covariance residual, or a
zero nonconstant mode is a different and explicitly reported obstruction.
The unchanged V13 source was evaluated above precisely to avoid silently
substituting it for this still separate incidence object.

No historical five-slot data have been populated here.  The tensor-port
comparison does not determine the named type/private rays, the charged matter
modules or their Yukawa coefficients, a compatible physical fermionic grading,
or an Einstein--Standard-Model continuum.  Nor does it show that the positive
comparison instrument is implemented by the old source.  These restrictions
preserve the original source-complete joint-packet requirement while reducing
the exact additional covariance/character information needed for its
determinant handoff.

\paragraph{Proof and artifact status.}
The complete V16 prefix is retained byte for byte.  The companion script checks
the native input/output charge ranges, the primitive and higher-harmonic
thresholds, the explicit alternating input/output embeddings, CP normalization,
mode reconstruction, the interference-readout formula, sharp diamond-norm
witnesses, the vanishing of all elementary proper marginal cross terms, and
the unchanged V13 nonabelian defect.  The all-channel, all-width and minimum
claims have analytic proofs above; finite enumeration is only a check of their
formulas.  The supplied V16 regression is rerun.  There is no
proof-assistant formalization, independent peer review, or mathematical
priority claim for the general harmonic-analysis methods.

\begin{thebibliography}{99}
\raggedright
\bibitem[28]{sel17-CDP}
G.~Chiribella, G.~M.~D'Ariano and P.~Perinotti,
\emph{Theoretical framework for quantum networks},
Phys.\ Rev.\ A \textbf{80} (2009), 022339;
\href{https://doi.org/10.1103/PhysRevA.80.022339}{doi:10.1103/PhysRevA.80.022339}.
Cited for the established distinction between complete quantum networks and
their marginals; the source-specific slot bounds and separating instruments
are proved in this block.  The symmetry-mode reference \cite{sel16-MS} is
retained rather than duplicated.
\bibitem[29]{sel17-Spacetime}
A.~P\'elissier,
\emph{Renewal Geometry and the Emergence of Lorentzian Spacetime}, supplied
manuscript, Section~2, \src{mt:grand-tensor} and
\src{prop:readable-relational-completion}.  Used for the declared
positive multitime entrance, not for occurrence of the tensor-port
comparison constructed here.
\end{thebibliography}
\addtocontents{toc}{\protect\endgroup}
% END OF SOURCE-SELECTION V17 DEVELOPMENT BLOCK.
% Preserve V1--V17 and append subsequent work before the final terminator.


\clearpage
\begin{center}
{\Large\bfseries Source-selection continuation V18}\\[.4em]
{\large Deriving the central direction from a finite tensor budget,\\
with exact alternative groups and forced determinant alternation}
\end{center}
\addtocontents{toc}{\protect\begingroup}
\addtocontents{toc}{\protect\renewcommand{\protect\numberline}{\protect\makebox[2.1em][l]}}

\section{Go upstream of the assumed hypercharge direction}
\label{sel18:context}

V17 establishes the first tensor order at which a recorded CP operation can
separate $S(\mathrm U(3)\times\mathrm U(2))$ from the full product group.
Its bandwidth theorem assumes covariance under that determinant subgroup,
including the central generator $-2\widehat N_C+3\widehat N_W$.
That is appropriate for a global-group refinement, but it does not derive
the central direction.  The present continuation removes this premise.
Only covariance under the two independently identified nonabelian factors
is imposed initially; the remaining central symmetry is reconstructed from
a two-by-two positive response Gram.

The main theorem has a general $\mathrm U(n)\times\mathrm U(m)$ form.
At the first mixed tensor orders, one surviving continuous central direction
with opposite signs on the defining factors forces the determinant relation.
For $n=3,m=2$ and five or six total native input/output factors, the exact
criterion is rank one of that response Gram and positivity of its mixed
entry.  It derives both the ratio $-2:3$ and the connected global subgroup,
without assuming a hypercharge coefficient, determinant tensor, charged
matter list, or anomaly equation.

\paragraph{Audit and nonduplication.}
The general marked integer-lattice stabilizer, fixed-amplitude/CP distinction,
and CP normalization examples of V16 are inputs.  V17 already proves the
bandwidth on the \emph{prescribed} determinant line, the five-factor
comparison channel and the insufficiency of its lower marginals.  None is
claimed again as a new result.  The older anomaly calculation
\cite{Attack}, \src{thm:V22-anomaly-hypercharge}, determines hypercharge
\emph{after} specifying charged matter types.  Here the direction is obtained
from the native tensor representation and measured symmetry response instead.
The additions are the two-centre charge polygon, the converse central selector,
the complete first-order lattice alternatives, and the tensor-factorization
forced at the first same-sign mode.  Harmonic decomposition of quantum
operations is established background \cite{sel16-MS}; all specialized bounds
and converse implications used here are proved below.

\paragraph{What remains physical.}
The tensor ports, their defining representations, and nonabelian covariance
must occur on the actual source, with each old label fixed.  A centre-response
rank is a property to test, not a consequence of source normalization.
No five-factor historical process table is supplied by these calculations.
The precise balance between mathematical reconstruction and occurring mixed
operations remains the joint-packet requirement of \cite{Duality},
\src{def:joint-packet}.  A covariance group of an incidence port is not the
commutant of all raw source coefficients.

\section{The two-centre charge polygon and its sharp operational cost}
\label{sel18:polygon}

For integers $n,m\ge2$ let
\[
 H=C\oplus(F\otimes W)\oplus N_0,\qquad
 \dim C=n,\quad\dim W=m,\quad \dim F,\dim N_0\ge1,
\]
with candidate action
\[
 R_0(U,V)=U\oplus(I_F\otimes V)\oplus I_{N_0},
 \quad G_0=\mathrm U(n)\times\mathrm U(m),
 \quad K=\mathrm{SU}(n)\times\mathrm{SU}(m).
\]
The specialization retained from V15--V17 is $(n,m,\dim F,\dim N_0)
=(3,2,3,5)$.  Original scalar provenance labels inside $N_0$ are not erased
by this charge notation.  Let $H_a,H_b$ be invariant subspaces of
$H^{\otimes a},H^{\otimes b}$ and put $s=a+b$.
Independent finite gauge-neutral ancillary factors are permitted and do not
contribute to $s$.  A complete original instrument has positive Choi
operators $J_y$ on $H_b\otimes\overline{H_a}$, with output-first convention
and representation $\mathcal R=R_b\otimes\overline{R_a}$.
The only group hypothesis in this section is
\begin{equation}
 \mathcal R(k)J_y\mathcal R(k)^*=J_y
 \quad (k\in K,\text{ every retained }y).
 \label{sel18:eq:Kcov}
\end{equation}
Write $\chi_{p,q}(U,V)=(\det U)^p(\det V)^q$ and define
\begin{equation}
 \ell_{n,m}(p,q)
 =\max\{|np|,|mq|,|np+mq|\}
 =\frac{|np|+|mq|+|np+mq|}{2}.
 \label{sel18:eq:hexnorm}
\end{equation}

\begin{theorem}[Complete two-centre bandwidth]
\label{sel18:thm:polygon}
Every $J_y$ satisfying \eqref{sel18:eq:Kcov} has an orthogonal expansion
\begin{equation}
 \boxed{J_y=\sum_{\ell_{n,m}(p,q)\le s}J_y^{[p,q]},\qquad
 \mathcal R(g)J_y^{[p,q]}\mathcal R(g)^*
       =\chi_{p,q}(g)J_y^{[p,q]}.}
 \label{sel18:eq:expansion}
\end{equation}
The adjoint exchanges $(p,q)$ and $(-p,-q)$.  An invariant port restriction
can delete allowed modes but cannot add modes outside this polygon.
Let $L_J$ be the integer span of the indices of nonzero marked modes.  Then
\begin{equation}
 \boxed{G_J=\pi^{-1}(\operatorname{Ann}L_J),\qquad
 \pi(U,V)=(\det U,\det V).}
 \label{sel18:eq:whole-lattice}
\end{equation}
This concerns modes of the \emph{whole Choi operator}; no character-scalar
Kraus-space assumption is required.
\end{theorem}
\begin{proof}
The subspace of operators fixed by the normal subgroup $K$ carries a unitary
representation of $G_0/K\cong\mathbb T^2$.  Simultaneous diagonalization
of commuting unitaries decomposes it into integer torus characters, giving
the expansion before its bound.  On one native Hilbert factor, the central
weights for $(e^{i\alpha}I_n,e^{i\beta}I_m)$ are
$(1,0),(0,1),(0,0)$.  Conjugation on its endomorphisms has weights
\[
 (0,0),\ (\pm1,0),\ (0,\pm1),\ (1,-1),\ (-1,1).
\]
Conjugating an input factor reverses these weights and hence leaves their
set unchanged.  On $s$ factors, every weight $(x,y)$ is a sum of $s$ members
of that set.  Each summand has $|x|,|y|,|x+y|\le1$, so the sum has
$|x|,|y|,|x+y|\le s$.  The character $\chi_{p,q}$ has central weight
$(np,mq)$, proving \eqref{sel18:eq:hexnorm}.  The equality between the two
expressions there follows by considering whether $np$ and $mq$ have equal
or opposite signs.  Neutral ancillas add only weight zero.
Orthogonality of distinct torus characters and Hermiticity of $J_y$ prove
the remaining mode statements.  A group element fixes all modes precisely
when its determinant pair annihilates every occurring index.  This is the
same lattice argument as Theorem~\ref{sel16:thm:lattice}, now applied to
the complete $K$-fixed Choi space rather than a scalar amplitude subspace.
\end{proof}

\begin{theorem}[Sharp cost of any determinant-character constraint]
\label{sel18:thm:cost}
For $(p,q)\ne(0,0)$, the least total native tensor order, minimized over
input/output splits, of a normalized CP source with exact covariance group
$\ker\chi_{p,q}$ is
\begin{equation}
 \boxed{s_{\min}(p,q)=\ell_{n,m}(p,q).}
 \label{sel18:eq:cost}
\end{equation}
The bound is attained by a state-preparation channel on that many output
factors.  This is a sharp comparison construction, not an assertion that
such a preparation is historically admitted.
\end{theorem}
\begin{proof}
Equality of annihilators determines the original integer lattice, for example
by Smith normal form as in Theorem~\ref{sel16:thm:lattice}.  Thus a source
with the stated exact group has $L_J=\mathbb Z(p,q)$.  Some nonzero mode is
$k(p,q)$ with $k\in\mathbb Z\setminus\{0\}$, and its cost is
$|k|\ell_{n,m}(p,q)$.  The bandwidth theorem proves the lower bound.

Choose a unit neutral vector $r$, one weak multiplicity copy, and normalized
alternating vectors $\Omega_C\in C^{\otimes n}$ and
$\Omega_W\in W^{\otimes m}$.  If $p,q\ge0$, with at least one positive,
use the superposition of
$\Omega_C^{\otimes p}\otimes\Omega_W^{\otimes q}$ and
$r^{\otimes(np+mq)}$.  These are orthogonal unit vectors with relative
character $(p,q)$.  If $p>0,q<0$, put $s=\max(np,m|q|)$ and superpose
\[
 \Omega_C^{\otimes p}\otimes r^{\otimes(s-np)},\qquad
 \Omega_W^{\otimes|q|}\otimes r^{\otimes(s-m|q|)}.
\]
Again their relative character is $(p,q)$.  Padding placements are fixed.
Changing the overall sign of $(p,q)$ exchanges the two vectors and does not
change the group.  The normalized pure-state density of either balanced
superposition has exactly the three modes $0,\pm(p,q)$ and is a normalized
channel from a scalar input.  Its stabilizer is exactly $\ker\chi_{p,q}$,
including all components when $p,q$ have a common divisor.  This proves
attainability in every case.  It does not assert that every fixed port split
attains the same minimum.
\end{proof}

\begin{remark}[Mixed signs can occur earlier]
For $n=3,m=2$, the characters $(1,1),(1,-1),(1,-2)$ have exact costs
$5,3,4$, respectively.  Nonabelian covariance alone does not privilege the
first of these.  Extending V17's bound $5|q|$ to all mixed determinant
characters would be incorrect: it holds on its prescribed diagonal line,
not on the full two-centre lattice.
\end{remark}

\section{A positive central response derives the first mixed direction}
\label{sel18:central}

On the Choi Hilbert space define the Hermitian charge operators
\[
 \widehat N_C=N_{C,b}\otimes I-I\otimes\overline{N_{C,a}},\qquad
 \widehat N_W=N_{W,b}\otimes I-I\otimes\overline{N_{W,a}},
\]
where $N_{C,a},N_{W,a}$ count the indicated defining factors in the
input representation, and similarly on output.  These are reconstructed
representation operators, not newly admitted controls.  Define the real
symmetric response Gram
\begin{equation}
 \mathsf F_{ij}
 =\operatorname{Re}\sum_y
 \Tr\bigl([\widehat N_i,J_y]^*[\widehat N_j,J_y]\bigr),
 \qquad i,j\in\{C,W\}.
 \label{sel18:eq:F}
\end{equation}
All original probabilities are retained in the unnormalized $J_y$.

\begin{proposition}[Positive mode decomposition and the surviving centre]
\label{sel18:prop:F}
Set $a_{p,q}=\sum_y\|J_y^{[p,q]}\|_{\HS}^2$ and
$v_{p,q}=(np,mq)^{\mathsf T}$.  Then
\begin{equation}
 \boxed{\mathsf F=\sum_{p,q}a_{p,q}v_{p,q}v_{p,q}^{\mathsf T}\succeq0.}
 \label{sel18:eq:F-modes}
\end{equation}
Its kernel is the Lie algebra of the surviving central subgroup, expressed
in the defining angles $(\alpha,\beta)$.  Its rank equals
$\operatorname{rank}_{\mathbb Z}L_J$.  In particular,
\begin{equation}
 \det\mathsf F
 =\frac12\sum_{\omega,\nu}a_\omega a_\nu
          \det(v_\omega,v_\nu)^2.
 \label{sel18:eq:F-det}
\end{equation}
A nonzero rank-one Gram therefore means that every occurring nonzero mode
lies on one rational line.  Its integer index is not determined by the
Gram in general.
\end{proposition}
\begin{proof}
Differentiating the mode character gives
$[\widehat N_C,J_y^{[p,q]}]=npJ_y^{[p,q]}$ and the corresponding formula
with $mq$.  Distinct mode spaces are orthogonal.  Substitution gives
\eqref{sel18:eq:F-modes}.  For a real vector $t$ its quadratic form is
$\sum_y\|[t_C\widehat N_C+t_W\widehat N_W,J_y]\|_{\HS}^2$.
It vanishes exactly when exponentiation of that generator fixes every
$J_y$ for all real times.  A positive sum of rank-one matrices has range
the real span of their vectors.  Integer vectors have the same real and
integer-span ranks.  Expanding the two-by-two determinant proves
\eqref{sel18:eq:F-det}; equivalently this is Cauchy--Binet.
\end{proof}

\begin{theorem}[First mixed-centre selection without a prescribed ratio]
\label{sel18:thm:first-mixed}
Assume only \eqref{sel18:eq:Kcov}, the native port specification, and
\begin{equation}
 s<n+m+\min(n,m).
 \label{sel18:eq:first-window}
\end{equation}
The following conditions are equivalent:
\begin{enumerate}[label=\textup{(\roman*)},leftmargin=2.8em]
\item $\rank\mathsf F=1$ and $\mathsf F_{CW}>0$;
\item the nonzero nonconstant modes are exactly $(1,1)$ and $(-1,-1)$;
\item $G_J=S(\mathrm U(n)\times\mathrm U(m))$.
\end{enumerate}
On this branch $s\ge n+m$, and for
$h=\sum_y\|J_y^{[1,1]}\|_{\HS}^2>0$,
\begin{equation}
 \boxed{\mathsf F=2h
  \begin{pmatrix}n^2&nm\\nm&m^2\end{pmatrix},\qquad
 \Ker\mathsf F=\mathbb R(-m,n).}
 \label{sel18:eq:derived-centre}
\end{equation}
Thus the relative central coefficients are an output of the source test,
not an assumed hypercharge equation.  The group is connected; no extra
character-power component survives at this order.
\end{theorem}
\begin{proof}
If (i) holds, Proposition~\ref{sel18:prop:F} makes all nonzero mode indices
collinear.  The positive mixed entry means their two nonzero coordinates
have the same sign: along a fixed line the product $np\,mq$ has a constant
sign, independent of the orientation or magnitude.  By Hermiticity an
occurring mode can be chosen with $p,q>0$.  These are positive integers,
and its cost is $np+mq$.  The smallest is $n+m$, attained only by $(1,1)$;
the next possible value is at least $n+m+\min(n,m)$.  Hence
\eqref{sel18:eq:first-window} forces $p=q=1$, and all other nonzero modes
are its adjoint pair.  This proves (ii).  That pair generates the primitive
lattice $\mathbb Z(1,1)$, so \eqref{sel18:eq:whole-lattice} proves (iii).
Conversely, (iii) forces that exact lattice.  Each nonconstant mode is
$k(1,1)$ and the port bound excludes $|k|\ge2$.  The lattice is nonzero,
so (ii) holds.  The adjoint modes have equal norms, giving
\eqref{sel18:eq:derived-centre} and hence (i).  Connectedness follows, for
example, from the connected determinant fibres or the standard connected
cover of the determinant subgroup.  Primitivity has already been proved by
the occurring index $(1,1)$, not inferred from connected Lie algebras.
\end{proof}

\begin{corollary}[Five or six native factors fix the structural direction]
\label{sel18:cor:SM}
For the retained $\mathbb C^3\oplus(\mathbb C^3\otimes\mathbb C^2)\oplus\mathbb C^5$ native action and
$s\le6$, exact nonabelian covariance together with
\begin{equation}
 \boxed{\det\mathsf F=0,\qquad \mathsf F_{CW}>0}
 \label{sel18:eq:SM-test}
\end{equation}
forces $s\ge5$ and
\[
 \boxed{3\alpha+2\beta=0,\qquad
 G_J=S(\mathrm U(3)\times\mathrm U(2)).}
\]
Equivalently, within this port range it suffices to establish that exactly
one central direction survives and that its defining charges on $C,W$
have opposite signs.  Their magnitudes need not be specified in advance.
The familiar $\mathbb Z_6$ covering description follows from the inherited
connected determinant-group theorem.
\end{corollary}
\begin{proof}
The Gram is positive.  A positive off-diagonal entry makes it nonzero and
excludes a zero diagonal entry.  Its determinant zero is therefore exactly
rank one.  Apply Theorem~\ref{sel18:thm:first-mixed}; its kernel gives the
opposite-sign formulation.  The coordinates $(\alpha,\beta)$ are angles
of the independently identified defining actions, not a sign convention
chosen after observing the response.
\end{proof}

This is weaker than imposing $-2\widehat N_C+3\widehat N_W$ as a conserved
source generator, but is not a proof that a central direction survives in
every source.  Rank two, rank zero, and a negative mixed entry are permitted
and are classified next.  The result also does not determine the numerical
normalization of a gauge kinetic term or the charged matter representations.

\section{All first-five-factor central groups, including the rejected branches}
\label{sel18:alternatives}

\begin{theorem}[Complete five-factor lattice alternatives]
\label{sel18:thm:lattices}
For $n=3,m=2$ and $s\le5$, all nonzero possible mode indices belong to
\begin{equation}
 \boxed{\pm(0,1),\ \pm(0,2),\ \pm(1,-2),\
        \pm(1,-1),\ \pm(1,0),\ \pm(1,1).}
 \label{sel18:eq:five-indices}
\end{equation}
There are exactly eleven possible embedded lattices $L_J$, and each is
attained by a normalized marked source on at most five output factors.
They are the zero lattice; the six rank-one lattices generated by the six
listed representatives; and the four rank-two lattices
\begin{equation}
 \begin{split}
 &\mathbb Z^2,\qquad
 \{(p,q):q\equiv0\pmod2\},\\
 &\{(p,q):q\equiv p\pmod2\},\qquad
 \{(p,q):q\equiv p\pmod3\}.
 \end{split}
 \label{sel18:eq:rank-two-list}
\end{equation}
The corresponding rank-two annihilators in determinant coordinates are,
respectively,
\[
 \{(1,1)\},\quad
 \{(1,1),(1,-1)\},\quad
 \{(1,1),(-1,-1)\},\quad
 \{(\zeta^{-1},\zeta):\zeta^3=1\}.
\]
Every group is the inverse image of its indicated annihilator under $\pi$.
\end{theorem}
\begin{proof}
The polygon bounds give $|p|\le1$, $|q|\le2$ and $|3p+2q|\le5$,
which gives exactly \eqref{sel18:eq:five-indices}.  Signs of generators
do not affect a lattice.  The only distinct collinear nonzero representatives
are $(0,1),(0,2)$, so the rank-one list follows.
For rank two, at least one generator has first coordinate one.  Such a
lattice has generators $(1,r),(0,d)$, where $d>0$ and $r$ is taken modulo
$d$.  The vertical number $d$ is the gcd of any available vertical generators
$1,2$ and the differences of selected values from $\{-2,-1,0,1\}$.
If $d=1$, the lattice is $\mathbb Z^2$.  If $d=2$, the selected values
have a common parity, giving the two displayed congruences.  A value $d=3$
is possible only with the two values $-2,1$ and no vertical generator;
then $r\equiv1\pmod3$.  No larger positive $d$ is possible.  Solving
$z_1z_2^r=1$, $z_2^d=1$ gives the annihilators.

To see attainability, use the pure preparation construction of
Theorem~\ref{sel18:thm:cost} for each desired representative, padding with
neutral factors to five when needed.  For any selected collection, choose
its preparations with strictly positive classical probabilities and retain
the choice as the original outcome label.  The sum of the output traces is
one.  Its marked mode lattice is precisely the sum of their individual
lattices.  This produces all eleven possibilities.  It constructs comparison
sources, not conditional probabilities for the historical source.
\end{proof}

\begin{center}
\small
\begin{tabular}{L{.22\textwidth}L{.30\textwidth}L{.30\textwidth}}
\toprule
Rank-one generator & Central equation & Defining surviving direction\\
\midrule
$(0,1)$ & $\det V=1$ & $(1,0)$\\
$(0,2)$ & $(\det V)^2=1$ & $(1,0)$, with two components\\
$(1,0)$ & $\det U=1$ & $(0,1)$\\
$(1,-1)$ & $\det U=\det V$ & $(2,3)$\\
$(1,-2)$ & $\det U=(\det V)^2$ & $(4,3)$\\
$(1,1)$ & $\det U\det V=1$ & $(-2,3)$\\
\bottomrule
\end{tabular}
\end{center}

\begin{proposition}[The sign and tensor window are essential]
\label{sel18:prop:counterexamples}
Normalized nonabelian-covariant sources with one surviving central direction
need not satisfy the determinant relation.  On three output factors,
$(\Omega_C+\Omega_W\otimes r)/\sqrt2$ selects $\ker\chi_{1,-1}$;
on four, $(\Omega_C\otimes r+\Omega_W^{\otimes2})/\sqrt2$ selects
$\ker\chi_{1,-2}$.  They can both be padded neutrally to five factors.
Their response Grams have negative mixed entries.
At seven output factors, the normalized superposition of
$\Omega_C\otimes\Omega_W^{\otimes2}$ and $r^{\otimes7}$ instead has
rank-one response with \emph{positive} mixed entry but group
$\ker\chi_{1,2}$, not $G_{\det}$.  Thus the five-or-six-factor window in
Corollary~\ref{sel18:cor:SM} cannot be extended to seven on these hypotheses.
\end{proposition}
\begin{proof}
These are the normalized states in the sharp-cost theorem.  Each balanced
pure state has cross-mode squared norm $1/4$, so its response is
$\tfrac12v_{p,q}v_{p,q}^{\mathsf T}$.  The signs and distinct groups follow
by substitution.  Their orthogonal summands have different determinant
characters and are each nonabelian singlets.
\end{proof}

\begin{remark}[An infinitesimal Gram does not always fix the global group]
A two-ray weak-determinant packet of visibility $v$ at index $(0,1)$ has
$\mathsf F_{WW}=2v^2$.  At index $(0,2)$, visibility $v/2$ gives the same
value and the same zero other entries.  Both are positive normalized
preparations, but only the second has a two-component covariance group.
The target branch avoids this ambiguity by its mixed positive signs and
strict tensor budget, not because every central Gram determines an integer
lattice.
\end{remark}

\section{At first occurrence the alternating tensor is forced}
\label{sel18:alternation}

The preceding selector does not insert a chosen determinant vector.  At its
first allowed order even the nonabelian tensor symmetry is forced.  We state
the full ambient result; an invariant port restriction retains its intersection
with this space.

\begin{lemma}[The first fundamental singlet is the alternating line]
\label{sel18:lem:alternation}
The invariant subspace of $(\mathbb C^n)^{\otimes n}$ under
$\mathrm{SU}(n)$ is $\Lambda^n\mathbb C^n$, of dimension one.
\end{lemma}
\begin{proof}
Invariance under the determinant-one diagonal torus requires every basis
label to occur the same number of times in a nonzero tensor coefficient.
There are $n$ slots, so each label occurs exactly once.  The tensor is
therefore a sum over permutations of $e_1\otimes\cdots\otimes e_n$.
For $i\ne j$, invariance under the complexified Lie algebra implies that
the sum of $E_{ij}$ acting in all slots annihilates it.  For any resulting
ordered tensor with two $i$ labels and no $j$ label, there are exactly two
contributing permutations, differing by interchange of $i,j$.  Their
coefficients must sum to zero.  Such interchanges generate all permutations,
so every coefficient is its permutation sign times one constant.  The
alternating tensor is indeed invariant because the determinant is one.
\end{proof}

\begin{theorem}[Structure of every first mixed determinant mode]
\label{sel18:thm:alternation}
Let $s=a+b=n+m$ and use the full native ports with no extra ancillary
factors.  The entire character-$(1,1)$ subspace of
$\End(H^{\otimes b}\otimes\overline{H^{\otimes a}})$ is an orthogonal
sum over choices of the $n$ colour positions among the $s$ slots.  On each
choice its charged tensor factor is exactly
\begin{equation}
 \boxed{\Lambda^n C\otimes\Lambda^m W.}
 \label{sel18:eq:forced-alt}
\end{equation}
Every slot pairs a defining charged sector with a neutral sector; no neutral
spectator slot or colour--weak exchange slot can occur in this mode.
If $f=\dim F$ and $\nu=\dim N_0$, its complex dimension is
\begin{equation}
 \boxed{\binom{n+m}{n} f^m\nu^{n+m}.}
 \label{sel18:eq:mode-dim}
\end{equation}
For the retained native representation this is
$\binom53\,3^2\,5^5=281250$, independently of the split $(a,b)$.
The determinant tensor factor is compulsory, but neither a unique source ray
nor a unique assignment of typed matter ports is selected by this count.
\end{theorem}
\begin{proof}
Under the common central phase $(e^{it}I_n,e^{it}I_m)$ the required mode
has weight $n+m=s$.  Each local endomorphism factor has maximum such weight
one.  Equality therefore forces every slot to have weight one.  On an output
slot that space is
$\Hom(N_0,C\oplus(F\otimes W))$; on a conjugated input slot it is
$\Hom(\overline{C\oplus(F\otimes W)},\overline{N_0})$.
After reordering tensor factors each contributes one defining $C$ or $W$
factor and a neutral factor of dimension $\nu$, with an additional $F$
factor for each weak choice.  A colour--weak exchange has common weight zero
and cannot enter a maximal-weight summand.

The individual central weights $(n,m)$ force exactly $n$ colour and $m$
weak choices.  Different choices occupy orthogonal block patterns.  The
nonabelian invariant space of their charged factors is one-dimensional by
Lemma~\ref{sel18:lem:alternation} in each group.  All multiplicity and neutral
factors transform trivially and remain arbitrary.  There are
$\binom{n+m}{n}$ patterns, $f^m$ weak multiplicity coordinates and
$\nu^{n+m}$ neutral coordinates.  This proves the dimension formula as
well as the factorization.  Input conjugation replaces dual fundamental
factors in the domain by fundamental factors in this operator mode; hence
the count is independent of the input/output split.  Restriction to invariant
physical port subspaces can only remove summands or directions.
\end{proof}

At $s=5$, this theorem reconstructs a completely alternating three-colour,
two-weak \emph{CP correlation type} from a positive central selector.  It does
not identify that operator mode with the single up-type amplitude
$\Theta:C\otimes W^{\otimes2}\to\Lambda^2C^*$ of V110: a choice of
physical input/output typing and coherent amplitude remains additional data.
The theorem provides an exact projection/factorization test on a supplied
joint Choi panel, rather than a licence to execute its isolated mode.

\begin{corollary}[A compatible parity may exclude the primitive mode]
\label{sel18:cor:parity}
Suppose, as an additional actual source condition, the native grading is
\[
 \Gamma=(-1)^{\epsilon_C}I_C\oplus
       (-1)^{\epsilon_W}I_{F\otimes W}\oplus I_{N_0},
 \qquad \epsilon_C,\epsilon_W\in\{0,1\},
\]
and every original CP branch is covariant under its tensor input/output
action.  A mode $(p,q)$ must then satisfy
\begin{equation}
 n\epsilon_Cp+m\epsilon_Wq\equiv0\pmod2.
 \label{sel18:eq:parity}
\end{equation}
For $n=3,m=2$ and $\epsilon_C=1$, the primitive determinant mode is excluded
at every tensor order; at $s\le6$ the positive selector
\eqref{sel18:eq:SM-test} is impossible.
\end{corollary}
\begin{proof}
The displayed grading is the represented central element
$((-1)^{\epsilon_C}I_n,(-1)^{\epsilon_W}I_m)$ of $G_0$.
Its action on mode $(p,q)$ is
$(-1)^{n\epsilon_Cp+m\epsilon_Wq}$.  Fixed-label covariance requires this
to be one whenever the mode is nonzero.  Substitute $(p,q)=(1,1)$ and the
stated dimensions.  Equivalently the displayed grading has $\chi=-1$
and cannot belong to a group equal to $\ker\chi$.
\end{proof}

This is a conditional compatibility check, not a fermion-statistics
assignment to $C,W$.  A chirality operator which an odd Dirac amplitude flips
is not automatically a superselection symmetry required to preserve every
CP branch.  The supplied Clifford/finite grading must be compared on its
actual carrier before invoking this corollary.  An odd neutral sector or a
different graded matter inventory changes the displayed premise and is not
a counterexample to it.  No parity condition is silently added to V17.

\section{Finite reconstruction and honest error tests}
\label{sel18:tests}

The exact nonabelian condition is already a finite source comparison.  With
any Hermitian bases of $\mathfrak{su}(n),\mathfrak{su}(m)$ acting on the
actual Choi port, put
\begin{equation}
 \Delta_K=\sum_{y,T}\|[\widehat T,J_y]\|_{\HS}^2.
 \label{sel18:eq:K-test}
\end{equation}
It vanishes if and only if \eqref{sel18:eq:Kcov} holds, since the two
special unitary groups are connected and generated by these one-parameter
subgroups.  There is no central generator in \eqref{sel18:eq:K-test}.

\begin{proposition}[Finite phase grid for the unconstrained central lattice]
\label{sel18:prop:DFT}
Under exact \eqref{sel18:eq:Kcov}, let
$P_s=\lfloor s/n\rfloor$, $Q_s=\lfloor s/m\rfloor$,
$u=2P_s+1$, $v=2Q_s+1$ and
\[
 g_{jk}=\bigl(e^{2\pi i j/(nu)}I_n,
                 e^{2\pi i k/(mv)}I_m\bigr).
\]
Then every mode is reconstructed by
\begin{equation}
 \boxed{J_y^{[p,q]}=
 \frac1{uv}\sum_{j=0}^{u-1}\sum_{k=0}^{v-1}
 e^{-2\pi i(pj/u+qk/v)}\mathcal R(g_{jk})J_y\mathcal R(g_{jk})^*.}
 \label{sel18:eq:DFT}
\end{equation}
For five native factors this is a $3\times5$ grid; for six it is a
$5\times7$ grid.  No limiting derivatives are required.
\end{proposition}
\begin{proof}
On the $K$-fixed operator space the determinant phases are $2\pi$-periodic,
even though a chosen central lift need not be periodic on the full Hilbert
space.  The expansion has $|p|\le P_s$, $|q|\le Q_s$.
Discrete character orthogonality on the two grids separates all these
indices without aliasing.  Applying it to \eqref{sel18:eq:expansion} proves
the formula.  Polygon-forbidden modes within this rectangular grid vanish.
\end{proof}

The transformed Choi matrices can be calculated from physical process
tomography and independently known representation matrices.  Interpreting
the grid as a set of executed phase interventions additionally requires
admission of those operations.  Neither interpretation assigns physical
Choi coordinates to a scalar Hankel matrix without the V15 product audit.

\begin{proposition}[Response-error bound and a one-sided selection test]
\label{sel18:prop:error}
Suppose the marked Choi estimates obey
$\sum_y\|\widehat J_y-J_y\|_{\HS}^2\le\varepsilon^2$ on the same
identified representation.  Let $\mathscr D:\mathbb R^2\to
\bigoplus_y\End(H_b\otimes\overline{H_a})$ send $t$ to
$([t_C\widehat N_C+t_W\widehat N_W,J_y])_y$, with the target treated as a
real HS space.  Denote the estimated map by $\widehat{\mathscr D}$ and set
$e=\sqrt2\,s\varepsilon$.  Then
\begin{equation}
 \boxed{\|\widehat{\mathscr D}-\mathscr D\|\le e,\qquad
 \|\widehat{\mathsf F}-\mathsf F\|
 \le (2\|\widehat{\mathscr D}\|+e)e=:b_F.}
 \label{sel18:eq:error}
\end{equation}
If exact nonabelian covariance and $\rank\mathsf F=1$ have independent
justification, $\widehat{\mathsf F}_{CW}>b_F$ certifies the first-mixed
selector in its stated tensor window.  Conversely, a second singular value
of $\widehat{\mathscr D}$ greater than $e$ rules out any surviving continuous
central symmetry for the true packet.
\end{proposition}
\begin{proof}
Each count operator on the Choi carrier has spectrum in $[-a,b]$.  Subtract
its scalar midpoint, which does not change commutators.  Its operator norm
is then at most $s/2$, so
$\|[\widehat N_i,X]\|_{\HS}\le s\|X\|_{\HS}$.
The two column errors of $\mathscr D$ therefore have combined squared norm
at most $2s^2\varepsilon^2$.  The operator norm is bounded by this column
HS norm.  Expanding the difference of the two Gram operators gives the
second inequality.  Every entry error is bounded by $b_F$; the positive
mixed entry then follows.  Singular values change by at most the operator
norm perturbation, proving the rank-two obstruction.
\end{proof}

Small $\det\widehat{\mathsf F}$ does not prove exact rank one.  Likewise
small $\Delta_K$ does not prove exact nonabelian covariance.  There is no
unconditional noisy equality test for these exact symmetries in an arbitrary
continuous source family.  The positive certificate above explicitly retains
that logical boundary.  A deterministic weight $w_y$ in an event multiplies
$J_y$ and its mode amplitudes by $w_y$, and its contributions to $\mathsf F$
by $w_y^2$; locked opportunity factors are not cancelled.  Time or kinetic
normalizations are not inferred from this response metric.

\section{Evaluation on the existing tensor comparison, not a replacement source}
\label{sel18:evaluation}

Retain exactly V17's two-input, three-output instrument
\eqref{sel17:eq:instrument}, with the same $B_0,B_1$, success weight
$\lambda$, visibility $\eta$, and failure output.  The input rays have
central charges $(0,0),(0,2)$; the output rays have charges $(3,0),(0,0)$.
Nonabelian covariance follows from their alternating representation types,
not from an imposed hypercharge equation.

\begin{proposition}[The unchanged comparison passes the weaker selector]
\label{sel18:prop:witness}
For that complete marked source,
\begin{equation}
 \boxed{\Delta_K=0,\qquad
 \mathsf F=2\lambda^2|\eta|^2
       \begin{pmatrix}9&6\\6&4\end{pmatrix}.}
 \label{sel18:eq:witness}
\end{equation}
At $\eta\ne0$ the central direction and primitive global group now follow
from Corollary~\ref{sel18:cor:SM} without assuming that direction in the
covariance audit.  At $\eta=0$ the response is zero and the full product
group remains.  No coefficient, initial port, outcome label, or failure map
has changed from V17.
\end{proposition}
\begin{proof}
The success Choi cross mode is
$\lambda\eta|B_0\rangle\!\rangle\langle\!\langle B_1|$, with
squared norm $\lambda^2|\eta|^2$.  Its weight is $(1,1)$ since the arm
weights are $(1,0),(0,-1)$.  The diagonal and failure contributions are
$G_0$-invariant.  The two adjoint modes give
\eqref{sel18:eq:witness}.  The normalized source and the arm transformation
laws were verified in Theorem~\ref{sel17:thm:sharp}; the new point is the
converse direction supplied by the first-mixed theorem.
\end{proof}

The six-outcome V13 benchmark is \emph{not} substituted for this comparison.
Its previously computed nonabelian Choi defect in
Proposition~\ref{sel17:prop:old-defect} remains positive.
Accordingly it fails the weaker nonabelian entrance as well; removing an
assumed central equation does not repair that separate source defect.

\section{The remaining source question and verification scope}
\label{sel18:status}

This continuation replaces the prescribed hypercharge generator by a
source-response criterion.  On a native five- or six-factor port,
nonabelian covariance, one surviving continuous central symmetry, and opposite
signs of its defining charges force the determinant direction and its
primitive global form.  The first same-sign mode also has a forced alternating
three-colour/two-weak tensor factor.  A determinant line was not loaded into
the calculation to obtain either conclusion.

The sign criterion is still a real source condition.  The exact eleven-lattice
classification supplies the rejected alternatives, including wrong mixed
directions and extra disconnected groups.  Thus the result is not that
nonabelian covariance and a small tensor budget \emph{alone} select the
Standard Model.  It is a stronger conditional theorem with a directly
computable two-by-two test and a complete first-order alternative list.
Neither an algebraic determinant factor nor its response selects one of the
many neutral/multiplicity coordinate directions in the full first-mode space.

The next historical input is the same-cut complete tensor-process row,
its nonabelian comparison, and its central response, rather than a specified
hypercharge coefficient.  The physical tensor-port occurrence, product/CP
identification, and any additional grading restriction must first be verified.
The finite tomographic entrance of \cite{sel17-Spacetime} reconstructs a
\emph{supplied} exposed process; it does not assert admission of every
comparison port.  The named type/private states, charged matter modules,
anomaly cancellation on those modules, odd Dirac incidence, and continuum
measure questions remain outside this selector.  Historical structural
Standard-Model selection is not declared closed.

\paragraph{Exact checks and preservation.}
The complete V17 source prefix is retained byte for byte.  The companion
script enumerates the native charge polygons and all 64 first-five-factor
mode-support subsets; it finds precisely the eleven proved lattices.  It
checks the first-mixed window for multiple defining dimensions, the explicit
wrong-direction and seven-factor counterexamples, the central response of
the unchanged V17 normalized instrument, discrete Fourier nonaliasing,
uniqueness of the three-colour and two-weak alternating singlets, the full
first-mode multiplicity formula, the optional parity congruence, and physical
weight scaling.  The supplied V17 checks are rerun separately.  Analytic
all-source statements are proved above, not extrapolated from these finite
enumerations.  No historical probabilities are invented, and neither
proof-assistant certification nor independent peer review is claimed.

\addtocontents{toc}{\protect\endgroup}
% END OF SOURCE-SELECTION V18 DEVELOPMENT BLOCK.
% Preserve V1--V18 and append later work before the final terminator.


\clearpage
\begin{center}
{\Large\bfseries Source-selection continuation V19}\\[.4em]
{\large Source-locked symmetry, exact action compilation,\\
and the sharp native-memory cost of the determinant group}
\end{center}
\addtocontents{toc}{\protect\begingroup}
\addtocontents{toc}{\protect\renewcommand{\protect\numberline}{\protect\makebox[2.1em][l]}}

\section{Cross-program audit and a change of target}
\label{sel19:context}

V18 supplies a central-direction selector for an arbitrary nonabelian-covariant
CP tensor port.  It does not show that the port belongs to the independent,
coherently identity-calibrated six-outcome source class.  The newly supplied
regenerative notes contain an explicit action-to-instrument construction in
precisely that class.  This is the useful transfer: first determine its exact
symmetry and its restrictions, then ask whether its action and auxiliary
operations occur.  Repeating a generic five-slot CP comparison would miss this
compatibility question.

\paragraph{What the four supplied sources actually provide.}
The audit concerns their relevant constructions, proofs, and terminal scope
statements; it is not an independent verification of every archived result.
\begin{center}\small
\begin{tabular}{L{.28\textwidth}L{.64\textwidth}}
\toprule
Source & Retained result and limit of the transfer\\
\midrule
Native regenerative stationarity V21
& Versions 12 and 15 give the exact six-outcome action compiler and its
  gauge-commuting specialization; Version 20 supplies neutral multiplicity in
  a pure-gauge plaquette.  Gauge action, action coefficients, contexts and
  preparation are declared there.  They do not select an unknown gauge group
  or populate our historical tensor port.\\[.45em]
Yang--Mills mass-gap V43
& Conditional centering has unit Jacobian and an exact telescoping action.
  The zero-background innovation Hessian has a uniform free floor, but the
  original Haar/pivot density, mixed Wilson interaction and transported wall
  remain.  No symmetry of a truncated or reweighted source is imported.\\[.45em]
Exact-action Einstein bridge V16
& The actual source excites a diagonal soft sector, while the retained flux
  excludes its nonzero limiting vectors from the finite-energy domain.
  The compression estimate is not a bound for the full inverse.  It supplies
  no internal tensor incidence or operator realization of that incidence.\\[.45em]
Yang--Mills shell mixing V16
& The measured one-loop coefficient includes normal coarea, the constrained
  Hessian, gauge determinant and fine/coarse Haar factors.  Its flat-root
  evaluation is complete; the noncommuting source calculation is separately
  open.  This scalar determinant is not the determinant-character CP coherence
  of Versions 16--18.\\
\bottomrule
\end{tabular}
\end{center}
The exact source locations are recorded in
\cite{sel19-Reg,sel19-YM,sel19-Ein,sel19-Shell}.  These distinctions change
our task rather than silently strengthening a companion hypothesis.

\paragraph{Inherited mathematics, not new claims.}
We use the phase-locking implication of Lemma~\ref{sel7:lem:phase-lock}
and Corollary~\ref{sel17:cor:old-covariance}, the character-weight geometry of
Theorem~\ref{sel18:thm:polygon}, and the compiler of
\cite{sel19-Reg}, \src{v12:thm:locked-bank} and \src{v15:thm:locked}.
The general mode and coherent-implementation distinctions are established
background \cite{sel16-MS,sel16-AFCB}.  The additions are the stronger
\emph{locked-source} memory bound, a symmetry-exact converse for the
compiler, a linear-Choi source test, an explicit attaining native realization,
and an exact comparison showing that its limiting action does not certify
its finite recorded symmetry.  We make no priority claim for the underlying
finite spectral calculus.

\section{The linear Choi symmetry test on a locked source}
\label{sel19:linear}

Let $\mathcal H$ be one finite acted-on memory.  Its original efficient
instrument has maps $\mathcal J_e(X)=K_eXK_e^*$, with six independent
coefficients satisfying
\begin{equation}
 \sum_e K_e^*K_e=I,\qquad \sum_eK_e=\sqrt2 I.
 \label{sel19:eq:lock}
\end{equation}
The second equality uses the supplied coherent coefficient convention.  It
is not inferred from the classical outcome law.  With original labels fixed,
the inherited phase-locking result is
\begin{equation}
 \operatorname{Cov}_{G_0}(\mathcal J_e:e)
 =\{g\in G_0:[R(g),K_e]=0\text{ for every }e\}.
 \label{sel19:eq:phase-lock}
\end{equation}
Indeed covariance initially gives $R(g)K_eR(g)^*=z_eK_e$; conjugating the
independent identity expansion forces all $z_e=1$.  This use of the earlier
lemma is recalled to fix the meaning of symmetry, not asserted as a new lemma.
Without independence, charged terms could cancel in the identity expansion.

Use output-first vectorization $|K\rangle\!\rangle$ and
$J_e=|K_e\rangle\!\rangle\langle\!\langle K_e|$.  For Hermitian represented
Lie generators $T_a$ put $\widehat T_a=T_a\otimes I-I\otimes\overline{T_a}$.

\begin{proposition}[An intrinsic source response linear in the Choi data]
\label{sel19:prop:linear}
The real symmetric matrix
\begin{equation}
 \mathsf L_{ab}
 =\operatorname{Re}\sum_e\Tr([T_a,K_e]^*[T_b,K_e])
 =\frac12\sum_e\Tr\bigl(J_e\{\widehat T_a,\widehat T_b\}\bigr)
 \succeq0
 \label{sel19:eq:L}
\end{equation}
has kernel exactly the infinitesimal fixed-label covariance directions of
\eqref{sel19:eq:lock}.  It is unchanged by independent phase choices for the
six coefficients and is reconstructed linearly from their physical Choi
operators.  No purity comparison of two copies of $J_e$ is required.
\end{proposition}
\begin{proof}
Vectorization sends $[T_a,K_e]$ to $\widehat T_a|K_e\rangle\!\rangle$.
Taking the real part of its inner product with the corresponding $b$ vector
gives the anticommutator formula.  For real $x$ the quadratic form is
$\sum_e\|[\sum_a x_aT_a,K_e]\|_{\HS}^2$.  It vanishes exactly when all those
commutators vanish.  Exponentiation and \eqref{sel19:eq:phase-lock} give
both directions of the infinitesimal symmetry statement.  Multiplying one
coefficient by a phase leaves each inner product and its Choi operator
unchanged.
\end{proof}

The response differs from V18's \emph{quadratic} Choi covariance Gram
$\mathsf F$.  Their kernels agree on this verified locked source, but their
nonzero eigenvalues, event-weight scaling, and measurement errors need not.
The extra coherent calibration licenses the linear test; a general CP port
cannot be assigned that calibration for free.

\section{Identity locking changes the sharp native-memory bound}
\label{sel19:band}

Keep exactly the native action
\begin{equation}
 R_0(U,V)=U\oplus(I_3\otimes V)\oplus I_2\oplus I_3,
 \qquad G_0=\mathrm U(3)\times\mathrm U(2).
 \label{sel19:eq:native}
\end{equation}
A $k$-copy memory is an invariant subspace of $\mathcal H_{14}^{\otimes k}$,
possibly with finite neutral multiplicity.  Input and output are this same
memory, as required by \eqref{sel19:eq:lock}.  Let $N_C,N_W$ count its defining
colour and weak factors.  Suppose the complete efficient instrument is
covariant, with labels fixed, under
$K_0=\mathrm{SU}(3)\times\mathrm{SU}(2)$.

\begin{theorem}[Coefficient bandwidth and the locked-source group]
\label{sel19:thm:band}
Each coefficient has an orthogonal decomposition
\[
 K_e=\sum_{p,q}K_e^{(p,q)},\qquad
 R(g)K_e^{(p,q)}R(g)^*=(\det U)^p(\det V)^qK_e^{(p,q)},
\]
and every occurring index obeys
\begin{equation}
 \boxed{\max\{3|p|,\,2|q|,\,|3p+2q|\}\le k.}
 \label{sel19:eq:band}
\end{equation}
If $L_K$ is the integer span of the indices of nonzero coefficient components,
then
\begin{equation}
 \boxed{G_{\rm lock}=\pi^{-1}(\operatorname{Ann}L_K),
 \qquad \pi(U,V)=(\det U,\det V).}
 \label{sel19:eq:lattice}
\end{equation}
In particular, an exactly $S(\mathrm U(3)\times\mathrm U(2))$-covariant
locked instrument requires $k\ge5$.  A group exactly $\ker(\det U\det V)^d$
requires $k\ge5d$.
\end{theorem}
\begin{proof}
Equation~\eqref{sel19:eq:phase-lock} makes every coefficient commute with
$K_0$.  Conjugation on this invariant operator space factors through the
determinant torus.  Simultaneous unitary diagonalization gives its integer
characters.  A native Hilbert vector has central weight $(a,b)$ with
$a,b\ge0$, $a+b\le k$.  A matrix coefficient has target-minus-source weight
$(x,y)$, for which $|x|,|y|,|x+y|\le k$.  A determinant character has
$(x,y)=(3p,2q)$, proving \eqref{sel19:eq:band}.  Invariant subspaces and
neutral multiplicity cannot enlarge this range.

Orthogonality implies that a coefficient is fixed if and only if each of its
nonzero character components is fixed.  Intersecting these equations gives
\eqref{sel19:eq:lattice}, by the integer-annihilator argument already used in
Theorem~\ref{sel16:thm:lattice}.  The annihilator of the connected determinant
kernel is $\mathbb Z(1,1)$, and that of $\ker\chi^d$ is
$d\mathbb Z(1,1)$.  Equality with either proper group requires a nonzero
index in its respective lattice.  Its minimum cost in
\eqref{sel19:eq:band} is five or $5d$.
\end{proof}

This is a bound on \emph{coefficients}, not on their pairwise Choi differences.
V17's arbitrary equal-port CP selector has minimum $k=3$; it is not claimed to
satisfy \eqref{sel19:eq:lock}.  Nor does a two-input/three-output port have an
identity endomorphism with which to satisfy that equation.  The two bounds
refer to different admissible source classes, not competing estimates of the
same object.

\begin{corollary}[A locked-source central-direction selector]
\label{sel19:cor:selector}
Let $\mathsf L$ in \eqref{sel19:eq:L} use $N_C,N_W$.  Under the preceding
hypotheses, if $k\le6$, then
\begin{equation}
 \boxed{\det\mathsf L=0,\quad \mathsf L_{CW}>0
 \quad\Longleftrightarrow\quad
 G_{\rm lock}=S(\mathrm U(3)\times\mathrm U(2)).}
 \label{sel19:eq:selector}
\end{equation}
On this branch $k\ge5$ and $\Ker\mathsf L=\mathbb R(-2,3)$.
More generally, for defining dimensions $n,m\ge2$, the same assertion holds
in $k<n+m+\min(n,m)$, with cost $n+m$ and null direction $(-m,n)$.
\end{corollary}
\begin{proof}
Orthogonal coefficient components give
\[
 \mathsf L=\sum_{p,q}b_{pq}(3p,2q)^{\mathsf T}(3p,2q),
 \qquad b_{pq}=\sum_e\|K_e^{(p,q)}\|_{\HS}^2\ge0.
\]
A nonzero rank-one positive matrix has all its contributing vectors collinear.
A positive mixed entry forces that line to have both coordinates nonzero of
the same sign.  A nonzero index on it has cost $3|p|+2|q|$; below seven the
only possibility is $(1,1)$ or its negative.  Its lattice is primitive,
proving the forward implication.  Conversely the stated exact group forces
all indices into $\mathbb Z(1,1)$ and at least one to be nonzero.  The
bandwidth excludes every nonprimitive multiple at $k\le6$.  The same positive
sum gives rank one, positive mixed entry and the displayed kernel.  The proof
with $n,m$ replaces five and seven by $n+m$ and $n+m+\min(n,m)$.
Individual coefficients need not be Hermitian; no equality between the two
opposite mode masses is used.
\end{proof}

\begin{corollary}[Four-copy locked steps have no temporal workaround]
\label{sel19:cor:network}
If every memory-bearing primitive on at most four native copies satisfies
\eqref{sel19:eq:lock} and fixed-label determinant-group covariance, it is
already $G_0$-covariant.  Any finite network of these steps, $G_0$-invariant
preparations, $G_0$-covariant terminal operations and classical decisions remains
$G_0$-covariant.  Neutral ancillary memory does not remove the obstruction.
\end{corollary}
\begin{proof}
Theorem~\ref{sel19:thm:band} excludes every nonzero index on the determinant
line below five.  The resulting $G_0$ covariance is preserved under composition,
tensor product and nonnegative classical sums, as in
Proposition~\ref{sel16:prop:no-free-phase}.  Charged hidden memory would change
the counted representation and is not included in the statement.
\end{proof}

\section{Transfer the existing action compiler, with a converse}
\label{sel19:compiler}

Let $H=H^*$ be an independently specified finite action generator on the same
memory, with its represented candidate action denoted by $R_0(g)$ here.
Assume there is a five-dimensional reducing subspace $A$ on which
\emph{all of $G_0$ acts trivially}.  Choose five orthonormal eigenvectors
$u_0,\ldots,u_4$ of $H|_A$ and define
\[
 C_\mu=|u_0\rangle\langle u_{\mu-1}|\quad(2\le\mu\le5),\qquad
 R=\sum_{\mu=2}^5C_\mu^*C_\mu.
\]
These are the operations of the supplied compiler, now required to be neutral
under the larger comparison group, not merely under a desired subgroup.  No
new Hilbert directions are needed when $A$ is already present.  Their physical
admission is nonetheless additional source content.

Let $O:\mathbb R^5\to\mathbf1^\perp\subset\mathbb R^6$ be a real isometry
with first column $(1,1,1,-1,-1,-1)^{\mathsf T}/\sqrt6$.  For $0<\tau<1$ put
\begin{align}
 t_\tau&=3\tau/2,& D_\tau&=(I-\tau^4R)^{1/2},& U_\tau&=e^{-it_\tau H},
 \nonumber\\
 B_1&=U_\tau D_\tau,& B_\mu&=\tau^2C_\mu\quad(\mu\ge2),&
 K_e(\tau)&=I/\sqrt{18}+\sqrt{2/3}\sum_\mu O_{e\mu}B_\mu.
 \label{sel19:eq:compiler}
\end{align}
The form and locking property of this construction are inherited from
\cite{sel19-Reg}.  The next theorem adds an exact \emph{converse} for its
symmetry and an explicit, dimension-free invertibility margin.

\begin{theorem}[Symmetry-exact action compilation below the phase-alias threshold]
\label{sel19:thm:compiler}
The six maps in \eqref{sel19:eq:compiler} satisfy \eqref{sel19:eq:lock} and
are independent.  For $0<\tau\le1/4$ every coefficient has
\begin{equation}
 s_{\min}(K_e)\ge
 \delta_*:=\frac{\sqrt{255}}{48}-\frac1{\sqrt{18}}-\frac1{24}>0.
 \label{sel19:eq:floor}
\end{equation}
This margin is independent of the spectrum of $H$ and of the total dimension.
If $\Delta_H=\max\operatorname{spec} H-\min\operatorname{spec} H$ and
\begin{equation}
                     t_\tau\Delta_H<2\pi,
 \label{sel19:eq:noalias}
\end{equation}
then the complete fixed-label covariance group inside $G_0$ is exactly
\begin{equation}
 \boxed{G_\tau=\{g\in G_0:[R_0(g),H]=0\}.}
 \label{sel19:eq:compiler-group}
\end{equation}
Thus the compiler does not impose an additional symmetry restriction beyond
the actual generator on this branch.
\end{theorem}
\begin{proof}
The neutral $A$ and its eigenbasis give $[H,R]=0$ and make every $C_\mu,R$
commute with $G_0$.  Moreover $\sum_\mu B_\mu^*B_\mu=I$.  Orthogonality
of $O$ and $O^{\mathsf T}\mathbf1=0$ give the two locking identities.  In
an eigenbasis extending $u_0,\ldots,u_4$, $I,B_1$ are diagonal, whereas the
four $C_\mu$ have distinct off-diagonal entries.  The moduli of $B_1$ on
$u_0$ and $u_1$ are respectively one and $\sqrt{1-\tau^4}$, so $B_1$ is
not scalar.  This proves independence of $I,B_1,C_2,\ldots,C_5$, and the
invertible coefficient change proves independence of the six $K_e$.

Each row of $O$ has squared norm $5/6$, with first entry squared $1/6$.
Since the $C_\mu$ have one common target and orthogonal input rays,
\[
 \left\|\sqrt{2/3}\,\tau^2\sum_{\mu=2}^5O_{e\mu}C_\mu\right\|
                         =2\tau^2/3.
\]
The first two terms of $K_e$ are $I/\sqrt{18}\pm B_1/3$.
The reverse triangle inequality bounds their least singular value below by
$\sqrt{1-\tau^4}/3-1/\sqrt{18}$.  Subtract the displayed correction and
use monotonicity on $[0,1/4]$ to obtain \eqref{sel19:eq:floor}.  For positivity,
$\sqrt{255}>15$ and $1/\sqrt{18}<1/4$ already give
$\delta_*>15/48-1/4-1/24=1/48$.

Commutation with $H$ implies commutation with every compiler operation, since
the helpers are $G_0$-neutral.  Conversely a symmetry of the six CP branches
commutes with their coefficients by \eqref{sel19:eq:phase-lock}, hence with
$B_1$.  It already commutes with $D_\tau$, which is invertible, so it commutes
with $U_\tau$.  Under \eqref{sel19:eq:noalias} the exponential is injective
on the actual spectrum of $H$.  Its spectral projections therefore determine
those of $H$, or equivalently $H$ is a polynomial in $U_\tau$ by finite
interpolation.  Commutation with $U_\tau$ then implies commutation with $H$.
\end{proof}

Only the converse in \eqref{sel19:eq:compiler-group} requires the phase-alias
bound.  It is not a new restriction on the companion's finite-observable
approximation theorem, which expressly allows unbounded generator norms.
Nor is \eqref{sel19:eq:compiler} inferred from a scalar stationary action:
the complete amplitude implementation and its neutral helpers must be part
of the source for the theorem to apply.

\begin{theorem}[Two actual step settings remove the spectral-width assumption]
\label{sel19:thm:two-times}
Use the same $H,C_\mu,R$ at two admitted settings $\tau_1,\tau_2\in(0,1)$, and retain
the setting together with the six outcome labels.  If $\tau_1/\tau_2$ is
irrational, then
\begin{equation}
 \boxed{G_{\tau_1}\cap G_{\tau_2}
           =\{g\in G_0:[R_0(g),H]=0\},}
 \label{sel19:eq:two-times}
\end{equation}
without any bound on $\Delta_H$.  There is no spectrum-independent positive
comparison of the resulting response Gram with the energy-generator Gram on
unbounded spectral families at fixed settings.
\end{theorem}
\begin{proof}
The same argument extracts $U_{\tau_1}$ and $U_{\tau_2}$.  Their joint
spectral character at energy $E$ is $(e^{-it_1E},e^{-it_2E})$.  Equal joint
characters at distinct $E,F$ would require $t_1(E-F),t_2(E-F)\in2\pi\mathbb Z$,
contradicting the irrational ratio.  Their joint spectral projections are
therefore those of $H$, proving \eqref{sel19:eq:two-times}.  For the final
statement a response across an energy gap $\Delta$ has multiplier bounded by
$8$, while its generator squared-commutator has multiplier $\Delta^2$.
Allowing arbitrarily large such gaps rules out a common positive lower
comparison constant.  The exact spectral multiplier is given next.
\end{proof}

These are two genuinely available settings of one action, not two proof-only
choices of time or two samples assumed from one observed setting.  A rational
ratio does not provide the stated all-spectrum guarantee.

\section{A calibrated action-to-source Ward comparison}
\label{sel19:response}

For Hermitian generators $T_a$ of the candidate group define
\[
 (\mathsf W_H)_{ab}=\operatorname{Re}\Tr([T_a,H]^*[T_b,H]).
\]
They annihilate the neutral helper space $A$.  The source response
$\mathsf L_\tau$ is that of Proposition~\ref{sel19:prop:linear}.

\begin{theorem}[Exact response and spectral sinc bounds]
\label{sel19:thm:response}
For the compiler \eqref{sel19:eq:compiler},
\begin{equation}
 \boxed{(\mathsf L_\tau)_{ab}
        =\frac23\operatorname{Re}\Tr([T_a,U_\tau]^*[T_b,U_\tau]).}
 \label{sel19:eq:response}
\end{equation}
Under \eqref{sel19:eq:noalias}, put
$b_\tau=\operatorname{sinc}(t_\tau\Delta_H/2)$, with
$\operatorname{sinc}(0)=1$.  In the order of positive quadratic forms,
\begin{equation}
 \boxed{\frac32\tau^2 b_\tau^2\mathsf W_H
      \preceq\mathsf L_\tau
      \preceq\frac32\tau^2\mathsf W_H.}
 \label{sel19:eq:sinc}
\end{equation}
In particular their ranks and kernels agree.  When the central rank is one,
the sign of the mixed central entry also agrees.  Without the spectral bound,
for $T=\sum_a x_aT_a$ and $H=\sum_E E P_E$ the exact identity is
\begin{equation}
 x^{\mathsf T}\mathsf L_\tau x
 =\frac23\sum_{E,F}4\sin^2\!\left(\frac{t_\tau(E-F)}2\right)
                         \|P_ETP_F\|_{\HS}^2.
 \label{sel19:eq:multiplier}
\end{equation}
\end{theorem}
\begin{proof}
The neutral helper commutators vanish.  Orthogonality of the five coefficient
columns leaves $(2/3)$ times the Gram of $[T_a,U_\tau]D_\tau$.
Since $A$ reduces $H$ and $T_a|_A=0$, $[T_a,U_\tau]$ vanishes on $A$.
Thus multiplication by $D_\tau$ does not change it, proving
\eqref{sel19:eq:response}.  In energy coordinates the $(E,F)$ block is
$(e^{-it_\tau F}-e^{-it_\tau E})P_ETP_F$.  Orthogonality of blocks gives
\eqref{sel19:eq:multiplier}, whereas $\|[T,H]\|_{\HS}^2$ has multiplier
$(E-F)^2$.  On $[0,\pi)$ the positive function $\sin x/x$ is decreasing:
its derivative has numerator $x\cos x-\sin x$, whose derivative is
$-x\sin x\le0$.  Hence every ratio of these squared multipliers lies
between $t_\tau^2b_\tau^2$ and $t_\tau^2$.  This proves
\eqref{sel19:eq:sinc}, since $(2/3)t_\tau^2=3\tau^2/2$.  Equal kernels
follow from the positive lower constant.  Two nonzero rank-one positive
matrices with the same kernel have proportional forms, giving the sign claim.
\end{proof}

At two irrationally related settings the sum of the sine-squared multipliers
is positive at every actual nonzero gap.  This is an exact finite certificate,
not a regulator-uniform lower bound.  A single setting can alias all relevant
gaps and return zero response despite $\mathsf W_H\ne0$.

\begin{corollary}[A source-action determinant handoff]
\label{sel19:cor:action-handoff}
On five or six native copies, let the actual $H$ commute with the nonabelian
action and satisfy $\rank\mathsf W_H^{\rm cent}=1$ and
$(\mathsf W_H^{\rm cent})_{CW}>0$.  A neutral-helper compiler satisfying
\eqref{sel19:eq:noalias}, or its two-setting version in
Theorem~\ref{sel19:thm:two-times}, has exactly the connected determinant
covariance group.  Conversely its linear-Choi test
\eqref{sel19:eq:selector} identifies that same group without prescribing
the central ratio.  No charged matter list or anomaly equation enters this
handoff.
\end{corollary}
\begin{proof}
The coefficient-space bandwidth proof applies also to the Hermitian $H$.
Its nonabelian invariance and central rank-and-sign condition force its
nonzero character support into the primitive $(1,1)$ line in the stated
copy window.  Thus its stabilizer is the connected determinant group.
Theorems~\ref{sel19:thm:compiler} and~\ref{sel19:thm:two-times} transfer it.
Alternatively use the source test of Corollary~\ref{sel19:cor:selector};
Theorem~\ref{sel19:thm:response} identifies its rank and sign below aliasing.
\end{proof}

\section{An exact five-copy source attains every locking requirement}
\label{sel19:witness}

This is an explicit \emph{constructed realization} using the newly supplied
compiler, not a numerical assignment to the historical Store.  Keep the
native representation \eqref{sel19:eq:native} and a fixed defining weak copy.
Choose orthonormal vectors $r_0,\ldots,r_4$ in its already present neutral
five-dimensional sector.  On five copies put
\[
 n=r_0^{\otimes5},\qquad \omega=\Omega_C\otimes\Omega_W,
\]
where $\Omega_C,\Omega_W$ are the normalized alternating vectors already used
in V16.  They satisfy $R(g)n=n$ and $R(g)\omega=\chi(g)\omega$.
Five orthonormal neutral helper vectors, orthogonal to $n,\omega$, are
\[
 u_0=r_1^{\otimes5},\quad u_1=r_2^{\otimes5},\quad
 u_2=r_3^{\otimes5},\quad u_3=r_4^{\otimes5},\quad
 u_4=r_0^{\otimes4}\otimes r_1.
\]
Their provenance may be retained; selecting these vectors does not identify
or erase the original two neutral scalar labels.

Define operators on the full five-copy space, extended by zero outside
$\operatorname{span}\{n,\omega\}$,
\begin{equation}
 P=|n\rangle\langle n|+|\omega\rangle\langle\omega|,\qquad
 S=|n\rangle\langle\omega|+|\omega\rangle\langle n|,\qquad
 H_s=s(P+S)\succeq0\quad(s>0).
 \label{sel19:eq:H}
\end{equation}
Its eigenvalues are $0,2s$, and its symmetry inside $G_0$ is exactly
$\ker\chi$.  The whole action, not just a transition square, is supplied.

\begin{theorem}[Sharp normalized native determinant source]
\label{sel19:thm:witness}
Apply \eqref{sel19:eq:compiler} to $H_s$ and the five neutral helpers.  All
six coefficients are independent, satisfy the two original locking equations,
and are invertible for $\tau\le1/4$.  Whenever $\sin(3\tau s/2)\ne0$,
\begin{equation}
 G_\tau=\ker\chi,\qquad
 \mathsf L_\tau^{\rm cent}
 =\frac43\sin^2(3\tau s/2)
       \begin{pmatrix}9&6\\6&4\end{pmatrix}.
 \label{sel19:eq:witness}
\end{equation}
At the exact choice $\tau=1/8$, $s=4\pi/3$,
\begin{equation}
 U_\tau=I-\frac{1+i}{2}(P+S),\qquad
 D_\tau=I+(\sqrt{4095}/64-1)R,
 \qquad
 \boxed{\mathsf L_\tau^{\rm cent}
      =\frac23\begin{pmatrix}9&6\\6&4\end{pmatrix}.}
 \label{sel19:eq:exact-point}
\end{equation}
Its nonzero central eigenvalue is $26/3$ and its null direction is $(-2,3)$.
Thus five native copies are both necessary and sufficient for an independent
six-outcome identity-calibrated source with this exact group.  Every original
forward word in the displayed realization is invertible.
\end{theorem}
\begin{proof}
The helper subspace reduces $H_s$, which is zero there.  The compiler theorem
proves all normalization, independence and invertibility assertions on the
entire space.  On its two-dimensional active subspace,
\[
 U_\tau=I-P+e^{-it_\tau s}
                         (\cos(t_\tau s)P-i\sin(t_\tau s)S).
\]
The $P$ term and complement are invariant under $G_0$.  A nonzero $S$
coefficient commutes with $R(g)$ exactly when $\chi(g)=1$.  The coefficient
phase-locking and invertible $D_\tau$ identify this with $G_\tau$, even at
nonalias-free times.  In this block $N_C=\diag(0,3)$ and
$N_W=\diag(0,2)$.  The squared commutator Gram with $S$ is twice the outer
product $(3,2)^{\mathsf T}(3,2)$.  Apply \eqref{sel19:eq:response} and
$|e^{-it_\tau s}\sin(t_\tau s)|^2=\sin^2(t_\tau s)$ to obtain
\eqref{sel19:eq:witness}.  The exact choices give $t_\tau s=\pi/4$ and
\eqref{sel19:eq:exact-point}.  The rank-one matrix has trace $26/3$.
The lower bound on copy number is Theorem~\ref{sel19:thm:band}.  Products
of invertible coefficients remain invertible, proving the last assertion.
\end{proof}

No sharp readout or rank-one original history is needed for this result.
All calculations can be performed on the seven-dimensional span of
$n,\omega,u_0,\ldots,u_4$ plus its scalar complement.  The physical native
carrier is still the full $14^5$ representation: a seven-dimensional working
block must not be relabelled as a new seven-dimensional fundamental source.
The original tensor ports and the admitted operations have not been shown to
occur historically by this realization.

The same construction with $\omega=(\Omega_C\otimes\Omega_W)^{\otimes d}$
and $5d$ native factors attains the $\ker\chi^d$ lower bound.  Helpers again
lie in the existing neutral tensor sector.  These are extra component groups,
not a change of the inherited sixfold covering kernel for $\ker\chi$.

\begin{example}[An exact single-setting alias]
\label{sel19:ex:alias}
For the same $H_s$, if $3\tau s=2\pi$, then $U_\tau=I$.  Every compiler
coefficient is $G_0$-invariant, although $H_s$ has only determinant-group
symmetry.  This is an exact finite failure of the converse outside its
spectral qualification, not an arbitrarily small numerical residual.
Two irrationally related admitted settings cannot both alias this nonzero gap.
\end{example}

\section{Action convergence versus finite recorded symmetry}
\label{sel19:counter}

The companions' dynamics must not be used to infer unrecorded implementation
properties.  This has a direct source-level comparison for the same compiler.
Let $\Phi_\tau$ be its forgotten accepted channel and
$\mathcal U_t(X)=e^{-itH}Xe^{itH}$.

\begin{proposition}[A completely bounded action approximation]
\label{sel19:prop:approx}
For every real scalar $c$ and $0<\tau<1$,
\begin{equation}
 \boxed{\|\Phi_\tau-\mathcal U_\tau\|_\diamond
       \le \tau^2\|H-cI\|^2+2\tau^4.}
 \label{sel19:eq:approx}
\end{equation}
Consequently $n$ accepted steps obey the upper bound $n$ times its right-hand
side at time $n\tau$.  Two choices of helper operations with the same $H$
have forgotten-channel difference at most $4\tau^4$, independently of $H$,
and $n$-step difference at most $4n\tau^4$.
\end{proposition}
\begin{proof}
Orthogonality of $O$ gives the exact channel, including its residual term,
\[
 \Phi_\tau(X)=\tfrac13X+\frac23\left[
 U_\tau D_\tau X D_\tau U_\tau^*+
                   \tau^4\sum_{\mu=2}^5 C_\mu XC_\mu^*\right].
\]
The second CP map in brackets has diamond norm $\|R\|=1$.  The sandwich
$D_\tau(\cdot)D_\tau$ differs from the identity in diamond norm by at most
$2\|I-D_\tau\|$, since $\|D_\tau\|\le1$.  Hence $\Phi_\tau$ differs
from $M_\tau=\frac13\mathcal U_0+\frac23\mathcal U_{3\tau/2}$ by at most
$(2/3)[2(1-\sqrt{1-\tau^4})+\tau^4]\le2\tau^4$.

Taylor's formula about time $\tau$ for the unitary channels has a remainder
bounded by $(t-\tau)^2\|\operatorname{ad}_H\|_\diamond^2/2$.
The weights of $M_\tau$ give zero first centered moment and second centered
moment $\tau^2/2$.  Thus
$\|M_\tau-\mathcal U_\tau\|_\diamond
 \le\tau^2\|\operatorname{ad}_H\|_\diamond^2/4
 \le\tau^2\|H-cI\|^2$.
All these bounds hold after every ancillary amplification.  Telescoping
compositions of channels, whose diamond norms are one, gives the $n$-step
bounds.  Comparing both helper choices with the same $M_\tau$ gives the
last assertion.  This is a finite-generator norm estimate; it does not replace
the stronger fixed-observable, large-generator result of \cite{sel19-Reg}.
\end{proof}

\begin{proposition}[The same limiting action allows a symmetry-breaking refinement]
\label{sel19:prop:bad-helper}
On the unchanged full five-copy carrier and with the same $H_s$, replace just
$u_1$ in the helper bank by
\[
 v=e_1\otimes r_0^{\otimes4},
\]
where $e_1$ is a unit colour vector.  Keep $u_0,u_2,u_3,u_4$ and the same
original six labels and normalization.  This is another valid independent,
invertible six-outcome compiler.  For a colour generator $T$ with eigenvalue
one on $e_1$ and zero on the neutral factor,
\begin{equation}
 \boxed{\sum_e\|[T,K_e^{\rm bad}(\tau)]\|_{\HS}^2
                         =\frac23\tau^4>0.}
 \label{sel19:eq:bad-helper}
\end{equation}
Thus it fails even the nonabelian covariance entrance at every finite
$\tau>0$, while its forgotten accepted evolution and that of the neutral
compiler differ by at most $4n\tau^4$ after $n$ steps with the same $H_s$.
\end{proposition}
\begin{proof}
The five helper vectors are still orthonormal zero-energy eigenvectors of
$H_s$.  The inherited locking and independence calculation therefore applies,
as does the explicit singular-value bound.  Their $R$ commutes with the
chosen diagonal colour generator, and $U_\tau$ commutes with it because
$H_s$ is an $\mathrm{SU}(3)$ singlet.  Only
$C_2^{\rm bad}=|u_0\rangle\langle v|$ has a nonzero commutator, namely
$[T,C_2^{\rm bad}]=-C_2^{\rm bad}$.  Orthogonality of $O$ leaves its
squared norm $(2/3)\tau^4$, proving \eqref{sel19:eq:bad-helper}.
The phase-locking result turns that nonzero coefficient defect into failure of
fixed-label CP covariance.  Proposition~\ref{sel19:prop:approx} proves the
comparison of the two unchanged-action evolutions.
\end{proof}

These are different complete sources, not identical old histories with two
interpretations.  Their difference is deliberately retained, rather than
removed by twirling.  Convergence to a stationary or coherent limit cannot
certify exact finite recorded gauge symmetry.  This is why the companion's
explicit neutral-multiplicity construction is substantive and not decorative.

\section{Physical errors, calibration and the next source row}
\label{sel19:status}

If the reconstructed Choi bank satisfies
$\sum_e\|\widehat J_e-J_e\|_1\le\epsilon$, the central entries of the linear
response obey
\begin{equation}
 |\widehat{\mathsf L}_{ij}-\mathsf L_{ij}|\le k^2\epsilon,
 \qquad \|\widehat{\mathsf L}-\mathsf L\|_\op\le2k^2\epsilon.
 \label{sel19:eq:errors}
\end{equation}
Indeed $N_C,N_W$ commute, and their Choi differences have norm at most $k$;
trace duality in \eqref{sel19:eq:L} gives the first inequality, and the
maximum absolute row sum gives the second.  A positive lower mixed endpoint
is therefore certifiable.  A small determinant or a small nonabelian residual
is not an exact rank-one or exact symmetry certificate.  The source identities
or an independently established conservation law must supply those zeros.

At a fixed locked opportunity the accepted coefficients are $\sqrt\vartheta
K_e$ and the no-response coefficient is $\sqrt{1-\vartheta}I$.  The latter
contributes zero to \eqref{sel19:eq:L}, so the physical response is
$\vartheta\mathsf L$, not $\vartheta^2\mathsf L$.  V18's squared-Choi
response instead has the latter scaling.  Tensoring the four inherited
protected record sectors keeps their labels; preparing one such sector gives
the displayed slice.  It does not license deleting random initial sector data.
The calibrated non-Poisson clock of \cite{sel19-Reg} is retained when the
compiler is used in that larger source.  Our accepted-step diamond estimate
is not an estimate of an uncalibrated physical elapsed time.

There is also a cofinal limitation visible within the positive example itself:
\[
 \lambda^+\bigl(\vartheta\mathsf L_\tau^{\rm cent}\bigr)
 =\frac{52\vartheta}{3}\sin^2(3\tau s/2)
                  \sim39\vartheta s^2\tau^2.
\]
Its exact group is constant at all small positive $\tau$, while this raw
response gap vanishes.  Neither exact source symmetry nor a pointwise
positive form is a uniform physical mass gap.  The spectral sinc factor,
operator domain, actual observation and time calibration must remain in any
further estimate, just as the new companions retain their source and density
terms.

\paragraph{Changed direction.}
The useful next object is not another arbitrary five-slot CP channel.  It is
the \emph{actual locked action implementation}: its common memory, its full
$H$, the helper operations and their neutral subspace, and either a verified
spectral window or two admitted noncommensurate step settings.  On five or six
native copies, its linear-Choi central test gives the determinant direction
without a preassigned hypercharge ratio.  If an available three-copy CP port
is the object instead, it belongs to V17's broader category and cannot be
identified with this six-outcome class by discarding the coherent identity
condition.

The comparison above proves a nonempty source-compatible active realization,
not that the historical Store has selected $H_s$ or its tensor port.  The
newly supplied pure-gauge and gauge--Higgs dynamics specify their gauge group
and action before constructing the outcomes.  Their results do not populate
our historical $\mathrm U(3)\times\mathrm U(2)$ source response, choose
charged matter modules, or establish the named type/private and Clifford
incidences.  The scalar shell determinant and the Einstein lapse inverse
cannot stand in for those missing operator data.  Full microscopic
Standard-Model selection and the Einstein--SM continuum are not claimed.

\paragraph{Verification and preservation.}
The complete V18 prefix is retained byte for byte.  The companion exact-check
script verifies the stronger coefficient weight budget, both normalizations,
independence, the seven-dimensional block and its scalar complement, the
linear Choi identity, the exact central response and alias, the nonneutral
helper defect and the event-weight scaling.  It checks the existing V18
regression separately.  The analytic general proofs above, not finite sample
ranks, establish the all-source and all-depth implications.  The four new
companion proofs are audited at the cited interfaces; their absent external
verification scripts are not claimed to have been rerun.  There is no
proof-assistant certification or independent peer review.

\begin{thebibliography}{99}
\raggedright
\bibitem[30]{sel19-Reg}
\emph{Renewal Native Regenerative Stationarity Closure Notes V21}, supplied
cumulative LaTeX.  Used at \src{v12:thm:locked-bank},
\src{v15:thm:locked}, \src{v20:prop:multiplicity}, and
\src{v21:sec:ledger}; these are constructed action/source and scope results,
not a historical Standard-Model selection theorem.
\bibitem[31]{sel19-YM}
\emph{Renewal Yang--Mills Mass-Gap Research Notes V43}, supplied cumulative
LaTeX.  Used at \src{thm:v43-triangular}, \src{thm:v43-free-history}, and
\src{prop:v43-wall}; the surviving density, nonlinear interactions and
transported domains are retained.
\bibitem[32]{sel19-Ein}
\emph{Renewal Exact-Action Einstein Bridge Research Notes V16}, supplied
cumulative LaTeX.  Used at \src{v16:no-soft-flux},
\src{v16:lift-theorem}, and \src{v16:status}; no full inverse or
internal incidence is inferred from its compression estimate.
\bibitem[33]{sel19-Shell}
\emph{Renewal Yang--Mills Shell Mixing Research Notes V16}, supplied
cumulative LaTeX.  Used at \src{thm:v16-assembly}, \src{thm:v16-flat}, and
\src{sec:v16-next}; all factors of its measured scalar coefficient remain
separate from the determinant-character source of this lineage.
\end{thebibliography}
\addtocontents{toc}{\protect\endgroup}
% END OF SOURCE-SELECTION V19 DEVELOPMENT BLOCK.
% Preserve V1--V19 and append subsequent work before the final terminator.


\clearpage
\begin{center}
{\Large\bfseries Source-selection continuation V20}\\[.4em]
{\large Inverting the recorded source: identity participation,\\
neutral implementation tests, and the action-logarithm fibre}
\end{center}
\addtocontents{toc}{\protect\begingroup}
\addtocontents{toc}{\protect\renewcommand{\protect\numberline}{\protect\makebox[2.1em][l]}}

\section{Go upstream from a supplied implementation to the actual marked maps}
\label{sel20:context}

V19 proves a symmetry-exact action compiler, but its positive use starts with
an independently specified Hermitian action, a neutral auxiliary eigenbasis,
and an implementation step.  The present question is inverse: can the original
\emph{separately recorded CP maps} decide whether that implementation class is
present, and reconstruct its data without adding an operation?  This is more
specific than another existence example and does not require choosing a new
Hamiltonian with the desired stabilizer.

The answer separates three levels.  First, an intrinsic identity-participation
test fixes compatible coefficient representatives and, more generally, locks
all symmetry phases even when their identity coefficients are unequal.
Second, a finite inverse test recognizes exactly the neutral-helper compiler
at one step, reconstructing its dimensionless step, helper rays, and unitary
factor.  Third, the unitary factor determines an explicit fibre of Hermitian
logarithms, not a unique physical action without an additional spectral or
multi-setting condition.  Finally, the complete six-direction Choi spectrum
identifies the exact small-step conditioning cost of this inverse problem.

\paragraph{Audit and inherited results.}
The linear Choi commutator form is already proved in
Proposition~\ref{sel19:prop:linear} and, for typed incidence banks, in
\cite{Rigidity}, \src{v3:thm:Choi-Howe}.  We do not reprove it as a new
certificate.  Coefficient phase locking, the determinant bandwidth, the
first-mixed-direction theorem, and forward action compilation are respectively
\ref{sel7:lem:phase-lock}, \ref{sel19:thm:band},
\ref{sel19:cor:selector}, and \ref{sel19:thm:compiler}.  The inverse conditions
below have been checked against those sections and against
\cite{sel19-Reg}, \src{v12:thm:locked-bank} and \src{v15:thm:locked}.
This is a targeted overlap audit, not a priority claim or an independent
verification of every supplied archive.  Hamiltonian identification from
physical process tomography is an established setting \cite{sel20-Wang};
the special inverse problem here retains the original marked implementation
and its identity normalization.

\paragraph{What is and is not reconstructed.}
Every input Choi operator below belongs to an actual outcome on one identified
input/output memory, with its original probability weight.  The formulae are
calculations on those maps.  They do not implement coherent sums of outcomes,
make a Kraus component separately observable, or infer the phase of a different
controlled Stinespring experiment.  Such an experiment can require additional
implementation data \cite{sel20-Abbott}.  Equality of the reconstructed maps
with the supplied maps does, however, preserve all of their original sequential
cylinders exactly.  The scalar variational action reconstructed from a reward
process is not identified with a Hermitian logarithm solely because both are
called actions.

\section{An intrinsic identity anchor from the separate Choi operators}
\label{sel20:anchor}

Let $\mathcal J_e$ be a normalized instrument on $M_N$, with nonzero
rank-one Choi operators $J_e$ and fixed labels $1\le e\le s$.  Use the
output-first vectorization already fixed in V19, and put
\begin{equation}
 i=|I_N\rangle\!\rangle,\qquad S=\sum_eJ_e,\qquad \Pi=SS^\dagger.
 \label{sel20:eq:sum}
\end{equation}
Assume $\rank S=s$.  Equivalently, any representatives
$J_e=|L_e\rangle\!\rangle\langle\!\langle L_e|$ are linearly independent.
For each label define the intrinsic numbers and vectors
\begin{equation}
 a_e=\langle i,S^\dagger J_eS^\dagger i\rangle\ge0,
 \qquad x_e=J_eS^\dagger i,\qquad
 r_I=\|(I-\Pi)i\|^2.
 \label{sel20:eq:anchor-data}
\end{equation}
No between-label coherent Choi block is part of these data.

\begin{theorem}[Identity participation and unique balanced calibration]
\label{sel20:thm:anchor}
Under the preceding hypotheses the following statements hold.
\begin{enumerate}[label=\textup{(\roman*)},leftmargin=2.8em]
\item $r_I=0$ exactly when $I_N$ lies in the coefficient span.  In that
case the unique expansion $I_N=\sum_e\alpha_eL_e$ has
$|\alpha_e|^2=a_e$ and $x_e=\alpha_e|L_e\rangle\!\rangle$.
Thus $a_e>0$ means that label $e$ actually participates in that expansion.
\item For a fixed positive real $c$, representatives $K_e$ of the
\emph{same} marked CP maps with $\sum_eK_e=cI_N$ exist exactly when
\begin{equation}
 \boxed{r_I=0,\qquad c^2a_e=1\quad\text{for every }e.}
 \label{sel20:eq:balanced}
\end{equation}
When they exist they are unique and satisfy
\begin{equation}
 \boxed{|K_e\rangle\!\rangle=cJ_eS^\dagger i.}
 \label{sel20:eq:canonical-K}
\end{equation}
In particular, the six-outcome calibration $c=\sqrt2$ is tested by
$a_1=\cdots=a_6=1/2$.
\item If $c$ has not been prescribed, a positive common coefficient exists
exactly when $r_I=0$ and all $a_e$ equal the same positive number $a$.
It is then determined as $c=a^{-1/2}$.
\end{enumerate}
\end{theorem}
\begin{proof}
Choose representatives temporarily and let $W:\C^s\to M_N$ have columns
$|L_e\rangle\!\rangle$.  It is injective and $S=WW^*$.
If $G=W^*W$, then $S^\dagger W=WG^{-1}$ and
$W^*S^\dagger W=I_s$; both identities follow by applying the operators
on $\Ran W$ and its orthogonal complement.  Thus $\Pi$ is the projection
onto $\Ran W$, which proves the first equivalence.  When $i=W\alpha$,
\[
 \langle\!\langle L_e|S^\dagger i\rangle=\alpha_e,
 \qquad J_eS^\dagger i=\alpha_e|L_e\rangle\!\rangle.
\]
This proves (i), including its phase independence.

Representatives of the same nonzero rank-one Choi ray have the form
$K_e=z_eL_e$ with $|z_e|=1$.  Their identity sum is equivalent, by
independence, to $z_e=c\alpha_e$ for every $e$.  These phases exist exactly
under \eqref{sel20:eq:balanced}, and then give
\eqref{sel20:eq:canonical-K} uniquely.  Their Choi matrices are unchanged,
so their trace-preserving normalization is unchanged as well.
Finally $c^2a_e=1$ determines the common number and $c$, proving (iii).
All formulas have been expressed using $J_e,S,i$, eliminating the temporary
representatives.
\end{proof}

Balanced calibration is stronger than necessary for symmetry phase locking.
This distinction removes a numerical normalization from the selection
hypotheses without changing the original maps.

\begin{theorem}[Full identity participation locks all recorded symmetry phases]
\label{sel20:thm:phase-lock}
Suppose only
\begin{equation}
                 r_I=0,\qquad a_e>0\quad(1\le e\le s).
 \label{sel20:eq:full-participation}
\end{equation}
For every unitary representation $R(g)$ on the same memory,
\begin{equation}
 \boxed{\operatorname{Cov}_{G}(\mathcal J_e:e)
       =\{g:[R(g),L_e]=0\ \text{for all }e\}.}
 \label{sel20:eq:phase-lock}
\end{equation}
The right side is independent of the individual representative phases.
Neither equal $a_e$, a prescribed coherent sum, nor execution of an
identity-interference experiment is required.
\end{theorem}
\begin{proof}
Fixed-label covariance gives $R(g)L_eR(g)^*=z_e(g)L_e$ by equality of
rank-one Choi rays.  Conjugate the unique identity expansion from
Theorem~\ref{sel20:thm:anchor}.  Its difference from the original expansion
is $\sum_e\alpha_e(z_e(g)-1)L_e=0$.  Independence and
$\alpha_e\ne0$ imply every $z_e(g)=1$.  The converse is immediate.
This is the inherited phase-locking argument applied to coefficients
reconstructed intrinsically, rather than to a supplied equal-coefficient
identity.
\end{proof}

\begin{example}[The identity ray alone does not suffice]
\label{sel20:ex:missing-participation}
On a qubit take $L_0=I/\sqrt2$, $L_1=E_{01}/\sqrt2$ and
$L_2=E_{10}/\sqrt2$.  Their squared input effects sum to $I$ and they are
independent.  Here $r_I=0$, but $(a_0,a_1,a_2)=(2,0,0)$.
Every diagonal unitary fixes all three CP branches, while it rotates the
last two coefficients by opposite phases.  Hence the conclusion
\eqref{sel20:eq:phase-lock} fails without participation of those labels.
This is not a counterexample to the balanced or full-participation theorem.
\end{example}

\begin{corollary}[A determinant selector without a prescribed coherent sum]
\label{sel20:cor:selection}
Let the independently verified memory be an invariant subspace of at most six
copies of the native representation \eqref{sel19:eq:native}, with any
retained neutral multiplicity.  Suppose the actual efficient, independent,
marked instrument satisfies \eqref{sel20:eq:full-participation} and is
fixed-label covariant under $\mathrm{SU}(3)\times\mathrm{SU}(2)$.
Compute the linear Choi response $\mathsf L$ of
Proposition~\ref{sel19:prop:linear} for $N_C,N_W$.  Then
\begin{equation}
 \boxed{\det\mathsf L=0,\quad\mathsf L_{CW}>0
       \quad\Longleftrightarrow\quad
       G_{\mathcal J}=S(\mathrm U(3)\times\mathrm U(2)).}
 \label{sel20:eq:selection}
\end{equation}
The positive branch has at least five native copies and null direction
$\mathbb R(-2,3)$.  The result needs neither the number six of outcome
labels nor the equal-sum value $\sqrt2$.
\end{corollary}
\begin{proof}
Theorem~\ref{sel20:thm:phase-lock} turns covariance into coefficient
commutation.  Each coefficient is therefore nonabelian invariant.  The proof
of Theorem~\ref{sel19:thm:band} now applies with no change: its bandwidth
uses only the native representation.  The proof of
Corollary~\ref{sel19:cor:selector} uses that bandwidth, the positive
coefficient-mode sum, and full phase locking, not the magnitude of the
identity coefficients.  It gives \eqref{sel20:eq:selection}, the copy lower
bound and the kernel.  The linear Choi expression for the coefficient form
is the inherited identity; the new participation test licenses its
interpretation as the actual covariance kernel in this larger source class.
\end{proof}

The input is still the \emph{output-sensitive} marked Choi panel, not scalar
outcome counts.  Also, an independently retained idle identity outcome can
make the enlarged Kraus list dependent.  It must not be appended to an
independent accepted list and then silently treated as one more independent
ray.  One may test the actual accepted slice and retain its original
acceptance weight separately, as in V19.

\section{A finite inverse theorem for the neutral-helper compiler}
\label{sel20:inverse}

Now require six labels and \eqref{sel20:eq:balanced} with $c=\sqrt2$.
The preceding theorem fixes $K_e$ from the original maps.  Keep the real
record isometry $O$ declared by the implementation class in V19; it is not
fitted to a desired group.  Its columns are orthonormal, orthogonal to
$\mathbf1$, and its first column is
$(1,1,1,-1,-1,-1)^{\mathsf T}/\sqrt6$.  Calculate
\begin{equation}
 B_\mu=\sqrt{3/2}\sum_eO_{e\mu}K_e,\qquad 1\le\mu\le5.
 \label{sel20:eq:B}
\end{equation}
The inherited normal-form identity gives
\begin{equation}
 \sum_\mu B_\mu^*B_\mu=I,\qquad
 K_e=I/\sqrt{18}+\sqrt{2/3}\sum_\mu O_{e\mu}B_\mu.
 \label{sel20:eq:normal-form}
\end{equation}
This follows also by projecting the six coefficient columns onto the identity
column and $\mathbf1^\perp$; it is not a new normalization construction.

For the four tail columns calculate
\begin{equation}
 q=\frac14\sum_{\mu=2}^5\|B_\mu\|_{\HS}^2,
 \qquad Z=\sum_{\mu=2}^5B_\mu B_\mu^*.
 \label{sel20:eq:tail}
\end{equation}
The following tests are evaluated in order, so undefined quantities are never
used after a failed support test.
\begin{enumerate}[label=\textup{(I\arabic*)},leftmargin=3em]
\item $0<q<1$ and
$\Tr(B_\mu^*B_\nu)=q\delta_{\mu\nu}$ for $2\le\mu,\nu\le5$.
\item $\rank Z=1$.  Put $p_0=Z/(4q)$ and
$p_j=B_{j+1}^*B_{j+1}/q$ for $1\le j\le4$.
\item With $R=\sum_{j=1}^4p_j$, one has $p_0R=0$.
\item Put $D=(I-qR)^{1/2}$ and $U=B_1D^{-1}$.
For $0\le j\le4$, $[U,p_j]=0$.
\item The projection $A=p_0+R$ has neutral range under the full comparison
group: $(R_0(g)-I)A=0$ for all $g\in G_0$.
\end{enumerate}
For the connected $G_0$ a finite Hermitian Lie generating set $T_a$ tests the
last condition by $\sum_a\|T_aA\|_{\HS}^2=0$.  This is \emph{trivial action}
on $A$, stronger than merely $[T_a,A]=0$.

\begin{theorem}[Necessary and sufficient inverse compiler test]
\label{sel20:thm:inverse}
The original six marked maps have a neutral-helper representation of the
precise form \eqref{sel19:eq:compiler}, with the fixed $O$, some
$0<\tau<1$, and a finite Hermitian generator, if and only if tests
\textup{(I1)}--\textup{(I5)} pass.  On this branch
\begin{equation}
 \boxed{\tau=q^{1/4},\qquad C_\mu=B_\mu/\sqrt q\ (\mu\ge2),
       \qquad U=B_1(I-qR)^{-1/2}}
 \label{sel20:eq:inverse}
\end{equation}
are reconstructed, rather than independently supplied.  The projections
$p_0,\ldots,p_4$ are mutually orthogonal of rank one, so $A$ has rank five.
Every Hermitian logarithm $H$ of $U$ at time $t=3\tau/2$ that commutes with
these five projections gives exactly the old maps.  Such a logarithm exists.
The sole vector freedom in reconstructing $C_\mu=|u_0\rangle\langle
u_{\mu-1}|$ is a common phase on the five reconstructed vectors.
\end{theorem}
\begin{proof}
Necessity follows by substituting V19's exact compiler.  Its four tail
operators are $\tau^2|u_0\rangle\langle u_{\mu-1}|$, giving (I1)--(I3)
and $q=\tau^4$.  Its polar unitary is $e^{-itH}$, giving (I4) because the
helper vectors are eigenvectors of $H$; their stated neutrality is (I5).

Conversely choose a unit vector $u_0$ spanning $\Ran Z$.  Positivity of the
sum defining $Z$ implies $\Ran B_\mu\subseteq\C u_0$ for every tail
column: for $v\perp u_0$, $\sum_\mu\|B_\mu^*v\|^2=0$.
Write $B_\mu=\sqrt q\,|u_0\rangle\langle u_{\mu-1}|$.
The tail Gram in (I1) makes $u_1,\ldots,u_4$ orthonormal.  Thus $p_0$ and
all $p_j$ are rank-one projections; (I3) makes all five orthogonal, and $R$
is a rank-four projection.  Formula \eqref{sel20:eq:normal-form} gives
$B_1^*B_1=I-qR=D^2\succ0$.  In finite equal input/output dimensions,
$U=B_1D^{-1}$ is unitary.

Choose any spectral arguments $\phi_z$ with $z=e^{-i\phi_z}$ for the finitely
many eigenvalues of $U=\sum_z zP_z$, and put
$H_0=t^{-1}\sum_z\phi_zP_z$.  It is Hermitian, $e^{-itH_0}=U$, and (I4)
implies that every $p_j$ commutes with $H_0$.  Being rank one, each of those
ranges is an eigenspace ray.  Condition (I5) supplies precisely the neutral
reducing helper subspace required by V19.  Substituting $\tau=q^{1/4}$ and
these data into that compiler gives $B_1$ and all four original tail columns,
hence every $K_e$ by \eqref{sel20:eq:normal-form}.  Their marked CP maps
are unchanged.  Any other logarithm commuting with the five $p_j$ has the
same property.  Finally $u_{\mu-1}=C_\mu^*u_0$; changing the phase of $u_0$
changes all five vectors by that same phase and leaves every $C_\mu$ fixed.
\end{proof}

\begin{remark}[Meaning of implementation recognition]
The theorem recognizes equality with a specified \emph{recorded instrument
class}.  It does not show that the laboratory or microscopic source executes
$C_\mu$, $D$, or $U$ as individual gates.  None is inserted in an original
history.  Equality of the full marked maps is enough to apply the symmetry
conclusions of that class, and preserves all original sequential probabilities
and conditioned states for every input and untouched ancillary reference.
The inferred $\tau$ is the dimensionless parameter of this exact normal form;
physical elapsed time and the acceptance clock require their original
calibration.  Changing the given $O$ to make a failed test pass would change
which implementation class is being tested.
\end{remark}

\begin{corollary}[The action-implementation premise can be replaced by the inverse panel]
\label{sel20:cor:inverse-selection}
If tests \textup{(I1)}--\textup{(I5)} pass, the fixed-label symmetry inside
$G_0$ is exactly $\{g:[R_0(g),U]=0\}$.  For a branch logarithm scalar on
each distinct $U$-eigenspace it is also the stabilizer of that logarithm.
On a verified five- or six-copy native memory, nonabelian covariance and the
linear central test \eqref{sel20:eq:selection} consequently identify the
connected determinant group without a separately supplied $H$, helper
basis, or $\tau$.
\end{corollary}
\begin{proof}
The neutral helper operators and $D$ commute with every $R_0(g)$.  The
identity-participation theorem and \eqref{sel20:eq:normal-form} show that
source covariance is equivalent to commutation with all $B_\mu$, hence to
commutation with $B_1=UD$, and hence with $U$.  Distinct spectral values of
$U$ have distinct chosen real arguments modulo $2\pi$, so such a logarithm
has exactly the same spectral projections.  The last assertion is
Corollary~\ref{sel20:cor:selection}.
\end{proof}

\section{The exact action ambiguity and a record-sensitive energy origin}
\label{sel20:logs}

The inverse theorem does not select the historical Hamiltonian from one
sampling setting.  The remaining ambiguity is explicitly reconstructible.

\begin{theorem}[All compatible Hermitian action generators]
\label{sel20:thm:log-fibre}
Let the inverse tests pass and write $U=\sum_z zP_z$, $z=e^{-i\phi_z}$,
with one fixed real argument $\phi_z$ for each distinct eigenvalue.
All compatible Hermitian generators, and only those, are
\begin{equation}
 \boxed{H=\bigoplus_z\left(\frac{\phi_z}{t}I_{P_z}
                      +\frac{2\pi}{t}N_z\right),\qquad t=3\tau/2,}
 \label{sel20:eq:log-fibre}
\end{equation}
where $N_z=N_z^*$ has integer spectrum and commutes with the restrictions of
$p_0,\ldots,p_4$ to $P_z\mathcal H$.
Every such $H$ has stabilizer contained in the recorded source symmetry.
Equality holds for scalar $N_z$ on every $P_z$; non-scalar choices can make
the inclusion strict while leaving \emph{all original marked maps} unchanged.
If an independently calibrated energy interval of width less than $2\pi/t$
is given, there is at most one compatible $H$ with spectrum in that interval.
\end{theorem}
\begin{proof}
A Hermitian $H$ commutes with its exponential and hence with each $P_z$.
On $P_z$ all its eigenvalues solve $e^{-itE}=z$, which is exactly
$E=(\phi_z+2\pi n)/t$, $n\in\mathbb Z$.  The finite spectral theorem gives
\eqref{sel20:eq:log-fibre}.  Commutation with the helper projections is
equivalent to the stated block condition and makes them eigenvector rays.
These observations prove necessity and sufficiency.

If $R_0(g)$ commutes with $H$ it commutes with $U$, proving containment.
If every $N_z$ is scalar, the real eigenvalues on different $P_z$ cannot
coincide, since their exponentials differ.  Thus the spectral projections,
and therefore stabilizers, coincide.  An interval of width less than one
period contains at most one real preimage of each $z$; it forces scalar
$N_z$ with a uniquely fixed integer when a solution exists.  This proves
uniqueness in the independently specified window.  It does not assert that
a window chosen after reconstruction identifies a physical energy lift.
\end{proof}

\begin{example}[An invisible charged action on the unchanged source]
\label{sel20:ex:charged-alias}
Keep V19's exact five-copy source with $H_s$, $\tau=1/8$, and its neutral
helpers.  The native colour vector $v=e_1\otimes r_0^{\otimes4}$ is
orthogonal to the determinant active block and all helpers.  Put
\[
              H'=H_s+\frac{2\pi}{t}|v\rangle\langle v|,
              \qquad t=3\tau/2.
\]
Both generators are positive, have the same helper eigenvectors, and give
\emph{identical} $U$, $B_\mu$, and six marked CP maps at this setting.
Every finite original recorded history is identical, including its action on
inputs entangled with a reference.  But $H_s$ commutes with the whole
nonabelian group, whereas $H'$ fails to commute with a colour generator
mixing $e_1$ and $e_2$.  This is the non-scalar integer freedom on a degenerate
$U$-eigenspace.  The source at a second physical step may distinguish them;
longer histories of the unchanged first setting do not.
\end{example}

\begin{proposition}[A constant energy shift is visible in the marked compiler]
\label{sel20:prop:energy-origin}
Keep fixed $\tau$, $O$ and the helper operators, and replace $H$ by $H+cI$.
Their forgotten accepted channels are identical for every real $c$.
Their complete marked instruments are identical if and only if
\begin{equation}
                              tc\in2\pi\mathbb Z.
 \label{sel20:eq:energy-period}
\end{equation}
For V19's zero-energy helper input $u_0$, shifting by $c=\pi/t$ changes the
six-label probability distribution by total variation exactly $2\sqrt2/3$,
although the forgotten output channel is unchanged on every input.
\end{proposition}
\begin{proof}
The shift multiplies $U$ and $B_1$ by $e^{-itc}$, while the tail columns and
$D$ remain unchanged.  In the forgotten channel
\eqref{sel19:eq:approx}'s proof, $U$ appears only in a conjugation, so its
common phase cancels.  If the marked maps agree, their balanced canonical
representatives agree by Theorem~\ref{sel20:thm:anchor}.  Formula
\eqref{sel20:eq:B} then gives the same $B_1$.  Its invertibility forces
$e^{-itc}=1$.  The converse follows by substitution.

On the helper $u_0$ the tail amplitudes vanish and $D u_0=u_0$.
Write $s_e=1$ for the first three labels and $-1$ for the last three.
If $Hu_0=E_0u_0$, the original probability is
\begin{equation}
              p_e(E_0)=\frac16+s_e\frac{\sqrt2}{9}\cos(tE_0).
 \label{sel20:eq:helper-probability}
\end{equation}
For $E_0=0$ and $t(E_0+c)=\pi$ each of the six differences has magnitude
$2\sqrt2/9$.  Half their sum gives $2\sqrt2/3$.
\end{proof}

This is not a claim that an arbitrary classical action constant has been
measured.  It specifies the energy-origin information carried by this
\emph{particular marked implementation}.  The identity amplitude in every
$K_e$ makes a phase that is invisible in the forgotten channel visible in
the retained label.  It is therefore unsafe to identify the complete compiler
from its limiting or even its exact forgotten dynamics alone.

\section{The inverse chart has an exact small-step conditioning spectrum}
\label{sel20:conditioning}

The previous tests are exact finite statements.  Their inverse does not have
a regulator-independent condition number.  This cost is visible in the
original source rather than in an unrelated counterexample.

For an actual compiler on an $N$-dimensional memory let
$\mathcal W_\tau:\C^6\to M_N$ have columns
$|K_e(\tau)\rangle\!\rangle$, and let
$\Gamma_\tau=\mathcal W_\tau^*\mathcal W_\tau$.
The nonzero spectrum of the full summed Choi operator $S_\tau$ is the same
as that of $\Gamma_\tau$; in particular these eigenvalues do not require a
between-label phase convention.

\begin{theorem}[The complete six-dimensional source Gram]
\label{sel20:thm:gram}
In the orthonormal record basis $(\mathbf1/\sqrt6,O_1,\ldots,O_5)$,
\begin{equation}
 \boxed{\Gamma_\tau\simeq
 \begin{pmatrix}
 N/3&\frac{\sqrt2}{3}\Tr B_1\\
 \frac{\sqrt2}{3}\overline{\Tr B_1}&\frac23(N-4\tau^4)
 \end{pmatrix}
 \oplus\frac23\tau^4 I_4.}
 \label{sel20:eq:gram}
\end{equation}
Here the displayed two off-diagonal entries follow the convention that the
first-row inner product is linear in its second argument.
Its two-by-two determinant is
\begin{equation}
 d_2(\tau)=\frac29\{N(N-4\tau^4)-|\Tr B_1|^2\}>0.
 \label{sel20:eq:det2}
\end{equation}
For fixed non-scalar $H$, put
$V_H=N\Tr H^2-(\Tr H)^2>0$.  As $\tau\downarrow0$ the other two
eigenvalues are $N+O(\tau^2)$ and
$\frac{V_H}{2N}\tau^2+O(\tau^4)$, while the four displayed small
eigenvalues are exactly $2\tau^4/3$.  Consequently
\begin{equation}
 \boxed{\det\Gamma_\tau=\frac8{81}V_H\tau^{18}+O(\tau^{20}),
 \qquad \|\mathcal W_\tau^\dagger\|=\sqrt{3/2}\,\tau^{-2}}
 \label{sel20:eq:nonscalar-conditioning}
\end{equation}
for all sufficiently small positive $\tau$ in the norm equality.
For a scalar $H$, instead,
\begin{align}
 d_2(\tau)&=\frac{8(N-4)}9(1-\sqrt{1-\tau^4})^2,\nonumber\\
 \lambda_{\min}(\Gamma_\tau)
    &=\frac{2(N-4)}{9N}\tau^8+O(\tau^{12}),\label{sel20:eq:scalar-conditioning}\\
 \det\Gamma_\tau
    &=\frac{32(N-4)}{729}\tau^{24}+O(\tau^{28}).\nonumber
\end{align}
Thus the full coefficient inverse grows as $\tau^{-4}$ on that branch.
\end{theorem}
\begin{proof}
The identity record column is mapped to $I/\sqrt3$, the next to
$\sqrt{2/3}B_1$, and the last four to
$\sqrt{2/3}\tau^2C_\mu$.  In the common eigenbasis of $H$ and the helper
projections, each $C_\mu$ has a strictly off-diagonal entry.  It is therefore
Hilbert--Schmidt orthogonal to both $I$ and the diagonal $B_1$.  The $C_\mu$
are orthonormal.  Also $\|B_1\|_{\HS}^2=N-4\tau^4$.
These statements prove \eqref{sel20:eq:gram} and its determinant.
Strict positivity follows because $B_1$ is not scalar: its moduli are one
on $u_0$ and $\sqrt{1-\tau^4}$ on the four other helper rays.

For fixed $H$, in its eigenbasis
$B_1=e^{-itH}(I-\tau^4R)^{1/2}$, with $t=3\tau/2$.
Expanding the finite sum of exponentials gives
\[
 N(N-4\tau^4)-|\Tr B_1|^2=t^2V_H+O(\tau^4).
\]
There is no cubic term: the squared trace is a sum of cosines of energy
differences with real even-in-$\tau$ amplitudes.  Hence
$d_2(\tau)=V_H\tau^2/2+O(\tau^4)$.  The trace of the two-by-two block is
$N-8\tau^4/3$, which gives its stated eigenvalue asymptotics.
For sufficiently small $\tau$ its smaller eigenvalue is larger than
$2\tau^4/3$.  Multiplication by the four tail eigenvalues gives
\eqref{sel20:eq:nonscalar-conditioning}; the pseudoinverse norm is the
reciprocal square root of the least Gram eigenvalue.

When $H$ is scalar, set $d=\sqrt{1-\tau^4}$.  Its irrelevant common phase
has modulus one and $|\Tr B_1|=N-4+4d$.  Direct expansion gives
$N(N-4\tau^4)-(N-4+4d)^2=4(N-4)(1-d)^2$.
Use $1-d=\tau^4/2+O(\tau^8)$, the same two-by-two trace, and the four
exact tail eigenvalues to obtain \eqref{sel20:eq:scalar-conditioning}.
The source has $N\ge5$, so the leading constant is positive.
\end{proof}

These are conditioning statements for the full coefficient synthesis, not
physical relaxation rates or mass gaps.  They show why inversion of a
numerically almost rank-one summed Choi packet needs explicit accuracy.
A specially chosen right-hand side can be better conditioned than the
worst-case inverse, so the theorem is not a claim that every scalar statistic
has that error amplification.  The helper reconstruction itself contains the
explicit division $C_\mu=B_\mu/\tau^2$.

\begin{proposition}[Certified coefficient errors propagate with their actual inverse scale]
\label{sel20:prop:errors}
If the calibrated representatives have certified errors
$\sum_e\|\widehat K_e-K_e\|_{\HS}^2\le\epsilon^2$, then
\begin{equation}
 \left(\sum_\mu\|\widehat B_\mu-B_\mu\|_{\HS}^2\right)^{1/2}
          \le\sqrt{3/2}\,\epsilon.
 \label{sel20:eq:B-error}
\end{equation}
At an independently fixed $\tau>0$, the analogous tail-helper error is at
most $\sqrt{3/2}\epsilon/\tau^2$.  More generally, for fixed synthesis
representatives with $\sigma_{\min}(\mathcal W)\ge s_*>0$, an estimate
$\widehat{\mathcal W}=\mathcal W+E$ with $\|E\|\le\delta<s_*$ and a
consistent target $b=\mathcal W\alpha$ with $\|\widehat b-b\|\le\eta$
satisfy
\begin{equation}
 \boxed{\|\widehat{\mathcal W}^{\dagger}\widehat b-\alpha\|
       \le\frac{\eta+\delta\|\alpha\|}{s_*-\delta}.}
 \label{sel20:eq:inverse-error}
\end{equation}
\end{proposition}
\begin{proof}
Orthogonal projection by the columns of $O$ is contractive, which proves
\eqref{sel20:eq:B-error}; division by $\tau^2$ gives the next assertion.
The variational characterization of the least singular value gives
$\sigma_{\min}(\widehat{\mathcal W})\ge s_*-\delta>0$.
Since $\widehat{\mathcal W}^{\dagger}\widehat{\mathcal W}=I$, the
coordinate error equals
$\widehat{\mathcal W}^{\dagger}(\widehat b-b-E\alpha)$.
Taking norms proves \eqref{sel20:eq:inverse-error}.
\end{proof}

This proposition starts with certified representatives or coordinates; it does
not silently turn a noisy positive Choi estimate into an exact rank-one ray.
Rank, neutrality, and commutation zeros in the inverse test remain exact
hypotheses.  A nonzero certified residual excludes the corresponding class;
small residuals alone do not prove membership.  The raw accepted Choi sum
is $\vartheta S$, whose nonzero eigenvalues acquire that same factor.
Neither the old event weights nor their conditioning are normalized away.

\section{Apply the inverse to the unchanged source and to its existing obstruction}
\label{sel20:evaluation}

\begin{proposition}[Exact recovered data for the V19 benchmark]
\label{sel20:prop:benchmark}
Apply the inverse to the unchanged V19 maps at $\tau=1/8$, $s=4\pi/3$.
Their data satisfy
\begin{equation}
 r_I=0,\qquad a_e=\frac12\ (1\le e\le6),\qquad
 q=\frac1{4096},\qquad
 \rank A=5.
 \label{sel20:eq:benchmark-anchor}
\end{equation}
All helper directions are neutral, and the recovered unitary is precisely
\eqref{sel19:eq:exact-point}.  On the independently declared energy interval
$[0,10]$ the unique compatible generator is $H_s$.
The four exact tail eigenvalues of the full summed Choi operator are
$1/6144$.  Writing $N=14^5$, the remaining two are those of
\eqref{sel20:eq:gram} with
\begin{equation}
              \Tr B_1=N-5+\frac{\sqrt{4095}}{16}-i,
              \qquad \|B_1\|_{\HS}^2=N-\frac1{1024}.
 \label{sel20:eq:benchmark-trace}
\end{equation}
The recovered source central response is the unchanged
$\frac23\left(\begin{smallmatrix}9&6\\6&4\end{smallmatrix}\right)$.
\end{proposition}
\begin{proof}
V19's independent representatives have the exact identity coefficients
$1/\sqrt2$.  Theorem~\ref{sel20:thm:anchor} gives the first two values
without choosing new phases.  The tail columns in that source are
$\tau^2C_\mu$, giving $q=\tau^4$ and the five original helper rays.
The inverse polar calculation therefore gives its original $U$.
Its two unitary eigenvalues are $1,-i$, so with $t=3/16$ their only energy
lifts in $[0,10]$ are $0,8\pi/3$.  Indeed $8\pi/3<10<32\pi/3=2\pi/t$.
The rank-one $-i$ eigenspace is the positive eigenline of $H_s$.
This proves the assertion about that \emph{declared} interval, not its
historical availability.

The determinant block has unitary trace $1-i$, the undamped helper contributes
one, its four input helpers contribute $4\sqrt{4095}/64$, and the scalar
complement contributes $N-7$.  Summing proves
\eqref{sel20:eq:benchmark-trace}.  The tail eigenvalues are
$(2/3)\tau^4=1/6144$.  Substituting the recovered matrices in the linear
source response returns the exact V19 value.  All these identities hold on
the full native representation, even though they can be checked using its
seven-dimensional working block and the retained scalar complement.
\end{proof}

The source is the earlier active benchmark, not an identified historical
Store.  Reconstructing its own data is a test of the inverse, not additional
historical evidence.

\begin{proposition}[The old charged-helper counterexample is rejected at unit strength]
\label{sel20:prop:bad-helper}
Keep exactly the bad-helper instrument of
Proposition~\ref{sel19:prop:bad-helper}.  Its identity anchor, recovered
$\tau$, tail Gram, common range, orthogonal rays and polar compatibility all
pass.  Neutrality fails: for its stated colour generator $T$,
\begin{equation}
                       \boxed{\|TA\|_{\HS}^2=1.}
 \label{sel20:eq:bad-neutral}
\end{equation}
This defect is independent of $\tau$, although the corresponding original
source response is $2\tau^4/3$.
\end{proposition}
\begin{proof}
The reconstruction returns the same four helper input rays as the actual
source.  Exactly one is the colour ray $v$ on which $T$ has eigenvalue one;
the other four rays of $A$ are neutral.  Thus $TA=|v\rangle\langle v|$,
giving \eqref{sel20:eq:bad-neutral}.  All other tests are the same
orthogonality and eigenvector statements already checked in V19.
\end{proof}

\begin{example}[Full participation is strictly weaker than balanced locking]
\label{sel20:ex:unbalanced}
For comparison only, mix the first two coefficients of the positive V19 bank
by the real orthogonal matrix
$\left(\begin{smallmatrix}3/5&4/5\\-4/5&3/5\end{smallmatrix}\right)$,
leaving the other four unchanged and retaining six output labels.
This creates a different marked instrument with the same forgotten channel.
Its coefficients remain independent and its identity participation numbers are
\begin{equation}
                 \left(\frac{49}{50},\frac1{50},
                              \frac12,\frac12,\frac12,\frac12\right).
 \label{sel20:eq:unbalanced}
\end{equation}
Every number is positive, but there is no common positive $c$ for a balanced
identity sum.  Nevertheless its coefficient span, linear commutator response,
and phase-locked covariance group are unchanged.  Thus
Corollary~\ref{sel20:cor:selection} applies beyond the exactly calibrated
compiler.  This algebraic comparison is not a permitted reinterpretation or
classical relabelling of the old source: coherent Kraus mixing changes its
resolved maps.
\end{example}

\section{What has been removed, and what remains a source input}
\label{sel20:status}

The positive advance has two parts.  First, one can certify phase rigidity of
a complete efficient source from its separate quantum maps, without supplying
a coherent coefficient sum.  Full identity participation, not equal identity
weights, is the essential condition for this argument.  Second, when testing
the companion's exact compiler, the generator need no longer be chosen first:
the marked Choi panel reconstructs the helper support, dimensionless step and
unitary implementation, or returns a failed finite condition.  The returned
spectral fibre records exactly what is left of the action ambiguity.

The distinction between these two results matters.  A source can pass the
more general symmetry selector but fail the particular compiler test, as the
unbalanced comparison shows.  Failure of one inverse implementation class is
not a no-go theorem for every source capable of the same determinant group.
Conversely, passing a compiler membership test does not force its central
response to have rank one or the desired sign.  Those are still evaluated
rather than presumed.  A reconstructed branch logarithm is not an independent
proof of the physical variational action or of a unique continuous-time law.

The next historical input is therefore the \emph{same-cut marked CP panel}
on an independently verified native memory, including its original record
partition and calibrations.  It is not an assumed auxiliary eigenbasis.  This
panel may already be part of a supplied tomographically complete physical
process table; the theorem does not invent that table or new tensor ports.
The manuscript's requirement that mixed maps occur on one actual joint carrier
is retained \cite{Duality}.  No historical five-copy panel has been populated
here, and no Standard-Model particle list, fermionic grading, determinant
matter incidence, or Einstein--SM continuum is asserted from these inverse
identities.

\paragraph{Verification and preservation.}
The complete V19 source preceding its final end-document command is preserved
byte for byte.  The new checker evaluates the canonical phase recovery under
independent representative rephasings, all inverse tests, the exact recovered
benchmark and bad-helper obstruction, the unbalanced full-participation
comparison, the identity-only counterexample, the six-direction spectrum,
and both small-step orders.  Its explicit calculations retain the scalar
complement rather than replacing the physical five-copy representation by
the working block.  The full V19 checker is rerun separately.  The all-source
claims above have written proofs; the finite checks are not a proof-assistant
formalization, an independent peer review, or a historical source measurement.

\begin{thebibliography}{99}
\raggedright
\bibitem[34]{sel20-Wang}
Y.~Wang, D.~Dong, B.~Qi, J.~Zhang, I.~R.~Petersen and H.~Yonezawa,
\emph{A Quantum Hamiltonian Identification Algorithm: Computational Complexity
and Error Analysis}, arXiv:1610.08841 (2016).
\href{https://arxiv.org/abs/1610.08841}{arXiv:1610.08841}.
Cited for the established tomography-based Hamiltonian-identification setting,
not as a source of the inverse locking or neutral-helper theorems proved here.
\bibitem[35]{sel20-Abbott}
A.~A.~Abbott, J.~Wechs, D.~Horsman, M.~Mhalla and C.~Branciard,
\emph{Communication through coherent control of quantum channels},
Quantum \textbf{4} (2020), 333.
\href{https://doi.org/10.22331/q-2020-09-24-333}{doi:10.22331/q-2020-09-24-333}.
Used only for the distinction between separate channel data and a physically
coherently controlled implementation.  No such control is executed by the
inverse calculations in this block.
\end{thebibliography}
\addtocontents{toc}{\protect\endgroup}
% END OF SOURCE-SELECTION V20 DEVELOPMENT BLOCK.
% Preserve V1--V20 and append subsequent work before the final terminator.



\clearpage
\begin{center}
{\Large\bfseries Source-selection continuation V21}\\[.4em]
{\large Mixed outcomes, identity-generated Choi anchors,\\
and symmetry retained after the original labels are forgotten}
\end{center}
\addtocontents{toc}{\protect\begingroup}
\addtocontents{toc}{\protect\renewcommand{\protect\numberline}{\protect\makebox[2.1em][l]}}

\section{The next entrance is the actual mixed quantum operation}
\label{sel21:context}

V20 removes a prescribed coherent coefficient sum from its symmetry theorem,
but still assumes an efficient instrument: each original outcome has Choi rank
one and the coefficient rays are independent.  A source with unresolved noise,
a retained coarse record, or an idle outcome need not satisfy either condition.
Refining its Kraus decomposition on paper would not create an actual finer
record.  The present continuation works directly with the supplied positive
Choi operators and permits arbitrary ranks, dependencies, and numbers of labels.

There are two different advances.  An identity-cyclicity test on the Choi
operator algebra promotes covariance to pointwise invariance of the entire
Kraus span.  A smaller anchored test can determine the desired subgroup even
when that promotion is impossible.  The second test is applied to the unchanged
V19/V20 source after \emph{all six accepted labels are forgotten}.  Its quantum
output channel still has exactly the determinant subgroup.  The source is not
changed to obtain this result.  Separate comparison channels then show what
is, and is not, robust under explicitly declared additional noise.

\paragraph{Audit and established inputs.}
The Choi-derived common Kraus commutant is already
\cite{Rigidity}, \src{v3:thm:Choi-Howe}; it is not a new theorem here.
V20's efficient identity anchor and its participation condition are
Theorems~\ref{sel20:thm:anchor}--\ref{sel20:thm:phase-lock}.
The native coefficient bandwidth and first-mixed character classification
are \ref{sel19:thm:band} and \ref{sel18:thm:first-mixed}.
The neutral-helper compiler and its inverse are
\ref{sel19:thm:compiler} and \ref{sel20:thm:inverse}.
The distinction between channel covariance and stronger conditions on its
Kraus space is established in the literature
\cite{sel21-Haapasalo,sel21-deGroot}.  The cyclic-subspace argument is elementary
finite-dimensional operator algebra.  The additions here are its intrinsic
identity anchor, its record and error tests, the mixed-outcome determinant
selector, and exact evaluations on the existing source.  No priority is
claimed for the general covariance or cyclic-vector methods.

\paragraph{Fixed physical meaning.}
Every Choi operator refers to the same identified input/output memory and its
original probability weight.  The identity vector below represents the
identity map on that memory; it is not a freely chosen source vector fitted to
a desired symmetry.  Products of Choi matrices are \emph{data calculations},
not temporal products of the original quantum maps.  They require neither
coherent erasure of records nor separate access to a hidden Kraus component.
The tensor representation and the nonabelian symmetry tests remain actual
source conditions.  No historical five-copy process panel is supplied by this
continuation.

\section{Identity cyclicity removes efficiency and independence}
\label{sel21:cyclic}

Let $\mathcal J_y:M_N\to M_N$ be the actual CP outcomes of a normalized
instrument, with Choi operators $J_y\succeq0$, $1\le y\le s$.  Zero outcomes
may be removed.  Fix output-first vectorization and put
\begin{equation}
 i=|I_N\rangle\!\rangle,\quad e=i/\sqrt N,\quad
 S=\sum_yJ_y,\quad
 \mathscr W=\Span\{e\}+\Ran S,\quad d=\dim\mathscr W.
 \label{sel21:eq:W}
\end{equation}
All $J_y$ leave $\mathscr W$ invariant: positivity gives
$\Ran J_y\subseteq\Ran S$.  Define, on this space,
\begin{equation}
 \mathscr D=C^*(I_{\mathscr W},J_y|_{\mathscr W}:y),\qquad
 \mathscr V=\mathscr D e.
 \label{sel21:eq:cyclic-space}
\end{equation}
The original system algebra is instead
$\mathcal A_S=C^*(I_N,\operatorname{unvec}\Ran S)$.
These two algebras act on different spaces and are not identified.

\begin{theorem}[Finite identity-cyclicity and the exact residual rotation space]
\label{sel21:thm:cyclic}
Set
\begin{equation}
 Q_0=|e\rangle\langle e|,\qquad
 Q_{\ell+1}=Q_\ell+\sum_yJ_yQ_\ell J_y.
 \label{sel21:eq:recursion}
\end{equation}
Then $\Ran Q_{d-1}=\mathscr V$.  A first equality of two successive ranges
already certifies saturation at every longer depth.  In particular
\begin{equation}
 Q_{d-1}\succ0\text{ on }\mathscr W
 \quad\Longleftrightarrow\quad \mathscr V=\mathscr W.
 \label{sel21:eq:cyclic-test}
\end{equation}
For the unrestricted Choi-space unitary group
\[
 \mathcal U_e=\{V\in\mathrm U(\mathscr W):[V,J_y]=0\ (\forall y),\ Ve=e\},
\]
its common fixed-vector space is exactly $\mathscr V$.  Thus
$\mathscr V=\mathscr W$ if and only if $\mathcal U_e=\{I_{\mathscr W}\}$.
If it fails, $2P_{\mathscr V}-I_{\mathscr W}$ is an explicit nonidentity
member of $\mathcal U_e$.
\end{theorem}
\begin{proof}
The range of a sum of positive matrices is the span of their ranges.  Therefore
\eqref{sel21:eq:recursion} gives
$\Ran Q_{\ell+1}=\Ran Q_\ell+\sum_yJ_y\Ran Q_\ell$.
Induction identifies this with the span of $J_{y_k}\cdots J_{y_1}e$ of
length at most $\ell+1$.  A plateau makes this space invariant under each
$J_y$, hence under their unital algebra.  Starting at dimension one, at most
$d-1$ strict increases are possible.  The Hermitian generators make the
saturated space reducing, and it equals $\mathscr De$.

Every $V\in\mathcal U_e$ fixes each displayed word vector.  Conversely,
the orthogonal projection onto the reducing space $\mathscr V$ commutes
with every $J_y$, and $e$ lies in its range.  The displayed reflection is
therefore in $\mathcal U_e$ and negates every vector in $\mathscr V^\perp$.
A vector fixed by the whole group must lie in $\mathscr V$.  This proves all
assertions, including the case $d=1$.
\end{proof}

The residual reflection is a unitary on an \emph{operator Hilbert space}.
It need not have the form $R\otimes\overline R$ for a physical system unitary.
Failure of \eqref{sel21:eq:cyclic-test} is therefore not, by itself, a proof
of an extra physical gauge symmetry.

\begin{theorem}[Mixed-outcome covariance becomes pointwise coefficient symmetry]
\label{sel21:thm:rigidity}
Let $R(g)$ be any independently given unitary action on the same memory.
Every fixed-label covariance of the actual instrument fixes
$\operatorname{unvec}\mathscr V$ pointwise.  If
$\mathscr V=\mathscr W$, then
\begin{equation}
 \boxed{\operatorname{Cov}_{G}(\mathcal J_y:y)
       =\{g:R(g)\in\mathcal A_S'\}.}
 \label{sel21:eq:rigidity}
\end{equation}
The same conclusion holds under the weaker algebraic condition
\begin{equation}
 \Ran S\subseteq\operatorname{vec}
       C^*(\operatorname{unvec}\mathscr V).
 \label{sel21:eq:algebraic-completion}
\end{equation}
Neither Choi rank one, independent coefficient rays, nor a coherent-sum
calibration is assumed.
\end{theorem}
\begin{proof}
Write $\mathcal R(g)=R(g)\otimes\overline{R(g)}$.  Choi covariance is
$[\mathcal R(g),J_y]=0$ for every original label.  Always
$\mathcal R(g)e=e$.  It preserves $\mathscr W$ and its restriction belongs
to $\mathcal U_e$, so Theorem~\ref{sel21:thm:cyclic} makes it the identity
on $\mathscr V$.  If this space is $\mathscr W$, it fixes every vectorized
Kraus coefficient.  Their unvectorized span generates $\mathcal A_S$.
Under \eqref{sel21:eq:algebraic-completion}, fixing the anchor operators
also fixes their products and adjoints and hence the same Kraus span.
Conversely, a system unitary commuting with that span fixes each Kraus
coefficient in any factorization of each $J_y$, so it is a covariance of
each actual map.  This proves \eqref{sel21:eq:rigidity} without selecting
a Kraus refinement.
\end{proof}

\begin{corollary}[Recovery of the efficient criterion and of the idle outcome]
\label{sel21:cor:efficient}
V20's full-participation condition implies $\mathscr V=\mathscr W$.
Appending an actual idle identity outcome, with a nonzero common rescaling
of the accepted bank when required for normalization, does not change this
conclusion, even though the enlarged coefficient list is dependent.  More generally,
appending positive multiples of existing Choi matrices or of
$|i\rangle\langle i|$ leaves $\mathscr V$ unchanged.
\end{corollary}
\begin{proof}
On the independent efficient bank, $i\in\Ran S$ and $S^\dagger$ is a
polynomial in $S$: interpolate $x\mapsto1/x$ on the finite positive spectrum,
with value zero at zero.  Thus $S^\dagger e\in\mathscr De$.  V20 gives
$J_yS^\dagger i=\alpha_y|K_y\rangle\!\rangle$ with every $\alpha_y\ne0$.
Hence every coefficient lies in $\mathscr V$, which proves cyclicity.
The projector $|i\rangle\langle i|$ maps every vector into $\mathbb Ce$;
adding it or a repeated existing generator cannot enlarge the cyclic space.
It also does not enlarge $\mathscr W$.  Positive common rescalings do not
change any of these spans.
\end{proof}

\begin{corollary}[The identity-cyclic mixed determinant selector]
\label{sel21:cor:cyclic-selector}
On an invariant subspace of at most six parallel native copies, let the actual
instrument be nonabelian covariant and identity cyclic as above.  Its
Kraus-independent linear Choi form $\mathsf L$ from
Proposition~\ref{sel19:prop:linear} satisfies
\[
 \det\mathsf L=0,\quad\mathsf L_{CW}>0
 \quad\Longleftrightarrow\quad
 G_{\mathcal J}=S(\mathrm U(3)\times\mathrm U(2)).
\]
The positive branch has at least five copies and central null direction
$(-2,3)$.  The same holds under \eqref{sel21:eq:algebraic-completion}.
\end{corollary}
\begin{proof}
Theorem~\ref{sel21:thm:rigidity} makes every supported Kraus vector
nonabelian invariant and identifies its common pointwise stabilizer with
the actual covariance group.  Apply the inherited coefficient bandwidth
and first-mixed theorem to that finite span, using any factorization only
to expand the intrinsic positive form.  The character and positivity
arguments do not depend on the number or independence of its vectors.
\end{proof}

\section{Record loss and a quantitative cyclicity margin}
\label{sel21:record}

\begin{proposition}[Classical record processing has a monotone anchor]
\label{sel21:prop:record}
For an actual stochastic record processing
$J'_z=\sum_yP(z\mid y)J_y$, one has the same $S,\mathscr W$ and
$\mathscr V'\subseteq\mathscr V$.  If $P$ has full column rank, then
$\mathscr V'=\mathscr V$ and the two fixed-label covariance groups are
identical.  The linear Kraus-commutant form obtained by summing over all
outcomes is unchanged under \emph{every} stochastic record processing.
\end{proposition}
\begin{proof}
Column sums one give $\sum_zJ'_z=S$.  Each new Choi generator is a linear
combination of the old ones, proving cyclic-space inclusion.  Full column
rank gives a scalar left inverse and hence equality of generated unital
algebras.  It also makes $[\mathcal R(g),J'_z]=0$ equivalent to all the old
commutators being zero.  The inverse is a data identity, not an executed
stochastic inverse.  Finally the linear form depends only on $S$, which
proves its stated invariance even when actual covariance enlarges.
\end{proof}

For instance, on a qubit the two original coefficients $P_0,P_1$ form a
normalized efficient instrument and are identity cyclic.  After forgetting
their labels, $S$ is the identity on the diagonal coefficient space, so its
identity-cyclic subspace is only $\mathbb CI$.  The old fixed-label group is
the diagonal unitary group.  The forgotten channel also permits the unitary
swapping the two rays.  This verifies that an unchanged summed Kraus form
does not, in general, preserve the global covariance group.

\begin{proposition}[Approximate covariance controls coefficient motion on a cyclic branch]
\label{sel21:prop:quantitative}
Let
\[
 F_L=\sum_{|w|\le L}|J_we\rangle\langle J_we|,
 \qquad F_L\succeq\gamma P_{\mathscr W},\quad\gamma>0,
\]
where $J_w$ is a word in the actual $J_y$ and the empty word is the identity.
Take $\mu\ge\max(1,\max_y\|J_y\|)$ and
\[
 C_L=\sum_{\ell=1}^L\ell^2s^{\ell-1}\mu^{2\ell-2}.
\]
For every physical $\mathcal R=R\otimes\overline R$,
\begin{equation}
 \boxed{\sum_{y,\alpha}\|RK_{y\alpha}R^*-K_{y\alpha}\|_{\HS}^2
 \le\frac{\|S\|C_L}{\gamma}
                 \sum_y\|[\mathcal R,J_y]\|_{\op}^2.}
 \label{sel21:eq:quantitative}
\end{equation}
Here any Kraus decomposition of each original outcome gives the same left
side.  The analogous bound holds with $[T,K_{y\alpha}]$ on the left and
$[\widehat T,J_y]$ on the right for any Hermitian system $T$.
\end{proposition}
\begin{proof}
For a word of length $\ell$, expand $[\mathcal R,J_w]$ as the sum of
$\ell$ terms.  Since $\mathcal Re=e$, Cauchy--Schwarz bounds
$\| (\mathcal R-I)J_we\|^2$ by
$\ell\mu^{2\ell-2}$ times the sum of the squared commutator norms at its
letters.  Summing over all $s^\ell$ words counts each letter at each
position $s^{\ell-1}$ times, giving
\[
 \Tr[(\mathcal R-I)F_L(\mathcal R-I)^*]
 \le C_L\sum_y\|[\mathcal R,J_y]\|^2.
\]
The assumed form lower bound controls
$\gamma\|(\mathcal R-I)P_{\mathscr W}\|_{\HS}^2$.
The left side of \eqref{sel21:eq:quantitative} is exactly
$\Tr[(\mathcal R-I)S(\mathcal R-I)^*]$, at most $\|S\|$ times that
squared norm.  No invariance of $\mathscr W$ under an approximately
covariant $\mathcal R$ is required.  For the Lie assertion use
$\widehat Te=0$ and the same word-commutator expansion.
\end{proof}

This is a fixed-source bound.  A small or vanishing $\gamma$ is paid
explicitly; the result asserts no regulator-uniform rigidity or mass gap.
The recursion \eqref{sel21:eq:recursion} computes the range without enumerating
all words.  Positive reweightings may condition this calculation, but their
constants must be retained in quantitative comparisons.

\section{A smaller mixed-outcome selector from the identity Choi moment}
\label{sel21:moment}

Complete cyclicity can be unnecessarily strong.  Define the first identity
Choi moments
\begin{equation}
 A_y=\frac1N\operatorname{unvec}(J_y i),\qquad
 \mathsf M_{ab}=\operatorname{Re}\sum_y
          \Tr([N_a,A_y]^*[N_b,A_y]),\quad a,b\in\{C,W\}.
 \label{sel21:eq:moment}
\end{equation}
The factor $1/N$ is a declared diagnostic normalization.  $A_y$ is generally
not Hermitian or positive and is not $\mathcal J_y(I)$.
For a factorization of a Choi matrix it equals
$N^{-1}\sum_\alpha\overline{\Tr K_{y\alpha}}K_{y\alpha}$, independently
of that factorization.  More generally any finite list of vectors in
$\mathscr V$, unvectorized with fixed coefficients, can be used as anchors.
Every actual covariance fixes these anchors pointwise.

Let $R_k$ be the independently verified native action on an invariant subspace
of $H_{14}^{\otimes k}$ with retained neutral multiplicity, and retain the
original $N_C,N_W$ on that space.  Assume the actual marked Choi matrices are
covariant under $K=\mathrm{SU}(3)\times\mathrm{SU}(2)$.  This is the full
Choi test, not just invariance of the $A_y$.

\begin{theorem}[Mixed-outcome determinant selection without a Kraus-rank hypothesis]
\label{sel21:thm:selector}
Suppose $k\le6$ and
\begin{equation}
 \det\mathsf M=0,\qquad \mathsf M_{CW}>0.
 \label{sel21:eq:moment-test}
\end{equation}
Let $v$ be a nonzero vector in $\Ker\mathsf M$, derived from this Gram.
Then $k\ge5$, $v\in\mathbb R(-2,3)$, and the pointwise stabilizer of the
anchor packet inside $G_0=\mathrm U(3)\times\mathrm U(2)$ is exactly
$G_{\det}=S(\mathrm U(3)\times\mathrm U(2))$.  Furthermore
\begin{equation}
 \boxed{\operatorname{Cov}_{G_0}(\mathcal J_y:y)=G_{\det}
 \quad\Longleftrightarrow\quad
 r_v:=\sum_y\|[\widehat T_v,J_y]\|_{\HS}^2=0,}
 \label{sel21:eq:selector}
\end{equation}
where $T_v=v_CN_C+v_WN_W$ and
$\widehat T_v=T_v\otimes I-I\otimes\overline{T_v}$.
No Choi rank, coefficient independence, calibrated identity expansion, or
identity-cyclicity assumption enters this theorem.
\end{theorem}
\begin{proof}
Nonabelian Choi covariance and $\mathcal R(g)i=i$ make each $A_y$
nonabelian invariant.  Consequently its Fourier decomposition consists of
determinant characters $(\det U)^p(\det V)^q$.  The inherited
\emph{endomorphism} bandwidth on $k$ native copies is
\[
 \max\{3|p|,2|q|,|3p+2q|\}\le k.
\]
Orthogonality of the modes expresses $\mathsf M$ as a positive sum of
outer products $(3p,2q)(3p,2q)^{\mathsf T}$, with coefficients the squared
HS norms of all occurring anchor modes.  Rank one forces collinearity; its
positive mixed entry forces same-sign nonzero coordinates.  At $k\le6$,
only $(p,q)=\pm(1,1)$ is possible.  At least one occurs, even though a
non-Hermitian anchor need not contain both signs.  Its stabilizer is
$\ker\chi$, with $\chi=\det U\det V$.  This proves the claimed copy
floor, primitive global group, and derived null direction.

Every actual covariance fixes every anchor, so it is contained in
$G_{\det}$.  If $r_v=0$, the Choi matrices commute with the derived central
one-parameter group.  Together with their assumed nonabelian covariance,
this gives covariance under the connected group $G_{\det}$.  The two
inclusions give equality.  Conversely equality implies that commutation
and hence $r_v=0$.  No statement about pointwise invariance of the remaining
Kraus span was needed.
\end{proof}

This is a sufficient selector with an exact last alternative on its entrance
branch; some determinant-covariant sources have scalar first anchors and do
not pass \eqref{sel21:eq:moment-test}.  The copy floor is not a universal
mixed-channel lower bound: V17 attains three copies in its wider CP class.
The theorem instead detects an endomorphism character already anchored by
the supplied channel.  It derives the central ratio before testing its
survival in the full operation.

For $S$ of a single forgotten CPTP channel, $S/N$ is its normalized Choi
state.  Thus $\Tr(B^*A)=\langle\!\langle B|(S/N)i\rangle$ can be obtained
from real and imaginary parts of the corresponding rank-one comparison
observable, or from physical process tomography.  This observation specifies
data requirements; it does not admit an absent entangled preparation or
comparison Read into the historical protocol.

\begin{proposition}[A perturbation bound for the smaller response]
\label{sel21:prop:moment-error}
If $\sum_y\|\widehat J_y-J_y\|_{\HS}^2\le\epsilon^2$, then
$\sum_y\|\widehat A_y-A_y\|_{\HS}^2\le\epsilon^2/N$.
Let $D$ send $t\in\mathbb R^2$ to
$([t_CN_C+t_WN_W,A_y])_y$, and define $\widehat D$ similarly.  Then
\begin{equation}
 \|D-\widehat D\|\le\delta:=\frac{2\sqrt2 k\epsilon}{\sqrt N},\qquad
 \|\mathsf M-\widehat{\mathsf M}\|
     \le\delta(2\|\widehat D\|+\delta).
 \label{sel21:eq:moment-error}
\end{equation}
\end{proposition}
\begin{proof}
$\|i\|=\sqrt N$ and $\|\Delta Ji\|\le\|\Delta J\|_{\HS}\sqrt N$
give the first bound.  Since $\|N_C\|,\|N_W\|\le k$, the elementary
commutator norm inequality and Cauchy--Schwarz give the displayed bound
on $D-\widehat D$.  Expand $D^*D-\widehat D^*\widehat D$ to obtain the
last estimate.
\end{proof}
Small errors can certify a positive mixed entry or exclude rank one.  They do
not prove exact nonabelian covariance, $\det\mathsf M=0$, or $r_v=0$.
An original opportunity weight multiplies $A_y$ by that weight and its
quadratic Gram by its square.  This differs from the linear Kraus form of
V19; the two calibrations are not interchanged.

\section{The original neutral compiler retains its symmetry with no accepted labels}
\label{sel21:compiler}

Keep exactly a compiler recognized in V20.  Write $q=\tau^4$,
$D=(I-qR)^{1/2}$, $B=UD$, and $C_\mu=|u_0\rangle\langle u_{\mu-1}|$
for its four neutral helper maps.  Their traces are zero.  The accepted
channel obtained by \emph{forgetting}, not coherently recombining, its six
labels has the inherited exact form
\begin{equation}
 \Phi(X)=\frac13X+\frac23\left(BXB^*+
                      q\sum_{\mu=2}^5C_\mu XC_\mu^*\right).
 \label{sel21:eq:forgotten}
\end{equation}
Its Choi operator and first anchor are
\begin{align}
 S&=\frac13|I\rangle\!\rangle\langle\!\langle I|
    +\frac23|B\rangle\!\rangle\langle\!\langle B|
    +\frac{2q}{3}\sum_\mu|C_\mu\rangle\!\rangle\langle\!\langle C_\mu|,
                                                        \label{sel21:eq:S}\\
 A&=\frac13I+\frac{2\overline{b}}{3N}B,\qquad b=\Tr B.
                                                        \label{sel21:eq:A}
\end{align}

\begin{theorem}[Exact symmetry retention after all original accepted labels are forgotten]
\label{sel21:thm:forget}
Assume the whole helper support is $G_0$-neutral, $0<q<1$, and $b\ne0$.
Then
\begin{equation}
 \boxed{\operatorname{Cov}_{G_0}(\Phi)
       =\{g:[R_k(g),U]=0\}
       =\operatorname{Cov}_{G_0}(\mathcal J_e:e).}
 \label{sel21:eq:forget}
\end{equation}
On the inherited no-alias branch this group is also the stabilizer of the
original Hermitian $H$.  The first identity holds without a no-alias
assumption: it then identifies only the symmetry of $U$.
The identity-cyclic space of the single $S$ is exactly
$\Span\{|I\rangle\!\rangle,|B\rangle\!\rangle\}$, of dimension two.
Its Kraus support has dimension six.
\end{theorem}
\begin{proof}
Orthogonality of the compiler's record columns proves
\eqref{sel21:eq:forgotten} and \eqref{sel21:eq:S} without changing its
outcomes.  Contract \eqref{sel21:eq:S} against $i/N$; the helper traces
vanish, giving \eqref{sel21:eq:A}.  A covariance of $\Phi$ fixes $A$.
Since $b\ne0$, it commutes with $B$.  Every $R_k(g)$ commutes with $D$
and $C_\mu$ by neutrality, and $D$ is invertible.  Hence it commutes with
$B$ if and only if it commutes with $U$.  Conversely such commutation fixes
all terms of \eqref{sel21:eq:forgotten}, and all six original coefficients.
The no-alias identification is the already proved spectral criterion of V19.

$B$ is diagonal in the helper eigenbasis, so it is HS orthogonal to all
$C_\mu$, as is $I$.  Thus $S$ preserves $\Span\{i,|B\rangle\!\rangle\}$.
The $b\ne0$ contraction generates $B$ from $i$ in one step.  $B$ cannot be
scalar: its modulus is $\sqrt{1-q}$ on the four helper input rays and one
on their orthogonal complement.  The four helpers are independent and
orthogonal to the first two directions.  This proves the final dimensions.
\end{proof}

The incomplete identity-cyclic span therefore need not be an obstruction to
\emph{relative} physical symmetry.  The four missing directions here are
known to be neutral; the anchored polar factor already decides the group.
This result does not reconstruct the individual six recorded maps from the
forgotten channel, or identify its scalar energy origin.  V20's distinct
marked implementations with the same forgotten channel remain distinct.

\begin{theorem}[Invariant isotropic noise preserves the group, not the Kraus commutant]
\label{sel21:thm:noise}
For any CPTP map $\Phi$ on $M_N$, let
$\mathcal T(X)=\Tr(X)I_N/N$ and
$\Phi_\epsilon=(1-\epsilon)\Phi+\epsilon\mathcal T$, $0<\epsilon<1$.
For every comparison unitary group,
\begin{equation}
 \operatorname{Cov}(\Phi_\epsilon)=\operatorname{Cov}(\Phi),\qquad
 S_\epsilon=(1-\epsilon)S+\frac\epsilon N I_{N^2}\succ0,
 \label{sel21:eq:noise}
\end{equation}
while $\mathcal A_{S_\epsilon}=M_N$.  The normalized first anchor and its
central response obey
\begin{equation}
 A_\epsilon=(1-\epsilon)A+\frac\epsilon{N^2}I_N,\qquad
 \mathsf M_\epsilon=(1-\epsilon)^2\mathsf M.
 \label{sel21:eq:noise-anchor}
\end{equation}
For any Hermitian system $T$, its linear Kraus form is instead
\begin{equation}
 q_{S_\epsilon}(T)=(1-\epsilon)q_S(T)
      +2\epsilon\left(\Tr T^2-\frac{(\Tr T)^2}{N}\right).
 \label{sel21:eq:noise-Kraus}
\end{equation}
It is positive on every nonscalar $T$, even one generating an exact symmetry
of $\Phi_\epsilon$.
\end{theorem}
\begin{proof}
$\mathcal T$ is covariant under every unitary.  The covariance difference
of the mixture is therefore exactly $1-\epsilon$ times that of $\Phi$,
proving the group equality.  Its Choi matrix is the displayed strictly
positive matrix, so its Kraus support is all of $M_N$.  Contracting against
$i/N$ proves \eqref{sel21:eq:noise-anchor}; the added identity commutes with
both central generators.  Finally the depolarizing Kraus operators are
$E_{ab}/\sqrt N$.  Diagonalizing Hermitian $T$ shows
\[
 \frac1N\sum_{a,b}\|[T,E_{ab}]\|_{\HS}^2
 =\frac1N\sum_{a,b}(t_a-t_b)^2
 =2\Tr T^2-2(\Tr T)^2/N.
\]
Linearity of the intrinsic form in $S$ gives \eqref{sel21:eq:noise-Kraus}.
\end{proof}

This theorem describes a \emph{declared comparison family}, not a permission
to add noise to the old source.  It shows why a positive linear Choi--Howe
form cannot generally be called a covariance detector for mixed outcomes.
That form still correctly detects the common coefficient commutant.  The
smaller anchored test can detect the actual covariance even when the
coefficient algebra has become the full matrix algebra.

\section{Exact evaluation of the unchanged source and separating comparisons}
\label{sel21:benchmark}

Keep the full native dimension $N=14^5=537824$, the original setting
$\tau=1/8$, $s=4\pi/3$, and the V19 vectors $n,\omega,u_0,\ldots,u_4$.
Set
\[
 d_0=\sqrt{1-1/4096}=\frac{\sqrt{4095}}{64},\qquad
 x=N-5+4d_0.
\]
On the determinant pair, $U=I-(1+i)\left(\begin{smallmatrix}1&1\\1&1\end{smallmatrix}\right)/2$;
on the helpers and the orthogonal complement it is the identity.

\begin{proposition}[One identity Choi moment certifies the forgotten benchmark]
\label{sel21:prop:benchmark}
For exactly the unchanged source,
\begin{equation}
 b=x-i\ne0,\qquad
 \boxed{\mathsf M=\frac{4(x^2+1)}{9N^2}
                 \begin{pmatrix}9&6\\6&4\end{pmatrix}.}
 \label{sel21:eq:benchmark}
\end{equation}
The full forgotten channel is nonabelian covariant and has $r_{(-2,3)}=0$.
Thus its covariance group is exactly $G_{\det}$, despite its Choi rank six
and identity-cyclic dimension two.  The first-anchor test does not require
retaining any of the six accepted labels.
\end{proposition}
\begin{proof}
The determinant block has trace $1-i$, the five helpers contribute
$1+4d_0$, and the $N-7$ scalar complement contributes $N-7$.
Their sum is $x-i$.  The two off-diagonal entries of $B$ on $n,\omega$
have squared modulus $1/2$, and their charge differences are
$\pm(3,2)$.  All other entries commute with both counting operators.  Hence
the coefficient commutator Gram of $B$ is
$\left(\begin{smallmatrix}9&6\\6&4\end{smallmatrix}\right)$.
Equation~\eqref{sel21:eq:A} gives \eqref{sel21:eq:benchmark}.
Every neutral helper and the determinant pair action commute with the
nonabelian group and $-2N_C+3N_W$.  Therefore the complete channel passes
both covariance tests.  Apply Theorem~\ref{sel21:thm:selector}, or directly
Theorem~\ref{sel21:thm:forget} and the inherited stabilizer of $H_s$.
\end{proof}

\begin{proposition}[An unresolved-rank test on the same record family]
\label{sel21:prop:blur}
For $0<\zeta<1$ set
\[
 J'_z=(1-\zeta)J_z+\frac\zeta6 S,\qquad z=1,\ldots,6.
\]
Each actual processed outcome has Choi rank six, but its identity-cyclic
space has dimension six already at depth one.  At $\zeta=1$ all six
operators equal $S/6$, and the cyclic dimension drops to two.  The covariance
group is $G_{\det}$ throughout $0\le\zeta\le1$ for this benchmark.
\end{proposition}
\begin{proof}
Positive coefficients on all six old rays give rank six for every
$\zeta>0$.  The old coefficients are independent and
\[
 \Tr K_e=\frac{N}{\sqrt{18}}\pm\frac b3.
\]
Its imaginary part is nonzero, so each $J_ei$ is a nonzero multiple of
$|K_e\rangle\!\rangle$.  These six vectors span the full support.
The stochastic mixing matrix has eigenvalues $1,1-\zeta$ and is invertible
for $\zeta<1$, hence its first-anchor span is the same.  At $\zeta=1$
Theorem~\ref{sel21:thm:forget} gives the cyclic dimension and the group.
For $\zeta<1$ the group assertion is
Proposition~\ref{sel21:prop:record}.  No hidden Kraus component was declared
to be a processed outcome.
\end{proof}

For the isotropic-noise comparison of this same forgotten benchmark,
$\rank S_\epsilon=N^2$ but the identity-cyclic space stays two-dimensional:
$S_\epsilon$ is an affine function of $S$ with nonzero slope.  Formula
\eqref{sel21:eq:benchmark} is multiplied by $(1-\epsilon)^2$ and still
certifies the exact determinant group.  In contrast, put
$Y=-2N_C+3N_W$.  On one native factor,
\[
 \frac{\Tr Y}{14}=\frac67,\qquad
 \frac{\Tr Y^2}{14}=\frac{33}{7},\qquad
 \operatorname{Var}_{I/14}(Y)=\frac{195}{49}.
\]
On five factors its normalized trace variance is $975/49$.  Therefore the
linear Kraus form on the \emph{exactly conserved} generator is
\begin{equation}
 \boxed{q_{S_\epsilon}(Y)=\frac{1950\epsilon N}{49}>0.}
 \label{sel21:eq:noise-Y}
\end{equation}
No physical charge or source weight has been divided out of this form.

\begin{proposition}[The anchored direction need not survive the whole channel]
\label{sel21:prop:counter}
The last covariance residual in Theorem~\ref{sel21:thm:selector} cannot be
removed.  On the same five-copy native carrier choose the colour determinant
ray $c=\Omega_C\otimes r_0^{\otimes2}$, which is orthogonal to $n,\omega$
and all neutral helpers.  Put
\[
 P=|n\rangle\langle n|+|c\rangle\langle c|,\quad
 X=|n\rangle\langle c|+|c\rangle\langle n|,\quad Q=I-P,
 \qquad \mathcal N(Z)=XZX+QZQ.
\]
For $0<\epsilon<1$ let
$\widetilde\Phi_\epsilon=(1-\epsilon)\Phi+\epsilon\mathcal N$.
This is a normalized nonabelian-covariant channel, and its anchor response
is still $(1-\epsilon)^2\mathsf M$.  But its full Choi residual satisfies
\begin{equation}
 \boxed{\|[\widehat Y,J(\widetilde\Phi_\epsilon)]\|_{\HS}^2
           =288\epsilon^2>0.}
 \label{sel21:eq:counter}
\end{equation}
It consequently has no surviving continuous central direction.
\end{proposition}
\begin{proof}
$X^*X=P$ and $Q^*Q=Q$, so $\mathcal N$ is trace preserving.  The three
rays $n,c,\omega$ are nonabelian singlets and their projections are
$G_0$ invariant.  Thus $X,Q$ commute with the nonabelian group.
$\Tr X=0$, so the first identity Choi moment of $\mathcal N$ is
$(N-2)Q/N$, a $G_0$-invariant operator.  Its addition has no central
commutator response.  This proves the asserted anchored Gram.

On $n,c$, the eigenvalues of $Y$ are $0,-6$.  The two vectorized arms of
$X$ therefore have $\widehat Y$ charges $6,-6$.  Their cross blocks in
$|X\rangle\!\rangle\langle\!\langle X|$ have commutator coefficients
$12,-12$ and are HS orthogonal.  The squared commutator norm is $288$.
$Q$ commutes with $Y$, and the original $S$ also does, proving
\eqref{sel21:eq:counter}.  Any continuous central covariance of the mixture
must fix its first anchor and hence lie in $\mathbb R(-2,3)$; the nonzero
residual excludes that remaining line.  The comparison changes the original
source.  It does not assert that the new $X$ operation historically occurs.
\end{proof}

\section{What is now decidable without a resolved Kraus record}
\label{sel21:status}

The efficient-outcome hypothesis is no longer necessary for a sufficient
source-level symmetry selector.  The finite Choi-space cyclicity test
licenses the whole Kraus-commutant test when it passes.  When it does not,
the first identity moment can already derive the primitive central direction,
while the full-operation residual decides whether that direction survives.
Neither route invents a Kraus measurement.  The two routes are different:
the strictly positive noisy benchmark passes the smaller symmetry selector
while its complete coefficient commutant is only scalar.

Most importantly, an actual calculation rather than a new favourable
realization reduces the benchmark's data demand.  The unchanged neutral-helper
source retains the determinant group after all six accepted labels are
forgotten.  Its first identity moment suffices; complete reconstruction of
the original helper rays or of a Hamiltonian logarithm is unnecessary for
this \emph{group} conclusion.  The recorded implementation and the
continuous-time action are not reconstructed uniquely from that smaller
object.  V20's energy aliases and distinct marked refinements remain in force.

The next historical source row can therefore be the same-cut \emph{forgotten
quantum output channel}, provided that its actual representation and covariance
are verified.  Classical outcome frequencies alone still do not supply it.
Evaluate its identity moment and the derived central residual before asking
for a resolved five-copy implementation.  If this smaller test fails, the
full marked Choi cyclic space is a second finite route, not an instruction
to modify the record policy.  The original joint-packet requirement remains
that all mixed data belong to one actual carrier \cite{Duality}.

Nothing here populates that missing historical channel, supplies the charged
matter modules, identifies the geometrically named private and type rays,
proves a physical fermionic grading, or constructs an Einstein--SM continuum.
The result concerns covariance of a specified operation inside its verified
$G_0$.  It is not a mass-gap estimate, a proof of gauge selection from every
renewal law, or a claim that the internal algebra alone fixes all particle
incidences.

\paragraph{Verification and preservation.}
All V20 bytes preceding its final end-document command are retained.  The new
script checks cyclic ranges, the residual reflection, dependent idle and
mixed-outcome record examples, normalized complete channels, the unchanged
seven-dimensional working block with its full native scalar complement,
the first Choi moment and both central responses, isotropic-noise support
and variance formulas, and the exact $288\epsilon^2$ obstruction.
The V20 checker is rerun separately.  Finite computations verify the displayed
identities; all-source statements have the proofs above.  These are not
proof-assistant formalizations, independent peer review, or historical source
measurements.

\begin{thebibliography}{99}
\raggedright
\bibitem[36]{sel21-Haapasalo}
E.~Haapasalo, \emph{The Choi--Jamiolkowski isomorphism and covariant quantum
channels}, arXiv:1906.11442v2 (2019).
\href{https://arxiv.org/abs/1906.11442}{arXiv:1906.11442}.
Background for covariance expressed in Choi data, not a source of the
identity-cyclicity or benchmark-retention claims proved here.
\bibitem[37]{sel21-deGroot}
C.~de Groot, A.~Turzillo and N.~Schuch,
\emph{Symmetry Protected Topological Order in Open Quantum Systems},
Quantum \textbf{6} (2022), 856.
\href{https://doi.org/10.22331/q-2022-11-10-856}{doi:10.22331/q-2022-11-10-856}.
Background for distinct weak and strong channel symmetries.  No topological
phase or many-body anomaly theorem from that work is asserted here.
\end{thebibliography}
\addtocontents{toc}{\protect\endgroup}
% END OF SOURCE-SELECTION V21 DEVELOPMENT BLOCK.
% Preserve V1--V21 and append subsequent work before the final terminator.


\clearpage
\hypersetup{pdftitle={Historical Standard-Model Source Selection: V1--V22}}
\begin{center}
{\Large\bfseries Source-selection continuation V22}\\[.4em]
{\large What incomplete measurements actually certify:\\
positive completion fibres and a two-scalar native symmetry witness}
\end{center}
\addtocontents{toc}{\protect\begingroup}
\addtocontents{toc}{\protect\renewcommand{\protect\numberline}{\protect\makebox[2.1em][l]}}

\section{Go upstream from a reconstructed channel to its actual measurements}
\label{sel22:context}

V21 removes efficiency, coefficient independence, and even the accepted labels
from a sufficient symmetry test.  It still requires a quantum output channel,
or enough physical comparisons to evaluate its anchor and its full covariance
residual.  The next question is not another favourable action or Kraus bank.
It is whether a smaller \emph{incomplete measurement record} already forces
the desired symmetry for every normalized positive channel consistent with it.
Choosing one symmetric reconstruction from an underdetermined record would
not establish that conclusion.

We prove a finite necessary-and-sufficient test for symmetry shared by the
entire positive completion fibre.  At a faithful Choi point this gives a sharp
linear-information lower bound.  At the opposite boundary, a positive null
measurement can confine all completions to an operator-support face.  Applied
to the \emph{unchanged} V19--V21 benchmark, that mechanism gives a sufficient
certificate consisting of two real linear expectations, and we prove that one
such expectation cannot suffice.  The certificate determines the covariance
group without determining the channel.  It requires two explicitly specified
collective comparison effects on a physical Choi state; their availability is
not inferred from the original six outcome names.

\paragraph{Audit and inherited inputs.}
The first identity moment, its covariance, the exact forgotten compiler, and
its evaluated native trace are respectively
\ref{sel21:thm:selector}, \eqref{sel21:eq:forgotten},
\eqref{sel21:eq:A}, and Proposition~\ref{sel21:prop:benchmark}.
Their formulas are retained.  We do not reprove the determinant-character
bandwidth, the action inverse, the mixed-source cyclicity theorem, or the
support-polar duality in \cite{Rigidity}.  Targeted searches and literal
inspection of the mounted cumulative source and rigidity notes found those
results, but not the positive-completion symmetry criterion or the two-scalar
minimality theorem below.  This is a limited nonduplication audit, not a
verification of every archived proof.  Positivity-assisted tomography and the
minimal-face geometry of semidefinite feasible sets are established methods
\cite{sel22-BDK,sel22-SWW}; the specialized claims below have direct proofs.

\paragraph{Scope of data and of exactness.}
Fix one physically identified memory of dimension $N$ and its actual channel
$\Phi:M_N\to M_N$.  Let
\[
 J=J(\Phi),\qquad \rho_J=J/N,\qquad
 \Tr_{\rm out}J=I_N.
\]
Output-first vectorization is retained.  A one-use comparison prepares a
maximally entangled input and reference, applies the unknown channel once,
and measures a declared effect on the joint output, thereby producing a
linear expectation in $\rho_J$.  No joint measurement of two independent
Choi copies is included in the information counts below.  Equivalent
independently validated process data can supply these same expectations.
Known reference frames, tensor-factor identifications, and admitted comparison
effects are part of the protocol, not consequences of its two scalar values.
All exact zero statements below concern exact data or exact source identities;
finite statistical error is treated separately.

\section{An exact test on the whole positive completion fibre}
\label{sel22:fibre}

Let $\mathcal L$ be a real linear map on Hermitian $N^2$-by-$N^2$ matrices.
It includes the trace-preserving constraints and all declared real comparison
expectations, with measured vector $b$.  Write
\begin{equation}
 \mathcal C_b=\{J\succeq0:\mathcal L(J)=b\},\qquad
 \Tr_{\rm out}J=I_N\quad(J\in\mathcal C_b).
 \label{sel22:eq:fibre}
\end{equation}
We initially assume feasibility.  This compact convex set retains all positive
completions, not just completions of a preferred rank or Kraus type.  The
physical action $R(g)$ induces
$\mathcal R(g)=R(g)\otimes\overline{R(g)}$ on the Choi Hilbert space.

\begin{lemma}[Maximal supported completion and its exact tangent space]
\label{sel22:lem:face}
Let $P$ be the projection onto the span of the ranges of all $J\in\mathcal C_b$.
There exists $J_\circ\in\mathcal C_b$ strictly positive on $\Ran P$ and zero
on its orthogonal complement.  Define
\begin{equation}
 \mathcal N_P=\{X=X^*:X=PXP,\ \mathcal L(X)=0\}.
 \label{sel22:eq:null}
\end{equation}
Then
\begin{equation}
 \boxed{\operatorname{aff}\mathcal C_b=J_\circ+\mathcal N_P.}
 \label{sel22:eq:affine}
\end{equation}
A maximal-support point can be found by at most $N^2-1$ strict support
enlargements from any nonzero feasible point, using linear semidefinite
maximizations over the same completion fibre.
\end{lemma}
\begin{proof}
The ambient space is finite dimensional.  Choose finitely many feasible
matrices whose ranges span $\Ran P$.  Their positive convex average is
feasible and has precisely that range: the kernel of a sum of positive
matrices is the intersection of their kernels.  It is therefore strictly
positive on $P$.  Every difference of feasible matrices lies in $\mathcal N_P$.
Conversely, for $X\in\mathcal N_P$ and
$|t|\|X\|<\lambda_{\min}(J_\circ|_P)$, both $J_\circ\pm tX$ are feasible.
They are zero off $P$ and remain positive on it.  This proves the affine hull.

For the last assertion, from a feasible $J_0$ with support $P_0$ maximize
$\Tr((I-P_0)J)$ over $\mathcal C_b$.  A maximum exists by compactness.  If
it is zero, positivity implies every feasible range lies in $P_0$, so the
support is maximal.  If it is positive, average a maximizing matrix with
$J_0$.  The support strictly increases.  At most $N^2-1$ such increases are
possible.  This is an exact finite procedure, not a claim of good numerical
conditioning, a one-step facial reduction, or polynomial computational cost.
\end{proof}

\begin{theorem}[Symmetry forced by incomplete physical process data]
\label{sel22:thm:universal}
Let $X_1,\ldots,X_d$ be any real basis of $\mathcal N_P$.  The covariance
group common to every compatible channel is exactly
\begin{equation}
 \boxed{
 G_{\rm forced}(b)=
 \{g:[\mathcal R(g),J_\circ]=0,
          \ [\mathcal R(g),X_a]=0\ (1\le a\le d)\}.}
 \label{sel22:eq:forced}
\end{equation}
For a connected proposed subgroup with Hermitian system generators $T_j$,
its inclusion in every compatible covariance group is equivalent to
\begin{equation}
 \boxed{
 \sum_j\left(\|[\widehat T_j,J_\circ]\|_{\HS}^2+
                    \sum_a\|[\widehat T_j,X_a]\|_{\HS}^2\right)=0,
 \quad \widehat T_j=T_j\otimes I-I\otimes\overline T_j.}
 \label{sel22:eq:forced-Lie}
\end{equation}
If either condition fails, an explicit compatible channel fails the asserted
symmetry: use $J_\circ$, or one of $J_\circ\pm tX_a$ at a sufficiently small
nonzero real $t$.
\end{theorem}
\begin{proof}
Covariance is the linear equation
$\mathcal R(g)J\mathcal R(g)^*=J$.  If it holds on the fibre, it holds on
its affine hull, and taking the two signs in the lemma forces it separately
on $J_\circ$ and every null direction.  Conversely those equalities imply it
on every $J_\circ+X$, hence on the whole fibre.  Differentiation and
exponentiation give the connected-group version.  A sum of squared norms is
zero exactly when every commutator is zero.  If the central point is fixed
but a null direction is not, either sign of a nonzero sufficiently small
perturbation fails.  If the central point already fails, no perturbation is
needed.  The chosen basis or maximal-support average affects the displayed
coordinates but not the resulting group.
\end{proof}

The theorem certifies a \emph{common subgroup}, not equality of the covariance
groups of all completions.  An additional data-certified symmetry-breaking
signal is needed to exclude larger groups.  The next sections give such a
signal for the native benchmark.

\begin{proposition}[Dual test and a positive null measurement]
\label{sel22:prop:dual}
For a Hermitian comparison $F$, its expectation is constant on $\mathcal C_b$
exactly when
\[
 PFP\in\Ran\mathcal L_P^*,\qquad
 \mathcal L_P:\operatorname{Herm}(P)\longrightarrow\operatorname{Ran}\mathcal L.
\]
If $PFP=\mathcal L_P^*c$, that constant is $c\cdot b$.
In particular, if a declared positive tester $Z\succeq0$ has measured
expectation zero, then every feasible $J$ has range in $\Ker Z$.
\end{proposition}
\begin{proof}
Constancy is orthogonality of $PFP$ to $\Ker\mathcal L_P$ by
Lemma~\ref{sel22:lem:face}.  In finite real inner-product spaces,
$(\Ker\mathcal L_P)^\perp=\Ran\mathcal L_P^*$, proving both assertions.
For the last statement,
$\Tr(ZJ)=\|Z^{1/2}J^{1/2}\|_{\HS}^2=0$ forces $Z^{1/2}J^{1/2}=0$.
\end{proof}

\section{Faithful noise changes the linear information needed for exact symmetry}
\label{sel22:interior}

Let
\[
 \mathcal T_N=\{X=X^*:\Tr_{\rm out}X=0\},\qquad
 \dim_{\mathbb R}\mathcal T_N=N^4-N^2.
\]
For a subgroup $H$, write $\mathcal F_H$ for its fixed Hermitian Choi space.
The next bound counts independent real linear observations in addition to
normalization.  It permits arbitrary collective one-use testers; restricting
the actual experimental alphabet can only reduce their power.

\begin{theorem}[Sharp linear codimension bound at a faithful covariant point]
\label{sel22:thm:count}
Suppose a completion fibre contains an $H$-covariant $J_\circ\succ0$.
For every compatible channel to be $H$-covariant, its $p$ independent real
measurement constraints must satisfy
\begin{equation}
 \boxed{p\ge c_H:=\dim\mathcal T_N-
                         \dim(\mathcal T_N\cap\mathcal F_H).}
 \label{sel22:eq:count}
\end{equation}
The bound is sharp for testing inclusion of $H$: $c_H$ linear expectations
can be chosen to force exactly the vanishing of all noninvariant tangent
coordinates.  They need not reconstruct the invariant coordinates or the
full channel.
\end{theorem}
\begin{proof}
Here $P=I$.  If $\mathcal M$ is the measurement map restricted to
$\mathcal T_N$, the invisible tangent is $\Ker\mathcal M$.
Theorem~\ref{sel22:thm:universal} requires
$\Ker\mathcal M\subseteq\mathcal T_N\cap\mathcal F_H$.
Thus $\rank\mathcal M\ge c_H$, proving necessity.  Conversely, choose a
real basis of the orthogonal complement of $\mathcal T_N\cap\mathcal F_H$
inside $\mathcal T_N$ and measure its $c_H$ coordinates.  Relative to the
invariant normalized depolarizing Choi matrix, setting them to zero leaves
only invariant tangents.  Each Hermitian linear functional can be shifted
and positively rescaled to an effect without changing its information content.
\end{proof}

\begin{corollary}[Exact central count on the unchanged native inventory]
\label{sel22:cor:central-count}
For a system generator $Y$ with distinct eigenvalues $y$ of multiplicities
$d_y$, let
\[
 h_r=\sum_{y-z=r}d_yd_z,\qquad N=\sum_y d_y.
\]
For covariance under the circle generated by $Y$,
\begin{equation}
 \boxed{c_Y=N^4-N^2-\sum_r h_r^2+\sum_y d_y^2.}
 \label{sel22:eq:charge-count}
\end{equation}
For $Y=-2N_C+3N_W$ on the full five-copy native carrier,
\[
 \sum_y d_y z^y=(3z^{-2}+5+6z^3)^5,\qquad N=14^5=537824,
\]
this gives
\begin{equation}
 \begin{aligned}
 \dim\mathcal T_N&=83\,668\,255\,424\,995\,546\,905\,600,\\
 c_Y&=79\,711\,967\,200\,263\,072\,203\,952.
 \end{aligned}
 \label{sel22:eq:large-count}
\end{equation}
The second number is about $95.27\%$ of the first.  It is a necessary count
for \emph{all} completions of a faithful channel to preserve this circle
when no further exact source restrictions are imposed.
\end{corollary}
\begin{proof}
The Choi representation has eigenvalue differences $r=y-z$ with multiplicities
$h_r$.  Its fixed Hermitian space has real dimension $\sum_rh_r^2$.
Partial trace maps this space onto the system's $Y$-fixed Hermitian space,
of dimension $\sum_y d_y^2$: the invariant preimage of $B$ is
$I_N/N\otimes\overline B$.  Its kernel therefore has the difference of
these dimensions.  Apply Theorem~\ref{sel22:thm:count}.  Multiplying the
five finite Laurent polynomials gives
\[
 \sum_y d_y^2=24643380100,\qquad
 \sum_rh_r^2=3956288224757118081748,
\]
which proves the exact integers.  The companion script performs these
integer convolutions, not a rank calculation on $N^2$-dimensional arrays.
\end{proof}

V21's isotropically noisy benchmark has $J_\epsilon\succ0$ at every
$\epsilon>0$ and is exactly $G_{\det}$-covariant.  The bound applies when
\emph{inferring} its symmetry from otherwise unrestricted linear data, even
though its group is known if that entire noise model is independently supplied.
It does not apply to quadratic Choi-purity comparisons, prior exact covariance,
a known Kraus-support face, or a restricted independently justified source
family.  It is not a sample-complexity or experimental running-time estimate.

\section{Two scalar comparisons certify the unchanged native source}
\label{sel22:two}

Keep the actual V19--V21 native five-copy representation $R_5$ and the seven
orthonormal vectors
\[
 n,\ \omega,\ u_0,\ldots,u_4.
\]
The six vectors other than $\omega$ are neutral, while
$R_5(U,V)\omega=\chi(U,V)\omega$, with
$\chi(U,V)=\det U\det V$.  These are the same identified vectors as in the
benchmark, not fitted eigenvectors of an unknown channel.  Let $Q$ project
onto their seven-dimensional span and put $Q_b=I-Q$, $m=N-7>0$.
Consider the unital represented algebra
\begin{equation}
 \mathcal A_Q=\mathcal B(Q\mathcal H)\oplus\C Q_b,
 \qquad \dim_\C\mathcal A_Q=50.
 \label{sel22:eq:AQ}
\end{equation}
Every element of $\mathcal A_Q$ commutes with $G_{\det}$ on the full native
carrier.  The scalar bulk is retained rather than replaced by one physical
state.

Let $P_Q$ be the Hilbert--Schmidt projection onto
$\operatorname{vec}\mathcal A_Q$.  In any orthonormal basis $(v_a)_{a=1}^7$
of $Q\mathcal H$, write $F_{ab}=|v_a\rangle\langle v_b|$.  Then
\begin{equation}
 P_Q=\sum_{a,b=1}^7|F_{ab}\rangle\!\rangle\langle\!\langle F_{ab}|
       +\frac{|Q_b\rangle\!\rangle\langle\!\langle Q_b|}{N-7}.
 \label{sel22:eq:PQ}
\end{equation}
This projection is independent of that basis.  Put
\[
 e=|I_N\rangle\!\rangle/\sqrt N,\quad
 f=|\,|\omega\rangle\langle n|\,\rangle\!\rangle,\quad
 X_\circ=|f\rangle\langle e|+|e\rangle\langle f|.
\]
The vectors $e,f$ are orthonormal, so $\|X_\circ\|=1$ and
$(I+X_\circ)/2$ is an effect.  Define two real scalar data of an otherwise
unknown normalized channel:
\begin{equation}
 \ell=\Tr[(I-P_Q)\rho_J],\qquad
 \mu=\Tr(X_\circ\rho_J).
 \label{sel22:eq:two-data}
\end{equation}
Thus $\ell$ is a probability and $(1+\mu)/2$ is a probability in a second
binary comparison.  No Kraus outcomes of the unknown channel are resolved.

\begin{theorem}[Two-scalar native determinant certificate]
\label{sel22:thm:two}
On the preceding independently identified native carrier and comparison bank,
every CPTP map satisfying
\begin{equation}
                    \boxed{\ell=0,\qquad \mu\ne0}
 \label{sel22:eq:two-test}
\end{equation}
has covariance group inside $G_0$ exactly $G_{\det}$.
No nonabelian or central covariance of the unknown map is an additional
premise, and no upper bound on its Choi rank is assumed.
\end{theorem}
\begin{proof}
The positive null measurement confines $J$ to $P_Q$.  In every Kraus
factorization, each coefficient consequently belongs to $\mathcal A_Q$.
It therefore commutes with every $R_5(h)$, $h\in G_{\det}$.  This proves
inclusion of that entire group in the channel covariance, without twirling
or changing the channel.

Every covariance fixes the identity Choi moment
$A=N^{-1}\operatorname{unvec}(Ji)$, by V21's argument or directly from
$\mathcal R(g)i=i$.  The two data satisfy
\begin{equation}
 \mu=\frac{2}{\sqrt N}\operatorname{Re}\langle\omega|A|n\rangle.
 \label{sel22:eq:mu-anchor}
\end{equation}
Thus this matrix element is nonzero.  Under $R_5(g)$ it is multiplied by
$\chi(g)$, so invariance of $A$ forces $\chi(g)=1$.  This excludes every
larger covariance group inside $G_0$.  Both implications concern every
positive channel with the same two data, not just the old benchmark.
\end{proof}

The ratio $-2:3$ follows here from the \emph{independently identified}
three-colour/two-weak alternating tensor $\omega$.  This is a data-reduction
result, not a new derivation of that tensor inventory.  The support tester
is a strong physical comparison: it certifies a whole-domain amplitude
restriction from a null probability.  Merely knowing that $P_Q$ is a
mathematical projector does not admit its measurement.

\begin{proposition}[The original benchmark passes without further reconstruction]
\label{sel22:prop:benchmark}
For exactly the V21 forgotten source, with $N=537824$, $\tau=1/8$, and
\[
 d_0=\sqrt{4095}/64,\qquad x=N-5+4d_0,
\]
the two comparisons are
\begin{equation}
 \boxed{\ell_*=0,\qquad
       \mu_*=-\frac{2(x-1)}{3N\sqrt N}\ne0.}
 \label{sel22:eq:benchmark}
\end{equation}
The second signal is about $-9.09\,10^{-4}$; its displayed exact expression,
not this rounded value, is used in the proof.
\end{proposition}
\begin{proof}
Each original $K_e$ acts by a seven-dimensional matrix on $Q\mathcal H$
and by a scalar on $Q_b\mathcal H$, so its Choi vector lies in $P_Q$.
This proves $\ell_*=0$ without changing the coefficients.  The inherited
formula $A=I/3+2\overline b B/(3N)$ has $b=x-i$ and
$\langle\omega|B|n\rangle=-(1+i)/2$.  Consequently
\[
 \langle\omega|A|n\rangle
       =-\frac{(x-1)+i(x+1)}{3N}.
\]
Equation~\eqref{sel22:eq:mu-anchor} proves the result.  In particular the
full $N-7$ bulk trace has not been dropped from $x$ or from normalization.
\end{proof}

\begin{proposition}[The certificate leaves many channels undetermined]
\label{sel22:prop:nonunique}
The face of normalized channels with $\ell=0$ has affine real dimension
$50^2-50=2450$.  The fibre with the additional exact benchmark value
$\mu=\mu_*$ has affine real dimension $2449$.  Every one of its channels
nevertheless has the group specified in Theorem~\ref{sel22:thm:two}.
\end{proposition}
\begin{proof}
Hermitian matrices supported on $P_Q$ have real dimension $50^2$.
Their output partial traces lie in $\overline{\mathcal A_Q}$, of real
Hermitian dimension $50$, and the partial trace onto it is surjective:
for Hermitian $A\in\mathcal A_Q$, the Hermitian Choi matrix
$(|A\rangle\!\rangle\langle\!\langle I|+
|I\rangle\!\rangle\langle\!\langle A|)/2$ has partial trace
$\overline A$.  A strictly supported normalized point exists, with Kraus
coefficients $|v_a\rangle\langle v_b|/\sqrt7$ and $Q_b$.  Hence
Lemma~\ref{sel22:lem:face} gives dimension $2450$.

Take the unitaries equal to $\pm\sigma_x$ on $\Span\{n,\omega\}$ and
identity elsewhere.  Their channels lie on this face and have
$\mu=\pm2(N-2)/(N\sqrt N)$.  The value \eqref{sel22:eq:benchmark} lies
strictly between them, since $0<x-1<N-2$.  Mix a small positive amount of
the strictly supported point into an appropriate convex combination of
these two unitary channels, adjusting the combination to keep $\mu=\mu_*$.
This gives a full-support point of the constrained face.  The second
functional is nonconstant there and removes exactly one affine dimension.
\end{proof}

\section{Two scalars are minimal, and faithful noise exposes unseen alternatives}
\label{sel22:minimal}

\begin{theorem}[One linear scalar cannot certify the benchmark's exact group]
\label{sel22:thm:one-no-go}
Let $J_*$ be the unchanged V21 Choi matrix.  With normalization and the native
representation fixed, no single real linear expectation $f(J_*)$ has the
property that \emph{all} CPTP completions with that value have covariance
group exactly $G_{\det}$.  Thus two independent real linear observations
are necessary and sufficient at this particular benchmark, in the unrestricted
one-use comparison class of Section~\ref{sel22:context}.
\end{theorem}
\begin{proof}
The possible values of $f$ on the compact convex CPTP set form a closed
interval.  If $f(J_*)$ is an endpoint, use
$J_{\id}=|I\rangle\!\rangle\langle\!\langle I|$.
The identity vector belongs to $\Ran J_*$ by the inherited formula for $S$.
Thus $J_*$ is strictly positive on its support, both matrices obey the same
partial trace, and $J_*\pm t(J_{\id}-J_*)$ remain CPTP for small positive
$t$.  Extremality of the value implies $f(J_{\id})=f(J_*)$.
The identity channel has the whole $G_0$ covariance and is a compatible
counterexample.  This also covers a constant functional.

If the value lies strictly inside the interval, a strictly positive CPTP
matrix with the same value exists.  Indeed, first mix a sufficiently small
amount of the strictly positive depolarizing Choi matrix into a convex
combination of two feasible matrices with values on opposite sides, then
adjust their mixture weights to match the given value.  Denote that point
by $J_\circ$.  If it is not $G_{\det}$-covariant there is already a
counterexample.  Otherwise the one-scalar invisible tangent has codimension
at most one in $\mathcal T_N$.

The $G_{\det}$-fixed tangent has codimension at least two.  To see this
without a large dimension calculation, choose a native defining colour copy
in the bulk, two vectors $v,w$ in that copy, and the neutral input $n$.
The two Hermitian tangents
\[
 (|v\rangle\langle w|+|w\rangle\langle v|)\otimes
                  |\overline n\rangle\langle\overline n|,
 \qquad
 i(|v\rangle\langle w|-|w\rangle\langle v|)\otimes
                  |\overline n\rangle\langle\overline n|
\]
have zero output partial trace and form a rotating real plane under an
$\mathrm{SU}(3)$ diagonal subgroup.  Their span meets the fixed space only
at zero.  Thus some invisible tangent is not fixed.  A sufficiently small
perturbation of $J_\circ$ along it remains positive, preserves the one scalar,
and fails covariance by Theorem~\ref{sel22:thm:universal}.  This proves
impossibility in both cases.  Sufficiency of two scalars is already proved.
\end{proof}

\begin{theorem}[Exactly identical two-scalar data can hide broken symmetry after noise]
\label{sel22:thm:noise-completion}
Keep V21's explicitly declared isotropic-noise comparison
$\Phi_\epsilon=(1-\epsilon)\Phi_*+\epsilon\mathcal T$, $0<\epsilon<1$.
Its two data are
\begin{equation}
 \ell_\epsilon=\epsilon(1-50/N^2),\qquad
 \mu_\epsilon=(1-\epsilon)\mu_*.
 \label{sel22:eq:noisy-data}
\end{equation}
Choose orthogonal $v,w$ in one native colour copy orthogonal to $Q$ and put
\[
 A_b=|v\rangle\langle w|+|w\rangle\langle v|,\qquad
 \Delta(X)=\langle n|X|n\rangle A_b.
\]
For every $0<|t|\le\epsilon/N$, the map
$\widetilde\Phi_{\epsilon,t}=\Phi_\epsilon+t\Delta$ is CPTP and has
\emph{exactly} the same two data.  It fails $\mathrm{SU}(3)$ covariance and
satisfies the sharp distance identity
\begin{equation}
 \boxed{
 \inf_{\Psi\ {G_{\det}\text{-covariant CPTP}}}
 \|\widetilde\Phi_{\epsilon,t}-\Psi\|_\diamond=2|t|.}
 \label{sel22:eq:noise-distance}
\end{equation}
\end{theorem}
\begin{proof}
The noisy normalized Choi state is
$(1-\epsilon)\rho_*+\epsilon I_{N^2}/N^2$.  Since $\rank P_Q=50$ and
$\Tr X_\circ=0$, this gives \eqref{sel22:eq:noisy-data}.
The perturbation Choi matrix is
$Z=A_b\otimes|\overline n\rangle\langle\overline n|$, of operator norm
one and zero output partial trace.  The faithful Choi floor $\epsilon/N$
therefore makes $J_\epsilon+tZ$ positive at the stated values of $t$.
Every vector in its support has bulk output and the active input $n$,
which is orthogonal to $\operatorname{vec}\mathcal A_Q$.  Thus $P_QZ=0$;
also $Z e=Zf=0$ and $\Tr Z=0$.  Both measured expectations are unchanged.

Choose $h\in\mathrm{SU}(3)\subset G_{\det}$ acting on $v,w$ by $i,-i$.
It fixes $n$ and sends $A_b$ to $-A_b$.  The original noisy channel is
covariant, so the conjugation defect of the new channel is $2t\Delta$.
The diamond norm of $\Delta$ is exactly $\|A_b\|_1=2$: the upper bound
follows from the norm-one input functional and preparation map, and input
$|n\rangle\langle n|$ attains it.  Hence the defect is $4|t|$.
Distance to any covariant channel is at least half this defect by the
triangle inequality and unitary invariance of the diamond norm.  The original
$\Phi_\epsilon$ is a covariant comparison at distance $2|t|$, proving equality.
The perturbation is a comparison of different positive channels with identical
partial data, not a perturbation assigned to the historical source.
\end{proof}

\section{A controlled comparison for a small support error}
\label{sel22:robust}

Exact symmetry is a closed constraint and cannot be inferred by rounding a
positive leakage to zero.  A quantitatively nearby symmetric channel is a
different conclusion, for which a complete normalization argument is needed.
Simply projecting a Choi matrix does not generally preserve trace.

\begin{theorem}[Support leakage admits a normalized algebraic comparison]
\label{sel22:thm:repair}
Let $\mathcal A\subset M_N$ be any represented unital $C^*$-algebra, let
$P_{\mathcal A}$ be the HS projection onto its vectorized space, and set
$\ell=\Tr[(I-P_{\mathcal A})J(\Phi)/N]$ for an actual CPTP channel.
There exists a CPTP comparison $\widehat\Phi$ with all Kraus coefficients
in $\mathcal A$ such that
\begin{equation}
 \boxed{\|\Phi-\widehat\Phi\|_\diamond
                           \le\min\{2,4\sqrt{N\ell}\}.}
 \label{sel22:eq:repair}
\end{equation}
If $\mathcal A=\mathcal A_Q$ and
$|\mu|>4\sqrt{N\ell}$, such a comparison can be chosen with covariance
group exactly $G_{\det}$.  This certifies proximity to a channel with that
group, not equality of the actual source's group.
\end{theorem}
\begin{proof}
Choose any Kraus list $K_j$ and let $\widetilde K_j$ be its HS projections
onto $\mathcal A$.  Form the coefficient columns
$V\psi=(K_j\psi)_j$ and $\widetilde V\psi=(\widetilde K_j\psi)_j$.
Then $V^*V=I$ and
\[
 \|V-\widetilde V\|\le
  \left(\sum_j\|K_j-\widetilde K_j\|_{\HS}^2\right)^{1/2}
  =\sqrt{N\ell}=:e.
\]
If $e<1$, the singular values of $\widetilde V$ lie between $1-e$ and
$1+e$.  Put $B=\widetilde V^*\widetilde V\in\mathcal A$ and
$\widehat V=\widetilde V B^{-1/2}$.  Inverse functional calculus inside
the finite unital algebra gives
$\widehat K_j=\widetilde K_jB^{-1/2}\in\mathcal A$ and
$\sum_j\widehat K_j^*\widehat K_j=I$.
The polar singular values give
$\|\widehat V-\widetilde V\|\le e$, hence
$\|V-\widehat V\|\le2e$.  Expanding the two isometric output maps on a
system and arbitrary reference gives diamond distance at most
$2\|V-\widehat V\|\le4e$; tracing the Kraus register is contractive.
For $e\ge1/2$, the identity channel in $\mathcal A$ also gives the general
bound two.  These cases prove the minimum in \eqref{sel22:eq:repair}.

For the last assertion, $|\mu|\le1$ makes its strict inequality imply
$e<1/4$, so the polar comparison is available.  Normalized Choi-state trace
distance is bounded by channel diamond distance.  Therefore
$|\widehat\mu-\mu|\le4e$, and $\widehat\mu\ne0$.
Its support leakage is zero.  Apply Theorem~\ref{sel22:thm:two}.
The comparison coefficients depend on a Kraus frame only up to the usual
isometric change, because HS projection is linear and $B$ is frame independent.
No normalization filter is asserted to be executed by the original source.
\end{proof}

\begin{proposition}[A vanishing normalized leakage need not control an active input]
\label{sel22:prop:dimension}
Keep $Q$ as above and let $m=N-7\ge3$.  For a bulk defining-colour unit vector
$v$ put $U_v=I-2|v\rangle\langle v|$.  Its unitary channel has
\begin{equation}
 \ell_v=\frac4N(1-1/m),\qquad
 \inf_{\Psi\ {G_{\det}\text{-covariant CPTP}}}
       \|\operatorname{Ad}_{U_v}-\Psi\|_\diamond\ge1.
 \label{sel22:eq:rare}
\end{equation}
Thus a dimension-free modulus tending to zero with normalized Choi leakage
cannot replace the whole-domain normalization cost in
Theorem~\ref{sel22:thm:repair} over increasing native carriers or neutral
spectator extensions.
\end{proposition}
\begin{proof}
The HS projection of $U_v$ onto $\mathcal A_Q$ is
$Q+(1-2/m)Q_b$.  The squared norm of its orthogonal residual is
$4(1-1/m)$, giving the displayed leakage.  Choose a colour rotation sending
$v$ to an orthogonal bulk vector $w$ and fixing the neutral vector $n$.
It conjugates $U_v$ to $U_w$.  The two unitary channels have diamond distance
two: on $(v+n)/\sqrt2$ their outputs are orthogonal pure states.
As above, half that covariance defect bounds the distance to every covariant
channel.  Enlarging by neutral spectators retains this test while $\ell_v$
tends to zero.  No averaging over a rare input is substituted for a
worst-case channel norm.
\end{proof}

If data bounds give $\ell\le\bar\ell$ and
$|\mu-\widehat\mu|\le\eta$, the strict condition
$|\widehat\mu|>\eta+4\sqrt{N\bar\ell}$ gives the preceding
comparison certificate.  At exact $\ell=0$, a certified nonzero $\mu$
instead gives the actual exact group.  Finite sampling of a null event alone
cannot certify its probability is exactly zero.  Acceptance normalization and
its estimation error must be retained: for a verified state-independent
acceptance weight $\vartheta$, unnormalized data scale by $\vartheta$;
conditioning does not remove that physical occurrence cost.

\section{Result and the next actual source row}
\label{sel22:status}

The missing object need not always be a complete quantum process table.  On
the verified benchmark inventory, one positive support-null comparison and
one interference expectation force the primitive determinant group for every
CPTP completion.  This result includes nonabelian covariance instead of
assuming it after an incomplete reconstruction.  It leaves a $2449$-dimensional
family of channels at the exact benchmark data, so it is a symmetry
certificate rather than disguised full tomography.  The two-scalar minimum
is specific to the unchanged rank-six benchmark and the stated one-use linear
comparison class; a different pure unitary target, prior constraints, or
nonlinear multi-copy measurements is not covered by that lower bound.

The cost has not disappeared.  The two effects are collective tests on the
whole native input/reference space.  Their construction uses the already
identified neutral vectors and the actual alternating tensor ray; those
identifications and the ability to perform the effects must not be inferred
from two classical scalar records.  Without such an admitted positive null
comparison, the completion-fibre theorem determines which additional physical
testers are needed.  A small symmetric fit does not replace that audit.
Faithful noise makes this distinction especially sharp: the same two exact
expectations have explicit positive symmetry-breaking completions.

No historical comparison outcomes have been supplied here.  The direction for
a subsequent source audit is to identify the \emph{admitted tester span} and
its positive support constraints on the same carrier, then compute its entire
compatible fibre or exhibit a separating completion.  That is a smaller and
more definite question than proposing another action with the target symmetry.
It does not identify the charged matter modules, their named type/private and
Clifford incidences, a fermionic grading, a common variational action, or a
four-dimensional continuum.  The manuscript's original same-history
source-complete requirement remains in force \cite{Duality}.

\paragraph{Verification and append-only status.}
The complete V21 source prefix is retained byte for byte.  The V22 script
checks finite positive support geometry and covariance tangents, the native
charge convolutions, the unchanged compiler's seven-dimensional working block
with its full scalar complement, the two exact expectations, the dimensions
of the supported channel face, and the explicit noisy same-data perturbation.
Small finite exact matrices verify polar normalization and the rare-input
obstruction; these are checks of the written general proofs, not numerical
rank extrapolations.  The complete supplied V21 checker is rerun separately.
No proof-assistant formalization, independent peer review, historical
measurement, or general mathematical-priority claim is made.

\begin{thebibliography}{99}
\raggedright
\bibitem[38]{sel22-BDK}
C.~H.~Baldwin, I.~H.~Deutsch and A.~Kalev,
\emph{Strictly-complete measurements for bounded-rank quantum-state tomography},
Phys.\ Rev.\ A \textbf{93} (2016), 052105.
\href{https://doi.org/10.1103/PhysRevA.93.052105}{doi:10.1103/PhysRevA.93.052105}.
Background for using positivity against completions of arbitrary rank;
the present target is a covariance group rather than full state reconstruction.
\bibitem[39]{sel22-SWW}
S.~Sremac, H.~Woerdeman and H.~Wolkowicz,
\emph{Complete Facial Reduction in One Step for Spectrahedra},
arXiv:1710.07410 (2017).
\href{https://arxiv.org/abs/1710.07410}{arXiv:1710.07410}.
Background for minimal faces and relative interior of semidefinite feasible
sets; no algorithmic complexity or one-step reduction result is imported.
\end{thebibliography}
\addtocontents{toc}{\protect\endgroup}
% END OF SOURCE-SELECTION V22 DEVELOPMENT BLOCK.
% Preserve V1--V22 and append subsequent work before the final terminator.


\clearpage
\begin{center}
{\Large\bfseries Source-selection continuation V23}\\[.4em]
{\large Native preparation--survival certificates, optimal null-test strength,\\
and a same-measure Yang--Mills fluctuation audit}
\end{center}
\addtocontents{toc}{\protect\begingroup}
\addtocontents{toc}{\protect\renewcommand{\protect\numberline}{\protect\makebox[2.1em][l]}}

\section{The new source changes the observation question}
\label{sel23:context}

V22 reduces the determinant-group certificate to a positive support-null
comparison and a nonzero interference comparison.  Its remaining measurement
hypothesis is substantial: both effects act on the full input--reference
Choi space.  We now ask whether the same conclusion can be certified by
preparing states of the original native memory, using the source once, and
measuring that memory alone.  No new source action is selected in answering
this question.

The new supplied Yang--Mills file \cite{sel23-YM86} changes the emphasis in a
useful way.  Its V85 transports the \emph{complete restricted measure}, with
its gauge determinant, Haar density, chart boundary and normalizer.  V86
uses a literal plaquette second difference to prove nonconcentration of the
nonlinear Wilson observable, not of a substitute linear field.  Its result
is at a specified logarithmic scale; the thermal-scale improvement and a
physical continuum interpretation remain separate.  We transfer that
observable-level discipline, not an alleged Yang--Mills mass gap or a
Standard-Model tensor port.

\paragraph{Audited interfaces and nonduplication.}
The retained inputs are V22's support-face theorem, its one-scalar
impossibility result, and its normalized polar comparison
(Theorems~\ref{sel22:thm:two}, \ref{sel22:thm:one-no-go},
\ref{sel22:thm:repair}).  The native seven-ray implementation is exactly
V19--V21.  The supplied YM V86 is read at its complete-density definition
\src{eq:v75-density}, global-colour equivariance in V62, the V84--V85
normalization/transport interfaces, and
\src{lem:v86-three-events}, \src{thm:v86-joint},
\src{prop:v86-thermal}, \src{lem:v86-doubling}.
None of these gives the preparation--survival certificate below.  The
existing duality notes' nondisturbance, positive realization and calibrated
commutator results are not claimed anew.  This is a targeted source audit,
not independent verification of all 86 YM bodies or of their external scripts.

Complementary-basis process tests and projective-design verification are
established methods \cite{sel23-Hofmann,sel23-ZZ,sel23-ZH}.  The results here
concern a \emph{whole channel face}, not a single target unitary: an
unrestricted active matrix block and a scalar complement must both be
retained.  We prove the associated constants, its optimal reference-free
null-test strength, the direct determinant signal, and an exact transfer of
the supplied YM fluctuation bound to an openly specified probe.

\paragraph{Scope of the access claim.}
``Without a reference ancilla'' means that no extra quantum system remains
entangled with the input while the source is used.  Classical random choices
of preparations and measurements are allowed.  States within the original
five-copy native memory may still be entangled across those five factors.
The preparation frame and its phase-sensitive output Read must actually be
admitted.  Neither a preferred basis nor a calibrated phase reference is
created by a mathematical formula.  All norms below retain the usual diamond
norm without a factor of one half, and the original channel is CPTP.

\section{A positive test made only of native preparations and survival Reads}
\label{sel23:survival}

First work on a general finite memory $\mathcal H=\C^N$ with an
independently identified projection $Q$ of rank $q\ge1$.  Write
$Q_b=I-Q$, $d=N-q\ge2$, and
\[
 \mathcal A_Q=\mathcal B(Q\mathcal H)\oplus\C Q_b.
\]
Let $P_{\mathcal A}$ be the HS projection onto its vectorized space.  For
any Kraus family $K_j$ of the actual channel, define its unnormalized
support leakage by
\begin{equation}
 L_{\mathcal A}(\Phi)=\Tr[(I-P_{\mathcal A})J_\Phi]
 =\sum_j\left\|K_j-QK_jQ-
       \frac{\Tr(Q_bK_jQ_b)}dQ_b\right\|_{\HS}^2.
 \label{sel23:eq:leakage}
\end{equation}
It is independent of the factorization of $J_\Phi$; V22's normalized leakage
is $L_{\mathcal A}/N$.

Choose an orthonormal basis $e_0,\ldots,e_{d-1}$ of the \emph{same} bulk
$Q_b\mathcal H$, and its Fourier basis
\[
 f_k=d^{-1/2}\sum_{j=0}^{d-1}e^{2\pi i jk/d}e_j.
\]
These bases are mutually unbiased in every integer dimension $d$.  Denote
$p_v=|v\rangle\langle v|$, and the HS projections onto their diagonal
operator spaces by $\mathcal D_E,\mathcal D_F$.  Directly,
\begin{equation}
 \mathcal D_E\mathcal D_F=\mathcal D_F\mathcal D_E=P_c,
 \qquad P_c(B)=\frac{\Tr B}{d}Q_b.
 \label{sel23:eq:MUB}
\end{equation}
Indeed, every matrix element $|\langle e_j,f_k\rangle|^2$ is $1/d$.
Consequently the traceless parts of the two diagonal spaces are orthogonal.

Use three kinds of one-use experiments.  In the first prepare $Q/q$ and
record whether the output is in $Q_b$.  In the other two, choose uniformly a
vector of the indicated bulk basis, prepare it, and record failure of the
\emph{same} vector's survival Read:
\begin{align}
 r_Q&=\Tr[Q_b\Phi(Q/q)],\
 r_E&=\frac1d\sum_j\Tr[(I-p_{e_j})\Phi(p_{e_j})],\
 r_F&=\frac1d\sum_j\Tr[(I-p_{f_j})\Phi(p_{f_j})].
 \label{sel23:eq:three-rates}
\end{align}
The identity in the last two effects is the full $I_N$; escape to $Q$ is a
failure, not an outcome discarded before normalization.  Put $D=2d+q$ and
randomize the three contexts with respective probabilities $q/D,d/D,d/D$.
The single observed failure probability is
\begin{equation}
                  p_{\rm sb}=\frac{q r_Q+d r_E+d r_F}{D}.
 \label{sel23:eq:failure}
\end{equation}
The mixture probabilities are part of the specified experiment.  Computing
this weighted average from separately measured rates gives the same datum.

\begin{theorem}[Support exactness and the complete two-basis strength]
\label{sel23:thm:survival}
For every CPTP map on this identified memory,
\begin{equation}
 \boxed{p_{\rm sb}=0\quad\Longleftrightarrow\quad
       \Ran J_\Phi\subseteq\operatorname{vec}\mathcal A_Q.}
 \label{sel23:eq:null}
\end{equation}
No Choi-rank, unitality or coefficient-independence assumption is needed.
Quantitatively,
\begin{equation}
 \boxed{\frac{L_{\mathcal A}(\Phi)}D\le p_{\rm sb}
                   \le\frac{2L_{\mathcal A}(\Phi)}D.}
 \label{sel23:eq:strength}
\end{equation}
As a positive tester on coefficient space its eigenvalues are zero on
$\operatorname{vec}\mathcal A_Q$, $1/D$ with multiplicity
$qd+2(d-1)$, and $2/D$ with multiplicity $qd+(d-1)^2$.
\end{theorem}
\begin{proof}
For any prepared unit vector $v$,
\[
 \Tr[(I-p_v)\Phi(p_v)]=\sum_j\|(I-p_v)K_jv\|^2.
\]
Every term is nonnegative.  Vanishing of the two bulk-basis rates says that
each $K_j$ takes every $e_k$ and every $f_k$ into its own line.  In particular
there is no $QK_jQ_b$ block.  Its bulk restriction belongs to both diagonal
algebras, whose intersection by \eqref{sel23:eq:MUB} is $\C Q_b$.
Vanishing of $r_Q$ separately gives $Q_bK_jQ=0$.  This proves membership
in $\mathcal A_Q$, with no condition on $QK_jQ$.  The converse follows by
substitution.  The range statement is equivalent to this coefficient
membership for any Kraus factorization.

For the constants, fix a coefficient $K$ and put $B=Q_bKQ_b$ and
$B_0=B-(\Tr B/d)Q_b$.  Its contribution to $D p_{\rm sb}$ is exactly
\begin{equation}
 \|Q_bKQ\|_{\HS}^2+2\|QKQ_b\|_{\HS}^2+
 2\|B_0\|_{\HS}^2-\|\mathcal D_EB_0\|_{\HS}^2
                         -\|\mathcal D_FB_0\|_{\HS}^2.
 \label{sel23:eq:exact-form}
\end{equation}
This follows by summing $\|Kv\|^2-|\langle v,Kv\rangle|^2$ in each
basis.  The scalar part cancels exactly.  On traceless bulk operators,
$\mathcal D_E+\mathcal D_F$ is an orthogonal projection of rank $2(d-1)$.
Thus the last three terms have eigenvalue one on its range and two on its
orthogonal complement.  The two rectangular blocks are orthogonal to these
spaces and to one another.  Their dimensions are each $qd$.  This proves
the stated spectrum and both inequalities, after summing the coefficients.
\end{proof}

A single bulk basis would not suffice: a nonscalar diagonal unitary on that
basis, extended by the identity on $Q$, passes its survival test and the
$Q$-escape test exactly.  Its Fourier-basis failure is nonzero.  This is an
access distinction, not another invocation of source recurrence.

\section{The optimal reference-free null-test budget}
\label{sel23:optimal}

A general randomized preparation--measurement failure test, without a retained
quantum reference, has the positive tester
\begin{equation}
 T=\sum_\alpha p_\alpha M_\alpha\otimes\overline{\rho_\alpha},
 \quad p_\alpha\ge0,\quad\sum_\alpha p_\alpha=1,\quad
 \rho_\alpha\succeq0,\quad\Tr\rho_\alpha=1,\quad0\preceq M_\alpha\preceq I.
 \label{sel23:eq:tester}
\end{equation}
Here $\Tr(TJ_\Phi)$ is the failure probability, in the output-first Choi
convention.  Impose $T\operatorname{vec}\mathcal A_Q=0$.  Its uniform
leakage strength is the largest $\gamma$ such that
$T\succeq\gamma(I-P_{\mathcal A})$.  This is a coefficient-space comparison
constant, not a physical contraction rate or a sample-complexity optimum.

\begin{theorem}[Sharp strength among all such positive null tests]
\label{sel23:thm:optimal}
For $d\ge2$ the maximal possible strength is
\begin{equation}
                         \boxed{\gamma_{\rm opt}=\frac1{N+1}.}
 \label{sel23:eq:optimum}
\end{equation}
It is attained by using the $Q$-input escape test with probability $q/(N+1)$
and a projective two-design bulk survival test with probability
$(d+1)/(N+1)$.  The two-basis test of
Theorem~\ref{sel23:thm:survival} achieves $1/(2d+q)$, within a factor of two
of this optimum, using only two bulk bases.
\end{theorem}
\begin{proof}
Let $r_b=\sum_\alpha p_\alpha\Tr(Q_b\rho_\alpha)$ and $r_Q=1-r_b$.
Compress $T\succeq\gamma(I-P_{\mathcal A})$ to output-bulk/input-active
space and take traces.  Its dimension is $dq$, whereas
\[
 \Tr[(Q_b\otimes\overline Q)T]
 =\sum_\alpha p_\alpha\Tr(Q_bM_\alpha)\Tr(Q\rho_\alpha)
 \le d r_Q.
\]
Hence $\gamma q\le r_Q$.

Since $T$ kills $|Q_b\rangle\!\rangle$, positivity makes each summand in
\eqref{sel23:eq:tester} have zero expectation there:
$\Tr(M_\alpha Q_b\rho_\alpha Q_b)=0$.
When $\Tr(Q_b\rho_\alpha)>0$, the positive contraction
$Q_bM_\alpha Q_b$ must vanish on a nonzero subspace and has trace at most
$d-1$.  When that weight is zero, its contribution is zero.  Therefore
\[
 \Tr[(Q_b\otimes\overline{Q_b})T]\le(d-1)r_b.
\]
On the traceless bulk coefficient space the lower bound has trace
$\gamma(d^2-1)$, so $\gamma(d+1)\le r_b$.
Adding gives $\gamma(N+1)\le1$.  This upper bound allows arbitrary mixed
preparations, output effects and finite classical randomization.

For attainment take an ensemble of bulk pure states with second moment
$\mathbb E(p_\psi\otimes p_\psi)=(I+\mathrm{Swap})/[d(d+1)]$.
Its survival-failure tester is
\[
 \frac1d Q\otimes\overline{Q_b}
 +\frac1{d+1}(I_{bb}-P_c),
 \quad I_{bb}=Q_b\otimes\overline{Q_b},\quad
 P_c=|Q_b\rangle\!\rangle\langle\!\langle Q_b|/d.
\]
Combining with the indicated active-input probability gives
\begin{equation}
 T_{\rm opt}=\frac1{N+1}\left[
 Q_b\otimes\overline Q+\frac{d+1}{d}Q\otimes\overline{Q_b}
                            +I_{bb}-P_c\right].
 \label{sel23:eq:Topt}
\end{equation}
Its kernel is precisely $\operatorname{vec}\mathcal A_Q$, and its least
complementary eigenvalue is $1/(N+1)$.

An exact finite weighted ensemble exists in every dimension, without assuming
a complete collection of mutually unbiased bases.  With total weight
$1/(d+1)$ use the computational basis uniformly; with total weight
$d/(d+1)$ use uniformly the vectors
$d^{-1/2}\sum_j z_j e_j$, where the $z_j$ are independent fourth roots of
unity.  A common phase is redundant.  The latter second moment is
$[I+\mathrm{Swap}-\sum_j|jj\rangle\langle jj|]/d^2$; the basis contribution
cancels its last term with the stated weights.  This proves the design
identity directly.  The ensemble can be very large; no efficient preparation
claim is made.  Finally $(N+1)/(2d+q)>1/2$ proves the two-basis comparison.
\end{proof}

The optimum does not remove the dimension cost.  It reallocates the input
sampling mass optimally between the active block and the bulk.  In particular
it does not identify a small mean failure with a small worst-input error
without paying a factor depending on $N$.

\section{A direct native interference replaces the diluted Choi signal}
\label{sel23:group}

Return to the unchanged five-copy native inventory of V19--V22.  Thus $q=7$,
$N=14^5$, the orthonormal vectors $n,u_0,\ldots,u_4$ are neutral under
$G_0=\mathrm U(3)\times\mathrm U(2)$, and $R(g)\omega=\chi(g)\omega$ with
$\chi(g)=\det U\det V$.  Prepare $p_n$ on the native memory and measure
\begin{equation}
 Y_\circ=i(|n\rangle\langle\omega|-|\omega\rangle\langle n|),
 \qquad \xi=\Tr[Y_\circ\Phi(p_n)].
 \label{sel23:eq:xi}
\end{equation}
This Hermitian operator has norm one.  Thus $(1+\xi)/2$ is a binary output
probability, including the complement where the effect is $I/2$.
It is a specified phase-sensitive Read on the native tensor sector, not a
between-outcome Kraus coherence experiment.

\begin{theorem}[Reference-ancilla-free determinant certificate]
\label{sel23:thm:group}
For every CPTP channel on the independently identified native carrier, either
of the preceding support tests gives
\begin{equation}
 \boxed{p_{\rm sb}=0,\quad\xi\ne0
       \quad\Longrightarrow\quad
       \operatorname{Cov}_{G_0}(\Phi)=\Ker\chi.}
 \label{sel23:eq:group}
\end{equation}
The same statement holds with $p_{\rm opt}=\Tr(T_{\rm opt}J_\Phi)$ in place
of $p_{\rm sb}$.  It holds for \emph{every} positive trace-preserving
completion of the two observed expectations.  Nonabelian covariance is not
a separate premise, and the individual source labels need not be recovered.
\end{theorem}
\begin{proof}
The support null places all coefficients in $\mathcal A_Q$.  Every such
operator commutes with $R(h)$ for $h\in\Ker\chi$: that group acts as the
identity on $Q\mathcal H$, while the coefficient is scalar on the retained
bulk.  Hence $\Ker\chi\subseteq\operatorname{Cov}(\Phi)$.

If $g$ is any covariance of the whole channel, the invariant input $p_n$
gives $R(g)\Phi(p_n)R(g)^*=\Phi(p_n)$.  Since $\xi\ne0$, the matrix element
$\langle\omega|\Phi(p_n)|n\rangle$ is nonzero.  Conjugation multiplies it
by $\chi(g)$, so $\chi(g)=1$.  This proves the opposite inclusion.
No representative of a Kraus ray has been chosen in either direction.
\end{proof}

\begin{proposition}[The same source has a dimension-independent signal]
\label{sel23:prop:benchmark}
For exactly the old benchmark, without changing its channel or native dimension,
\begin{equation}
 \boxed{p_{\rm sb}=p_{\rm opt}=0,\qquad \xi_*=\frac23.}
 \label{sel23:eq:benchmark}
\end{equation}
The binary interference outcome has probability $5/6$.
The support-compatible channel face still has affine dimension $2450$, and
fixing $\xi=2/3$ leaves dimension $2449$.
\end{proposition}
\begin{proof}
Support exactness is inherited from the original coefficient blocks; both new
null testers have the same kernel as $P_{\mathcal A}$.  On $(n,\omega)$ the
unchanged polar unitary is
\[
 U_*|_{\Span\{n,\omega\}}=
 \frac12\begin{pmatrix}1-i&-1-i\\-1-i&1-i\end{pmatrix}.
\]
The helper terms annihilate $n$ and the positive polar factor is the identity
on this two-dimensional sector.  Therefore the original forgotten channel gives
\begin{equation}
 \Phi_*(p_n)|_{\Span\{n,\omega\}}
 =\frac13p_n+\frac23U_*p_nU_*^*
 =\begin{pmatrix}2/3&i/3\\-i/3&1/3\end{pmatrix}.
 \label{sel23:eq:output}
\end{equation}
There is zero output elsewhere.  Pairing with
$Y_\circ=\left(\begin{smallmatrix}0&i\\-i&0\end{smallmatrix}\right)$ gives
$2/3$.  Neither the $N-7$ bulk nor its trace was reassigned; it simply is not
populated by this particular input experiment.

The face dimension $2450$ is V22's proved partial-trace rank count, unchanged
by replacing its positive null effect by one with the same kernel.  A strictly
supported channel on this face has Kraus operators
$|v_a\rangle\langle v_b|/\sqrt7$ and $Q_b$ and has $\xi=0$.  A unitary on
$Q$ can take $n$ to an eigenvector of $Y_\circ$ of eigenvalue one.  Mixing
that channel with weight $2/3$ and the strictly supported channel with weight
$1/3$ gives a faithful point of the constrained face at $\xi=2/3$.
The functional is nonconstant on the face and removes exactly one dimension.
\end{proof}

The signal improvement from V22's $O(N^{-1/2})$ normalized identity-Choi
comparison to $2/3$ is obtained by changing the \emph{probe}, not by rescaling
a recorded probability.  The inherited one-scalar theorem still applies to
this benchmark, since preparation--measurement expectations are linear Choi
functionals.  Thus two independent real expectations remain minimal in that
stated one-use linear-data class.  This reuses the old minimality proof.

\section{What accuracy is paid, and which reference is still required}
\label{sel23:accuracy}

\begin{corollary}[A normalized comparison from the observed native failure rate]
\label{sel23:cor:repair}
For the two-basis test there is a CPTP comparison $\widehat\Phi$ with all
coefficients in $\mathcal A_Q$ and
\begin{equation}
 \|\Phi-\widehat\Phi\|_\diamond
       \le\delta:=\min\{2,4\sqrt{(2d+q)p_{\rm sb}}\}.
 \label{sel23:eq:repair}
\end{equation}
If $|\xi|>\delta$, that comparison has covariance exactly $\Ker\chi$.
For the optimal tester replace $(2d+q)p_{\rm sb}$ by $(N+1)p_{\rm opt}$.
These are proximity results, not exact-group conclusions for $\Phi$ at a
positive failure rate.
\end{corollary}
\begin{proof}
The strength inequalities bound $L_{\mathcal A}$ by the indicated measured
quantity.  Apply the already proved normalized polar-comparison theorem
\ref{sel22:thm:repair}, whose right side is $4\sqrt{L_{\mathcal A}}$.
A norm-one output observable on a state changes by at most the diamond
error.  Thus $|\widehat\xi|\ge|\xi|-\delta>0$; the exact support of the
comparison and Theorem~\ref{sel23:thm:group} give its group.
\end{proof}

\begin{proposition}[The old rare-input obstruction has a small actual test rate]
\label{sel23:prop:rare}
Let $v$ be the first computational bulk vector and a native defining-colour
ray.  Postcompose the unchanged source by
$\operatorname{Ad}_{V_v}$, $V_v=I-2|v\rangle\langle v|$.
The new channel still has $\xi=2/3$, but
\begin{equation}
 p_{\rm sb}=\frac{4(1-1/d)}{2d+q},\qquad
 p_{\rm opt}=\frac{4(1-1/d)}{N+1}.
 \label{sel23:eq:rare}
\end{equation}
Its diamond distance from every $\Ker\chi$-covariant channel is at least one.
\end{proposition}
\begin{proof}
The original channel is the identity on every bulk input; its helper Kraus maps vanish there.
The reflection fixes every computational bulk ray; on any Fourier ray its
survival probability is $(1-2/d)^2$.  There is no active-to-bulk escape.
The randomized protocol gives the first formula.  Its only coefficient
leakage is the bulk traceless part of $V_v$, of squared norm $4(1-1/d)$,
and \eqref{sel23:eq:Topt} gives the second.

Choose orthogonal native colour rays $w,z$ in the same defining copy as
$v$, and a colour special unitary exchanging the rays $v,w$.  Use the bulk
input $(v+z)/\sqrt2$.  The original channel acts identically there, while
the outputs of $V_v$ and $V_w$ are orthogonal.  Thus the two conjugate
channels have diamond distance two.  Every $\Ker\chi$-covariant channel is
fixed by this conjugation; the triangle inequality gives distance at least
one.  The determinant Read remains
unchanged because the reflection is the identity on $Q$.
\end{proof}

This example is a probe evaluation of the already retained rare-input
mechanism, not a new reason to infer exact symmetry from small error.
For a fixed i.i.d. channel, $m$ independent null-test trials with zero failures
exclude a failure probability $p\ge a$ with error probability at most
$(1-a)^m\le e^{-ma}$.  Thus a conservative confidence upper bound is
$\log(1/\varepsilon)/m$.  To use \eqref{sel23:eq:repair} with tolerance
$\delta_0$, one sufficient number of zero-failure trials is
\[
       m\ge16(2d+q)\delta_0^{-2}\log(1/\varepsilon),
\]
with a separately certified lower bound $|\xi|>\delta_0$.
This is not a proof of an exact zero, an adversarial-channel guarantee, or
an efficient implementation of the preparation bank.  A verified
state-independent acceptance probability $\vartheta$ gives expected raw
opportunity cost $m/\vartheta$; it must not be dropped by conditioning.

\begin{proposition}[Neutral classical data cannot provide the missing phase Read]
\label{sel23:prop:neutral}
Suppose a one-use tester bank $T_a$ is invariant under
$\mathcal R(g)=R(g)\otimes\overline{R(g)}$ for every $g\in G_0$.
For any channel $\Phi$, its $G_0$-twirl is a CPTP completion with exactly the
same measured data.  Hence this bank cannot certify a proper covariance
subgroup of $G_0$ among otherwise unrestricted CPTP completions.
\end{proposition}
\begin{proof}
The twirled Choi matrix is
$\overline J=\int\mathcal R(g)J\mathcal R(g)^*\,dg$ and is positive and
trace preserving.  Invariance and cyclicity of trace give
$\Tr(T_a\overline J)=\Tr(T_aJ)$ for each $a$.  The twirled channel has the
entire $G_0$ covariance, which proves the claim.
If preparations, controls and output effects are themselves covariant with
invariant initial data, composition, conditioning on fixed invariant labels
and classical randomization preserve this tester invariance.  Adding a
neutral classical reward does not change that argument.
\end{proof}

In the new certificate $Y_\circ$ is expressly \emph{not} invariant under the
whole $G_0$.  Its admission remains the coherent reference requirement.
Removing an entangled reference ancilla does not remove this operational
requirement.  In particular, nonconcentration of a scalar Wilson trace does
not manufacture that quantum Read.  This is an explicit scope boundary of
the transfer from the new YM source, not a defect of its fluctuation theorem.

\section[A same-measure Yang--Mills fluctuation probe]{The complete YM law gives a different, symmetry-preserving support failure}
\label{sel23:YM}

We next apply the supplied result to the literal nonlinear holonomy, with no
Gaussian replacement.  In this section $\mu=\mu_{0,\beta}$ is exactly
\cite{sel23-YM86}'s fixed-radius, mean-zero Coulomb law,
\begin{equation}
 d\mu(a)=Z_\beta^{-1}\mathbf1_{\Omega_R}(a)
       \det F(a)e^{-\beta S(a)-\mathcal H(a)}\,da.
 \label{sel23:eq:YMlaw}
\end{equation}
All determinant, Haar, domain and normalization factors are retained.
Let $U_p(a)\in\mathrm{SU}(2)$ be its actual plaquette product and
$s_p=1-\tfrac12\operatorname{Re}\Tr U_p\in[0,2]$.
For fixed $0<R\le1/4$, V86 proves, along
\begin{equation}
 M,\beta_M\longrightarrow\infty,\qquad
 \frac{\sqrt{\beta_M}}{M^2\sqrt{\log(2+\beta_M)}}\longrightarrow0,
 \label{sel23:eq:YMwindow}
\end{equation}
with $H_\beta=\log(2+\beta)$ and its specified $\sigma_R>0$,
\begin{equation}
 \mathbb E s_p\ge\frac{\sigma_R^2}{2\beta H_\beta},\quad
 \operatorname{Var}s_p\ge\frac{\sigma_R^4}{4\beta^2H_\beta^2},\quad
 \sup_{|I|\le2\sigma_R^2/(\beta H_\beta)}\mu\{s_p\in I\}\le\frac34.
 \label{sel23:eq:inherited-YM}
\end{equation}
These are inherited same-law estimates, not new YM assertions.  They do not
provide thermal-scale tightness or a physical Hamiltonian gap.

\paragraph{The added probe is declared, not attributed to the old source.}
Choose an isometric native weak doublet $\iota:\C^2\to Q_b\mathcal H$ on
which $G_0$ acts as its defining $V$, with the other native factors neutral.
Such a subrepresentation exists in the retained five-copy inventory and is
orthogonal to $Q$, whose vectors are nonabelian singlets.  Denote its
projection by $P_w$.  Consider the \emph{specified comparison probe}
\begin{equation}
 V_p(a)=I-P_w+\iota U_p(a)\iota^*,\quad
 \Gamma_\mu(X)=\int V_p(a)XV_p(a)^*\,d\mu(a),\quad
 \Phi_\mu=\Gamma_\mu\Phi_*.
 \label{sel23:eq:YMprobe}
\end{equation}
The field sample is independent of the quantum input.  Its coupling to the
native quantum memory is additional declared access.  The supplied YM file does not assert its
execution.  This definition changes neither the classical input law nor the
old $\Phi_*$; it specifies a new composite comparison channel and does not
claim equality with either original process.  Statements below concern its
field-forgotten quantum channel, not fixed-label covariance of each field
sample.

\begin{theorem}[Actual plaquette fluctuations can preserve the determinant group]
\label{sel23:thm:YMgroup}
The law \eqref{sel23:eq:YMlaw} at zero background makes $\Gamma_\mu$
$G_0$-covariant.  Consequently
\begin{equation}
 \boxed{\operatorname{Cov}_{G_0}(\Phi_\mu)=\Ker\chi,
                                  \qquad \xi(\Phi_\mu)=2/3.}
 \label{sel23:eq:YMgroup}
\end{equation}
Nevertheless its leakage outside the support face of V22 is exactly
\begin{equation}
 \boxed{L_{\mathcal A}(\Phi_\mu)
                 =4\mathbb E_\mu\bigl[s_p(1-s_p/d)\bigr].}
 \label{sel23:eq:YMleak}
\end{equation}
For the optimal native null test the conditional failure probability at field
$a$ is
\begin{equation}
             \pi(a)=\frac{4s_p(a)(1-s_p(a)/d)}{N+1}.
 \label{sel23:eq:pi}
\end{equation}
In particular positive support leakage need not mean loss of the actual
covariance group.
\end{theorem}
\begin{proof}
A constant colour conjugation preserves the Coulomb and mean-zero linear
space, the radius constraints, and its Lebesgue measure.  The supplied gauge
matrix transforms by orthogonal conjugation, so its positivity domain and
determinant are preserved.  The Haar factors depend on link norms and the
Wilson action is invariant.  Thus the \emph{complete} normalized law is
invariant, not only a quadratic approximation to it.  The plaquette product
transforms as $U_p\mapsto gU_pg^*$.

For a physical weak unitary $V\in\mathrm U(2)$, its scalar phase cancels in
conjugation and its special-unitary part induces exactly this transformation.
All other comparison-group actions commute with the specified weak-doublet
probe.  Change variables in \eqref{sel23:eq:YMprobe} to obtain $G_0$
covariance of $\Gamma_\mu$.  It fixes every operator supported on $Q$ because
each $V_p(a)$ is the identity there.  Thus it retains \eqref{sel23:eq:output}
and its nonzero $\xi$.  Composition preserves the $\Ker\chi$ covariance of
$\Phi_*$.  The same invariant-input argument as in
Theorem~\ref{sel23:thm:group} excludes any larger group.

Every coefficient of $\Phi_*$ is an arbitrary $Q$ block plus $k_jQ_b$.
Trace preservation on a bulk input gives $\sum_j|k_j|^2=1$.
After the probe its only component outside $\mathcal A_Q$ is
$k_j[V_p|_{Q_b}-(\Tr V_p|_{Q_b}/d)Q_b]$.
Now $V_p|_{Q_b}$ is unitary and
\[
              \Tr V_p|_{Q_b}=d-2+\Tr U_p=d-2s_p.
\]
Its squared traceless HS norm is $d-(d-2s_p)^2/d$.
Sum $|k_j|^2$ and integrate to prove \eqref{sel23:eq:YMleak}.
There are no active--bulk rectangular coefficients in this probe, so
\eqref{sel23:eq:Topt} pairs with this leakage by exactly $1/(N+1)$.
This proves \eqref{sel23:eq:pi}, without changing the sampling law.
\end{proof}


\begin{proposition}[Two literal Wilson moments determine the entire probe channel]
\label{sel23:prop:YMchannel}
Let $m_1=\mathbb E_\mu W_p$, $m_2=\mathbb E_\mu W_p^2$, and $P=P_w$,
$C=I-P$.  The complete field-forgotten probe is exactly
\begin{equation}
 \Gamma_\mu(X)=CXC+m_1(CXP+PXC)
       +\frac{4m_2-1}{3}PXP
       +\frac{2(1-m_2)}3\Tr(PX)P.
 \label{sel23:eq:YMchannel}
\end{equation}
In particular, for the unchanged base source,
\begin{equation}
 \boxed{\mathbb E s_p\le\|\Phi_\mu-\Phi_*\|_\diamond
       \le\|\Gamma_\mu-\id\|_\diamond
       \le\min\{2,5\mathbb E s_p\}.}
 \label{sel23:eq:YMnorm}
\end{equation}
The constants do not contain the native dimension $N$.
\end{proposition}
\begin{proof}
Write the actual holonomy as $U_p=wI+i\sum_{a=1}^3x_a\sigma_a$ with
$w=W_p$ real and $|x|^2=1-w^2$.  Simultaneous colour conjugation rotates
$x$ and leaves $w$ fixed.  It implies
$\mathbb Ex_a=\mathbb E(wx_a)=0$ and
$\mathbb Ex_ax_b=(1-m_2)\delta_{ab}/3$.
Expansion of $V_pXV_p^*$ gives the cross-block factor $m_1$.
On $PXP$ use $\sum_a\sigma_a Z\sigma_a=2\Tr(Z)I-Z$ to obtain
\eqref{sel23:eq:YMchannel}.  This is an evaluation of the exact group
moments of the original law, not its Gaussian approximation.

Let $\mathcal D(X)=CXC+PXP$ and
$\mathcal R_P(X)=PXP-\Tr(PX)P/2$.  The difference is
\[
 \Gamma_\mu-\id=(m_1-1)(\id-\mathcal D)
                   -\frac{4(1-m_2)}3\mathcal R_P.
\]
Here $\|\id-\mathcal D\|_\diamond\le1$, since it is half the difference
of the identity channel and conjugation by $C-P$.
Also $\|\mathcal R_P\|_\diamond\le3/2$: compression to the weak qubit is
completely contractive in trace norm, and on that qubit identity minus
complete depolarization is $(3\id-\sum_a\operatorname{Ad}_{\sigma_a})/4$.
The channel triangle inequality gives $3/2$.
Consequently the norm is at most
$(1-m_1)+2(1-m_2)\le5\mathbb E(1-W_p)$ and also at most two.
Postcomposition by a channel proves the middle bound.

For the lower bound choose a weak unit vector $v\in P\mathcal H$ and a
bulk vector $z\perp P\mathcal H$, still in $Q_b\mathcal H$.  On the input
$(v+z)/\sqrt2$ the old source is the identity.  The norm-one observable
$|v\rangle\langle z|+|z\rangle\langle v|$ has original expectation one
and new expectation $m_1$.  This proves the claimed lower bound for the
actual composite-channel difference.
\end{proof}

This dimension-free comparison uses the \emph{proved transformation law and
specific coupling} of the fluctuation.  It does not contradict V22's lower
bound for unrestricted noisy completions, or identify every perturbation
with this covariant noise model.

\begin{corollary}[A same-measure nonconcentration scale for the actual probe probability]
\label{sel23:cor:YMscale}
Along \eqref{sel23:eq:YMwindow}, eventually
\begin{align}
 \mathbb E\pi&\ge
 \frac{2(1-2/d)\sigma_R^2}{(N+1)\beta H_\beta}>0,
                              \label{sel23:eq:pi-mean}\\
 \operatorname{Var}\pi&\ge
 \frac{4(1-4/d)^2\sigma_R^4}{(N+1)^2\beta^2H_\beta^2},
                              \label{sel23:eq:pi-var}\\
 \sup_{|I|\le8(1-4/d)\sigma_R^2/[(N+1)\beta H_\beta]}
                          \mu\{\pi\in I\}&\le\frac34.
                              \label{sel23:eq:pi-nonconc}
\end{align}
For the finite two-basis protocol the weaker but explicit mean bound is
\[
 \mathbb E p_{\rm sb}(a)\ge
           \frac{2(1-2/d)\sigma_R^2}{(2d+7)\beta H_\beta}.
\]
Here the variance is over the original field law of the conditional
probability $\pi(a)$, not the shot-noise variance of one binary outcome.
\end{corollary}
\begin{proof}
The native $d=N-7$ is greater than four.  On $0\le s\le2$,
$f(s)=4s(1-s/d)/(N+1)$ satisfies
\[
 f(s)\ge\frac{4(1-2/d)}{N+1}s,\qquad
 f'(s)\ge c_d:=\frac{4(1-4/d)}{N+1}>0.
\]
Use the mean bound in \eqref{sel23:eq:inherited-YM} for the first assertion.
For independent copies $S,S'$ of $s_p$,
$\operatorname{Var}f(S)=\mathbb E|f(S)-f(S')|^2/2
\ge c_d^2\operatorname{Var}S$, proving the second.
The preimage under $f$ of any interval of width at most
$2c_d\sigma_R^2/(\beta H_\beta)$ has width at most
$2\sigma_R^2/(\beta H_\beta)$; apply the inherited concentration-function
bound.  Finally combine \eqref{sel23:eq:strength} and
\eqref{sel23:eq:YMleak} for the two-basis mean.
\end{proof}

The lower bound shows that enforcing the old \emph{strong} scalar-bulk support
would reject this symmetry-preserving fluctuating comparison at every
sufficiently large finite pair $(M,\beta_M)$.  It does not show that the
historical source contains the probe.  Nor does it upgrade the logarithmic
scale to $\beta^{-1}$: the dyadic free-energy gate left open in YM V86
remains open here.  All source and clock normalizations remain explicit.

\section{Result, verification scope, and the next useful comparison}
\label{sel23:status}

The measurement assumption has been reduced in a concrete way.  On the same
identified native carrier, an entangled input--reference state and a global
Choi interference measurement are unnecessary for the sufficient group
certificate.  A finite pair of complementary bulk frames, active-to-bulk
escape testing, and one actual neutral-input interference Read suffice.  The
unchanged source yields a probability $5/6$ in the latter experiment rather
than a signal diluted by the total native dimension.  A sharp minimax bound
records the remaining null-test strength; a small average error is not
confused with a uniform actual-input guarantee.

The new YM file also prevents an unproductive continuation: exact support on
one chosen small operator face should not be mistaken for necessary
covariance.  Its literal nonlinear holonomy supplies a same-law fluctuating
probe for which that face condition fails by a proved margin while the
covariance group remains correct.  When the source contains such fluctuations,
one should retain their transformation law or test the full compatible
covariance fibre, not project them away in order to pass the older null test.

No historical native preparation frame, phase Read, or composite YM probe has
been identified by these calculations.  Scalar Wilson rewards alone do not
supply the noninvariant quantum tester.  The next source-level task is to
identify which of these \emph{reference-ancilla-free} native preparations and
survival/phase Reads are actually admitted, and which covariant fluctuations
are part of their complete same-cut process.  This is distinct from another
construction of a target-shaped action.  Charged matter, named type/private
and Clifford incidences, grading, continuum existence and a physical mass gap
are not concluded.

\paragraph{Proof and artifact status.}
The entire V22 prefix is preserved before its final document terminator.
The new checker verifies finite complementary-frame tester spectra, the
finite two-design identity and optimal allocation, exact source output and
normalization, the rare-input rates, and literal quaternion probe formulas.
The YM fluctuation inequalities are inherited at their stated same-law
scope; their external YM checker is not supplied or claimed to have been
rerun.  The complete supplied V22 checker is run separately.  Exact finite
checks supplement, rather than replace, the all-dimension and all-channel
proofs.  There is no proof-assistant formalization, independent peer review,
new historical measurement, or general mathematical-priority claim.

\begin{thebibliography}{99}
\raggedright
\bibitem[40]{sel23-YM86}
\emph{Renewal Yang--Mills Mass-Gap Research Notes, Post-Thirteenth-Archive
Compact Lineage, V86}, supplied file
\src{Renewal_Yang_Mills_Mass_Gap_Research_Notes_V86(2).tex}.
Used at the exact-density and global-colour interfaces and at the
V84--V86 normalized transport and plaquette nonconcentration results.
The original law, logarithmic scale and remaining thermal/continuum
conditions are retained.
\bibitem[41]{sel23-Hofmann}
H.~F.~Hofmann, \emph{Complementary classical fidelities as an efficient
criterion for the evaluation of experimentally realized quantum operations},
Phys.\ Rev.\ Lett.\ \textbf{94} (2005), 160504;
\href{https://arxiv.org/abs/quant-ph/0409083}{arXiv:quant-ph/0409083}.
Background for complementary-input process verification, not a claimed new
origin of the two-basis method.
\bibitem[42]{sel23-ZZ}
H.~Zhu and H.~Zhang, \emph{Efficient verification of quantum gates with local
operations}, Phys.\ Rev.\ A \textbf{101} (2020), 042316;
\href{https://arxiv.org/abs/1910.14032}{arXiv:1910.14032}.
Background for preparation--measurement verification; no locality or
sample-efficiency theorem for the native tensor basis is imported.
\bibitem[43]{sel23-ZH}
H.~Zhu and M.~Hayashi, \emph{Optimal verification and fidelity estimation of
maximally entangled states}, Phys.\ Rev.\ A \textbf{99} (2019), 052346;
\href{https://arxiv.org/abs/1901.09772}{arXiv:1901.09772}.
Background for two-design verification.  The active-block/scalar-complement
minimax bound above is proved in the stated native tester class.
\end{thebibliography}
\addtocontents{toc}{\protect\endgroup}
% END OF SOURCE-SELECTION V23 DEVELOPMENT BLOCK.
% Preserve V1--V23 and append subsequent work before the final terminator.


\clearpage
\providecommand{\Herm}{\operatorname{Herm}}
\hypersetup{pdftitle={Renewal SM Source Selection Research Notes: V1--V24}}
\begin{center}
{\Large\bfseries Source-selection continuation V24}\\[.35em]
{\large Operational-frame closure: reconstructing the symmetry certificate\\
from the already admitted tomography rather than adding testers}
\end{center}
\addtocontents{toc}{\protect\begingroup}
\addcontentsline{toc}{part}{V24: Operational-frame closure and tester-span identifiability}

\section{Go upstream from tester availability to the declared intervention frame}
\label{sel24:context}

V22--V23 reduced the determinant-group audit to a small number of positive or
interference comparisons on one independently identified native carrier.  Their
logical boundary was still physical access to those comparisons.  The present
continuation asks a more upstream question: if the operational entrance already
contains an actual tomographically complete intervention frame on this same
input/output port, must the V22--V23 comparison effects themselves be executable,
or can their expectation values be reconstructed from the old probability table?

The answer is the latter.  Exact execution and exact statistical
\emph{identifiability} are different notions.  A signed linear combination of
probabilities is not a new quantum operation, but it can determine exactly the
expectation of an unexecuted Hermitian tester.  This distinction was deliberately
kept open in the earlier source audit: tomographic span was not treated as a
license to execute every completely positive subchannel.  We retain that
restriction.  The result below uses only classical postprocessing after the
admitted experiments have occurred.

The motivation is source-specific.  The spacetime entrance already requires
physical intervention frames that are tomographically complete on each
\emph{exposed finite operator port}.  This does not by itself prove that the
five-copy native memory of V19--V23 is such an exposed port.  V24 therefore
removes the separate tester-availability premise \emph{conditional on that
same-cut port identification}; it does not manufacture the port.

\paragraph{Nonduplication.}
The positive-completion theorem of V22 already characterizes symmetry forced by
an arbitrary linear data map.  We do not repeat it.  V23 already gives an
ancilla-free physical realization of one sufficient comparison bank.  The new
work is the exact bridge from an actual preparation/effect frame to the V22
linear functionals, the normalization-aware identifiability criterion, its
optimal error amplification constant, and the resulting source-level reduction.
No target effect is inserted into the operation alphabet.

\section{Exact reconstruction from an admitted preparation--effect frame}
\label{sel24:frame}

Let $\mathcal H_{\rm in}\cong\C^{d_{\rm in}}$ and
$\mathcal H_{\rm out}\cong\C^{d_{\rm out}}$.  Let
$\rho_a\succeq0$, $\Tr\rho_a=1$, $1\le a\le A$, be actual admitted input
preparations, and let $0\preceq E_b\preceq I$, $1\le b\le B$, be actual admitted
output effects.  For an unknown CPTP map $\Phi$ use the output--input Choi
convention
\[
 J_\Phi=\sum_{i,j}\Phi(|i\rangle\langle j|)\otimes|i\rangle\langle j|,
 \qquad \Tr_{\rm out}J_\Phi=I_{\rm in}.
\]
The physically observed one-use probabilities are
\begin{equation}
 p_{ba}=\Tr\!\left[E_b\Phi(\rho_a)\right]
       =\Tr\!\left[(E_b\otimes\rho_a^{\mathsf T})J_\Phi\right].
 \label{sel24:eq:data}
\end{equation}
No linear combination on the right is interpreted as an executable effect.

Put
\[
 \mathscr P=\Span_{\R}\{\rho_a\}\subseteq\Herm(\mathcal H_{\rm in}),
 \qquad
 \mathscr E=\Span_{\R}\{E_b\}\subseteq\Herm(\mathcal H_{\rm out}).
\]

\begin{theorem}[Operational-frame reconstruction without target-tester execution]
\label{sel24:thm:dual-frame}
Assume $\mathscr P=\Herm(\mathcal H_{\rm in})$ and
$\mathscr E=\Herm(\mathcal H_{\rm out})$.  Choose any Hilbert--Schmidt dual
frames $\widetilde\rho_a,\widetilde E_b$ satisfying
\[
 X=\sum_a\Tr(\widetilde\rho_aX)\rho_a,
 \qquad
 Y=\sum_b\Tr(\widetilde E_bY)E_b
\]
for all Hermitian $X,Y$.  Then every Hermitian Choi functional $F$ has the exact
expansion
\begin{equation}
 F=\sum_{a,b}c_{ba}(F)\,E_b\otimes\rho_a^{\mathsf T},
 \qquad
 c_{ba}(F)=
 \Tr\!\left[(\widetilde E_b\otimes\widetilde\rho_a^{\mathsf T})F\right],
 \label{sel24:eq:dual-expansion}
\end{equation}
and consequently
\begin{equation}
 \boxed{\Tr(FJ_\Phi)=\sum_{a,b}c_{ba}(F)p_{ba}.}
 \label{sel24:eq:reconstruct}
\end{equation}
Thus the value of $F$ is determined by classical postprocessing of the actual
frame probabilities.  Neither $F$ nor its positive and negative parts need be
admitted as physical effects.
\end{theorem}

\begin{proof}
The tensor products of two spanning real Hermitian frames span the Hermitian
operators on $\mathcal H_{\rm out}\otimes\mathcal H_{\rm in}$.  Applying the two
dual-frame identities successively to each tensor factor gives
\eqref{sel24:eq:dual-expansion}.  Pairing with $J_\Phi$ and using
\eqref{sel24:eq:data} gives \eqref{sel24:eq:reconstruct}.  All coefficients are
real because the paired operators are Hermitian.  The reconstruction is a
linear identity between observed numbers; no signed combination is promoted to
a quantum operation.
\end{proof}

Full tomography is stronger than necessary.  The next theorem gives the exact
normalization-aware criterion for one proposed witness.

\begin{theorem}[Exact tester-span criterion modulo trace preservation]
\label{sel24:thm:span}
Let
\[
 \mathscr T
 =\Span_{\R}\{E_b\otimes\rho_a^{\mathsf T}:a,b\}
 \subseteq\Herm(\mathcal H_{\rm out}\otimes\mathcal H_{\rm in}),
\]
and let
\[
 \mathscr N
 =\{I_{\rm out}\otimes Y^{\mathsf T}:Y=Y^*\}
\]
be the adjoint space of the trace-preserving normalization constraint.  For a
Hermitian tester $F$, the following are equivalent.
\begin{enumerate}[label=\textup{(\roman*)},leftmargin=2.8em]
\item $\Tr(FJ)$ has the same value for all Hermitian $J$ with
      $\Tr_{\rm out}J=I$ and the same frame data \eqref{sel24:eq:data};
\item $F\in\mathscr T+\mathscr N$;
\item there exist real $c_{ba}$ and Hermitian $Y$ such that
\begin{equation}
 \boxed{
 F=\sum_{a,b}c_{ba}E_b\otimes\rho_a^{\mathsf T}
       +I_{\rm out}\otimes Y^{\mathsf T}.}
 \label{sel24:eq:modnorm}
\end{equation}
\end{enumerate}
When these conditions hold,
\begin{equation}
 \Tr(FJ_\Phi)=\sum_{a,b}c_{ba}p_{ba}+\Tr Y.
 \label{sel24:eq:modnorm-value}
\end{equation}
If they fail and the compatible fibre contains a faithful Choi point
$J_0\succ0$, there are two CPTP Choi matrices $J_\pm$ with identical admitted
frame probabilities but different values of $\Tr(FJ_\pm)$.
\end{theorem}

\begin{proof}
The tangent space to the affine trace-preserving slice is
\[
 \mathscr V_{\rm TP}
 =\{X=X^*: \Tr_{\rm out}X=0\}.
\]
The data-null tangent is
\[
 \mathscr K=\mathscr V_{\rm TP}\cap\mathscr T^\perp.
\]
Condition (i) is exactly $F\perp\mathscr K$.  In finite dimension,
\[
 \mathscr K^\perp
 =\mathscr V_{\rm TP}^\perp+\mathscr T
 =\mathscr N+\mathscr T,
\]
because the adjoint of partial trace maps $Y$ to
$I_{\rm out}\otimes Y^{\mathsf T}$ in the present convention.  This proves the
equivalence and \eqref{sel24:eq:modnorm-value}.

If $F\notin\mathscr T+\mathscr N$, orthogonal projection onto $\mathscr K$
gives an $X\in\mathscr K$ with $\Tr(FX)\ne0$.  Since $J_0\succ0$, for all
sufficiently small $\varepsilon>0$ the matrices
$J_\pm=J_0\pm\varepsilon X$ remain positive.  They remain trace preserving and
have the same admitted data because $X\in\mathscr K$, while their $F$ values
differ by $2\varepsilon\Tr(FX)$.
\end{proof}

\begin{remark}[Why informational completeness is not effect completeness]
\label{sel24:rem:not-execution}
Theorem~\ref{sel24:thm:dual-frame} does not contradict the earlier no-go results
for executable refinements.  A dual coefficient $c_{ba}$ may be negative or
larger than one.  Such coefficients are permitted in offline statistical
postprocessing and forbidden as probabilities of a randomized physical effect.
Tomographic completeness therefore determines an expectation without licensing
the corresponding operation.
\end{remark}

\section{Optimal error amplification for a source-native frame}
\label{sel24:conditioning}

Exact source identities are useful only if their finite conditioning cost is
visible.  For a tester satisfying Theorem~\ref{sel24:thm:span}, define its
normalization-aware $\ell^1$ reconstruction cost
\begin{equation}
 \kappa_1(F\mid\mathscr T)
 :=\inf\left\{\sum_{a,b}|c_{ba}|:
 F-\sum_{a,b}c_{ba}E_b\otimes\rho_a^{\mathsf T}\in\mathscr N\right\}.
 \label{sel24:eq:kappa}
\end{equation}
The infimum is attained because this is a finite-dimensional linear program.

\begin{theorem}[Sharp worst-case frame-error constant]
\label{sel24:thm:error}
Suppose estimates satisfy
$|\widehat p_{ba}-p_{ba}|\le\epsilon$ for every admitted frame probability.
For any exact representation \eqref{sel24:eq:modnorm},
\[
 \left|\sum_{a,b}c_{ba}\widehat p_{ba}+\Tr Y-\Tr(FJ_\Phi)\right|
 \le\epsilon\sum_{a,b}|c_{ba}|.
\]
Optimizing over all representations gives
\begin{equation}
 \boxed{|\widehat f-f|\le\epsilon\,\kappa_1(F\mid\mathscr T).}
 \label{sel24:eq:error}
\end{equation}
This coefficient is the smallest possible uniform constant for arbitrary
independent entrywise errors bounded by $\epsilon$.
\end{theorem}

\begin{proof}
The upper bound is the triangle inequality.  Conversely, for any fixed
coefficient representation the adversarial signs
$\widehat p_{ba}-p_{ba}=\epsilon\operatorname{sgn}(c_{ba})$ attain its
$\ell^1$ norm at the linear level.  Minimizing over equivalent representations
modulo normalization gives the quotient norm \eqref{sel24:eq:kappa}.  Finite
linear-program duality gives attainment and shows no smaller uniform constant
can represent the same functional on the data quotient.
\end{proof}

There is a corresponding dual obstruction.  The linear-program dual is a
maximization over $X\in\mathscr V_{\rm TP}$ satisfying
$|\Tr[(E_b\otimes\rho_a^{\mathsf T})X]|\le1$.  Hence a large
$\kappa_1$ is accompanied by a normalized unresolved tangent direction on which
$F$ is large.  Poor tomography conditioning is therefore itself a finite source
witness, not merely a numerical inconvenience.

\section{The V22--V23 determinant certificate is frame reconstructible}
\label{sel24:application}

Return to the independently identified V19--V23 five-copy native carrier.  Keep
its seven named vectors $n,\omega,u_0,\ldots,u_4$, the active projection $Q$, and
\[
 \mathcal A_Q=\mathcal B(Q\mathcal H)\oplus\C Q_b.
\]
V22 uses the two Hermitian Choi functionals
\begin{equation}
 F_{\rm leak}=\frac{I-P_Q}{N},
 \qquad
 F_{\rm int}=\frac{X_\circ}{N},
 \label{sel24:eq:twoF}
\end{equation}
whose values are respectively $\ell$ and $\mu$.  V23 supplies a different native
preparation/readout implementation of the same symmetry conclusion.  We now
ask only whether the old operational frame determines the two values in
\eqref{sel24:eq:twoF}.

\begin{theorem}[Operational-frame closure of the two-scalar determinant selector]
\label{sel24:thm:closure}
Let $\{\rho_a\},\{E_b\}$ be the actual preparation/effect frame exposed on the
same identified native input/output port.  Assume either
\begin{enumerate}[label=\textup{(A\arabic*)},leftmargin=3em]
\item both frames are tomographically complete on their Hermitian operator
      spaces; or
\item more minimally,
      $F_{\rm leak},F_{\rm int}\in\mathscr T+\mathscr N$.
\end{enumerate}
Then $\ell$ and $\mu$ are exact signed linear functions of the already observed
frame probabilities.  No V22 support projector, Choi interference effect, V23
complementary-basis survival Read, or native phase Read has to be added to the
physical operation alphabet in order to \emph{infer} these two numbers.

If their reconstructed values satisfy
\begin{equation}
 \boxed{\ell=0,\qquad\mu\ne0,}
 \label{sel24:eq:closure-test}
\end{equation}
then every CPTP completion of the physical frame data has covariance group
inside $G_0=\mathrm U(3)\times\mathrm U(2)$ exactly
\[
 \boxed{S(\mathrm U(3)\times\mathrm U(2)).}
\]
With entrywise frame error at most $\epsilon$, the signs and zero/nonzero
margins are certified whenever the corresponding bounds from
Theorem~\ref{sel24:thm:error}, together with the positive-completion residuals,
separate zero from the measured intervals.
\end{theorem}

\begin{proof}
Under (A1), Theorem~\ref{sel24:thm:dual-frame} reconstructs both functionals.
Under (A2), Theorem~\ref{sel24:thm:span} does.  The values are properties of the
same unknown physical Choi operator; reconstructing them does not replace that
operator by a fitted channel.  V22's two-scalar theorem then applies verbatim:
$\ell=0$ confines every positive completion to the represented algebra
$\mathcal A_Q$, which contains $G_{\det}$ in its covariance, and $\mu\ne0$
forces the determinant character to be trivial on every additional covariance.
The error statement is the quotient-norm bound above combined with the V22
positive-completion logic; no rounding of a positive residual to zero is used.
\end{proof}

\begin{corollary}[What the operational entrance would remove]
\label{sel24:cor:entrance}
Suppose the historical five-copy native input/output memory is proved to be one
of the \emph{exposed finite operator ports} on which the Renewal operational
entrance supplies tomographically complete physical intervention frames.  Then,
for purposes of the V22 symmetry selector, tester availability is no longer an
independent source premise.  The remaining source rows are:
\begin{enumerate}[label=\textup{(\roman*)},leftmargin=2.8em]
\item same-cut identification of that exposed port with the native colour--weak
      tensor carrier;
\item population of its actual frame probability table;
\item the reconstructed numerical conditions $\ell=0$ and $\mu\ne0$ (or another
      forcing panel from the V22 completion theorem).
\end{enumerate}
Physical execution of the V23 native phase/survival protocol remains an optional
direct verification, not a logical requirement for reconstructing the two Choi
functionals.
\end{corollary}

The premise in this corollary is intentionally precise.  Tomographic
completeness on \emph{some} upstream port cannot be transported to a different
five-copy memory by matching dimensions or representation labels.  The
same-carrier condition of the duality manuscript remains in force.

\section{Incomplete frames return an explicit unresolved quantum direction}
\label{sel24:failure}

The operational-frame reduction is falsifiable.  It is not a declaration that
every finite intervention bank is sufficient.

\begin{proposition}[Frame defect and faithful positive counter-completions]
\label{sel24:prop:defect}
Let $F$ be either tester in \eqref{sel24:eq:twoF}.  If
$F\notin\mathscr T+\mathscr N$, define the normalized frame defect
\begin{equation}
 \delta_F
 =\left\|P_{\mathscr K}F\right\|_{\HS}>0,
 \qquad
 \mathscr K=\mathscr V_{\rm TP}\cap\mathscr T^\perp.
 \label{sel24:eq:defect}
\end{equation}
Then $X=P_{\mathscr K}F/\delta_F$ is a unit-Hilbert--Schmidt, trace-preserving
unresolved direction satisfying
\[
 \Tr[(E_b\otimes\rho_a^{\mathsf T})X]=0\quad\forall a,b,
 \qquad
 \Tr(FX)=\delta_F.
\]
At every faithful compatible Choi point $J_0\succ0$, the perturbations
$J_0\pm\varepsilon X$ are distinct CPTP completions with the same physical frame
data for all sufficiently small $\varepsilon>0$, and their $F$ values differ by
$2\varepsilon\delta_F$.
\end{proposition}

\begin{proof}
Orthogonal projection onto $\mathscr K$ gives all displayed identities.  Strict
positivity of $J_0$ is open in operator norm and
$\|X\|_{\op}\le\|X\|_{\HS}=1$, so
$0<\varepsilon<\lambda_{\min}(J_0)$ suffices for positivity.  The trace and data
constraints are unchanged because $X\in\mathscr K$.
\end{proof}

This is the correct upstream stopping rule.  If the historical frame fails to
identify one of the two witnesses, adding the desired tester would be a new
physical operation.  The programme must instead enlarge the actually observed
frame, exploit a positive boundary condition that shrinks the completion face,
or retain the unresolved direction as a genuine source ambiguity.

\begin{example}[Informational span and executable cone are different]
\label{sel24:ex:qubit}
On a qubit, the four physical states
\[
 \rho_0=|0\rangle\langle0|,
 \quad\rho_1=|1\rangle\langle1|,
 \quad\rho_+=|+\rangle\langle+|,
 \quad\rho_{+i}=|+i\rangle\langle+i|
\]
span all Hermitian matrices.  Their real linear span therefore reconstructs the
expectation of the opposite phase state
$\rho_{-i}=|-i\rangle\langle-i|$.  But $\rho_{-i}$ is not in their convex hull:
its Bloch $y$ coordinate is $-1$, whereas every convex combination of the four
listed states has nonnegative $y$ coordinate.  Signed tomography can therefore
recover a target functional that randomized execution of the supplied
preparations cannot realize.  This is exactly the distinction used in
Theorem~\ref{sel24:thm:closure}.
\end{example}

\section{Result and next upstream source row}
\label{sel24:status}

V24 changes the remaining source-selection question.  The phase/survival effects
introduced as direct physical tests in V22--V23 are not necessarily new source
premises.  If the same native memory is already an exposed port with an actual
tomographically complete preparation/effect frame, their Choi functionals are
exactly reconstructible from the existing probability table by signed classical
postprocessing.  The operational frame need not contain those effects in its
positive executable cone.

Thus the next historical audit should not ask first for an exotic determinant
measurement.  It should ask:
\[
 \boxed{
 \text{Is the selected native colour--weak memory an actual exposed operator port,}
 \quad
 \text{and what is its admitted intervention-frame span?}}
\]
If that frame is complete, evaluate the two reconstructed scalars.  If it is
partial, run the finite membership test
$F_{\rm leak},F_{\rm int}\in\mathscr T+\mathscr N$; on failure the projection
\eqref{sel24:eq:defect} is the explicit missing quantum direction.  This is more
upstream and less target-shaped than adding another active measurement.

Historical Standard-Model selection is still not declared closed.  The
same-carrier port identification and the historical frame values have not been
populated here.  Nor does process tomography derive the charged matter modules,
fermionic grading, Clifford/Yukawa incidence, or continuum dynamics.  It only
removes a separate tester-execution premise when the already-declared
operational tomography truly lives on the required port.

\paragraph{Proof and artifact status.}
The entire V23 source prefix is retained before its final terminator.  The V24
checker verifies dual-frame reconstruction, the normalization quotient, the
qubit cone/span separation, and finite faithful counter-completions.  The all-
dimension theorems are the written linear-algebra proofs above; finite scripts
are regressions, not proof-assistant formalization or independent peer review.
No historical probability table is invented.

\addtocontents{toc}{\protect\endgroup}
% END OF SOURCE-SELECTION V24 DEVELOPMENT BLOCK.
% Preserve V1--V24 and append subsequent work before the final terminator.


\clearpage
\begin{center}
{\Large\bfseries Source-selection continuation V25}\\[.4em]
{\large Contextual tomography of hidden memories, Hankel-saturated landing,\\
and the irreducible joint-five-factor occurrence row}
\end{center}
\addtocontents{toc}{\protect\begingroup}
\addtocontents{toc}{\protect\renewcommand{\protect\numberline}{\protect\makebox[2.1em][l]}}

\section{Move upstream from an exposed port to the original cylinder table}
\label{sel25:context}

V24 shows that a tomographically complete physical intervention frame on the
\emph{same exposed native port} determines the determinant witnesses by signed
classical postprocessing.  The remaining word ``exposed'' is stronger than is
needed for inference.  The operational Grand-Tensor construction already gives
support-minimal sequential memories and says that they are unique up to
record-preserving memory unitaries.  An internal memory need not itself admit
arbitrary physical interventions for the old prefixes and futures to probe it.
The present continuation makes this distinction exact.

At one cut, an actual past history prepares an unnormalized positive memory
state and an actual future word defines a positive memory effect.  Their pairing
is an ordinary cylinder probability.  If those past states and future effects
span the Hermitian memory spaces, then any \emph{actual event inserted between
them} is process-tomographed from the original cylinder table, even if the
memory is hidden and never exposed as an intervention port.  More generally,
only the span of the particular target Choi functional is needed.  This is the
hidden-memory analogue of V24's exposed-frame theorem.

There is a sharp boundary.  V13's exact rank-$196$ original-record panel already
saturates the Hermitian coordinates of one fourteen-dimensional memory.  It can
therefore tomograph any actual same-cut event on that memory.  It does not,
however, create a five-factor joint source operation.  If five copies are
physically co-instantiated and the product prefix--suffix contexts are admitted,
the single-copy basis tensorizes and tomographs an actual joint event on
$14^5$ dimensions.  Five independent repetitions without such a joint event
are insufficient: V17 already gives channels with identical every-subcritical
marginal and different determinant covariance groups.

\paragraph{Audited inputs and nonduplication.}
The finite process-comb construction, its support-minimal memory and its unitary
uniqueness are inherited from the operational reconstruction used by the
spacetime programme.  V13's Theorem~\ref{sel13:thm:full-panel} supplies the
exact rank-$196$ probability witness and explicitly proves that its selected
past states and future effects each span $M_{14}$.  V17's
Theorem~\ref{sel17:thm:marginals} is retained as the joint-order no-go.  V24's
normalization-aware span theorem is the exposed-port precursor of the contextual
theorem below.  None of these statements is reproved as a new source-selection
result.  The new step is to replace direct intervention tomography of a hidden
memory by prefix--suffix tomography of an actually occurring event and then to
identify exactly which part does \emph{not} tensorize without a physical joint
source row.

\paragraph{Conventions.}
All finite memories below are support-minimal.  The transpose in a Choi tensor
is taken in an arbitrary temporary memory basis.  No conclusion depends on that
basis.  Real spans of Hermitian matrices are used for tomography.  Cylinder
probabilities are unnormalized when the prefix itself has probability below
one; its original weight is retained.  A ``context'' means an actually
occurring prefix or future word, not an independently admitted control on the
hidden memory.

\section{Prefix--suffix tomography of an internal event}
\label{sel25:contextual}

Fix an actual event $a$ between two cuts of a support-minimal sequential
realization.  Let the incoming and outgoing memories be
$A_-\cong\C^{d_-}$ and $A_+\cong\C^{d_+}$.  An admitted prefix $p$ ending at the
incoming cut determines an unnormalized state $\sigma_p\succeq0$ on $A_-$.
An admitted future $f$ beginning after the outgoing cut determines a positive
Heisenberg effect $E_f\succeq0$ on $A_+$.  If the event acts by the actual CP
map $\Phi_a$, then
\begin{equation}
 p(fap)=\Tr\!\left[E_f\Phi_a(\sigma_p)\right]
       =\Tr\!\left[(E_f\otimes\sigma_p^{\mathsf T})J_a\right].
 \label{sel25:eq:cylinder-Choi}
\end{equation}
The first equality is the sequential realization of the original cylinder;
the second fixes the Choi convention.  No $\sigma_p$ or $E_f$ is thereby
promoted to an executable preparation or Read on the hidden memory.

Put
\begin{equation}
 \mathscr R_-:=\Span_{\R}\{\sigma_p:p\in\mathcal P\},\qquad
 \mathscr O_+:=\Span_{\R}\{E_f:f\in\mathcal F\},
 \label{sel25:eq:RO}
\end{equation}
and define the contextual tester space
\begin{equation}
 \mathscr T_{\rm ctx}
 :=\Span_{\R}\{E_f\otimes\sigma_p^{\mathsf T}:f,p\}.
 \label{sel25:eq:Tctx}
\end{equation}
For trace-preserving events retain the normalization-annihilator complement in
exactly V24's convention,
\begin{equation}
 \mathscr N_-:=\{I_{+}\otimes Y^{\mathsf T}:Y=Y^*\}.
 \label{sel25:eq:Nctx}
\end{equation}

\begin{theorem}[Contextual identifiability on a hidden memory]
\label{sel25:thm:contextual}
Let $F=F^*$ be a Choi functional of the actual event $a$.
\begin{enumerate}[label=\textup{(\roman*)},leftmargin=2.8em]
\item Its value $\Tr(FJ_a)$ is determined by the cylinder probabilities
      \eqref{sel25:eq:cylinder-Choi} and trace preservation whenever
      \begin{equation}
       \boxed{F\in\mathscr T_{\rm ctx}+\mathscr N_-.}
       \label{sel25:eq:ctx-span}
      \end{equation}
      In that case there are real coefficients $c_{fp}$ and a Hermitian $Y$
      such that
      \begin{equation}
       \boxed{\Tr(FJ_a)=\sum_{f,p}c_{fp}\,p(fap)+\Tr Y.}
       \label{sel25:eq:ctx-reconstruction}
      \end{equation}
\item Conversely, if $F\notin\mathscr T_{\rm ctx}+\mathscr N_-$ and the
      compatible event fibre contains a faithful Choi point $J_0\succ0$, there
      are two CPTP maps with all the same contextual cylinder probabilities
      but different values of $\Tr(FJ)$.  The invisible tangent may be chosen
      as the normalized projection of $F$ onto
      \begin{equation}
       \mathscr K_{\rm ctx}
       :=\{X=X^*: \Tr_{+}X=0\}\cap\mathscr T_{\rm ctx}^{\perp}.
       \label{sel25:eq:ctx-kernel}
      \end{equation}
\item If $\mathscr R_-=\mathrm{Herm}(M_{d_-})$ and
      $\mathscr O_+=\mathrm{Herm}(M_{d_+})$, then
      $\mathscr T_{\rm ctx}$ is the full Hermitian Choi space and the complete
      $J_a$ is determined from the old cylinders.  Direct access to the memory
      as an exposed port is unnecessary for this inference.
\end{enumerate}
At a positive boundary point, failure of \eqref{sel25:eq:ctx-span} is only a
linear obstruction; the smaller positive completion face must be used exactly
as in V22 before asserting non-identifiability.
\end{theorem}

\begin{proof}
If
$F=\sum_{f,p}c_{fp}E_f\otimes\sigma_p^{\mathsf T}
   +I_+\otimes Y^{\mathsf T}$,
then \eqref{sel25:eq:cylinder-Choi} and
$\Tr_+J_a=I_-$ give
\[
 \Tr(FJ_a)=\sum_{f,p}c_{fp}p(fap)+\Tr Y,
\]
which proves (i).  For (ii), orthogonal projection onto
$\mathscr K_{\rm ctx}$ gives a nonzero $X$ with all contextual pairings zero
and $\Tr_+X=0$, but with $\Tr(FX)\ne0$.  For sufficiently small
$\varepsilon>0$, $J_0\pm\varepsilon X$ remain positive and trace preserving.
They therefore give the required two physical completions.  Finally,
Hermitian bases on the output and input have tensor products forming a real
basis of the Hermitian Choi space, proving (iii).
\end{proof}

\begin{remark}[What has been removed]
The theorem removes \emph{direct tester execution on the hidden memory}, not
physical occurrence of the event $a$.  The probability $p(fap)$ exists only
when the prefix, the event, and the suffix occur in the declared grammar on
one sequential carrier.  Inserting a mathematically desired $a$ between old
contexts would change the source.
\end{remark}

\section{A Hankel-rank criterion for contextual completeness}
\label{sel25:hankel}

At a stationary cut with one $d$-dimensional memory, choose finite prefix and
future lists and form the ordinary probability Hankel matrix
\begin{equation}
 H_{f,p}=p(fp)=\Tr(E_f\sigma_p).
 \label{sel25:eq:Hankel}
\end{equation}
Its factorization goes through the real Hilbert space
$\mathrm{Herm}(M_d)$ of dimension $d^2$.

\begin{theorem}[Hankel saturation is hidden-memory tomography]
\label{sel25:thm:Hankel}
For every such support-minimal realization,
\begin{equation}
 \rank H\le
 \min\{\dim\mathscr R,\dim\mathscr O\}\le d^2.
 \label{sel25:eq:Hankel-bound}
\end{equation}
If a finite subpanel has rank $d^2$, then
\begin{equation}
 \boxed{
 \mathscr R=\mathscr O=\mathrm{Herm}(M_d).}
 \label{sel25:eq:Hankel-saturation}
\end{equation}
Choose $d^2$ independent prefixes $p_i$ and futures $f_j$, and let
$\widehat\sigma_i,\widehat E_j$ be the Hilbert--Schmidt dual bases.  Every
actual event $a$ between these cuts is then reconstructed by
\begin{equation}
 \boxed{
 J_a=\sum_{i,j=1}^{d^2}p(f_jap_i)\,
       \widehat E_j\otimes\widehat\sigma_i^{\mathsf T}.}
 \label{sel25:eq:event-reconstruction}
\end{equation}
The dual matrices need not be positive; the formula is statistical
reconstruction, not an executable hidden-memory protocol.
\end{theorem}

\begin{proof}
Let $R$ have columns the coordinate vectors of the $\sigma_p$ in a Hermitian
basis and let $O$ have columns those of the $E_f$.  Then, up to the harmless
choice of row/column orientation, $H=O^{\mathsf T}R$.  The rank inequalities
follow.  Equality with $d^2$ forces both spans to have their maximum possible
dimension.  For the last formula, the product duals
$\widehat E_j\otimes\widehat\sigma_i^{\mathsf T}$ are dual to the contextual
testers $E_j\otimes\sigma_i^{\mathsf T}$, so expansion of $J_a$ in that dual
basis gives \eqref{sel25:eq:event-reconstruction}.
\end{proof}

\begin{corollary}[The V13 rank-$196$ witness already tomographs its hidden memory]
\label{sel25:cor:V13}
At the exact V13 source point of Theorem~\ref{sel13:thm:full-panel},
$d=14$ and the selected original-record panel has rank
$196=14^2$.  Hence its 196 past states and 196 future effects are complete
Hermitian frames on the support-minimal memory.  Any \emph{actual} same-cut
CP event inserted between those selected contexts is determined by its
original cylinder probabilities $p(f_jap_i)$, without exposing that memory as
a physical intervention port.

The same conclusion holds for any historical source on which the corresponding
rank-$196$ condition is actually observed.  The V13 nonzero comparison point
itself is not reclassified as historical data.
\end{corollary}

\begin{proof}
Theorem~\ref{sel13:thm:full-panel} proves both the probability rank and the
past-state/future-effect span assertions directly.  Apply
Theorem~\ref{sel25:thm:Hankel}.  The final sentence retains V13's own
historical-population firewall.
\end{proof}

\section{Memory-gauge invariance: no preferred hidden basis is needed}
\label{sel25:gauge}

Support-minimal sequential memories are unique only up to memory unitaries.
That ambiguity is not a new port-identification problem for contextual
tomography.

\begin{proposition}[Record-preserving unitary covariance of contextual landing]
\label{sel25:prop:gauge}
Let another support-minimal realization use memory unitaries
$U_-:A_-\to A'_-$ and $U_+:A_+\to A'_+$.  Then
\begin{align}
 \sigma'_p&=U_-\sigma_pU_-^*,&
 E'_f&=U_+E_fU_+^*,\nonumber
 J'_a&=(U_+\otimes\overline{U_-})J_a
       (U_+\otimes\overline{U_-})^*.
 \label{sel25:eq:gauge-transform}
\end{align}
The contextual probabilities and all span ranks are unchanged.  If an
intrinsic target functional is transported by the same rule,
$F'=(U_+\otimes\overline{U_-})F(U_+\otimes\overline{U_-})^*$, then
\begin{equation}
 \boxed{\Tr(F'J'_a)=\Tr(FJ_a).}
 \label{sel25:eq:gauge-invariant}
\end{equation}
Thus a contextual certificate does not require choosing the canonical memory
unitary left unfixed by process-comb tomography.
\end{proposition}

\begin{proof}
These are exactly the inter-realization transformations of a sequential CP
realization.  Cyclicity of trace proves every cylinder equality and
\eqref{sel25:eq:gauge-invariant}.  Conjugation is an invertible real-linear map
on Hermitian matrices, so all span dimensions and membership tests are
unchanged.
\end{proof}

The proposition removes a \emph{basis} ambiguity, not a \emph{carrier}
ambiguity.  A hidden memory at one cut cannot be identified with a different
physical tensor product merely because the two Hilbert spaces have the same
dimension.

\section{When the contextual frame tensorizes to a real five-factor event}
\label{sel25:tensor}

The determinant selector from V16--V23 lives at a higher tensor order.  The
right upstream question is therefore whether the historical grammar contains
one actual joint event with enough old contexts around it.

Suppose five copies of the same $d$-dimensional support-minimal memory are
physically co-instantiated at a cut.  Assume the operational grammar admits the
product prefix states and product future effects
\begin{equation}
 \sigma_{\mathbf i}=\sigma_{i_1}\otimes\cdots\otimes\sigma_{i_5},
 \qquad
 E_{\mathbf j}=E_{j_1}\otimes\cdots\otimes E_{j_5},
 \label{sel25:eq:product-contexts}
\end{equation}
where each one-copy family is a Hermitian basis.  This is a physical
\emph{parallel context-closure} condition.  It is stronger than the ability to
repeat one experiment five times on five unrelated runs.

\begin{theorem}[Tensorized contextual landing of a joint event]
\label{sel25:thm:tensor}
Under \eqref{sel25:eq:product-contexts}, let $A$ be one actual CP event acting
jointly on the five co-instantiated memories.  Put $N=d^5$.  Then the
$N^2$ product prefixes and $N^2$ product futures span the complete Hermitian
input and output spaces on $M_N$, and the probabilities
\begin{equation}
 P^A_{\mathbf j,\mathbf i}
 =\Tr\!\left[E_{\mathbf j}A(\sigma_{\mathbf i})\right]
 \label{sel25:eq:joint-panel}
\end{equation}
uniquely determine the full joint Choi operator $J_A$.
Explicitly,
\begin{equation}
 \boxed{
 J_A=\sum_{\mathbf i,\mathbf j}
 P^A_{\mathbf j,\mathbf i}
 \left(\bigotimes_{r=1}^{5}\widehat E_{j_r}\right)
 \otimes
 \left(\bigotimes_{r=1}^{5}\widehat\sigma_{i_r}\right)^{\mathsf T}.}
 \label{sel25:eq:joint-reconstruction}
\end{equation}
Consequently every V22 Choi functional, including the determinant-support and
interference functionals when defined on this same native joint carrier, is a
signed linear function of the original joint cylinder table.  The five-copy
port itself need not be directly exposed to new tomography operations.
\end{theorem}

\begin{proof}
Tensor products of real Hermitian bases form a real Hermitian basis of the
fivefold matrix algebra.  Their product duals are the displayed tensor products
of one-copy dual bases.  Apply Theorem~\ref{sel25:thm:Hankel} in the resulting
$N$-dimensional memory, or expand the Choi operator directly against the dual
product testers.
\end{proof}

If a square one-copy context Hankel $H$ is used on each factor, the no-event
product context matrix is $H^{\otimes5}$.  Hence
\begin{equation}
 \rank(H^{\otimes5})=(\rank H)^5,
 \qquad
 \det(H^{\otimes5})=\det(H)^{5n^4}
 \quad(n=\dim H).
 \label{sel25:eq:Kronecker-det}
\end{equation}
At V13's saturated size $n=196$, one nonzero one-copy determinant therefore
certifies the tensor-product context rank algebraically; one need not perform a
separate gigantic rank proof.  Physical occurrence of the product contexts and
of $A$ is still required.

\begin{corollary}[What V24's exposed-port premise can be weakened to]
\label{sel25:cor:weaken}
For the determinant selector, direct exposure of the five-copy native memory
may be replaced by the conjunction of:
\begin{enumerate}[label=\textup{(\roman*)},leftmargin=2.8em]
\item an actually co-instantiated five-factor memory on the selected native
      representation;
\item an actual joint event $A$ whose covariance is being tested;
\item parallel prefix--suffix context closure as in
      \eqref{sel25:eq:product-contexts}; and
\item single-factor contextual completeness, for example a measured
      rank-$196$ Hankel certificate on each fourteen-dimensional factor.
\end{enumerate}
Under these conditions V22--V24's determinant functionals are reconstructible
from historical cylinders without any direct intervention on the hidden
five-copy memory.
\end{corollary}

\section{Why independent repetitions cannot supply the missing joint source}
\label{sel25:no-go}

The preceding theorem makes the remaining physical row narrower, but it does
not remove it.  A tensor product of \emph{coordinate frames} is free once the
parallel contexts occur; a tensor product of \emph{unknown dynamics} is not.

\begin{theorem}[The joint-event row is irreducible to subcritical marginals]
\label{sel25:thm:joint-no-go}
Knowing all single-copy contextual process data, and even all covariantly
formed native marginals with fewer than five total input/output factors, does
not determine an unknown five-slot joint event.  In particular V17's explicit
family $\Lambda_\eta$ satisfies
\begin{equation}
 \mathfrak F(\Lambda_\eta)=\mathfrak F(\Lambda_0)
 \label{sel25:eq:all-marginals}
\end{equation}
for every $G_0$-covariant marginal construction with fewer than five retained
native slots, while
\begin{equation}
 \operatorname{Cov}(\Lambda_0)=G_0,
 \qquad
 \operatorname{Cov}(\Lambda_\eta)=G_{\det}\quad(\eta\ne0).
 \label{sel25:eq:group-separation}
\end{equation}
Therefore five statistically independent repetitions of a one-copy source, or
a complete archive of their separate marginal tables, cannot be relabelled as
the historical five-factor determinant event.  An actual joint source row or an
equivalent same-history joint Choi packet remains necessary.
\end{theorem}

\begin{proof}
This is the operational consequence of
Theorem~\ref{sel17:thm:marginals}.  That theorem gives equality of the reduced
\emph{quantum channels on all retained inputs}, not merely equality of selected
scalar populations, and gives the two different full covariance groups for the
same family.  If the subcritical data determined the complete joint event, the
two family members would have to coincide, contradicting their nonzero diamond
distance and different covariance groups.
\end{proof}

This also identifies the precise point at which the original rank-$196$ source
stops helping.  It saturates the operator coordinates of one memory and makes
hidden single-copy event tomography possible.  It does not provide the
cross-copy Choi coherence that V17 proves invisible to every subcritical
marginal.

\section{Conditioning cost of contextual and parallel tomography}
\label{sel25:conditioning}

Removing an execution assumption does not remove numerical conditioning.  Let
$\mathcal A_R$ and $\mathcal A_O$ be the real analysis maps taking an incoming
Hermitian matrix to its prefix-frame coordinates and an outgoing Hermitian
matrix to its future-effect coordinates, restricted to their full spans.  The
contextual Choi analysis map is their Kronecker product, up to transpose and a
permutation of coordinates.  Therefore
\begin{equation}
 \boxed{
 \kappa_2(\mathcal A_{\rm ctx})
 =\kappa_2(\mathcal A_O)\,\kappa_2(\mathcal A_R).}
 \label{sel25:eq:condition-one}
\end{equation}
For five identical parallel context frames,
\begin{equation}
 \boxed{
 \kappa_2(\mathcal A_{\rm ctx}^{\otimes5})
 =\kappa_2(\mathcal A_{\rm ctx})^5.}
 \label{sel25:eq:condition-five}
\end{equation}
These are exact singular-value identities.  They show why algebraic
identifiability of the five-factor event does not promise practical precision.
A source-native frame that is only mildly ill-conditioned on one factor can be
severely ill-conditioned after tensorization.  V24's target-specific
normalization-aware $\ell^1$ reconstruction remains preferable whenever only a
few determinant functionals are needed.

\begin{proof}
The nonzero singular values of $A\otimes B$ are all products
$s_i(A)s_j(B)$.  Their largest-to-smallest ratio is therefore the product of
the two condition numbers.  Iterate five times.
\end{proof}

\section{Result and next upstream source row}
\label{sel25:status}

V25 removes a stronger premise than V24.  A hidden support-minimal memory does
not have to be promoted to a directly controllable exposed port before its
actual transition can be tomographed.  Prefix histories already prepare
positive memory states, future words already define positive effects, and a
full Hankel rank proves that these contexts span the complete hidden operator
space.  The reconstruction is invariant under the memory unitary left
undetermined by minimal process-comb realization.

For the fourteen-dimensional source class the exact rank-$196$ certificate is
therefore an operator-tomography certificate as well as a memory-dimension
certificate.  On any historical source where that rank is actually populated,
all actual same-cut single-memory events can be reconstructed from ordinary
cylinders without adding hidden-memory testers.

The determinant frontier is correspondingly sharper:
\[
 \boxed{
 \begin{gathered}
 \text{the missing object is not direct five-copy port exposure,}\\
 \text{but one actual co-instantiated five-factor joint source row.}
 \end{gathered}}
\]
If that row occurs and the product prefix--suffix contexts occur around it,
V13-style single-factor saturation tensorizes and V22--V24's determinant
functionals become historical-cylinder observables.  If the joint row does not
occur, V17 proves that no archive of subcritical marginals can manufacture it.

Thus the next historical audit should search the event grammar and same-history
readable packet for a genuine five-factor joint incidence/transition, or for a
source-complete mixed word whose represented Choi packet is equivalent to one.
Only after such occurrence is established should its contextual panel be
populated.  Repeating the one-copy source five times is not a substitute.

Historical Standard-Model selection is still not declared closed.  The exact
V13 point is a mathematical witness, not historical population; parallel
context closure and the five-factor event have not been supplied.  Charged
matter naming, fermionic grading and continuum dynamics remain independent
interfaces.  The gain is a more upstream stopping rule: failure now returns
\emph{joint-event absence or contextual-rank deficiency}, rather than an
artificial demand for direct access to a hidden memory.

\paragraph{Proof and artifact status.}
The complete V24 source prefix is retained before its final terminator.  The
V25 checker verifies exact hidden-qubit contextual tomography, memory-unitary
covariance, Kronecker rank/conditioning identities, and a faithful invisible
TP tangent for an incomplete context bank.  The rank-$196$ and five-slot
physical statements use the already proved V13 and V17 theorems and are not
replaced by small numerical examples.  No historical joint-event table is
invented, and no proof-assistant formalization or independent peer review is
claimed.

\addtocontents{toc}{\protect\endgroup}
% END OF SOURCE-SELECTION V25 DEVELOPMENT BLOCK.
% Preserve V1--V25 and append subsequent work before the final terminator.


\clearpage
\begin{center}
{\Large\bfseries Source-selection continuation V26}\\[.4em]
{\large Returning from five-copy proxies to the historical typed incidence:\\
five-slot compression, coherent Choi anchoring, and a finite archaeology test}
\end{center}
\addtocontents{toc}{\protect\begingroup}
\addtocontents{toc}{\protect\renewcommand{\protect\numberline}{\protect\makebox[2.1em][l]}}

\section{Amend the V25 frontier: five representation slots do not require five memory copies}
\label{sel26:context}

V25 gives a correct sufficient route: if five copies of a saturated hidden
memory are physically co-instantiated, product contexts tomograph any actual
five-copy joint event, while V17 proves that separate lower-order marginals do
not manufacture such an event.  The historical audit in V110 of
\cite{Attack}, however, shows that this route is not source minimal for the
Standard-Model determinant question.  The determinant character is already the
fully alternating component of one \emph{typed up-incidence operator}
\[
 T:C\otimes W\otimes W\longrightarrow\Lambda^2C^*,
 \qquad \dim C=3,\quad\dim W=2,
\]
carried by one same-history Store--Clifford word.  V16, independently, already
identified the extra coherence needed when that represented operator is
forgotten down to a phase-blind CP packet.  The present continuation joins
these two facts and corrects the interpretation of V25's ``five-factor row.''

The inherited ingredients are not reproved as new statements.  In particular,
we retain the determinant-incidence isomorphism and finite flat-word
alternative \cite{Attack},
\src{thm:V110-det-incidence-iso}, \src{thm:V110-flat-det}; V16's neutral--
determinant Choi criterion, Theorem~\ref{sel16:thm:anchor}; and V25's
subcritical-marginal no-go, Theorem~\ref{sel25:thm:joint-no-go}.  The new point
is the operational synthesis: the five defining representation slots may live
inside one rectangular typed incidence word.  Five copies of the fourteen-
dimensional comparison memory are only one possible realization of the same
character order.

\begin{proposition}[One typed incidence already has the sharp five-slot determinant order]
\label{sel26:prop:slots}
Let
\[
 \mathscr U=\Hom(C\otimes W\otimes W,\Lambda^2C^*).
\]
Under the central action
\[
 (e^{i\alpha}I_C,e^{i\beta}I_W)\in U(3)\times U(2),
\]
the determinant-incidence line
\[
 \mathscr U_{\det}\simeq(\det C\otimes\det W)^*
\]
has character
\begin{equation}
 \boxed{e^{-i(3\alpha+2\beta)}.}
 \label{sel26:eq:character}
\end{equation}
It uses exactly five defining representation slots: one colour and two weak
slots on the input, and two dual colour slots in the output
$\Lambda^2C^*$.  Hence a nonzero vector on this line has stabilizer
\begin{equation}
 \boxed{3\alpha+2\beta=0,\qquad
        S(U(3)\times U(2)).}
 \label{sel26:eq:stabilizer}
\end{equation}
The same sharp total slot order that appears in V17--V18 is therefore realized
by one $12\to3$ typed operator.  It does not require a Hilbert memory
$H_{14}^{\otimes5}$.
\end{proposition}

\begin{proof}
The input $C\otimes W\otimes W$ has central weight
$\alpha+2\beta$, while $\Lambda^2C^*$ has weight $-2\alpha$.
The induced Hom character is target minus source, namely
$-3\alpha-2\beta$.  Equivalently, complete alternation produces one copy of
$(\det C\otimes\det W)^*$.  The line is fixed exactly when its character is
one.  Its three colour slots consist of the input $C$ and the two dual colour
slots in the exterior-square output; its two weak slots are the two input
copies of $W$.
\end{proof}

\begin{remark}[Explicit amendment to the V25 stopping rule]
V25's five-copy joint event remains a valid sufficient comparison and its
marginal no-go remains unchanged.  What is amended is the claim that such a
five-copy event is the \emph{next necessary historical object}.  For the
structural determinant packet, an actual typed $12\to3$ up-incidence word is a
strictly smaller candidate: it already carries the five defining slots on one
same-history operation.
\end{remark}

\section{The operator-word route: one alternating shadow is enough}
\label{sel26:word}

Let $\mathfrak A_{\det}:\mathscr U\to
(\det C\otimes\det W)^*$ be the complete alternation from V110, and let
$P_{\det}$ be the Hilbert--Schmidt orthogonal projection onto
$\mathscr U_{\det}$.  The inherited normalization is
\begin{equation}
 \boxed{\|P_{\det}T\|_{\HS}^2
       =6\|\mathfrak A_{\det}T\|^2.}
 \label{sel26:eq:norm}
\end{equation}
Thus
\begin{equation}
 m_{\det}(T):=\|\mathfrak A_{\det}T\|^2
              =\frac16\|P_{\det}T\|_{\HS}^2
 \label{sel26:eq:mdet}
\end{equation}
is phase independent.

\begin{theorem}[Historical typed-word determinant test without a five-copy port]
\label{sel26:thm:word}
Assume the already reconstructed internal carriers $C,W$ and let
$\mathfrak b$ be an actual same-history represented operation word on the
geometry-shorted packet.  Suppose its typed up-incidence projection
\[
 T_u(\mathfrak b)\in\mathscr U
\]
is part of the reconstructed history action algebra.  Then
\begin{equation}
 \boxed{m_{\det}(\mathfrak b)
 :=\|\mathfrak A_{\det}T_u(\mathfrak b)\|^2>0}
 \label{sel26:eq:word-positive}
\end{equation}
if and only if that same historical word contains a nonzero protected
determinant-incidence shadow.  On the positive branch the determinant tensor is
reconstructed, up to its irrelevant nonzero scalar, as
\[
 \tau_{\mathfrak b}=\mathfrak A_{\det}T_u(\mathfrak b),
\]
and its stabilizer is \eqref{sel26:eq:stabilizer}.  No parallel five-copy
memory, product context bank, or independently added determinant word is
required.
\end{theorem}

\begin{proof}
This is the V110 determinant-incidence isomorphism applied to an actually
occurring represented word rather than to an added comparison tensor.  The
alternation is a canonical equivariant linear projection.  Its norm is zero
exactly when the determinant component vanishes, by
\eqref{sel26:eq:norm}.  When it is nonzero, Proposition~\ref{sel26:prop:slots}
gives its stabilizer.  The statement uses the word as an element of the
reconstructed action algebra; it does not replace that operator by its
phase-blind CP square.
\end{proof}

The last sentence is essential.  V16 already showed that an isolated
character-valued amplitude and its CP square have different stabilizers.  The
operator-word route is legitimate only when the Grand-Tensor reconstruction
retains the represented operation itself as part of the history algebra.
When only a CP outcome packet is physically retained, the next section gives
the intrinsic replacement.

\begin{proposition}[A basis-free scalar flat-word archaeology test]
\label{sel26:prop:flat}
Let $\mathcal W_r^{\rm odd}$ be the finite labelled typed operation-word span
through depth $r$, with its intrinsic Hilbert--Schmidt Gram, and define
\[
 L_r:\mathcal W_r^{\rm odd}\longrightarrow
 (\det C\otimes\det W)^*,
 \qquad L_rw=\mathfrak A_{\det}T_u(w).
\]
Put
\begin{equation}
 \boxed{M_r:=\Tr(L_r^*L_r).}
 \label{sel26:eq:Mr}
\end{equation}
Then $M_r\ge0$ and is independent of the orthonormal word basis.  Moreover:
\begin{enumerate}[label=\textup{(\roman*)},leftmargin=2.8em]
\item $M_r>0$ exactly when some historical word of depth at most $r$ has a
      nonzero determinant alternation;
\item if
      $\rank\mathcal W_r^{\rm odd}=\rank\mathcal W_{r+1}^{\rm odd}$ and
      $M_{r+1}=0$, then $m_{\det}(w)=0$ for every word in the complete generated
      historical operation algebra.
\end{enumerate}
Thus the V110 flat-zero alternative can be evaluated as one basis-invariant
nonnegative Gram trace instead of a basis-by-basis search.
\end{proposition}

\begin{proof}
For any orthonormal basis $(w_j)$,
$M_r=\sum_j\|L_rw_j\|^2$, proving positivity, basis independence and (i).
At the first equality of consecutive word ranks the represented word span is
flat and is a unital $*$-algebra containing every later word.  If $M_{r+1}=0$,
then $L_{r+1}=0$; linearity of the typed projection and alternation therefore
makes $L$ vanish on the entire generated algebra.  This is the inherited V110
flat-depth argument in basis-free Gram form.
\end{proof}

\section{If only the CP packet survives, diagonal determinant mass is not enough}
\label{sel26:choi}

We now formulate the exact operational distinction that reconciles V110 with
V16.  Work after the inherited spacetime and nuisance shorts.  Let
$\mathscr K_{\rm Cl}$ be the one-dimensional normalized Store--parity Clifford
shadow line with unit vector $\widehat B$, and let
$\mathscr K_{\det}=\mathscr U_{\det}$ have normalized unit vector
$\widehat\Theta$.  Write
\[
 P_B=|\widehat B\rangle\!\rangle\langle\!\langle\widehat B|,
 \qquad
 P_D=|\widehat\Theta\rangle\!\rangle
              \langle\!\langle\widehat\Theta|.
\]
The first line is neutral under $G_0=U(3)\times U(2)$, while the second carries
$\chi^{-1}$, $\chi(U,V)=\det U\det V$.

For any positive marked Choi packet $J$ on the direct sum of these two
coefficient sectors define
\begin{equation}
 \mu_J=\Tr(P_BJ),\qquad
 \nu_J=\Tr(P_DJ),\qquad
 C_J=P_BJP_D,
 \qquad
 h_J=\|C_J\|_{\HS}^2.
 \label{sel26:eq:choi-data}
\end{equation}

\begin{theorem}[Factorization-free neutral--determinant coherence test]
\label{sel26:thm:choi}
For every $J\succeq0$ in the preceding two-sector packet,
\begin{equation}
 \boxed{0\le h_J\le\mu_J\nu_J.}
 \label{sel26:eq:coh-bound}
\end{equation}
Furthermore, inside $G_0$:
\begin{enumerate}[label=\textup{(\roman*)},leftmargin=2.8em]
\item if $h_J>0$, every fixed-label covariance of $J$ satisfies
      $\chi=1$;
\item if the packet contains only these two sectors, then
      \begin{equation}
       \boxed{h_J>0\iff\operatorname{Cov}_{G_0}(J)=\ker\chi,
       \qquad
       h_J=0\iff\operatorname{Cov}_{G_0}(J)=G_0.}
       \label{sel26:eq:choi-group}
      \end{equation}
\item for the coherent rank-one word
      \[
       |K\rangle\!\rangle
       =\sqrt\mu\,|\widehat B\rangle\!\rangle
        +e^{i\theta}\sqrt\nu\,
          |\widehat\Theta\rangle\!\rangle,
      \]
      one has $h_J=\mu\nu$;
\item the incoherent packet
      \[
       J_{\rm mix}=\mu P_B+\nu P_D
      \]
      has the \emph{same} two diagonal masses but $h_{J_{\rm mix}}=0$.
\end{enumerate}
With additional already reconstructed typed sectors, $h_J>0$ still forces the
covariance group to be contained in $\ker\chi$; equality additionally requires
those sectors to satisfy their inherited determinant-compatible residuals.
\end{theorem}

\begin{proof}
Compress positivity of $J$ to the two one-dimensional lines.  Its scalar block
matrix is
\[
 \begin{pmatrix}\mu_J&c\\\bar c&\nu_J\end{pmatrix}\succeq0,
 \qquad h_J=|c|^2,
\]
so $|c|^2\le\mu_J\nu_J$.  Under $g\in G_0$, the upper-right block transforms
as $c\mapsto\chi(g)c$.  Thus invariance with $c\ne0$ forces $\chi(g)=1$.
When no other sectors are present, every element of $\ker\chi$ acts trivially
on both displayed coefficient lines, proving \eqref{sel26:eq:choi-group}.
The coherent rank-one and incoherent formulas follow by direct block
multiplication.
\end{proof}

\begin{corollary}[Operational meaning of the inherited one-ray construction]
\label{sel26:cor:one-ray}
The source-minimal V110 comparison ray
\[
 z\longmapsto z\left(\sqrt\mu\,\widehat B
                 \oplus\sqrt\nu\,\widehat\Theta\right)
\]
has, when it is realized as one actual coherent marked quantum word,
\begin{equation}
 \boxed{h_J=\mu\nu>0.}
 \label{sel26:eq:one-ray-h}
\end{equation}
Thus it lands exactly on V16's positive neutral--determinant coherence branch.
If the two shadows are instead retained only as classically distinguished or
incoherently mixed alternatives, their marginal masses remain $\mu,\nu$ but
$h_J=0$.  Source-rank counting alone therefore does not establish historical
coherence; the actual operation-word algebra or its Choi cross block must do so.
\end{corollary}

\section{Historical cylinders can reconstruct the Choi coherence directly}
\label{sel26:contextual}

V25 now becomes useful in a smaller way than its five-copy construction.  Let
one actual historical word $\mathfrak b$ have an input memory and output memory
with context families $(\sigma_p)$ and $(E_f)$.  Its original cylinders are
\[
 p(f\mathfrak b p)=\Tr[E_f\Phi_{\mathfrak b}(\sigma_p)].
\]
If the context Hankel is saturated, or merely if the target functionals below
lie in the normalization-aware contextual tester span, V25 reconstructs the
required Choi blocks without exposing that memory as a new port.

\begin{theorem}[Cylinder-level determinant test for a historical typed word]
\label{sel26:thm:cylinder}
Assume the typed Clifford and up-incidence projections of one actual marked
word $\mathfrak b$ are identified on its support-minimal memory.
\begin{enumerate}[label=\textup{(\alph*)},leftmargin=2.8em]
\item If the reconstructed history algebra retains $\mathfrak b$ as a coherent
      represented operator, then the scalar
      $m_{\det}(\mathfrak b)$ of \eqref{sel26:eq:word-positive} is the direct
      historical occurrence test.
\item If only its positive Choi packet $J_{\mathfrak b}$ is retained, then the
      factorization-free scalar
      \[
       h_{\mathfrak b}
       =\|P_BJ_{\mathfrak b}P_D\|_{\HS}^2
      \]
      is reconstructible from the old cylinder table whenever the relevant
      contextual tester span is complete for these block functionals.  A full
      $d^2$ context Hankel rank is sufficient.
\item On a rank-one efficient branch
      $J_{\mathfrak b}=|K_{\mathfrak b}\rangle\!\rangle
      \langle\!\langle K_{\mathfrak b}|$, the diagonal determinant Choi mass
      obeys
      \begin{equation}
       \boxed{\Tr(P_DJ_{\mathfrak b})
              =\|P_{\det}T_u(K_{\mathfrak b})\|_{\HS}^2
              =6m_{\det}(K_{\mathfrak b}),}
       \label{sel26:eq:rankone-mass}
      \end{equation}
      independent of the unresolved global Kraus phase.
\end{enumerate}
Hence no $14^5$ contextual panel is required when the historical determinant
is carried by one typed operation word.  What must occur physically is that
word, together with enough old contexts to determine the relevant operator or
Choi block.
\end{theorem}

\begin{proof}
Part (a) is Theorem~\ref{sel26:thm:word}.  For (b), each matrix element of the
Choi packet is a real-linear functional of $J_{\mathfrak b}$ and V25's
contextual theorem reconstructs every functional in the contextual tester span;
$h_{\mathfrak b}$ is then computed from the reconstructed cross block.  Full
Hankel rank $d^2$ makes the context states and effects Hermitian bases and is
therefore sufficient.  For (c), a rank-one Choi operator determines its
coefficient up to one global phase.  Orthogonal projection onto the determinant
line commutes with that phase, and the inherited normalization
\eqref{sel26:eq:norm} gives the factor six.
\end{proof}

\begin{remark}[Where the V13 rank-$196$ certificate can and cannot be used]
If the typed word in question acts on the same saturated fourteen-dimensional
memory of V13, its $196=14^2$ original context rank is already sufficient for
part (b).  If the Store--Clifford incidence lives on a different typed quotient,
that identification must be proved and its own context span audited.  Equality
of Hilbert dimensions is not a carrier-identification theorem.
\end{remark}

\section{A sharper finite historical stopping rule}
\label{sel26:stopping}

The preceding results replace V25's parallel-five-copy frontier by a smaller
branching audit.

\begin{theorem}[Operator-versus-CP historical determinant alternative]
\label{sel26:thm:alternative}
After the inherited internal-algebra, weak-reset and typed-carrier steps, audit
the actual odd labelled word hierarchy as follows.
\begin{enumerate}[label=\textup{(S\arabic*)},leftmargin=3em]
\item If the Grand-Tensor history algebra reconstructs coherent represented
      words, evaluate the positive Gram trace $M_r$ of
      Proposition~\ref{sel26:prop:flat} while the word span grows.
      The first $M_r>0$ lands a determinant-incidence word.  A flat span with
      $M_{r+1}=0$ is a complete historical determinant no-go on that algebra.
\item If the physical record retains only CP packets for the relevant words,
      reconstruct their Choi packets from the old contexts and test
      $h_{\mathfrak b}$.  A nonzero $h_{\mathfrak b}$ lands the primitive
      determinant phase relative to the neutral Clifford shadow.  Positive
      diagonal determinant mass with $h_{\mathfrak b}=0$ is not enough.
\item Invoke a parallel five-copy event only if such a joint event is itself an
      actual historical row.  It is no longer a prerequisite of the determinant
      archaeology.
\end{enumerate}
No step is licensed to replace a classicalized CP record by an unobserved
coherent operator, or to replace a missing typed word by five independent
repetitions.
\end{theorem}

\begin{proof}
(S1) is the flat algebra argument in Proposition~\ref{sel26:prop:flat}.
(S2) is Theorem~\ref{sel26:thm:choi} combined with V25 contextual tomography.
(S3) follows from Proposition~\ref{sel26:prop:slots} and preserves the V17/V25
no-go against manufacturing an unknown joint event from marginals.
\end{proof}

\begin{example}[Same shadow masses, opposite historical answers]
\label{sel26:ex:coherent-mix}
Take $\mu,\nu>0$ and the two normalized coefficient rays above.  The coherent
packet
\[
 J_{\rm coh}
 =\left|\sqrt\mu\,\widehat B+\sqrt\nu\,\widehat\Theta\right\rangle\!\rangle
  \left\langle\!\left\langle
    \sqrt\mu\,\widehat B+\sqrt\nu\,\widehat\Theta
  \right|\right.
\]
and the incoherent packet
\[
 J_{\rm mix}=\mu P_B+\nu P_D
\]
have identical Clifford and determinant diagonal masses.  Yet
\[
 h(J_{\rm coh})=\mu\nu,
 \qquad h(J_{\rm mix})=0.
\]
Thus marginal shadow occurrence does not decide historical determinant phase
selection.  The distinguishing datum is precisely the cross block retained by
one coherent word or by its output-sensitive Choi packet.
\end{example}

\section{Result and next upstream row}
\label{sel26:status}

V26 changes the source-selection frontier materially.  The number five that
appears in the determinant character is a \emph{representation-slot order}, not
a theorem that five copies of the fourteen-dimensional memory must coexist.
The inherited up-type incidence
\[
 C\otimes W\otimes W\longrightarrow\Lambda^2C^*
\]
already realizes all five defining slots in one rectangular typed operation.
Its complete alternation is the determinant line.

Accordingly, the historical programme now has two legitimate entrances.
If the reconstructed Grand Tensor contains the actual coherent typed operation
word, its alternating mass $m_{\det}$ is the direct finite test, with the flat
word hierarchy giving a terminating positive-or-no-go alternative.  If the
record retains only an outcome-sensitive CP packet, the intrinsic
neutral--determinant Choi cross block $h$ is the required operational phase
anchor.  V25's hidden-memory contextual tomography can reconstruct this block
from ordinary old cylinders whenever the relevant context span is saturated.

This synthesis also sharpens the status of the V110 active one-ray completion.
One source ray carrying nonzero Clifford and determinant shadows is indeed the
source-minimal coherent realization, and on that realization
$h=\mu\nu>0$.  But the same two marginal shadow masses can be classically
separated, giving $h=0$ and retaining the larger product-group CP covariance.
Historical occurrence therefore means occurrence of the coherent word or its
cross block, not merely occurrence of both marginals.

The next upstream audit should therefore search the \emph{actual typed odd
operation-word hierarchy} on the already reconstructed internal carriers,
apply the canonical up-incidence projection and complete alternation, and stop
at the first flat depth.  Only if the operational formalism exposes no coherent
word algebra should one fall back to Choi-block reconstruction.  A five-copy
parallel source should be audited only when it actually occurs in the event
grammar; it is no longer the default missing object.

Historical Standard-Model selection is still not declared closed.  V26 does
not prove that the locked historical Store contains the requisite typed odd
word, nor does it populate its alternating mass or Choi cross block.  It
removes a stronger and partly artificial premise: five-memory co-instantiation
is not required merely because the determinant character has total order five.
The remaining obstruction is one same-history typed operation occurrence on
the already selected carriers, exactly the finite historical row isolated by
V110.

\paragraph{Proof and artifact status.}
The V25 prefix is retained before its final terminator.  The V26 checker
verifies the explicit $3\times12$ determinant-incidence matrix and its norm six,
the five-slot central character, the coherent-versus-incoherent Choi cross
blocks, the positivity bound $h\le\mu\nu$, and basis invariance of the flat-word
Gram trace.  The historical flat-depth implication itself is the written
finite-algebra proof above and the inherited V110 theorem, not a numerical
extrapolation.  No historical typed word or probability value is invented, and
no proof-assistant formalization or independent peer review is claimed.

\addtocontents{toc}{\protect\endgroup}
% END OF SOURCE-SELECTION V26 DEVELOPMENT BLOCK.
% Preserve V1--V26 and append subsequent work before the final terminator.



\clearpage
\begin{center}
{\Large\bfseries Source-selection continuation V27}\\[.4em]
{\large Resolved efficiency, automatic determinant coherence,\\
and the exact label-erasure boundary}
\end{center}
\addtocontents{toc}{\protect\begingroup}
\addtocontents{toc}{\protect\renewcommand{\protect\numberline}{\protect\makebox[2.1em][l]}}

\section{Go upstream from a phase test to the resolution of the historical branch}
\label{sel27:context}

V26 separates two historical regimes.  A coherent represented operator word can
be tested directly by its alternating up-incidence mass.  If that operator has
been forgotten down to a positive Choi packet, V26 instead retains the
neutral--determinant cross block
\[
 h_J=\|P_BJP_D\|_{\rm HS}^2.
\]
This is exact but can still be downstream.  A resolved quantum outcome is often
\emph{efficient}: one labelled branch has one Kraus coefficient.  Compositions
of such labelled branches remain efficient.  On that branch a positive neutral
shadow and a positive determinant shadow cannot be mutually incoherent; their
cross block is forced by rank one.

This continuation proves that statement without assuming a phase convention,
gives a factorization-free Choi-purity test for efficiency, and quantifies how
much mixing is allowed before a separate cross-block audit becomes necessary.
It also isolates the exact record operation that can destroy the conclusion:
coarse graining several individually coherent labelled branches can cancel the
relative determinant phase while preserving both diagonal shadow masses.

The retained inputs are V26's two-sector Choi theorem and its historical
operator-versus-CP alternative, together with the inherited rule that a
historical claim must use the actually retained outcome labels.  No claim is
made here that every primitive Store outcome is efficient.  Efficiency is an
independently checkable property of the physical marked branch, not a
replacement for its occurrence.

\section{Resolved efficient words stay rank one}
\label{sel27:efficient}

\begin{definition}[Resolved efficient branch]
A marked completely positive branch $e$ is \emph{efficient} when there is one
operator $K_e$ such that
\[
 \Phi_e(X)=K_eXK_e^*.
\]
The branch is \emph{resolved} when the label $e$ is retained rather than summed
with other alternatives.  For a labelled word
$w=e_m\cdots e_1$ put
\[
 K_w=K_{e_m}\cdots K_{e_1}.
\]
\end{definition}

\begin{theorem}[Efficiency is closed under labelled composition]
\label{sel27:thm:composition}
If every branch in the resolved word $w=e_m\cdots e_1$ is efficient, then
\begin{equation}
 \boxed{\Phi_w(X)=K_wXK_w^*,\qquad
 J_w=|K_w\rangle\!\rangle\langle\!\langle K_w|.}
 \label{sel27:eq:word-rankone}
\end{equation}
Thus either $K_w=0$ or the Choi rank of the complete labelled word is exactly
one.  No relative phase between two linear shadows of $K_w$ is an additional
physical datum.
\end{theorem}

\begin{proof}
Composition gives
$\Phi_{e_m}\circ\cdots\circ\Phi_{e_1}(X)
 =K_{e_m}\cdots K_{e_1}XK_{e_1}^*\cdots K_{e_m}^*$.
Vectorization gives the Choi formula.  Multiplying $K_w$ by one global phase
leaves $J_w$ unchanged and multiplies every linear projection of $K_w$ by the
same phase; it cannot change a relative phase between two projections of that
same coefficient.
\end{proof}

\begin{proposition}[A factorization-free purity scalar]
\label{sel27:prop:purity}
For every nonzero positive Choi matrix $J$ define
\begin{equation}
 \boxed{\Pi(J):=(\Tr J)^2-\Tr(J^2)\ge0.}
 \label{sel27:eq:purity-defect}
\end{equation}
Then
\begin{equation}
 \boxed{\Pi(J)=0\iff\rank J=1.}
 \label{sel27:eq:purity-rank}
\end{equation}
The scalar is invariant under all input/output memory unitaries and under every
change of Kraus factorization.  Hence efficiency of a reconstructed nonzero
historical CP packet can be certified without selecting a Kraus frame.
\end{proposition}

\begin{proof}
If $\lambda_i\ge0$ are the eigenvalues of $J$, then
\[
 \Pi(J)=2\sum_{i<j}\lambda_i\lambda_j.
\]
It vanishes exactly when at most one eigenvalue is nonzero.  The remaining
claims follow because the expression depends only on the Choi spectrum.
\end{proof}

\section{Purity forces neutral--determinant coherence}
\label{sel27:purity-coherence}

Let $P_B,P_D$ be orthogonal coefficient-sector projections, with $P_D$ the
determinant-incidence sector and $P_B$ the retained neutral Store--Clifford
sector.  For an arbitrary positive Choi packet put
\begin{equation}
 \mu=\Tr(P_BJ),\qquad
 \nu=\Tr(P_DJ),\qquad
 h=\|P_BJP_D\|_{\rm HS}^2.
 \label{sel27:eq:masses}
\end{equation}
V26 proves $h\le\mu\nu$.  The missing converse is controlled by the total
Choi mixedness.

\begin{theorem}[Global purity controls the missing determinant coherence]
\label{sel27:thm:purity-bound}
For every $J\succeq0$ and orthogonal projections $P_B,P_D$,
\begin{equation}
 \boxed{0\le \mu\nu-h\le\frac12\Pi(J).}
 \label{sel27:eq:purity-bound}
\end{equation}
Consequently:
\begin{enumerate}[label=\textup{(\roman*)},leftmargin=2.8em]
\item if $J$ is nonzero and rank one, then
      \begin{equation}
       \boxed{h=\mu\nu;}
       \label{sel27:eq:rankone-coherence}
      \end{equation}
\item more generally,
      \begin{equation}
       \boxed{2\mu\nu>\Pi(J)\quad\Longrightarrow\quad h>0.}
       \label{sel27:eq:mixed-sufficient}
      \end{equation}
\end{enumerate}
The sufficient inequality is frame independent and uses only the two diagonal
sector masses and the Choi purity defect.
\end{theorem}

\begin{proof}
Choose an orthonormal basis adapted to
$P_B\oplus P_D\oplus(I-P_B-P_D)$.  For every positive matrix,
\begin{equation}
 \frac12\Pi(J)
 =\sum_{i<j}\left(J_{ii}J_{jj}-|J_{ij}|^2\right),
 \label{sel27:eq:minor-sum}
\end{equation}
where every summand is a nonnegative $2\times2$ principal minor.  Summing only
the pairs with $i$ in $P_B$ and $j$ in $P_D$ gives exactly
\[
 \sum_{i\in B,j\in D}
 \left(J_{ii}J_{jj}-|J_{ij}|^2\right)
 =\mu\nu-h.
\]
This proves both inequalities.  If $J$ has rank one, Proposition~\ref{sel27:prop:purity}
gives $\Pi(J)=0$, so equality follows.  The strict sufficient condition is the
same bound rearranged.
\end{proof}

\begin{corollary}[No independent phase row on a resolved efficient determinant word]
\label{sel27:cor:no-phase-row}
Let $\mathfrak b$ be one actual resolved historical word whose CP branch is
efficient.  If its neutral Store--Clifford shadow and determinant-incidence
shadow both have positive Choi masses,
\[
 \mu_{\mathfrak b}>0,
 \qquad
 \nu_{\mathfrak b}>0,
\]
then
\begin{equation}
 \boxed{h_{\mathfrak b}
 =\mu_{\mathfrak b}\nu_{\mathfrak b}>0.}
 \label{sel27:eq:automatic-h}
\end{equation}
Therefore V26's determinant covariance conclusion follows without executing or
reconstructing a separate phase-sensitive cross-block observable.  In the V110
one-word branch, positive Store--Clifford occurrence and
$m_{\det}(\mathfrak b)>0$ suffice for this conclusion whenever the same marked
word is independently certified efficient.
\end{corollary}

\begin{remark}[What is and is not removed]
The corollary removes a \emph{relative phase test}; it does not infer that the
historical word occurs, that its two typed shadows are positive, or that its
CP branch is efficient.  If any one of those properties is unresolved, it
remains an actual source row.  Additional typed sectors must still satisfy the
inherited determinant-compatible residuals before the complete covariance
group can be identified with $S(U(3)\times U(2))$ rather than a subgroup.
\end{remark}

\section{A quantitative fallback for a mixed historical packet}
\label{sel27:mixed}

Theorem~\ref{sel27:thm:purity-bound} is useful even when the branch is not
exactly efficient.  Suppose historical contextual tomography supplies certified
bounds
\[
 \mu\ge\mu_->0,
 \qquad
 \nu\ge\nu_->0,
 \qquad
 \Pi(J)\le\Pi_+.
\]
Then
\begin{equation}
 \boxed{
 2\mu_-\nu_->\Pi_+
 \quad\Longrightarrow\quad
 h\ge\mu_-\nu_- -\frac{\Pi_+}{2}>0.}
 \label{sel27:eq:robust-purity}
\end{equation}
Thus a sufficiently pure packet can establish determinant coherence from
phase-insensitive scalar information alone.  Exact zero purity is not required.

\begin{proposition}[Sharpness on the two-sector packet]
\label{sel27:prop:sharp}
If $P_B$ and $P_D$ are one-dimensional and $J$ has support only on their sum,
then
\begin{equation}
 \boxed{\Pi(J)=2(\mu\nu-h).}
 \label{sel27:eq:two-sector-exact}
\end{equation}
Hence
\[
 2\mu\nu>\Pi(J)\iff h>0
\]
is exact in the two-sector model.  The incoherent diagonal packet saturates the
boundary $\Pi(J)=2\mu\nu$, while every rank-one coherent packet has
$\Pi(J)=0$.
\end{proposition}

\begin{proof}
In an adapted basis the packet is
$\left(\begin{smallmatrix}\mu&c\\\bar c&\nu\end{smallmatrix}\right)$,
with $h=|c|^2$.  Direct expansion gives
$(\mu+\nu)^2-(\mu^2+\nu^2+2|c|^2)=2(\mu\nu-h)$.
\end{proof}

\section{Coarse graining can erase coherence that every resolved branch possesses}
\label{sel27:coarse}

The remaining distinction is record theoretic, not algebraic.  Rank-one
closure applies to each \emph{resolved} word.  Summing different words can
increase Choi rank and cancel their determinant cross blocks.

\begin{theorem}[Exact label-erasure boundary]
\label{sel27:thm:erasure}
Let $\widehat B,\widehat\Theta$ be orthonormal coefficient rays and let
$\mu,\nu>0$.  Define two efficient labelled success branches
\[
 K_+=\sqrt\mu\,\widehat B+\sqrt\nu\,\widehat\Theta,
 \qquad
 K_-=\sqrt\mu\,\widehat B-\sqrt\nu\,\widehat\Theta.
\]
Each resolved Choi packet satisfies
\[
 h_+=h_-=\mu\nu>0.
\]
If the two labels are forgotten, their summed success packet is
\begin{equation}
 \boxed{
 J_{\rm forget}=J_++J_-
 =2\mu P_B+2\nu P_D,\qquad h_{\rm forget}=0.}
 \label{sel27:eq:erasure}
\end{equation}
After multiplying both success branches by a sufficiently small common scalar,
one may add one failure operation and obtain a normalized physical instrument
with exactly this resolved-versus-forgotten distinction.
\end{theorem}

\begin{proof}
The two rank-one cross blocks are respectively
$+\sqrt{\mu\nu}\,|\widehat B\rangle\!\rangle
\langle\!\langle\widehat\Theta|$ and its negative.  They cancel in the sum,
while both diagonal blocks add.  Scaling the success effects until their sum is
bounded by the identity leaves a positive remainder effect; a standard reset
failure branch completes the instrument without changing the displayed success
Choi relation.
\end{proof}

\begin{corollary}[The historical question is same-label provenance]
\label{sel27:cor:provenance}
When the primitive outcome bank is efficient, the correct determinant audit is:
\begin{enumerate}[label=\textup{(P\arabic*)},leftmargin=3em]
\item keep the original resolved labels through the complete word;
\item test whether one and the same labelled word has positive neutral and
      determinant masses;
\item only after this test may labels be forgotten for downstream averaged
      dynamics.
\end{enumerate}
If the two shadows occur only in different alternatives, their separate
positivity does not establish determinant phase selection.  Conversely, on one
resolved efficient word no additional phase calibration is required.
\end{corollary}

\section{Result and next upstream row}
\label{sel27:status}

V27 moves the historical frontier from ``measure the neutral--determinant
phase'' to the more primitive question ``what did the original record resolve?''
For one labelled efficient word, rank-one Choi structure forces maximal
coherence between every pair of its nonzero linear shadows.  The determinant
cross block is then not an independent physical datum.  The V110 one-word
conditions reduce to occurrence of the same resolved branch, positive
Store--Clifford and determinant shadows, and an efficiency certificate.

When efficiency is not exact, the scalar purity defect
\[
 \Pi(J)=(\Tr J)^2-\Tr J^2
\]
provides a phase-free fallback.  The universal bound
\[
 h\ge\mu\nu-\Pi(J)/2
\]
can certify a nonzero determinant cross block without reconstructing its
complex phase.  In the isolated two-sector problem this criterion is sharp.

The negative result is equally structural.  Coarse graining can sum two
individually efficient coherent branches and erase the cross block exactly.
Therefore an averaged CP map cannot be promoted back to an efficient historical
word merely because it admits a one-Kraus decomposition after fitting.  The
original outcome resolution is part of the source provenance.

The next upstream audit should therefore classify the \emph{actual primitive
label bank} feeding the V110 typed odd hierarchy: determine which primitive
outcomes are efficient, which labels survive through the Store--Clifford word,
and whether the positive Clifford and determinant shadows land in the same
resolved word before any coarse graining.  If they do, V27 removes the separate
phase row.  If they do not, V26's cross-block test remains irreducible.  No new
carrier or determinant tensor should be introduced to repair a failure of label
provenance.

Historical Standard-Model selection is still not declared closed.  The result
is a reduction in source cost and a sharper stopping rule, not evidence that the
required resolved word already occurs in the historical Store.

\paragraph{Proof and artifact status.}
The complete V26 prefix is retained before its final terminator.  The V27
checker verifies rank-one closure under labelled composition, the exact purity
identity, the global principal-minor bound, sharpness on the two-sector packet,
and complete cancellation under the explicit two-label coarse graining.  These
finite checks support the written general proofs; they are not a
proof-assistant formalization or independent peer review.

\addtocontents{toc}{\protect\endgroup}
% END OF SOURCE-SELECTION V27 DEVELOPMENT BLOCK.
% Preserve V1--V27 and append subsequent work before the final terminator.



\clearpage
\begin{center}
{\Large\bfseries Source-selection continuation V28}\\[.4em]
{\large Contextual efficiency discriminants, composite Kraus collapse,\\
and a direct historical-cylinder determinant test}
\end{center}
\addtocontents{toc}{\protect\begingroup}
\addtocontents{toc}{\protect\renewcommand{\protect\numberline}{\protect\makebox[2.1em][l]}}

\section{Go upstream from primitive efficiency to the candidate historical word}
\label{sel28:context}

V27 shows that a resolved efficient word has Choi rank one and therefore cannot
hide an incoherent separation between two nonzero linear shadows of the same
coefficient.  Its proposed next audit was to classify which primitive labelled
outcomes are efficient.  That is sufficient but, in general, stronger than
necessary.  A primitive CP outcome may have several Kraus coefficients and yet
become efficient after composition with the actual incoming or outgoing branch.
Conversely, an efficient prefix may become mixed after a later labelled
operation.  Choi rank of a \emph{composite word} is therefore the exact object;
it is not a monotone count of the ranks of its primitive letters.

This continuation replaces the primitive-efficiency premise by two intrinsic
objects reconstructed on the candidate word itself.  First, a Kraus-product
Gram identifies its exact Choi rank without assigning physical meaning to an
unobserved Kraus label.  Second, V25 contextual tomography turns the Choi-purity
defect and the two typed shadow masses into a quadratic polynomial of ordinary
historical cylinder probabilities.  Combining this with V27 gives one finite,
phase-free historical coherence margin.  No preferred Kraus frame, direct
hidden-memory intervention, or new phase-sensitive Read is required.

The distinction between an operationally primitive label and an extreme ray of
the CP cone is retained throughout.  A mathematical spectral decomposition of
a mixed Choi matrix is not an admitted refinement of the historical record.

\section{The composite Kraus span, not the primitive bank, decides efficiency}
\label{sel28:kraus-span}

Let a resolved word be
\[
 w=e_m\cdots e_1,
\]
and choose arbitrary Kraus representations
\[
 \Phi_{e_j}(X)=\sum_{\alpha_j}
 K^{(j)}_{\alpha_j}XK^{(j)*}_{\alpha_j}.
\]
For a multi-index
$\boldsymbol\alpha=(\alpha_m,\ldots,\alpha_1)$ put
\begin{equation}
 L_{\boldsymbol\alpha}
 :=K^{(m)}_{\alpha_m}\cdots K^{(1)}_{\alpha_1}.
 \label{sel28:eq:products}
\end{equation}
Zero products are retained in formulas but may be discarded from spans.

\begin{theorem}[Exact composite Kraus-span theorem]
\label{sel28:thm:kraus-span}
The resolved composite word has
\begin{equation}
 \boxed{
 J_w=\sum_{\boldsymbol\alpha}
 |L_{\boldsymbol\alpha}\rangle\!\rangle
 \langle\!\langle L_{\boldsymbol\alpha}|.}
 \label{sel28:eq:word-choi}
\end{equation}
If
\begin{equation}
 (\Gamma_w)_{\boldsymbol\alpha,\boldsymbol\beta}
 =\Tr(L_{\boldsymbol\alpha}^*L_{\boldsymbol\beta}),
 \label{sel28:eq:kraus-gram}
\end{equation}
then
\begin{equation}
 \boxed{
 \rank J_w
 =\rank\Gamma_w
 =\dim\Span\{L_{\boldsymbol\alpha}\}.}
 \label{sel28:eq:rank-equality}
\end{equation}
Consequently a nonzero resolved word is efficient if and only if all nonzero
composite coefficients in \eqref{sel28:eq:products} are collinear.  No
factor-wise efficiency assumption is necessary.
\end{theorem}

\begin{proof}
Let $V$ be the matrix whose columns are the vectors
$|L_{\boldsymbol\alpha}\rangle\!\rangle$.  Composition of Kraus forms gives
$J_w=VV^*$, while \eqref{sel28:eq:kraus-gram} is $\Gamma_w=V^*V$.
The two matrices have the same nonzero singular spectrum and hence the same
rank, equal to the column-space dimension of $V$.  Rank one is exactly
collinearity of all nonzero columns.
\end{proof}

\begin{example}[A mixed primitive can collapse to an efficient word]
\label{sel28:ex:collapse}
On a qubit let
\[
 \Delta(X)=P_0XP_0+P_1XP_1,
 \qquad P_j=|j\rangle\langle j|,
\]
so $\Delta$ has Choi rank two.  Precede it by the resolved efficient branch
\[
 \mathcal P_0(X)=P_0XP_0.
\]
The two composite Kraus coefficients are
\[
 P_0P_0=P_0,
 \qquad
 P_1P_0=0,
\]
so
\begin{equation}
 \boxed{\Delta\circ\mathcal P_0=\mathcal P_0}
 \label{sel28:eq:collapse}
\end{equation}
and the composite word has Choi rank one.  Thus requiring every primitive
outcome to be efficient would reject a genuinely efficient historical word.
In the opposite order, $\Delta$ itself remains rank two; composite Choi rank is
not monotone along a word.
\end{example}

\begin{remark}[Primitive does not mean Kraus-extremal]
A primitive label is primitive in the declared event grammar: the physical
record contains no finer admitted outcome at that cut.  It need not be an
extreme ray of the cone of CP maps.  Every positive Choi matrix has a spectral
sum of rank-one positive terms, but regarding those terms as historical
sub-outcomes would append a new record.  Theorem~\ref{sel6:thm:efficient-refinement}
proves rigidity in the opposite direction: when the old branch is already
rank one, every exact record-respecting refinement is only classical.
\end{remark}

\section{The word-efficiency discriminant is a quadratic function of old cylinders}
\label{sel28:cylinders}

Fix one actual candidate typed word $w$ on a support-minimal input/output
memory pair.  By V25, choose historical prefix/future contexts whose Hermitian
Choi testers
\[
 F_a
\]
span the relevant Choi space.  Write the observed contextual cylinders as
\begin{equation}
 p_a(w)=\Tr(F_aJ_w).
 \label{sel28:eq:cylinder-vector}
\end{equation}
Choose any Hermitian dual frame $D^a$ with
\[
 J_w=\sum_a p_a(w)D^a.
\]
This dual is a reconstruction coordinate system; its members need not be
physical effects.

Define the real frame tensors
\begin{equation}
 t_a=\Tr D^a,
 \qquad
 G_{ab}=\Tr(D^aD^b),
 \qquad
 Q:=tt^{\mathsf T}-G.
 \label{sel28:eq:frame-Q}
\end{equation}
Let $P_B$ and $P_D$ be the already identified neutral Store--Clifford and
determinant-incidence coefficient-sector projections and put
\begin{equation}
 b_a=\Tr(P_BD^a),
 \qquad
 d_a=\Tr(P_DD^a).
 \label{sel28:eq:sector-vectors}
\end{equation}
All of these coefficients are fixed once the historical context frame and the
typed sector identification are fixed.

\begin{theorem}[Contextual efficiency discriminant]
\label{sel28:thm:context-purity}
For the actual cylinder vector $p=p(w)$,
\begin{align}
 \boxed{\Pi(J_w)=p^{\mathsf T}Qp,}
 \label{sel28:eq:Pi-cylinder}\\
 \boxed{\mu_w=b^{\mathsf T}p,
 \qquad \nu_w=d^{\mathsf T}p.}
 \label{sel28:eq:masses-cylinder}
\end{align}
Therefore, for a nonzero compatible CP word,
\begin{equation}
 \boxed{
 p^{\mathsf T}Qp=0
 \iff
 \rank J_w=1.}
 \label{sel28:eq:rank-cylinder}
\end{equation}
The value is independent of the chosen spanning context basis and its dual.
It is computed from ordinary historical cylinders; no Kraus decomposition and
no additional quantum experiment is involved.
\end{theorem}

\begin{proof}
Substituting $J_w=\sum_ap_aD^a$ gives
\[
 \Tr J_w=t^{\mathsf T}p,
 \qquad
 \Tr J_w^2=p^{\mathsf T}Gp.
\]
Their difference is \eqref{sel28:eq:Pi-cylinder}.  The two sector masses are
linear contractions of the same reconstruction.  V27's intrinsic purity
proposition gives \eqref{sel28:eq:rank-cylinder}.  Any other spanning context
basis reconstructs the same $J_w$, so the displayed scalars are invariant even
though $Q,b,d$ individually change coordinates.
\end{proof}

\begin{definition}[Historical determinant-coherence margin]
\label{sel28:def:margin}
For one resolved candidate word define
\begin{equation}
 \boxed{
 \mathfrak c(w)
 :=2\mu_w\nu_w-\Pi(J_w)
 =2(b^{\mathsf T}p)(d^{\mathsf T}p)-p^{\mathsf T}Qp.}
 \label{sel28:eq:margin}
\end{equation}
This is a quadratic polynomial of its contextual cylinder probabilities.
\end{definition}

\begin{theorem}[Direct cylinder-level determinant-coherence certificate]
\label{sel28:thm:direct}
Let $h_w=\|P_BJ_wP_D\|_{\rm HS}^2$.  Then
\begin{equation}
 \boxed{h_w\ge\frac12\mathfrak c(w).}
 \label{sel28:eq:margin-bound}
\end{equation}
Hence
\begin{equation}
 \boxed{
 \mathfrak c(w)>0
 \quad\Longrightarrow\quad
 h_w>0.}
 \label{sel28:eq:positive-margin}
\end{equation}
In particular, if
\[
 p^{\mathsf T}Qp=0,
 \qquad b^{\mathsf T}p>0,
 \qquad d^{\mathsf T}p>0,
\]
then
\begin{equation}
 \boxed{h_w=\mu_w\nu_w>0.}
 \label{sel28:eq:efficient-cylinder}
\end{equation}
On a packet supported only on $P_B\oplus P_D$,
\begin{equation}
 \boxed{\mathfrak c(w)=2h_w,}
 \label{sel28:eq:margin-exact}
\end{equation}
so the same quadratic statistic is exact rather than merely sufficient.
\end{theorem}

\begin{proof}
V27 proves
$\mu_w\nu_w-h_w\le\Pi(J_w)/2$.  Rearranging gives
\eqref{sel28:eq:margin-bound}.  The efficient case uses
$\Pi(J_w)=0$.  V27's two-sector equality gives
\eqref{sel28:eq:margin-exact}.
\end{proof}

\begin{corollary}[No primitive-efficiency audit is needed for the V110 word]
\label{sel28:cor:v110}
For any actually occurring V110 candidate typed odd word $\mathfrak b$ whose
historical prefix/future contexts span its Choi space, compute only its cylinder
vector $p(\mathfrak b)$ and the fixed quadratic data
$Q,b,d$.  If
\begin{equation}
 \boxed{\mathfrak c(\mathfrak b)>0,}
 \label{sel28:eq:v110-direct}
\end{equation}
then the neutral--determinant Choi coherence is nonzero.  On the efficient
branch, the simpler conditions
\[
 p^{\mathsf T}Qp=0,
 \qquad
 \mu_{\mathfrak b}>0,
 \qquad
 m_{\det}(\mathfrak b)>0
\]
are sufficient, with the fixed normalization between $m_{\det}$ and the
projected determinant Choi mass inherited from V26.  The efficiencies of the
primitive letters composing $\mathfrak b$ need not be established separately.
\end{corollary}

\section{Conditioning and exact scope of the finite historical test}
\label{sel28:conditioning}

The quadratic form is exact on the ideal historical table.  For an estimated
context vector $\widehat p=p+\delta$ with
$\|\delta\|_2\le\varepsilon$, suppose a known bound
$\|p\|_2\le R$ is retained.  Write
$B=\|b\|_2\|d\|_2$ and $q=\|Q\|_{\op}$.  Direct expansion gives
\begin{equation}
 \boxed{
 |\mathfrak c(\widehat p)-\mathfrak c(p)|
 \le
 (4B+2q)R\varepsilon+(2B+q)\varepsilon^2.}
 \label{sel28:eq:conditioning}
\end{equation}
Thus a measured lower margin exceeding the displayed error certifies positive
historical determinant coherence.  If only componentwise probability errors
$|\delta_a|\le\epsilon$ are known, one may take
$\varepsilon\le\sqrt{M}\epsilon$ for $M$ retained contexts.  These constants
belong to the chosen historical frame and are not discarded by calling that
frame tomographically complete.

\begin{proof}[Proof of \eqref{sel28:eq:conditioning}]
Expand the bilinear term
$2(b^{\mathsf T}p)(d^{\mathsf T}p)$ at $p+\delta$ and use
Cauchy--Schwarz.  Its change is at most
$4BR\varepsilon+2B\varepsilon^2$.  For the quadratic purity term,
$|(p+\delta)^{\mathsf T}Q(p+\delta)-p^{\mathsf T}Qp|
\le2qR\varepsilon+q\varepsilon^2$.  Add the bounds.
\end{proof}

\begin{remark}[The finite negative stopping rule remains asymmetric]
The inherited V110 flat-depth theorem is still a complete historical no-go when
the determinant alternation vanishes on the flat typed-word algebra: then
$\nu_w=0$ for every later historical word, so no word can satisfy
\eqref{sel28:eq:positive-margin}.  The converse is deliberately not asserted.
Once the flat algebra contains a determinant direction, a favourable
\emph{linear combination} of word operators is not automatically an actually
resolved word.  Efficiency and same-label provenance are nonlinear properties
of the physical labelled cylinder.  A positive repaired algebra therefore
cannot replace the actual word by a fitted linear combination.
\end{remark}

\section{Result and next upstream row}
\label{sel28:status}

V28 removes an unnecessary premise from V27.  Primitive-outcome efficiency is
not the historical invariant.  The exact object is the Choi rank of the
resolved candidate word itself, and that rank may decrease or increase under
composition.  The composite Kraus-span theorem identifies the intrinsic
criterion, while contextual tomography turns it into the quadratic historical
statistic
\[
 p(w)^{\mathsf T}Qp(w).
\]

Together with the two typed shadow masses, the single scalar
\[
 \boxed{
 \mathfrak c(w)
 =2(b^{\mathsf T}p(w))(d^{\mathsf T}p(w))
  -p(w)^{\mathsf T}Qp(w)}
\]
is a direct finite sufficient certificate for determinant coherence.  On a
rank-one word it is automatically positive whenever both shadows occur, and on
the isolated two-sector packet it equals twice the exact cross-block mass.
Thus the historical question can be asked directly of ordinary labelled
cylinders without a Kraus model, a primitive-efficiency census, or a separate
phase experiment.

The remaining occurrence boundary is now sharper.  One must identify an actual
V110 typed odd word in the physical grammar and retain enough old prefix/future
contexts to evaluate its fixed quadratic certificate.  A flat zero determinant
alternation terminates the historical search negatively.  A positive word
margin closes the determinant-coherence row positively, subject to the already
inherited Clifford and finite-Dirac residuals.  If the margin is nonpositive on
a mixed packet, that failure is not a no-go: the exact V26 cross-block
reconstruction remains the stronger test.

No new colour, weak, generation, determinant, or memory carrier is introduced.
Historical Standard-Model selection is still not declared closed because no
historical cylinder vector for the requisite typed odd word has been populated
here.  The next genuinely upstream task is therefore its \emph{same-history
word landing}: relate the abstract V110 typed projection to a concrete labelled
word of the reconstructed operation grammar and evaluate the finite quadratic
cylinder panel there.  Further manipulation of hypothetical Kraus factors would
be downstream of that missing occurrence row.

\paragraph{Proof and artifact status.}
The complete V27 prefix is retained before its final terminator.  The V28
checker verifies the composite Kraus-span rank identity on exact finite
examples, a rank-two primitive collapsing to a rank-one composite, the
quadratic contextual purity formula, the coherent/incoherent determinant
margins, and the stated conditioning inequality.  These checks supplement the
written finite-dimensional proofs; they are not a proof-assistant
formalization, historical data population, or independent peer review.

\addtocontents{toc}{\protect\endgroup}
% END OF SOURCE-SELECTION V28 DEVELOPMENT BLOCK.
% Preserve V1--V28 and append subsequent work before the final terminator.


\clearpage
\begin{center}
{\Large\bfseries Source-selection continuation V29}\\[.4em]
{\large Liouville word landing, same-label flatness,\\
and a finite historical provenance algorithm}
\end{center}
\addtocontents{toc}{\protect\begingroup}
\addtocontents{toc}{\protect\renewcommand{\protect\numberline}{\protect\makebox[2.1em][l]}}

\section{Go upstream from a named candidate to the complete labelled grammar}
\label{sel29:context}

V28 reduces the determinant-coherence audit of one already identified typed odd
word to a quadratic function of its ordinary historical cylinders.  One source
assumption nevertheless remains in that formulation: the candidate word has
already been named.  Searching an unrestricted list of longer words would be a
poor stopping rule, and ordinary flatness of the represented operator span does
not solve the problem because same-label coexistence of two shadows is
nonlinear on that span.

The correct linear object is one layer higher.  A completely positive labelled
word has a Liouville matrix; composition of labelled CP maps is ordinary matrix
multiplication, while the Choi matrix is obtained from the Liouville matrix by
a fixed linear reshuffling.  The neutral--determinant Choi cross block is
therefore a \emph{linear} functional of the Liouville word.  Finite-dimensional
flatness can consequently be applied at this level.  On a coherent one-Kraus
word this Liouville representation is precisely the doubled amplitude
$K_w\otimes\overline{K_w}$, explaining why the lift restores the phase
information which is nonlinear on the ordinary operator span.

The result below gives both a positive landing algorithm and a negative stopping
rule.  It uses only actual labelled words of the declared grammar.  No linear
combination is promoted to an event label, and no favourable Kraus refinement
of a coarse outcome is inserted.

\section{The Liouville word span linearizes same-label Choi coherence}
\label{sel29:liouville}

For notational clarity first fix a recurrent cut type with finite memory
$\mathcal H\cong\C^d$.  The typed case is obtained by taking the direct sum of
all source--target Liouville spaces and declaring incompatible compositions to
be zero.  Let $\Sigma$ be the finite alphabet of actually resolved CP labels.
For $a\in\Sigma$ write
\[
 \Phi_a:M_d\to M_d,
 \qquad
 \operatorname{vec}(\Phi_a(X))=S_a\operatorname{vec}(X),
\]
where $S_a\in M_{d^2}(\C)$ is its Liouville matrix.  For the word
$w=a_m\cdots a_1$ put
\[
 S_w=S_{a_m}\cdots S_{a_1},
 \qquad S_{\varnothing}=I_{d^2}.
\]
Let $\mathfrak R$ denote the fixed Liouville-to-Choi reshuffling, so
$J_w=\mathfrak R(S_w)$.

Retain the already reconstructed neutral Store--Clifford and
 determinant-incidence Choi-sector projections $P_B,P_D$.  Define the linear
cross-block map
\begin{equation}
 \boxed{
 \mathcal C(S):=P_B\,\mathfrak R(S)\,P_D.}
 \label{sel29:eq:crossmap}
\end{equation}
Thus the exact historical cross mass is
\begin{equation}
 h_w=\|\mathcal C(S_w)\|_{\rm HS}^2.
 \label{sel29:eq:hword}
\end{equation}

Vectorize the Liouville matrices and set
\begin{equation}
 z_w=\operatorname{vec}(S_w),
 \qquad
 \mathcal Z_r=\operatorname{span}_{\C}\{z_w:|w|\le r\}.
 \label{sel29:eq:Zr}
\end{equation}
Left composition by one physical label is linear:
\[
 z_{aw}=L_a z_w,
 \qquad
 L_a=I_{d^2}\otimes S_a
\]
for column-major vectorization (with the harmless transpose modification under
the opposite convention).

\begin{theorem}[Same-label Liouville flatness and exact historical stopping]
\label{sel29:thm:flat}
Suppose for some $r$ that
\begin{equation}
 \boxed{\mathcal Z_r=\mathcal Z_{r+1}.}
 \label{sel29:eq:flat}
\end{equation}
Then $\mathcal Z_s=\mathcal Z_r$ for every $s\ge r$, and:
\begin{enumerate}[label=\textup{(\roman*)},leftmargin=2.8em]
\item if $\mathcal C$ vanishes on $\mathcal Z_r$, then
      \begin{equation}
       \boxed{h_w=0\quad\text{for every later historical word }w;}
       \label{sel29:eq:no-go}
      \end{equation}
\item if $\mathcal C$ does not vanish on $\mathcal Z_r$, then at least one
      actual labelled word $w$ with $|w|\le r$ satisfies
      \begin{equation}
       \boxed{h_w>0.}
       \label{sel29:eq:witness}
      \end{equation}
\end{enumerate}
If $R=\dim\operatorname{span}\{z_w:w\text{ historical}\}$, the first flat
step occurs by $r\le R-1$ for a single nonzero seed.  In particular one never
needs an unbounded word search.
\end{theorem}

\begin{proof}
Equation~\eqref{sel29:eq:flat} implies invariance under every $L_a$: each
$L_a z_w=z_{aw}$ for $|w|\le r$ lies in $\mathcal Z_{r+1}=\mathcal Z_r$, and
linearity extends this to the whole span.  Induction gives flatness at all
later depths.  If $\mathcal C$ vanishes on the flat span it vanishes on every
later $z_w$, which proves (i).  Conversely $\mathcal Z_r$ is spanned by actual
word vectors.  A nonzero linear map on that span cannot vanish on every
spanning word, proving (ii).  Finally the nondecreasing integer sequence
$\dim\mathcal Z_r$ is bounded by $R$; if it increased strictly through depth
$R-1$, it would already have dimension $R$ there and the next step would be
flat.
\end{proof}

\begin{corollary}[Basis-witness algorithm without exponential word retention]
\label{sel29:cor:algorithm}
Start from the empty-word Liouville vector and keep a list of actual words whose
$z_w$ form a basis of the current reachable span.  In one sweep append every
admissible primitive label to each retained basis word.  Keep a new word only
when its Liouville vector increases the span, and test
$\mathcal C(S_w)$ when that word is retained.  Repeat until a complete sweep
adds no vector.

The algorithm terminates after at most $R$ retained basis words.  It returns an
actual positive-coherence word whenever one exists; if it terminates with
$\mathcal C=0$ on the retained basis, Theorem~\ref{sel29:thm:flat} proves a
complete same-label historical no-go for that grammar.
\end{corollary}

\begin{proof}
Applying all primitive left-composition maps to a basis spans the next
reachable layer, so discarding dependent word vectors loses no reachable
Liouville direction.  If a dependent actual word had nonzero cross block while
all retained basis words had zero cross block, linearity of $\mathcal C$ would
be contradicted.  The dimension bound gives termination.
\end{proof}

\section{Why ordinary operator flatness is not the right stopping rule}
\label{sel29:ordinary}

On an efficient coherent branch $\Phi_w(X)=K_wXK_w^*$ and
\begin{equation}
 \boxed{S_w=K_w\otimes\overline{K_w}.}
 \label{sel29:eq:doubled}
\end{equation}
Thus Liouville reachability is exactly the doubled operator-word reachability.
The following two-dimensional examples isolate the distinction.

\begin{proposition}[Ordinary flatness can occur before same-label coherence]
\label{sel29:prop:late}
Let
\[
 A=\begin{pmatrix}0&1\\1&1\end{pmatrix},
 \qquad x_0=\binom10,
 \qquad x_n=A^n x_0,
\]
and let the two linear shadows be the first and second coordinates.  Then the
ordinary word span is already flat at depth one:
\[
 \operatorname{span}\{x_0,x_1\}=\C^2
 =\operatorname{span}\{x_0,x_1,x_2\}.
\]
Neither $x_0$ nor $x_1$ has both shadows, but
\[
 \boxed{x_2=\binom11}
\]
does.  By contrast the doubled vectors
$y_n=x_n\otimes\overline{x_n}$ have ranks
\[
 1,\ 2,\ 3,\ 3
\]
through depths $0,1,2,3$; the doubled span does not become flat until it has
acquired the cross direction carried by $x_2$.
\end{proposition}

\begin{proposition}[A flat doubled zero is a genuine provenance no-go]
\label{sel29:prop:swap}
Replace $A$ by
\[
 R=\begin{pmatrix}0&1\\1&0\end{pmatrix}.
\]
Then the actual words alternate between $(1,0)^{\mathsf T}$ and
$(0,1)^{\mathsf T}$.  The ordinary span again contains both shadow directions,
but no historical word contains both.  The doubled reachable span is already
flat at depth one and is generated by
\[
 e_1\otimes e_1,
 \qquad
 e_2\otimes e_2,
\]
on which the cross-shadow map vanishes.  Hence
Theorem~\ref{sel29:thm:flat} returns the permanent same-label no-go.
\end{proposition}

These examples show exactly why V110's linear determinant-alternation flatness
and V28's same-label question are different.  Linear flatness detects whether a
determinant direction exists in the historical algebra.  Liouville flatness
detects whether an \emph{actual CP-labelled word} ever carries the required
coherence with the neutral Store--Clifford sector.

\section{The flatness test is itself computable from historical cylinders}
\label{sel29:cylinders}

Under the V25 contextual rank condition, choose historical testers $F_a$ and a
Hermitian dual frame $D^a$ so that for every resolved word
\begin{equation}
 J_w=\sum_a p_a(w)D^a,
 \qquad
 p_a(w)=\Tr(F_aJ_w).
 \label{sel29:eq:Jfromp}
\end{equation}
Define the fixed positive semidefinite context matrix
\begin{equation}
 \boxed{
 H_{ab}
 :=\operatorname{Re}\Tr\!\left[
   (P_BD^aP_D)^*(P_BD^bP_D)
 \right].}
 \label{sel29:eq:H}
\end{equation}

\begin{proposition}[Exact cross mass as one positive quadratic cylinder form]
\label{sel29:prop:H}
For every historical word,
\begin{equation}
 \boxed{
 h_w=p(w)^{\mathsf T}H\,p(w)\ge0.}
 \label{sel29:eq:hquadratic}
\end{equation}
The equality is frame-coordinate invariant.  Consequently the basis-witness
algorithm of Corollary~\ref{sel29:cor:algorithm} may test its retained words
using only their old contextual cylinder vectors.  Full display of $J_w$ is
optional.
\end{proposition}

\begin{proof}
Insert~\eqref{sel29:eq:Jfromp} into $P_BJ_wP_D$, take the Hilbert--Schmidt
norm squared, and expand.  Hermiticity of $J_w$ makes the result real; the
Gram construction makes $H$ positive semidefinite.  Changing contextual frame
changes $p$ and $H$ contragrediently while leaving the reconstructed cross
block fixed.
\end{proof}

The transition side is equally source native.  The table-native Hankel
realization already supplies finite linear transition maps for primitive
labels, and the support-minimal CP realization supplies the same Liouville word
semigroup up to the inherited memory-unitary gauge.  Under a memory-unitary
change all Liouville vectors are conjugated by one invertible linear map and the
sector projections are transported with it.  Reachable ranks, flat depth, and
vanishing or nonvanishing of the physical cross block are therefore invariant.
No preferred hidden-memory basis is required.

\section{Application to the V110 historical Store--Clifford question}
\label{sel29:v110}

There are now two nested finite historical stopping rules.

First, the inherited V110 determinant-alternation map is linear on the typed
operator-word span.  If the typed word hierarchy is flat and its determinant
alternation is zero, then no determinant incidence occurs at all and the
historical branch terminates negatively.

Second, suppose determinant incidence does occur somewhere in the historical
algebra.  That does not yet prove that one resolved word carries both the
Store--Clifford and determinant shadows.  Apply the Liouville construction to
the actual resolved labelled CP grammar and the cross map
\eqref{sel29:eq:crossmap}.  Then:
\begin{equation}
 \boxed{
 \begin{array}{c}
 \mathcal Z_r=\mathcal Z_{r+1},\\[0.15em]
 \mathcal C|_{\mathcal Z_r}=0
 \end{array}
 \quad\Longrightarrow\quad
 \text{no historical resolved word ever closes the same-label determinant row}.}
 \label{sel29:eq:v110-nogo}
\end{equation}
This conclusion can hold even when determinant-only and Clifford-only words
exist separately.  It is therefore stronger than the flat-zero determinant
alternative but does not contradict it.

On the positive branch the basis-witness algorithm returns a concrete actual
labelled word $\mathfrak b$.  Its cylinder vector then enters V28's quadratic
margin, or equivalently the exact V26 cross mass
\eqref{sel29:eq:hquadratic}.  The inherited V110 Clifford-shape and finite-Dirac
residuals are evaluated on that same returned word.  No fitted linear
combination of source words is used at any stage.

\begin{remark}[Finite ambient bounds versus actual historical rank]
On a single $d$-dimensional memory the ambient Liouville space has complex
dimension $d^4$, so the theorem gives the crude universal bound
$r\le d^4-1$.  For $d=14$ this is $38415$.  This is a termination bound, not a
recommendation to enumerate that many depths.  The basis algorithm stops at the
actual reachable Liouville rank, which may be far smaller, and retains at most
that many witness words.  On an efficient amplitude-resolved subbranch the
same construction is the doubled operator lift of a word algebra of dimension
$R_{\rm op}$ and therefore has rank at most $R_{\rm op}^2$.
\end{remark}

\section{Result and next upstream row}
\label{sel29:status}

V29 removes the final open-ended search implicit in V28.  The historical
question ``does some longer actual word carry both shadows?'' is not an
unbounded language problem once the finite operational realization is retained.
The correct state variable is the Liouville/Choi word, not the ordinary operator
span.  Same-label determinant coherence is a linear cross-block condition on
that finite word space, and equality of two consecutive reachable spans is a
complete stopping rule.

This gives a source-native algorithm with two possible terminal outputs:
\[
 \boxed{
 \text{concrete historical witness word}
 \quad\text{or}\quad
 \text{finite flat same-label no-go}.}
\]
The negative output remains valid even when the ordinary historical algebra
contains the determinant and Clifford directions separately.  The positive
output supplies the actual word needed by V28; its old cylinder vector then
certifies the exact cross mass without a Kraus model or a new phase experiment.

The remaining upstream issue is therefore no longer word length.  It is the
\emph{typed ancestry of the primitive CP labels themselves}: whether the
historical grammar's resolved labels and the V110 Store--Clifford/up-incidence
typed projections are represented on the same geometry-shorted physical
quotient.  If that ancestry identification is available, V29 provides a finite
complete word search.  If it is not, no amount of further multiplication of
abstract source words can establish the missing common carrier.

Historical Standard-Model selection is still not declared closed.  V29 proves
finite decidability of the same-label occurrence question conditional on the
already reconstructed typed projections and resolved physical grammar; it does
not populate those historical ancestry identifications.

\paragraph{Proof and artifact status.}
The complete V28 prefix is retained before its final terminator.  The V29
checker verifies the Liouville flatness mechanism, the late-coherence example,
the permanent swap no-go, the exact positive contextual cross form, and the
append-only prefix identity.  These are exact finite-dimensional checks of the
written proofs; they are not a proof-assistant formalization, historical data
population, or independent peer review.

\addtocontents{toc}{\protect\endgroup}
% END OF SOURCE-SELECTION V29 DEVELOPMENT BLOCK.
% Preserve V1--V29 and append subsequent work before the final terminator.


\clearpage
\begin{center}
{\Large\bfseries Source-selection continuation V30}\\[.4em]
{\large Mixed-Gram ancestry landing, source-sufficient labels,\\
and a finite common-carrier audit}
\end{center}
\addtocontents{toc}{\protect\begingroup}
\addtocontents{toc}{\protect\renewcommand{\protect\numberline}{\protect\makebox[2.1em][l]}}

\section{Go upstream from a common-carrier name to a measured ancestry relation}
\label{sel30:context}

V29 gives a finite complete search once the resolved CP grammar and the
Store--Clifford/up-incidence sector projections have been placed on one
geometry-shorted physical coefficient carrier.  The remaining phrase ``placed
on one carrier'' must not be treated as a naming convention.  Equal dimensions,
equal marginal Grams, or a formal identification of two abstract source spaces
do not establish common ancestry.

The appropriate datum is a same-history \emph{mixed Gram}.  This continuation
specializes the programme's source-ancestry short to the historical CP-label
problem.  It proves three statements.  First, one positive Schur residual is
necessary and sufficient for a proposed typed shadow bank to lie in the
actually occurring operational source range.  Second, even after this landing,
a separate positive conditional-source variance decides whether the retained
primitive label is sufficient for that ancestry or whether an unrecorded
sub-label is still required.  Third, marginal source Grams alone cannot decide
either issue: there are aligned and orthogonal joint realizations with identical
marginals.

The result does not infer a new source from a zero residual.  It says that a
previously reconstructed typed shadow is already a deterministic component of
the occurring operational source.  Conversely, a positive residual counts the
minimum additional source dimension required to realize the proposed ancestry.

\section{The mixed Gram lands a typed shadow bank on the operational source}
\label{sel30:schur}

Let
\[
 S:E\longrightarrow\mathcal K
\]
be the source synthesis reconstructed from the resolved historical CP grammar
on one geometry-shorted cut, and let
\[
 T:F\longrightarrow\mathcal K
\]
be a proposed typed shadow synthesis on the \emph{same-history joint panel}.
Here $F$ may contain the Store--Clifford line, the up-incidence determinant
line, or a larger typed bank.  Put
\begin{equation}
 G=S^*S,
 \qquad H=T^*T,
 \qquad C=S^*T.
 \label{sel30:eq:grampieces}
\end{equation}
The complete joint source Gram is
\begin{equation}
 \mathbb G=
 \begin{pmatrix}
 G&C\\ C^*&H
 \end{pmatrix}\succeq0.
 \label{sel30:eq:blockgram}
\end{equation}
Define the operational source-range projector
\begin{equation}
 P_S:=SG^\dagger S^*
 \label{sel30:eq:PS}
\end{equation}
and the typed-ancestry residual
\begin{equation}
 \boxed{
 R_{T\mid S}:=H-C^*G^\dagger C.}
 \label{sel30:eq:R}
\end{equation}

\begin{theorem}[Exact mixed-Gram ancestry landing]
\label{sel30:thm:landing}
For the same-history syntheses above,
\begin{equation}
 \boxed{
 R_{T\mid S}
 =T^*(I-P_S)T\succeq0.}
 \label{sel30:eq:Rprojector}
\end{equation}
Moreover the following are equivalent:
\begin{enumerate}[label=\textup{(\roman*)},leftmargin=2.8em]
\item $R_{T\mid S}=0$;
\item $\Ran T\subseteq\Ran S$;
\item with
      \begin{equation}
       \Theta:=G^\dagger C,
       \label{sel30:eq:Theta}
      \end{equation}
      one has the exact ancestry factorization
      \begin{equation}
       \boxed{T=S\Theta.}
       \label{sel30:eq:factor}
      \end{equation}
\end{enumerate}
If the residual is nonzero, then
\begin{equation}
 \boxed{
 \rank R_{T\mid S}
 =\dim\Ran((I-P_S)T)}
 \label{sel30:eq:rankcost}
\end{equation}
is the minimum number of additional orthogonal source coordinates required by
any enlargement of the operational source range that contains the typed bank.
\end{theorem}

\begin{proof}
The Moore--Penrose formula makes $P_S$ the orthogonal projection onto
$\Ran S$.  Hence
\[
 T^*P_ST
 =T^*SG^\dagger S^*T
 =C^*G^\dagger C,
\]
which proves \eqref{sel30:eq:Rprojector}.  Positivity is immediate.  Its
vanishing is equivalent to $(I-P_S)T=0$, which is equivalent to range
inclusion.  On that branch,
\[
 S\Theta=SG^\dagger S^*T=P_ST=T.
\]
Conversely a factorization through $S$ gives range inclusion.  Finally,
$R_{T\mid S}$ is the Gram of the orthogonal residual $(I-P_S)T$; Gram
minimality gives \eqref{sel30:eq:rankcost}.
\end{proof}

The theorem distinguishes a \emph{deterministic typed follower} from a new
source birth.  The follower can have nonzero norm and nontrivial determinant
character; zero ancestry residual means only that its source has already been
paid for by the historical operational bank.

\begin{corollary}[Whitened ancestry contraction]
\label{sel30:cor:white}
On the supports of $G$ and $H$, define
\[
 K:=G^{\dagger/2}CH^{\dagger/2}.
\]
Then $K$ is a contraction and
\begin{equation}
 \boxed{
 R_{T\mid S}
 =H^{1/2}(I-K^*K)H^{1/2}}
 \label{sel30:eq:whiteR}
\end{equation}
with all powers restricted to $\supp H$.  Thus exact ancestry is equivalent to
$K^*K=I$ on the typed source support.  Equality of the two source ranges would
add the reverse condition $KK^*=I$ on the operational support; that stronger
condition is not required merely to descend the typed shadows from the
historical source.
\end{corollary}

\begin{proof}
Write the polar source syntheses as
$S=V_GG^{1/2}$ and $T=V_HH^{1/2}$ on their supported coefficient spaces.
Then $K=V_G^*V_H$.  It is therefore a contraction, and substitution in
\eqref{sel30:eq:R} gives \eqref{sel30:eq:whiteR}.
\end{proof}

\section{For the V110 common ray the ancestry audit is one scalar}
\label{sel30:rankone}

The V110 source-minimal determinant--Clifford packet is one-dimensional at its
source entrance.  In that case take $F=\C$ and write
\[
 H=\nu>0,
 \qquad C=c\in E.
\]
The full mixed-Gram audit collapses to
\begin{equation}
 \boxed{
 \delta_{\rm anc}
 :=\nu-c^*G^\dagger c\ge0.}
 \label{sel30:eq:deltaanc}
\end{equation}

\begin{corollary}[One-scalar landing of the common determinant--Clifford ray]
\label{sel30:cor:oneray}
The proposed V110 common source ray belongs to the historical operational
source range if and only if
\begin{equation}
 \boxed{\delta_{\rm anc}=0.}
 \label{sel30:eq:delta0}
\end{equation}
On that branch its coefficient vector in the minimum operational source
coordinates is $G^\dagger c$.  If $\delta_{\rm anc}>0$, one additional source
coordinate is necessary and sufficient to carry the unexplained orthogonal
part of that ray.
\end{corollary}

This is weaker than asking the whole typed Standard-Model packet to equal the
operational source.  Only the single same-history source ray needed for the
Store--Clifford/determinant handoff is tested.

\section{Primitive labels must also be source sufficient}
\label{sel30:sufficiency}

Range landing alone does not prove that a retained primitive label carries a
well-defined typed ancestry.  The same displayed label can occur on several
histories while hiding different source profiles.  V29 may use a primitive
symbol as one transition only after that unresolved variation has been tested.

Let $(\Omega,\nu)$ be the faithful finite set of actual occurrences of one
source/target cut type, let
\[
 \lambda:\Omega\to\Sigma
\]
be the resolved primitive label including all retained cut and predictive-state
marks, and let
\[
 T_\omega:F\to\mathcal K
\]
be the landed typed profile on occurrence $\omega$.  For a label $a$ put
\[
 \nu_a=\sum_{\lambda(\omega)=a}\nu_\omega,
 \qquad
 \overline T_a
 =\nu_a^{-1}\sum_{\lambda(\omega)=a}\nu_\omega T_\omega.
\]
Define the conditional typed-source innovation
\begin{equation}
 \boxed{
 \mathbb I_{T\mid\lambda}
 :=\frac12\sum_{a\in\Sigma}\frac1{\nu_a}
 \sum_{\substack{\omega,\eta\\
                  \lambda(\omega)=\lambda(\eta)=a}}
 \nu_\omega\nu_\eta
 (T_\omega-T_\eta)^*(T_\omega-T_\eta)\succeq0.}
 \label{sel30:eq:variance}
\end{equation}

\begin{theorem}[Exact source-sufficiency test for primitive labels]
\label{sel30:thm:sufficient}
One has the positive variance decomposition
\begin{equation}
 \boxed{
 \sum_{\omega}\nu_\omega T_\omega^*T_\omega
 =\sum_a\nu_a\overline T_a^*\overline T_a
 +\mathbb I_{T\mid\lambda}.}
 \label{sel30:eq:vardec}
\end{equation}
Moreover
\begin{equation}
 \boxed{
 \mathbb I_{T\mid\lambda}=0
 \iff
 T_\omega=T_\eta
 \text{ whenever }\lambda(\omega)=\lambda(\eta).}
 \label{sel30:eq:sufficientiff}
\end{equation}
Hence a zero innovation makes the typed ancestry a deterministic function of
the retained primitive label.  If it is positive, the present label grammar is
not ancestry sufficient; the joint level-set partition of the profiles
$T_\omega$ is the unique coarsest finite refinement that is sufficient, and
\begin{equation}
 \rank\mathbb I_{T\mid\lambda}
 \label{sel30:eq:innovationrank}
\end{equation}
counts the minimum orthogonal typed-source dimension erased by the proposed
coarse record.
\end{theorem}

\begin{proof}
Expand the weighted variance separately inside each label fibre.  Every term in
\eqref{sel30:eq:variance} is positive.  Its vanishing forces
$(T_\omega-T_\eta)^*(T_\omega-T_\eta)=0$ for every pair in a fibre, hence
$T_\omega=T_\eta$.  The converse is immediate.  Refining by equality of the
finite profiles is sufficient and is coarsest by construction.  The rank
statement is the Gram dimension of the omitted orthogonal innovation.
\end{proof}

A useful consequence is that hidden Kraus refinements and hidden ancestry
refinements are different issues.  V28 does not require a physically resolved
Kraus index.  The present theorem asks only whether two occurrences carrying
the same \emph{physical primitive label} have the same typed source ancestry.

\section{Marginal source Grams cannot establish common ancestry}
\label{sel30:nogo}

\begin{theorem}[Same marginals, opposite ancestry]
\label{sel30:cth:marginals}
Knowledge of $G=S^*S$ and $H=T^*T$ without the mixed Gram $C=S^*T$ cannot
decide whether the typed source is historically carried by the operational
source.
\end{theorem}

\begin{proof}
Take $E=F=\C$ and $\mathcal K=\C^2$.  In both models let
\[
 S1=e_1,
 \qquad \|T1\|=1.
\]
For the aligned model take $T1=e_1$; then
\[
 G=H=1,
 \qquad C=1,
 \qquad R_{T\mid S}=0.
\]
For the orthogonal model take $T1=e_2$; then
\[
 G=H=1,
 \qquad C=0,
 \qquad R_{T\mid S}=1.
\]
Thus every separate source energy is identical while the ancestry conclusion
is opposite.
\end{proof}

The missing object is therefore specifically a same-history cross Gram or an
equivalent cross-Hankel panel.  A formal unitary identification chosen after
seeing the two marginals would merely replace the absent ancestry datum by a
fit.

\begin{corollary}[The full block Gram reconstructs the minimum joint carrier]
\label{sel30:cor:jointcarrier}
Once the physical block Gram \eqref{sel30:eq:blockgram} is available, its Gram
quotient has minimum joint carrier dimension
\begin{equation}
 \boxed{\rank\mathbb G,}
 \label{sel30:eq:jointdim}
\end{equation}
and every source-minimal joint realization differs by one global unitary.
Consequently $R_{T\mid S}$ and the landed typed subspaces are invariant under
the remaining joint-carrier gauge.
\end{corollary}

\section{The ancestry scalars are themselves historical frame functionals}
\label{sel30:cylinders}

Suppose the V25 contextual tomography condition holds on the relevant source
port.  Any fixed Hermitian same-history mixed-Gram entry is then a signed linear
functional of the old contextual cylinder table.  In particular, after choosing
finite operational and typed source bases, the entries of $G,H,C$ can be
reconstructed without physically executing the Moore--Penrose coefficient map
$\Theta$.  The nonlinear ancestry residual is formed only after those measured
linear coordinates have been assembled.

This distinction is important.  The signed dual-frame coefficients used to
reconstruct $C$ are statistical postprocessing coefficients; they are not
claimed to be executable intervention probabilities.  If the admitted
historical tester span does not contain the necessary mixed-Gram functional,
V24's completion-fibre argument applies: a compatible positive completion can
move the missing cross coordinate while preserving all admitted marginal data.
The ancestry result must then remain unresolved rather than being fixed by a
preferred completion.

For estimated blocks $\widehat G,\widehat H,\widehat C$, the most stable audit
uses the whitened correlation $K$ of Corollary~\ref{sel30:cor:white} on a
source support whose positive spectral floor has been independently certified.
Without such a floor, small numerical Schur residuals do not prove exact range
inclusion because the pseudoinverse is unstable at a changing source rank.

\section{Application to the V29 historical search}
\label{sel30:v29}

The finite same-label search of V29 is historically justified once two
ancestry checks pass on the same geometry-shorted cut:
\begin{enumerate}[label=\textup{(A\arabic*)},leftmargin=3em]
\item the mixed-Gram landing residual for the required Store--Clifford and
      up-incidence typed shadow bank vanishes;
\item the retained primitive labels are source sufficient for that landed bank,
      so $\mathbb I_{T\mid\lambda}=0$.
\end{enumerate}
On this branch the V110 typed sectors are deterministic shadows of the actual
resolved CP grammar rather than independently installed coordinates.  The
Liouville basis search may then be run on those historical labels, and every
returned word has a well-defined typed ancestry on the same carrier.

If the first test fails, $\rank R_{T\mid S}$ is the exact minimum new source
rank required to add the proposed typed packet.  If the second test fails, the
primitive label must be refined to the coarsest source-sufficient physical
record before a same-label search is meaningful.  Neither failure is repaired
by multiplying more abstract words.

For the source-minimal V110 common ray the first audit needs only
$\delta_{\rm anc}$ from \eqref{sel30:eq:deltaanc}.  Thus the remaining common-
carrier question has been reduced from an unspecified carrier identification
to one mixed overlap vector and one scalar Schur defect, followed by the finite
record-sufficiency variance.

\section{Result and next upstream row}
\label{sel30:status}

V30 removes the phrase ``the primitive CP labels and V110 projections are on
the same quotient'' as an untested identification.  The historical common
carrier is now a finite positive reconstruction problem:
\[
 \boxed{
 \begin{aligned}
 &\text{marginal source Grams + same-history mixed Gram}\\
 &\qquad\longrightarrow R_{T\mid S}
 \longrightarrow \mathbb I_{T\mid\lambda}
 \longrightarrow \text{ancestry-landed CP grammar}.
 \end{aligned}}
\]
Only after the two zero tests pass does V29's finite Liouville search apply.
The output is then still either one concrete historical Store--Clifford/
determinant word or a flat same-label no-go.

The next genuinely upstream question is therefore even more specific: whether
the original historical cylinder table contains the \emph{mixed cross-Hankel
coordinates} needed to reconstruct $C=S^*T$ for the V110 common source ray.
If those coordinates are in the admitted tester span, V24--V25 make the scalar
ancestry audit table native.  If not, the same-marginal countertheorem proves
that no source-only argument can identify the missing carrier alignment.

Historical Standard-Model selection is still not declared closed.  V30 proves
that common-carrier ancestry and label sufficiency are finitely decidable once
the appropriate same-history mixed coordinates occur; it does not populate
those historical cross coordinates.

\paragraph{Proof and artifact status.}
The complete V29 prefix is retained before its final terminator.  The V30
checker verifies the Schur ancestry identity, the whitened contraction, the
rank-one scalar reduction, conditional-source variance, the same-marginal
counterexample, and append-only prefix identity.  These are exact
finite-dimensional checks of the written proofs, not proof-assistant
formalization, historical data population, or independent peer review.

\addtocontents{toc}{\protect\endgroup}
% END OF SOURCE-SELECTION V30 DEVELOPMENT BLOCK.
% Preserve V1--V30 and append subsequent work before the final terminator.



\clearpage
\begin{center}
{\Large\bfseries Source-selection continuation V31}\\[.4em]
{\large Word-Gram ancestry compilation, inserted shadow moments,\\
and the exact boundary of cross-Hankel data}
\end{center}
\addtocontents{toc}{\protect\begingroup}
\addtocontents{toc}{\protect\renewcommand{\protect\numberline}{\protect\makebox[2.1em][l]}}

\section{Go upstream from a missing cross panel to the historical word Gram}
\label{sel31:context}

V30 ended with a deliberately conservative question: does the historical
cylinder table contain the mixed cross-Hankel coordinates needed to reconstruct
$C=S^*T$?  That question is too large whenever the two source families are
already represented by admitted histories on one Grand-Tensor carrier.  The
programme established much earlier that in this case a separate mixed-source
experiment is unnecessary: if
\[
 W_r:\mathcal W_{\le r}\longrightarrow\mathcal K
\]
is the physical word synthesis and $\mathbb K_r=W_r^*W_r$ its word Gram, then
source families of the form $S_i=W_rB_i$ satisfy
\[
 S_i^*S_j=B_i^*\mathbb K_rB_j.
\]
The purpose of the present continuation is to make the consequence for the
V30 ancestry audit explicit and to extend it to the typed compressions actually
used by the V110 Store--Clifford/determinant packet.

There are two logically different cases.  A \emph{word-native shadow} is a
linear operation on represented history operators built from already occurring
history-algebra elements, such as source/target compression, multiplication,
adjoint, or a finite group average.  Its mixed Grams are ordinary inserted word
moments and hence belong to the Grand-Tensor moment hierarchy.  An
\emph{external shadow} uses a projector, router, tensor-slot identification, or
coherent link not represented in that historical algebra.  Its cross block is
not determined by the old word Gram; V30's same-marginal no-go then remains in
force.

Thus V31 does not declare a new physical operation.  It proves an exact
compiler deciding when the V30 mixed ancestry coordinates are already present
in the historical positive moment object.

\section{Word-native shadows and inserted moment compilation}
\label{sel31:inserted}

Let $\mathcal A\subseteq\mathcal B(\mathcal H)$ be the finite represented
history algebra with normalized regular trace $\tau$, and write
\[
 \langle X,Y\rangle_{\rm HS}:=\tau(X^*Y).
\]
For represented words $x,y$, the Grand-Tensor Gram is
\[
 \mathbb K(y,x)=\tau(y^*x).
\]

\begin{definition}[Word-native linear shadow]
\label{sel31:def:wordnative}
A linear map $\mathscr L:\mathcal A\to\mathcal A$ is called word native if it
has a finite represented form
\begin{equation}
 \boxed{
 \mathscr L(X)=\sum_{\alpha=1}^m c_\alpha A_\alpha X B_\alpha,}
 \label{sel31:eq:wordnative}
\end{equation}
where $A_\alpha,B_\alpha\in\mathcal A$ are already reconstructed history-
algebra elements and the coefficients $c_\alpha$ are fixed by the declared
shadow, not fitted from the target conclusion.  Finite sums of compressions
$PXQ$, conjugations, and finite group averages are included.
\end{definition}

This definition is representation covariant: under one common memory unitary,
all $A_\alpha,B_\alpha,X$ and the shadow transform together.  No preferred
matrix basis is part of the criterion.

\begin{theorem}[Inserted word-moment compiler]
\label{sel31:thm:inserted}
Let $x,y$ be admitted represented words and let $\mathscr L,\mathscr M$ be
word-native maps as in \eqref{sel31:eq:wordnative}, with
\[
 \mathscr M(Z)=\sum_\beta d_\beta C_\beta ZD_\beta.
\]
Then every ordinary--shadow and shadow--shadow inner product is an ordinary
history moment:
\begin{align}
 \langle y,\mathscr L(x)\rangle_{\rm HS}
 &=\sum_\alpha c_\alpha\,
   \tau(y^*A_\alpha xB_\alpha),
 \label{sel31:eq:oneinsert}\\
 \langle\mathscr L(x),\mathscr M(z)\rangle_{\rm HS}
 &=\sum_{\alpha,\beta}\overline{c_\alpha}d_\beta\,
   \tau(B_\alpha^*x^*A_\alpha^*C_\beta zD_\beta).
 \label{sel31:eq:twoinsert}
\end{align}
Consequently these blocks are finite linear combinations of entries of the
same Grand-Tensor word-Gram hierarchy.

If $x,y,z$ have depth at most $r$ and
\[
 m_{\mathscr L}=\max_\alpha(|A_\alpha|+|B_\alpha|),
 \qquad
 m_{\mathscr M}=\max_\beta(|C_\beta|+|D_\beta|),
\]
then a symmetric word Gram through depth
\begin{equation}
 \boxed{R\ge r+m_{\mathscr L}+m_{\mathscr M}}
 \label{sel31:eq:depthbound}
\end{equation}
is sufficient for all blocks in \eqref{sel31:eq:twoinsert}.  After finite
flatness, no deeper physical data are required: products are reduced by the
already reconstructed multiplication table of $\mathcal A$.
\end{theorem}

\begin{proof}
The first identity is the definition of the Hilbert--Schmidt product.  By
cyclicity of the regular trace,
\[
 \tau(y^*A_\alpha xB_\alpha)
 =\tau\!\left((A_\alpha^*yB_\alpha^*)^*x\right),
\]
so each term is one ordinary word-Gram entry after expanding the represented
product $A_\alpha^*yB_\alpha^*$ in the word basis.  Likewise
\[
 \tau(B_\alpha^*x^*A_\alpha^*C_\beta zD_\beta)
 =\tau\!\left(
 (C_\beta^*A_\alpha xB_\alpha D_\beta^*)^*z
 \right),
\]
which proves the second assertion and the stated depth bound.  At flatness the
finite word span is already a unital $*$-algebra, so every inserted product
reduces in its multiplication table and no longer raises the required physical
word depth.
\end{proof}

\begin{corollary}[The old same-history Schur compiler is the zero-insertion case]
\label{sel31:cor:V13}
Let $W_r:\mathcal W_{\le r}\to\mathcal K$ be the physical word synthesis and
let $B_0:E_0\to\mathcal W_{\le r}$ and
$B_1:E_1\to\mathcal W_{\le r}$ be fixed word coefficient maps.  For
$S_i=W_rB_i$,
\begin{equation}
 \boxed{S_i^*S_j=B_i^*\mathbb K_r^{\rm phys}B_j.}
 \label{sel31:eq:V13}
\end{equation}
Hence V30's mixed block $C=S_0^*S_1$ is not an independent cross-Hankel row on
this branch.  It is already an entry of the physical word Gram, in agreement
with the earlier same-history Schur compiler.
\end{corollary}

\section{Compile the entire V30 ancestry audit from one positive word Gram}
\label{sel31:ancestry}

Allow now a typed family to be a word-native shadow of admitted histories.
Let
\[
 S=W_rB_S,
 \qquad
 T=\mathscr L W_rB_T,
\]
where $B_S,B_T$ are fixed coefficient maps and $\mathscr L$ is word native.
Define the inserted matrices on the formal word space by
\[
 K^{00}=W_r^*W_r,
 \qquad
 K^{0L}=W_r^*\mathscr L W_r,
 \qquad
 K^{LL}=W_r^*\mathscr L^*\mathscr L W_r.
\]
Theorem~\ref{sel31:thm:inserted} says that all three are reconstructed from
ordinary Grand-Tensor moments.

\begin{theorem}[Table-native ancestry short]
\label{sel31:thm:ancestry}
The V30 ancestry blocks are exactly
\begin{equation}
 \boxed{
 G=B_S^*K^{00}B_S,
 \quad
 C=B_S^*K^{0L}B_T,
 \quad
 H=B_T^*K^{LL}B_T.}
 \label{sel31:eq:GCH}
\end{equation}
Therefore
\begin{equation}
 \boxed{
 R_{T\mid S}
 =B_T^*K^{LL}B_T
 -B_T^*(K^{0L})^*B_S
  (B_S^*K^{00}B_S)^\dagger
  B_S^*K^{0L}B_T}
 \label{sel31:eq:compiledshort}
\end{equation}
is a deterministic function of the historical word-Gram hierarchy and the
declared word-native compiler.  No independently executed mixed ancestry panel
is required.
\end{theorem}

\begin{proof}
Equation~\eqref{sel31:eq:GCH} follows by substitution into the three inner
products.  The Schur formula is V30's exact ancestry identity applied to those
blocks.  Every matrix on its right-hand side is fixed by the word moments and
the predetermined compiler coefficients.
\end{proof}

The same observation removes a second apparent datum from V30.  If occurrences
$\omega$ carrying one retained label have represented word vectors $x_\omega$
then
\[
 \|T_\omega-T_\eta\|^2
 =\|\mathscr L(x_\omega-x_\eta)\|^2
\]
is an inserted shadow--shadow moment.  Hence the complete conditional source
innovation $\mathbb I_{T\mid\lambda}$ of V30 is also word-Gram native whenever
the occurrence-specific word vectors and $\mathscr L$ are historical.

\begin{corollary}[No new ancestry data after flatness]
\label{sel31:cor:flat}
Suppose the physical word hierarchy is flat and the operational source,
typed shadow compiler, and retained label fibres are all represented in that
flat algebra.  Then both
\[
 R_{T\mid S}
 \qquad\text{and}\qquad
 \mathbb I_{T\mid\lambda}
\]
are fixed by the existing finite Grand Tensor.  Further histories cannot alter
either ancestry conclusion without enlarging the represented history algebra
or adding a non-word-native operation.
\end{corollary}

\section{Application to the Store--Clifford and up-incidence shadows}
\label{sel31:V110}

The V110 packet uses two kinds of linear extraction from one same-history odd
word: a Store--parity Clifford compression and the raw up-type incidence
compression, followed by complete determinant alternation.  The alternation
itself is a fixed finite linear projection on the typed Hom space and therefore
requires no new physical intervention.  The only issue for ancestry is whether
the source/target projections and routers used to form those typed
compressions are represented on the same historical word algebra.

Write schematically
\[
 \mathscr L_{\rm Cl}(X)=P_-XP_+,
 \qquad
 \mathscr L_u(X)=P_u X P_{QH},
\]
with the understood already-reconstructed typed identifications of the V110
packet.  These formulas are not used to assert that the projectors occur.  They
state the compiler once they do.

\begin{theorem}[Word-native V110 ancestry compiler]
\label{sel31:thm:V110compiler}
Assume on the strict historical branch that:
\begin{enumerate}[label=\textup{(N\arabic*)},leftmargin=3em]
\item the grading, type, weak-leg and incidence source/target projectors entering
      the V110 Store--Clifford and up-incidence definitions are represented
      elements (or finite functional-calculus elements) of the same flat
      historical $*$-algebra as the resolved CP words;
\item every router or finite symmetry average used in those compressions is a
      fixed word-native operation on that algebra;
\item the coefficient maps selecting the historical word bank are fixed before
      evaluating determinant occurrence.
\end{enumerate}
Then the complete V30 common-ancestry Gram for the Store--Clifford/up-incidence
bank, its scalar common-ray defect $\delta_{\rm anc}$, and the label-sufficiency
innovation are all deterministic functions of the existing Grand-Tensor word
Gram and multiplication table.  In particular, the ``missing cross-Hankel
coordinate'' of V30 is not an additional physical datum on this branch.

After this compilation, the determinant alternation
$\mathfrak A_{\det}T_u(w)$ and the same-label cross mass of V26--V29 are
computed on the returned historical word without adding an ancestry port.
\end{theorem}

\begin{proof}
Conditions (N1)--(N2) put the typed compressions in the class of
Definition~\ref{sel31:def:wordnative}; finite functional calculus remains
inside a finite $C^*$-algebra.  Apply Theorem~\ref{sel31:thm:inserted} and then
Theorem~\ref{sel31:thm:ancestry}.  Complete determinant alternation is a fixed
linear postprocessing of the already reconstructed up-type coordinates, so it
does not modify the ancestry conclusion.
\end{proof}

This theorem matches the logical content of the V110 historical formulation:
its word span is explicitly a labelled typed operation-word span on the
geometry-shorted physical quotient.  If that phrase is realized literally by
the Grand-Tensor word algebra, V30's separate common-carrier panel is redundant.
If instead it refers only to the repaired active compiler, the theorem does not
promote it to historical occurrence.

\section{Sharp boundary: an external shadow is not recoverable from the word Gram}
\label{sel31:boundary}

The previous reduction must not be applied to a projector or coherent router
that has merely been named in an external model.  If $\mathscr L$ is not a
reconstructed operation on the historical word algebra, formulas
\eqref{sel31:eq:oneinsert}--\eqref{sel31:eq:twoinsert} contain matrix elements
not fixed by $\mathbb K$.

\begin{theorem}[Word moments cannot manufacture an external shadow]
\label{sel31:cth:external}
There are conservative extensions with identical complete historical word
algebra, multiplication table, word-Gram hierarchy, and resolved CP cylinder
probabilities but different values of a proposed external grading-changing or
typed-incidence cross block.  Therefore no function of the old word Gram alone
can determine V30's ancestry residual for a non-word-native shadow.
\end{theorem}

\begin{proof}
Use the inherited no-free-matter construction: tensor any faithful realization
of the declared history algebra with an unread two-level spectator and let all
old histories act trivially on that factor.  Every old word moment is unchanged.
One extension retains no cross-support link; another admits a future-visible
coherent spectator link and its typed compression.  The old word algebra and
its entire Gram hierarchy agree, while the proposed external cross block is
zero in the first extension and nonzero in the second.  Hence the cross is not
a function of the old Grand Tensor unless the shadow operation itself has been
landed in that tensor.
\end{proof}

The countertheorem identifies the exact remaining source row.  It is not a
numerical cross-Hankel coordinate once the shadow is word native.  It is the
\emph{historical word-nativity} of the grading/type/incidence compiler itself.

\section{Revised finite pipeline and next upstream row}
\label{sel31:status}

The V30--V31 ancestry logic is now
\[
 \boxed{
 \begin{aligned}
 &\text{historical word Gram + multiplication table}\\
 &\quad+\ \text{verified word-native typed compiler}\\
 &\qquad\Longrightarrow (G,C,H)\\
 &\qquad\Longrightarrow R_{T\mid S},\ \mathbb I_{T\mid\lambda}\\
 &\qquad\Longrightarrow \text{V29 finite same-label Liouville search}.
 \end{aligned}}
\]
On this branch there is no independent mixed ancestry measurement to populate.
The word Gram already contains it.

The next genuinely upstream question is therefore not whether a missing cross
number can be measured.  It is whether the historical Grand Tensor itself
contains the \emph{typing operations} used by the V110 compiler: the protected
grading projectors, the relevant internal type/source/target projections, and
the required router identifications on the same cut.  If they are reconstructed
as elements or finite functional-calculus operations of the historical flat
algebra, V31 closes the ancestry interface.  If any is only part of the repaired
active completion, the external-shadow countertheorem prevents its historical
promotion.

Historical Standard-Model selection is still not declared closed.  V31 removes
an unnecessary cross-Hankel premise conditional on word-native typing; it does
not prove historical occurrence of those typing operators.

\paragraph{Proof and artifact status.}
The complete V30 prefix is retained before its final terminator.  The V31
checker verifies the inserted-moment identities on an exact finite matrix
algebra, the compiled $G,C,H$ blocks and Schur residual, finite-flat reduction,
and append-only prefix identity.  These are exact finite-dimensional checks of
the written proofs, not proof-assistant formalization, historical data
population, or independent peer review.

\addtocontents{toc}{\protect\endgroup}
% END OF SOURCE-SELECTION V31 DEVELOPMENT BLOCK.
% Preserve V1--V31 and append subsequent work before the final terminator.


\clearpage
\hypersetup{pdftitle={Historical Standard-Model Source Selection: V1--V32}}
\begin{center}
{\Large\bfseries Source-selection continuation V32}\\[.4em]
{\large Complete compiler recognition, historical grading actions,\\
and the exact cost of directed typing}
\end{center}
\addtocontents{toc}{\protect\begingroup}
\addtocontents{toc}{\protect\renewcommand{\protect\numberline}{\protect\makebox[2.1em][l]}}

\section[Whole-compiler audit and inherited grading results]{Audit the compiler rather than demand every implementing projector}
\label{sel32:context}

V31 reduces the ancestry Gram to old word moments once the typing compiler is
word native.  Its last proposed task was to establish membership of each
grading projector, type projection, and router in the historical algebra.
That sufficient entrance is stronger than necessary for some outputs and
insufficient to identify others.  The present continuation separates three
questions: whether a \emph{specified whole map} has an elementary historical
formula, whether its action has actually been identified from the old marks,
and whether its \emph{directed} physical components are determined by that
action.  None is equated with execution of a new intervention.

The source audit found an important inherited result that must not be counted
again.  Grand Tensor V076, Theorem
\src{thm:SMST-grading-commutant}, already proves that two grading implementers
induce the same action precisely when their relative unitary belongs to the
commutant.  It also proves equality of their even/odd splittings.  Its
Theorem \src{thm:SMST-exact-locked-sign} reconstructs the locked sign exactly;
Theorem \src{thm:SMST-record-grading-rigidity} separates this action from an
oriented record label \cite{Grand}.  Here the new consequences are a complete
finite recognition test for V31's compiler class, an exact blockwise directed-
compression defect, and the smallest algebraic enlargement required to retain
that directed information.  The raw locked bank is evaluated without choosing
any of its unknown edge matrices.  The normalizer facts used in the proof are
finite matrix-algebra facts, not a new Skolem--Noether theorem
\cite{sel32-SN}.  The operational distinction between an algebraic tensor factor
and an identified physical one is also established background
\cite{sel32-ZLL}.

\paragraph{Two explicit qualifications to V31.}
Definition~\ref{sel31:def:wordnative} defines a \emph{complex-linear}
elementary map.  The operation $X\mapsto X^*$ is conjugate linear and is not
of that form; it is separately determined by the reconstructed involution.
Also, an externally implemented automorphism which permutes the centre need
not be elementary even when its action is completely known.  Neither point
invalidates V31's inserted-moment formulas for maps satisfying its displayed
definition.  They limit the informal list of examples and motivate the
larger, action-based inference below.

Throughout this block $\A$ is one fixed, faithfully represented, finite
historical $C^*$-algebra.  A flat word table, its involution, and its faithful
tracial Gram are treated as inputs, not reconstructed anew.  A proposed
typing is never selected by minimizing a determinant or matter residual.

\section[Exact recognition and minimum compiler length]{An exact recognition and minimum-length theorem for elementary compilers}
\label{sel32:elementary}

Write the abstract and represented decompositions as
\begin{equation}
 \A\simeq\bigoplus_{\lambda=1}^k M_{n_\lambda}(\C),\qquad
 \Hh\simeq\bigoplus_\lambda\C^{n_\lambda}\otimes\C^{m_\lambda},\qquad
 \A=\bigoplus_\lambda M_{n_\lambda}(\C)\otimes I_{m_\lambda}.
 \label{sel32:eq:wedderburn}
\end{equation}
The $m_\lambda$ are represented multiplicities, not additional matrix blocks
of $\A$.  Let $z_\lambda$ be its minimal central units, let
$\mathscr W=L^2(\A,\tau)$, and let $\mathsf P_\lambda x=z_\lambda x$.
These are orthogonal projections for any faithful tracial $\tau$.
Use $\|\cdot\|_{\mathrm{HS}(\mathscr W)}$ for the Hilbert--Schmidt norm of
maps on this word space.  Define
\[
 \operatorname{Elem}(\A)
 =\left\{x\mapsto\sum_{j=1}^r a_jxb_j:a_j,b_j\in\A,\ r<\infty\right\}.
\]
The coefficients of a proposed map must be known or supplied by identified
source marks.  The following theorem decides its \emph{representability},
not that prior identification.

\begin{theorem}[Complete elementary-compiler recognition]
\label{sel32:thm:elementary}
For a complex-linear map $\mathscr L:\A\to\A$, the following are equivalent:
\begin{enumerate}[label=\textup{(\roman*)},leftmargin=2.8em]
\item $\mathscr L\in\operatorname{Elem}(\A)$;
\item $\mathscr L(zx)=z\mathscr L(x)$ for every $z\in Z(\A)$ and $x\in\A$;
\item $\mathsf P_\lambda\mathscr L\mathsf P_\mu=0$ whenever $\lambda\ne\mu$.
\end{enumerate}
Consequently
\begin{equation}
 \boxed{\operatorname{Elem}(\A)
 =\bigoplus_\lambda\End_\C(M_{n_\lambda}),\qquad
 \dim_\C\operatorname{Elem}(\A)=\sum_\lambda n_\lambda^4.}
 \label{sel32:eq:elem-dimension}
\end{equation}
On each block choose matrix units $e_{ij}^\lambda$ and write
$\mathscr L(e_{kl}^\lambda)=\sum_{ij}c^\lambda_{ij;kl}e_{ij}^\lambda$.
An explicit historical formula is
\begin{equation}
 \boxed{\mathscr L(x)=
 \sum_{\lambda,i,j,k,l}c^\lambda_{ij;kl}
            e_{ik}^\lambda x e_{lj}^\lambda.}
 \label{sel32:eq:compiler-formula}
\end{equation}
\end{theorem}
\begin{proof}
Every central $z$ commutes with $a_j,b_j$, so (i) implies (ii).
The central units sum to the identity, which proves equivalence of (ii) and
(iii).  Under (iii) the map acts separately on each simple ideal.
The matrix-unit identity
$e_{ik}^\lambda e_{ab}^\mu e_{lj}^\lambda
=\delta_{\lambda\mu}\delta_{ka}\delta_{bl}e_{ij}^\lambda$
proves \eqref{sel32:eq:compiler-formula}.  Thus every block endomorphism is
elementary, giving (i) and the dimension.  The matrix units are coordinates
inside the already reconstructed algebra; changing them changes the formula
but not the map or the criterion.
\end{proof}

\begin{proposition}[Exact distance and a finite obstruction]
\label{sel32:prop:distance}
The orthogonal projection of a specified map onto the elementary class is
\[
 \mathfrak P_{\mathrm{el}}\mathscr L
       =\sum_\lambda\mathsf P_\lambda\mathscr L\mathsf P_\lambda.
\]
Its squared distance is
\begin{equation}
 \boxed{\delta_{\mathrm{el}}(\mathscr L)^2
 =\sum_{\lambda\ne\mu}
      \|\mathsf P_\lambda\mathscr L\mathsf P_\mu\|_{\mathrm{HS}(\mathscr W)}^2
 =\frac12\sum_\lambda
      \|[\mathscr L,\mathsf P_\lambda]\|_{\mathrm{HS}(\mathscr W)}^2.}
 \label{sel32:eq:elementary-defect}
\end{equation}
A nonzero summand returns a source-block vector whose image enters a different
central ideal.  If $\|\widehat{\mathscr L}-\mathscr L\|_{\mathrm{HS}(\mathscr W)}
\le\varepsilon$ and the central projections are exact, then
$|\delta_{\mathrm{el}}(\widehat{\mathscr L})-
\delta_{\mathrm{el}}(\mathscr L)|\le\varepsilon$.
\end{proposition}
\begin{proof}
The spaces $\mathsf P_\lambda\End(\mathscr W)\mathsf P_\mu$ are mutually
orthogonal.  The elementary maps are exactly the diagonal spaces by
Theorem~\ref{sel32:thm:elementary}.  Each off-diagonal block occurs in two
commutators with central projections, yielding the factor $1/2$.
A nonzero block has a nonzero column in any complete basis, supplying the
witness.  The final bound is contractivity of the orthogonal projection onto
the complementary map space and the reverse triangle inequality.
\end{proof}

For a nonorthogonal flat word basis with Gram $G$, the same calculation is
made in orthonormal coordinates.  A map matrix $L$ becomes
$G^{1/2}LG^{-1/2}$; no empirical reweighting of the source is performed.  If a
proposed ambient map does not even map $\A$ into $\A$, its range defect must
first be tested on a basis using the actual same-carrier representation.
Neither a zero centre defect nor a convenient extension supplies that missing
range statement.

\begin{proposition}[Minimum algebraic number of left--right terms]
\label{sel32:prop:length}
For a nonzero elementary map form, on each block, the matrix
\[
 (\mathsf R_\lambda)_{(i,k),(l,j)}=c^\lambda_{ij;kl}.
\]
The least $r$ for which $\mathscr L(x)=\sum_{j=1}^r a_jxb_j$ is
\begin{equation}
 \boxed{r_{\min}(\mathscr L)=\max_\lambda\rank\mathsf R_\lambda.}
 \label{sel32:eq:length}
\end{equation}
This is an algebraic compiler length, not a count of executable CP outcomes.
\end{proposition}
\begin{proof}
One elementary term has
$c_{ij;kl}=a_{ik}b_{lj}$ and therefore rank one after this reshuffling.
Every $r$-term formula gives block ranks at most $r$.  Conversely factor each
$\mathsf R_\lambda$ as a sum of its rank-many outer products, reshape the
factors into $a_j^\lambda,b_j^\lambda$, and pad the shorter decompositions
with zeros.  Taking the direct sums of corresponding factors produces a
formula with the maximum rank, proving both bounds.
\end{proof}

\begin{corollary}[What membership means on the previously saturated carriers]
\label{sel32:cor:saturated}
On the actual fourteen-dimensional full-matrix witness of
Theorem~\ref{sel13:thm:full-panel}, every \emph{specified} complex-linear
endomorphism of $M_{14}$ is elementary; this space has dimension $196^2=38416$.
On the abstract internal algebra $M_3\oplus M_2\oplus\C\oplus\C$, the
elementary space instead has dimension $81+16+1+1=99$, inside a
$225$-dimensional endomorphism space.
\end{corollary}
\begin{proof}
Apply \eqref{sel32:eq:elem-dimension} to the two algebra types.  The full
matrix witness itself is inherited, not reconstructed here.
\end{proof}

Thus on a verified full $M_{14}$ carrier, asking for algebra membership alone
cannot identify the physically named typing: every operator on that same
carrier already belongs to the algebra.  The source must still identify which
map is the typed one.  This corollary does not insert the distinct
$C\otimes W\otimes W\to\Lambda^2C^*$ incidence port into $\End(H_{14})$;
its typed domains and the earlier character obstruction remain unchanged.

\section[Historical grading action and the full odd generator]{A historical grading action suffices for the full odd generator}
\label{sel32:grading}

An actual relabelling or source comparison may determine a $*$-automorphism
$\alpha:\A\to\A$ without supplying an implementing unitary as an executable
word.  To make ``determines'' finite, fix a word basis $a_1,\ldots,a_D$ and
its multiplication and involution tables.  Proposed images $y_i\in\A$,
obtained from the claimed same-source action, specify a linear map.
Its unit, product, involution and order-two residuals are the finite equations
\begin{equation}
 \alpha(1)=1,\qquad
 \alpha(a_i a_j)=y_i y_j,\qquad
 \alpha(a_i^*)=y_i^*,\qquad
 \alpha(y_i)=a_i.
 \label{sel32:eq:grading-audit}
\end{equation}
These are tested using the old faithful Gram.  Zero residual means exact
identity in $\A$; a small residual does not mean exact identity.

\begin{theorem}[Action-based odd compiler]
\label{sel32:thm:odd}
If \eqref{sel32:eq:grading-audit} holds, the specified action is an involutive
$*$-automorphism.  It preserves the normalized regular trace, and
\begin{equation}
 \boxed{\mathscr O_\alpha=\frac12(\id-\alpha)}
 \label{sel32:eq:odd}
\end{equation}
is the orthogonal projection onto its $(-1)$-eigenspace in
$L^2(\A,\tau_{\rm reg})$.  For every self-adjoint historical generator $A$,
$D_\alpha=\mathscr O_\alpha(A)$ is reconstructed from that action and the
old multiplication table.  If a physical involution $J$ implements the same
action, then
\begin{equation}
 \boxed{D_\alpha=\frac12(A-JAJ)
       =P_-^JAP_+^J+P_+^JAP_-^J.}
 \label{sel32:eq:DF}
\end{equation}
No membership assertion about $J$ or its two projections is required for
this \emph{summed} odd operator.

The map $\mathscr O_\alpha$ is elementary in V31's sense exactly when
$\alpha$ fixes every central unit.  If it exchanges equal-size simple blocks
in pairs, then
\begin{equation}
 \boxed{\delta_{\mathrm{el}}(\mathscr O_\alpha)^2
       =\frac12\sum_{\{\lambda,\mu\}\ \mathrm{exchanged}} n_\lambda^2.}
 \label{sel32:eq:odd-centre-defect}
\end{equation}
It is still table-determined when this defect is nonzero, provided its
permutation action is part of the actual identified source data.
\end{theorem}
\begin{proof}
The equations extend by linearity and conjugate linearity to the complete
product and involution rules.  The last equation makes the map invertible.
A finite $*$-automorphism permutes equal-size simple ideals and acts by
unitary conjugation on each transported block, so it preserves the regular
trace.  Its order two makes it self-adjoint as a unitary on word space,
proving the projection assertion.  Expanding $P_\pm^J=(I\pm J)/2$ proves
\eqref{sel32:eq:DF}.  For a fixed central unit all cross-blocks vanish; for
an exchanged pair the two cross-blocks of $\mathscr O_\alpha$ equal
$-\alpha/2$.  Each has squared Hilbert--Schmidt norm $n_\lambda^2/4$.
Proposition~\ref{sel32:prop:distance} proves the last assertions.
\end{proof}

For any table-determined linear map with coordinate matrix $L$, the inserted
Grams are simply $GL$ and $L^*GL$.  Thus V31's ancestry formulas remain
available even for the known centre-moving action, without pretending that
such an action has a left--right formula.  Knowledge of $L$ is an essential
premise, not a consequence of the existence of an abstract linear map.

\begin{proposition}[Inner form and exact odd dimension]
\label{sel32:prop:inner}
On a simple block fixed by $\alpha$ there is a self-adjoint unitary $j_\lambda$
inside $M_{n_\lambda}$ with
$\alpha(x)=j_\lambda xj_\lambda$.  It is determined up to sign on that block.
If its two ranks are $r_\lambda^+,r_\lambda^-$, then
\[
 \dim_\C(\A_\lambda^-) =2r_\lambda^+r_\lambda^-,\qquad
 \Tr_{\End(\A_\lambda)}\alpha=(r_\lambda^+-r_\lambda^-)^2.
\]
An exchanged pair of $n\times n$ blocks contributes $n^2$ odd dimensions.
All these dimensions are fixed by the specified action, not by a chosen
physical grading orientation.
\end{proposition}
\begin{proof}
For a matrix $*$-automorphism, the images of a matrix-unit system are another
matrix-unit system.  Choosing an orthonormal vector in the range of the first
minimal projection and applying the remaining matrix units constructs a
unitary $v$ implementing it.  Since $\alpha^2=\id$, $v^2$ is scalar.
Multiplication by a scalar phase gives $j^2=I$, and a unitary with this
property is self-adjoint.  Two such implementers differ by a central scalar,
which is $\pm1$.  The odd blocks are the two rectangular corners between
the $j$ eigenspaces, of total dimension $2r^+r^-$.  Conjugation has trace
$|\Tr j|^2$.  For an exchanged pair the odd elements form the graph
$(x,-\beta x)$ of the block isomorphism $\beta$, giving dimension $n^2$.
\end{proof}

An abstract centre-exchanging action need not have a unitary implementer on
an arbitrarily repeated physical representation.  For an exchanged pair in
\eqref{sel32:eq:wedderburn}, such an implementer requires
$m_\lambda=m_\mu$; equality suffices by an orthonormal block identification,
with its inverse on the paired block.  The action-based reconstruction does
not assert a physical $J$ when this representation condition fails.

A two-dimensional counterexample to ``inner implementer required'' is already
$\A=\C\oplus\C$ with its known swap.  The odd map is
$\frac12\left(\begin{smallmatrix}1&-1\\-1&1\end{smallmatrix}\right)$ on
coefficient space.  Its elementary distance squared is $1/2$, but all its
inserted moments are determined once the actual swap is identified.

\section[Directed typing and multiplicity ambiguity]{Directed typing retains a precisely classified multiplicity ambiguity}
\label{sel32:directed}

Equality of odd parts, which is inherited from the monograph, must not be
strengthened to equality of the individual directed compressions.  The next
theorem gives the complete obstruction on a fixed simple block.  Here norms
and partial traces are the \emph{unnormalized physical} Hilbert--Schmidt ones;
they are not replaced by the regular trace of an abstractly isomorphic algebra.

\begin{theorem}[Directed-compression fibre and exact distance]
\label{sel32:thm:directed}
Let $\A=M_n\otimes I_m$ and let a physical self-adjoint unitary $J$ normalize
$\A$, inducing the specified involution $\alpha=\operatorname{Ad}_j$ with
$j=j^*=j^{-1}\in M_n$.  Then necessarily
\begin{equation}
 \boxed{J=j\otimes c,\qquad c=c^*=c^{-1}\in M_m.}
 \label{sel32:eq:fibre}
\end{equation}
Put $p_\pm=(I\pm j)/2$, $e_\pm=(I\pm c)/2$, and
$s_\pm=\rank e_\pm$, $t=(s_+-s_-)/m$.  For $X=x\otimes I_m$,
\begin{align}
 P_-^JXP_+^J
 &=p_-xp_+\otimes e_+ + p_+xp_-\otimes e_-,
 \label{sel32:eq:directed}\\
 \nabla_J(x)
 &:=\inf_{a\in\A}\|P_-^JXP_+^J-a\|_{\HS}^2\notag\\
 &=\frac{s_+s_-}{m}
   \bigl(\|p_-xp_+\|_{\HS}^2+\|p_+xp_-\|_{\HS}^2\bigr)\notag\\
 &=\frac{1-t^2}{4}\|\mathscr O_\alpha(X)\|_{\HS}^2.
 \label{sel32:eq:direction-distance}
\end{align}
The nearest element is
\begin{equation}
 \boxed{\mathbb E_\A(P_-^JXP_+^J)
 =\left(\frac{s_+}{m}p_-xp_++\frac{s_-}{m}p_+xp_-\right)\otimes I_m.}
 \label{sel32:eq:direction-expectation}
\end{equation}
If both $p_+,p_-$ are nonzero, the directed map sends all of $\A$ into $\A$
if and only if $c=I_m$ or $c=-I_m$.  The two choices exchange the directions.
If $\alpha=\id$, both directed compressions vanish on $\A$, irrespective
of $c$.
\end{theorem}
\begin{proof}
The unitary $(j\otimes I)J$ commutes with $\A$, hence equals $I\otimes c$.
It commutes with $j\otimes I$; self-adjointness and squaring give
$c^*=c=c^{-1}$.  The grading projections are
$P_+^J=p_+\otimes e_+ +p_-\otimes e_-$ and
$P_-^J=p_-\otimes e_+ +p_+\otimes e_-$, proving
\eqref{sel32:eq:directed}.
The physical trace-preserving conditional expectation onto $\A$ is
$x\otimes y\mapsto (\Tr y/m)x\otimes I_m$, yielding
\eqref{sel32:eq:direction-expectation}.  It is the orthogonal projection.
The two rectangular $p$ corners are Hilbert--Schmidt orthogonal, and
\[
 \|e_+-(s_+/m)I\|_{\HS}^2
 =\|e_--(s_-/m)I\|_{\HS}^2=\frac{s_+s_-}{m}.
\]
This proves the distance and its last expression, since
$\|\mathscr O_\alpha(X)\|_{\HS}^2
=m(\|p_-xp_+\|^2+\|p_+xp_-\|^2)$.
If both $p$ corners are available, zero distance for every $x$ is equivalent
to $s_+s_-=0$.  If one $p$ corner is absent, $j$ is scalar and both terms
vanish.
\end{proof}

\begin{theorem}[Smallest algebra retaining all directed compressions]
\label{sel32:thm:enlargement}
Under Theorem~\ref{sel32:thm:directed}, assume $p_+\ne0\ne p_-$.
The smallest unital $*$-algebra containing $\A$ and all $P_-^JXP_+^J$,
$X\in\A$, is
\begin{equation}
 \boxed{\A_{\rm dir}=M_n\otimes C^*(c)=C^*(\A,J).}
 \label{sel32:eq:min-enlargement}
\end{equation}
For a nonscalar $c$ it has abstract type $M_n\oplus M_n$, with represented
multiplicities $s_+,s_-$ and dimension $2n^2$.  Thus the exact additional
\emph{algebra} dimension is $n^2$, with no added ambient Hilbert direction.
If $\alpha=\id$, the directed outputs are zero and force no enlargement.
\end{theorem}
\begin{proof}
All displayed directed outputs lie in $M_n\otimes C^*(c)$.
Choose a matrix unit $e_{ij}$ from the $p_+$ to the $p_-$ sector.
Its directed output is $e_{ij}\otimes e_+$.  Multiplication on both sides
by matrix units in $\A$ gives every $e_{kl}\otimes e_+$ and their sum
$I_n\otimes e_+$.  Subtraction from $I$ gives $I_n\otimes e_-$.
Hence the lower bound equals the upper bound.  Also
$(j\otimes I)J=I\otimes c$, proving equality with $C^*(\A,J)$.
The scalar and trivial-action cases follow directly from the formulas.
\end{proof}

\begin{example}[Same historical action, different physical directed shadows]
\label{sel32:ex:multiplicity}
Let $\A=M_2\otimes I_2$, take the Pauli matrices $Z,X$, and compare
$J_0=Z\otimes I_2$ and $J_1=Z\otimes Z$ on the \emph{same} carrier.
Both induce $x\mapsto ZxZ$ on every old word, so the complete old word Gram
and the grading action agree.  On $X\otimes I_2$, both odd parts equal
$X\otimes I_2$.  But
\[
 P_-^{J_0}(X\otimes I_2)P_+^{J_0}=e_{10}\otimes I_2,\qquad
 P_-^{J_1}(X\otimes I_2)P_+^{J_1}
 =e_{10}\otimes e_{00}+e_{01}\otimes e_{11}.
\]
The first lies in $\A$.  The second has squared distance $1$ from $\A$,
with nearest element $X\otimes I_2/2$, and generates the enlargement
$M_2\otimes\operatorname{diag}(\C^2)$ of dimension eight.
This is a comparison of two unpopulated physical grading extensions, not an
assertion that an already measured $J$ can be changed without changing its
measurements.
\end{example}

The ambiguity theorem of the monograph identifies the residual commutant.
Theorems~\ref{sel32:thm:directed}--\ref{sel32:thm:enlargement} determine
\emph{which part of it directed typing needs}, its exact reconstruction
error, and its minimal algebraic support.  These do not turn the enlarged
algebra into a bank of performed operations.

\section[Locked record signs versus pointer reversal]{Evaluation on the literal locked bank: sign and reversal give different answers}
\label{sel32:locked}

Return now to coefficients already supplied, rather than another chosen
matter realization.  On the literal locked writer retain
\[
 A_{e,\sigma,r}=K_e\otimes\Pi_\sigma\otimes P_r,
 \qquad \sum_eK_e^*K_e=I,
 \qquad \sum_eK_e=\sqrt2 I.
\]
Here $\Pi_\pm$ and $P_c,P_\ell$ are the nonzero orthogonal rank-one record
projections.  The accepted scaling $\sqrt\vartheta$ and the identity
no-response coefficient do not change the generated algebra.  Its already
proved factorization is
\[
 \A_{\rm lock}=\A_K\otimes\mathcal D_2\otimes\mathcal D_2,
 \qquad \A_K=C^*(I,K_e),
\]
as in V9 and Proposition~\ref{sel11:prop:local-lift}.
The exact sign formula in \cite{Grand} is retained.

\begin{proposition}[The two actual record signs are native and centrally even]
\label{sel32:prop:locked-sign}
For every sign--type record,
\[
 Q_{\sigma r}=I\otimes\Pi_\sigma\otimes P_r
     =\sum_e A_{e,\sigma,r}^*A_{e,\sigma,r}\in Z(\A_{\rm lock}).
\]
Consequently both recorded sign involutions, and every binary function of
these record projections, belong to the old algebra but induce the identity
grading action on it.  For every original labelled history $w$, and indeed
every element of this generated algebra,
\begin{equation}
 \boxed{\frac12(w-Z_{\rm rec}wZ_{\rm rec})=0,\qquad
        P_-^{Z_{\rm rec}}wP_+^{Z_{\rm rec}}=0.}
 \label{sel32:eq:locked-even}
\end{equation}
Thus for these \emph{particular source-native signs}, membership is not the
missing condition: their nontrivial odd action on the literal locked bank is
absent.
\end{proposition}
\begin{proof}
Use $\sum_eK_e^*K_e=I$ and the tensor product of the two record projections.
The resulting $Q_{\sigma r}$ commutes with every coefficient and its adjoint,
hence is central in the complete generated algebra.  Centrality gives both
zero identities, including products of arbitrary length.
\end{proof}

This is not a no-go for the separately loaded full-cell action or for an
actually occurring grading-changing Store route.  Such an operation would
change the action algebra being tested and must enter under its original
source label.  It is also not a new proof of the inherited sharp-sign formula;
it is its application to the compiler membership question.

\begin{proposition}[A known pointer-swap action is not a directed historical router]
\label{sel32:prop:swap}
Fix one retained type and use the compatible local reversal of
Proposition~\ref{sel11:prop:local-lift},
$J_R=I\otimes X$ on $H_K\otimes\C^2$.  It exchanges the two central
pointer blocks of $\A_K\otimes\mathcal D_2$ and is therefore not an
elementary inner action there.  Nevertheless its specified action determines
\[
 \mathscr O_R(B\otimes\Pi_+)=\tfrac12B\otimes Z,
 \qquad B\in\A_K.
\]
The directed compression is instead
\begin{equation}
 \boxed{P_-^{J_R}(B\otimes\Pi_+)P_+^{J_R}
       =\frac14 B\otimes(Z+iY),}
 \label{sel32:eq:swap-directed}
\end{equation}
and has squared distance $\|B\|_{\HS}^2/8$ from
$\A_K\otimes\mathcal D_2$ for $B\ne0$.
The smallest algebra containing the old algebra and these directed outputs
for all $B$ is $\A_K\otimes M_2$.  With both old types retained this is
$\A_K\otimes M_2\otimes\mathcal D_2$, not a deletion of a type copy.
\end{proposition}
\begin{proof}
The swap sends $\Pi_+$ to $\Pi_-$; its odd difference is as stated.
In the $Z$ basis,
$P_-^X\Pi_+P_+^X=(Z+iY)/4$.
The diagonal expectation removes $iY/4$, whose squared norm is $1/8$.
Tensor multiplication gives the distance.  Setting $B=I$, the directed
operator and its adjoint recover $Y$; the old algebra contains $Z$, and
these generate $M_2$ on the pointer.  Conversely every displayed operator
belongs to $\A_K\otimes M_2$, proving minimality.
\end{proof}

This calculation keeps three notions separate: a sharp record sign, an
identified symmetry that swaps records, and an actually present cross-record
operator.  Nonzero reversal-odd \emph{source differences} do not by themselves
supply the last object.  The pointer-swap remains the compatible local action
already proved in V11, not a chronological oriented-edge identification.

\section[What algebra membership does not identify]{Typing that is not fixed by the represented algebra}
\label{sel32:identification}

\begin{proposition}[A bare matrix block does not canonically select a ray]
\label{sel32:prop:no-canonical-ray}
A projection assigned to an unmarked finite tracial algebra in a manner
invariant under all its inner automorphisms must be central.  On $M_n$ its
only possibilities are $0$ and $I$.  Therefore neither a proper type line,
a private plane, nor a nontrivial grading orientation is selected merely by
knowing the full matrix algebra and its trace.
\end{proposition}
\begin{proof}
Invariance means $upu^*=p$ for every unitary $u\in\A$.  Every element of
$\A$ is a complex linear combination of its unitaries, so $p$ commutes with
all of $\A$ and is central.  A simple matrix algebra has only two central
projections.  The assertion concerns the unmarked algebra; an actual
labelled effect or history can break this freedom and is not erased by this
argument.
\end{proof}

There is a positive source route that must be tried before adding such a
projection.  If an identified bin of actual outcomes has effect
$E=\sum_{a\in B,\nu}K_{a\nu}^*K_{a\nu}$, then $E\in\A$ whenever those
coefficient words belong to $\A$.  For a normalized instrument,
$0\preceq E\preceq I$, and faithfulness gives
\[
 \tau(E-E^2)=0\quad\Longleftrightarrow\quad E^2=E.
\]
In that case the bin itself supplies a historical input projection.
Its sharp probability effect does not prescribe a L\"uders state update,
and an output record projection cannot be substituted for this input effect.
The preceding locked calculation is an evaluated instance: its projections
pass this test but are central, so they supply no odd route within that bank.

\begin{proposition}[Native router existence is blockwise, not a total-rank test]
\label{sel32:prop:router}
For specified $p,q\in\A$ write their ranks in each abstract simple block as
$r_\lambda,s_\lambda$.  Then
\[
 \dim_\C(q\A p)=\sum_\lambda r_\lambda s_\lambda.
\]
A native partial isometry with $v^*v=p$ and $vv^*=q$ exists exactly when
$r_\lambda=s_\lambda$ for every $\lambda$.  Such an existence result does
not select a named router or prove that it occurs as a resolved word.
\end{proposition}
\begin{proof}
The rectangular corner in a block is $M_{s_\lambda\times r_\lambda}$.
Initial and final projections of a partial isometry have equal ranks in each
block.  Conversely orthonormal identifications of these blockwise ranges
give a direct-sum partial isometry.  Independent unitary changes on these
ranges preserve the algebra and ranks while changing the router; the old
labelled source must fix any required choice.
\end{proof}

For instance $p=(e_{11},0)$ and $q=(0,e_{11})$ in $M_2\oplus M_2$ have
equal total rank but $q\A p=0$.  On a simple factor the same total-rank
obstruction disappears, yet selecting one of the available identifications
still requires the original labelled information.

\section[V110 consequences and proof-status ledger]{Consequences for the V110 handoff and proof-status ledger}
\label{sel32:status}

The useful replacement for V31's factor-by-factor membership demand is to
identify and test the \emph{whole compiler needed by the claimed output}.
For a known $\mathscr L:\A\to\A$, the centre defect decides whether the
exact V31 elementary formulas exist and constructs them when they do.
For a separately identified historical grading action, the summed odd
operator is table-determined even if the implementer is outside $\A$ or
moves its centre.  Neither route requires making that reconstructed linear
map into a physical intervention.

For the manuscript's full odd generator, \eqref{sel32:eq:DF} removes the
individual-projector premise once the correct grading action has been
established.  This is compatible with the inherited monograph grading
ambiguity theorem.  For a directed Store--Clifford compression or a typed
incidence port the conclusion is narrower: residual multiplicity orientation
can remain, and \eqref{sel32:eq:direction-distance} measures precisely when
it lies outside the old algebra.  The minimal enlargement theorem states its
algebraic cost without declaring an extra source operation to have occurred.

The literal locked signs provide a resolved audit, not another hypothetical
positive model.  They are already native and central, and their odd compiler
vanishes on that bank.  Its separately compatible pointer reversal has a
known nontrivial action but needs a larger algebra for directed compression.
Neither observation identifies the V110 matter grading with the wrong one
merely because the labels are binary.

\paragraph{Proved in this continuation.}
Complete elementary-map recognition and exact map distance; minimum
left--right compiler length; action-based odd reconstruction including
centre exchange; the blockwise directed distance and minimal enlargement;
the literal locked-sign and pointer-swap evaluations; and the distinctions
between algebra membership, marked identification, and native router existence.
The supporting finite matrix classification is standard and is proved here
for completeness, not claimed as a new general classification theorem.

\paragraph{Inherited inputs.}
V31's flat word Gram and inserted-moment ancestry formulas, V13's full-matrix
witness, V11's local pointer reversal and source algebra, and Grand Tensor
V076's grading-automorphism/commutant and exact locked-sign theorems.
The V110 determinant-incidence and same-word coherence requirements are not
reproved or weakened.

\paragraph{Still unresolved on the historical Standard-Model branch.}
The supplied tables have not identified a nontrivial V110 matter-grading
action and the directed type/weak-leg/incidence identifications on the actual
common historical carrier.  No missing numerical matrix has been replaced
by one of the present comparison examples.  A full matrix-algebra membership
check alone cannot close that identification.  Once the identified compiler
and its actual primitive grammar are supplied, V31 compiles ancestry and
V29 runs its finite same-label search.  The flavour, grading and tensor-port
requirements for the complete matter handoff remain separate.  No historical
Standard-Model selection or interacting continuum is declared proved here.

\paragraph{Verification scope.}
The new script checks the finite identities and counterexamples exactly over
rational and Gaussian-rational matrices, including basis changes and the
physical multiplicity factors.  The V31 mathematical prefix is preserved
byte for byte before its terminal document command.  The supplied V31
checker is rerun separately; its success is not an independent verification
of all inherited mathematics.  No proof-assistant certification, physical
measurement, or priority claim follows from these calculations.

\addtocontents{toc}{\protect\endgroup}
\renewcommand{\refname}{Additional background for V32}
\begin{thebibliography}{99}
\raggedright
\bibitem[44]{sel32-SN}
J.~Szigeti and L.~van Wyk,
\emph{A constructive elementary proof of the Skolem--Noether theorem for
matrix algebras}, American Mathematical Monthly \textbf{124} (2017),
966--968. \href{https://doi.org/10.4169/amer.math.monthly.124.10.966}
{doi:10.4169/amer.math.monthly.124.10.966};
\href{https://arxiv.org/abs/1810.08368}{arXiv:1810.08368}.
Background for constructive inner implementation of matrix automorphisms.
\bibitem[45]{sel32-ZLL}
P.~Zanardi, D.~A.~Lidar and S.~Lloyd,
\emph{Quantum tensor product structures are observable induced},
Physical Review Letters \textbf{92} (2004), 060402.
\href{https://doi.org/10.1103/PhysRevLett.92.060402}
{doi:10.1103/PhysRevLett.92.060402}.
Background for the operational interpretation of represented subsystem
algebras; no historical source data are imported from this reference.
\end{thebibliography}

% END OF SOURCE-SELECTION V32 DEVELOPMENT BLOCK.
% Preserve V1--V32 and append subsequent work before the final terminator.


\clearpage
\hypersetup{pdftitle={Historical Standard-Model Source Selection: V1--V33}}
\begin{center}
{\Large\bfseries Source-selection continuation V33}\\[.4em]
{\large Reconstructing the finite family of record-preserving gradings\\
from efficient word rays and one-step moments}
\end{center}
\addtocontents{toc}{\protect\begingroup}
\addtocontents{toc}{\protect\renewcommand{\protect\numberline}{\protect\makebox[2.1em][l]}}

\section[From an unknown grading to a finite source test]{Go upstream from a supplied grading action to its recorded source constraints}
\label{sel33:context}

V32 determines what follows from an \emph{identified} historical grading
and quantifies the remaining directed multiplicity ambiguity.  Its literal
locked-sign calculation also prevents a central record sign from being
renamed an odd matter operation.  This continuation asks a prior question:
can the original marked quantum maps determine the \emph{family of possible
grading actions}, before one supplies a grading matrix?  On the class below
the answer is a finite binary-linear reconstruction, with a separate test
for its implementation on the represented Hilbert space.

\paragraph{Exact scope.}
Fix one finite historical carrier and retain its original CP outcome maps,
including their labels.  A \emph{record-preserving grading} means an
involutive $*$-automorphism implemented by a self-adjoint unitary $Z$ whose
conjugation preserves each of those maps separately.  This is a property to
be tested, not an assumption that every physical matter grading must have.
A grading which permutes outcomes, or which sends an action pulse to its
inverse pulse, belongs to a different covariance problem.  Failure of the
present test is not a theorem excluding all nonzero odd matter operators.
An explicit example below makes this boundary quantitative.

We assume that some finite bank of \emph{actually occurring resolved
histories} has nonzero Choi rank one and that its coefficient rays generate
the historical algebra being tested.  The bank may contain compound
histories: primitive outcomes need not all be efficient.  V28 already proves
that a mixed primitive can become efficient after a preceding resolved
history.  The remaining mixed primitive maps will be tested directly, not
replaced by hypothetical Kraus-resolved labels.  The same-cut identification
and physical Hilbert representation of the CP maps remain inputs.

The new results concern a finite algebra of quantum operations and its
involutive automorphisms.  They are not a classification of arbitrary
fermionic gradings, nor a derivation of charged tensor slots.  The general
relation between algebra gradings and their automorphism groups is
established mathematics \cite{sel33-gradings}.  The contributions here are
the phase-independent binary moment test, its complete mixed-outcome
filter, the reconstructible commutator pairing of its implementers, and
its application to the current source-selection boundary.  V32's
Skolem--Noether, directed-distance and locked-sign results are retained.

\section[Binary moment constraints and completeness]{A finite binary moment test for all homogeneous involutions}
\label{sel33:binary}

Let $k_1,\ldots,k_m$ be nonzero coefficient representatives of the selected
efficient word rays, and put $\A=C^*(I,k_1,\ldots,k_m)$.  No linear
independence of this bank is assumed.  Use the canonical normalized regular
trace
\[
 \tau(x)=\frac{\Tr_{\A}(L_x)}{\dim_\C\A},\qquad L_xy=xy.
\]
It is positive, faithful, tracial and invariant under every algebra
automorphism.  For $\A=\bigoplus_\lambda M_{n_\lambda}$ its value is
$\sum_\lambda n_\lambda\Tr(x_\lambda)/\sum_\lambda n_\lambda^2$.
This algebraic trace is computed from the flat multiplication table.  It is
not substituted for a physical state or for the represented Hilbert trace;
those enter the implementation test separately.

Choose a basis $b_1=I,b_2,\ldots,b_d$ of $\A$ consisting of labelled words
in the $k_a,k_a^*$, of length at most $r$.  One obtains it by retaining only
linearly independent words until left multiplication by every $k_a,k_a^*$
closes.  This uses at most $d=\dim_\C\A$ retained basis words.  Formal
adjoints enter the reconstructed algebra, not the executed event alphabet.
Let $v_i\in\mathbb F_2^m$ count the occurrences, modulo two, of each
$k_a$ \emph{or} $k_a^*$ in $b_i$.  In particular $v_1=0$.

Form the finite matrices
\begin{equation}
 G_{ij}=\tau(b_i^*b_j),\qquad
 M^a_{ij}=\tau(b_i^*k_ab_j),\qquad
 M^{a*}_{ij}=\tau(b_i^*k_a^*b_j).
 \label{sel33:eq:moments}
\end{equation}
Here $G\succ0$ and $M^{a*}=(M^a)^*$.  Keeping both displayed moment banks
is harmless; one is redundant.  Define a binary matrix $\mathsf R$ by
retaining the following rows whenever the indicated scalar is nonzero:
\begin{equation}
 \begin{array}{c|c}
 \text{nonzero scalar}&\text{row of }\mathsf R\\\hline
 G_{ij}&v_i+v_j\\
 M^a_{ij}\text{ or }M^{a*}_{ij}&v_i+v_j+e_a.
 \end{array}
 \label{sel33:eq:rows}
\end{equation}
All additions in this table are in $\mathbb F_2^m$.  Repeated or zero rows
can be discarded.  The required moments have word degree at most $2r+1$;
a flat historical multiplication table already determines them.

\begin{theorem}[Complete binary reconstruction]
\label{sel33:thm:binary}
For $s\in\mathbb F_2^m$, the assignment
\[
 k_a\longmapsto(-1)^{s_a}k_a
\]
extends to a unital complex-linear $*$-automorphism $\alpha_s$ of $\A$
if and only if
\begin{equation}
 \boxed{\mathsf R s=0.}
 \label{sel33:eq:code}
\end{equation}
The extension is unique and involutive.  Its coordinate matrix on the
labelled basis is
\[
 D_s=\diag\bigl((-1)^{s\cdot v_i}\bigr),\qquad
 \alpha_s\!\left(\sum_i x_ib_i\right)=\sum_i(D_sx)_ib_i.
\]
The complete family $\mathcal C=\Ker_{\mathbb F_2}\mathsf R$ is an
elementary abelian group of distinct actions:
$\alpha_s\alpha_t=\alpha_{s+t}$ and
$|\mathcal C|=2^{m-\rank_{\mathbb F_2}\mathsf R}$.
\end{theorem}
\begin{proof}
If an extension exists, multiplicativity gives
$\alpha_s(b_i)=(-1)^{s\cdot v_i}b_i$.  Invariance of the regular trace
follows because $L_{\alpha_s(x)}=\alpha_sL_x\alpha_s^{-1}$.
Applying trace and involution preservation to \eqref{sel33:eq:moments}
gives
\begin{equation}
 D_sGD_s=G,\qquad
 D_sM^aD_s=(-1)^{s_a}M^a,
 \label{sel33:eq:moment-identities}
\end{equation}
and the same identity for $M^{a*}$.  Entrywise these are exactly
\eqref{sel33:eq:rows}--\eqref{sel33:eq:code}.

Conversely assume those binary equations.  Define $\alpha_s$ by $D_s$ on
the basis.  It is a unital linear involution, and $D_sGD_s=G$ makes it
unitary on $L^2(\A,\tau)$.  The coordinate matrix of left multiplication
by $k_a$ is $G^{-1}M^a$.  Since $D_s$ commutes with $G$, the second
identity in \eqref{sel33:eq:moment-identities} gives
\[
 \alpha_s(k_ax)=(-1)^{s_a}k_a\alpha_s(x),\qquad
 \alpha_s(k_a^*x)=(-1)^{s_a}k_a^*\alpha_s(x).
\]
Iteration proves this identity for every word, and linear extension in the
left factor proves $\alpha_s(xy)=\alpha_s(x)\alpha_s(y)$ on all of
$\A$.  The generator and adjoint formulas then give
$\alpha_s(x^*)=\alpha_s(x)^*$.  Uniqueness follows from generation.
Composition adds the signs.  Because each $k_a$ is nonzero, distinct sign
assignments induce distinct actions.  This proves every assertion.
\end{proof}

\begin{corollary}[No coherent cross-label phase is required]
\label{sel33:cor:phase}
The code and its reconstructed actions are independent of the chosen Kraus
phases and of the labelled word basis.  They are therefore determined by
the selected same-cut rank-one Choi maps and their represented operator
algebra, without a coherent comparison between different outcome labels.
\end{corollary}
\begin{proof}
Replacing $k_a$ by $e^{i\theta_a}k_a$ multiplies each word by a scalar
phase (with the opposite phase for an adjoint).  Every entry in
\eqref{sel33:eq:moments} is multiplied by a unit-modulus scalar.  Its zero
or nonzero status is unchanged, as are all parity counts.  The assignment
$\alpha_s(k_a)=(-1)^{s_a}k_a$ itself is phase independent.  Basis
independence follows from the necessary-and-sufficient theorem, not an
assumption that two bases have the same zero pattern.
\end{proof}

\begin{proposition}[Pair overlaps alone do not suffice]
\label{sel33:prop:cubic}
On $M_2$ take the normalized three-outcome comparison instrument with
coefficients $X/3,2Y/3,2Z/3$, where $X,Y,Z$ are the Pauli matrices.
Its basis $I,X,Y,Z$ has diagonal trace Gram, allowing every choice of three
signs if only pair overlaps are checked.  The one-step moments impose
\begin{equation}
 \boxed{s_X+s_Y+s_Z=0.}
 \label{sel33:eq:pauli-constraint}
\end{equation}
There are exactly four actions, not eight.  In particular the assignment
which negates all three axes fails to be a complex-linear automorphism.
\end{proposition}
\begin{proof}
The relation $XY=iZ$ requires $(-1)^{s_X+s_Y}=(-1)^{s_Z}$.
Equivalently $\tau(Z^*XY)=i\ne0$ supplies the row $(1,1,1)$.
The four allowed assignments are conjugation by $I,X,Y,Z$, proving
sufficiency and completeness.  The normalization is
$(1+4+4)I/9=I$.  This is a declared comparison instrument, not a change to
any historical source.
\end{proof}

\section[Physical implementation and mixed outcomes]{Which reconstructed actions are physical grading symmetries?}
\label{sel33:physical}

Write the known representation as
\[
 \A=\bigoplus_\lambda M_{n_\lambda}\otimes I_{m_\lambda}
 \quad\text{on}\quad
 \Hh=\bigoplus_\lambda\C^{n_\lambda}\otimes\C^{m_\lambda}.
\]
An $\alpha_s$ permutes the minimal central units by an involution
$\pi_s$.  Isomorphic ideals have equal $n_\lambda$; their represented
multiplicities need not agree.

\begin{theorem}[Represented implementation and complete mixed-map filter]
\label{sel33:thm:physical}
An action $\alpha_s$ from Theorem~\ref{sel33:thm:binary} is implemented
on this \emph{same} Hilbert space by a self-adjoint unitary if and only if
\begin{equation}
 \boxed{m_\lambda=m_{\pi_s(\lambda)}\quad\text{for every }\lambda.}
 \label{sel33:eq:multiplicity-test}
\end{equation}
On a full matrix carrier $\A=M_n$, every action passes, and its
implementing involution $Z_s$ is unique up to a common sign.  The set of
record-preserving grading actions of an efficient generating instrument is
then exactly $\mathcal C$.

More generally, suppose the selected efficient histories generate $M_n$
but some primitive outcomes $\Phi_b$ are mixed.  Reconstruct $Z_s$ for
$s\in\mathcal C$ and retain exactly those satisfying
\begin{equation}
 \boxed{
 (Z_s\otimes\overline Z_s)J_b
 (Z_s\otimes\overline Z_s)^*=J_b
 \quad\text{for every original primitive label }b.}
 \label{sel33:eq:mixed-filter}
\end{equation}
The survivors are precisely all record-preserving grading actions of the
\emph{complete} instrument.  They form a binary subspace of $\mathcal C$.
No refinement of a mixed label is needed or asserted.
\end{theorem}
\begin{proof}
A spatial unitary must map the represented support of $z_\lambda$ onto
that of $\alpha_s(z_\lambda)$; dimensions give necessity of
\eqref{sel33:eq:multiplicity-test}.  On a fixed simple block, an inner
unitary implementer $u$ exists.  Since the action squares to the identity,
$u^2$ is scalar, so a phase makes $u$ a self-adjoint involution.  Tensor it
with $I_{m_\lambda}$.  On a two-cycle of simple blocks, choose a unitary
intertwiner $v$ between their matrix representations, tensored with an
identification of their equal multiplicity spaces.  The off-diagonal
operator with blocks $v,v^*$ is a self-adjoint involution implementing the
exchange.  Their direct sum proves sufficiency.  This uses the finite
matrix implementation described in V32; it does not select the physical
multiplicity orientation when that orientation is unresolved.

On $M_n$, two implementing unitaries differ by a scalar; self-adjoint
involutions differ by a sign.  Conversely, if $Z$ preserves a rank-one
outcome with nonzero coefficient $k_a$, Choi equality gives
$Zk_aZ=\zeta_a k_a$, $|\zeta_a|=1$.  Applying conjugation twice forces
$\zeta_a^2=1$.  Hence $Z$ induces one of the sign actions of the code.
The converse follows by squaring the sign inside each CP outcome.

A symmetry of all primitive maps preserves every actually composed
resolved history.  It therefore belongs to the code of any efficient
historical generating bank.  For each candidate on $M_n$,
\eqref{sel33:eq:mixed-filter} is exactly covariance of the remaining
maps and is sufficient.  These covariance symmetries form a group, and the
candidate actions commute and have exponent two; their intersection is
therefore a binary subspace.  The sign of $Z_s$ cancels in the Choi action.
\end{proof}

The full-matrix assumption in the last paragraph is substantive.  On
$M_n\otimes I_m$ an action can have different implementers
$Z_s=j_s\otimes c$ on the unresolved multiplicity.  V32 computes their
directed ambiguity.  If a remaining outcome acts outside the generated
subalgebra, testing just one chosen $c$ need not classify the other
implementations.  The present full-matrix filter avoids that hidden choice.

\begin{proposition}[Constructive implementer from a positive linear system]
\label{sel33:prop:laplacian}
On $M_n$, for $s\in\mathcal C$ the simultaneous solution space
\begin{equation}
 \mathcal N_s=\{T:Tk_a=(-1)^{s_a}k_aT,
       \ Tk_a^*=(-1)^{s_a}k_a^*T\ \forall a\}
 \label{sel33:eq:intertwiner}
\end{equation}
is one-dimensional.  Any nonzero solution is a scalar multiple of a
unitary; a scalar phase makes its normalized version a self-adjoint
involution.  It is computable as the kernel of the positive quadratic form
\[
 q_s(T)=\sum_a\bigl(\|Tk_a-(-1)^{s_a}k_aT\|_{\HS}^2
             +\|Tk_a^*-(-1)^{s_a}k_a^*T\|_{\HS}^2\bigr).
\]
\end{proposition}
\begin{proof}
Existence follows from Theorem~\ref{sel33:thm:physical}.  For any solution,
$T^*T$ commutes with all $k_a,k_a^*$ and is scalar.  A nonzero solution is
therefore invertible and a scalar multiple of a unitary.  Comparing it with
one implementer shows that their relative operator lies in $M_n'=\C I$.
Also $T^2$ commutes with the bank; its scalar value is nonzero.  Phase
normalization of the unitary produces square $I$, hence self-adjointness.
The positive form has exactly the displayed simultaneous kernel.
\end{proof}

\begin{example}[An intrinsic code need not have a spatial lift]
\label{sel33:ex:multiplicity}
The bank $k_0=3I/5,k_1=4z/5$ on $\C\oplus\C$, $z=(1,-1)$,
has the two actions $(s_0,s_1)=(0,0),(0,1)$.  The latter exchanges the
simple blocks.  In the representation with multiplicities $(1,2)$ it
cannot be implemented by any unitary: the $+1$ and $-1$ eigenspaces of
$z$ have different dimensions.  With multiplicities $(1,1)$ the swap
matrix implements it.  The intrinsic algebra and code agree in both
cases; the physical trace and spatial feasibility do not.
\end{example}

\section[Commutator pairing and classical-record limits]{The compatible actions need not have commuting binary implementers}
\label{sel33:pairing}

On $M_n$ reconstruct $Z_s$ as above.  The ambiguity $Z_s\mapsto-Z_s$
does not affect the following invariant.

\begin{theorem}[Reconstructible commutator pairing]
\label{sel33:thm:pairing}
There is a well-defined alternating bilinear form
$\beta:\mathcal C\times\mathcal C\to\mathbb F_2$ such that
\begin{equation}
 \boxed{Z_sZ_t=(-1)^{\beta(s,t)}Z_tZ_s.}
 \label{sel33:eq:pairing}
\end{equation}
It is determined by the recorded bank and is unchanged by simultaneous
memory-unitary conjugation.  If $\dim\mathcal C=k$ and
$\rank\beta=2q$, then
\begin{equation}
 \boxed{2^q\text{ divides }n,\qquad
 \max\{\dim V:V\subseteq\mathcal C,\ \beta|_V=0\}=k-q.}
 \label{sel33:eq:pairing-bound}
\end{equation}
Thus at most $k-q$ independent candidate actions admit mutually commuting
involutive implementers.  Algebraic existence of these implementers does
not establish physical sharp readout of them.
\end{theorem}
\begin{proof}
The actions $\alpha_s,\alpha_t$ commute, so their implementer commutator
is a scalar.  Write $Z_sZ_t=\zeta Z_tZ_s$.  Conjugation by the involution
$Z_s$ twice gives $\zeta^2=1$.  This defines $\beta$, independent of
implementer phases.  An implementer of $\alpha_{s+t}$ is a scalar
multiple of $Z_sZ_t$, so commuting that product past $Z_u$ multiplies the
two signs.  This proves bilinearity, and $\beta(s,s)=0$ proves
alternation.  The scalar can be read from the trace of
$Z_sZ_tZ_sZ_t$.

An alternating form over $\mathbb F_2$ has a symplectic basis of $q$
anticommuting pairs plus a radical of dimension $k-2q$.  The $2q$
corresponding Hermitian involutions generate a full $M_{2^q}$ factor:
for one pair, diagonalize the first; the second exchanges its two
eigenspaces unitarily, giving the Pauli factor.  Subsequent pairs commute
with that factor and repeat this construction on its multiplicity.
Hence $n$ is a multiple of $2^q$.  In a $2q$-dimensional symplectic
space an isotropic subspace has dimension at most $q$; adjoining the
radical gives $k-q$, and a symplectic basis attains it.  Commuting
involutions give commuting sharp projectors mathematically.  Their
admission as measurements is a separate source condition.
\end{proof}

This distinction matters for the current programme.  Several commuting
actions on the old algebra need not be several jointly sharp classical
sign records on its Hilbert representation.  Anticommuting implementers
can induce commuting automorphisms.  Treating every admissible grading as
an extra simultaneously readable classical mark would erase the pairing
\eqref{sel33:eq:pairing}.

\section[Exact qubit and Clifford evaluations]{Positive reconstruction and exact ambiguity examples}
\label{sel33:examples}

\begin{theorem}[A labelled one-way qubit operation determines a nontrivial grading]
\label{sel33:thm:damping}
Consider the normalized, explicitly declared comparison bank
\begin{equation}
 k_0=\begin{pmatrix}1&0\\0&3/5\end{pmatrix},\qquad
 k_1=\frac45|0\rangle\langle1|.
 \label{sel33:eq:damping}
\end{equation}
Its generated algebra is $M_2$ and its code is
$\{(0,0),(0,1)\}$.  Hence it has a unique nonidentity
record-preserving grading action.  The input and output supports of the
\emph{same recorded map} $\Phi_1$ reconstruct
\begin{equation}
 \boxed{
 p_{\rm in}=\supp\Phi_1^*(I)=|1\rangle\langle1|,\quad
 p_{\rm out}=\supp\Phi_1(I)=|0\rangle\langle0|,\quad
 Z=p_{\rm in}-p_{\rm out}.}
 \label{sel33:eq:oriented-support}
\end{equation}
They sum to $I$ and are orthogonal.  Declaring the directed input to have
the positive sign, this fixes the implementing involution and gives
$k_1=p_{\rm out}k_1p_{\rm in}$ with squared norm $16/25$.
Neither $Z$ nor a separate grading action is supplied to the construction.
\end{theorem}
\begin{proof}
The two normalization effects sum to $I$.  The nonzero off-diagonal
matrix unit and its adjoint, together with $I$, generate $M_2$.
Because $\tau(k_0)=4/5\ne0$, its sign must be positive.
Conjugation by $\diag(-1,1)$ fixes $k_0$ and negates $k_1$, so both
listed assignments occur and no others can.  The two support formulas
follow directly from $k_1^*k_1=(16/25)|1\rangle\langle1|$ and
$k_1k_1^*=(16/25)|0\rangle\langle0|$.  They and the CP map are
independent of its unresolved common Kraus phase.  The sign convention
orients the already reconstructed action; it is not an extra matter
representation or determinant relation.
\end{proof}

This proves a positive source-level grading reconstruction, not that the
historical Store equals \eqref{sel33:eq:damping}.  On a larger carrier a
one-way rank-one jump may leave an untested spectator subspace.  The full
support and $M_2$ generation used above cannot be omitted.

\begin{proposition}[An actual mixed axis reduces the candidate family]
\label{sel33:prop:axis}
The comparison bank $3X/5,4Z/5$ has code $\mathbb F_2^2$, with
nonidentity actions implemented by $X,Y,Z$ and a pairing of rank two.
The different normalized bank
\[
 X/3,\qquad 2Z/3,\qquad 2(3X+4Z)/15
\]
has code $\{(0,0,0),(1,1,1)\}$ and unique nonidentity action
$\operatorname{Ad}_Y$.
\end{proposition}
\begin{proof}
For the first bank all four sign choices are Pauli conjugations.  For the
second, the nonzero trace overlaps of its third axis with each of the first
two force their three signs to agree.  Conjugation by $Y$ reverses each,
so this constraint is sufficient.  Each unscaled axis is a Hermitian
unitary, and the squared prefactors sum to one.  These are distinct
comparison instruments; the third axis is not silently executed as a
linear combination of the first instrument's old outcomes.
\end{proof}

\begin{proposition}[The represented four-axis Clifford bank has sixteen actions]
\label{sel33:prop:clifford}
Take the represented Hermitian matrices
\[
 \sigma_0=X\otimes I,\quad \sigma_1=Z\otimes Z,\quad
 \sigma_2=Z\otimes Y,\quad \sigma_3=Y\otimes I.
\]
Their normalized comparison outcomes have coefficients $\sigma_a/2$.
The word algebra is $M_4$; the sixteen ordered square-free Clifford
monomials form a basis through length four.  The binary code is all of
$\mathbb F_2^4$.  In the basis of single-axis sign assignments,
\begin{equation}
 \boxed{\beta_{ab}=1\ (a\ne b),\qquad\beta_{aa}=0,\qquad
       \rank_{\mathbb F_2}\beta=4.}
 \label{sel33:eq:clifford-pairing}
\end{equation}
There are fifteen nonidentity actions, but at most two independent ones
have commuting implementers.  For a sign assignment reversing $j$ of the
four axes, the manuscript's Clifford occurrence margin is $j/4$.
\end{proposition}
\begin{proof}
The axes anticommute and square to $I$; direct matrix multiplication gives
the sixteen independent monomials and the normalization.
For any target sign vector $s$, let
$c=\sum_a s_a\pmod2$ and choose the subset with indicator $s_a+c$.
If its cardinality is $t$, the product of its axes, multiplied by
$i^{t(t-1)/2}$, is a Hermitian involution with exactly the required
conjugation signs.  Thus all sixteen assignments occur.  For a single-axis
sign assignment the implementer is, up to phase, the product of the
other three axes.  These four implementers pairwise anticommute, giving
\eqref{sel33:eq:clifford-pairing}; its binary matrix is nonsingular.
Each reversed unitary axis exchanges two equal-dimensional grading
fibres, so its directed corner has squared norm $4/2=2$.
In $p_{\rm Cl}=(2\cdot4)^{-1}\sum_a\|P_-\sigma_aP_+\|_{\HS}^2$
this contributes $1/4$.  Unreversed axes contribute zero.
\end{proof}

The Clifford matrices are inherited represented data, not newly derived
spacetime axes.  Making all four into normalized executable outcomes is
explicitly a comparison instrument.  The calculation does not choose one
of its fifteen nonidentity grading actions as the historical matter
grading, and it does not add four commuting orientation records.

\section[No-go scope and inherited locked source]{What a zero code does and does not exclude}
\label{sel33:boundary}

\begin{proposition}[No homogeneous grading is not no odd matter]
\label{sel33:prop:nonhomogeneous}
Let
\[
 R_X=(3I+4iX)/5,\quad R_Z=(3I+4iZ)/5,\qquad
 k_X=3R_X/5,\quad k_Z=4R_Z/5.
\]
This normalized two-outcome comparison bank generates $M_2$ and has
trivial binary code.  Nevertheless the valid grading $Z$ gives
\begin{equation}
 \boxed{\frac{k_X-Zk_XZ}{2}=\frac{12i}{25}X\ne0,\qquad
 \left\|\frac{k_X-Zk_XZ}{2}\right\|_{\HS}^2=\frac{288}{625}.}
 \label{sel33:eq:nonhomogeneous}
\end{equation}
It is not a record-preserving symmetry of this instrument.
\end{proposition}
\begin{proof}
$R_X,R_Z$ are unitary; the squares of $3/5,4/5$ sum to one.
Their nonzero traces force both signs to be positive, so the code is
trivial.  They generate $X,Z$ algebraically and hence $M_2$.
Conjugation by $Z$ sends $R_X$ to $R_X^*$, which is not a scalar
multiple of $R_X$ because both its identity and $X$ coefficients are
nonzero.  The rank-one CP outcome is therefore not preserved.  Direct
subtraction proves \eqref{sel33:eq:nonhomogeneous}.
\end{proof}

The literal locked source of V32 remains a separate inherited zero case.
Its retained coefficient rays are independent on each record block, and
\[
 \sum_e k_e=\sqrt2 I.
\]
Their full identity participation already forces every fixed-label
covariance to fix each coefficient, by V20's phase-rigidity theorem.
Thus no nontrivial action on that generated algebra is obtained merely by
assigning different homogeneous parities to the same locked labels.
The present code must return that conclusion wherever its full rays are
specified; no unknown edge coefficient is assigned a favourable value.
V32's central record signs implement the trivial action, and its
pointer swap permutes labels rather than belonging to the fixed-label
class.  These are retained results, not counted as new source closures.

A similar limitation concerns mixed labels without a generating efficient
bank.  A single completely depolarizing outcome on $M_2$ has full Choi
rank and admits every unitary conjugation, including a continuum of
involutive actions.  There is no finite independent sign per Kraus ray of
that \emph{one} recorded outcome.  Applying the binary theorem to an
arbitrarily chosen Kraus refinement would change the retained record
problem.  The mixed filter in Theorem~\ref{sel33:thm:physical} avoids
this mistake precisely when an actual efficient generating bank is present.

\section[Finite calibration and error certificates]{Identification cost and one-sided stability of the finite test}
\label{sel33:calibration}

\begin{proposition}[Exact sign-calibration cost]
\label{sel33:prop:calibration}
Let $\mathcal C_0$ be the surviving binary space on a full matrix carrier,
of dimension $k_0$.  An actually established homogeneous sign on a nonzero old
word with parity vector $v$ imposes the equation $s\cdot v=b$.
A collection of such calibrations identifies a unique action exactly when
its row map is injective on $\mathcal C_0$ and its values are consistent.
At least $k_0$ independent binary scalar calibrations are required in
this restricted query model; $k_0$ suitably chosen primitive-coordinate
signs suffice algebraically.  With only $k_0=1$, independent knowledge
that the action is nonidentity selects its unique nonzero element.
\end{proposition}
\begin{proof}
The compatible assignments form an affine fibre of the row map on
$\mathcal C_0$.  It is a singleton exactly when its kernel is zero.
A map into $\mathbb F_2^\ell$ cannot be injective on a $k_0$-dimensional
space when $\ell<k_0$.  Coordinate functionals span the dual of any
subspace of $\mathbb F_2^m$, so one may choose $k_0$ independent ones.
This is an information count for the displayed parity queries.  It does
not claim physical access to those signs from outcome names alone, or a
lower bound for arbitrary quantum experiments.
\end{proof}

The action reconstructed in the qubit one-way example has an additional
source-derived orientation through \eqref{sel33:eq:oriented-support}.
In general identifying $\alpha_s$ does not by itself identify an
oriented writer, and on a representation with residual multiplicity it
does not remove V32's directed-typing ambiguity.

\begin{proposition}[Certified exclusions without assuming exact zeros]
\label{sel33:prop:noise}
Suppose each scalar in \eqref{sel33:eq:moments} has an estimate $\hat z$
with certified error $|z-\hat z|\le\varepsilon_z$.
Rows for which $|\hat z|>\varepsilon_z$ are certainly valid constraints.
Their binary kernel contains every exact homogeneous grading of the true
bank.  If these certified rows have rank $m$, no nontrivial action in this
class exists.  More generally, for a proposed $s$ and an exactly
identified labelled basis, define
\[
 \Delta_s^2=\|D_sGD_s-G\|_F^2+
       \sum_a\|D_sM^aD_s-(-1)^{s_a}M^a\|_F^2.
\]
Then $\Delta_s=0$ exactly for the code assignments, and certified errors
$\|\widehat G-G\|_F\le\epsilon_G$,
$\|\widehat M^a-M^a\|_F\le\epsilon_a$ imply
\begin{equation}
 \boxed{|\widehat\Delta_s-\Delta_s|
 \le 2\sqrt{\epsilon_G^2+\sum_a\epsilon_a^2}.}
 \label{sel33:eq:noise}
\end{equation}
\end{proposition}
\begin{proof}
The reverse triangle inequality proves the certified nonzero test.  Each
retained row is therefore satisfied by every true action, giving the
kernel inclusion and full-rank exclusion.  The residual expression is
exactly \eqref{sel33:eq:moment-identities}; the adjoint moments are
redundant.  Each difference map on a matrix has Frobenius operator norm
at most two, and the reverse triangle inequality in their Hilbert direct
sum gives \eqref{sel33:eq:noise}.
\end{proof}

Small estimated residuals do not prove exact covariance.  Nor do the bounds
certify a proposed word basis, memory dimension, or efficient Choi rank
unless those have been separately established.  Near a vanishing overlap
the binary code can enlarge discontinuously; a convention of rounding
small overlaps to zero would manufacture a symmetry.

\section[Consequences and proof-status ledger]{Consequences for the historical source programme}
\label{sel33:status}

The upstream task is now finite on a stated source class:
\[
 \boxed{
 \begin{gathered}
 \text{actual efficient generating histories on one identified carrier}\\
 \longrightarrow\text{binary word-moment code and spatial test}\\
 \longrightarrow\text{mixed-primitive covariance filter}\\
 \longrightarrow\text{identified grading action or remaining finite family}.
 \end{gathered}}
\]
No grading projector is inserted in the entrance of this reconstruction.
When a unique nonidentity action is selected by actual source information,
V32 determines its full odd compiler.  Its directed implementation is then
tested on the same carrier, V31 compiles the ancestry, and V29 searches
for an actual same-label determinant--Clifford word.  These later stages
are not automatic consequences of a nontrivial binary code.

\paragraph{Proved in this continuation.}
The complete finite binary moment classification; independence from Kraus
phases and labelled basis; a constructive spatial implementation test and
full-matrix mixed-outcome filter; the canonical commutator pairing and its
commuting-record bound; the positive qubit directed reconstruction;
the cubic Pauli obstruction, sixteen-action Clifford evaluation,
representation-multiplicity obstruction and nonhomogeneous counterexample;
and finite calibration and one-sided error certificates.  The methods of
finite algebra grading, trace invariance and inner implementation are
standard; the source-level tests and their stated applications are the
relative development.

\paragraph{Inherited inputs.}
V32's action/implementer and directed-typing distinction, V31's flat
historical moment calculus, V28's efficiency-of-composite observation,
V25's same-cut contextual tomography, V29's same-label search, and V20's
identity-participation obstruction.  The displayed Clifford matrices are
already part of the companion realization.  None of those results is
reproved as a new historical Standard-Model selection.

\paragraph{Unresolved physical rows.}
The actual V110 matter grading has not been shown to be a fixed-label
symmetry of an efficiently generated same-cut instrument.  The supplied
historical matrices have not been populated to select a nontrivial code
element or the required directed matter compiler.  A grading that permutes
labels or has both even and odd components in one coefficient is not ruled
out by a trivial code.  The qubit and Clifford evaluations are explicit
comparison calculations, not historical substitutions.  Charged tensor
slots, determinant incidence, finite-Dirac completeness and a physical
Einstein--Standard-Model continuum remain separate.

\paragraph{Verification scope.}
The companion checker builds exact word bases, moments and binary kernels;
checks every surviving sign assignment against multiplication; reconstructs
implementers and their pairing; and verifies the stated examples with
rational or Gaussian-rational arithmetic.  The V32 prefix is retained
unchanged before its terminal document command.  Regression tests are
computational checks, not independent verification of the full preceding
programme or proof-assistant certification.

\addtocontents{toc}{\protect\endgroup}
\renewcommand{\refname}{Additional background for V33}
\begin{thebibliography}{99}
\raggedright
\bibitem[46]{sel33-gradings}
A.~Gordienko and O.~Schnabel,
\emph{Categories and weak equivalences of graded algebras},
\href{https://arxiv.org/abs/1805.04073}{arXiv:1805.04073} (2018).
Background on gradings and universal grading groups.  No physical source
matrix or recorded grading is inferred from this reference.
\end{thebibliography}

% END OF SOURCE-SELECTION V33 DEVELOPMENT BLOCK.
% Preserve V1--V33 and append subsequent work before the final terminator.


\clearpage
\hypersetup{pdftitle={Historical Standard-Model Source Selection: V1--V34}}
\begin{center}
{\Large\bfseries Source-selection continuation V34}\\[.4em]
{\large Outcome permutations, adjoint reversal, and\protect\\
phase-independent reconstruction of grading actions}
\end{center}
\addtocontents{toc}{\protect\begingroup}
\addtocontents{toc}{\protect\renewcommand{\protect\numberline}{\protect\makebox[2.1em][l]}}

\section[The interface beyond fixed labels]{Go upstream from fixed-label symmetry to the actual reversal rule}
\label{sel34:context}

V33 reconstructs fixed-label involutions by a finite binary code.  It
explicitly leaves two different interfaces: a grading may permute the
actual labels, or it may carry a forward operation to its Hilbert--Schmidt
adjoint without that adjoint being an admitted operation.  A trivial
fixed-label code answers neither question.  This continuation gives a
finite moment classification for both interfaces, followed by two simpler
phase-independent tests on important source classes.  The forward-only
qubit bank of V33 is evaluated without changing its coefficients: its
fixed-label action remains trivial, but it has a unique unitary
adjoint-reversal action.

\paragraph{Source and scope.}
All maps in this continuation act on one independently identified finite
carrier.  A permutation of histories must come from the declared grammar;
it is not permission to add new labels.  Formal adjoints in the flat
historical algebra are permitted for calculation, not automatically for
execution.  The source must distinguish a fixed-label symmetry, an outcome
permutation, an adjoint-reversal relation, and reversal of the temporal
order.  None is identified with a physical matter grading solely because
it exists mathematically.  The results use V33's actual efficient
\emph{generating history bank}; primitive outcomes need not all be
efficient.  Missing mixed outcomes are tested at the end.

Covariant instruments and their pointwise Kraus descriptions are established
objects \cite{sel34-covariant}.  Likewise, an adjoint of a bistochastic
channel is a normalized channel, whereas a general trace-preserving
channel's adjoint need not be trace preserving \cite{sel34-reverse}.
The contribution here is the inverse, same-source classification by finite
moments, its identity-anchor reduction, and its application to the
historical grading interface.  These references supply background, not
new physical records.  V33's fixed-label code, V32's spatial multiplicity
criterion, and V31's inserted moment calculus remain inherited inputs.

\section[Signed permutations and finite moment tests]{A complete finite test for permuted or adjointed coefficient rays}
\label{sel34:moments}

Let the nonzero matrices $k_1,\ldots,k_m$ represent actual efficient
history maps.  Set $\A=C^*(I,k_1,\ldots,k_m)$ and equip this algebra
with the normalized regular trace $\tau$ used in V33.  Choose a labelled
star-word basis $b_1=I,b_2,\ldots,b_d$, where $d=\dim_\C\A$.
For each basis word, its vector $v_i\in\mathbb Z^m$ counts occurrences
of $k_a$ minus occurrences of $k_a^*$.  The adjoints are formal algebra
letters, not new physical events.

Fix an involution $\pi$ of the generating labels and signs
$\epsilon_a\in\{1,-1\}$ such that $\epsilon_{\pi(a)}=\epsilon_a$.
Set
\[
 h_a=\begin{cases}k_{\pi(a)},&\epsilon_a=1,\\
                   k_{\pi(a)}^*,&\epsilon_a=-1.\end{cases}
\]
For a word $b_i$, let $b_i^{\pi,\epsilon}$ be its evaluation after
replacing $k_a$ by $h_a$ and $k_a^*$ by $h_a^*$, \emph{without
reversing the order of the letters}.  Define
\begin{align}
 G_{ij}&=\tau(b_i^*b_j),&
 G'_{ij}&=\tau((b_i^{\pi,\epsilon})^*b_j^{\pi,\epsilon}),\nonumber\\
 M^a_{ij}&=\tau(b_i^*k_ab_j),&
 M'^a_{ij}&=\tau((b_i^{\pi,\epsilon})^*h_ab_j^{\pi,\epsilon}).
 \label{sel34:eq:moments}
\end{align}
The target words need not be postulated to be a basis: the test below
will establish this when it succeeds.

\begin{theorem}[Finite signed-permutation reconstruction]
\label{sel34:thm:moment}
For $z=(z_1,\ldots,z_m)\in\mathrm U(1)^m$, the assignment
$k_a\mapsto z_ah_a$ extends to a unital complex-linear involutive
$*$-automorphism of $\A$ if and only if
\begin{align}
 z^{v_j-v_i}G'_{ij}&=G_{ij},\nonumber\\
 z^{v_j-v_i+e_a}M'^a_{ij}&=M^a_{ij},\label{sel34:eq:moment-test}\\
 z_a z_{\pi(a)}^{\epsilon_a}&=1\qquad(1\le a\le m).\nonumber
\end{align}
Here $z^v=\prod_a z_a^{v_a}$.  The action is unique and is reconstructed by
\begin{equation}
 \alpha_z\Big(\sum_i x_ib_i\Big)
   =\sum_i x_i z^{v_i}b_i^{\pi,\epsilon}.
 \label{sel34:eq:action}
\end{equation}
If the basis words have length at most $r$, only moments of degree at most
$2r+1$ are used.  A flat historical multiplication table already determines
all of them.
\end{theorem}
\begin{proof}
An automorphism preserves the regular trace because left multiplication
by $\alpha(x)$ is conjugate, as a linear operator on $\A$, to left
multiplication by $x$.  Multiplicativity and preservation of adjoints give
$\alpha(b_i)=z^{v_i}b_i^{\pi,\epsilon}$.  Trace invariance applied to
$G$ and $M^a$ gives the first two equalities.  On a generator,
\[
 \alpha^2(k_a)=z_a z_{\pi(a)}^{\epsilon_a}k_a,
\]
so the last equality is necessary.

Conversely use \eqref{sel34:eq:action} to define a linear map $T$ on the
basis.  The first equality makes $T$ an isometry of $L^2(\A,\tau)$.
It is therefore invertible, the target words are independent, and $T(I)=I$.
For every $i,j$, the second equality says
\[
 \langle Tb_i,z_ah_aTb_j\rangle
      =\langle b_i,k_ab_j\rangle
      =\langle Tb_i,T(k_ab_j)\rangle.
\]
Thus $T(k_ax)=z_ah_aT(x)$ for all $x$.  Taking adjoints of the moment
matrices supplies the analogous identity
$T(k_a^*x)=\overline z_a h_a^*T(x)$.  Iterate these identities for star
words and extend linearly in the first factor.  This proves
$T(xy)=T(x)T(y)$ and $T(x^*)=T(x)^*$.  The last equality in
\eqref{sel34:eq:moment-test} makes $T^2$ the identity on generators and
therefore on $\A$.  Generation proves uniqueness.  The degree count is
immediate from \eqref{sel34:eq:moments}.
\end{proof}

When all $\epsilon_a=1$, the theorem classifies outcome permutations of
these efficient maps.  When $\epsilon_a=-1$, the corresponding target map
is the Hilbert--Schmidt adjoint of the indicated CP map.  This is a
relation to an already computable map; it is not proof that this map is
an executable reverse event.  In the fixed-label case $\pi=\id$ and
$\epsilon_a=1$, the final equations impose $z_a^2=1$, and the theorem
reduces exactly to V33's binary code.

\begin{corollary}[Kraus-phase and basis independence]
\label{sel34:cor:phase}
Changing representatives to $k'_a=\lambda_a k_a$, $|\lambda_a|=1$,
changes the phase coordinates by
\begin{equation}
 z'_a=\frac{\lambda_a z_a}{\lambda_{\pi(a)}^{\epsilon_a}}.
 \label{sel34:eq:phase-gauge}
\end{equation}
It leaves the reconstructed automorphism family unchanged.  Changing the
labelled word basis likewise leaves that family unchanged.
\end{corollary}
\begin{proof}
Apply the same action to $k'_a$ and express the result in the new target
representative.  This gives \eqref{sel34:eq:phase-gauge}; it maps the
moment equations bijectively.  Basis independence follows from the
necessary-and-sufficient characterization.
\end{proof}

\section[Integer phase equations and physical filters]{The solution space can be a torus, not a binary code}
\label{sel34:torus}

\begin{theorem}[Exact phase-fibre classification]
\label{sel34:thm:torus}
For a prescribed $(\pi,\epsilon)$, first reject any moment equation in
\eqref{sel34:eq:moment-test} whose two scalar coefficients have unequal
modulus, including a mismatch of zeros.  Omit pairs which are both zero.
Every remaining equation is
\[
 z^{R_\ell}=\eta_\ell,\qquad R_\ell\in\mathbb Z^m,
 \quad |\eta_\ell|=1,
\]
including the order-two rows $e_a+\epsilon_a e_{\pi(a)}$ with right-hand
side one.  Let $URV=D$ be an integer Smith normal form, of rank $q$, with
nonzero diagonal entries $d_1\mid\cdots\mid d_q$.  Transform the targets
by $\eta'_j=\prod_\ell\eta_\ell^{U_{j\ell}}$.
The equations are feasible exactly when $\eta'_j=1$ for every zero row
of $D$.  If feasible, their solution set is a translate, as a phase space,
of
\begin{equation}
 \boxed{\mathrm U(1)^{m-q}\times
              \mu_{d_1}\times\cdots\times\mu_{d_q},}
 \label{sel34:eq:phase-fibre}
\end{equation}
where $\mu_d$ denotes the $d$th roots of unity.  Every phase point yields
one distinct action.  This description does not assert that the actions
with fixed nontrivial permutation form a group under composition.
\end{theorem}
\begin{proof}
Division by a nonzero target moment gives the monomial equation.  Integer
row operations multiply equations or invert them; integer column operations
are invertible monomial coordinate changes on the torus.  In Smith
coordinates the nonzero rows are $y_j^{d_j}=\eta'_j$.  Every point of the
unit circle has exactly $d_j$ such roots; a zero row is consistent exactly
as stated.  The remaining coordinates are free.  Distinct phase points
act differently on some nonzero generator.  This proves the claim,
including disconnected and positive-dimensional fibres.
\end{proof}

This is a finite exact test.  For algebraic input data the modulus, root,
and consistency comparisons are exact algebraic operations.  The zero
conditions cannot be established by rounding small numerical moments.
Several admissible outcome permutations give a finite union of these
families.  Efficient composite histories may be used, but their target
histories must be those induced by the actual primitive record rule.

\begin{proposition}[Same-carrier implementation and mixed-label filter]
\label{sel34:prop:spatial}
An involutive algebra action obtained above has a self-adjoint unitary
implementation on
$\A=\bigoplus_\lambda M_{n_\lambda}\otimes I_{m_\lambda}$ exactly
when the multiplicities agree along the action's central-block permutation.
On a full matrix carrier $M_N$, its implementer is unique up to sign.
For each phase point it is obtained from
\begin{equation}
 Tk_a=z_ah_aT,\qquad
 Tk_a^*=\overline z_a h_a^*T.
 \label{sel34:eq:intertwiner}
\end{equation}
The solution space is one-dimensional.  A nonzero solution can be normalized
and multiplied by a phase to give a self-adjoint unitary $Z$.

On $M_N$, retain a phase point only if every remaining original mixed
primitive label satisfies the intended relation
\[
 \operatorname{Ad}_Z\,\Phi_b\,\operatorname{Ad}_Z
     =\Phi_{\pi(b)}
 \quad\hbox{or}\quad
 \operatorname{Ad}_Z\,\Phi_b\,\operatorname{Ad}_Z
     =\Phi_{\pi(b)}^\dagger,
\]
as actually prescribed.  This is a direct Choi equality.  It neither
requires nor creates a Kraus refinement of that label.
\end{proposition}
\begin{proof}
The spatial multiplicity assertion is V33's same-carrier criterion and is
retained here, not replaced by an abstract dimension match.  On $M_N$,
inner implementation of the automorphism gives a nonzero solution of
\eqref{sel34:eq:intertwiner}.  Schur's lemma makes the solution line
one-dimensional.  Its adjoint product is scalar, and involutivity makes
the square of a unitary representative scalar; a phase gives $Z^2=I$.
The final equalities are exactly the required relations for the untested
maps.  On a represented algebra with residual multiplicity, one must
retain and test the implementation fibre rather than select an arbitrary
representative.  V32 describes that additional directed freedom.
\end{proof}

\section[Identity anchoring removes permutation phases]{The original locked class has a simpler, finite permutation audit}
\label{sel34:anchor}

Take linearly independent efficient coefficient vectors $|k_a\rangle\!\rangle$
and their actual Choi matrices $J_a$.  Put
\[
 S=\sum_aJ_a,\quad i=|I\rangle\!\rangle,
 \qquad i\in\Ran S,
\]
and assume full identity participation in the V20 sense.  Define directly
from the CP maps
\begin{equation}
 a_a=\langle i,S^\dagger J_a S^\dagger i\rangle>0,
 \qquad
 \ell_a=\operatorname{unvec}(J_aS^\dagger i).
 \label{sel34:eq:anchor}
\end{equation}
The $a_a$ are participation weights, not branch probabilities.

\begin{theorem}[Phase-free outcome-permutation criterion]
\label{sel34:thm:anchor}
The anchors in \eqref{sel34:eq:anchor} are independent and satisfy
$\sum_a\ell_a=I$.  A spatial automorphism permutes the actual efficient
maps by $\pi$ if and only if
\begin{equation}
 \boxed{a_{\pi(a)}=a_a,\qquad
               \alpha(\ell_a)=\ell_{\pi(a)}\quad\hbox{for all }a.}
 \label{sel34:eq:anchor-perm}
\end{equation}
For a fixed permutation there is at most one algebra action on the generated
algebra.  Its existence is tested by Theorem~\ref{sel34:thm:moment} with
$z_a=1$, $\epsilon_a=1$, applied to the anchors.  If $\pi^2=\id$,
that action is automatically involutive.
\end{theorem}
\begin{proof}
Let $W$ have columns $|k_a\rangle\!\rangle$ and write $i=Wr$ in its
unique coefficient expansion.  Then $W^*S^\dagger i=r$,
$\ell_a=r_ak_a$, and $a_a=|r_a|^2$.  Nonzero participation proves
independence and $\sum_a\ell_a=I$.

A spatial automorphism permuting $J_a$ fixes $i$, preserves $S$ and
$S^\dagger$, and therefore permutes both anchors and weights.
Conversely
$J_a=|\ell_a\rangle\!\rangle\langle\!\langle\ell_a|/a_a$,
so \eqref{sel34:eq:anchor-perm} gives the required Choi equalities.
The anchors generate the original algebra.  Hence an action on them is
unique; its square fixes the generators when $\pi^2=\id$.
\end{proof}

\begin{corollary}[The six-letter identity-calibrated bank]
\label{sel34:cor:76}
For six independent coefficients with $\sum_a k_a=\sqrt2 I$, one has
$a_a=1/2$ and $\ell_a=k_a/\sqrt2$.  All permutation phases are fixed:
$\alpha(k_a)=k_{\pi(a)}$.  There are at most
\begin{equation}
 \boxed{1+15+45+15=76}
 \label{sel34:eq:76}
\end{equation}
involutive edge permutations to test on one fixed record/type slice, each
with at most one action on its generated algebra.  Typed admissibility can
only reduce this list.  The identity permutation yields the trivial action,
as already proved in V20 and V33; nontrivial record permutations are not
excluded by that fixed-label conclusion.
\end{corollary}
\begin{proof}
The identity coefficients are all $1/\sqrt2$, giving the anchors.
The number of involutions on six points is
$\sum_{j=0}^3 6!/(2^j j!(6-2j)!)=76$.  The preceding theorem gives all
remaining assertions.  No unspecified historical edge matrix is populated
by this combinatorial count.
\end{proof}

\section[Adjoint reversal from one linear kernel]{A smaller phase-free test for trace-visible efficient histories}
\label{sel34:reverse}

On a full matrix carrier let the actual efficient generating histories have
$\Tr k_a\ne0$.  V21's intrinsic identity Choi moment is
\begin{equation}
 A_a=\frac1N\operatorname{unvec}(J_a|I\rangle\!\rangle)
     =\frac{\overline{\Tr k_a}}{N}k_a
     =H_a+iQ_a,
 \label{sel34:eq:trace-anchor}
\end{equation}
where $H_a,Q_a$ are Hermitian.  Nonzero trace ensures that these anchors
still generate $M_N$; it is an explicit entrance condition of the smaller
test, not of the general moment theorem.

\begin{theorem}[Unique adjoint-reversal or an exact kernel obstruction]
\label{sel34:thm:reverse}
Define the complex-linear operator
\begin{equation}
 \mathscr D_\#(T)=([T,H_a],\{T,Q_a\})_{a=1}^m,
 \qquad
 \mathscr L_\#=\mathscr D_\#^*\mathscr D_\#.
 \label{sel34:eq:reverse-kernel}
\end{equation}
There exists a self-adjoint unitary $Z$ with
\begin{equation}
 \operatorname{Ad}_Z\,\Phi_a\,\operatorname{Ad}_Z=\Phi_a^\dagger
 \quad\hbox{for every selected history}
 \label{sel34:eq:reverse-map}
\end{equation}
if and only if $\Ker\mathscr D_\#\ne0$.  A nonzero kernel is exactly
one-dimensional; normalizing any nonzero member reconstructs $Z$ up to sign.
The corresponding grading action is unique, and may be the identity.
The complete criterion depends only on the original CP maps, not on any
cross-label phase or admitted inverse operation.
\end{theorem}
\begin{proof}
For an efficient map, \eqref{sel34:eq:reverse-map} is equivalent to
$Zk_aZ=z_ak_a^*$ with $|z_a|=1$.  Taking the trace gives
$z_a=\Tr k_a/\overline{\Tr k_a}$.  By
\eqref{sel34:eq:trace-anchor}, this is equivalent to $ZA_aZ=A_a^*$,
or $[Z,H_a]=0$ and $\{Z,Q_a\}=0$.

Conversely let $T\ne0$ satisfy the linear equations.  Its adjoint obeys
the same equations, so $T^*T$ commutes with every $H_a,Q_a$.  They generate
$M_N$, hence $T^*T=cI$ with $c>0$, and $T$ is invertible.  If $T_1,T_2$
are two nonzero solutions, $T_1^{-1}T_2$ commutes with the generators and
is scalar.  Thus the kernel is a line.  After normalizing $T$ to a unitary,
$T^*$ lies on that same line.  Multiplication by a scalar phase makes it
self-adjoint; its square is then $I$.  The first part of the proof yields
\eqref{sel34:eq:reverse-map}.  Phase independence follows directly from
\eqref{sel34:eq:trace-anchor}.
\end{proof}

\begin{proposition}[Normalization and temporal-order safeguards]
\label{sel34:prop:reverse-scope}
If \eqref{sel34:eq:reverse-map}, possibly with a permutation of all
primitive labels, holds for a normalized instrument whose total map is
$\Phi$, then $\Phi$ must also be unital.  Unitality alone is not
sufficient for this simultaneous branch reversal.

If $\alpha\Phi_a\alpha=\Phi_a^\dagger$, then for an original word
$w=a_m\cdots a_1$,
\begin{equation}
 \alpha\Phi_w\alpha
 =\Phi_{a_m}^\dagger\cdots\Phi_{a_1}^\dagger
 =(\Phi_{a_1}\cdots\Phi_{a_m})^\dagger.
 \label{sel34:eq:reverse-order}
\end{equation}
This is generally not $\Phi_w^\dagger$.  The action remains
multiplicative and complex linear; it is neither the adjunction map on
all products nor an antiunitary time-reversal operator.
\end{proposition}
\begin{proof}
Summing all branch relations gives $\alpha\Phi\alpha=\Phi^\dagger$.
The left side is trace preserving, so the adjoint map is trace preserving;
equivalently $\Phi(I)=I$.  The three-axis example below proves
insufficiency.  Conjugation of a composition preserves its order, whereas
taking its Hilbert--Schmidt adjoint reverses order.  This proves
\eqref{sel34:eq:reverse-order} and the distinction.
\end{proof}

\begin{proposition}[One-sided errors and the reconstruction margin]
\label{sel34:prop:reverse-errors}
For Hermitian Choi estimates with
$\|\widehat J_a-J_a\|_{\HS}\le\delta_a$, the operators reconstructed
by \eqref{sel34:eq:trace-anchor}--\eqref{sel34:eq:reverse-kernel} satisfy
\begin{equation}
 \|\widehat{\mathscr D}_\#-\mathscr D_\#\|_{\op}
 \le\frac2{\sqrt N}\Big(\sum_a\delta_a^2\Big)^{1/2}=:\varepsilon_\#.
 \label{sel34:eq:reverse-error}
\end{equation}
Consequently $\sigma_{\min}(\widehat{\mathscr D}_\#)>\varepsilon_\#$
excludes an exact reverser.  On a verified nonzero-kernel branch, let
$\gamma_\#$ be the least positive eigenvalue of $\mathscr L_\#$.  Then
\[
 \operatorname{dist}_{\HS}(T,\mathbb CZ)
       \le\|\mathscr D_\#T\|_{\HS}/\sqrt{\gamma_\#}.
\]
Small estimated singular values do not prove exact existence.
\end{proposition}
\begin{proof}
The contraction defining $A_a$ has norm at most $1/\sqrt N$ from Choi
Hilbert--Schmidt norm to operator Hilbert--Schmidt norm.  If
$\Delta A=\Delta H+i\Delta Q$, then
$\|\Delta H\|_{\HS}^2+\|\Delta Q\|_{\HS}^2=\|\Delta A\|_{\HS}^2$.
Both commutator and anticommutator with a matrix $B$ have norm at most
$2\|B\|_{\op}$ on Hilbert--Schmidt space.  Sum the squared bounds to
obtain \eqref{sel34:eq:reverse-error}.  The singular-value perturbation
bound gives the exclusion.  Spectral decomposition of the positive form on
the complement of its one-dimensional kernel gives the last estimate.
\end{proof}

\section[Exact source evaluations and separating examples]{What the tests return on unchanged and calibrated comparison banks}
\label{sel34:examples}

\begin{theorem}[The unchanged V33 forward-pulse bank]
\label{sel34:thm:forward-example}
Retain exactly the two coefficients of
Proposition~\ref{sel33:prop:nonhomogeneous}:
\[
 k_X=(9I+12iX)/25,\qquad k_Z=(12I+16iZ)/25.
\]
Their fixed-label grading action is trivial, and no action swaps their two
recorded maps.  Nevertheless they have a unique adjoint-reversal action,
implemented by $Y$ up to sign.  Their intrinsic anchors are
\begin{equation}
 A_X=(81I+108iX)/625,\qquad
 A_Z=(144I+192iZ)/625.
 \label{sel34:eq:forward-anchors}
\end{equation}
On the orthonormal basis $(I,X,Y,Z)/\sqrt2$, the reversal form is diagonal:
\begin{equation}
 \boxed{\mathscr L_\#=
 \frac1{390625}\diag(194112,46656,0,147456),\qquad
 \gamma_\#=\frac{46656}{390625}.}
 \label{sel34:eq:forward-spectrum}
\end{equation}
Its source-derived Hermitian operators $Q_X=108X/625$ and
$Q_Z=192Z/625$ are both odd for this uniquely reconstructed action.
\end{theorem}
\begin{proof}
Normalization follows from
$k_X^*k_X=9I/25$ and $k_Z^*k_Z=16I/25$.  The nonzero traces force
positive phases for every fixed-label action, so the generating $X,Z$
must be fixed.  A swap is impossible because the two Choi traces,
$18/25$ and $32/25$, differ.  Formula
\eqref{sel34:eq:trace-anchor} gives \eqref{sel34:eq:forward-anchors}.
Their Hermitian parts are scalar, so only the two anticommutators with
$108X/625$ and $192Z/625$ enter.  The Pauli anticommutation identities
give \eqref{sel34:eq:forward-spectrum} and its one-dimensional kernel
$\mathbb CY$.  Theorem~\ref{sel34:thm:reverse} finishes the proof.
\end{proof}

No reverse pulse is inserted in this evaluation.  It proves an exact
same-cut relation to the calculable adjoint maps.  Identifying that relation
as the \emph{physical protected matter action}, or executing its reverse
maps, remains additional source content.  The adjoint map is not an inverse
map of a general channel.  Even here the original unequal branch weights
are retained.  The two primitive matrices do not commute:
\begin{equation}
 [k_X,k_Z]=\frac{384i}{625}Y,
 \qquad\|[k_X,k_Z]\|_{\HS}^2=\frac{294912}{390625}.
 \label{sel34:eq:order-example}
\end{equation}
Thus the ordering warning is visible on this unchanged bank.

\begin{example}[A continuous family for one outcome swap]
\label{sel34:ex:circle}
The normalized bank $k_1=E_{01}$, $k_2=E_{10}$ generates $M_2$.
For the swap $\pi=(12)$, every allowed action has
\[
 k_1\mapsto z k_2,\qquad k_2\mapsto z^{-1}k_1,
 \quad z\in\mathrm U(1).
\]
It is implemented by
$Z_\phi=\left(\begin{smallmatrix}0&e^{i\phi}\\e^{-i\phi}&0\end{smallmatrix}\right)$
with $z=e^{-2i\phi}$.  Conversely an outcome swap must exchange the two
rank-one support projections, so its self-adjoint implementer is of this
form up to sign.  The moment exponent matrix has Smith diagonal $(1,0)$:
one circle remains.  The fixed-label involutive actions, in contrast, are
only $\id$ and $\operatorname{Ad}_Z$.  Finiteness in V33 cannot therefore
be carried over to label permutations without an anchoring condition.
\end{example}

\begin{example}[A finite phase fibre and a rate obstruction]
\label{sel34:ex:pauli-swap}
For $k_1=X/\sqrt2$, $k_2=Z/\sqrt2$ and their swap, the moment system has
Smith diagonal $(1,2)$ and two solutions $z_1=z_2=\pm1$.  The
implementers are $(X+Z)/\sqrt2$ and $(X-Z)/\sqrt2$.  The identities
$k_a^2=I/2$ force $z_a^2=1$, and order two forces the signs to agree,
which proves completeness.  Replacing the coefficients by $3X/5,4Z/5$
preserves normalization but makes the swap impossible, by unequal Choi
traces.  Dividing away those weights would change the recorded source.
\end{example}

\begin{proposition}[The calibrated six-root projector family]
\label{sel34:prop:root}
In an orthonormal realization of the rank-minimal $A_3$ projector branch,
let
\[
 \begin{gathered}
 (v_a)_{a=1}^6=\frac1{\sqrt2}\bigl((1,1,0),(1,-1,0),(1,0,1),
                    (1,0,-1),(0,1,1),(0,1,-1)\bigr),\\
 P_a=v_av_a^*,\qquad k_a=P_a/\sqrt2.
 \end{gathered}
\]
These satisfy the two locked normalizations and generate $M_3$.
Their identity anchors are $\ell_a=P_a/2$, all with weight $1/2$.
Of the 76 involutive label permutations, exactly ten extend to algebra
actions: the identity and nine nonidentity actions.

Pairwise projector Grams alone admit twenty involutive permutations.  The
remaining ten are excluded by multiplication.  For example, swapping only
$P_1,P_2$ preserves the pair Gram, but
\begin{equation}
 \Tr(P_1P_3P_5)=\frac18,\qquad
 \Tr(P_2P_3P_5)=-\frac18.
 \label{sel34:eq:root-cubic}
\end{equation}
\end{proposition}
\begin{proof}
Summing the six projectors gives $2I_3$, so both normalization identities
follow.  Their diagonal sums span diagonal matrices and their differences
span the three symmetric off-diagonal directions.  Multiplication generates
all matrix units.  Independence follows from those six directions.  The
anchor theorem gives $\ell_a=P_a/2$.

Their three orthogonal pairs are $(1,2),(3,4),(5,6)$; all other pair overlaps
have square $1/4$.  Thus the pair-Gram permutations form the pair-permuting
group of order 48.  It has 20 involutions: eight preserving all three pairs
and twelve exchanging two pairs.  The actual unitary actions are the signed
coordinate permutations modulo their common sign, a group of order 24
isomorphic to $S_4$.  To check exhaustiveness, an action permuting the
projectors is inner; each orthogonal pair has common orthogonal complement
one coordinate axis.  It therefore permutes the coordinate axes and is a
monomial unitary.  Preservation of the six real root lines forces the three
coordinate phases to be equal up to sign.  The action group is consequently
exactly the stated quotient.  Its involutions are the identity, six
transpositions, and three double transpositions.  Finally, the product of
the three consecutive root overlaps gives \eqref{sel34:eq:root-cubic}.
\end{proof}

This evaluates the already classified projector \emph{source class}, in a
specified root frame.  It is not assigned to the historical undetermined
edge coefficients, and does not make its external triplet an internal
colour representation.  The ten-candidate conclusion is a grading-action
calculation on that comparison class.

\begin{proposition}[Two necessary safeguards for adjoint reversal]
\label{sel34:prop:no-reverse}
The normalized unital bank
$k_a=(3I+4i\sigma_a)/(5\sqrt3)$ for $\sigma_a=X,Y,Z$ has no simultaneous
unitary adjoint reverser.  Its anchors have scalar Hermitian parts and
imaginary parts $4\sigma_a/25$.  The form
$\mathscr L_\#$ has spectrum
\[
 \{192/625,64/625,64/625,64/625\},
\]
so its kernel is zero.  A unitary grading cannot negate all three Pauli axes.

The one-way bank of Theorem~\ref{sel33:thm:damping} has no complete
Hilbert--Schmidt adjoint-reversal instrument either: its total unitality
defect is $\diag(16/25,-16/25)$, with squared norm $512/625$.
Its different, fixed-label grading result in V33 is unchanged.
\end{proposition}
\begin{proof}
The first anchors follow from their traces.  Summing the three
anticommutator forms gives the displayed spectrum.  The total map is
unital because every Kraus coefficient is a scalar multiple of a unitary;
this proves that unitality is not sufficient.  For the second bank sum
$k_ak_a^*$ and subtract $I$.  Proposition~\ref{sel34:prop:reverse-scope}
then excludes adjoint reversal, independently of coefficient phases.
\end{proof}

\section[Historical consequences and status]{What is decided, and which physical row remains}
\label{sel34:status}

The current source audit has three distinct finite branches:
\[
 \boxed{
 \begin{gathered}
 \text{fixed labels}\longrightarrow\text{V33 binary code},\\
 \text{permuted labels}\longrightarrow\text{V34 moment/phase equations},\\
 \text{adjoint reversal}\longrightarrow\text{V34 signed equations or linear kernel}.
 \end{gathered}}
\]
The independent identity-anchored bank needs no continuous phase search for
outcome permutations.  The trace-visible efficient bank needs only one
linear kernel for same-label adjoint reversal.  Both are tests on original
quantum maps.  Neither replaces a failed fixed-label test by a conveniently
chosen, different physical symmetry.  A symmetry of the complete cylinder law additionally requires
the specified initial preparations and terminal Reads to transform
consistently.  Those finite boundary tests are not consequences of
operator-map covariance alone.

A returned nonidentity action determines V32's full odd-part compiler.
Directed multiplicity information, the correct input/output tensor types,
and the same-label determinant--Clifford coherence remain to be tested.
An outcome-permuting action may make an odd \emph{linear combination} of
two original coefficients available in the reconstructed algebra.  That
combination is not automatically one executed outcome.  V29's search must
still use the original resolved CP labels, not the symmetrized differences.
Likewise an algebraic inverse or adjoint word is not inserted into the
forward grammar.  This retains the source-complete distinction at the
centre of V26--V33.

All step weights are physical.  Adding the original idle opportunity and
scaling accepted Choi maps by $\vartheta$ leaves the reversal kernel
unchanged but scales \eqref{sel34:eq:reverse-kernel}'s quadratic form by
$\vartheta^2$; the idle identity anchor adds no constraint.  The positive
number in \eqref{sel34:eq:forward-spectrum} is a finite reconstruction
margin, not a physical mass gap or relaxation rate.

\paragraph{Proved in V34.}
The signed-permutation moment equivalence and its phase-coordinate
invariance; the finite Smith-form phase-fibre classification; the
same-carrier/mixed-map testing interface; the phase-free identity-anchor
permutation theorem and 76-permutation bound; the trace-anchor
adjoint-reversal kernel theorem and one-sided calibration bounds; the
unique reverser of the unchanged V33 forward bank; the continuous and
finite permutation fibres; the complete ten-involution root-projector
classification and its pair-versus-cubic distinction; and the unitality
and temporal-order obstructions.  The finite algebra and Smith-normal-form
methods are standard.  The source-specific inverse tests and evaluated
interfaces are the development claimed here.

\paragraph{Inherited.}
V33's full-matrix and represented-multiplicity implementation results,
V32's action/implementer and directed-typing distinction, V31's flat
historical moments, V20's identity participation, V21's intrinsic Choi
anchor, and V29's same-label search.  The fixed-label zero for the
identity-calibrated source is unchanged, as are the earlier determinant
and type-occurrence conditions.

\paragraph{Unresolved.}
The actual V110 matter grading has not been identified as a fixed-label,
record-permuting, or adjoint-reversing source action.  The unknown historical
edge matrices are not assigned the root-projector or qubit values used in
the exact comparisons.  No historical determinant--Clifford candidate word
has been populated by this continuation.  The source must identify its
physical record rule and then supply the appropriate mixed-map and directed
type rows.  No claim is made of historical Standard-Model selection,
fermionic measure control, or an Einstein--Standard-Model continuum.

\paragraph{Verification.}
The accompanying checker constructs the signed word bases and moment
exponent matrices, evaluates Smith invariants, independently solves
matrix intertwiner kernels, verifies Choi relations and retained
normalizations, and exhausts all 76 root-label involutions.  It also runs
the unmodified V33 regression.  These are exact arithmetic checks
supplementing the written proofs, not proof-assistant certification,
physical measurements, or independent peer review.

\addtocontents{toc}{\protect\endgroup}
\renewcommand{\refname}{Additional background for V34}
\begin{thebibliography}{99}
\raggedright
\bibitem[47]{sel34-covariant}
E.~Haapasalo and J.-P.~Pellonp\"a\"a,
\emph{Optimal covariant quantum measurements},
\href{https://arxiv.org/abs/2009.14080}{arXiv:2009.14080} (2020).
Background on finite-outcome covariant instruments and their Kraus
operators.  The present inverse tests do not import any source data.
\bibitem[48]{sel34-reverse}
E.~Aurell, J.~Zakrzewski and K.~\.{Z}yczkowski,
\emph{Time reversals of irreversible quantum maps},
\href{https://arxiv.org/abs/1505.02259}{arXiv:1505.02259} (2015).
Background for the distinction between bistochastic adjoints and other
notions of reverse dynamics; no antiunitary or detailed-balance assumption
is imposed here.
\end{thebibliography}

% END OF SOURCE-SELECTION V34 DEVELOPMENT BLOCK.
% Preserve V1--V34 and append subsequent work before the final terminator.


\clearpage
\hypersetup{pdftitle={Historical Standard-Model Source Selection: V1--V35}}
\begin{center}
{\Large\bfseries Source-selection continuation V35}\\[.4em]
{\large The stationary source fixes the reverse comparison:\\
nonunital grading reconstruction and its exact boundaries}
\end{center}
\addtocontents{toc}{\protect\begingroup}
\addtocontents{toc}{\protect\renewcommand{\protect\numberline}{\protect\makebox[2.1em][l]}}

\section[Which adjoint belongs to the source?]{Go upstream from the trace adjoint to the actual stationary source}
\label{sel35:context}

V34 correctly proves that Hilbert--Schmidt adjoint reversal of a normalized
instrument requires unitality.  This is a boundary of that particular
reverse comparison, not of every stationary reversal.  The companion
feedback/return analysis already retains faithful, possibly nontracial
stationary states and explicitly distinguishes its weighted observable
adjoints from ordinary trace adjoints; see \cite{Rigidity},
\src{v6:eq:faithfulness}, \src{v7:eq:occupation-metric}, and
\src{v7:thm:occupation}.  Those occupation and faithfulness theorems are
inherited.  They do not themselves identify a physical matter grading.

The present continuation uses the source's own invariant density to define
the stationary Petz reverse, and develops an inverse test for an involutive
\emph{unitary parity} relating each original branch to that reverse.  The
reverse construction and its trajectory interpretation are established
methods \cite{sel35-crooks,sel35-transition}.  What is developed here is the
source-state anchor, the unique-candidate kernel including mixed original
outcomes, its nonunital exact evaluation, and the separation from the
companion's GNS-metric adjoint.  This is not an assertion of physical
antiunitary time reversal, of an executable reverse gate, or of detailed
balance of an unspecified historical source.

\paragraph{Same-cut and data convention.}
All maps below act on the same independently identified $M_n(\C)$.
An actual finite instrument is $\{\Phi_a\}_{a\in\Sigma}$, with
$\Phi=\sum_a\Phi_a$ trace preserving.  The maps and the stationary
state are reconstructed on that carrier, for example by the contextual
procedure of V25.  Retained labels and their probabilities are not
renormalized individually.  We write $\Phi_a^\dagger$ for the ordinary
Hilbert--Schmidt adjoint.  A faithful stationary density means
\begin{equation}
 \sigma\succ0,\qquad \Tr\sigma=1,\qquad\Phi(\sigma)=\sigma.
 \label{sel35:eq:stationary}
\end{equation}
Existence is the inherited finite positive feasibility question; it is not
obtained by adding an arbitrarily small reset.  The uniqueness conditions
below make $\sigma$ a determined source output, rather than a density chosen
to make a desired symmetry pass.

\section[The normalized stationary reverse]{An exact reverse instrument without a unitality assumption}
\label{sel35:reverse}

Put $\Gamma_\sigma(X)=\sigma^{1/2}X\sigma^{1/2}$ and define
\begin{equation}
 \boxed{\Phi_a^{\mathrm R,\sigma}
       =\Gamma_\sigma\Phi_a^\dagger\Gamma_\sigma^{-1}.}
 \label{sel35:eq:petz}
\end{equation}
This equation refers to Schr\"odinger maps throughout.

\begin{proposition}[Stationary reversal preserves the complete normalization]
\label{sel35:prop:reverse}
Under \eqref{sel35:eq:stationary}, \eqref{sel35:eq:petz} is a CP instrument
whose total map is trace preserving and fixes $\sigma$.  If
$\Phi_a(X)=\sum_\mu k_{a\mu}Xk_{a\mu}^*$, its reverse coefficients are
\begin{equation}
 r_{a\mu}=\sigma^{1/2}k_{a\mu}^*\sigma^{-1/2}.
 \label{sel35:eq:reverse-kraus}
\end{equation}
Reversal is involutive and reverses composition order.  In particular,
for $\Phi_w=\Phi_{a_m}\cdots\Phi_{a_1}$,
\begin{equation}
 \Phi_w^{\mathrm R,\sigma}
  =\Phi_{a_1}^{\mathrm R,\sigma}\cdots\Phi_{a_m}^{\mathrm R,\sigma}.
 \label{sel35:eq:reverse-order}
\end{equation}
The stationary word probabilities obey
\begin{equation}
 \Tr[\Phi_{a_m}\cdots\Phi_{a_1}(\sigma)]
 =\Tr[\Phi_{a_1}^{\mathrm R,\sigma}\cdots
             \Phi_{a_m}^{\mathrm R,\sigma}(\sigma)].
 \label{sel35:eq:probability-reverse}
\end{equation}
\end{proposition}
\begin{proof}
The displayed coefficients give complete positivity and
\[
 \sum_{a,\mu}r_{a\mu}^*r_{a\mu}
 =\sigma^{-1/2}\Phi(\sigma)\sigma^{-1/2}=I.
\]
The total map fixes $\sigma$, since
$\Gamma_\sigma\Phi^\dagger\Gamma_\sigma^{-1}(\sigma)
 =\Gamma_\sigma\Phi^\dagger(I)=\sigma$.
Each positive branch is trace nonincreasing because their sum is trace
preserving.  The superoperators $\Gamma_\sigma$ and its inverse are
self-adjoint in the Hilbert--Schmidt metric.  Substitution proves
$(\Phi_a^{\mathrm R,\sigma})^{\mathrm R,\sigma}=\Phi_a$;
taking the adjoint of a product proves \eqref{sel35:eq:reverse-order}.
Finally,
$\Tr[\Gamma_\sigma\Phi_w^\dagger(I)]
 =\Tr[\sigma\Phi_w^\dagger(I)]=\Tr[\Phi_w(\sigma)]$.
The result is independent of the chosen Kraus list.
\end{proof}

\begin{lemma}[State uniqueness and preservation are source tests]
\label{sel35:lem:state}
If \eqref{sel35:eq:stationary} holds and the complete coefficient algebra
of $\Phi$ is $M_n$, its stationary density is unique.  More generally,
uniqueness is decided by the finite fixed-space equation
$\dim\ker(\Phi-I)=1$.  The reverse total map has the same fixed-space
dimension.  If the stationary density is unique and an involutive unitary
conjugation $\alpha$ satisfies
\begin{equation}
 \alpha\Phi_a\alpha=\Phi_{\pi(a)}^{\mathrm R,\sigma}
 \label{sel35:eq:parity-permutation}
\end{equation}
for a permutation $\pi$ of the retained labels, then $\alpha(\sigma)=\sigma$.
\end{lemma}
\begin{proof}
For a Hermitian fixed observable $h=\Phi^\dagger(h)$, normalization gives
\[
 \Phi^\dagger(h^2)-h^2
 =\sum_{a,\mu}[h,k_{a\mu}]^*[h,k_{a\mu}]\succeq0.
\]
Its trace against $\sigma$ is zero by stationarity, and faithfulness forces
every commutator to vanish.  A full coefficient algebra therefore has
only scalar fixed observables.  A linear map and its adjoint have the same
nullity at eigenvalue one, giving uniqueness of the invariant state and
the stated fixed-space test.  Equation \eqref{sel35:eq:petz} makes the
reverse similar to $\Phi^\dagger$, so its fixed-space dimension is the same.
Summing \eqref{sel35:eq:parity-permutation} shows that $\alpha(\sigma)$ is
stationary for the reverse.  Its unique stationary density is $\sigma$.
\end{proof}

In particular, a self-adjoint unitary implementer $Z$ of an actual
reversal must commute with the reconstructed $\sigma$.  This commutation
need not be guessed as an independent grading axiom when the source state
is unique.

\begin{proposition}[A stationary Choi reflection and a complete label test]
\label{sel35:prop:choi}
Use vectorization $|k\rangle\!\rangle$ in output--input order, and let
$\mathcal C|k\rangle\!\rangle=|k^*\rangle\!\rangle$ be the antiunitary
adjunction map.  Define the positive stationary Choi packet
\begin{equation}
 \Omega_a=(I\otimes\overline{\sigma}^{1/2})J_a
                         (I\otimes\overline{\sigma}^{1/2}).
 \label{sel35:eq:stationary-choi}
\end{equation}
Then $\Omega_a^{\mathrm R,\sigma}=\mathcal C\Omega_a\mathcal C$.
For $Z=Z^*=Z^{-1}$ commuting with $\sigma$,
\eqref{sel35:eq:parity-permutation} is equivalent to
\begin{equation}
 (Z\otimes\overline Z)\Omega_a(Z\otimes\overline Z)^*
                 =\mathcal C\Omega_{\pi(a)}\mathcal C
 \quad\text{for every original label }a.
 \label{sel35:eq:stationary-residual}
\end{equation}
It follows that the original stationary law is invariant under
$ a_1\cdots a_m\mapsto\pi(a_m)\cdots\pi(a_1)$.
\end{proposition}
\begin{proof}
Equation \eqref{sel35:eq:stationary-choi} is the sum of the projectors on
$|k_{a\mu}\sigma^{1/2}\rangle\!\rangle$.  By
\eqref{sel35:eq:reverse-kraus},
$r_{a\mu}\sigma^{1/2}=(k_{a\mu}\sigma^{1/2})^*$, proving the reflection
identity.  Commutation of $Z$ with $\sigma$ makes the stationary congruence
intertwine its system conjugation.  This congruence is invertible, so Choi
equality is equivalent to equality of the original maps.  Insert
\eqref{sel35:eq:parity-permutation} in
\eqref{sel35:eq:probability-reverse}, cancel intermediate $\alpha$'s, and
use trace invariance and $\alpha(\sigma)=\sigma$.
\end{proof}

\begin{corollary}[The exact two-sided contextual identity]
\label{sel35:cor:boundary}
For positive operators $E,F$ and every actual word map $\Phi_w$,
\begin{equation}
 \Tr[E\Phi_w(\Gamma_\sigma(F))]
 =\Tr[F\Phi_w^{\mathrm R,\sigma}(\Gamma_\sigma(E))].
 \label{sel35:eq:boundary}
\end{equation}
If \eqref{sel35:eq:parity-permutation} holds, the right side equals
$\Tr[\alpha(F)\Phi_{w^{\leftarrow\pi}}
       (\Gamma_\sigma(\alpha(E)))]$, where
$w^{\leftarrow\pi}=\pi(a_m)\cdots\pi(a_1)$ in chronological order.
\end{corollary}
\begin{proof}
Move $\Phi_w$ to its trace adjoint, then use self-adjointness of
$\Gamma_\sigma$ and
$\Phi_w^{\mathrm R,\sigma}\Gamma_\sigma
 =\Gamma_\sigma\Phi_w^\dagger$.
For the second assertion apply \eqref{sel35:eq:reverse-order}, substitute
\eqref{sel35:eq:parity-permutation}, and use commutation of $\alpha$
with $\Gamma_\sigma$.  The boundary weights are not divided out.
\end{proof}

The operators $\Gamma_\sigma(F)$ are unnormalized states; their physical
preparation or contextual reconstruction must be established before
\eqref{sel35:eq:boundary} is interpreted as an executed comparison.
The formula does not add a reverse word to the original grammar.

The antiunitary $\mathcal C$ here represents adjunction of data matrices;
it is not a newly postulated physical time-reversal operation.  The
normalized positive packet $\sum_a\Omega_a$ has trace one.  Its
reconstruction is a deterministic transformation of the original Choi
maps and the source density, not a new interference experiment.

\section[A phase-free stationary grading kernel]{A unique candidate from stationary anchors, with an exact mixed-outcome filter}
\label{sel35:kernel}

For each original map define its stationary anchor
\begin{equation}
 \boxed{A_a^\sigma=\operatorname{unvec}(J_a|\sigma\rangle\!\rangle)
  =\sum_\mu\overline{\Tr(\sigma k_{a\mu})}\,k_{a\mu}.}
 \label{sel35:eq:anchor}
\end{equation}
It is independent of the Kraus frame.  For $\sigma=I/n$ it is exactly the
identity Choi moment used in V21 and V34.  Introduce the algebraic whitening
\begin{equation}
 B_a=\sigma^{-1/4}A_a^\sigma\sigma^{1/4}=H_a+iQ_a,
 \qquad H_a=H_a^*,\quad Q_a=Q_a^*,
 \label{sel35:eq:whiten}
\end{equation}
and the linear space $\mathcal E_\sigma=\{T:[T,\sigma]=0\}$.  On it put
\begin{equation}
 \mathscr D_\sigma T=([T,H_a],\{T,Q_a\})_{a\in\Sigma},\qquad
 \mathscr L_\sigma=(\mathscr D_\sigma|_{\mathcal E_\sigma})^*
                         (\mathscr D_\sigma|_{\mathcal E_\sigma}).
 \label{sel35:eq:kernel}
\end{equation}
Whitening is used to solve equations; it is not an executed channel or a
change to the physical branch normalization.

\begin{theorem}[Stationary anchor reconstruction]
\label{sel35:thm:kernel}
Suppose \eqref{sel35:eq:stationary} holds and
\begin{equation}
 C^*(I,A_a^\sigma:a\in\Sigma)=M_n(\C).
 \label{sel35:eq:anchor-generation}
\end{equation}
Then $\sigma$ is the unique stationary density.  Any unitary involutive
\emph{same-label} stationary reverser
\begin{equation}
 \operatorname{Ad}_Z\Phi_a\operatorname{Ad}_Z
                      =\Phi_a^{\mathrm R,\sigma}\quad(a\in\Sigma)
 \label{sel35:eq:same-label}
\end{equation}
lies in $\ker\mathscr D_\sigma$.  This kernel is either zero or one
dimensional.  In the second case it reconstructs one self-adjoint unitary
$Z$, unique up to sign, as the only possible action satisfying
\eqref{sel35:eq:same-label}.

For arbitrary mixed original outcomes the candidate succeeds exactly when
the finite Choi residuals \eqref{sel35:eq:stationary-residual}, with
$\pi=\id$, all vanish.  If every tested outcome is efficient and
$\Tr(\sigma k_a)\ne0$, the residual test is automatic: a nonzero kernel
is necessary and sufficient for \eqref{sel35:eq:same-label}.
\end{theorem}
\begin{proof}
The anchors lie in the original coefficient span, so their full generated
algebra implies a full coefficient algebra.  Lemma~\ref{sel35:lem:state}
gives uniqueness and shows that any reverser commutes with $\sigma$.
The weighted trace of \eqref{sel35:eq:reverse-kraus} is
$\Tr(\sigma r_{a\mu})=\overline{\Tr(\sigma k_{a\mu})}$.
Consequently the reverse anchor is
\[
 (A_a^\sigma)^{\mathrm R}=\sigma^{1/2}(A_a^\sigma)^*\sigma^{-1/2}.
\]
The same-label equation therefore implies
$ZA_a^\sigma Z=(A_a^\sigma)^{\mathrm R}$.
After whitening, this is $ZB_aZ=B_a^*$, equivalently the two equations in
\eqref{sel35:eq:kernel}.

Conversely take nonzero $T$ in that kernel.  The equations imply
$TB_a=B_a^*T$ and $TB_a^*=B_aT$; their adjoints show that $T^*$ satisfies
the same equations.  Hence $T^*T$ commutes with all $B_a,B_a^*$ and with
$\sigma$.  Since $A_a^\sigma=\sigma^{1/4}B_a\sigma^{-1/4}$, it commutes
with all anchors and their adjoints, so it is a positive scalar multiple
of $I$.  Thus $T$ is invertible and its polar normalization $U$ is unitary.
The same argument shows that any two solutions differ by a scalar, and
$U^2$ commutes with the anchors and $\sigma$.  It is scalar.  Choosing a
scalar phase makes $U^2=I$, which also makes $U=U^*$.  This constructs $Z$
and proves uniqueness up to sign.  It establishes anchor reversal, not yet
reversal of a mixed map; Proposition~\ref{sel35:prop:choi} gives the exact
remaining test.

For an efficient outcome put $t_a=\Tr(\sigma k_a)\ne0$.
The anchor equality reads
$\overline{t_a}Zk_aZ=t_a\sigma^{1/2}k_a^*\sigma^{-1/2}$.
The two coefficients therefore differ by the unit-modulus scalar
$t_a/\overline{t_a}$, proving equality of their CP maps.  No coefficient
phase is a new source datum.
\end{proof}

The action can be the identity: nonidentity must be checked after the
kernel is solved.  If \eqref{sel35:eq:anchor-generation} fails, the theorem
does not assert uniqueness or nonexistence.  In that case the full Choi
equations remain valid, and the previous permutation/moment methods or a
larger actual anchor bank may be needed.  A zero stationary trace is not a
license to split a mixed physical outcome into unrecorded Kraus labels.

\begin{corollary}[Existing contextual data suffice]
\label{sel35:cor:contextual}
On a same-cut contextually tomographic branch, the finite original cylinder
table reconstructs $\Phi_a$, the fixed-space equation for $\sigma$, the
anchors \eqref{sel35:eq:anchor}, and the kernel and Choi residuals above.
No direct execution of $\Gamma_\sigma^{-1}$, a reverse branch, or $Z$ is
required for the inference.  Unitary changes of the reconstructed memory
coordinates conjugate every answer and preserve all tests.
\end{corollary}
\begin{proof}
Use V25 for the original Choi maps.  Stationarity is a finite linear
system followed by the actual positivity test.  All remaining expressions
are determined matrix operations.  Under $k\mapsto VkV^*$ and
$\sigma\mapsto V\sigma V^*$, functional calculus and
\eqref{sel35:eq:anchor} give $B_a\mapsto VB_aV^*$, and the kernel
transports by the same conjugation.  No signed statistical reconstruction
is promoted to a physical operation.
\end{proof}

\section[An exact nonunital source]{A source-determined state and grading with no trace-adjoint reverse}
\label{sel35:example}

\begin{theorem}[A normalized nonunital comparison with a unique branchwise reverser]
\label{sel35:thm:example}
Consider the declared two-outcome comparison instrument
\begin{equation}
 k_+=\frac1{\sqrt{10}}\begin{pmatrix}2&2\\-1&1\end{pmatrix},\qquad
 k_-=\frac1{\sqrt{10}}\begin{pmatrix}2&-2\\1&1\end{pmatrix}.
 \label{sel35:eq:example}
\end{equation}
Its coefficients are invertible and generate $M_2$.  Its unique stationary
state, channel spectrum, and unitality defect are
\begin{equation}
 \boxed{\sigma=\diag(4/5,1/5),\quad
 \operatorname{spec}\Phi=\{1,4/5,0,0\},\quad
 \Phi(I)-I=\frac35 Z.}
 \label{sel35:eq:example-state}
\end{equation}
There is no complete unitary Hilbert--Schmidt adjoint reversal.  Nevertheless
its unique same-label stationary reversal action is $\operatorname{Ad}_Z$:
\begin{equation}
 \sigma^{1/2}k_\pm^*\sigma^{-1/2}=k_\mp=Zk_\pm Z.
 \label{sel35:eq:example-reverse}
\end{equation}
The stationary anchors and their whitened values are
\begin{equation}
 A_\pm^\sigma=\frac9{50}\begin{pmatrix}2&\pm2\\\mp1&1\end{pmatrix},
 \qquad B_\pm=\frac9{50}\begin{pmatrix}2&\pm\sqrt2\\\mp\sqrt2&1\end{pmatrix}.
 \label{sel35:eq:example-anchor}
\end{equation}
On $\mathcal E_\sigma=\Span\{I,Z\}$, with basis $(I,Z)/\sqrt2$,
\begin{equation}
 \boxed{\mathscr L_\sigma=\diag(324/625,0),\qquad
           \ker\mathscr D_\sigma=\C Z.}
 \label{sel35:eq:example-gap}
\end{equation}
The original one-step probabilities at stationarity are $1/2,1/2$;
the two-step probabilities are
\begin{equation}
 p(++ )=p(--)=\frac{89}{500},\qquad
 p(+-)=p(-+)=\frac{161}{500}.
 \label{sel35:eq:example-words}
\end{equation}
Thus the retained word law is not an independently resampled uniform law.
\end{theorem}
\begin{proof}
Direct multiplication gives
$k_\pm^*k_\pm=I/2\pm3X/10$, proving normalization.
The difference of the two matrices is nonzero off diagonal; their sum has
two distinct diagonal entries.  Its spectral projections and the off-diagonal
corners therefore generate $M_2$.  In the Pauli coordinates,
\[
 \Phi(I)=I+3Z/5,\quad \Phi(X)=0,\quad
 \Phi(Y)=4Y/5,\quad \Phi(Z)=0.
\]
This proves \eqref{sel35:eq:example-state}, including uniqueness of the
stationary state.  The unitality defect has squared Hilbert--Schmidt norm
$18/25$, so the V34 obstruction to trace-adjoint reversal applies.
Equation \eqref{sel35:eq:example-reverse} is direct.  The stationary trace
of either coefficient is $9/(5\sqrt{10})$, giving
\eqref{sel35:eq:example-anchor}.  Their Hermitian parts are
$27I/100+9Z/100$ and their imaginary parts are
$\pm9\sqrt2Y/50$.  The commutator vanishes on $\mathcal E_\sigma$;
the two anticommutator squares give \eqref{sel35:eq:example-gap}.
Theorem~\ref{sel35:thm:kernel} proves completeness.  Finally
$\Phi_\pm(\sigma)=\sigma/2\mp3X/25$; taking the next outcome effect
gives \eqref{sel35:eq:example-words}.
\end{proof}

This instrument is an exact comparison, not the unknown historical Store.
The same $Z$ can also be recognized by V34 as an outcome exchange.  The new
point is that this exchange equals the \emph{normalized stationary reverse}
while the ordinary trace-adjoint reverse is excluded.  Here the reverse
maps are already the two other retained branches; no added physical gate
is needed even for that realization.  The total channel satisfies
$\Phi^{\mathrm R,\sigma}=\Phi$, so forgetting the labels admits the
identity parity as well.  Branchwise reversal is the condition selecting
the nonidentity action.

\begin{proposition}[The unchanged V34 example and the singular one-way boundary]
\label{sel35:prop:inherited}
For the unchanged V34 bank
$k_X=(9I+12iX)/25$, $k_Z=(12I+16iZ)/25$, the unique stationary state is
$I/2$.  The stationary anchors equal V34's original anchors; the reverse,
its unique action $\operatorname{Ad}_Y$, and its response margin
$46656/390625$ are unchanged.

For the unchanged one-way V33 bank
$k_0=\diag(1,3/5)$, $k_1=(4/5)E_{01}$, the only stationary density is
$E_{00}$.  It therefore has no faithful stationary reverse on the stated
$M_2$ carrier.  Restricting to its stationary support gives $M_1$, on which
there is no nonidentity grading action.
\end{proposition}
\begin{proof}
The first channel is unital and its coefficients generate $M_2$; apply
Lemma~\ref{sel35:lem:state}.  At $\sigma=I/2$ all formulae reduce to V34.
For the one-way bank stationarity of the lower population gives
$x_{11}=9x_{11}/25$, hence $x_{11}=0$.  Positivity forces its off-diagonal
entries to vanish.  Trace one gives $E_{00}$.  The support statement follows.
Its fixed-label grading from V33 is not contradicted: that is a different
source relation.
\end{proof}

\section[Three boundaries for reverse inference]{Observable adjoints, mixed anchors, and scalar word symmetry are different tests}
\label{sel35:boundaries}

\begin{proposition}[The GNS adjoint need not be a reverse CP map]
\label{sel35:prop:gns}
Let $\mathsf T=\Phi^\dagger$ be the Heisenberg map and use the
companion's one-sided metric
$\langle X,Y\rangle_\sigma=\Tr(\sigma X^*Y)$.
Its Hilbert adjoint is
\begin{equation}
 \mathsf T^{\ddagger_{\rm GNS}}(Y)=\Phi(Y\sigma)\sigma^{-1}.
 \label{sel35:eq:gns}
\end{equation}
This need not preserve Hermiticity, even when $\sigma$ is faithful.
For \eqref{sel35:eq:example},
\begin{equation}
 \mathsf T^{\ddagger_{\rm GNS}}(X)
       =\begin{pmatrix}0&-6/5\\3/10&0\end{pmatrix}.
 \label{sel35:eq:gns-counter}
\end{equation}
By contrast, the adjoint for the symmetric metric
$\langle X,Y\rangle_{\sigma,1/2}
 =\Tr(\sigma^{1/2}X^*\sigma^{1/2}Y)$ is
$\Gamma_\sigma^{-1}\Phi\Gamma_\sigma$, a unital CP map.  Its
ordinary trace dual is \eqref{sel35:eq:petz}.
\end{proposition}
\begin{proof}
The one-sided metric is the Hilbert--Schmidt metric with weight
$R_\sigma:Y\mapsto Y\sigma$.  Its adjoint is
$R_\sigma^{-1}\mathsf T^\dagger R_\sigma$, proving
\eqref{sel35:eq:gns}.  Substitute $X\sigma$ and use the Pauli action in
Theorem~\ref{sel35:thm:example} to obtain
\eqref{sel35:eq:gns-counter}; a positive complex-linear map must preserve
Hermiticity.  The symmetric metric has weight $\Gamma_\sigma$ instead,
so its adjoint is the stated CP sandwich.  Stationarity gives its value
$I$ at $I$, and taking the trace dual gives \eqref{sel35:eq:petz}.
\end{proof}

This does not change the companion's variational circulation identities:
they use the GNS Hilbert adjoint as such.  It prevents interpreting that
adjoint as an available or even positive Schr\"odinger reverse instrument
without an additional theorem.  The present stationary reverse is one
specified CP-compatible comparison, not a claim that all quantum notions
of detailed balance are equivalent.

\begin{proposition}[A mixed original label can pass the anchors and fail reversal]
\label{sel35:prop:mixed}
Retain V34's $k_X,k_Z$ above, take $0<\epsilon<1$, and put
$V=(X+Y)/\sqrt2$.  Define two original mixed-label maps by
\[
 \Psi_1=(1-\epsilon)\Phi_X+\epsilon\operatorname{Ad}_V,
 \qquad \Psi_2=(1-\epsilon)\Phi_Z.
\]
Their sum is a normalized unital channel with unique stationary state
$I/2$.  The stationary anchors are $(1-\epsilon)$ times the original
V34 anchors, so the unique candidate remains $Y$.  But, with the ordinary
unnormalized Choi convention,
\begin{equation}
 \boxed{\|J(\operatorname{Ad}_Y\Psi_1\operatorname{Ad}_Y)
                 -J(\Psi_1^{\mathrm R,I/2})\|_{\HS}^2
                    =8\epsilon^2>0.}
 \label{sel35:eq:mixed-residual}
\end{equation}
No simultaneous same-label stationary reverser exists.  The corresponding
difference of the first branch maps has diamond norm $2\epsilon$.
\end{proposition}
\begin{proof}
Normalization and unitality are preserved by the displayed positive
combination.  Its coefficient algebra is still $M_2$, so the state is
unique.  Since $\Tr V=0$, this term contributes zero to the anchor.
The two original efficient terms obey V34 reversal.  The remaining Choi
difference is $\epsilon$ times the difference of the projectors on
$|YVY\rangle\!\rangle$ and $|V\rangle\!\rangle$.
These two vectors are orthogonal and each has squared norm two, giving
\eqref{sel35:eq:mixed-residual}.  The maximally entangled input produces
orthogonal output pure states for the two unitary channels, giving diamond
norm two before multiplication by $\epsilon$.  The unique-candidate
theorem excludes a different hidden reverser.
\end{proof}

\begin{proposition}[Perfect scalar reversal does not force quantum reversal]
\label{sel35:prop:scalar}
For the unchanged three-axis comparison of V34,
$k_a=(3I+4i\sigma_a)/(5\sqrt3)$, $a=X,Y,Z$, every original record word
has probability $3^{-|w|}$ from every initial state.  Thus its complete
stationary record law is exactly invariant under order reversal, yet it
has no same-label unitary stationary reverser.
\end{proposition}
\begin{proof}
Each coefficient is a unitary times $1/\sqrt3$, so every branch has
state-independent probability $1/3$, at every history.  The channel is
unital with full coefficient algebra, hence its unique stationary state
is $I/2$.  Its stationary reverse is the trace adjoint, and V34's positive
reversal form, with eigenvalues $192/625,64/625,64/625,64/625$, has zero
kernel.  Equivalently a unitary would have to negate all three Pauli axes,
contradicting preservation of their ordered products.
\end{proof}

The probability identity \eqref{sel35:eq:probability-reverse} is therefore
not a substitute for quantum output-sensitive reconstruction.  This
boundary is all-depth, not a failure to sample enough long words.

\section[Finite calibration and physical frontier]{Source-normalized error control and the remaining occurrence question}
\label{sel35:status}

\begin{proposition}[A one-sided stable kernel test]
\label{sel35:prop:error}
Fix an exactly identified faithful $\sigma$ and its commutant
$\mathcal E_\sigma$.  If
$\|\widehat A_a^\sigma-A_a^\sigma\|_\op\le e_a$, put
\[
 d_a=\left(\frac{\lambda_{\max}(\sigma)}
                    {\lambda_{\min}(\sigma)}\right)^{1/4}e_a,
 \qquad \delta=2\sqrt{2\sum_a d_a^2}.
\]
Then $\|\widehat{\mathscr D}_\sigma-\mathscr D_\sigma\|\le\delta$ and
\begin{equation}
 \|\widehat{\mathscr L}_\sigma-\mathscr L_\sigma\|
       \le\delta(2\|\widehat{\mathscr D}_\sigma\|+\delta).
 \label{sel35:eq:error}
\end{equation}
A least eigenvalue of the estimated form exceeding this bound certifies
nonexistence of a stationary reverser under the theorem's hypotheses.
If the exact kernel is known to be one dimensional with positive spectral
gap $\gamma$, then for any unit Hilbert--Schmidt vector $T\in\mathcal E_\sigma$,
\begin{equation}
 \operatorname{dist}(T,\ker\mathscr D_\sigma)
 \le\frac{\|\widehat{\mathscr D}_\sigma T\|+\delta}{\sqrt\gamma}.
 \label{sel35:eq:kernel-stability}
\end{equation}
\end{proposition}
\begin{proof}
Submultiplicativity bounds the whitening error by $d_a$.  Taking Hermitian
and imaginary parts does not increase that bound.  A commutator or
anticommutator by an operator of norm at most $d_a$ has Hilbert--Schmidt
operator norm at most $2d_a$.  Summing their squared estimates proves the
bound on $\mathscr D$.  Expand the difference of the two Gram operators to
obtain \eqref{sel35:eq:error}.  Weyl's inequality proves the exclusion.
On the orthogonal complement of the exact kernel,
$\|\mathscr D T\|\ge\sqrt\gamma\|T\|$; the triangle inequality then
proves \eqref{sel35:eq:kernel-stability}.
\end{proof}

The state in this bound is fixed, not estimated silently.  If its exact
spectral projectors are unresolved, the equation $[T,\sigma]=0$ must also
be tested with its own certified state error; its kernel dimension cannot
be inferred from nearly repeated eigenvalues.  Small residuals do not
establish exact reversal.  Faithfulness enters the fourth-root whitening
through its measured eigenvalue floor.

For the original idle opportunity with state-independent acceptance
$\vartheta>0$, the total map is
$(1-\vartheta)\id+\vartheta\Phi$ and has the same stationary states.
Its stationary reverse retains exactly these two weights.  Every accepted
anchor is multiplied by $\vartheta$, the idle anchor is scalar, and
\begin{equation}
 \mathscr L_{\sigma,\mathrm{opportunity}}
                  =\vartheta^2\mathscr L_{\sigma,\mathrm{accepted}}.
 \label{sel35:eq:opportunity}
\end{equation}
The commutation condition with $\sigma$ is imposed on the same space
$\mathcal E_\sigma$ rather than assigned an arbitrary weighted penalty.
The null direction is unchanged; this response is not a dynamical mass gap.
State-dependent postselection does not admit this simplification.

\paragraph{Proved in V35.}
A source-state-normalized reverse instrument and its complete temporal
order rule; the state-preservation implication and stationary Choi
reflection; the phase-free stationary-anchor kernel, uniqueness of its
candidate, and the exact mixed-outcome filter; an invertible two-outcome
nonunital example with recovered state, parity, response and record
probabilities; the GNS-versus-CP-reverse counterexample; unchanged-source
checks and the singular one-way boundary; a false-positive mixed anchor
with an exact Choi and diamond defect; the all-depth scalar-record
obstruction; and the one-sided conditioning and opportunity bounds.
The reverse construction is standard.  The inverse-source interface and
its evaluated boundaries are the relative development.

\paragraph{Inherited inputs.}
The companion's faithful-state and occupation reconstruction,
V34's trace-adjoint and outcome-permutation results, V25's contextual
process tomography, V32's odd-part compiler, and V29's actual-label
search.  Nothing here changes V34's unitality theorem: it concerns a
different reverse, recovered here only when $\sigma=I/n$.  No old
coefficient, physical rate, or cutoff is changed in the inherited examples.

\paragraph{Unresolved physical rows.}
The historical V110 grading has not been identified with this stationary
reverse comparison.  A uniquely reconstructed $Z$ would determine
$(A-ZAZ)/2$ for an already identified self-adjoint full-cell generator $A$;
it does not supply that generator, its directed type maps, or its
same-label Clifford--determinant word.  The historical CP table and its
stationary state have not been numerically populated here.  Where the
source has no faithful stationary density on the retained carrier, the
present inverse cannot be used by deleting transient matter directions or
adding noise.  No Einstein--Standard-Model continuum or fermionic measure
claim is made.

\paragraph{Next source-level fork.}
If the physical comparison is the trace adjoint, V34 applies.  If the source
specifies stationary reversal, first reconstruct its actual stationary
density and apply the V35 kernel and full Choi test.  If neither relation
is physically identified, neither mathematical reverser is to be renamed
the protected matter grading.  This fork prevents another cycle of changing
the reversal convention merely because one candidate has a nonzero kernel.

\clearpage
\begin{samepage}
\paragraph{Verification.}
The accompanying script checks exact matrix normalization, stationary
fixed spaces, every matrix-unit CP identity, anchor and Kraus-frame
invariance, kernel dimensions and spectra, stationary word reversal,
noncommuting temporal order, the GNS defect, mixed residuals, and physical
opportunity scaling.  It can rerun the unchanged V34 checker.  These
computational certificates supplement the proofs; they are not physical
measurements, proof-assistant formalization, or independent peer review.
\end{samepage}

\addtocontents{toc}{\protect\endgroup}
\renewcommand{\refname}{Additional background for V35}
\begin{thebibliography}{99}
\raggedright
\bibitem[49]{sel35-crooks}
G.~E.~Crooks, \emph{Quantum operation time reversal},
Phys. Rev. A \textbf{77}, 034101 (2008).
\href{https://arxiv.org/abs/0706.3749}{arXiv:0706.3749}.
Stationary reverse maps and trajectory reversal are background, not new
claims of the present source reconstruction.
\bibitem[50]{sel35-transition}
R.~Duvenhage, K.~Oerder, and K.~van den Heuvel,
\emph{Quantum detailed balance via elementary transitions},
Quantum \textbf{9}, 1743 (2025).
\href{https://arxiv.org/abs/2411.02339}{arXiv:2411.02339}.
Background for stationary Choi states, parity, and the distinction between
adjoint/reversing conventions.  No historical source data are imported.
\end{thebibliography}

% END OF SOURCE-SELECTION V35 DEVELOPMENT BLOCK.
% Preserve V1--V35 and append subsequent work before the final terminator.


\clearpage
\hypersetup{pdftitle={Historical Standard-Model Source Selection: V1--V36}}
\begin{center}
{\Large\bfseries Source-selection continuation V36}\\[.4em]
{\large A reconstructed reverse is not a coherent action:\\
transient source support, original-record recovery, and local-law comparison}
\end{center}
\addtocontents{toc}{\protect\begingroup}
\addtocontents{toc}{\protect\renewcommand{\protect\numberline}{\protect\makebox[2.1em][l]}}

\section[Audit of the four companion interfaces]{Stop changing reversal conventions: test the actual source and its action claim}
\label{sel36:context}

V35 constructs a stationary reverse and a unique candidate involution under
explicit source-state and generation hypotheses. It does not identify
that involution with the protected matter grading. The present audit
does not introduce another convention chosen to produce a nonzero kernel.
Instead it asks two prior questions: does the concrete action compiler
satisfy the faithful-stationary entrance at all, and does a recovered
reverse imply preservation of an unknown coherent input? The answers
separate sharply on the sources already constructed.

\paragraph{The supplied regenerative continuation.}
The cumulative \emph{Native Regenerative Stationarity Closure V38}
\cite{sel36-native} contains two useful, distinct interfaces. Its
\src{v12:eq:B-family} and \src{v12:thm:locked-bank} give the actual
one-context six-outcome compiler with four one-way auxiliary maps.
Its \src{v25:lem:polar} and \src{v25:thm:word} give the exact optimal
terminal recovery formula for a qubit, retaining the original outcome
history. These normalization and recovery results are inherited below.
The new V38 calculation, \src{v38:thm:isolation}, excludes every nearby
phase-complete restoration of a missing profile at its specified
two-profile root: the two actual residuals are comparable to the
missing amplitude squared. Its \src{v38:prop:mean} also quantifies
the diverging preparation inverse. This does not populate a historical
matter-grading table or an interior three-profile root.

\paragraph{The supplied Einstein bridge.}
The \emph{Exact Action Einstein Bridge V28} \cite{sel36-einstein},
\src{v28:quadratic}, removes the earlier Poisson cokernel on a
larger mixed-mode stratum and solves the original quadratic initial
tilt. Its \src{v28:cubic-theorem} controls the nonconstant cubic
rows but leaves three constant-shift means, which are nonzero at
the displayed witness. Its differential lift \src{v28:lift} is
not identified with the physical multiplier rate. It would therefore
be incorrect either to import the old axial null tests into that
new stratum or to read its coefficient solution as propagated
source stationarity. Neither inference is used here.

\paragraph{The supplied perturbative continuation.}
The \emph{Perturbative ESM Continuum V23} \cite{sel36-esm} proves
a derivative floor for the particular combined local counterterm and
removes a putative finite cubic-artifact return at its stated one-loop,
fixed-positive-lower-scale orders. The first-background, higher-antifield,
and mixed matter/chiral obligations remain explicit in its final ledger.
These coefficientwise local identities are not positive Choi matrices
or Kraus operations. They supply neither a quantum stationary density
nor the determinant--Clifford incidence of the historical instrument.
The lesson used here is to calculate the actual coefficient before
concluding from a conservative scaling budget.

\paragraph{The supplied Yang--Mills continuation.}
The \emph{Yang--Mills Mass Gap V32} \cite{sel36-ym},
\src{r32:kernels}, compares normalized conditional laws at the
\emph{same exterior}, with explicit local oscillation and total-variation
bounds. It expressly does not compare unconditional marginals or global
spectral gaps. Section~\ref{sel36:local} transfers precisely that local
result to an already identified quantum output bank, without constructing
a new YM--SM coupling.

\paragraph{Status of the additional mathematics.}
Stationary reversal and environment-assisted correction are established,
different notions \cite{sel35-crooks,sel36-gw}. The new work evaluates
their difference on the unchanged V35 source, classifies the stationary
support of the inherited action compiler, gives an exact two-parameter
interpolation and its physical-clock retention limit, and derives a
spectral compatibility test for a proposed coherent odd action.
This last test is conditional on an \emph{all-input coherent-memory}
interpretation. It is not a new prerequisite for the finite algebraic
duality, a classical field reader, or the structural determinant theorem.

All diamond norms below have no factor \(1/2\).
No optimizing recovery is inserted in the source chronology.
``Terminal reader'' means a CPTP map allowed to use every originally
retained classical label and the final quantum output, but not a
bypass copy of the unknown logical input.

\section[The one-context compiler has transient helpers]{Stationary support of the inherited locked action compiler}
\label{sel36:compiler}

Retain one finite action \(H=H^\dagger\), five orthonormal eigenvectors
\(u_0,u_1,\ldots,u_4\), and the original auxiliary operators
\[
 C_j=|u_0\rangle\langle u_j|,\qquad
 R=\sum_{j=1}^4|u_j\rangle\langle u_j|,\qquad P=I-R .
\]
For \(0<\epsilon<1\), put
\[
 q=\epsilon^4,\quad D=(I-qR)^{1/2},\quad
 U=e^{-3i\epsilon H/2},\quad B=UD .
\]
The unchanged six coefficients are
\(K_e=I/\sqrt{18}+\sqrt{2/3}\sum_{\mu=1}^5O_{e\mu}B_\mu\),
where \(B_1=B\), \(B_{j+1}=\epsilon^2C_j\), and \(O\) is the
retained real isometry onto \(\mathbf1^\perp\).
Their already proved normalization gives the forgotten channel
\begin{equation}
 \Phi_\epsilon(X)=\frac13X+
 \frac23\left(BXB^\dagger+q\sum_{j=1}^4C_jXC_j^\dagger\right).
 \label{sel36:eq:compiler}
\end{equation}
This formula is independent of a new Kraus choice: it is the sum of
the six original marked maps.

\begin{theorem}[Exact stationary face and finite-time helper survival]
\label{sel36:thm:stationary-face}
For the source \eqref{sel36:eq:compiler},
\begin{equation}
 \Phi_\epsilon^\dagger(R)=(1-2q/3)R .
 \label{sel36:eq:helper}
\end{equation}
Its stationary densities are exactly
\begin{equation}
 \left\{\sigma\succeq0:\Tr\sigma=1,\quad
 \sigma=P\sigma P,\quad[U|_{P\Hh},\sigma]=0\right\}.
 \label{sel36:eq:stationary-face}
\end{equation}
Its maximal stationary support is \(P\), of rank \(N-4\).
In particular there is no faithful stationary density on the original
\(N\)-dimensional carrier, for any fixed \(\epsilon>0\).

After \(n\) actual acceptances,
\begin{equation}
 \Tr(R\Phi_\epsilon^n(\rho))
    =(1-2\epsilon^4/3)^n\Tr(R\rho).
 \label{sel36:eq:survival}
\end{equation}
On the inherited independent renewal clock, if \(A_T\) is the accepted
count and \(\mathbb E A_T\le T/\epsilon+1\), then
\begin{equation}
 0\le \Tr(R\rho)-\mathbb E\Tr(R\rho_{A_T})
       \le\frac23(T\epsilon^3+\epsilon^4)\Tr(R\rho).
 \label{sel36:eq:physical-survival}
\end{equation}
\end{theorem}
\begin{proof}
The eigenvector choice gives \([U,R]=0\), while \(RC_j=0\).
Hence \(B^\dagger RB=(1-q)R\) and
\(C_j^\dagger RC_j=0\), proving \eqref{sel36:eq:helper}.
Stationarity implies
\(\Tr(R\sigma)=(1-2q/3)\Tr(R\sigma)\), so \(\Tr(R\sigma)=0\).
Positivity then implies \(\sigma=P\sigma P\), including vanishing
off-diagonal corners: \(R\sigma^{1/2}=0\).
On this support all \(C_j\sigma\) vanish and \(D\sigma=\sigma\);
the fixed-point equation is
\(\sigma=\frac13\sigma+\frac23U\sigma U^\dagger\).
This is equivalent to the commutation in
\eqref{sel36:eq:stationary-face}. Conversely it makes a supported
density stationary. The density \(P/(N-4)\) belongs to that set,
so its support is maximal.

Iteration of the adjoint identity proves
\eqref{sel36:eq:survival}, with no assumption on the initial density.
Conditioning on the independent clock keeps this same identity.
The inequality \(1-(1-a)^n\le na\), followed by the inherited
accepted-count bound, proves \eqref{sel36:eq:physical-survival}.
\end{proof}

Thus the V35 stationary-reverse entrance fails on this specific full
compiler carrier, including the V19--V22 benchmark. Deleting \(R\)
would change the retained preparation and history problem; adding a
reset to make the state faithful would change the source.

This failure does not imply a failure of the compiler's finite-time
action approximation. The long-time and small-event limits are different.

\begin{proposition}[Coherent finite-time evolution despite the stationary loss]
\label{sel36:prop:unitary-approx}
For fixed finite \(H\), let
\(h_H=\inf_{c\in\R}\|H-cI\|_\op\) and
\(\mathcal U_t=\operatorname{Ad}_{e^{-itH}}\).
At \(n\) accepted events,
\begin{equation}
 \|\Phi_\epsilon^n-\mathcal U_{n\epsilon}\|_\diamond
 \le \min\{2,\ n(\epsilon^2h_H^2+2\epsilon^4)\}.
 \label{sel36:eq:unitary-approx}
\end{equation}
In particular it tends to zero for \(n=\lfloor T/\epsilon\rfloor\).
Nevertheless the stationary support has rank \(N-4\) at every
nonzero \(\epsilon\).
\end{proposition}
\begin{proof}
Let \(\mathcal Q_\epsilon=\frac13\id+\frac23\mathcal U_{3\epsilon/2}\).
Since \(\|D\|\le1\), \(\|I-D\|=1-\sqrt{1-q}\le q\), and
\(\|\sum_j C_j(\cdot)C_j^\dagger\|_\diamond=\|R\|=1\),
\(\|\Phi_\epsilon-\mathcal Q_\epsilon\|_\diamond\le2q\).
For completeness, the diamond norm of a CP map is the operator norm
of its effect: its Stinespring factorization gives the upper bound and
a positive maximizing input gives equality.

Put \(\mathcal L=-i[H,\cdot]\), so \(\|\mathcal L\|_\diamond\le2h_H\).
Expand the isometric group \(e^{t\mathcal L}\) about \(t=\epsilon\)
with the integral second-order remainder. At times \(0\) and
\(3\epsilon/2\), with weights \(1/3,2/3\), the linear terms cancel,
and their weighted squared displacement is \(\epsilon^2/2\).
The remainder is bounded by
\(\epsilon^2\|\mathcal L^2\|_\diamond/4\le\epsilon^2h_H^2\).
Telescoping powers of trace-preserving CP maps proves
\eqref{sel36:eq:unitary-approx}.
\end{proof}

The estimate is at accepted-event time; it does not replace the original
random clock by a deterministic one. Equation
\eqref{sel36:eq:physical-survival} is the actual-clock helper statement.
The companion's separate renewal estimates govern its field-reader
limit. Also, this stationary-face theorem concerns \emph{one fixed
context}. The native spatial construction retains a label
\((\text{context},e)\); different local contexts can refill one
another's helper sectors. No common Lyapunov projection \(R\) for
that entire mixed instrument has been proved here.

\section[What V35's reverser does not recover]{Exact optimal recovery of the unchanged nonunital source}
\label{sel36:recovery}

Retain V35's original two coefficients and stationary density:
\begin{equation}
 k_\pm=\frac1{\sqrt{10}}
 \begin{pmatrix}2&\pm2\\\mp1&1\end{pmatrix},
 \qquad \sigma=\diag(4/5,1/5).
 \label{sel36:eq:old-bank}
\end{equation}
V35 proves their unique stationary-reversal action
\(\operatorname{Ad}_Z\), \(Z=\diag(1,-1)\).
For all length-\(n\) original words, let
\[
 \mathcal T_n(\rho)=\bigoplus_{|w|=n}k_w\rho k_w^\dagger,\qquad
 e_n=\inf_{\mathcal R\ {\rm CPTP}}
       \|\mathcal R\mathcal T_n-\id\|_\diamond .
\]
The infimum allows every terminal reader of the actual flags and
terminal system, not just the stationary reverse.

\begin{theorem}[Unique stationary reversal with irreducible record-retention loss]
\label{sel36:thm:old-retention}
For \eqref{sel36:eq:old-bank},
\begin{equation}
 \boxed{e_n=1-(4/5)^n,\qquad
 \sup_{\mathcal R}F_{\rm e}(\mathcal R\mathcal T_n)
                   =\frac{1+(4/5)^n}{2}.}
 \label{sel36:eq:old-retention}
\end{equation}
The first equality is unchanged if the target is any prescribed
logical unitary channel. No reader can undo the information loss
by knowing all original labels.
\end{theorem}
\begin{proof}
The imported polar-recovery theorem
\cite{sel36-native}, \src{v25:lem:polar}, states for an efficient
flagged qubit-input instrument that
\(e=1-\sum_w\sqrt{\det(k_w^\dagger k_w)}\), with the corresponding
entanglement-fidelity maximum. Here \(\det k_+=\det k_-=2/5>0\).
Multiplicativity of the determinant gives
\[
 \sum_{|w|=n}\sqrt{\det(k_w^\dagger k_w)}
       =\sum_{|w|=n}|\det k_w|=(2\cdot2/5)^n .
\]
This proves both assertions. A reader composed with the desired
unitary, or its inverse, proves the unitary-target equality.
The imported theorem already allows all terminal CPTP readers;
there is no limitation to a convenient Kraus inverse.
\end{proof}

This does not contradict the efficient same-label coherence theorem
of V27. Rank-one Choi evolution preserves purity on a conditioned
pure input but can still disclose information about that input in
its outcome probabilities. A coherent coefficient, a recoverable
unknown input, and a stationary reverse are different objects.

\section[The original renewal clock fixes the retention rate]{Exact retention on the actual two-phase timing law}
\label{sel36:clock}

Retain \cite{sel36-native}, \src{v10:eq:full-table}:
the phase transition is \(H\to(H,P)\) with probabilities
\((1/5,4/5)\), and \(P\to(H,P)\) with probabilities \((2/3,1/3)\).
A \(P\to H\) completion is accepted with probability \(\vartheta\).
The physical tick is
\(d=4\vartheta\epsilon/11\), so \(\epsilon\) is the mean time
between accepted operations. Nonaccepted ticks act as the identity.
The clock is independent of the unknown quantum input.

For any efficient qubit instrument on this clock put
\(s=\sum_a|\det k_a|\in[0,1]\), and keep all its word and clock
records through tick \(j\). Write \(e_j^{\rm clk}\) for the optimal
terminal recovery error. Let \(\zeta\) be the original initial phase
distribution, considered as a row vector.

\begin{theorem}[Exact clock transfer and a uniform physical rate]
\label{sel36:thm:clock}
At the actual tick time \(t_j=jd\),
\begin{equation}
 \boxed{e_j^{\rm clk}=1-\zeta B_\vartheta(s)^j\mathbf1,\qquad
 B_\vartheta(s)=
 \begin{pmatrix}
 1/5&4/5\\ (2/3)(1-\vartheta+\vartheta s)&1/3
 \end{pmatrix}.}
 \label{sel36:eq:clock}
\end{equation}
For fixed \(s<1\), put
\[
 \lambda_\vartheta(s)=
 \frac{4+\sqrt{121-120\vartheta(1-s)}}{15},
 \qquad C_s=\frac{12}{1+\sqrt{1+120s}} .
\]
Then, for every initial phase and \(0<\vartheta\le1\),
\begin{equation}
 1-e_j^{\rm clk}\le
 C_s\exp\{- (1-s)t_j/\epsilon\}.
 \label{sel36:eq:clock-rate}
\end{equation}
For the unchanged V35 bank \(s=4/5\), this is
\(C_s=12/(1+\sqrt{97})\) and exponent \(-t_j/(5\epsilon)\).
\end{theorem}
\begin{proof}
A complete clock path of probability \(p_c\) and \(A_j=n\)
acceptances has composite coefficients \(\sqrt{p_c}\,k_w\).
For qubit input their determinant magnitudes are
\(p_c|\det k_w|\). The same imported terminal-recovery identity,
now applied to all classical clock and word flags, gives
\(e_j^{\rm clk}=1-\mathbb E s^{A_j}\).
The original tick table weights precisely one accepted
\(P\to H\) transition by \(s\), proving the matrix formula.

The positive eigenvector for \(\lambda_\vartheta(s)\) is
\(v=(1,(1+\sqrt{121-120\vartheta(1-s)})/12)^T\).
Its entries lie between \(1/C_s\) and \(1\), so positivity gives
\(B_\vartheta(s)^j\mathbf1\le C_s\lambda_\vartheta(s)^j\mathbf1\).
Rationalizing the root gives
\[
 1-\lambda_\vartheta(s)=
 \frac{8\vartheta(1-s)}
      {11+\sqrt{121-120\vartheta(1-s)}}
 \ge\frac{4\vartheta(1-s)}{11}.
\]
Use \(1-x\le e^{-x}\) and the original value of \(d\).
\end{proof}

This is an application of the inherited count transform, not a new
clock. In particular rare acceptance does not suppress the fixed
per-acceptance loss when the original mean accepted time tends to zero.

\section[A family separates parity, information loss, and the generator]{An exact nonunital interpolation and its continuous-time limit}
\label{sel36:family}

For \(0<a,b<1\) define the declared comparison family
\begin{equation}
 k_\pm(a,b)=\frac1{\sqrt2}
 \begin{pmatrix}\sqrt{1-a}&\pm\sqrt b\\
                 \mp\sqrt a&\sqrt{1-b}\end{pmatrix},
 \qquad
 s(a,b)=\sqrt{(1-a)(1-b)}+\sqrt{ab}.
 \label{sel36:eq:family}
\end{equation}
It specializes exactly to \eqref{sel36:eq:old-bank} at
\((a,b)=(1/5,4/5)\). Varying \(a,b\) changes the comparison source;
it is not an allowed reinterpretation of its old recorded maps.

\begin{theorem}[Exact state, reversal, identification margin, and recovery]
\label{sel36:thm:family}
If \(a\ne b\), this normalized instrument has full coefficient
algebra \(M_2\), unique faithful stationary density
\[
 \sigma_{a,b}=\diag\left(\frac b{a+b},\frac a{a+b}\right),
\]
and unique stationary-reversal action \(\operatorname{Ad}_Z\).
At \(n\) acceptances its optimal terminal recovery error is
\begin{equation}
 \boxed{e_n(a,b)=1-s(a,b)^n.}
 \label{sel36:eq:family-retention}
\end{equation}
Let
\[
 c_\sigma=\frac{b\sqrt{1-a}+a\sqrt{1-b}}{a+b}.
\]
On the source-determined space
\(\{T:[T,\sigma_{a,b}]=0\}=\Span\{I,Z\}\), V35's
stationary-anchor form has the exact matrix
\begin{equation}
 \boxed{\mathscr L_\sigma
       =\diag(2c_\sigma^2\sqrt{ab},0)}
 \label{sel36:eq:family-gap}
\end{equation}
in the orthonormal basis \((I,Z)/\sqrt2\).
\end{theorem}
\begin{proof}
Summing \(k_\pm^\dagger k_\pm\) cancels their off-diagonal entries
and gives \(I\). On populations, the channel has transition matrix
\(\left(\begin{smallmatrix}1-a&b\\a&1-b\end{smallmatrix}\right)\).
Its fixed density is as displayed. On Pauli coherences its multipliers
are \(c-d,c+d\), where \(c=\sqrt{(1-a)(1-b)}\), \(d=\sqrt{ab}\);
the remaining population eigenvalue is \(1-a-b\).
Writing \(a=\sin^2\theta,b=\sin^2\phi\), with
\(0<\theta,\phi<\pi/2\), gives
\(c+d=\cos(\theta-\phi)<1\) when \(a\ne b\) and
\(|c-d|=|\cos(\theta+\phi)|<1\).
Thus the stationary density is unique.

The sum of the coefficients is diagonal with distinct eigenvalues,
and their difference has both off-diagonal entries nonzero.
Their spectral projections and corners generate all matrix units.
The stationary reverse coefficients are \(k_\mp=Zk_\pm Z\).
Their weighted traces equal \(c_\sigma/\sqrt2>0\), so their
stationary anchors generate the same full algebra. V35's theorem
therefore gives uniqueness.

More explicitly the whitened anchors are
\[
 \sigma^{-1/4}A_\pm^\sigma\sigma^{1/4}
 =\frac{c_\sigma}{2}
 \begin{pmatrix}
 \sqrt{1-a}&\pm(ab)^{1/4}\\
 \mp(ab)^{1/4}&\sqrt{1-b}
 \end{pmatrix}.
\]
Their imaginary Hermitian parts are opposite multiples of \(Y\).
On \(\Span\{I,Z\}\) the commutators with the diagonal parts vanish.
The two anticommutators have squared norm \(2c_\sigma^2\sqrt{ab}\)
on \(I/\sqrt2\), and zero on \(Z/\sqrt2\).
This proves \eqref{sel36:eq:family-gap}.
Finally \(\det k_\pm=s(a,b)/2>0\); the determinant argument of
Theorem~\ref{sel36:thm:old-retention} proves
\eqref{sel36:eq:family-retention}.
\end{proof}

At \(a=b\) every effect is \(I/2\), the original branches are
scaled unitaries, and recovery is perfect. The coefficient algebra
is then abelian and the stationary state is not unique; this boundary
does not meet the unique-grading hypothesis.

\begin{theorem}[A unique grading can persist in a dissipative physical limit]
\label{sel36:thm:weak-limit}
Fix distinct \(u,v>0\), take \(a=u\epsilon,b=v\epsilon\), and retain
the original mean-event calibration. Put
\[
 c_{u,v}=\tfrac12(\sqrt u-\sqrt v)^2>0,\qquad
 L=\begin{pmatrix}0&\sqrt v\\-\sqrt u&0\end{pmatrix}.
\]
For sufficiently small \(\epsilon\),
\begin{align}
 s(u\epsilon,v\epsilon)&=1-c_{u,v}\epsilon+O(\epsilon^2),\notag\\
 \gamma_\sigma(\epsilon)&=2\sqrt{uv}\epsilon+O(\epsilon^2),\notag\\
 \Phi_\epsilon&=\id+\epsilon\mathcal D_L+O_\diamond(\epsilon^2),
 \quad
 \mathcal D_L(X)=LXL^\dagger-\tfrac12\{L^\dagger L,X\}.
 \label{sel36:eq:weak-limit}
\end{align}
For \(n=\lfloor T/\epsilon\rfloor\),
\(\Phi_\epsilon^n\to e^{T\mathcal D_L}\) and
\(e_n\to1-e^{-c_{u,v}T}\).
On the actual two-phase clock, at ticks tending to \(T>0\),
\begin{equation}
 \boxed{e_j^{\rm clk}\longrightarrow1-e^{-c_{u,v}T},}
 \label{sel36:eq:clock-weak-limit}
\end{equation}
also for \(\vartheta=\vartheta_\epsilon\in(0,1]\).
The recovered grading is \(Z\) at every \(\epsilon\), but the
first-order generator here is dissipative, not a nonzero
Hamiltonian commutator.
\end{theorem}
\begin{proof}
Taylor expansion of the two positive square roots gives the first
two expansions. Write \(D_\epsilon=\diag(\sqrt{1-u\epsilon},
\sqrt{1-v\epsilon})\). Cancellation between the two original labels
gives exactly
\(\Phi_\epsilon(X)=D_\epsilon XD_\epsilon+\epsilon LXL^\dagger\).
Since \(D_\epsilon=I-\epsilon L^\dagger L/2+O(\epsilon^2)\),
the channel expansion follows in finite dimension.
The semigroup limit follows by telescoping against
\(e^{\epsilon\mathcal D_L}\); the latter is CPTP, or equivalently
is the norm limit of these normalized finite channels.
The determinant formula gives the accepted-time recovery limit.

For the actual clock, the eigenvector in the proof of
Theorem~\ref{sel36:thm:clock} differs from \(\mathbf1\) by
\(O(\vartheta_\epsilon\epsilon)\). Uniform Taylor expansion of its
positive eigenvalue gives
\[
 \lambda_{\vartheta_\epsilon}(s_\epsilon)
 =1-\frac{4\vartheta_\epsilon}{11}c_{u,v}\epsilon
      +O(\vartheta_\epsilon\epsilon^2),
\]
uniformly for \(0<\vartheta_\epsilon\le1\). Positivity sandwiches
\(B^j\mathbf1\) between scalar multiples of \(\lambda^jv\) whose
ratio tends to one. Dividing \(\log\lambda\) by
\(d=4\vartheta_\epsilon\epsilon/11\) gives \(-c_{u,v}+O(\epsilon)\).
This proves \eqref{sel36:eq:clock-weak-limit}, including arbitrary
initial phase.

The expansion has the displayed jump dissipator. It is not zero and
is not a Hamiltonian commutator: for example it changes the diagonal
population \(E_{00}\) at rate \(u(E_{11}-E_{00})\).
Also \(ZLZ=-L\). Thus an odd jump and a recovered parity are present,
but neither is an independently identified odd Hamiltonian.
\end{proof}

\section[The stationary spectrum constrains a coherent odd action]{A same-source spectral obstruction to an odd Hamiltonian interpretation}
\label{sel36:action}

Let a proposed physical family have stationary densities \(\sigma_\epsilon\)
and an independently identified self-adjoint generator \(A_\epsilon\).
A stationary reverse relation alone does not imply the unitary
comparison in the following theorem.

\begin{theorem}[Stationary-state cost of a coherent action approximation]
\label{sel36:thm:action}
Suppose \(\Phi_\epsilon(\sigma_\epsilon)=\sigma_\epsilon\),
\(\Tr\sigma_\epsilon=1\), and put
\[
 \delta_\epsilon=
 \|\Phi_\epsilon-\operatorname{Ad}_{e^{-i\epsilon A_\epsilon}}\|_\diamond,
 \qquad
 h_\epsilon=\inf_{c\in\R}\|A_\epsilon-cI\|_\op.
\]
Then
\begin{equation}
 \|[A_\epsilon,\sigma_\epsilon]\|_1
 \le\frac{\delta_\epsilon}{\epsilon}+2\epsilon h_\epsilon^2.
 \label{sel36:eq:stationary-action}
\end{equation}
Consequently, bounded \(A_\epsilon\to A\),
\(\sigma_\epsilon\to\sigma\), and
\(\delta_\epsilon=o(\epsilon)\) imply \([A,\sigma]=0\).

If \(Z=Z^\dagger=Z^{-1}\) commutes with a fixed \(\sigma\), the
self-adjoint operators satisfying
\(ZAZ=-A\) and \([A,\sigma]=0\) have real dimension
\begin{equation}
 \boxed{2\sum_\lambda
 \rank(P_\lambda^\sigma P_+^Z)\,
 \rank(P_\lambda^\sigma P_-^Z).}
 \label{sel36:eq:odd-dim}
\end{equation}
In particular an odd coherent action requires a stationary
eigenvalue present in both grading sectors.
\end{theorem}
\begin{proof}
Stationarity bounds
\(\|e^{-i\epsilon A_\epsilon}\sigma_\epsilon
 e^{i\epsilon A_\epsilon}-\sigma_\epsilon\|_1
 \le\delta_\epsilon\).
The integral Taylor remainder is at most
\(\epsilon^2\|[A_\epsilon,[A_\epsilon,\sigma_\epsilon]]\|_1/2
 \le2\epsilon^2h_\epsilon^2\).
The triangle inequality proves \eqref{sel36:eq:stationary-action}
and passage to the limit gives the commutator identity.

Diagonalize \(\sigma\) and \(Z\) simultaneously. Commutation with
\(\sigma\) allows matrix entries only within a common eigenvalue.
Oddness allows only entries between opposite signs. On the
\(\lambda\)-eigenspace a self-adjoint odd matrix is
\(\left(\begin{smallmatrix}0&B_\lambda^\dagger\\
B_\lambda&0\end{smallmatrix}\right)\), with arbitrary complex
\(B_\lambda\). Counting its real parameters proves
\eqref{sel36:eq:odd-dim}.
\end{proof}

If the spectra of \(\sigma\) on opposite grading sectors are disjoint,
let \(\Delta\) be their minimum separation. For
\(A_{\rm odd}=(A-ZAZ)/2\), entrywise calculation gives
\[
 \Delta\|A_{\rm odd}\|_{\HS}
 \le\|[A,\sigma]\|_{\HS}\le\|[A,\sigma]\|_1 .
\]
Thus \eqref{sel36:eq:stationary-action} gives the explicit bound
\begin{equation}
 \|A_{\rm odd}\|_{\HS}
 \le\Delta^{-1}
       \left(\delta_\epsilon/\epsilon+2\epsilon h_\epsilon^2\right)
 \label{sel36:eq:odd-bound}
\end{equation}
when the exact stationary state and grading in this expression are fixed.

For V35's \(\sigma=\diag(4/5,1/5)\), \(Z=\diag(1,-1)\), the
dimension in \eqref{sel36:eq:odd-dim} is zero and
\(\Delta=3/5\). More explicitly, for the independently proposed
\(A=\Omega X\),
\begin{equation}
 \boxed{
 \|\Phi-\operatorname{Ad}_{e^{-i\epsilon\Omega X}}\|_\diamond
 \ge\frac65|\sin(\epsilon\Omega)|
 }
 \label{sel36:eq:qubit-action}
\end{equation}
for \emph{every} channel \(\Phi\) fixing that \(\sigma\).
Indeed the Bloch vector of \(\sigma\) has length \(3/5\);
rotation by angle \(2\epsilon\Omega\) changes it by
\(2(3/5)|\sin(\epsilon\Omega)|\). This is an exact stationary-input
lower bound, not a formal expansion.

These statements concern a specified all-input unitary action on the
same quantum memory. They do not preclude the manuscript's algebraic
finite Dirac operator, a different occupation state, an action on a
larger correlated carrier, or a classical decoded field evolution.
They prevent an unproved identification of those different targets.

\section[What the Yang--Mills local comparison actually transports]{Conditional source transfer without changing the exterior law}
\label{sel36:local}

The supplied YM V32 gives, at the same exterior \(\xi\) and retained
region \(D\),
\begin{equation}
 e^{-\omega_D}\le
 \frac{dK_D^\xi}{d\overline K_D^\xi}\le e^{\omega_D},
 \quad
 d_{\rm TV}(K_D^\xi,\overline K_D^\xi)
 \le\min\{1,\omega_D/4\},
 \quad
 \omega_D=\frac{2Cr^2n_D}{\gamma},\quad n_D\le120|D|.
 \label{sel36:eq:YM-input}
\end{equation}
Its hybrid keeps the full original Haar space, the coherent scalar
normalization, all rough factors, and the original exterior condition.

\begin{proposition}[Local positive-source and quantum-output comparison]
\label{sel36:prop:local}
Let \(V(\eta):E\to F\) be any fixed bounded measurable finite-dimensional
source response on that same conditional region and exterior. Then
\[
 G=\int V^\dagger V\,dK_D^\xi,\qquad
 \overline G=\int V^\dagger V\,d\overline K_D^\xi
\]
satisfy
\begin{equation}
 e^{-\omega_D}\overline G\preceq G
                      \preceq e^{\omega_D}\overline G,
 \qquad \Ker G=\Ker\overline G .
 \label{sel36:eq:local-Gram}
\end{equation}
If an already identified family of CPTP output maps
\(\{\mathcal V_\eta\}\) is the same for both laws, its two conditional
averages obey
\begin{equation}
 \left\|\int\mathcal V_\eta\,dK_D^\xi
            -\int\mathcal V_\eta\,d\overline K_D^\xi
 \right\|_\diamond
 \le\min\{2,\omega_D/2\}.
 \label{sel36:eq:local-diamond}
\end{equation}
The same estimate holds with the sampled classical variable retained.
For any two such flagged output channels \(\mathcal T,\overline{\mathcal T}\),
their optimal terminal unitary-recovery errors satisfy
\[
 |e(\mathcal T)-e(\overline{\mathcal T})|
                   \le\|\mathcal T-\overline{\mathcal T}\|_\diamond.
\]
\end{proposition}
\begin{proof}
Apply the density bounds to
\(\|V(\eta)x\|^2\ge0\) for every \(x\).
This proves both Loewner inequalities and hence equality of kernels.
For the quantum statement, integration against the signed measure
\(K_D^\xi-\overline K_D^\xi\), using
\(\|\mathcal V_\eta\|_\diamond=1\), bounds the difference by
\(2d_{\rm TV}\). The same estimate in the classical direct-integral
output uses the norm of the signed measure, so retaining the sample
does not worsen it. Finite recorded partitions suffice when the
grammar has a finite alphabet.

A CPTP terminal reader is a diamond contraction. For each fixed
reader the two distances to a target unitary differ by no more
than the source-channel distance. Taking infima, first in one
direction and then the other, proves the last assertion.
\end{proof}

For \(n\) successive uses from a common initial preparation, with the
same actually implemented conditional update policy, uniform local error \(\delta\) gives full-history
error at most \(\min\{2,n\delta\}\). This follows by telescoping
the adaptive channels including their memory and flags, not by assuming
independent quantum inputs. On the independent clock it gives
\(\delta\,\mathbb E A_T\le\delta(T/\epsilon+1)\).
At the supplied thermal window this is
\(O(|D|T\log\gamma/(\epsilon\gamma^2))\), not automatically a
vanishing bound. An actual joint scale must pay for the accepted count.

Two limitations are necessary. First, the response \(V\) or quantum
bank \(\mathcal V_\eta\) has to be identified independently; the
conditional field law does not manufacture an SM operator.
Second, \eqref{sel36:eq:local-Gram} preserves kernels of positive
source forms, not every cancellation in an averaged quantum map.
For example mix the two fixed qubit reset maps to the \(X\)-eigenstates.
Equal positive weights give output \(I/2\), while weights
\((1+t)/2,(1-t)/2\), \(0<|t|<1\), give \((I+tX)/2\).
The two Choi supports are both full and their density ratios are bounded,
but the first channel is covariant under all \(U(2)\), whereas the second
has only the stabilizer of its output density. Equivalence of sampling
measures therefore cannot replace a covariance or reverse-instrument
test.

Finally, integrating the local kernels against the \emph{same} exterior
distribution preserves \eqref{sel36:eq:local-diamond}. The actual and
hybrid global laws generally have different exterior distributions.
No estimate for that missing difference, a global susceptibility,
or a continuum gap is inferred.

\section[What has been proved and what should be tested next]{Proof-status ledger and the noncircular source target}
\label{sel36:status}

\paragraph{Proved in V36.}
The inherited one-context action compiler has an exactly identified
transient helper sector and no faithful stationary state on its full
carrier, despite its coherent fixed-action finite-time approximation.
The unchanged V35 source has exact all-history optimal recovery
error \(1-(4/5)^n\), with the original-clock formula and rate above.
The declared two-parameter family determines its stationary state,
unique parity, identification margin, retention exponent, and
dissipative weak limit. The stationary spectrum supplies a sharp
cross-grading support condition for a proposed coherent action.
The same-exterior YM comparison transports positive source forms
and conditional quantum outputs at its original, local scope.

\paragraph{Inherited rather than redone.}
The six-coefficient normalization and finite-observable action
approximation, the rectangular-code polar-recovery theorem, and
the two-phase timing law come from the regenerative notes.
The V35 reverse and inverse-grading theorems are unchanged.
The YM logarithmic remainder, support counts, and conditional density
bounds are imported, not reproved as lattice estimates.
The new exact checks test the algebraic transfers and source examples,
not the companion numerical compilers or their external certificates.

\paragraph{What is not established.}
No classical profile root, propagated Einstein tangency, completed
ESM background identity, historical matter-grading label, or
determinant--Clifford cylinder has been populated. The original
local multi-context action instrument has not been assigned the
single-context stationary face. No optimizing terminal decoder or
hybrid kernel is silently executed. None of the finite identification
margins is interpreted as a Yang--Mills mass gap.

\paragraph{Next source-level target.}
The outstanding issue is not another mathematical reverse convention.
For an actual protected grading, determine whether its physical
interpretation is a source covariance, a jump parity, or the odd
component of an independently identified action. Only for an asserted
all-input coherent action should the recovery and stationary-spectrum
tests above be required. For the structural duality alone, continue
with the same-carrier directed typing and the V29 search on original
labels; do not impose thermodynamic faithfulness as a new universal
axiom. The three unresolved classical cubic means, the phase-isolated
regenerative boundary, and the perturbative descendant equations
remain in their separate source programmes.

\paragraph{Verification scope.}
The executable checks accompanying this append use exact rational and
algebraic matrices for normalization, stationary support, determinant
products, clock recurrences, parity reconstruction, source spectra,
weak-generator coefficients, and conditional quantum comparisons.
Analytic limit and norm proofs are the arguments above, not an
inference from sampled parameter values. Prior companion external
verification packages were not supplied with the four TeX files and
are not claimed to have been rerun.

\addtocontents{toc}{\protect\endgroup}
\renewcommand{\refname}{Additional sources and background for V36}
\begin{thebibliography}{99}
\raggedright
\bibitem[51]{sel36-native}
Renewal Geometry research programme,
\emph{Renewal Native Regenerative Stationarity Closure Notes V38},
supplied source \src{Renewal_Native_Regenerative_Stationarity_Closure_Notes_V38(2).tex}.
Relevant labels: \src{v12:thm:locked-bank}, \src{v20:prop:multiplicity},
\src{v25:lem:polar}, \src{v25:thm:word}, \src{v10:eq:full-table},
\src{v38:thm:isolation}, \src{v38:prop:mean}.
\bibitem[52]{sel36-einstein}
Renewal Geometry research programme,
\emph{Renewal Exact Action Einstein Bridge Research Notes V28},
supplied source \src{Renewal_Exact_Action_Einstein_Bridge_Research_Notes_V28(2).tex}.
Relevant labels: \src{v28:quadratic}, \src{v28:cubic-theorem},
\src{v28:lift}, \src{v28:status}.
\bibitem[53]{sel36-esm}
Renewal Geometry research programme,
\emph{Renewal Perturbative ESM Continuum Research Notes V23},
supplied source \src{Renewal_Perturbative_ESM_Continuum_Research_Notes_V23(3).tex}.
See the V23 derivative-floor development and closing ledgers,
\src{sec:v23-ledgers}.
\bibitem[54]{sel36-ym}
Renewal Geometry research programme,
\emph{Renewal Yang--Mills Mass Gap Research Notes V32},
supplied source \src{Renewal_Yang_Mills_Mass_Gap_Research_Notes_V32(2).tex}.
Relevant labels: \src{r32:derivatives}, \src{r32:native},
\src{r32:kernels}, \src{r32:thermal-rates}.
\bibitem[55]{sel36-gw}
M.~Gregoratti and R.~F.~Werner,
\emph{On quantum error-correction by classical feedback in discrete time},
J. Math. Phys. \textbf{45}, 2600--2612 (2004).
\href{https://arxiv.org/abs/quant-ph/0403092}{arXiv:quant-ph/0403092}.
Environment-assisted optimal recovery is background; its qubit
diamond-norm form is imported from the supplied regenerative theorem.
\end{thebibliography}

% END OF SOURCE-SELECTION V36 DEVELOPMENT BLOCK.
% Preserve V1--V36 and append subsequent work before the final terminator.


\clearpage
\hypersetup{pdftitle={Historical Standard-Model Source Selection: V1--V37}}
\begin{center}
{\Large\bfseries Source-selection continuation V37}\\[.4em]
{\large Stationary-reversal circulation:\\
exact Hamiltonian recognition and the asymmetric-noise obstruction}
\end{center}
\addtocontents{toc}{\protect\begingroup}
\addtocontents{toc}{\protect\renewcommand{\protect\numberline}{\protect\makebox[2.1em][l]}}

\section[The existing Hamiltonian projection and the missing comparison]{Use the reconstructed reverse; do not choose another convention}
\label{sel37:context}

V36 distinguishes a source-compatible reversal, an odd jump, and coherent
motion of an unknown input. Its two unchanged examples leave a concrete
question: which part of a source's asymmetry under its already fixed
stationary reverse is actually a Hamiltonian commutator? This continuation
answers that question without choosing another reverse, replacing an
instrument, or assuming that its full dynamics are unitary.

The Grand Tensor already defines the orthogonal projection of a
Hermiticity-preserving map onto Hamiltonian derivations and its squared
residual; see \cite{sel37-grand},
\src{def:SMST-Hamiltonian-tangent} and
\src{thm:SMST-Hamiltonian-tangent-certificate}. The duality development
already gives the conditional Choi cone for a quantum Markov generator;
see \cite{Rigidity}, \src{v12:lem:generator-cone}. Those existence and
projection theorems are inherited. We give an explicit conditional-Choi
formula for their residual on a forward--reverse difference, prove its
stationary and grading consequences, and evaluate the original V34 and
V35 sources. A qutrit comparison establishes the exact obstruction and
its sharp diamond distance. Canonical Hamiltonians and quantum detailed
balance are established subjects \cite{sel37-hs,sel37-fu}; no claim of
priority is made for their general frameworks.

\paragraph{Source and normalization.}
All maps in this block act on one already identified matrix memory
$M_n$. The matrix trace is ordinary, not divided by $n$. For a
Hermiticity-preserving map $\mathcal K$ we use the output-first convention
\begin{equation}
 J(\mathcal K)=\sum_{i,j}\mathcal K(E_{ij})\otimes E_{ij},
 \quad \iota=|I_n\rangle\!\rangle,
 \quad P_0=\frac{\iota\iota^*}{n},\quad Q_0=I-P_0.
 \label{sel37:eq:conventions}
\end{equation}
Thus $J(\operatorname{Ad}_k)=|k\rangle\!\rangle\langle\!\langle k|$,
and $\operatorname{Tr}_{\rm out}J(\mathcal K)=0$ exactly when
$\mathcal K$ annihilates the trace. Choi reshuffling preserves the
Hilbert--Schmidt norm of superoperators. A forward--reverse difference is
not a CP map; its signed conditional Choi block must not be interpreted
as a positive event or as a new physical outcome.

\section[An explicit formula for the inherited Hamiltonian residual]{Conditional Choi coordinates of a trace-annihilating map}
\label{sel37:normal}

For a Hermiticity-preserving, trace-annihilating map $\mathcal K$ define
\begin{align}
 C_{\mathcal K}&=Q_0J(\mathcal K)Q_0,
 & B_{\mathcal K}&=\operatorname{unvec}(J(\mathcal K)\iota),\notag\\
 H_{\mathcal K}&=\frac{i}{2n}(B_{\mathcal K}-B_{\mathcal K}^*),
 & M_C&=(\operatorname{Tr}_{\rm out}C_{\mathcal K})^{\mathsf T}.
 \label{sel37:eq:normal-data}
\end{align}
Let $\Psi_C$ be the unique Hermiticity-preserving map with Choi matrix
$C=C_{\mathcal K}$, and set
\begin{equation}
 \mathcal D_C(X)=\Psi_C(X)-\tfrac12(M_CX+XM_C).
 \label{sel37:eq:signed-dissipator}
\end{equation}
Here $\Psi_C^*(I)=M_C$. Positivity of $C$ is not assumed.

\begin{proposition}[Explicit projection and exact orthogonal residual]
\label{sel37:prop:normal}
The inherited Hamiltonian projection has the closed form
\eqref{sel37:eq:normal-data}. More precisely, $H_{\mathcal K}$ is Hermitian
and traceless, and
\begin{equation}
 \boxed{\mathcal K=-i[H_{\mathcal K},\,\cdot\,]+\mathcal D_C.}
 \label{sel37:eq:normal-form}
\end{equation}
The two terms are orthogonal in the real Hilbert--Schmidt superoperator
metric. Consequently its exact squared distance from all Hamiltonian
commutators is
\begin{equation}
 \boxed{
 \Delta_{\rm Ham}(\mathcal K)
 =\|C\|_{\HS}^2+\frac n2\|M_C\|_{\HS}^2
                      +\frac12(\operatorname{Tr}M_C)^2.}
 \label{sel37:eq:distance}
\end{equation}
In particular,
\begin{equation}
 \boxed{\mathcal K=-i[H,\,\cdot\,]\text{ for some }H=H^*
             \quad\Longleftrightarrow\quad C_{\mathcal K}=0.}
 \label{sel37:eq:derivation-criterion}
\end{equation}
\end{proposition}
\begin{proof}
$C=C^*$ and $C\iota=0$. Its associated $\Psi_C$ preserves Hermiticity,
and $M_C=M_C^*$, so \eqref{sel37:eq:signed-dissipator} preserves
Hermiticity and annihilates the trace. Its Choi matrix is
\[
 J(\mathcal D_C)=C-\tfrac12\bigl(
       |M_C\rangle\!\rangle\iota^*+
       \iota\langle\!\langle M_C|\bigr).
\]
The difference $J(\mathcal K-\mathcal D_C)$ has zero $Q_0$--$Q_0$ corner.
Every such Hermitian matrix has the form
$|A\rangle\!\rangle\iota^*+\iota\langle\!\langle A|$.
The corresponding map is $X\mapsto AX+XA^*$; trace annihilation forces
$A+A^*=0$. It is therefore a Hamiltonian commutator.

For a traceless Hermitian $H$,
\[
 J(-i[H,\cdot])=-i|H\rangle\!\rangle\iota^*
                       +i\iota\langle\!\langle H|.
\]
Its product with $\iota$ is $-in|H\rangle\!\rangle$. On the other hand,
$\operatorname{unvec}(J(\mathcal D_C)\iota)
=-(nM_C+(\operatorname{Tr}M_C)I)/2$ is Hermitian. Taking the
anti-Hermitian part gives exactly \eqref{sel37:eq:normal-data}.
The trace of $B_{\mathcal K}$ is the real number
$\iota^*J(\mathcal K)\iota$, so $H_{\mathcal K}$ is traceless.

$C$ is orthogonal to every matrix with an identity row or column.
The real inner product of the displayed Hamiltonian Choi matrix with
$|M\rangle\!\rangle\iota^*+\iota\langle\!\langle M|$ is zero for
Hermitian $H,M$. Thus $\mathcal D_C$ is orthogonal to every Hamiltonian
commutator, identifying the old orthogonal projection. Finally,
\[
 \bigl\||M\rangle\!\rangle\iota^*+
       \iota\langle\!\langle M|\bigr\|_{\HS}^2
       =2n\|M\|_{\HS}^2+2(\operatorname{Tr}M)^2.
\]
Pythagoras proves \eqref{sel37:eq:distance}. Its right side vanishes
exactly when $C=0$.
\end{proof}

\begin{corollary}[A norm obstruction that does not use a fitted action]
\label{sel37:cor:norm}
For the same map,
\begin{equation}
 \frac{\|C_{\mathcal K}\|_1}{n}
 \le\inf_{H=H^*}\|\mathcal K+i[H,\cdot]\|_\diamond
 \le\|\mathcal D_C\|_\diamond\le2\|C_{\mathcal K}\|_1.
 \label{sel37:eq:diamond}
\end{equation}
\end{corollary}
\begin{proof}
The normalized maximally entangled input gives
$\|\mathcal T\|_\diamond\ge\|J(\mathcal T)\|_1/n$.
Compressing its Choi matrix by $Q_0$ removes any Hamiltonian term and is
trace-norm contractive. For the upper bound diagonalize the Hermitian
$C=\sum_j c_j|L_j\rangle\!\rangle\langle\!\langle L_j|$ with
$\|L_j\|_{\HS}=1$. Then
$\mathcal D_C=\sum_jc_j\mathcal D[L_j]$ and
$\|\mathcal D[L_j]\|_\diamond\le2\|L_j\|_{\op}^2\le2$.
Summing gives the claim. This bounds the unchanged residual; deleting it
would be a different source.
\end{proof}

\section[Stationary reversal: noise matching or a non-Hamiltonian witness]{The complete finite source criterion}
\label{sel37:reverse}

Let $\Phi$ be the actual CPTP map and let $\sigma\succ0$ be an
independently identified stationary density. Use exactly V35's reverse,
not the one-sided weighted adjoint excluded there:
\[
 \Phi^{\mathrm R}=\Gamma_\sigma\Phi^*\Gamma_\sigma^{-1},
 \qquad \Gamma_\sigma(X)=\sigma^{1/2}X\sigma^{1/2}.
\]
Both maps are CPTP, fix $\sigma$, and reversal is involutive. Define
\begin{equation}
 \mathcal A_\sigma=\tfrac12(\Phi-\Phi^{\mathrm R}),\qquad
 \mathcal S_\sigma=\tfrac12(\Phi+\Phi^{\mathrm R})-\id,\qquad
 \mathcal E_\sigma=Q_0(J(\Phi)-J(\Phi^{\mathrm R}))Q_0.
 \label{sel37:eq:reverse-data}
\end{equation}
The word ``circulation'' below means the antisymmetric response
$\mathcal A_\sigma$; it is not an assertion of a Hamiltonian, a particle
current, or physical time reversal.

\begin{theorem}[Exact recognition of Hamiltonian stationary circulation]
\label{sel37:thm:reverse}
There is a Hermitian $H$ with $\mathcal A_\sigma=-i[H,\cdot]$ if and
only if
\begin{equation}
 \boxed{\mathcal E_\sigma=0.}
 \label{sel37:eq:matching}
\end{equation}
On this branch the unique traceless coefficient is
\begin{equation}
 \boxed{H_\circ=\frac{i}{4n}(B_\Delta-B_\Delta^*),\qquad
 B_\Delta=\operatorname{unvec}\bigl((J(\Phi)-J(\Phi^{\mathrm R}))\iota\bigr),}
 \label{sel37:eq:Hcircle}
\end{equation}
and $[H_\circ,\sigma]=0$. The exact decomposition is
\begin{equation}
 \Phi-\id=\mathcal S_\sigma-i[H_\circ,\cdot],\qquad
 \mathcal S_\sigma^{\mathrm R}=\mathcal S_\sigma,\qquad
 \mathcal S_\sigma(\sigma)=0.
 \label{sel37:eq:finite-split}
\end{equation}
If $\mathcal E_\sigma\ne0$, the exact residual is
\eqref{sel37:eq:distance} with $C=\mathcal E_\sigma/2$.
This returns a specified non-Hamiltonian source direction, not merely a
failure to guess $H$.
\end{theorem}
\begin{proof}
$\mathcal A_\sigma$ preserves Hermiticity and annihilates the trace, and
$C_{\mathcal A_\sigma}=\mathcal E_\sigma/2$. Apply
Proposition~\ref{sel37:prop:normal}. Stationarity gives
$\mathcal A_\sigma(\sigma)=0$, which on the zero branch is exactly
$[H_\circ,\sigma]=0$. The remaining identities follow by adding and
subtracting $\Phi$ and $\Phi^{\mathrm R}$.
\end{proof}

The symmetric part is legitimate positive comparison dynamics:
if $M=(\Phi+\Phi^{\mathrm R})/2$, then
$e^{t\mathcal S_\sigma}=e^{-t}\sum_{k\ge0}t^kM^k/k!$ is CPTP and fixes
$\sigma$. This mathematical Poisson expansion is not substituted for the
original renewal clock. Nor does \eqref{sel37:eq:finite-split} imply
$\Phi=\operatorname{Ad}_{e^{-iH_\circ}}$ or a composition of two channels.
For nontracial $\sigma$, an R-symmetric generator can itself have a
nonzero canonical Hamiltonian counterterm. Accordingly
$\mathcal S_\sigma$ is called R-symmetric dynamics, not automatically
``pure dissipation.''

\begin{corollary}[The recovered odd coefficient and its word-native status]
\label{sel37:cor:odd}
Suppose the already verified physical involution satisfies
$Z\sigma Z=\sigma$ and
$\operatorname{Ad}_Z\Phi\operatorname{Ad}_Z=\Phi^{\mathrm R}$.
Even before \eqref{sel37:eq:matching} is known, the projection
$H_\circ=H_{\mathcal A_\sigma}$ obeys
\begin{equation}
 ZH_\circ Z=-H_\circ,\qquad
 H_\circ=\tfrac12\bigl(H_{\Phi-\id}-ZH_{\Phi-\id}Z\bigr).
 \label{sel37:eq:odd-H}
\end{equation}
If the forward Kraus coefficients lie in a fixed represented unital
$C^*$-algebra $\mathcal B$ and $\sigma\in\mathcal B$ is invertible,
then $H_\circ\in\mathcal B$. Thus its word-native algebra membership
requires no additional generator atom.
\end{corollary}
\begin{proof}
Reversal commutes with conjugation by $Z$ because $Z$ commutes with
$\sigma$. Thus $\operatorname{Ad}_Z\mathcal A_\sigma
\operatorname{Ad}_Z=-\mathcal A_\sigma$.
The Hilbert--Schmidt projection, or its explicit identity moment, is linear
and unitarily equivariant. Tracelessness removes any additive scalar and
proves both identities. Functional calculus gives
$\sigma^{\pm1/2}\in\mathcal B$, so the reverse Kraus coefficients
$\sigma^{1/2}k_j^*\sigma^{-1/2}$ also lie in $\mathcal B$.
For any Kraus bank,
$\operatorname{unvec}(J\iota)=\sum_j\overline{\operatorname{Tr}k_j}\,k_j$.
The forward and reverse moments, and hence $H_\circ$, belong to
$\mathcal B$.
\end{proof}

When $\mathcal E_\sigma\ne0$, \eqref{sel37:eq:odd-H} still reconstructs
the canonical odd Hamiltonian \emph{component}; the odd non-Hamiltonian
residual must also be retained. Conversely, a zero matching residual
with $H_\circ=0$ certifies no nonzero odd coherent coefficient. A canonical
linear combination is not necessarily an original resolved word, so this
result does not bypass V29's same-label occurrence test or identify the
V110 up-incidence port by dimension alone.

\section[The two unchanged source-selection examples]{A nonzero qubit circulation and an exactly zero one}
\label{sel37:examples}

\begin{theorem}[Automatic matching on a tracial qubit]
\label{sel37:thm:qubit}
For every unital qubit channel and $\sigma=I_2/2$,
$\mathcal E_\sigma=0$. Thus its Hilbert--Schmidt antisymmetric response
is always a Hamiltonian commutator. No efficiency or generating-bank
assumption is needed.
\end{theorem}
\begin{proof}
On the orthonormal Hermitian basis $(I,X,Y,Z)/\sqrt2$, unitality and
trace preservation make the channel matrix $1\oplus T$ with $T$ real.
Its adjoint has matrix $1\oplus T^{\mathsf T}$, so its antisymmetric
part is $0\oplus(T-T^{\mathsf T})/2$. Every real antisymmetric
three-by-three matrix acts as cross product with a unique real vector.
The Pauli relation $[\sigma_i,\sigma_j]=2i\epsilon_{ijk}\sigma_k$
identifies exactly these maps with $-i[H,\cdot]$ for traceless
Hermitian $H$. Apply \eqref{sel37:eq:derivation-criterion}.
\end{proof}

\begin{proposition}[The unchanged V34 pulse bank]
\label{sel37:prop:old-forward}
Keep the original coefficients
\[
 k_X=(9I+12iX)/25,\qquad k_Z=(12I+16iZ)/25.
\]
Its forgotten channel has the exact decomposition
\begin{equation}
 \boxed{\Phi-\id=-i[H_\circ,\cdot]
                  +\frac{144}{625}\mathcal D[X]
                  +\frac{256}{625}\mathcal D[Z],\qquad
 H_\circ=-\frac{108X+192Z}{625}.}
 \label{sel37:eq:old-forward}
\end{equation}
Its already reconstructed reverser is $Y$;
$YH_\circ Y=-H_\circ$ and
\begin{equation}
 \|H_\circ\|_{\HS}^2=\frac{97056}{390625},\qquad
 \operatorname{spec}H_\circ=
       \left\{-\frac{12\sqrt{337}}{625},
                   \frac{12\sqrt{337}}{625}\right\}.
 \label{sel37:eq:old-forward-spectrum}
\end{equation}
\end{proposition}
\begin{proof}
For real $a,b$,
$(aI+ibX)\rho(aI-ibX)=a^2\rho+b^2X\rho X+iab[X,\rho]$.
Apply this identity to the two original coefficients. Their scalar-square
sum is $225/625$, and the two jump-square coefficients sum to $400/625$,
so subtracting the identity channel gives
\eqref{sel37:eq:old-forward}. The Pauli anticommutation relations give
$H_\circ^2=(108^2+192^2)I/625^2$ and the displayed spectrum.
The dissipators are unchanged under Hilbert--Schmidt reversal, whereas
the commutator changes sign. This proves the exact circulation statement.
\end{proof}

No inverse pulse is added. The result identifies an odd coefficient
from the old quantum maps, not a logarithm of their finite channel.
In particular the two positive dissipators in
\eqref{sel37:eq:old-forward} do not disappear. Per accepted event this
coefficient is dimensionless; a physical Hamiltonian rate needs the
actual time calibration.

\begin{proposition}[The unchanged V35/V36 nonunital family]
\label{sel37:prop:old-nonunital}
For the original family
\[
 k_\pm(a,b)=\frac1{\sqrt2}
 \begin{pmatrix}\sqrt{1-a}&\pm\sqrt b\\
                  \mp\sqrt a&\sqrt{1-b}\end{pmatrix},
 \qquad 0<a,b<1,
\]
with $\sigma=\operatorname{diag}(b,a)/(a+b)$, its forgotten channel obeys
\begin{equation}
 \boxed{\Phi^{\mathrm R}=\Phi,\quad
         \mathcal A_\sigma=0,\quad H_\circ=0.}
 \label{sel37:eq:zero-family}
\end{equation}
This includes exactly the V35 point $(a,b)=(1/5,4/5)$ and the V36
weak-event family $a=u\epsilon$, $b=v\epsilon$.
\end{proposition}
\begin{proof}
The unchanged stationary reverse calculation gives
$\sigma^{1/2}k_\pm^*\sigma^{-1/2}=k_\mp$.
Summing the two original branches proves the first identity and hence the
others. This does not erase the marked exchange of the outcomes or the
nonzero odd jump established in V36.
\end{proof}

Thus the previous existence and uniqueness of a reverser does not itself
supply a nonzero canonical odd action. The new distinction is computed
on the same original channels, rather than on another convention chosen
to produce an answer.

\section[A smallest tracial counterexample and an exact diamond obstruction]{Classical circulation inside a qutrit is not coherent motion}
\label{sel37:qutrit}

The qubit theorem has a sharp dimension boundary. On $\C^3$ use indices
modulo three and define two fixed CPTP maps
\[
 F(\rho)=\sum_{j=0}^2E_{j+1,j}\rho E_{j,j+1},\qquad B=F^*.
\]
For $p,q>0$, $p+q<1$, consider the \emph{declared comparison channel}
\begin{equation}
 \Phi_{p,q}=(1-p-q)\id+pF+qB.
 \label{sel37:eq:qutrit-channel}
\end{equation}
Its original Kraus bank can retain the identity outcome and the six
individual directed transitions. This is not assigned to a historical
SM source.

\begin{theorem}[Reversible parity with purely non-Hamiltonian circulation]
\label{sel37:thm:qutrit}
The channel \eqref{sel37:eq:qutrit-channel} has unique stationary density
$I_3/3$. The involution $Ze_j=e_{-j}$ satisfies
$\operatorname{Ad}_Z\Phi_{p,q}\operatorname{Ad}_Z=\Phi_{p,q}^*$.
Nevertheless, if $p\ne q$, its antisymmetric response has
\begin{equation}
 \boxed{\mathcal A=\frac{p-q}{2}(F-B),\qquad H_{\mathcal A}=0,\qquad
 \Delta_{\rm Ham}(\mathcal A)=\frac32(p-q)^2>0.}
 \label{sel37:eq:qutrit-defect}
\end{equation}
More strongly,
\begin{equation}
 \boxed{\inf_{H=H^*}\|\mathcal A+i[H,\cdot]\|_\diamond=|p-q|.}
 \label{sel37:eq:qutrit-diamond}
\end{equation}
Thus dimension three is the smallest matrix dimension in which a
tracial-stationary channel can have non-Hamiltonian antisymmetric response.
\end{theorem}
\begin{proof}
On diagonal matrices the map is the irreducible aperiodic three-cycle
Markov kernel with strictly positive clockwise and counterclockwise
rates and a positive holding probability. Its unique stationary vector
is uniform; directly its other two eigenvalues are
$1-3(p+q)/2\pm i\sqrt3(p-q)/2$, strictly inside the unit circle.
On the six off-diagonal matrix units the eigenvalue is $1-p-q<1$.
This proves the full stationary assertion without an unobserved-sector
compression. Reflection exchanges $F$ and $B$.

Let $P_F=\sum_j|E_{j+1,j}\rangle\!\rangle
\langle\!\langle E_{j+1,j}|$, and similarly define $P_B$ from the
opposite transitions. These are orthogonal rank-three projections,
both orthogonal to $\iota$. Hence
\[
 J(\mathcal A)=C_{\mathcal A}
       =\tfrac{p-q}{2}(P_F-P_B),\qquad
 M_{C_{\mathcal A}}=0.
\]
Its identity moment is zero and its squared norm is
$6(p-q)^2/4$, proving \eqref{sel37:eq:qutrit-defect}.
For every Hermitian $H$, the diagonal part of $[H,E_{00}]$ vanishes,
whereas
$\mathcal A(E_{00})=(p-q)(E_{11}-E_{22})/2$.
Diagonal pinching is trace-norm contractive, so this single input gives
lower bound $|p-q|$ against \emph{every} Hamiltonian candidate.
The choice $H=0$ attains it, because
$\|F-B\|_\diamond\le2$ and the same input gives equality.
Dimension one has zero antisymmetric response, and
Theorem~\ref{sel37:thm:qubit} covers dimension two.
\end{proof}

At $p=1/3$, $q=1/6$, the exact comparison has
$\Delta_{\rm Ham}=1/24$ and diamond obstruction $1/6$.
The conditional signed Choi eigenvalues are $+1/12$ and $-1/12$, each
three times, with three zeros. They are eigenvalues of a difference,
not negative probabilities in either physical channel.

The obstruction persists beside a genuine odd Hamiltonian.
For $a,b>0$ and $H_h=h\operatorname{diag}(0,1,-1)$, the valid generator
\begin{equation}
 \mathcal G=-i[H_h,\cdot]+a(F-\id)+b(B-\id)
 \label{sel37:eq:qutrit-generator}
\end{equation}
fixes $I_3/3$ and obeys the same parity--adjoint relation. Its half-difference
is $-i[H_h,\cdot]+(a-b)(F-B)/2$. The inherited projection returns
$H_h$ exactly, while its residual has squared norm $3(a-b)^2/2$ and
optimal diamond distance $|a-b|$. The proof above applies unchanged,
since diagonal pinching kills every Hamiltonian difference. When $a=b$
there is a nonzero Hamiltonian circulation with strictly positive
R-symmetric dissipation. Equality of forward and reverse noise is
therefore weaker than absence of noise.

\section[When a finite circulation becomes a physical tangent]{A source-family transfer with the original time retained}
\label{sel37:tangent}

A finite-step circulation must not be relabelled a continuous-time action.
Suppose an \emph{actually specified} family on one fixed memory satisfies
\begin{equation}
 \Phi_\epsilon=\id+\epsilon\mathcal L+o_\diamond(\epsilon),\qquad
 \Phi_\epsilon(\sigma)=\sigma,\qquad \sigma\succ0.
 \label{sel37:eq:tangent-family}
\end{equation}
The parameter $\epsilon$ is the declared physical step, or the independently
calibrated accepted-time parameter; it is not inferred from the channel.
Suppose also that the same verified $Z$ preserves $\sigma$ and relates
each family member to its stationary reverse.

\begin{theorem}[A derived odd Hamiltonian within R-symmetric dynamics]
\label{sel37:thm:tangent}
The conditional-noise matching condition
\begin{equation}
 Q_0\bigl(J(\Phi_\epsilon)-J(\Phi_\epsilon^{\mathrm R})\bigr)Q_0
                                      =o_{\HS}(\epsilon)
 \label{sel37:eq:first-noise-match}
\end{equation}
is equivalent to the following limiting decomposition:
\begin{equation}
 \boxed{
 \mathcal L=\mathcal S-i[H,\cdot],\quad
 \mathcal S^{\mathrm R}=\mathcal S,\quad
 \mathcal S(\sigma)=0,\quad
 [H,\sigma]=0,\quad ZHZ=-H,
 }
 \label{sel37:eq:tangent-split}
\end{equation}
where $\mathcal S=(\mathcal L+\mathcal L^{\mathrm R})/2$ generates a
CPTP semigroup and $H$ is traceless and unique. It is reconstructed by
\begin{equation}
 H=\lim_{\epsilon\downarrow0}\frac{H_\circ(\Phi_\epsilon)}{\epsilon}.
 \label{sel37:eq:Hlimit}
\end{equation}
No condition $\mathcal S=0$ is imposed.
\end{theorem}
\begin{proof}
Differentiating trace preservation and Hermiticity gives the respective
linear identities for $\mathcal L$. For each $\epsilon>0$, positivity
of $J(\Phi_\epsilon)$ gives
$Q_0J(\mathcal L)Q_0\succeq0$ in the limit. The inherited finite
conditional-Choi criterion makes $\mathcal L$ a quantum Markov generator.
Stationarity gives $\mathcal L(\sigma)=0$. Reversal is a fixed continuous
linear operation at faithful $\sigma$, so
$\Phi_\epsilon^{\mathrm R}=\id+\epsilon\mathcal L^{\mathrm R}
+o(\epsilon)$. The reverse generator also generates a CPTP semigroup,
as follows either by reversing $e^{t\mathcal L}$ or by the same cone test.
Their average $\mathcal S$ is a generator and is R-symmetric.

Condition \eqref{sel37:eq:first-noise-match} is exactly
$C_{(\mathcal L-\mathcal L^{\mathrm R})/2}=0$.
Apply Theorem~\ref{sel37:thm:reverse} at the trace-annihilating level and
Corollary~\ref{sel37:cor:odd}. Linearity and continuity of the inherited
Hamiltonian projection give \eqref{sel37:eq:Hlimit}.
Conversely, a decomposition with an R-symmetric part and $[H,\sigma]=0$
has $(-i[H,\cdot])^{\mathrm R}=+i[H,\cdot]$; conditional compression
annihilates this commutator and proves
\eqref{sel37:eq:first-noise-match}.
\end{proof}

A trace-free source Hamiltonian can be recovered even when the symmetric
part is noisy. This is a weaker action claim than an all-input coherent
channel approximation. In particular, the V36 family
$a=u\epsilon,b=v\epsilon$ has $H=0$ by
Proposition~\ref{sel37:prop:old-nonunital}, while its original nonzero Lindblad
jump and dissipative generator remain. The one-context action compiler
of V36 has no faithful stationary state on its full memory, so the
stationary part of this theorem does not apply to it. Its already
established ordinary channel tangent and canonical Hamiltonian
reconstruction remain valid independently of stationarity.

The commutation $[H,\sigma]=0$ retains V36's stationary-spectrum boundary:
a nonzero $Z$-odd $H$ can connect the grading sectors only inside a common
stationary eigenspace. The dimension formula and the explicit
$\operatorname{diag}(4/5,1/5)$ obstruction are inherited, not new
assumptions or repeated theorems here.

\section[Finite data, retained error budgets, and original-clock scaling]{Operational reconstruction without executing the subtraction}
\label{sel37:data}

\begin{proposition}[Stable non-Hamiltonian and coefficient tests]
\label{sel37:prop:errors}
Let $\widehat{\mathcal K}$ and $\mathcal K$ preserve Hermiticity and
annihilate the trace, with
$\|J(\widehat{\mathcal K})-J(\mathcal K)\|_{\HS}\le\eta$.
Then
\begin{align}
 \|C_{\widehat{\mathcal K}}-C_{\mathcal K}\|_{\HS}&\le\eta,\notag\\
 \bigl|\sqrt{\Delta_{\rm Ham}(\widehat{\mathcal K})}
              -\sqrt{\Delta_{\rm Ham}(\mathcal K)}\bigr|&\le\eta,\notag\\
 \|H_{\widehat{\mathcal K}}-H_{\mathcal K}\|_{\HS}
                                      &\le\eta/\sqrt{2n}.
 \label{sel37:eq:error-bounds}
\end{align}
For a forward--reverse difference one may use
$\eta=(\eta_{\rm f}+\eta_{\rm R})/2$, where the second budget includes
any uncertainty in the stationary density used for reversal.
\end{proposition}
\begin{proof}
Compression and orthogonal projection are contractions in
Hilbert--Schmidt norm. Distance to the Hamiltonian subspace is
one-Lipschitz. The inherited identity
$\|\operatorname{ad}_H\|_{\rm sup,HS}^2=2n\|H\|_{\HS}^2$ for
traceless Hermitian $H$ gives the last bound. The difference estimate is
the triangle inequality. If $\sigma$ is exact, reversal on superoperator
Hilbert--Schmidt space has norm at most
$\lambda_{\max}(\sigma)/\lambda_{\min}(\sigma)$, which is one explicit
way to obtain $\eta_{\rm R}$ from a forward estimate. It follows from
$\|\Gamma_\sigma\|=\lambda_{\max}(\sigma)$ and
$\|\Gamma_\sigma^{-1}\|=1/\lambda_{\min}(\sigma)$.
\end{proof}

Thus $\sqrt{\Delta_{\rm Ham}(\widehat{\mathcal K})}>\eta$ certifies a
non-Hamiltonian obstruction. A small residual is not a proof of its
exact vanishing. Approximate covariance, uncertain stationary supports,
or changing physical type identifications require their own budgets.

There is no additional coherent intervention in these tests. At a fixed
identified $\sigma$, let V25's actual contextual data reconstruct
$J(\Phi)=\sum_a p_aD^a$. Reversal, identity contraction, conditional
compression, and the orthogonal Hamiltonian projection are fixed real
linear maps on these data. If $R_a$ is the Choi matrix of the
orthogonal residual, using the inherited projection on all
Hermiticity-preserving maps, of
$\tfrac12(\operatorname{unvec}_{\rm Choi}D^a-
(\operatorname{unvec}_{\rm Choi}D^a)^{\mathrm R})$, define
$W_{ab}=\operatorname{Re}\operatorname{Tr}(R_a^*R_b)$. Then
\begin{equation}
 \boxed{W\succeq0,\qquad
 \Delta_{\rm Ham}(\mathcal A_\sigma)=p^{\mathsf T}Wp.}
 \label{sel37:eq:context-statistic}
\end{equation}
On the actual trace-preserving, $\sigma$-stationary data fibre, the
combined difference is trace annihilating and its norm is given by
\eqref{sel37:eq:distance}. The individual dual-frame matrices need not
satisfy either normalization condition. The inverse-Choi symbol in this formula denotes only the linear
matrix-to-map identification, not a physical decoder. If $\sigma$ is
reconstructed from the same $p$, $W$ depends on that reconstructed state;
the statement is not a claim of unconditional quadratic dependence on
unknown stationary data. This is an application of the earlier contextual
reconstruction, not a new tomography theorem.

For the original state-independent opportunity mixture
$\Phi_\vartheta=(1-\vartheta)\id+\vartheta\Phi$, the same $\sigma$
is stationary and its reverse has the same mixture. Therefore
\begin{equation}
 H_\circ(\Phi_\vartheta)=\vartheta H_\circ(\Phi),\quad
 \mathcal E_\sigma(\Phi_\vartheta)=\vartheta\mathcal E_\sigma(\Phi),\quad
 \Delta_{\rm Ham}(\mathcal A_{\sigma,\vartheta})
                 =\vartheta^2\Delta_{\rm Ham}(\mathcal A_\sigma).
 \label{sel37:eq:clock}
\end{equation}
The physical duration must still be supplied before a coefficient is
converted to a rate. These finite identification norms are not
Yang--Mills spectral gaps or continuum mass estimates.

\section[Proof status and the remaining same-history action identification]{What the present calculation closes}
\label{sel37:status}

\paragraph{Proved in this continuation.}
The reverse-noise difference supplies an exact test for whether the
antisymmetric response of the actual stationary pair is a Hamiltonian
derivation. Its full orthogonal residual and a diamond-norm obstruction
are explicit. A verified parity makes its canonical Hamiltonian component
odd, and on the matching branch that coefficient commutes with the actual
stationary state. Tracial qubits always satisfy matching. The unchanged
V34 pulse bank has the nonzero coefficient
\eqref{sel37:eq:old-forward}; the unchanged V35/V36 nonunital family has
zero circulation. A normalized qutrit supplies the smallest tracial
counterexample, with an exact optimum over every Hamiltonian comparison.
An actual weak-event family gives the tangent transfer without requiring
vanishing R-symmetric dynamics. Error and clock factors are retained.

\paragraph{Inherited rather than repeated.}
The canonical Hamiltonian projection, finite inner-derivation theorem,
conditional Choi generator cone, V35 stationary reverse, V36 recovery and
stationary-spectrum theorems, and V25 contextual tomography are inputs.
The present explicit residual identity and evaluations refine their
interface; they do not discover canonical GKSL coordinates or quantum
detailed balance. The four companion-interface boundaries audited in V36
are unchanged. No new scalar Einstein mean, regenerative profile root,
perturbative ESM identity, or Yang--Mills global estimate is claimed.

\paragraph{Not established.}
The actual historical SM channel family, its physical clock, and its
identification with a protected matter grading are not populated by these
calculations. Even a nonzero recovered $H_\circ$ is not automatically the
manuscript's full common action or its determinant-incidence operator.
The nonzero V34 coefficient is a qubit coefficient, not a colour or
three-generation construction. A generator's canonical odd component is
not an executed odd outcome. Full-record covariance and the directed
V110 typing still have to be audited where they are claimed.

\paragraph{Next source-level target.}
For a \emph{named} source compatible with the already identified parity,
evaluate the forward--reverse conditional Choi difference first. If it
vanishes, compute the derived $H_\circ$ rather than assume another action
matrix. If it does not, retain the explicit asymmetric-noise residual;
a different adjoint convention is not a repair. To use a tangent as the
physical odd action, compare it with the separately reconstructed
common-action variation on the same carrier and clock. For the purely
structural claim, retain directed typing and V29's finite same-label
search. This continuation does not add stationarity or detailed balance
as universal axioms of the duality.

\paragraph{Verification scope.}
The accompanying executable checks use rational and algebraic matrices,
exact ranks, exact Choi reshuffling and partial traces. They test the
normal form against full superoperator matrices, the unchanged examples,
the three-cycle stationary law and reflection, the exact conditional
spectrum, the sharp norm witnesses, and the data/clock identities.
The analytic equivalences and norm optima are established by the written
proofs, not by numerical rank thresholds or a sampled optimization.
Prior companion packages are not claimed to have been independently
verified. The V36 text is preserved apart from its terminal document
delimiter.

\addtocontents{toc}{\protect\endgroup}
\renewcommand{\refname}{Additional sources and background for V37}
\begin{thebibliography}{99}
\raggedright
\bibitem[56]{sel37-grand}
Renewal Geometry research programme,
\emph{Renewal Grand Tensor Monograph V076}, supplied source
\src{Renewal_Grand_Tensor_Monograph_V076(1).tex}.
Section \src{sec:SMST-channel-native-tangent}; in particular
\src{def:SMST-Hamiltonian-tangent},
\src{thm:SMST-Hamiltonian-tangent-certificate}, and
\src{thm:SMST-channel-native-tangent}.
\bibitem[57]{sel37-hs}
P.~Hayden and J.~Sorce,
\emph{A canonical Hamiltonian for open quantum systems},
J. Phys. A: Math. Theor. \textbf{55}, 225302 (2022).
\href{https://arxiv.org/abs/2108.08316}{arXiv:2108.08316}.
Canonical Hamiltonian projection and the traceless-jump convention
are established background, not new claims of this note.
\bibitem[58]{sel37-fu}
F.~Fagnola and V.~Umanit\`a,
\emph{Generators of KMS symmetric Markov semigroups on B(h):
symmetry and quantum detailed balance},
\href{https://arxiv.org/abs/0908.0967}{arXiv:0908.0967}.
Cited for the established symmetric stationary duality and detailed-balance
framework. The source's unitary grading here is not identified with an
antiunitary time-reversal operator.
\end{thebibliography}

% END OF SOURCE-SELECTION V37 DEVELOPMENT BLOCK.
% Preserve V1--V37 and append subsequent work before the final terminator.


\clearpage
\hypersetup{pdftitle={Historical Standard-Model Source Selection: V1--V38}}
\begin{center}
{\Large\bfseries Source-selection continuation V38}\\[.4em]
{\large Event circulation versus physical-time generators:\\
an all-branch embedding obstruction and a finite-time action audit}
\end{center}
\addtocontents{toc}{\protect\begingroup}
\addtocontents{toc}{\protect\renewcommand{\protect\numberline}{\protect\makebox[2.1em][l]}}

\section[The unresolved time interface, not another reversal convention]{What a finite source map does and does not identify}
\label{sel38:context}

V37 reconstructs the Hamiltonian projection of the actual finite
forward--reverse difference. Its conclusion is deliberately not that this
coefficient is a logarithm of the source map. The next comparison with a
common action must respect this distinction. This continuation evaluates
that interface before asking for another grading, action, or typed source.

Two results point in opposite directions. The unchanged V34 two-pulse
channel has the nonzero Hamiltonian circulation computed in
Proposition~\ref{sel37:prop:old-forward}, but is not the finite-time value
of \emph{any} time-homogeneous CPTP semigroup on that qubit. Conversely,
an explicit primitive qutrit semigroup with a genuine odd Hamiltonian
and reverse-symmetric depolarization fails V37's exact finite-step
noise-matching test at almost every short positive time. Neither result
contradicts V37: the map at one time and its infinitesimal generator are
different objects.

The general channel-embedding problem, logarithm branches and conditional
Choi tests are established subjects \cite{sel38-wolf}. The canonical
Hamiltonian projection is already supplied by the monograph
\cite{sel37-grand,sel37-hs}. We do not claim either theory as new.
The developments here are a finite isometric-sector obstruction evaluated
on the unchanged source, the opposite qutrit calculation, and their
combination into a clock-preserving test against an independently
reconstructed action. Every time parameter below is in a fixed declared
unit. No exponential waiting law is substituted for the original
renewal clock.

Use V37's output-first Choi convention
\[
 J(\mathcal K)=\sum_{i,j}\mathcal K(E_{ij})\otimes E_{ij},
 \qquad \iota=|I_n\rangle\!\rangle,\qquad
 Q_0=I-\iota\iota^*/n.
\]
Write $\|\cdot\|_{\mathrm F}$ for the Hilbert--Schmidt norm of a
superoperator in an orthonormal Hilbert--Schmidt operator basis; Choi
reshuffling preserves this norm. Matrix-operator norm on that space
is denoted $\|\cdot\|_{2\to2}$. The physical diamond norm is not divided
by two. The stationary reverse remains exactly
$\mathcal K^{\mathrm R}=\Gamma_\sigma\mathcal K^\dagger
\Gamma_\sigma^{-1}$, with
$\Gamma_\sigma(X)=\sigma^{1/2}X\sigma^{1/2}$.

\section[An exact finite obstruction to every homogeneous embedding]{A contraction semigroup cannot lose an equality sector later}
\label{sel38:obstruction}

The following elementary Hilbert-space lemma gives a branch-independent
test. It does not compute a matrix logarithm.

\begin{lemma}[Finite-time equality sectors of contraction semigroups]
\label{sel38:lem:equality}
Let $M$ act on a $d$-dimensional complex Hilbert space and suppose
$M+M^*\preceq0$. Put $D=-(M+M^*)\succeq0$ and $T_t=e^{tM}$.
For every $t>0$,
\begin{equation}
 \boxed{\ker(I-T_t^*T_t)
 =\bigcap_{j=0}^{d-1}\ker(D^{1/2}M^j)=:\mathcal K.}
 \label{sel38:eq:equality-space}
\end{equation}
The space $\mathcal K$ reduces $M$ and every $T_s$. The restriction
$M|_{\mathcal K}$ is skew-adjoint; in particular $T_s|_{\mathcal K}$
is unitary. Thus the equality sector is independent of the positive
observation time.
\end{lemma}
\begin{proof}
Differentiate $T_s^*T_s$ and integrate:
\[
 I-T_t^*T_t=\int_0^t e^{sM^*}D e^{sM}\,ds.
\]
A vector $v$ lies in its kernel exactly when
$D^{1/2}e^{sM}v=0$ for all $s\in[0,t]$, because the integrand's
quadratic form is nonnegative and continuous. Analyticity then gives
$D^{1/2}M^jv=0$ for every $j\ge0$. Cayley--Hamilton reduces this
to $0\le j<d$, proving \eqref{sel38:eq:equality-space}.
It also proves $M\mathcal K\subseteq\mathcal K$.
For $v\in\mathcal K$, $Dv=0$, hence $M^*v=-Mv\in\mathcal K$.
The space therefore reduces $M$ and the restriction is skew-adjoint.
\end{proof}

Let $\Phi$ be an actually identified CPTP map with a unique faithful
stationary density $\sigma$. Its Heisenberg map $\mathsf T=\Phi^\dagger$
is a contraction for
\[
 \langle X,Y\rangle_\sigma=\Tr(\sigma X^*Y),
\]
by the Schwarz inequality and stationarity. Its Hilbert adjoint
$\mathsf T^{\ddagger_\sigma}$ below is used only for this norm.
It is \emph{not} substituted for the stationary CP reverse; V35's
distinction between those operations remains in force. Let $P_\Phi$
be the orthogonal projection, in this metric, onto
\[
 \mathcal K_\Phi=\ker(I-\mathsf T^{\ddagger_\sigma}\mathsf T).
\]

\begin{theorem}[A source-native obstruction to all homogeneous logarithm branches]
\label{sel38:thm:obstruction}
If $\Phi=e^{t_0\mathcal L}$ for some $t_0>0$ and a time-homogeneous
CPTP generator $\mathcal L$ on the same $M_n$, then
\begin{equation}
 \boxed{(I-P_\Phi)\mathsf T P_\Phi=0.}
 \label{sel38:eq:invariance}
\end{equation}
Consequently the strictly positive finite residual
\[
 \eta_\Phi
 =\|(I-P_\Phi)\mathsf T P_\Phi\|_{\mathrm F,\sigma}^2>0
\]
excludes every such generator, at every positive duration and on every
logarithm branch.
\end{theorem}
\begin{proof}
The maps $\mathcal L$ and $\Phi$ commute. Therefore
$\Phi(\mathcal L(\sigma))=\mathcal L(\sigma)$.
The latter matrix is Hermitian and traceless. A nonzero such fixed
matrix would, by faithfulness, make $\sigma\pm\epsilon\mathcal L(\sigma)$
two distinct stationary densities for sufficiently small $\epsilon$.
Uniqueness gives $\mathcal L(\sigma)=0$.
Every $e^{s\mathcal L^\dagger}$ is consequently a contraction in the
same $\sigma$-metric. Differentiation at zero gives dissipativity of
$\mathcal L^\dagger$ in that Hilbert space. Apply
Lemma~\ref{sel38:lem:equality}. Its equality sector at $t_0$ reduces the
semigroup, proving \eqref{sel38:eq:invariance}.
\end{proof}

The theorem is necessary, not sufficient for embedding. If the actual
stationary state is not unique, the argument applies to any proposed
embedding that preserves the specified faithful state; uniqueness is
what removes that extra proviso. Exact kernel assertions require exact
data or an independent structural certificate. Small estimated
singular-value deficits are not promoted to exact equality sectors.

\begin{theorem}[The unchanged forward-pulse source has no homogeneous embedding]
\label{sel38:thm:old-source}
Retain exactly
\[
 k_X=(9I+12iX)/25,\qquad k_Z=(12I+16iZ)/25,
 \qquad \Phi(\rho)=k_X\rho k_X^*+k_Z\rho k_Z^*.
\]
This CPTP map is not $e^{t\mathcal L}$ for any $t>0$ and any
time-homogeneous CPTP generator on its original qubit.
\end{theorem}
\begin{proof}
Normalization and unitality follow from the two original coefficients.
In the orthonormal Pauli coordinates of the traceless Hermitian space,
its state-map matrix is
\begin{equation}
 T=\frac1{625}
 \begin{pmatrix}
 113&384&0\\
 -384&-175&216\\
 0&-216&337
 \end{pmatrix}.
 \label{sel38:eq:old-bloch}
\end{equation}
Direct multiplication gives
\[
 \operatorname{spec}(T^{\mathsf T}T)
 =\{1,77281/390625,77281/390625\},\qquad
 \det(I-T)=294912/390625\ne0.
\]
Thus $I/2$ is the unique stationary density. Define
\[
 v=\frac{(-3,-4,3)^{\mathsf T}}{\sqrt{34}},\qquad
 w=\frac{(-3,4,3)^{\mathsf T}}{\sqrt{34}}.
\]
Then $Tv=w$, $T^{\mathsf T}w=v$, and $v^{\mathsf T}w=1/17$.
The Heisenberg equality sector is $\C I$ together with the Pauli
direction $w$. Its image has nonzero component outside that sector.
The exact noninvariance residual, equivalently computed for the state
map and the projection onto $\C v$, is
\begin{equation}
 \boxed{\eta_\Phi=1-\frac1{17^2}=\frac{288}{289}>0.}
 \label{sel38:eq:old-obstruction}
\end{equation}
Apply Theorem~\ref{sel38:thm:obstruction}.
\end{proof}

There is also a direct state-level interpretation. The pure input with
Bloch vector $v$ is sent to the pure state with vector $w$. A second
unchanged application yields
\[
 \|T^2v\|^2=\frac{1332209}{6640625},\qquad
 \Tr\Phi^2(\rho_v)^2=\frac{3986417}{6640625}<1.
\]
An exact norm-equality interval inside a contraction semigroup would
belong to its invariant equality sector, not to a direction that loses
norm on the next interval.

This does not alter V37's computed
$H_\circ=-(108X+192Z)/625$ or its two dissipators. Indeed
$\Phi-\id=\sum_a\mathcal D[k_a]$ is always a valid generator
for a normalized Kraus bank. Executing it at rate $\nu$ would describe
the \emph{separate} Poisson-jump evolution
$e^{\nu t(\Phi-\id)}$, not the identity
$e^{t\mathcal L}=\Phi$ just excluded. Neither that Poisson clock nor a
continuous embedding is imposed on the original renewal source.
Discrete iteration $\Phi^m$ remains valid, as do time-dependent or
larger-carrier realizations when separately specified.
The obstruction is to a homogeneous semigroup on this same memory,
not to quantum dynamics or to a structural finite Dirac operator.


\begin{proposition}[A second certificate: constant positive-time Choi rank]
\label{sel38:prop:rank}
For every finite-dimensional, time-homogeneous CPTP semigroup
$e^{t\mathcal L}$, the Choi rank is constant for all $t>0$.
In particular, for an observed channel $\Phi$,
\begin{equation}
 \boxed{\rank J(\Phi^2)\ne\rank J(\Phi)
 \ \Longrightarrow\
 \Phi\ne e^{t\mathcal L}\text{ for every }t>0.}
 \label{sel38:eq:rank-test}
\end{equation}
This obstruction does not require unitality or a faithful stationary state.
\end{proposition}
\begin{proof}
Write a finite GKSL representation
$\mathcal L(X)=GX+XG^*+\sum_jL_jXL_j^*$,
where $G=-iH-\frac12\sum_jL_j^*L_j$. The positive quantum-trajectory
expansion has no-jump coefficient $e^{tG}$ and trajectory coefficients
\[
 e^{tG}\widetilde L_{j_k}(s_k)\cdots
       \widetilde L_{j_1}(s_1),\qquad
 \widetilde L_j(s)=e^{-sG}L_je^{sG},\qquad
 0<s_1<\cdots<s_k<t .
\]
The series of their positive Choi integrals converges in finite
dimension. A vector is in its kernel exactly when it is orthogonal
to the no-jump coefficient and to every trajectory coefficient:
zero of the nonnegative integrals implies zero of their continuous
integrands on the open ordered simplex.

Let $\mathfrak B$ be the unital algebra generated by
$(\operatorname{ad}_G)^rL_j$ for $r\ge0$. All trajectory coefficients
belong to $e^{tG}\mathfrak B$. Conversely, a functional vanishing
on their coefficient span vanishes, for every sequence of labels,
on an analytic function of $s_1,\ldots,s_k$ on an open set.
Its analytic continuation and all derivatives at zero vanish.
Those derivatives are exactly the products of
$(\operatorname{ad}_G)^rL_j$, up to signs.
The coefficient span is therefore $e^{tG}\mathfrak B$.
The no-jump coefficient supplies its identity element.
As $e^{tG}$ is invertible, its dimension is
$\dim\mathfrak B$, independent of $t>0$.
Only $r<n^2$ are needed by Cayley--Hamilton for
$\operatorname{ad}_G$ on $M_n$.
Finally, $\Phi^2=e^{2t\mathcal L}$ would have the same rank.
\end{proof}

For the unchanged two-pulse source,
\[
 \rank J(\Phi)=2,\qquad \rank J(\Phi^2)=4.
\]
Indeed, the four original composite coefficients
$k_Xk_X,k_Xk_Z,k_Zk_X,k_Zk_Z$ have vectorization determinant
\[
 -\frac{24461180928}{152587890625}\ne0
\]
in that order and the output-first row vectorization convention.
This independently certifies the same all-branch obstruction.
No hypothetical refinement of either outcome is introduced:
these are the actual two-event coefficient products.
The rank criterion concerns the forgotten map; a homogeneous
realization of a more finely marked process on the same memory
would also have to pass it.

\section[A coherent generator can have non-Hamiltonian finite circulation]{An opposite qutrit calculation with unchanged time units}
\label{sel38:sampling}

A failed finite matching test can also be too strong for an infinitesimal
action claim. This is demonstrated without unequal clockwise and
counterclockwise noise.

\begin{theorem}[Finite sampling creates a non-Hamiltonian antisymmetric map]
\label{sel38:thm:qutrit}
On $M_3$, let
\[
 \begin{gathered}
 D_0=\diag(-1,0,1),\qquad H=\omega D_0,\qquad
 \mathcal T(\rho)=\Tr(\rho)I_3/3,\\
 \mathcal L=-i[H,\cdot]+\gamma(\mathcal T-\id),
 \qquad \gamma>0,\quad\omega>0 .
 \end{gathered}
\]
This is a CPTP generator with unique stationary density $I_3/3$.
Let $Z$ exchange the first and third basis vectors. Then
$ZHZ=-H$ and $\operatorname{Ad}_Z\mathcal L\operatorname{Ad}_Z
=\mathcal L^\dagger$.
The exact channel at physical duration $t$ is
\begin{equation}
 \Phi_t=e^{-\gamma t}\operatorname{Ad}_{e^{-itH}}
                    +(1-e^{-\gamma t})\mathcal T .
 \label{sel38:eq:qutrit-channel}
\end{equation}
For $\mathcal A_t=(\Phi_t-\Phi_t^\dagger)/2$, its canonical coefficient
and full Hamiltonian residual are
\begin{align}
 H_\circ(t)
 &=\frac{e^{-\gamma t}}3
       \bigl(\sin(\omega t)+\sin(2\omega t)\bigr)D_0,\notag\\
 \Delta_{\rm Ham}(\mathcal A_t)
 &=\frac23 e^{-2\gamma t}
       \bigl(2\sin(\omega t)-\sin(2\omega t)\bigr)^2.
 \label{sel38:eq:qutrit-residual}
\end{align}
In particular the residual is strictly positive for
$0<\omega t<\pi$, although
$(\mathcal L-\mathcal L^\dagger)/2=-i[H,\cdot]$ exactly.
\end{theorem}
\begin{proof}
The depolarizing generator has jumps
$\sqrt{\gamma/3}E_{ij}$, and commutes with the Hamiltonian
superoperator. This gives \eqref{sel38:eq:qutrit-channel}, stationarity
and uniqueness: every traceless component has eigenvalue of modulus
$e^{-\gamma t}<1$ at each $t>0$. Reflection changes $H$ to $-H$ and
preserves depolarization, proving the parity relation.

Put $a=e^{-\gamma t}\sin(\omega t)$ and
$b=e^{-\gamma t}\sin(2\omega t)$. The map $\mathcal A_t$ vanishes
on diagonal matrix units. It sends $E_{01},E_{12},E_{02}$ to,
respectively, $iaE_{01},iaE_{12},ibE_{02}$, with the conjugate
relations on reverse matrix units. Its identity Choi contraction is
diagonal, so the inherited Hamiltonian projection is diagonal.
Equivalently, every off-diagonal Hamiltonian direction is orthogonal
to this matrix-unit multiplier. For a diagonal candidate, the normal
equations give the unique traceless minimizer $hD_0$ with
$h=(a+b)/3$. Its squared residual is
\[
 4(a-h)^2+2(b-2h)^2=\frac23(2a-b)^2.
\]
This proves \eqref{sel38:eq:qutrit-residual}. Finally
$2\sin x-\sin2x=2\sin x(1-\cos x)>0$ for $0<x<\pi$.
\end{proof}

This is a declared comparison semigroup, not a new historical source.
The same finite-sampling phenomenon remains at $\gamma=0$, although
the unique-stationary-state assertion then need not hold.
It is the nonlinear sine of different energy gaps, not a negative
physical noise probability, that prevents a finite antisymmetric map
from satisfying the additive Hamiltonian gap relations.

\begin{proposition}[An exact positive-data point]
\label{sel38:prop:qutrit-point}
At the declared time $t=1$, with $\gamma=\log2$ and $\omega=\pi/3$,
put
\[
 U=\diag(e^{i\pi/3},1,e^{-i\pi/3}),\qquad
 \Phi_*=\tfrac12\operatorname{Ad}_U+\tfrac12\mathcal T .
\]
Then
\begin{equation}
 \boxed{J(\Phi_*)=\tfrac12|U\rangle\!\rangle\langle\!\langle U|
                    +\tfrac16 I_9\succ0,\quad
 H_\circ=\tfrac{\sqrt3}{6}D_0,\quad
 \Delta_{\rm Ham}=\tfrac18.}
 \label{sel38:eq:qutrit-point}
\end{equation}
Nevertheless its principal superoperator logarithm is exactly
\[
 \log\Phi_*=-i[(\pi/3)D_0,\cdot]+(\log2)(\mathcal T-\id).
\]
The conditional Choi block of this generator is
$(\log2)Q_0/3\succeq0$, and its forward--reverse half-difference
is the Hamiltonian commutator with $(\pi/3)D_0$.
\end{proposition}
\begin{proof}
Equation \eqref{sel38:eq:qutrit-point} follows by substitution in the
previous theorem. On $I$, $\Phi_*$ has eigenvalue one. On traceless
diagonals it has eigenvalue $1/2$; on $E_{ij}$, $i\ne j$, it has
eigenvalue $e^{-i(h_i-h_j)}/2$. All their arguments lie strictly
between $-\pi$ and $\pi$, because the maximal gap is $2\pi/3$.
The principal scalar logarithms on these invariant spaces give the
displayed superoperator logarithm. Finally
$Q_0J(\mathcal T-\id)Q_0=Q_0/3$, whereas Hamiltonian
conditional compression is zero.
\end{proof}

The input channel at this point is algebraic: its entries use only
rational numbers, $i$ and $\sqrt3$. The recovered rates $\log2$
and $\pi/3$ follow from its spectral logarithms in the declared time
unit. A nonzero finite-step residual $1/8$ is therefore not a witness
of asymmetric \emph{generator} noise in this example. It remains an
exact witness about the finite map to which V37 applies it.

\section[A positive logarithm test, with its branch and physical scope stated]{Reconstruct a compatible generator before comparing its action}
\label{sel38:logarithm}

The following packages standard logarithm and generator criteria with
V37's exact reverse comparison. It is a reconstruction on a stated
branch, not a general new solution of the embedding problem.

\begin{theorem}[Principal-branch stationary generator and action test]
\label{sel38:thm:log}
Let $\Phi$ be an invertible CPTP map fixing an identified
$\sigma\succ0$, and suppose
$\operatorname{spec}\Phi\cap(-\infty,0]=\varnothing$.
For a declared duration $t_0>0$, define
\[
 \mathcal L_0=t_0^{-1}\operatorname{Log}\Phi,\quad
 C_0=Q_0J(\mathcal L_0)Q_0,\quad
 \mathcal A_0=(\mathcal L_0-\mathcal L_0^{\mathrm R})/2.
\]
Then $\mathcal L_0$ preserves Hermiticity, annihilates the trace and
fixes $\sigma$. Moreover
\begin{equation}
 \boxed{C_0\succeq0
 \iff \mathcal L_0\text{ generates a CPTP semigroup}.}
 \label{sel38:eq:log-positive}
\end{equation}
Every compatible generator whose spectral imaginary parts lie in
$(-\pi/t_0,\pi/t_0)$ equals $\mathcal L_0$. If
\eqref{sel38:eq:log-positive} passes, then
\begin{equation}
 \boxed{
 Q_0J(\mathcal A_0)Q_0=0
 \iff
 \mathcal L_0=\mathcal S_0-i[H_0,\cdot],\
 \mathcal S_0^{\mathrm R}=\mathcal S_0,\
 [H_0,\sigma]=0,}
 \label{sel38:eq:log-matching}
\end{equation}
where $\mathcal S_0$ is a CPTP generator and $H_0$ is the unique
traceless Hermitian coefficient. A verified involution $Z$ fixing
$\sigma$ and satisfying
$\operatorname{Ad}_Z\Phi\operatorname{Ad}_Z=\Phi^{\mathrm R}$
also satisfies $ZH_0Z=-H_0$.
\end{theorem}
\begin{proof}
In a Hermitian operator basis $\Phi$ has real entries. The principal
logarithm commutes with complex conjugation on the indicated domain,
so its matrix is real in that basis. The trace functional is a left
eigenvector of $\Phi$ at one; $\sigma$ is a right eigenvector there.
Analytic functional calculus and $\operatorname{Log}1=0$ give the
trace and stationarity assertions. The finite conditional-Choi
generator theorem, already inherited in V37, gives
\eqref{sel38:eq:log-positive}; the positive conditional block supplies
the traceless jumps, and the remaining trace-annihilating part has
the canonical Hamiltonian normal form.

The principal logarithm commutes with similarity, and
$\operatorname{Log}(\Phi^\dagger)=(\operatorname{Log}\Phi)^\dagger$.
Consequently
\[
 \operatorname{Log}(\Phi^{\mathrm R})
 =\Gamma_\sigma(\operatorname{Log}\Phi)^\dagger\Gamma_\sigma^{-1}.
\]
If $e^{t_0\mathcal L}=\Phi$ and the spectrum of $t_0\mathcal L$
lies in the open principal strip, the analytic identity
$\operatorname{Log}(e^z)=z$ in that strip, applied to $\mathcal L$,
gives $\mathcal L=\mathcal L_0$, including nondiagonalizable cases.
This proves the stated uniqueness.

When $\mathcal L_0$ is a generator, reversal of $e^{s\mathcal L_0}$
shows that $\mathcal L_0^{\mathrm R}$ is also a generator.
Their average $\mathcal S_0$ is therefore a generator. V37's
derivation criterion applied to $\mathcal A_0$ gives the matching
equivalence, stationarity gives $[H_0,\sigma]=0$, and functional
calculus transports the given parity to $\mathcal L_0$.
The oddness conclusion is the same canonical uniqueness argument
as Corollary~\ref{sel37:cor:odd}.
\end{proof}

Failure of $C_0\succeq0$ excludes this principal branch, not all
logarithms. Theorem~\ref{sel38:thm:old-source}, in contrast, excludes
all homogeneous branches for its particular source without logarithms.
The original V35 point $(a,b)=(1/5,4/5)$, retained in V36, has a
singular superoperator and is not eligible for a finite matrix
logarithm at all. Its eigenvalue-zero modes are retained rather than
replaced by an invertible fit; this assertion is not made for the
entire two-parameter family.

Even a positive logarithm establishes \emph{consistency} with one
homogeneous dynamics, not that the historical source implements it
between observations. Neither the scalar physical duration nor a
subdivision law follows from a single channel matrix.

\begin{proposition}[All grid-time cylinders can hide different physical tangents]
\label{sel38:prop:clock-ambiguity}
Let $\mathcal L$ be any nonzero CPTP generator and let $t_0>0$
be a fixed observation interval. For every $|a|<1$ the family
\[
 f_a(s)=s+\frac{a t_0}{2\pi}\sin(2\pi s/t_0),\qquad
 \Phi^{(a)}_{s_2,s_1}
   =e^{(f_a(s_2)-f_a(s_1))\mathcal L},\quad s_2\ge s_1,
\]
is a normalized time-dependent, CP-divisible dynamics.
Every map between consecutive grid times $jt_0,(j+1)t_0$ equals
$e^{t_0\mathcal L}$. Hence all grid-time cylinder probabilities,
including arbitrary identical interventions at those grid times,
are independent of $a$. But its initial physical generator is
\[
 \left.\frac{d}{ds}\Phi^{(a)}_{s,0}\right|_{s=0}
 =(1+a)\mathcal L.
\]
If $\mathcal L$ has a nonzero Hamiltonian coefficient, that coefficient
also changes by $1+a$.
\end{proposition}
\begin{proof}
$f_a'(s)=1+a\cos(2\pi s/t_0)>0$. Positive increments give CPTP
propagators and their additive exponents give composition. At grid
times $f_a(jt_0)=jt_0$, proving equality of every propagated interval
and therefore of every intervened cylinder built from them.
Differentiation at zero proves the tangent assertion.
\end{proof}

The proposition concerns the specified unmarked propagators and identical
external interventions. It does not assert equality of an additional,
continuously monitored jump record. The physical time $s$ is the same
in all these comparison sources. They have different rates between the
retained samples; this is not a permitted relabelling of a measured clock. The example says precisely
which information the grid fails to contain. Observing a finer actual
interval or imposing an independently justified homogeneous generator
changes the data. Selecting a convenient logarithm does not.

\section[Quantitative sampling correction and a held-out action comparison]{Finite source families can test an action without executing a signed map}
\label{sel38:action}

Let an actually specified homogeneous family be
$\Phi_t=e^{t\mathcal L}$, fixing the same faithful $\sigma$, and put
$\mathcal L^{\mathrm R}=\Gamma_\sigma\mathcal L^\dagger
\Gamma_\sigma^{-1}$. Suppose
\[
 \mathcal K=\frac{\mathcal L-\mathcal L^{\mathrm R}}2
        =-i[H,\cdot],\qquad
 \mathcal S=\frac{\mathcal L+\mathcal L^{\mathrm R}}2 .
\]
The generators and their clock are fixed; they are not changed in
the following estimates.

\begin{proposition}[The exact first sampling correction]
\label{sel38:prop:sampling}
For $\mathcal A_t=(\Phi_t-\Phi_t^{\mathrm R})/2$,
\begin{equation}
 \mathcal A_t=t\mathcal K+
        \frac{t^2}{2}(\mathcal S\mathcal K+\mathcal K\mathcal S)
        +O_{\mathrm F}(t^3).
 \label{sel38:eq:sampling-expansion}
\end{equation}
Without unspecified remainder constants,
\begin{align}
 \|\mathcal A_t-t\mathcal K\|_{\mathrm F}
 \le\frac{t^2}{4}\bigl(
 &\|\mathcal L^2\|_{\mathrm F}e^{t\|\mathcal L\|_{2\to2}} \notag\\
 &+\|(\mathcal L^{\mathrm R})^2\|_{\mathrm F}
                    e^{t\|\mathcal L^{\mathrm R}\|_{2\to2}}\bigr).
 \label{sel38:eq:first-budget}
\end{align}
In particular, $\sqrt{\Delta_{\rm Ham}(\mathcal A_t)}$ is at most
the right side. Nonzero finite residuals of this size do not exclude
a Hamiltonian antisymmetric generator.
\end{proposition}
\begin{proof}
Expand both exponentials; their quadratic difference is
$2(\mathcal S\mathcal K+\mathcal K\mathcal S)$.
For every finite matrix $M$,
\[
 \|e^{tM}-I-tM\|_{\mathrm F}
 \le\frac{t^2}{2}\|M^2\|_{\mathrm F}e^{t\|M\|_{2\to2}},
\]
by the exponential series and
$\|M^{j+2}\|_{\mathrm F}\le\|M^2\|_{\mathrm F}\|M\|_{2\to2}^j$.
Apply this to the two generators and divide their difference by two.
The distance to the Hamiltonian subspace is no larger than the
distance from $t\mathcal K$.
\end{proof}

\begin{theorem}[Two actual durations remove the quadratic sampling bias]
\label{sel38:thm:richardson}
If both $\Phi_t$ and $\Phi_{t/2}$ are actually determined on the
same port and clock, define the signed data estimator
\[
 \mathcal R_t=\frac{4\mathcal A_{t/2}-\mathcal A_t}{t}.
\]
Then
\begin{align}
 \|\mathcal R_t-\mathcal K\|_{\mathrm F}
 \le\frac{t^2}{8}\bigl(
 &\|\mathcal L^3\|_{\mathrm F}e^{t\|\mathcal L\|_{2\to2}}\notag\\
 &+\|(\mathcal L^{\mathrm R})^3\|_{\mathrm F}
                   e^{t\|\mathcal L^{\mathrm R}\|_{2\to2}}\bigr)
 =:B_t.
 \label{sel38:eq:second-budget}
\end{align}
If the estimated $\mathcal A_t,\mathcal A_{t/2}$ have Frobenius
errors at most $\eta_t,\eta_{t/2}$, the total budget is
\[
 B_t^{\rm data}=B_t+(4\eta_{t/2}+\eta_t)/t.
\]
Consequently
\[
 \|H_{\widehat{\mathcal R}_t}-H\|_{\HS}
       \le B_t^{\rm data}/\sqrt{2n},\qquad
 \sqrt{\Delta_{\rm Ham}(\widehat{\mathcal R}_t)}
       \le B_t^{\rm data}.
\]
\end{theorem}
\begin{proof}
The linear term in $4\mathcal A_{t/2}-\mathcal A_t$ is
$t\mathcal K$, and its quadratic term cancels.
The cubic exponential remainders are bounded by
$t^3\|M^3\|_{\mathrm F}e^{t\|M\|}/6$.
The additional factor $1/2$ in $\mathcal A_t$, the four half-step
terms and division by $t$ give the constant
$4/(8\cdot12)+1/12=1/8$ in \eqref{sel38:eq:second-budget}.
The noise estimate is the triangle inequality.
Projection onto Hamiltonian derivations and the exact norm
$\|-i[H,\cdot]\|_{\mathrm F}^2=2n\|H\|_{\HS}^2$
give the last claims.
\end{proof}

The estimator is generally not positive. It is classical postprocessing
of identified maps, as in V24--V28, and is not an executed source repair.
A formal square root of $\Phi_t$ does not substitute for the actual
half-duration source. The bound is useful only with a genuine rate
budget; no cutoff-uniform conclusion follows if its constants diverge.

\begin{proposition}[A held-out common-action comparison]
\label{sel38:prop:action}
Let $A_{\rm com}=A_{\rm com}^*$ be independently reconstructed on
the same carrier in the same time unit, and put
$A_{\rm com}^0=A_{\rm com}-\Tr(A_{\rm com})I/n$.
For any Hermiticity-preserving trace-annihilating comparison
$\mathcal K_0$,
\begin{equation}
 \boxed{
 \|\mathcal K_0+i[A_{\rm com},\cdot]\|_{\mathrm F}^2
 =\Delta_{\rm Ham}(\mathcal K_0)
           +2n\|H_{\mathcal K_0}-A_{\rm com}^0\|_{\HS}^2.}
 \label{sel38:eq:heldout}
\end{equation}
Thus exact agreement requires both a vanishing non-Hamiltonian
remainder and agreement of the two independently obtained
Hamiltonian coefficients modulo a scalar.
\end{proposition}
\begin{proof}
Subtract the canonical Hamiltonian projection from $\mathcal K_0$.
Its residual is orthogonal to every Hamiltonian commutator by V37.
Pythagoras and the inner-derivation norm identity give
\eqref{sel38:eq:heldout}.
\end{proof}

Equation \eqref{sel38:eq:heldout} is an application of the inherited
projection theorem, not a new action-occurrence theorem.
It applies to the positive logarithm's $\mathcal A_0$, or to a
physically calibrated tangent. For the corrected qutrit logarithm,
$A_{\rm com}=(\pi/3)D_0$ makes both terms zero. Using the finite
$\mathcal A_1$ instead gives
\[
 \frac18+12\left(\frac{\pi}{3}-\frac{\sqrt3}{6}\right)^2>0,
\]
a mismatch caused by comparing different source objects.
The same numerical coefficient on a different typed memory or a
different clock would not pass the common-carrier premise.

\section[Error control, clock retention, and the next source decision]{What this continuation changes}
\label{sel38:status}

\begin{proposition}[A local logarithm error budget]
\label{sel38:prop:log-error}
Let $\Phi,\widehat\Phi$ be trace-preserving Hermiticity-preserving
maps fixing the same exact $\sigma\succ0$. Suppose
\[
 \|\Phi-\id\|_{2\to2}\le r<1,\quad
 \|\widehat\Phi-\id\|_{2\to2}\le r,\quad
 \|\widehat\Phi-\Phi\|_{\mathrm F}\le\delta .
\]
Let $\mathcal A_0,\widehat{\mathcal A}_0$ be the antisymmetric
parts of their duration-$t_0$ principal logarithms. Then
\begin{equation}
 \boxed{
 \|\widehat{\mathcal A}_0-\mathcal A_0\|_{\mathrm F}
 \le\frac{1+\kappa_\sigma}{2t_0(1-r)}\,\delta,\qquad
 \kappa_\sigma=\frac{\lambda_{\max}(\sigma)}
                         {\lambda_{\min}(\sigma)}.}
 \label{sel38:eq:log-error}
\end{equation}
The Hamiltonian coefficient error is at most this bound divided
by $\sqrt{2n}$; the error in
$\sqrt{\Delta_{\rm Ham}}$ is at most the bound itself.
\end{proposition}
\begin{proof}
Use the convergent power series for $\operatorname{Log}(I+X)$.
The telescoping identity for $X^k-Y^k$ gives a Frobenius
bound $kr^{k-1}\|X-Y\|_{\mathrm F}$, so the logarithm
difference is at most $\delta/(1-r)$.
Reversal has Frobenius operator norm at most $\kappa_\sigma$.
The half-difference, division by the physical duration and the
inherited orthogonal-projection bounds prove the result.
\end{proof}

Uncertainty in $\sigma$, a spectrum near the logarithm cut, or an
unverified rank/equality sector needs additional error information.
No small numerical defect is declared to be an exact zero.
The local bound need not apply to the algebraic qutrit point; that
point is evaluated by its exact spectral decomposition instead.

\paragraph{Proved in this continuation.}
The finite equality sector of a contraction semigroup is a reducing
unitary sector, giving a branch-independent homogeneous-embedding
obstruction. The unchanged V34 pulse bank fails it with exact residual
$288/289$, even though its V37 circulation is Hamiltonian.
The explicit positive qutrit semigroup exhibits the converse:
Hamiltonian generator asymmetry and reversible depolarization yield
finite circulation defect $1/8$ at the stated point.
Its true coefficient is recovered by the positive principal logarithm.
A fixed grid cannot identify the initial physical rate without a
subdivision law. Two actual durations and explicit generator bounds
supply a signed second-order estimator. A held-out action comparison
separates non-Hamiltonian response from an independent coefficient
mismatch.

\paragraph{Inherited rather than repeated.}
The conditional-Choi generator criterion, canonical Hamiltonian
projection, stationary reverse, source-compatible parity and contextual
tomography are retained. General embedding and principal-logarithm
machinery is background \cite{sel38-wolf}; the same-memory source
obstruction and the evaluated distinctions above are the new interfaces.
V19/V20's coherent action aliases and V36's original two-phase clock
are not replaced by a homogeneous-time assumption.

\paragraph{Not established.}
No historical Standard-Model channel, generator or typed odd
determinant--Clifford word is populated here. The qubit circulation is
not promoted to colour or a generation carrier. An embeddable channel
does not establish that its historical implementation is homogeneous,
and a mathematical logarithm is not automatically the common
variational action. Exact identification with a physical grading and
the directed typing remain source conditions. Stationarity and
semigroup embeddability are not new axioms of the finite commutant
duality.

\paragraph{Next source-level decision.}
Before comparing any recovered coefficient with a common action,
identify what the physical table describes: one renewal event, a
fixed-duration propagator, or a calibrated small-event family.
For a proposed same-memory homogeneous interpretation, test the
equality-sector obstruction before solving for a logarithm. If an
independently admissible logarithm branch passes the generator cone,
compare its antisymmetric generator and its canonical coefficient
with the actual action. For a measured family, use the physical
derivative or a certified subdivision estimate. For the purely finite
structural claim, preserve the original labels and use V29's
same-label search; none of these time-embedding requirements replaces
that task.

\paragraph{Verification scope.}
The executable checks verify the exact Pauli transfer matrix,
stationary uniqueness, singular subspace and its noninvariance,
two-step purity, the algebraic qutrit channel and its canonical
residual, the analytic logarithm's explicit invariant-space action,
conditional positivity, parity, the sampling coefficients and the
held-out Pythagorean identity. They use no floating-point logarithm
or numerical rank threshold to prove nonembedding.
The written proofs establish the all-branch and all-time assertions.
No physical data collection or proof-assistant certification is claimed.

\addtocontents{toc}{\protect\endgroup}
\renewcommand{\refname}{Additional background for V38}
\begin{thebibliography}{99}
\raggedright
\bibitem[59]{sel38-wolf}
M.~M.~Wolf, J.~Eisert, T.~S.~Cubitt and J.~I.~Cirac,
\emph{Assessing non-Markovian quantum dynamics},
Phys. Rev. Lett. \textbf{101}, 150402 (2008).
\href{https://doi.org/10.1103/PhysRevLett.101.150402}
{doi:10.1103/PhysRevLett.101.150402};
\href{https://arxiv.org/abs/0711.3172}{arXiv:0711.3172}.
Cited for the established channel-embedding and logarithm/conditional-
positivity framework, not as a proof of the source-specific calculations
in this continuation.
\end{thebibliography}

% END OF SOURCE-SELECTION V38 DEVELOPMENT BLOCK.
% Preserve V1--V38 and append subsequent work before the final terminator.


\clearpage
\hypersetup{pdftitle={Historical Standard-Model Source Selection: V1--V39}}
\begin{center}
{\Large\bfseries Source-selection continuation V39}\\[.4em]
{\large Target-visible historical search and the limit of flatness:\\
minimal response coordinates, actual-word witnesses, and joint constraints}
\end{center}
\addtocontents{toc}{\protect\begingroup}
\addtocontents{toc}{\protect\renewcommand{\protect\numberline}{\protect\makebox[2.1em][l]}}

\section[Return to the finite structural occurrence problem]{The structural search does not require another continuous-time assumption}
\label{sel39:context}

V38 separates a finite event, a fixed-duration propagator, and a calibrated
family. None of its additional dynamical tests is needed for the finite
commutant duality. We therefore return to the actual-word problem of
Theorem~\ref{sel29:thm:flat}. Its output is a word with a nonzero specified
Choi cross block, not necessarily a word satisfying every additional typing,
covariance, or finite-Dirac condition.

Two questions must now be separated. First, how much of the represented
source must be retained to decide whether that cross block ever occurs?
Second, does flatness also bound the first word at which independent exact
zero constraints hold together with the nonzero cross block? The first has
a positive minimal-response answer. The second has a negative answer even
for one unitary event on a six-dimensional carrier: the Liouville span can
be flat at depth twelve while the first word with the required exact group
appears at an arbitrarily larger depth.

This is a scope refinement, not a reversal of V29. A linear output is nonzero
on a span exactly when it is nonzero on a spanning word. A vector satisfying
several exact equalities need not be one of those spanning words. V29's proof
uses the former fact and does not establish the latter. We give a finite
positive selection theorem when the linear side constraints hold throughout
the actual admitted language, and state explicitly what is not decided when
they do not.

Reachable/observable quotients and Hankel minimality are established linear
realization theory \cite{sel39-kiefer}. They already occur in the Grand
Tensor and the duality companion's response-minimality results. The new
applications here are the Choi-cross target quotient, an exact six-coordinate
calculation on the unchanged pulse bank, a complete-source covariance gate,
and the fixed-dimension delayed-selection family. No minimal linear
coordinate space is assigned an unproved positive cone or a quantum-memory
interpretation.

\section[The smallest linear description of the specified Choi cross]{Discard only what every target continuation fails to see}
\label{sel39:minimal}

Fix an actually identified recurrent memory $\C^d$, a finite bank of resolved
CP trace-nonincreasing maps $\Phi_a$, and the actual admissible word language.
For the moment every concatenation is admitted. Use V29's Liouville matrices
$S_w$, vectors $z_w=\operatorname{vec}S_w$, seed $z_0=\operatorname{vec}I_{d^2}$,
and append maps $L_a z_w=z_{aw}$. Retain the same identified Choi-sector
projections $P_B,P_D$. In fixed orthonormal output coordinates put
\[
 C z=\operatorname{vec}\bigl(P_B\mathfrak R(S)P_D\bigr),
 \qquad z=\operatorname{vec}S,
 \qquad y(w)=Cz_w,\qquad h_w=\|y(w)\|_2^2.
\]
The map $C$ is an analysis map. Its application is not a newly admitted
physical operation. Define
\begin{align}
 \mathcal Z&=\operatorname{span}\{z_w:w\text{ admitted}\},
 &R&=\dim_{\C}\mathcal Z,\notag\\
 \mathcal U_C&=\{z\in\mathcal Z:C L_vz=0\text{ for every }v\},
 &\mathcal X_C&=\mathcal Z/\mathcal U_C,
 &r_C&=\dim_{\C}\mathcal X_C.
 \label{sel39:eq:quotient}
\end{align}
Here $L_v=L_{a_m}\cdots L_{a_1}$ for chronological append order
$v=a_1\cdots a_m$. Accordingly $y(v\mid u)=CL_vL_uz_0$ means that the
history $u$ is followed by $v$.

\begin{theorem}[Minimal linear realization of the same-label cross response]
\label{sel39:thm:minimal}
The subspace $\mathcal U_C$ is invariant under every $L_a$. The quotient
$\mathcal X_C$ has well-defined induced append maps, seed and output, and
reproduces $y(w)$ for every original word. With output coordinate $j$, form
\[
 \mathbb H_C[(v,j),u]=(CL_vL_uz_0)_j.
\]
Then
\begin{equation}
 \boxed{\operatorname{rank}\mathbb H_C=r_C\le R.}
 \label{sel39:eq:hankel}
\end{equation}
Every complex-linear realization of this complete vector-valued response has
dimension at least $r_C$. The quotient is unique up to invertible linear
coordinates among reachable, observable realizations. If the response is not
identically zero, an actual word with $h_w>0$ has length at most $r_C-1$.
\end{theorem}
\begin{proof}
If $z\in\mathcal U_C$, then $CL_vL_az=0$ for every $v$, since $a$ followed
by $v$ is an admitted continuation. This proves invariance. The output is
well defined because the empty continuation is included. Quotienting loses
no output and is reachable by construction.

Define $\mathcal O:z\mapsto(CL_vz)_{v}$. Its kernel on $\mathcal Z$ is
exactly $\mathcal U_C$. The columns of $\mathbb H_C$ are $\mathcal O(z_u)$,
and the $z_u$ span $\mathcal Z$, proving \eqref{sel39:eq:hankel}. For any
other realization, the same matrix factors through its reachable state
space, giving the lower bound. Sending a reachable vector to its complete
future-output function identifies the reachable/observable quotient with
$\operatorname{Ran}\mathcal O$, proving uniqueness.

If $r_C>0$, the quotient seed is nonzero. Its reachable word spans start at
dimension one and either grow or become permanently invariant. They reach
the whole $r_C$-dimensional quotient by depth $r_C-1$. If $C$ vanished on
all those spanning words, it would vanish on the whole reachable quotient,
contradicting a nonzero response. Thus an actual spanning word has nonzero
output. This argument seeks a nonzero output only; it imposes no additional
zero constraint on that word.
\end{proof}

\begin{proposition}[Finite construction and a positive all-depth statistic]
\label{sel39:prop:gramian}
Restrict $L_a$ to $\mathcal Z$ and choose an orthonormal coordinate basis there.
Set
\[
 W_0=C^*C,\qquad W_{r+1}=C^*C+\sum_aL_a^*W_rL_a.
\]
Then
\begin{equation}
 W_r=\sum_{|v|\le r}L_v^*C^*CL_v\succeq0,
 \qquad \mathcal U_C=\ker W_{R-1}.
 \label{sel39:eq:gramian}
\end{equation}
A plateau in the associated observation row span proves permanence; a small
numerical eigenvalue is not such a proof. If
$0<q<1/M$ with $M=\sum_a\|L_a\|_{2\to2}^2$, the convergent equation
\[
 W(q)=C^*C+q\sum_aL_a^*W(q)L_a
\]
has the unique solution
\begin{equation}
 \boxed{z_0^*W(q)z_0=\sum_wq^{|w|}h_w\ge0.}
 \label{sel39:eq:generating}
\end{equation}
It is zero exactly when every actual word has zero cross block. The tail
beyond length $r$ is at most
$\|C\|^2\|z_0\|^2(qM)^{r+1}/(1-qM)$.
\end{proposition}
\begin{proof}
Expanding the recursion gives \eqref{sel39:eq:gramian}; positivity identifies
its kernel with the common kernel of all the listed rows. The row spaces
$\operatorname{span}\{\text{rows of }CL_v:|v|\le r\}$ form an increasing
chain in an $R$-dimensional dual space. A plateau is invariant under right
multiplication by each $L_a$. If $C\ne0$ the initial row rank is at least
one, so permanence occurs by $R-1$; if $C=0$ the conclusion is immediate.
The map $W\mapsto q\sum_aL_a^*WL_a$ has operator norm at most $qM<1$.
Its geometric series proves uniqueness, positivity and
\eqref{sel39:eq:generating}. Bound the length-$j$ terms using
$\sum_{|v|=j}\|L_v\|^2\le M^j$ and sum the geometric tail.
\end{proof}

The damping variable $q$ is a freely chosen convergence parameter in an
analysis formula, not a new waiting-time law. Neither $W(q)$ nor the
quotient predictor is asserted to be an executed instrument. No physical
acceptance weight is removed from $S_a$.

\paragraph{Actual finite-state typing.}
If the retained grammar includes a deterministic finite classical state,
use $e_x\otimes z_w$, append only through its actual transitions, and set the
output to zero off the specified terminal types. An incompatible word then
has zero vector. The same proofs apply on this enlarged linear space and
return admissible words. The classical state is counted. An arbitrary
history-dependent admission rule is not silently replaced by a finite-state
rule; the relevant finite representation must be established from the source.

\section[An exact six-coordinate audit of the unchanged pulse bank]{A smaller target model does not change the instrument}
\label{sel39:pulse}

Retain precisely the original coefficients of
Proposition~\ref{sel37:prop:old-forward} and Theorem~\ref{sel38:thm:old-source}, written
explicitly to avoid any source-identification ambiguity:
\[
 k_X=(9I+12iX)/25,\qquad k_Z=(12I+16iZ)/25.
\]
For this computational audit alone choose the two coefficient rays
$b=I/\sqrt2$ and $d=Y/\sqrt2$. They are test sectors, not a physical
Clifford/determinant identification. Let
$\chi(w)=\langle\!\langle b|J_w|d\rangle\!\rangle$.
In the orthonormal Pauli operator basis, each $S_a$ has the form
$\operatorname{diag}(p_a,M_a)$, where
\begin{align}
 p_X&=9/25,&
 M_X&=\frac1{625}
 \begin{pmatrix}225&0&0\\0&-63&216\\0&-216&-63\end{pmatrix},\notag\\
 p_Z&=16/25,&
 M_Z&=\frac1{625}
 \begin{pmatrix}-112&384&0\\-384&-112&0\\0&0&400\end{pmatrix}.
 \label{sel39:eq:pulses}
\end{align}
These factors include the actual branch probabilities.

\begin{theorem}[Exact target rank six, with a two-event witness]
\label{sel39:thm:pulse}
The unchanged instrument's full Liouville word span has dimension ten. Its
complete $\chi$-response has minimal complex-linear dimension six. It has
the explicit realization
\begin{equation}
 x_w=(M_we_1,M_we_3),\quad
 x_0=(e_1,e_3),\quad
 x_{aw}=\operatorname{diag}(M_a,M_a)x_w,
 \label{sel39:eq:sixmodel}
\end{equation}
with output
\begin{equation}
 y(w)=\tfrac12\bigl(e_1^{\mathsf T}M_we_3-e_3^{\mathsf T}M_we_1\bigr),
 \qquad \chi(w)=i\,y(w),\qquad h_w=y(w)^2.
 \label{sel39:eq:pulseout}
\end{equation}
The empty and one-event words have zero output. In chronological notation,
\begin{equation}
 \boxed{
 y(XZ)=\frac{41472}{390625},\qquad
 y(ZX)=-\frac{41472}{390625},\qquad
 h_{XZ}=h_{ZX}=\frac{1719926784}{152587890625}.}
 \label{sel39:eq:pulse-witness}
\end{equation}
Thus two events are the shortest witnesses for these specified test rays.
\end{theorem}
\begin{proof}
Every word coefficient is a positive real weight times an $SU(2)$ matrix,
so it has the quaternion form
$aI+i(xX+yY+zZ)$ with real coefficients. Direct multiplication gives
$\chi=-2iay$ and
$(M_{13}-M_{31})/2=-2ay$, proving \eqref{sel39:eq:pulseout} without a
chosen outcome phase. Appending a map multiplies both columns in
\eqref{sel39:eq:sixmodel} by the same $M_a$, giving the dimension-six upper
bound.

For completeness the lower bound is an exact finite probability-response
certificate, not a floating-point rank estimate. Use the chronological
prefix list and continuation list
\[
 \mathcal P=(\varnothing,X,Z,XX,XZ,ZX),\qquad
 \mathcal V=(\varnothing,X,Z,XX,ZX,XZ).
\]
The matrix $H_{v,u}=y(uv)$, with $u$ followed by $v$, has determinant
\begin{equation}
 \det H=
 \frac{136728756748667408445475367864500224}
 {867361737988403547205962240695953369140625}\ne0.
 \label{sel39:eq:sixminor}
\end{equation}
The exact checker exports the matrix and verifies this determinant by rational
arithmetic. Theorem~\ref{sel39:thm:minimal} proves minimality. For the full
Liouville span, every word belongs to $\C\oplus M_3$, giving the upper
bound ten. The ten chronological words
\[
 \varnothing,X,Z,XX,XZ,ZX,ZZ,XXZ,XZX,XZZ
\]
have linearly independent Liouville matrices, as certified by the exported
exact rational row reduction. This gives the matching lower bound.

Finally,
$k_Zk_X=(108I+144iX-192iY+144iZ)/625$,
while the opposite order changes the sign of its $Y$ term. Their $I$ and
$Y$ components give \eqref{sel39:eq:pulse-witness}. Neither primitive
coefficient has a $Y$ component, proving shortest length two.
\end{proof}

The new rank is not a six-dimensional quantum memory, an internal algebra,
or a particle count. It is the minimum linear predictor of one specified
coefficient cross response. The original instrument and all source weights
are unchanged. Its V38 homogeneous-embedding obstruction remains valid but
plays no role in this finite calculation.

\section[Nonzero coherence and exact typing have different stopping rules]{Why a flat span need not bound the first fully compatible word}
\label{sel39:delay}

The following family is a separate comparison, not a historical source.
Its purpose is to test an algorithmic inference. The determinant character
is explicitly present in its declared representation; no emergence of that
representation is claimed.

Let $G_0=U(3)\times U(2)$, $\chi(A,B)=\det A\det B$, and
$H=\ker\chi$. On a six-dimensional carrier set
\[
 \mathcal H=(\C n\oplus\C\omega)\oplus(C\oplus\C r),
 \qquad R(A,B)=1\oplus\chi(A,B)\oplus A\oplus1.
\]
Write $(c_1,c_2,c_3)$ for an orthonormal basis of $C$. For any integer
$m\ge7$, put $\alpha=\pi/(4m)$ and define the one-event unitary
\begin{equation}
 U_m=e^{-i\alpha X_{n\omega}}\oplus
       e^{-4i\alpha X_{c_1r}}\oplus I_{\operatorname{span}\{c_2,c_3\}},
 \qquad \Phi_m=\operatorname{Ad}_{U_m}.
 \label{sel39:eq:delayU}
\end{equation}
The direct sums denote the specified orthogonal subspaces. Let
$b=I_{C\oplus\C r}/2$ and $d=|\omega\rangle\langle n|$, both unit
Hilbert--Schmidt vectors. They transform by the trivial character and
$\chi$, respectively.

\begin{theorem}[Fixed-dimensional flatness, arbitrarily delayed exact selection]
\label{sel39:thm:delay}
For the single repeated event \eqref{sel39:eq:delayU}:
\begin{enumerate}[label=\textup{(\roman*)},leftmargin=2.8em]
\item The ordinary coefficient word span has dimension five. The Liouville
word span has dimension thirteen and is already flat at depth twelve.
\item The first strictly positive word length $k$ for which
$\Phi_m^k$ is covariant under $H$ is exactly $k=2m$.
\item At that length,
\[
 \operatorname{Cov}_{G_0}(\Phi_m^{2m})=H,
 \qquad |\langle\!\langle b|J(\Phi_m^{2m})|d\rangle\!\rangle|^2=4.
\]
The cross block is already nonzero at $k=1$, while $H$ covariance fails.
\end{enumerate}
Consequently neither the Hilbert dimension, alphabet size nor flat Liouville
rank bounds the first word simultaneously satisfying exact covariance and
nonzero coherence. This assertion does not exclude bounds that use the
actual eigenvalue orders or other additional source data.
\end{theorem}
\begin{proof}
Let $\zeta=e^{i\alpha}$, of order $8m$. The distinct eigenvalues of $U_m$
are $\zeta^j$ for $j\in\{0,1,-1,4,-4\}$. They are distinct, so its
minimal polynomial has degree five. The distinct eigenvalues of
$U_m\otimes\overline U_m$ have exponents
\[
 \mathcal D=\{0,\pm1,\pm2,\pm3,\pm4,\pm5,\pm8\}.
\]
For $m\ge7$ these thirteen exponents are distinct modulo $8m$. Since the
superoperator is diagonalizable, its minimal polynomial has degree thirteen.
Thus $I,S,\ldots,S^{12}$ are independent and $S^{13}$ lies in their span.
This proves (i) by an exact spectral argument.

If a unitary channel $\operatorname{Ad}_V$ is covariant under the connected
subgroup $SU(3)\subset H$, then
$R(g)VR(g)^*=c(g)V$ for a continuous character $c:SU(3)\to U(1)$.
This character is trivial: its Lie-algebra derivative vanishes on
$[\mathfrak{su}(3),\mathfrak{su}(3)]=\mathfrak{su}(3)$, and the group is
connected. Hence $V$ must commute with this $SU(3)$.

In $U_m^k$, the $c_1,r$ subblock has angle $k\pi/m$, while $c_2,c_3$
are fixed. Commutation with all $SU(3)$ requires its $c_1,r$ off-diagonal
entry to vanish and its $c_1$ diagonal entry to equal one. The exact
condition is therefore $e^{ik\pi/m}=1$, or $2m\mid k$. At $k=2m$ the
whole complementary block is the identity and the $n,\omega$ block is
$-iX_{n\omega}$. Thus it commutes with $H$.

For $g\in G_0$, a covariance can change $U_m^{2m}$ only by a common scalar.
Its identity action on the invariant four-dimensional complement forces that
scalar to be one. The off-diagonal $n,\omega$ entries then force $\chi(g)=1$.
This proves exact group equality. The neutral coefficient on $b$ is two and
the coefficient on $d$ is $-i$, giving cross mass four.

At arbitrary $k$, these two coefficients are
$1+\cos(k\pi/m)$ and $-i\sin(k\pi/(4m))$. Both are nonzero at $k=1$.
Choosing $m$ arbitrarily large proves the claimed lack of a rank-only
simultaneous-occurrence bound. No linear combination of powers is promoted
to an actual event.
\end{proof}

At $m=8$, the Liouville ranks are $1,2,\ldots,13,13,\ldots$, but the
first exact-group word has length sixteen. The checker verifies the spectral
orders in the exact cyclotomic field $\mathbb Q(\zeta_{64})$, not by numerical
angle rounding. In particular, existence of a symmetry-compatible vector
in the flat span is not a certificate that any chosen spanning word is
symmetry compatible. The nonzero-only statement of
Theorem~\ref{sel29:thm:flat} is unaffected.

\section[A valid finite selector for a source-certified covariant language]{Certify the side conditions on the language, not on a fitted word}
\label{sel39:safety}

Let $R(g)$ be the already identified comparison representation on the same
input/output memory and let $H\subset G_0$ be a connected compact subgroup.
For a Hermitian Lie-algebra generating bank $T_\ell$ on that memory, write
$\widehat T_\ell=T_\ell\otimes I-I\otimes\overline T_\ell$.
Define the linear defect output
\[
 D_Hz=\bigl([\widehat T_\ell,\mathfrak R(S)]\bigr)_\ell,
 \qquad z=\operatorname{vec}S.
\]
It vanishes exactly when the corresponding map is $H$ covariant. If required,
a fixed linear Choi-support condition $(I-P_{\rm type})J=0$ may be added as
another output. This does not encode nonlinear rank, isometry or numerical
coupling conditions.

\begin{theorem}[Universal side-condition gate and an actual-word selector]
\label{sel39:thm:selector}
Let $\mathcal Z$ be the reachable span for the actual finite-state typed
language, with defect and cross outputs read only at its designated terminal
types. The following statements are finite and exact.
\begin{enumerate}[label=\textup{(\roman*)},leftmargin=2.8em]
\item Every admitted terminal word passes the linear side conditions if and
only if $D_H$ (and the added linear defect rows) vanishes on a retained
actual-word basis of $\mathcal Z$. A failure supplies an actual violating
word, not merely an unphysical linear combination.
\item Suppose the gate passes. If the cross response is nonzero, a word of
length at most $r_C-1$ has a nonzero cross block and also passes all the
certified linear side conditions.
\item In particular let $H=\ker\chi$, and suppose the identified Choi-sector
projections $P_B,P_D$ transform pointwise by the trivial character and
$\chi^{\pm1}$. Then every returned word in (ii) has
\[
 \operatorname{Cov}_{G_0}(\Phi_w)=\ker\chi.
\]
\end{enumerate}
The subgroup-covariance gate is a sufficient entrance to this finite selector,
not a necessary property of a source that happens to have one suitable word.
\end{theorem}
\begin{proof}
The first assertion is linearity on a span of actual words. Finite reachability
and observation closure are those of Theorem~\ref{sel39:thm:minimal}. The
second assertion follows from its word bound because the gate holds for every
admitted terminal word. For (iii), covariance under $H$ is already certified.
Every covariance under $G_0$ must preserve the specified nonzero cross block.
It transforms by the scalar $\chi^{\pm1}$, so preservation forces $\chi=1$.
The two group inclusions give equality. No other source components are dropped;
the full-Choi covariance gate accounts for them.
\end{proof}

For an unrestricted language, covariance of every primitive CP map under the
same group already implies covariance of every composition. The typed-language
form also covers sources for which only designated complete words have that
covariance. Known symmetry identities may prove the gate symbolically; it
need not be estimated from a large table if it is an actual exact source law.

If the gate fails, neither its violating word nor an arbitrary first
coherence word settles the existential joint problem. The delayed family
proves that terminating V29 at its first flat layer would then be an
invalid completeness claim for exact-group selection. The current theorem
adds no such claim. An independently verified restricted language, an exact
symmetry law, or an additional constructive search argument is needed.

\section[Historical data and numerical margins retain their scope]{What can be reconstructed without full channel tomography}
\label{sel39:data}

Suppose the particular real and imaginary entries of $y(w)$ belong to the
admitted contextual tester span of V24/V25. Then there is a fixed matrix $B_C$
with
\[
 y(w)=B_Cp(w),\qquad h_w=\|B_Cp(w)\|_2^2.
\]
Only these target functionals are needed; the full Choi matrix need not be
reconstructed. Their contextual availability and same-carrier meaning still
have to be established. The entries of $\mathbb H_C$ are obtained by applying
the same postprocessing to the actual concatenated words, not by multiplying
their marginal outcome probabilities.

\begin{proposition}[A finite data realization and its exact-data boundary]
\label{sel39:prop:data}
If the complete target Hankel rank is independently certified as $r_C$, choose
$r_C$ independent prefix columns and $r_C$ independent continuation-coordinate
rows with invertible intersection $H_0$. Let $H_a$ be the same table with one
actual event $a$ between each chosen prefix and continuation. In the prefix
coordinate basis the transition is
\[
 T_a=H_0^{-1}H_a.
\]
The corresponding seed and output coordinates reconstruct $y(w)$ for every
word. These are generally not positive coordinates. An invertible finite
minor alone proves only a lower bound on the complete rank; the source's
finite closure or an independent dimension bound is needed for completeness.
\end{proposition}
\begin{proof}
Let $B$ collect the chosen quotient prefix states and $O$ the chosen future
functionals. Then $H_0=OB$ and $H_a=O\overline L_aB$. Independence makes
$B$ a basis and $O$ invertible on that basis, so
$\overline L_a B=BT_a$. The empty-prefix and empty-future equations give the
seed and output. Repeated multiplication proves the complete response
identity. A finite minor does not restrict unobserved reachable directions,
which explains the stated exact-data boundary.
\end{proof}

For $\|\widehat p-p\|_2\le\epsilon$, an actual word has the one-sided bound
\begin{equation}
 h_w\ge
 \left(\max\{0,\|B_C\widehat p(w)\|_2-\|B_C\|\epsilon\}\right)^2.
 \label{sel39:eq:error}
\end{equation}
This certifies a positive cross, not an exact vanishing of the side-condition
residual. Noisy Hankel plateaux likewise do not prove exact infinite-depth
zeros. All quantities use unnormalized branch maps. If every accepted event
carries an actual scalar factor $\vartheta$, then a word of length $k$ has
$y_{\rm raw}(w)=\vartheta^k y_{\rm acc}(w)$ and
$h_{\rm raw}(w)=\vartheta^{2k}h_{\rm acc}(w)$. The nonzero ranks are unchanged
for $\vartheta>0$, but numerical margins are not.

\section[Source status after the target-specific audit]{What is proved and what still needs the historical source}
\label{sel39:status}

The finite structural search has acquired a target-minimal realization. It
retains exactly the coordinates required for all future values of the
specified same-label Choi cross. The unchanged pulse bank gives an exact
six-coordinate realization and a shortest two-event witness. This is a new
calculation on the existing maps, not a new physical instrument or a claim
that its Pauli test rays are determinant incidence.

The delayed six-dimensional unitary family proves a different limitation:
flatness decides a nonzero linear output, but does not bound when independent
exact zero conditions and that nonzero output first occur on the same actual
word. This limits a possible extrapolation of V29; it does not invalidate its
stated cross-nonvanishing theorem. The universal-covariance/typing gate gives
a correct finite positive selector where the whole source language satisfies
the relevant linear conditions. It can be checked by old moment/cylinder data
or proved from an actual symbolic source law.

For the historical Standard-Model programme, the grading action, represented
type maps, original CP grammar and their common-cut ancestry remain required.
Once identified, compute the target quotient and audit the full-source
linear side conditions before promoting a cross witness to an exact-group
word. Nonlinear Clifford-shape, finite-Dirac and numerical normalization
conditions remain separately evaluated on the same word. No historical
cylinder vector for that word has been populated in this continuation.

No additional stationarity, homogeneous embedding or continuum assumption has
been added to the finite duality. The previous dynamical results retain their
own scope. The exact scripts below certify rational ranks and moments for the
unchanged bank, the cyclotomic delay example, target-observability closure,
and the stated group/word distinctions. They supplement the written proofs;
they are neither physical measurements nor proof-assistant certification.

\addtocontents{toc}{\protect\endgroup}
\renewcommand{\refname}{Background for V39}
\begin{thebibliography}{99}
\raggedright
\bibitem[60]{sel39-kiefer}
S.~Kiefer, \emph{Notes on equivalence and minimization of weighted automata}
(2020), \href{https://arxiv.org/abs/2009.01217}{arXiv:2009.01217}.
Cited for the established reachable/observable and Hankel-minimality framework,
not as a source of the Choi-cross calculations or delayed-selection example.
\end{thebibliography}

% END OF SOURCE-SELECTION V39 DEVELOPMENT BLOCK.
% Preserve V1--V39 and append subsequent work before the final terminator.


\clearpage
\hypersetup{pdftitle={Historical Standard-Model Source Selection: V1--V40}}
\begin{center}
{\Large\bfseries Source-selection continuation V40}\\[.4em]
{\large Exact word constraints beyond flatness:\\
arithmetic exclusion, finite symmetry-return languages, and retained records}
\end{center}
\addtocontents{toc}{\protect\begingroup}
\addtocontents{toc}{\protect\renewcommand{\protect\numberline}{\protect\makebox[2.1em][l]}}

\section[The next finite task is a language, not another coherence witness]{Move upstream to the exact constraints on actual words}
\label{sel40:context}

V39 establishes two facts that must be used together. Its target quotient
finds a nonzero same-label cross in finitely many linear coordinates, but
its delayed example excludes any bound, based only on flat Liouville rank,
on the first word passing an additional exact covariance constraint.
The universal side-condition gate of Theorem~\ref{sel39:thm:selector} is
sufficient but can fail even when a suitable word exists. The missing object
is therefore the language of original words passing that constraint, not
another choice of a favourable vector in the flat span.

This continuation gives two concrete resolutions. On the unchanged two-pulse
bank of Theorem~\ref{sel39:thm:pulse}, reduction of its exact numerator
matrices modulo five proves that \emph{no nonempty word} preserves the
already reconstructed $Y$-grading as a fixed-label CP operation. It also
excludes covariance under the full circle generated by $Y$. Thus the
nonzero two-event cross in that example cannot be promoted to that symmetry
claim at any later word length. This is a source-specific all-depth result,
not a claim about the unknown historical matter alphabet.

In the positive direction, an exactly closed finite orbit of transported
Lie generators recognizes all covariance-compatible words. Intersecting
that language with V39's target response gives a complete finite test for
simultaneous covariance and nonzero coherence. The construction is applied
to the unchanged delayed family: its covariance language has $2m$ states,
and the covariance-filtered cross has exact Hankel rank $4m$, while the
unfiltered Liouville span still has rank thirteen. This identifies the
additional record complexity hidden by an ordinary rank-only test.

The finite-automaton product and linear minimization are established
realization methods \cite{sel39-kiefer}. All specialized arithmetic,
source-symmetry, rank and probability statements below are proved here.
The automaton states are classical analysis coordinates computed from the
retained labels; no new physical controller, pulse, inverse operation or
clock is added. The V39 mathematical body is retained unchanged.

\section[Finite residue certificates for exact word constraints]{A finite necessary language can prove an all-depth exclusion}
\label{sel40:modular}

Let $\mathcal R$ be a commutative subring of $\C$, and let
$\varphi:\mathcal R\to\mathbb F$ be a unital homomorphism to a finite field.
Suppose exact source matrices have the form
\[
 S_a=c_a A_a,\qquad c_a\in\C\setminus\{0\},\qquad
 A_a\in M_N(\mathcal R).
\]
The matrices $S_a$ may be Liouville matrices, or coefficient matrices when
an independently proved covariance criterion is expressed on coefficients.
For chronological $w=a_1\cdots a_k$ put
$A_w=A_{a_k}\cdots A_{a_1}$. Suppose the required side condition implies
\[
 f_j(A_w)=0\quad(1\le j\le s),
\]
where the $f_j$ are fixed homogeneous matrix-entry polynomials over
$\mathcal R$. A finite disjunction of such systems is allowed. Homogeneity
makes nonzero rescaling of a word irrelevant to these exact zeros.

\begin{theorem}[Finite-field necessary-language certificate]
\label{sel40:thm:modular}
Let $\mathcal M$ be the finite matrix monoid generated by
$\overline A_a=\varphi(A_a)$ and the identity. The deterministic update
\[
 q\longmapsto\overline A_a q,\qquad q_0=I,
\]
recognizes the regular language on which the reduced polynomial conditions
hold. Every actual word passing the characteristic-zero side condition
belongs to that language. In particular, if no reachable nonempty word has
an accepted residue, the side condition is impossible for all nonempty
original words.

A modular zero is only a necessary condition. It does not certify the
original equality, positivity, normalization, or physical covariance.
\end{theorem}
\begin{proof}
Entrywise reduction respects sums and products, so the state after $w$ is
$\varphi(A_w)$ and $\varphi(f_j(A_w))=\overline f_j(\varphi(A_w))$.
A zero in $\mathcal R$ has zero image. This proves the implication for each
system and for a finite disjunction. There are at most $|\mathbb F|^{N^2}$
matrices, so exact reachability closure terminates. A separate two-state
empty/nonempty flag can be included when the identity residue is reached
again by a nonempty word. No converse to reduction is used.
\end{proof}

One may clear the denominators of each primitive matrix \emph{before}
reduction. The chosen denominator may vanish in the residue field: it is
never inverted there. One applies the homomorphism to $A_a$, not to an
undefined reduction of $S_a$. Likewise, a homomorphism sending $i$ to a
finite-field square root of $-1$ need not preserve complex conjugation.
Adjoints appearing in a source identity are expanded algebraically before
reduction. The residue matrices themselves are not quantum operations.

\begin{lemma}[Covariance tests for one invertible coefficient]
\label{sel40:lem:coefficient}
Let $\Psi_K(X)=KXK^*$ with $K$ invertible.
For a self-adjoint unitary $Z$,
\[
 \operatorname{Ad}_Z\Psi_K\operatorname{Ad}_Z=\Psi_K
 \quad\Longleftrightarrow\quad ZKZ=K\ \hbox{or}\ ZKZ=-K.
\]
Let $R:G\to U(d)$ be a continuous representation of a connected compact
group, with Hermitian Lie generators $T_\ell$. Then
\[
 \Psi_K\text{ is }G\text{-covariant}
 \quad\Longleftrightarrow\quad [K,T_\ell]=0\quad\hbox{for all }\ell.
\]
\end{lemma}
\begin{proof}
Equality of the nonzero rank-one Choi matrices implies
$ZKZ=zK$, $|z|=1$. Applying the involution again gives $z^2=1$.
For the connected group, covariance implies
$R(g)KR(g)^*=z(g)K$. The phase is continuous, as follows by taking its
Hilbert--Schmidt inner product with $K$ and dividing by $\|K\|_{\HS}^2$.
Taking determinants gives $z(g)^d=1$. Connectedness and $z(1)=1$ force
$z(g)=1$. Commutation with the group is equivalent to commutation with its
Lie generators. The converses follow by substitution.
\end{proof}

The invertibility assumption in the connected-group statement is material:
a singular coefficient can transform by a nontrivial character while its
CP square is covariant. It is not discarded merely because all source
matrices act on a finite-dimensional carrier.

\section[An all-depth result for the unchanged two-pulse bank]{Its reconstructed reverser never becomes a fixed-label word symmetry}
\label{sel40:pulses}

Keep exactly the original instrument
\[
 k_X=(9I+12iX)/25,\qquad k_Z=(12I+16iZ)/25,
 \qquad k_X^*k_X+k_Z^*k_Z=I.
\]
Write $k_X=(3/5)U_X$, $k_Z=(4/5)U_Z$, with
\[
 U_X=(3I+4iX)/5,\qquad U_Z=(3I+4iZ)/5.
\]
For a word of length $k$ define its integer Gaussian numerator
\[
 M_w=A_{a_k}\cdots A_{a_1},\qquad
 A_X=3I+4iX,\quad A_Z=3I+4iZ,
 \qquad U_w=5^{-k}M_w.
\]
Its original, input-independent probability is
\[
 p_w=(9/25)^{n_X(w)}(16/25)^{n_Z(w)},\qquad
 \Phi_w=p_w\operatorname{Ad}_{U_w}.
\]
All these factors remain in every CP word.

\begin{theorem}[No nonempty $Y$-homogeneous word]
\label{sel40:thm:pulse-nogo}
For every nonempty word in the original alphabet $\{X,Z\}$,
\begin{equation}
 \boxed{YU_wY\ne U_w,\qquad YU_wY\ne-U_w.}
 \label{sel40:eq:nohomogeneous}
\end{equation}
Consequently no such resolved word is preserved by
$\operatorname{Ad}_Y$, and none is covariant under
$\{e^{-itY}:t\in\R\}$. The empty word has both symmetries but zero
cross between $I/\sqrt2$ and $Y/\sqrt2$.

This is compatible with the uniquely reconstructed \emph{adjoint-reversal}
of V34: an adjoint-reversal relation does not say that an original CP label
is fixed. It also leaves V39's two-event coherence value unchanged.
\end{theorem}
\begin{proof}
Reduce $\mathbb Z[i]$ to $\mathbb F_5$ by $i\mapsto2$. Then
\[
 \overline A_X=\begin{pmatrix}3&3\\3&3\end{pmatrix}=3uu^{\mathsf T},
 \quad
 \overline A_Z=\begin{pmatrix}1&0\\0&0\end{pmatrix}=ee^{\mathsf T},
 \quad
 \overline Y=\begin{pmatrix}0&3\\2&0\end{pmatrix},
 \quad u=\binom11,\ e=\binom10.
\]
Every scalar joining two consecutive rank-one factors is nonzero:
$u^{\mathsf T}u=2$, $u^{\mathsf T}e=e^{\mathsf T}u=e^{\mathsf T}e=1$.
Thus every nonempty numerator product has nonzero rank-one reduction,
with image either $\mathbb F_5u$ or $\mathbb F_5e$.
Neither line is invariant under $\overline Y$:
$\overline Yu=(3,2)^{\mathsf T}$ and
$\overline Ye=(0,2)^{\mathsf T}$.

If $YM_w=\pm M_wY$ held in characteristic zero, its reduction would imply
that the image of $\overline M_w$ is invariant under $\overline Y$.
The contradiction proves \eqref{sel40:eq:nohomogeneous}. Apply
Lemma~\ref{sel40:lem:coefficient}, first to the involution and then to the
circle. For the identity coefficient,
$\langle\!\langle I/\sqrt2|J(\mathrm{id})|Y/\sqrt2\rangle\!\rangle=0$
because $\Tr Y=0$.
\end{proof}

\begin{proposition}[The exact finite residue monoid]
\label{sel40:prop:seventeen}
The two reduced numerator matrices generate exactly seventeen matrices:
\[
 \boxed{
 \{I\}\ \sqcup\
 \{\lambda ab^{\mathsf T}:\lambda\in\mathbb F_5^*,\ a,b\in\{u,e\}\}.}
\]
The sixteen nonidentity matrices all have rank one and none obeys
$\overline YM=\pm M\overline Y$. Only the empty word reaches the identity.
\end{proposition}
\begin{proof}
The rank-one calculation proves containment and excludes return to the
identity. Put $A=\overline A_X$, $B=\overline A_Z$.
Direct multiplication gives
$A^2=A$, $B^2=B$, $ABA=3A$, $BAB=3B$.
The element $3$ generates $\mathbb F_5^*$, so these products produce all
nonzero scalar multiples of $A$ and $B$, and multiplication by the other
letter produces all scalar multiples of $ue^{\mathsf T}$ and
$eu^{\mathsf T}$. The sixteen outer products are distinct. The last
assertion is the image-line argument in the theorem.
\end{proof}

\begin{proposition}[Exact finite-word separation, without a uniform all-depth gap]
\label{sel40:prop:distance}
Write
\[
 M_w=a_0I+i(a_1X+a_2Y+a_3Z),\qquad a_j\in\mathbb Z,
\]
and put $A_w=a_0^2+a_2^2$, $B_w=a_1^2+a_3^2$. For every nonempty word
of length $k$,
\[
 A_w,B_w\ge1,\qquad A_w+B_w=25^k,
\]
and the actual unnormalized branch has parity defect
\begin{equation}
 \boxed{
 \|\operatorname{Ad}_Y\Phi_w\operatorname{Ad}_Y-\Phi_w\|_\diamond
 =\frac{4p_w\sqrt{A_wB_w}}{25^k}
 \ge\frac{4p_w\sqrt{25^k-1}}{25^k}>0.}
 \label{sel40:eq:paritydistance}
\end{equation}
Nevertheless, after division by the known word probability, these defects
have infimum zero over all nonempty words.
\end{proposition}
\begin{proof}
The quaternion form is closed under multiplication with integer
coefficients. The determinant identity $\det A_a=25$ gives
$\sum a_j^2=25^k$. Commutation with $Y$ would mean $B_w=0$ and
anticommutation would mean $A_w=0$; both are excluded by the theorem.
The relative unitary $U_w^*YU_wY$ has determinant one and trace
$2(A_w-B_w)/25^k$. For two qubit unitary channels whose relative eigenvalues
are $e^{\pm i\theta}$, the diamond distance is $2|\sin\theta|$.
Indeed, purification reduces the pure-output overlap to
$|\Tr(\rho U_w^*YU_wY)|$; its minimum over densities is the distance of
zero from the line segment joining the two eigenvalues, namely
$|\cos\theta|$. Convexity of trace norm and purification justify taking
pure inputs with an ancilla. Substituting the trace and restoring $p_w$
gives the equality. Positive integers with sum $25^k$ have product at least
$25^k-1$, giving the bound.

For the final statement it suffices to repeat $X$. Let
$e^{i\theta}=(3+4i)/5$. This is not a root of unity: its degree-two
monic minimal polynomial over $\mathbb Q$ is
$x^2-(6/5)x+1$, so it is not an algebraic integer. Thus
$\theta/(2\pi)$ is irrational. The pigeonhole principle applied to the
first $M+1$ points on the unit circle gives a positive $k\le M$ with
$k\theta$ within $2\pi/M$ of an integer multiple of $2\pi$.
Along an unbounded subsequence, $U_X^k\to I$, and the normalized parity
defect tends to zero. No measured approximate matrix has been rounded to
an exact residue class.
\end{proof}

This result is about the already reconstructed $Y$ action and the unchanged
pulse alphabet. It does not exclude odd algebra elements, parity-changing
linear combinations, or a differently identified historical grading. In
particular V39's $I/Y$ test rays still have nonzero cross at $XZ$ and $ZX$;
they have not been identified as physical Standard-Model incidence rays.
The new conclusion is that searching further in this exact alphabet cannot
turn that particular cross witness into a fixed-$Y$-label symmetry witness.

\section[Selecting a word from a certified finite language]{Intersect the response with the side-condition record}
\label{sel40:language}

Let $y(w)=Cx_w$, $x_{aw}=T_ax_w$, be any finite-dimensional exact
realization of the desired vector-valued cross response, of dimension $r$.
Use original branch weights, not normalized conditional amplitudes.
Let a finite deterministic automaton have state set $Q$, initial state
$q_0$, transition $\delta(q,a)$, and accepting set $F\subseteq Q$.
All original admission constraints may first be included by a product with
the actual finite-state typed grammar. The automaton may be calculated from
source matrices; it need not be a separately executed physical memory.

\begin{theorem}[Finite selection on an exactly recognized side-condition language]
\label{sel40:thm:gated}
The gated response
\[
 y_F(w)=\mathbf1_{\{q(w)\in F\}}y(w)
\]
has a finite linear realization on $\C^Q\otimes\C^r$:
\begin{align*}
 \widetilde x_0&=e_{q_0}\otimes x_0,\\
 \widetilde T_a&=\sum_{q\in Q}
 |e_{\delta(q,a)}\rangle\langle e_q|\otimes T_a,\\
 \widetilde C(e_q\otimes x)&=\mathbf1_F(q)Cx.
\end{align*}
Its complete target Hankel rank $r_F$ satisfies $r_F\le |Q|r$.
If $y_F$ is nonzero, an actual accepted word with $y(w)\ne0$ exists at
length at most $r_F-1$ and is returned by reachable-word basis search.
If $y_F$ is zero after certified reachability and observation closure, no
accepted original word has the cross response.

If acceptance is \emph{equivalent} to the exact side condition, this is a
complete finite decision procedure for that condition together with a
nonzero cross. It does not require the condition to hold for every primitive
letter or for every word.
\end{theorem}
\begin{proof}
Induction gives
$\widetilde T_w\widetilde x_0=e_{q(w)}\otimes T_wx_0$.
Applying the output proves the formula. Reachable/observable reduction and
the actual-word length bound are exactly the elementary argument in
Theorem~\ref{sel39:thm:minimal}, applied to these product matrices.
Finite dimension gives the rank bound. Only original label sequences are
retained as basis witnesses; the tensor notation introduces neither a
coherent superposition of words nor a physical copy of the quantum system.
\end{proof}

\begin{corollary}[One-sided filters keep one-sided conclusions]
\label{sel40:cor:filters}
If every genuinely compatible word is accepted, a zero gated response is a
valid all-depth no-go; a nonzero gated word must still be checked against
the original side condition. If every accepted word is compatible, a
nonzero gated response certifies a valid witness, while a zero does not
exclude compatible words outside that language. A finite-field necessary
filter from Theorem~\ref{sel40:thm:modular} has the first interpretation,
not the second.
\end{corollary}
\begin{proof}
These are respectively the two inclusions between the accepted language and
the side-condition language. Apply the pointwise definition of $y_F$.
\end{proof}

For determinant selection use the \emph{complete} $H$-covariance condition,
$H=\ker(\det U\det V)$, not only the vanishing of a chosen matrix entry.
If the specified Choi sectors transform pointwise by $1$ and
$(\det U\det V)^{\pm1}$, any $H$-covariant word with nonzero cross has
covariance group exactly $H$ inside $U(3)\times U(2)$, by the two
inclusions in Theorem~\ref{sel39:thm:selector}. Other linear or nonlinear
conditions may be recognized by a separate exact finite monitor when the
source supplies one. This theorem does not prove such a monitor always
exists.

\section[An exact covariance monitor from the original coefficient rays]{Finite transport orbits remove the universal-covariance premise}
\label{sel40:orbit}

Let $k_a$ be the invertible efficient coefficients of actual resolved labels
on one fixed memory. Let $T_1,\ldots,T_s$ be an independently identified
Hermitian generating bank for the represented Lie algebra of a connected
compact comparison subgroup $H$. Define the ordered tuple
\[
 \mathcal F_0=(T_1,\ldots,T_s),\qquad
 \mathcal F_w=(k_wT_\ell k_w^{-1})_{\ell=1}^s.
\]
The inverses are used only to calculate this tuple. They are not admitted
inverse operations. Individual Kraus phases and nonzero scalar branch
weights cancel from $\mathcal F_w$.

\begin{theorem}[Finite symmetry-transport monitor]
\label{sel40:thm:orbit}
Suppose the exact orbit
$\mathcal Q=\{\mathcal F_w:w\text{ a word}\}$ is finite.
Then the original append rules
\[
 \mathcal F\longmapsto(k_aF_\ell k_a^{-1})_\ell
\]
define a finite deterministic automaton. With initial and sole accepting
state $\mathcal F_0$, it recognizes \emph{exactly} the words whose original
CP map is $H$-covariant. Every state and transition can be certified by a
finite list of exact tuple equalities and closure identities.

Combined with Theorem~\ref{sel40:thm:gated}, this yields an exact finite
same-word covariance--coherence selector even when the universal gate of
V39 fails. Its raw witness-length bound is $|\mathcal Q|r-1$; minimizing
the gated response can give a smaller bound.
\end{theorem}
\begin{proof}
The update follows from $k_{aw}=k_ak_w$. A return to $\mathcal F_0$ is
exactly $[k_w,T_\ell]=0$ for every $\ell$. Each word coefficient is
invertible, so Lemma~\ref{sel40:lem:coefficient} identifies this with
$H$ covariance of its CP map. Finite orbit closure gives a complete
transition table. The product realization proves the remaining claims.
Equality of the span of tuple coordinates is not orbit closure; the actual
finite list of transported tuples must close.
\end{proof}

It is essential to retain an \emph{ordered} generating tuple, not just its
linear span or generated algebra. A coefficient can normalize a subgroup
while acting nontrivially on it. For example $X$ sends $Y$ to $-Y$ and
therefore preserves $\R Y$, but $\operatorname{Ad}_X$ is not covariant
under the circle generated by $Y$. A span-only return would confuse
normalization with commutation.

Finite orbit is an additional source property, not a consequence of finite
Hilbert dimension. The $X$ pulse above has infinitely many distinct
transported $Y$ matrices: a return of two such matrices would give
$U_X^kYU_X^{-k}=Y$ for a nonzero integer $k$, contrary to its irrational
rotation angle. Thus an exact arithmetic exclusion can terminate a source
audit even when this particular positive finite-orbit method does not.

\section[The delayed family now has a complete exact language]{The missing resource is symmetry-return complexity}
\label{sel40:delay}

Retain without change the source of Theorem~\ref{sel39:thm:delay}:
\[
 R(A,B)=1\oplus\chi(A,B)\oplus A\oplus1,\qquad
 U_m=e^{-i\alpha X_{n\omega}}\oplus e^{-4i\alpha X_{c_1r}}
       \oplus I_{\operatorname{span}\{c_2,c_3\}},
 \quad \alpha=\frac{\pi}{4m},\quad m\ge7.
\]
Here $\chi(A,B)=\det A\det B$ is part of the comparison representation,
not a derived historical charge. There is one actual unitary label.
Set $b=I_{C\oplus\C r}/2$ and $d=|\omega\rangle\langle n|$,
as in V39, and define the real scalar response by
\[
 \langle\!\langle b|J(\operatorname{Ad}_{U_m^k})|d\rangle\!\rangle
 =i\,y_m(k).
\]

\begin{theorem}[Exact return language and its minimal gated rank]
\label{sel40:thm:delay}
For this unchanged family:
\begin{enumerate}[label=\textup{(\roman*)},leftmargin=2.8em]
\item The complete $H=\ker\chi$-covariance language is
$\{a^{2mj}:j\ge0\}$. Its minimal deterministic automaton has exactly
$2m$ states. The ordered $H$-generator transport orbit has exactly $2m$
elements and supplies such a monitor.
\item The ungated cross is
\[
 y_m(k)=\bigl(1+\cos(k\pi/m)\bigr)\sin(k\pi/(4m)),
\]
with minimal Hankel rank six. The covariance-gated cross is
\begin{equation}
 \boxed{
 y_{m,H}(k)=
 \begin{cases}
 2\sin(j\pi/2),&k=2mj,\\
 0,&2m\nmid k,
 \end{cases}
 \qquad \operatorname{rank}\mathbb H_{m,H}=4m.}
 \label{sel40:eq:gatedrank}
\end{equation}
\item The exact-group language is
\[
 \boxed{
 \operatorname{Cov}_{G_0}(\operatorname{Ad}_{U_m^k})=H
 \quad\Longleftrightarrow\quad k\equiv2m\pmod{4m}.}
\]
It has minimal deterministic automaton size $4m$. At every such word the
cross mass is four. At the other covariance-return words the group is
$G_0$. At nonreturn words the group does not contain $H$.
\item At $m=8$, the ordinary Liouville span has rank thirteen and is flat at
depth twelve, the covariance monitor has sixteen states, and the gated
cross has rank thirty-two. Its first nonzero word has length sixteen.
\end{enumerate}
\end{theorem}
\begin{proof}
V39 proves that $H$ covariance is equivalent to $2m\mid k$ by testing
the actual $SU(3)$ action on $c_1,c_2,c_3$. At such a length the entire
$C\oplus\C r$ block is the identity, while the $n,\omega$ block is a
rotation through $j\pi/2$. Since $H$ acts trivially on that block,
$2m\mid k$ is also sufficient. The same criterion says that the tuple
returns for the first time after $2m$ steps. Unitary similarity is
invertible, so no earlier pair of orbit states coincides. This proves the
orbit size. The residues modulo $2m$ are inequivalent states for the
language: append the suffix reaching residue zero from one of two distinct
residues, and only that one is accepted. This proves minimality.

Taking Hilbert--Schmidt inner products with $b$ and $d$ gives the displayed
$y_m$. With $\zeta=e^{i\pi/(4m)}$ its nonzero Fourier exponents are
$\pm1,\pm3,\pm5$, all distinct modulo $8m$. A Vandermonde factorization
therefore gives exact ungated Hankel rank six. At $k=2mj$ the cosine is
one, proving \eqref{sel40:eq:gatedrank}'s formula. The gated sequence obeys
$y_{m,H}(k+4m)=-y_{m,H}(k)$, so its rank is at most $4m$.

For the matching lower bound take the square Hankel matrix with row and
column indices $0,\ldots,4m-1$. Its only nonzero entries have index sum
$2m$, where the value is $2$, or index sum $6m$, where the value is $-2$.
It is the direct sum of two antidiagonal blocks of lengths $2m+1$ and
$2m-1$. Their determinant signs multiply to $-1$, and their entry-product
sign is also $-1$. Hence
\begin{equation}
 \boxed{\det[\,y_{m,H}(r+c)\,]_{r,c=0}^{4m-1}=2^{4m}\ne0.}
 \label{sel40:eq:hankeldet}
\end{equation}
This proves the exact rank. No modular upper-bound inference is used.

If $j$ is odd, the active rotation has a nonzero off-diagonal coefficient
and the other four directions are fixed. A covariance can multiply the
unitary only by a common phase; its identity complement fixes that phase
to one. The off-diagonal active coefficient then forces $\chi=1$.
If $j$ is even, the active block is $I$ or $-I$ and the whole operator
commutes with $G_0$. This proves (iii). The residues modulo $4m$ are
pairwise distinguishable by suffixes reaching residue $2m$, so its stated
finite-state cost is minimal. Part (iv) specializes these exact formulas
and the unchanged rank-thirteen calculation of V39.
\end{proof}

The covariance monitor is not a new experimentally executed operation. It
is an exact classification of the retained lengths of the original label.
It neither averages over a return nor conditions away a waiting time. If
an original scalar occurrence factor $\vartheta$ multiplies this event,
then the gated cross at length $k$ is multiplied by $\vartheta^k$ and the
mass by $\vartheta^{2k}$. The accepted language and exact Hankel rank remain
the same for $\vartheta>0$; the $4m$-square determinant in
\eqref{sel40:eq:hankeldet} is multiplied by
$\vartheta^{(4m)(4m-1)}$. The event count $2m$ is not identified with
physical elapsed time without the actual waiting-time law.

\section[Source status after the exact-language audit]{What is now decided and what is not}
\label{sel40:status}

\paragraph{Proved in this continuation.}
The unchanged forward-pulse source has no nonempty fixed-$Y$-covariant word,
as certified by a seventeen-element residue monoid and its two possible
image lines. This remains an exact all-depth exclusion even though its
coherence response is nonzero and its normalized covariance defects can
approach zero. A general finite-residue test supplies a sound necessary
language for homogeneous algebraic side constraints. A finite ordered
symmetry-transport orbit supplies an exact covariance language for
invertible efficient labels. Gating a target response by such a language
restores a complete finite simultaneous-occurrence selector. On the
unchanged delayed family the covariance monitor has size $2m$, the gated
Hankel rank is $4m$, and exact determinant-group words are precisely
$2m$ modulo $4m$.

\paragraph{Inherited inputs, not new derivations.}
The pulse coefficients, their adjoint-reversal $Y$, V39's nonzero cross,
its rank-thirteen delayed family, the Choi-sector interpretation and the
reachable/observable minimality method are retained. The determinant
character is explicitly supplied in the delayed comparison. Neither the
finite-field alphabet nor the automaton product is assigned a positive
quantum-state interpretation. No new physical pulse has been introduced.

\paragraph{Still unresolved for the historical source.}
The actual V110 directed incidence compiler and primitive same-cut source
maps have not been populated by these calculations. No historical
Standard-Model selection is claimed. For an exact algebraic source with
identified side constraints, first seek a finite necessary residue filter;
for a source with a certified finite symmetry-transport orbit, use its exact
return language. If neither applies, the existential joint-word problem
remains open for that source. In particular, linear flatness alone does not
settle nonlinear Clifford-isometry or finite-Dirac-shape conditions.

These results narrow the structural audit without importing faithful
stationarity, reversibility, homogeneous time embedding, or a continuum
limit as additional duality axioms. Finite-data uncertainty must be kept
separate: an approximate numerical orbit closure is not exact closure and
a rounded algebraic model is not the physical source. At a fixed tested
word a certified matrix error can be propagated through the parity or Choi
residual; the vanishing all-depth infimum in
Proposition~\ref{sel40:prop:distance} excludes a uniform word-independent
numerical separation claim. The accompanying exact computations supplement
the written proofs and do not constitute proof-assistant certification or
physical measurements.

\addtocontents{toc}{\protect\endgroup}
% END OF SOURCE-SELECTION V40 DEVELOPMENT BLOCK.
% Preserve V1--V40 and append subsequent work before the final terminator.


\clearpage
\hypersetup{pdftitle={Historical Standard-Model Source Selection: V1--V41}}
\begin{center}
{\Large\bfseries Source-selection continuation V41}\\[.4em]
{\large Return to the retained record policy:\\
positive classical filters, exact covariance, and a four-event source certificate}
\end{center}
\addtocontents{toc}{\protect\begingroup}
\addtocontents{toc}{\protect\renewcommand{\protect\numberline}{\protect\makebox[2.1em][l]}}

\section[The original word and its classically processed outcome]{The next upstream distinction is the retained record}
\label{sel41:context}

V40 proves that no nonempty resolved word of the unchanged two-pulse source
is fixed by its reconstructed $Y$ parity. That theorem is explicitly about
an individual word with its label retained. It does not decide the covariance
of an outcome obtained by classical selection and coarse graining of several
original words. The operational entrance is affine under admitted classical
randomization and additive under admitted coarse graining \cite{sel41-spacetime}.
Accordingly, these two event classes must be separated before another
word-length search is attempted.

This continuation proves a finite positive-filter criterion and evaluates it
on exactly the V34--V40 pulse source. A four-event classical acceptance rule
produces a mixed CP outcome covariant under the whole $Y$ circle, with
nonzero $I/Y$ coefficient coherence, despite the all-depth single-word
exclusion. The rule changes neither quantum pulse and preserves every old
branch after forgetting its additional acceptance flag. Its covariance
requires forgetting the fine word inside the selected outcome; the complete
old instrument is not made covariant. Moreover, positive filtering cannot
remove a wrong typed support or turn distinct source rays into one efficient
ray. Those facts preserve the stricter historical V110 incidence boundary.

Classical and conditional post-processing of quantum instruments, and the
convex duality used below, are established frameworks
\cite{sel41-instruments,sel41-convex}. The contribution here is the exact
source-constrained filter test, the unchanged-source calculation, its
four-event minimality and optimal conditional cross, and the separation of
record cancellation from typed-source occurrence. Neither the qubit rays nor
its circle are identified as historical Standard-Model representations.

All positive implementation statements below are relative to an explicitly
admitted classical acceptance policy on the same input and output cuts.
Affineness of a probability table alone does not certify that a particular
policy was historically executed. The matrices determining that policy are
not additional quantum operations. Exact finite models and measured
finite-precision data are kept separate.

\section[Classical filtering and the finite positive cone]{The admissible operation is a positive sum, not an amplitude sum}
\label{sel41:cone}

Let $\{\Phi_j\}_{j=1}^M$ be a finite original instrument on $M_n$,
possibly the complete fixed-depth word instrument, and let $J_j\succeq0$
be its Choi operators. A restricted bank is permitted when its untouched
complement is retained as an additional failure outcome. For an admitted
classical rule $0\le f_j\le1$, set
\begin{equation}
 \Phi_f=\sum_j f_j\Phi_j,\qquad J_f=\sum_j f_jJ_j.
 \label{sel41:eq:filter}
\end{equation}
The acceptance probability is $\Tr\Phi_f(\rho)$ and is not generally
state independent. In particular, $\Phi_f/\Tr\Phi_f(\rho)$ is not to be
called a linear channel unless $\Phi_f^*(I)=sI$ with $s>0$.

\begin{proposition}[Conservative flag refinement and exact positive support]
\label{sel41:prop:physical}
The maps $f_j\Phi_j$ and $(1-f_j)\Phi_j$, with both $j$ and the acceptance
bit retained, form a normalized classical refinement. Forgetting only the
acceptance bit recovers every original map $\Phi_j$ exactly. Forgetting
$j$ inside the accepted branch gives \eqref{sel41:eq:filter}, with no
cross terms between different Kraus coefficients.

For every positive filter,
\begin{equation}
 \boxed{\Ran J_f=\sum_{j:f_j>0}\Ran J_j.}
 \label{sel41:eq:support}
\end{equation}
Hence zero leakage outside a specified coefficient support $P$ requires
$\Ran J_j\subseteq\Ran P$ for every selected $j$. A nonzero $J_f$ has
rank one exactly when all selected nonzero $J_j$ lie on the same rank-one
positive ray. For any specified coefficient projections $P_B,P_D$,
\begin{equation}
 P_BJ_jP_D=0\ (\text{all }j)\quad\Longrightarrow\quad P_BJ_fP_D=0.
 \label{sel41:eq:nocreate}
\end{equation}
Thus classical filtering cannot create a missing same-word coherence,
repair an incompatible typed support, or purify distinct coefficient rays.
\end{proposition}
\begin{proof}
Normalization follows by summing the two positive scalar multiples for
each $j$. For $v$ in the Choi space,
$\langle v,J_fv\rangle=\sum_jf_j\|J_j^{1/2}v\|^2$.
It vanishes exactly when $v\in\Ker J_j$ for every selected $j$.
Orthogonal complements prove \eqref{sel41:eq:support}. The rank and
support statements follow immediately. The final assertion is linearity,
not positivity of a cross block. In particular a nonzero filtered cross
implies at least one original selected cross was nonzero, though other
crosses may cancel.
\end{proof}

Fix a connected compact comparison subgroup $H$ with an independently
identified representation on the input and output system. Let
$\mathscr D_H$ be a finite real linear covariance-defect map on Hermitian
Choi matrices. For example its coordinates can be the real and imaginary
parts of $[\widehat T_\ell,J]$, where $T_\ell$ generate the Lie algebra
and $\widehat T_\ell=T_\ell\otimes I-I\otimes\overline{T_\ell}$.
Its zero is exact $H$ covariance. More generally use the correct distinct
input and output representations in this formula.

Let $\mathscr C(J)$ list the real and imaginary coordinates of $P_BJP_D$.
Write
\[
 d_j=\mathscr D_H(J_j),\qquad c_j=\mathscr C(J_j),\qquad
 D=[d_1\ \cdots\ d_M],\quad C=[c_1\ \cdots\ c_M].
\]
Additional \emph{linear} side conditions can be included in $D$. In
particular, state-independent success is imposed by adjoining the
traceless coordinates of $(\Tr_{\rm out}J_j)^{\mathsf T}$ to $d_j$.
Nonlinear partial-isometry or finite-Dirac-shape conditions are not silently
included in this construction.

\begin{theorem}[Complete finite classical-filter test and dual exclusion]
\label{sel41:thm:LP}
Assume the full box of classical policies $f\in[0,1]^M$ is admitted.
An $H$-covariant selected outcome with nonzero specified cross exists
if and only if
\begin{equation}
 \boxed{\exists f\in[0,1]^M:\quad Df=0,\qquad Cf\ne0.}
 \label{sel41:eq:feasible}
\end{equation}
This is decided by finitely many ordinary linear programs: maximize both
signs of each coordinate of $Cf$ subject to $Df=0$ and $0\le f\le1$.
A positive optimum returns an admissible rule, not a new quantum pulse.
If all optima vanish, no admitted rule in this bank supplies the cross.

For a specified real cross functional $c^{\mathsf T}f$, the absence of
any positive value has the exact conic dual certificate
\begin{equation}
 \boxed{\exists u:\quad D^{\mathsf T}u\ge c\quad\text{componentwise}.}
 \label{sel41:eq:dual}
\end{equation}
A certificate for both signs of every cross coordinate proves the full
no-go. If a positive filter exists, one exists using at most
$\rank D+1$ original nonzero outcome maps, after possibly reducing its
success probability. This support bound does not assert an optimum
throughput at a fixed requested success probability.
\end{theorem}
\begin{proof}
Linearity of the Choi representation proves \eqref{sel41:eq:feasible}.
A nonzero real vector has a coordinate with a nonzero sign. Zero is always
feasible and each box-constrained program is compact, so the programs give
both implications. The upper bounds $f_j\le1$ do not affect existence
of a positive homogeneous objective: any vector in
$K=\{f\ge0:Df=0\}$ can be scaled into the box. The polar cone is
$K^\circ=\Ran D^{\mathsf T}-\R_+^M$, a closed polyhedral cone.
Thus $c^{\mathsf T}f\le0$ on $K$ is equivalent to
$c=D^{\mathsf T}u-v$ for $v\ge0$, which is
\eqref{sel41:eq:dual}. Alternatively this is the finite-dimensional
separating-hyperplane form of Farkas' lemma.

Discard zero $J_j$ and let $t_j=\Tr J_j>0$. Normalize a positive witness
by $t^{\mathsf T}f=1$. The set
$\{f\ge0:Df=0,t^{\mathsf T}f=1\}$ is compact. An extreme maximizer
of the selected positive coordinate has at most $\rank D+1$ positive
entries: otherwise a nonzero null vector of the equality rows on its
support would permit a small perturbation of both signs, contradicting
extremality. Scaling this finite vector down makes it a physical filter.
\end{proof}

If the source admits only prescribed deterministic bins or a finite set of
classical probabilities, those restrictions must replace the box.
The theorem does not manufacture a randomizer from linear reconstruction.
At the other extreme, allowing arbitrary nonnegative coefficients rather
than the box is legitimate only for existence \emph{up to scale}; the
normalization and original success factors must be restored afterward.

\begin{corollary}[A filtered determinant-group certificate]
\label{sel41:cor:det}
Suppose $H=\ker\chi$ inside $G_0=\mathrm U(3)\times\mathrm U(2)$,
$\chi(U,V)=\det U\det V$, and the actual identified coefficient sectors
$P_B,P_D$ transform pointwise by $1,\chi^{\pm1}$. A feasible filter
with $Df=0$ and $Cf\ne0$ has
$\operatorname{Cov}_{G_0}(\Phi_f)=H$. This concludes covariance of the
selected CP outcome, not of the complete original instrument, and does
not prove one efficient source ray or the other V110 typing conditions.
\end{corollary}
\begin{proof}
$Df=0$ gives $H$ covariance of the complete selected Choi matrix.
Every additional covariance must fix its nonzero $P_B/P_D$ cross, so its
relative character is one. Thus it belongs to $H$. This is the same
pair of group inclusions used in V39, now applied to a positive classical
sum of original outcome maps.
\end{proof}

\section[Why the first three original horizons fail]{Short fixed-depth exclusions use positivity, not a numerical search}
\label{sel41:short}

Retain the original source exactly:
\[
 k_X=\frac{9I+12iX}{25}=\frac35U_X,\qquad
 k_Z=\frac{12I+16iZ}{25}=\frac45U_Z,\qquad
 U_X=\frac{3I+4iX}{5},\quad U_Z=\frac{3I+4iZ}{5}.
\]
For chronological $w=a_1\cdots a_n$, set
\[
 M_w=A_{a_n}\cdots A_{a_1}
     =a_wI+i(b_wX+c_wY+d_wZ),\quad
 A_X=3I+4iX,\quad A_Z=3I+4iZ.
\]
All four coordinates are integers and
$a_w^2+b_w^2+c_w^2+d_w^2=25^n$. The actual word channel is
$\Phi_w=p_w\operatorname{Ad}_{U_w}$, where
\[
 U_w=M_w/5^n,\qquad
 p_w=(9/25)^{n_X(w)}(16/25)^{n_Z(w)}.
\]
Since $p_w>0$, a normalized mixture at this horizon can be realized by a
classical filter at some positive success probability. Its conditional
weights are denoted $r_w\ge0$, $\sum_wr_w=1$.

Let $H_Y=\{e^{-itY}:t\in\R\}\subset SU(2)$. Put
\begin{equation}
 F(w)=(a_wb_w,a_wd_w,b_wc_w,c_wd_w,b_wd_w,b_w^2-d_w^2)^{\mathsf T}.
 \label{sel41:eq:quaterniondefect}
\end{equation}

\begin{lemma}[Six exact circle-covariance equations]
\label{sel41:lem:six}
A fixed-depth mixture $\sum_wr_w\operatorname{Ad}_{U_w}$ is
$H_Y$ covariant exactly when
\begin{equation}
 \boxed{\sum_wr_wF(w)=0.}
 \label{sel41:eq:six}
\end{equation}
In the Hilbert--Schmidt orthonormal coefficient basis
$(I,X,Y,Z)/\sqrt2$, its $I/Y$ Choi cross is
\[
 \langle\!\langle I/\sqrt2|J|Y/\sqrt2\rangle\!\rangle
 =-i\gamma,\qquad
 \gamma=\sum_wr_w\frac{2a_wc_w}{25^n}.
\]
\end{lemma}
\begin{proof}
The Choi coordinate of $U_w$ is
$\sqrt2(a_w,ib_w,ic_w,id_w)^{\mathsf T}/5^n$.
Conjugation by $e^{-itY}$ fixes $I,Y$ and acts by all real planar
rotations on $X,Z$. The four cross entries between these two subspaces
vanish exactly when the first four entries in \eqref{sel41:eq:six}
vanish. The $X/Z$ diagonal block is a real symmetric matrix; it commutes
with all rotations exactly when its off-diagonal entry is zero and its
two diagonal entries agree. These are the last two equations. The
remaining $I/Y$ block is unrestricted by the circle. Multiplying its
coordinates gives the cross formula.
\end{proof}

\begin{theorem}[No nonzero circle-covariant filter at horizons one to three]
\label{sel41:thm:short}
For each fixed $n=1,2,3$, no nonzero nonnegative mixture of the original
length-$n$ outcomes is $H_Y$ covariant. Thus no classical acceptance
rule at any of these horizons supplies such an outcome.
\end{theorem}
\begin{proof}
For each horizon a linear functional of $F(w)$ is strictly positive on
all its words. Explicit separating rows and their values are
\[
\begin{array}{c|c|c}
 n&v_n&\{v_n\cdot F(w):|w|=n\}\\\hline
1&(1,1,0,0,0,0)&\{12\}\\
2&(-3,-3,0,0,8,0)&\{504\}\\
3&(-24,-24,0,0,-7,0)&\{36288,123552\}.
\end{array}
\]
These values follow by multiplying the two integer quaternion
numerators; at $n=3$ the larger value occurs only on $XXX$ and $ZZZ$.
Taking a nonzero nonnegative weighted sum cannot give zero, proving
\eqref{sel41:eq:six} impossible. This is a strict dual certificate, not
a floating-point infeasibility report.
\end{proof}

The theorem concerns a fixed number of actual events. Combining branches
from different observation times is a different policy and can change the
answer. In particular no claim of a three-event exclusion for arbitrary
variable-duration protocols is made.

\section[An exact four-event filter on the unchanged source]{Classical record processing restores a selected covariance}
\label{sel41:positive}

\begin{theorem}[Six-word filter with state-independent success]
\label{sel41:thm:four}
At horizon four, accept the following six original word records with the
listed probabilities $f_w$, and reject every other word:
\begin{equation}
\begin{array}{c|r|c|c}
 w&m_w&(a_w,b_w,c_w,d_w)&f_w\\\hline
XXXX&714375&(-527,-336,0,0)&317500/365211\\
XXXZ&1195236&(-351,132,-176,-468)&9/11\\
XXZZ&939114&(49,-168,-576,-168)&81/224\\
XZXZ&2597056&(-463,216,-288,216)&1\\
XZZZ&1195236&(-351,-468,-176,132)&729/2816\\
ZZZZ&714375&(-527,0,0,-336)&57864375/664846336
\end{array}
\label{sel41:eq:filtertable}
\end{equation}
All $f_w$ are in $[0,1]$. Their success probability is independent of
\emph{every} input density:
\begin{equation}
 \boxed{
 s_* =\frac{216649728}{1441015625}>0,\qquad
 \Phi_f^*(I)=s_*I.}
 \label{sel41:eq:success}
\end{equation}
The conditioned channel $\Psi_*=\Phi_f/s_*$ has weights
$r_w=m_w/7355392$ and Choi matrix, in the coefficient basis
$(I,X,Y,Z)/\sqrt2$,
\begin{equation}
 \boxed{
 J_* =\frac1{408125}
 \begin{pmatrix}
 355198&0&-132804i&0\\
 0&145152&0&0\\
 132804i&0&170748&0\\
 0&0&0&145152
 \end{pmatrix}.}
 \label{sel41:eq:choi}
\end{equation}
It is strictly positive, has trace two, and satisfies
\begin{equation}
 \boxed{
 \operatorname{Cov}_{SU(2)}(\Psi_*)=H_Y,\qquad
 \left|\langle\!\langle I/\sqrt2|J_*|Y/\sqrt2\rangle\!\rangle\right|
 =\frac{132804}{408125}>0.}
 \label{sel41:eq:exactgroup}
\end{equation}
No original word in this filter is $Y$-parity covariant. Theorem
\ref{sel40:thm:pulse-nogo} continues to exclude every individual nonempty
word at every depth.
\end{theorem}
\begin{proof}
The integer columns in \eqref{sel41:eq:filtertable} are literal products
of the original numerators, in the stated chronological convention.
They give $\sum_wm_w=7355392$ and
$\sum_wm_wF(w)=0$. Direct rational multiplication gives
$f_wp_w=s_*m_w/7355392$ on the six selected words, proving success and
normalization without dividing out an input-dependent event probability.
Using Lemma~\ref{sel41:lem:six} to form the coefficient outer products
gives \eqref{sel41:eq:choi}. Its $I/Y$ principal determinant numerator is
\[
 355198\cdot170748-132804^2=43012445688>0,
\]
and its two other diagonal entries are positive, proving strict positivity.
Trace preservation also follows from the mixture of original unitary
channels. Its Bloch matrix in the ordered axes $(X,Y,Z)$ is
\begin{equation}
 B_* =\frac1{408125}
 \begin{pmatrix}
 92225&0&-132804\\
 0&117821&0\\
 132804&0&92225
 \end{pmatrix}.
 \label{sel41:eq:bloch}
\end{equation}
It commutes with every rotation about $Y$. Conversely its nonzero real
antisymmetric part is a nonzero cross-product matrix with axis $Y$.
Every $SO(3)$ rotation commuting with $B_*$ commutes with that part and
therefore fixes $Y$, not just the unoriented line $\R Y$. It is a rotation
about $Y$. Lifting to $SU(2)$ proves the exact covariance group. Scalar
system phases, if the comparison group is $U(2)$, are automatically added
and have no action on density matrices. The final assertion is the
unchanged V40 theorem.
\end{proof}

For this displayed mixture, $s_*$ is the largest allowed success scalar:
$r_wp_w^{-1}s\le1$ for every selected word gives
$s\le\min_w p_w/r_w=s_*$, attained on $XZXZ$. It is \emph{not} claimed
to be the globally maximal success probability over all possible filters.
The source's actual coefficient algebra and the Choi rank are not reduced:
$J_*$ has rank four and its Kraus span is all of $M_2$.

\begin{theorem}[Optimal conditional cross at the first feasible horizon]
\label{sel41:thm:optimal}
Among all $H_Y$-covariant nonnegative mixtures of the sixteen original
four-event unitary channels, the maximum absolute $I/Y$ Choi cross is
exactly $132804/408125$. Thus \eqref{sel41:eq:choi} attains the largest
\emph{conditional} cross at that horizon. Together with
Theorem~\ref{sel41:thm:short}, four is the least fixed horizon with a
nonzero circle-covariant classically selected outcome.
\end{theorem}
\begin{proof}
For each original four-event quaternion let $f(w)=F(w)/25^4$ and
$c(w)=2a_wc_w/25^4$. Set
\[
 g=\frac{132804}{408125},\qquad
 u=\left(-\frac{1875}{2612},-\frac{1875}{2612},
          -\frac{1875}{2612},-\frac{1875}{2612},-2,0\right).
\]
The exact dual inequalities $g+u\cdot f(w)-c(w)\ge0$ hold on all
sixteen original words. For a transparent integer verification their
left sides, multiplied by $255078125$, are zero except as follows:
\[
\begin{array}{c|r}
\text{words}&255078125\,[g+u\cdot f(w)-c(w)]\\\hline
XZXX,
XZZX,
ZXXZ,
ZZXZ&204724800\\
ZXXX,
ZZZX&216798912\\
ZXZX&231654528\\
ZZXX&107718912
\end{array}
\]
Weighting by $r_w$, using $\sum r_w=1$ and $\sum r_wf(w)=0$, gives
$\gamma\le g$. Every word used in \eqref{sel41:eq:filtertable} has zero
dual slack, so it attains equality. Transposing the numerator reverses
the chronological word, preserves $a_w,b_w,d_w$, and changes the sign
of $c_w$. Reversal preserves word probabilities and covariance
constraints. Applying the same upper bound to the reversed mixture gives
$-\gamma\le g$. The cross is always purely imaginary, so the absolute
bound follows. All inequalities are rational equalities or inequalities
on a finite original bank; no numerical optimizer is part of the proof.
\end{proof}

The unnormalized success cross and its mass retain the probability budget:
\begin{equation}
 \boxed{
 s_*\frac{132804}{408125}=\frac{11943936}{244140625},\qquad
 h_f=\left(\frac{11943936}{244140625}\right)^2.}
 \label{sel41:eq:rawcross}
\end{equation}
These are not the conditional values in \eqref{sel41:eq:exactgroup}.

\section[The full record is not made covariant]{Record cancellation does not repair the whole historical source}
\label{sel41:record}

\begin{proposition}[Fine records, coarse bins, and the total-map obstruction]
\label{sel41:prop:records}
For a fixed-label group action, the fine accepted instrument
$\rho\mapsto\bigoplus_j f_j\Phi_j(\rho)$ is $H$ covariant exactly when
every selected nonzero $\Phi_j$ is $H$ covariant. Thus the positive pulse
filter above is not covariant when its original words are also retained.

For any complete classical post-processing into bins, covariance of every
bin implies covariance of the original total map. If a selected success
bin is covariant but the original total map is not, its complement must
fail covariance. For the unchanged pulse bank at horizon four the total
map $\Phi^4$ is not even fixed by $Y$ parity: its Bloch $(X,Y)$ entry is
\begin{equation}
 \boxed{(B_{\Phi^4})_{12}=\frac{4196550144}{152587890625}\ne0.}
 \label{sel41:eq:totaldefect}
\end{equation}
Therefore the complete two-outcome success/failure instrument in
Theorem~\ref{sel41:thm:four} is not $H_Y$ covariant.
\end{proposition}
\begin{proof}
The fine Choi matrix is a classical direct sum. Equality under a group
that leaves its labels fixed is equality of every nonzero summand; a
positive scalar does not change that equality. Summing the covariance
relations over any complete collection of bins proves the total-map
claim. The one-event Bloch matrix of the unchanged channel is
\[
 B_\Phi=\frac1{625}
 \begin{pmatrix}113&384&0\\-384&-175&216\\0&-216&337\end{pmatrix}.
\]
Raising it to the fourth power gives \eqref{sel41:eq:totaldefect}.
Conjugation by $Y$ acts on Bloch space as $\diag(-1,1,-1)$ and reverses
that entry, so covariance fails. Since success is covariant, failure
carries the nonzero remaining defect.
\end{proof}

This distinction is decisive for interpretation. The acceptance flag can be
adjoined conservatively, but erasing the fine word from the success output
is a genuine change of the retained record policy. A theorem about that
selected outcome is not a symmetry theorem for the entire original
history algebra or all recorded operations. There is no contradiction with
V40's fixed-word no-go and no common-ray V110 conclusion: the selected
packet has rank four, not one. The typed and directed conditions must still
be tested on the actual source to which such a conclusion is attached.

An additional finite check in the accompanying certificate enumerates all
$2^{16}$ deterministic subsets at this horizon and finds that only the
empty subset is circle covariant with the \emph{original} probabilities.
This observation is a finite exhaustive integer check, not used in the
proof of the positive theorem or its optimality. In particular the
explicit nontrivial probabilities $f_w$ must not be omitted from the
implementation statement.

\section[Conic closure, inherited zeros, and the historical interface]{The broader search needs positive closure, not merely a flat span}
\label{sel41:closure}

On one fixed input memory let $J_w$ denote all actual endomorphism-word
Choi matrices and set
\[
 \mathcal K_r=\operatorname{cone}\{J_w:|w|\le r\}.
\]
Each finite cone is polyhedral. Appending a primitive map acts linearly
by $J_w\mapsto(\Phi_a\otimes\id)(J_w)$ in the output--input Choi
convention. Typed finite-state admission can be retained by separate
cones at each grammar state; no inadmissible concatenation is added.

\begin{proposition}[Conic flatness is sufficient, but is not linear flatness]
\label{sel41:prop:conic}
If $\mathcal K_{r+1}=\mathcal K_r$ and all admissible append operations
are included, then $\mathcal K_s=\mathcal K_r$ for every $s\ge r$.
This equality is a finite family of nonnegative membership tests, one
for each appended generating ray. The finite filter test can then decide
all covariance--cross possibilities in this complete cone.

A negative answer excludes every finite nonnegative combination of
original words, hence every actual classical policy that uses them.
A positive cone element gives a physical selected outcome only if the
finite combination of horizons and record erasures it uses is admitted;
otherwise it remains an algebraic recipe rather than an executed protocol.
At a fixed common horizon the implementation is exactly
Proposition~\ref{sel41:prop:physical}.

For the unchanged two-pulse source the increasing cones never stabilize,
even though the ordinary Liouville span has finite rank ten. The positive
four-event result does not require conic flatness.
\end{proposition}
\begin{proof}
Cone equality makes $\mathcal K_r$ invariant under every append map:
apply the append map to its original generators and use nonnegative
linearity. Induction proves all-depth containment, and the reverse
containment is automatic. Membership of a vector in a finitely generated
cone is a finite system of linear equations with nonnegative variables.
The filter test and scaling argument therefore apply on the closed cone.
The implementation qualification retains the distinction between a finite
algebraic cone and the source's permitted record/time policies.

For nonstabilization, $J_{X^k}$ are rank-one positive matrices on infinitely
many distinct rays. Indeed $U_X$ has eigenvalues $(3\pm4i)/5$, whose
ratio is not a root of unity, as already proved in V40. No two different
positive powers are proportional as unitary matrices. If a rank-one
positive matrix is a sum of positive matrices, every nonzero summand has
the same support line. Thus a cone generated by finitely many original
rank-one word Choi matrices cannot contain all these distinct rays.
No finite $\mathcal K_r$ is therefore the complete cone.
\end{proof}

In contrast, the V29 all-depth zero $P_BJ_wP_D=0$ remains conclusive without
conic flatness: by \eqref{sel41:eq:nocreate}, it excludes every finite
classical filter. It is the \emph{covariance} exclusion of individual
words, not the \emph{cross-zero} exclusion, that may fail to survive coarse
graining. Similarly, every exact positive support test remains raywise;
covariance is a signed matrix equality and can cancel across records.

\section[Calibrated errors and source status]{What the new branch establishes and what it leaves open}
\label{sel41:status}

For fixed physical maps and two classical rules $f,\widehat f$,
\begin{equation}
 \|\Phi_f-\Phi_{\widehat f}\|_\diamond
 \le\sum_j|f_j-\widehat f_j|\,\|\Phi_j^*(I)\|_{\op}.
 \label{sel41:eq:policyerror}
\end{equation}
This is the triangle inequality and the CP-map identity
$\|\Phi_j\|_\diamond=\|\Phi_j^*(I)\|_{\op}$.
For the original random-unitary words the weights on the right are the
unchanged $p_w$. Errors in the original maps additionally contribute
$\sum_j f_j\|\Phi_j-\widehat\Phi_j\|_\diamond$. A policy solved only
against rounded matrices need not satisfy exact covariance of the true
source. A positive cross can be certified from a norm margin, but a small
covariance defect is not an exact zero.

At four consecutive raw opportunities with an independent acceptance
factor $\vartheta$, rejecting all records except the four-accepted words
multiplies $s_*$ and the unnormalized cross by $\vartheta^4$, and $h_f$
by $\vartheta^8$. The conditional channel remains \eqref{sel41:eq:choi}.
Four accepted events and four raw opportunities are different experiments;
no waiting-time law is silently replaced. The numerical success
$s_*\simeq0.15034516$ is per four accepted original events, not a
mass gap, a continuum rate, or a historical Standard-Model measurement.

\paragraph{Proved.}
A finite original same-cut bank admits a complete linear-programming test
for classically filtered covariance with nonzero cross, with explicit dual
exclusions and a sparse-existence bound. Positive support and efficient-ray
requirements do not admit cancellations. On the unchanged two-pulse bank,
all fixed horizons below four are excluded by strict exact separating rows.
A stated six-word classical rule at horizon four has state-independent
success, a strictly positive Choi matrix, exact circle covariance, and an
optimal conditional $I/Y$ cross. Neither the fine accepted record nor the
complete success/failure instrument has that covariance. Positive cone
closure gives a sufficient all-depth stopping rule, but linear flatness
does not imply it.

\paragraph{Inherited and conditional inputs.}
The original pulse coefficients, all-depth fixed-word exclusion, actual
branch probabilities, and target-sector conventions are unchanged.
Classical post-processing and finite-dimensional convex duality are
standard tools. The operational rule explicitly includes its acceptance
coins and erased record; those are not inferred to have occurred merely
because their probability formula is possible. The qubit circle and Pauli
rays are computational comparison data, not named matter modules.

\paragraph{Still unresolved.}
The historical primitive CP maps, protected record policy, and directed
V110 typing on the same physical cut have not been populated. For a
strictly resolved-word claim, V40's exclusions remain in force. For an
actually admitted classically processed CP outcome, the positive filter
cone is the correct search object. For the full original source, covariance
of one selected outcome is insufficient. For a one-ray finite-Dirac
handoff, rank four is insufficient. No historical Standard-Model selection,
new grading, actual determinant incidence, or continuum dynamics is declared.
The next source audit is which of these record policies the claimed packet
actually retains, followed by the appropriate word or positive-cone test on
that source, not a favourable change of quantum matrices.

\begingroup
\renewcommand{\refname}{References for the V41 development}
\begin{thebibliography}{3}
\bibitem{sel41-spacetime}
A. P\'elissier, \emph{Renewal Geometry and the Emergence of Lorentzian
Spacetime}, supplied manuscript
\src{Pasted text (2)(20260912-175458).txt},
Definition \src{def:supp-operational-datum} and its finite intervention
and coarse-graining assumptions.
\bibitem{sel41-instruments}
L. Lepp\"aj\"arvi and M. Sedl\'ak,
\emph{Post-processing of quantum instruments},
Phys. Rev. A \textbf{103}, 022615 (2021),
\href{https://arxiv.org/abs/2010.15816}{arXiv:2010.15816}.
\bibitem{sel41-convex}
S. Boyd and L. Vandenberghe, \emph{Convex Optimization},
Cambridge University Press (2004), Chapters 2 and 5;
\href{https://stanford.edu/~boyd/cvxbook/}{author-hosted edition}.
\end{thebibliography}
\endgroup
\addtocontents{toc}{\protect\endgroup}
% END OF SOURCE-SELECTION V41 DEVELOPMENT BLOCK.
% Preserve V1--V41 and append subsequent work before the final terminator.


\clearpage
\hypersetup{pdftitle={Historical Standard-Model Source Selection: V1--V42}}
\begin{center}
{\Large\bfseries Source-selection continuation V42}\\[.4em]
{\large Causal stopping without postselection:\\
a normalized four-event output and the retained-clock obstruction}
\end{center}
\addtocontents{toc}{\protect\begingroup}
\addtocontents{toc}{\protect\renewcommand{\protect\numberline}{\protect\makebox[2.1em][l]}}

\section[From successful filtering to a complete stopped output]{The source question is whether every run returns an output}
\label{sel42:context}

V41 constructs a circle-covariant selected outcome of the unchanged two-pulse
source, with success probability strictly below one. Its complete
success/failure instrument is not covariant. Two different ways of addressing
this failure must be distinguished. Repeating the same filter on the memory
left by failure does not restart the original input. Alternatively, a causal
rule can decide when to stop an ordinary original pulse sequence, without
rejecting any run. The first construction requires a renewal resolvent; the
second requires a prefix-consistent stopping tree, not the unconstrained
positive cone of V41.

The main positive calculation below retains the original quantum coefficients
and returns an output on every input after one to four original events. Its
coarse output is exactly circle covariant, with a nonzero $I/Y$ coefficient
cross. The displayed rule is classical and rational. There is no inverse
pulse, state reset, quantum correction, amplitude recombination, or failure
postselection. Nevertheless, the stopping duration and fine history must be
forgotten from the claimed covariant output. Retaining even the event count
invalidates that claim. A variable-duration output is not the unchanged
fixed-duration source or its full clocked instrument.

Classical postprocessing of instruments and the semantics of quantum loops
are established frameworks \cite{sel42-postprocessing,sel42-loops}. The
companion duality notes already supply a first-return Schur construction
\cite{Rigidity}, theorem \src{v5:thm:return}; no general priority is claimed
for the geometric-series formula. Here the new results are a causal stopping
polytope for the fixed source, a probability-one four-event certificate,
a sharp upper-horizon exclusion, and an exact covariance criterion for
repetition of the existing selected map. Source-dependent records and clocks
are retained in each statement.

All implementation conclusions assume that the stated classical stopping
coins and the chosen terminal cut are admitted. The new policy has not been
identified in a historical table. Its qubit representation is the existing
comparison bank, not a reconstruction of the Standard-Model matter modules.
The V1--V41 mathematical prefix is unchanged.

\section[Causal stopping is a finite affine channel problem]{A prefix-consistent tree, rather than arbitrary word weights}
\label{sel42:tree}

Let $\{\Phi_a:a\in\Sigma\}$ be the fixed original instrument on $M_n$, with
$\mathcal T=\sum_a\Phi_a$ trace preserving. Write
$\Phi_w=\Phi_{a_k}\cdots\Phi_{a_1}$ for chronological
$w=a_1\cdots a_k$, and $w^-$ for its immediate prefix. Fix $L\ge1$.
The source is required to execute at least one original event. Define
\begin{equation}
 c_\emptyset=1,\qquad c_w=0\quad(|w|=L),\qquad
 0\le c_w\le c_{w^-}\quad(1\le |w|<L).
 \label{sel42:eq:treepoly}
\end{equation}
Here $c_w$ is the product of continuation-coin probabilities along the
prefix through $w$; it is \emph{not} the probability of the original word.
After a reached prefix $w$ with $c_{w^-}>0$, continue with probability
$c_w/c_{w^-}$ and otherwise return the current memory. No subsequent quantum
operation is applied to that returned output. Choices at unreachable
prefixes have no effect. All length-$L$ runs terminate.

\begin{theorem}[Complete finite stopping-channel polytope]
\label{sel42:thm:tree}
Every rule \eqref{sel42:eq:treepoly} defines a normalized stopped instrument
\begin{equation}
 \mathcal S_w=(c_{w^-}-c_w)\Phi_w,\qquad 1\le |w|\le L,
 \qquad \sum_w\mathcal S_w^*(I)=I.
 \label{sel42:eq:stoppedmaps}
\end{equation}
Conversely, every causal classical stopping rule on these original
histories, with no extra retained quantum state and with $1\le\tau\le L$,
has such coefficients after marginalizing its private classical coins.
Its output with the whole stopping record discarded is
\begin{equation}
 \boxed{\mathcal G_c
 =\sum_{1\le |w|\le L}(c_{w^-}-c_w)\Phi_w
 =\mathcal T+\sum_{1\le |u|<L}c_u(\mathcal T-\id)\Phi_u.}
 \label{sel42:eq:telescoping}
\end{equation}
In particular, the set of complete stopped output channels at this bounded
horizon is an affine image of a compact polytope. Full-Choi covariance and
any additional specified linear source constraints are linear constraints
on $c$. Nonzero specified Choi cross is decided by maximizing both signs
of its finitely many real coordinates over that constrained polytope.
An infeasible system, or zero for every such optimum, excludes the
corresponding stopped protocol at this horizon.
\end{theorem}
\begin{proof}
Conditioned on an original prefix, the continuation decisions multiply as
ordinary conditional classical probabilities. Hence the stopping multiplier
is exactly $c_{w^-}-c_w$, which is nonnegative. Conversely, the cumulative
probability of continuing along a fixed prefix of any causal policy gives
these $c_w$, including policies that first sample a random deterministic
stopping tree. Bayes conditioning on the previously used private coins
changes the next conditional coin but not the product formula. Those coins
do not control the choice of quantum pulse; the next event is always the
same original instrument.

Expand the sum in \eqref{sel42:eq:stoppedmaps}. For every internal prefix
$u$, the positive contribution from its children is
$c_u\sum_a\Phi_a\Phi_u=c_u\mathcal T\Phi_u$, and its negative contribution
is $-c_u\Phi_u$. The root contributes $\mathcal T$, and the leaves have
$c_w=0$. This proves \eqref{sel42:eq:telescoping}. Its trace adjoint at $I$
is $I$ because $(\mathcal T-\id)^*(I)=0$. Thus normalization holds on all
inputs, including arbitrary reference-entangled inputs, without any
input-dependent division. Compactness follows from the finite inequalities
\eqref{sel42:eq:treepoly}. The Choi map is affine in $c$, giving the claimed
linear programs. A complex cross is nonzero exactly when one of its real
or imaginary coordinates is nonzero, which gives the two signed tests.
\end{proof}

The theorem supplies a stopping policy only within the class of admitted
coins and terminal cuts. If a source permits only deterministic stopping,
its continuation constraints are discrete; the full polytope is not then
an available implementation class. Finite typed grammars can be treated
with a separate tree of actually admitted prefixes and a complete outgoing
instrument at every nonterminal node. An inadmissible continuation is never
inserted.

\begin{corollary}[The original record law and the stopping budget]
\label{sel42:cor:budget}
For an input $\rho$, the probability of the terminal word $w$ is
$(c_{w^-}-c_w)\Tr\Phi_w(\rho)$. If
$\Phi_a=p_a\operatorname{Ad}_{U_a}$ with $\sum_ap_a=1$, these probabilities
are independent of $\rho$, and
\begin{equation}
 \Pr(\tau=j)=\sum_{|w|=j}p_w(c_{w^-}-c_w),\qquad
 \mathbb E\tau=1+\sum_{1\le |u|<L}p_uc_u,
 \quad p_w=\prod_{a\text{ in }w}p_a.
 \label{sel42:eq:budget}
\end{equation}
No new pulse distribution has been sampled: the stopping weights are
causally constrained multiples of the original word probabilities.
\end{corollary}
\begin{proof}
The terminal trace of \eqref{sel42:eq:stoppedmaps} proves the first assertion.
The random-unitary branch effect is $p_wI$. The tail-sum identity for a
positive bounded integer-valued random variable gives
$\mathbb E\tau=1+\sum_{j=1}^{L-1}\Pr(\tau>j)$, with
$\Pr(\tau>j)=\sum_{|u|=j}p_uc_u$.
\end{proof}

V41 permits failure because it selects from one fixed horizon. The present
prefix constraints instead require every run to return an output. They
cannot be dropped merely because a desired channel belongs to the cone of
all original word maps.

\section[A normalized four-event output from the unchanged pulses]{An exact no-postselection covariance certificate}
\label{sel42:positive}

Retain, without alteration,
\[
 k_X=\frac{9I+12iX}{25},\qquad
 k_Z=\frac{12I+16iZ}{25},\qquad
 p_X=\frac9{25},\quad p_Z=\frac{16}{25}.
\]
Words and multiplication conventions are exactly those of V41. Put
\begin{equation}
 D=6806449369850284704006144.
 \label{sel42:eq:policyden}
\end{equation}
The following table specifies $c_w=N_w/D$. Equal entries are grouped; the
empty prefix has $c_\emptyset=1$, and every length-four prefix has $c_w=0$.
The actual next continuation probability is $c_w/c_{w^-}$, not $c_w$.
\begin{equation}
\begin{array}{c|r}
 w&N_w\\\hline
 X,\ XZ&5421887613571925318400000\\
 Z,\ ZZ&5629987213880255514000000\\
 XX&3414004387497460550400000\\
 ZX,\ ZXZ&4298383783552331250000000\\
 XXX,\ XXZ&0\\
 XZX&3000486567096695875000000\\
 XZZ&1706587774631868480468750\\
 ZXX&2834965798898116125000000\\
 ZZX&3301256036512058800781250\\
 ZZZ&2392476915364671181640625
\end{array}
\label{sel42:eq:policy}
\end{equation}
The large integers are an exact rational policy encoding, not numbers
assigned to the quantum source. The companion certificate exports each
conditional coin separately and verifies all prefix inequalities by integer
arithmetic. This policy is not asserted to optimize coherence or expected
duration; its clean final channel is useful for an exact source comparison.

\begin{theorem}[Probability-one stopped channel with exact circle covariance]
\label{sel42:thm:positive}
The policy \eqref{sel42:eq:policy} is admissible in
\eqref{sel42:eq:treepoly} with $L=4$. Every input, including an unknown
reference-entangled input, returns an output with probability one after
at most four original pulses. In the Bloch axes $(X,Y,Z)$ its coarse output
has matrix
\begin{equation}
 \boxed{B_{\mathcal G}=
 \begin{pmatrix}1/8&0&-1/40\\0&1/8&0\\1/40&0&1/8\end{pmatrix}.}
 \label{sel42:eq:bloch}
\end{equation}
In the Hilbert--Schmidt orthonormal coefficient basis $(I,X,Y,Z)/\sqrt2$,
its Choi matrix is
\begin{equation}
 \boxed{J_{\mathcal G}=
 \begin{pmatrix}
 11/16&0&-i/40&0\\
 0&7/16&0&0\\
 i/40&0&7/16&0\\
 0&0&0&7/16
 \end{pmatrix}.}
 \label{sel42:eq:choi}
\end{equation}
It has trace two, Choi rank four, eigenvalues
\begin{equation}
 \frac7{16},\quad\frac7{16},\quad
 \frac9{16}-\frac{\sqrt{26}}{40},\quad
 \frac9{16}+\frac{\sqrt{26}}{40},
 \label{sel42:eq:eigen}
\end{equation}
and determinant $94129/1638400$. Its exact covariance group inside $SU(2)$
is
\begin{equation}
 \boxed{\operatorname{Cov}_{SU(2)}(\mathcal G)=H_Y
 :=\{e^{-itY}:t\in\mathbb R\},\qquad
 |\langle\!\langle I/\sqrt2|J_{\mathcal G}|Y/\sqrt2\rangle\!\rangle|
 =\frac1{40}.}
 \label{sel42:eq:group}
\end{equation}
The stopping distribution is independent of the quantum input. With
$D_\tau=10257597550531510272$, its probabilities for lengths $1,2,3,4$
are respectively the following numerators divided by $D_\tau$:
\begin{equation}
 (1885877839387158897,\ 854526994789864825,\
 3774534528676713675,\ 3742658187677772875).
 \label{sel42:eq:duration}
\end{equation}
Consequently
\begin{equation}
 \boxed{\mathbb E\tau=\frac{1868073010356757567}{641099846908219392}
 \simeq2.913856584689207.}
 \label{sel42:eq:mean}
\end{equation}
\end{theorem}
\begin{proof}
Each numerator in \eqref{sel42:eq:policy} lies between zero and its parent
numerator (with $D$ at the root). This gives nonnegative terminal
coefficients and the exact normalized instrument of Theorem
\ref{sel42:thm:tree}. To verify the channel, use the original conditional
unitary Bloch matrices
\[
 R_X=\begin{pmatrix}1&0&0\\0&-7/25&24/25\\0&-24/25&-7/25\end{pmatrix},
 \qquad
 R_Z=\begin{pmatrix}-7/25&24/25&0\\-24/25&-7/25&0\\0&0&1\end{pmatrix}.
\]
Their original average is
\begin{equation}
 T=\frac1{625}
 \begin{pmatrix}113&384&0\\-384&-175&216\\0&-216&337\end{pmatrix}.
 \label{sel42:eq:originalbloch}
\end{equation}
For $R_w=R_{a_k}\cdots R_{a_1}$, the exact identity to verify is the
small rational matrix equation
\begin{equation}
 T+\sum_{1\le|u|<4}p_uc_u(T-I_3)R_u=B_{\mathcal G}.
 \label{sel42:eq:smallcertificate}
\end{equation}
Substitution of the fourteen coefficients in \eqref{sel42:eq:policy}
proves each of its nine entries. Equivalently, form the original integer
quaternions $M_w=a_wI+i(b_wX+c_w'Y+d_wZ)$ as in V41 and sum their thirty
normalized Choi outer products with terminal weights
$p_w(c_{w^-}-c_w)$. This gives \eqref{sel42:eq:choi}, independently of the
Bloch calculation. The eigenvalues follow from its $I/Y$ block and its
other two diagonal entries; their positivity is explicit.

The symmetric part of $B_{\mathcal G}$ is $(1/8)I_3$ and its antisymmetric
part is a nonzero axial rotation generator about $Y$. An orthogonal
conjugation preserving $B_{\mathcal G}$ must preserve that oriented axial
vector. In $SO(3)$ these are exactly the rotations about $Y$, whose preimage
in $SU(2)$ is $H_Y$. This also directly verifies full-channel covariance,
not merely the selected cross. Finally substitute the same policy into
\eqref{sel42:eq:budget}. The four numerators in
\eqref{sel42:eq:duration} sum to $D_\tau$, and their weighted sum gives
\eqref{sel42:eq:mean}.
\end{proof}

No run is discarded. The policy decides whether to apply the next
\emph{original} pulse, not which pulse to apply, and no operation is
performed after returning the memory. The output is not the fixed-time
map $\Phi^4$; it is an explicitly specified stopped map. Its Choi rank four
also precludes interpreting it as the single efficient V110 source ray.

\section[Four is the sharp upper horizon for causal stopping]{Early stopping does not circumvent the three-event obstruction}
\label{sel42:minimal}

V41's separate exclusions at fixed depths one, two, and three do not by
themselves exclude a mixture of those depths. The following calculation
retains the stronger causal constraints and supplies the missing proof.

\begin{theorem}[No normalized covariant stopped channel by event three]
\label{sel42:thm:minimal}
For the unchanged pulse bank, no causal classical stopping rule with
$1\le\tau\le3$ has $H_Y$-covariant coarse quantum output. Therefore four
is the least possible upper bound on the number of original events for a
probability-one circle-covariant output. The rule in
Theorem~\ref{sel42:thm:positive} attains this bound with nonzero cross.
\end{theorem}
\begin{proof}
Every shorter rule is included at $L=3$ by setting its unused continuation
coefficients to zero. For a real Bloch matrix define
\[
 d(B)=(B_{12},B_{21},B_{23},B_{32},B_{11}-B_{33},B_{13}+B_{31})^{\mathsf T}.
\]
This is exactly the six-dimensional circle-covariance defect for unital
qubit channels. At horizon three put
$\mathbf c=(c_X,c_Z,c_{XX},c_{XZ},c_{ZX},c_{ZZ})^{\mathsf T}$.
By \eqref{sel42:eq:telescoping}, covariance is
\begin{equation}
 A\mathbf c=-d(T),\qquad
 A=\big[p_ud((T-I_3)R_u)\big]_{u=X,Z,XX,XZ,ZX,ZZ}.
 \label{sel42:eq:threeeq}
\end{equation}
The matrices are explicitly given by \eqref{sel42:eq:originalbloch} and
the two rotations above. Direct rational elimination gives
\begin{equation}
 \det A=-\frac{2^{71}3^{18}11^2\cdot47\cdot97\cdot7919}{5^{60}}
 \ne0,
 \label{sel42:eq:threedet}
\end{equation}
and its unique solution is
\begin{equation}
 \boxed{\mathbf c=\frac1{372193}
 (218750,218750,390625,390625,390625,390625)^{\mathsf T}.}
 \label{sel42:eq:threebad}
\end{equation}
The depth-two entries exceed one and exceed their parents, violating
\eqref{sel42:eq:treepoly}. Thus even this unique algebraic solution is
not a stopping policy. The determinant and the six substituted equations
are independent small exact certificates of exclusion; a floating-point
infeasibility report is not used as the proof.
\end{proof}

The minimum concerns a bounded number of original accepted pulses and all
causal stopping policies of the stated kind. It does not optimize expected
duration, the number of classical coin operations, or physical elapsed
time. In particular, the policy \eqref{sel42:eq:policy} was chosen to give
a transparent exact output, not an extremal duration or coherence value.

\section[Repeating the existing success filter is a different process]{Failure propagation survives the first-success sum}
\label{sel42:retry}

Let $\mathcal S,\mathcal F$ be the success and failure CP maps of one fixed
trial, with $\mathcal T_0=\mathcal S+\mathcal F$ trace preserving. Repeating
the trial on the \emph{same memory after each failure}, without correction
or reset, returns at attempt $j+1$ the map $\mathcal S\mathcal F^j$.
This is the usual quantum-loop composition, not an independent resampling
of the original quantum input.

\begin{theorem}[First-success channel and the exact covariance obstruction]
\label{sel42:thm:retry}
Assume $\rho(\mathcal F)<1$ as a finite superoperator. Then
\begin{equation}
 \boxed{\mathcal R=\sum_{j\ge0}\mathcal S\mathcal F^j
 =\mathcal S(\id-\mathcal F)^{-1}}
 \label{sel42:eq:retry}
\end{equation}
is CPTP and is the output with the stopping count discarded. If
$\mathcal F^*(I)=(1-s)I$ for $0<s\le1$, then the tail after $j$ trials is
$(1-s)^j$ for every input, convergence holds in diamond norm, and the mean
trial count is $1/s$.

Let $\mathcal U_g=\operatorname{Ad}_{R(g)}$, assume $\mathcal S$ is
$H$ covariant, and put $\mathcal F_g=\mathcal U_g\mathcal F\mathcal U_g^{-1}$.
For every $g\in H$,
\begin{equation}
 \boxed{\mathcal R-\mathcal U_g\mathcal R\mathcal U_g^{-1}
 =\mathcal R(\mathcal F-\mathcal F_g)(\id-\mathcal F_g)^{-1}.}
 \label{sel42:eq:retrydefect}
\end{equation}
Thus $g$ is a covariance of $\mathcal R$ exactly when
$\mathcal R(\mathcal F-\mathcal F_g)=0$. If $\mathcal S$ is injective as a
linear map on $M_n$, then
\begin{equation}
 \boxed{\mathcal R\text{ is }H\text{-covariant}
 \iff\mathcal F\text{ is }H\text{-covariant}
 \iff\mathcal T_0\text{ is }H\text{-covariant}.}
 \label{sel42:eq:retryiff}
\end{equation}
No CP inverse of $\mathcal S$ is assumed or executed.
\end{theorem}
\begin{proof}
Finite-dimensional spectral-radius convergence gives the resolvent series.
Every partial sum is CP. Since
$\mathcal S^*(I)=I-\mathcal F^*(I)$, its terminal effect is
$I-(\mathcal F^*)^{j+1}(I)$, tending to $I$. This proves complete
positivity and trace preservation of the limit. Under the scalar failure
effect, the $j$th failure map has CP diamond norm $(1-s)^j$, and the
remaining total stopped map has the same effect. This gives the uniform
tail and geometric mean count.

Covariance of $\mathcal S$ gives
$\mathcal U_g\mathcal R\mathcal U_g^{-1}
=\mathcal S(\id-\mathcal F_g)^{-1}$.
Apply the resolvent identity to obtain \eqref{sel42:eq:retrydefect}; its
rightmost factor is invertible. If $\mathcal S$ is injective, so is
$\mathcal R$, and the remaining factor vanishes precisely when
$\mathcal F=\mathcal F_g$. Adding the covariant success map proves the
last equivalence.
\end{proof}

\begin{proposition}[The unchanged V41 filter cannot be amplified this way]
\label{sel42:prop:retryexample}
Take exactly V41's four-event trial:
\[
 \mathcal S=s_*\Psi_*,\qquad
 \mathcal F=\Phi^4-\mathcal S,\qquad
 s_*=\frac{216649728}{1441015625}.
\]
The normalized $\Psi_*$ is circle covariant and injective. Repetition on
the failure output terminates almost surely, but its complete coarse
first-success channel is not even $Y$-parity covariant. Its Bloch entry
\begin{equation}
 \boxed{(B_{\mathcal R})_{12}
 =\frac{78438144473398561074048}{91995168929550527571306611}\ne0}
 \label{sel42:eq:retryentry}
\end{equation}
changes sign under the parity. Its mean number of original events is
\begin{equation}
 \boxed{\frac4{s_*}=\frac{1441015625}{54162432}
 \simeq26.605445357401972.}
 \label{sel42:eq:retrycost}
\end{equation}
\end{proposition}
\begin{proof}
The failure effect is $(1-s_*)I$, so the convergence statements follow
without resetting the state. V41 gives the Bloch matrix
\[
 P=\frac1{408125}
 \begin{pmatrix}92225&0&-132804\\0&117821&0\\132804&0&92225\end{pmatrix}.
\]
Its determinant is
$3080118177643661/67979755126953125>0$, and the identity component of
$\Psi_*$ is fixed. Hence it is an invertible linear superoperator. V41
already proves that $\Phi^4$ is not parity covariant, and
\eqref{sel42:eq:retryiff} excludes covariance of the repeated output.
The independent rational check
\[
 B_{\mathcal R}=s_*P(I_3-T^4+s_*P)^{-1}
\]
gives \eqref{sel42:eq:retryentry}. Multiplying the geometric mean trial
count by four gives \eqref{sel42:eq:retrycost}.
\end{proof}

The injectivity qualification is substantive. A reset success map can erase
some failure defects, so \eqref{sel42:eq:retryiff} is not asserted for
singular success maps. In general the exact criterion is
\eqref{sel42:eq:retrydefect}. Resetting an \emph{unknown input} after every
failure is also an additional operation, not a classical retry convention.
For example, a CPTP map on $M_n$ has a CPTP left inverse on the same $M_n$
only if it is unitary: every Kraus product in a left-inverse composition
must be scalar $I$; one nonzero product is invertible and forces all Kraus
coefficients of the channel to be proportional. The normalized failure
map in this example has Choi rank four, so no such same-memory recovery
exists once its fine record has been discarded. Keeping a fine unitary
word record can instead specify an inverse, but execution of that inverse
is a separate control assumption, not one of the original forward pulses.

\section[The clock distinguishes the complete instrument from its coarse output]{Stopping is not permission to discard a protected duration}
\label{sel42:clock}

Let $\mathcal G_j=\sum_{|w|=j}\mathcal S_w$ be the original stopped map
conditioned only on its event count. With the count retained the output is
$\bigoplus_j\mathcal G_j$, not $\mathcal G=\sum_j\mathcal G_j$.
The comparison group below leaves the classical time labels fixed.

\begin{theorem}[A sharp retained-duration obstruction at the minimal horizon]
\label{sel42:thm:clock}
No causal stopping policy for the unchanged pulses with $1\le\tau\le4$
can have an $H_Y$-covariant complete output if its stopping event count is
retained. In particular, the probability-one channel in
Theorem~\ref{sel42:thm:positive} owes its covariance to discarding that
count as well as its finer word record.

For the displayed policy put
$q_X=(9/25)(1-c_X)$ and $q_Z=(16/25)(1-c_Z)$. They are strictly positive.
The parity defect of its length-retaining instrument has the lower bound
\begin{equation}
 \boxed{
 \left\|\bigoplus_{j=1}^4
 (\operatorname{Ad}_Y\mathcal G_j\operatorname{Ad}_Y-\mathcal G_j)
 \right\|_\diamond
 \ge\frac{48}{25}\sqrt{q_X^2+q_Z^2}>0.}
 \label{sel42:eq:clockdefect}
\end{equation}
Its finely word-resolved instrument likewise is not parity covariant.
\end{theorem}
\begin{proof}
Fixed-label covariance of a direct sum is covariance of every nonzero
summand. Each $\mathcal G_j$ is a nonnegative mixture of the original
length-$j$ maps. V41's strict fixed-horizon separators exclude every
nonzero circle-covariant such mixture for $j=1,2,3$. Thus a covariant
length-retaining stopped instrument by horizon four must stop with
probability zero at those lengths. Every original word has positive
probability on every input, so each early terminal multiplier vanishes.
Induction forces $c_w=1$ for all $|w|<4$, leaving the deterministic map
$\Phi^4$ at length four. Its parity defect is nonzero by V41, a
contradiction.

For the displayed policy, the length-one map is
$q_X\operatorname{Ad}_{U_X}+q_Z\operatorname{Ad}_{U_Z}$.
Its difference from its parity conjugate, evaluated on
$\rho_Y=(I+Y)/2$, has a Bloch-output difference of length
$(48/25)\sqrt{q_X^2+q_Z^2}$. A traceless Hermitian qubit matrix with
Bloch vector $v$ has trace norm $\|v\|_2$. The corresponding block is a
contractive extraction of the complete flagged output, proving
\eqref{sel42:eq:clockdefect}. Equivalently, this one-block difference is
a Hamiltonian commutator whose diamond norm is that same number. Finally,
V40 excludes parity covariance of every nonempty individual original word;
any nonzero fine stopped branch therefore excludes fixed-word covariance
of the whole fine record.
\end{proof}

This proves a precise three-way distinction: all-input normalization at a
bounded stopping time, covariance of the quantum output after classical
record discard, and covariance of the duration-resolved process. The first
two hold in the positive example, while the third fails. The policy is a
new classical terminal-read rule, not a replacement of the protected
physical clock. It cannot establish a symmetry of the original infinitely
continued operational law.

\begin{corollary}[The original two-phase clock, without a Poisson replacement]
\label{sel42:cor:physicalclock}
Place the stopping policy on the inherited independent clock of
Section~\ref{sel36:clock}. Its phases are $H,P$; transitions have rows
$(1/5,4/5)$ and $(2/3,1/3)$, an accepted event is a $P\to H$ completion
marked with probability $\vartheta$, and the physical tick is
$d=4\vartheta\epsilon/11$. Let the initial phase law be
$\zeta=(\zeta_H,\zeta_P)$, retained rather than fitted.
For the number $K$ of clock ticks to the stopped output, define
\[
 h_H(z)=\frac{8\vartheta z^2}{15-8z+(8\vartheta-7)z^2},\qquad
 h_P(z)=\frac{2\vartheta z(5-z)}{15-8z+(8\vartheta-7)z^2}.
\]
Writing $\pi_j=\Pr(\tau=j)$ from \eqref{sel42:eq:duration}, one has
\begin{equation}
 \boxed{\mathbb E z^K=
 (\zeta_Hh_H(z)+\zeta_Ph_P(z))
 \sum_{j=1}^4\pi_jh_H(z)^{j-1},\qquad
 \mathbb E(dK)=\epsilon\left(\mathbb E\tau-
                   \frac{5\vartheta\zeta_P}{11}\right).}
 \label{sel42:eq:physicalclock}
\end{equation}
In particular, starting at phase $H$ gives exactly
$\mathbb E(dK)=\epsilon\mathbb E\tau$.
Keeping the elapsed tick count also destroys parity covariance of the
positive example, even if its pulse count is discarded.
\end{corollary}
\begin{proof}
Before an acceptance, the killed phase transition is
\[
 Q_\vartheta=\begin{pmatrix}1/5&4/5\\(2/3)(1-\vartheta)&1/3\end{pmatrix},
 \qquad b=\binom{0}{2\vartheta/3}.
\]
The two components of $z(I-zQ_\vartheta)^{-1}b$ are $h_H,h_P$.
Each acceptance returns the phase to $H$, and the clock is independent of
the pulse labels and stopping coins. Conditional on $\tau=j$, multiply
one initial waiting transform by $j-1$ copies of $h_H$, proving the first
identity. Differentiation at $z=1$ gives waiting means
$11/(4\vartheta)$ and $(11-5\vartheta)/(4\vartheta)$; multiply their sum
by the unchanged physical tick to obtain the second identity.

If $\zeta_P>0$, the earliest possible output occurs at tick one and its
map is $(2\vartheta\zeta_P/3)\mathcal G_1$. If $\zeta_P=0$, the earliest
output is at tick two and its map is $(8\vartheta/15)\mathcal G_1$.
No two-pulse output is possible at these respective earliest times.
Equation~\eqref{sel42:eq:clockdefect} gives a nonzero parity defect in
that elapsed-time block, proving the claim without a time rescaling.
\end{proof}

For another actual waiting law its own joint label/time distribution must
replace these formulas. A stopping bound of four accepted pulses is not a
four-tick bound; at fixed $\vartheta>0$ this clock has unbounded waiting
support. No readable time variable is silently dropped from a covariance
claim about the complete source.

\section[Source transfer, stability, and the remaining physical row]{What has advanced and what has not been selected}
\label{sel42:status}

For a general independently identified native carrier and comparison group,
Theorem~\ref{sel42:thm:tree} gives a finite, complete no-postselection
stopping test. If its whole stopped Choi matrix is $H$ covariant and has a
nonzero cross between identified $1$ and $\chi^{\pm1}$ coefficient sectors,
with $H=\ker\chi\subset\mathrm U(3)\times\mathrm U(2)$, the two-inclusion
argument of V41 gives exact covariance group $H$ for that coarse stopped
output. This is a conditional transfer of the stopping theorem, not an
identification of the qubit $I/Y$ comparison with those physical sectors.

The positive-support restriction is inherited without weakening:
\[
 \Ran J_{\mathcal G_c}
 =\sum_{w:c_{w^-}>c_w}\Ran J_w.
\]
Thus stopping cannot cancel a wrong typed support, manufacture a cross that
vanishes for every original word, or merge distinct Choi rays into one
efficient ray. No determinant tensor, grading matrix, directed matter map,
or finite-Dirac-shaped coherent word has been supplied by the new
rank-four output. V29's all-depth cross-zero exclusion still applies to
every such stopping policy.

At a fixed source and horizon, the affine formula gives the calibrated error
bound
\begin{equation}
 \|\mathcal G_c-\mathcal G_{\widehat c}\|_\diamond
 \le2\sum_{1\le|u|<L}|c_u-\widehat c_u|\,\|\Phi_u^*(I)\|_{\op}.
 \label{sel42:eq:stability}
\end{equation}
It follows from \eqref{sel42:eq:telescoping}, the CP diamond-norm identity,
and $\|\mathcal T-\id\|_\diamond\le2$. For the present source the weight
is exactly $p_u$. Both rules must satisfy the causal tree inequalities.
A rounded table can preserve normalization while breaking exact covariance;
a small calculated defect is not an exact symmetry certificate. Positive
coherence admits a one-sided norm margin, with the independently identified
coefficient sectors and original occurrence weights retained.

\paragraph{Proved.}
A finite causal stopping problem is an affine CP-channel polytope with
exact prefix inequalities. On the unchanged pulse source an explicit rational
policy stops every run within four events and gives the clean channel
\eqref{sel42:eq:choi}, with exact $H_Y$ covariance and cross $1/40$.
No normalized causal policy ending by event three is covariant. The mean
event count of the displayed policy is about $2.914$, without an optimality
claim for that mean. Repeating V41's old successful filter is a different
process: its geometric first-success resolvent fails even parity covariance,
as follows from a general injective-success criterion. Keeping the stopping
count rules out covariance for every bounded-four protocol in this source
class.

\paragraph{Inherited and explicitly conditional.}
The original two pulses, their probabilities, V40's resolved-word exclusion,
V41's filter and short-horizon separators, the comparison axes, and the
positive-support theorem are retained. Classical processing, finite linear
programming, and quantum-loop resolvents are standard methods. Numerical
linear programming was used only to discover a candidate stopping policy;
its theorem is certified by the displayed integer policy and independent
exact rational identities. No optimizer or numerical rank test is required
to reproduce the certificate. The original quantum matrices were not varied.

\paragraph{Still unresolved for the historical packet.}
The historical primitive quantum maps, directed V110 typing, actual terminal
policy, and which clock/record marks remain readable have not been populated.
The policy above is a possible classical protocol on the specified comparison
source, not its demonstrated historical selection. A no-postselection return
is stronger than one successful branch but weaker than covariance of the
complete time-resolved source. The next same-source audit must decide whether
its claimed operation is a resolved word, a fixed-horizon filter, a stopped
coarse output, or a protected-duration instrument. No amount of optimizing
one class can silently replace another. Historical Standard-Model selection
and continuum dynamics are not declared closed.

\begingroup
\renewcommand{\refname}{References for the V42 development}
\begin{thebibliography}{2}
\bibitem{sel42-postprocessing}
L. Lepp\"aj\"arvi and M. Sedl\'ak,
\emph{Post-processing of quantum instruments},
Phys. Rev. A \textbf{103}, 022615 (2021),
\href{https://arxiv.org/abs/2010.15816}{arXiv:2010.15816}.
\bibitem{sel42-loops}
M. Ying and Y. Feng, \emph{Quantum loop programs},
\href{https://arxiv.org/abs/quant-ph/0605218}{arXiv:quant-ph/0605218}.
\end{thebibliography}
\endgroup
\addtocontents{toc}{\protect\endgroup}
% END OF SOURCE-SELECTION V42 DEVELOPMENT BLOCK.
% Preserve V1--V42 and append subsequent work before the final terminator.


\clearpage
\hypersetup{pdftitle={Historical Standard-Model Source Selection: V1--V43}}
\begin{center}
{\Large\bfseries Source-selection continuation V43}\\[.4em]
{\large Exact covariance with the duration retained:\\
a seven-event stopping rule and renewal-clock identifiability}
\end{center}
\addtocontents{toc}{\protect\begingroup}
\addtocontents{toc}{\protect\renewcommand{\protect\numberline}{\protect\makebox[2.1em][l]}}

\section[From a four-event obstruction to a duration-resolved construction]{The retained clock changes the minimum horizon, not the quantum source}
\label{sel43:context}

V42 proves that the unchanged pulse source admits a normalized circle-covariant
coarse output with a stopping bound of four, but that no such bound preserves
circle covariance when the stopping count is retained. That is a
\emph{bounded-four} obstruction. It is not an all-horizon theorem. Rather
than discard another record, this continuation keeps the count and asks
whether the same source can meet the stronger requirement at a longer,
still bounded, stopping time.

The answer is positive at seven original events. Every run returns; every
nonzero count-conditioned map is circle covariant and has nonzero oriented
coherence. The full output retaining the independent original elapsed time
has exactly the same covariance group. An algebraic renewal identity gives
a sharp lower bound: no count-resolved covariant stopping rule can end by
five, and any such rule ending by six has fully isotropic count bins. Hence
seven is the least upper horizon for the \emph{exact circle group}, not just
for an unconstrained coherent output.

The original two coefficient maps, their probabilities, and the two-phase
clock are unchanged. Only a declared classical stopping policy is constructed.
It uses the past pulse labels, not elapsed waiting times, to decide whether
to continue. The next pulse is never controlled by that decision. The fine
word is forgotten within a count bin; the count and physical elapsed time
need not be forgotten. Admission of the displayed random coins and terminal
cuts is explicit. No historical Standard-Model source is inferred from this
qubit comparison, and its $I/Y$ coefficient rays are not assigned matter
semantics.

The inherited inputs are V41's short fixed-horizon separators and positive
support theorem, V42's causal-tree parameterization, and its literal clock
law. They are not reproved as new source constructions. Classical instrument
postprocessing is established machinery \cite{sel42-postprocessing}; the new
calculations are the duration-resolved renewal constraint, its exact
low-dimensional obstruction, the seven-event rational certificate, and the
clock-injectivity theorem below. The complete V1--V42 mathematical prefix is
retained without modification.

\section[A renewal identity for every bounded stopping rule]{A source-level constraint before searching the stopping tree}
\label{sel43:renewal}

Let $\{\Phi_a:a\in\Sigma\}$ be the fixed instrument on $M_n$, with
$\Phi=\sum_a\Phi_a$ trace preserving. Retain V42's chronological convention
$\Phi_w=\Phi_{a_j}\cdots\Phi_{a_1}$ and causal continuation weights
$c_\emptyset=1$, $0\le c_w\le c_{w^-}$, $c_w=0$ at $|w|=L$.
The map returning at event count $j$ is
\[
 \mathcal G_j=\sum_{|w|=j}(c_{w^-}-c_w)\Phi_w.
\]
It is unnormalized; the complete output is the CPTP instrument
$\rho\mapsto\bigoplus_{j=1}^L\mathcal G_j(\rho)$. The group acts on the
quantum input and output and leaves the count label fixed.

\begin{theorem}[Duration-resolved renewal constraint]
\label{sel43:thm:renewal}
Every bounded causal policy on the unchanged source satisfies
\begin{equation}
 \boxed{\Phi^L=\sum_{j=1}^L\Phi^{L-j}\mathcal G_j.}
 \label{sel43:eq:renewal}
\end{equation}
Consequently, if the first nonzero terminal count is at least $r$ and the
count-conditioned maps must lie in a specified linear covariance space
$\mathcal C$, then
\begin{equation}
 \Phi^L\in\operatorname{Ran}\mathcal M_{L,r},\qquad
 \mathcal M_{L,r}(B_r,\ldots,B_L)
       =\sum_{j=r}^L\Phi^{L-j}B_j,\quad B_j\in\mathcal C.
 \label{sel43:eq:aggregate}
\end{equation}
A linear functional annihilating this range but not $\Phi^L$ excludes every
such policy. If $\mathcal M_{L,r}$ is injective, the unnormalized terminal
maps are fixed by the original source, before the stopping coefficients are
chosen. Membership is necessary, not sufficient for a positive causal policy.
\end{theorem}
\begin{proof}
Let $\mathcal R_j=\sum_{|w|=j}c_w\Phi_w$ be the remaining, not-yet-returned
map, with $\mathcal R_0=\id$. Summing the unchanged next instrument gives
\[
 \mathcal R_j=\Phi\mathcal R_{j-1}-\mathcal G_j.
\]
Iterating and using $\mathcal R_L=0$ yields
\eqref{sel43:eq:renewal}. Equivalently, analytically continue each terminal
branch with the original instrument to count $L$. On any full length-$L$
word its terminal weights telescope to one. This continuation is a proof
of the identity, not an added operation on the returned output.
The range and injectivity statements are finite-dimensional linear algebra.
They do not impose positivity on a formal solution; that must be certified
by actual prefix weights.
\end{proof}

For a random-unitary source $\Phi_a=p_a\operatorname{Ad}_{U_a}$, each
$\mathcal G_j$ is trace-scaling and unital-scaling by the same constant
$p_j=\Pr(\tau=j)$. On traceless qubit coordinates let $T$ be the original
Bloch matrix and let $B_j$ be that of $\mathcal G_j$. Then
\begin{equation}
 \boxed{T^L=\sum_{j=1}^LT^{L-j}B_j.}
 \label{sel43:eq:blochrenewal}
\end{equation}
This smaller necessary system retains the complete covariance equations
for this random-unitary class. It does not replace the trace probabilities
$p_j$ or the causal feasibility conditions.

\section[A sharp seven-event lower bound]{Six events cannot retain the oriented circle response}
\label{sel43:lower}

Keep precisely
\[
 k_X=(9I+12iX)/25,\qquad k_Z=(12I+16iZ)/25,
 \qquad p_X=9/25,\quad p_Z=16/25.
\]
The conditional Bloch rotations and total matrix are
\begin{align}
 R_X&=\begin{pmatrix}1&0&0\\0&-7/25&24/25\\0&-24/25&-7/25\end{pmatrix},&
 R_Z&=\begin{pmatrix}-7/25&24/25&0\\-24/25&-7/25&0\\0&0&1\end{pmatrix},\nonumber\\
 T&=\frac1{625}\begin{pmatrix}113&384&0\\-384&-175&216\\0&-216&337\end{pmatrix}.
 \label{sel43:eq:originalT}
\end{align}
Put $H_Y=\{e^{-itY}:t\in\R\}\subset SU(2)$, and write
\[
 E_1=\diag(1,0,1),\quad E_2=\diag(0,1,0),\quad
 K_Y=\begin{pmatrix}0&0&-1\\0&0&0\\1&0&0\end{pmatrix}.
\]
A real Bloch matrix commutes with this circle precisely when it belongs to
$\mathcal C_Y=\Span_\R\{E_1,E_2,K_Y\}$.

\begin{theorem}[Lower horizons force failure or full isotropy]
\label{sel43:thm:lower}
For causal stopping on this unchanged source, with at least one original
event and no quantum operation after return:
\begin{enumerate}[label=\textup{(\roman*)},leftmargin=2.5em]
\item no rule with $\tau\le5$ has an $H_Y$-covariant count-retaining output;
\item if $\tau\le6$ and that output is $H_Y$ covariant, then every terminal
bin is in fact $SU(2)$ covariant, and its $I/Y$ Choi cross vanishes;
\item therefore a count-retaining output with exact covariance group $H_Y$
requires $L\ge7$.
\end{enumerate}
\end{theorem}
\begin{proof}
Fixed-count covariance is covariance of each CP summand $\mathcal G_j$.
The inherited Theorem~\ref{sel41:thm:short} excludes every nonzero positive
covariant mixture at counts $1,2,3$. Hence $\mathcal G_1=\mathcal G_2=
\mathcal G_3=0$. Since every original word has strictly positive scalar
effect, no branch can stop at these counts, so all corresponding $c_w=1$.

Use row-major vectorization of real $3\times3$ matrices. At $L=5$, the six
columns $\operatorname{vec}(TE_i),\operatorname{vec}(E_i)$, $i=1,2,3$,
span the possible right side of \eqref{sel43:eq:blochrenewal}. The row
\[
 \ell^{\mathsf T}=(0,0,1011,-252,0,-448,1011,0,0)
\]
annihilates every column, whereas
\begin{equation}
 \boxed{\ell^{\mathsf T}\operatorname{vec}(T^5)
       =-\frac{89351992774656}{3814697265625}\ne0.}
 \label{sel43:eq:fivewitness}
\end{equation}
This also excludes every smaller bound, by setting later terminal bins to
zero.

At $L=6$, form the square matrix $A_6$ whose columns, in order, are
$\operatorname{vec}(T^2E_i),\operatorname{vec}(TE_i),\operatorname{vec}(E_i)$,
with $i=1,2,3$ within each group. Its determinant is
\begin{equation}
 \boxed{\det A_6=
 \frac{224497403084481308393472}{14551915228366851806640625}\ne0.}
 \label{sel43:eq:sixdet}
\end{equation}
The unique solution is
\begin{equation}
 B_4=d_4I_3,\quad B_5=d_5I_3,\quad B_6=d_6I_3,
 \qquad
 \begin{aligned}
 d_4&=\frac{21016857999}{152587890625},\\
 d_5&=-\frac{56462835672}{3814697265625},\\
 d_6&=-\frac{47284070726}{3814697265625}.
 \end{aligned}
 \label{sel43:eq:isotropic}
\end{equation}
Direct rational multiplication verifies
$T^6=d_4T^2+d_5T+d_6I_3$. Invertibility of $A_6$ proves uniqueness,
including zero coefficients of $K_Y$. Each terminal map has scalar action
on $I$ and on the full traceless Bloch space and is therefore $SU(2)$
covariant. Its Choi matrix is diagonal in the Pauli coefficient basis, so
its $I/Y$ cross is zero. The scalar probabilities do not change this
conclusion. This is a necessary classification at six, not an assertion
that arbitrary choices of its trace probabilities are causally feasible.
\end{proof}

The calculation goes upstream of the exponentially large stopping tree:
a $9\times9$ exact linear system rules out the desired orientation at
all six-event policies. It does not infer an all-horizon obstruction from
the four-event result of V42.

\section[A rational seven-event policy]{Exact covariance in every nonzero duration bin}
\label{sel43:positive}

Set $r=77281$ and choose the following unnormalized count maps as a target
for a classical policy, not as new quantum operations:
\begin{equation}
 \begin{gathered}
 (p_4,p_5,p_6,p_7)=\left(\frac38,\frac38,\frac15,\frac1{20}\right),
 \qquad B_7=\frac1{100}I_3-\frac1{200}K_Y,\\
 B_4=d_4I_3-\frac{390625}{r}B_7,\qquad
 B_5=d_5I_3+\frac{171875}{r}B_7,\qquad
 B_6=d_6I_3-\frac{153443}{r}B_7.
 \end{gathered}
 \label{sel43:eq:targets}
\end{equation}
The $d_j$ are exactly \eqref{sel43:eq:isotropic}. More generally, for any
$B_7\in\mathcal C_Y$ these formulas give the unique $B_4,B_5,B_6$ satisfying
$T^7=T^3B_4+T^2B_5+TB_6+B_7$. This follows from the same invertible $A_6$
after left multiplication by $T$, or by direct rational substitution;
$T$ is invertible. At horizon seven a freely chosen terminal circle block
can therefore carry the orientation excluded at horizon six. Positivity
and causal realization still require the next certificate.

\subsection*{An exact rational policy certificate}

All prefixes of lengths at most three have $c_w=1$, and all length-seven
prefixes have $c_w=0$. At lengths four, five, and six use the following two
sets. A prefix in $\mathsf I$ inherits $c_w=c_{w^-}$; a prefix in
$\mathsf S$ has its own unknown $\xi_w=c_w$; every other prefix has $c_w=0$.
The order within each displayed $\mathsf S$ row is lexicographic, and the
unknowns are ordered by length and then that order.
\begin{center}
\small
\begin{tabular}{c L{.60\textwidth} L{.25\textwidth}}
\toprule
Depth & $\mathsf S$: independent continuation variables & $\mathsf I$: inherit parent\\
\midrule
4 & \texttt{XXXZ, XXZZ, XZXX, XZXZ, XZZX, XZZZ, ZXXZ, ZXZX, ZXZZ, ZZZX, ZZZZ}
  & \texttt{ZXXX, ZZXX, ZZXZ}\\[.4em]
5 & \texttt{XZXZZ, ZXXXX, ZXXXZ, ZXXZZ, ZZXXZ, ZZXZX, ZZXZZ, ZZZXZ, ZZZZX, ZZZZZ}
  & \texttt{XXZZX, XZZZX, XZZZZ}\\[.4em]
6 & \texttt{XXZZXX, XZZZZX, ZXXXZX, ZXXZZX, ZXXZZZ, ZZXZZX, ZZZZXZ, ZZZZZX, ZZZZZZ}
  & \texttt{ZXXXZZ, ZZZXZX}\\
\bottomrule
\end{tabular}
\end{center}
There are thirty independent variables. For a chronological word put
$p_w=\prod_i p_{a_i}$ and $R_w=R_{a_j}\cdots R_{a_1}$. Define the ten-entry
column, affine in $\xi$,
\begin{equation}
 F_j(\xi)=\sum_{|w|=j}p_w(c_{w^-}-c_w)
       \binom{1}{\operatorname{vec}(R_w)},\qquad j=4,5,6.
 \label{sel43:eq:Fj}
\end{equation}
The policy is \emph{defined exactly} by the thirty rational equations
\begin{equation}
 \boxed{F_j(\xi)=\binom{p_j}{\operatorname{vec}(B_j)},
                  \qquad j=4,5,6.}
 \label{sel43:eq:policylinear}
\end{equation}
Thus no rounding convention is part of its definition. If the coefficient
matrix in \eqref{sel43:eq:policylinear} is $M$ and the right-hand side
minus its constant part is $b$, the policy is $\xi=M^{-1}b$.
All source denominators in $M$ are powers of five; its determinant satisfies
\begin{equation}
 \boxed{\det M\equiv72\pmod{101}.}
 \label{sel43:eq:policydet}
\end{equation}
The reduction is the homomorphism on rationals with denominator prime to
$101$. It proves nonsingularity in characteristic zero, not a numerical
conditioning bound.

The following compact table certifies every strict conditional coin. If
$w\in\mathsf S$, write $q_w=c_w/c_{w^-}$. Its parent is nonzero, and the
listed integer $k_w$ satisfies the \emph{exact rational inequalities}
\begin{equation}
 \boxed{\frac{k_w}{1000}<q_w<\frac{k_w+1}{1000}.}
 \label{sel43:eq:coinbounds}
\end{equation}
The table is an interval certificate, not the implemented policy; the latter
is \eqref{sel43:eq:policylinear}. The other reachable conditional coins are
exactly zero or one.
\begin{center}
\small
\begin{tabular}{lr@{\qquad}lr@{\qquad}lr}
\toprule
$w$&$k_w$&$w$&$k_w$&$w$&$k_w$\\
\midrule
\texttt{XXXZ}&824&\texttt{XZXZZ}&917&\texttt{XXZZXX}&870\\
\texttt{XXZZ}&751&\texttt{ZXXXX}&658&\texttt{XZZZZX}&186\\
\texttt{XZXX}&141&\texttt{ZXXXZ}&464&\texttt{ZXXXZX}&160\\
\texttt{XZXZ}&846&\texttt{ZXXZZ}&800&\texttt{ZXXZZX}&42\\
\texttt{XZZX}&631&\texttt{ZZXXZ}&696&\texttt{ZXXZZZ}&387\\
\texttt{XZZZ}&469&\texttt{ZZXZX}&80&\texttt{ZZXZZX}&162\\
\texttt{ZXXZ}&849&\texttt{ZZXZZ}&644&\texttt{ZZZZXZ}&237\\
\texttt{ZXZX}&140&\texttt{ZZZXZ}&168&\texttt{ZZZZZX}&412\\
\texttt{ZXZZ}&455&\texttt{ZZZZX}&364&\texttt{ZZZZZZ}&687\\
\texttt{ZZZX}&842&\texttt{ZZZZZ}&796&&\\
\texttt{ZZZZ}&483&&&&\\
\bottomrule
\end{tabular}
\end{center}
Equations \eqref{sel43:eq:Fj}--\eqref{sel43:eq:policylinear}, the original
rational rotations, and the finite prefix table suffice to reproduce the
policy independently of an optimizer. The accompanying exact checker
constructs $M,b$, verifies \eqref{sel43:eq:policydet}, computes the rational
inverse solution, and checks every inequality in \eqref{sel43:eq:coinbounds}
by integer numerator signs. Its certificate also exports every exact
conditional coin, including unreachable prefixes.

\begin{theorem}[A normalized duration-retaining seven-event realization]
\label{sel43:thm:positive}
The preceding rational policy is causal. Every input, including an unknown
reference-entangled input, returns after four to seven original pulses.
Its count-conditioned maps have exactly the probabilities and Bloch blocks
in \eqref{sel43:eq:targets}. Each nonzero bin and the entire count-retaining
instrument have covariance group exactly $H_Y$ inside $SU(2)$.

Write $B_j=a_jI_3+b_jK_Y$. Their raw oriented coefficients are
\begin{equation}
 \boxed{(b_4,b_5,b_6,b_7)=
 \left(\frac{15625}{618248},-\frac{6875}{618248},
       \frac{153443}{15456200},-\frac1{200}\right).}
 \label{sel43:eq:crosses}
\end{equation}
In the Pauli coefficient basis the unnormalized Choi matrices are
\begin{equation}
 J_j=\begin{pmatrix}
 (p_j+3a_j)/2&0&-ib_j&0\\
 0&(p_j-a_j)/2&0&0\\
 ib_j&0&(p_j-a_j)/2&0\\
 0&0&0&(p_j-a_j)/2
 \end{pmatrix}\succ0.
 \label{sel43:eq:countchoi}
\end{equation}
In particular, all four bins have Choi rank four and nonzero cross.
Seven is the least possible upper horizon for a count-retaining output
with this exact circle group.
\end{theorem}
\begin{proof}
The nonsingularity certificate gives a unique rational solution, and the
strict bounds \eqref{sel43:eq:coinbounds}, together with the inherited and
zero assignments, prove every prefix inequality. The causal-tree theorem
\ref{sel42:thm:tree} therefore gives complete all-input normalization.
There is no stopping before event four. Equations
\eqref{sel43:eq:policylinear} give the first three terminal bins. Trace
normalization fixes $p_7=1/20$, and \eqref{sel43:eq:blochrenewal} fixes
$B_7$ to the value in \eqref{sel43:eq:targets}.

The Choi formula follows directly from the Pauli basis and is checked
independently by summing the 240 original length-four through length-seven
Kraus squares with their terminal weights. Each $p_j-a_j>0$ and
$(p_j+3a_j)(p_j-a_j)/4-b_j^2>0$, verified as exact rational inequalities;
hence \eqref{sel43:eq:countchoi} is strictly positive. The sign and nonzero
values of \eqref{sel43:eq:crosses} retain the original stopping probabilities.
No conditional division is needed for the complete instrument.

The symmetric part of $B_j$ is scalar, while its antisymmetric part is
$b_jK_Y\ne0$. A proper rotation commutes with this matrix exactly when it
fixes the oriented $Y$ axis. Lifting from $SO(3)$ to $SU(2)$ gives exactly
$H_Y$, including the central $\{\pm I\}$. Fixed-label covariance of the
count direct sum is the intersection of the bin covariance groups, again
$H_Y$. The lower bound is Theorem~\ref{sel43:thm:lower}.
\end{proof}

The conditional channel at the last count is especially simple:
\begin{equation}
 \boxed{B_{\mathcal G_7/p_7}=\frac15I_3-\frac1{10}K_Y.}
 \label{sel43:eq:lastconditional}
\end{equation}
The original coarse sum also has exact group $H_Y$, since
\begin{equation}
 \boxed{\sum_{j=4}^7B_j
 =\frac{21337430402174137}{294803619384765625}I_3
       +\frac{36864}{1932025}K_Y.}
 \label{sel43:eq:coarse}
\end{equation}
Neither this coarse sum nor a conditional bin is an efficient source ray.
Their positive rank does not supply the V110 one-ray incidence.

\begin{corollary}[The accepted-event budget is exact]
\label{sel43:cor:budget}
The rule has
\begin{equation}
 \boxed{\mathbb E\tau=\frac{197}{40}=4.925,\qquad
 \mathbb E\tau^2=\frac{1001}{40},\qquad
 \operatorname{Var}\tau=\frac{1231}{1600}.}
 \label{sel43:eq:countbudget}
\end{equation}
These expectations are independent of the input state. No claim is made
that the mean, variance, throughput, or cross magnitude is optimal over
seven-event rules; the optimality statement concerns only the upper horizon
for the exact group.
\end{corollary}
\begin{proof}
Insert the four probabilities of \eqref{sel43:eq:targets}. Their input
independence follows from the scalar effects of every original word.
\end{proof}

\section[Elapsed time retains the covariance information]{An independent renewal clock is injective on the count instrument}
\label{sel43:clockinjective}

The following result concerns the complete output instrument, not merely its
mean waiting time. Let $\{\mathcal G_j:j\ge1\}$ be CP maps whose sum is
trace preserving (or trace nonincreasing). Let $W_0,W_1,W_2,\ldots$ be
independent nonnegative waiting times, independent of the quantum process,
pulse labels, and stopping coins. The $W_i$, $i\ge1$, have one common law,
with $\Pr(W_1>0)>0$; $W_0$ can have a different initial law. Assume finite
waiting times almost surely. A return at count $j$ is reported at
$T_j=W_0+\cdots+W_{j-1}$. The internal comparison group leaves both count
and elapsed time fixed. Stopping decisions cannot use the realized waits.

\begin{theorem}[Equality of count and elapsed-time covariance groups]
\label{sel43:thm:clock}
Under these hypotheses the covariance group of the complete elapsed-time
instrument is exactly that of the complete count-retaining instrument:
\begin{equation}
 \boxed{G_{\rm elapsed}=\bigcap_{j\ge1}\operatorname{Cov}(\mathcal G_j)
                         =G_{\rm count}.}
 \label{sel43:eq:clockgroup}
\end{equation}
The conclusion includes bounded stopping and any almost-surely finite
stopping with summable terminal CP maps. An independent nondegenerate clock
does not erase an exact count-conditioned covariance defect, even though
several counts may contribute to one elapsed-time interval.
\end{theorem}
\begin{proof}
Let $h_0(s)=\mathbb E e^{-sW_0}$ and $h(s)=\mathbb E e^{-sW_1}$, for
$s>0$. The Laplace transform of the operator-valued elapsed-time measure is
\begin{equation}
 \widehat{\mathcal G}(s)
       =h_0(s)\sum_{j\ge1}h(s)^{j-1}\mathcal G_j.
 \label{sel43:eq:laplace}
\end{equation}
Absolute convergence follows from complete positivity: if
$E_j=\mathcal G_j^*(I)$, then
$\sum_j\|\mathcal G_j\|_\diamond
 =\sum_j\|E_j\|\le\sum_j\Tr E_j\le n$.
Fix a group element $g$ and write
$D_j=\operatorname{Ad}_{R(g)}\mathcal G_j
\operatorname{Ad}_{R(g)}^{-1}-\mathcal G_j$.
If every $D_j=0$, every elapsed-time bin is covariant by integration.
Conversely, elapsed-time covariance gives
$\sum_jh(s)^{j-1}D_j=0$ for all $s>0$, since $h_0(s)>0$.
The matrix-valued power series $F(z)=\sum_jz^{j-1}D_j$ is analytic for
$|z|<1$. The assumption $\Pr(W_1>0)>0$ makes $h(s)$ nonconstant and
strictly between zero and one on $s>0$. Its values contain an interval with
an accumulation point inside the disk. Applying the identity theorem to
every matrix entry gives $F\equiv0$, hence $D_j=0$ for every $j$.
No physical measurement implementing an inverse Laplace transform is assumed.
\end{proof}

This is an identifiability statement, not an assertion of well-conditioned
experimental inversion. If the waiting law depends on pulse labels, if it
has extra quantum correlations, or if the stopping policy uses measured
waits, the factorization \eqref{sel43:eq:laplace} need not hold and the
theorem cannot be invoked. A deterministic zero waiting time is also
excluded: it would erase the count into a single output time.

\begin{proposition}[Finite triangular inversion for the inherited clock]
\label{sel43:prop:triangular}
For the two-phase clock of V42 let
\[
 D(z)=15-8z+(8\vartheta-7)z^2,\qquad
 h(z)=\frac{8\vartheta z^2}{D(z)},\qquad
 h_P(z)=\frac{2\vartheta z(5-z)}{D(z)}.
\]
With initial law $\zeta$, put $h_0=\zeta_Hh+\zeta_Ph_P$. Write
$h(z)=az^2+O(z^3)$ and $h_0(z)=a_0z^{r_0}+O(z^{r_0+1})$. Then
\[
 a=8\vartheta/15,\qquad
 (r_0,a_0)=
 \begin{cases}(1,2\vartheta\zeta_P/3),&\zeta_P>0,\\
 (2,a),&\zeta_P=0.
 \end{cases}
\]
For a horizon $L$, the elapsed coefficients at times
$r_0,r_0+2,\ldots,r_0+2(L-1)$ form a triangular linear system for the
count maps $\mathcal G_1,\ldots,\mathcal G_L$, with determinant
\begin{equation}
 \boxed{a_0^L a^{L(L-1)/2}>0.}
 \label{sel43:eq:timedet}
\end{equation}
In particular, at initial phase $H$ and $L=7$ it is
$(8\vartheta/15)^{28}$.
\end{proposition}
\begin{proof}
The contribution of count $j$ has generating function
$h_0(z)h(z)^{j-1}\mathcal G_j$, whose first term is
$a_0a^{j-1}z^{r_0+2(j-1)}\mathcal G_j$. Later counts cannot contribute to
that power. The stated diagonal entries and determinant follow. Higher
powers give overidentifying observations. Small $\vartheta$ makes these
leading coefficients very small; exact injectivity must not be confused
with a uniform statistical margin.
\end{proof}

\section[The unchanged physical clock now passes]{The count and elapsed time are both retained in the positive example}
\label{sel43:physicalclock}

\begin{theorem}[The seven-event instrument with its original duration]
\label{sel43:thm:physical}
Place the rule of Theorem~\ref{sel43:thm:positive} on the original independent
two-phase clock, retaining $\vartheta>0$, initial law $\zeta$, and physical
tick $d=4\vartheta\epsilon/11$. The complete output retains either the
accepted count, the elapsed tick count, or both. In all three cases its
covariance group inside $SU(2)$ is exactly $H_Y$.

If $K$ is the number of ticks, its scalar generating function is
\begin{equation}
 \boxed{\mathbb E z^K=h_0(z)
    \left(\frac38h(z)^3+\frac38h(z)^4
                     +\frac15h(z)^5+\frac1{20}h(z)^6\right).}
 \label{sel43:eq:timepgf}
\end{equation}
Its original physical mean is
\begin{equation}
 \boxed{\mathbb E(dK)=\epsilon
           \left(\frac{197}{40}-\frac{5\vartheta\zeta_P}{11}\right).}
 \label{sel43:eq:physicalmean}
\end{equation}
For initial phase $H$, its physical variance is
\begin{equation}
 \boxed{\operatorname{Var}(dK)=
       \epsilon^2\frac{1102431-819520\vartheta}{193600}.}
 \label{sel43:eq:physicalvariance}
\end{equation}
The seven-event bound is not a seven-tick bound; the waiting law retains
unbounded support.
\end{theorem}
\begin{proof}
Every accepted event returns the clock to phase $H$, independently of the
quantum label. Conditional on $\tau=j$, the elapsed transform is
$h_0h^{j-1}$, giving \eqref{sel43:eq:timepgf}. Each count map is $H_Y$
covariant, so every time/count bin and every elapsed-time marginal is.
The equality of groups follows from Theorem~\ref{sel43:thm:clock}; it also
has a direct earliest-time witness here. At initial $H$, the earliest return
is tick eight, with map $(8\vartheta/15)^4\mathcal G_4$. This bin has
exactly $H_Y$, since $b_4\ne0$. With $\zeta_P>0$, the earliest return is
tick seven with multiplier $(2\vartheta\zeta_P/3)(8\vartheta/15)^3$, again
on $\mathcal G_4$. Therefore elapsed-time coarse graining introduces no
extra covariance.

The first waiting means are $11/(4\vartheta)$ and
$(11-5\vartheta)/(4\vartheta)$, exactly as in V42. Combine them with
\eqref{sel43:eq:countbudget} and multiply by the original $d$ to obtain
\eqref{sel43:eq:physicalmean}. Starting at $H$, each inter-acceptance wait
has variance $(121-104\vartheta)/(16\vartheta^2)$. Independence and the
random-sum variance identity give
\[
 \operatorname{Var}K
 =\mathbb E\tau\,\operatorname{Var}W_H
       +\operatorname{Var}\tau\,(\mathbb EW_H)^2.
\]
Substitute \eqref{sel43:eq:countbudget} and $d$ to obtain
\eqref{sel43:eq:physicalvariance}. These are fluctuations of the unchanged
clock, not a Poisson replacement or a rescaling chosen to conceal a delay.
\end{proof}

\section[Remaining fine records and structural scope]{What the longer stopping rule does not supply}
\label{sel43:status}

Retaining duration was not the final obstruction for this source. V42's
four-event coarse construction fails that stronger test, but the seven-event
policy passes it. The general clock theorem also explains why: it is the
count-conditioned maps that must be repaired, not their waiting distribution.
An independent original clock cannot hide a failed count bin.

The fine pulse word remains different. The unchanged all-depth result
\ref{sel40:thm:pulse-nogo} excludes $Y$-parity covariance for every nonempty
individual word. Our policy selects only such words. Hence the instrument
retaining that fine word is still not fixed-label parity covariant, even
though the output with count and elapsed time retained is. Classical coins
may be marginalized; a claim that also protects them or their original word
must be tested on that larger record. The new policy has not been identified
as a historical choice.

The inherited positive-support identity remains exact:
\[
 \Ran J_{\mathcal G_j}
 =\sum_{\substack{|w|=j\\c_{w^-}>c_w}}\Ran J_w.
\]
No wrong typed support is cancelled. A cross absent in every original word
cannot be created by stopping, and distinct coefficient rays are not merged
into one efficient ray. All four evaluated terminal bins have rank four;
none establishes the V110 common coherent determinant--Clifford word,
its directed matter typing, or its finite-Dirac shape. The conditional
transfer to a genuine $\mathrm U(3)\times\mathrm U(2)$ packet still requires
its actual native representation and source maps. The Pauli comparison is
not such an identification.

\paragraph{Proved.}
Every bounded stopping policy obeys the exact source renewal identity.
For the unchanged two-pulse source, that identity excludes count-resolved
circle covariance by five events and forces full isotropy at six. A fully
specified rational seven-event policy returns every input with exact circle
covariance in each nonzero duration bin and mean count $197/40$. Its complete
original elapsed-time instrument has the same exact group. Independent
renewal-time marking is injective on count instruments, with an explicit
finite triangular certificate for the original clock.

\paragraph{Inherited and explicitly conditional.}
The two original pulses, their probabilities, the short-horizon separators,
V42's causal policy class, the independent physical clock, and the
fine-word and positive-support exclusions are retained. Numerical linear
programming was used only to discover a policy support pattern. The final
policy is defined by rational equations, and all feasibility and covariance
claims have optimizer-free exact certificates. The minimum is for the upper
accepted-event horizon and the exact circle group; no mean-time or
coherence optimum is asserted.

\paragraph{Still unresolved for the historical packet.}
The actual directed V110 compiler, primitive quantum maps, and historically
admitted record policy have not been populated. The result removes an overly
broad interpretation of the bounded-four duration obstruction, not the
same-carrier occurrence requirement. The finite duality acquires no new
stationarity or continuum axiom. Further optimization of this comparison
should not substitute for identifying the historical typed source.

\addtocontents{toc}{\protect\endgroup}
% END OF SOURCE-SELECTION V43 DEVELOPMENT BLOCK.
% Preserve V1--V43 and append subsequent work before the final terminator.


\clearpage
\hypersetup{pdftitle={Historical Standard-Model Source Selection: V1--V44}}
\begin{center}
{\Large\bfseries Source-selection continuation V44}\\[.4em]
{\large Nonabelian noise and efficient recorded outcomes:\\
a sharp refinement obstruction before the one-ray handoff}
\end{center}
\addtocontents{toc}{\protect\begingroup}
\addtocontents{toc}{\protect\renewcommand{\protect\numberline}{\protect\makebox[2.1em][l]}}

\section[From stopped covariance to efficient source occurrence]{Why the circle example does not settle the nonabelian record question}
\label{sel44:context}

V43 supplies exact circle covariance in every nonzero duration bin, with
all inputs returned. Each bin nevertheless has Choi rank four. Its final
section correctly retains the distinction between that mixed operation and
a common coherent determinant--Clifford ray. Further stopping-policy
optimization does not answer whether a given covariant channel admits
\emph{any} efficient covariant recorded realization, still less whether
such a realization is the original one.

The relevant distinction is representation theoretic. For an Abelian group,
every irreducible noise representation is one dimensional. A covariant
channel therefore has a mathematical efficient fixed-label covariant
refinement. This fails in general for a nonabelian group: a nontrivial
irreducible noise sector cannot be split into efficient outcomes while
keeping every outcome fixed by the symmetry. We make this statement exact,
compute the largest covariant part that admits such a refinement, and prove
a sharp, probability-weighted lower bound on the covariance defect of
\emph{every} efficient realization. The bound cannot be reduced by merely
splitting outcome labels.

V5 already established rank constraints for an independently specified
readout algebra, in Theorem~\ref{sel5:thm:pure-capacity}. We do not claim a
new abstract rank-one effect criterion. Here the constraint algebra is
reconstructed from the \emph{actual channel symmetry}; the new quantitative
minimum, its equality characterization, and the native representation
calculations address a different question. Covariant Stinespring theory
and instrument postprocessing are established frameworks
\cite{sel10-Verdon,sel42-postprocessing}. All specialized implications below
are proved directly. They are developments relative to the audited notes,
not claims of literature priority for the representation decomposition.

An efficient outcome means a nonzero \emph{Choi-rank-one} CP map. Its Kraus
coefficient need not have operator rank one. A fixed-label covariant
refinement is an alternative instrument with the same forgotten channel,
not necessarily a classical refinement of the original records. No
covariance or native representation is inferred from dimension matching.

\section[Covariant refinements and their irreducible ranks]{The complete invariant refinement cone}
\label{sel44:cone}

Let $H$ be a compact group with identified continuous unitary input and
output representations $U_g,V_g$. On coefficient space
$\mathcal W=\mathcal H_{\rm out}\otimes\overline{\mathcal H_{\rm in}}$,
write $R_g=V_g\otimes\overline U_g$. We use the unnormalized Choi convention
$J_\Phi=\sum_a|k_a\rangle\!\rangle\langle\!\langle k_a|$, so
$\Tr_{\rm out}J_\Phi=I$ for a trace-preserving map. Fix a nonzero
$H$-covariant CP map with Choi matrix $J$, hence $[R_g,J]=0$.

In a verified isotypic decomposition, Schur's lemma gives
\begin{equation}
 \mathcal W=\bigoplus_\lambda V_\lambda\otimes M_\lambda,
 \qquad R_g=\bigoplus_\lambda u_\lambda(g)\otimes I,
 \qquad J=\bigoplus_\lambda I_{d_\lambda}\otimes A_\lambda,
 \quad A_\lambda\succeq0,
 \label{sel44:eq:isotypic}
\end{equation}
where the $u_\lambda$ are inequivalent complex irreducibles and
$d_\lambda=\dim V_\lambda$. Set $r_\lambda=\rank A_\lambda$, omit zero
blocks when taking maxima, and put $r=\rank J=\sum_\lambda d_\lambda r_\lambda$.

\begin{theorem}[Efficient refinability and irreducible outcome rank]
\label{sel44:thm:refine}
A fixed-label $H$-covariant CP decomposition $J=\sum_yJ_y$ is exactly a
family
\begin{equation}
 J_y=\bigoplus_\lambda I_{d_\lambda}\otimes A_{\lambda y},
 \qquad A_{\lambda y}\succeq0,
 \qquad\sum_yA_{\lambda y}=A_\lambda.
 \label{sel44:eq:allrefinements}
\end{equation}
Consequently:
\begin{enumerate}[label=\textup{(\roman*)},leftmargin=2.8em]
\item A complete efficient fixed-label covariant decomposition exists if
and only if $d_\lambda=1$ for every occupied block. When it exists, its
minimum number of nonzero efficient outcomes is $r$.
\item The extreme rays of the cone of covariant CP maps are
$I_{d_\lambda}\otimes|v\rangle\langle v|$, each supported in one
isotypic block. Their Choi ranks are $d_\lambda$. Decomposing the given
$J$ into these invariant-cone extreme rays needs exactly
$\sum_\lambda r_\lambda$ rays at minimum.
\item Among all complete covariant decompositions, the smallest possible
largest nonzero outcome Choi rank is
\begin{equation}
 \boxed{\min_{J=\sum_yJ_y\ {
m covariant}}\max_y\rank J_y
       =\max_{\lambda:r_\lambda>0}d_\lambda.}
 \label{sel44:eq:minmaxrank}
\end{equation}
\end{enumerate}
Every outcome is trace nonincreasing when $J$ is a channel Choi matrix.
These are existence statements for instruments with the same forgotten
map; they do not assert access to the required environment measurements.
\end{theorem}
\begin{proof}
An invariant Hermitian operator has the block form in
\eqref{sel44:eq:isotypic}; positivity and summation give
\eqref{sel44:eq:allrefinements}. Its rank is
$\sum_\lambda d_\lambda\rank A_{\lambda y}$. It equals one precisely
for a rank-one matrix in a single $d_\lambda=1$ block. A complete sum of
such matrices cannot cover an occupied block with $d_\lambda>1$.
Conversely, spectral decompositions of all the $A_\lambda$ suffice when
all occupied $d_\lambda=1$. A sum of $m$ rank-one matrices has rank at
most $m$, proving the minimum $r$.

The invariant cone is the direct sum of positive matrix cones on the
$M_\lambda$. Its extreme rays have exactly the stated form. At least
$r_\lambda$ rank-one multiplicity matrices are needed in block $\lambda$,
and its spectral decomposition attains that bound. Any covariant
outcome contributing to an occupied $\lambda$ has rank at least
$d_\lambda$, while the extreme-ray decomposition attains the minimum
in \eqref{sel44:eq:minmaxrank}. Finally $0\preceq J_y\preceq J$
implies $0\preceq\Tr_{\rm out}J_y\preceq I$.
\end{proof}

\begin{corollary}[The canonical efficiently refinable part]
\label{sel44:cor:charpart}
Let $P_{\rm char}$ project onto the sum of the $d_\lambda=1$ isotypic
spaces. Then
\begin{equation}
 \boxed{J_{\rm char}=P_{\rm char}JP_{\rm char}\preceq J}
 \label{sel44:eq:charpart}
\end{equation}
is the largest aggregate of simultaneously selected efficient
fixed-label $H$-covariant suboperations of $J$. It has an efficient
covariant decomposition, and $J-J_{\rm char}$ remains covariant and
positive. For a specified input state $\rho$, the maximum total success
probability in this mathematical class is
\begin{equation}
 \boxed{\Tr\!\left[\rho\,(\Tr_{\rm out}J_{\rm char})^{\mathsf T}\right].}
 \label{sel44:eq:charprob}
\end{equation}
The same statement holds for any selected union of character sectors,
including only the trivial character.
\end{corollary}
\begin{proof}
Every efficient invariant ray belongs to a one-dimensional group type.
Their sum $E$ is supported on $P_{\rm char}$, so $E\preceq J$ gives
$E\preceq P_{\rm char}JP_{\rm char}$. The latter is a direct sum of
positive matrices in character multiplicities and hence has the required
spectral decomposition. The input-effect formula follows from
$\Phi^*(I)=(\Tr_{\rm out}J_\Phi)^{\mathsf T}$ in this Choi convention.
\end{proof}

\begin{remark}[Connected symmetries cannot continuously permute finitely many labels]
\label{sel44:rem:labels}
For connected $H$, a continuous action on a finite outcome set is trivial.
Thus allowing continuous symmetry relabelling does not bypass the theorem
for the finite alphabets used here. A disconnected group or a continuous
outcome space has a different covariance question. Also, different efficient
outcomes may carry different characters; no common Kraus phase is required.
We therefore use the explicit phrase ``efficient fixed-label covariance''
rather than an ambiguous strong-symmetry convention.
\end{remark}

\section[A sharp defect over all efficient realizations]{An outcome-splitting-invariant quantitative obstruction}
\label{sel44:defect}

The rank obstruction is exact but can be made quantitative without imposing
a lower bound on individual outcome probabilities. For any efficient
CP decomposition $J=\sum_yJ_y$, define
\begin{equation}
 \mathfrak d_H(\{J_y\})
 =\frac1{2\Tr J}\sum_{y:J_y\ne0}
   \frac{\displaystyle\int_H
   \|R_gJ_yR_g^*-J_y\|_{\HS}^2\,dg}{\Tr J_y},
 \label{sel44:eq:defect}
\end{equation}
with Haar measure of total mass one. For a trace-preserving map on $M_N$,
$\Tr J_y/N$ is the outcome probability on the maximally mixed input.
The denominator in each term is essential. Replacing $J_y$ by $a_1J_y,\ldots,a_kJ_y$, $a_i\ge0$,
$\sum_i a_i=1$, leaves \eqref{sel44:eq:defect} unchanged. This is a
normalized quadratic covariance diagnostic, not a diamond distance or a
physical mass gap.

\begin{theorem}[Exact least covariance defect of an efficient realization]
\label{sel44:thm:sharpdefect}
For the invariant $J$ in \eqref{sel44:eq:isotypic},
\begin{equation}
 \boxed{
 \inf_{J=\sum_yJ_y,\ \rank J_y=1}\mathfrak d_H(\{J_y\})
 =\frac1{\Tr J}\sum_\lambda
       \left(1-\frac1{d_\lambda}\right)\Tr(P_\lambda J).
 }
 \label{sel44:eq:sharpdefect}
\end{equation}
The infimum is attained by a finite decomposition. Equality occurs exactly
when every nonzero normalized coefficient vector is supported on a single
isotypic block and is a product vector in $V_\lambda\otimes M_\lambda$.
In particular, the minimum is zero exactly under the efficient-refinability
condition of Theorem~\ref{sel44:thm:refine}.
\end{theorem}
\begin{proof}
Write $J_y=t_y|\psi_y\rangle\langle\psi_y|$, $\|\psi_y\|=1$.
The Haar twirl $\mathcal P_H$ is the orthogonal projection onto the
invariant operators. Unitary invariance of the Hilbert--Schmidt norm gives
\[
 \int_H\|R_gJ_yR_g^*-J_y\|_{\HS}^2dg
 =2t_y^2\left(1-\|\mathcal P_H(|\psi_y\rangle\langle\psi_y|)\|_{\HS}^2\right).
\]
For $\psi=\bigoplus_\lambda\psi_\lambda$, set
$p_\lambda=\|\psi_\lambda\|^2$ and
$\rho_\lambda=\Tr_{V_\lambda}|\psi_\lambda\rangle\langle\psi_\lambda|$.
Schur averaging yields
\[
 \mathcal P_H(|\psi\rangle\langle\psi|)
 =\bigoplus_\lambda\frac{I_{d_\lambda}}{d_\lambda}\otimes\rho_\lambda,
 \qquad
 \|\mathcal P_H(|\psi\rangle\langle\psi|)\|_{\HS}^2
 =\sum_\lambda\frac{\Tr\rho_\lambda^2}{d_\lambda}
 \le\sum_\lambda\frac{p_\lambda^2}{d_\lambda}
 \le\sum_\lambda\frac{p_\lambda}{d_\lambda}.
\]
Multiply by $t_y/\Tr J$, sum over $y$, and use
$\sum_y t_yp_{\lambda y}=\Tr(P_\lambda J)$. This proves the lower bound.
For attainability, diagonalize each $A_\lambda$ and decompose
$I_{d_\lambda}\otimes A_\lambda$ in the product basis of an orthonormal
basis of $V_\lambda$ and those multiplicity eigenvectors. Each resulting
pure vector has twirled purity $1/d_\lambda$, so equality holds.

Equality in the second inequality requires every nonzero $p_\lambda$ to
be one, since all $d_\lambda$ are finite and positive. Equality in the
first then requires the one nonzero reduced $\rho_\lambda$ to have rank
one, equivalently a product vector. Every $t_y$ is positive, so these
conditions must hold for every outcome of a minimizing decomposition.
\end{proof}

\begin{corollary}[Faithful isotropic noise gives a calibrated positive floor]
\label{sel44:cor:white}
On an $N$-dimensional input/output system with the same representation,
let $\mathcal T(X)=\Tr(X)I/N$ and
$J_\epsilon=(1-\epsilon)J_0+\epsilon I_{N^2}/N$, where $J_0$ is an
$H$-covariant channel. If
$\mathcal W=\bigoplus_\lambda V_\lambda\otimes\C^{m_\lambda}$, then
\begin{equation}
 \boxed{\mathfrak d_H^{\min}(J_\epsilon)
 =(1-\epsilon)\mathfrak d_H^{\min}(J_0)
 +\epsilon\left(1-\frac{\sum_\lambda m_\lambda}{N^2}\right).}
 \label{sel44:eq:noisebound}
\end{equation}
For any $\epsilon>0$, an occupied $d_\lambda>1$ precludes a complete
efficient fixed-label covariant realization, however many labels are used.
This does not preclude nonzero efficient covariant \emph{suboperations}.
\end{corollary}
\begin{proof}
The right-hand side of \eqref{sel44:eq:sharpdefect} is linear in $J$ when
$\Tr J=N$ is fixed. For $I_{N^2}/N$, the $\lambda$-weight is
$d_\lambda m_\lambda/N^2$. Substitute and use
$\sum_\lambda d_\lambda m_\lambda=N^2$.
\end{proof}

\section[Finite tests and the physical readout distinction]{Recover the symmetry restriction before testing source access}
\label{sel44:access}

The refinement question has a finite test even before an entire irreducible
inventory is computed. Suppose $H$ is connected. Let
$\widehat T_1,\ldots,\widehat T_s$ be the Hermitian generators of the
\emph{derived} Lie algebra $[\mathfrak h,\mathfrak h]$ in the Choi
representation and define
\begin{equation}
 L_{\rm der}=\sum_{\ell=1}^s\widehat T_\ell^2\succeq0,
 \qquad P_{\rm char}=P_{\Ker L_{\rm der}},
 \qquad \eta_H(J)=\frac{\Tr[(I-P_{\rm char})J]}{\Tr J}.
 \label{sel44:eq:finite}
\end{equation}

\begin{proposition}[A source-relative finite efficiency screen]
\label{sel44:prop:finite}
For an exactly $H$-covariant nonzero $J$, the $P_{\rm char}$ in
\eqref{sel44:eq:finite} is precisely the character-sector projection in
Corollary~\ref{sel44:cor:charpart}. In particular,
\begin{equation}
 \boxed{\eta_H(J)=0
 \iff J\text{ admits a complete efficient fixed-label covariant refinement}.}
 \label{sel44:eq:zeroscreen}
\end{equation}
When noncharacter sectors are present, with dimensions between
$d_{\min}>1$ and $d_{\max}$,
\[
 (1-d_{\min}^{-1})\eta_H(J)
 \le\mathfrak d_H^{\min}(J)
 \le(1-d_{\max}^{-1})\eta_H(J).
\]
At a fixed identified representation, $\eta_H$ is a linear Choi functional
up to its known normalization. It can be evaluated from V24/V25 contextual
testers whenever that functional belongs to their actual span.
\end{proposition}
\begin{proof}
The kernel of the positive sum consists exactly of the vectors annihilated
by all derived-algebra generators. They are fixed by the connected derived
subgroup. The action on this fixed space factors through the compact Abelian
quotient, whose irreducibles are characters. Conversely every character is
trivial on the derived subgroup. This identifies the projection. Positivity
of $J$ makes its zero leakage equivalent to support in that projection;
Theorem~\ref{sel44:thm:refine} applies. The bounds follow from
\eqref{sel44:eq:sharpdefect} term by term.
\end{proof}

\begin{proposition}[The induced readout algebra is not the original record bank]
\label{sel44:prop:dilation}
Let $J=WW^*$ with $W:E\to\mathcal W$ injective. Covariance induces the
unique unitary representation
\begin{equation}
 v_g=W^\dagger R_gW,\qquad R_gW=Wv_g.
 \label{sel44:eq:environment}
\end{equation}
A coefficient-side effect $0\preceq A\preceq I_E$ gives an invariant
suboperation $WAW^*$ exactly when $[A,v_g]=0$ for every $g$. Thus the
symmetry-allowed effect algebra is $v(H)'$, with the same irreducible
multiplicity obstruction as above. The actual physical environment effect
is $A^{\mathsf T}$ in the corresponding Kraus convention.
\end{proposition}
\begin{proof}
The support of $J$ is $R$-invariant. Hence $R_gW$ is in $\Ran W$, giving
the intertwining identity. Moreover
\[
 v_gv_g^*=W^\dagger R_gJR_g^*(W^\dagger)^*
          =W^\dagger J(W^\dagger)^*=I_E.
\]
The group law and uniqueness follow from injectivity. Invariance of
$WAW^*$ is equivalent to $W(v_gAv_g^*-A)W^*=0$, and applying the
left inverse and its adjoint gives the condition. The CP-effect
parameterization is the inherited finite Radon--Nikodym correspondence.
\end{proof}

This proposition identifies the algebra to which V5's access theorem can
be applied; it does not declare every effect in it physically available.
A minimal dilation can be a smaller mathematical presentation than the
original word-and-clock environment. Replacing the latter by the former
cannot erase a protected record in an occurrence proof. Diagonal classical
postprocessing of the actual recorded branches and a new coherent
environment measurement are different operations.

\section[The native nonabelian comparison]{A fourteen-dimensional noise calculation with exact resource counts}
\label{sel44:native}

For this comparison retain the specified V15/V16 native block-unitary
representation, restricted to $H=S(\mathrm U(3)\times\mathrm U(2))$:
\begin{equation}
 \mathcal H_{14}=C\oplus(\C^3\otimes W)\oplus\C^5,
 \qquad U(g_3,g_2,z)=z^{-2}g_3\oplus(I_3\otimes z^3g_2)\oplus I_5.
 \label{sel44:eq:native}
\end{equation}
Here $C\cong\C^3$ and $W\cong\C^2$. The two previously displayed
neutral blocks of sizes two and three are retained as the five-dimensional
neutral subspace. This is an independently declared representation, not a
new derivation of a matter carrier. The integer subscript below is the
weight of the covering $\mathrm U(1)$ parameter $z$; all types descend to
$H$. These are \emph{noise representations in operator space}, not new
particle multiplets.

\begin{theorem}[Native nonabelian noise and the efficient-record obstruction]
\label{sel44:thm:native}
For \eqref{sel44:eq:native}, the Choi representation has the following
irreducibles and multiplicities:
\begin{center}
\small
\begin{tabular}{lrr}
\toprule
Irreducible type & dimension $d_\lambda$ & multiplicity $m_\lambda$\\
\midrule
$(1,1)_0$ & 1 & 35\\
$(8,1)_0$ & 8 & 1\\
$(1,3)_0$ & 3 & 9\\
$(3,1)_{-2}$ and $(\overline3,1)_{2}$ & 3 each & 5 each\\
$(1,2)_{3}$ and $(1,2)_{-3}$ & 2 each & 15 each\\
$(3,2)_{-5}$ and $(\overline3,2)_{5}$ & 6 each & 3 each\\
\bottomrule
\end{tabular}
\end{center}
In particular
\begin{equation}
 \boxed{\sum_\lambda d_\lambda m_\lambda=196,
 \qquad\sum_\lambda m_\lambda=91,
 \qquad\dim\mathcal W_{\rm char}=35.}
 \label{sel44:eq:counts}
\end{equation}
Every strictly positive $H$-covariant Choi matrix on this carrier has no
complete efficient fixed-label $H$-covariant realization. In any covariant
realization some outcome must have Choi rank at least eight; eight is
attainable as the largest rank. An extreme-ray covariant decomposition has
minimum 91 outcomes.

For the specifically declared depolarizing comparison $J_{\mathcal T}=I_{196}/14$,
\begin{equation}
 \boxed{\eta_H(J_{\mathcal T})=\frac{23}{28},
 \qquad\mathfrak d_H^{\min}(J_{\mathcal T})=\frac{15}{28}.}
 \label{sel44:eq:nativedefects}
\end{equation}
Its maximal efficient covariant suboperation has input effect
\begin{equation}
 \boxed{E_{\rm char}
 =\frac1{42}I_C\oplus\frac3{28}I_{\C^3\otimes W}
                    \oplus\frac5{14}I_{\C^5}.}
 \label{sel44:eq:nativeeffect}
\end{equation}
Its success on the maximally mixed native input is $5/28$, and its
worst-input guaranteed success is $1/42$.
\end{theorem}
\begin{proof}
Use $\End C=1\oplus8$, $\End W=1\oplus3$, and expand every ordered
Hom block of \eqref{sel44:eq:native}. Diagonal blocks contribute
$1+9+25=35$ trivial copies, one adjoint eight, and nine weak triplets.
Hom blocks between the charged sectors and the five neutral directions
give the two conjugate colour types with multiplicity five and the two
weak types with multiplicity fifteen. Hom blocks between colour and weak
give the two six-dimensional types with multiplicity three. Target-minus-
source weights are respectively $\pm2,\pm3,\pm5$. The only characters
are the 35 trivial copies. This proves the table and \eqref{sel44:eq:counts}.
Full Choi support makes every $r_\lambda=m_\lambda$, so
Theorem~\ref{sel44:thm:refine} gives the rank and count assertions.

For depolarization, the character trace fraction is $35/196=5/28$ and
\eqref{sel44:eq:sharpdefect} is
$1-91/196=15/28$. To compute the effect, the invariant operator algebra is
\[
 \C I_3\oplus(I_2\otimes M_3)\oplus M_5
\]
up to the fixed tensor-factor ordering. On a block
$\C^d\otimes\C^m$, an orthonormal operator basis is
$I_d/\sqrt d\otimes E_{ab}$. The sum of its $K^*K$ matrices is
$(m/d)I_{dm}$. Multiplying by $1/14$ gives
\eqref{sel44:eq:nativeeffect}. Its trace divided by fourteen is $5/28$;
its least eigenvalue is $1/42$.
\end{proof}

The 35-dimensional algebra here is the commutant of the \emph{gauge action
on this carrier}; it is not the earlier 15-dimensional internal algebra
obtained as an intersection with the external tetrahedral action.
Depolarization has a larger covariance group than $H$. None of these
calculations selects $H$ as its exact group or inserts a determinant
incidence into the fundamental-only memory. The result is a diagnostic of
refinability under the stated subgroup. It applies without change to an
actual typed port once that port's own representation decomposition is
used instead of this comparison table.

For any $H$-covariant channel on \eqref{sel44:eq:native}, an explicit
$\epsilon$ depolarizing component gives
\[
 \mathfrak d_H^{\min}(J_\epsilon)
 =(1-\epsilon)\mathfrak d_H^{\min}(J_0)+\frac{15\epsilon}{28}.
\]
An arbitrarily weak positive component therefore obstructs a complete
symmetric efficient record; it does not exclude one selected coherent
suboperation. Requiring every source outcome to be efficient would be
strictly stronger than the V110 one-ray occurrence condition.

\section[The unchanged final duration bin]{An efficient symmetric realization exists, but not as a classical refinement}
\label{sel44:pulse}

We evaluate only V43's already specified count-seven conditional channel;
no new stopping coefficients are optimized. Its probability is $p_7=1/20$
and its conditional Bloch matrix is $I_3/5-K_Y/10$. In the same
orthonormal Pauli coefficient basis as V43,
\begin{equation}
 J_7^{\rm cond}=
 \begin{pmatrix}
 4/5&0&i/10&0\\
 0&2/5&0&0\\
 -i/10&0&2/5&0\\
 0&0&0&2/5
 \end{pmatrix}.
 \label{sel44:eq:J7}
\end{equation}
The $H_Y$ coefficient representation has the trivial type on
$\Span\{I,Y\}$ and characters $e^{\pm2it}$ on the lines
$\C(X\pm iZ)$. All irreducibles have dimension one, but their characters
need not agree.

\begin{proposition}[Exact efficient decomposition of the existing duration channel]
\label{sel44:prop:pulse}
Put
\begin{equation}
 \lambda_\pm=\frac{6\pm\sqrt5}{10},\qquad
 t_\pm=2\mp\sqrt5,\qquad
 K_\pm=\sqrt{\frac{\lambda_\pm}{2}}\,
            \frac{I+it_\pm Y}{\sqrt{1+t_\pm^2}},\qquad
 L_\pm=\sqrt{\frac25}\,\frac{X\pm iZ}{2}.
 \label{sel44:eq:fourkraus}
\end{equation}
Then these four coefficients realize exactly \eqref{sel44:eq:J7} and
\begin{equation}
 \boxed{\sum_{\pm}(K_\pm^*K_\pm+L_\pm^*L_\pm)=I.}
 \label{sel44:eq:normalized}
\end{equation}
Each efficient outcome is fixed-label $H_Y$-covariant. Four is the minimum
number of efficient outcomes for this conditional channel.

The largest aggregate of covariant efficient suboperations whose
coefficients \emph{commute} with $H_Y$ is the neutral $I/Y$ Choi block,
and its effect is $3I/5$. Thus both neutral outcomes can jointly be selected
with conditional success $3/5$; the largest success of one neutral
unitary outcome is
\begin{equation}
 \boxed{\frac{\lambda_+}{2}=\frac3{10}+\frac{\sqrt5}{20}.}
 \label{sel44:eq:singleprob}
\end{equation}
Before conditioning on count seven, the combined neutral success is
$3/100$. These attainability statements require the appropriate new
readout of a dilation of the count channel.
\end{proposition}
\begin{proof}
The neutral block in \eqref{sel44:eq:J7} is
$3I_2/5+M/10$, where
$M=\left(\begin{smallmatrix}2&i\\-i&-2\end{smallmatrix}\right)$ and
$M^2=5I_2$. Its unit eigenvectors are proportional to
$(1,it_\pm)^{\mathsf T}$ with eigenvalues $\lambda_\pm$.
Their unvectorizations give $K_\pm$. On the $X/Z$ plane the matrix is
$2I_2/5$; its character eigenvectors give $L_\pm$. This proves the Choi
identity. The $K$ maps are scaled unitary conjugations and have effects
$\lambda_\pm I/2$, whose sum is $3I/5$. The $L$ effects sum to $2I/5$.
Their transformation laws are $K_\pm\mapsto K_\pm$ and
$L_\pm\mapsto e^{\pm2it}L_\pm$ under $e^{-itY}$. Choi rank four gives
the minimum four. Corollary~\ref{sel44:cor:charpart}, restricted to the
trivial character, gives the aggregate maximum.

For one neutral unitary coefficient $K=\sqrt s\,V$, its Choi trace is
$2s$. The inequality $|K\rangle\!\rangle\langle\!\langle K|
\preceq J_7^{\rm cond}$ implies $2s\le\lambda_+$ by its top eigenvalue.
The coefficient $K_+$ attains it. Multiplication by the original $p_7$
gives the unconditioned probability without changing the clock.
\end{proof}

\begin{corollary}[No classical recovery of this efficient symmetric record]
\label{sel44:cor:no-classical}
No nonzero efficient $H_Y$-covariant suboperation can be produced by
nonnegative classical processing of the original nonempty fine pulse
words selected in V43, even with additional splitting or postselection of
those fine records. Hence \eqref{sel44:eq:fourkraus} is not a recovered
historical fine-word instrument.
\end{corollary}
\begin{proof}
The original word maps are efficient. If a positive sum of their Choi
matrices has rank one, every selected nonzero Choi matrix is on that same
ray, by V41's positive-support theorem. If the sum is $H_Y$-covariant,
every such ray is $H_Y$-covariant. V40's unchanged all-depth exclusion
rules this out for every nonempty original pulse word. Extra classical
splitting only changes nonnegative weights. It cannot alter the rays.
\end{proof}

This is a concrete separation between three claims: the coarse duration
channel is symmetric; it admits a mathematical symmetric efficient
realization; its original fine recorded realization is not symmetric.
The Abelian representation makes the second claim automatic in principle.
That implication cannot be transferred to a nonabelian matter packet
without the noise-sector test above. Nor does the pair of invariant Pauli
operators identify the V110 Clifford and determinant sectors.

\section[Status and source-native next test]{A noise-representation test, not another stopping rule}
\label{sel44:status}

For the finite duality, no general efficiency assumption has been added.
For a claimed common coherent source ray, the relevant questions remain
which typed ray occurs and which original records and quantum outputs
retain it. The new results separate the mathematical availability of a
covariant efficient suboperation from both a complete efficient covariant
realization and its physical access.

\paragraph{Proved.}
An exact invariant Choi decomposition classifies all fixed-label covariant
refinements. Nonabelian noise irreducibles impose minimal outcome Choi
ranks. The sharp normalized defect \eqref{sel44:eq:sharpdefect} is the
minimum over all efficient realizations and is invariant under trivial
outcome splitting. The character-supported compression is the exact
maximum efficiently refinable covariant part. On the declared native
fourteen-dimensional representation, faithful noise requires at least one
rank-eight covariant outcome; the depolarizing comparison has minimum
normalized fine-record defect $15/28$. The unchanged V43 last count channel
has the explicit four-outcome efficient symmetric realization above, but
none of those new rays is obtainable by classical processing of its old
fine records.

\paragraph{Inherited and explicitly conditional.}
V5 supplies the effect-algebra access distinction, V40 the original
fine-word exclusion, V41 positivity of coarse supports, and V43 the actual
evaluated count map. Their proofs and source are unchanged. Covariant
representation theory and dilation are established machinery. The proposed
refinements are mathematical factorizations until their readout effects
are admitted. Haar averaging in the diagnostic is statistical calculation,
not a twirling operation imposed on the source. The noise-sector table uses
its stated native action and is not the external/internal duality census.

\paragraph{Still unresolved for the historical packet.}
The actual V110 typed CP packet, its same-cut representation and its
admitted environment or record access have not been populated. The useful
upstream calculation is its nonabelian noise support and its
character-supported component, followed by the existing typed-cross and
same-word tests on that component. A zero noncharacter mass makes a
complete efficient covariant realization possible; it does not identify
the historical realization. Positive noncharacter mass obstructs that
complete refinement but need not obstruct a single common ray. Neither
branch supplies a new particle inventory, a grading, a determinant
incidence, or a continuum theorem by itself.

\addtocontents{toc}{\protect\endgroup}
% END OF SOURCE-SELECTION V44 DEVELOPMENT BLOCK.
% Preserve V1--V44 and append subsequent work before the final terminator.


\clearpage
\hypersetup{pdftitle={Historical Standard-Model Source Selection: V1--V45}}
\begin{center}
{\Large\bfseries Source-selection continuation V45}\\[.4em]
{\large Analytic fluctuation records, exact support rigidity,\\
and the sharp coherence available to a quantum readout}
\end{center}
\addtocontents{toc}{\protect\begingroup}
\addtocontents{toc}{\protect\renewcommand{\protect\numberline}{\protect\makebox[2.1em][l]}}

\section[Companion audit and the original-law interface]{What the four related developments do and do not supply}
\label{sel45:context}

V44 separates an efficiently refinable covariant Choi component from its
physical accessibility through the original records. The newly supplied
companions suggest testing that distinction on an actual fluctuating source
law, rather than designing another pulse filter. We audit their terminal
arguments and the interfaces used below; we do not independently reprove
all their inherited developments.

The Yang--Mills mass-gap file V79 retains the exact compact spin measure of
its V78, equation \src{r78:weighted}. Its new zero-momentum orientation
contact is not scalar. Its finite-range realization and volume-uniform
quadratic Fourier residual concern the \emph{first} covariance coefficient
(\src{r79:local}); the original nonlinear remainder in
\src{r79:unpaid} remains unbounded uniformly in volume. This positive
finite-graph measure, not a signed truncated covariance expansion, is the
useful transfer here.

The exact-action Einstein bridge V36 retains constant-translation rows in
its full spectral pencil. Theorem \src{v36:slow-real} finds an additional
frozen real pair of order $a^{-1}$, distinct from the rapid modes. It does
not show that every selected history excites that pair, nor give a new
nonautonomous propagation estimate. Native stationarity V46 identifies two
neutral forcing channels and obtains original constrained histories through
$s_*a^2$ (section \src{v46:sec:status}). The invariant graph is still
normally repelling; its preparation-accuracy condition is necessary, not a
new selection mechanism. Neither finite constrained system supplies a
quantum environment readout or a determinant-incidence operation.

Perturbative ESM V51 supplies an all-positive-scalar-degree local metric
current section. At the quartic stage, its exact BV subtraction retains the
scalar-quintic and metric-sextic partners
(\src{thm:v51-BV}, \src{eq:v51-next}). That primitive has ghost number
minus two. It is not the missing invariant quantum action counterterm,
a convergent infinite field series, or a positive Choi packet. The finite
matrix bounds in that file are not converted into a probability measure.

These distinctions motivate two new questions. Can ordinary postselection
of an absolutely continuous analytic fluctuation ever isolate an efficient
source ray? When a coherent environment readout is instead admitted, what
is the largest same-ray cross that the \emph{unchanged} Choi packet permits?
The first has a strong support-rigidity answer, and the second has a sharp
closed-form optimum. A compact-spin diagnostic evaluates both access
regimes without using the unpaid YM remainder. The compact diagnostic is
explicitly conditional on coupling a native relative spin to a qubit; the
YM source file does not itself assert that quantum coupling.

The real-analytic zero-set theorem and the numerical radius of a rank-one
operator are established results \cite{sel45-Mityagin,sel45-numerical}.
The support and optimization applications below, their original-law
calculations, and their physical qualifications are proved explicitly.
No priority claim is made for those underlying analytic or operator
inequalities. All V1--V44 source bytes are retained.

\section[Analytic filtering preserves the complete support]{A positive-probability record cannot isolate a proper analytic source face}
\label{sel45:rigidity}

Let $\Omega$ be a connected real-analytic manifold with its smooth volume
class. Let $\mu$ be a probability measure absolutely continuous with
respect to that class. Let
\[
 J:\Omega\longrightarrow\operatorname{Herm}(\mathcal W),\qquad
 J(x)\succeq0,
\]
be real analytic, with integrable trace, and set
$\overline J=\int J(x)\,d\mu(x)\ne0$. The matrices can describe
conditional Choi densities, including densities of mixed CP operations.
A classical acceptance policy is a measurable $f:\Omega\to[0,1]$;
its unnormalized accepted packet is
\begin{equation}
 J_f=\int f(x)J(x)\,d\mu(x).
 \label{sel45:eq:filter}
\end{equation}
No amplitude sum or environment interference is contained in this rule.

\begin{theorem}[Analytic support rigidity under classical acceptance]
\label{sel45:thm:rigidity}
Every policy with $J_f\ne0$ satisfies
\begin{equation}
 \boxed{\Ker J_f=\Ker\overline J,\qquad
        \Ran J_f=\Ran\overline J.}
 \label{sel45:eq:rigidity}
\end{equation}
In particular no positive-probability classical selection decreases Choi
rank or selects a proper coefficient-support subspace. A finite or
countable classical output bin can be efficient only if the whole
analytic source is already supported on one ray.
\end{theorem}
\begin{proof}
If $v\in\Ker\overline J$, positivity gives
$v^*J(x)v=0$ for $\mu$-almost every $x$, hence $v\in\Ker J_f$.
Conversely let $v\in\Ker J_f$. The real-valued function
$q_v(x)=v^*J(x)v$ is analytic and nonnegative. It vanishes almost
everywhere on $\{f>0\}$ with respect to $\mu$. Since $J_f\ne0$,
this set has positive $\mu$ measure. After removing the null set where
$q_v$ is nonzero it still has positive $\mu$ measure, and therefore
positive smooth volume by absolute continuity. A nontrivial real-analytic
function on a connected analytic manifold has a zero set of smooth
measure zero \cite{sel45-Mityagin}. Applied in coordinate charts and
then by analytic continuation, this gives $q_v\equiv0$ on $\Omega$.
Thus $v\in\Ker\overline J$. In finite dimension equality of kernels
is equality of ranges. Apply the same argument separately to every
nonzero classical bin.
\end{proof}

The conclusion is stronger than the inherited identity that the support
of a positive sum is the sum of its supports. Analyticity shows that
\emph{every} positive-probability subset already spans the full source.
The argument neither requires a covariance assumption nor chooses a
preferred Kraus representation. If analytic coefficients $k_j(x)$ are
supplied, $J(x)=\sum_j|k_j(x)\rangle\!\rangle\langle\!\langle k_j(x)|$
is one sufficient way to verify the analytic entrance.

\begin{corollary}[History-wise and stopping-wise form]
\label{sel45:cor:histories}
For a fixed original finite word or path stratum $w$, suppose its
conditional source has an analytic Choi density $J_w(x)$ on a connected
parameter manifold and an absolutely continuous original law. Any nonzero
classically selected part of that stratum retains
$\Ran\int J_w\,d\mu_w$. If several strata are combined with positive
weights, their retained ranges add. In particular, if every selectable
stratum has rank at least two, no nonzero finite accepted output of an
almost-surely finite stopping policy can have rank one.
\end{corollary}
\begin{proof}
Apply Theorem~\ref{sel45:thm:rigidity} within each nonzero stratum. For
a finite or convergent countable positive sum the kernel is the
intersection of the summand kernels, since its nonnegative quadratic
forms sum. At least one nonzero summand of rank two already precludes
rank one of the sum.
\end{proof}

This statement does not assume that repeatedly using one frozen field is
equivalent to drawing fresh fields. Either model must supply its own
joint path density and analytic map. The theorem is reapplied to the
composite word, not inferred from its primitive ranks: analytic factors
can multiply to a constant coefficient. A continuum-valued fine record may
have efficient \emph{densities} at individual points of measure zero.
The theorem concerns nonzero integrated outcomes; it does not deny the
existence of such densities or a covariance that moves their continuous
labels. Singular delta conditioning, quantum feedback depending on $x$,
or coherent measurement of the field dilation are different operations.

Connectedness and analyticity matter. On two disconnected components,
two different constant rays can be selected separately. Smoothness alone
also fails: on $[-1,1]$ let
$\varphi(x)=0$ for $x\le0$ and $\varphi(x)=e^{-1/x^2}$ for $x>0$.
The unit vector $(1,\varphi(x))/\sqrt{1+\varphi(x)^2}$ is smooth,
its complete preparation channel has rank two, and selecting $x<0$
returns one exact ray with probability $1/2$. It is not analytic at zero.

\section[The actual compact spin law as an access test]{A finite-graph quantum diagnostic without a covariance expansion}
\label{sel45:compact}

Retain the exact finite connected graph $G=(V,E)$ and positive
conductances $\omega_e$ from the supplied YM file. Identify $S^3$ with
$SU(2)$ and keep its normalized Haar measure. For every finite
$\gamma\ge0$ put
\begin{equation}
 d\mu_{\gamma,\omega}(u)=Z_{\gamma}(\omega)^{-1}
 e^{-\gamma\sum_{e=xy}\omega_e(1-u_x\cdot u_y)}
 \prod_{x\in V}dH(u_x).
 \label{sel45:eq:compactlaw}
\end{equation}
For $\gamma>0$ this is precisely the measure in YM V78,
\src{r78:weighted}, retained in V79. The value $\gamma=0$ below is
its explicitly declared Haar boundary comparison, not an extrapolation of
a large-activity estimate. There is no Faddeev--Popov or Wilson-exterior
factor in this particular compact spin integral; it is not relabelled as
the general Wilson law. The analytic theorem is applied to compact spin
variables, not to the discrete wire or representation-count marginal, nor
to a complex Casimir source.

Fix an edge $xy$ and define its actual relative spin
\[
 g(u)=u_xu_y^{-1}=q_0(u)I+i\sum_{j=1}^3q_j(u)\sigma_j,
 \qquad q_0=u_x\cdot u_y,\quad \sum_{j=0}^3q_j^2=1.
\]
To discuss a CP operation, declare the diagnostic coupling
$\rho\mapsto g(u)\rho g(u)^*$. It is an added interpretation of the
native relative spin on an identified qubit, not a coupling supplied by
the YM file. Averaging that specified bank uses exactly
\eqref{sel45:eq:compactlaw}:
\begin{equation}
 \Phi_{\gamma,e}(\rho)=\int g(u)\rho g(u)^*\,d\mu_{\gamma,\omega}(u),
 \qquad m_{\gamma,e}=\mathbb E(q_0^2).
 \label{sel45:eq:probe}
\end{equation}

\begin{theorem}[Full analytic support and a positive coherent-readout capacity]
\label{sel45:thm:spin}
For every fixed graph above and every finite $\gamma\ge0$,
$0<m_{\gamma,e}<1$. In the orthonormal coefficient basis
$(I,X,Y,Z)/\sqrt2$ the exact Choi matrix is
\begin{equation}
 \boxed{J_{\gamma,e}=
 \diag\left(2m_{\gamma,e},
       \frac{2(1-m_{\gamma,e})}{3}I_3\right).}
 \label{sel45:eq:spinchoi}
\end{equation}
Every nonzero accepted outcome obtained by a measurable classical rule
on the \emph{entire original spin configuration} has Choi rank four.
Nonetheless the largest $SU(2)$-covariant efficient suboperation allowed
by CP order is
\begin{equation}
 \boxed{\rho\longmapsto m_{\gamma,e}\rho.}
 \label{sel45:eq:identitypart}
\end{equation}
It has input-independent success $m_{\gamma,e}$ and is realized by a
specified coherent effect on the unmeasured field dilation. The V44
minimum normalized efficient-record defect and the exact channel distance
are
\begin{equation}
 \boxed{\mathfrak d_{SU(2)}^{\min}=\frac23(1-m_{\gamma,e}),
 \qquad \|\Phi_{\gamma,e}-\id\|_\diamond=2(1-m_{\gamma,e}).}
 \label{sel45:eq:spindiagnostics}
\end{equation}
\end{theorem}
\begin{proof}
The density in \eqref{sel45:eq:compactlaw} is strictly positive on the
connected compact manifold $(S^3)^V$. The relative spin takes every
value in $SU(2)$, so $q_0^2$ is neither identically zero nor one and
both inequalities follow. Common conjugation of the original spins
leaves their dot products, Haar measures and normalizer unchanged.
It rotates $(q_1,q_2,q_3)$ by every spatial rotation and fixes $q_0$.
Consequently
\[
 \mathbb E(q_0q_j)=0,\qquad
 \mathbb E(q_iq_j)=\frac{1-m_{\gamma,e}}3\delta_{ij}\quad(i,j>0).
\]
Expanding $|g\rangle\!\rangle\langle\!\langle g|$ gives
\eqref{sel45:eq:spinchoi}. Its support is all $M_2$: alternatively,
the values $I,iX,iY,iZ$ span that space. The density is analytic, so
Theorem~\ref{sel45:thm:rigidity} applies to every nonzero classical
accepted bin, irrespective of whether that bin is itself symmetric.

The coefficient representation is the trivial line on $I$ plus the
three-dimensional adjoint representation. An efficient fixed-label
$SU(2)$-covariant coefficient must therefore be scalar. In
\eqref{sel45:eq:spinchoi}, $s|I\rangle\!\rangle\langle\!\langle I|
\preceq J$ holds exactly for $s\le m_{\gamma,e}$, proving maximality.
For attainability on the unmeasured dilation, write
$(V\psi)(u)=g(u)\psi$ in
$L^2(\mu_{\gamma,\omega})\otimes\C^2$. It is an isometry. The unit
environment wavefunction $f_0=q_0/\sqrt{m_{\gamma,e}}$ gives
\[
 (\langle f_0|\otimes I)V
 =\frac1{\sqrt{m_{\gamma,e}}}\int q_0g\,d\mu
 =\sqrt{m_{\gamma,e}}I.
\]
The effect $|f_0\rangle\langle f_0|$ and its complement form a
normalized environment measurement. This effect has a non-diagonal
integral kernel. It is not multiplication by a classical acceptance
function on the field record. If that field has been irreversibly recorded
classically, this construction is not an available reprocessing of its
record.

The defect formula follows from V44's sharp theorem with irreducible
dimensions one and three. Finally the channel is the Pauli mixture with
weights $m_{\gamma,e},(1-m_{\gamma,e})/3$ on each Pauli. Its difference
from the identity has diamond norm at most $2(1-m_{\gamma,e})$ by the
triangle inequality. Applying it to one normalized Bell state gives four
orthogonal Bell components and attains that trace norm.
\end{proof}

At the Haar boundary $m_{0,e}=1/4$, the channel is completely depolarizing.
Its sharp efficient-record defect is $1/2$, its distance from the identity
is $3/2$, and the allowed coherent identity suboperation has success
$1/4$. No positive-probability classical field bin realizes that identity
suboperation. This comparison therefore separates CP-order availability
from the classical access regime on a full, analytic source law.

\begin{corollary}[A fixed-graph asymptotic with its correct scope]
\label{sel45:cor:laplace}
Let $r_e=b_e^{\mathsf T}L_\omega^\dagger b_e$ be the effective resistance
in the supplied compact integral. Its inherited coupling-differentiable
Laplace theorem, YM \src{r78:laplace}, implies at fixed graph
\begin{equation}
 \begin{split}
 m_{\gamma,e}&=1-\frac{3r_e}{\gamma}+O_{G,\omega}(\gamma^{-2}),\\
 \mathfrak d_{SU(2)}^{\min}
     &=\frac{2r_e}{\gamma}+O_{G,\omega}(\gamma^{-2}),\\
 \|\Phi_{\gamma,e}-\id\|_\diamond
     &=\frac{6r_e}{\gamma}+O_{G,\omega}(\gamma^{-2}).
 \end{split}
 \label{sel45:eq:asymptotic}
\end{equation}
At each finite $\gamma$, however, all nonzero classical bins still have
rank four. No volume-uniform error or physical-time estimate is claimed.
\end{corollary}
\begin{proof}
Writing $s_e=1-q_0$, exact differentiation of the original integral gives
$\mathbb E s_e=-\gamma^{-1}\partial_{\omega_e}\log Z$ and
$\operatorname{Var}s_e=\gamma^{-2}\partial_{\omega_e}^2\log Z$.
The supplied $C^2$ Laplace result and
$\partial_{\omega_e}\log\det L_\omega^o=r_e$ give
$\mathbb E s_e=3r_e/(2\gamma)+O(\gamma^{-2})$ and
$\mathbb E s_e^2=O(\gamma^{-2})$. Insert these in
$m=1-2\mathbb E s_e+\mathbb E s_e^2$ and use
\eqref{sel45:eq:spindiagnostics}. The error constants are those of the
fixed-graph theorem. The unpaid V79 original-law remainder is not used.
\end{proof}

\section[Sharp probability versus approximate efficiency]{What classical acceptance can approximate at nonzero success}
\label{sel45:tradeoff}

The exact obstruction does not imply a uniform lower bound on approximate
efficiency. In the compact diagnostic define
$\ell(u)=1-q_0(u)^2$. For a classical policy with success
$s=\int f\,d\mu>0$, the success is input independent because each
conditional operation is unitary. The normalized accepted channel has
identity entanglement-fidelity loss
\[
 \epsilon_f=\frac1s\int f\ell\,d\mu.
\]
Let $Q_\ell(v)=\inf\{t:\mu(\ell\le t)\ge v\}$ be the lower quantile.

\begin{theorem}[Optimal classical approximation of the identity branch]
\label{sel45:thm:quantile}
For every $0<s\le1$,
\begin{equation}
 \boxed{
 \inf_{0\le f\le1,\ \int f=s}
 \left\|\frac{\Phi_f}{s}-\id\right\|_\diamond
 =\frac2s\int_0^s Q_\ell(v)\,dv.}
 \label{sel45:eq:quantile}
\end{equation}
A threshold rule accepts the smallest values of $\ell$, with a fractional
coin on a boundary atom if necessary. The minimum is strictly positive
for every $s>0$ under \eqref{sel45:eq:compactlaw}.
\end{theorem}
\begin{proof}
If $f_*$ is the threshold rule at threshold $t$, then
$(f-f_*)(\ell-t)\ge0$ pointwise for every competing $f$.
The equal-success constraint removes the $t$ term, proving
$\int f\ell\ge\int f_*\ell=\int_0^sQ_\ell(v)dv$.
For any channel the trace distance on a Bell input is at least twice
its loss of overlap with that Bell state, giving the lower bound
$2\epsilon_f$. The threshold rule is conjugation invariant, since
$\ell$ is; the conditional channel is consequently a Pauli-isotropic
channel of the form used in Theorem~\ref{sel45:thm:spin}. Its diamond
distance equals $2\epsilon_{f_*}$, attaining the bound. Finally the
analytic nonconstant function $\ell$ has a zero set of measure zero,
so its integral on any set of positive acceptance mass is positive.
\end{proof}

This is an optimization over scalar field-record acceptance only. An
inverse-$g$ feedback would solve a different recovery problem by changing
the allowed quantum operations.

\begin{proposition}[Fixed acceptance mass gives a positive spectral floor]
\label{sel45:prop:floor}
In Theorem~\ref{sel45:thm:rigidity}, let $P_0=\supp\overline J$ and
fix $0<s\le1$. With quantiles computed under the unchanged $\mu$, define
\begin{equation}
 \lambda_s=\min_{v\in P_0\mathcal W,\ \|v\|=1}
       \int_0^s Q_{v^*J(\cdot)v}(u)\,du.
 \label{sel45:eq:spectral-floor}
\end{equation}
Then $\lambda_s>0$, and every classical policy with $\int f\,d\mu=s$
satisfies $J_f\succeq\lambda_sP_0$. This is the largest uniform lower
coefficient in that operator inequality. No uniformity as the original
source law or cutoff changes is asserted.
\end{proposition}
\begin{proof}
For each unit $v$ on $P_0$, its analytic quadratic response is nonzero
and thus strictly positive almost everywhere. Its lower quantile integral
is positive at every $s>0$. That integral is the infimum of
$\int f v^*Jv$ over the specified policies, by the same threshold
argument as above. As a function of $v$ it is continuous: the difference
for two unit vectors is bounded by
$2\|v-w\|\int\|J(x)\|_{\rm op}\,d\mu(x)$, which is finite.
Compactness of the unit sphere of $P_0\mathcal W$ gives a strictly
positive minimum. The variational characterization of the least
eigenvalue gives the operator bound. Conversely the minimizing vector
and its quantile policy attain the corresponding Rayleigh value, proving
optimality of the uniform coefficient.
\end{proof}

For random-unitary densities the mass $s$ is the actual input-independent
success probability. Hence approaching an efficient ray while keeping
that probability bounded below is impossible on one fixed analytic
source of rank greater than one. The floor can collapse when the original
law changes, as it does in the large-activity compact comparison; no
uniform cutoff conclusion is hidden in \eqref{sel45:eq:spectral-floor}.

\begin{proposition}[Exact Haar tradeoff and its small-success order]
\label{sel45:prop:haar}
At $\gamma=0$, $\ell$ has density
$\frac2\pi\sqrt{\ell/(1-\ell)}$ on $(0,1)$. With threshold
$\delta=\sin^2\theta$, $0<\theta\le\pi/2$, put
\begin{equation}
 s(\delta)=\frac{2\theta-\sin2\theta}{\pi},\qquad
 M(\delta)=\frac{12\theta-8\sin2\theta+\sin4\theta}{8\pi}.
 \label{sel45:eq:haar}
\end{equation}
Then the optimal distance at success $s(\delta)$ is $2M(\delta)/s(\delta)$.
In particular
\begin{equation}
 \begin{split}
 s(\delta)&=\frac4{3\pi}\delta^{3/2}(1+O(\delta)),\\
 \frac{2M(\delta)}{s(\delta)}&=\frac65\delta(1+O(\delta)).
 \end{split}
 \label{sel45:eq:haar-asymptotic}
\end{equation}
Thus the largest success compatible with diamond error at most $d_*$
is asymptotic to
$\frac4{3\pi}(5d_*/6)^{3/2}$ as $d_*\downarrow0$.
At $\delta=1/2$, the exact pair is
\begin{equation}
 \boxed{s=\frac12-\frac1\pi,\qquad
 d_{\min}=\frac{3\pi-8}{2(\pi-2)}.}
 \label{sel45:eq:haar-half}
\end{equation}
\end{proposition}
\begin{proof}
A uniform point on $S^3$ has first-coordinate density
$2\sqrt{1-q_0^2}/\pi$. The two inverse images of
$\ell=1-q_0^2$ give the stated density. Substitute
$\ell=\sin^2\theta$: its mass and first-moment integrands are
$4\sin^2\theta/\pi$ and $4\sin^4\theta/\pi$.
Their elementary integrals are \eqref{sel45:eq:haar}. Expansion near
zero gives $M(\delta)=4\delta^{5/2}/(5\pi)(1+O(\delta))$ and the
asserted rate. The quantile average is strictly increasing, so inverting
the error relation gives the maximal small-error success. Evaluating at
$\theta=\pi/4$ proves \eqref{sel45:eq:haar-half}.
\end{proof}

\section[Optimal coherent selection of the two typed shadows]{A sharp one-ray coherence formula on the original positive packet}
\label{sel45:quantum}

Return now to an independently identified finite typed Choi packet $J$,
not to the compact qubit diagnostic. Let $b,d$ be orthonormal coefficient
vectors for the desired neutral and determinant-incidence lines; put
$P=|b\rangle\langle b|+|d\rangle\langle d|$ and $Q=I-P$.
The largest positive suboperation supported on $P$ is the inherited
operator short
\begin{equation}
 S_P(J)=PJP-PJQ(QJQ)^\dagger QJP
       =\begin{pmatrix}\mu&c\\\bar c&\nu\end{pmatrix}_{b,d}.
 \label{sel45:eq:short}
\end{equation}
Using $PJP$ alone would ignore coupled source directions and can give a
false success budget. The zero/nonzero pattern of the true short is
retained, including singular cases.

\begin{theorem}[Largest same-ray cross permitted by CP order]
\label{sel45:thm:coherent}
For the unchanged $J$,
\begin{equation}
 \boxed{
 \sup_{v\in P\mathcal W,\ |v\rangle\langle v|\preceq J}
 |\langle b,v\rangle|^2|\langle d,v\rangle|^2
 =\frac14\bigl(\sqrt{\mu\nu}+|c|\bigr)^2.}
 \label{sel45:eq:optimum}
\end{equation}
The supremum is attained. It is positive exactly when $\mu\nu>0$,
even if the original short has $c=0$. If these two coefficient lines
transform by $1$ and $(\det U\det V)^{-1}$, every maximizing nonzero
cross ray is covariant exactly under $S(\mathrm U(3)\times\mathrm U(2))$
inside the specified comparison group.

Attainability is through an arbitrary allowed effect of a minimal
Stinespring environment. With a restricted original readout algebra,
\eqref{sel45:eq:optimum} is only an upper bound; in the analytic classical
regime of Theorem~\ref{sel45:thm:rigidity} it can be positive while the
accessible efficient yield is zero.
\end{theorem}
\begin{proof}
Positivity of $J$ ensures the range condition for the generalized Schur
complement. Therefore $|v\rangle\langle v|\preceq J$, for $v\in P$,
is equivalent to $|v\rangle\langle v|\preceq S_P(J)$. Write
$v=S_P(J)^{1/2}u$, $\|u\|\le1$. With
$a=S_P(J)^{1/2}b$ and $z=S_P(J)^{1/2}d$, the unsquared cross is
$|\langle a,u\rangle\langle u,z\rangle|$.
Its maximum is the numerical radius of the rank-one operator
$|z\rangle\langle a|$:
\[
 \frac{\|a\|\|z\|+|\langle a,z\rangle|}{2}
 =\frac{\sqrt{\mu\nu}+|c|}{2}.
\]
For completeness, normalize nonzero $a,z$ to $e,f$ and choose their
relative phase so $\langle e,f\rangle\ge0$. The bound follows from
\[
 |\langle e,u\rangle\langle u,f\rangle|
 \le\tfrac12\bigl(|\langle e,u\rangle|^2+|\langle f,u\rangle|^2\bigr)
 \le\tfrac12(1+|\langle e,f\rangle|).
\]
The unit vector parallel to $e+f$ attains it; zero and collinear cases
follow directly. This is the standard rank-one numerical-radius identity
\cite{sel45-numerical}. Squaring proves \eqref{sel45:eq:optimum}.

More explicitly, when $\mu,\nu>0$ a maximizing vector has squared
components
\[
 x=\frac\mu2\left(1+\frac{|c|}{\sqrt{\mu\nu}}\right),\qquad
 y=\frac\nu2\left(1+\frac{|c|}{\sqrt{\mu\nu}}\right),
\]
with its off-diagonal phase equal to that of $c$ (arbitrary if $c=0$).
The residual short is positive and has rank at most one. A ray with both
components nonzero is fixed up to scalar phase under the two-character
action precisely when their characters agree, namely when
$\det U\det V=1$.

Finally let $J=WW^*$ be a minimal injective Kraus synthesis. The selected
rank-one suboperation corresponds to an effect
$|W^\dagger v\rangle\langle W^\dagger v|\preceq I$ in the coefficient
convention, or its transpose in the physical Kraus-environment convention.
Its complement retains normalization. This is the already established
CP Radon--Nikodym correspondence, not proof that this effect is present
in the original accessible record algebra.
\end{proof}

A zero original cross can therefore coexist with a positive mathematical
coherent selection budget. The compensating remainder then has an opposite
cross; their forgotten sum is unchanged. This does not create an efficient
ray through positive classical addition.

\begin{proposition}[New optimum on the unchanged V43 last duration channel]
\label{sel45:prop:oldbin}
Keep exactly $J_7^{\rm cond}$ from \eqref{sel44:eq:J7}, with its
original probability $p_7=1/20$. On $P=\Span\{I,Y\}$ in the orthonormal
coefficient convention, its short is
$\left(\begin{smallmatrix}4/5&i/10\\-i/10&2/5\end{smallmatrix}\right)$.
The largest same-ray cross mass is
\begin{equation}
 \boxed{h_{\max}=\frac{33+8\sqrt2}{400}.}
 \label{sel45:eq:oldopt}
\end{equation}
It is attained by
\begin{equation}
 K_*=(\sqrt{x}\,I-i\sqrt{y}\,Y)/\sqrt2,\qquad
 x=\frac{8+\sqrt2}{20},\quad y=\frac{8+\sqrt2}{40},
 \label{sel45:eq:oldkraus}
\end{equation}
with input-independent conditional success
$3(8+\sqrt2)/80$. Its original unconditioned success is
$3(8+\sqrt2)/1600$, and its unconditioned cross mass is
$(33+8\sqrt2)/160000$.
\end{proposition}
\begin{proof}
There are no cross blocks from $P$ to its complement in the retained
matrix, so the short is its literal neutral block. Substitute
$\mu=4/5,\nu=2/5,|c|=1/10$ in
\eqref{sel45:eq:optimum}. Equation \eqref{sel45:eq:oldkraus} has the
maximizing squared components and phase. Its effect is
$(x+y)I/2$, and the neutral residual is positive rank one. The complete
residual has rank three. Multiplication by $p_7$, and by $p_7^2$ for the
squared cross, proves the raw budgets.
\end{proof}

This objective differs from V44's maximum single-unitary \emph{success}
$\lambda_+/2$. Neither optimum improves the other objective. The original
pulse records still cannot access this new ray by classical processing,
by Corollary~\ref{sel44:cor:no-classical}. There is no new pulse,
stopping rule, elapsed-time rescaling or historical matter identification.

\section[Formal jets do not certify exact support]{Why the perturbative companion does not resolve the access question}
\label{sel45:formal}

\begin{proposition}[Equal perturbative data, different finite-rank refinability]
\label{sel45:prop:flat}
For $a>0$ set $t(a)=e^{-1/a^2}$, and $t(0)=0$. On a defining qubit let
\[
 \Phi_a^{(0)}=\id,\qquad
 \Phi_a^{(1)}=(1-t(a))\id+t(a)\mathcal T_2,
 \quad \mathcal T_2(\rho)=\Tr(\rho)I_2/2.
\]
These are normalized $SU(2)$-covariant channels, smooth at $a=0$, with
identical Taylor coefficients there to every order. At every $a>0$ the
first Choi matrix has rank one and the second rank four. V44 gives
\begin{equation}
 \boxed{\mathfrak d_{SU(2)}^{\min}(\Phi_a^{(0)})=0,\qquad
 \mathfrak d_{SU(2)}^{\min}(\Phi_a^{(1)})=t(a)/2>0.}
 \label{sel45:eq:flat}
\end{equation}
Their diamond distance is exactly $3t(a)/2$. For any prescribed finite
jet order $r$, replacing $t(a)$ by $a^{r+1}$ on $0\le a<1$ gives the
same obstruction with families analytic at zero and agreeing through
order $r$.
\end{proposition}
\begin{proof}
The mixed Choi matrix is
$(1-t)|I\rangle\!\rangle\langle\!\langle I|+(t/2)I_4$.
Its eigenvalues are $2-3t/2$ on the identity line and $t/2$ on each
adjoint axis. All are positive for $0<t<1$. Its normalized nontrivial
weight is $3t/4$, giving \eqref{sel45:eq:flat} by the adjoint dimension
three. The Bell-input calculation in
Theorem~\ref{sel45:thm:spin} gives distance $3t/2$. Every derivative
of $e^{-1/a^2}$ tends to zero at the origin since it is that exponential
times a polynomial in $1/a$. The finite-order analytic example is immediate.
\end{proof}

The all-order comparison is deliberately smooth, not analytic at zero.
An actual convergent analytic germ with all its coefficients fixed would
remove that particular freedom. Neither coefficientwise perturbation
nor the local current theorem in ESM V51 asserts such a germ. No claim is
made that its actual action contains the displayed noise. The proposition
only prevents perturbative identities, without a justified positive
reconstruction and remainder, from deciding an exact Choi-support claim.
It neither refutes nor strengthens the companion's Noether/BV theorem.

\section[Status and next source row]{A sharp distinction between fluctuation records and coherent access}
\label{sel45:status}

\paragraph{Proved.}
An absolutely continuous analytic source on a connected parameter manifold
has a rigid Choi support under every nonzero classical acceptance policy.
The result also applies stratum by stratum to finite histories and stopping
policies. On the original compact-spin law, a specified relative-spin
qubit diagnostic has an exact scalar-plus-adjoint Choi matrix and full
rank at every finite activity. Its coherent identity suboperation has
success $m_{\gamma,e}>0$, but no nonzero classical field bin is efficient.
The exact quantile theorem gives the optimal success--error tradeoff for
classical approximation, including the Haar $3/2$ success exponent. A
separate short-operator optimization determines the largest same-ray
neutral--determinant cross and is evaluated on the unchanged V43 duration
matrix. Equal formal jets do not determine exact finite-source rank.

\paragraph{Inherited and conditional.}
The positive compact measure and its fixed-graph differentiable Laplace
estimate are source inputs from YM V78 retained in V79. No volume-uniform
remainder is inferred. The graph law is not a general Wilson exterior.
The diagnostic qubit coupling is declared here and is not a historical YM
or SM operation. The character/refinement theorem, accessible-effect
distinction and old pulse-record exclusion are inherited from V44 and its
specified predecessors. Coherent environment effects are mathematical
possibilities until actually admitted; none is substituted for a recorded
classical field. The other companions supply constraint, spectral and
local-cohomological information, not those effects.

\paragraph{Still unresolved for the historical packet.}
No literal V110 typed source or joint quantum readout was populated.
The next question is whether its source variability belongs to an
absolutely continuous analytic branch, a discrete or singular branch,
or an accessible coherent dilation. A positive analytic classical branch
of rank greater than one cannot be made into the required single ray by
finer postselection; a zero-support jet or signed covariance coefficient
does not change this. If coherent access is independently present, the
actual short \eqref{sel45:eq:short}, not marginal masses or a fitted
reference packet, fixes the attainable cross. The common typing, grading,
record policy and physical clock remain separate source rows. Complete
efficient covariance is not added as an axiom of the finite duality.

\begingroup
\renewcommand{\refname}{Additional primary-source context for V45}
\begin{thebibliography}{9}
\small\raggedright
\bibitem{sel45-Mityagin}
B.~S. Mityagin, \emph{The zero set of a real analytic function},
Mathematical Notes \textbf{107} (2020), 529--530;
\url{https://arxiv.org/abs/1512.07276}.
Used only for the analytic zero-set property in the support proof.
\bibitem{sel45-numerical}
D.~Den\v ci\'c, H.~Stankovi\'c, M.~Krsti\'c and I.~Damnjanovi\'c,
\emph{$q$-numerical radius of rank-one operators and the generalized
Buzano inequality} (2025),
\url{https://arxiv.org/abs/2503.05036}.
Primary-source context for the standard rank-one numerical-radius identity;
the $q=1$ case used here is proved directly.
\end{thebibliography}
\endgroup
\addtocontents{toc}{\protect\endgroup}
% END OF SOURCE-SELECTION V45 DEVELOPMENT BLOCK.
% Preserve V1--V45 and append subsequent work before the final terminator.


\clearpage
\hypersetup{pdftitle={Historical Standard-Model Source Selection: V1--V46}}
\begin{center}
{\Large\bfseries Source-selection continuation V46}\\[.4em]
{\large Partial readout, inaccessible records,\\
and the exact boundary of coherent source selection}
\end{center}
\addtocontents{toc}{\protect\begingroup}
\addtocontents{toc}{\protect\renewcommand{\protect\numberline}{\protect\makebox[2.1em][l]}}

\section[Audit the joint readout, not another source completion]{What the unchanged system channel does not determine}
\label{sel46:context}

V45 proves both analytic classical-support rigidity and the largest
neutral--determinant cross available under unrestricted coherent readout.
Those are different access classes. The remaining step is not to choose one
of them by terminology: part of a dilation may be quantum-accessible while
another part has already recorded information irreversibly relative to the
admitted operations. We therefore keep the system channel fixed and include
its \emph{actual accessible output} in the source data.

The CP-order correspondence, the algebra-constrained common-input maximum
and the positive-sum support theorem are inherited from
Theorems~\ref{sel5:thm:access-short}, \ref{sel5:thm:pure-capacity} and the
later refinement developments. The new statements permit an arbitrary
mixed system--readout extension, not merely effects in a represented
subalgebra of a noiseless environment. In particular the pullbacks of all
readout effects may span the full matrix space while their positive cone
still excludes every desired efficient outcome. The tests below distinguish
these possibilities without changing any old system map.

Environment-assisted operations and CP Radon--Nikodym theory are established
frameworks \cite{sel46-Raginsky,sel46-Gregoratti}. The kernel test has the
same positivity mechanism as unambiguous discrimination of mixed states
\cite{sel46-Feng}. We prove the finite formulas needed here and make no
priority claim for those general frameworks. All readout-channel comparisons
below are declared comparisons, not reconstructions of an unprovided
historical environment. The V45 compact field law and all four companion
interfaces remain unchanged; no new companion theorem is required.

\section[The access-constrained maximum from one positive matrix]{An exact target-support test on a mixed joint output}
\label{sel46:access}

Let $\mathcal W=H_{\rm out}\otimes\overline{H_{\rm in}}$ be the
identified system coefficient space. Write the actual joint Choi operator
of the system and its accessible readout $A$ as
\begin{equation}
 \Omega\succeq0\quad\text{on }\mathcal W\otimes A,
 \qquad J=\Tr_A\Omega,
 \qquad \Tr_{\rm out}J=I_{\rm in}.
 \label{sel46:eq:joint}
\end{equation}
An inaccessible environment has already been traced out in $\Omega$.
The same formulas apply to a subnormalized CP packet, retaining its original
normalization. A terminal effect $0\preceq F\preceq I_A$ selects
\begin{equation}
 J_F=\Tr_A\bigl[(I\otimes F^{1/2})\Omega
                    (I\otimes F^{1/2})\bigr],\qquad J_{I-F}=J-J_F.
 \label{sel46:eq:selection}
\end{equation}
The equality can be evaluated linearly in $F$ by partial-trace cyclicity.
The maximizations below assume all effects on the identified $A$ are
admitted. Otherwise their values are upper bounds and the displayed effects
must separately pass the actual access restriction. Arbitrary local
readout processing, ancillary systems and terminal measurements reduce to
such an effect. Conditional feedback on the system itself is not included.

Put $R_A=\Tr_{\mathcal W}\Omega$ and work on $A_0=\supp R_A$.
Positivity makes every block of $\Omega$ outside $A_0$ zero, so this
removes only an identically unused readout direction. For a specified
coefficient-support projector $P$, let $Q=I-P$ and define
\begin{equation}
 D_P=\Tr_{\mathcal W}\bigl[(Q\otimes I)\Omega(Q\otimes I)\bigr],
 \qquad q_P=P_{\Ker D_P}\quad\text{on }A_0.
 \label{sel46:eq:bad-record}
\end{equation}
Then $0\preceq D_P\preceq R_A$, and $R_A$ is faithful on $A_0$.

\begin{theorem}[Maximal accessible CP packet on a specified support]
\label{sel46:thm:max}
For every effect on $A_0$,
\begin{equation}
 \Ran J_F\subseteq P\mathcal W
 \quad\Longleftrightarrow\quad
 0\preceq F\preceq q_P.
 \label{sel46:eq:feasible}
\end{equation}
Consequently
\begin{equation}
 \boxed{J^{\rm acc}_P:=J_{q_P}}
 \label{sel46:eq:max}
\end{equation}
is the unique largest accessible positive suboperation supported on $P$.
It satisfies $J^{\rm acc}_P\preceq S_P(J)$, with $S_P(J)$ the full
operator short used in V45. A nonzero such suboperation exists exactly when
$q_P\ne0$.
\end{theorem}
\begin{proof}
For positive $J_F$, its support lies in $P$ exactly when
$\Tr(QJ_F)=0$. By partial-trace cyclicity,
\[
 \Tr(QJ_F)=\Tr(FD_P)
          =\|D_P^{1/2}F^{1/2}\|_{\HS}^2.
\]
This vanishes exactly when $\Ran F\subseteq\Ker D_P$.
For an effect this is equivalent to $F\preceq q_P$. The map
$F\mapsto J_F$ preserves order, proving maximality. Every $J_F$ is below
$J$, hence a supported one is below its operator short. Finally
$\Tr J_{q_P}=\Tr(q_PR_A)>0$ when $q_P\ne0$, since $R_A$ is faithful
on $A_0$. The converse is immediate.
\end{proof}

\begin{corollary}[Exact capacity of one proposed coherent ray]
\label{sel46:cor:ray}
For a unit $v\in\mathcal W$, put $P_v=|v\rangle\langle v|$ and
$q_v=P_{\Ker D_{P_v}}$ on $A_0$. The largest accessible coefficient
mass on that ray is
\begin{equation}
 \boxed{\alpha_v=\Tr(q_vR_A),\qquad J^{\rm acc}_{P_v}=\alpha_vP_v.}
 \label{sel46:eq:ray-capacity}
\end{equation}
Thus a nonzero efficient suboperation on $v$ exists if and only if the
accessible readout has a nonzero effect annihilating the entire unwanted
coefficient sector. Its maximum mean success on the maximally mixed
system input is $\alpha_v/\dim H_{\rm in}$.
\end{corollary}
\begin{proof}
Apply the theorem to a one-dimensional support and take the trace.
The mean-success formula uses the original Choi convention
$\Tr J_F/\dim H_{\rm in}$; it does not assert state-independent success.
\end{proof}

This is an unambiguous-exclusion test: a conclusive readout must have zero
probability on every unwanted coefficient direction. It does not ask for
an inverse channel or for recovery of all quantum information. If
$\Omega=|\Psi\rangle\langle\Psi|$ and the whole purifying output is
accessible, write $|\Psi\rangle=\sum_jw_j\otimes|j\rangle$.
Then $J_F=WF^{\mathsf T}W^*$, $J=WW^*$, and the standard CP-order
correspondence gives $J^{\rm acc}_P=S_P(J)$. An inaccessible output makes
$\Omega$ mixed; that equality is no longer automatic.

\begin{proposition}[Further loss cannot improve exact accessible support]
\label{sel46:prop:monotone}
Let a specified readout channel $\mathcal N:A\to A'$ act after
\eqref{sel46:eq:joint}. The set of selectable system suboperations can
only shrink. In particular its new maximum on $P$ is at most
$J^{\rm acc}_P$, and its capacity on every fixed ray is at most $\alpha_v$.
\end{proposition}
\begin{proof}
A terminal effect $F'$ pulls back to the effect $\mathcal N^*(F')$ on
$A$, and selects exactly the same system suboperation. This inclusion of
feasible sets proves both assertions.
\end{proof}

\section[An exact approximation floor at any success]{The same test gives a generalized eigenvalue certificate}
\label{sel46:floor}

\begin{theorem}[Sharp conditional leakage fraction]
\label{sel46:thm:floor}
On the identified support $A_0$, define
\begin{equation}
 \delta_P=\lambda_{\min}\bigl(R_A^{-1/2}D_PR_A^{-1/2}\bigr).
 \label{sel46:eq:delta}
\end{equation}
Then
\begin{equation}
 \boxed{\inf_{F\succeq0,\ F\preceq I,\ J_F\ne0}
 \frac{\Tr(QJ_F)}{\Tr J_F}=\delta_P.}
 \label{sel46:eq:floor}
\end{equation}
The infimum is attained by a scaled rank-one effect. Moreover
$\delta_P=0$ if and only if an exactly $P$-supported nonzero outcome is
accessible. When $P=P_v$, every normalized accepted Choi state satisfies
\begin{equation}
 \left\|\frac{J_F}{\Tr J_F}-P_v\right\|_1\ge2\delta_{P_v}.
 \label{sel46:eq:trace-floor}
\end{equation}
\end{theorem}
\begin{proof}
For $Y=R_A^{1/2}FR_A^{1/2}$, the ratio is
$\Tr(YR_A^{-1/2}D_PR_A^{-1/2})/\Tr Y$, bounded below by the least
eigenvalue. If $u$ is a minimizing unit eigenvector, choose
$F=tR_A^{-1/2}|u\rangle\langle u|R_A^{-1/2}$ with
$0<t\le\|R_A^{-1/2}u\|^{-2}$. This is a nonzero effect attaining the
bound. The zero equivalence follows by invertibility of $R_A$ and
Theorem~\ref{sel46:thm:max}. For a normalized Choi state $\rho$, the
two-outcome measurement $\{P_v,I-P_v\}$ gives
$\|\rho-P_v\|_1\ge2\Tr[(I-P_v)\rho]$.
\end{proof}

The floor concerns normalized \emph{Choi states}. It is not, without
additional normalization information, a diamond distance between
conditioned system channels. It holds at arbitrarily small acceptance:
postselecting more rarely cannot beat the fixed accessible-output floor.
This differs from V45's classical quantile optimization, where the
conditional error can approach zero at vanishing success.

\begin{proposition}[One-sided certificate under reconstruction error]
\label{sel46:prop:error}
Suppose the readout support is independently fixed, the reconstructed
matrices obey $0\preceq\widehat D\preceq\widehat R$, and
$\|D_P-\widehat D\|\le\epsilon_D$,
$\|R_A-\widehat R\|\le\epsilon_R<r$, where
$r=\lambda_{\min}(\widehat R)>0$. Put
$\widehat\delta=\lambda_{\min}(\widehat R^{-1/2}\widehat D
\widehat R^{-1/2})$. Then
\begin{equation}
 \boxed{\delta_P\ge
 \max\left\{0,\widehat\delta-
       \frac{\epsilon_D+\widehat\delta\epsilon_R}{r-\epsilon_R}
       \right\}.}
 \label{sel46:eq:error}
\end{equation}
A strictly positive lower bound excludes exact access to that target ray.
\end{proposition}
\begin{proof}
Since $\widehat D\succeq\widehat\delta\widehat R$,
$D_P\succeq\widehat\delta R_A-
(\epsilon_D+\widehat\delta\epsilon_R)I$. Use
$R_A\succeq(r-\epsilon_R)I$ and conjugate by $R_A^{-1/2}$.
\end{proof}

Exact zero cannot be certified by a small numerical eigenvalue alone.
Changing the unverified accessible support or adding a readout degree of
freedom would change the source question, not improve this error bound.

\section[Partial loss of the original fine record]{An invertible readout map can still exclude every new efficient ray}
\label{sel46:dephasing}

Keep a finite original efficient bank, possibly redundant,
\[
 J=\sum_{j=1}^m|w_j\rangle\langle w_j|,
 \qquad |\Psi\rangle=\sum_jw_j\otimes|j\rangle_A.
\]
Its indices can be original word labels, including their actual probability
weights. The normalized system channel is unchanged if only its readout is
processed. Consider the declared readout family
\begin{equation}
 \mathcal D_\kappa(X)=\kappa X+(1-\kappa)\Delta(X),
 \quad 0\le\kappa\le1,
 \quad \Delta(X)=\sum_j|j\rangle\langle j|X|j\rangle\langle j|.
 \label{sel46:eq:dephasing}
\end{equation}
It can be realized by inaccessible marker vectors
$e_j=\sqrt\kappa e_0+\sqrt{1-\kappa}f_j$, where $e_0,f_1,\ldots,f_m$
are orthonormal, through $|j\rangle\mapsto|j\rangle_A\otimes e_j$.
The inaccessible markers have pairwise overlap $\kappa$.
The retained joint packet is
$\Omega_\kappa=(\id\otimes\mathcal D_\kappa)(|\Psi\rangle\langle\Psi|)$.
Every member has marginal $J$ and selects
\begin{equation}
 J_F^{(\kappa)}
 =\kappa WF^{\mathsf T}W^*
  +(1-\kappa)\sum_jF_{jj}|w_j\rangle\langle w_j|.
 \label{sel46:eq:dephased-selection}
\end{equation}

\begin{theorem}[Any nonzero fine-record leakage fixes the accessible rays]
\label{sel46:thm:record}
For every $0\le\kappa<1$, the accessible nonzero rank-one system rays
are exactly the original rays $\C w_j$ with $w_j\ne0$.
For a fixed unit ray $v$, its largest accessible coefficient mass is
\begin{equation}
 \boxed{\alpha_v^{(\kappa)}=
   \sum_{j:w_j\in\C v}\|w_j\|^2,\qquad 0\le\kappa<1.}
 \label{sel46:eq:class-capacity}
\end{equation}
In particular a new coherent ray allowed by $0\preceq J'\preceq J$
can be inaccessible for every $\kappa<1$, even arbitrarily close to one.
No injectivity or linear independence of the original coefficient bank is
assumed.
\end{theorem}
\begin{proof}
Both terms in \eqref{sel46:eq:dephased-selection} are positive. If their
sum has range $\C v$, every $j$ with $F_{jj}>0$ and $w_j\ne0$ must
satisfy $w_j\in\C v$. Positivity of $F$ gives
$|F_{ij}|^2\le F_{ii}F_{jj}$, so all its nonzero rows and columns lie
among these indices and indices with $w_j=0$. On that coordinate support,
$F\preceq I$; order preservation bounds the selected packet by
$\sum_{j:w_j\in\C v}|w_j\rangle\langle w_j|$. The diagonal
projection onto those indices attains this bound. Selecting a single
original index proves the existence assertion. The proof does not use an
inverse of $W$ or discard an old label.
\end{proof}

\begin{proposition}[Full linear tomography does not restore positive access]
\label{sel46:prop:span}
For $0<\kappa<1$, $\mathcal D_\kappa$ is an invertible linear map on
all matrices. Its accessible effect pullbacks span the entire Hermitian
space and generate its full matrix algebra. Its inverse is not positive.
Thus informational completeness, and even full generated readout algebra,
do not imply the unrestricted effect interval used by V45.
\end{proposition}
\begin{proof}
The map fixes diagonal entries and multiplies every off-diagonal entry by
$\kappa$, so it is invertible. Effects span Hermitian matrices and an
invertible linear map preserves that span. On the two-index effect
$|+\rangle\langle+|$, the inverse has eigenvalues
$(1\pm\kappa^{-1})/2$, one negative. It cannot implement the missing
positive effect by inverse postprocessing. This is the access distinction
behind the theorem, not a failure of statistical reconstruction.
\end{proof}

An arbitrary sequence of operations on the accessible register alone before
a final acceptance flag still pulls back to an effect $F$. Ancillary
processing, adaptation and a rarer outcome do not evade
Theorem~\ref{sel46:thm:record}. Access to the inaccessible marker or
feedback on the system is a different operation class.

\section[Evaluation on the unchanged count-seven channel]{The V45 optimum requires more than nearly coherent fine-record access}
\label{sel46:old}

Use exactly the original V43 count-seven terminal coefficients. If its
terminal word $w$ has stopping coefficient $r_w=c_{w^-}-c_w>0$, let
\[
 w_w=\operatorname{vec}\bigl(\sqrt{r_w/p_7}\,k_w\bigr),
 \qquad p_7=1/20.
\]
They give the unchanged conditional packet $J_7^{\rm cond}$ in
\eqref{sel44:eq:J7}. Consider only the readout family
\eqref{sel46:eq:dephasing} on this \emph{actual fine-word coefficient
bank}. The family changes how much of that record is irreversibly copied;
it does not change the system channel, pulse probabilities, stopping coins,
count, or physical clock. Its $\kappa>0$ members are comparison
implementations, not permissions to reverse a classical record already
written in the original process.

\begin{corollary}[Sharp discontinuity of efficient symmetric access]
\label{sel46:cor:old}
In this family, the maximal squared $I/Y$ cross of an accessible efficient
$H_Y$-covariant suboperation is
\begin{equation}
 \boxed{
 \mathfrak h_{\rm eff}(\kappa)=
 \begin{cases}
 0,&0\le\kappa<1,\\[1mm]
 (33+8\sqrt2)/400,&\kappa=1.
 \end{cases}}
 \label{sel46:eq:jump}
\end{equation}
At full access the maximizing ray and its success are exactly those of
Proposition~\ref{sel45:prop:oldbin}; no new maximizing channel is fitted.
Its unconditional cross budget retains the factor $p_7^2=1/400$.
\end{corollary}
\begin{proof}
The inherited V40 all-depth exclusion says that no nonempty original pulse
word is even parity-covariant under $Y$, hence none is $H_Y$-covariant.
Theorem~\ref{sel46:thm:record} excludes every nonzero efficient
$H_Y$-covariant selection for $\kappa<1$. At $\kappa=1$ the full pure
dilation realizes the entire CP-order interval, including in a redundant
fine-word dilation after restricting its actual Schmidt support. The
unchanged neutral short is
$\left(\begin{smallmatrix}4/5&i/10\\-i/10&2/5\end{smallmatrix}\right)$;
V45's exact optimization gives the second value. Multiplying an original
outcome by $p_7$ multiplies its squared cross by $p_7^2$.
\end{proof}

The quantity in \eqref{sel46:eq:jump} restricts to \emph{exactly efficient}
selected maps. There is no discontinuity of the coarse channel and no claim
of a discontinuity of an approximate optimal fidelity. The coarse channel
remains exactly circle covariant for all these implementations. The theorem
strengthens V44's classical fine-word obstruction to every partially leaked
record in this specified family. It is not a universal claim that all
small environment noises destroy all pure selections; the general test is
Theorem~\ref{sel46:thm:max}.

\section[Two typed shadows with an inaccessible marker]{Exact purity and surviving determinant covariance are different questions}
\label{sel46:two}

Let $b,d$ be two independently identified orthonormal coefficient vectors.
For a declared two-sector packet, take the joint vector
\begin{equation}
 |\Psi\rangle
 =\sqrt\mu\,b\otimes|0\rangle_A\otimes e_b
  +\sqrt\nu\,d\otimes|1\rangle_A\otimes e_d,
 \quad \mu,\nu>0,\quad \langle e_d,e_b\rangle=\zeta,
 \quad r=|\zeta|\le1.
 \label{sel46:eq:markers}
\end{equation}
The marker vectors are unit vectors and are inaccessible. The source is
$\Omega=\Tr_{\rm marker}|\Psi\rangle\langle\Psi|$ and its marginal is
always $J=\mu P_b+\nu P_d$. These are coefficient masses, not assumed
input-independent probabilities. A fully normalized special case is a
state-preparation channel with scalar input and $\mu+\nu=1$. A typed
subpacket of a larger instrument retains its original positive remainder.

For $F=\left(\begin{smallmatrix}a&z\\\bar z&b_0\end{smallmatrix}\right)$,
\begin{equation}
 J_F=
 \begin{pmatrix}
 \mu a&\sqrt{\mu\nu}\,\zeta\bar z\\
 \sqrt{\mu\nu}\,\bar\zeta z&\nu b_0
 \end{pmatrix},
 \quad 0\le a,b_0\le1,
 \quad |z|^2\le\min\{ab_0,(1-a)(1-b_0)\}.
 \label{sel46:eq:marker-packet}
\end{equation}

\begin{theorem}[Sharp cross and purity limits under incomplete erasure]
\label{sel46:thm:two}
The following statements hold for the \emph{whole} accessible effect interval.
\begin{equation}
 \boxed{\max_F\|P_bJ_FP_d\|_{\HS}^2=\frac{\mu\nu r^2}{4}.}
 \label{sel46:eq:marker-cross}
\end{equation}
A nonzero efficient selected ray having both shadows exists exactly when
$r=1$. For any selected packet with both shadow masses positive, define
$t=\Tr J_F$ and $p=\mu a/t$. Then
\begin{equation}
 \boxed{1-\Tr[(J_F/t)^2]
 \ge2(1-r^2)p(1-p).}
 \label{sel46:eq:marker-purity}
\end{equation}
This bound is attained at every $0<p<1$ by a sufficiently scaled rank-one
readout effect. For any prescribed balanced target
$v_\theta=(b+e^{i\theta}d)/\sqrt2$, when the readout phases are unrestricted,
\begin{equation}
 \boxed{\inf_{F:J_F\ne0}
 \|J_F/\Tr J_F-P_{v_\theta}\|_1=1-r.}
 \label{sel46:eq:marker-distance}
\end{equation}
If $b,d$ transform by $1,\chi^{-1}$ with
$\chi(U,V)=\det U\det V$, every selected packet with $\zeta z\ne0$
has covariance group exactly $\ker\chi$ inside the comparison group.
For $0<r<1$ this gives exact determinant covariance but \emph{not} an
efficient common source ray.
\end{theorem}
\begin{proof}
The effect constraints imply $|z|\le1/2$: for $s=a+b_0$,
$ab_0\le s^2/4$ and $(1-a)(1-b_0)\le(2-s)^2/4$; their minimum is
at most $1/4$. A balanced rank-one projector attains it, proving
\eqref{sel46:eq:marker-cross}. The determinant is
$\mu\nu(ab_0-r^2|z|^2)$. If $ab_0>0$, it is zero precisely when
$r=1$ and $|z|^2=ab_0$. For a normalized two-dimensional positive
matrix, $1-\Tr\rho^2=2\det\rho$. The inequality
$|z|^2\le ab_0$ proves \eqref{sel46:eq:marker-purity}.
For any required $p$, choose $a,b_0>0$ in the required ratio and
$|z|^2=ab_0$, then scale the effect to have largest eigenvalue at most one.

For the balanced target, the normalized unwanted-sector matrix in
\eqref{sel46:eq:delta} is unitarily equivalent to
$\frac12\left(\begin{smallmatrix}1&-r\\-r&1\end{smallmatrix}\right)$;
its least eigenvalue is $(1-r)/2$. Theorem~\ref{sel46:thm:floor}
gives the lower bound $1-r$. Choose a rank-one effect with
$\mu a=\nu b_0$, and choose its phase so the output cross has the
target phase. The normalized selected packet is
$\frac12\left(\begin{smallmatrix}1&r e^{-i\theta}\\r e^{i\theta}&1\end{smallmatrix}\right)$,
which attains that distance. Finally the two diagonal blocks are invariant,
while a nonzero off-diagonal block is multiplied by the nontrivial relative
character. It is fixed precisely when $\chi=1$.
\end{proof}

For example, $\mu=\nu=1/2$, $r=3/5$ and a balanced effect give
$\Tr J_F=1/2$, $h=9/400$, normalized purity $17/25$, and best balanced
pure-target distance $2/5$. For $r=1$ the same marginal admits a pure
coherent selection; for $r=0$ none of its accessible effects produces a
cross. Thus the marginal $J$ alone cannot determine any of these access
answers. The familiar marker-overlap mechanism is not new
\cite{sel46-Neves}; the source claims being distinguished here are its
exact CP efficiency, symmetry and inaccessible-record implications.

\section[Loss that makes every accepted packet full rank]{A readout-noise floor without altering the original system channel}
\label{sel46:faithful}

\begin{proposition}[A faithful replacement component in the readout]
\label{sel46:prop:replacement}
Let $\Omega_0$ have system marginal $J\ne0$, let $\tau_A\succ0$ be a
density on the identified accessible support, and set
\begin{equation}
 \Omega_\eta=(1-\eta)\Omega_0+\eta J\otimes\tau_A,
 \qquad 0<\eta\le1.
 \label{sel46:eq:replacement}
\end{equation}
This is the result of a specified local readout-replacement channel and
still has system marginal $J$. Every nonzero accepted output has
$\Ran J_F=\Ran J$. No proper supported selection exists. More
quantitatively, for any $P$ with $\Tr(QJ)>0$,
\begin{equation}
 \boxed{\delta_P\ge
 \frac{\eta\Tr(QJ)\lambda_{\min}(\tau_A)}
      {\lambda_{\max}(R_{A,\eta})}>0.}
 \label{sel46:eq:replacement-floor}
\end{equation}
\end{proposition}
\begin{proof}
The selected operator obeys
$J_F\succeq\eta\Tr(F\tau_A)J$. A nonzero effect has strictly
positive $\tau_A$ expectation, and $J_F\preceq J$, proving equality
of supports. Likewise
$D_{P,\eta}\succeq\eta\Tr(QJ)\tau_A$. Compare this with
$R_{A,\eta}\preceq\lambda_{\max}(R_{A,\eta})I$ and use
\eqref{sel46:eq:delta}.
\end{proof}

In particular, a channel can have a full-rank system Choi matrix admitting
many mathematical pure suboperations while an arbitrarily weak faithful
replacement component on its \emph{readout} makes all of them inaccessible.
This is not the noise model of the original YM law or the historical Store.
It is a separating comparison at fixed $J$, and prevents physical access
from being inferred from that marginal alone.

\section[Status and the next historical readout test]{A smaller joint source row, not an assumption of perfect erasure}
\label{sel46:status}

\paragraph{Proved.}
An actual mixed joint system--readout packet determines the maximal
accessible supported suboperation by one positive unwanted-sector matrix.
Its kernel gives exact ray availability; a generalized eigenvalue gives the
sharp conditional leakage fraction and a certified one-sided obstruction
under errors. Any nonzero irreversible fine-record component in the
specified dephasing family restricts efficient outcomes to original rays,
although its effect pullbacks remain informationally complete. Applied to
the original V43 fine-word bank, it excludes every new efficient symmetric
ray for all $\kappa<1$ and retains the unchanged V45 optimum at full
access. A two-shadow marker calculation separates surviving exact group
covariance from exact one-ray occurrence and gives sharp cross, purity and
approximation bounds. Faithful readout replacement preserves the entire
Choi support of every accepted output.

\paragraph{Inherited and conditional.}
V45's CP-order optimum and analytic classical theorem remain unchanged.
V40's original pulse-word exclusion and V43's exact stopping bank are inputs
to the evaluated access-family corollary, not new historical matter data.
All formulas keep the marginal channel and original source weights fixed.
Changing the joint readout extension is explicitly a different access
implementation. The general maximum assumes that all effects on the
\emph{actually accessible} output are permitted; smaller banks need their
own feasibility restriction. A classical record cannot be reinterpreted as
an unmeasured coherent register. Finite symbolic checks supplement the proofs
and do not establish physical availability or a new grading.

\paragraph{Still unresolved for the historical packet.}
The next concrete row is the same-cut joint Choi packet of the typed source
and its actual accessible readout, with inaccessible records traced out.
For a proposed V110 ray, only $R_A$ and its unwanted-sector matrix $D_{P_v}$
are needed for its exact capacity test; full recovery of the environment is
not required. These contractions may be inferred from admitted joint
contexts when their tester span suffices, but system-only tomography does
not supply them. A nonzero kernel must still correspond to an admitted
readout effect. No literal historical V110 packet, marker overlap, or
readout channel is populated here. One efficient common source ray remains
a stronger statement than determinant covariance of a mixed outcome; neither
is replaced by a formal choice of a Kraus decomposition.

\begingroup
\renewcommand{\refname}{Additional primary-source context for V46}
\begin{thebibliography}{9}
\small\raggedright
\bibitem{sel46-Raginsky}
M.~Raginsky, \emph{Radon--Nikodym derivatives of quantum operations},
J. Math. Phys. \textbf{44} (2003), 5003--5020;
\url{https://arxiv.org/abs/math-ph/0303056}.
\bibitem{sel46-Gregoratti}
M.~Gregoratti and R.~F. Werner, \emph{Quantum lost and found},
J. Mod. Opt. \textbf{50} (2003), 915--933;
\url{https://arxiv.org/abs/quant-ph/0209025}.
\bibitem{sel46-Feng}
Y.~Feng, R.~Duan and M.~Ying,
\emph{Unambiguous discrimination between quantum mixed states} (2004);
\url{https://arxiv.org/abs/quant-ph/0403147}.
\bibitem{sel46-Neves}
L.~Neves, G.~Lima, J.~Aguirre, F.~A. Torres-Ruiz, C.~Saavedra and A.~Delgado,
\emph{Control of quantum interference in the quantum eraser},
New J. Phys. \textbf{11} (2009), 073035;
\url{https://arxiv.org/abs/0904.4242}.
\end{thebibliography}
\endgroup
\addtocontents{toc}{\protect\endgroup}
% END OF SOURCE-SELECTION V46 DEVELOPMENT BLOCK.
% Preserve V1--V46 and append subsequent work before the final terminator.


\clearpage
\hypersetup{pdftitle={Historical Standard-Model Source Selection: V1--V47}}
\begin{center}
{\Large\bfseries Source-selection continuation V47}\\[.4em]
{\large Symmetry-accounted readout, finite determinant references,\\
and the cost of extracting a common coherent ray}
\end{center}
\addtocontents{toc}{\protect\begingroup}
\addtocontents{toc}{\protect\renewcommand{\protect\numberline}{\protect\makebox[2.1em][l]}}

\section[The readout may contain the missing symmetry resource]{Audit the selection operation before assigning its symmetry to the source}
\label{sel47:context}

V46 determines exact ray accessibility from an actual joint system--readout
packet, assuming the full effect interval on its identified accessible output.
Its kernel and leakage theorems remain unchanged. The next question is not
another unconstrained Kraus decomposition. It is whether the proposed
\emph{readout operation itself} requires a phase resource under the comparison
group. A symmetry-invariant joint source and a symmetry-breaking readout can
produce a selected packet with a smaller group without showing that the
unmeasured source intrinsically selected that group.

We make this distinction operational on the \emph{same two-shadow packet}
used in Theorem~\ref{sel46:thm:two}. No historical reference state is supplied
by this calculation. The new input, when used, is explicitly a finite quantum
reference carrying specified determinant characters. All joint effects must
commute with the total comparison action. We derive the exact cross budget,
the success of an efficient common ray, a distinction between these two
optimizations, and an obstruction to returning that finite reference exactly
and uncorrelated on a successful branch.

The harmonic-mode framework and limitations on broadcasting finite quantum
references are established mathematics \cite{sel16-MS,sel16-GS,sel47-LM,sel47-MS}.
No general priority claim is made for asymmetry as a resource, phase-reference
bounds, steering or the no-broadcasting principle. The contribution relative
to this lineage is the explicit reference-accounted selector, its exact
one-ray cost on the existing determinant--Clifford comparison, the gapped-mode
activation example, and the heralded product-return test below. The
all-effects entrance of V46 and the already coherent historical-word branch
of V27 are not retroactively given a symmetry restriction. In particular, an
\emph{already occurring} efficient word with both shadows still has the
coherence proved there without a newly added reference.

\section[Covariance of the complete source and readout]{A symmetry of the selected packet need not be a symmetry selected by the source}
\label{sel47:covariance}

Let $G$ be a compact group, with identified unitary actions $R_g$ on the
system coefficient space $\mathcal W$ and $U_g$ on the accessible output $A$.
Assume the \emph{actual joint} Choi packet is invariant:
\begin{equation}
 (R_g\otimes U_g)\Omega(R_g\otimes U_g)^*=\Omega.
 \label{sel47:eq:joint-invariant}
\end{equation}
For an accessible effect $F$, retain V46's convention
$J_F=\Tr_A[(I\otimes F^{1/2})\Omega(I\otimes F^{1/2})]$.

\begin{proposition}[The readout-to-source covariance map]
\label{sel47:prop:steering}
For every $g\in G$,
\begin{equation}
 \boxed{R_gJ_FR_g^*=J_{U_gFU_g^*}.}
 \label{sel47:eq:intertwining}
\end{equation}
Thus an invariant effect selects a $G$-invariant packet. More generally, the
stabilizer of the effect is contained in the stabilizer of its selected
packet. If a reference $R$ in state $\tau$ is appended independently and
$F$ is invariant under $U_g\otimes V_g$, then its effective effect
\begin{equation}
 E_\tau=\Tr_R[(I_A\otimes\tau^{1/2})F(I_A\otimes\tau^{1/2})]
 \label{sel47:eq:effective}
\end{equation}
satisfies $0\preceq E_\tau\preceq I_A$. Every symmetry fixing $\tau$ fixes
$J_{E_\tau}$. All statements concern fixed outcome labels.
\end{proposition}
\begin{proof}
Joint invariance gives
$(R_g\otimes I)\Omega(R_g^*\otimes I)
=(I\otimes U_g^*)\Omega(I\otimes U_g)$.
Partial-trace cyclicity proves \eqref{sel47:eq:intertwining}.
Equation~\eqref{sel47:eq:effective} is the pullback of a positive effect
through preparation of $\tau$, hence is an effect. If $V_g\tau V_g^*=\tau$,
conjugate its defining expression by $U_g$, use invariance of $F$, and
change the traced reference basis by $V_g$. This proves invariance of
$E_\tau$ and then of $J_{E_\tau}$.
\end{proof}

The proposition is an inclusion, not an equality: a source can be insensitive
to some noninvariant readout directions. It applies only after the joint
symmetry has been verified, not from invariance of the marginal $J$ alone.
Finite classical outcomes do not evade it by a continuous permutation action
of a connected group, since every such action is constant.

For the determinant comparison let
\begin{equation}
 G_0=\mathrm U(3)\times\mathrm U(2),\quad
 \chi(U,V)=\det U\det V,\quad
 H=\ker\chi.
 \label{sel47:eq:group}
\end{equation}
Take the already identified orthonormal coefficient vectors $b,d$ with
characters $1,\chi^{-1}$, and give the accessible flags $|0\rangle,|1\rangle$
characters $1,\chi$. In V46's marker packet this makes the full joint source
invariant, if the inaccessible marker is neutral. After tracing the marker,
\begin{equation}
 \begin{split}
 \Omega={}&\mu|b0\rangle\langle b0|+
             \nu|d1\rangle\langle d1|\\
 &+\sqrt{\mu\nu}\bigl(\zeta|b0\rangle\langle d1|
                     +\bar\zeta|d1\rangle\langle b0|\bigr),
 \qquad r=|\zeta|\le1.
 \end{split}
 \label{sel47:eq:source}
\end{equation}
Here $\mu,\nu>0$ and the original weights are not changed. Its system marginal
is $\mu P_b+\nu P_d$. With scalar input and $\mu+\nu=1$ this is a normalized
state-preparation source; for a typed Choi subpacket, all other positive
source components and the input normalization remain separately retained.
The new \emph{character assignment on the readout} is a declared comparison
hypothesis, not a deduction from the marker overlap or marginal channel.

Without an asymmetric reference, every allowed invariant effect on this
readout is diagonal, so it produces zero $b/d$ cross even when $r=1$.
The unrestricted balanced eraser from V46 is not invariant under this
comparison action. This does not prohibit that eraser if it is independently
admitted; it identifies the resource hidden in such admission.

\section[Exact determinant coherence from a finite reference]{A sharp first-character budget with arbitrary charge multiplicities}
\label{sel47:budget}

Let the independently supplied finite reference have representation
\begin{equation}
 V_g=\bigoplus_{n\in\mathbb Z}\chi(g)^n I_{R_n},\qquad
 \tau\succeq0,\quad\Tr\tau=1,
 \quad\tau_{nm}=P_n\tau P_m.
 \label{sel47:eq:reference}
\end{equation}
Only finitely many $R_n$ are nonzero. No physical energy spacing, waiting time,
or gauge-charge reference is inferred from an ordinary clock. Allow the
\emph{whole invariant effect interval} on $A\otimes R$; a smaller physical
bank can only lower the following maxima. Define
\begin{equation}
 \boxed{v_1(\tau)=\sum_n\|\tau_{n,n+1}\|_1.}
 \label{sel47:eq:v1}
\end{equation}

\begin{theorem}[Sharp reference-limited cross]
\label{sel47:thm:cross}
For the unchanged source \eqref{sel47:eq:source},
\begin{equation}
 \boxed{
 \max_{\substack{0\preceq F\preceq I\\{}[F,U_g\otimes V_g]=0}}
 \|P_bJ_{E_\tau}P_d\|_{\HS}^2
 =\frac{\mu\nu r^2}{4}\,v_1(\tau)^2.}
 \label{sel47:eq:cross-max}
\end{equation}
The maximum is attained. A nonzero determinant cross is selectable exactly
when $r>0$ and $v_1(\tau)>0$. On the stated two-character packet, every such
selected cross has covariance group exactly $H$. When $v_1(\tau)>0$, the
reference density itself has stabilizer exactly $H$ inside $G_0$.
\end{theorem}
\begin{proof}
The total-character-$n$ subspace is
$(|0\rangle\otimes R_n)\oplus(|1\rangle\otimes R_{n-1})$.
An invariant effect is block diagonal in these subspaces. Let $Z_n$ be its
off-diagonal block from $|1\rangle\otimes R_{n-1}$ to
$|0\rangle\otimes R_n$. Since $2F-I$ is a contraction, $\|Z_n\|\le1/2$.
Taking the partial trace in \eqref{sel47:eq:effective} gives
\[
 (E_\tau)_{01}=\sum_n\Tr(\tau_{n-1,n}Z_n),\qquad
 |(E_\tau)_{01}|\le\tfrac12v_1(\tau).
\]
Each bound can be attained with a common phase: choose twice $Z_n$ to be an
appropriate polar partial isometry of $\tau_{n-1,n}$, and set the two
corresponding diagonal blocks to half their identities. The block
$\frac12\left(\begin{smallmatrix}I&V\\V^*&I\end{smallmatrix}\right)$ and
its complement are positive for every contraction $V$. Different total
characters occupy orthogonal blocks, so all choices are simultaneous.
Unused boundary blocks may be set to zero. V46's selection formula now gives
$J_{bd}=\sqrt{\mu\nu}\,\zeta\overline{(E_\tau)_{01}}$, proving the maximum.
The two diagonal blocks are invariant and the nonzero cross transforms by
$\chi$, so its stabilizer is exactly $H$.
Finally a nonzero $\tau_{n,n+1}$ transforms by $\chi^{-1}$; fixing it forces
$\chi=1$. Conversely $H$ acts trivially on every $R_n$, proving the last claim.
\end{proof}

Thus in this invariant-source implementation the selected determinant group
is \emph{not inferred without a determinant-sensitive input resource}.
The theorem states exactly which mode of that resource is used. It does not
show that the historical source lacks such a resource, nor that an initially
coherent historical determinant--Clifford word must obtain it from a new
apparatus.

\begin{proposition}[Finite-reference bound and a direct error estimate]
\label{sel47:prop:path}
Suppose the allowed reference charges have no run of more than $L$ consecutive
integers. Then
\begin{equation}
 \boxed{v_1(\tau)\le\cos\frac{\pi}{L+1}<1.}
 \label{sel47:eq:path}
\end{equation}
The maximum is attained on a longest run by a pure reference with normalized
amplitudes
$x_j=\sqrt{2/(L+1)}\sin((j+1)\pi/(L+1))$, $0\le j<L$.
If $\widehat\tau$ is a density on the same charge space and the adjacent
block errors have sum at most $\epsilon$, then
\begin{equation}
 |v_1(\tau)-v_1(\widehat\tau)|\le\epsilon,
 \qquad
 |h_{\max}(\tau)-h_{\max}(\widehat\tau)|
 \le\frac{\mu\nu r^2}{2}\epsilon.
 \label{sel47:eq:error}
\end{equation}
\end{proposition}
\begin{proof}
Write $\tau=XX^*$ and $X_n=P_nX$. The Schatten Cauchy--Schwarz inequality gives
$\|\tau_{n,n+1}\|_1\le\|X_n\|_{\HS}\|X_{n+1}\|_{\HS}
=\sqrt{p_np_{n+1}}$, where $p_n=\Tr\tau_{nn}$ and $\sum_np_n=1$.
With $x_n=\sqrt{p_n}$, the resulting sum is the Rayleigh form of the direct
sum of path matrices having $1/2$ on adjacent off-diagonals. A run of length
$L$ has largest eigenvalue $\cos(\pi/(L+1))$: substituting the displayed sine
vector proves its eigenvalue and normalization, and the remaining sine
vectors give the full spectrum. A pure reference on that vector saturates
both inequalities. The trace-norm triangle inequality proves the error in
$v_1$; use $0\le v_1(\tau),v_1(\widehat\tau)\le1$ to compare their squares.
\end{proof}

This is a bound on an unnormalized selected cross, not on its conditional
purity or success. A state whose stabilizer is $H$ can still have $v_1=0$;
the exact counterexample in Section~\ref{sel47:examples} makes the distinction
operational.

\section[Efficient-ray success is a different optimization]{Resolve total character before asking for a pure common source}
\label{sel47:pure}

Retain $r=1$ in \eqref{sel47:eq:source}; the constant marker phase can be
absorbed into the target phase. Let the reference be a \emph{specified pure}
state $\psi$, with $p_n=\|P_n\psi\|^2$. Consider the prescribed normalized ray
\[
 v_q=\sqrt q\,b+e^{i\theta}\sqrt{1-q}\,d,\qquad 0<q<1.
\]

\begin{theorem}[Exact common-ray capacity with a finite pure reference]
\label{sel47:thm:pure}
Under invariant readout effects, the largest coefficient mass $\alpha_q$ with
$J_{E_\tau}=\alpha_qP_{v_q}$ is
\begin{equation}
 \boxed{
 \alpha_q(\psi)=
 \sum_n\frac{\mu\nu p_np_{n-1}}
 {q\nu p_{n-1}+(1-q)\mu p_n}.}
 \label{sel47:eq:pure-capacity}
\end{equation}
Terms with a zero numerator are zero. The maximum is attained by one effect
which is a sum of rank-one projectors in different total-character blocks.
No fine total-character label needs to be coherently superposed: every
nonzero selected block gives the \emph{same} system ray.
For $r<1$, the capacity for every ray having both shadows is zero, whatever
finite reference is supplied under this readout-only protocol.
\end{theorem}
\begin{proof}
In total character $n$, the source's two occupied readout directions have
coefficient map
$W_n=\diag(\sqrt\mu\,\|P_n\psi\|,
\sqrt\nu\,\|P_{n-1}\psi\|)$, after retaining their known phases in the
readout basis. Selecting its $2\times2$ effect produces a positive packet
below
$J_n=\diag(\mu p_n,\nu p_{n-1})$.
Different total-character blocks add positively. If their sum has the desired
rank-one support, each nonzero summand must have that support. Its largest
mass is
\[
 \alpha_{q,n}=\bigl(\langle v_q,J_n^{-1}v_q\rangle\bigr)^{-1}
 =\frac{\mu\nu p_np_{n-1}}
 {q\nu p_{n-1}+(1-q)\mu p_n}.
\]
For positive diagonal entries, the effect transpose
$\alpha_{q,n}|W_n^{-1}v_q\rangle\langle W_n^{-1}v_q|$
is a rank-one projector and attains this value. If either entry is zero,
a ray with both shadows is impossible on that block. Summing these effects
is allowed because the blocks are orthogonal; summing their output masses
proves the formula. Arbitrary charge multiplicities cause no change, since a
pure reference occupies only the vector $P_n\psi$ within each sector.
For $r<1$, every reference-assisted effect still pulls back to an ordinary
effect on $A$ by \eqref{sel47:eq:effective}. The exact determinant calculation
in Theorem~\ref{sel46:thm:two} excludes a two-shadow pure output even in that
larger unrestricted effect interval.
\end{proof}

With scalar input and $\mu+\nu=1$, $\alpha_q$ is the physical success
probability. For a typed channel coefficient, the selected input effect is
$\alpha_q K_{v_q}^*K_{v_q}$, with $K_{v_q}=\operatorname{unvec}v_q$;
only its mean on the maximally mixed input is $\alpha_q/\dim H_{\rm in}$.
State-independent success is not implicit.

For a uniform reference on $L$ consecutive charges,
\begin{equation}
 \boxed{
 \alpha_q=\frac{L-1}{L}\,
 \frac{\mu\nu}{q\nu+(1-q)\mu}.}
 \label{sel47:eq:uniform}
\end{equation}
The second factor is the unrestricted V45 CP-order mass on this ray.
Thus a finite reference can enable an \emph{exact} efficient branch with
positive success, although it cannot attain the unrestricted maximal cross.
The reference is consumed by this readout; exact return is examined below.

\begin{proposition}[Sharp two- and three-level balanced-ray optima]
\label{sel47:prop:three}
For $\mu=\nu=q=1/2$, optimize also over every reference density on the given
consecutive charges. On two levels the maximum is $\alpha_{1/2}=1/4$.
On three levels,
\begin{equation}
 \boxed{
 \max_\psi\alpha_{1/2}=6-4\sqrt2,
 \qquad
 (p_0,p_1,p_2)=
 (1-1/\sqrt2,\sqrt2-1,1-1/\sqrt2)
 }
 \label{sel47:eq:three-pure}
\end{equation}
is a maximizing population vector. Its selected same-ray cross is
$17-12\sqrt2$. In contrast, the maximal cross over all references and
invariant effects is $1/32$, attained by reference populations
$(1/4,1/2,1/4)$ and a \emph{mixed} selected packet.
\end{proposition}
\begin{proof}
First, a mixed reference is a convex sum of pure references. If one effect
selects a given pure ray from that sum, positivity forces every nonzero
component output onto the same ray. Its total mass cannot exceed the best
pure-reference mass. Thus optimizing over pure references loses nothing.
Equation~\eqref{sel47:eq:pure-capacity} reduces to
$\sum_np_np_{n-1}/(p_n+p_{n-1})$.
On two levels this is $p_0p_1\le1/4$.
The function $f(x,y)=xy/(x+y)$ is jointly concave on $x,y>0$, since
\[
 \nabla^2f=-\frac2{(x+y)^3}
 \begin{pmatrix}y^2&-xy\\-xy&x^2\end{pmatrix}\preceq0,
\]
and extends continuously at zero. On three levels, average a population
vector with its reversal; concavity does not decrease the objective.
It therefore suffices to set $p_0=p_2=a$, $p_1=1-2a$,
$0\le a\le1/2$. The objective is $2a(1-2a)/(1-a)$.
Its only interior maximum is $a=1-1/\sqrt2$, giving
$6-4\sqrt2$; its endpoint values are zero. For a balanced selected pure ray,
the squared cross is $\alpha_{1/2}^2/4=17-12\sqrt2$.
For the cross optimum apply \eqref{sel47:eq:path} with $L=3$.
The attaining reference has amplitudes $(1/2,1/\sqrt2,1/2)$.
Balanced projectors on its two occupied two-dimensional total-character
blocks, with zero boundary effects, give
\begin{equation}
 J_{\rm cross}=
 \begin{pmatrix}3/16&\sqrt2/8\\\sqrt2/8&3/16\end{pmatrix},
 \quad \Tr J_{\rm cross}=3/8,\quad
 \det J_{\rm cross}=1/256>0.
 \label{sel47:eq:three-mixed}
\end{equation}
This attains cross $1/32$ but has rank two. Its normalized purity is $17/18$.
\end{proof}

The two objectives must not be substituted for one another. In particular,
a small mixedness at the maximum cross does not certify V110's efficient
common ray.

\section[Evaluated reference and marker comparisons]{The old system packet is unchanged while the reference condition is exposed}
\label{sel47:examples}

\begin{proposition}[A gapped reference has the right stabilizer but the wrong mode]
\label{sel47:prop:gapped}
Let $\psi=(|0\rangle+|2\rangle+|5\rangle)/\sqrt3$ with the determinant
charge representation \eqref{sel47:eq:reference}. Its density has stabilizer
exactly $H$, but $v_1(|\psi\rangle\langle\psi|)=0$.
One copy therefore cannot select a determinant cross from
\eqref{sel47:eq:source} under invariant effects. Two independently prepared
copies, with the product state and an invariant \emph{joint} readout, have
\begin{equation}
 \boxed{v_1(\psi^{\otimes2})=\frac{\sqrt2}{9}.}
 \label{sel47:eq:gapped}
\end{equation}
For $\mu=\nu=1/2$ the maximal cross is $r^2/648$.
At $r=1$, the maximum success on the prescribed balanced efficient ray is
$2/27$.
\end{proposition}
\begin{proof}
Fixing the one-copy density forces both $\chi^2=1$ and $\chi^5=1$, hence
$\chi=1$, but there are no adjacent occupied charges. The two-copy total
charge probabilities, retaining their degeneracies, are
\[
 \begin{array}{c|rrrrrr}
 n&0&2&4&5&7&10\\\hline
 p_n&1/9&2/9&1/9&2/9&2/9&1/9.
 \end{array}
\]
The reference remains pure; its only adjacent pair is $(4,5)$.
Thus $v_1=\sqrt{p_4p_5}=\sqrt2/9$.
Theorems~\ref{sel47:thm:cross} and~\ref{sel47:thm:pure} give respectively
$r^2/648$ and $(p_4p_5)/(p_4+p_5)=2/27$.
The positive two-copy result uses an actual product reference and a joint
invariant effect. Repeating a one-copy experiment without that joint access
is not the same operation.
\end{proof}

This example also shows that the stabilizer of a resource does not determine
which modes it can transfer in one use. Conversely, the two copies do not
create asymmetry from invariant states: their original nonzero modes add in
the tensor product to produce the required adjacent mode.

\begin{proposition}[The unchanged V46 marker point with a two-level reference]
\label{sel47:prop:old}
Retain V46's packet $\mu=\nu=1/2$, $r=3/5$, and supply the declared pure
reference $(|0\rangle+|1\rangle)/\sqrt2$. With the invariant projector onto
$(|0,1\rangle+|1,0\rangle)/\sqrt2$ the selected packet is
\begin{equation}
 \boxed{
 J_{\rm ref}=\begin{pmatrix}1/8&3/40\\3/40&1/8\end{pmatrix},
 \quad\Tr J_{\rm ref}=1/4,
 \quad h=9/1600.}
 \label{sel47:eq:old}
\end{equation}
Its normalized purity remains $17/25$ and its distance from the balanced
pure target remains $2/5$. It has the exact determinant group but is not an
efficient ray. With the \emph{same reference and effect}, setting $r=1$
gives the efficient packet $\frac14P_{(b+d)/\sqrt2}$; setting $r=0$ gives
zero cross. These are the same source/marker comparisons as in V46, with the
reference restriction now explicit.
\end{proposition}
\begin{proof}
The effective effect is $E_\tau=\frac12P_+=
\frac14\left(\begin{smallmatrix}1&1\\1&1\end{smallmatrix}\right)$.
Insert it into V46's two-sector selection formula. Normalization divides
\eqref{sel47:eq:old} by $1/4$ and recovers exactly the old normalized marker
matrix $\frac12\left(\begin{smallmatrix}1&3/5\\3/5&1\end{smallmatrix}\right)$.
Its two eigenvalues, purity and distance give the displayed values.
The reference has $v_1=1/2$, so Theorem~\ref{sel47:thm:cross} also proves
optimality of this \emph{raw cross} under the stated reference restriction.
\end{proof}

A finite reference does not erase information already lost into the
inaccessible marker. Nor does a reference-free theorem about a direct
historical efficient word require a reference to be supplied afterward.
These are different source classes. The comparison here uses no original
pulse-stopping policy, alters no V43 duration probability, and attributes no
gauge charge to the companion's physical renewal clock.

\section[No exact product return of a finite reference]{Heralding does not make the determinant reference catalytic}
\label{sel47:catalysis}

For a finite integer-charge representation, define the Fourier modes
\begin{equation}
 X^{(k)}=\frac1{2\pi}\int_0^{2\pi}
 e^{-ik\phi}e^{i\phi N}Xe^{-i\phi N}\,d\phi.
 \label{sel47:eq:modes}
\end{equation}
They satisfy $\|X^{(k)}\|_1\le\|X\|_1$.
Every covariant linear map preserves each mode, and a finite reference
density has a largest nonnegative occupied mode
$m=\max\{k\ge0:\tau^{(k)}\ne0\}$.
These are the established mode-selection rules \cite{sel16-MS}; the
following restricted product-return consequence has an elementary proof.

\begin{theorem}[No nonzero heralded product-return branch]
\label{sel47:thm:return}
Let the non-reference input packet be invariant under the full comparison
circle. Start with its product with a finite reference density $\tau$.
No covariant trace-nonincreasing operation with a fixed success label can
produce
\begin{equation}
 p\,\rho\otimes\tau,\qquad p>0,
 \quad \rho\text{ supported on }\Span\{b,d\},
 \quad P_b\rho P_d\ne0.
 \label{sel47:eq:forbidden-return}
\end{equation}
The statement includes mixed $\rho$ and arbitrarily small $p$.
For any normalized actual output $\omega$ of such a successful branch,
\begin{equation}
 \boxed{
 \|\omega-\rho\otimes\tau\|_1
 \ge \|\rho^{(1)}\|_1\,\|\tau^{(m)}\|_1>0.}
 \label{sel47:eq:return-error}
\end{equation}
In particular, for a balanced pure target and a uniform $L$-level pure
reference, the bound is $1/(2L)$.
\end{theorem}
\begin{proof}
The initial invariant packet times $\tau$ has no mode above $m$.
Covariance preserves this fact for a selected, unnormalized output and hence
also for its nonzero scalar normalization. But $\rho$ has modes only
$-1,0,1$, with $\rho^{(1)}\ne0$. The proposed product output has mode $m+1$
exactly $\rho^{(1)}\otimes\tau^{(m)}\ne0$, a contradiction.
Applying the mode-$(m+1)$ contraction to the difference gives
\eqref{sel47:eq:return-error}, since the trace norm multiplies on tensor
products. For the uniform reference, the top mode is
$|L-1\rangle\langle0|/L$; for the balanced target the cross has trace norm
$1/2$. The formula also covers $L=1$, with $m=0$ and $\|\tau^{(0)}\|_1=1$.
\end{proof}

The result concerns a reference returned \emph{exactly and uncorrelated on
the individual success branch}. It does not analyze all protocols permitting
output--reference correlations, an unchanged reference only after averaging
over labels, asymmetric success flags, infinite references, or feedback
using an additional asymmetry resource. The general no-broadcasting literature
addresses broader but differently quantified settings \cite{sel47-LM,sel47-MS}.
No claim about those variants is obtained by removing the stated premises.

\begin{corollary}[A finite exact independent-output budget]
\label{sel47:cor:budget}
With the same invariant non-reference entrance, suppose a covariant selected
branch produces $k$ independent two-shadow output states, each with nonzero
cross, and a reference state $\tau'$ uncorrelated with them. Let $m'$ be the
largest occupied mode of $\tau'$. Then necessarily
\begin{equation}
 \boxed{k+m'\le m.}
 \label{sel47:eq:budget}
\end{equation}
Thus a reference occupying $L$ consecutive charges can supply at most
$L-1$ such independent coherent outputs on an exact branch even if it is
consumed completely. Returning it exactly would force $k=0$.
\end{corollary}
\begin{proof}
The product output has a nonzero highest mode $k+m'$: it is the tensor
product of every output's nonzero mode-one block and $\tau'^{(m')}$.
The initial input has no modes above $m$. Covariant processing cannot create
that highest mode if $k+m'>m$. The consecutive-charge bound uses $m\le L-1$.
\end{proof}

This is a necessary bandwidth bound, not a sufficient implementation theorem
or a probability estimate. It applies to independent \emph{coherent packets},
not to particle number, generations, pulse count, or physical clock ticks.

\section[Status and the next same-source test]{Distinguish an occurring ray from a ray prepared with a new reference}
\label{sel47:status}

\paragraph{Proved.}
A symmetry-invariant joint packet has an intertwining steering map.
Invariant readout supplemented by a specified finite determinant reference
has the sharp cross budget \eqref{sel47:eq:cross-max}; arbitrary charge
multiplicities give no improvement over the finite-path bound. A pure
reference gives the exact common-ray success \eqref{sel47:eq:pure-capacity},
with different optimal states for raw cross and efficient balanced-ray
success already on three levels. A gapped reference can have the target
stabilizer yet no transferable first mode; two actual copies activate that
mode in the displayed example. The unchanged V46 marker point retains its
purity obstruction after reference assistance. A finite reference cannot be
returned exactly and independently while a nonzero determinant cross is
created on any fixed heralded branch.

\paragraph{Inherited and conditional.}
The V46 joint packet, complete-effect maximum and inaccessible-marker bounds
are inherited. V27's same-word efficiency theorem is unchanged: a ray already
occurring in the source is not required to undergo this invariant erasure
protocol. The comparison group, coefficient sectors and reference charge
action are independently specified. Maximizing over invariant effects is
still an upper bound for a smaller actually admitted bank. In the positive
comparisons the finite reference is a declared additional preparation, not a
resource inferred from an ordinary time calibration. Success weights refer
to normalized preparation sources where stated; channel subpacket weights
must retain their original input effects and physical clock. No dynamical
rate or mass gap is computed.

\paragraph{Still unresolved for the historical source.}
The next source row is the joint action of the comparison group on the
\emph{actual readout implementation}, including any reference preparation
and retained success/failure record. If the source is already asymmetric on
the required common word, its old typed cross is tested directly. If it is
invariant and the proposed readout is invariant with an invariant reference,
selection cannot supply the missing determinant cross. If a reference is
actually present, its first-character block and any irreversible marker
must be evaluated, and a positive capacity still needs an admitted effect.
No historical readout state, determinant-charged reference, V110 word, or
new physical operation has been populated here. A selected subgroup obtained
by importing a reference with that subgroup is a compatibility result, not
an intrinsic historical selection theorem.

\begingroup
\renewcommand{\refname}{Additional primary-source context for V47}
\begin{thebibliography}{9}
\small\raggedright
\bibitem{sel47-LM}
M.~Lostaglio and M.~P. M\"uller,
\emph{Coherence and asymmetry cannot be broadcast},
Phys. Rev. Lett. \textbf{123} (2019), 020403;
\url{https://arxiv.org/abs/1812.08214}.
\bibitem{sel47-MS}
I.~Marvian and R.~W. Spekkens,
\emph{A no-broadcasting theorem for quantum asymmetry and coherence and a
trade-off relation for approximate broadcasting},
Phys. Rev. Lett. \textbf{123} (2019), 020404;
\url{https://arxiv.org/abs/1812.08766}.
The harmonic-mode context is the already cited \cite{sel16-MS}; the
finite product-return proof above is self-contained. These references do not
supply the source with a new quantum reference or a physical readout.
\end{thebibliography}
\endgroup
\addtocontents{toc}{\protect\endgroup}
% END OF SOURCE-SELECTION V47 DEVELOPMENT BLOCK.
% Preserve V1--V47 and append subsequent work before the final terminator.


\clearpage
\hypersetup{pdftitle={Historical Standard-Model Source Selection: V1--V48}}
\begin{center}
{\Large\bfseries Source-selection continuation V48}\\[.4em]
{\large Mixed-reference support, exact common-ray capacity,\\
and the finite cost of a pure reference}
\end{center}
\addtocontents{toc}{\protect\begingroup}
\addtocontents{toc}{\protect\renewcommand{\protect\numberline}{\protect\makebox[2.1em][l]}}

\section[Reference modes do not certify efficient-ray access]{Audit the positive reference state before enlarging its charge inventory}
\label{sel48:context}

V47 distinguishes an already occurring coherent history from an invariant
source read with a separately supplied determinant reference. Its first-mode
norm gives the exact raw cross, but its closed efficient-ray capacity is for
a \emph{pure} specified reference. A nonzero first mode is not a purity
certificate. Enlarging a mixed reference or retrying a rare effect need not
repair this missing condition. We keep the source and operation class fixed
and replace that pure-state premise by an exact test on adjacent positive
reference blocks, including their multiplicities.

The inputs retained without rederivation are V46's accessible-support kernel
criterion, V47's invariant two-shadow packet and first-mode optimum, and its
pure-reference formula. The new results are an explicit mixed-reference
capacity, a sharp conditional purity floor, an all-finite-copy obstruction on
invariant faithful support, equal-summary separating references, and the
all-length optimum of the inherited pure-reference functional. The latter
replaces the uniform-reference linear loss by a quadratically decaying loss;
it is not an instruction to add a reference to the historical source.

The general mode-preservation framework is established
\cite{sel16-MS}. Rank-sensitive coherence distillation and obstructions for
mixed quantum references are also established subjects
\cite{sel48-LRA,sel48-Marvian,sel48-LZS}. The operation classes in those
works are not silently identified with our invariant \emph{readout-only}
protocol. All claims used below are proved in that explicitly stated class;
no general priority or exhaustive literature claim is made.

\paragraph{Fixed source and access convention.}
Retain orthonormal coefficient vectors $b,d$ transforming by $1,\chi^{-1}$,
where $\chi(U,V)=\det U\det V$ and $H=\ker\chi$. Their accessible flags
$|0\rangle,|1\rangle$ transform by $1,\chi$. With identical inaccessible
markers, the joint source is the unnormalized pure packet
\begin{equation}
 |\Omega\rangle=\sqrt\mu\,b\otimes|0\rangle
                  +\sqrt\nu\,d\otimes|1\rangle,
 \qquad \mu,\nu>0.
 \label{sel48:eq:source}
\end{equation}
The reference has a \emph{given} state $\tau$ on the given representation
$R=\bigoplus_nR_n$, with character $\chi^n$ on $R_n$; $\Tr\tau=1$.
It is independent of the source. We optimize over $0\preceq F\preceq I$
on the accessible flag times reference with $[F,N_A+N_R]=0$.
Thus the sole new resource being audited is $\tau$, not a new system map,
charged clock or unconstrained phase-sensitive effect. A smaller admitted
effect bank can only decrease the capacities below.
The target is $v_q=\sqrt q\,b+e^{i\theta}\sqrt{1-q}\,d$, $0<q<1$.
Coefficient mass is denoted by $\alpha$ in $J_F=\alpha P_{v_q}$.
It is a success probability only for a normalized scalar-input preparation
source. For a channel coefficient, its input effect remains
$\alpha K_{v_q}^*K_{v_q}$, as in V47.

\section[Exact capacity from adjacent mixed blocks]{The missing condition is saturation of a positive contraction}
\label{sel48:capacity}

In total readout charge $n$, only $|0\rangle\otimes R_n$ and
$|1\rangle\otimes R_{n-1}$ interact. Write
\begin{equation}
 A_n=P_n\tau P_n,\quad B_n=P_{n-1}\tau P_{n-1},\quad
 C_n=P_n\tau P_{n-1},\qquad
 M_n=\begin{pmatrix}A_n&C_n\\C_n^*&B_n\end{pmatrix}\succeq0.
 \label{sel48:eq:blocks}
\end{equation}
Restrict $A_n,B_n$ to their supports; inverses below refer to those supports.
Positivity implies that $C_n$ has the corresponding source and target supports.
Zero diagonal blocks cannot contribute a target with both shadows and are
omitted. Define the contraction
\begin{equation}
 K_n=A_n^{-1/2}C_nB_n^{-1/2},\qquad \|K_n\|\le1,
 \qquad \mathcal E_n=\ker(I-K_n^*K_n).
 \label{sel48:eq:contraction}
\end{equation}
Let $V_n:\C^{e_n}\to\supp B_n$ be an isometry onto $\mathcal E_n$, and
$U_n=K_nV_n$, also an isometry. Empty spaces contribute zero.

\begin{theorem}[Exact mixed-reference common-ray capacity]
\label{sel48:thm:capacity}
For the fixed source \eqref{sel48:eq:source}, arbitrary finite reference
$\tau$, and invariant readout effects, the maximal mass on $P_{v_q}$ is
\begin{equation}
 \boxed{
 \alpha_q(\tau)=\sum_n\Tr\left[
 \frac q\mu U_n^*A_n^{-1}U_n
 +\frac{1-q}\nu V_n^*B_n^{-1}V_n
 \right]^{-1}.}
 \label{sel48:eq:capacity}
\end{equation}
The trace and inverse in each summand act on $\C^{e_n}$. The maximum is
attained by one invariant effect; it is independent of $\theta$.
For every $0<q<1$, positivity of this capacity is equivalent to
\begin{equation}
 \boxed{
 \exists n:\ e_n>0
 \quad\Longleftrightarrow\quad
 \exists n:\ \rank M_n<\rank A_n+\rank B_n.}
 \label{sel48:eq:rank-test}
\end{equation}
Thus no global purity premise is needed, but an adjacent block must saturate
its Cauchy--Schwarz constraint in a nonzero direction.
\end{theorem}
\begin{proof}
The invariant effect is a direct sum $F=\bigoplus_nF_n$; the selected system
packets add positively. Put $u_q=\sqrt{1-q}\,b-e^{i\theta}\sqrt q\,d$.
The unwanted target mass from block $n$ is $\Tr(F_nD_n)$, where
\begin{equation}
 D_n=\begin{pmatrix}
 (1-q)\mu A_n&-e^{i\theta}\sqrt{q(1-q)\mu\nu}\,C_n\\
 -e^{-i\theta}\sqrt{q(1-q)\mu\nu}\,C_n^*&q\nu B_n
 \end{pmatrix},\qquad
 R_n=\diag(\mu A_n,\nu B_n).
 \label{sel48:eq:unwanted}
\end{equation}
Both are positive and $\Tr(F_nR_n)$ is the selected total mass.
A positive sum is supported on $\C v_q$ exactly when each summand has zero
$u_q$ mass. Since $\Tr(F_nD_n)=0$ iff $F_n$ is supported on $\ker D_n$,
V46's criterion, or this elementary trace identity alone, gives the maximizing
effect $F_n=P_{\ker D_n}$ on the support of $R_n$.

Conjugating $D_n$ by the inverse of
$\diag(\sqrt{\mu(1-q)}A_n^{1/2},\sqrt{\nu q}B_n^{1/2})$
gives $\left(\begin{smallmatrix}I&-e^{i\theta}K_n\\
-e^{-i\theta}K_n^*&I\end{smallmatrix}\right)$.
Its kernel is exactly the graph of $e^{i\theta}K_n$ over $\mathcal E_n$.
A full-column-rank synthesis for $\ker D_n$ is consequently
\[
 Z_n=\begin{pmatrix}
 e^{i\theta}A_n^{-1/2}U_n/\sqrt{\mu(1-q)}\\
 B_n^{-1/2}V_n/\sqrt{\nu q}
 \end{pmatrix}.
\]
Then $F_n=Z_n(Z_n^*Z_n)^{-1}Z_n^*$ and
$Z_n^*R_nZ_n=I/[q(1-q)]$. Cyclicity of trace gives exactly the $n$th term
of \eqref{sel48:eq:capacity}. All blocks are orthogonal, so their projectors
sum to an effect bounded by $I$. The congruence also gives
$e_n=\rank A_n+\rank B_n-\rank M_n$ and proves
\eqref{sel48:eq:rank-test}. Directions in the kernel of $R_n$ add no output;
removing them has not postselected any nonzero source component.
\end{proof}

\begin{corollary}[Nondegenerate charges and the inherited pure formula]
\label{sel48:cor:scalar}
When each occupied charge is one-dimensional, put $p_n=\tau_{nn}$ and
$c_n=\tau_{n,n-1}$. Then
\begin{equation}
 \boxed{
 \alpha_q(\tau)=
 \sum_{\substack{n:p_np_{n-1}>0\\|c_n|^2=p_np_{n-1}}}
 \frac{\mu\nu p_np_{n-1}}
 {q\nu p_{n-1}+(1-q)\mu p_n}.}
 \label{sel48:eq:scalar}
\end{equation}
Only saturated adjacent pairs contribute. For a pure reference every occupied
adjacent pair is saturated, recovering Theorem~\ref{sel47:thm:pure} exactly.
If the inaccessible marker overlap is $r<1$, no two-shadow pure ray is
selectable, regardless of these reference blocks.
\end{corollary}
\begin{proof}
A scalar contraction has an isometric direction precisely when its modulus
is one, proving the formula. A pure reference has
$|c_n|=\sqrt{p_np_{n-1}}$. The marker statement is the inherited V46 bound;
reference preparation pulls every joint effect back to an ordinary effect
on the accessible flag and cannot improve that unrestricted bound.
\end{proof}

\section[The exact conditional-purity floor]{Postselection cannot exceed the best normalized adjacent coherence}
\label{sel48:floor}

Define $\kappa(\tau)=\max_n\|K_n\|$, setting it to zero if no two nonzero
adjacent sectors occur. Retain a marker overlap $r\in[0,1]$, with its phase
incorporated into the chosen target phase. The source and reference supports
are unchanged.

\begin{theorem}[Sharp conditional coherence and balanced-target error]
\label{sel48:thm:floor}
For every nonzero invariant selected packet
$J_F=\left(\begin{smallmatrix}a&z\\\bar z&b'\end{smallmatrix}\right)$,
\begin{equation}
 |z|\le r\kappa(\tau)\sqrt{ab'}.
 \label{sel48:eq:normalized-cross}
\end{equation}
If $p=a/(a+b')$, then
\begin{equation}
 \boxed{
 1-\Tr\left[(J_F/\Tr J_F)^2\right]
 \ge2(1-r^2\kappa(\tau)^2)p(1-p).}
 \label{sel48:eq:purity}
\end{equation}
For any balanced target $v_{1/2}$ with its specified phase,
\begin{equation}
 \boxed{
 \inf_{F:J_F\ne0}
 \left\|J_F/\Tr J_F-P_{v_{1/2}}\right\|_1
 =1-r\kappa(\tau).}
 \label{sel48:eq:distance}
\end{equation}
The infimum is attained with positive success by a rank-one effect in a
maximizing total-charge block, with arbitrarily reduced success also allowed.
These are normalized coefficient-state bounds, not conditioned-channel
diamond distances.
\end{theorem}
\begin{proof}
For one block of $F_n\succeq0$, its off-diagonal entry has a contraction
factorization $F_{10}=F_{11}^{1/2}LF_{00}^{1/2}$ with $\|L\|\le1$.
Using $C_n=A_n^{1/2}K_nB_n^{1/2}$ and Hilbert--Schmidt Cauchy--Schwarz gives
\[
 |\Tr(F_{10}C_n)|
 \le\|K_n\|
 \sqrt{\Tr(F_{00}A_n)\Tr(F_{11}B_n)}.
\]
Include the source weights and marker factor; sum over $n$ and apply
Cauchy--Schwarz again. This proves \eqref{sel48:eq:normalized-cross}.
The purity identity for a two-by-two density gives
\eqref{sel48:eq:purity}. Its expectation on the balanced target is at most
$(1+r\kappa)/2$, so the two-outcome measurement containing that target
bounds the trace distance below by $1-r\kappa$.

For attainment, choose unit singular vectors $u,v$ of some maximizing $K_n$.
An effect vector with components proportional to
$A_n^{-1/2}u/\sqrt\mu$ and $B_n^{-1/2}v/\sqrt\nu$, with the appropriate
relative phase, gives equal selected diagonal masses and normalized cross
$r\kappa/2$. Scale its projector to be bounded by $I$. It has strictly
positive success and normalized packet
$\frac12\left(\begin{smallmatrix}1&r\kappa e^{-i\theta}\\
r\kappa e^{i\theta}&1\end{smallmatrix}\right)$, whose trace distance is
exactly $1-r\kappa$. If there are no adjacent occupied sectors, the
single-sector choices can instead be mixed to equalize the two diagonal
masses; both source weights are positive, and the same conclusion holds
with $\kappa=0$.
\end{proof}

Equations~\eqref{sel48:eq:rank-test} and \eqref{sel48:eq:distance} distinguish
three questions: whether any raw cross is possible ($v_1(\tau)>0$), whether
an exact two-shadow pure ray is possible ($\kappa=1$ and $r=1$), and the
probability at which a prescribed ray can occur (\eqref{sel48:eq:capacity}).
Neither the stabilizer nor the unnormalized first-mode norm determines all
three answers.

\begin{proposition}[Finite-copy obstruction on invariant support]
\label{sel48:prop:faithful}
If $\supp\tau$ is invariant under the reference group, no covariant
trace-nonincreasing operation, with a fixed success label and an independent
invariant non-reference input, can produce a nonzero pure noninvariant
output from $\tau^{\otimes k}$ for any finite $k\ge1$.
In particular every faithful reference on its declared charge carrier has
zero exact common-ray capacity at every finite copy number. A sequence of
invariant readout-only operations, ancillary invariant preparations and
classical retries cannot evade this conclusion on a finite success branch.
\end{proposition}
\begin{proof}
Write $S=\supp\tau$. The state is strictly positive on $S$ and, for finite
$k$, on $S^{\otimes k}$. After fixing the invariant non-reference input and
discarding any unwanted output, a proposed successful protocol defines a
covariant CP map $\Lambda$ on $\B(S^{\otimes k})$. If
$\Lambda(\tau^{\otimes k})=\alpha P_v$, $\alpha>0$, positivity and
faithfulness imply that every Kraus coefficient of $\Lambda$ has range in
$\C v$. Hence $\Lambda(I_{S^{\otimes k}})=cP_v$ for $c>0$.
The input identity is invariant, so covariance forces $P_v$ to be invariant,
a contradiction. Composition and fixed-branch classical selection remain CP
and covariant. An arbitrary number of finite success branches also cannot
sum to an exact pure asymmetric packet, since positivity would force each
nonzero summand onto that ray. No claim is made about a limiting infinite
reference, approximate outputs, asymmetric flags, or noninvariant ancillary
resources.
\end{proof}

This proposition is a direct source-specific use of the known full-support
purification obstruction \cite{sel48-Marvian,sel48-LZS}, not a new general
no-distillation principle. The reference support itself need not be invariant
for a singular mixed reference, which is why a mixed state can still pass
\eqref{sel48:eq:rank-test}.

\section[Equal summaries, different exact capacities]{The first-mode norm and the reference stabilizer do not settle purity}
\label{sel48:examples}

\begin{proposition}[Two three-level references with opposite answers]
\label{sel48:prop:pair}
On charges $0,1,2$, compare the specified states
\begin{equation}
 \tau_A=\frac16\begin{pmatrix}1&1&1\\1&1&1\\1&1&4\end{pmatrix},
 \qquad
 \tau_B=\begin{pmatrix}1/6&1/12&0\\1/12&1/6&1/4\\0&1/4&2/3\end{pmatrix}.
 \label{sel48:eq:pair}
\end{equation}
They have the same populations $(1/6,1/6,2/3)$, the same first-mode norm
$v_1=1/3$, and the same stabilizer $H$. For the normalized balanced source
$\mu=\nu=q=1/2$ and $r=1$, both have maximal raw cross $1/144$.
Nevertheless,
\begin{equation}
 \boxed{
 \alpha_{1/2}(\tau_A)=1/12,\qquad
 \alpha_{1/2}(\tau_B)=0.}
 \label{sel48:eq:pair-capacity}
\end{equation}
The first state has rank two and $\kappa=1$. The second has determinant
$1/288>0$, $\kappa=3/4$, and optimal balanced-target trace distance $1/4$.
No finite number of independent copies of $\tau_B$ permits an exact pure
asymmetric success branch under the stated covariant rules.
\end{proposition}
\begin{proof}
$\tau_A=(\mathbf1\mathbf1^*+3E_{22})/6\succeq0$ has rank two.
The leading principal minors of $\tau_B$ are $1/6,1/48,1/288$, so it is
strictly positive. The adjacent moduli sum to $1/3$ in both cases, and a
nonzero adjacent coherence forces $\chi=1$ for their stabilizers.
For $\tau_A$, only the $0,1$ block saturates, contributing
$(1/6)^2/(1/3)=1/12$ in \eqref{sel48:eq:scalar}; its other adjacent
normalized modulus is $1/2$. For $\tau_B$ these moduli are $1/2,3/4$,
so neither saturates. Apply Theorem~\ref{sel48:thm:floor} and
Proposition~\ref{sel48:prop:faithful}. V47's unchanged first-mode formula
gives $(\mu\nu/4)(1/3)^2=1/144$ for the raw cross.
\end{proof}

This is a comparison of \emph{two specified reference preparations}; their
full matrices are not claimed equal or statistically indistinguishable under
tomography. The system source, character assignment and first-mode summary
are held fixed. The example isolates the exact reference information missing
from those coarser summaries.

\begin{proposition}[Multiplicities contribute through an isometric direction]
\label{sel48:prop:degenerate}
On two adjacent charges of multiplicity two, take
\[
\begin{aligned}
 A&=\diag(1/10,1/5),\qquad B=\diag(3/10,2/5),\\
 K&=\frac15\begin{pmatrix}3&-4\\4&3\end{pmatrix}\diag(1,1/3),\\
 C&=A^{1/2}KB^{1/2},\qquad
 \tau=\begin{pmatrix}A&C\\C^*&B\end{pmatrix}.
\end{aligned}
\]
This is a normalized rank-three reference. For $\mu=2/3$, $\nu=1/3$ and
$q=2/5$, at every specified target phase,
\begin{equation}
 \boxed{\alpha_q(\tau)=25/252.}
 \label{sel48:eq:degenerate}
\end{equation}
\end{proposition}
\begin{proof}
The contraction has singular values $1,1/3$; hence positivity and rank follow
from the block congruence in Theorem~\ref{sel48:thm:capacity}.
Take $V=e_1$, $U=(3,4)^{\mathsf T}/5$. Then
$U^*A^{-1}U=34/5$ and $V^*B^{-1}V=10/3$.
The denominator in \eqref{sel48:eq:capacity} is $252/25$.
The kernel projector constructed in that proof attains its reciprocal and
preserves the full complementary failure outcome.
\end{proof}

\section[Noise and copying: exact and approximate statements differ]{Small reference noise closes the exact branch without destroying its signal}
\label{sel48:noise}

\begin{theorem}[Faithful dephasing of a fixed pure reference]
\label{sel48:thm:noise}
Let $\psi$ have strictly positive populations $p_n$ on $L\ge2$ consecutive
nondegenerate charges. Keep the original reference carrier and source, but
consider the declared noisy reference
\begin{equation}
 \tau_\eta=(1-\eta)|\psi\rangle\langle\psi|
                +\eta\,\Delta(|\psi\rangle\langle\psi|),
 \qquad 0<\eta\le1.
 \label{sel48:eq:noise}
\end{equation}
It is faithful, $\kappa(\tau_\eta)=1-\eta$, and no finite number of copies
can give an exact noninvariant pure output under the stated covariant rules.
At $r=1$, the best normalized balanced-ray error in one use is exactly
$\eta$, at a positive success probability. In fact, applying the
\emph{unchanged optimal pure-reference balanced-ray effect} gives
\begin{equation}
 \boxed{
 J_{\eta}=\frac{\alpha_{1/2}(\psi)}2
 \begin{pmatrix}1&(1-\eta)e^{-i\theta}\\
 (1-\eta)e^{i\theta}&1\end{pmatrix}.}
 \label{sel48:eq:same-effect}
\end{equation}
Its success mass is unchanged and its balanced-target trace distance is
$\eta$. For $\eta<1$ the reference stabilizer and the selected mixed-packet
covariance are still exactly $H$.
\end{theorem}
\begin{proof}
Conjugating $\tau_\eta$ by $\diag(p_n^{-1/2})$ gives, up to the pure
reference phases, $\eta I+(1-\eta)\mathbf1\mathbf1^*$, which is strictly
positive. Every normalized adjacent entry is $1-\eta$.
The full-support and distance statements follow from the preceding results.
The old effect's output diagonals depend only on $p_n$, whereas every cross
entry is linear in $\tau_{n,n-1}$ and is multiplied by $1-\eta$.
At zero noise the output was $\alpha_{1/2}(\psi)P_{v_{1/2}}$; this proves
\eqref{sel48:eq:same-effect} and its trace distance.
The stabilizer assertion uses the nonzero adjacent mode; the two-character
selected packet then has the already proved exact determinant covariance.
\end{proof}

Thus the exact capacity is discontinuous at a pure reference, even though
$\|\tau_\eta-P_\psi\|_1\le2\eta$ and the same-policy success is unchanged.
This does not mean an observable finite discontinuity in the noisy packet:
its trace distance tends to zero. It is the \emph{exact rank-one statement}
that is unstable. Increasing $L$ with fixed $\eta$ does not reduce the
one-use normalized error floor of this specified noise family.

\begin{proposition}[Two copies improve approximation without exact purification]
\label{sel48:prop:copies}
For $\tau_c=\frac12\left(\begin{smallmatrix}1&c\\c&1\end{smallmatrix}\right)$,
$0<c<1$, one has
\begin{equation}
 \boxed{
 \kappa(\tau_c)=c,\qquad
 \kappa(\tau_c^{\otimes2})=\frac{\sqrt2c}{\sqrt{1+c^2}}<1.}
 \label{sel48:eq:copies}
\end{equation}
At $c=3/5$, the optimal balanced-target errors are $2/5$ and
$1-3/\sqrt{17}$, respectively. Both exact capacities are zero, as are those
at every finite copy number.
\end{proposition}
\begin{proof}
The two-copy charge sectors have dimensions $1,2,1$.
The middle diagonal block is
$\frac14\left(\begin{smallmatrix}1&c^2\\c^2&1\end{smallmatrix}\right)$,
the endpoint blocks are $1/4$, and each adjacent block is
$c(1,1)^{\mathsf T}/4$ or its transpose. Whitening gives the single nonzero
singular value in \eqref{sel48:eq:copies} for either adjacent pair.
Theorem~\ref{sel48:thm:floor} proves the sharp errors; faithfulness proves the
all-finite-copy exact exclusion. The joint two-copy invariant measurement
is an explicitly different permitted resource from two separate one-copy
measurements. No uniform error floor as the copy number tends to infinity
is claimed.
\end{proof}

\section[All-length optimal pure-reference success]{A finite recurrence and a quadratic rather than linear size penalty}
\label{sel48:length}

The previous exclusions concern a \emph{specified} imperfect reference.
For comparison, the best exact ray permitted by an independently available
pure reference of given size can also be determined beyond V47's two- and
three-level calculations. Keep the normalized balanced source and target,
$\mu=\nu=q=1/2$, $r=1$. On $L$ consecutive nondegenerate charges define
\begin{equation}
 a_L=\max_{p_j\ge0,\ \sum_{j=0}^{L-1}p_j=1}
       \sum_{j=0}^{L-2}\frac{p_jp_{j+1}}{p_j+p_{j+1}},
 \label{sel48:eq:variational}
\end{equation}
with $0/0=0$. By V47's positive-decomposition argument, the same maximum
holds against all mixed-reference competitors.

\begin{theorem}[Unique optimal populations and an exact finite dual certificate]
\label{sel48:thm:length}
For $L\ge2$ the maximizing population vector is unique, strictly positive
and reflection symmetric. Its value $a_L<1/2$ is determined by the unique
solution, with $0<t_j<1$, of
\begin{align}
 t_0&=1-\sqrt{a_L},\nonumber\\
 t_j&=1-\sqrt{a_L-t_{j-1}^2}\quad(1\le j\le L-2),\nonumber\\
 t_{L-2}&=\sqrt{a_L},\qquad
 \frac{p_{j+1}}{p_j}=\frac{1-t_j}{t_j}.
 \label{sel48:eq:recurrence}
\end{align}
Normalization recovers the populations; $t_j=p_j/(p_j+p_{j+1})$.
Equivalently,
\begin{equation}
 \boxed{
 a_L=\min_{0\le t_j\le1}\max\{
 (1-t_0)^2,\ t_0^2+(1-t_1)^2,\ldots,
 t_{L-3}^2+(1-t_{L-2})^2,\ t_{L-2}^2\}.}
 \label{sel48:eq:dual}
\end{equation}
For $L=2$ the list contains only its two endpoints.
In particular the new four-level optimum is
\begin{equation}
 \boxed{
 a_4=25/64,\qquad p=(3,5,5,3)/16,\qquad
 t=(3/8,1/2,5/8).}
 \label{sel48:eq:L4}
\end{equation}
\end{theorem}
\begin{proof}
The objective is continuous on the simplex. Its value is below $1/2$, since
$4xy/(x+y)\le x+y$ and equality throughout would require constant positive
populations with zero endpoints. If a maximizing vector has a zero component
adjacent to a positive one, scale its positive components by $1-\epsilon$
and place mass $\epsilon$ on that zero. Each adjacent positive component
adds $\epsilon+O(\epsilon^2)$ to the objective, whereas rescaling loses
only $\epsilon a_L$. Since $a_L<1/2$, this strictly improves it.
A zero exists next to a positive component whenever the support is proper,
so the maximizer is interior.

For an interior variation $h$, the Hessian quadratic form is
\[
 -2\sum_{j=0}^{L-2}
 \frac{(p_{j+1}h_j-p_jh_{j+1})^2}{(p_j+p_{j+1})^3}.
\]
Its null space consists of $h$ proportional to $p$; intersecting with the
simplex tangent gives zero. Joint concavity, already verified in V47,
therefore gives uniqueness; reflection gives symmetry. The Lagrange
multiplier equations, and degree-one homogeneity, set every first derivative
equal to $a_L$. Their endpoint and interior expressions are exactly
$(1-t_0)^2$, $t_{j-1}^2+(1-t_j)^2$, and $t_{L-2}^2$, proving the recurrence.
Conversely, a solution in the stated interval gives positive populations
with constant gradient and hence the unique maximum.

For any $t\in[0,1]$, $xy/(x+y)\le(1-t)^2x+t^2y$, with equality when
$t=x/(x+y)$. Summing proves the upper bound in \eqref{sel48:eq:dual};
the just constructed maximizing populations attain it. In \eqref{sel48:eq:L4}
all four dual coefficients equal $25/64$, and direct substitution of the
populations attains that value. This proves the global four-level optimum,
not merely a stationary point.
\end{proof}

\begin{proposition}[Certified inverse-square finite-size loss]
\label{sel48:prop:scaling}
Let $\theta_L=\pi/(L+1)$. Then
\begin{equation}
 \boxed{
 \frac{1-\cos\theta_L}{2}
 \le\frac12-a_L\le1-\cos\theta_L.}
 \label{sel48:eq:scaling}
\end{equation}
Thus the optimal success deficit is of order $L^{-2}$, whereas a uniform
reference has deficit exactly $1/(2L)$. No sharp leading asymptotic constant
is claimed in \eqref{sel48:eq:scaling}.
\end{proposition}
\begin{proof}
Put $p_{-1}=p_L=0$ and $u_j=\sqrt{p_j}$. An exact rearrangement gives
\[
 \frac12-\sum_j\frac{p_jp_{j+1}}{p_j+p_{j+1}}
 =\frac14\sum_{j=-1}^{L-1}
 \frac{(p_{j+1}-p_j)^2}{p_{j+1}+p_j}.
\]
Each summand lies between $(u_{j+1}-u_j)^2$ and twice that quantity.
The minimum Dirichlet path energy on $\sum_j u_j^2=1$ is
$2(1-\cos\theta_L)$, attained by the normalized sine vector.
The lower inequality applies to every population vector. Testing the sine
vector and the upper summand inequality gives the upper inequality.
The uniform value is the inherited formula with balanced weights.
\end{proof}

At four levels the exact optimal balanced-ray mass is $25/64$, with squared
same-ray cross $625/16384$. Applying any nonzero dephasing fraction $\eta$ of
Theorem~\ref{sel48:thm:noise} leaves that same-policy success unchanged, but
makes its output rank two with best one-use balanced-target error $\eta$.
The optimized size and the reference's positive support are therefore
independent requirements; increasing the former does not establish the latter.

\section[Status and the actual next source row]{Separate character content, saturated support, and physical reference provenance}
\label{sel48:status}

\paragraph{Proved.}
The mixed-reference efficient-ray capacity is an explicit sum of inverses
on the isometric directions of adjacent whitened blocks. The corresponding
maximum normalized coherence gives a sharp purity and balanced-target error
bound. Invariant faithful support excludes exact asymmetric pure selection
at every finite copy number. Two explicit references share populations,
first-mode norm and stabilizer but have opposite exact-ray answers.
Multiplicities, a noisy pure reference and its unchanged readout effect, and
a collective two-copy approximation are evaluated without altering the
system packet. The best balanced exact-ray success on $L$ charges has a
unique positive optimizer, a finite dual/recurrence certificate, the exact
four-level value $25/64$, and a certified $L^{-2}$ deficit.

\paragraph{Inherited and conditional.}
V47's raw cross formula and its pure-reference special case remain unchanged;
V46 supplies the support-selection principle. Theorem~\ref{sel48:thm:capacity}
is its symmetry-resolved evaluation, not a new general operator-short
principle. The source weights, reference charge action, marker and all
allowed effects are stated before optimization. The positive examples are
reference-assisted compatibility constructions, not a mechanism selecting
the reference or its stabilizer. Neither the ordinary clock nor a
coefficientwise continuum identity supplies a pure determinant reference.
An already occurring efficient historical word still needs no auxiliary
reference from this protocol. No new particle, generation, coupling,
continuum measure or mass-gap statement is made.

\paragraph{Still unresolved for the historical source.}
The next finite row on a reference-assisted branch is the actual adjacent
reference block $M_n$, including its support and physical ancestry. A
nonzero first mode only establishes possible interference; exact efficiency
requires the rank defect in \eqref{sel48:eq:rank-test} and then an admitted
kernel effect. Strict positivity of a reconstructed adjacent block is a
one-sided exclusion when certified above its error. A near-zero eigenvalue
is not an exact saturation certificate. System-only data do not identify
this reference row. No historical $\tau$, readout effect or V110 word has
been populated. The remaining historical question is its actual occurrence,
not whether a larger mathematically chosen reference has a favourable
capacity.

\begingroup
\renewcommand{\refname}{Additional primary-source context for V48}
\begin{thebibliography}{9}
\small\raggedright
\bibitem{sel48-LRA}
L.~Lami, B.~Regula and G.~Adesso,
\emph{Generic bound coherence under strictly incoherent operations},
Phys. Rev. Lett. \textbf{122} (2019), 150402;
\url{https://arxiv.org/abs/1809.06880}.
\bibitem{sel48-Marvian}
I.~Marvian, \emph{Coherence distillation machines are impossible in quantum
thermodynamics}, Nature Communications \textbf{11} (2020), 25;
\url{https://doi.org/10.1038/s41467-019-13846-3}.
\bibitem{sel48-LZS}
C.~L. Liu, D.~L. Zhou and C.~P. Sun,
\emph{Criterion for a state to be distillable via stochastic incoherent
operations} (2021);
\url{https://arxiv.org/abs/2112.06168}.
These sources provide context for support-sensitive coherence conversion,
not permission to identify their operation classes with the present
invariant readout. The capacity and comparison proofs above are self-contained.
\end{thebibliography}
\endgroup
\addtocontents{toc}{\protect\endgroup}
% END OF SOURCE-SELECTION V48 DEVELOPMENT BLOCK.
% Preserve V1--V48 and append subsequent work before the final terminator.


\clearpage
\hypersetup{pdftitle={Historical Standard-Model Source Selection: V1--V49}}
\begin{center}
{\Large\bfseries Source-selection continuation V49}\\[.4em]
{\large Exact finite-copy access from the actual reference support\\
Exterior-power selection, charge periods, and mixed-reference activation}
\end{center}
\addtocontents{toc}{\protect\begingroup}
\addtocontents{toc}{\protect\renewcommand{\protect\numberline}{\protect\makebox[2.1em][l]}}

\section[The finite-copy question is a support question]{Move upstream of one-copy saturation without supplying a new reference}
\label{sel49:context}

V48 gives the exact one-copy capacity from adjacent positive blocks. It also
excludes exact asymmetric pure selection at every finite copy number when
the reference support is invariant under the whole charge group. Between
these statements lies a genuine question: can a singular mixed reference,
with no saturated adjacent one-copy block, produce an exact common ray by
joint use of its own copies? We do not improve the reference by assumption.
Instead we determine the symmetry of its actual support, prove a finite-copy
existence criterion, and give a constructive invariant selection that attains
positive success precisely on the primitive support branch.

The inherited inputs are Theorem~\ref{sel48:thm:capacity}, the pure-reference
capacity already proved in V47, and the support-sensitive obstruction in
V48. The exterior-power identity used below is ordinary finite linear
algebra, not a new fermionic postulate. Pure-state asymmetry conversions and
full-support distillation obstructions are established subjects
\cite{sel49-MS,sel48-Marvian}. The finite-copy support-period criterion and
its source-specific realization are proved here; no claim of general
priority or equivalence with other incoherent-operation classes is made.

\paragraph{Fixed physical entrance.}
Keep the actual reference representation
$R=\bigoplus_n R_n$, with integer charge operator
$N=\sum_n n\Pi_n$ and $V_g|_{R_n}=\chi(g)^nI$, where
$\chi(g)=\det U\det V$ and $H=\ker\chi$. The reference is an arbitrary
specified density $\tau$, not a fitted pure state. Write
\begin{equation}
 P=\supp\tau,\qquad r=\rank\tau,\qquad D=\dim R.
 \label{sel49:eq:entrance}
\end{equation}
We use the same invariant source
$|\Omega\rangle=\sqrt\mu\,b\otimes|0\rangle+
\sqrt\nu\,d\otimes|1\rangle$, with $\mu,\nu>0$, coefficient characters
$1,\chi^{-1}$, and flag characters $1,\chi$. Inaccessible source markers
are identical. The target is the specified ray
$v_q=\sqrt q\,b+e^{i\theta}\sqrt{1-q}\,d$, $0<q<1$.
For $k$ \emph{actually independently prepared} reference copies, effects
$0\preceq F\preceq I$ on flag times $R^{\otimes k}$ must commute with
$N_A+\sum_{j=1}^kN_j$. The full interval of such invariant effects is the
stated access class. No system correction, extra asymmetric ancilla,
reference purification supplied from outside, or recovered inaccessible
marker is allowed. Success mass has the same meaning as in V48; it is an
ordinary probability only for a normalized preparation source. Restricting
the actual effect bank can invalidate attainability, but not the exclusions.

\section[Support symmetry survives pure selection]{The stabilizer of the support is stronger than the stabilizer of the state}
\label{sel49:support}

For any compact-group representation on a finite reference, define
\begin{equation}
 G_P=\{g:V_gPV_g^*=P\}.
 \label{sel49:eq:Gp}
\end{equation}
This need not equal $\operatorname{Stab}(\tau)$: an element can preserve
the occupied subspace while rotating a nonuniform density inside it.

\begin{lemma}[Pure outputs retain every input-support symmetry]
\label{sel49:lem:necessary}
Let $\Lambda$ be a fixed-label covariant CP map from a finite reference to
an output representation. If
$\Lambda(\tau^{\otimes k})=\alpha P_v$ with $\alpha>0$ and $P_v$ a pure
state, then every $g\in G_P$ fixes $P_v$. This remains true with independent
invariant ancillary inputs, discarded outputs, and covariant adaptive
processing, when the complete success branch is retained as one CP map.
\end{lemma}
\begin{proof}
There are positive constants $a,b$ with
$aP\preceq\tau\preceq bP$. Thus
$a^k\Lambda(P^{\otimes k})\preceq\alpha P_v
\preceq b^k\Lambda(P^{\otimes k})$.
Positivity forces $\Lambda(P^{\otimes k})$ to be a nonzero multiple of
$P_v$. Since $P^{\otimes k}$ is invariant under $G_P$, covariance forces
that same invariance of $P_v$. Appending an invariant input and taking a
partial trace are covariant CP operations, so the identical argument
applies to the composed branch. No assertion about a most likely outcome
or its asymptotic rate is used.
\end{proof}

For the fixed two-shadow source, inserting an invariant effect induces a
covariant CP map from the reference to the coefficient output. The lemma
therefore applies directly to the readout-only protocol, even though the
invariant source contains system--flag correlations.

Define the integer support period by
\begin{equation}
 g(P)=\gcd\{|n-m|:\ n\ne m,\ \Pi_nP\Pi_m\ne0\},
 \qquad \gcd\varnothing=0.
 \label{sel49:eq:period}
\end{equation}
The convention $g=0$ means that all charge blocks of $P$ are diagonal.

\begin{proposition}[A finite period test and fixed-support capacity comparison]
\label{sel49:prop:period}
For the determinant-character representation,
\begin{equation}
 G_P=
 \begin{cases}
 G_0,&g(P)=0,\\
 \{g\in G_0:\chi(g)^{g(P)}=1\},&g(P)>0.
 \end{cases}
 \label{sel49:eq:stabilizer}
\end{equation}
Consequently an exact target $P_{v_q}$ in the fixed protocol requires
$g(P)=1$, at every finite copy number.
Moreover, if $a\tau\preceq\tau'\preceq b\tau$ for $a,b>0$, the same
invariant effects select exact rays from either state, and their optimal
$k$-copy target masses obey
\begin{equation}
 a^k\alpha_q(\tau^{\otimes k})
 \le\alpha_q((\tau')^{\otimes k})
 \le b^k\alpha_q(\tau^{\otimes k}).
 \label{sel49:eq:comparison}
\end{equation}
Thus exact feasibility depends on support, whereas the positive yield
retains the reference's actual eigenvalues.
\end{proposition}
\begin{proof}
The $(n,m)$ block of $V_gPV_g^*$ is
$\chi(g)^{n-m}\Pi_nP\Pi_m$. All its nonzero blocks are fixed exactly when
all the corresponding integer powers equal one. B\'ezout's identity gives
\eqref{sel49:eq:stabilizer}. Since $P_{v_q}$ has stabilizer exactly $H$,
Lemma~\ref{sel49:lem:necessary} gives the necessary period one. Tensoring
the operator inequalities and applying the positive readout map gives
$a^kJ_F(\tau^{\otimes k})\preceq J_F((\tau')^{\otimes k})
\preceq b^kJ_F(\tau^{\otimes k})$. The supports of the two nonzero outputs
therefore agree. Optimize their scalar masses on the fixed target ray to
obtain \eqref{sel49:eq:comparison}.
\end{proof}

This strengthens the invariant-support obstruction without changing it.
A faithful reference has $P=I$ and $g(P)=0$. A singular reference with
$g(P)=2$ also fails the primitive test, even if its density has a nonzero
first character mode and stabilizer exactly $H$.

\section[Invariant selection of the support exterior line]{A mixed reference can yield a pure reference without an imported phase}
\label{sel49:exterior}

Let $\mathsf A_r=r!^{-1}\sum_{\pi\in S_r}\operatorname{sgn}(\pi)U_\pi$
be the orthogonal antisymmetrizer on $R^{\otimes r}$. It commutes with
$V_g^{\otimes r}$ for \emph{every} unitary representation $V$.
Let $\lambda_1,\ldots,\lambda_r>0$ be the nonzero eigenvalues of $\tau$.
Choose an orthonormal support basis $u_1,\ldots,u_r$ and denote its normalized
exterior vector by $\Psi_P=u_1\wedge\cdots\wedge u_r$. Its ray is independent
of that basis; it is a state on the retained exterior-power carrier, not
an assertion that the ray alone is an invariant representation.

\begin{theorem}[Exact invariant exterior selection]
\label{sel49:thm:exterior}
The original $r$-copy reference satisfies
\begin{equation}
 \boxed{
 \mathsf A_r\tau^{\otimes r}\mathsf A_r
 = e_r(\tau)P_{\Psi_P},\qquad
 e_r(\tau)=\prod_{j=1}^r\lambda_j>0.}
 \label{sel49:eq:exterior}
\end{equation}
Its conditional pure output has exactly the support stabilizer:
\begin{equation}
 \operatorname{Stab}(P_{\Psi_P})=G_P.
 \label{sel49:eq:exterior-stab}
\end{equation}
The success is $e_r(\tau)\le r^{-r}$. The construction consumes $r$
reference copies and is covariant; it does not use their preparation
purification or an asymmetric auxiliary state.
\end{theorem}
\begin{proof}
In an eigenbasis of $\tau$, its $r$th exterior power has the products of
$r$ distinct eigenvalues as eigenvalues. Only the product of all $r$
nonzero eigenvalues survives. Its vector is $\Psi_P$, proving
\eqref{sel49:eq:exterior}; the same identity follows directly by summing
the permutation matrix elements. A unitary preserves the line
$\C\Psi_P$ exactly when it preserves $\Ran P$. For example,
$\Ran P=\{x:x\wedge\Psi_P=0\}$ when $r<D$, and the $r=D$ case is
immediate. This proves \eqref{sel49:eq:exterior-stab}. The arithmetic--geometric
mean inequality and $\sum_j\lambda_j=1$ give the probability bound.
\end{proof}

The antisymmetrizer is a collective operation on distinguishable prepared
reference copies. It is not an identification of those copies as physical
fermions, a CAR postulate, the matter determinant incidence, or an admitted
history of the Store. In the stated full invariant-effect class it can also
be absorbed into a single final readout effect, as proved next.

\begin{proposition}[The selected charge law is a polynomial of the old density]
\label{sel49:prop:minors}
Choose any orthonormal charge eigenbasis indexed by $i$ with charges $n_i$.
For each $r$-element index set $I$, let $\tau_I$ be its principal submatrix.
The unnormalized charge weights of \eqref{sel49:eq:exterior} are
\begin{equation}
 w_N=\sum_{|I|=r,\ \sum_{i\in I}n_i=N}\det\tau_I,\qquad
 p_N=\frac{w_N}{e_r(\tau)},\qquad \sum_Nw_N=e_r(\tau).
 \label{sel49:eq:minors}
\end{equation}
They are independent of basis changes inside the charge sectors.
If $\mathcal S_P=\{N:w_N>0\}$, then
\begin{equation}
 \gcd\{N-M:N,M\in\mathcal S_P\}=g(P).
 \label{sel49:eq:wedge-period}
\end{equation}
Equivalently, $w_N$ is the coefficient of $t^rz^N$ in the finite Laurent
polynomial $\det(I+t z^N\tau)$, where here $z^N$ denotes the matrix
$\sum_nz^n\Pi_n$.
\end{proposition}
\begin{proof}
The compound matrix of $\tau$ in the normalized exterior basis has diagonal
entries $\det\tau_I$. Taking its charge-resolved trace gives
\eqref{sel49:eq:minors}. Expanding a determinant by principal minors gives
the polynomial formula. In a pure state with occupied charge set
$\mathcal S_P$, an element fixes the density exactly when all occupied
relative charge phases agree. Its stabilizer is thus the subgroup determined
by the gcd of these differences. Compare with
\eqref{sel49:eq:exterior-stab} and \eqref{sel49:eq:stabilizer}.
\end{proof}

\section[Complete finite-copy existence criterion]{The support period, not a search through successively larger copies}
\label{sel49:finite}

\begin{lemma}[Finite convolution produces an adjacent pair exactly at period one]
\label{sel49:lem:bezout}
Let $\mathcal S\subset\mathbb Z$ be a nonempty finite set. Some finite
sumset $\ell\mathcal S$ contains two consecutive integers exactly when
$\gcd(\mathcal S-\mathcal S)=1$.
More explicitly, if integers $c_s$ satisfy
\begin{equation}
 \sum_{s\in\mathcal S}c_s=0,\qquad
 \sum_{s\in\mathcal S}s c_s=1,
 \label{sel49:eq:bezout}
\end{equation}
then $\ell=\frac12\sum_s|c_s|$ suffices.
\end{lemma}
\begin{proof}
Every difference of two $\ell$-fold sums is a multiple of the gcd, proving
necessity. For sufficiency choose a base charge $s_0$ and a B\'ezout
combination of $s-s_0$ equal to one. Setting $c_{s_0}$ to the negative sum
of the other coefficients gives \eqref{sel49:eq:bezout}. Its positive and
negative parts each use $\ell$ charges, and their sums differ by one.
\end{proof}

\begin{theorem}[Exact finite-copy determinant-ray alternative]
\label{sel49:thm:main}
For the fixed source, reference and invariant readout class of
Section~\ref{sel49:context}, the following are equivalent:
\begin{enumerate}[label=\textup{(\roman*)},leftmargin=2.4em]
\item some finite number $k$ of the original independent reference copies
permits $J_F=\alpha P_{v_q}$ with $\alpha>0$;
\item $G_P=H$;
\item $g(P)=1$.
\end{enumerate}
No global purity or one-copy adjacent-saturation assumption is required.
A constructive bound is $k=r\ell$, with $\ell$ from
Lemma~\ref{sel49:lem:bezout} applied to \eqref{sel49:eq:minors}.
Write $p^{*\ell}$ for the convolution of the normalized exterior charge law.
The construction achieves at least
\begin{equation}
 \boxed{
 \alpha_q(\tau^{\otimes r\ell})\ \ge\
 e_r(\tau)^\ell
 \sum_N
 \frac{\mu\nu\,p_N^{*\ell}p_{N-1}^{*\ell}}
 {q\nu p_{N-1}^{*\ell}+(1-q)\mu p_N^{*\ell}}
 >0,}
 \label{sel49:eq:yield}
\end{equation}
with zero terms omitted. This is a guaranteed yield, not a general optimal
copy count or a sharp capacity formula for all $r\ell$-copy references.
\end{theorem}
\begin{proof}
Necessity is Lemma~\ref{sel49:lem:necessary} and
Proposition~\ref{sel49:prop:period}. For sufficiency, apply
$\mathsf A_r$ to each of $\ell$ disjoint groups of $r$ copies. The event has
probability $e_r(\tau)^\ell$ and produces
$\Psi_P^{\otimes\ell}$. Its occupied charges are $\ell\mathcal S_P$,
with positive probabilities $p^{*\ell}$. Distinct charge patterns belong
to orthogonal tensor subspaces, so no amplitude cancellation removes a sum charge. Period one and
Lemma~\ref{sel49:lem:bezout} supply an occupied adjacent pair. V47's
pure-reference formula, or the scalar saturated case of
Theorem~\ref{sel48:thm:capacity}, gives the sum in \eqref{sel49:eq:yield}.

For a readout-only implementation, let
$B=\mathsf A_r^{\otimes\ell}$ and let $F_0$ be the invariant final effect
on the filtered support, extended by zero. Use the single effect
$F=(I_A\otimes B)F_0(I_A\otimes B)$. It lies in $[0,I]$ and commutes with
the full charge. Contracting it against the original reference produces
exactly the stated sequential result. Thus no independent noninvariant gate
has been added by using the filtering proof. The complementary effect
$I-F$ retains the rest of the original source; forgetting the selection
still gives $\mu P_b+\nu P_d$.
\end{proof}

There is also a general version for a target with characters $1,\chi^{-\Delta}$
and compensating flag $1,\chi^\Delta$: a nonzero integer $\Delta$ is
accessible by this finite-copy exact protocol precisely when
$\Delta\in g(P)\mathbb Z$, with $0\mathbb Z=\{0\}$. The proof replaces
one by $\Delta$ in \eqref{sel49:eq:bezout}. This describes different
specified target gaps, not an unannounced change to the primitive V110
packet. In particular a support period two can carry a pure second harmonic
without supplying the primitive determinant coherence.

\section[An exact two-copy mixed-reference activation]{The complete rank-two qutrit alternative and a normalized attaining effect}
\label{sel49:qutrit}

\begin{theorem}[One copy, two copies, or never on three consecutive charges]
\label{sel49:thm:qutrit}
Let $\tau$ have rank two on charges $0,1,2$, with all three diagonal
populations positive, and put
$d_{ij}=\det\tau_{\{i,j\}}\ge0$. For every prescribed target $v_q$:
\begin{enumerate}[label=\textup{(\roman*)},leftmargin=2.4em]
\item if $d_{01}=0$ or $d_{12}=0$, exact positive success is already possible
with one copy;
\item if all three $d_{01},d_{02},d_{12}$ are positive, the minimum number of
copies is exactly two;
\item if $d_{02}=0$, primitive exact success is impossible at every finite
copy number.
\end{enumerate}
These cases are disjoint and exhaustive. If a diagonal population vanishes,
the rank-two support is a charge-coordinate plane and primitive exact
success is likewise impossible.
\end{theorem}
\begin{proof}
Write $\tau_{ij}=\langle x_i,x_j\rangle$ for three nonzero vectors spanning
a two-dimensional space. A zero $d_{ij}$ means those two vectors are
parallel. Two zero minors would make all three parallel, contradicting rank
two. Thus at least two minors are positive. V48 gives one-copy success exactly
at a zero adjacent minor. For two copies, the exterior vector has occupied
charges $1,2,3$, with unnormalized weights $d_{01},d_{02},d_{12}$.
All-positive weights therefore give adjacent coherence and two-copy success,
while both original adjacent minors are positive and exclude one-copy
success. If only the middle weight is zero, the occupied exterior charges
are $1,3$, so $g(P)=2$ by \eqref{sel49:eq:wedge-period} and
Theorem~\ref{sel49:thm:main} excludes every finite copy number. A zero
diagonal forces its entire row and column to vanish by positivity; rank
two fills the remaining coordinate plane.
\end{proof}

\begin{proposition}[A rank-two reference activated exactly at two copies]
\label{sel49:prop:activation}
Consider the declared reference preparation
\begin{equation}
 \tau_*=\frac{1}{4}\begin{pmatrix}1&0&1\\0&1&1\\1&1&2\end{pmatrix},
 \qquad
 P_*=\frac{1}{3}\begin{pmatrix}2&-1&1\\-1&2&1\\1&1&2\end{pmatrix}.
 \label{sel49:eq:activation}
\end{equation}
Its spectrum is $(3/4,1/4,0)$, all three $d_{ij}=1/16$, and its one-copy
capacity is zero. For the normalized balanced source and target,
\begin{equation}
 \boxed{\alpha_{1/2}(\tau_*)=0,\qquad
 \alpha_{1/2}(\tau_*^{\otimes2})=\frac{1}{16}.}
 \label{sel49:eq:activation-cap}
\end{equation}
The second value is the full two-copy optimum, not only a lower bound.
\end{proposition}
\begin{proof}
In the normalized exterior basis
$a_{ij}=(|ij\rangle-|ji\rangle)/\sqrt2$,
\begin{equation}
 \mathsf A_2\tau_*^{\otimes2}\mathsf A_2
 =\frac{1}{16}
 \begin{pmatrix}1&1&-1\\1&1&-1\\-1&-1&1\end{pmatrix}_{a_{01},a_{02},a_{12}}
 =\frac{3}{16}P_\psi,
 \quad \psi=\frac{a_{01}+a_{02}-a_{12}}{\sqrt3}.
 \label{sel49:eq:wedge-star}
\end{equation}
Its three charges are $1,2,3$. The uniform pure-reference balanced capacity
is $1/3$, giving $(3/16)(1/3)=1/16$ without discarding this preparation cost.
An explicit single invariant effect on flag times the \emph{original two
copies} is
\begin{align}
 F&=|f_2\rangle\langle f_2|+|f_3\rangle\langle f_3|,\nonumber\\
 f_2&=(|0\rangle a_{02}+|1\rangle a_{01})/\sqrt2,\qquad
 f_3=(-|0\rangle a_{12}+|1\rangle a_{02})/\sqrt2.
 \label{sel49:eq:effect}
\end{align}
The vectors are orthonormal and have total charges two and three. Direct
contraction gives
$J_F=\frac1{32}\left(\begin{smallmatrix}1&1\\1&1\end{smallmatrix}\right)$.
For optimality, apply the unwanted-output kernel formula in the proof of
Theorem~\ref{sel48:thm:capacity} to $\tau_*^{\otimes2}$. The total-charge
one and four paired blocks have zero kernel; the paired blocks at charges
two and three have one-dimensional kernels, exactly $\C f_2$ and
$\C f_3$. Boundary blocks cannot support a target with both shadows.
Each of the two kernel projectors selects mass $1/32$, and no admissible
effect can exceed its projector. Their sum is therefore optimal.
\end{proof}

The normalized selected output is an exact balanced ray; its unnormalized
squared cross is $1/1024$. Replacing identical inaccessible markers by overlap
$r_{\rm mark}<1$ gives instead
$J_F=\frac1{32}\left(\begin{smallmatrix}1&r_{\rm mark}\\
r_{\rm mark}&1\end{smallmatrix}\right)$. The independent marker obstruction
has not been removed by processing the reference.

\begin{proposition}[A singular coherent state whose support blocks every primitive ray]
\label{sel49:prop:period-two}
On the same three charges take
\begin{equation}
 \tau_{\rm per}=
 \begin{pmatrix}
 1/4&\sqrt2/8&1/4\\
 \sqrt2/8&1/2&\sqrt2/8\\
 1/4&\sqrt2/8&1/4
 \end{pmatrix},\qquad
 P_{\rm per}=\begin{pmatrix}1/2&0&1/2\\0&1&0\\1/2&0&1/2\end{pmatrix}.
 \label{sel49:eq:period-two}
\end{equation}
It has spectrum $(3/4,1/4,0)$ and state stabilizer exactly $H$, with
$v_1(\tau_{\rm per})=\sqrt2/4>0$. Its support has period two, so primitive
exact capacity is zero at \emph{every} finite copy number.
The two-copy exterior event succeeds with probability $3/16$ but prepares
$(a_{01}-a_{12})/\sqrt2$, whose charge difference is two, not one.
\end{proposition}
\begin{proof}
Use the support basis $(|0\rangle+|2\rangle)/\sqrt2,|1\rangle$.
In that basis the state is
$\left(\begin{smallmatrix}1/2&1/4\\1/4&1/2\end{smallmatrix}\right)$.
This proves the spectrum and support. Its adjacent density entries are
nonzero, giving state period one, but its support projector has only the
nonzero off-diagonal gap two. Its minors are
$(d_{01},d_{02},d_{12})=(3/32,0,3/32)$; the compound matrix gives the stated
exterior vector. Apply Theorem~\ref{sel49:thm:main}.
\end{proof}

The unchanged V48 reference
$\tau_A=\frac16\left(\begin{smallmatrix}1&1&1\\1&1&1\\1&1&4\end{smallmatrix}\right)$
lies in the first case: its minors are $(0,1/12,1/12)$, and its already
proved one-copy balanced capacity is $1/12$. Its exterior event has mass
$1/6$ and prepares $(a_{02}+a_{12})/\sqrt2$. That particular two-copy
route has total target mass $1/24$ and is not claimed optimal. The new
criterion therefore retains the old example without imposing a pure
reference or replacing its density.

\section[Copy costs and exact-rank limitations]{Period one is finite feasibility, not a dimension-only resource bound}
\label{sel49:cost}

\begin{proposition}[Fixed dimension can require arbitrarily many reference copies]
\label{sel49:prop:delay}
For $m\ge2$, the pure reference
$\psi_m=(|0\rangle+|2\rangle+|2m+1\rangle)/\sqrt3$ has support period
one. Its minimum copy number for a primitive exact target in the stated
readout class is exactly $m$.
\end{proposition}
\begin{proof}
A charge of $\psi_m^{\otimes k}$ is $2a+(2m+1)b$ with nonnegative
$a,b$ and $a+b\le k$. Two such charges differ by one only if
$2x+(2m+1)y=\pm1$, with $|x|\le k$. Necessarily $y$ is odd and nonzero,
so $|x|\ge m$. No $k<m$ works. At $k=m$, $m$ copies of charge two have
total $2m$, whereas one charge $2m+1$ and $m-1$ zero charges have total
$2m+1$. Both have positive amplitude. The pure-reference capacity is then
positive. The state dimension remains three but its specified charge range
grows; no bounded-charge counterexample is asserted.
\end{proof}

\begin{proposition}[Rank-filling noise changes exact feasibility, not discontinuously the output]
\label{sel49:prop:noise}
For $r<D$, put $\tau_\eta=(1-\eta)\tau+\eta I_D/D$, $0<\eta\le1$.
Its support is $I_D$, hence primitive exact capacity is zero at every finite
copy number. Nevertheless the \emph{same} $r$-copy exterior filter has
success $e_r(\tau_\eta)$ and fidelity with its former conditional pure output
\begin{equation}
 F_\eta=\frac{\prod_{i=1}^{r}((1-\eta)\lambda_i+\eta/D)}
 {e_r(\tau_\eta)},\qquad
 \left\|\frac{\mathsf A_r\tau_\eta^{\otimes r}\mathsf A_r}
 {e_r(\tau_\eta)}-P_{\Psi_P}\right\|_1=2(1-F_\eta).
 \label{sel49:eq:noise}
\end{equation}
Both the selected unnormalized state and its nonzero normalized output
converge to the old ones as $\eta\downarrow0$.
\end{proposition}
\begin{proof}
The old eigenbasis diagonalizes $\tau_\eta$, with the displayed shifted
nonzero eigenvalues and $D-r$ further eigenvalues $\eta/D$. The exterior
spectrum consists of all products of $r$ distinct such eigenvalues.
The old support exterior vector is one eigenvector, of the numerator weight
in \eqref{sel49:eq:noise}; all others are orthogonal. This proves the exact
fidelity and trace distance. The full-support exclusion follows from
Proposition~\ref{sel49:prop:period}, not an approximation argument.
\end{proof}

For $\tau_*$, the same two-copy effect \eqref{sel49:eq:effect} under this
specified noise gives the directly checked packet
\begin{equation}
 J_\eta=\begin{pmatrix}
 1/32+\eta/16-11\eta^2/288&(\eta-3)(\eta-1)/96\\
 (\eta-3)(\eta-1)/96&1/32+\eta/24-5\eta^2/288
 \end{pmatrix}.
 \label{sel49:eq:noise-packet}
\end{equation}
It is rank two for every $0<\eta\le1$ and approaches the exact ray
continuously. The unknown historical source is not assigned this noise.
The point is that a small measured eigenvalue is not a certificate of the
exact rank needed by the positive theorem.

\section[Status and the source-level stopping rule]{One support-period computation replaces an unbounded exact-copy search}
\label{sel49:status}

\paragraph{Proved.}
On the stated finite integer-charge representation and full invariant
readout class, finite-copy primitive efficient selection is possible
exactly at support period one. The exterior-power event constructs a pure
reference with the support's exact stabilizer, including a quantified
preparation cost and a finite B\'ezout copy bound. The occupied exterior
charges are reconstructed from principal minors of the \emph{original}
reference density. Rank-two qutrit references on three consecutive charges
have a complete one-copy/two-copy/never alternative. A mixed reference with
zero one-copy capacity attains an exact balanced ray at two copies with
optimal success $1/16$. Another singular reference has the correct state
stabilizer and nonzero first mode but is excluded at every finite copy
number by its support period two. Charge width, eigenvalue yields and
rank-filling noise retain separate resource costs.

\paragraph{Inherited and conditional.}
The one-copy capacity and marker obstruction of V48, the pure-reference
formula of V47 and the accessible-effect distinctions of V46 remain
unchanged. The new sufficient operation is a collective invariant effect
on physically available independent copies. It is not classical processing
of fine records, a reusable unmodified reference, or a free preparation of
copies. The support symmetry is a \emph{resource already present in the
reference}; an invariant source and invariant reference cannot generate it.
The exterior antisymmetrizer does not identify matter statistics or infer a
new Standard-Model representation. No historical word, coupling, generation,
physical clock or continuum measure is constructed.

\paragraph{Still unresolved for the historical source.}
The next datum on this reference-assisted route is the actual reference
support projector and its identified charge action, together with independent
copy preparation and the admitted invariant readout bank. If $g(P)\ne1$, no
finite-copy search in this class can supply the primitive exact ray. If
$g(P)=1$, the theorem supplies an explicit positive construction and yield
bound; whether its effect occurs in the old physical grammar remains a
separate test. There is no need to keep guessing increasingly large adjacent
blocks merely to decide \emph{existence}. An already occurring coherent
V110 word remains a different direct-history route and is not required to
pass this reference protocol. Historical Standard-Model selection is not
asserted closed.

\begingroup
\renewcommand{\refname}{Additional primary-source context for V49}
\begin{thebibliography}{9}
\small\raggedright
\bibitem{sel49-MS}
I.~Marvian and R.~W.~Spekkens,
\emph{The theory of manipulations of pure state asymmetry: basic tools and
equivalence classes of states under symmetric operations},
New J. Phys. \textbf{15} (2013), 033001;
\url{https://arxiv.org/abs/1104.0018}.
The pure-state symmetry framework is background. The support-period,
exterior selection and exact finite-copy statements used here are proved
above. The support-invariant obstruction is also related to the already
cited \cite{sel48-Marvian}; no result from a different operation class is
used as a substitute for the stated proofs.
\end{thebibliography}
\endgroup
\addtocontents{toc}{\protect\endgroup}
% END OF SOURCE-SELECTION V49 DEVELOPMENT BLOCK.
% Preserve V1--V49 and append subsequent work before the final terminator.


\clearpage
\hypersetup{pdftitle={Historical Standard-Model Source Selection: V1--V50}}
\begin{center}
{\Large\bfseries Source-selection continuation V50}\\[.4em]
{\large Prepare the reference from the original source\\
Record-dependent support, exact capacity, and irreversible label loss}
\end{center}
\addtocontents{toc}{\protect\begingroup}
\addtocontents{toc}{\protect\renewcommand{\protect\numberline}{\protect\makebox[2.1em][l]}}

\section[Audit the preparation, not another reference ansatz]{The actual source output precedes the support-period test}
\label{sel50:context}

V49 decides finite-copy feasibility once a reference density and its charge
representation are given. It does not identify their preparation. The next
useful step is therefore to examine the original instrument's \emph{joint
classical--quantum output}, rather than choose another reference with favorable
support. We derive the exact support-period rule after classical processing
of the old records and evaluate the unchanged determinant-action compiler of
Theorem~\ref{sel19:thm:witness}. Its known invariant input produces usable pure
references when one original bit is retained. Its forgotten output is faithful
on an invariant two-dimensional support and cannot produce the same exact ray
at any finite copy number. The difference persists at every positive event
count and has an explicit classical-noise test.

\paragraph{Inherited inputs and new scope.}
The compiler, its full native carrier, its helpers and its determinant action
are inherited from V19; its response on the invariant input is also the
previously evaluated V23 probe. The one-copy pure-reference formula of
Theorem~\ref{sel47:thm:pure}, the mixed-reference saturation criterion of
Theorem~\ref{sel48:thm:capacity}, and the finite-copy support-period theorem of
Theorem~\ref{sel49:thm:main} are inputs, not new results. The additions are the
record-resolved preparation criterion, its behavior under an arbitrary
classical confusion matrix, and the exact all-count evaluation of this old
source. Classical instrument postprocessing and symmetry-constrained state
conversion are established frameworks \cite{sel50-instruments,sel50-asymmetry}.
The specialized proofs below do not claim priority for positive-support joins,
block-diagonal classical registers, or likelihood overlap estimates.

\paragraph{Fixed protocol.}
Write $G_0=\mathrm U(3)\times\mathrm U(2)$,
$\chi(U,V)=\det U\det V$, and $H=\ker\chi$. The independently identified
reference carrier is a finite integer-charge representation
$\mathcal R=\bigoplus_n\mathcal R_n$, with $V_g=\sum_n\chi(g)^n\Pi_n$.
The source used in the \emph{readout} stage is the same normalized balanced
preparation packet
\begin{equation}
 |\Omega\rangle=(b\otimes|0\rangle_A+d\otimes|1\rangle_A)/\sqrt2,
 \qquad b\mapsto b,\quad d\mapsto\chi^{-1}d,
 \label{sel50:eq:target-source}
\end{equation}
where the two flags carry charges zero and one. Its inaccessible markers are
identical. The target is $(b+d)/\sqrt2$; a prescribed relative phase is handled
by the corresponding invariant effect. Full invariant effect access on the
flag and the available quantum reference is a stated assumption. References
are consumed. No coherent access to the \emph{classical preparation record}
is assumed. All probabilities below are for this normalized preparation
source; for a general channel coefficient packet the same calculations give
coefficient mass, and its actual input effect must still be retained.

\section[Exact feasibility with a retained classical record]{The record changes the support before any distillation is attempted}
\label{sel50:record}

Let $\{\Phi_a\}$ be an actual recorded preparation instrument on an actual
specified input $\rho$. Its unnormalized output states are
\begin{equation}
 R_a=\Phi_a(\rho)\succeq0,\qquad
 \sum_a\Tr R_a=1.
 \label{sel50:eq:prep}
\end{equation}
No efficiency assumption is made. Suppose the retained classical output is
$y$, with the actual stochastic matrix $M_{y a}\ge0$,
$\sum_y M_{ya}=1$. Its complete output is
\begin{equation}
 T_y=\sum_a M_{ya}R_a,\qquad
 \tau_M=\bigoplus_y T_y,\qquad P_y=\supp T_y.
 \label{sel50:eq:flagged}
\end{equation}
The record transforms trivially under $G_0$. Zero-probability bins are omitted.
For a support projection $P$, use the V49 convention
\begin{equation}
 g(P)=\gcd\{|n-m|:\ n\ne m,\ \Pi_nP\Pi_m\ne0\},
 \qquad \gcd\varnothing=0.
 \label{sel50:eq:g}
\end{equation}

\begin{theorem}[A source-output criterion for all finite copy numbers]
\label{sel50:thm:flagged}
Under the fixed protocol, independent preparations of the \emph{actual}
classical--quantum output $\tau_M$ can supply the exact target ray at some
finite copy number with positive success if and only if
\begin{equation}
 \boxed{g_M:=\gcd_y g(P_y)=1.}
 \label{sel50:eq:flagged-period}
\end{equation}
The full invariant-effect optimum is unchanged when effects are required to
be diagonal in every retained classical record. In particular, attainability
does not require a coherent eraser of those classical labels.
Furthermore,
\begin{equation}
 \boxed{\Ran P_y=\sum_{a:M_{ya}>0}\Ran R_a.}
 \label{sel50:eq:join}
\end{equation}
Thus exact feasibility depends only on the output supports and the zero
pattern of the classical channel, while success weights retain its actual
positive entries. Additional classical processing can only enlarge the
support stabilizer, not create finite-copy exact feasibility that was absent
before processing.
\end{theorem}
\begin{proof}
For any $x$,
$\langle x,T_yx\rangle=\sum_a M_{ya}\|R_a^{1/2}x\|^2$.
Its kernel is exactly the intersection of the selected kernels, proving
\eqref{sel50:eq:join}. The support of $\tau_M$ is $\bigoplus_yP_y$.
Its $(n,m)$ charge block is nonzero exactly when one of the corresponding
blocks of $P_y$ is nonzero. Its V49 support period is therefore $g_M$.
Theorem~\ref{sel49:thm:main}, including its positive constructive direction,
applies to this finite representation with charge multiplicities.

Let $\mathsf P$ pinch an effect in all classical record projectors. The source
state and $\tau_M^{\otimes k}$ are already block diagonal in those records, so
$F$ and $\mathsf P(F)$ induce identical system outputs. Pinching is unital and
positive; hence $0\le\mathsf P(F)\le I$. It preserves invariance because the
record projectors commute with the charge action. This proves that the V49
construction can be realized by record-conditioned quantum effects, without
coherent operations on the record itself. It is an existence statement under
the fixed full-effect access, not admission of every conditional effect into
the historical grammar.

Finally a unitary preserving every $\Ran R_a$ preserves each of their sums.
Thus the intersection of the original support stabilizers is a subgroup of
the intersection after coarse processing. Equivalently, the latter cannot
lose a support symmetry that formerly obstructed the target. This also
follows by composing a putative successful protocol with the classical
channel, which is a covariant operation.
\end{proof}

\begin{remark}[One good bin is sufficient, not necessary]
The gcd in \eqref{sel50:eq:flagged-period} is over \emph{all} retained bins.
For example, on charges $0,2,3$, a labelled equal mixture of the pure states
$(|0\rangle+|2\rangle)/\sqrt2$ and
$(|0\rangle+|3\rangle)/\sqrt2$ has individual support periods two and three,
but joint period one. Two independent samples with different labels have
charges $0,2,3,5$ with equal probabilities and conditional balanced-ray
capacity $1/8$. The different-label event has probability $1/2$, so the full
two-sample capacity is $1/16$. Equal-label pairs have zero capacity. This
example only illustrates the theorem; it is not a source assigned to the
Store, and neutral record labels do not themselves carry the missing phase.
\end{remark}

\begin{corollary}[A covariant preparation cannot generate its own missing resource]
\label{sel50:cor:origin}
If $\rho$ is $G_0$-invariant and every original preparation branch is
fixed-label $G_0$-covariant, then $g_M=0$ for every classical record policy.
\end{corollary}
\begin{proof}
Every $R_a$ is invariant, hence commutes with the charge action. Its support
and every sum in \eqref{sel50:eq:join} do also. Apply \eqref{sel50:eq:g}.
\end{proof}
This is the source-provenance safeguard: a positive preparation result below
must use asymmetry of an already specified source operation, not infer it
from neutral classical labels.

\section[The unchanged compiler prepares the reference]{One old preparation and one original bit give exact success}
\label{sel50:native}

Keep the \emph{entire} V19 native carrier $\mathcal H_{14}^{\otimes5}$, its
five neutral helpers, its original six labels, and the exact point
$\tau=1/8$, $s=4\pi/3$ of \eqref{sel19:eq:exact-point}. The known invariant
input is $p_n=|n\rangle\langle n|$, where $n=r_0^{\otimes5}$. The previously
identified $\omega=\Omega_C\otimes\Omega_W$ has character $\chi$.
Their two-dimensional span $Q\mathcal H$ reduces every original coefficient.
It is used only to calculate this preparation, not substituted for the full
native representation.

All helper terms vanish on $Q\mathcal H$, so the restriction of the original
six coefficients is
\begin{equation}
 K_e|_Q=I_2/\sqrt{18}+\epsilon_e U/3,
 \quad \epsilon_e=+1\ (e=1,2,3),\ -1\ (e=4,5,6),
 \quad U=I_2-\frac{1+i}{2}\begin{pmatrix}1&1\\1&1\end{pmatrix}.
 \label{sel50:eq:restriction}
\end{equation}
Retaining only this original sign bit loses no quantum information on this
prepared branch. Its restricted coefficients are
\begin{equation}
 L_\pm=\frac1{\sqrt3}(I_2/\sqrt2\pm U),\qquad
 R_\pm=L_\pm p_nL_\pm^*.
 \label{sel50:eq:L}
\end{equation}
The grouped maps need not be efficient on the rest of the full native
carrier. No such claim is used.

\begin{theorem}[Exact source-prepared reference capacity]
\label{sel50:thm:one-step}
The unchanged source produces
\begin{equation}
 R_\pm=\begin{pmatrix}
 \frac13\pm\frac{\sqrt2}{6}&
 \mp\frac{\sqrt2}{12}+i(\frac16\pm\frac{\sqrt2}{12})\\
 \mp\frac{\sqrt2}{12}-i(\frac16\pm\frac{\sqrt2}{12})&\frac16
 \end{pmatrix},
 \quad p_\pm=\Tr R_\pm=\frac12\pm\frac{\sqrt2}{6}.
 \label{sel50:eq:R}
\end{equation}
Both conditional references are pure and occupy both charges zero and one.
Using one prepared output with the retained sign and invariant readout, the
maximum exact balanced-target success is
\begin{equation}
 \boxed{A_1=\frac4{21}.}
 \label{sel50:eq:A1}
\end{equation}
For the normalized target preparation \eqref{sel50:eq:target-source}, the
success packet and its cross are
\begin{equation}
 J_{\rm succ}=\frac2{21}\begin{pmatrix}1&1\\1&1\end{pmatrix},
 \qquad \Tr J_{\rm succ}=\frac4{21},\qquad h=\frac4{441}.
 \label{sel50:eq:success}
\end{equation}
Forgetting the preparation bit instead gives
\begin{equation}
 \bar\rho_1=R_++R_-=
 \begin{pmatrix}2/3&i/3\\-i/3&1/3\end{pmatrix},
 \qquad \det\bar\rho_1=\frac19.
 \label{sel50:eq:forgotten}
\end{equation}
Its support is the full invariant $Q$; no finite number of its independent
copies can supply the exact target in the fixed covariant protocol.
\end{theorem}
\begin{proof}
Equations~\eqref{sel50:eq:restriction}--\eqref{sel50:eq:L} follow by evaluating
the \emph{unchanged} compiler, since its first contrast column is
$(1,1,1,-1,-1,-1)^T/\sqrt6$. Multiplication gives \eqref{sel50:eq:R}.
Write $v_\pm=L_\pm n=(a_\pm,b_\pm)^T$. The weights are positive and both
amplitudes are nonzero. The V47 formula on a pure two-charge reference,
including its original preparation probability, gives
\begin{equation}
 A_{1,\pm}=
 \frac{|a_\pm|^2|b_\pm|^2}{|a_\pm|^2+|b_\pm|^2}
 =\frac2{21}\pm\frac{\sqrt2}{42}.
 \label{sel50:eq:individual}
\end{equation}
Classical pinching shows that the optima for the two labels add. Their sum
is $4/21$; no unrecorded coherent recombination of those labels is used.

For an explicit attaining effect, inside the total-charge-one sector with
basis $|0\rangle_A|1\rangle_R,|1\rangle_A|0\rangle_R$, use
\begin{equation}
 f_\pm=(\overline a_\pm,\overline b_\pm)^T/\sqrt{p_\pm},
 \qquad F_\pm=|f_\pm\rangle\langle f_\pm|,
 \label{sel50:eq:effect}
\end{equation}
and zero outside that sector. It selects the coefficient
$a_\pm b_\pm(b+d)/\sqrt{2p_\pm}$.
The combined effect is diagonal in the sign record and invariant on flag
and reference. Its complement is positive and restores the original
unconditioned target marginal on summation. This proves attainability,
normalization and \eqref{sel50:eq:success}. Finally
\eqref{sel50:eq:forgotten} is faithful on $Q$. V49 gives the all-finite-copy
exclusion because $g(Q)=0$.
\end{proof}

The exact probability $4/21$ includes the actual sign probabilities. It is
per prepared source output, not conditioned on choosing its more favorable
sign. At one fixed locked opportunity its unconditional value includes the
unchanged acceptance factor $\vartheta$: it is $4\vartheta/21$. Waiting until
one genuine acceptance gives the displayed accepted-event distribution, not
a different pulse or reference. No elapsed-time recalibration is performed.

The output reference and its charge-zero/one support are supplied by this
specified compiler rather than a new reference ansatz. \emph{The compiler
already contains the symmetry-breaking action $H_s=s(P+S)$.} Thus this
calculation transfers a known source resource; it does not derive its
initial determinant incidence or prove historical selection. The joint
invariant effect \eqref{sel50:eq:effect} still needs physical admission.

\section[All original event counts can be evaluated]{One retained count suffices; forgetting it gives an all-depth obstruction}
\label{sel50:counts}

The two restricted coefficients $L_+,L_-$ commute. Let $m\ge1$ be the number
of original accepted events and $r$ the number with positive sign. Define
\begin{equation}
 v_{m r}=L_+^rL_-^{m-r}n,
 \qquad R_{m r}=\binom{m}{r}|v_{m r}\rangle\langle v_{m r}|.
 \label{sel50:eq:count-state}
\end{equation}
Their sum is normalized. The count $r$ is a deterministic coarse function
of the existing fine record. It is not a newly controlled pulse schedule.

\begin{theorem}[Exact support and minimum ray-preserving record]
\label{sel50:thm:counts}
For every $m\ge1$, all $m+1$ vectors $v_{m r}$ have two nonzero charge
components and define pairwise distinct rays. Every retained count bin is
therefore a primitive pure reference. Any positive bin mixing two different
counts has full invariant support $Q$, and cannot supply an exact asymmetric
ray at any finite copy number on its own.

The count partition is the coarsest deterministic partition preserving the
conditional output ray of \emph{every} original word at that horizon. A weaker
policy retaining only the existence of some usable ray need not retain every
count: one exclusive nonzero count bin already suffices.

In contrast, with all pulse records forgotten the output is
\begin{equation}
 \boxed{
 \bar\rho_m=\frac12\begin{pmatrix}
 1+\Re z^m&i\Im z^m\\-i\Im z^m&1-\Re z^m
 \end{pmatrix},\quad
 z=\frac{1+2i}{3},\quad
 \det\bar\rho_m=\frac{1-(5/9)^m}{4}>0.}
 \label{sel50:eq:all-time}
\end{equation}
No finite collection of independently prepared, unrecorded outputs at
positive event counts can supply the exact target in the covariant protocol.
\end{theorem}
\begin{proof}
Use the fixed eigenvectors $f_0=(n-\omega)/\sqrt2$ and
$f_1=(n+\omega)/\sqrt2$ of $U$, with eigenvalues $1,-i$. Put
\begin{align}
 a_{m r}&=\left(\frac{1/\sqrt2+1}{\sqrt3}\right)^r
          \left(\frac{1/\sqrt2-1}{\sqrt3}\right)^{m-r},\
 b_{m r}&=\left(\frac{1/\sqrt2-i}{\sqrt3}\right)^r
          \left(\frac{1/\sqrt2+i}{\sqrt3}\right)^{m-r}.
 \label{sel50:eq:eigen-amplitudes}
\end{align}
Then $v_{m r}=(a_{m r}f_0+b_{m r}f_1)/\sqrt2$ and
\begin{equation}
 \frac{|a_{m r}|^2}{|b_{m r}|^2}
 =\frac{(3+2\sqrt2)^{2r-m}}{3^m}.
 \label{sel50:eq:ratio}
\end{equation}
If a charge component vanished, $a_{m r}=\pm b_{m r}$ would force this ratio
to equal one. Taking the algebraic norm in $\mathbb Q(\sqrt2)$ would then give
$1=3^{2m}$, impossible for $m\ge1$. Different $r$ give different positive
ratios in \eqref{sel50:eq:ratio}, so they cannot define the same ray.
Two different rays span $Q$; positivity makes the support of their nonzero
mixture exactly $Q$. This proves the record and support statements without
searching through $6^m$ histories.

On this reducing block, forgetting the bit at each event gives exactly
$\Lambda=\frac13\id+\frac23\operatorname{Ad}_U$.
Its Bloch multiplier on $z_{
m Bloch}-iy_{
m Bloch}$ is $(1+2i)/3$.
The initial value for $p_n$ is one. This gives \eqref{sel50:eq:all-time} and
its determinant. Any finite tensor product of such unrecorded states is
faithful on an invariant tensor support. The support-symmetry obstruction
of V49 excludes the pure asymmetric target, also after covariant processing.
\end{proof}

The limiting forgotten state is $Q/2$. This calculation concerns the
prepared reducing two-dimensional branch, not every stationary density of
the full helper-bearing compiler. It is consistent with V36's separate
stationary-support calculation. A nonzero off-diagonal density element at a
finite time does not replace the support condition needed for exact selection.

\begin{proposition}[Nonvanishing yield from the original count record]
\label{sel50:prop:yield}
The exact one-use balanced-ray capacity of the count-retaining reference is
\begin{equation}
 A_m=\sum_{r=0}^m\binom{m}{r}
 \frac{|(v_{m r})_0|^2|(v_{m r})_1|^2}
 {\|v_{m r}\|^2}.
 \label{sel50:eq:Am}
\end{equation}
It satisfies
\begin{equation}
 \boxed{
 \frac{1-(\sqrt{2/3})^m}{4}\le A_m\le\frac14,
 \qquad \lim_{m\to\infty}A_m=\frac14.}
 \label{sel50:eq:Am-bound}
\end{equation}
The first exact values are
\begin{equation}
 A_1=\frac4{21},\qquad A_2=\frac{934}{4365},\qquad
 A_3=\frac{245074}{1306557}.
 \label{sel50:eq:first-values}
\end{equation}
In particular, the finite-count capacity is \emph{not} monotone in $m$.
\end{proposition}
\begin{proof}
The state in each count bin is pure, so the inherited two-charge formula
gives \eqref{sel50:eq:Am} with its unchanged binomial probability.
Writing $a=a_{m r}$, $b=b_{m r}$, its two charge populations give
\begin{equation}
 \frac14-A_m=
 \sum_r\binom{m}{r}
 \frac{(\Re(a\bar b))^2}{2(|a|^2+|b|^2)}\ge0.
 \label{sel50:eq:deficit}
\end{equation}
Use $(\Re(a\bar b))^2\le |ab|^2$ and
$|a|^2+|b|^2\ge2|ab|$. The right side is at most
$\frac14\sum_r\binom{m}{r}|a_{m r}b_{m r}|$.
The binomial sum equals
\[
 \left[
 \sqrt{\frac{3+2\sqrt2}{6}\frac12}+
 \sqrt{\frac{3-2\sqrt2}{6}\frac12}
 \right]^m=(\sqrt{2/3})^m.
\]
This proves \eqref{sel50:eq:Am-bound}. Direct substitution at $m=1,2,3$
gives \eqref{sel50:eq:first-values}, including the strict decrease from two
to three. This is the elementary likelihood-overlap mechanism already used
in the earlier nondemolition analysis, now evaluated for this unchanged
source and this exact-ray objective.
\end{proof}

The index $m$ counts the original accepted source events, not independent
reference copies and not physical ticks. Preparing several such references
requires independent repetitions of the known input preparation. On a fixed
sequence of $m$ raw opportunities, the all-accepted branch retains the factor
$\vartheta^m$; its target success is $\vartheta^m A_m$. Allowing the original
idle waits changes the duration law, not the conditional reference formulas.

\section[Classical errors can remove every exact ray]{A positive confusion entry is not inverted by statistical reconstruction}
\label{sel50:confusion}

At a fixed $m$, let $M_{y r}$ be any actual classical readout channel on the
count record. A bin $y$ is \emph{exclusive} for $r$ when $M_{y r}>0$ and
$M_{y s}=0$ for every $s\ne r$. The unnormalized count states in
\eqref{sel50:eq:count-state} already include their original probabilities.

\begin{theorem}[Exact record-loss criterion on the unchanged source]
\label{sel50:thm:confusion}
The processed classical--quantum reference permits the primitive exact ray
at some finite copy number if and only if at least one nonzero bin is
exclusive. When no exclusive bin exists, no finite-copy covariant protocol
can supply it. When exclusive bins exist, the exact one-use capacity is
\begin{equation}
 \boxed{
 A_m(M)=\sum_{r=0}^m\ \sum_{y\text{ exclusive for }r}M_{yr}\binom{m}{r}
 \frac{|(v_{m r})_0|^2|(v_{m r})_1|^2}{\|v_{m r}\|^2}.}
 \label{sel50:eq:conf-capacity}
\end{equation}
Thus, if every readout bin confuses at least two count values, arbitrarily
small but strictly positive mixing can eliminate exact finite-copy capacity.
\end{theorem}
\begin{proof}
Each exclusive bin has a primitive rank-one support; every other nonzero bin
has full invariant support $Q$ by Theorem~\ref{sel50:thm:counts}. The global
period of the flagged support is one exactly when some exclusive bin occurs;
otherwise it is zero. Apply Theorem~\ref{sel50:thm:flagged}.
For a single use, a positive sum can have the prescribed rank-one output only
if every nonzero selected bin has that same pure output. A full-support bin
cannot contribute such an output under the invariant protocol. Exclusive
bins independently attain their inherited pure-reference optima, which add
as \eqref{sel50:eq:conf-capacity}. There is no signed cancellation of mixed
states and no coherent erasure of classical record values.
\end{proof}

\begin{proposition}[Exact binary-error floor]
\label{sel50:prop:BSC}
At one event, let the old sign pass through a binary symmetric classical
channel with error $0\le\delta\le1/2$, and put $t=1-2\delta$.
Its two unnormalized output blocks are
$T_\pm=(R_++R_-)/2\pm t(R_+-R_-)/2$. They obey
\begin{equation}
 \det T_\pm=\frac{1-t^2}{36},\qquad
 \kappa_\pm^2=
 \frac{t^2\pm\sqrt2t+1}{2\pm\sqrt2t}.
 \label{sel50:eq:noise}
\end{equation}
For $0<\delta\le1/2$, the exact ray capacity vanishes at every finite copy
number. In a single use, the best normalized balanced-target trace distance,
allowing arbitrarily rare success, is
\begin{equation}
 \boxed{
 d_{\min}(\delta)=
 1-\sqrt{\frac{t^2+\sqrt2t+1}{2+\sqrt2t}}.}
 \label{sel50:eq:noise-floor}
\end{equation}
For $0<\delta<1/2$ the classical confusion matrix is linearly invertible,
so the two old output matrices are statistically reconstructible, despite
the exact physical selection obstruction.
\end{proposition}
\begin{proof}
Substitute \eqref{sel50:eq:R}; the two determinants give full support for
$\delta>0$. Each charge sector is one-dimensional, so its normalized adjacent
coherence is $|T_{01}|/\sqrt{T_{00}T_{11}}$, giving
\eqref{sel50:eq:noise}. For $0\le t<1$, the plus value is maximal, since
$\kappa_\pm^2=1-(1-t^2)/(2\pm\sqrt2t)$.
V48's sharp one-use error theorem applies blockwise; selecting the better
classical bin attains $1-\kappa_+$. The trace-norm lower bound applies also to
mixtures of bins by $|z|\le\max\kappa_\pm\sqrt{ab}$. Finally the confusion
matrix has determinant $1-2\delta=t$, but its inverse has negative entries
for $0<\delta<1/2$. It is a statistical inverse, not a positive operation on
the unknown reference.
\end{proof}

At $\delta=0$, the capacity is $4/21$. At every positive error it is zero,
while the best one-use approximation error in \eqref{sel50:eq:noise-floor}
tends continuously to zero as $\delta\downarrow0$. This is a change only to
the declared classical readout channel; the quantum instrument and its
forgotten marginal are unchanged. At positive error the recovered coarse
density may still have exact stabilizer $H$. That fact cannot certify the
perfect support required for an efficient common ray. Approximate multi-copy
improvements are not excluded or evaluated by the one-use error formula.
With uncertain data, strict positivity can certify exclusion; a small
estimated determinant cannot certify the exact zero required for success.

\section[Status and the actual source packet]{Use the old preparation and its retained record before supplying a reference}
\label{sel50:status}

\paragraph{Proved.}
For an actual recorded preparation, finite-copy exact feasibility after any
specified classical readout channel is determined by the gcd of its
record-conditioned support periods. The criterion remains attainable with
only classical access to the preparation labels and full invariant quantum
readout of the conditional reference. On the unchanged V19 source and its
known invariant input, one original bit yields exact balanced-ray success
$4/21$. At every positive event count, the count of signs gives distinct pure
primitive reference rays; its exact one-use capacity converges to $1/4$ with
an explicit exponential bound but is not monotone. Forgetting that count
always gives faithful invariant support and an all-finite-copy exclusion.
A complete confusion-channel test identifies which original bins retain
exact capacity. Binary record error has an explicit sharp one-use
approximation floor even when its statistical inverse exists.

\paragraph{Inherited and conditional.}
The conversion criteria, readout normalization, and support obstruction of
V47--V49 are retained, not reproved as independent discoveries. The known
native preparation, action $H_s$, charge assignment, helpers, and original
coefficients are those of V19. There is no new reference density selected to
fit a desired answer and no change to the quantum instrument or its clock.
The two-dimensional calculations describe a prepared reducing branch of the
full $14^5$ carrier. Independent repetitions and invariant joint effect
access remain physical conditions. A smaller admitted effect bank can reduce
attainable capacity. Classical labels are used as controls, not as coherent
quantum references or as new gauge-charged states.

\paragraph{Still unresolved for the historical source.}
The positive benchmark has a determinant-sensitive action at its entrance;
therefore its prepared reference does not constitute an intrinsic derivation
of that action or of the Standard Model. For the historical programme the
next packet is the actual normalized preparation, its original joint
record--quantum output, the actual record-confusion/forgetting map, and the
identified charge action. Those objects determine the relevant support before
V49 is invoked. A source that is invariant on every retained branch cannot
bootstrap the missing reference from neutral records. Conversely, an already
occurring reference-bearing branch need not be replaced by an independently
postulated pure reference. The V110 directed typing and one-history incidence
remain separate; none is populated here. No stationarity, reference purity,
recoverability, or new readout is promoted to a universal axiom of the finite
duality. Historical Standard-Model selection is not asserted closed.

\begingroup
\renewcommand{\refname}{Primary-source context for V50}
\begin{thebibliography}{9}
\small\raggedright
\bibitem{sel50-instruments}
L.~Lepp\"aj\"arvi and M.~Sedl\'ak,
\emph{Post-processing of quantum instruments},
Phys. Rev. A \textbf{103} (2021), 022615;
\url{https://arxiv.org/abs/2010.15816}.
\bibitem{sel50-asymmetry}
I.~Marvian and R.~W.~Spekkens,
\emph{The theory of manipulations of pure state asymmetry: basic tools and
equivalence classes of states under symmetric operations},
New J. Phys. \textbf{15} (2013), 033001;
\url{https://arxiv.org/abs/1104.0018}.
These sources provide established operation-class context. The complete
source-specific record-period criterion and calculations used in this layer
are proved above, with the explicit V47--V49 inherited inputs. No claim of
general priority, physical measurement, or equivalence with other classes
of incoherent operations is made.
\end{thebibliography}
\endgroup
\addtocontents{toc}{\protect\endgroup}
% END OF SOURCE-SELECTION V50 DEVELOPMENT BLOCK.
% Preserve V1--V50 and append subsequent work before the final terminator.


\clearpage
\hypersetup{pdftitle={Historical Standard-Model Source Selection: V1--V51}}
\begin{center}
{\Large\bfseries Source-selection continuation V51}\\[.4em]
{\large From a fixed prepared state to the original channel\\
Support symmetries, finite-query obstructions, and invariant neutral probes}
\end{center}
\addtocontents{toc}{\protect\begingroup}
\addtocontents{toc}{\protect\renewcommand{\protect\numberline}{\protect\makebox[2.1em][l]}}

\section[Audit the preparation restriction before demanding a record]{A fixed output obstruction is not a channel obstruction}
\label{sel51:context}

Theorem~\ref{sel50:thm:one-step} prepares a reference using the specified
invariant input $n$, and proves that forgetting its sign record gives an
invariant full support on the $n/\omega$ plane. That exact statement concerns
independent copies of that prepared output. It is not a theorem about every
invariant preparation or every use of the forgotten channel. The next
upstream question is whether the original channel itself has a larger
support symmetry, and whether another \emph{invariant} preparation on its
already present neutral sector can avoid the obstruction.

This layer answers both questions. A support-preserving group of the Choi
matrix obstructs every finite covariant adaptive query, not only copies of
one output. On the character-supported branch, the V49 criterion lifts to a
necessary-and-sufficient finite-query criterion. The unchanged V19 source
then provides two sharply different restrictions: its two-level forgotten
channel has an all-finite-query primitive-ray obstruction, but access to one
existing neutral helper gives a one-query, record-free positive construction.

\paragraph{Inherited mathematics and the operation class.}
The source coefficients and vectors are those of
\eqref{sel19:eq:compiler}--\eqref{sel19:eq:H}; the evaluated time and action
are unchanged. The pure-reference capacity
(Theorem~\ref{sel47:thm:pure}), mixed-reference support test
(Theorem~\ref{sel48:thm:capacity}), and finite-copy support-period theorem
(Theorem~\ref{sel49:thm:main}) are inherited, not rederived here.
Quantum-channel resource-purification obstructions, including adaptive
protocols, are established \cite{sel51-FangLiu}. The support comparison below
is a self-contained symmetry specialization, not a general priority claim.
The covariant-operation and character-mode framework is likewise established
\cite{sel51-Marvian,sel51-Modes}.

The positive constructions permit newly specified \emph{invariant}
preparations and readout effects on the original carrier. Their admission is
an explicit capability hypothesis, not a consequence of algebra membership.
In particular, neutral vectors in different classical superselection sectors
cannot be coherently combined unless that preparation is genuinely admitted.
No old effect is inverted and no inaccessible environment is measured.
All fine preparation outcomes are forgotten before the new quantum readout.

Write $G_0=\mathrm U(3)\times\mathrm U(2)$,
$\chi(g)=\det U\det V$, and $H=\ker\chi$. As in V50, the target readout
source is the specified normalized invariant packet
\begin{equation}
 |\Xi\rangle=(b|0\rangle_A+d|1\rangle_A)/\sqrt2,
 \quad b\mapsto b,\quad d\mapsto\chi^{-1}d,
 \quad |0\rangle_A\mapsto|0\rangle_A,\quad
 |1\rangle_A\mapsto\chi|1\rangle_A.
 \label{sel51:eq:target}
\end{equation}
Its markers are identical and the target is $v=(b+d)/\sqrt2$.
An effect on $A$ and the available reference must be invariant, and the
reference is consumed. The stated success optima concern this \emph{fixed
steering protocol}, not arbitrary state-conversion operations. For a general
channel coefficient packet the same output weights are coefficient masses;
its input effect must also be retained.

\section[Channel-support symmetries obstruct every finite query]{A finite-query support comparison with a quantitative remainder}
\label{sel51:queries}

Let $\Phi:\B(H_{\rm in})\to\B(H_{\rm out})$ be a CPTP map with Choi
matrix $J$, and let $P=\supp J$. The Choi representation is
$R_g=U_{\rm out}(g)\otimes\overline{U_{\rm in}(g)}$. Define the closed
subgroup
\begin{equation}
 G_P=\{g:R_gPR_g^*=P\}.
 \label{sel51:eq:GP}
\end{equation}
A finite-query protocol uses independent memoryless calls to $\Phi$,
invariant ancillary preparations, and $G_0$-covariant intervening operations.
It may keep quantum memory, use classical feed-forward, and herald a fixed
success label. That label is not allowed to encode an uncounted asymmetric
reference. The queried primitive is the \emph{forgotten channel}; granting
access to its original environment or fine instrument changes the primitive.

\begin{lemma}[Same Choi support gives the same success support at finite depth]
\label{sel51:lem:comparison}
If two channels have the same Choi support, there exist constants
$0<c\le1\le C<\infty$ such that
\begin{equation}
 c\Psi\preceq_{\rm CP}\Phi\preceq_{\rm CP}C\Psi.
 \label{sel51:eq:CPorder}
\end{equation}
For the unnormalized successful states $X_k,Y_k$ of the same protocol using
at most $k$ calls to these channels, respectively,
\begin{equation}
 c^kY_k\preceq X_k\preceq C^kY_k.
 \label{sel51:eq:queryorder}
\end{equation}
Consequently their success supports coincide, including whether success is
zero. No division by a small success probability is needed for this statement.
\end{lemma}
\begin{proof}
Both Choi matrices are positive definite on their common support. The least
and greatest eigenvalues of
$J_\Psi^{-1/2}J_\Phi J_\Psi^{-1/2}$ on that support give $c,C$.
Trace preservation implies $c\le1\le C$. A fixed finite adaptive experiment
is linear in each channel slot and sends a CP insertion to a positive
unnormalized output, even when the other slots are non-trace-preserving CP
maps. Replace one slot at a time in \eqref{sel51:eq:CPorder}. Multilinearity
and positivity give \eqref{sel51:eq:queryorder}. Classical sums and different
terminal branches preserve the inequalities. A branch with fewer calls can
be padded with discarded calls on invariant inputs; alternatively its bounds
are stronger since $c\le1\le C$. Finally, comparable positive matrices have
the same kernel.
\end{proof}

\begin{theorem}[Every exact pure output inherits the channel-support symmetry]
\label{sel51:thm:queryobstruction}
Every nonzero pure state heralded from finitely many calls to $\Phi$ by the
stated covariant protocol has a density matrix invariant under $G_P$.
For a quantitative form, let
\begin{equation}
 \Psi=\int_{G_P}g\cdot\Phi\,dg,
 \qquad J_\Psi=\int_{G_P}R_gJR_g^*\,dg,
 \label{sel51:eq:twirl}
\end{equation}
where $(g\cdot\Phi)(X)=U_{\rm out}(g)\Phi(U_{\rm in}(g)^*XU_{\rm in}(g))
U_{\rm out}(g)^*$. Choose $c,C$ as in \eqref{sel51:eq:CPorder}.
For any normalized target ray $v$, put
$f_P(v)=\lambda_{\max}(\int_{G_P}R'_gP_vR_g^{\prime *}\,dg)$,
where $R'_g$ is its output representation. Every normalized successful
output $\rho_k$ obeys
\begin{equation}
 \boxed{\|\rho_k-P_v\|_1
 \ge 2(c/C)^k\bigl(1-f_P(v)\bigr).}
 \label{sel51:eq:error}
\end{equation}
The exact pure-output exclusion also holds for an almost-surely finite
protocol with classical termination, even if its call count is unbounded.
No uniform positive error bound over unbounded call counts is asserted.
\end{theorem}
\begin{proof}
Every integrand of \eqref{sel51:eq:twirl} has support $P$, so $J_\Psi$ has
that same support. The channel $\Psi$ is $G_P$-covariant. With invariant
inputs and covariant controls, $Y_k$ is $G_P$-invariant. By
Lemma~\ref{sel51:lem:comparison}, a nonzero pure $X_k$ forces $Y_k$ onto the
same ray; that ray must be invariant.

For a general $X_k\ne0$, normalize $Y_k$ to $\sigma_k$. Taking traces of
\eqref{sel51:eq:queryorder} and its lower inequality gives
$\rho_k\succeq(c/C)^k\sigma_k$. Since $\sigma_k$ is invariant,
$\Tr(P_v\sigma_k)\le f_P(v)$. Testing $I-P_v$ and using the trace-zero
Helstrom inequality yields \eqref{sel51:eq:error}. Finally, a positive sum
of nonzero terminal states can be rank one only when every summand is on
that ray. Apply the finite-depth conclusion to each classical terminal
count; at least one finite count has positive success mass.
\end{proof}

This theorem uses the stabilizer of the \emph{support}, not the covariance
group of the channel density. A channel may have the desired covariance
but a larger support stabilizer. The latter can obstruct pure extraction
at every finite query depth.

\begin{corollary}[Complete finite-query criterion on the determinant-character branch]
\label{sel51:cor:queryperiod}
Suppose $\Ran J$ lies in the subspace fixed pointwise by $H$ of the Choi
representation. This subspace decomposes into integer powers $\chi^n$,
with projectors $\Pi_n$. Put
\begin{equation}
 g(P)=\gcd\{|n-m|:\Pi_nP\Pi_m\ne0,\ n\ne m\},
 \qquad \gcd\varnothing=0.
 \label{sel51:eq:period}
\end{equation}
Assume the protocol additionally admits an invariant maximally entangled
input with a retained conjugate input representation, independent channel
calls, and the invariant effects of Theorem~\ref{sel49:thm:main}.
Then a primitive exact balanced target ray is attainable with positive
success at some finite query count if and only if $g(P)=1$.
Adaptive querying does not enlarge this existence region beyond independent
parallel Choi preparations, although their success probabilities may differ.
\end{corollary}
\begin{proof}
Since $\chi$ maps $G_0$ onto the circle with kernel $H$, the support
stabilizer inside this character representation is $G_0$ for $g(P)=0$,
$H$ for $g(P)=1$, and $\{g:\chi(g)^{g(P)}=1\}$ otherwise. The target
has stabilizer exactly $H$, so necessity follows from
Theorem~\ref{sel51:thm:queryobstruction}.
For sufficiency prepare
$|I\rangle\!\rangle/\sqrt{\dim H_{\rm in}}$ on
$H_{\rm in}\otimes\overline{H_{\rm in}}$. It is invariant. One call,
with the conjugate reference retained, produces $J/\dim H_{\rm in}$.
The inherited V49 theorem applies to this actual prepared state. Its
exterior-projection and finite charge-sum construction give positive success
from finitely many independent calls. The required operations are invariant;
no source environment is accessed. This is a preparation protocol, not a
claim that the Choi state simulates the channel on arbitrary inputs.
\end{proof}

The conjugate reference and its entangled preparation are explicit access
requirements, not automatic consequences of tomography. The negative
Theorem~\ref{sel51:thm:queryobstruction} does not require this extra access.

\section[The unchanged source has two different access restrictions]{A two-level obstruction and one existing neutral direction}
\label{sel51:source}

Keep V19's full native space and all six coefficients at its exact point.
The original vectors $n,\omega,u_0$ are mutually orthogonal;
$n,u_0$ are neutral, $\omega$ has character $\chi$, and all helper maps
$C_\mu$ annihilate their span. The span
$H_3=\Span\{n,\omega,u_0\}$ is invariant under every original coefficient.
The forward channel restricted to inputs supported there is normalized and
stays there. It is exactly
\begin{equation}
 \Phi_3=\tfrac13\id+\tfrac23\operatorname{Ad}_{U_3},\qquad
 U_3=\begin{pmatrix}a&h&0\\h&a&0\\0&0&1\end{pmatrix},
 \quad a=\frac{1-i}{2},\quad h=-\frac{1+i}{2}.
 \label{sel51:eq:U3}
\end{equation}
Indeed each original restriction is $I_3/\sqrt{18}\pm U_3/3$, three of
each sign; summing their CP maps gives \eqref{sel51:eq:U3}. No quantum
coefficient is changed. $H_3$ is a prepared invariant sector, not a replacement
for the complete $14^5$-dimensional native carrier. It is not generally
reducing for the coefficient algebra: $C_\mu^*u_0=u_{\mu-1}$ leaves it.
No coefficient-adjoint descent is used here.

\begin{theorem}[The forgotten two-level source cannot supply the primitive ray]
\label{sel51:thm:twolevel}
Restrict both input and output to $H_2=\Span\{n,\omega\}$, and forget all
original outcomes. The resulting channel $\Phi_2$ has covariance group $H$,
but its Choi support has stabilizer
\begin{equation}
 G_{P_2}=\{g:\chi(g)^2=1\}.
 \label{sel51:eq:periodtwo}
\end{equation}
No finite covariant adaptive query protocol using only this channel can
herald a primitive pure target. For at most $k$ calls the normalized
balanced-target trace error is at least
\begin{equation}
 \boxed{(3-2\sqrt2)^k.}
 \label{sel51:eq:twolevelerror}
\end{equation}
In contrast, the support stabilizer of the unchanged restriction $\Phi_3$
is exactly $H$.
\end{theorem}
\begin{proof}
For $\Phi_2$, its Kraus span is $\Span\{I_2,X\}$. In the orthonormal
coefficient basis $I_2/\sqrt2,X/\sqrt2$, the Choi matrix is
\begin{equation}
 J_2=\frac23\begin{pmatrix}2&i\\-i&1\end{pmatrix}.
 \label{sel51:eq:J2}
\end{equation}
Conjugation by $\diag(1,\chi)$ preserves the traceless line $\C X$
exactly at $\chi=\pm1$, proving \eqref{sel51:eq:periodtwo}.
The density \eqref{sel51:eq:J2} is not invariant at $\chi=-1$, so its
channel covariance is $H$. Twirling over that sign yields
$J_{\Psi_2}=\diag(4/3,2/3)$. The relative eigenvalues are
$1\pm1/\sqrt2$, whence $c/C=3-2\sqrt2$.
The balanced target twirl over this extra sign has largest eigenvalue $1/2$;
Theorem~\ref{sel51:thm:queryobstruction} proves the claims.

For $\Phi_3$, the Kraus span is $\Span\{I_3,U_3\}$. The neutral diagonal
part $\diag(a,a,1)$ of $U_3$ is not scalar. If the span were preserved at
some $\chi$, its transformed $U_3$ would be $xI_3+yU_3$. Comparing the
neutral diagonals forces $x=0,y=1$. Comparing either nonzero charged
entry then forces $\chi=1$. The whole span is fixed at $H$, proving the
last assertion.
\end{proof}

The exclusion already allows arbitrary invariant ancillary memory and
adaptive controls, but only calls to the restricted two-level primitive.
The next construction accesses $u_0$ through the original full source, and
therefore lies outside that restriction. It does not regain an erased
record. This is a change of admitted preparation, not a change of the source.

\section[One invariant preparation gives record-free success]{Exact extraction using an existing neutral helper}
\label{sel51:neutral}

Choose the explicitly invariant pure preparation and its orthogonal neutral
partner
\begin{equation}
 \psi_z=\sqrt{1-z}\,n+\sqrt z\,u_0,\qquad
 e_z=\sqrt z\,n-\sqrt{1-z}\,u_0,
 \qquad 0<z<1.
 \label{sel51:eq:psi}
\end{equation}
This requires coherent preparation within the indicated neutral sector.
Let $Q_z=P_{H_3}-P_{\psi_z}$. It is invariant, and can be included in one
final invariant effect rather than separately executed.

\begin{lemma}[An invariant contamination must be excluded on a pure asymmetric branch]
\label{sel51:lem:kill}
Let $\rho=tP_\psi+\rho_1$, $t>0$, with invariant $P_\psi$. In the fixed
steering protocol \eqref{sel51:eq:target}, any invariant effect $F$ producing
a nonzero multiple of a pure noninvariant target must annihilate
$|0\rangle_A\otimes\psi$ and $|1\rangle_A\otimes\psi$.
It therefore gives the same output from $\rho$ as from
$(I-P_\psi)\rho(I-P_\psi)$.
\end{lemma}
\begin{proof}
Positive summands of a rank-one output have that same ray or are zero.
The summand from $P_\psi$ is invariant, since both the source and effect
are invariant. It cannot have the noninvariant target ray, so it is zero.
Its trace equals one half the sum of the two diagonal expectations of $F$
on the displayed vectors. Both are nonnegative. Each is zero, and positivity
of $F$ then forces both vectors into its kernel. Compression of $F$ proves
the final statement.
\end{proof}

\begin{theorem}[Exact one-query record-free capacity on the neutral probe family]
\label{sel51:thm:neutral}
For the unchanged channel \eqref{sel51:eq:U3} prepared in
$P_{\psi_z}$, forgetting every original outcome, the invariant filtered
reference is pure. Its unnormalized charge populations are
\begin{equation}
 A_0=\frac{z(1-z)}3,\qquad A_1=\frac{1-z}3,
 \qquad s_z=A_0+A_1=\frac{1-z^2}3.
 \label{sel51:eq:neutralweights}
\end{equation}
The largest exact balanced-target success in the fixed one-reference
invariant readout protocol is
\begin{equation}
 \boxed{\alpha(z)=\frac{z(1-z)}{3(1+z)}.}
 \label{sel51:eq:alphaz}
\end{equation}
This is optimal for that fixed preparation even without requiring the
preliminary projection $Q_z$. Within the family \eqref{sel51:eq:psi}, its
maximum is
\begin{equation}
 \boxed{z_* =\sqrt2-1,\qquad
 \alpha(z_*)=\frac{3-2\sqrt2}{3}.}
 \label{sel51:eq:neutralopt}
\end{equation}
At the simple preparation $z=1/2$, the exact success is $1/18$ and the
selected packet is
\begin{equation}
 \boxed{J_{\rm succ}=\frac1{36}
 \begin{pmatrix}1&1\\1&1\end{pmatrix},\qquad h=\frac1{1296}.}
 \label{sel51:eq:one18}
\end{equation}
\end{theorem}
\begin{proof}
The state before filtering is
$\rho_z=P_{\psi_z}/3+(2/3)P_{U_3\psi_z}$.
Since $a-1=h$,
\[
 Q_zU_3\psi_z=h\sqrt{1-z}(\sqrt z\,e_z+\omega).
\]
Thus $Q_z\rho_zQ_z$ is rank one, with populations
\eqref{sel51:eq:neutralweights}. The inherited pure-reference formula gives
$A_0A_1/(A_0+A_1)$, proving attainability and \eqref{sel51:eq:alphaz}.
Lemma~\ref{sel51:lem:kill} forces every successful effect to exclude
$P_{\psi_z}$; after that exclusion the reference is exactly the computed
pure state. The same formula is therefore an upper bound for arbitrary
admitted invariant effects on this fixed preparation. Differentiation gives
$\alpha'(z)=(1-2z-z^2)/(3(1+z)^2)$, proving the unique maximum.

At $z=1/2$, write $e=(n-u_0)/\sqrt2$. Directly,
\begin{equation}
 \rho_{1/2}=\begin{pmatrix}
 1/3&i/6&1/3-i/6\\
 -i/6&1/6&-1/6-i/6\\
 1/3+i/6&-1/6+i/6&1/2
 \end{pmatrix}_{n,\omega,u_0},\qquad
 Q\rho_{1/2}Q=\frac1{12}
 \begin{pmatrix}1&\sqrt2\\\sqrt2&2\end{pmatrix}_{e,\omega}.
 \label{sel51:eq:referenceexplicit}
\end{equation}
The reference-selection probability is $1/4$ and its normalized state is
$(e+\sqrt2\omega)/\sqrt3$. The single final effect
\begin{equation}
 F=|f\rangle\langle f|,\qquad
 f=\frac{|0\rangle_A\omega+\sqrt2|1\rangle_Ae}{\sqrt3}
 \label{sel51:eq:effect18}
\end{equation}
lies in one total-charge sector, hence is invariant. It already vanishes
on both flag copies of $\psi_{1/2}$. Contraction of
$|\Xi\rangle\langle\Xi|\otimes\rho_{1/2}$ gives
\eqref{sel51:eq:one18}. Its complement $I-F$ is positive and returns the
remaining part of the target marginal $I_2/2$.
\end{proof}

No reference ancilla is retained during this preparation, and the original
six source labels are unavailable. The known neutral helper is used as an
input component, not as an extra asymmetric resource. Its \emph{coherent
preparability} and the invariant effect \eqref{sel51:eq:effect18} remain
explicit access requirements. This is not claimed to maximize success over
all invariant preparations on the complete native space, multiple queries,
or arbitrary covariant conversion channels.

The unchanged idle/accepted opportunity map is
$(1-\vartheta)\id+\vartheta\Phi_3$. The same effect kills its idle output
$P_{\psi_z}$ as well. The raw opportunity success is therefore
$\vartheta\alpha(z)$, including $\vartheta/18$ at $z=1/2$.
Neither the acceptance record nor a new waiting-time normalization is needed
for that probability. No elapsed-time estimate is inferred from it.

\section[Invariant input correlations give a second exact benchmark]{The Choi preparation improves success without environment access}
\label{sel51:bell}

For comparison, explicitly admit a retained conjugate three-dimensional
reference, and prepare the invariant vector
\begin{equation}
 \Omega_{x,y,z}=\sqrt x\,n\bar n+\sqrt y\,\omega\bar\omega
                  +\sqrt z\,u_0\bar u_0,
 \quad x,y,z\ge0,\quad x+y+z=1.
 \label{sel51:eq:Bell}
\end{equation}
All terms have total character one. This entangled preparation is not
asserted to follow from classical tomography or from an ordinary clock.
Let $\xi=(U_3\otimes I)\Omega_{x,y,z}$. The one-query output is
$P_\Omega/3+2P_\xi/3$. The invariant projection $I-P_\Omega$ leaves a pure
unnormalized reference, with charge masses
\begin{equation}
 A_{-1}=y/3,\qquad A_0=z(1-z)/3,\qquad A_{+1}=x/3.
 \label{sel51:eq:Bellweights}
\end{equation}

\begin{proposition}[Exact optimum in the correlated-probe family]
\label{sel51:prop:Bell}
For the preparation \eqref{sel51:eq:Bell}, the best exact balanced-target
success in the same invariant readout protocol is
\begin{equation}
 \beta(x,y,z)=\frac13\left(
 \frac{z(1-z)y}{z(1-z)+y}+
 \frac{z(1-z)x}{z(1-z)+x}\right),
 \label{sel51:eq:beta}
\end{equation}
with zero terms omitted. Over this preparation family,
\begin{equation}
 \boxed{\beta_{\max}=\frac{2-\sqrt3}{3},\quad
 z=\frac{\sqrt3-1}{2},\quad x=y=\frac{3-\sqrt3}{4}.}
 \label{sel51:eq:Bellopt}
\end{equation}
For $x=y=z=1/3$, the reference projection succeeds with $8/27$, and the
complete target success is $4/45$. Its selected packet is
$2\begin{psmallmatrix}1&1\\1&1\end{psmallmatrix}/45$.
\end{proposition}
\begin{proof}
One has
$\langle\Omega,\xi\rangle=(1+z)/2-i(1-z)/2$.
The two charged components have squared norms $x/2,y/2$; after removing
$\Omega$, the neutral squared norm is $z(1-z)/2$. Multiplication by $2/3$
gives \eqref{sel51:eq:Bellweights}. Lemma~\ref{sel51:lem:kill} again
forces every efficient asymmetric success effect to annihilate both flag
copies of $\Omega$. The remaining pure reference has two adjacent pairs.
Their independent total-charge blocks give exactly the sum
\eqref{sel51:eq:beta}, by the inherited pure-reference capacity.
For fixed $z\in(0,1)$, concavity of $t\mapsto at/(a+t)$ with
$a=z(1-z)>0$ gives the maximum at $x=y=(1-z)/2$.
The resulting function is $2z(1-z)/(3(2z+1))$; differentiation proves
\eqref{sel51:eq:Bellopt}. Substitution gives $4/45$ for the uniform probe.
The pure conditional reference there has populations $(3/8,1/4,3/8)$,
whose conditional capacity is $3/10$, and $(8/27)(3/10)=4/45$.
\end{proof}

All these probabilities include the success of removing the invariant
contamination. They are not conditional probabilities with an omitted
preparation cost. The construction consumes a retained quantum reference
and is not a refinement of an erased classical source record. It leaves
the original source channel unchanged.

\section[The channel-level noise obstruction survives every finite query]{Exact symmetry can persist while exact pure extraction fails}
\label{sel51:noise}

\begin{proposition}[A sharp CP comparison for a specified noisy restriction]
\label{sel51:prop:noise}
On the identified $H_3$, consider the declared comparison family
\begin{equation}
 \Phi_{3,\eta}=(1-\eta)\Phi_3+\eta\mathcal T_3,
 \qquad\mathcal T_3(X)=\Tr(X)I_3/3,\quad 0<\eta\le1.
 \label{sel51:eq:noise}
\end{equation}
For $0<\eta<1$, its channel covariance inside $G_0$ remains exactly $H$,
but its Choi support is the entire coefficient space. No finite covariant
adaptive query protocol on this restricted primitive can herald an exact
primitive pure target. For at most $k$ calls,
\begin{equation}
 \boxed{\|\rho_k-P_v\|_1\ge
 \left(\frac{\eta}{8-7\eta}\right)^k.}
 \label{sel51:eq:noisebound}
\end{equation}
This does not rule out approximation with increasing query number.
\end{proposition}
\begin{proof}
The nonzero eigenvalues of $J_3$ are $8/3$ and $1/3$. Indeed its two
weighted Kraus vectors have Gram
$\begin{psmallmatrix}1&\sqrt2(2-i)/3\\
\sqrt2(2+i)/3&2\end{psmallmatrix}$, with trace three and determinant $8/9$.
Thus $J_{3,\eta}$ has smallest eigenvalue $\eta/3$ and largest
$(8-7\eta)/3$. Relative to $J_{\mathcal T_3}=I_9/3$,
\[
 \eta\mathcal T_3\preceq_{\rm CP}\Phi_{3,\eta}
       \preceq_{\rm CP}(8-7\eta)\mathcal T_3.
\]
The comparison is $G_0$-covariant. The balanced target has maximum overlap
$1/2$ with an invariant density. Apply
Lemma~\ref{sel51:lem:comparison} and the proof of
Theorem~\ref{sel51:thm:queryobstruction}. The positive bound also proves
the exact exclusion. For $\eta<1$, covariance of the mixture is equivalent
to covariance of $\Phi_3$, since $\mathcal T_3$ is invariant and can be
subtracted from the equality. The latter is $H$, for example because its
output on invariant $n$ has nonzero $n/\omega$ coherence.
\end{proof}

The noise in \eqref{sel51:eq:noise} is a new, explicitly distinguished
comparison channel, not noise assigned to the historical Store or silently
added to the benchmark. The all-query conclusion here confines each source
call to $H_3$. For a noisy \emph{full native} channel
$(1-\eta)\Phi_{\rm native}+\eta\mathcal T_N$, the same theorem gives a
full-carrier exclusion, with CP constants $c=\eta$ and
$C=\eta+(1-\eta)N\|J_{\rm native}\|_{\op}$.
Access to a different noise-free source port is not included in either
exclusion. Nor do these arguments allow an arbitrary enlargement of the
invariant operation bank without recording that assumption.

\section{Proof status and the source boundary after V51}
\label{sel51:status}

\paragraph{Proved here.}
A channel-support symmetry propagates through every finite covariant
adaptive source query and gives \eqref{sel51:eq:error}. Combined with the
inherited reference theorem it is a complete finite-query existence test
on the specified determinant-character branch with invariant Choi-preparation
access. The original two-level forgotten source has a period-two obstruction
and explicit error bound. Its unchanged three-level restriction, using the
already present neutral helper, admits the record-free preparations and
optimal fixed-probe readouts computed above. The declared noisy channel has
an all-finite-query obstruction while retaining its exact covariance group.

\paragraph{Inherited, not newly established.}
The native tensor carrier, its invariant and determinant vectors, its
symmetry-breaking action, neutral helpers and six coefficients are the V19
constructed source. The V47--V49 reference conversion theorems and V50's
$4/21$ bit-retaining preparation are inputs. V50's forgotten fixed-input
state still has zero finite-copy capacity: changing the input to
\eqref{sel51:eq:psi} does not reverse that old preparation or recover its bit.
The general relation between CP order and finite-query resource obstruction
belongs to the established purification framework cited above.

\paragraph{Unresolved historical inputs.}
The benchmark already contains determinant-sensitive $H_s$; no invariant
preparation here derives that initial resource. Historical admission of a
coherent neutral preparation, the same-cut charge identification, invariant
joint effects, and any conjugate-reference port must be independently shown.
They are not consequences of the commutant's dimension or of classical frame
spanning. The V110 directed odd incidence and its original word remain
unpopulated. The present protocols produce selected \emph{coefficient-state}
rays in the fixed comparison source, not a new physical Yukawa operator.

The productive next audit therefore distinguishes the support of a specified
prepared state from the support of the primitive channel available for new
queries. Requiring retention of a particular preparation bit is sufficient
on V50's chosen experiment, but is not necessary once the actual neutral
preparation bank is enlarged in the explicitly stated way. No stationarity,
continuous-time embeddability, or full SM--Einstein continuum is imposed as a
new condition of the finite duality.

\paragraph{Verification boundary.}
The accompanying script uses exact rational/algebraic matrices, independently
reconstructs all six original seven-dimensional working-block coefficients,
and checks their three-level restrictions, the positive comparison spectra,
the invariant preparation/readout effects, and the stated optimizations.
The finite-query and all-copy assertions rest on the written proofs, not
on testing a finite list of protocols. The unchanged V50 checker is rerun;
this is not an independent audit of all fifty preceding layers. No physical
measurement, proof-assistant formalization, or independent peer review is
claimed.

\begingroup
\renewcommand{\refname}{Primary-source context for V51}
\begin{thebibliography}{9}
\small\raggedright
\bibitem{sel51-FangLiu}
K.~Fang and Z.-W.~Liu,
\emph{No-Go Theorems for Quantum Resource Purification: New Approach and
Channel Theory}, PRX Quantum \textbf{3} (2022), 010337;
\url{https://arxiv.org/abs/2010.11822}.
\bibitem{sel51-Marvian}
I.~Marvian and R.~W.~Spekkens,
\emph{The theory of manipulations of pure state asymmetry: basic tools and
equivalence classes of states under symmetric operations},
New J. Phys. \textbf{15} (2013), 033001;
\url{https://arxiv.org/abs/1104.0018}.
\bibitem{sel51-Modes}
I.~Marvian and R.~W.~Spekkens,
\emph{Modes of asymmetry: the application of harmonic analysis to symmetric
quantum dynamics and quantum reference frames},
Phys. Rev. A \textbf{90} (2014), 062110;
\url{https://arxiv.org/abs/1312.0680}.
\end{thebibliography}
\endgroup
\addtocontents{toc}{\protect\endgroup}
% END OF SOURCE-SELECTION V51 DEVELOPMENT BLOCK.
% Preserve V1--V51 and append subsequent work before the final terminator.


\clearpage
\hypersetup{pdftitle={Historical Standard-Model Source Selection: V1--V52}}
\begin{center}
{\Large\bfseries Source-selection continuation V52}\\[.4em]
{\large Exact one-query preparation limits on the full native source\\
Minimal invariant reference and the preparation-support boundary}
\end{center}
\addtocontents{toc}{\protect\begingroup}
\addtocontents{toc}{\protect\renewcommand{\protect\numberline}{\protect\makebox[2.1em][l]}}

\section[Optimize the original preparation class, not another reference]{The full native one-query problem}
\label{sel52:context}

V51 constructs record-free exact rays on an invariant three-dimensional
forward sector of the original V19 source. Its probabilities are optimized
within two specified preparation families, not over the full native input
space or arbitrary invariant ancillary probes. The present layer closes
that finite optimization problem in the same steering protocol. It proves
that the unused native directions and the four damped helper inputs cannot
improve either optimum. A two-dimensional conjugate reference attains the
larger value; the three-dimensional reference in V51 is unnecessary.
The smaller value is the optimum over all invariant preparations without a
retained quantum reference, not merely the displayed neutral pure family.

This changes an access question, not the historical source. The action,
coefficients, helper directions, carrier and acceptance factor are those of
V19 and V51. No source matrix is chosen anew. Exact preparation of an
invariant state is still a physical capability: it does not follow from
knowing its coordinates, from algebra membership, or from a spanning
statistical frame. The inherited reference capacities of V47--V49 are
used rather than claimed again. Symmetric-operation and channel-purification
frameworks are established \cite{sel52-Marvian,sel52-FangLiu}; the
full-source optimization and access comparisons below have direct proofs.

\paragraph{The unchanged channel on its full carrier.}
Let $\mathcal H=\mathcal H_{14}^{\otimes5}$, $N=14^5$. The orthonormal
vectors $n,u_0,u_1,\ldots,u_4$ are neutral, and $\omega$ has character
$\chi(g)=\det U\det V$ under the declared $G_0=\mathrm U(3)\times\mathrm U(2)$
action. Put
\begin{equation}
 \begin{gathered}
 P_{\rm act}=P_n+P_\omega,\qquad R=\sum_{j=1}^4P_{u_j},
 \qquad C_j=|u_0\rangle\langle u_j|,\\
 q=\frac1{4096},\quad d=\sqrt{1-q}=\frac{\sqrt{4095}}{64},
 \quad D=I+(d-1)R,\\
 U=I-\frac{1+i}{2}|n+\omega\rangle\langle n+\omega|.
 \end{gathered}
 \label{sel52:eq:source}
\end{equation}
The original six-outcome compiler, with its outcomes forgotten, is exactly
\begin{equation}
 \boxed{\Phi(X)=\frac13X+\frac23\left(UDXD U^*
                +q\sum_{j=1}^4 C_jXC_j^*\right).}
 \label{sel52:eq:fullchannel}
\end{equation}
It is normalized since $D^2+qR=I$. The operator $U$ is the identity on
$P_{\rm act}^\perp$, not an unspecified unitary on that complement. Every
native direction is retained in the optimization. The notation $d$ in
\eqref{sel52:eq:source} is a scalar damping factor; the target coefficient
vector below is denoted $d_{\rm tar}$ to avoid ambiguity.

\paragraph{The precise operation class.}
Retain V51's independently prepared normalized invariant target source
\begin{equation}
 |\Xi\rangle=\frac{b|0\rangle_A+d_{\rm tar}|1\rangle_A}{\sqrt2},
 \qquad v=\frac{b+d_{\rm tar}}{\sqrt2},
 \label{sel52:eq:target}
\end{equation}
where $b,d_{\rm tar}$ have characters $1,\chi^{-1}$ and the flag states have
characters $1,\chi$. Their inaccessible markers are identical. Prepare an
invariant density $\rho$ on $\mathcal H\otimes B$, independently of
\eqref{sel52:eq:target}; make one call to $\Phi\otimes\id_B$; discard all
original source labels; and apply an invariant effect $0\le F\le I$ on
$A\otimes\mathcal H\otimes B$. The desired output is $\alpha P_v$ on the
untouched target coefficient space. The reference is consumed. Source-side
invariant postprocessing and classical decisions after the call can be
absorbed into $F$. Classical choices before the call can be included as
neutral flags in $\rho$.

Write $A_{\rm bare}$ for the supremum without a retained quantum reference,
allowing classical preparation tags, and $A_{\rm ref}$ for the supremum over
arbitrary finite $B$, arbitrary unitary representations on $B$, and all
jointly invariant preparations. All success probabilities include every
projection and filter; in particular $\alpha\le1/2$ because the original
target marginal is $I_2/2$. These are optima for this one-call,
independent-target \emph{preparation--steering class}, not for arbitrary
covariant channel supermaps, multiple calls, coherent control of the query,
or operations acting directly on the target output.

\section[A normal form for all invariant probes]{Invariant contamination and the helper penalty}
\label{sel52:normal}

For an invariant pure probe $\Psi\in\mathcal H\otimes B$, let
\begin{equation}
 \begin{gathered}
 x=\|(P_n\otimes I)\Psi\|^2,\quad
 y=\|(P_\omega\otimes I)\Psi\|^2,\quad s=x+y,\quad z=1-s,\\
 r=\|(R\otimes I)\Psi\|^2,\qquad 0\le r\le z.
 \end{gathered}
 \label{sel52:eq:weights}
\end{equation}
A probe transforming by one character is covered by factoring out that
common character; its density and all the support projections used below
remain invariant. Define
\[
 \eta_j=(C_j\otimes I)\Psi,\qquad
 E=\Span\{\Psi,\eta_1,\ldots,\eta_4\},\qquad Q_E=I-P_E,
 \qquad \xi=Q_E(UD\otimes I)\Psi.
\]
All vectors in $E$ have the same total character as $\Psi$.

\begin{lemma}[Exact reduction to three relative charges]
\label{sel52:lem:normal}
Every effect producing a nonzero multiple of $P_v$ must annihilate
$\C^2_A\otimes E$. It has exactly the same target output when the queried
reference is replaced by $2P_\xi/3$. Relative to the probe's character,
$\xi$ has only charges $-1,0,1$. Put $g=\|\xi_0\|^2$. Its unnormalized
charge masses and optimal target success are
\begin{align}
 A_{-1}&=\frac y3,\qquad A_0=\frac{2g}{3},\qquad A_{+1}=\frac x3,
 \label{sel52:eq:masses}\\
 \boxed{\alpha(\Psi)=\frac{2g}{3}
       \left(\frac{x}{2g+x}+\frac{y}{2g+y}\right),}
 \label{sel52:eq:capacity}
\end{align}
with zero terms omitted. If $\varepsilon_d=1-d$, then
\begin{equation}
 \boxed{g\le \frac{s(1-s)}2
                  +r\varepsilon_d\bigl(\varepsilon_d(1-r)-s\bigr).}
 \label{sel52:eq:helperbound}
\end{equation}
For $r=0$ equality holds in \eqref{sel52:eq:helperbound}.
\end{lemma}
\begin{proof}
The query output is
\[
 \tfrac13P_\Psi+\tfrac23P_{(UD\otimes I)\Psi}
                  +\tfrac{2q}{3}\sum_jP_{\eta_j}.
\]
Each first or last summand is invariant. Its selected target contribution is
positive and invariant, and lies below the desired rank-one noninvariant
output. It must be zero. Its trace is one half the sum of the two positive
diagonal expectations of $F$ on the flag copies of its vector. Positivity
then puts both copies in $\ker F$. This is the inherited
Lemma~\ref{sel51:lem:kill}, applied to every invariant contamination vector.
Compression by $Q_E$ therefore loses no possible exact success.

Write $a=(1-i)/2$ and $h=-(1+i)/2$. The only character-changing terms of
$UD$ are $h|\omega\rangle\langle n|$ and
$h|n\rangle\langle\omega|$. They are orthogonal to $E$, have squared norms
$x/2,y/2$, and give \eqref{sel52:eq:masses}. The neutral component of the
remaining pure reference has squared norm $g$. The inherited
pure-reference formula, or its two independent total-charge-block
Cauchy--Schwarz proofs, gives
$A_0A_{-1}/(A_0+A_{-1})+A_0A_{+1}/(A_0+A_{+1})$.
This is \eqref{sel52:eq:capacity}; the two attaining rank-one effect
projectors occupy orthogonal total-charge blocks and sum to an effect.

For completeness retain the helper inputs before their further removal.
Let $\Psi_s=P_{\rm act}\Psi$, $\Psi_r=R\Psi$ and
$\Psi_z=(I-P_{\rm act}-R)\Psi$, with tensor identities suppressed.
The neutral component of $(UD\otimes I)\Psi$ is
$a\Psi_s+d\Psi_r+\Psi_z$. Projecting only off $\Psi$ gives squared norm
\begin{align*}
 g_{\rm pre}
 &=s(z-r)|a-1|^2+sr|a-d|^2+r(z-r)|1-d|^2\\
 &=\frac{sz}{2}-sr d\varepsilon_d+r(z-r)\varepsilon_d^2
  =\frac{sz}{2}+r\varepsilon_d
                         \bigl(\varepsilon_d(1-r)-s\bigr).
\end{align*}
This is the elementary variance identity for three orthogonal vectors with
weights $s,r,z-r$ summing to one. Removing the other vectors of $E$ can only
decrease the norm, so $g\le g_{\rm pre}$. If $r=0$, all $\eta_j$ vanish,
$E=\C\Psi$, and equality follows.
\end{proof}

The term in \eqref{sel52:eq:helperbound} retains the original damping and
all four original jumps. Treating the helpers as absent from an arbitrary
input would not justify the bound. The result also handles every unaffected
native representation sector, not just the seven-dimensional working block.

\section[Global optima on the unchanged full native carrier]{The preparation families already attain the full one-query limits}
\label{sel52:optima}

\begin{theorem}[Exact unassisted and invariant-reference optima]
\label{sel52:thm:global}
For \eqref{sel52:eq:fullchannel} and the precisely stated steering class,
\begin{equation}
 \boxed{A_{\rm bare}=\frac{3-2\sqrt2}{3},\qquad
        A_{\rm ref}=\frac{2-\sqrt3}{3}.}
 \label{sel52:eq:global}
\end{equation}
These are maxima over the full original native space. In the second class
arbitrary finite ancillary representations and mixed invariant probes are
allowed. The same values bound arbitrary finite classical preparation
ensembles with their tags retained. An optimal pure probe does not populate
the four damped helper inputs.
\end{theorem}
\begin{proof}
An invariant mixed source--reference density has an invariant purification:
if its representation is $V$, the vector $|\sqrt\rho\rangle\!\rangle$ in
its space tensored with the conjugate space is fixed by $V\otimes\bar V$.
Retaining the extra reference can only enlarge the attainable effect bank;
the old effect tensored with identity reproduces the old output. Thus an
upper bound for arbitrary invariant pure probes bounds the entire assisted
class. This is an upper-bound reduction, not a claim that the purifying
register was present in a given experiment.

At the original parameter, $d>63/64$, hence $0<\varepsilon_d<1/64$.
The capacity in \eqref{sel52:eq:capacity} is increasing in $g$ and is at
most $(x+y)/3=s/3$. If $s<\varepsilon_d$, it is less than $1/192$;
the already attainable value $1/18$ excludes this region from either global
maximum. If $s\ge\varepsilon_d$, \eqref{sel52:eq:helperbound} gives
$g\le s(1-s)/2$. For fixed $s$ and $g>0$, strict concavity of
$t/(2g+t)$ gives the maximum at $x=y=s/2$. Consequently
\[
 \alpha(\Psi)\le\frac{2z(1-z)}{3(1+2z)},\qquad z=1-s.
\]
Its unique maximizer is $z=(\sqrt3-1)/2$, with value
$(2-\sqrt3)/3$. The correlated probe in
Proposition~\ref{sel51:prop:Bell} attains it without helper-input weight,
so the assisted upper bound is exact.

For the unassisted class, decompose an invariant input density into its
representation blocks. Outside the trivial and $\chi$ isotypic sectors,
$U=D=I$ and $C_j=0$, so the channel acts identically and that invariant
part cannot contribute to a pure asymmetric success. In each of the two
character sectors an invariant density is a convex combination of pure
vectors transforming by that character. No such vector can have both
$x>0$ and $y>0$. For each, \eqref{sel52:eq:capacity} and the same helper
bound therefore give
\[
 \alpha\le\frac{z(1-z)}{3(1+z)}.
\]
This is maximized at $z=\sqrt2-1$ with value $(3-2\sqrt2)/3$.
The neutral probe $\sqrt{1-z}\,n+\sqrt z\,u_0$ attains the value by
Theorem~\ref{sel51:thm:neutral}. For any mixed input and any one fixed
successful effect, positivity forces every nonzero component contribution
onto the same target ray. Summing their weighted individual upper bounds
therefore gives the same bound without assuming an executable spectral
refinement. Retaining classical tags cannot exceed the largest component
bound either.

At either maximum, $s>\varepsilon_d$. The helper correction is then
strictly negative when $r>0$. The capacity is strictly increasing in its
nonzero neutral mass, so a maximizing pure probe has $r=0$.
\end{proof}

For the assisted pure optimum the necessary populations are
\begin{equation}
 x=y=\frac{3-\sqrt3}{4},\qquad z=\frac{\sqrt3-1}{2}.
 \label{sel52:eq:bestweights}
\end{equation}
The proof does not uniquely name the spectator vector: many actual native
neutral directions could represent it if they were preparable. The source
chooses neither a preferred input nor an optimizing readout merely because
the optimum is known.

\section[Two reference levels suffice; neutral multiplicity does not]{The minimal symmetry-carrying ancillary probe}
\label{sel52:reference}

\begin{theorem}[A two-dimensional invariant probe attains the assisted optimum]
\label{sel52:thm:qubit}
Let $B=\Span\{|0\rangle,|1\rangle\}$ carry characters $1,\chi^{-1}$.
The invariant vector
\begin{equation}
 \boxed{\Psi_{xyz}=
 (\sqrt x\,n+\sqrt z\,u_0)\otimes|0\rangle
       +\sqrt y\,\omega\otimes|1\rangle,\qquad x+y+z=1,}
 \label{sel52:eq:twolevelprobe}
\end{equation}
with \eqref{sel52:eq:bestweights} attains $A_{\rm ref}$.
Two is the smallest ancillary Hilbert dimension attaining that value.
An ancillary register of any dimension on which $G_0$ acts by one scalar
character, in particular a completely neutral register, cannot improve
$A_{\rm bare}$.
\end{theorem}
\begin{proof}
The two terms in \eqref{sel52:eq:twolevelprobe} have total character one;
the state contains no initial asymmetry. It has zero helper-input weight.
Write $s=x+y$, $\Psi_s=\sqrt x\,n0+\sqrt y\,\omega1$ and
$\Psi_z=\sqrt z\,u_00$. For $sz>0$, the three orthonormal relative-charge
vectors are
\[
 e_-=n1,\qquad
 e_0=\frac{z\Psi_s-s\Psi_z}{\sqrt{sz}},\qquad
 e_+=\omega0.
\]
Since $a-1=h$, direct application of the unchanged $U$ gives
\[
 (I-P_\Psi)(U\otimes I)\Psi
       =h\bigl(\sqrt y\,e_-+\sqrt{sz}\,e_0+\sqrt x\,e_+\bigr).
\]
After multiplying by $2/3$, its charge masses are
$A_-=y/3$, $A_0=sz/3$, $A_+=x/3$. Define
\begin{align}
 f_0&=\frac{\sqrt{A_-}\,|0\rangle_Ae_0+
                 \sqrt{A_0}\,|1\rangle_Ae_-}{\sqrt{A_-+A_0}},\\
 f_1&=\frac{\sqrt{A_0}\,|0\rangle_Ae_++
                 \sqrt{A_+}\,|1\rangle_Ae_0}{\sqrt{A_0+A_+}},
 \qquad F=P_{f_0}+P_{f_1}.
 \label{sel52:eq:effect}
\end{align}
The summands have different total charges, so $0\le F\le I$ and $F$ is
invariant. They both annihilate all flag copies of $\Psi$.
Contraction with \eqref{sel52:eq:target} gives
\begin{equation}
 J_F=\frac{\alpha}{2}
 \begin{pmatrix}1&1\\1&1\end{pmatrix},\qquad
 \alpha=\frac{sz}{3}
       \left(\frac{x}{sz+x}+\frac{y}{sz+y}\right).
 \label{sel52:eq:qubitoutput}
\end{equation}
At \eqref{sel52:eq:bestweights}, this is
$\alpha=(2-\sqrt3)/3$. Its unnormalized squared cross is
$(7-4\sqrt3)/36$.

A one-dimensional ancillary representation is scalar and adds no charged
correlation. More generally, with a scalar-character ancillary action,
an invariant state decomposes into source isotypic sectors as in the
unassisted proof. Pure probes in the two affected sectors still have
$xy=0$, irrespective of neutral ancillary multiplicities. The same
unassisted bound follows. Since $A_{\rm ref}>A_{\rm bare}$, an ancillary
dimension smaller than two cannot attain $A_{\rm ref}$.
\end{proof}

Discarding the ancillary qubit before the readout changes this exact
capability. The source marginal of \eqref{sel52:eq:twolevelprobe} is
\[
 \rho_{\mathcal H}=
 \begin{pmatrix}x&0&\sqrt{xz}\\0&y&0\\\sqrt{xz}&0&z\end{pmatrix}_{n,\omega,u_0}.
\]
For $xyz>0$ its forgotten-source output obeys
\begin{equation}
 \det\Phi_3(\rho_{\mathcal H})=\frac{yz(x+y)}9>0.
 \label{sel52:eq:discardedreference}
\end{equation}
At the assisted optimum this is $(3\sqrt3-5)/24>0$. Its support is the
entire invariant $H_3$, so no finite number of independent copies of that
reduced output can supply an exact primitive ray. The joint output with the
qubit retained attains $A_{\rm ref}$ instead. This is a statement about the
specified prepared output, not an obstruction to new source queries.

The source marginal of \eqref{sel52:eq:twolevelprobe} and the ancillary
marginal are both invariant. The probe is nevertheless entangled when
$x+z>0$ and $y>0$. The useful extra capability is a coherent invariant
correlation with a \emph{compensating} character sector, not an arbitrarily
large neutral register. No claim is made that every separable invariant
probe is equivalent to a neutral register, or that every optimal probe has
been classified. The conjugate character and the joint preparation are
physical admissions, not consequences of their two-dimensional cost.

The entire protocol uses one effect, with complement $I-F$; its failure
output plus \eqref{sel52:eq:qubitoutput} is $I_2/2$. There is no separate
postselection cost omitted in either global value. On the original raw
opportunity channel $(1-\vartheta)\id+\vartheta\Phi$, the same effect
annihilates the idle contribution. Both exact maxima therefore scale as
\begin{equation}
 \boxed{A_{\rm bare}^{\rm opp}=\vartheta A_{\rm bare},\qquad
        A_{\rm ref}^{\rm opp}=\vartheta A_{\rm ref}.}
 \label{sel52:eq:opportunity}
\end{equation}
An event count, a fixed opportunity, and elapsed physical time are not
identified by this equality.

\section[Which three-level preparations actually pass?]{A complete unassisted preparation-support test}
\label{sel52:preptest}

Return to the original forward-invariant sector
$H_3=\Span\{n,\omega,u_0\}$. Every invariant density has the form
\begin{equation}
 \rho=\begin{pmatrix}a&0&c\\0&b&0\\\bar c&0&z\end{pmatrix},\qquad
 a,b,z\ge0,\quad a+b+z=1,\quad |c|^2\le az.
 \label{sel52:eq:input}
\end{equation}
The original source outcomes are forgotten and the prepared output
$\sigma=\Phi_3(\rho)$ is used as a reference.

\begin{theorem}[All exact-capable invariant inputs on the three-level port]
\label{sel52:thm:prepface}
In the inherited invariant steering protocol, the following are equivalent:
\begin{enumerate}[label=\textup{(\roman*)},leftmargin=2.5em]
\item a finite number of independent copies of $\sigma$ permits a nonzero
      exact primitive target ray;
\item one copy of $\sigma$ permits such a ray;
\item $b=0$ and $|c|^2=az>0$.
\end{enumerate}
In the positive case $\rho$ is a pure coherent neutral state, and the exact
one-copy balanced-target capacity is
\begin{equation}
 \frac{z(1-z)}{3(1+z)}.
 \label{sel52:eq:prepfacecapacity}
\end{equation}
Every mixed invariant density on $H_3$ gives an output with invariant support
and therefore zero exact primitive-ray capacity at every finite copy number.
These statements concern copies of this prepared output, not further calls
to $\Phi_3$ with a different preparation policy.
\end{theorem}
\begin{proof}
Positivity of the two terms in $\Phi_3=\id/3+2\operatorname{Ad}_U/3$ gives
$\supp\sigma=\Span(\supp\rho,U\supp\rho)$. If the neutral block in
\eqref{sel52:eq:input} has rank two, its support contains both $n,u_0$;
adding $Un$ fills $H_3$. If that block has rank one and $b>0$, write its
ray as $\psi=\sqrt a\,n+e^{i\phi}\sqrt z\,u_0$ after normalization.
For $z>0$, $\Span\{\psi,\omega,U\omega\}$ fills $H_3$.
For $z=0$ the support is the invariant $n/\omega$ plane. The remaining
rank-one boundary inputs $P_n,P_\omega,P_{u_0}$ have output support
respectively the same invariant plane, that plane, and $\C u_0$.
Every listed support is $G_0$ invariant. The inherited support theorem
excludes an asymmetric pure success from any finite number of copies.

The only remaining possibility is $b=0$, $|c|^2=az>0$. It is precisely
$\rho=P_{\psi_z}$ up to a phase on $u_0$. That phase is an invariant
unitary commuting with $U$; transporting the readout preserves its optimum.
The unchanged positive construction of V51 gives
\eqref{sel52:eq:prepfacecapacity}, so one copy suffices. This proves both
necessity and sufficiency without assigning a purity condition to other
reference protocols or to the finite duality.
\end{proof}

\begin{proposition}[Neutral preparation noise changes support, not the source]
\label{sel52:prop:prepnoise}
For $0<z<1$ and $0\le\delta\le1$, dephase only the neutral preparation:
\[
 \rho_{z,\delta}=
 \begin{pmatrix}
 1-z&0&(1-\delta)\sqrt{z(1-z)}\\
 0&0&0\\
 (1-\delta)\sqrt{z(1-z)}&0&z
 \end{pmatrix}.
\]
All these inputs are invariant and have the same populations. The unchanged
channel gives
\begin{equation}
 \boxed{\det\Phi_3(\rho_{z,\delta})
       =\frac{\delta(2-\delta)z(1-z)^2}{9}.}
 \label{sel52:eq:noisedet}
\end{equation}
Thus every $\delta>0$ makes the prepared output faithful and excludes exact
primitive-ray extraction from every finite number of its independent copies.
For $z=1/2$, the \emph{unchanged} effect
\eqref{sel51:eq:effect18} returns
\begin{equation}
 J_\delta=\frac1{36}
 \begin{pmatrix}
 1&1-\delta-i\delta\\
 1-\delta+i\delta&1+4\delta
 \end{pmatrix},\quad
 \Tr J_\delta=\frac{1+2\delta}{18},\quad
 \det J_\delta=\frac{\delta(3-\delta)}{648}.
 \label{sel52:eq:noisyoutput}
\end{equation}
Its normalized distance from $P_v$ is exactly
\begin{equation}
 \boxed{\left\|\frac{J_\delta}{\Tr J_\delta}-P_v\right\|_1
             =\frac{\sqrt{14}\,\delta}{1+2\delta}.}
 \label{sel52:eq:noisyerror}
\end{equation}
\end{proposition}
\begin{proof}
Insert the input into $\rho/3+2U\rho U^*/3$ and expand its determinant.
The rank conclusion follows from positivity, trace one, and the strictly
positive determinant. For the last formulas contract the resulting density
with the already specified six-dimensional effect on flag and $H_3$.
The normalized difference is Hermitian of trace zero with eigenvalues
$\pm\sqrt{14}\,\delta/[2(1+2\delta)]$, proving
\eqref{sel52:eq:noisyerror}.
\end{proof}

This is preparation noise, not channel noise or fine-record error. The
selected mixed output approaches the pure one continuously; only the
\emph{exact} rank-one assertion has a boundary. Its growing success
probability in \eqref{sel52:eq:noisyoutput} is not evidence of improved
exact-ray production. No multi-copy approximate optimum is claimed.

\section[A spanning preparation bank can have zero exact capacity]{The actual preparation bank remains the historical input}
\label{sel52:bank}

\begin{proposition}[An exact source-relative tomography--preparation separation]
\label{sel52:prop:bank}
In $H_3$, let
\[
 Z_0=P_n-P_{u_0},\quad X_0=|n\rangle\langle u_0|+|u_0\rangle\langle n|,
 \quad Y_0=i(|u_0\rangle\langle n|-|n\rangle\langle u_0|),
 \quad T_0=P_n+P_{u_0}-2P_\omega.
\]
The five invariant states
\begin{equation}
 \rho_0=I_3/3,\qquad
 \rho_j=I_3/3+B_j/12,\quad
 (B_1,B_2,B_3,B_4)=(Z_0,X_0,Y_0,T_0)
 \label{sel52:eq:bank}
\end{equation}
span the complete five-dimensional real invariant Hermitian space. All
satisfy $\rho_j\succeq I_3/6$; their outputs satisfy
$\Phi_3(\rho_j)\succeq I_3/18$. Every randomized preparation from this bank,
even with its classical tag retained, has zero finite-copy exact
primitive-ray capacity in the fixed covariant conversion class. Nevertheless
\begin{equation}
 \boxed{P_{(n+u_0)/\sqrt2}=2\rho_4+6\rho_2-7\rho_0.}
 \label{sel52:eq:signedbank}
\end{equation}
\end{proposition}
\begin{proof}
Each $B_j$ is traceless and invariant. Their real linear independence and
independence from identity prove spanning. Their least eigenvalues are at
least $-2$, giving the input floor. The identity component of
\eqref{sel52:eq:fullchannel} gives the output floor. Every nonzero tagged
block therefore has full invariant support; so do finite tensor products
and positive classical mixtures. The support obstruction applies.
Expanding the right side of \eqref{sel52:eq:signedbank} gives
$I_3/3+T_0/6+X_0/2=P_{(n+u_0)/\sqrt2}$.
\end{proof}

This is a new evaluation on the unchanged source of the inherited
V24 distinction between signed inference and positive executability, not a
new general tomography theorem. The five preparations reconstruct that pure
probe's hypothetical response statistically, but do not execute it by
randomization. If independently admitted quantum preprocessing can create
the pure probe, the preparation bank has been enlarged and the stated zero
no longer applies.

For a finite bank of actual invariant $H_3$ preparations $\rho_a$, their
one-copy capacities are now explicit: zero unless the test in
Theorem~\ref{sel52:thm:prepface} passes, and otherwise
\eqref{sel52:eq:prepfacecapacity}. At fixed probabilities $w_a$ and with
neutral tags retained, the exact success is $\sum_aw_a\alpha_a$; if the
preparation may be chosen it is $\max_a\alpha_a$. These statements concern
one-query, one-reference steering. Multiple prepared references have their
own inherited conversion problem. Noise-free equality of a preparation
minor cannot be certified from a merely small measured residual.

\section{Proof status and next source calculation}
\label{sel52:status}

\paragraph{Proved in this layer.}
The full-native helper bound proves both one-query optima over the stated
preparation classes. A two-dimensional conjugate reference attains the
assisted value; scalar-character ancillary multiplicities cannot improve
the unassisted value. The three-level preparation face, its exact noise
boundary, and the positive spanning-bank example are fully evaluated.

\paragraph{Inherited rather than newly supplied.}
V19 supplies the channel, action, native representation, helpers and clock.
V51 already achieves both numerical values in restricted families. V52
proves the global upper bounds in the declared one-call class and reduces
the sufficient reference dimension. V47--V51 supply the reference-capacity
and support-obstruction inputs. No historical instrument is inferred.

\paragraph{Unresolved historical input.}
The actual invariant preparation bank, compensating ancillary character,
and readout effects still require same-cut identification. The source
already contains determinant-sensitive $H_s$; the protocol transfers this
resource rather than deriving its occurrence. Neither coherent neutral
preparation nor conjugate-reference access follows from a signed tomography
expansion. The V110 word, directed grading and finite-Dirac incidence remain
unpopulated. No continuum, stationarity, preparation purity, or ancillary
access is added as a universal axiom of the finite duality.

The next source row is therefore the executable preparation and retained
reference, followed by its admitted positive readout bank. Another
optimization over an unnamed preparation family is unnecessary. Multiple
queries and direct target processing remain different operation classes.

\paragraph{Verification boundary.}
The checker reconstructs the original working-block coefficients and checks
the formulas, optima, effects and examples. The full-native and all-copy
claims use the written proofs, not numerical enumeration. V51's unchanged
checker is rerun. No physical measurement, complete audit of the immutable
prefix, proof-assistant certification, or independent peer review is claimed.

\begingroup
\renewcommand{\refname}{Primary-source context for V52}
\begin{thebibliography}{9}
\small\raggedright
\bibitem{sel52-Marvian}
I.~Marvian and R.~W.~Spekkens,
\emph{The theory of manipulations of pure state asymmetry: basic tools and
equivalence classes of states under symmetric operations},
New J. Phys. \textbf{15} (2013), 033001;
\url{https://arxiv.org/abs/1104.0018}.
\bibitem{sel52-FangLiu}
K.~Fang and Z.-W.~Liu,
\emph{No-Go Theorems for Quantum Resource Purification: New Approach and
Channel Theory}, PRX Quantum \textbf{3} (2022), 010337;
\url{https://arxiv.org/abs/2010.11822}.
\end{thebibliography}
\endgroup
\addtocontents{toc}{\protect\endgroup}
% END OF SOURCE-SELECTION V52 DEVELOPMENT BLOCK.
% Preserve V1--V52 and append subsequent work before the final terminator.


\clearpage
\hypersetup{pdftitle={Historical Standard-Model Source Selection: V1--V53}}
\begin{center}
{\Large\bfseries Source-selection continuation V53}\\[.4em]
{\large Correlated invariant preparations and the source-reuse law\\
Exact joint-ray capacity with locally unusable outputs}
\end{center}
\addtocontents{toc}{\protect\begingroup}
\addtocontents{toc}{\protect\renewcommand{\protect\numberline}{\protect\makebox[2.1em][l]}}

\section[Audit the preparation correlations and the joint query]{The upstream datum is the complete preparation--query pair}
\label{sel53:context}

V52 completely solves its one-query preparation--steering class. Its
three-level mixed-input exclusion concerns independent copies of a
\emph{specified reduced output}. Neither that exclusion nor its optimum
settles a correlated multi-register preparation followed by queries on all
registers. The present layer evaluates that distinction on the unchanged
V19/V51 source. It does not optimize another one-query family, introduce a
new determinant-sensitive action, or identify a tensor product of marginal
channels with an unmeasured joint source.

The result is a source-relative positive alternative. In a specified
one-active-slot neutral preparation space, the exact joint-ray capacity is
computed for \emph{every} density matrix, including mixed ones. Each local
prepared input can be mixed and each local output can have full invariant
support, while the joint output permits an exact efficient ray. Classical
correlation alone cannot do this in this space. The preparation therefore
has an actual entanglement cost, even though it carries no asymmetry.
A second calculation keeps the same one-register channel but changes its
\emph{joint} two-register implementation. The exact capacity then depends
on a cross-query coefficient invisible to every one-register experiment.

\paragraph{Audit and nonduplication.}
The interior-CP-frame obstruction is already
Theorem~\ref{sel5:thm:interior-frame}; it is not reproved here. V49 supplies
the support-period obstruction, V51 the original three-level channel, V48
the general mixed-reference conversion setting, and V52 the one-query
optima and local preparation boundary. The new computations are the
correlated-input normal form, its exact supported readout and optimum,
the source-reuse comparison, and their original-opportunity scaling.
Relational encodings and the distinction between correlated and memoryless
channel use are established settings
\cite{sel53-Bartlett,sel53-Buscemi,sel53-Modes}. The specialized formulas
below have direct proofs.

\paragraph{Unchanged one-register source.}
Write $u=u_0$ for the already present neutral helper output and retain
$H_3=\Span\{n,\omega,u\}$, with characters $1,\chi,1$, where
$\chi(A,B)=\det A\det B$. All helper jumps annihilate this forward input
space. Its original forgotten accepted channel is $\Phi_{2/3}$ in
\begin{equation}
 \Phi_p=(1-p)\id+p\operatorname{Ad}_U,\qquad
 U=\begin{pmatrix}(1-i)/2&-(1+i)/2&0\\
 -(1+i)/2&(1-i)/2&0\\0&0&1\end{pmatrix},\qquad 0<p<1.
 \label{sel53:eq:source}
\end{equation}
The variable $p$ is used only to retain the raw opportunity calculation:
there $p=2\vartheta/3$. No adjustable source probability is inferred.
The same full native carrier remains in force; $H_3$ is a forward
restriction, not an asserted reducing subalgebra under all adjoints.

\paragraph{The new, explicitly admitted operation class.}
Co-prepare $k$ source registers in the invariant space defined below; apply
$\Phi_p^{\otimes k}$, without retaining any source outcome; retain the
\emph{joint} quantum output as the reference. Independently supply the
same normalized target packet as V52,
\begin{equation}
 |\Xi\rangle=\frac{b|0\rangle+d_{\rm tar}|1\rangle}{\sqrt2},
 \quad v_{\beta,\theta}=\sqrt\beta\,b+
           e^{i\theta}\sqrt{1-\beta}\,d_{\rm tar},\quad 0<\beta<1,
 \label{sel53:eq:target}
\end{equation}
where the target vectors have characters $1,\chi^{-1}$ and the flag
vectors $1,\chi$. Their inaccessible markers are identical. Optimize only
over invariant effects $0\le F\le I$ on the flag and the entire queried
reference. No operation acts directly on the target. This assumes actual
coherent access within equal-total-charge readout sectors. The success
probability includes all filters and query outcomes. A correlated
preparation is not supplied by randomization of its local marginals.

\section[An exact correlated-output normal form]{The complete output of one-active-slot preparations}
\label{sel53:normal}

Let $N_j$ denote the $k$-register vector with $n$ in slot $j$ and $u$
elsewhere, and let $W_j$ replace that $n$ by $\omega$. Thus each $N_j$ is
neutral and each $W_j$ has character $\chi$. For a density $R$ on
$\C^k$ prepare
\begin{equation}
 \rho_R=\sum_{j,l}R_{jl}|N_j\rangle\langle N_l|,
 \qquad D_R=\diag(R_{11},\ldots,R_{kk}).
 \label{sel53:eq:prep}
\end{equation}
Every such preparation is invariant. Remove indices with $R_{jj}=0$
from the coordinate calculation; positivity makes their entire rows and
columns zero, and their physical registers stay in $u$. In the formulas
below $k$ is the number of occupied indices and $D_R\succ0$. No unused
native direction is removed from the underlying source.

Put
\begin{equation}
 h=-\frac{1+i}{2},\quad c=1+ph,\quad d=ph,\quad
 \eta=\frac{p(1-p)}2,\qquad c-d=1.
 \label{sel53:eq:constants}
\end{equation}
The names $c,d$ here are scalar coefficients, not the target vector
$d_{\rm tar}$. Define $Vx=(cx,dx)$ and $Wx=(x,x)$ into the ordered
reference space $\Span\{N_j\}\oplus\Span\{W_j\}$.

\begin{theorem}[Joint normal form and its exact kernel]
\label{sel53:thm:normal}
The original product query has output
\begin{equation}
 \boxed{\sigma_R=\Phi_p^{\otimes k}(\rho_R)
       =VRV^*+\eta WD_RW^*.}
 \label{sel53:eq:normal}
\end{equation}
In its two charge blocks this is
\begin{equation}
 \sigma_R=\begin{pmatrix}A&C\\C^*&B\end{pmatrix},\quad
 A=|c|^2R+\eta D_R,\quad B=|d|^2R+\eta D_R,\quad
 C=c\bar d R+\eta D_R.
 \label{sel53:eq:blocks}
\end{equation}
Both diagonal blocks are positive definite on the occupied coordinates,
$\Tr\sigma_R=1$, and
\begin{equation}
 \boxed{\ker\sigma_R=\{(x,-x):Rx=0\},\qquad
        \rank\sigma_R=k+\rank R.}
 \label{sel53:eq:kernel}
\end{equation}
\end{theorem}
\begin{proof}
Use the random-unitary expression in \eqref{sel53:eq:source} to evaluate
this \emph{specified product channel}. If $\xi_j\in\{0,1\}$ denotes its
independent mixture variable with mean $p$, a vector $N_j$ becomes
$N_j+h\xi_j(N_j+W_j)$. The mean map is $V$. For distinct indices the
centered mixture variables have zero covariance; the diagonal variance is
$p(1-p)$. Multiplying by $|h|^2=1/2$ gives precisely the second term of
\eqref{sel53:eq:normal}. This is a calculation of the channel, not a claim
that the original six physical labels resolve these two mixture terms.
Linearity gives the formula for every mixed $R$.

The trace is $\bigl(|c|^2+|d|^2+2\eta\bigr)\Tr R=1$.
Since $D_R\succ0$ and $\eta>0$, positivity of both terms shows that a
kernel vector $(x,y)$ must satisfy $x+y=0$ and
$R^{1/2}(\bar c x+\bar d y)=0$. The second equation is $Rx=0$ because
$c-d=1$. The dimension formula follows. The same positive diagonal term
makes $A,B$ faithful.
\end{proof}

\begin{corollary}[All one-register data omit the decisive preparation]
\label{sel53:cor:marginals}
Writing $r_j=R_{jj}$, the local input and output are
\begin{align}
 \rho_j&=r_jP_n+(1-r_j)P_u,\\
 \sigma_j&=r_j\begin{pmatrix}1-p/2&ip/2&0\\
                 -ip/2&p/2&0\\0&0&0\end{pmatrix}
                 +(1-r_j)P_u,\qquad
 \det\sigma_j=\eta r_j^2(1-r_j).
 \label{sel53:eq:marginals}
\end{align}
For $0<r_j<1$, every local output has full invariant $H_3$ support and
has zero primitive exact-ray capacity at every finite independent-copy
number. The complete local input and output tomographies depend only on
$D_R$, not on the off-diagonal preparation coordinates.
\end{corollary}
\begin{proof}
Tracing all but one register kills $|N_j\rangle\langle N_l|$ for $j\ne l$.
Applying the unchanged local channel gives the displayed matrix and
determinant. The all-copy conclusion is the inherited invariant-support
obstruction, applied to these specified reduced outputs. It is not applied
to the correlated state $\sigma_R$.
\end{proof}

\section[Exact capacity is a preparation-kernel weight]{A complete invariant readout solution on the joint archive}
\label{sel53:capacity}

Let $P_0=P_{\ker R}$ on the occupied slot coordinates. In total charge
one, define an isometry from $\C^k$ into flag--reference space by
\begin{equation}
 L_{\beta,\theta}x=
 \sqrt\beta\,|0\rangle\otimes\mathsf W(x)
 +e^{-i\theta}\sqrt{1-\beta}\,|1\rangle\otimes\mathsf N(x),
 \label{sel53:eq:L}
\end{equation}
where $\mathsf N(x)=\sum_jx_jN_j$ and
$\mathsf W(x)=\sum_jx_jW_j$ are the separate charge embeddings.

\begin{theorem}[Exact capacity for every correlated density]
\label{sel53:thm:capacity}
In the stated invariant-effect protocol, the maximum success for any
specified two-shadow pure ray \eqref{sel53:eq:target} is
\begin{equation}
 \boxed{\alpha_p(R)=\frac{\eta}{2}\Tr(P_0D_R)
              =\frac{p(1-p)}4\Tr(P_0D_R).}
 \label{sel53:eq:capacity}
\end{equation}
It is attained by the invariant projector
\begin{equation}
 F_{\beta,\theta}=L_{\beta,\theta}P_0L_{\beta,\theta}^*,
 \qquad J_F=\alpha_p(R)P_{v_{\beta,\theta}}.
 \label{sel53:eq:effect}
\end{equation}
The effect acts as zero on all other physical readout sectors; $I-F$ is
its admitted positive complement. In particular,
\begin{equation}
 \boxed{\alpha_p(R)>0\iff\rank R<k.}
 \label{sel53:eq:criterion}
\end{equation}
This is also equivalent to existence of a successful conversion from some
finite number of independent copies of the whole prepared joint output.
\end{theorem}
\begin{proof}
An invariant effect is block diagonal in total character. Total charges
zero and two give a target supported only on $b$ or only on $d_{\rm tar}$.
A positive contribution on either ray cannot lie below a nonzero pure
packet with both components. Their supported effects must therefore vanish.
Faithfulness of $A,B$ makes this a zero on their whole physical supports.
Only total charge one remains, with order
$(|0\rangle\otimes W_j,|1\rangle\otimes N_j)$.

Zero expectation on the target vector orthogonal to $v_{\beta,\theta}$
forces the effect into the kernel of the corresponding positive unwanted
matrix. That matrix is an invertible diagonal congruence, with a block
swap and phases, of $\sigma_R$. By \eqref{sel53:eq:kernel} its kernel
is exactly $\Ran(L_{\beta,\theta}P_0)$. Thus every feasible effect on
this sector obeys $0\le F\le L_{\beta,\theta}P_0L_{\beta,\theta}^*$.
The upper endpoint maximizes its positive success functional.

On $\ker R$, all compressions of $A,B,C$ in \eqref{sel53:eq:blocks}
are the same matrix $\eta P_0D_RP_0$. Direct contraction of the original
normalized target packet gives diagonal entries
$\eta\beta\Tr(P_0D_R)/2$ and
$\eta(1-\beta)\Tr(P_0D_R)/2$, with the stated common phase in the
off-diagonal entry. This proves both attainment and maximality.
Since $D_R\succ0$, its trace on a nonzero $P_0$ is positive. If $R$ is
faithful, \eqref{sel53:eq:kernel} makes the whole output faithful on its
invariant two-charge support. The inherited support theorem then excludes
all finite-copy exact conversion. If $R$ is singular, the displayed
one-output effect already succeeds.
\end{proof}

The independence used in the last sentence concerns repeated preparations
of the \emph{whole correlated archive}; it is not independence of its
register marginals. The target source remains independent of that archive.
For $\beta=0$ or $1$ the target is invariant and the pure-asymmetric
constraint changes; \eqref{sel53:eq:capacity} is not its optimum.

\begin{proposition}[The optimum in the complete one-active-slot class]
\label{sel53:prop:optimum}
For $k$ occupied slots,
\begin{equation}
 \boxed{\alpha_p(R)\le\frac{p(1-p)}4\left(1-\frac1k\right).}
 \label{sel53:eq:optimum}
\end{equation}
Equality holds precisely for a pure $R=|a\rangle\langle a|$ with
$|a_j|^2=1/k$. For the unchanged accepted channel, the maximum is
$(k-1)/(18k)$. A genuinely mixed example is
$R=(I-|a\rangle\langle a|)/(k-1)$ for the flat vector $a$; its success
is $p(1-p)/(4k)$, or $1/(18k)$ at $p=2/3$.
\end{proposition}
\begin{proof}
A density satisfies $R\preceq P_{\supp R}$. Thus, writing $r_j=R_{jj}$,
\[
 \Tr(P_0D_R)\le\sum_j r_j(1-r_j)
              =1-\sum_jr_j^2\le1-1/k.
\]
Equality in the first bound, since $D_R\succ0$, requires
$P_{\supp R}=R$, hence rank one. Equality in the second requires all
$r_j=1/k$. Both conditions suffice. For the mixed example $D_R=I/k$
and $P_0=P_a$, giving the stated trace $1/k$.
\end{proof}

This is an optimum \emph{within the declared correlated preparation class}
with every register queried. It is not an optimum over all $k$-query
strategies and does not improve or contradict V52's distinct one-query
bounds. In particular it is not necessary to query every member of an
entangled preparation if a retained neutral bypass register is independently
admitted; that is another access policy.

\section[Two locally unusable outputs can be jointly useful]{Exact same-marginal comparisons and the entanglement cost}
\label{sel53:examples}

\begin{theorem}[Two unchanged source calls, one exact joint ray]
\label{sel53:thm:two}
For $k=2$, prepare the invariant vector
\begin{equation}
 \Psi=\frac{n\otimes u+u\otimes n}{\sqrt2}
 \label{sel53:eq:Bellneutral}
\end{equation}
and apply two independent accepted source channels. In order
$(N_1,N_2,W_1,W_2)$, the reference is
\begin{equation}
 \sigma_*=
 \begin{pmatrix}
 1/3&5/18&i/6&-1/18+i/6\\
 5/18&1/3&-1/18+i/6&i/6\\
 -i/6&-1/18-i/6&1/6&1/9\\
 -1/18-i/6&-i/6&1/9&1/6
 \end{pmatrix},\qquad \rank\sigma_*=3.
 \label{sel53:eq:sigma2}
\end{equation}
Let $z=(1,-1)/\sqrt2$. The invariant rank-one effect onto
$(|0\rangle\mathsf W(z)+|1\rangle\mathsf N(z))/\sqrt2$ returns
\begin{equation}
 \boxed{J_*^{\rm sel}=\frac1{72}
 \begin{pmatrix}1&1\\1&1\end{pmatrix},\qquad
 \alpha_*=\frac1{36},\qquad h_*=\frac1{5184}.}
 \label{sel53:eq:tworesult}
\end{equation}
This is the exact optimum for this prepared joint reference. Each local
output has determinant $1/72$ and zero exact finite-copy capacity.
\end{theorem}
\begin{proof}
Insert $R=\tfrac12\begin{psmallmatrix}1&1\\1&1\end{psmallmatrix}$
into \eqref{sel53:eq:blocks}, using
$c=(2-i)/3$, $d=-(1+i)/3$, $\eta=1/9$. Its kernel is generated by
$(1,-1,-1,1)^T$. The capacity theorem gives
$\Tr(P_0D_R)=1/2$ and hence $1/36$. The selected matrix and cross
are direct substitutions. The marginal determinant is
$(1/9)(1/2)^2(1/2)=1/72$.
\end{proof}

Replacing this preparation by
$\tfrac12(P_{n\otimes u}+P_{u\otimes n})$ preserves every local input
and output state and every one-register channel, but its $R=I_2/2$ makes
the entire joint output faithful on its invariant four-dimensional support.
Its finite-copy exact capacity is zero. Retaining the \emph{classical}
slot label does not fix this: each labelled reference occupies its full
invariant $N_j/W_j$ plane. Thus no amount of separate marginal tomography
can decide between these two preparation policies.

There is also a mixed positive example. For $k=3$, take
$R=(I_3-P_{(1,1,1)/\sqrt3})/2$. It has rank two, $D_R=I_3/3$, and
reference rank five. The exact success is $1/54$. Every local input is
mixed, every local output is faithful with determinant $2/243$, and the
complete preparation is mixed. This does not repeat V52's pure-local-input
criterion on a different name: the physically retained object is now the
correlated multi-register preparation.

\begin{proposition}[Classical correlation is insufficient in this class]
\label{sel53:prop:entanglement}
A fully separable state supported on $\Span\{N_1,\ldots,N_k\}$ is a
diagonal mixture of the $P_{N_j}$. Therefore every positive exact capacity
in \eqref{sel53:eq:criterion} requires an entangled preparation in this
class, despite full $G_0$ invariance.
\end{proposition}
\begin{proof}
Positivity forces each pure product vector in a separable decomposition
to belong to the same support. Every local factor then lies in
$\Span\{n,u\}$. The sum of the local $n$-occupation projections has
sharp value one. On a product vector its variance is the sum of the
nonnegative local variances, so every local occupation is zero or one,
and exactly one is one. The product vector is therefore one of the
$N_j$. A diagonal density is faithful on its occupied indices and has
zero capacity by Theorem~\ref{sel53:thm:capacity}.
\end{proof}

The word ``occupation'' is a coordinate description of these neutral
slots, not a particle-number or generation assertion. Neither entanglement
nor an invariant joint effect follows from the single-register physical
frame merely because it spans the appropriate matrices.

\section[The same single-use source can have different joint capacities]{The source-reuse correlation must also be retained}
\label{sel53:reuse}

The product-query assumption is independently substantive. To separate it
from preparation, retain \eqref{sel53:eq:Bellneutral} and define the
following \emph{distinct joint-channel comparison family}:
\begin{align}
 \mathcal C_t={}&(t-1/3)\id\otimes\id
 +(2/3-t)(\operatorname{Ad}_U\otimes\id+
                    \id\otimes\operatorname{Ad}_U)
 +t\operatorname{Ad}_{U\otimes U},\qquad 1/3\le t\le2/3.
 \label{sel53:eq:jointfamily}
\end{align}
All coefficients are nonnegative and sum to one. For every joint input,
including entangled inputs, either marginal channel is exactly the old
$\Phi_{2/3}$. At $t=4/9$, this is the independent product query; at
$t=2/3$, one common mixture variable acts on both registers.

\begin{theorem}[A joint coefficient invisible to all local channel tests]
\label{sel53:thm:correlation}
For \eqref{sel53:eq:jointfamily} on the same preparation and target
protocol, the exact capacity is
\begin{equation}
 \boxed{\alpha(t)=\frac{2/3-t}{8}.}
 \label{sel53:eq:jointcapacity}
\end{equation}
The same invariant effect as in Theorem~\ref{sel53:thm:two} is optimal
for $t<2/3$. At $t=2/3$ the reference support is invariant and no finite
number of its independent copies can supply an exact primitive ray.
Consequently the identical local channel and identical preparation give
success $1/24$, $1/36$, or zero at $t=1/3,4/9,2/3$, respectively.
\end{theorem}
\begin{proof}
Represent the two mixture indicators by $\xi_1,\xi_2$ only to calculate
the displayed channel. Their means are $2/3$ and covariance matrix is
\[
 \Gamma_t=\begin{pmatrix}2/9&t-4/9\\t-4/9&2/9\end{pmatrix}.
\]
For the flat preparation vector $a=(1,1)/\sqrt2$, the output is
$VP_aV^*+W(\Gamma_t/4)W^*$, with $c,d$ as above.
The covariance eigenvalues are $t-2/9>0$ and $2/3-t$.
For $t<2/3$ the noise matrix is faithful, so the kernel and maximal-effect
proof of Theorem~\ref{sel53:thm:capacity} apply with
$\eta D_R$ replaced by $\Gamma_t/4$. The success is one half its
trace on $a^\perp$, namely $(2/3-t)/8$.
At $t=2/3$ both terms have their slot range in $\C a$. Their charge
vectors span the two-dimensional invariant plane
$\Span\{\mathsf N(a),\mathsf W(a)\}$, with strictly positive density there. The
invariant-support exclusion gives the final conclusion.
\end{proof}

The mixture indicators in this proof are \emph{not} asserted to be the
original six recorded outcomes. A random-unitary formula for a forgotten
channel does not identify a physical dilation or a cross-query record.
The family \eqref{sel53:eq:jointfamily} exhibits different possible
\emph{joint} implementations with the same local map. Determining $t$
requires the actual joint channel or an independently calibrated joint
implementation, not the marginals or a chosen unravelling. No value other
than $4/9$ is assigned to the explicitly independent benchmark.

\section[Unchanged local data can hide a preparation-support boundary]{Correlation noise and physical opportunity normalization}
\label{sel53:noiseclock}

\begin{proposition}[Dephasing the preparation, not the source]
\label{sel53:prop:noise}
Keep two independent original accepted queries and replace only the joint
neutral preparation matrix by
\[
 R_\delta=\frac12\begin{pmatrix}1&1-\delta\\1-\delta&1\end{pmatrix},
 \qquad0\le\delta\le1.
\]
All local input and output states are unchanged. The joint output has
\begin{equation}
 \boxed{\det\sigma_{R_\delta}=\frac{\delta(2-\delta)}{1296}.}
 \label{sel53:eq:noisedet}
\end{equation}
Thus every $\delta>0$ excludes exact primitive-ray extraction from every
finite number of independent joint archives. The \emph{unchanged} effect
of Theorem~\ref{sel53:thm:two} nevertheless gives
\begin{equation}
 J_\delta=\frac1{72}
 \begin{pmatrix}1+2\delta&1-\delta-3i\delta\\
 1-\delta+3i\delta&1+5\delta\end{pmatrix},\qquad
 \Tr J_\delta=\frac{2+7\delta}{72},\quad
 \det J_\delta=\frac{\delta}{576}.
 \label{sel53:eq:noiseselected}
\end{equation}
Its normalized balanced-target trace distance is
\begin{equation}
 \boxed{\left\|\frac{J_\delta}{\Tr J_\delta}-P_v\right\|_1
       =\frac{3\sqrt{14}\,\delta}{2+7\delta}.}
 \label{sel53:eq:noiseerror}
\end{equation}
This is an evaluation of that effect, not an optimal approximation bound.
\end{proposition}
\begin{proof}
Diagonalize the two slot coordinates into their symmetric and antisymmetric
vectors. The preparation eigenvalues are $1-\delta/2$ and $\delta/2$.
Each corresponding two-charge block of the output has determinant equal
to its preparation eigenvalue divided by $18$. Their product gives
\eqref{sel53:eq:noisedet}. Direct contraction gives
\eqref{sel53:eq:noiseselected}; the traceless normalized difference has
opposite eigenvalues of modulus
$3\sqrt{14}\delta/[2(2+7\delta)]$. The all-copy exclusion follows from
faithfulness on the invariant four-dimensional support.
\end{proof}

\begin{proposition}[The raw source clock is not conditioned away]
\label{sel53:prop:raw}
If each occupied register instead receives one \emph{independent raw locked
opportunity}, with original acceptance factor $\vartheta>0$, its forgotten
map on $H_3$ is $\Phi_p$ with $p=2\vartheta/3$. The exact capacity is
\begin{equation}
 \boxed{\alpha_{\rm raw}(R)=
 \frac{\vartheta}{6}\left(1-\frac{2\vartheta}{3}\right)
 \Tr(P_{\ker R}D_R).}
 \label{sel53:eq:raw}
\end{equation}
For the two-register flat pure preparation this is
$\vartheta(1-2\vartheta/3)/12$. It is not
$\vartheta^2/36$: the latter is the distinct policy that retains only the
branch on which \emph{both} opportunities accept.
\end{proposition}
\begin{proof}
The idle map is identity on the input and the forgotten accepted map is
$\id/3+2\operatorname{Ad}_U/3$. Their original mixture is
$(1-2\vartheta/3)\id+(2\vartheta/3)\operatorname{Ad}_U$.
Insert this $p$ into \eqref{sel53:eq:capacity}. With no acceptance
postselection, trajectories in which different registers undergo different
mixture components also contribute and must not be deleted. Keeping both
acceptances instead multiplies the normalized accepted-product probability
$1/36$ by $\vartheta^2$.
\end{proof}

This is a probability per specified set of raw opportunities. It is not an
elapsed-time estimate, and independence of their clock/source blocks is
explicit. If each register is first run until one accepted event, the
accepted-product result applies with its separately retained waiting law.
Co-instantiation of the registers, or their sequential use with a retained
quantum memory, is a physical source requirement rather than a consequence
of counting independent experimental repetitions.

\section{Proof status and the next source row}
\label{sel53:status}

\paragraph{Proved in this layer.}
The exact product-query output, kernel, optimal invariant pure-ray effect,
and success formula hold for every density in the stated correlated
one-active-slot preparation space. One joint archive suffices whenever any
finite number can suffice. Mixed, locally unusable preparations can be
positive globally; all fully separable preparations in this space are
negative. The identical-local-channel family isolates a separate joint
source coefficient. Preparation dephasing and raw opportunities are
calculated without changing their marginals or suppressing costs.

\paragraph{Inherited, not newly established.}
The original native source, its determinant-sensitive action, the forward
three-level restriction and the target packet are V19/V51/V52 inputs.
V48--V51 supply the finite-copy invariant-support exclusions. V52's
one-query optimum and its local mixed-preparation classification remain
unchanged. The present results do not improve those optima; they identify
a different executable-preparation and source-reuse interface. No numerical
optimization or genericity claim is used in the new capacity formulas.

\paragraph{Unresolved historical input.}
The actual correlated neutral preparation, the joint query law and the
invariant collective readout must be independently admitted on the same
carrier. Their local marginals do not determine them. The entangled
preparation carries no determinant asymmetry, but it is still a coherent
preparation resource; the asymmetry is transferred from the already supplied
source action. No historical V110 word, matter type, grading or finite-Dirac
incidence is populated here. A single efficient \emph{target coefficient
state} is not identified with an executed microscopic Yukawa word.

The next useful audit is therefore the original \emph{joint}
preparation--process table, including retained quantum correlations and its
actual replica/source-reuse policy. Neither a new favourable local pure
state nor a tensor product silently substituted for marginal channel data
answers that question. The direct historical-word route remains separate;
correlated preparation is not imposed as an axiom of finite duality.

\paragraph{Verification boundary.}
The accompanying checker evaluates the reduced normal form against the
full two- and three-register source action on the stated forward sector,
checks the kernel, selected effects and exact optima on independent
rational examples, and verifies the joint-correlation and opportunity
polynomials symbolically. General rank and maximality statements rest on
the written proofs. The unchanged V52 checker is rerun; this is not a
complete independent re-audit of the preceding fifty-two layers. No physical
measurements, proof-assistant formalization, or independent peer review are
claimed.

\begingroup
\renewcommand{\refname}{Primary-source context for V53}
\begin{thebibliography}{9}
\small\raggedright
\bibitem{sel53-Bartlett}
S.~D.~Bartlett, T.~Rudolph and R.~W.~Spekkens,
\emph{Reference frames, superselection rules, and quantum information},
Rev. Mod. Phys. \textbf{79} (2007), 555;
\url{https://arxiv.org/abs/quant-ph/0610030}.
\bibitem{sel53-Buscemi}
F.~Buscemi and N.~Datta,
\emph{The quantum capacity of channels with arbitrarily correlated noise},
IEEE Trans. Inf. Theory \textbf{56} (2010), 1447--1460;
\url{https://arxiv.org/abs/0902.0158}.
\bibitem{sel53-Modes}
I.~Marvian and R.~W.~Spekkens,
\emph{Modes of asymmetry: the application of harmonic analysis to symmetric
quantum dynamics and quantum reference frames},
Phys. Rev. A \textbf{90} (2014), 062110;
\url{https://arxiv.org/abs/1312.0680}.
\end{thebibliography}
\endgroup
\addtocontents{toc}{\protect\endgroup}
% END OF SOURCE-SELECTION V53 DEVELOPMENT BLOCK.
% Preserve V1--V53 and append subsequent work before the final terminator.


\clearpage
\hypersetup{pdftitle={Historical Standard-Model Source Selection: V1--V54}}
\begin{center}
{\Large\bfseries Source-selection continuation V54}\\[.4em]
{\large All marginal-compatible two-register channels\\
An exact joint-correlation disk and a sharp coherent-ray capacity}
\end{center}
\addtocontents{toc}{\protect\begingroup}
\addtocontents{toc}{\protect\renewcommand{\protect\numberline}{\protect\makebox[2.1em][l]}}

\section[Replace the assumed joint mixture by its complete positive fibre]{The upstream object is the entire fixed-marginal channel fibre}
\label{sel54:context}

V53 isolates a genuine joint-source datum: its family
\eqref{sel53:eq:jointfamily} has identical local channels but different
exact-ray capacities. That family assumes a classical correlation between
the particular two terms in one random-unitary formula. A representation
of a local channel is not a specification of its joint use. This layer
removes that restriction: it classifies \emph{every} bipartite CPTP map on
the verified three-level forward port whose two marginal maps equal the
original source on \emph{all} joint inputs. The original correlated
preparation and invariant target readout class are retained.

The conclusion is stronger than a further separating example. Four real
parameters describe the complete channel fibre; one complex parameter
determines the selected preparation's entire output and optimal
balanced-ray success. Its exact feasible set is a disk. The product-query
point keeps V53's success $1/36$. The maximum over all compatible joint
maps is $3/14$, attained on a specified different implementation with
shared classical randomness. That implementation is not assigned to the
historical source. In particular its alternative local unitary
unravelling is not declared to be the original recorded coefficient bank.

\paragraph{Inherited inputs and nonduplication.}
The local source is \eqref{sel53:eq:source} at $p=2/3$, the preparation is
\eqref{sel53:eq:Bellneutral}, and the independent target packet and allowed
readout effects are \eqref{sel53:eq:target}, with balanced target.
V53's product-channel formula and its one-parameter correlated comparison
are recovered, not reclassified as new results. The positive unwanted-output
kernel method is inherited from V46/V48/V53. What is new is the exhaustive
fixed-marginal channel classification, its feasible joint disk, the
capacity at every point, and the global maximum on that fibre.
The distinction between nonsignalling and a chosen joint implementation,
and correlation-matrix descriptions of dephasing, are established
frameworks \cite{sel54-Beckman,sel54-Eggeling,sel54-Puchala}. All specialized
formulas and bounds below are proved directly.

\paragraph{Actual scope of the local constraint.}
Let $H_3=\Span\{n,\omega,u\}$ and retain its characters $1,\chi,1$,
where $\chi(A,B)=\det A\det B$. Write
\[
 h=(n+\omega)/\sqrt2,\qquad d=(n-\omega)/\sqrt2,\qquad
 P_1=P_h,\quad P_0=I-P_1,\quad U=P_0-iP_1.
\]
Here $d$ denotes the local dark vector, not the target coefficient
$d_{\rm tar}$. Set $\Phi=\id/3+2\operatorname{Ad}_U/3$.
A candidate joint map $\mathcal E$ must satisfy
\begin{equation}
 \boxed{\Tr_2\mathcal E(X)=\Phi(\Tr_2X),\qquad
 \Tr_1\mathcal E(X)=\Phi(\Tr_1X)\quad\text{for every }X\in\End(H_3^{\otimes2}).}
 \label{sel54:eq:marginals}
\end{equation}
These are complete linear channel identities, not just equality on the
one preparation used below. They imply the same identities with an
untouched reference. They do not establish independence, localizability,
a joint law for the old six labels, or a temporal memory model. Only the
verified forward port is classified; an extension of a new comparison
map to the full helper-bearing historical instrument is not inferred.

\section[Four real parameters classify every compatible joint map]{A finite correlation-matrix normal form}
\label{sel54:normal}

Put $\Pi_{ab}=P_a\otimes P_b$ for $a,b\in\{0,1\}$, in order
$00,01,10,11$, and set
\[
 z=\frac{1+2i}{3},\qquad q_0=|z|^2=\frac59,\qquad
 \delta_0=1-q_0=\frac49.
\]
The symbol $q_0$ here is a scalar spectral coefficient, not a physical
preparation or a probability assigned to an original label.

\begin{theorem}[Complete nonsignalling fibre of the fixed local source]
\label{sel54:thm:fibre}
Every CPTP map satisfying \eqref{sel54:eq:marginals} has the unique form
\begin{equation}
 \mathcal E_{a,b}(X)=\sum_{s,t\in\{00,01,10,11\}}C(a,b)_{st}\Pi_sX\Pi_t,
 \qquad
 C(a,b)=\begin{pmatrix}
 1&z&z&a\\
 \bar z&1&b&z\\
 \bar z&\bar b&1&z\\
 \bar a&\bar z&\bar z&1
 \end{pmatrix}\succeq0.
 \label{sel54:eq:C}
\end{equation}
Conversely every such positive matrix defines a CPTP map satisfying
\eqref{sel54:eq:marginals}. Its Choi rank is $\rank C(a,b)$.
Thus the unconstrained part of the fixed-marginal fibre has four real
coordinates, not the one coordinate of V53's particular mixture family.
\end{theorem}
\begin{proof}
Let $\mathcal D=\End(P_0H_3)\oplus\End(P_1H_3)$. The local Heisenberg
channel fixes $\mathcal D$ pointwise. By the marginal identities,
$\mathcal E^*$ fixes $\mathcal D\otimes I$ and $I\otimes\mathcal D$.
For every $a$ in either algebra it also fixes $a^*a$ and $aa^*$.
If $V$ is a Stinespring isometry for the unital map $\mathcal E^*$, the
squared norm of $(a\otimes I)V-Va$ is the Schwarz equality residual and
vanishes. Accordingly every Kraus coefficient commutes with both local
algebras. Their joint commutant is exactly
$\Span\{\Pi_{00},\Pi_{01},\Pi_{10},\Pi_{11}\}$. Write the coefficients
as $L_j=\sum_s\ell_{js}\Pi_s$.

The channel therefore has \eqref{sel54:eq:C} with
$C_{st}=\sum_j\ell_{js}\bar\ell_{jt}$. Positivity is automatic and trace
preservation gives $C_{ss}=1$. The first local marginal fixes the two
entries for changing the first binary sector, independently of the
second; the other marginal fixes the corresponding second-sector entries.
Their values are $z$ and $\bar z$. Only $C_{00,11}$ and $C_{01,10}$
remain free. Conversely a Gram factorization of $C$ supplies the displayed
Kraus coefficients. Terms with different traced-out sectors have zero
partial trace, and terms with equal traced-out sectors have precisely the
local coefficient, proving both marginal identities. Linear independence
of the four orthogonal nonzero projections proves uniqueness and equality
of the Choi and correlation-matrix ranks.
\end{proof}

\begin{theorem}[The exact disk of the relevant joint coherence]
\label{sel54:thm:disk}
A complex number $b$ occurs in \eqref{sel54:eq:C} if and only if
\begin{equation}
 \boxed{|b-5/9|\le4/9.}
 \label{sel54:eq:disk}
\end{equation}
For $b\ne1$ in this disk, define
\[
 s_b=\frac{2(1-\Re b)}{1-|b|^2}.
\]
The complete compatible set of $a$ is
\begin{equation}
 \boxed{|a-z^2s_b|\le1-\frac59s_b.}
 \label{sel54:eq:adisk}
\end{equation}
For $b=1$ the same formula holds with $s_b=1$.
At every boundary point of \eqref{sel54:eq:disk} other than $b=1$,
$a=(-3+4i)/5$ is forced and $C$ has rank two.
\end{theorem}
\begin{proof}
The principal block on $00,01,10$ has Schur complement
\[
 \begin{pmatrix}4/9&b-5/9\\\bar b-5/9&4/9\end{pmatrix}
\]
over its first entry. This proves necessity of the disk. For $b\ne1$ in
it, $|b|<1$. Reorder $C$ into the middle pair $01,10$ followed by the
endpoints $00,11$. The middle block is
$M_b=\left(\begin{smallmatrix}1&b\\\bar b&1\end{smallmatrix}\right)$,
and $\mathbf1^*M_b^{-1}\mathbf1=s_b$. The endpoint Schur complement is
\[
 \begin{pmatrix}
 1-q_0s_b&a-z^2s_b\\
 \bar a-\bar z^2s_b&1-q_0s_b
 \end{pmatrix}.
\]
The disk is precisely $q_0s_b\le1$; the remaining positivity condition is
\eqref{sel54:eq:adisk}. Thus every $b$ in the disk extends, for example
by choosing its displayed centre. At $b=1$, $M_b$ has range
$\C\mathbf1$, all cross rows lie in that range, and its Moore--Penrose
inverse gives $\mathbf1^*M_b^\dagger\mathbf1=1$. The same argument applies.
For a nontrivial boundary point, $q_0s_b=1$, the endpoint complement is
zero, and the invertible middle block proves the rank and unique $a$.
\end{proof}

\section[One complex joint row determines the prepared reference]{The original neutral preparation throughout the complete fibre}
\label{sel54:output}

Retain exactly $\Psi=(N_1+N_2)/\sqrt2$ from V53, with
$N_1=n\otimes u$, $N_2=u\otimes n$, $W_1=\omega\otimes u$,
$W_2=u\otimes\omega$. It is invariant under $G_0$ and no source outcome
is read. In order $(N_1,N_2,W_1,W_2)$ write $b=x+iy$.

\begin{proposition}[Complete output and its rank stratification]
\label{sel54:prop:output}
The output is independent of $a$ and equals
\begin{equation}
 \sigma_b=\begin{pmatrix}
 1/3&\bar b/8+5/24&i/6&\bar b/8-1/8+i/6\\
 b/8+5/24&1/3&b/8-1/8+i/6&i/6\\
 -i/6&\bar b/8-1/8-i/6&1/6&\bar b/8+1/24\\
 b/8-1/8-i/6&-i/6&b/8+1/24&1/6
 \end{pmatrix}.
 \label{sel54:eq:output}
\end{equation}
Its rank is three in the open disk and two on its entire boundary. In the
open disk its kernel is exactly
$\C(N_1-N_2-W_1+W_2)$. Both charge-diagonal blocks are positive definite
except at $b=1$, where the support is the invariant plane
$\Span\{N_1+N_2,W_1+W_2\}$ and the density is faithful there.
\end{proposition}
\begin{proof}
In the joint eigenspaces the input has components
\[
 \frac1{\sqrt2}\,\frac{d\otimes u+u\otimes d}{\sqrt2},\qquad
 \frac12h\otimes u,\qquad\frac12u\otimes h,
\]
and no $11$ component. Its output is the corresponding diagonal
congruence of the $00,10,01$ principal block of $C$; this explains why
$a$ drops out. Returning to the $N/W$ basis gives \eqref{sel54:eq:output}.
The determinant of that weighted three-by-three block is
\[
 \frac{(4/9)^2-|b-5/9|^2}{32}.
\]
Its first two-by-two principal minor is nonzero, so the stated ranks
follow. The missing antisymmetric dark vector gives the universal kernel.
The determinants of the two charge blocks are, respectively,
$1/9-|\bar b/8+5/24|^2$ and
$1/36-|\bar b/8+1/24|^2$.
On the disk both are positive except at its rightmost point $b=1$.
There, each is supported on the symmetric slot vector. The remaining
normalized charge block has determinant $1/9$, proving faithfulness.
\end{proof}

\begin{corollary}[An invariant joint measurement identifies the missing row]
\label{sel54:cor:testers}
Let
$X_W=|W_1\rangle\langle W_2|+|W_2\rangle\langle W_1|$ and
$Y_W=-i|W_1\rangle\langle W_2|+i|W_2\rangle\langle W_1|$,
zero on the orthogonal complement. Both commute with $G_0$, have norm
one, and
\begin{equation}
 \boxed{b=4\Tr(X_W\sigma_b)-\frac13+4i\Tr(Y_W\sigma_b).}
 \label{sel54:eq:measuredrow}
\end{equation}
Thus two invariant expectation values on the unchanged preparation
determine this task completely, although they do not determine the joint
channel parameter $a$.
\end{corollary}
\begin{proof}
The two $W$ vectors carry the same character. Their off-diagonal matrix
element in \eqref{sel54:eq:output} is $\bar b/8+1/24$. Taking its two
quadratures gives the formula. Each expectation is equivalently a binary
probability for $(I+X_W)/2$ or $(I+Y_W)/2$ on the full output.
\end{proof}

These are collective \emph{invariant} effects, not separately guaranteed
historical Reads. Signed contextual reconstruction is another route when
the effects lie in the admitted tester span. Full one-register tomography
cannot supply this row: by \eqref{sel54:eq:marginals}, every local input
and output marginal remains exactly the V53 value.

\section[The entire disk has an exact pure-ray capacity]{The invariant readout optimum, including singular boundaries}
\label{sel54:capacity}

Keep the independent normalized target packet
$(b_{\rm tar}|0\rangle+d_{\rm tar}|1\rangle)/\sqrt2$, with target
characters $1,\chi^{-1}$ and flag characters $1,\chi$. The target is
$v=(b_{\rm tar}+d_{\rm tar})/\sqrt2$; this $b_{\rm tar}$ is distinct
from the scalar $b$ in the joint channel. All inaccessible target markers
are identical. Maximize only over invariant effects on the flag and the
whole retained reference. There is no operation on the target itself.

\begin{theorem}[Capacity throughout the full marginal-compatible fibre]
\label{sel54:thm:capacity}
The largest coefficient mass $\alpha(b)$ with selected output
$J_F=\alpha(b)P_v$ is
\begin{equation}
 \boxed{
 \alpha(b)=\begin{cases}
 (1-\Re b)/16,&|b-5/9|<4/9,\\[1mm]
 3(1-\Re b)/(13-5\Re b),&|b-5/9|=4/9,\ b\ne1,\\[1mm]
 0,&b=1.
 \end{cases}}
 \label{sel54:eq:capacity}
\end{equation}
Every optimum is attained. For this normalized preparation source, the
coefficient mass is its success probability, with all selection costs
included. In particular,
\begin{equation}
 \boxed{\max_{\mathcal E\text{ satisfying }\eqref{sel54:eq:marginals}}
 \alpha(\mathcal E)=\frac3{14},}
 \label{sel54:eq:max}
\end{equation}
attained only at $b=1/9$, which also forces $a=(-3+4i)/5$.
For every allowed $b$ there is one fixed invariant effect giving success
$(1-\Re b)/16$, irrespective of whether it is optimal. Hence one output
permits the exact ray if and only if $b\ne1$; if $b=1$, no finite number
of independent copies of that output permits it by the inherited covariant
protocol.
\end{theorem}
\begin{proof}
Write $\sigma_b=\left(\begin{smallmatrix}A&C\\C^*&B\end{smallmatrix}\right)$
in its two charge blocks. Only total readout charge one can contribute to
a pure balanced output with both target components. Contributions from
total charge zero or two occupy a single target ray and cannot be
cancelled in a positive sum. On the remaining sector, ordered as
$(|0\rangle W_j,|1\rangle N_j)$, the unwanted-target expectation is
one quarter the pairing with
\begin{equation}
 D_b=\begin{pmatrix}B&-C^*\\-C&A\end{pmatrix}\succeq0.
 \label{sel54:eq:bad}
\end{equation}
Therefore the maximal feasible effect is $P_{\ker D_b}$; success equals
$\frac12\Tr[P_{\ker D_b}\diag(B,A)]$. This uses actual readout effect
order, not an arbitrary suboperation below the system Choi matrix.

In the open disk the kernel has the single normalized vector
\[
 f_0=(|0\rangle W_-+|1\rangle N_-)/\sqrt2,\qquad
 N_-=(N_1-N_2)/\sqrt2,\quad W_-=(W_1-W_2)/\sqrt2.
\]
Its selected packet is $(1-x)\left(\begin{smallmatrix}1&1\\1&1\end{smallmatrix}\right)/32$.
The same effect remains feasible at the boundary, proving the universal
lower formula.

For a boundary point other than $1$, use the exhaustive parameterization
\[
 b_t=\frac{1+9t^2-8it}{9(1+t^2)},\qquad t\in\R.
\]
In symmetric/antisymmetric slot coordinates, the reference kernel is
$\{(T_ty,y):y\in\C^2\}$ with
\[
 T_t=\begin{pmatrix}(1-3i)/5&0\\-3it(2-i)/5&-1\end{pmatrix}.
\]
Thus the unwanted kernel is the range of $L_t=\binom{I}{-T_t}$.
Direct multiplication of \eqref{sel54:eq:output} gives
\begin{align*}
 G_t=L_t^*L_t&=\begin{pmatrix}
 (9t^2+7)/5&3t(1-2i)/5\\3t(1+2i)/5&2
 \end{pmatrix},\\
 R_t=\tfrac12L_t^*\diag(B,A)L_t
 &=\frac1{9(1+t^2)}
 \begin{pmatrix}3t^2+2&-it\\it&1\end{pmatrix}.
\end{align*}
The maximum is
$\Tr(G_t^{-1}R_t)=3/(14+9t^2)$, since
$\det G_t=(14+9t^2)/5$. Substitution of $x=\Re b_t$ proves the
second branch of \eqref{sel54:eq:capacity}. At $b=1$ the reference has
faithful invariant support by Proposition~\ref{sel54:prop:output}; its
all-finite-copy pure-asymmetry obstruction is inherited from V49.

The open-disk success is strictly below $1/18$. On the boundary,
$3(1-x)/(13-5x)$ is strictly decreasing for $1/9\le x<1$.
This proves the maximum, its unique $b$, and its forced $a$.
\end{proof}

The theorem optimizes the \emph{joint channel completion} and readout for
the fixed V53 preparation. It does not optimize all preparations, all
multi-query strategies, or full-native extensions of the old instrument.
It does not replace the product source's success $1/36$ by $3/14$.

\section[A fully normalized extremizer and the unchanged product point]{Exact benchmarks and a classical shared-phase realization}
\label{sel54:examples}

\begin{proposition}[An explicit extremal joint source and optimal effect]
\label{sel54:prop:extreme}
Let $\gamma=\bar z=(1-2i)/3$ and
\[
 \lambda_\pm=\gamma\left(1\pm\frac{2i}{\sqrt5}\right),\qquad
 V_\pm=P_0+\lambda_\pm P_1.
\]
Then $|\lambda_\pm|=1$, and the normalized shared-randomness map
\begin{equation}
 \boxed{\mathcal E_*=\tfrac12\operatorname{Ad}_{V_+\otimes V_-}
                    +\tfrac12\operatorname{Ad}_{V_-\otimes V_+}}
 \label{sel54:eq:extremechannel}
\end{equation}
has both local marginal channels exactly equal to $\Phi$. It realizes
$b=1/9$, $a=(-3+4i)/5$ and the maximum $3/14$.
Its prepared reference has eigenvalues $7/9,2/9,0,0$.
\end{proposition}
\begin{proof}
The phase mean is $(\lambda_++\lambda_-)/2=\gamma$, so either local
averaged channel agrees with $\Phi$ on the spectral off-diagonal blocks
and fixes the diagonal blocks. The mean relative phase gives $b=1/9$,
while $\lambda_+\lambda_-=\bar z/z$ gives the stated $a$.
This proves normalization and the exact local identities independently
of the correlation-matrix criterion.

In the slot basis $N_\pm=(N_1\pm N_2)/\sqrt2$,
$W_\pm=(W_1\pm W_2)/\sqrt2$, the reference is a direct sum of the
pure two-charge blocks
\[
 \sigma_+=\frac19\begin{pmatrix}5&-1+3i\\-1-3i&2\end{pmatrix},
 \qquad
 \sigma_-=\frac19\begin{pmatrix}1&1\\1&1\end{pmatrix}.
\]
The optimal effect is the sum of the orthogonal projectors onto
\[
 f_-=\frac{|0\rangle W_-+|1\rangle N_-}{\sqrt2},\qquad
 f_+=\frac{(-1-3i)|0\rangle W_++2|1\rangle N_+}{\sqrt{14}}.
\]
Both vectors have total character $\chi$, so the effect is invariant.
Their successes are $1/18$ and $10/63$, respectively. Both select the
same target ray, giving $J_F=\frac3{28}
\left(\begin{smallmatrix}1&1\\1&1\end{smallmatrix}\right)$ and
$h=9/784$. The positive complement restores the original target marginal.
\end{proof}

The construction needs no entangled noise environment or communication
between the two local operations: one shared classical bit suffices.
However $V_\pm$ are a different unitary realization of the forgotten local
channel. They are not original selectable pulses of the six-outcome
compiler. This is an explicit compatible joint implementation, not a
refinement or an inferred execution of the original labelled source.

\begin{proposition}[The inherited V53 points keep their values]
\label{sel54:prop:regression}
The independent product channel has
$(a,b)=(z^2,5/9)$ and success $1/36$. V53's correlated family is
\[
 a=-\frac13+i\left(\frac43-2t\right),\qquad
 b=2t-\frac13,\qquad \frac13\le t\le\frac23.
\]
It occupies only the real interval $1/3\le b\le1$, and
\[
 \alpha=\frac{2/3-t}{8}.
\]
Its maximum remains $1/24$, strictly below the full-fibre maximum $3/14$.
All these examples have the same complete local maps and the same local
prepared output with determinant $1/72$.
\end{proposition}
\begin{proof}
Evaluate the spectral coefficients of the displayed original mixtures.
Their $b$ values lie in the open disk except at $t=2/3$. Apply
Theorem~\ref{sel54:thm:capacity}; the endpoint is the invariant-support
zero. Marginal equality follows on all inputs from
Theorem~\ref{sel54:thm:fibre}, and its prepared-state specialization is
V53's original local density.
\end{proof}

\section[Boundary gains require exact joint information]{Stability, source constraints, and the historical next row}
\label{sel54:boundary}

\begin{proposition}[A compatible perturbation removes only the boundary bonus]
\label{sel54:prop:boundary}
Keep all local channels and the correlated preparation fixed. Define the
\emph{different} joint comparison
$\mathcal E_\epsilon=(1-\epsilon)\mathcal E_*+
\epsilon\Phi^{\otimes2}$, $0<\epsilon\le1$. Then
\begin{equation}
 \boxed{b_\epsilon=\frac{1+4\epsilon}{9},\qquad
 \alpha(\mathcal E_\epsilon)=\frac{2-\epsilon}{36},\qquad
 \|\mathcal E_\epsilon-\mathcal E_*\|_\diamond\le2\epsilon.}
 \label{sel54:eq:perturbation}
\end{equation}
The exact optimum has limit $1/18$ along this interior path, whereas
$\alpha(\mathcal E_*)=3/14$. The universal effect $P_{f_0}$ still selects
an exact ray at every point of the path.
\end{proposition}
\begin{proof}
The joint coefficient and channel depend affinely on the mixture. For
$0<\epsilon\le1$ the coefficient is in the open disk; substitution in
\eqref{sel54:eq:capacity} proves the success formula. The diamond estimate
uses the distance at most two between channels. The old kernel direction
is enforced by the exact local identities and original pure correlated
preparation, not by an approximate rank computation.
\end{proof}

The discontinuity concerns the \emph{optimal exact pure-output} condition,
not discontinuous channels or measured probabilities. A near-boundary
value does not prove the extra exact support kernel. On the other hand,
under independently verified exact marginal identities, a bound
$\Re b\le\widehat x+\varepsilon_x<1$ gives the robust one-sided
\emph{success} certificate
\[
 \alpha_{f_0}\ge\frac{1-\widehat x-\varepsilon_x}{16}>0.
\]
The real quadrature alone decides this one-copy existence question; the
imaginary quadrature is needed in general to distinguish boundary
optimality. For example $b=5/9$ and $b=5/9+4i/9$ have the same real
quadrature but optimal successes $1/36$ and $3/23$.

If the marginal identities themselves have only an error certificate, the
exact support conclusion cannot be retained. For a fixed preparation and
fixed admitted effect, a joint diamond error at most $\varepsilon$
perturbs the unnormalized target output by at most $\varepsilon$ in trace
norm. Relative to an ideal output $\alpha P_v$ with
$\varepsilon<\alpha$, its normalized target error is at most
$2\varepsilon/(\alpha-\varepsilon)$. This is an approximation statement,
not an exact-efficiency inference from small measured residuals.

\paragraph{Original labels and clock.}
All displayed successes are for two normalized uses of the accepted local
forward channel, with its source labels forgotten. A specified policy
requiring both raw opportunities to accept multiplies the corresponding
conditional success by its \emph{actual} joint acceptance probability;
under independent acceptances this is $\vartheta^2$. For the original
independent raw-opportunity source without that postselection, V53's
success $\vartheta(1-2\vartheta/3)/12$ remains unchanged. Fixed local
channel marginals alone identify neither that joint acceptance law nor an
elapsed-time instrument. No new comparison is assigned its clock by
suppressing those data.

\paragraph{Proved in this layer.}
Every CPTP extension of the fixed three-level local maps, with the
all-input marginal identities, is described by \eqref{sel54:eq:C}.
The relevant complex joint row lies in an exact disk, determines the
unchanged prepared output, and yields the exact readout maximum on every
rank stratum. One fixed invariant effect succeeds for every point other
than $b=1$. The global maximum $3/14$ is attained by a fully specified
shared-classical-phase implementation. V53's product and one-parameter
results are retained as strict subcases, not improved on the unchanged
product implementation.

\paragraph{Inherited and still conditional.}
The native determinant-sensitive action, correlated invariant preparation,
independent target source, target-character identification, identical
inaccessible target markers, and invariant joint effect access remain
inputs. They do not follow from marginal tomography. The complete
historical six-label bank and its full helper sector have not been
replaced by the alternative local unitary realization. Nonsignalling does
not assert a particular dilation, simultaneous physical availability, or
an independent-repetition policy.

\paragraph{Next historical calculation.}
On an independently verified same-cut port, check the \emph{complete local
channel identities} and evaluate the invariant two-register row
\eqref{sel54:eq:measuredrow} for the actual joint implementation. There is
no need to fit a two-mixture correlation model first. Conversely, absent
that row, neither $1/36$ nor $3/14$ may be assigned to the historical
source. The direct V110 typed-word route is separate; no selected
coefficient-state ray is identified with an executed microscopic Yukawa
word, fermionic grading, or finite-Dirac incidence. Historical
Standard-Model selection is not declared closed.

\paragraph{Verification boundary.}
Exact checks below compare the correlation-matrix construction with full
nine-dimensional maps, all local marginal identities on a matrix basis,
reference ranks, kernel-optimal effects, the boundary formula, and the
shared-phase realization. The V53 checker is rerun unchanged. General
classification and optimality rest on the written proofs; this is not an
independent re-audit of the entire preceding lineage, a physical
measurement, or proof-assistant certification.

\begingroup
\renewcommand{\refname}{Primary-source context for V54}
\begin{thebibliography}{9}
\small\raggedright
\bibitem{sel54-Beckman}
D.~Beckman, D.~Gottesman, M.~A.~Nielsen and J.~Preskill,
\emph{Causal and localizable quantum operations}, Phys. Rev. A
\textbf{64} (2001), 052309;
\url{https://arxiv.org/abs/quant-ph/0102043}.
\bibitem{sel54-Eggeling}
T.~Eggeling, D.~Schlingemann and R.~F.~Werner,
\emph{Semicausal operations are semilocalizable}, Europhys. Lett.
\textbf{57} (2002), 782--788;
\url{https://arxiv.org/abs/quant-ph/0104027}.
\bibitem{sel54-Puchala}
Z.~Pucha\l a, K.~Korzekwa, R.~Salazar, P.~Horodecki and K.~\.{Z}yczkowski,
\emph{Dephasing superchannels}, Phys. Rev. A \textbf{104} (2021), 052611;
\url{https://arxiv.org/abs/2107.06585}.
\end{thebibliography}
\endgroup
\addtocontents{toc}{\protect\endgroup}
% END OF SOURCE-SELECTION V54 DEVELOPMENT BLOCK.
% Preserve V1--V54 and append subsequent work before the final terminator.


\clearpage
\hypersetup{pdftitle={Historical Standard-Model Source Selection: V1--V55}}
\begin{center}
{\Large\bfseries Source-selection continuation V55}\\[.4em]
{\large Exchangeable source reuse, finite extension budgets,\\
and a joint-probability certificate for independent queries}
\end{center}
\addtocontents{toc}{\protect\begingroup}
\addtocontents{toc}{\protect\renewcommand{\protect\numberline}{\protect\makebox[2.1em][l]}}

\section[Joint-channel feasibility is not indefinitely exchangeable reuse]{The upstream question is consistency across numbers of source uses}
\label{sel55:context}

V54 classifies every two-register channel with the original local maps.
Its largest balanced-ray success, $3/14$, belongs to a different joint
implementation from the product source. A two-register completion need not
be the two-register marginal of one exchangeable source used an arbitrary
number of times. This layer tests that additional, independently meaningful
source condition. It does not alter the assigned product channel, optimize
another preparation, or impose exchangeability on the renewal grammar.

The local forward port, its correlated two-register preparation, and the
invariant target-readout class remain exactly those of V53--V54. The new
results give a sharp finite extension bound, classify the complete
infinitely exchangeable two-register fibre, and show when one measured
joint coefficient proves independence of an already verified exchangeable
hierarchy. A hierarchy is not inferred merely from the observation of
that coefficient. Chronologically ordered dynamics with memory can be
nonexchangeable and remain perfectly valid physical sources.

\paragraph{Inherited versus new.}
The normal form \eqref{sel54:eq:C}, the all-input marginal identities,
the output \eqref{sel54:eq:output}, and the capacity
\eqref{sel54:eq:capacity} are inherited, not new optimizations. The
quantum-operation de Finetti theorem is also an established input
\cite{sel55-Fuchs}; the reduction to the particular local source, the
exact moment region, finite extension witnesses, and capacity consequences
are proved below. No general novelty is claimed for exchangeability,
positive block matrices, or variance inequalities.

Keep $P_1=P_{(n+\omega)/\sqrt2}$, $P_0=I-P_1$ on $H_3$ and put
\begin{equation}
 \Phi_\zeta(X)=P_0XP_0+P_1XP_1+
       \zeta P_0XP_1+\bar\zeta P_1XP_0,\qquad |\zeta|\le1.
 \label{sel55:eq:local}
\end{equation}
These are CPTP channels, with the original source at
$z=(1+2i)/3$. Write $s=|z|^2=5/9$ and $r_0=1-s=4/9$.
The symbol $s$ here is not a stopping time or a selected probability.
When $|\zeta|=1$, $\Phi_\zeta=\operatorname{Ad}_{P_0+\bar\zeta P_1}$.
An alternative value of $\zeta$ in a comparison construction is not an
original selectable six-outcome source label.

\begin{definition}[Finite extension and a consistent exchangeable hierarchy]
\label{sel55:def:extension}
An $M$-register extension of the specified two-register channel $\E_2$
means a CPTP map $\E_M$ on $H_3^{\otimes M}$ which commutes with every
permutation of the registers and satisfies, for each pair $i,j$,
\[
 \Tr_{\{i,j\}^c}\E_M(X)=\E_2(\Tr_{\{i,j\}^c}X)
 \quad\text{for every joint operator }X.
\]
The natural relabelling of the retained pair is understood. The
one-register marginal is the original $\Phi_z$.
An infinitely exchangeable hierarchy is a family $(\E_k)_{k\ge1}$ of
permutation-covariant CPTP maps with $\E_1=\Phi_z$ and
\begin{equation}
 \Tr_{k+1}\E_{k+1}(X)=\E_k(\Tr_{k+1}X)
 \quad\text{for every }k\ge1\text{ and every }X.
 \label{sel55:eq:consistency}
\end{equation}
\end{definition}
The identities include entangled inputs and extend to untouched references.
A source can be symmetric for two uses without meeting either definition
at a larger number of uses. Repeating a special two-use block indefinitely
also need not produce an exchangeable hierarchy across all registers.
No timing record, joint acceptance law, or environment access is specified
by these channel identities.

\section[A finite positive minor bounds how many symmetric extensions exist]{The exact finite extension budget}
\label{sel55:finite}

\begin{lemma}[Spectral-sector normal form at any finite number of registers]
\label{sel55:lem:normal}
A CPTP map with the stated one-register marginals has the form
\[
 \E_M(X)=\sum_{\boldsymbol a,\boldsymbol b\in\{0,1\}^M}
 C^{(M)}_{\boldsymbol a,\boldsymbol b}
 \Pi_{\boldsymbol a}X\Pi_{\boldsymbol b},\qquad
 \Pi_{\boldsymbol a}=\bigotimes_jP_{a_j},
\]
where $C^{(M)}\succeq0$ and all diagonal entries are one. If the
pair marginals equal $\E_2=C(a,b)$ of \eqref{sel54:eq:C}, then the
principal block indexed by $0^M,e_1,\ldots,e_M$ is
\begin{equation}
 G_M=\begin{pmatrix}
 1&z\mathbf1^{\mathsf T}\\
 \bar z\mathbf1&(1-b)I_M+b\mathbf1\mathbf1^{\mathsf T}
 \end{pmatrix}.
 \label{sel55:eq:principal}
\end{equation}
For a permutation-covariant extension, $b$ is real.
\end{lemma}
\begin{proof}
The dual local channel fixes pointwise
$\mathcal D=\B(P_0H_3)\oplus\B(P_1H_3)$. The marginal identities
make the dual joint channel fix every embedded copy of $\mathcal D$.
If $A$ is a unitary in one such copy and $K_\ell$ are joint Kraus
coefficients, fixedness of $A,A^*A$ gives
\[
 \sum_\ell(AK_\ell-K_\ell A)^*(AK_\ell-K_\ell A)=0.
\]
Thus every $K_\ell$ commutes with all those local algebras. Their joint
commutant is the span of the orthogonal projectors $\Pi_{\boldsymbol a}$.
Expanding the coefficients in that basis gives a positive coefficient Gram
$C^{(M)}$; trace preservation gives unit diagonal. Conversely any such
matrix defines a CPTP map. All projectors have nonzero rank, so this
representation is faithful. On the displayed principal block, a single
changed bit is determined by $\Phi_z$, and two opposite changed bits are
determined by the $01,10$ entry of $\E_2$. Exchange of those two sites
sends $b$ to $\bar b$ and proves reality.
\end{proof}

\begin{theorem}[Sharp projected finite-extension interval]
\label{sel55:thm:finite}
For every $M\ge2$, the coefficient $b$ of an $M$-register extension
necessarily satisfies
\begin{equation}
 \boxed{b_M^{\min}:=s-\frac{1-s}{M-1}
       =\frac{5M-9}{9(M-1)}\ \le b\le1.}
 \label{sel55:eq:finitebound}
\end{equation}
Conversely, every real $b$ in this interval occurs for some compatible
$a$ and some permutation-covariant $M$-register channel satisfying all
pair marginal identities. The statement projects out $a$: it does not
assert that every $a$ allowed by V54 extends at that $b$.
\end{theorem}
\begin{proof}
The Schur complement of the first entry of \eqref{sel55:eq:principal}
is $(1-b)I_M+(b-s)\mathbf1\mathbf1^{\mathsf T}$. Its eigenvalues
are $1-b$ on $\mathbf1^\perp$ and $1+(M-1)b-Ms$ on
$\C\mathbf1$, giving the necessity.

For sufficiency at the lower endpoint, choose unit complex numbers
$\zeta_1,\ldots,\zeta_M$ with sum $Mz$. An explicit choice is as
follows. Put $u=z/|z|$. For even $M$, take equally many values
$u(|z|+i\sqrt{1-s})$ and $u(|z|-i\sqrt{1-s})$. For odd $M$, take
one $u$ and equally many values $u(c_M\pm i\sqrt{1-c_M^2})$, where
$c_M=(M|z|-1)/(M-1)\in[-1,1]$. Uniformly permute this fixed list
between the $M$ registers and apply the corresponding product unitary
channels. The resulting map is normalized, permutation covariant, and
nonsignalling on every subset. Its local coefficient is $z$, and
\[
 b=\frac1{M(M-1)}\sum_{i\ne j}\zeta_i\bar\zeta_j
   =\frac{M^2s-M}{M(M-1)}=b_M^{\min}.
\]
At the upper endpoint use one common unit phase at all sites, taking
$1$ with probability $1/3$ and $i$ with probability $2/3$. Its local
coefficient is also $z$, and $b=1$. Convex mixtures of these two
$M$-register maps attain the entire interval. Their value of $a$ is the
actual average of $\zeta_i\zeta_j$, not an independently assigned row.
\end{proof}

\begin{corollary}[A finite witness against further reuse]
\label{sel55:cor:budget}
If $b<s$, every extension satisfies
\begin{equation}
 \boxed{M\le\frac{1-b}{s-b}.}
 \label{sel55:eq:budget}
\end{equation}
In particular V54's maximizing channel with $b=1/9$ has no
three-register extension. The vector $v=(-3z,1,1,1)^{\mathsf T}$
in the required $G_3$ gives
\[
 v^*G_3v=-4/3,\qquad \|v\|^2=8,
\]
an explicit negative quadratic form rather than a failure to find an
extension. The original V53 classical endpoint $t=1/3$, with
$b=1/3$ and $a=(-1+2i)/3$, extends to exactly three registers and not
four: uniformly place one identity and two original $\operatorname{Ad}_U$
channels in the three sites.
\end{corollary}
\begin{proof}
Rearrange \eqref{sel55:eq:finitebound}. The displayed quadratic form is
$3[1+2b-3s]$. For the three-register construction the spectral phases
are $(1,i,i)$, so their pair moments are precisely the stated $a,b$;
\eqref{sel55:eq:budget} gives the upper limit of three.
\end{proof}
These comparison extensions require no entangled shared environment:
classical permutation of product unitary channels suffices to attain
the bound. This does not identify their hidden phase labels with the
original six recorded outcomes, or supply a joint law for those records.

\section[The exact success maximum decreases with exchangeable extension depth]{Capacity on the same two-register preparation}
\label{sel55:capacity}

The correlated invariant input, independent balanced target packet,
identical inaccessible target markers, and allowed invariant readout
effects remain those of V54. Only the source-reuse class is restricted.
Let $A_M$ be the largest exact balanced-ray success of its pair marginal
among $M$-register extensions.

\begin{theorem}[Sharp capacity under finite symmetric extension]
\label{sel55:thm:capacity}
One has
\begin{equation}
 \boxed{A_2=\frac3{14},\qquad
 A_M=\frac{M}{36(M-1)}\quad(M\ge3).}
 \label{sel55:eq:capacities}
\end{equation}
The bounds are attained. For $M\ge3$, the single V54 invariant effect
$f_0$ attains the bound at $b=b_M^{\min}$; no boundary-optimized effect
or new preparation is needed.
\end{theorem}
\begin{proof}
For $M=2$, permutation symmetry forces only $b$ real; V54's maximizer
already has real $b=1/9$ and is symmetric. For $M\ge3$, the interval
\eqref{sel55:eq:finitebound} lies strictly inside the real diameter of
V54's disk except at $b=1$. Thus V54's exact capacity is
$(1-b)/16$ at every point of this interval, with zero at $b=1$.
It is largest at the lower endpoint. The constructions of
Theorem~\ref{sel55:thm:finite} attain that endpoint, giving
$(1-s)M/[16(M-1)]$. V54 already proves that $f_0$ is optimal at
these interior points.
\end{proof}
For example, $A_3=1/24$, $A_4=1/27$, and $A_6=1/30$.
The sharp two-register value $3/14$ is not a capacity of a source that
can be extended symmetrically even once more. This is an additional
source-consistency test, not a correction to V54's two-register theorem.
Repeating the special pair implementation in distinct blocks remains
possible but does not have the stated permutation symmetry across blocks.

\section[An indefinitely exchangeable source is a mixture of identical dephasings]{The complete infinitely exchangeable marginal fibre}
\label{sel55:infinite}

\begin{theorem}[Source-specific de Finetti representation]
\label{sel55:thm:representation}
Every hierarchy in Definition~\ref{sel55:def:extension} has a unique
probability measure $\mu$ on the closed complex unit disk such that
\begin{equation}
 \boxed{\E_k=\int\Phi_\zeta^{\otimes k}\,d\mu(\zeta),\qquad
        \int\zeta\,d\mu=z.}
 \label{sel55:eq:definetti}
\end{equation}
Conversely every such measure defines a hierarchy. For its two-register
parameters,
\begin{equation}
 \boxed{a=\int\zeta^2d\mu,\qquad b=\int|\zeta|^2d\mu.}
 \label{sel55:eq:moments}
\end{equation}
\end{theorem}
\begin{proof}
Apply the established de Finetti theorem for finite-dimensional quantum
operations \cite[Definition 1 and Theorem 1]{sel55-Fuchs} to the exact
permutation and all-input discard identities. It gives a unique measure
over CPTP maps $\Lambda$ with
$\E_k=\int\Lambda^{\otimes k}d\nu(\Lambda)$ and
$\Phi_z=\int\Lambda d\nu$. This external theorem is the general
representation input, not a result claimed anew here.

The local Choi support is exactly
$\operatorname{vec}\Span\{P_0,P_1\}$, since $|z|<1$.
The integral of the nonnegative Choi mass on its orthogonal complement
is zero. Hence almost every $J_\Lambda$ has that same containing support.
Its Kraus coefficients therefore have the form $u_jP_0+v_jP_1$.
Trace preservation gives $\sum|u_j|^2=\sum|v_j|^2=1$, while
$\zeta=\sum u_j\bar v_j$ has modulus at most one. Thus
$\Lambda=\Phi_\zeta$. The map $\zeta\mapsto\Phi_\zeta$ is
injective, giving uniqueness of $\mu$ and the mean identity. Reading the
$00,11$ and $01,10$ coefficients of the tensor product gives
\eqref{sel55:eq:moments}. The converse follows directly from products
of normalized local channels and integration.
\end{proof}
This is an exact representation of the joint channel hierarchy. It does
not prove that its implementation exposes a physical classical label
$\zeta$, nor admit measurements conditioned on that label. Independence
conditional on a mathematical mixing variable is not unconditional
independence of the original source uses.

\begin{theorem}[Complete extendible two-register moment region]
\label{sel55:thm:momentregion}
A V54 pair channel belongs to an infinitely exchangeable hierarchy if and
only if
\begin{equation}
 \boxed{b\in[s,1]\subset\R,\qquad
        |a-z^2|\le b-s.}
 \label{sel55:eq:region}
\end{equation}
Every point has an explicit extension represented by a measure with at
most five atoms on the closed unit disk.
\end{theorem}
\begin{proof}
Put $\eta=\zeta-z$, $\delta=b-s$, and $\xi=a-z^2$.
The necessity follows from
$\delta=\int|\eta|^2d\mu\ge0$,
$|\xi|=|\int\eta^2d\mu|\le\delta$, and $b\le1$.
For sufficiency form the real covariance matrix
\[
 V=\frac12\begin{pmatrix}
 \delta+\Re\xi&\Im\xi\\
 \Im\xi&\delta-\Re\xi
 \end{pmatrix}\succeq0,
 \qquad\Tr V=\delta\le r_0.
\]
Write its spectral decomposition as $\sum_{j=1}^2\lambda_j u_ju_j^{\mathsf T}$,
identifying each real unit vector $u_j$ with a unit complex number.
The chord through $z$ in direction $u_j$ meets the unit circle at
$z+t_{j\pm}u_j$, where
\[
 t_{j\pm}=-c_j\pm\sqrt{c_j^2+r_0},\qquad
 c_j=\Re(\bar zu_j).
\]
Since $r_0>0$, these roots have opposite signs. Give the two endpoints
probabilities $-t_{j-}/(t_{j+}-t_{j-})$ and
$t_{j+}/(t_{j+}-t_{j-})$. This chord measure has mean $z$ and
covariance $r_0u_ju_j^{\mathsf T}$, since
$-t_{j+}t_{j-}=r_0$. Mix the two chord measures with weights
$\lambda_j/r_0$ and put the remaining mass $1-\delta/r_0$ at $z$.
The resulting measure has the required first and second moments and at
most five atoms. Theorem~\ref{sel55:thm:representation} completes the
extension.
\end{proof}
The construction is not a claim that five atoms determine the original
infinite hierarchy. They provide one compatible hierarchy for a specified
pair. Different measures can have the same first two moments.

\begin{corollary}[The independent source is the unique indefinitely exchangeable maximizer]
\label{sel55:cor:infcap}
For this unchanged correlated preparation and readout class,
\begin{equation}
 \boxed{\alpha(\E_2)=\frac{1-b}{16}\le\frac1{36}.}
 \label{sel55:eq:infcap}
\end{equation}
Equality is attained precisely when the entire hierarchy is
$\E_k=\Phi_z^{\otimes k}$ for every $k$. At $b=1$ the capacity is
zero. Every intermediate value occurs in an explicit consistent hierarchy.
\end{corollary}
\begin{proof}
Apply V54's capacity on the real interval $[s,1]$. Equality requires
$b=s$, which in \eqref{sel55:eq:moments} means
$\int|\zeta-z|^2d\mu=0$ and hence $\mu=\delta_z$.
To attain all intermediate values, take
\[
 \mu_t=(1-t)\delta_z+t\bigl(\tfrac13\delta_1+\tfrac23\delta_i\bigr),
 \qquad 0\le t\le1.
\]
It has $b=s+t(1-s)$ and capacity $(1-t)/36$.
\end{proof}
Thus the original product value $1/36$ is unchanged and becomes the sharp
maximum under indefinite exchangeability. It is not inferred from the
one-register marginal alone. The variables $\zeta$ used for these
comparison hierarchies are not new selectable source operations.

\section[One joint scalar certifies independence under the verified hierarchy]{Quantitative reconstruction and the limit of pair data}
\label{sel55:independence}

\begin{proposition}[Exact distance from the independent pair]
\label{sel55:prop:distance}
On the infinitely exchangeable class, let $\delta=b-s\ge0$.
Then
\begin{equation}
 \boxed{\|\E_2-\Phi_z^{\otimes2}\|_\diamond=\delta,
 \qquad\alpha(\E_2)=\frac1{36}-\frac{\delta}{16}.}
 \label{sel55:eq:pairdistance}
\end{equation}
At every higher order,
\begin{equation}
 \boxed{\|\E_k-\Phi_z^{\otimes k}\|_\diamond
 \le\min\left\{2,\ k\sqrt\delta,\ \binom{k}{2}\delta\right\}.}
 \label{sel55:eq:higherdistance}
\end{equation}
Diamond norms use no factor of $1/2$.
\end{proposition}
\begin{proof}
In the two-register correlation matrix the only differences from the
product map are $\xi=a-z^2$ between $00,11$ and $\delta$ between
$01,10$. Pinch into the two orthogonal sums of those sectors. On each,
the difference is a phase times its off-diagonal projection, multiplied
by $\xi$ or $\delta$. The off-diagonal projection is
$(\id-\operatorname{Ad}_{P_+-P_-})/2$ and has diamond norm at most
one. Thus the total norm is at most $\max\{|\xi|,\delta\}=\delta$.
A balanced input superposition of unit vectors in $01$ and $10$ gives
a traceless two-by-two difference with eigenvalues $\pm\delta/2$,
attaining the bound. The capacity follows from
\eqref{sel55:eq:infcap}.

For higher orders, $\|\Phi_\zeta-\Phi_z\|_\diamond=|\zeta-z|$
by the same two-sector argument. Tensor telescoping and Cauchy--Schwarz
give $k\sqrt\delta$. For the sharper quadratic estimate put
$F_\eta(t)=\Phi_{z+t\eta}^{\otimes k}$. The line segment remains in
the unit disk, so every nondifferentiated factor is CPTP of diamond norm
one. Hence
$\|F_\eta''(t)\|_\diamond\le k(k-1)|\eta|^2$.
The average of $F_\eta'(0)$ is zero, since $\int\eta d\mu=0$.
Taylor's integral formula and integration over $\mu$ give
$\binom{k}{2}\delta$. Finally the distance of two channels is at most two.
\end{proof}
The exact independence certificate is therefore $b=5/9$, \emph{after}
the consistent exchangeable hierarchy has been established. At known
$z$, an uncertainty $\delta\le\epsilon$ yields
$\|\E_k-\Phi_z^{\otimes k}\|_\diamond\le\binom{k}{2}\epsilon$.
This is an all-input operational estimate, not a noise-model fit or a
physical mass gap. A finite table cannot itself certify an infinite
extension hypothesis.

\begin{proposition}[Pair tomography does not generally determine the hierarchy]
\label{sel55:prop:third}
Let $r=1/6$, and take either the three or four equally weighted complex
numbers
\[
 \zeta_j=z+r e^{2\pi i j/q},\qquad j=0,\ldots,q-1,
 \qquad q=3\text{ or }4.
\]
Both sets lie strictly inside the unit disk. The two resulting hierarchies
have identical one- and two-register channels:
\[
 a=z^2,\qquad b=s+r^2=7/12,\qquad\alpha=5/192.
\]
Their three-register channels nevertheless satisfy
\begin{equation}
 \boxed{\|\E_3^{(3)}-\E_3^{(4)}\|_\diamond=r^3=1/216.}
 \label{sel55:eq:third}
\end{equation}
\end{proposition}
\begin{proof}
Here $|z|+r<1$. Both phase laws have zero first and second centered
complex moments and centered squared modulus $r^2$. The cubic law has
third centered moment $r^3$; the fourth-root law has zero third centered
moment. Mixed centered moments of total degree at most three agree.
Thus the only differing three-register coefficients are $000,111$ and
its adjoint, by $r^3$. The same pinching and balanced-input proof as in
Proposition~\ref{sel55:prop:distance} gives the exact diamond distance.
\end{proof}
The $b=5/9$ independence certificate is special because it forces zero
mixing variance. Away from it, even complete two-register tomography does
not identify higher process correlations. This example uses distinct
comparison source hierarchies; it does not modify the product benchmark.

\section[The original invariant tester now checks a reuse assumption]{Source-level tests, retained costs, and proof status}
\label{sel55:status}

The actual invariant observables of Corollary~\ref{sel54:cor:testers}
still give
\[
 b=4\langle X_W\rangle-\frac13+4i\langle Y_W\rangle.
\]
Exchangeability requires $\langle Y_W\rangle=0$. A proposed $M$-register
extension requires $b\ge b_M^{\min}$. If an estimate has certified
error $|\widehat b-b|\le\epsilon_b$ and the original $z$ is exact,
\[
 \Re\widehat b+\epsilon_b<b_M^{\min}
\]
is a one-sided exclusion certificate. It does not turn a small imaginary
part into proof of symmetry. Under an already justified infinite
hierarchy, $b=5/9$ is equivalent to
$\langle X_W\rangle=2/9$ and proves all-order independence; a positive
upper bound on $b-5/9$ gives \eqref{sel55:eq:higherdistance}.
These invariant expectations still require actual Reads or their admitted
contextual reconstruction. Their symmetry does not guarantee executability.

\paragraph{Proved in this layer.}
The finite extension bound holds for every positive exchangeable extension,
including quantum implementations, and is sharp after projecting out $a$.
The exact pair success is at most $M/[36(M-1)]$ for $M\ge3$ and $1/36$
for an infinite hierarchy. The full infinite pair region has a constructive
five-atom extension. Conditional on that hierarchy, one joint scalar
certifies all-order independence or supplies the quantitative bounds
\eqref{sel55:eq:pairdistance}--\eqref{sel55:eq:higherdistance}. Away from
the independence point, pair-identical higher-order alternatives remain.

\paragraph{Inherited inputs and historical limits.}
The determinant-sensitive action, forward port, correlated invariant
preparation, target-character assignment, identical inaccessible markers,
and invariant effect access remain V19/V51/V53/V54 inputs. The de Finetti
theorem does not establish its physical exchangeability hypothesis. These
results concern forgotten channels: they supply no joint six-label
instrument, acceptance law, elapsed-time record, or full-native extension.
$M$ counts exchangeable registers, not ticks or generations.

The $3/14$ channel remains a valid designated pair operation; it simply
cannot have the specified third exchangeable extension. Exchangeability
is not an axiom of finite duality or of ordered renewal dynamics.
Before another readout optimization, identify whether the actual reuse
protocol is independent, indefinitely exchangeable, exchangeable only
through a fixed register count, or nonexchangeable. The joint coefficient
$b$ constrains that choice but does not establish the reuse grammar.
No historical V110 word, grading, or directed finite-Dirac incidence is
populated here. Historical Standard-Model selection remains unclosed.

\paragraph{Verification boundary.}
The checker evaluates finite Gram obstructions, explicit extension maps,
all-input marginal identities, moment constructions, and the pair-identical
triple distinction, and reruns the unchanged immediate-parent checker.
General hierarchy assertions rest on the written proofs. These are not
physical measurements, proof-assistant formalization, or an independent
re-audit of the complete preceding manuscript.

\begingroup
\renewcommand{\refname}{Primary-source context for V55}
\begin{thebibliography}{9}
\small\raggedright
\bibitem{sel55-Fuchs}
C.~A.~Fuchs, R.~Schack and P.~F.~Scudo,
\emph{De Finetti representation theorem for quantum-process tomography},
Phys. Rev. A \textbf{69} (2004), 062305;
\url{https://arxiv.org/abs/quant-ph/0307198}.
The exact quantum-operation theorem is used at
Theorem~\ref{sel55:thm:representation}; its hypotheses are stated in
Definition~\ref{sel55:def:extension}.
\end{thebibliography}
\endgroup
\addtocontents{toc}{\protect\endgroup}
% END OF SOURCE-SELECTION V55 DEVELOPMENT BLOCK.
% Preserve V1--V55 and append subsequent work before the final terminator.


% BEGIN IMMUTABLE SOURCE-SELECTION V56 BODY
\clearpage
\hypersetup{pdftitle={Historical Standard-Model Source Selection: V1--V56}}
\begin{center}
{\Large\bfseries Source-selection continuation V56}\\[.4em]
{\large Causal reuse instead of exchangeability:\\
finite memory, graph-compatible pair saturation, and retained clock records}
\end{center}
\addtocontents{toc}{\protect\begingroup}
\addtocontents{toc}{\protect\renewcommand{\protect\numberline}{\protect\makebox[2.1em][l]}}

\section[Companion audit and the actual source-reuse question]{A chronological source need not be exchangeable}
\label{sel56:context}

V55 proves a sharp bound under exchangeable source reuse. It expressly does
not impose exchangeability on an ordered renewal grammar. The four newly
supplied companion sources make this distinction worth resolving before
another preparation or readout optimization. We retain the V53 correlated
preparation, the V54 all-input marginal condition and invariant-readout
capacity, and the V55 extension obstruction. We now ask which of those
pair channels admit an indefinitely consistent \emph{ordered} implementation,
with an explicitly counted causal memory and without future-dependent
preparation. The answer differs from exchangeability.

\paragraph{Audited companion interfaces.}
The following are source-derived statements, not theorems claimed anew here.
The audit file gives exact source-line ranges and SHA-256 fingerprints.
\begin{center}\small
\begin{tabular}{L{.24\textwidth}L{.67\textwidth}}
\toprule
Supplied source & Relevant result and retained boundary\\\midrule
Native stationarity V54\cite{sel56-Native} & The time-dependent covector response retains a
rank-two differentiated-source term and pays both fast scales. Its growing-mode
correction uses terminal data on the prescribed interval. The source explicitly
does not identify that prepared solution with a causal decoder or a nonlinear
original trajectory.\\[.3em]
Exact-action bridge V45\cite{sel56-Einstein} & The weak Poisson nullity and two secular chains are
certified; a positive rank-four intermediate pencil yields two oscillatory
pairs. A next-order skew coefficient and nonlinear propagation remain separate.
The established physical interval remains of order $a^2$.\\[.3em]
Perturbative ESM V64\cite{sel56-ESM} & A single counterterm action fixes its next-order insertion
and marked restrictions. A finite Hessian-integrability test retains a failed-input
residual. The populated critical common action and complete two-loop local
master/reference equation are not supplied.\\[.3em]
Yang--Mills V100 f17\cite{sel56-YM} & An exact nonzero resonant crossing obstructs a whole-head
electric inverse. A local resonance-preserving conjugation retains a spatially
summable remainder in its stated strong-electric regime. Neither the comparison
gap nor that conjugation proves the physical continuum gap.\\\bottomrule
\end{tabular}
\end{center}
These sources motivate retaining the complete consistency and causality
conditions, rather than completing each marginal independently. They do not
supply a quantum controller, its stationary law, or the historical determinant
operation. No descriptor variable, BV coefficient, or electric resonance is
identified with the two-state memory constructed below. The new constructions
are normalized comparison sources, not changes installed in any companion.

\paragraph{Inherited local interface.}
On the original forward port $H_3$, retain
\[
 P_1=P_{(n+\omega)/\sqrt2},\quad P_0=I-P_1,\quad
 z=(1+2i)/3,\qquad
 \Phi_z(X)=\sum_{a=0}^1P_aXP_a+zP_0XP_1+\bar zP_1XP_0.
\]
Put $s=|z|^2=5/9$ and $d=\dim H_3=3$. The designated extremal pair
$\E_*$ of Proposition~\ref{sel54:prop:extreme} uses
\begin{equation}
 \lambda_\pm=\bar z(1\pm2i/\sqrt5),\qquad
 V_\pm=P_0+\lambda_\pm P_1,\qquad
 \E_*=\tfrac12\operatorname{Ad}_{V_+\otimes V_-}
       +\tfrac12\operatorname{Ad}_{V_-\otimes V_+}.
 \label{sel56:eq:phases}
\end{equation}
Its pair parameters are $b_*=1/9$, $a_*=(-3+4i)/5$, and its optimal
balanced-ray success for the inherited preparation/readout task is $3/14$.
The $V_\pm$ are the already declared alternative unitary realization of the
\emph{forgotten} local channel. They are not the original six resolved
coefficients or available selectable pulses. Quantum memory channels are
established machinery \cite{sel56-KW,sel56-BM}; the specialized graph rigidity,
lag capacities and original-clock comparison are proved below.

\section[A two-state source realizes consistent ordered reuse]{An explicit causal family with unchanged local channels}
\label{sel56:causal}

Let $-1\le r\le1$ and let $S_j\in\{+1,-1\}$ be the stationary Markov
chain with uniform initial law and
\begin{equation}
 Q_r(t\mid s)=\frac{1+rst}{2}.
 \label{sel56:eq:kernel}
\end{equation}
Use $V_s$ on register $j$ and then update the memory to $t$. On the
classical memory embedded in $\C^2$ this is the CPTP map with coefficients
$\sqrt{Q_r(t\mid s)}V_s\otimes|t\rangle\langle s|$.
No input state or future outcome enters \eqref{sel56:eq:kernel}.

\begin{theorem}[Stationary causal implementation and all lagged pair maps]
\label{sel56:thm:memory}
The channels
\begin{equation}
 \E_k^{(r)}(X)=\sum_{s_1,\ldots,s_k}
 \frac12\prod_{j=1}^{k-1}Q_r(s_{j+1}\mid s_j)
 \operatorname{Ad}_{V_{s_1}\otimes\cdots\otimes V_{s_k}}(X)
 \label{sel56:eq:Ek}
\end{equation}
are normalized, stationary under shifts of the slot indices, and consistent
under discarding a final output. Their one-register marginals are $\Phi_z$
for all joint inputs, including reference-entangled inputs. The pair at
index separation $\ell\ge1$ is the V54 channel $C(a_\ell,b_\ell)$ with
\begin{equation}
 \boxed{b_\ell=\frac59+\frac49 r^\ell,\qquad
 a_\ell=z^2\left(1-\frac45r^\ell\right).}
 \label{sel56:eq:lag}
\end{equation}
For $|r|<1$ its $k$-register Choi rank is $2^k$; at $r=\pm1$ it is two
for every $k\ge1$. For $r\ne0$, the minimum persistent Hilbert dimension
of a causal realization is two, counting shared classical information as
memory. At $r=0$ the hierarchy is exactly $\Phi_z^{\otimes k}$.
\end{theorem}
\begin{proof}
The displayed coefficients obey $\sum_{s,t}K_{ts}^*K_{ts}=I\otimes I_2$.
Their repeated execution gives \eqref{sel56:eq:Ek}. Normalization and the
all-input marginal identities follow term by term from unitary conjugation
and the marginal probability law. Stationarity and prefix consistency are
those of the stationary chain. Tracing an intervening slot retains its
memory transition: pairs at separation $\ell$ use $Q_r^\ell$, not $Q_r$.
Set $\zeta_s=\bar\lambda_s=z(1-2is/\sqrt5)$. Then
$\mathbb E S_j=0$ and $\mathbb E S_jS_{j+\ell}=r^\ell$.
The two relevant correlation entries are
$\mathbb E(\zeta_j\bar\zeta_{j+\ell})$ and
$\mathbb E(\zeta_j\zeta_{j+\ell})$, giving \eqref{sel56:eq:lag}.

$V_+,V_-$ are linearly independent. Their $2^k$ tensor products are
therefore linearly independent. All their probabilities are positive when
$|r|<1$, proving the rank. At either endpoint exactly two distinct tensor
products have positive probability. A one-dimensional persistent memory,
with all other shared information included in that count, gives a product
of the single-use channels. Since $b_1\ne5/9$ for $r\ne0$, this is impossible.
The explicit two-state construction attains the lower bound. At $r=0$ the
independent signs give the original product channel.
\end{proof}
At $r=-1$ the controller alternates deterministically. Every adjacent pair
is exactly $\E_*$, at every depth. Pairs at even separation instead have
$b_\ell=1$; hence this is not an exchangeable extension of $\E_*$. The
three-register contradiction of V55 remains valid. There is no causal
contradiction: all memory is supplied by one initially sampled bit and the
fixed forward update. The process is stationary but not mixing in the
accepted-slot index. No future boundary condition is used.

\section[Pair support determines the entire extremal reuse graph]{Bipartiteness is necessary and sufficient, even for quantum extensions}
\label{sel56:graph}

A pair-marginal identity on an edge $ij$ means
$\Tr_{\{i,j\}^c}\E(X)=\E_*(\Tr_{\{i,j\}^c}X)$ for every joint $X$,
not only for the correlated preparation. The local channel is $\Phi_z$.

\begin{theorem}[Exact compatibility and uniqueness on a connected reuse graph]
\label{sel56:thm:graph}
Let $G$ be a finite connected simple graph with at least two vertices.
A CPTP map on $H_3^{\otimes |G|}$ with pair marginal $\E_*$ on every
edge exists if and only if $G$ is bipartite. On the bipartite branch it is
unique. For either two-colouring $c_j\in\{+1,-1\}$, it equals
\begin{equation}
 \boxed{\E_G=\tfrac12\operatorname{Ad}_{\bigotimes_jV_{c_j}}
                   +\tfrac12\operatorname{Ad}_{\bigotimes_jV_{-c_j}}.}
 \label{sel56:eq:graphchannel}
\end{equation}
Pairs in opposite parts have channel $\E_*$; pairs in the same part have
$\tfrac12\operatorname{Ad}_{V_+\otimes V_+}
 +\tfrac12\operatorname{Ad}_{V_-\otimes V_-}$, with $b=1$ and $a=z^2/5$.
The conclusion allows arbitrary quantum joint implementations; shared
classical randomness is an output of the rigidity statement, not an assumption.
\end{theorem}
\begin{proof}
Use the normalized Choi state $\rho_G=J_{\E}/d^{|G|}$. The all-input
marginal identities imply that its reduced state on each paired input/output
site is the corresponding normalized local or edge Choi state. Positivity
implies the support rule
\[
 \supp\Tr_B\rho\subseteq L\quad\Longrightarrow\quad
 \supp\rho\subseteq L\otimes B,
\]
since $\Tr[(P_{L^\perp}\otimes I)\rho]=0$.
The local support is $\mathscr S=\Span\{|V_+\rangle\!\rangle,
|V_-\rangle\!\rangle\}$ and the edge support is
\[
 \mathscr T=\Span\{|V_+\rangle\!\rangle\otimes|V_-\rangle\!\rangle,
                    |V_-\rangle\!\rangle\otimes|V_+\rangle\!\rangle\}.
\]
An invertible coordinate map on $\mathscr S$ sends these two independent
local vectors to the two coordinate vectors. It is a proof coordinate
change, not a physical filter. The intersection of the edge support
constraints is thereby the coordinate span of all sign strings satisfying
$s_i=-s_j$ on every edge. An odd cycle has no such string and would force
$J_{\E}=0$, contradicting trace preservation. A connected bipartite graph
has exactly the two strings in \eqref{sel56:eq:graphchannel}.

Let their unitary coefficients be $K_+,K_-$. Hence
$J_{\E}=\sum_{a,b=\pm}h_{ab}|K_a\rangle\!\rangle
\langle\!\langle K_b|$ for a positive $2\times2$ matrix $h$.
Set
\[
 q=\bar\lambda_-\lambda_+=\frac{1+4i\sqrt5}{9}.
\]
The spectrum of $K_-^*K_+$ contains $1,q,\bar q$: choose all local
spectral indices zero, then one index in each part equal to one.
The trace-preserving equation is
\[
 (h_{++}+h_{--})I+h_{+-}K_-^*K_++\bar h_{+-}K_+^*K_-=I.
\]
Its three scalar equations on $1,q,\bar q$ force $h_{+-}=0$ and
$h_{++}+h_{--}=1$, since $q$ is nonreal and $\Re q\ne1$.
The one-register marginal has mean coefficient $z$ only when
$h_{++}=h_{--}=1/2$. This proves uniqueness. Conversely the displayed
mixture is CPTP and has every required all-input edge marginal. Its other
pair marginals follow by tracing unitary factors.
\end{proof}

\begin{corollary}[Extremal neighbouring pairs force a persistent ordered source]
\label{sel56:cor:chain}
If every neighbouring pair in an indefinitely consistent ordered hierarchy
is $\E_*$, every finite interval is \eqref{sel56:eq:Ek} with $r=-1$.
Its next-neighbour pair has $b_2=1$, not $b_1=1/9$.
Every designated adjacent pair retains the inherited maximum success $3/14$;
even-separated pairs have zero exact capacity in that same preparation task.
\end{corollary}
\begin{proof}
Apply the theorem to each path. Its unique channels have consistent interval
restrictions. The pair parameters and capacities are the stated endpoint
values of \eqref{sel56:eq:lag} and Theorem~\ref{sel54:thm:capacity}.
\end{proof}
This does not assert independence of simultaneous successful readouts on
overlapping pairs. The same physical register cannot be consumed twice by
an independent-pair preparation. The statement concerns designated pair
marginals and their separately specified experiments. Likewise a reuse graph
is not identified with the spacetime lattice or its tetrahedral source graph.
A triangle (and therefore an all-pairs three-register extension) is excluded;
a chain or any connected bipartite reuse graph is not.

\section[Odd cycles have a quantitative original-channel obstruction]{A finite scalar consistency test away from exact saturation}
\label{sel56:robust}

\begin{proposition}[Odd-cycle lower bound]
\label{sel56:prop:cycle}
Suppose an all-input compatible joint channel has the original one-register
marginals, and let $b_{ij}$ be the opposite-bit coefficient of each pair.
For every odd cycle of length $L\ge3$,
\begin{equation}
 \boxed{\frac1L\sum_{ij\in\mathrm{cycle}}\Re b_{ij}
       \ge\frac59-\frac49\cos(\pi/L).}
 \label{sel56:eq:oddcycle}
\end{equation}
Consequently at least one edge obeys
\begin{equation}
 \boxed{\|\E_{ij}-\E_*\|_\diamond
       \ge\frac49\{1-\cos(\pi/L)\}.}
 \label{sel56:eq:oddnoise}
\end{equation}
These are necessary conditions for a joint channel. Sharpness at the level
of a vector Gram is not asserted to supply every higher channel coefficient.
\end{proposition}
\begin{proof}
The inherited all-input normal form, Lemma~\ref{sel55:lem:normal}, supplies
a positive coefficient Gram. Its principal block on the all-zero string
and the single excitations has diagonal one, first row $z$, and pair entries
$b_{ij}$. Taking its Schur complement gives vectors $u_i$ of norm one with
$\langle u_i,u_j\rangle=(b_{ij}-5/9)/(4/9)$, up to a consistent conjugation
convention that does not change real parts. The cycle adjacency matrix has
least eigenvalue $-2\cos(\pi/L)$; this follows by diagonalizing on the finite
Fourier vectors. Therefore
$2\sum_{ij}\Re\langle u_i,u_j\rangle\ge-2L\cos(\pi/L)$,
proving \eqref{sel56:eq:oddcycle}. For at least one edge the real deviation
from $1/9$ is at least the bound in \eqref{sel56:eq:oddnoise}.
A balanced superposition of one vector in spectral sector $01$ and one in
$10$ witnesses trace-distance change $|b_{ij}-1/9|$ between the two pair
channels. This lower bounds their diamond distance.
\end{proof}
For a triangle the scalar lower bound is $1/3$ and the necessary maximum
edge distance is $2/9$. Thus a certified near-extremal pair cannot be copied
onto all three edges with arbitrarily small original-channel error. The
corresponding invariant observations are the unchanged V54 ones:
$b_{ij}=4\langle X_W\rangle_{ij}-1/3+4i\langle Y_W\rangle_{ij}$.
No unmeasured zero imaginary part or exact marginal identity is inferred
from a small numerical residual.

\section[Mixing retains an advantage but removes the saturated bonus]{Exact lag capacities and a minimal memory witness}
\label{sel56:capacity}

\begin{theorem}[Capacity throughout the chronological Markov family]
\label{sel56:thm:capacity}
For the original two-register preparation and invariant target readout,
the pair at separation $\ell$ in \eqref{sel56:eq:Ek} has optimal success
\begin{equation}
 \boxed{\alpha_\ell(r)=
 \begin{cases}
 (1-r^\ell)/36,&|r|<1,\\
 3/14,&r=-1,\ \ell\text{ odd},\\
 0,&r=-1,\ \ell\text{ even},\text{ or }r=1.
 \end{cases}}
 \label{sel56:eq:capacity}
\end{equation}
For $-1<r<0$, every odd lag has larger capacity than the exchangeable
maximum $1/36$, while every even lag has smaller capacity. Both converge
to $1/36$ as $\ell\to\infty$. These are mixing, indefinitely causal sources,
not infinitely exchangeable ones.
For any $r,r'\in[-1,1]$, their adjacent pair channels satisfy
\begin{equation}
 \boxed{\|\E_2^{(r)}-\E_2^{(r')}\|_\diamond=\frac49|r-r'|.}
 \label{sel56:eq:pairdistance}
\end{equation}
\end{theorem}
\begin{proof}
For $|r|<1$, \eqref{sel56:eq:lag} lies in the open V54 disk. Its capacity
is $(1-b_\ell)/16$, giving the first line. The boundary cases follow from
the previously evaluated extremal and common-choice channels. Since
$Q_r^\ell$ converges to the uniform kernel at rate $|r|^\ell$, the stated
lag limits follow. This proves mixing of the controller and separation of
its distant output blocks; it does not identify a continuum mixing rate.

The adjacent pair difference has just two off-diagonal blocks in its
spectral-sector coefficient matrix: $00,11$ changes by
$-(4/5)z^2(r-r')$, and $01,10$ changes by $(4/9)(r-r')$.
Both have modulus $(4/9)|r-r'|$. Pinch onto those two paired subspaces.
On either subspace the difference is a phase rotation of that modulus
times the off-diagonal projection $\tfrac12(\id-\operatorname{Ad}_{P_+-P_-})$,
whose diamond norm is one. Pinching and the direct-sum triangle inequality
give the upper bound with that modulus. A balanced input on the $01,10$
subspace attains it. The argument is unchanged by spectral multiplicities
or an untouched reference.
\end{proof}

For example $r=-1/2$ gives $b_1=1/3$, $a_1=(-21+28i)/45$ and
$\alpha_1=1/24$, but $b_2=2/3$ and $\alpha_2=1/48$.
This continues at every depth with one two-state causal memory. It does
not violate V55's three-register exchangeability bound because its pair
marginals depend on separation.

\begin{proposition}[The enhanced boundary forces persistent memory]
\label{sel56:prop:boundary}
Let $r=-1+2\epsilon$, $0<\epsilon<1$, so that $\epsilon$ is the probability
of failing to flip the controller at each update. Then
\[
 \alpha_1(r)=\frac{1-\epsilon}{18},\qquad
 \|\E_2^{(r)}-\E_*\|_\diamond=\frac{8\epsilon}{9},
\]
and
\begin{equation}
 \|\E_k^{(r)}-\E_k^{(-1)}\|_\diamond
 \le2\{1-(1-\epsilon)^{k-1}\}.
 \label{sel56:eq:coupling}
\end{equation}
The value $3/14$ at $\epsilon=0$ is not the interior limit of the exact
capacity. Its extra yield requires exact pair-support saturation.
Within the exact all-input marginal class, attaining that value at every
neighbouring pair forces the nonmixing alternating hierarchy of
Corollary~\ref{sel56:cor:chain}.
\end{proposition}
\begin{proof}
The first assertions follow from the theorem. Couple the chains to the
same initial sign; the outputs agree whenever no failed flip occurs in
$k-1$ transitions. The failure probability is $1-(1-\epsilon)^{k-1}$,
and the diamond distance between any two channels is at most two.
The rigidity conclusion is the graph theorem, not an inference from a
small estimated eigenvalue.
\end{proof}

The lagged observable reconstructs $r=(9b_1-5)/4$ \emph{within the declared
model}. A necessary model check is
$9b_2-5=(9b_1-5)^2/4$. More generally $9b_\ell-5=4r^\ell$.
For $r\notin\{0,1\}$, the $2\times2$ Hankel matrix with rows
$(b_1,b_2)$ and $(b_2,b_3)$ has determinant $s(1-s)r(1-r)^2\ne0$;
the lag response therefore has two linear coordinates.
Finite agreement with these relations does not prove the Markov realization,
stationarity, or arbitrary-depth completeness. The original all-input maps
or a separately verified memory description remain the basis of such a claim.

\section[The original clock randomizes an unrecorded count, not the source law]{Physical elapsed time and retained memory records}
\label{sel56:clock}

For this subsection only, explicitly couple the comparison controller to the
inherited independent two-phase opportunity clock. It updates exactly when
an opportunity accepts, and is idle otherwise. Let $N_j$ be its acceptance
count after $j$ ticks. The initial phase law $\zeta$ is the old one. Summing
over the clock paths gives the tilted transition matrix
\begin{equation}
 B_\vartheta(r)=
 \begin{pmatrix}
 1/5&4/5\\
 (2/3)(1-\vartheta+\vartheta r)&1/3
 \end{pmatrix},\qquad
 \mathbb E(r^{N_j})=\zeta B_\vartheta(r)^j\mathbf1.
 \label{sel56:eq:clock}
\end{equation}
It is not a stochastic matrix when $r\ne1$; it is the exact generating
matrix for the original count.

\begin{proposition}[Clock averaging and accepted-count memory differ]
\label{sel56:prop:clock}
At the alternating endpoint, the controller correlation conditional on
$N_j=m$ is $(-1)^m$ and has unit magnitude. If that count is forgotten,
its correlation is \eqref{sel56:eq:clock} with $r=-1$. For the inherited
$0<\vartheta<1/2$, this tends to zero; its two characteristic roots are
\begin{equation}
 \lambda_\pm(r)=\frac{4\pm\sqrt{121-120\vartheta(1-r)}}{15},
 \label{sel56:eq:clockroots}
\end{equation}
whenever the square root is real. In particular $r=-1$ has
$|\lambda_-|<\lambda_+<1$. With original tick
$d_{\rm clk}=4\vartheta\epsilon_{\rm clk}/11$, the dominant decay rate is
$-d_{\rm clk}^{-1}\log\lambda_+(-1)$ and approaches $2/\epsilon_{\rm clk}$
as $\vartheta\downarrow0$.
\end{proposition}
\begin{proof}
On the completion transition, failure contributes weight $1-\vartheta$
and acceptance contributes $\vartheta r$; the other transitions do not
change $N_j$. Matrix multiplication proves the generating identity.
The conditional statement is the deterministic flip rule. The trace and
determinant of $B_\vartheta(r)$ give \eqref{sel56:eq:clockroots}; at $r=-1$
and $0<\vartheta<1/2$ both roots have modulus below one, and the displayed
ordering follows directly. Finally
$\lambda_+(-1)=1-8\vartheta/11+O(\vartheta^2)$, proving the rate limit.
\end{proof}
This is not a decoherence of the count-retaining controller. Given the
count, its bit can be transported back to the initial bit exactly.
Nor does an elapsed-time decay justify replacing the gap-dependent pair
channels by product queries. The coupling to this clock is stated, not
inferred from their local marginal channels. A joint acceptance record,
clock correlation or full six-label implementation needs its own source
identification. The companion's terminally prepared covector solution is
not invoked as a controller or a reset.

\section{Proof status and the next historical source test}
\label{sel56:status}

\paragraph{Proved in this layer.}
The complete stationary two-state chronological family is evaluated with
its exact lagged channels and capacities. The extremal V54 pair has a
quantum-compatible extension on a connected reuse graph exactly when that
graph is bipartite, and that extension is unique. An odd cycle gives an
original-channel distance obstruction. For mixing Markov reuse, negative
odd-lag correlations improve the selected pair capacity above the
exchangeable bound without inconsistency; the exact saturated bonus forces
persistent alternating memory. The independently coupled original clock
has a calculated generating matrix and does not erase retained-count memory.

\paragraph{Inherited inputs and what is not claimed.}
The local forgotten channel, determinant-sensitive action, pair preparation,
target characters and invariant readout remain V19/V51/V53/V54 inputs.
The pair capacity formula is inherited, not optimized again. All controller
families use the V54 alternative realization of the forgotten local channel;
none is identified with the historical six resolved coefficients. No
asymmetric reference, matter grading or new spacetime coordinate is derived.
The four companion sources supply none of those missing rows, no fixed
positive nonlinear physical lifespan, and no interacting continuum limit
through the present argument. The graph concerns the reuse of source slots,
not gauge-lattice edges. Finite memory is an explicit comparison cost, not a
proof of historical occurrence.

\paragraph{Next noncircular audit.}
Before applying the exchangeable or chronological formula, identify which
slot pairs in the actual grammar have the claimed all-input marginals and
which original memory/clock records are retained. A chain and an all-pairs
reuse protocol have different exact compatibility conditions. If the
extremal adjacent pair is actually established, the graph theorem determines
its full ordered channel hierarchy; no independent controller ansatz is
then needed at the channel level. Otherwise the lagged invariant observables
can test a proposed memory model but cannot supply it from local tomography.
Historical Standard-Model selection remains unclosed: no historical V110
same-word directed incidence has been populated.

\paragraph{Verification boundary.}
The new checker uses exact algebra for the phases, finite memory transitions,
full coefficient Grams, all-input matrix-basis marginal identities, graph
support intersections, normalizing cross constraints, odd-cycle tests,
capacity values and clock roots. The general graph and all-depth assertions
rest on their written proofs. The immediate parent checker is rerun; the
four companion analytical theorems are audited at named interfaces, not
independently re-proved or numerically certified here. These checks are not
physical measurements, a proof-assistant formalization or peer review.

\begingroup
\renewcommand{\refname}{Primary-source context and audited companions for V56}
\begin{thebibliography}{9}
\small\raggedright
\bibitem{sel56-KW} D.~Kretschmann and R.~F.~Werner,
\emph{Quantum channels with memory}, Phys. Rev. A \textbf{72} (2005), 062323;
\url{https://arxiv.org/abs/quant-ph/0502106}.
Cited for established causal memory-channel structure, not for the
source-specific pair-support theorem or the capacities above.
\bibitem{sel56-BM} G.~Bowen and S.~Mancini,
\emph{Quantum channels with a finite memory}, Phys. Rev. A \textbf{69} (2004),
012306; \url{https://doi.org/10.1103/PhysRevA.69.012306}.
\bibitem{sel56-Native} Renewal Geometry programme,
\emph{Native Regenerative Stationarity Closure Notes}, supplied V54,
\src{Renewal_Native_Regenerative_Stationarity_Closure_Notes_V54(2).tex};
\src{v54:thm:lift}, \src{v54:cor:initial}, \src{v54:eq:native-defect}.
\bibitem{sel56-Einstein} Renewal Geometry programme,
\emph{Exact Action Einstein Bridge Research Notes}, supplied V45,
\src{Renewal_Exact_Action_Einstein_Bridge_Research_Notes_V45(2).tex};
\src{v45:weak-rank}, \src{v45:positive-pencil}, \src{v45:status}.
\bibitem{sel56-ESM} Renewal Geometry programme,
\emph{Perturbative ESM Continuum Research Notes}, supplied V64,
\src{Renewal_Perturbative_ESM_Continuum_Research_Notes_V64(1).tex};
\src{thm:v64-hessian}, \src{sec:v64-matching}, \src{sec:v64-ledgers}.
\bibitem{sel56-YM} Renewal Geometry programme,
\emph{Yang--Mills Mass Gap Research Notes}, supplied V100 f17,
\src{Renewal_Yang_Mills_Mass_Gap_Research_Notes_V100_f17.tex};
\src{r100:resonance-theorem}, \src{r100:local-normal-theorem}.
\end{thebibliography}
\endgroup
\addtocontents{toc}{\protect\endgroup}
% END IMMUTABLE SOURCE-SELECTION V56 BODY
% APPEND FUTURE ITERATIONS HERE, BEFORE THE FINAL DOCUMENT COMMAND.


% BEGIN SOURCE-SELECTION V57 DEVELOPMENT BLOCK
\clearpage
\hypersetup{pdftitle={Historical Standard-Model Source Selection: V1--V57}}
\begin{center}
{\Large\bfseries Source-selection continuation V57}\\[.4em]
{\large Retained-record rigidity: the original instrument,\\
unique source reuse, and the separate acceptance-clock freedom}
\end{center}
\addtocontents{toc}{\protect\begingroup}
\addtocontents{toc}{\protect\renewcommand{\protect\numberline}{\protect\makebox[2.1em][l]}}

\section[The missing marginal is the original instrument]{Return to the original resolved source}
\label{sel57:context}

V54--V56 classify compatible joint channels after the original outcomes have
been forgotten. In particular, Theorem~\ref{sel56:thm:graph} supplies a unique
extremal channel on every connected bipartite reuse graph. Its final source
boundary explicitly leaves the original six recorded coefficients unassigned
to that alternative realization. Another optimization of the forgotten
channel would not settle this interface. We instead ask whether a joint
implementation can retain the \emph{original local instrument}, including its
state update and its classical outcome, for every joint input.

The distinction changes the answer on the unchanged benchmark. The original
six branch effects are linearly independent. That property, together with
Choi-rank-one branches, forces any all-input non-signalling extension to be a
product. On the three-dimensional forward port, retaining just its binary
sign instrument on one register already gives this conclusion. Adding the
actual idle outcome admits a different freedom: correlations in the acceptance
flags. We classify that freedom without dividing out the clock.

\paragraph{Source audit and scope.}
The normalized compiler, frame, helper vectors, and parameter point are
inherited from \eqref{sel19:eq:compiler}. The local forward restriction is
V50/V51's original sign instrument; the pair channel and its selected-ray
capacity are V54 inputs. No coefficient, label, preparation, or clock is
replaced. General Choi positivity and channel marginal problems are
established frameworks \cite{sel57-Choi,sel57-HLA}; quantum-instrument
compatibility retains both outcomes and postmeasurement states
\cite{sel57-LS}. The product-input, all-input marginal problem below is
stated directly; it is not identified with compatibility of two measurements
on one common input. The specialized rigidity and source evaluations have
self-contained proofs. No general priority claim is made.

\begin{definition}[Recorded all-input marginal]
\label{sel57:def:marginal}
Let $\mathcal I_a(X)=K_aXK_a^*$, $a\in\mathsf A$, be nonzero efficient
branches from $H_A$ to $H'_A$, with
$\sum_aE_a=I_A$, $E_a=K_a^*K_a$. Let $\Psi$ be a channel from $H_B$ to
$H'_B$. A joint implementation with the original classical $A$ record is
$\mathcal Q=\bigoplus_a\mathcal Q_a$, where the $\mathcal Q_a$ are CP and,
for every $X\in\End(H_A\otimes H_B)$,
\begin{equation}
\begin{aligned}
 \Tr_{B'}\mathcal Q_a(X)&=\mathcal I_a(\Tr_BX),\\
 \Tr_{A'}\sum_a\mathcal Q_a(X)&=\Psi(\Tr_AX).
\end{aligned}
 \label{sel57:eq:marginals}
\end{equation}
The identities are identities of linear maps, not tests on only a fixed
preparation or a maximally mixed input. $B'$ may include its own retained
classical record.
\end{definition}
The total map is normalized by the first identity. Unknown entangled inputs
and an untouched reference are included by complete positivity and equality
of the maps. Coherent off-diagonal output-record blocks are not classical
outcomes. Pinching such a register leaves its forgotten system channel
unchanged but does not license identifying the unpinched state with a classical
instrument.

\section[One recorded marginal forces product reuse]{Efficient branches and independent effects remove joint freedom}
\label{sel57:rigidity}

\begin{theorem}[One-sided recorded-source factorization]
\label{sel57:thm:oneside}
Every joint implementation satisfying the first identity of
\eqref{sel57:eq:marginals} has unique channels $\Psi_a$ such that
\begin{equation}
                 \mathcal Q_a=\mathcal I_a\otimes\Psi_a.
 \label{sel57:eq:conditionedproduct}
\end{equation}
If the Hermitian effects $E_a$ are linearly independent and the second
identity also holds, then
\begin{equation}
 \boxed{\Psi_a=\Psi\quad\text{for all }a,\qquad
 \mathcal Q=\mathcal I\otimes\Psi.}
 \label{sel57:eq:product}
\end{equation}
No efficiency, effect independence, unitary realization, or classical-memory
model is assumed for $\Psi$.
\end{theorem}
\begin{proof}
Use the unnormalized Choi convention and reorder tensor factors to put each
output beside its input. The first marginal identity implies
\[
 \Tr_{B'}J(\mathcal Q_a)=|K_a\rangle\!\rangle
       \langle\!\langle K_a|\otimes I_B.
\]
For a positive matrix $Z$, the support of $Z$ is contained in the support of
any marginal tensored with the traced space: a vector in the complementary
marginal kernel has zero positive quadratic form, so its tensor subspace is
in $\Ker Z$. It follows that
\[
 J(\mathcal Q_a)=|K_a\rangle\!\rangle
       \langle\!\langle K_a|\otimes R_a,\qquad R_a\succeq0.
\]
The same marginal equation gives $\Tr_{B'}R_a=I_B$. Thus $R_a$ is the Choi
matrix of a channel $\Psi_a$, proving \eqref{sel57:eq:conditionedproduct} and
uniqueness. Taking the other marginal yields
\[
 \sum_a E_a^{\mathsf T}\otimes
           \bigl(J(\Psi_a)-J(\Psi)\bigr)=0.
\]
Independence of the $E_a$ makes each coefficient vanish. Equivalently, apply
any trace-dual family $D_a=D_a^*$ with $\Tr(D_aE_b)=\delta_{ab}$ to the
first factor. This proves the assertion without assigning any hidden
variable to the source.
\end{proof}

The outcome probabilities need not be independent on an entangled input.
Equation~\eqref{sel57:eq:product} means that the \emph{operation} factors;
correlations already present in the input are retained. Nor does it exclude
chronological interactions on one repeatedly used memory whose later local
map genuinely depends on the earlier input. Those interactions do not satisfy
\eqref{sel57:eq:marginals} as a map on separate slots.

\begin{corollary}[All finite slot counts]
\label{sel57:cor:allcounts}
If every slot has an efficient classical instrument with linearly independent
effects, the only joint instrument having those all-input one-slot marginals
is their tensor product, at every finite slot count.
\end{corollary}
\begin{proof}
Apply the positive-support argument to every local flag. Each joint branch
has Choi matrix $t_{\boldsymbol a}\bigotimes_jJ(\mathcal I_{a_j}^{(j)})$,
with $t_{\boldsymbol a}\ge0$. For a fixed flag $a_1$, the first marginal
identity reads
\[
 \sum_{a_2,\ldots,a_k}t_{a_1\cdots a_k}
       E_{a_2}^{(2)}\otimes\cdots\otimes E_{a_k}^{(k)}=I.
\]
Products of independent effect families are independent; the identity
already has coefficient one on each product. Hence every $t_{\boldsymbol a}=1$.
This argument does not assume that higher subset marginals were supplied
independently of the one-slot identities.
\end{proof}

\section[The exact remaining coupling polytope]{When effects are dependent, the freedom is explicit}
\label{sel57:polytope}

\begin{theorem}[Complete coupling polytope for two efficient instruments]
\label{sel57:thm:polytope}
Let both local instruments be efficient, with nonzero coefficients $K_a,L_b$
and effects $E_a,F_b$. Put
\[
 \mathcal N_E=\{x\in\R^{|\mathsf A|}:\sum_ax_aE_a=0\},\qquad
 \mathcal N_F=\{y\in\R^{|\mathsf B|}:\sum_by_bF_b=0\}.
\]
Every compatible joint classical instrument is uniquely
\begin{equation}
 \boxed{\mathcal Q_{ab}=t_{ab}(\mathcal I_a\otimes\mathcal J_b),\qquad
 t\ge0,\qquad t-\mathbf1\mathbf1^{\mathsf T}
       \in\mathcal N_E\otimes\mathcal N_F.}
 \label{sel57:eq:polytope}
\end{equation}
Its real affine dimension is
$\dim\mathcal N_E\,\dim\mathcal N_F$. It is compact and contains the
product coupling in its relative interior. In particular one independent
effect family is enough to make the coupling unique.
\end{theorem}
\begin{proof}
The two local support conditions give
$J(\mathcal Q_{ab})=t_{ab}J(\mathcal I_a)\otimes J(\mathcal J_b)$.
The remaining conditions are exactly
\[
 \sum_b t_{ab}F_b=I_B\quad\text{for each }a,
 \qquad \sum_a t_{ab}E_a=I_A\quad\text{for each }b.
\]
Subtracting the product solution, every row belongs to $\mathcal N_F$ and
every column belongs to $\mathcal N_E$. Their intersection is the tensor
product displayed in \eqref{sel57:eq:polytope}. Positivity imposes just the
entrywise inequalities. The strictly positive product array is a relative
interior point. From $t_{ab}F_b\preceq I_B$ one gets
$t_{ab}\le\|F_b\|_{\op}^{-1}$, so the finite closed polytope is bounded.
\end{proof}

\begin{example}[Coefficient generation does not imply effect independence]
\label{sel57:ex:pulse}
The distinct original V34 qubit comparison has
$K_X=(9I+12iX)/25$ and $K_Z=(12I+16iZ)/25$. It generates $M_2$, but
$E_X=(9/25)I$, $E_Z=(16/25)I$. Its effect kernel is
$\R(16,-9)$. The exact array
\[
 t=\begin{pmatrix}0&25/16\\25/16&175/256\end{pmatrix},\qquad
 t-\mathbf1\mathbf1^{\mathsf T}
       =-\frac1{256}\binom{16}{-9}(16,-9),
\]
defines a compatible, nonproduct joint \emph{recorded} instrument. Its
joint record probabilities are $0,9/25,9/25,7/25$ for $XX,XZ,ZX,ZZ$.
Thus full operator-algebra generation and independent coefficient rays do
not replace the effect test. This is a comparison using that already stated
pulse bank, not the six-outcome action compiler.
\end{example}

\section[An exact effect certificate for the original compiler]{The six original records pass the rigidity test}
\label{sel57:compiler}

Keep the full native carrier of the V19 point, $N=14^5$, and its orthonormal
vectors $n,\omega,u_0,u_1,\ldots,u_4$. All unspecified directions retain their
original scalar action. Put
\[
 P_+=\tfrac12|n+\omega\rangle\langle n+\omega|,\quad
 R=\sum_{j=1}^4P_{u_j},\quad U=I-(1+i)P_+,\quad
 D=(I-\tau^4R)^{1/2},\quad \tau=1/8.
\]
The actual five-column frame is
\[
 O=\begin{pmatrix}
1&1&1&0&0\\ 1&-1&1&0&0\\ 1&0&-2&0&0\\
-1&0&0&1&1\\ -1&0&0&-1&1\\ -1&0&0&0&-2
\end{pmatrix}
\diag\left(\frac1{\sqrt6},\frac1{\sqrt2},\frac1{\sqrt6},
            \frac1{\sqrt2},\frac1{\sqrt6}\right).
\]
Thus the unchanged coefficients are
\begin{equation}
 K_e=\frac I{\sqrt{18}}+\frac{s_e}{3}UD
       +\sqrt{\frac23}\tau^2\sum_{j=1}^4O_{e,j+1}|u_0\rangle\langle u_j|,
 \qquad s_e=\begin{cases}1&e\le3,\\-1&e>3.\end{cases}
 \label{sel57:eq:bank}
\end{equation}
Their normalization and invertibility are inherited; the new calculation
concerns the six effects $E_e=K_e^*K_e$.

\begin{theorem}[The literal branch effects are independent]
\label{sel57:thm:effects}
Evaluate each $E_e$ on the six fixed linear functionals
\[
 \langle n|E_e|n\rangle,\quad \langle n|E_e|\omega\rangle,
 \quad \langle u_0|E_e|u_j\rangle\ (j=1,2,3,4).
\]
The resulting real $6\times6$ matrix $M$ has
\begin{equation}
               \boxed{\det M=\frac{\sqrt2}{220150628352}>0.}
 \label{sel57:eq:effectdet}
\end{equation}
In particular the complete native effects are linearly independent. A joint
channel preserving even one full original recorded marginal and an arbitrary
fixed other marginal must factor as in Theorem~\ref{sel57:thm:oneside}.
\end{theorem}
\begin{proof}
For a compact analytic check replace the phase by
$U_\phi=I+(e^{-i\phi}-1)P_+$ and write $c=\cos\phi$,
$\alpha_\pm=(1/\sqrt2\pm1)/3$. This is a calculation in a displayed
comparison family; the original point is $c=0$, $\tau=1/8$.
The six rows are
\begin{align*}
 M_{1e}&=1/6+s_e(1+c)/(9\sqrt2),&
 M_{2e}&=s_e(c-1)/(9\sqrt2),\\
 M_{3\bullet}&=\alpha_+\tau^2(1,-1,0,0,0,0)/\sqrt3,&
 M_{4\bullet}&=\alpha_+\tau^2(1,1,-2,0,0,0)/3,\\
 M_{5\bullet}&=\alpha_-\tau^2(0,0,0,1,-1,0)/\sqrt3,&
 M_{6\bullet}&=\alpha_-\tau^2(0,0,0,1,1,-2)/3.
\end{align*}
Subtracting the second row from the first removes $c$ from that row. The
remaining contrasts are the two independent zero-sum directions in each
three-label block. Direct block elimination gives
\[
 \det M=\frac{\tau^8(1-c)}{\sqrt2\,3^8}.
\]
At the original point this is \eqref{sel57:eq:effectdet}. These are entries
of the full effects, so no unsupported reduction of the ambient carrier is
used. Linear independence of these columns implies independence of the
operators themselves. The product conclusion follows from the earlier
proved theorem, not from a numerical rank threshold.
\end{proof}

The same result applies when the original separate order and type records
are retained. If their orthogonal record projections are $Q_\ell$, the
accepted effects are $E_e\otimes Q_\ell$. They are independent by compression
onto each record block. This statement retains those factors; it does not
rename an output sign as an internal matter grading.

The row determinant is a reconstruction certificate, not a probability,
gauge coupling, or mass gap. The comparison formula also shows why a
small-step limit needs conditioning control: its determinant can tend to zero
although exact independence holds at each nondegenerate positive step.

\section[One original sign bit already fixes the forgotten pair]{The minimum record needed on the forward three-level port}
\label{sel57:bit}

On $H_3=\Span\{n,\omega,u_0\}$ the helper inputs vanish and the sign-grouped
original branches are efficient:
\[
 L_\pm=\frac1{\sqrt3}\left(\frac I{\sqrt2}\pm U\right),\qquad
 \mathcal L_\pm(X)=L_\pm XL_\pm^*,\qquad
 \Phi=\mathcal L_++\mathcal L_-.
\]
This is the inherited \emph{forward} restriction. No reducing coefficient
algebra or full-carrier efficiency of the grouped branches is asserted.

\begin{theorem}[A single retained sign marginal forces product queries]
\label{sel57:thm:bit}
Let $P_1=P_+$ on $H_3$, $P_0=I_3-P_1$. The sign effects are
\begin{equation}
 E_\pm=\frac12I_3\pm\frac{\sqrt2}{3}P_0,\qquad
 \det\bigl[\Tr(E_sE_t)\bigr]_{s,t=\pm}=\frac49.
 \label{sel57:eq:sign-effects}
\end{equation}
Every joint implementation preserving the original sign instrument on one
register and the original forgotten channel $\Phi$ on the other, as all-input
marginals, is $\mathcal L\otimes\Phi$. In particular its forgotten quantum
channel is exactly $\Phi\otimes\Phi$.
\end{theorem}
\begin{proof}
Since $U=P_0-iP_1$, $U+U^*=2P_0$. Expanding $L_\pm^*L_\pm$ gives the
formula. $\rank P_0=2$, $\rank P_1=1$ yields
\[
 [\Tr(E_sE_t)]_{s,t}
 =\begin{pmatrix}
 43/36+2\sqrt2/3&11/36\\11/36&43/36-2\sqrt2/3
 \end{pmatrix}.
\]
Its positive determinant and Theorem~\ref{sel57:thm:oneside} prove the claim.
\end{proof}

Thus the enhanced V54 pair $\E_*$ cannot be obtained by correlating copies
of the original source \emph{while preserving that one local recorded
instrument}. It remains a legitimate joint implementation of the forgotten
local channel. Erasing all local outcome information admits both $\E_*$
and the product source; preserving the original sign on one register permits
only the latter. This is minimum positive classical alphabet size within this
specified test: one value gives the forgotten channel and does not force
independence; the original two-value sign does.

For the fixed V53 correlated preparation and its \emph{quantum-output-only}
readout, the inherited success is therefore $1/36$, not $3/14$. Conditioning
that readout on retained sign records is a different task whose optimum is
not computed here. In the shared spectral basis, $\E_*-\Phi^{\otimes2}$
has only its $00/11$ and $01/10$ coherence multipliers, both of modulus $4/9$.
Pinching onto those two orthogonal blocks and testing a balanced vector in
either block gives
\begin{equation}
              \boxed{\|\E_*-\Phi^{\otimes2}\|_\diamond=4/9.}
 \label{sel57:eq:distance}
\end{equation}
Thus the enhanced pair is separated in the original all-input channel norm
from every implementation preserving those exact marginals. No restriction
on the hidden implementation dimension removes this obstruction.

\section[Acceptance correlations are not quantum-source correlations]{Retain the actual idle branch}
\label{sel57:clock}

Let any normalized efficient accepted instrument $\{\mathcal I_a\}$ have
independent effects. At an original opportunity, retain the maps
\[
 \mathcal K_0=(1-\vartheta)\mathrm{id},\qquad
 \mathcal K_a=\vartheta\mathcal I_a,\qquad 0<\vartheta<1.
\]
Their effects are no longer independent: the idle effect is proportional to
the sum of all accepted effects. It must not be suppressed in the marginal
problem.

\begin{theorem}[Complete raw-opportunity coupling]
\label{sel57:thm:idle}
Every compatible two-opportunity classical instrument with these full
original local marginals has exactly one free parameter $q$:
\begin{align}
 \mathcal Q_{00}&=(1-2\vartheta+q)\,\mathrm{id}\otimes\mathrm{id},\nonumber\\
 \mathcal Q_{0b}&=(\vartheta-q)\,\mathrm{id}\otimes\mathcal I_b,\nonumber\\
 \mathcal Q_{a0}&=(\vartheta-q)\,\mathcal I_a\otimes\mathrm{id},\nonumber\\
 \mathcal Q_{ab}&=q\,\mathcal I_a\otimes\mathcal I_b,
 \qquad \max(0,2\vartheta-1)\le q\le\vartheta.
 \label{sel57:eq:raw}
\end{align}
Here $q$ is the joint acceptance probability and is independent of the input.
Whenever $q>0$, conditioning on two acceptances returns exactly the product
accepted instrument. In particular it cannot return $\E_*$ on the indicated
forward quantum port.
\end{theorem}
\begin{proof}
The effect kernel is one-dimensional, spanned by
$v=(\vartheta,-(1-\vartheta),\ldots,-(1-\vartheta))$. Indeed a vanishing
linear combination becomes
$\sum_a[(1-\vartheta)x_0+\vartheta x_a]E_a=0$, and effect independence
sets every bracket to zero. The coupling polytope is therefore
$t=\mathbf1\mathbf1^{\mathsf T}+u vv^{\mathsf T}$. Multiplying its entries
by the original idle/accepted weights and setting
$q=\vartheta^2[1+u(1-\vartheta)^2]$ gives \eqref{sel57:eq:raw}.
Its four coefficient classes are nonnegative exactly on the displayed
interval. Summing all $a,b$ in the final row yields effect $qI\otimes I$.
The conditional instrument is consequently the product, with no nonlinear
input-dependent normalization.
\end{proof}

For independent acceptance $q=\vartheta^2$; common acceptance allows
$q=\vartheta$. For the fixed output-only pair task, success at two raw
opportunities, retaining only double acceptances, is $q/36$. For the
historical two-phase clock the actual $q$ depends on which opportunities,
phase records and common randomness are identified. The theorem does not
supply that joint clock law. It separates its allowable acceptance
correlation from the accepted quantum dynamics. It is a two-slot theorem;
no unproved classification of arbitrary multi-slot clock correlations is
inferred from its single parameter.

\section[A stable one-sided version]{Marginal accuracy and the effect-dual cost}
\label{sel57:stability}

\begin{proposition}[Exact recorded marginal, approximate other marginal]
\label{sel57:prop:stability}
Assume the first identity of \eqref{sel57:eq:marginals} exactly, and assume
that the $E_a$ are independent. Let $D_a=D_a^*$ be a trace-dual family and
put $\kappa_E=\max_a\|D_a\|_1$. If
\[
 \delta=\left\|\Tr_{A'}\sum_a\mathcal Q_a-
                            \Psi\circ\Tr_A\right\|_\diamond,
\]
then
\begin{equation}
 \boxed{\|\mathcal Q-\mathcal I\otimes\Psi\|_\diamond
                         \le\min\{2,\kappa_E\delta\}.}
 \label{sel57:eq:robust}
\end{equation}
For the original three-level sign instrument one can take
$\kappa_E=1+3/\sqrt2$.
\end{proposition}
\begin{proof}
The first marginal gives $\mathcal Q_a=\mathcal I_a\otimes\Psi_a$ exactly.
Inserting $D_a\otimes X$ into the other marginal difference isolates
$(\Psi_a-\Psi)(X)$. The preparation map $X\mapsto D_a\otimes X$ has
complete trace norm $\|D_a\|_1$, including Hermitian $D_a$ not positive.
Thus $\|\Psi_a-\Psi\|_\diamond\le\|D_a\|_1\delta$.
For an arbitrary joint density and reference, first apply the instrument
$\mathcal I$ on $A$; its unnormalized positive outputs have total trace one.
Applying $\Psi_a-\Psi$ on the other system in each classical block gives
total trace norm at most $\max_a\|\Psi_a-\Psi\|_\diamond$. The same bound
holds in diamond norm. The two maps compared are channels, giving the
additional bound two. For the sign effects, put $u=3/(2\sqrt2)$ and choose
\[
 D_+=\tfrac u2P_0+(1-u)P_1,\qquad
 D_-=-\tfrac u2P_0+(1+u)P_1.
\]
They satisfy the trace-dual equations and have trace norms $2u-1,2u+1$.
\end{proof}

This estimate is not a theorem for arbitrary noisy approximations to both
marginals. Exact efficiency and the first all-input identity are retained.
Statistical errors cannot silently establish either premise. Independence
of a reconstructed finite effect frame can, however, be certified one-sided:
if a measured $6\times6$ witness $\widehat M$ has operator error at most
$\eta$ and $\sigma_{\min}(\widehat M)>\eta$, the true effect family is
independent. No Choi purification, rank truncation, or source projection is
performed by this test. The effect-dual norm displays the cost of poor
conditioning instead of hiding it in a qualitative completeness statement.

\section{Status and the next source-level audit}
\label{sel57:status}

\paragraph{Proved.}
The complete original-instrument marginal problem factors whenever one
local efficient branch bank has independent effects. For two efficient banks
all remaining freedom is a finite positive coupling polytope determined by
the two effect kernels. The literal V19 six-effect bank passes an exact
finite minor test on its unchanged full carrier. On its three-level forward
port one original sign marginal suffices. An actual idle branch leaves
precisely a two-slot acceptance correlation and cannot alter the conditional
accepted product. A one-sided diamond estimate keeps the dual-effect
conditioning explicit.

\paragraph{Inherited and not claimed.}
The compiler and the $1/36$ and $3/14$ pair-task capacities are inherited;
we do not optimize preparation or readout again. The source-sensitive
comparison with V56 is now decided: the extremal bipartite graph channel
cannot simultaneously preserve the original local sign instrument, even on
one of its endpoint registers. This does not invalidate its alternative
forgotten-channel realization, forbid correlated input states, or exclude
chronological memory that changes the all-input local marginals. The V34
scalar-effect bank supplies an explicit exception to effect rigidity.

\paragraph{Next noncircular audit.}
Identify the actual retained local instrument before declaring a reuse graph
or choosing a channel-only extension. If its original efficient effects are
independent, reuse is fixed at the channel level without a controller
ansatz, an exchangeability assumption, or more pair optimization. If the
actual effects are dependent, compute their kernel and the physical
coupling constraints; do not substitute coefficient-algebra rank. Clock
flags and fixed input states are distinct from all-input branch identities.
No historical V110 directed determinant--Clifford cylinder has been populated.
Historical Standard-Model selection remains unclosed, but the enhanced
channel-only construction is no longer an unresolved candidate for reuse of
the unchanged recorded compiler.

\paragraph{Verification boundary.}
The new checker verifies the exact six-effect minor, full normalization,
record-block independence, the binary effect Gram and its dual, compatible
and incompatible finite coupling arrays, idle-branch marginals, original
three-level matrix-unit maps and the extremal-pair distance witness. The
immediate V56 checker is rerun unchanged and its certificate compared with
the supplied baseline. The general positivity and all-depth statements rest
on their written proofs. These are not physical measurements, formal
proof-assistant certification, or independent peer review.

\begingroup
\renewcommand{\refname}{Primary-source context for V57}
\begin{thebibliography}{9}\small\raggedright
\bibitem{sel57-Choi} M.-D. Choi, Completely positive linear maps on complex
matrices, \emph{Linear Algebra Appl.} \textbf{10} (1975), 285--290.
\url{https://doi.org/10.1016/0024-3795(75)90075-0}.
Choi positivity is standard background; the support arguments are proved above.
\bibitem{sel57-HLA} C.-Y. Hsieh, M. Lostaglio, and A. Ac\'in,
Quantum Channel Marginal Problem, \emph{Phys. Rev. Research}
\textbf{4} (2022), 013249; \url{https://arxiv.org/abs/2102.10926}.
Cited for the established all-input marginal framework, not as a source of
the evaluated instrument or its unknown historical realization.
\bibitem{sel57-LS} L. Lepp\"aj\"arvi and M. Sedl\'ak,
Incompatibility of quantum instruments, \emph{Quantum} \textbf{8} (2024),
1246; \url{https://arxiv.org/abs/2212.11225}.
Cited for the distinction between outcome/state-update instruments and
their induced channels; our separate-slot marginal conditions are explicit.
\end{thebibliography}
\endgroup
\addtocontents{toc}{\protect\endgroup}
% END SOURCE-SELECTION V57 DEVELOPMENT BLOCK
% Append future development before the final document command.


% BEGIN SOURCE-SELECTION V58 DEVELOPMENT BLOCK
\clearpage
\hypersetup{pdftitle={Historical Standard-Model Source Selection: V1--V58}}
\begin{center}
{\Large\bfseries Source-selection continuation V58}\\[.4em]
{\large Noisy recorded-marginal rigidity, its optimal square-root order,\\
and a coherent-code obstruction to the enhanced pair}
\end{center}
\addtocontents{toc}{\protect\begingroup}
\addtocontents{toc}{\protect\renewcommand{\protect\numberline}{\protect\makebox[2.1em][l]}}

\section[From exact branch identity to a measured error budget]{The unresolved hypothesis in the recorded-source theorem}
\label{sel58:context}

Theorem~\ref{sel57:thm:oneside} fixes joint source reuse from one efficient
recorded marginal and an independent effect family. Its stability statement,
Proposition~\ref{sel57:prop:stability}, still keeps that first marginal
\emph{exact}. No finite noisy panel by itself proves exact Choi rank one.
This continuation treats an arbitrary actual joint instrument: neither its
recorded branches nor either of its marginals is assumed to have the exact
reference form. A dimension-independent comparison theorem then turns two
all-input marginal error bounds into a joint product bound. The reference
instrument is not substituted for the actual process.

The second calculation returns to the unchanged enhanced pair of V54. That
pair is exactly isometric on a two-dimensional code. The original sign bit
has different probabilities on the two code inputs. An application of the
already proved V6 record/coherence inequality therefore gives a considerably
stronger, source-specific incompatibility test than generic perturbation
continuity. The physical ingredients and the error conventions are retained
in each comparison.

\paragraph{Audit and mathematical inputs.}
V57's rank-one marginal factorization, effect-dual estimate, original sign
coefficients, and the distance $4/9$ to the enhanced forgotten pair are
inherited. The record-fidelity inequality in the proof of
Theorem~\ref{sel6:thm:information-disturbance} is inherited as well.
Continuity of Stinespring dilations is an established theorem
\cite{sel58-KSW}; its finite-dimensional channel form is stated explicitly
below. The new results are the two-noisy-marginal consequence, its
source-calibrated constants and sharp exponent example, and the evaluated
code/record obstruction. No general priority claim is made for dilation
continuity or information--disturbance. No new historical source table is
populated.

\paragraph{Actual maps and unhalved norms.}
Let $\mathcal I=\bigoplus_a\mathcal I_a$ be a specified normalized reference
instrument, with nonzero $\mathcal I_a(X)=K_aXK_a^*$ and independent effects
$E_a=K_a^*K_a$. Let $\Psi:B\to B'$ be any specified channel. The actual
channel $\mathcal Q=\bigoplus_a\mathcal Q_a$ has input $AB$ and outputs the
original classical flag $a$, $A'$, and $B'$. It need not have efficient
branches. Write
\begin{align}
 \mathcal M_A&=\bigoplus_a\Tr_{B'}\mathcal Q_a,
 &\mathcal M_B&=\Tr_{A'}\sum_a\mathcal Q_a,\nonumber\\
 \epsilon&=\|\mathcal M_A-\mathcal I\circ\Tr_B\|_\diamond,
 &\eta&=\|\mathcal M_B-\Psi\circ\Tr_A\|_\diamond.
 \label{sel58:eq:errors}
\end{align}
These are identities or distances of maps on the entire joint input, not
fits at one preparation. Diamond norms are \emph{not divided by two}.
Let $D_a=D_a^*$ satisfy $\Tr(D_aE_b)=\delta_{ab}$ and put
$\kappa_E=\max_a\|D_a\|_1$. One may minimize this constant over trace-dual
families. No lower bound on the probability of each label on every input
is imposed.

\section[Two noisy marginals control the whole joint channel]{A positive extension comparison with the flag retained}
\label{sel58:robust}

\begin{lemma}[Lifting a nearby marginal while retaining its extension]
\label{sel58:lem:lift}
Let $\mathcal Q:X\to YZ$ be a channel, $\mathcal M=\Tr_Z\mathcal Q$, and
let $\mathcal N:X\to Y$ be a channel with
$\|\mathcal M-\mathcal N\|_\diamond\le e$. There exists a channel
$\widetilde{\mathcal Q}:X\to YZ$ such that
\begin{equation}
 \Tr_Z\widetilde{\mathcal Q}=\mathcal N,
 \qquad
 \|\mathcal Q-\widetilde{\mathcal Q}\|_\diamond
                  \le\min\{2,2\sqrt e\}.
 \label{sel58:eq:lift}
\end{equation}
If $Y$ contains an original classical flag and both $\mathcal Q$ and
$\mathcal N$ are diagonal in that flag, the comparison can also be chosen
classical in the same flag, without introducing new labels.
\end{lemma}
\begin{proof}
In the finite-dimensional Schr\"odinger convention, the Stinespring
continuity theorem \cite[Theorem~1]{sel58-KSW} supplies a common environment
and isometries $V_{\mathcal M},V_{\mathcal N}$ with
$\|V_{\mathcal M}-V_{\mathcal N}\|_{\op}\le\sqrt e$. Every extension of
$\mathcal M$ is obtained by a channel on such a dilation environment.
For completeness, begin with a minimal dilation of $\mathcal M$.
Stinespring uniqueness embeds it both in the chosen common environment
and in the environment $Z\otimes E$ of a dilation of $\mathcal Q$.
On the embedded minimal subspace, the map to $Z$ is that isometry followed
by $\Tr_E$; extend it to a channel $\mathcal D$ on the orthogonal complement
by a fixed trace-one output. Then
$\mathcal Q=(\id_Y\otimes\mathcal D)(V_{\mathcal M}(\cdot)V_{\mathcal M}^*)$.
Define the comparison with $V_{\mathcal N}$ and the same $\mathcal D$.
It has marginal $\mathcal N$. Expanding the difference of the two isometric
channels and using complete trace-norm contraction gives the bound
$2\|V_{\mathcal M}-V_{\mathcal N}\|_{\op}$. Two channels have distance at
most two. Finally pinch the original flag in the comparison. This leaves its
marginal $\mathcal N$ and the actual $\mathcal Q$ unchanged, and cannot
increase their distance.
\end{proof}

\begin{theorem}[Recorded-source rigidity with two imperfect marginals]
\label{sel58:thm:robust}
Under \eqref{sel58:eq:errors}, with no rank restriction on the actual maps,
\begin{equation}
 \boxed{
 \|\mathcal Q-\mathcal I\otimes\Psi\|_\diamond
 \le B_E(\epsilon,\eta)
 :=\min\{2,\,2(1+\kappa_E)\sqrt\epsilon+\kappa_E\eta\}.}
 \label{sel58:eq:main}
\end{equation}
The bound is independent of the dimensions of the second slot and every
hidden implementation environment. At $\epsilon=0$ it recovers V57's
one-sided estimate. The square-root \emph{order} is optimal in this class;
no optimality is claimed for the prefactor in \eqref{sel58:eq:main}.
\end{theorem}
\begin{proof}
Apply Lemma~\ref{sel58:lem:lift} to $Y=(a,A')$, $Z=B'$ and
$\mathcal N=\mathcal I\circ\Tr_B$. It gives an original-flag comparison
$\widetilde{\mathcal Q}$ with the first marginal exact and joint error at
most $r=2\sqrt\epsilon$. V57's positive-support argument yields
$\widetilde{\mathcal Q}_a=\mathcal I_a\otimes\Psi_a$ for normalized
channels $\Psi_a$. Partial-trace contraction bounds the other marginal's
error by $\eta+r$. The unchanged effect-dual proof of
Proposition~\ref{sel57:prop:stability} gives
$\|\widetilde{\mathcal Q}-\mathcal I\otimes\Psi\|_\diamond
\le\kappa_E(\eta+r)$. A triangle inequality gives
\eqref{sel58:eq:main}, and channel normalization supplies the cap two.
\end{proof}

The dilation and the comparison are proof objects. They are not a physical
recovery map, an accessible environment, or a replacement of noisy measured
branches by exact pure branches. The conclusion is a bound on the
\emph{unchanged actual} $\mathcal Q$. In particular, independent effects
alone do not assert exact factorization when the branch maps are mixed.
The next example quantifies the failure without changing even the effects.

\section[The square-root order is necessary on the original sign bank]{Small marginal error can conceal larger joint correlations}
\label{sel58:sharp}

Retain the three-level V57 sign instrument
\[
 L_\pm=(I_3/\sqrt2\pm U)/\sqrt3,
 \qquad U=P_0-iP_1,\qquad \rank P_0=2,\quad\rank P_1=1.
\]
Choose orthonormal $e,f\in\Ran P_0$ and set
$G=P_e-P_f+P_1$, $R_t=e^{-itG}$. Thus $G=G^*=G^{-1}$ and $G$ commutes
with both original coefficients. As a \emph{declared perturbed comparison},
discard the second input and produce a classical qubit $B'$ with
\begin{equation}
 \mathcal Q_{t,a}(X)=\frac12\sum_{s=\pm1}
 R_{st}L_a(\Tr_BX)L_a^*R_{st}^*\otimes P_s^{B'}.
 \label{sel58:eq:sharpmap}
\end{equation}
This introduces explicitly specified correlated output rotations; it is not
claimed to be the unchanged physical compiler. The nominal reference
instrument and its effects, however, remain exactly the original ones.
Let $\Psi(X)=\Tr(X)I_2/2$.

\begin{proposition}[Exact two-sided errors and exact joint separation]
\label{sel58:prop:sharp}
For $0<t<\pi/2$, the actual first branches in
\eqref{sel58:eq:sharpmap} have Choi rank two, while their effects are exactly
$E_a=L_a^*L_a$. The second marginal is exactly $\Psi\circ\Tr_A$, and
\begin{equation}
 \boxed{
 \epsilon_t=2\sin^2t,\qquad \eta_t=0,\qquad
 \|\mathcal Q_t-\mathcal I\otimes\Psi\|_\diamond
      =2\sin t=\sqrt{2\epsilon_t}.}
 \label{sel58:eq:sharp}
\end{equation}
The reference bank and its effect Gram are fixed as $t$ varies. No
input-dependent or vanishing-probability outcome has been conditioned away. No $o(\sqrt\epsilon)$ joint estimate with fixed
$\kappa_E$ and $\eta=0$ can hold generally.
\end{proposition}
\begin{proof}
After summing $s$, the first marginal is
$(\cos^2t\,\id+\sin^2t\,\operatorname{Ad}_G)\circ\mathcal I$.
This proves the upper bound $2\sin^2t$. On
$q=(e+f)/\sqrt2$, every $L_a$ is scalar and $Gq\perp q$; the output
trace distance therefore attains that upper bound, with the original
outcome probabilities summing to one. Tracing out the first output and
summing $a$ proves the stated second marginal exactly. Because the $s$
blocks are orthogonal, the joint distance is at most the average of
$\|\operatorname{Ad}_{R_t}-\id\|_\diamond$ and
$\|\operatorname{Ad}_{R_{-t}}-\id\|_\diamond$, each equal to
$2\sin t$. The same input $q$ attains both pure-state distances and hence
their block sum. The effects are unchanged by the output unitaries.
The invertible $L_a$ and the two nonproportional matrices $R_t,R_{-t}$
give Choi rank two. All assertions are exact map statements, so arbitrary
untouched references are included.
\end{proof}

For example, $\sin t=1/10$ gives first-marginal error $1/50$, zero
second-marginal error, and joint error $1/5$. The side output carrying $s$
can itself be classical: output entanglement is not required for the
square-root obstruction. Checking only the original branch probabilities
would miss this family entirely, because every effect remains unchanged.

\section[A calibrated efficient comparison from positive data]{Near efficiency need not be declared exact}
\label{sel58:calibration}

A specified efficient reference can come from an independent exact model,
as above. There is also an optional positive-data comparison construction.
It is not a proof that the observed instrument is efficient.

\begin{proposition}[Rank-one component comparison with its full error paid]
\label{sel58:prop:calibration}
Let $\mathcal T=\bigoplus_a\mathcal T_a$ be any normalized instrument on
$A$, with Choi matrices $J_a\succeq0$. Choose nonzero rank-one components
$|K_a\rangle\!\rangle\langle\!\langle K_a|\preceq J_a$, for example the
largest spectral component. Put
\[
 Z=\sum_aK_a^*K_a=I-R,
 \quad
 R=\left[\sum_a\Tr_{A'}
 (J_a-|K_a\rangle\!\rangle\langle\!\langle K_a|)\right]^{\mathsf T},
 \quad r=\|R\|_{\op}<1.
\]
Then $\widehat K_a=K_a Z^{-1/2}$ defines a normalized efficient
comparison $\widehat{\mathcal I}$ on the same full input and original labels,
and
\begin{equation}
 \boxed{
 \|\mathcal T-\widehat{\mathcal I}\|_\diamond
 \le\Gamma(r):=r+2(1-\sqrt{1-r})\le3r.}
 \label{sel58:eq:calibration}
\end{equation}
Its effects are independent exactly when the retained $K_a^*K_a$ are.
If the measured first joint marginal is within $e$ of
$\mathcal T\circ\Tr_B$, Theorem~\ref{sel58:thm:robust} applies with
$\epsilon\le e+\Gamma(r)$ and the comparison effect-dual norm.
\end{proposition}
\begin{proof}
The residual instrument is CP with total effect $R$, so its diamond norm
is $r$. The retained Kraus column $V$ has $V^*V=Z$ and the normalized
column is $\widehat V=VZ^{-1/2}$. Polar decomposition gives
$\|V-\widehat V\|_{\op}=1-\sqrt{1-r}$; both column norms are at most
one. Expanding their induced CP maps, including their classical flags,
bounds their distance by $2(1-\sqrt{1-r})$. Add the residual norm $r$.
Finally $\widehat E_a=Z^{-1/2}(K_a^*K_a)Z^{-1/2}$, and this invertible
linear congruence preserves independence. The last assertion is a triangle
inequality on the first marginal.
\end{proof}

All residual noise and the normalization change have been included in
\eqref{sel58:eq:calibration}. If $r=1$, this full-domain construction is
not available; compressing the support would change the input problem.
The top eigenvector is only one convenient comparison choice, with no
optimality claim. Finite measurement intervals must certify the residual
and the effect frame, rather than assign them their nominal values.
For the Hilbert--Schmidt effect Gram $G_E$, the canonical trace-dual family
satisfies
\begin{equation}
 \kappa_E\le\max_a\sqrt{d_A(G_E^{-1})_{aa}}
 \le\sqrt{d_A/\lambda_{\min}(G_E)}.
 \label{sel58:eq:gramdual}
\end{equation}
Thus an operator-error interval smaller than the observed least Gram
eigenvalue certifies a finite constant. An ill-conditioned frame is not
replaced by an unqualified claim of tomographic completeness.

\section[The original sign has a sharp conditioning constant]{What the calibrated event scale does to rigidity}
\label{sel58:sign}

Retain the original sign construction but allow its specified action phase
$\phi$:
\[
 U_\phi=P_0+e^{-i\phi}P_1,\qquad
 L_\pm(\phi)=(I_3/\sqrt2\pm U_\phi)/\sqrt3,
 \qquad c=\cos\phi<1.
\]
No variation of this parameter is identified with a change of physical clock
unless that family has been independently declared.

\begin{proposition}[Optimal trace-dual cost for the sign effects]
\label{sel58:prop:dual}
The effects and the least possible maximum trace-dual norm are
\begin{equation}
 \boxed{
 E_\pm(\phi)=\tfrac12I_3\pm\tfrac{\sqrt2}{3}(P_0+cP_1),
 \qquad
 \kappa_{\min}(\phi)=\frac{3/\sqrt2+1+c}{1-c}.}
 \label{sel58:eq:dual}
\end{equation}
At the unchanged phase $\phi=\pi/2$ this is $1+3/\sqrt2$. The general
joint-source bound becomes
\begin{equation}
 \boxed{
 \|\mathcal Q-\mathcal L\otimes\Psi\|_\diamond
 \le\min\{2,(4+3\sqrt2)\sqrt\epsilon
                    +(1+3/\sqrt2)\eta\}.}
 \label{sel58:eq:signbudget}
\end{equation}
\end{proposition}
\begin{proof}
Average any Hermitian dual over the unitaries of $P_0$ and $P_1$. This
preserves its dual constraints and cannot increase its trace norm. Put
$D_\pm=x_\pm P_0/2+y_\pm P_1$ and $u=3/(2\sqrt2)$. The two dual equations
are $x_\pm+y_\pm=1$ and $x_\pm+cy_\pm=\pm u$, hence
\[
 x_\pm=\frac{\pm u-c}{1-c},\qquad
 y_\pm=\frac{1\mp u}{1-c}.
\]
Since $u>1$ and $-1\le c<1$, their norms are
$(2u-1-c)/(1-c)$ and $(2u+1+c)/(1-c)$. The second is the maximum.
The averaging argument proves optimality over all duals, not just block
scalar candidates. Substitution in \eqref{sel58:eq:main} gives
\eqref{sel58:eq:signbudget}.
\end{proof}

In a separately fixed-Hamiltonian family with
$\phi=3\tau\Delta_H/2$, the dual cost grows as $\tau^{-2}$ when
$\Delta_H\ne0$ is fixed. This sufficient rigidity budget tends to zero
if, for example, $\epsilon=o(\tau^4)$ and $\eta=o(\tau^2)$.
These are sufficient precision scalings for this bound, not necessary
physical laws. Exact effect independence at every nonzero scale does not
by itself give cutoff-uniform statistical rigidity. At $c=1$ the two
reference effects are scalar multiples of $I$ and the independence premise
fails exactly. This conditioning issue is separate from V57's independent
accepted effects at its specified nonzero phase.

\section[A stronger record test for the enhanced pair]{The unchanged extremal pair contains an isometric code}
\label{sel58:code}

Retain V54's exact enhanced pair on $H_3\otimes H_3$,
\[
 \mathcal E_*=\tfrac12\operatorname{Ad}_{V_+\otimes V_-}
             +\tfrac12\operatorname{Ad}_{V_-\otimes V_+},
 \quad
 V_\pm=P_0+\lambda_\pm P_1,
 \quad
 \lambda_\pm=\bar z(1\pm2i/\sqrt5),\quad z=(1+2i)/3.
\]
Choose any unit $e\in\Ran P_0$ and the unit $h\in\Ran P_1$, and let
$|0_c\rangle=e\otimes e$, $|1_c\rangle=h\otimes h$. Both Kraus unitaries
agree on this two-dimensional code:
\begin{equation}
 |0_c\rangle\longmapsto|0_c\rangle,\qquad
 |1_c\rangle\longmapsto\gamma|1_c\rangle,
 \qquad \gamma=\lambda_+\lambda_-=\frac{-3-4i}{5}.
 \label{sel58:eq:code}
\end{equation}
The restriction of $\mathcal E_*$ is therefore exactly isometric.

\begin{theorem}[Record visibility excludes an accurate enhanced-pair claim]
\label{sel58:thm:code}
Let any actual two-sign joint instrument $\mathcal Q$ have forgotten channel
$\mathcal E=\sum_\pm\mathcal Q_\pm$. Let
\[
 \gamma_E=\|\mathcal E-\mathcal E_*\|_\diamond,
 \qquad c_*=\sqrt2/3,
 \qquad p_*=(\tfrac12+c_*,\tfrac12-c_*),
 \qquad q_*=(\tfrac12,\tfrac12).
\]
If its first recorded marginal is within $\epsilon$ of the original sign
instrument in \eqref{sel58:eq:errors}, then
\begin{equation}
 \boxed{\gamma_E\ge1-F(\epsilon),}
 \label{sel58:eq:codebudget}
\end{equation}
where, for $0\le\epsilon\le c_*$,
\begin{align}
 F(\epsilon)={}&
 \sqrt{(\tfrac12+c_*-\epsilon/2)(\tfrac12+\epsilon/2)}\nonumber\\
 &+\sqrt{(\tfrac12-c_*+\epsilon/2)(\tfrac12-\epsilon/2)},
 \label{sel58:eq:fidelity}
\end{align}
and $F(\epsilon)=1$ for $\epsilon\ge c_*$. In particular,
\begin{equation}
 \boxed{
 \mathcal E=\mathcal E_*
       \ \Longrightarrow\ \epsilon\ge\frac{\sqrt2}{3},
 \qquad
 \epsilon=0
       \ \Longrightarrow\ \gamma_E\ge1-\sqrt{2/3}.}
 \label{sel58:eq:codeendpoints}
\end{equation}
No efficient-branch assumption on $\mathcal Q$, no independent-effect
assumption on its measured effects, and no second-marginal hypothesis is
needed for these exclusions.
\end{theorem}
\begin{proof}
Let $p,q$ be the actual sign distributions on $0_c,1_c$. Their target
values are $p_*,q_*$ by the original effects. Because the marginal outputs
are normalized, trace-distance contraction gives
$|p_+-p_{*,+}|\le\epsilon/2$ and
$|q_+-1/2|\le\epsilon/2$.

The visibility step in the proof of
Theorem~\ref{sel6:thm:information-disturbance}, applied to this input code
and its output isometry \eqref{sel58:eq:code}, gives
\[
 \gamma_E\ge1-\sum_{a=\pm}\sqrt{p_aq_a}.
\]
Explicitly, if $M_{a\beta}$ are actual Kraus coefficients and
$f_0=|0_c\rangle$, $f_1=\gamma|1_c\rangle$, then
$|\langle f_0,\mathcal E(|0_c\rangle\langle1_c|)f_1\rangle|
\le\sum_a\sqrt{p_aq_a}$ by Cauchy--Schwarz. A reference qubit and the
Hermitian off-diagonal code witness used in V6 turn this into the stated
channel-norm lower bound. This is the inherited information--disturbance
argument, not a new general principle.

For $\epsilon<c_*$ the two allowed binary-probability intervals are
disjoint. Their maximum classical fidelity is at their closest endpoints,
$p_+=1/2+c_*-\epsilon/2$, $q_+=1/2+\epsilon/2$: write the fidelity as
$\cos(\arccos\sqrt{p_+}-\arccos\sqrt{q_+})$ to see the monotonicity.
This gives \eqref{sel58:eq:fidelity}; when the intervals overlap the upper
bound is one. At zero error $F(0)^2=(1+\sqrt{1-4c_*^2})/2=2/3$.
For exact $\mathcal E_*$, every CP refinement of its rank-one code
restriction is proportional to that restriction. Thus $p=q$, which also
gives $\epsilon\ge c_*$ directly. No constraint on a hidden dilation was
introduced.
\end{proof}

The threshold $\sqrt2/3$ is already forced by two classical input/record
comparisons. It is not claimed to be the minimum full recorded-marginal
channel error over all realizations of $\mathcal E_*$. The fidelity bound
is sharp for instruments acting on the isolated code with specified record
laws: the diagonal coefficients
$M_a=\sqrt{p_a}|f_0\rangle\langle0_c|
     +\sqrt{q_a}|f_1\rangle\langle1_c|$
produce exactly the dephasing factor $\sum_a\sqrt{p_aq_a}$. Additional
full-channel and marginal requirements can increase the minimum error.

This exclusion is stronger, for this source, than the generic consequence
of \eqref{sel58:eq:main} and the inherited pair distance $4/9$. The generic
inequality would only give
\[
 \frac49-\gamma_E\le
        (4+3\sqrt2)\sqrt\epsilon+(1+3/\sqrt2)\eta.
\]
The isometric-code calculation uses more source information and requires
less tomography for exclusion. It does not prove general product reuse
from two classical frequencies. That positive conclusion still requires
the all-input conditions of Theorem~\ref{sel58:thm:robust}.

\section[What finite error does and does not certify]{Operational consequences and status}
\label{sel58:status}

\begin{proposition}[Fixed-protocol output and normalization budget]
\label{sel58:prop:postprocess}
If two joint channels differ by at most $\Delta$ in diamond norm, every
identical physical preparation, ancillary processing and terminal CP
success branch gives subnormalized outputs $X,Y$ with
$\|X-Y\|_1\le\Delta$. If both are nonzero and
$p=\Tr X\ge s>0$, then
\begin{equation}
 |\Tr X-\Tr Y|\le\Delta,\qquad
 \left\|\frac X{\Tr X}-\frac Y{\Tr Y}\right\|_1
       \le\min\{2,2\Delta/s\}.
 \label{sel58:eq:postprocess}
\end{equation}
For a circuit with at most $n$ uses, the same telescoping estimate is
$n\Delta$ whenever the per-use bound holds on the complete actual input,
including every memory crossing that interface.
\end{proposition}
\begin{proof}
Preparation and CP trace-nonincreasing output processing are completely
trace-norm contractions on Hermitian inputs. The normalized difference is
bounded by $(\|X-Y\|_1+|\Tr X-\Tr Y|)/p$. A finite telescoping replacement
of one complete interface at a time proves the last assertion. A local
marginal estimate alone is not the required complete-memory bound.
\end{proof}

This gives error control for a \emph{fixed} protocol, not continuity of an
optimized exact-rank-one success capacity. V54's boundary example already
shows why the latter inference would be false. Nor does a small right-hand
side of \eqref{sel58:eq:main} establish exact product structure, symmetry,
or a vanished typed source residual. The present theorem supplies an
operational neighbourhood with explicit constants.

\paragraph{Data and clock boundary.}
The errors in \eqref{sel58:eq:errors} can be certified from a sufficiently
complete original contextual frame and its error bounds. For a
Hermiticity-preserving map $\mathcal D$ with input dimension $d$,
$\|\mathcal D\|_\diamond\le\|J(\mathcal D)\|_1
=d\|J(\mathcal D)/d\|_1$ is a conservative conversion; no absence of
ancillary amplification is assumed. The smaller code exclusion needs only
the displayed record probabilities plus an independently tested coherent
code action. Statistical identifiability still does not create an
unavailable preparation or Read.

The normalized accepted-instrument statement is not applied to the raw bank
by deleting its idle branch. The raw effect family is dependent, and
Theorem~\ref{sel57:thm:idle} retains its independent acceptance parameter.
Conditioning noisy maps on acceptance requires the corresponding success
floor as in \eqref{sel58:eq:postprocess}; it may not be divided out without
cost. No joint acceptance probability or physical elapsed time is inferred
from the local accepted-channel errors.

\paragraph{Proved and inherited.}
V57's exact recorded-source theorem now has a two-noisy-marginal extension,
without exact rank assumptions on the measured joint source. The
square-root order is necessary even for a fixed original sign bank with
unchanged effects. A normalized efficient comparison can be certified from
positive data with its full error paid. The optimal sign dual cost shows
what happens at a weak-event limit. Finally, the unchanged enhanced pair
has a source-specific coherent-code obstruction of size $\sqrt2/3$ to an
exact original sign marginal, independent of any Kraus-rank census.

\paragraph{Next noncircular source audit.}
Populate the actual all-input recorded maps and their error budgets on the
identified carrier; use a certified effect condition number rather than
exact zeroes inferred from floating-point fits. For the enhanced-pair claim,
first test the two code record laws and the protected code coherence, since
they already give a stronger exclusion. Do not reopen a forgotten-channel
completion search when the original recorded effects and state updates
rule it out. The determinant-sensitive benchmark action, its typing and
clock remain inherited. No historical V110 determinant--Clifford cylinder
has been supplied; historical Standard-Model selection remains unclosed.

\paragraph{Verification boundary.}
The checker reconstructs the original sign maps and enhanced pair,
verifies all-input marginal identities for the sharpness comparison,
checks the code isometry, interval-fidelity formulas, efficient comparison
normalization, and phase-dependent dual constants. General dilation
continuity is the cited theorem; the applications above have written
proofs. Exact symbolic checks are not measurements, proof-assistant
certification, or independent peer review. No global numerical optimization
of a diamond norm is used.

\begingroup
\renewcommand{\refname}{Primary-source context for V58}
\begin{thebibliography}{9}\small\raggedright
\bibitem{sel58-KSW} D. Kretschmann, D. Schlingemann, and R. F. Werner,
A continuity theorem for Stinespring's dilation,
\emph{Journal of Functional Analysis} \textbf{255} (2008), 1889--1904.
\url{https://doi.org/10.1016/j.jfa.2008.07.023};
\url{https://arxiv.org/abs/0710.2495}.
Theorem 1 supplies the common-dilation distance at most the square root of
the completely bounded distance; the extension and recorded-instrument
consequences are proved above.
\end{thebibliography}
\endgroup
\addtocontents{toc}{\protect\endgroup}
% END SOURCE-SELECTION V58 DEVELOPMENT BLOCK
% Append future development before the final document command.


% BEGIN SOURCE-SELECTION V59 DEVELOPMENT BLOCK
\clearpage
\hypersetup{pdftitle={Historical Standard-Model Source Selection: V1--V59}}
\begin{center}
{\Large\bfseries Source-selection continuation V59}\\[.4em]
{\large Mixed-outcome marginal rigidity, complete coupling kernels,\\
and the dependence on the retained temporal interface}
\end{center}
\addtocontents{toc}{\protect\begingroup}
\addtocontents{toc}{\protect\renewcommand{\protect\numberline}{\protect\makebox[2.1em][l]}}

\section[The entrance is extremality, not efficient outcomes]{A source criterion that does not require an efficient comparison}
\label{sel59:context}

V57 proves exact product reuse from an efficient recorded marginal with
independent effects. V58 permits noisy actual marginals but compares them
with an efficient reference. Neither result says that a mixed-outcome
source cannot itself determine the joint operation. The more intrinsic
question is whether its positive Choi support has a nonzero
\emph{normalization-preserving perturbation}. This continuation uses that
kernel to classify the complete joint family, not merely to choose a
rank-one approximation to the source.

The general extremality criterion for instruments and the uniqueness of
causal joint channels with an extreme marginal are established results
\cite{sel59-DPS,sel59-HHP}. They are stated and proved in the present
finite conventions, not claimed as new general principles. The
programme-level developments are the explicit tensor-kernel
parameterization, its higher-order version, the completely bounded
conditioning extension of V58, and exact evaluations on unchanged sources.
In particular, the V35 nonunital channel is rigid after all labels are
forgotten, while its two-event endpoint channel is not. Thus record
resolution and temporal resolution are separate source conditions.

No new companion file was supplied for this iteration. The directly used
antecedents are \eqref{sel36:eq:family},
Theorems~\ref{sel57:thm:polytope}, \ref{sel57:thm:bit},
\ref{sel57:thm:idle}, and \ref{sel58:thm:robust}, and
Lemma~\ref{sel58:lem:lift}. Their source constructions, clocks, and
previous error statements are retained. All joint marginal identities below
hold as maps on \emph{every} input, including reference-entangled inputs.
They are not conditional probabilities at one preparation or an assumption
about successive uses of an unobserved memory.

\section[Normalization kernels of arbitrary recorded maps]{The finite support test and its established extremality meaning}
\label{sel59:kernel}

Let \(\mathcal I=\bigoplus_a\mathcal I_a:A\to\bigoplus_a A'_a\)
be a specified normalized instrument. Its output labels remain classical.
For each nonzero branch choose a \emph{minimal} Kraus list
\(K_{a\mu}\), \(1\le\mu\le r_a=\rank J_a\), and the injective map
\(W_a e_\mu=|K_{a\mu}\rangle\!\rangle\). Thus \(J_a=W_aW_a^*\).
Put
\[
 \mathcal D_{\mathcal I}=\bigoplus_aM_{r_a},\qquad
 q_{\mathcal I}=\sum_ar_a^2,
 \qquad
 \mathcal R_{\mathcal I}(T)
   =\sum_a\Tr_{A'_a}(W_aT_aW_a^*).
\]
We use the output-first Choi convention. Explicitly,
\begin{equation}
 \mathcal R_{\mathcal I}(T)
 =\sum_{a,\mu,\nu}(T_a)_{\mu\nu}
                  (K_{a\nu}^*K_{a\mu})^{\mathsf T},\qquad
 \mathcal R_{\mathcal I}(\mathbf1)=I_A.
 \label{sel59:eq:normalization}
\end{equation}
This is a completely positive linear map on the coefficient algebra. Define
\begin{equation}
 \mathcal N_{\mathcal I}=\ker\mathcal R_{\mathcal I},\qquad
 \nu_{\mathcal I}
   =q_{\mathcal I}-\rank\mathcal R_{\mathcal I}.
 \label{sel59:eq:nullity}
\end{equation}
The complex dimension of the kernel equals the real dimension of its
Hermitian part. Kraus indices are calculation coordinates, not added
physical outcome labels. Changing minimal Kraus bases conjugates these
spaces and leaves their dimensions and positivity statements unchanged.

\begin{proposition}[Finite extremality test; standard criterion in source coordinates]
\label{sel59:prop:extreme}
The following are equivalent: \(\nu_{\mathcal I}=0\); the products
\(\{K_{a\mu}^*K_{a\nu}\}_{a,\mu,\nu}\) are linearly independent;
and \(\mathcal I\) is extreme in the convex set of normalized instruments
with these original classical labels. Equivalently, their Hilbert--Schmidt
Gram is strictly positive. In particular,
\begin{equation}
 \nu_{\mathcal I}=0\quad\Longrightarrow\quad
                    \sum_ar_a^2\le (\dim A)^2.
 \label{sel59:eq:rankbound}
\end{equation}
\end{proposition}
\begin{proof}
Transpose and exchange of matrix indices identify the first two conditions.
If \(0\ne H=H^*\in\mathcal N_{\mathcal I}\), then
\(W_a(I\pm tH_a)W_a^*\) are distinct normalized instruments for sufficiently
small \(t>0\), with average \(J_a\). Conversely, in a convex decomposition
of \(J_a\), positivity forces both summands to have support inside
\(\Ran W_a\). Their difference is \(W_aH_aW_a^*\); normalization puts
\(H\) in the kernel. Injectivity of the synthesis gives the converse.
The dimension bound follows from the codomain dimension. This is the
finite-dimensional criterion in \cite[Theorem 5]{sel59-DPS}.
\end{proof}

There is also a Kraus-free formulation: test the kernel of
\(\sum_a\Tr_{A'_a}X_a\) with
\(X_a=X_a^*\), \(X_a=P_{\supp J_a}X_aP_{\supp J_a}\).
An injective kernel test is therefore intrinsic to the actual branch maps.
It is not a claim that a small fitted Choi eigenvalue is exactly zero.
An \emph{extreme instrument} need not have efficient branches, and an
extreme channel need not preserve pure states, be reversible, or have an
invertible Liouville matrix.

\section[All joint implementations form a kernel spectrahedron]{Complete fixed-marginal freedom, not a chosen controller model}
\label{sel59:coupling}

Take a second specified instrument \(\mathcal J:B\to\bigoplus_bB'_b\)
with minimal syntheses \(V_b\) and normalization map
\(\mathcal R_{\mathcal J}\). A compatible joint instrument
\(\mathcal Q=\bigoplus_{ab}\mathcal Q_{ab}\) must satisfy
\begin{align}
 \sum_b\Tr_{B'_b}\mathcal Q_{ab}
       &=\mathcal I_a\circ\Tr_B,\nonumber\\
 \sum_a\Tr_{A'_a}\mathcal Q_{ab}
       &=\mathcal J_b\circ\Tr_A.
 \label{sel59:eq:marginals}
\end{align}
A fixed permutation of Choi tensor factors groups each local output with
its own input throughout the next statement.

\begin{theorem}[Exact mixed-outcome coupling spectrahedron]
\label{sel59:thm:couplings}
Every compatible joint instrument, and only such an instrument, has
\begin{equation}
 J(\mathcal Q_{ab})=(W_a\otimes V_b)T_{ab}(W_a\otimes V_b)^*,\qquad
 \boxed{T=\mathbf1+X\succeq0,\quad
 X=X^*\in\mathcal N_{\mathcal I}\otimes\mathcal N_{\mathcal J}.}
 \label{sel59:eq:couplings}
\end{equation}
The coefficient \(T\) is unique. The set is compact and has real affine
dimension \(\nu_{\mathcal I}\nu_{\mathcal J}\). In particular,
\begin{equation}
 \boxed{\mathcal Q=\mathcal I\otimes\mathcal J
 \text{ is the only compatible instrument}
 \iff \nu_{\mathcal I}\nu_{\mathcal J}=0.}
 \label{sel59:eq:unique}
\end{equation}
No locality of a dilation, shared-randomness form, or rank assumption on
the joint branches is imposed.
\end{theorem}
\begin{proof}
A positive operator with a marginal supported on a subspace has its full
support in that subspace tensored with the other factors: zero expectation
on the complementary projection annihilates the positive operator.
Apply this to both identities \eqref{sel59:eq:marginals}, separately to
every positive joint branch. Their intersecting supports are
\(\Ran W_a\otimes\Ran V_b\). Injectivity gives the unique positive
coefficient \(T_{ab}\).

The first marginal identities, expressed in the linearly independent
operators \(W_aE_{\mu\nu}W_a^*\), give
\((\id\otimes\mathcal R_{\mathcal J})(T)=\mathbf1\otimes I_B\).
The other identities give
\((\mathcal R_{\mathcal I}\otimes\id)(T)=I_A\otimes\mathbf1\).
Subtract the product solution \(\mathbf1\). The intersection of the two
linear kernels is exactly
\(\mathcal N_{\mathcal I}\otimes\mathcal N_{\mathcal J}\): decompose
each coefficient space into its kernel and any complementary subspace,
on which the corresponding normalization map is injective. The converse
follows by the same equations and complete positivity.

The identity is strictly positive on every coefficient block, so all
Hermitian kernel directions are feasible for small coefficients. This
proves the affine dimension and also gives nonproduct solutions when both
kernels are nonzero. Compactness follows from normalization and the
injectivity of the fixed finite synthesis: each Choi matrix has bounded
trace, and the inverse synthesis on its support is bounded. The uniqueness
implication is the causal extreme-marginal principle
\cite[Section 5]{sel59-HHP}; the displayed coordinates also describe all
nonunique cases and show necessity in this separate-input product setting.
\end{proof}

When each local branch is efficient, \(\mathcal D_{\mathcal I}\) and
\(\mathcal D_{\mathcal J}\) are commutative. Positivity in
\eqref{sel59:eq:couplings} is entrywise nonnegativity, recovering V57's
scalar coupling polytope. Mixed branches replace that polytope by a
matrix positive-semidefinite set. A single channel is an instrument with
one label; hence either extreme \emph{forgotten channel} already forces
product reuse. This does not conflict with uniqueness phenomena for two
observables on the \emph{same} input, where a product joint observable
need not exist.

\begin{corollary}[A finite witness whenever universal rigidity fails]
\label{sel59:cor:coin}
If \(\nu_{\mathcal I}>0\), a classical two-state auxiliary output suffices
to exhibit a nonproduct extension with the exact first marginal and an
input-independent fair-bit second marginal.
\end{corollary}
\begin{proof}
Choose nonzero Hermitian \(H\in\mathcal N_{\mathcal I}\), and define
normalized instruments with Choi matrices
\(W_a(I\pm tH_a)W_a^*\). Output either instrument together with its
sign, with probabilities \(1/2\). The first marginal is unchanged, each
sign has total effect \(I/2\), and the two conditional instruments differ.
A second input, if present, can be discarded. This is an alternative
extension, not an assertion that the old source exposes the sign.
\end{proof}

\section[The unchanged mixed qubit is rigid without any record]{An evaluated criterion that is stronger than branch efficiency}
\label{sel59:qubit}

Retain precisely the V35 coefficients, now forgetting their two labels:
\[
 k_+=\frac1{\sqrt{10}}\begin{pmatrix}2&2\\-1&1\end{pmatrix},\qquad
 k_-=\frac1{\sqrt{10}}\begin{pmatrix}2&-2\\1&1\end{pmatrix},\qquad
 \Phi(X)=\sum_\pm k_\pm Xk_\pm^*.
\]
\begin{theorem}[Full rigidity of the original nonunital channel]
\label{sel59:thm:old}
The forgotten channel has Choi rank two and \(\nu_\Phi=0\).
Its four-product Gram, in the order \((++,+-,-+,--)\), is
\begin{equation}
 G_{\rm prod}=\frac1{25}
 \begin{pmatrix}17&0&0&8\\0&17&-8&0\\0&-8&17&0\\8&0&0&17\end{pmatrix},
 \qquad
 \boxed{\operatorname{spec}G_{\rm prod}=\{1,1,9/25,9/25\}.}
 \label{sel59:eq:oldgram}
\end{equation}
Its determinant is \(81/625\). Any non-signalling joint instrument with
this forgotten local channel and any specified other local instrument is
the product. No original sign record, stationary reverse, or recovery
hypothesis is needed for that conclusion.
\end{theorem}
\begin{proof}
The four products are
\[
 \frac12I+\frac3{10}X,\quad
 \frac3{10}Z-\frac i2Y,\quad
 \frac3{10}Z+\frac i2Y,\quad
 \frac12I-\frac3{10}X.
\]
Their Gram is \eqref{sel59:eq:oldgram}. Apply
Proposition~\ref{sel59:prop:extreme} and
Theorem~\ref{sel59:thm:couplings}. The coefficients are unchanged and
linearly independent, so their Choi rank is two.
\end{proof}

This is not a claim of coherent recoverability: V36's unchanged recovery
error \(1-(4/5)^n\) and V37's zero stationary Hamiltonian circulation remain
valid. Convex extremality constrains joint implementations; it does not
reverse information loss. The only effect of the forgotten channel is
\(I\), so independence of effects alone cannot distinguish it from any
other one-outcome channel.

\begin{proposition}[The whole inherited nonunital family]
\label{sel59:prop:family}
For the original family \eqref{sel36:eq:family}, \(0<a,b<1\), the
forgotten channel is extreme exactly when \(a\ne b\). In either minimal
Kraus basis its four-product Gram has determinant
\begin{equation}
 \boxed{\det G_{\rm prod}(a,b)=(a-b)^4.}
 \label{sel59:eq:familydet}
\end{equation}
At \(a=b\), its normalization map has rank two and nullity two.
\end{proposition}
\begin{proof}
Use the unitary Kraus change
\(A=\diag(\sqrt{1-a},\sqrt{1-b})\),
\(B=\left(\begin{smallmatrix}0&\sqrt b\\-\sqrt a&0\end{smallmatrix}\right)\).
The determinants for the two diagonal products and the two off-diagonal
products have absolute value \(|a-b|\) each. The coefficient determinant
therefore has absolute value \((a-b)^2\), proving
\eqref{sel59:eq:familydet}. At equality, both diagonal products are
scalar, and the two off-diagonal products span one direction; both
surviving directions are nonzero. A unitary change of minimal Kraus basis
acts unitarily on the matrix-coefficient Hilbert space, preserving the
Gram determinant.
\end{proof}

\begin{theorem}[Forgetting one temporal interface reopens the coupling family]
\label{sel59:thm:powers}
For the unchanged channel in Theorem~\ref{sel59:thm:old}, every integer
\(n\ge1\) has
\begin{equation}
 J(\Phi^n)=\frac12\left(I_4+\frac35 Z\otimes I
                       -(4/5)^nY\otimes Y\right),
 \qquad
 \lambda_\pm=\frac{1\pm\sqrt{9/25+(4/5)^{2n}}}{2},
 \label{sel59:eq:powers}
\end{equation}
each eigenvalue having multiplicity two. Thus \(\rank J(\Phi)=2\),
whereas \(J(\Phi^n)\succ0\) for every \(n\ge2\). For these endpoint
channels \(\nu_{\Phi^n}=12\); the complete pair-coupling family with both
local marginals \(\Phi^n\) has real affine dimension \(144\).
\end{theorem}
\begin{proof}
The Bloch map is \((x,y,z)\mapsto(0,4y/5,3/5)\). Iteration gives
\eqref{sel59:eq:powers}. The two traceless terms anticommute and square
to scalars, giving the eigenvalues. At \(n=1\) the square root is one;
it is smaller than one for \(n\ge2\). Full Choi support has coefficient
space dimension sixteen. Its normalization map has rank four, because
synthesis is invertible and partial trace maps onto \(M_2\). The
coupling dimension follows from Theorem~\ref{sel59:thm:couplings}.
\end{proof}

For a concrete endpoint counter-completion, let \(J_2=J(\Phi^2)\) and
\(D=Z\otimes Z\) in output-first local Choi coordinates. Then
\(\Tr_{\rm out}D=0\), and
\[
 J_2\otimes J_2+\frac1{1024}D\otimes D
\]
is positive and has exactly the two endpoint marginals. Indeed
\(\lambda_{\min}(J_2)=(25-\sqrt{481})/50>1/20\), so the positive product
floor exceeds the perturbation norm. It is not the product channel.
This does not provide an implementation with the same original one-event
interfaces: if each complete intermediate bipartite step satisfies those
exact one-event marginals, each step is a product. Endpoint tomography
alone omits that stronger temporal information.

\section[Recorded, forgotten, and imperfectly recorded sign sources]{The effect census must be replaced by a complete Kraus-product census}
\label{sel59:sign}

On the unchanged V57 three-level forward port, retain
\(L_\pm=(I/\sqrt2\pm U)/\sqrt3\), \(U=P_0-iP_1\),
\(\rank P_0=2\), \(\rank P_1=1\). Let \(\mathcal L_\pm\) denote
the two CP branches and \(\Phi_{\rm s}=\mathcal L_++\mathcal L_-\).
Define a \emph{declared} classical misidentification comparison by
\[
 \mathcal L^{(\delta)}_+=(1-\delta)\mathcal L_++\delta\mathcal L_-,
 \qquad
 \mathcal L^{(\delta)}_-=\delta\mathcal L_++(1-\delta)\mathcal L_-.
\]

\begin{proposition}[Exact dependence on the retained record policy]
\label{sel59:prop:sign}
The dimensions of the complete fixed-local-map pair families are as follows:
\begin{center}\small
\begin{tabular}{L{.40\textwidth}ccc}
\toprule
Original or explicitly processed object & \(q\) & \(\nu\) & Pair dimension\\
\midrule
Original sign instrument & 2&0&0\\
All signs forgotten, \(\Phi_{\rm s}\)&4&2&4\\
Misidentified signs, \(0<\delta<1\)&8&6&36\\
\bottomrule
\end{tabular}\end{center}
For \(\delta\ne1/2\), the two \emph{effects} of the misidentified
instrument are still independent. Thus effect independence is insufficient
for mixed branches, while the full product-kernel test remains exact.
\end{proposition}
\begin{proof}
The original effects are independent by Theorem~\ref{sel57:thm:bit}.
The forgotten channel has minimal Kraus span \(\Span\{P_0,P_1\}\),
and its products span exactly that two-dimensional algebra. For
\(0<\delta<1\), each recorded branch has that same two-dimensional
Kraus span, so \(q=2\cdot2^2=8\), and the normalization map still has
rank two. The effects are
\(I/2\pm(1-2\delta)\sqrt2 P_0/3\). Apply
Theorem~\ref{sel59:thm:couplings}.
\end{proof}

The four-dimensional forgotten-channel family is the family already
parameterized by V54. Its dimension is not a new optimization. Here it is
recovered from the normalization kernel without choosing a correlated
mixture model. The thirty-six-dimensional comparison keeps a \emph{noisy}
classical sign in the output, not the true original sign. A linearly
invertible confusion matrix does not physically restore the latter.
The family dimension jumps at exact zero confusion, while operational
errors need not jump. The source and marginal precision still constrain
how far feasible channels lie from the original product, as in V58 and the
next theorem. None of these comparisons is silently assigned to the
historical instrument.

\section[Mixed extreme references also have noisy product bounds]{The inverse complementary-observable norm controls stability}
\label{sel59:robust}

Define the unital completely positive normalization-adjoint map
\[
 \Gamma_{\mathcal I}(T)
   =\sum_{a,\mu,\nu}(T_a)_{\mu\nu}K_{a\mu}^*K_{a\nu}.
\]
It differs from \(\mathcal R_{\mathcal I}\) only by transposes on both
matrix spaces. Suppose \(\nu_{\mathcal I}=0\), and set
\begin{equation}
 \kappa_{\mathcal I}
 =\|\Gamma_{\mathcal I}^{-1}:\Ran\Gamma_{\mathcal I}
                       \to\mathcal D_{\mathcal I}\|_{\rm cb}.
 \label{sel59:eq:kappa}
\end{equation}
The range inherits its concrete matrix norms. The inverse is only a
bounded linear inference map, not a positive physical operation.
The constant is finite. For trace-duals
\(\Tr(D_{a\mu\nu}K_{b\kappa}^*K_{b\lambda})
=\delta_{ab}\delta_{\mu\kappa}\delta_{\nu\lambda}\), a computable bound is
\[
 \kappa_{\mathcal I}
 \le\max_a\sum_{\mu,\nu}\|D_{a\mu\nu}\|_1.
\]
For efficient branches this reduces to V57's effect-dual bound.

\begin{theorem}[Two noisy marginals with a mixed extreme reference]
\label{sel59:thm:robust}
Let the actual classical-output joint instrument \(\mathcal Q\) have
marginal errors \(\epsilon,\eta\) as in \eqref{sel58:eq:errors}, but let
\(\mathcal I\) be any extreme instrument, not necessarily efficient.
Then
\begin{equation}
 \boxed{
 \|\mathcal Q-\mathcal I\otimes\Psi\|_\diamond
 \le\min\{2,\ 2(1+\kappa_{\mathcal I})\sqrt\epsilon
                         +\kappa_{\mathcal I}\eta\}.}
 \label{sel59:eq:robust}
\end{equation}
There is no rank assumption on the actual maps and no dependence on the
second input, output, or hidden-environment dimension outside the
specified norm errors. No claim of an optimal prefactor is made.
\end{theorem}
\begin{proof}
First keep the first marginal exact. Minimal dilation uniqueness for
each branch gives channels
\(\Lambda_a:M_{r_a}\otimes B\to B'\) with
\[
 \mathcal Q_a(X)=(\id_{A'_a}\otimes\Lambda_a)
                 ((V_a\otimes I)X(V_a^*\otimes I)),
 \qquad V_a\xi=\sum_\mu K_{a\mu}\xi\otimes|\mu\rangle.
\]
One can also derive this directly from the Choi support: expansion in the
independent \(|K_{a\mu}\rangle\!\rangle\) makes the first marginal
normalization exactly the trace-preserving identity for \(\Lambda_a\).
The combined original dilation has \(\sum_aV_a^*V_a=I\).

In the Heisenberg picture let
\(\Lambda^*(Y)=\bigoplus_a\Lambda_a^*(Y)\), and
\(\Lambda_0^*(Y)=\mathbf1\otimes\Psi^*(Y)\). The second marginal defect is
\((\Gamma_{\mathcal I}\otimes\id_B)(\Lambda^*-\Lambda_0^*)\).
By \eqref{sel59:eq:kappa} and amplification, its inverse gives
\(\|\Lambda^*-\Lambda_0^*\|_{\rm cb}\le\kappa_{\mathcal I}\eta\).
Amplifying each block by \(A'_a\), then compressing by the normalized
original dilation, cannot increase this norm. Duality with the diamond
norm proves the exact-first-marginal bound.

For a noisy first marginal, use the unchanged
Lemma~\ref{sel58:lem:lift} and its classical-flag pinching. The comparison
has first marginal exact and lies at distance at most \(2\sqrt\epsilon\)
from the actual source. Its second error is at most
\(\eta+2\sqrt\epsilon\). Apply the preceding bound and the triangle
inequality. The cap two is channel normalization. The comparison and
inverse are not implemented on the actual source.
\end{proof}

\begin{proposition}[Exact conditioning for the unchanged mixed qubit]
\label{sel59:prop:kappa}
For Theorem~\ref{sel59:thm:old},
\begin{equation}
 \boxed{\kappa_\Phi=5/3,\qquad
 \|\mathcal Q-\Phi\otimes\Psi\|_\diamond
 \le\min\{2,(16/3)\sqrt\epsilon+(5/3)\eta\}.}
 \label{sel59:eq:oldrobust}
\end{equation}
\end{proposition}
\begin{proof}
In the original two-Kraus basis,
\(\Gamma(I)=I\), \(\Gamma(X)=3Z/5\),
\(\Gamma(Y)=-Y\), \(\Gamma(Z)=3X/5\).
If \(H=(X+Z)/\sqrt2\) and
\(\mathcal D_Y=(4/5)\id+(1/5)\operatorname{Ad}_Y\), then
\(\Gamma=\operatorname{Ad}_H\mathcal D_Y\).
The inverse of \(\mathcal D_Y\) is
\((4/3)\id-(1/3)\operatorname{Ad}_Y\), of completely bounded norm at most
\(5/3\). Its action on \(X\), of operator norm one, has norm \(5/3\),
proving equality. Insert the value into \eqref{sel59:eq:robust}.
\end{proof}

An unknown measured source does not become extreme because a fitted
reference is extreme. The fixed-rank product Gram can give a one-sided
conditioning certificate when its support is independently established;
otherwise the norm errors in \eqref{sel59:eq:robust} must pay for the
comparison. Full Choi rank filling can restore nontrivial kernels even
arbitrarily close to an extreme source. V58's square-root-order example
remains valid in this larger theorem class.

\section[Which higher-order dependence can lower-order marginals miss?]{A tensor-kernel hierarchy with exact source marginals}
\label{sel59:higher}

For \(n\) specified local instruments let
\(\mathcal D_i=\mathcal N_i\oplus\mathcal V_i\) be any
adjoint-stable linear splitting, with the normalization map injective on
\(\mathcal V_i\). Suppose every \(k\)-slot marginal is the specified
product instrument, on every global input, where \(1\le k<n\).

\begin{theorem}[Complete higher-order coupling space]
\label{sel59:thm:higher}
In the same minimal coefficient coordinates, all compatible joint maps
are exactly \(T=\mathbf1+X\succeq0\), where \(X=X^*\) belongs to
\begin{equation}
 \boxed{
 \bigoplus_{\substack{S\subseteq\{1,\ldots,n\}\\|S|\ge k+1}}
 \left(\bigotimes_{i\in S}\mathcal N_i\right)
 \otimes\left(\bigotimes_{i\notin S}\mathcal V_i\right).}
 \label{sel59:eq:higher}
\end{equation}
Its real affine dimension is
\(\sum_{|S|\ge k+1}\prod_{i\in S}\nu_i
                      \prod_{i\notin S}(q_i-\nu_i)\).
In particular, if at most \(k\) local instruments are nonextreme, the
stated \(k\)-slot marginal identities force the entire product channel.
\end{theorem}
\begin{proof}
Positive marginal support restricts the total Choi matrix to the tensor
product of the local supports, exactly as in
Theorem~\ref{sel59:thm:couplings}. Retaining slots \(A\) and tracing their
complement applies the identity on \(\mathcal D_i\), \(i\in A\), and
the normalization map elsewhere. A component whose kernel-factor set is
\(S\) survives this map exactly when \(S\subseteq A\). The surviving
components are independent under the chosen splittings, because the
normalization maps are injective on the complements. Requiring zero
defect for every \(|A|=k\) therefore removes exactly the components with
\(|S|\le k\). Conversely, every remaining component vanishes in every
such marginal. Positivity, the interior product point, and the finite
synthesis prove the remaining statements.
\end{proof}

For an explicit comparison, take the qubit dephasing channel
\(\Delta=(\id+\operatorname{Ad}_Z)/2\). For \(|t|\le1\), set
\[
 \mathcal E_t=\frac18\sum_{s\in\{0,1\}^3}
 [1+t(-1)^{s_1+s_2+s_3}]
 \operatorname{Ad}_{Z^{s_1}\otimes Z^{s_2}\otimes Z^{s_3}}.
\]
Every pair marginal is \(\Delta\otimes\Delta\) on every joint input,
yet \(\|\mathcal E_t-\Delta^{\otimes3}\|_\diamond=|t|\).
The upper bound is the sum of the absolute coefficients; the input
\(|+\rangle^{\otimes3}\) produces orthogonal \(X\)-basis outputs and
attains it. This elementary example illustrates a genuine three-kernel
component. It is not a correlated implementation assigned to either
unchanged benchmark, and pairwise marginal equality is not a surrogate
for the full source table.

\section[Source audit and proof status]{What has actually moved upstream}
\label{sel59:status}

\paragraph{Proved.}
The complete fixed-marginal joint family is a positive tensor-kernel set
for arbitrary branch ranks. Vanishing of either local normalization kernel
is exactly product uniqueness in the separate-input setting. The same
kernel describes which higher-order correlations evade lower-order map
data. V58's norm estimate extends to mixed extreme reference instruments
through an inverse complementary-observable norm. The unchanged V35
forgotten channel is extreme, with \(\kappa=5/3\); its two-event endpoint
channel has a nontrivial, explicitly evaluated coupling family. The
original, forgotten, and misidentified sign objects have distinct exact
kernel dimensions.

\paragraph{Inherited inputs.}
Instrument extremality, extreme-marginal uniqueness, positive support
factorization, and dilation continuity are established mathematics. The
original benchmark maps and their earlier stationary, recovery, and
coherence conclusions are unchanged. No positive statement here uses a
new grading, determinant tensor, controller, or reverse operation. Formal
Kraus changes and inverse normalization maps are inference coordinates,
not new physical refinements.

\paragraph{Remaining source question.}
For an actual all-input local map, test its full supported normalization
kernel rather than infer reuse freedom from its Choi rank or its effects
alone. State explicitly whether the retained interface is an original
event, a noisy recorded event, or a multi-event endpoint channel. An
exact instrument kernel cannot be certified by rounding small numerical
eigenvalues to zero. Where only a calibrated comparison is justified,
retain the errors in \eqref{sel59:eq:robust}, and retain success denominators
under subsequent conditioning. The raw idle branch is never deleted: its
positive-normalization dependencies and acceptance correlations remain
part of the same kernel calculation. No joint clock law follows from a
local marginal.

The historical primitive channel, its directed V110 typing, and its
same-history determinant--Clifford word have not been populated. Historical
Standard-Model selection is not declared closed. This continuation removes
an unnecessarily restrictive \emph{efficiency} premise from the reuse
criterion; it does not add extremality as an axiom of the finite duality.

\paragraph{Verification boundary.}
The script checks all newly displayed finite matrices, product-kernel
ranks, positive counter-completions, their complete matrix-unit marginals,
the nonunital-family determinant, the old channel powers, the exact
inverse-norm witness, and a pairwise-product/triple-distinct source. The
general theorems have the written proofs above. No physical measurements,
proof-assistant formalization, or independent peer review is claimed.

\begingroup
\renewcommand{\refname}{Primary-source context for V59}
\begin{thebibliography}{9}\small\raggedright
\bibitem{sel59-DPS} G. M. D'Ariano, P. Perinotti, and M. Sedl\'ak,
Extremal quantum protocols,
\emph{Journal of Mathematical Physics} \textbf{52} (2011), 082202.
\url{https://doi.org/10.1063/1.3610676};
\url{https://arxiv.org/abs/1101.4889}.
Theorem 5 supplies the established instrument Kraus-product extremality
criterion. The finite proof and the source calculations are given here.
\bibitem{sel59-HHP} E. Haapasalo, T. Heinosaari, and J.-P. Pellonp\"a\"a,
When do pieces determine the whole? Extreme marginals of a completely
positive map,
\emph{Reviews in Mathematical Physics} \textbf{26} (2014), 1450002.
\url{https://doi.org/10.1142/S0129055X14500020};
\url{https://arxiv.org/abs/1209.5933}.
Section 5 proves the established causal extreme-marginal product principle.
The tensor-kernel description and evaluated interfaces above retain this
context rather than claiming general priority.
\end{thebibliography}
\endgroup
\addtocontents{toc}{\protect\endgroup}
% END SOURCE-SELECTION V59 DEVELOPMENT BLOCK
% Append future development before the final document command.


% BEGIN SOURCE-SELECTION V60 DEVELOPMENT BLOCK
\clearpage
\hypersetup{pdftitle={Historical Standard-Model Source Selection: V1--V60}}
\begin{center}
{\Large\bfseries Source-selection continuation V60}\\[.4em]
{\large Choi-support extremality, basis-free reuse rigidity,\\
and the exact source cost of record erasure}
\end{center}
\addtocontents{toc}{\protect\begingroup}
\addtocontents{toc}{\protect\renewcommand{\protect\numberline}{\protect\makebox[2.1em][l]}}

\section[Move upstream from Kraus products to Choi support]{The reuse criterion is intrinsic to the exact branch supports}
\label{sel60:context}

V59 expresses source-reuse freedom through the normalization kernel of a
minimal Kraus synthesis.  That criterion is invariant under Kraus changes, but
the formulas can still suggest that one must reconstruct a preferred Kraus
frame.  This continuation removes that appearance completely.  The kernel,
its dimension, and the resulting product-rigidity alternative depend only on
the exact Choi support subspaces of the recorded branches.  Branch weights and
all nonzero eigenvalues inside those supports are irrelevant to the yes/no
extremality question.

This is useful for the historical programme for two reasons.  First, V24--V25
already reduce an actual same-cut process packet to its Choi matrices whenever
the admitted frame or contextual cylinder table is complete on the required
support.  No further Kraus reconstruction is needed.  Second, support loss or
support filling is a sharp source event: it can change reuse rigidity even when
the channel varies continuously.  A small fitted eigenvalue is therefore not
silently rounded to zero below.

No new companion file is used in V60.  The directly inherited inputs are
V57's exact independence of the six original compiler effects, V57's raw idle
coupling, V59's mixed-outcome coupling spectrahedron, and the V24--V25 Choi
reconstruction interfaces.  The established extremality criterion remains
credited to the sources cited in V59; the new statements are the support-only
compiler, its direct positive Laplacian, the one-sided historical obstruction,
and the exact evaluations below.

\section[The support-only normalization map]{A basis-free extremality and reuse certificate}
\label{sel60:support}

Let the specified normalized instrument have Choi matrices
\(
J_a\succeq0
\)
on output spaces \(A'_a\) and one common input space \(A\).  Put
\[
 P_a=P_{\supp J_a},\qquad r_a=\rank J_a,
\]
and define the finite operator space
\[
 \mathfrak S_{\mathbf P}
 =\bigoplus_a P_a\,\mathcal B(A'_a\otimes A)\,P_a.
\]
On it define
\begin{equation}
 \boxed{
 \Pi_{\mathbf P}\bigl((X_a)_a\bigr)
 =\sum_a\Tr_{A'_a}X_a.}
 \label{sel60:eq:Pi}
\end{equation}
The harmless transpose used by the output-first Choi convention in V59 is
suppressed here; transpose is an invertible Hilbert--Schmidt isometry and does
not change any kernel or rank statement.

\begin{theorem}[Choi-support sufficiency for normalization rigidity]
\label{sel60:thm:support}
Let \(\mathcal N_{\mathcal I}\) and \(\nu_{\mathcal I}\) be V59's
normalization kernel and its dimension.  Then there is a canonical linear
isomorphism
\[
 \boxed{
 \mathcal N_{\mathcal I}\cong\ker\Pi_{\mathbf P},
 \qquad
 \nu_{\mathcal I}=\dim\ker\Pi_{\mathbf P}.}
\]
Consequently the following data depend only on the exact support projectors
\(P_a\):
\begin{enumerate}[label=\textup{(\roman*)},leftmargin=2.8em]
\item whether the recorded instrument is extreme;
\item whether it forces product reuse against every specified second marginal;
\item the affine dimension of V59's compatible pair family;
\item the dimension of every higher-order normalization-kernel tensor sector.
\end{enumerate}
In particular, two normalized instruments with the same branch support
projectors have the same normalization nullity, even if their positive
spectra inside those supports are different.
\end{theorem}
\begin{proof}
Choose any minimal Kraus synthesis \(W_a:\mathbb C^{r_a}\to
A'_a\otimes A\), so \(J_a=W_aW_a^*\).  Its polar decomposition has the form
\[
 W_a=J_a^{1/2}U_a,
\]
where \(U_a:\mathbb C^{r_a}\to\Ran P_a\) is unitary.  Since
\(J_a^{1/2}\) is invertible on \(\Ran P_a\), the congruence
\[
 \Xi_a(T_a)=W_aT_aW_a^*
\]
is a linear isomorphism from \(M_{r_a}\) onto
\(P_a\mathcal B(A'_a\otimes A)P_a\).  The normalization map of V59 is
exactly \(\Pi_{\mathbf P}\circ\bigoplus_a\Xi_a\), up to the displayed
transpose convention.  The kernel statement follows.  V59's extremality and
coupling theorems then give the remaining conclusions.
\end{proof}

There is therefore no physical content in choosing an eigen-Kraus basis for
this audit.  One may diagonalize \(J_a\), use a square root, or use any other
minimal factorization; all are calculation gauges on one fixed support.

\begin{proposition}[Support Laplacian and a direct positive certificate]
\label{sel60:prop:laplacian}
With Hilbert--Schmidt inner products, the adjoint of
\(\Pi_{\mathbf P}\) is
\begin{equation}
 \boxed{
 \Pi_{\mathbf P}^*(Y)
 =\bigl(P_a(I_{A'_a}\otimes Y)P_a\bigr)_a.}
 \label{sel60:eq:adjoint}
\end{equation}
Hence the positive support Laplacian
\begin{equation}
 \boxed{
 \mathscr L_{\mathbf P}=\Pi_{\mathbf P}^*\Pi_{\mathbf P}\succeq0}
 \label{sel60:eq:laplacian}
\end{equation}
satisfies
\[
 \langle X,\mathscr L_{\mathbf P}X\rangle
 =\left\|\sum_a\Tr_{A'_a}X_a\right\|_{\HS}^2,
 \qquad
 \ker\mathscr L_{\mathbf P}=\ker\Pi_{\mathbf P}.
\]
Thus exact product rigidity is equivalent to strict positivity of
\(\mathscr L_{\mathbf P}\) on \(\mathfrak S_{\mathbf P}\).
Its least eigenvalue on that finite space is a basis-free conditioning margin.
\end{proposition}
\begin{proof}
For \(X_a=P_aX_aP_a\), cyclicity of trace gives
\[
 \langle Y,\Tr_{A'_a}X_a\rangle
 =\Tr\!\left[(I_{A'_a}\otimes Y)^*X_a\right]
 =\langle P_a(I_{A'_a}\otimes Y)P_a,X_a\rangle.
\]
Summing over \(a\) proves the adjoint formula.  The quadratic identity and
kernel statement follow from \(\mathscr L=\Pi^*\Pi\).
\end{proof}

The conditioning margin is an identification quantity.  It is not a physical
spectral gap, relaxation rate, or mass.

\begin{corollary}[Support-stratum invariance and boundary discontinuity]
\label{sel60:cor:stratum}
Along any family of normalized instruments for which all exact support
projectors \(P_a\) remain fixed, \(\nu_{\mathcal I}\), extremality, and product
uniqueness remain fixed.  They can change only when at least one exact Choi
support changes.
\end{corollary}

This explains why infinitesimal full-rank contamination can destroy an exact
extremality statement while changing the channel continuously.  No observable
discontinuity of the channel is implied.

\section[Rank alone gives one-sided source no-goes]{A support lower bound can exclude rigidity without reconstructing Kraus products}
\label{sel60:rank}

Let \(d=\dim A\) and \(q=\sum_ar_a^2=\dim\mathfrak S_{\mathbf P}\).
Since \(\Pi_{\mathbf P}\) maps into a \(d^2\)-dimensional space,
\begin{equation}
 \boxed{
 \nu_{\mathcal I}\ge q-d^2
 =\sum_ar_a^2-d^2.}
 \label{sel60:eq:ranklower}
\end{equation}
The positive part of the right-hand side is therefore an immediate lower bound
on the dimension of the joint-source ambiguity.

\begin{proposition}[Full-support channel]
\label{sel60:prop:full}
For a one-outcome channel from a \(d_{\rm in}\)-dimensional input to a
\(d_{\rm out}\)-dimensional output with strictly positive Choi matrix,
\begin{equation}
 \boxed{
 \nu=(d_{\rm out}^2-1)d_{\rm in}^2.}
 \label{sel60:eq:full}
\end{equation}
\end{proposition}
\begin{proof}
Now \(P=I_{A'\otimes A}\), so \(\Pi=\Tr_{A'}\) on the complete matrix
algebra.  It is onto: \(Y\mapsto |0\rangle\langle0|\otimes Y\) is a right
inverse.  The domain dimension is \(d_{\rm out}^2d_{\rm in}^2\) and the range
dimension is \(d_{\rm in}^2\).
\end{proof}

For a square qubit channel, every strictly positive Choi packet therefore has
\[
 \boxed{\nu=16-4=12.}
\]
This recovers V59's endpoint result for \(\Phi^n\), \(n\ge2\), from support
alone.  No four-product Gram needs to be recomputed at each time.

\begin{corollary}[Certified support lower bounds already give a robust exclusion]
\label{sel60:cor:lower}
Suppose an experiment or exact source theorem certifies subspaces
\(F_a\subseteq\Ran J_a\) of dimensions \(s_a\).  Then
\[
 \boxed{
 \nu_{\mathcal I}\ge\sum_as_a^2-d^2.}
\]
If the right-hand side is positive, product rigidity is impossible regardless
of every unresolved positive eigenvalue and every unresolved enlargement of the
supports.
\end{corollary}

This is one-sided.  Certifying \(\nu=0\) is stronger: exact support ranks and the
strict positivity of the support Laplacian are needed.  A numerically small
Choi eigenvalue or a numerically small Laplacian eigenvalue is not set to zero.

\section[Exact evaluations on unchanged source interfaces]{Record erasure, temporal coarse graining, and the raw idle branch}
\label{sel60:evaluations}

\subsection{The unchanged nonunital qubit and its endpoint channel}

For V35's forgotten channel, V59 already proves that its two-dimensional Choi
support has injective normalization map.  Theorem~\ref{sel60:thm:support}
therefore strengthens the interpretation:
\[
 \boxed{
 \text{every normalized channel with exactly that same Choi support is extreme.}}
\]
The conclusion does not depend on V35's particular nonzero Choi eigenvalues.

By contrast, V59's \(\Phi^n\), \(n\ge2\), has full qubit Choi support.  Hence
Proposition~\ref{sel60:prop:full} gives \(\nu=12\) immediately.  Temporal
coarse graining has changed the support stratum, not merely the weights inside
one stratum.

\subsection{The original sign record on the forward port}

Retain V57's two efficient sign branches \(L_+,L_-\) on
\(H_3=\operatorname{span}\{n,\omega,u_0\}\).  Their effects are independent,
so the recorded sign instrument has
\[
 \boxed{\nu_{\rm sign}=0.}
\]
Forgetting the sign merges the two coefficient rays into one Choi support of
dimension two.  V59 gives
\[
 \boxed{\nu_{\rm forgotten}=2.}
\]
The new theorem shows that these values are properties of the two exact support
arrangements.  They do not depend on how the positive weights are distributed
inside those supports.

Now retain the idle branch at a raw opportunity.  For \(0<\vartheta<1\), the
three efficient branch supports are the idle identity ray and the two sign
rays.  Their normalized effects lie in the two-dimensional span generated by
\(I_3\) and the inherited projector \(P_0\).  Hence
\begin{equation}
 \boxed{\nu_{\rm raw}=1.}
 \label{sel60:eq:rawnullity}
\end{equation}
This is the support-only origin of V57's single raw acceptance-correlation
parameter.  Changing \(\vartheta\) inside \((0,1)\) changes the physical branch
weights but not this nullity.  At \(\vartheta=1\), the idle support disappears
and the recorded accepted instrument returns to \(\nu=0\).

\subsection{The full six-outcome compiler has the same raw one-dimensional freedom}

V57 proves on the complete native carrier that the six accepted branch effects
are linearly independent and sum to identity.  Therefore the accepted
six-outcome instrument has \(\nu=0\).  Add the raw idle coefficient
\(\sqrt{1-\vartheta}\,I\), \(0<\vartheta<1\).  Its effect lies in the span of
the six accepted effects, and there is exactly one relation:
\[
 (1-\vartheta)I
 -\frac{1-\vartheta}{\vartheta}
   \sum_{e=1}^{6}\vartheta K_e^*K_e=0.
\]
Thus
\begin{equation}
 \boxed{
 \nu_{\rm six,accepted}=0,
 \qquad
 \nu_{\rm six,raw}=1.}
 \label{sel60:eq:sixraw}
\end{equation}
The one-dimensional joint-source freedom identified earlier is therefore not an
artifact of the three-level restriction.  It is already present on the complete
native raw opportunity, and it is caused by the extra idle support relation.

This statement does not determine the joint clock law or the value of the raw
acceptance-correlation parameter.  It identifies only the exact dimension of
the allowed fixed-marginal freedom.

\section[The Choi packet itself returns a counter-completion]{No Kraus frame is needed on the nonextreme branch}
\label{sel60:witness}

The support formulation also gives a direct witness when rigidity fails.
Let \(0\ne X=(X_a)_a=X^*\in\ker\Pi_{\mathbf P}\).  Define
\[
 \lambda_*
 =\min_{a:X_a\ne0}
 \frac{\lambda_{\min}(J_a|_{\Ran P_a})}{\|X_a\|_{\op}}.
\]
Every numerator is strictly positive on its exact support.

\begin{proposition}[Basis-free positive counter-decomposition]
\label{sel60:prop:counter}
For every \(0<t<\lambda_*\),
\begin{equation}
 \boxed{
 J_a^{\pm}=J_a\pm tX_a\succeq0,
 \qquad
 \sum_a\Tr_{A'_a}J_a^{\pm}=I_A.}
 \label{sel60:eq:counter}
\end{equation}
The two normalized instruments are distinct and average to the original one.
Thus a nonzero support-kernel vector is itself an explicit physical
nonextremality witness, without choosing a Kraus basis.
\end{proposition}
\begin{proof}
On \(\Ran P_a\),
\(J_a\succeq\lambda_{\min}(J_a|_{\Ran P_a})P_a\).  Hence
\(J_a\pm tX_a\succeq0\) for the displayed range of \(t\).  Outside the
support both terms vanish.  The trace-preserving identity follows from
\(X\in\ker\Pi_{\mathbf P}\).  Nonzero \(X\) makes the two instruments distinct.
\end{proof}

This is a source counter-completion, not an assertion that the original device
classically chooses between the two decomposed instruments on each run.

\section[Historical compilation and revised source frontier]{What the existing cylinder table has to determine}
\label{sel60:status}

V24--V25 already show how an actual same-cut Choi packet is reconstructed when
the historical intervention frame or prefix--suffix contextual frame is
complete on the relevant memory.  V60 removes the next apparent requirement:
\emph{after those Choi matrices are known, no Kraus-frame reconstruction is
needed for the reuse-rigidity test.}

The finite historical compiler is now
\[
 \boxed{
 \{J_a\}
 \longrightarrow
 \{P_{\supp J_a}\}
 \longrightarrow
 \mathscr L_{\mathbf P}
 \longrightarrow
 \begin{cases}
 \ker\mathscr L_{\mathbf P}=0:&\text{forced product reuse},\\
 \ker\mathscr L_{\mathbf P}\ne0:&\text{explicit coupling directions}.
 \end{cases}}
\]
When only support lower bounds are certified, the dimension inequality
\eqref{sel60:eq:ranklower} can already prove that rigidity is impossible.
Positive rigidity still requires the exact support projectors or another exact
source theorem that fixes them.

This also identifies what record loss actually changes.  Forgetting a label can
merge Choi supports and expose new off-diagonal normalization directions;
coarse temporal observation can fill the entire Choi space; adding an idle
branch can add one exact support relation.  These are different support-level
operations and should not be summarized merely as ``more noise.''

\paragraph{Proved.}
The normalization kernel and all V59 fixed-marginal reuse dimensions depend
only on the exact branch Choi supports.  The basis-free support Laplacian gives
the direct finite certificate.  Full Choi support gives an exact nullity
formula.  The original nonunital qubit, its multi-event endpoint, the retained
sign source, the forgotten sign channel, and the raw idle interfaces are all
re-evaluated in these coordinates.  A nonzero support-kernel vector directly
produces positive counter-completions.

\paragraph{Inherited inputs.}
V24--V25 supply the process-tomography interfaces; V57 supplies exact effect
independence for the original compiler; V59 supplies the normalization-kernel
and joint-coupling theorems.  No new grading, determinant tensor, readout,
controller, or clock is added.

\paragraph{Remaining source question.}
The historical primitive Choi packets themselves still have to be populated on
the correct typed carrier.  Exact support is a source statement.  Numerical
near-nullity does not establish an exact support drop, while a certified support
lower bound can already establish a one-sided nonrigidity obstruction.  The raw
idle branch and its acceptance relation remain part of the same calculation.
No joint clock law follows from the local supports.

Historical Standard-Model selection is not declared closed.  The directed V110
determinant--Clifford word and its same-cut typed incidence remain unpopulated.
The present continuation moves the reuse audit upstream from a chosen Kraus
representation to the Choi support already reconstructed by the operational
programme.

\paragraph{Verification boundary.}
The companion checker evaluates the support-only maps directly from Choi column
spaces, reproduces the unchanged extreme qubit and full-rank endpoint nullities,
checks the sign/forgotten/raw-idle hierarchy, and verifies positive
counter-decompositions.  The general theorems have the written proofs above.
No physical measurement, proof-assistant formalization, or independent peer
review is claimed.

\addtocontents{toc}{\protect\endgroup}
% END SOURCE-SELECTION V60 DEVELOPMENT BLOCK
% Append future development before the final document command.


% BEGIN SOURCE-SELECTION V61 DEVELOPMENT BLOCK
\clearpage
\hypersetup{pdftitle={Historical Standard-Model Source Selection: V1--V61}}
\begin{center}
{\Large\bfseries Source-selection continuation V61}\\[.4em]
{\large Rigid support envelopes, positive-zero witnesses,\\
and why approximate tomography cannot certify exact reuse}
\end{center}
\addtocontents{toc}{\protect\begingroup}
\addtocontents{toc}{\protect\renewcommand{\protect\numberline}{\protect\makebox[2.1em][l]}}

\section[Move upstream from exact support reconstruction]{Exact support projectors are stronger than the positive branch needs}
\label{sel61:context}

V60 moved the reuse audit from a chosen Kraus frame to the exact branch Choi
supports.  That is intrinsic, but it still leaves an apparently demanding
historical instruction: reconstruct every exact support projector and then test
the support Laplacian.  The present continuation separates two logically
different tasks.

For a \emph{positive} rigidity claim, the exact support need not be known.  It
is enough to know an exact \emph{support envelope} inside which the historical
Choi packets must lie.  If the normalization map is injective on that larger
envelope, normalization already fixes the entire instrument.  Positive
zero-witnesses can certify such an envelope without full process tomography.
For a \emph{negative} rigidity claim, by contrast, certified support lower
bounds remain enough through V60's dimension obstruction.

This asymmetry matters because exact extremality has no open robustness radius
in channel norm.  Arbitrarily small full-support contamination can turn a
rigid instrument into a nonextreme one.  Consequently a small numerical Choi
eigenvalue, or a small measured leakage, is not promoted below into an exact
zero.  Exact positive closure must come from an algebraic source identity, an
exact positive-zero witness, or an equivalent structural statement.  Finite
precision remains useful for one-sided nonrigidity and approximate leakage
bounds.

No new source, grading, determinant tensor, controller, or clock is introduced
in V61.  The unchanged examples are V35's nonunital qubit, V57's sign and raw
compiler interfaces, V59's normalization-kernel theorem, and V60's support
Laplacian.

\section[Rigid support envelopes identify the entire source]{Normalization plus an injective envelope is a singleton theorem}
\label{sel61:envelope}

Let \(Q_a\) be arbitrary orthogonal projectors on the branch Choi spaces
\(A'_a\otimes A\).  They are only assumed to be \emph{upper} support bounds;
they need not equal the eventual supports.  Put
\[
 \mathfrak S_{\mathbf Q}
 =\bigoplus_a Q_a\mathcal B(A'_a\otimes A)Q_a,
 \qquad
 \Pi_{\mathbf Q}((X_a)_a)=\sum_a\Tr_{A'_a}X_a.
\]

\begin{theorem}[Rigid support-envelope theorem]
\label{sel61:thm:envelope}
Assume
\begin{equation}
 \boxed{\ker\Pi_{\mathbf Q}=\{0\}.}
 \label{sel61:eq:envelopeinjective}
\end{equation}
Then there is \emph{at most one} normalized instrument \((J_a)_a\) satisfying
\[
 J_a\succeq0,
 \qquad
 \supp J_a\subseteq\Ran Q_a,
 \qquad
 \sum_a\Tr_{A'_a}J_a=I_A.
\]
Hence, if one normalized candidate \((J_a^{(0)})_a\) with those support bounds
is known to exist, every normalized instrument obeying the same bounds equals
that candidate branch by branch.

In particular, an injective support envelope determines the positive spectra
inside the envelope as well as the exact supports; neither has to be supplied
separately.
\end{theorem}
\begin{proof}
Let \(J\) and \(K\) be two normalized instruments obeying the support bounds.
Then \(X_a=J_a-K_a\) belongs to
\(Q_a\mathcal B(A'_a\otimes A)Q_a\) and
\[
 \Pi_{\mathbf Q}(X)=
 \sum_a\Tr_{A'_a}(J_a-K_a)=0.
\]
Injectivity gives \(X=0\).  Positivity is needed only to interpret the two
solutions as instruments; uniqueness itself is linear.
\end{proof}

The result is stronger than V60's support-stratum statement.  V60 shows that
all normalized instruments with one fixed exact support have the same
extremality status.  Theorem~\ref{sel61:thm:envelope} says that when the support
normalization map is injective, there cannot be two different normalized
instruments anywhere inside that support envelope at all.

\begin{corollary}[A structural support upper bound is enough for product reuse]
\label{sel61:cor:reuse}
Under the hypotheses of Theorem~\ref{sel61:thm:envelope}, every normalized
instrument supported in \(\mathbf Q\) is extreme.  Therefore V59's fixed-marginal
theorem forces product reuse against every specified second marginal.
\end{corollary}
\begin{proof}
The actual exact supports \(P_a\) satisfy \(P_a\le Q_a\).  Any normalization
perturbation supported in \(P_a\) is also an element of
\(\ker\Pi_{\mathbf Q}\), hence vanishes.  V60 then gives extremality, and V59
gives product reuse.
\end{proof}

This is the useful direction historically: an exact source theorem may provide
a simple forbidden subspace even when it does not diagonalize the nonzero
Choi packet.

\section[Positive zeroes collapse an incomplete completion fibre]{A finite exact witness can replace full Choi tomography}
\label{sel61:zeros}

Suppose the actual historical data do not reconstruct a complete Choi matrix,
but they do determine expectations of positive operators
\(H_{a\ell}\succeq0\).  Define
\[
 Q_a=P_{\bigcap_\ell\ker H_{a\ell}}.
\]

\begin{lemma}[Positive-zero support lemma]
\label{sel61:lem:zero}
For \(J\succeq0\) and \(H\succeq0\),
\begin{equation}
 \boxed{
 \Tr(HJ)=0
 \quad\Longleftrightarrow\quad
 \supp J\subseteq\ker H.}
 \label{sel61:eq:zerolemma}
\end{equation}
\end{lemma}
\begin{proof}
Because both operators are positive,
\[
 \Tr(HJ)=\|H^{1/2}J^{1/2}\|_{\HS}^2.
\]
The trace vanishes exactly when \(H^{1/2}J^{1/2}=0\), which is equivalent to
\(\Ran J\subseteq\ker H\).
\end{proof}

\begin{theorem}[Positive-zero collapse of a contextual completion set]
\label{sel61:thm:collapse}
Let \(\mathcal C\) be any positive trace-preserving completion set for an
incompletely observed instrument.  Suppose the observed or algebraically fixed
data include the exact identities
\begin{equation}
 \boxed{\Tr(H_{a\ell}J_a)=0\quad\text{for all }a,\ell}
 \label{sel61:eq:exactzeros}
\end{equation}
for every \(J\in\mathcal C\), with \(H_{a\ell}\succeq0\).  If
\(\ker\Pi_{\mathbf Q}=0\), then \(\mathcal C\) contains at most one instrument.
If a candidate satisfying the remaining measured data exists, it is the unique
positive completion.
\end{theorem}
\begin{proof}
Lemma~\ref{sel61:lem:zero} places every completion inside the same support
envelope \(\mathbf Q\).  Theorem~\ref{sel61:thm:envelope} then gives uniqueness.
\end{proof}

Thus V24's completion-fibre logic admits a sharp positive shortcut.  One need
not span every Hermitian Choi functional if the available data contain enough
\emph{positive exact zeroes} to place the packet in an injective support face.
The positive operator itself may be reconstructed by signed contextual
postprocessing when it lies in the admitted tester span; what matters for the
support inference is its positivity and the exact zero value, not whether it
was implemented as one primitive detector outcome.

A source grammar can provide the same information algebraically.  For example,
if an exact branch construction gives linear annihilators \(B_{a\ell}\) with
\(B_{a\ell}|K\rangle\!\rangle=0\) on every Kraus vector in branch \(a\), then
\[
 H_{a\ell}=B_{a\ell}^*B_{a\ell}\succeq0
\]
is an exact positive-zero witness.  This is a structural identity, not a
small-number threshold.

\begin{proposition}[Approximate zero gives leakage, not exact support]
\label{sel61:prop:leakage}
Let \(Q\) be a projector and suppose
\[
 H\succeq g(I-Q),\qquad g>0,
 \qquad
 J\succeq0,
 \qquad
 \Tr(HJ)\le\varepsilon.
\]
Then
\begin{equation}
 \boxed{\Tr[(I-Q)J]\le\varepsilon/g.}
 \label{sel61:eq:leakage}
\end{equation}
No conclusion \(\supp J\subseteq\Ran Q\) follows from a positive but nonzero
right-hand side.
\end{proposition}
\begin{proof}
Take the trace of the operator inequality \(H\succeq g(I-Q)\) against
\(J\succeq0\).
\end{proof}

\section[Exact rigidity has zero norm radius]{Why finite-precision tomography alone cannot prove an exact support drop}
\label{sel61:noradius}

The preceding distinction is not merely conservative bookkeeping.

\begin{theorem}[Dense full-support nonextreme contamination]
\label{sel61:thm:dense}
Let \(\mathcal I\) be any finite normalized instrument on input dimension
\(d\), with branch output dimensions \(d_a'\), and assume
\[
 \sum_a(d_a')^2>1.
\]
For every \(\epsilon\in(0,1)\) there exists a normalized instrument
\(\mathcal I_\epsilon\) such that
\begin{equation}
 \boxed{
 \|\mathcal I_\epsilon-\mathcal I\|_\diamond\le2\epsilon,}
 \label{sel61:eq:densebound}
\end{equation}
all branch Choi matrices of \(\mathcal I_\epsilon\) are strictly positive, and
\(\mathcal I_\epsilon\) is nonextreme.  Consequently no nontrivial exact
product-rigidity statement has an open diamond-norm neighbourhood consisting
only of product-rigid instruments.
\end{theorem}
\begin{proof}
Choose probabilities \(p_a>0\), \(\sum_ap_a=1\), and full-rank states
\(\tau_a>0\) on the branch outputs.  The instrument
\[
 \mathcal D_a(X)=p_a\Tr(X)\tau_a
\]
is normalized and each of its Choi matrices is strictly positive on
\(A'_a\otimes A\).  Set
\[
 \mathcal I_\epsilon=(1-\epsilon)\mathcal I+\epsilon\mathcal D.
\]
Every branch Choi matrix is then strictly positive.  Since the diamond distance
between two normalized instruments is at most two, the displayed bound follows.
For the full-support instrument the support-domain dimension is
\[
 \sum_a(d_a'd)^2=d^2\sum_a(d_a')^2>d^2.
\]
V60's dimension bound therefore gives a nonzero normalization kernel.
\end{proof}

The theorem does \emph{not} say that the physical source contains such
contamination.  It says that a finite error bar around a reconstructed extreme
instrument cannot, by itself, exclude it.  Exact support requires exact source
structure.  This is why the historical programme should seek annihilators,
selection rules, exact record identities, or other structural zeroes rather
than numerically thresholding small Choi eigenvalues.

The negative direction is different and remains robust.

\begin{proposition}[One-sided finite-precision nonrigidity certificate]
\label{sel61:prop:robustnegative}
Suppose a reconstructed branch Choi matrix \(\widehat J_a\) obeys a certified
operator-norm error bound
\[
 \|J_a-\widehat J_a\|_{\op}\le\eta_a.
\]
If \(\widehat J_a\) has \(s_a\) eigenvalues strictly larger than \(\eta_a\),
then \(\rank J_a\ge s_a\).  Hence
\begin{equation}
 \boxed{
 \sum_as_a^2>d^2
 \quad\Longrightarrow\quad
 \text{the actual instrument is nonextreme}.}
 \label{sel61:eq:robustnegative}
\end{equation}
\end{proposition}
\begin{proof}
Weyl's eigenvalue inequality implies that each of the first \(s_a\) true
eigenvalues remains positive.  Apply V60's rank lower bound.
\end{proof}

This gives a useful asymmetry: finite precision can robustly exclude rigidity,
but exact positive rigidity requires structural support information.

\section[Small support kernels need only a few local calibrations]{Identification and extremality are different questions}
\label{sel61:calibration}

An envelope can fail to be rigid while still requiring only a few scalar data
to identify the local instrument.  This should not be confused with product
reuse.

Let
\[
 \mathfrak N_{\mathbf Q}^{\rm sa}
 =\{X=X^*\in\mathfrak S_{\mathbf Q}:\Pi_{\mathbf Q}X=0\},
 \qquad
 \nu=\dim_{\mathbb R}\mathfrak N_{\mathbf Q}^{\rm sa}.
\]
Because the complex kernel is closed under adjoint, this real dimension equals
the complex nullity used in V59--V60.

\begin{theorem}[Minimal calibration theorem inside a support envelope]
\label{sel61:thm:calibration}
Let \(f_1,\ldots,f_m\) be real linear functionals on Hermitian branch packets,
and let \(F=(f_1,\ldots,f_m)\).  Among normalized Hermitian packets supported in
\(\mathbf Q\), the normalization and calibration values uniquely determine a
candidate exactly when
\begin{equation}
 \boxed{
 \mathfrak N_{\mathbf Q}^{\rm sa}\cap\ker F=\{0\}.}
 \label{sel61:eq:calibrationcriterion}
\end{equation}
If a normalized candidate is positive definite on every nonzero envelope
\(Q_a\), then at least \(\nu\) independent real scalar calibrations are locally
necessary, and \(\nu\) generic independent calibrations are sufficient.
\end{theorem}
\begin{proof}
Two normalized supported packets with the same calibration data differ by an
element of the displayed intersection, proving the first statement.  At a
relative-interior positive candidate, every sufficiently small Hermitian
normalization-kernel perturbation remains positive in both signs.  Therefore
positivity does not reduce the local affine dimension \(\nu\); an injective real
linear map on a \(\nu\)-dimensional space needs at least \(\nu\) scalar
coordinates.  Generic \(\nu\) functionals separate that finite-dimensional
space.
\end{proof}

The theorem concerns \emph{local source identification}.  If \(\nu>0\), adding
calibration values does not make the identified instrument extreme: a joint
source can still exploit its normalization kernel while preserving the now
fully known local marginal.

\section[Exact evaluations on the unchanged interfaces]{What the stronger entrance buys on the existing examples}
\label{sel61:evaluations}

\subsection{V35's forgotten nonunital qubit needs no Choi eigenvalue calibration}

Let \(Q_{35}\) be the two-dimensional Choi support already reconstructed for the
unchanged V35 channel.  V60 proves
\[
 \ker\Pi_{Q_{35}}=0.
\]
Therefore Theorem~\ref{sel61:thm:envelope} strengthens the earlier conclusion:
\begin{equation}
 \boxed{
 \text{trace preservation + the upper support bound }Q_{35}
 \text{ uniquely determine the entire channel}.}
 \label{sel61:eq:v35unique}
\end{equation}
The nonzero Choi eigenvalues need not be separately supplied.  If a historical
source identity provides a positive annihilator whose kernel is \(Q_{35}\), the
exact channel follows from that support identity and normalization alone.

This remains a statement about the unchanged finite qubit source.  Its
information loss, stationary reverse, and recovery obstructions are not altered.

\subsection{The recorded sign source and the six accepted compiler branches}

For V57's two recorded sign branches, take \(Q_+\) and \(Q_-\) to be their
one-dimensional Choi rays.  Their effects are independent, hence
\(\ker\Pi_{\mathbf Q}=0\).  Once those two exact branch rays are fixed, trace
preservation uniquely fixes their weights.  The same argument applies to the
six original accepted compiler rays on the complete native carrier: V57's exact
effect independence makes their support envelope injective.  Thus
\begin{equation}
 \boxed{
 \text{accepted branch rays + normalization}
 \Longrightarrow
 \text{the full accepted instrument}.}
 \label{sel61:eq:sixunique}
\end{equation}
This is a source-identification statement, not a claim that the historical Store
has already supplied those six rays.

\subsection{The raw idle source: one calibration fixes the local instrument but not reuse}

Add the actual idle ray.  V60 gives
\[
 \nu_{\rm raw}=1.
\]
Write the supported normalized family using the fixed accepted rays as
\[
 J_0=t\,J_{\rm idle},
 \qquad
 J_e=(1-t)J_e^{\rm acc},
 \qquad 0<t<1,
\]
with the obvious normalization convention for the ray representatives.  The
single acceptance scalar \(1-t\) therefore fixes the local raw instrument
inside its support envelope.  This is the \(\nu=1\) case of
Theorem~\ref{sel61:thm:calibration}.

But \(\nu_{\rm raw}=1\) still means that the now fully identified local
instrument is nonextreme.  V57's joint acceptance-correlation freedom remains.
Knowing the local acceptance probability does not force independent acceptance
across opportunities.

\subsection{The forgotten sign channel and the two-event endpoint}

The forgotten sign channel has support nullity two.  A relative-interior channel
inside that support envelope therefore needs two independent real local
calibrations in addition to trace preservation for unique identification.
That does not remove its four-dimensional compatible pair family.

The two-event qubit endpoint is more severe: its Choi support is full and V60
gives \(\nu=12\).  Support information contributes no positive rigidity at all;
twelve real local channel coordinates remain after trace preservation.  This is
exactly the regime in which a structural support shortcut cannot replace the
full endpoint process data.

\section[Revised historical compiler]{Seek exact forbidden directions before full tomography}
\label{sel61:status}

The source-selection entrance can now be organized asymmetrically.

For a candidate historical branch family, first seek exact positive forbidden
directions:
\[
 H_{a\ell}\succeq0,
 \qquad
 \Tr(H_{a\ell}J_a)=0.
\]
They define a support envelope \(\mathbf Q\).  Then compute the finite support
normalization map \(\Pi_{\mathbf Q}\).

If
\[
 \boxed{\ker\Pi_{\mathbf Q}=0,}
\]
normalization and those exact zeroes identify the entire instrument; full Choi
tomography is unnecessary.  If the kernel is small but nonzero, its dimension
tells exactly how many independent local scalar calibrations are required to
identify a relative-interior source within that envelope.  If certified support
\emph{lower} bounds already satisfy V60's dimension inequality, nonrigidity is
proved without resolving the remaining spectrum.

Approximate zeroes occupy a different line:
\[
 \Tr(HJ)\le\varepsilon
 \quad\Longrightarrow\quad
 \text{controlled leakage only}.
\]
Theorem~\ref{sel61:thm:dense} forbids upgrading such a statement into exact
extremality from norm closeness alone.  Any exact positive closure must retain
its structural-zero premise explicitly.

\paragraph{Proved.}
Injective support envelopes are singleton normalized faces; positive exact-zero
witnesses can collapse an incomplete contextual completion set to that
singleton; exact extremality has zero diamond-norm robustness radius against
full-support contamination; rank lower bounds give a robust one-sided
nonrigidity test; and the normalization-kernel dimension is exactly the number
of generic real local calibrations required inside a relative-interior support
envelope.

\paragraph{Inherited inputs.}
V24 supplies the distinction between contextual inference and physical
execution.  V35 supplies the unchanged nonunital qubit.  V57 supplies the exact
accepted compiler effect independence and the raw idle interface.  V59--V60
supply the normalization-kernel and support-only reuse theorems.

\paragraph{Remaining source question.}
The historical primitive typed branches still have to supply either their
actual Choi packets or exact structural annihilators on the correct same-cut
carrier.  A fitted near-zero eigenvalue is not such an annihilator.  The raw
idle relation and clock records remain part of the instrument, and local
calibration does not determine their joint reuse law.

Historical Standard-Model selection is not declared closed.  The V110 directed
determinant--Clifford word and its typed incidence remain unpopulated.  The gain
in V61 is narrower but upstream: a historical proof need not reconstruct every
positive Choi weight if it can prove the right exact forbidden directions.

\paragraph{Verification boundary.}
The companion checker verifies the singleton-envelope linear systems on the
unchanged qubit and compiler interfaces, the exact positive-annihilator
identities, the raw one-parameter family and its acceptance calibration, and
full-support nonextreme contaminations.  The general finite-dimensional
statements have the written proofs above.  No empirical exact zero,
proof-assistant formalization, or independent peer review is claimed.

\addtocontents{toc}{\protect\endgroup}
% END SOURCE-SELECTION V61 DEVELOPMENT BLOCK
% Append future development before the final document command.


% BEGIN SOURCE-SELECTION V62 DEVELOPMENT BLOCK
\clearpage
\hypersetup{pdftitle={Historical Standard-Model Source Selection: V1--V62}}
\begin{center}
{\Large\bfseries Source-selection continuation V62}\\[.4em]
{\large Fourth-moment rigidity from source-native amplitudes,\\
without Choi reconstruction or support projectors}
\end{center}
\addtocontents{toc}{\protect\begingroup}
\addtocontents{toc}{\protect\renewcommand{\protect\numberline}{\protect\makebox[2.1em][l]}}

\section[Move upstream from support envelopes]{When the source already supplies amplitude words, the fourth moment is the direct object}
\label{sel62:context}

V60--V61 moved the fixed-marginal reuse test from a chosen Kraus frame to exact
Choi supports and then further to exact support envelopes.  That is the correct
general route for a genuinely mixed primitive branch.  It is nevertheless
stronger than necessary on a different historical branch that has occurred
throughout the programme: the source grammar may already supply a finite
\emph{efficient amplitude refinement}.  In that case one should not first
reconstruct the complete Choi packet, diagonalize it, reconstruct its support,
and only then test normalization rigidity.

The present continuation identifies the direct source-native object.  For each
retained record $a$, assume that the historical grammar supplies actual
amplitude operations
\[
 L_{a1},\ldots,L_{am_a}:A\longrightarrow A'_a
\]
such that
\begin{equation}
 \Phi_a(X)=\sum_{\lambda=1}^{m_a}
 L_{a\lambda}X L_{a\lambda}^* .
 \label{sel62:eq:refinement}
\end{equation}
The amplitudes in \eqref{sel62:eq:refinement} need not be separately retained
classical labels, but they must be source-native operations on the same cut;
an arbitrary mathematical Kraus decomposition invented after the fact is not
promoted to historical words below.

The main result is that the normalization kernel of V59--V60 is already the
kernel of one positive \emph{fourth-word-moment Gram}.  When the amplitudes are
historical words, V31's inserted-moment compiler and finite flatness compute
this matrix from the old Grand Tensor.  Neither a Choi matrix nor a support
projector has to be displayed.

This gives a second, complementary entrance to V61.  The support-envelope route
remains the correct fallback when primitive mixed branches have no historical
amplitude refinement.  The two routes are not conflated.

\section[The fourth-moment product Gram is the normalization kernel]{A direct finite extremality and reuse certificate}
\label{sel62:fourth}

Begin with a support-minimal Kraus frame for each retained branch,
\[
 \Phi_a(X)=\sum_{\mu=1}^{r_a}K_{a\mu}XK_{a\mu}^*,
 \qquad
 r_a=\rank J_a .
\]
Define the input-side product bank
\begin{equation}
 P_{a\mu\nu}:=K_{a\mu}^*K_{a\nu}
 \label{sel62:eq:productbank}
\end{equation}
and its Hilbert--Schmidt Gram matrix
\begin{equation}
 \boxed{
 \mathbb G^{(4)}_{a\mu\nu,\,b\rho\sigma}
 :=\Tr\!\left(P_{a\mu\nu}^*P_{b\rho\sigma}\right).}
 \label{sel62:eq:fourthgram}
\end{equation}
The superscript records that a matrix entry contains four source amplitudes.

\begin{theorem}[Fourth-moment normalization theorem]
\label{sel62:thm:fourth}
For a normalized finite-dimensional instrument,
\begin{equation}
 \boxed{
 \dim\ker\mathbb G^{(4)}
 =\nu_{\mathcal I},}
 \label{sel62:eq:nullity}
\end{equation}
where $\nu_{\mathcal I}$ is the normalization-kernel dimension of V59--V60.
Consequently the following are equivalent:
\begin{enumerate}[label=\textup{(\roman*)},leftmargin=3em]
\item $\mathbb G^{(4)}$ is positive definite;
\item the operators $K_{a\mu}^*K_{a\nu}$ are linearly independent;
\item the instrument is extreme;
\item every joint instrument having this instrument as one all-input marginal
      and any specified second marginal is the product implementation.
\end{enumerate}
If $\mathbb G^{(4)}$ is singular, every vector in its kernel gives explicitly a
normalization-kernel direction.
\end{theorem}

\begin{proof}
Let
\[
 \Gamma_{\mathcal I}:\bigoplus_aM_{r_a}(\mathbb C)
 \longrightarrow\mathcal B(A),
 \qquad
 \Gamma_{\mathcal I}(T)
 =\sum_{a,\mu,\nu}T^{(a)}_{\mu\nu}
   K_{a\mu}^*K_{a\nu}.
\]
Use the matrix-unit basis on the domain and the Hilbert--Schmidt product on
$\mathcal B(A)$.  Equation~\eqref{sel62:eq:fourthgram} is exactly the matrix of
$\Gamma_{\mathcal I}^*\Gamma_{\mathcal I}$.  Hence
\[
 \ker\mathbb G^{(4)}=\ker\Gamma_{\mathcal I}.
\]
V59 identifies this kernel with the normalization kernel, proving
\eqref{sel62:eq:nullity}.  Positive definiteness is equivalent to injectivity.
The equivalence with extremality is the finite instrument extremality criterion
already audited in V59, and V59's fixed-marginal theorem gives the product
statement.  If $c\in\ker\mathbb G^{(4)}$, then the coefficient matrices obtained
by reshaping $c_{a\mu\nu}$ satisfy
$\Gamma_{\mathcal I}(c)=0$ and therefore are the promised normalization-kernel
direction.
\end{proof}

The theorem is basis free in the physically relevant sense.  A different
minimal Kraus frame in branch $a$ is related by a unitary $U_a$.  The product
bank transforms by $\overline U_a\otimes U_a$, so the whole fourth Gram changes
by unitary similarity.  Its spectrum, rank, nullity, and positive-definiteness
status are unchanged.

\begin{corollary}[Historical fourth moments can bypass Choi tomography]
\label{sel62:cor:historical}
Assume every minimal branch amplitude $K_{a\mu}$ is a fixed linear combination
of admitted historical words of depth at most $r$ on the same cut.  Then every
entry of $\mathbb G^{(4)}$ is a finite linear combination of ordinary
Grand-Tensor word moments.  A symmetric historical word Gram through depth
$2r$ is sufficient.  After finite flatness, the reconstructed multiplication
table is sufficient and no deeper physical histories are required.
\end{corollary}

\begin{proof}
Each product $K_{a\mu}^*K_{a\nu}$ lies in the span of represented words of
depth at most $2r$.  Equation~\eqref{sel62:eq:fourthgram} is the ordinary
Hilbert--Schmidt Gram of those product words.  Thus it is an inserted moment of
the kind compiled in V31.  A word Gram through depth $2r$ contains all pairings
of such product words.  At flatness all products reduce in the existing finite
multiplication table.
\end{proof}

This is an inference statement.  Taking linear combinations of historical
amplitudes in order to choose a convenient minimal frame is a calculation on
the reconstructed source span; it is not the declaration of a new executable
source outcome.

\section[Nonminimal source refinements are quotiented first]{Kraus dependencies must not be mistaken for normalization freedom}
\label{sel62:nonminimal}

A historical source may present more fine amplitudes than the Choi rank of the
retained coarse branch.  Directly forming all products from such a redundant
list creates null vectors that merely express Kraus-frame redundancy.

For one branch define the amplitude Gram
\begin{equation}
 A^{(a)}_{\lambda\rho}
 :=\Tr(L_{a\lambda}^*L_{a\rho}).
 \label{sel62:eq:ampgram}
\end{equation}
Its rank is the dimension of the operator span of the supplied amplitudes.  A
unitary rotation on the finite coefficient space can be chosen so that the
last $m_a-r_a$ rotated amplitudes vanish identically; deleting those zeros gives
a support-minimal Kraus frame for the same CP map.  This is ordinary Kraus
compression and leaves \eqref{sel62:eq:refinement} unchanged.

\begin{proposition}[Two-stage Gram compiler]
\label{sel62:prop:twostage}
Suppose the source supplies a finite amplitude refinement
\eqref{sel62:eq:refinement}.  Then the normalization nullity is computable from
historical moments by the following finite procedure:
\begin{enumerate}[label=\textup{(\arabic*)},leftmargin=3em]
\item compute each amplitude Gram $A^{(a)}$ and quotient its exact kernel;
\item choose any orthonormal basis of the resulting coefficient quotient and
      form the corresponding support-minimal amplitudes $K_{a\mu}$;
\item build $\mathbb G^{(4)}$ from \eqref{sel62:eq:fourthgram};
\item return $\dim\ker\mathbb G^{(4)}$.
\end{enumerate}
The returned number is independent of every basis choice in steps (1)--(2).
Exact zeroes in $A^{(a)}$ and in $\mathbb G^{(4)}$ are algebraic rank
statements.  Numerical near-nullity is not silently converted into an exact
quotient.
\end{proposition}

\begin{proof}
The amplitude Gram is the Gram of the operator vectors
$L_{a\lambda}$, so its kernel is exactly the coefficient kernel of their
synthesis map.  A unitary coefficient rotation separates that kernel and
preserves the CP sum.  Any two resulting minimal frames are related by a
branch unitary.  The preceding Kraus-frame covariance then proves that the
rank and nullity of the fourth Gram are independent of the chosen quotient
basis.  Theorem~\ref{sel62:thm:fourth} identifies the returned nullity with the
intrinsic normalization kernel.
\end{proof}

This two-stage compiler is especially useful when a coarse historical record
is known to arise by erasing a finer efficient record.  The erased fine labels
need not remain experimentally readable, but their amplitude refinement must
be independently established by the source grammar.  If no such historical
refinement exists, an arbitrary Kraus decomposition of the coarse Choi matrix
is only a mathematical representation; V60--V61 remain the correct route.

\section[Second moments are not enough]{Two channels with identical amplitude Gram can have opposite reuse answers}
\label{sel62:second}

The fourth moment cannot in general be replaced by the ordinary amplitude
Gram.  There is an exact separating pair already inside the examples used in
this programme.

Let the unchanged V35 nonunital channel have
\[
 k_+=\frac1{\sqrt{10}}
 \begin{pmatrix}2&2\\-1&1\end{pmatrix},
 \qquad
 k_-=\frac1{\sqrt{10}}
 \begin{pmatrix}2&-2\\1&1\end{pmatrix}.
\]
Compare it with the qubit dephasing channel
\[
 d_0=\frac{I}{\sqrt2},
 \qquad
 d_1=\frac{Z}{\sqrt2}.
\]
Both amplitude families have the identical second Gram
\begin{equation}
 \boxed{
 [\Tr(k_i^*k_j)]_{ij}
 =[\Tr(d_i^*d_j)]_{ij}
 =I_2.}
 \label{sel62:eq:same2gram}
\end{equation}

\begin{proposition}[Identical second moments, opposite fourth-moment rigidity]
\label{sel62:prop:separate}
In the product ordering $(00,01,10,11)$, the V35 fourth Gram is
\begin{equation}
 \boxed{
 \mathbb G^{(4)}_{\rm V35}
 =\frac1{25}
 \begin{pmatrix}
 17&0&0&8\\
 0&17&-8&0\\
 0&-8&17&0\\
 8&0&0&17
 \end{pmatrix},}
 \label{sel62:eq:v35gram}
\end{equation}
with eigenvalues $1,1,9/25,9/25$ and determinant $81/625$.
Thus the channel is product rigid.

For the dephasing channel,
\begin{equation}
 \boxed{
 \mathbb G^{(4)}_{\rm deph}
 =\frac12
 \begin{pmatrix}
 1&0&0&1\\
 0&1&1&0\\
 0&1&1&0\\
 1&0&0&1
 \end{pmatrix},}
 \label{sel62:eq:dephgram}
\end{equation}
whose spectrum is $\{1,1,0,0\}$.  Its normalization nullity is two.
Therefore no criterion depending only on the pairwise amplitude Gram can decide
fixed-marginal reuse rigidity.
\end{proposition}

\begin{proof}
Equation~\eqref{sel62:eq:same2gram} is direct.  For V35 the four products are
those already evaluated in V59; their Hilbert--Schmidt Gram is
\eqref{sel62:eq:v35gram}.  For dephasing the products are
$I/2,Z/2,Z/2,I/2$, giving \eqref{sel62:eq:dephgram}.  The displayed spectra
then follow exactly.
\end{proof}

Operationally, the ordinary source Gram knows that the two amplitude vectors
are orthonormal in both examples.  The fourth Gram additionally knows how their
input-side products multiply.  That multiplication information is precisely
what normalization rigidity needs.

\section[Exact evaluations on the existing source interfaces]{The direct fourth-moment route reproduces the prior census}
\label{sel62:examples}

The new compiler changes the route, not the already established answers.

\paragraph{Unchanged V35 forgotten channel.}
Its two source-native amplitude words satisfy
\eqref{sel62:eq:v35gram}; hence the old Grand-Tensor fourth moments already
certify $\nu=0$.  On this source one may bypass both explicit Choi support
reconstruction and V61's positive support annihilator, provided the two
amplitudes themselves have been historically landed.

\paragraph{Retained three-level sign instrument.}
Each retained sign branch is efficient, so the product bank consists only of
its two effects $E_+,E_-$.  V57 computed
\[
 \boxed{
 \det[\Tr(E_sE_t)]_{s,t=\pm}=\frac49>0.}
\]
This is exactly the two-dimensional fourth-moment criterion of
Theorem~\ref{sel62:thm:fourth}; it proves $\nu_{\rm sign}=0$ directly from
source products.

\paragraph{Forgotten sign channel.}
After erasing the sign, the two amplitudes become one channel Kraus frame.  In
the spectral decomposition $U=P_0-iP_1$, all four products
$L_s^*L_t$ lie in the two-dimensional span of $P_0,P_1$.  Hence the fourth
Gram has rank two and nullity two:
\[
 \boxed{\nu_{\rm forgotten}=4-2=2,}
\]
in agreement with V59--V61.

\paragraph{Six accepted compiler branches.}
The six branches are efficient.  Their product bank is exactly the six effect
operators tested by V57's invertible six-functional evaluation matrix.  Thus
the same historical inserted moments certify
\[
 \boxed{\nu_{\rm six,accepted}=0.}
\]
No support projector on the full native Choi carrier is required for this
positive branch.

\paragraph{Raw idle plus accepted branches.}
The idle product is proportional to $I$, while the six accepted effects sum to
$I$.  There is exactly one product-bank relation, so
\[
 \boxed{\nu_{\rm six,raw}=1.}
\]
The fourth-moment route therefore recovers the single acceptance-correlation
direction without first constructing the raw Choi supports.

These evaluations also clarify the hierarchy of historical inputs.  An exact
amplitude refinement is stronger than a coarse branch label but weaker than a
complete Choi tomography.  When that refinement is genuinely source native,
the fourth moment is the shortest rigidity certificate.  When it is not, the
support-envelope route remains indispensable.

\section[Revised historical compiler]{Ask first whether the amplitude refinement is historical}
\label{sel62:status}

The source-reuse entrance now branches before any Choi reconstruction:
\[
 \boxed{
 \begin{array}{c}
 \text{historical efficient amplitude refinement available}\\[.15em]
 \Downarrow\\[-.1em]
 \text{amplitude Gram quotient}\\[-.05em]
 \Downarrow\\[-.1em]
 \mathbb G^{(4)}\text{ from old word moments}
 \end{array}}
\]
If no historical amplitude refinement is available, one instead uses
\[
 \boxed{
 \text{V60--V61 Choi-support / support-envelope route}.}
\]
On the first branch, finite flatness and V31's inserted-moment compiler mean
that no independent process-tomography panel is needed solely for the reuse
question.  The needed information is already a finite fourth moment of the
historical word algebra.  A nonzero kernel gives the complete local
normalization ambiguity used by V59's joint-coupling theorem; a positive
definite matrix forces product reuse.

\paragraph{Proved.}
The normalization nullity equals the nullity of the minimal Kraus-product
fourth Gram.  Its rank and spectrum are invariant under every minimal Kraus
basis change.  A redundant historical amplitude refinement is reduced by its
ordinary amplitude Gram before the fourth-moment test.  Historical words of
depth $r$ require only the existing word Gram through depth $2r$, and no deeper
histories after flatness.  Pairwise amplitude moments alone are insufficient,
as shown by an exact V35/dephasing pair with identical second Gram and opposite
fourth-moment rigidity.  The V35, sign, forgotten-sign, six-accepted, and raw
idle interfaces reproduce their previous nullities directly.

\paragraph{Inherited inputs.}
V31 supplies the inserted-word-moment compiler and the strict distinction
between word-native and external shadows.  V57 supplies the compiler effect
independence.  V59 supplies the mixed-outcome normalization-kernel theorem.
V60--V61 remain the general support-based route when source-native amplitude
words are unavailable.

\paragraph{Sharp boundary.}
A generic Kraus decomposition of an observed mixed Choi packet is not a
historical amplitude refinement.  Its Kraus vectors may be used internally to
prove the basis-free V59--V60 theorem, but they cannot be inserted into the old
word Gram unless their source ancestry is independently established.  Likewise,
exact zeros in the amplitude and fourth Grams remain structural statements;
finite numerical near-nullity is not promoted into an exact quotient.

\paragraph{Remaining source question.}
The unresolved historical row moves one step upstream: determine whether the
primitive typed Store events admit the required \emph{source-native efficient
amplitude refinement} on the same geometry-shorted cut, or instead only a mixed
coarse CP packet.  In the former case V62 compiles reuse rigidity directly from
old fourth moments.  In the latter case V61's exact annihilator/support-envelope
programme remains necessary.

Historical Standard-Model selection is not declared closed.  In particular,
none of these fourth-moment identities populate the V110 directed
Store--Clifford/up-incidence word or manufacture its typing operations.  They
only shorten the source-reuse audit once the relevant amplitudes have actually
landed in the historical word algebra.

\paragraph{Verification boundary.}
The companion checker reconstructs the fourth Grams from exact matrices,
verifies Kraus-frame covariance, checks the redundant-frame quotient, proves
the identical-second-Gram separating example, and reproduces the unchanged
nullity census.  The general word-depth and historical-ancestry statements have
the written proofs above.  No empirical exact rank, proof-assistant
formalization, physical source refinement, or independent peer review is
claimed.

\addtocontents{toc}{\protect\endgroup}
% END SOURCE-SELECTION V62 DEVELOPMENT BLOCK
% Append future development before the final document command.


% BEGIN SOURCE-SELECTION V63 DEVELOPMENT BLOCK
\clearpage
\hypersetup{pdftitle={Historical Standard-Model Source Selection: V1--V63}}
\begin{center}
{\Large\bfseries Source-selection continuation V63}\\[.4em]
{\large Envelope fourth moments from historical selection rules,\\
without efficient refinements or exact Choi supports}
\end{center}
\addtocontents{toc}{\protect\begingroup}
\addtocontents{toc}{\protect\renewcommand{\protect\numberline}{\protect\makebox[2.1em][l]}}

\section[Move upstream from amplitude refinements]{A support envelope, not a historical Kraus refinement, is enough on the positive branch}
\label{sel63:context}

V62 deliberately imposed a strong ancestry firewall: its fourth-moment
compiler was invoked only after a source-native efficient amplitude refinement
had been established.  That firewall is necessary if one wants to identify the
\emph{actual normalization nullity} from a displayed Kraus family.  It is
stronger than necessary for the one-sided positive conclusion needed for
historical source selection.

Indeed V61 already showed that an exact Choi \emph{support envelope} can force a
unique normalized source even when the actual support is a strict subspace of
that envelope.  The present continuation combines that observation with V62's
fourth-moment compiler.  The result is a direct test on an operator-space
\emph{selection rule}.  The operators spanning the envelope need not be actual
Kraus amplitudes and need not be executable source outcomes.

For each retained branch $a$, let
\[
 \mathsf S_a\subset\mathcal B(A,A'_a)
\]
be a finite-dimensional operator space known exactly to contain every possible
Kraus operator of that branch.  Equivalently, if $Q_a$ denotes the orthogonal
projector onto $\operatorname{vec}(\mathsf S_a)$, the historical statement is
only
\begin{equation}
 \operatorname{supp}J_a\subseteq Q_a.
 \label{sel63:eq:envelope}
\end{equation}
Such an envelope can come from exact typing, grading, source/target, incidence,
or word-algebra selection rules.  No decomposition of $J_a$ is assumed.

The main theorem below shows that V61's normalization map on the Choi envelope
is exactly the product-bank fourth moment of \emph{any basis of
$\mathsf S_a$}.  Hence a positive-definite envelope fourth Gram proves the
whole local instrument and product reuse without first proving an efficient
historical refinement.

\section[Envelope fourth moments equal the support-envelope kernel]{The V61 and V62 entrances are the same finite linear map}
\label{sel63:envelope-fourth}

Choose a linearly independent basis
\[
 B_{a1},\ldots,B_{am_a}
\]
of $\mathsf S_a$ and define the product bank
\begin{equation}
 C_{a\mu\nu}:=B_{a\mu}^*B_{a\nu}.
 \label{sel63:eq:products}
\end{equation}
Its Hilbert--Schmidt Gram is
\begin{equation}
 \boxed{
 \mathbb E^{(4)}_{a\mu\nu,\,b\rho\sigma}
 :=\Tr\!\left(C_{a\mu\nu}^*C_{b\rho\sigma}\right).}
 \label{sel63:eq:envelopegram}
\end{equation}
The superscript again records four operator amplitudes, but the $B_{a\mu}$ are
now merely coordinates on an exact support envelope.

\begin{theorem}[Envelope fourth-moment theorem]
\label{sel63:thm:envelope}
Let $Q_a=P_{\operatorname{vec}(\mathsf S_a)}$ and let
$\Pi_{\mathbf Q}$ be the V61 support-envelope normalization map.  Then
\begin{equation}
 \boxed{
 \dim\ker\mathbb E^{(4)}
 =\dim\ker\Pi_{\mathbf Q}.}
 \label{sel63:eq:kernelidentity}
\end{equation}
In particular, if $\mathbb E^{(4)}$ is positive definite, then:
\begin{enumerate}[label=\textup{(\roman*)},leftmargin=3em]
\item there is at most one normalized instrument satisfying
      \eqref{sel63:eq:envelope};
\item if such an instrument exists, its actual branch supports may be strict
      subspaces of the $Q_a$ and the instrument is nevertheless extreme;
\item every joint implementation having that instrument as one all-input
      marginal and any specified second marginal is the product
      implementation.
\end{enumerate}
No efficient refinement of the historical branch is required.
\end{theorem}

\begin{proof}
Let $W_a:\mathbb C^{m_a}\to A'_a\otimes A$ be the injective synthesis map
$W_ae_\mu=|B_{a\mu}\rangle\!\rangle$.  Every Hermitian operator supported in
$Q_a$ has a unique form
\[
 X_a=W_aT_aW_a^*,\qquad T_a=T_a^*.
\]
Moreover $X_a\succeq0$ if and only if $T_a\succeq0$, because $W_a^*$ is
surjective onto the coefficient space.  With the standard vectorization
convention,
\[
 \Tr_{A'_a}X_a
 =\left(\sum_{\mu,\nu}(T_a)_{\mu\nu}
 B_{a\nu}^*B_{a\mu}\right)^{\!\mathsf T}.
\]
Coefficient transposition and the final matrix transpose are invertible linear
operations, so the kernel of $\Pi_{\mathbf Q}$ is isomorphic to the kernel of
\[
 \Gamma_{\mathbf S}((T_a)_a)
 :=\sum_{a,\mu,\nu}(T_a)_{\mu\nu}B_{a\mu}^*B_{a\nu}.
\]
Equation~\eqref{sel63:eq:envelopegram} is the matrix of
$\Gamma_{\mathbf S}^*\Gamma_{\mathbf S}$ in matrix-unit coordinates.  Hence
$\ker\mathbb E^{(4)}=\ker\Gamma_{\mathbf S}$ and
\eqref{sel63:eq:kernelidentity} follows.  V61 then gives uniqueness of a
normalized instrument inside an injective envelope.  Restricting an injective
normalization map to the smaller actual supports remains injective, so the
actual instrument is extreme.  V59's fixed-marginal theorem gives the product
statement.
\end{proof}

A change of envelope basis by an invertible matrix acts on product coordinates
by an invertible congruence.  Thus the rank, nullity, and nonsingularity of
$\mathbb E^{(4)}$ are basis independent.  If the envelope basis is chosen
Hilbert--Schmidt orthonormally, any further orthonormal basis change is unitary
and the spectrum itself is invariant, exactly as in V62.

\begin{corollary}[Historical word-span envelope compiler]
\label{sel63:cor:word}
Suppose the source grammar proves only that
\[
 \mathsf S_a\subseteq
 \operatorname{span}\{w_{a1},\ldots,w_{am_a}\},
\]
where the $w_{a\mu}$ are admitted historical words on the same cut, of depth
at most $r$.  Quotient exact linear dependencies in this word list using its
ordinary Gram, choose any basis of the resulting envelope, and form
\eqref{sel63:eq:envelopegram}.  The matrix is compiled from the old word
moments through depth $2r$, or from the finite multiplication table after
flatness.  If it is positive definite, the historical source and its product
reuse are fixed even though none of the basis words has been asserted to be an
actual Kraus amplitude.
\end{corollary}

\begin{proof}
The only physical premise is the support inclusion
\eqref{sel63:eq:envelope}.  Once a basis of the exact word span has been
obtained, every $B_{a\mu}^*B_{a\nu}$ is a linear combination of represented
words of depth at most $2r$.  Their Hilbert--Schmidt Gram is therefore an old
inserted word moment, exactly as in V62.  Theorem~\ref{sel63:thm:envelope}
then applies without interpreting the basis vectors as primitive source
outcomes.
\end{proof}

This removes V62's efficient-refinement premise only for the \emph{positive
rigidity} route.  If the envelope Gram is singular, its kernel means only that
the current structural envelope is too broad.  It does not prove that the
actual, potentially smaller, Choi supports are nonextreme.

\section[A strict envelope can already be a singleton]{Exact support equality is genuinely unnecessary}
\label{sel63:strict}

The preceding distinction is not merely formal.  There are normalized sources
for which the exact structural envelope is strictly larger than the actual
support and nevertheless contains only one normalized instrument.

Let $A=A'=\mathbb C^3$ and write $E_{ij}=|i\rangle\langle j|$.  Consider two
branch envelopes
\[
 \mathsf S_1=\operatorname{span}\{E_{11},E_{12}\},
 \qquad
 \mathsf S_2=\operatorname{span}\{P_{23}\},
 \qquad
 P_{23}=E_{22}+E_{33}.
\]
A normalized candidate inside them is
\begin{equation}
 J_1=|E_{11}\rangle\!\rangle\langle\!\langle E_{11}|,
 \qquad
 J_2=|P_{23}\rangle\!\rangle\langle\!\langle P_{23}|,
 \label{sel63:eq:strictcandidate}
\end{equation}
because $E_{11}+P_{23}=I_3$.  The actual first-branch support is only the ray
$\operatorname{span}\{E_{11}\}$, a strict subspace of $\mathsf S_1$.

\begin{proposition}[Strict-envelope singleton]
\label{sel63:prop:strict}
For the envelopes above, the product bank is
\[
 E_{11},\ E_{12},\ E_{21},\ E_{22},\ E_{22}+E_{33}.
\]
Its fourth Gram is
\begin{equation}
 \boxed{
 \mathbb E^{(4)}=
 \begin{pmatrix}
 1&0&0&0&0\\
 0&1&0&0&0\\
 0&0&1&0&0\\
 0&0&0&1&1\\
 0&0&0&1&2
 \end{pmatrix},
 \qquad
 \det\mathbb E^{(4)}=1.}
 \label{sel63:eq:strictgram}
\end{equation}
Consequently \eqref{sel63:eq:strictcandidate} is the unique normalized
instrument in the stated envelopes, despite the strict support inclusion in
branch one.
\end{proposition}

\begin{proof}
The five displayed products are linearly independent, and their
Hilbert--Schmidt products give \eqref{sel63:eq:strictgram}.  Theorem
\ref{sel63:thm:envelope} gives uniqueness.  Directly, writing the branch-one
coefficient matrix as $T=(t_{ij})$ and the branch-two weight as $s$, trace
preservation reads
\[
 t_{11}E_{11}+t_{12}E_{12}+t_{21}E_{21}
 +t_{22}E_{22}+s(E_{22}+E_{33})=I_3.
\]
Hence $s=1$, $t_{22}=0$, $t_{11}=1$, and $t_{12}=t_{21}=0$, which is exactly
\eqref{sel63:eq:strictcandidate}.
\end{proof}

Thus one need not prove that every allowed structural direction is occupied.
On an injective envelope, normalization itself can force the unused directions
to have zero weight.

\section[Dimension screens and structural narrowing]{A broad word algebra is usually too large; typing must do real work}
\label{sel63:dimension}

The envelope theorem gives an immediate necessary screen.  Put
$m_a=\dim\mathsf S_a$ and $d=\dim A$.  Since
\[
 \Gamma_{\mathbf S}:\bigoplus_aM_{m_a}(\mathbb C)
 \longrightarrow M_d(\mathbb C),
\]
injectivity requires
\begin{equation}
 \boxed{
 \sum_a m_a^2\le d^2.}
 \label{sel63:eq:dimension}
\end{equation}
If this fails, the current selection rules cannot by themselves prove a unique
normalized source.

\begin{proposition}[Monotonicity under exact structural refinement]
\label{sel63:prop:monotone}
Suppose $\mathsf S'_a\subseteq\mathsf S_a$ for every branch.  Then the
normalization map for $\mathbf S'$ is the restriction of the map for
$\mathbf S$.  In particular:
\begin{enumerate}[label=\textup{(\roman*)},leftmargin=3em]
\item if the larger envelope is injective, every smaller exact envelope is
      injective;
\item if the larger envelope is singular, no conclusion about the actual
      source follows until the envelope is narrowed or the actual supports are
      identified;
\item successive exact typing or annihilator constraints can only remove
      normalization-kernel directions, never create a failure of injectivity
      that was absent before.
\end{enumerate}
\end{proposition}

\begin{proof}
Choose bases of $\mathsf S'_a$ and extend them to bases of $\mathsf S_a$.
The coefficient domain for the smaller envelope is then a linear subspace of
the larger coefficient domain and its normalization synthesis map is the
restriction of the larger map.  The three statements are immediate.
\end{proof}

This explains why algebra membership alone was insufficient in V32.  On a
single full $M_{14}$ operator envelope one has
\[
 m=14^2=196,
 \qquad
 m^2-d^2=196^2-196=38220.
\]
So the normalization map has a kernel of dimension at least $38220$ before any
typing information is used.  Merely showing that a primitive operation lies
somewhere in the full historical $M_{14}$ algebra cannot prove source
rigidity.  Exact source/target, grading, incidence, or character constraints
must narrow the operator module substantially.

More generally, if a typed branch is known to lie in an exact bimodule such as
\[
 \mathsf S_a=P_{t(a)}\,\mathcal A\,P_{s(a)},
\]
then its dimension is computed directly from the finite historical algebra,
and the product bank
$\mathsf S_a^*\mathsf S_a\subseteq P_{s(a)}\mathcal A P_{s(a)}$ is compiled by
its multiplication table.  Equation~\eqref{sel63:eq:dimension} is therefore a
cheap first test before any Choi reconstruction or optimization.

\section[Existing examples and the revised fork]{What V63 changes and what it does not}
\label{sel63:examples}

\paragraph{V35 forgotten channel.}
Taking the exact two-dimensional operator envelope
$\mathsf S=\operatorname{span}\{k_+,k_-\}$ reproduces V62's positive fourth
Gram and therefore rigidity.  The statement no longer depends on declaring
$k_\pm$ to be retained fine outcomes; it is enough that the Choi support is
known to lie in their historical operator span.

\paragraph{Retained sign and six accepted compiler branches.}
Their one-dimensional branch envelopes reproduce the V57/V62 positive effect
Grams.  Hence their product reuse follows equally from exact branch selection
rules plus normalization.  No additional amplitude interpretation is needed.

\paragraph{Raw idle source.}
The idle ray together with the six accepted rays gives a singular envelope
with exactly the one known relation proportional to
$I-\sum_eE_e$.  Here the envelope equals the actual support family, so the
single kernel direction is physical and reproduces the acceptance-correlation
parameter.  This is stronger information than singularity of an arbitrary
broad envelope.

\paragraph{Forgotten sign channel.}
Its exact two-dimensional Choi support gives a four-dimensional coefficient
domain but only a two-dimensional product span, hence nullity two.  If one knew
only a still larger structural envelope, singularity of that larger envelope
would not by itself establish the same nullity.

These examples sharpen the logical fork after V62:
\[
 \boxed{
 \begin{array}{c}
 \text{exact historical operator-space envelope}\\[.15em]
 \Downarrow\\[-.2em]
 \text{envelope fourth Gram from old word moments}\\[.15em]
 \Downarrow\\[-.2em]
 \mathbb E^{(4)}>0\;\Longrightarrow\;
 \text{unique source and product reuse}
 \end{array}}
 \label{sel63:eq:newpipeline}
\]
No historical efficient refinement is needed on this positive branch.

If the envelope fourth Gram is singular, there are two genuinely different
possibilities.  The actual support may be smaller and still rigid, in which
case one must add exact structural annihilators or identify the actual support
using V60--V61.  Or the actual support may occupy the singular directions, in
which case V59's coupling freedom is real.  The envelope alone does not decide
between them.

This is also the correct firewall for the directed V110 incidence.  A typed
module or character line may provide a narrow envelope, and its fourth moment
may show that normalization would force a unique branch \emph{if a normalized
historical branch is known to lie there}.  The calculation does not establish
that the Store actually contains that directed incidence, nor does it create
the required same-history word.

\section[Revised historical target]{Search first for exact selection rules, not hidden Kraus labels}
\label{sel63:status}

The remaining source audit can now be ordered more efficiently:
\begin{enumerate}[label=\textup{(\arabic*)},leftmargin=3em]
\item use the historical grammar, type data, and old word algebra to construct
      the narrowest exact operator-space envelopes $\mathsf S_a$ available on
      the same geometry-shorted cut;
\item apply the dimension screen \eqref{sel63:eq:dimension};
\item compile the envelope fourth Gram from the existing multiplication table
      or inserted word moments;
\item if it is positive definite, stop: normalization already fixes the local
      instrument and V59 fixes reuse;
\item if it is singular, narrow the envelope with additional exact
      annihilators or pass to the actual support programme of V60--V61.
\end{enumerate}

This is upstream of searching for a physically retained Kraus refinement.  It
asks only for exact \emph{selection rules} on the primitive branch.  Those
rules can be easier to inherit from source/target typing and finite word
algebra than a microscopic environment label.

\paragraph{Proved.}
The fourth product Gram of any operator-space envelope has kernel exactly equal
to V61's support-envelope normalization kernel.  Positive definiteness forces
a unique normalized instrument even when the actual support is strict.  The
criterion is basis independent, compiles from historical word moments, obeys
the dimension screen $\sum_am_a^2\le d^2$, and is monotone under exact
structural narrowing.  A strict qutrit envelope was exhibited for which
normalization removes an allowed but unoccupied Kraus direction uniquely.

\paragraph{Inherited inputs.}
V31 supplies finite multiplication after flatness.  V59 supplies the
fixed-marginal coupling theorem.  V60 identifies normalization rigidity with
Choi support.  V61 proves the support-envelope singleton theorem.  V62 supplies
the fourth-product Gram compiler.  V63 identifies these last two linear maps
without assuming that the envelope basis consists of historical Kraus
amplitudes.

\paragraph{Sharp boundary.}
A singular broad envelope proves only that the current structural information
is insufficient.  It does not prove nonextremality of the actual source.
Conversely, positive rigidity still requires exact support containment; a
finite-precision near-zero leakage estimate is not silently promoted into the
exact statement \eqref{sel63:eq:envelope}.  V61's zero-robustness boundary
therefore remains unchanged.

\paragraph{Remaining historical question.}
Determine the narrow exact typed operator-space envelope of the primitive
Store branches on the same cut, especially the branch that would have to carry
the V110 directed Store--Clifford/up-incidence structure.  If its envelope
fourth Gram is injective, no hidden efficient refinement is needed for the
reuse proof.  What is still missing is the historical occurrence and typing of
that branch itself.

Historical Standard-Model selection is not declared closed.

\paragraph{Verification boundary.}
The companion checker verifies the envelope-kernel identity on exact finite
examples, the strict-envelope singleton, basis covariance, the $M_{14}$
dimension obstruction, and the unchanged V35/sign/raw-idle census.  The
general finite-dimensional statements have the written proofs above.  No
empirical exact support containment, proof-assistant formalization, or
independent peer review is claimed.

\addtocontents{toc}{\protect\endgroup}
% END SOURCE-SELECTION V63 DEVELOPMENT BLOCK
% Append future development before the final document command.



\clearpage
\begin{center}
{\Large\bfseries Source-selection continuation V64}\\[.4em]
{\large Typed-bimodule obstruction and Schur compression:\\
why source/target labels alone are too broad, and how finite nonabelian covariance plus one swap parity isolates the determinant incidence line}
\end{center}
\addtocontents{toc}{\protect\begingroup}
\addtocontents{toc}{\protect\renewcommand{\protect\numberline}{\protect\makebox[2.1em][l]}}

\section[Move upstream of the envelope basis]{From arbitrary operator envelopes to typed representation modules}
\label{sel64:context}

V63 reduces the positive source-rigidity problem to an exact historical
operator-space envelope $\mathsf S_a$ and the injectivity of its normalization
map.  The remaining upstream question is how much of that envelope follows
from the type data themselves.  There are two sharply different answers.
Source/target projectors alone usually leave a full rectangular bimodule and
are far too weak.  Exact covariance and finite tensor parity can collapse that
same rectangle to a one-dimensional intertwiner line.

The results below are finite-dimensional.  They do not assume that the
candidate continuous group is already a physical gauge symmetry of the
historical Store.  Whenever an exact covariance action is invoked, its
historical identification remains a separate input.  The purpose is to
identify precisely which structural constraints would be sufficient to turn
the V63 envelope theorem into a source-rigidity certificate.

\section[Exact rectangular-bimodule nullity]{Source and target projectors alone almost never force rigidity}
\label{sel64:rect}

Let the input carrier split orthogonally as
\[
 A=\bigoplus_{s\in\mathcal S}A_s,
 \qquad P_s=P_{A_s},\qquad p_s=\dim A_s.
\]
For each recorded branch $a$, fix a source type $s(a)$ and a target space
$T_a$ of dimension $q_a\ge1$.  Suppose the only exact structural restriction
is
\begin{equation}
 \boxed{\mathsf S_a=\Hom(A_{s(a)},T_a).}
 \label{sel64:eq:rect-envelope}
\end{equation}
Thus the branch is typed, but otherwise every rectangular amplitude with those
source and target types is admitted.

\begin{theorem}[Exact normalization kernel of full typed rectangles]
\label{sel64:thm:rect-kernel}
Assume that every input sector $A_s$ occurs in at least one branch.  For the
envelopes \eqref{sel64:eq:rect-envelope}, V63's normalization synthesis map is,
after canonical identifications,
\begin{equation}
 \boxed{
 \Gamma_{\rm rect}\bigl((X_a)_a\bigr)
 =\bigoplus_{s\in\mathcal S}
   \sum_{a:\,s(a)=s}\Tr_{T_a}X_a,}
 \label{sel64:eq:partialtrace}
\end{equation}
where $X_a\in M_{q_ap_s}(\C)\simeq
M_{q_a}(\C)\otimes M_{p_s}(\C)$.  It is surjective onto
$\bigoplus_sM_{p_s}(\C)$ and has complex nullity
\begin{equation}
 \boxed{
 \nu_{\rm rect}
 =\sum_{s\in\mathcal S}p_s^2
   \left(\sum_{a:\,s(a)=s}q_a^2-1\right).}
 \label{sel64:eq:rect-nullity}
\end{equation}
Consequently the rectangular envelope is injective if and only if, for every
source sector $s$, exactly one branch leaves $s$ and its target dimension is
$q_a=1$.
\end{theorem}

\begin{proof}
Choose orthonormal bases $\{e_\alpha\}_{\alpha=1}^{p_s}$ of $A_s$ and
$\{f_i\}_{i=1}^{q_a}$ of $T_a$.  The matrix units
$B_{a,i\alpha}=|f_i\rangle\langle e_\alpha|$ form an orthonormal basis of
$\mathsf S_a$.  Their products are
\[
 B_{a,i\alpha}^*B_{a,j\beta}
 =\delta_{ij}|e_\alpha\rangle\langle e_\beta|.
\]
Hence the coefficient matrix $X_a$ contributes exactly its partial trace over
the target index.  This proves \eqref{sel64:eq:partialtrace}.  For each fixed
$s$, any one nonzero target sector makes the sum of partial traces surjective
onto $M_{p_s}$, so the rank is $p_s^2$.  The coefficient-domain dimension on
that source block is
$p_s^2\sum_{a:s(a)=s}q_a^2$.  Subtracting the rank and summing over orthogonal
source blocks proves \eqref{sel64:eq:rect-nullity}.  Since the $q_a$ are
positive integers, zero nullity is equivalent to a single $q_a=1$ branch for
each $s$.
\end{proof}

\begin{corollary}[The bare V110 up-incidence type is far too broad]
\label{sel64:cor:up-rect}
For the typed space
\[
 \mathscr U=\Hom(C\otimes W\otimes W,\Lambda^2C^*),
 \qquad \dim C=3,\quad\dim W=2,
\]
one has input dimension $p=12$ and target dimension $q=3$.  If the historical
information is only ``a $12\to3$ up-incidence map'', then the corresponding
rectangular envelope has
\begin{equation}
 \boxed{
 \nu_{\mathscr U}=(3^2-1)12^2=1152.}
 \label{sel64:eq:up-nullity}
\end{equation}
Thus source/target typing alone cannot identify or rigidify the determinant
incidence.
\end{corollary}

This strengthens the V32 warning that membership in the full historical
algebra is too weak.  Even after the correct rectangular source and target
ports are known, the normalization kernel remains enormous unless additional
selection rules remove most of the rectangle.

\section[Schur compression of a typed rectangle]{Exact covariance converts type data into a much smaller intertwiner envelope}
\label{sel64:schur}

Let a compact group $G$ act unitarily on a source representation $H_s$ and a
target representation $H_t$.  For a unitary character $\chi:G\to U(1)$ define
the semi-intertwiner envelope
\begin{equation}
 \mathsf S_\chi
 =\{K:H_s\to H_t:
 R_t(g)K R_s(g)^*=\chi(g)K\ \text{for all }g\in G\}.
 \label{sel64:eq:semi-intertwiner}
\end{equation}
Equivalently, $K$ intertwines $R_s$ with the twisted target
$\chi^{-1}R_t$.

Write the isotypic decompositions
\[
 H_s\simeq\bigoplus_{\lambda}V_\lambda\otimes\C^{m_\lambda},
 \qquad
 (\chi^{-1}H_t)\simeq
 \bigoplus_{\lambda}V_\lambda\otimes\C^{n_\lambda}.
\]

\begin{theorem}[Schur-compressed envelope and symmetry-reduced dimension screen]
\label{sel64:thm:schur}
The envelope \eqref{sel64:eq:semi-intertwiner} is canonically isomorphic to
\begin{equation}
 \boxed{
 \mathsf S_\chi\simeq
 \bigoplus_\lambda I_{V_\lambda}\otimes
 M_{n_\lambda\times m_\lambda}(\C),
 \qquad
 \dim\mathsf S_\chi=\sum_\lambda m_\lambda n_\lambda.}
 \label{sel64:eq:schur-envelope}
\end{equation}
Moreover every product $K^*L$, with $K,L\in\mathsf S_\chi$, belongs to the
source commutant
\begin{equation}
 \boxed{
 R_s(G)'\simeq
 \bigoplus_\lambda I_{V_\lambda}\otimes M_{m_\lambda}(\C),
 \qquad
 \dim R_s(G)'=\sum_\lambda m_\lambda^2.}
 \label{sel64:eq:commutant}
\end{equation}
Hence for several semi-intertwiner branch envelopes on the same source
representation, a necessary condition for V63 injectivity is
\begin{equation}
 \boxed{
 \sum_a\left(\sum_\lambda m_\lambda n_{a,\lambda}\right)^2
 \le \sum_\lambda m_\lambda^2.}
 \label{sel64:eq:schur-screen}
\end{equation}
\end{theorem}

\begin{proof}
After twisting the target by $\chi^{-1}$, the condition is ordinary
$G$-equivariance.  Schur's lemma makes an intertwiner scalar on each irreducible
factor $V_\lambda$ and arbitrary only on the multiplicity spaces, giving
\eqref{sel64:eq:schur-envelope}.  The product of two such intertwiners is
$G$-invariant on the source and therefore lies in the commutant
\eqref{sel64:eq:commutant}.  V63's normalization map is built entirely from
these products.  Its coefficient-domain dimension is the left side of
\eqref{sel64:eq:schur-screen}, while its image lies in a vector space of the
dimension on the right.  Injectivity requires the displayed inequality.
\end{proof}

\begin{corollary}[Multiplicity-free irreducible source]
\label{sel64:cor:schur-rays}
Suppose the source representation is irreducible and $k$ recorded branches each
have a nonzero one-dimensional semi-intertwiner envelope.  Then
$K_a^*K_a=c_aI$ with $c_a>0$, and the normalization kernel on that source
sector has dimension exactly
\begin{equation}
 \boxed{k-1.}
 \label{sel64:eq:kminusone}
\end{equation}
In particular one such branch is normalization-rigid on the sector, whereas
several symmetry-equivalent ray branches require $k-1$ additional independent
record calibrations to separate their weights.
\end{corollary}

\begin{proof}
By Schur's lemma each $K_a^*K_a$ is a positive scalar multiple of the identity.
The normalization map is therefore the nonzero linear functional
$(x_1,\ldots,x_k)\mapsto\sum_ac_ax_a$ into the one-dimensional commutant.  Its
kernel has dimension $k-1$.
\end{proof}

Thus symmetry can compress a large typed rectangle drastically, but symmetry
alone does not distinguish several recorded copies of the same irreducible
intertwiner type.

\section[Finite derivation of the determinant line]{Nonabelian covariance plus one weak-slot swap parity fixes the V110 character line}
\label{sel64:detline}

The V110 incidence is a particularly sharp instance of Schur compression.
Let
\[
 G_{\rm na}=SU(C)\times SU(W),
 \qquad \dim C=3,\quad\dim W=2,
\]
and let $\tau_W$ swap the two weak tensor factors.  Define
\begin{align}
 \mathsf S_{\rm alt}^{\rm na}
 :=\{T\in\Hom(C\otimes W\otimes W,\Lambda^2C^*):\ &
 T(I_C\otimes\tau_W)=-T,\nonumber\\[-.2em]
 &R_{\rm out}(U)T R_{\rm in}(U,V)^*=T
 \ \forall(U,V)\in G_{\rm na}\}.
 \label{sel64:eq:na-alt-envelope}
\end{align}

\begin{theorem}[Finite nonabelian selection of the determinant-incidence line]
\label{sel64:thm:detline}
The envelope \eqref{sel64:eq:na-alt-envelope} is one-dimensional.  Every
nonzero $T$ in it factors through
\[
 C\otimes\Lambda^2W
 \simeq C
\]
and is an isomorphism onto $\Lambda^2C^*$ up to an overall scalar.  When the
same tensor is regarded under the full central torus of $U(C)\times U(W)$, its
weight is automatically
\begin{equation}
 \boxed{
 R_{\rm out}(e^{i\alpha}I_C)
 T
 R_{\rm in}(e^{i\alpha}I_C,e^{i\beta}I_W)^*
 =e^{-i(3\alpha+2\beta)}T.}
 \label{sel64:eq:derived-character}
\end{equation}
Thus the primitive determinant character is derived from the tensor type,
nonabelian covariance, and weak-slot antisymmetry; it need not be separately
inserted as a central selection rule.
\end{theorem}

\begin{proof}
The swap relation kills $C\otimes\operatorname{Sym}^2W$ and factors $T$ through
$C\otimes\Lambda^2W$.  Since $\Lambda^2W$ is the trivial one-dimensional
representation of $SU(2)$, the surviving source is the defining $SU(3)$
representation $C$.  The standard exterior-volume isomorphism identifies
$\Lambda^2C^*$ with the same irreducible $SU(3)$ representation.  Schur's
lemma therefore makes the intertwiner space one-dimensional.

For central phases, the source $C\otimes W\otimes W$ has weight
$e^{i(\alpha+2\beta)}$, while $\Lambda^2C^*$ has weight $e^{-2i\alpha}$.
The induced Hom weight is target times the inverse source weight, yielding
$e^{-i(3\alpha+2\beta)}$.
\end{proof}

\begin{remark}[Finite linear compiler]
Because $SU(3)\times SU(2)$ is connected, the covariance equations in
\eqref{sel64:eq:na-alt-envelope} are equivalent to the finitely many Lie-algebra
intertwining equations for a basis of $\mathfrak{su}(3)\oplus\mathfrak{su}(2)$,
together with the single swap equation.  Hence the one-dimensional envelope
can be recovered as the kernel of a finite exact linear system on the
$36$-dimensional space $\mathscr U$; no group integration is required.
\end{remark}

Choose the canonical generator $T_0$ normalized so that
\begin{equation}
 \boxed{T_0^*T_0=P_{\rm alt},}
 \label{sel64:eq:T0effect}
\end{equation}
where $P_{\rm alt}$ projects onto
$C\otimes\Lambda^2W\subset C\otimes W\otimes W$.  Equation
\eqref{sel64:eq:T0effect} follows again from Schur's lemma; an explicit volume
form fixes the normalization.

\section[What normalization can now fix]{Isolation of the alternating input sector versus unresolved competing branches}
\label{sel64:normalization}

The determinant line being one-dimensional does not by itself prove that a
historical branch exists, nor that its weight is separated from every other
branch.  Normalization supplies the missing scalar only under an additional
support-isolation condition.

\begin{proposition}[Alternating-sector isolation fixes the incidence weight]
\label{sel64:prop:isolation}
Let one branch have envelope $\mathsf S_D=\C T_0$ with
$T_0^*T_0=P_{\rm alt}$.  Suppose every other branch amplitude $K$ obeys
\begin{equation}
 \boxed{K P_{\rm alt}=0.}
 \label{sel64:eq:other-annihilate}
\end{equation}
Then in every normalized instrument contained in these envelopes, the
coefficient of $|T_0\rangle\!\rangle\langle\!\langle T_0|$ is exactly one.
All remaining normalization freedom is confined to the orthogonal symmetric
input sector $P_{\rm sym}=I-P_{\rm alt}$.
\end{proposition}

\begin{proof}
Compress trace preservation
$\sum_aK_a^*K_a=I$ on both sides by $P_{\rm alt}$.  Every non-determinant
branch vanishes by \eqref{sel64:eq:other-annihilate}, while the determinant
branch contributes $xP_{\rm alt}$.  Therefore
$xP_{\rm alt}=P_{\rm alt}$ and $x=1$.  Compression by $P_{\rm sym}$ contains
all remaining equations and is independent of this scalar.
\end{proof}

If another branch also acts on $C\otimes\Lambda^2W$, the conclusion no longer
follows from normalization alone.  This is exactly the multiplicity warning in
Corollary~\ref{sel64:cor:schur-rays}: several ray branches on the same
irreducible source sector share a scalar effect direction and require extra
record information.

The historical advantage is therefore precise.  To fix the determinant
branch weight one need not reconstruct its complete Choi packet.  It is enough
to establish:
\begin{enumerate}[label=\textup{(\roman*)},leftmargin=3em]
\item the branch is nonzero and obeys the weak-swap and nonabelian covariance
      equations of \eqref{sel64:eq:na-alt-envelope};
\item no competing historical branch acts on the same alternating input
      sector, or the competing weights are separately calibrated.
\end{enumerate}
Neither statement is supplied merely by the abstract existence of the V110
intertwiner line.

\section[Revised historical target]{Search for finite tensor selection rules before Choi support}
\label{sel64:status}

V64 changes the upstream order again.  The cheapest exact test is now:
\begin{enumerate}[label=\textup{(\arabic*)},leftmargin=3em]
\item identify the actual source/target projectors of a primitive typed branch;
\item apply Theorem~\ref{sel64:thm:rect-kernel}.  If the resulting full
      rectangular envelope already has large nullity, do not waste effort on
      its fourth Gram;
\item impose exact finite tensor selection rules---swap parity, grading,
      incidence, and any historically justified nonabelian covariance---and
      compute the resulting intertwiner envelope by a finite linear nullspace;
\item apply the Schur dimension screen \eqref{sel64:eq:schur-screen} and then
      V63's envelope fourth Gram only on that compressed module;
\item if the determinant incidence is isolated on the alternating input
      sector, use Proposition~\ref{sel64:prop:isolation} to fix its normalization
      weight without full process tomography.
\end{enumerate}

For the V110 up-incidence, this calculation draws a sharp line.  The bare
$12\to3$ typed rectangle has nullity $1152$ and is useless for rigidity.  Exact
weak-slot antisymmetry together with exact $SU(3)\times SU(2)$ covariance
collapses that rectangle to the unique determinant-incidence line and derives
the central weight $-(3,2)$.  What remains missing is not the abstract line but
its \emph{historical occurrence on one same-history Store word}, together with
the required covariance/typing provenance and the absence or calibration of
competing alternating-sector branches.

\paragraph{Proved.}
The normalization nullity of arbitrary full rectangular source/target
bimodules is given exactly by \eqref{sel64:eq:rect-nullity}.  Exact covariance
compresses a typed rectangle to the Schur multiplicity envelope
\eqref{sel64:eq:schur-envelope}.  For the up-incidence type, weak-slot
antisymmetry plus nonabelian covariance gives a one-dimensional envelope and
automatically yields the determinant character $e^{-i(3\alpha+2\beta)}$.
If this branch is the sole branch acting on the alternating input sector,
trace preservation fixes its coefficient exactly.

\paragraph{Inherited inputs.}
V26 supplies the already audited determinant-incidence interpretation and
five-slot character count.  V31 supplies finite multiplication after
flatness.  V59--V63 supply the normalization-kernel, support-envelope, and
operator-envelope rigidity theorems.  V64 does not re-prove historical
occurrence of any V110 word.

\paragraph{Sharp boundary.}
The $SU(3)\times SU(2)$ covariance used in
Theorem~\ref{sel64:thm:detline} is an exact structural input.  V18 already
showed that nonabelian covariance cannot be silently inferred from the desired
central answer.  Likewise, a formal tensor swap in the representation space is
useful only after the historical weak-factor ordering and its action have been
identified on the same cut.  Approximate covariance or approximate
antisymmetry yields approximate envelope leakage, not the exact one-dimensional
support statement used here.

\paragraph{Remaining historical question.}
Land one actual primitive Store branch on the same geometry-shorted cut and
verify the finite linear constraints defining
$\mathsf S_{\rm alt}^{\rm na}$.  If that succeeds, the determinant character
and one-dimensional operator envelope follow without assuming a central
$U(1)$ direction in advance.  One must then test whether any competing branch
also acts on the alternating sector and, separately, whether the neutral
Store--Clifford shadow is coherent with this incidence on the same resolved
history.

Historical Standard-Model selection is not declared closed.

\paragraph{Verification boundary.}
The companion checker verifies the rectangular nullity formula on exact finite
examples, the $12\to3$ V110 count, the one-dimensional finite Lie/swap kernel,
the canonical effect $T_0^*T_0=P_{\rm alt}$, and the normalization-isolation
claim.  It also reruns the V63 regression checker.  No empirical historical
covariance, exact swap action, or Store occurrence is claimed.

\addtocontents{toc}{\protect\endgroup}
% END SOURCE-SELECTION V64 DEVELOPMENT BLOCK
% Append future development before the final document command.



\clearpage
\begin{center}
{\Large\bfseries Source-selection continuation V65}\\[.4em]
{\large Finite matrix-unit naturality and a single positive Choi witness:\\
removing the weak-swap premise and replacing continuous nonabelian covariance by six exact historical defect identities}
\end{center}
\addtocontents{toc}{\protect\begingroup}
\addtocontents{toc}{\protect\renewcommand{\protect\numberline}{\protect\makebox[2.1em][l]}}

\section[Move upstream of covariance and swap]{Replace continuous covariance and weak-slot parity by finite matrix-unit identities}
\label{sel65:context}

V64 compressed the bare $12\to3$ up-incidence rectangle to one determinant
line by assuming two exact structural inputs: $SU(3)\times SU(2)$ covariance
and odd parity under interchange of the two weak slots.  The remaining
upstream question is whether these conditions can themselves be replaced by
finite identities internal to the already identified matrix factors.

The answer is stronger than the V64 entrance.  Once the color and weak
factors are represented on the typed legs, six off-diagonal matrix-unit
naturality equations leave exactly the same one-dimensional incidence line.
No weak-swap equation is needed, no central generator is used, and no
continuous group integration appears.  Moreover the six equations combine
into one positive Choi witness.  Exact vanishing of that witness forces an
\emph{a priori mixed} branch onto the determinant line, so an efficient Kraus
refinement need not be supplied in advance.

The same-cut representation data remain essential.  An abstract algebra
isomorphism $M_3\oplus M_2\oplus\C^2$ does not by itself identify which
historical source and target legs carry the defining actions used below.
Likewise the output type $\Lambda^2C^*$ and the two weak tensor legs must be
actual typed carrier data, not labels assigned after seeing the desired
answer.

\section[Six finite naturality defects]{The determinant line without a weak-swap assumption}
\label{sel65:six}

Let $C\cong\C^3$ and $W\cong\C^2$, and keep the V64 typed amplitude space
\[
 \mathscr U=\Hom(C\otimes W\otimes W,\Lambda^2C^*).
\]
For $X\in\End(C)$ let $d\rho_C^{\rm in}(X)=X\otimes I_W\otimes I_W$.
On $C^*$ the infinitesimal dual action is $-X^{\mathsf T}$, and its natural
exterior-square action on $\Lambda^2C^*$ will be denoted
$d\rho_C^{\rm out}(X)$.  Explicitly,
\begin{equation}
 d\rho_C^{\rm out}(X)(\varphi\wedge\psi)
 =-(\varphi\circ X)\wedge\psi
   -\varphi\wedge(\psi\circ X).
 \label{sel65:eq:color-out}
\end{equation}
For $Y\in\End(W)$ put
\begin{equation}
 d\rho_W^{\rm in}(Y)
 =I_C\otimes Y\otimes I_W+I_C\otimes I_W\otimes Y.
 \label{sel65:eq:weak-in}
\end{equation}
The output has no weak factor.  Define the linear naturality defects
\begin{align}
 \mathcal R_X(T)&=d\rho_C^{\rm out}(X)T-Td\rho_C^{\rm in}(X),
 \label{sel65:eq:R}\\
 \mathcal S_Y(T)&=-T d\rho_W^{\rm in}(Y).
 \label{sel65:eq:S}
\end{align}
Choose any matrix-unit systems $E_{ij}$ of the identified color factor and
$F_{ab}$ of the identified weak factor, and use the six off-diagonal
generators
\begin{equation}
 \mathcal G_C=\{E_{12},E_{21},E_{23},E_{32}\},\qquad
 \mathcal G_W=\{F_{12},F_{21}\}.
 \label{sel65:eq:gens}
\end{equation}
They involve only the nonabelian factors.  No central $U(1)$ direction and
no tensor-slot swap occurs in their definition.

\begin{theorem}[Six-generator naturality compiler]
\label{sel65:thm:six}
The common kernel
\begin{equation}
 \mathsf S_{\rm nat}
 =\bigcap_{X\in\mathcal G_C}\ker\mathcal R_X
  \cap
  \bigcap_{Y\in\mathcal G_W}\ker\mathcal S_Y
 \subset\mathscr U
 \label{sel65:eq:natline}
\end{equation}
is one-dimensional.  Every nonzero $T\in\mathsf S_{\rm nat}$ satisfies:
\begin{enumerate}[label=\textup{(\roman*)},leftmargin=2.8em]
\item exact $SU(3)\times SU(2)$ covariance;
\item automatic weak-slot antisymmetry
\begin{equation}
 \boxed{T(I_C\otimes\tau_W)=-T;}
 \label{sel65:eq:auto-swap}
\end{equation}
\item the central character
\begin{equation}
 \boxed{e^{-i(3\alpha+2\beta)}}
 \label{sel65:eq:auto-char}
\end{equation}
under $(e^{i\alpha}I_C,e^{i\beta}I_W)\in U(3)\times U(2)$.
\end{enumerate}
Thus the separate weak-swap hypothesis in V64 is redundant once the finite
weak matrix-unit naturality equations are established.
\end{theorem}

\begin{proof}
The four color matrix units in \eqref{sel65:eq:gens} generate
$\mathfrak{sl}_3(\C)$ under complex linear combinations and Lie brackets;
the two weak units generate $\mathfrak{sl}_2(\C)$.  If the displayed
residuals vanish on the generators, they also vanish on commutators because
\[
 [d\rho(X),d\rho(Z)]T-T[d\rho_{\rm in}(X),d\rho_{\rm in}(Z)]
\]
is obtained by composing the two corresponding intertwining identities.
Hence \eqref{sel65:eq:natline} is the full
$\mathfrak{sl}_3\oplus\mathfrak{sl}_2$ intertwiner space.

For the weak factor,
\[
 W\otimes W=\operatorname{Sym}^2W\oplus\Lambda^2W,
\]
where $\operatorname{Sym}^2W$ is the nontrivial three-dimensional irreducible
$\mathfrak{sl}_2$ module and $\Lambda^2W$ is the trivial line.  Since the
output is weak-trivial, every weak intertwiner annihilates the symmetric
summand and factors through $\Lambda^2W$.  Therefore the swap acts by $-1$
on every nonzero survivor, proving \eqref{sel65:eq:auto-swap}; no separate
swap datum was used.

After that factorization the color problem is
\[
 \Hom_{\mathfrak{sl}_3}(C,\Lambda^2C^*).
\]
Both factors are the same irreducible $\mathfrak{sl}_3$ module, so this Hom
space is one-dimensional by Schur's lemma.  This proves the dimension claim.
The compact groups $SU(3)$ and $SU(2)$ are connected and generated by
exponentials of their Lie algebras, so the infinitesimal identities exponentiate
to exact group covariance.  Finally the central phase count depends only on
the tensor type: the output $\Lambda^2C^*$ has color weight $-2\alpha$ and
the inverse input contributes $-(\alpha+2\beta)$, giving
$-(3\alpha+2\beta)$ exactly as in V64.
\end{proof}

The two nonabelian factors play complementary roles.  The weak equations
alone leave a nine-dimensional space
\begin{equation}
 \ker\mathcal S\cong\Hom(C,\Lambda^2C^*)\otimes\Lambda^2W^*,
 \qquad \dim=9,
 \label{sel65:eq:weakonly}
\end{equation}
and already force weak antisymmetry.  The color equations alone leave a
four-dimensional space, corresponding to the free $W\otimes W$ coefficient.
Together they leave one line.  Thus the disappearance of the swap premise is
not achieved by discarding the weak factor; it is a consequence of its exact
matrix-unit action.

\section[A single positive Choi witness]{Mixed primitive branches can be certified without a Kraus refinement}
\label{sel65:witness}

Equip $\mathscr U$ with the Hilbert--Schmidt inner product and vectorize it as
a $36$-dimensional coefficient space.  Let $L_X$ and $M_Y$ denote the matrices
of \eqref{sel65:eq:R}--\eqref{sel65:eq:S}.  Define the positive naturality
operator
\begin{equation}
 \boxed{
 H_{\rm nat}
 =\sum_{X\in\mathcal G_C}L_X^*L_X
  +\sum_{Y\in\mathcal G_W}M_Y^*M_Y\succeq0.}
 \label{sel65:eq:Hnat}
\end{equation}
Its kernel is exactly $\mathsf S_{\rm nat}$.  For the standard normalized
matrix units and the orthonormal wedge basis, the exact spectrum is
\begin{equation}
 \boxed{
 \operatorname{spec}H_{\rm nat}
 =\{0^{(1)},2^{(5)},3^{(4)},4^{(7)},5^{(8)},
      6^{(4)},7^{(4)},8^{(2)},10^{(1)}\}.}
 \label{sel65:eq:spectrum}
\end{equation}
Let $P_D$ be the orthogonal projector onto the unique determinant-incidence
line.  Equation \eqref{sel65:eq:spectrum} gives the operator inequality
\begin{equation}
 \boxed{H_{\rm nat}\succeq2(I-P_D).}
 \label{sel65:eq:gap}
\end{equation}
A different orthonormal matrix-unit system is obtained by unitary conjugation
inside the color and weak factors; it unitarily conjugates the entire defect
system.  Hence the zero set and the spectrum in \eqref{sel65:eq:spectrum} do
not depend on a preferred internal basis.

\begin{theorem}[One positive defect scalar certifies the exact incidence ray]
\label{sel65:thm:positive}
Let $J_D\succeq0$ be the Choi operator of an arbitrary recorded primitive
branch on the V110 typed port.  No rank assumption is made.  Put
\begin{equation}
 d_{\rm nat}=\Tr(H_{\rm nat}J_D).
 \label{sel65:eq:defect}
\end{equation}
Then
\begin{equation}
 \boxed{
 \Tr[(I-P_D)J_D]\le\frac12d_{\rm nat}.}
 \label{sel65:eq:leakage}
\end{equation}
In particular,
\begin{equation}
 \boxed{
 d_{\rm nat}=0
 \quad\Longrightarrow\quad
 J_D=x\,|T_0\rangle\!\rangle\langle\!\langle T_0|
 \quad\text{for some }x\ge0,}
 \label{sel65:eq:rayforced}
\end{equation}
where $T_0$ is any normalized generator of $\mathsf S_{\rm nat}$.
If additionally $\Tr J_D>0$, then $x>0$: the determinant incidence actually
occurs and is automatically efficient, nonabelian-covariant, and weak-odd.
\end{theorem}

\begin{proof}
Equation \eqref{sel65:eq:gap}, positivity of $J_D$, and trace monotonicity give
\eqref{sel65:eq:leakage}.  If $d_{\rm nat}=0$, the positive-zero identity from
V61 gives
$\operatorname{supp}J_D\subseteq\ker H_{\rm nat}=\operatorname{Ran}P_D$.
That support is one-dimensional by Theorem~\ref{sel65:thm:six}, proving
\eqref{sel65:eq:rayforced}.  A positive operator on a one-dimensional support
is a nonnegative scalar multiple of its rank-one projector.  Positive trace
excludes the zero scalar.
\end{proof}

This is stronger than first selecting a favourable Kraus decomposition of a
mixed map.  The witness acts on the primitive branch Choi packet itself.  A
zero value forces every possible Kraus vector into the same one-dimensional
line and therefore proves efficiency rather than assuming it.

The exact-zero premise remains essential for the exact ray statement.  A
small positive value gives the quantitative leakage bound
\eqref{sel65:eq:leakage}; it is not rounded to zero.  This retains the V61
support firewall while making the relevant gap explicit for the present typed
problem.

\section[Historical compilation of the witness]{No continuous group action has to be measured}
\label{sel65:historical}

The witness \eqref{sel65:eq:Hnat} is finite once the typed matrix-factor
representation has been identified.  It can enter the historical programme in
two different ways.

\paragraph{Contextual-Choi route.}
$H_{\rm nat}$ is one Hermitian positive functional on the primitive branch
Choi packet.  V24's tester-span criterion decides whether its expectation is
already identifiable by signed postprocessing of the actual contextual frame.
If it is, the single scalar $d_{\rm nat}$ can be reconstructed without full
process tomography.  Exact zero then invokes Theorem~\ref{sel65:thm:positive}
and V61's positive-annihilator logic.  Statistical identifiability of the
functional is not the same as physical execution of $H_{\rm nat}$ as an
effect.

\paragraph{Word-native route.}
If the branch amplitudes and the identified matrix units are themselves
historical-algebra operations, each residual in
\eqref{sel65:eq:R}--\eqref{sel65:eq:S} is a finite linear combination of
historical products.  Its squared Hilbert--Schmidt norm is therefore an old
word moment.  After V31 flatness, the finite multiplication table computes all
six residual norms and their sum exactly.  No new continuous twirl, phase
reference, or tensor-swap operation is introduced.

The second route is a convenience only when the amplitudes are already
word-native.  The first route remains valid for an a priori mixed branch.  In
neither route is a determinant character inserted: all six generators are
traceless off-diagonal matrix units of the two nonabelian factors.

\begin{corollary}[Two-scalar historical incidence certificate]
\label{sel65:cor:two-scalar}
Suppose the same-cut typed port and its matrix-factor actions have been
identified, and the actual source data determine
\begin{equation}
 \boxed{d_{\rm nat}=0,\qquad \Tr J_D>0.}
 \label{sel65:eq:two-scalars}
\end{equation}
Then a nonzero primitive determinant-incidence branch occurs on that cut.
No weak-slot swap, central $U(1)$ generator, continuous covariance test, Kraus
refinement, or full Choi reconstruction is additionally required to identify
its operator ray.
\end{corollary}

The word ``two-scalar'' refers only to the final certification once the typed
matrix-factor representation has already been landed.  Establishing that
representation can require much more historical data; the corollary does not
hide that cost.

\section[Normalization and the remaining source rows]{What V65 removes, and what it does not}
\label{sel65:normalization}

If the V64 alternating-sector isolation hypothesis also holds, choose $T_0$
so that
\[
 T_0^*T_0=P_{\rm alt}.
\]
Then \eqref{sel65:eq:rayforced} and trace preservation compressed to
$P_{\rm alt}$ fix $x=1$ exactly.  Thus the combined route is
\begin{equation}
 \boxed{
 d_{\rm nat}=0
 \ +\ \Tr J_D>0
 \ +\ \text{alternating-sector isolation}
 \ \Longrightarrow\ 
 J_D=|T_0\rangle\!\rangle\langle\!\langle T_0|.}
 \label{sel65:eq:full-cert}
\end{equation}
If another branch acts on the same alternating source sector, V64's scalar
multiplicity ambiguity remains and must be calibrated from the retained
record.

Several historical rows also remain untouched:
\begin{enumerate}[label=\textup{(\roman*)},leftmargin=2.8em]
\item the internal $M_3$ and $M_2$ factors and their matrix-unit systems must
      act on the correct same-cut typed carrier;
\item the two weak tensor legs and the output type $\Lambda^2C^*$ must be
      historical typing data.  V65 removes the need to identify their swap
      as an independent operation, but it does not manufacture the two legs;
\item an abstract isomorphism of the internal algebra does not prove that a
      primitive Store branch has the Choi functional tested by
      $H_{\rm nat}$;
\item the neutral Store--Clifford shadow still has to be coherent with this
      determinant incidence on the same resolved history.
\end{enumerate}

In particular the V32 warning survives intact: membership in a saturated full
matrix algebra is too weak to identify a physical typing action.  The present
naturality relations become historical only after the represented matrix
factors and their tensor-leg actions have been identified.

\section[Revised historical target]{Search for finite naturality defects before continuous symmetry}
\label{sel65:status}

The cheapest exact V110 incidence audit is now:
\begin{enumerate}[label=\textup{(\arabic*)},leftmargin=3em]
\item identify the same-cut $C$, $W$, $W$ source legs and $\Lambda^2C^*$
      target together with the already reconstructed $M_3$ and $M_2$ actions;
\item form the six matrix-unit residuals in
      \eqref{sel65:eq:R}--\eqref{sel65:eq:S};
\item compile or reconstruct the single positive scalar
      $d_{\rm nat}=\Tr(H_{\rm nat}J_D)$;
\item if $d_{\rm nat}=0$ and the branch is nonzero, conclude the exact
      determinant-incidence ray, automatic weak antisymmetry, and exact
      nonabelian covariance;
\item only then use V64 normalization isolation and test the same-history
      neutral Store--Clifford coherence.
\end{enumerate}

This removes two external-looking inputs from V64.  The weak swap is no longer
an independent premise, and the continuous nonabelian covariance statement is
replaced by finitely many matrix-unit identities inside the identified source
algebra.  The central determinant direction is still obtained only afterwards
from the tensor type.

\paragraph{Proved.}
The six off-diagonal matrix-unit naturality equations have a one-dimensional
common kernel in the V110 incidence space.  Weak antisymmetry follows from the
weak equations and requires no separate swap condition.  The finite equations
exponentiate to exact $SU(3)\times SU(2)$ covariance.  Their positive Choi
defect operator has the exact spectrum \eqref{sel65:eq:spectrum} and gap two,
which yields the leakage estimate \eqref{sel65:eq:leakage}.  Exact zero forces
an arbitrary positive primitive Choi packet onto the unique determinant line.

\paragraph{Inherited inputs.}
V15--V16 supply the represented internal $M_3,M_2$ factors and the typed
incidence context; V24 supplies contextual functional identifiability; V31
supplies finite multiplication after flatness; V61 supplies the positive-zero
support principle; and V64 supplies the normalization-isolation handoff.  No
historical Store occurrence is imported from those results.

\paragraph{Sharp boundary.}
The finite defects still presuppose the correct represented matrix-factor
actions on the typed legs.  Approximate residuals yield only
\eqref{sel65:eq:leakage}.  The zero map also has zero defect, so historical
occurrence requires a separate nonzero branch mass.  A different output type
or different tensor-leg action defines a different naturality problem and
cannot be silently projected onto the one solved here.

\paragraph{Remaining historical question.}
Populate one actual primitive Store branch on the V110 same-cut typed port and
certify the finite matrix-unit defect scalar \eqref{sel65:eq:defect} together
with nonzero branch mass.  If that succeeds, the nonabelian covariance, weak
antisymmetry, determinant character and one-dimensional operator ray no longer
need separate historical assumptions.  One must still exclude or calibrate
competing alternating-sector branches and establish same-history coherence
with the neutral Store--Clifford component.

Historical Standard-Model selection is not declared closed.

\paragraph{Verification boundary.}
The companion checker constructs the intrinsic $\Lambda^2C^*$ action, verifies
the six-generator kernel and automatic swap parity, evaluates the exact
positive-defect spectrum and gap, checks the color-only/weak-only dimensions,
and reruns the V64 regression checker.  It does not certify that the historical
Store supplies the required typed matrix units, exact defect zero, nonzero
branch occurrence, or same-history Clifford coherence.

\addtocontents{toc}{\protect\endgroup}
% END SOURCE-SELECTION V65 DEVELOPMENT BLOCK
% Append future development before the final document command.

\end{document}
